\documentclass[11pt,oneside]{book}
\usepackage[margin=1in]{geometry}
\usepackage{amsmath,amssymb,amsthm,mathtools}
\usepackage{booktabs,longtable,array}
\usepackage{enumitem}
\usepackage[colorlinks=true,linkcolor=blue,urlcolor=blue,citecolor=blue]{hyperref}
\usepackage[nameinlink,noabbrev]{cleveref}

\newtheorem{theorem}{Theorem}[chapter]
\newtheorem{lemma}[theorem]{Lemma}
\newtheorem{proposition}[theorem]{Proposition}
\newtheorem{corollary}[theorem]{Corollary}
\newtheorem{assessment}[theorem]{Assessment}
\newtheorem{firewall}[theorem]{Claim firewall}
\newtheorem{openproblem}[theorem]{Open problem}
\newtheorem{record}[theorem]{Record}
\newtheorem{programme}[theorem]{Programme}
\theoremstyle{definition}
\newtheorem{definition}[theorem]{Definition}
\theoremstyle{remark}
\newtheorem{remark}[theorem]{Remark}

\newcommand{\MGclosed}{\textsc{closed}}
\newcommand{\MGopen}{\textsc{open}}
\newcommand{\MGpartial}{\textsc{partial}}
\newcommand{\opnorm}[1]{\lVert#1\rVert_{\mathrm{op}}}
\newcommand{\trnorm}[1]{\lVert#1\rVert_{1}}
\newcommand{\Tr}{\operatorname{Tr}}
\newcommand{\cB}{\mathcal B}
\newcommand{\cU}{\mathcal U}
\newcommand{\cV}{\mathcal V}

\title{Renewal Yang--Mills Mass-Gap Research Notes\\
Third Series, Version V1}
\author{Proof-indexed continuation after archives f1 and f2}
\date{September 2026}

\begin{document}
\maketitle
\tableofcontents

\chapter{Context, archive map, and exact claim boundary}
\label{ch:V1-context}

\section{Purpose of this restarted notebook}

This notebook continues a constructive Yang--Mills programme whose
ultimate target is a four-dimensional continuum gauge theory satisfying
the Osterwalder--Schrader axioms and possessing a strictly positive
physical spectral gap.  The work does not postulate the principal
uniform cluster estimate.  It attempts to derive that estimate from
exact Wilson heat-kernel algebra, physical Fourier resolution,
product-state capacity, probability marks, and a finite connected-tree
compiler.

Two completed source sequences are now frozen:
\begin{align*}
 {\cal A}_1&=
 \texttt{Renewal\_Yang\_Mills\_Mass\_Gap\_Research\_Notes\_V135\_f1.tex},\\
 {\cal A}_2&=
 \texttt{Renewal\_Yang\_Mills\_Mass\_Gap\_Research\_Notes\_V100\_f2.tex}.
\end{align*}
The present file starts a new active sequence at V1 because appending to
the second archive had again made navigation and proof auditing
unreliable.  It does not restart the mathematics.  It gives a compact
map of proved results, states the exact archive interfaces used below,
and continues at the final f2 bottleneck: a missing directed response
at an endpoint of a matching-breaking covariance.

\begin{record}[Exact archive identities]\label{rec:V1-archives}
The f1 archive has \(88\,517\) source lines, \(3\,083\,467\) bytes, and
SHA--256
\[
 \texttt{82f6bbffeecdad3c5c63551660b19a0bdc7f3d6dffe4c68b4e9bba3a783a2b71}.
\]
The f2 archive has \(49\,581\) source lines, \(1\,706\,928\) bytes, and
SHA--256
\[
 \texttt{6444a1044e74dfcb1a72ed9c7248bc0f234efc4a0e4248c736cd7db212dfca9a}.
\]
Every archive label quoted in typewriter font refers to these exact
files.  Such a quotation is a proof index, not a local cross-reference.
\end{record}

\begin{record}[Version and companion-audit protocol]\label{rec:V1-protocol}
The next note is obtained by copying the complete active file, appending
one new chapter, validating the result, and retaining only that newest
active version.  Frozen archives are retained.  At every active version
divisible by five, the newest Standard-Model--Einstein and
Navier--Stokes notes in \texttt{latex\_notes} are inspected read-only.
At active V100 the sequence is frozen with suffix \(\texttt{\_f3}\),
and another contextual V1 is started.
\end{record}

\section{Proof route and logical order}\label{sec:V1-route}

The intended implication chain is
\begin{equation}
\begin{split}
\text{exact Wilson cumulants}
&\Longrightarrow \text{marked-tree atlas (MTA)}
\Longrightarrow \text{weighted uniform fusion control (WUFC)}\\
&\Longrightarrow \text{connected source/Perron bounds (CPM)}
\Longrightarrow \text{uniform finite-edge estimates}\\
&\Longrightarrow \text{constructive cutoff removal and reflection positivity}\\
&\Longrightarrow \text{Osterwalder--Schrader reconstruction}
\Longrightarrow \text{positive physical spectral gap}.
\end{split}
\label{eq:V1-route}
\end{equation}
The programme is still inside the first implication.

\section{Major mathematical results obtained so far}\label{sec:V1-summary}

\subsection{Harmonic normalization, probability marks, and topology}

For irreducible inputs \(R_1,\ldots,R_m\), f1 V46 proves
\[
 \widehat h(U)=
 \frac{\prod_i d_{R_i}}{d_U}
 \Tr_{M_U}\!\left[
 W_U^*\Bigl(\bigotimes_i A_i\Bigr)W_U
 (I_{H_U}\otimes K_U)\right],
\]
together with complete isotypic pinching
\[
 \sum_U\trnorm{W_U^*(\bigotimes_i A_i)W_U}
 \le \prod_i\trnorm{A_i}.
\]
Thus the output dimension cancels and multiplicity copies are not
summed.  F1 V31 proves binary MTA using the exact probability marks
\(\vartheta_{x,e}=a_{x,e}/\Xi_x\), and the ternary row is closed after
the V46 correction.

Every minimal connected four-row coordinate skeleton has topology
\[
 [4],\qquad[3,2],\qquad[3,3],\qquad[2,2,2].
\]
F1 V41 controls all sixteen labelled four-vertex trees.  The
exclusive-pair \([2,2,2]\) sector is closed in f1 V53, and the
noncollision \([3,2]\) sector in f1 V56.

\subsection{Collision algebra and the numerator boundary}

F1 V57--V129 keeps the noncommuting collision orientations separate,
constructs positive product-state capacities, and closes each fixed
completely regrouped collision monomial.  V130 gives an exact
Boolean--Möbius selector and V131 an exact sixteen-tree representation.
V133 disproves the V132 product-bridge interpolation: several nominal
bridges can collapse to one quotient covariance.

Returning upstream, f1 V134--V135 proves
\[
 \opnorm{\mathbf K^{\mathrm{conn}}_{\{1,2,3,4\}}}
 \le C\sum_{\tau\in\mathfrak T_4}
 \prod_{ij\in E(\tau)}s_\alpha H_{ij}.
\]
This is an unnormalized numerator; taking its norm has already
forgotten the repeated probability marks and final resolvent quotient.

\subsection{First-connectivity and the one-payer calculus}

The f2 archive develops the normalized source before norming it:
\begin{enumerate}[leftmargin=3em]
\item V61 proves exact first-connectivity factorization and the
      linked-word criterion for arbitrary proper row partitions.
\item V65--V78 proves complete composite two-block covariance minima
      \(C\min\{1,s_\alpha H\}\) unpaid and
      \(C\min\{H,s_\alpha^{-1}\}\) paid.
\item V81 and V96 allocate the sole final resolvent once, including an
      arbitrarily supported common suffix without coordinate loss.
\item V83--V85 proves a higher-incidence square tail, gives an explicit
      \(SU(2)\) counterexample to the false crossing square tail, and
      isolates the sharp output-channel obstruction.
\item V92--V94 proves literal tree marking, payer-compatible
      normalization, and exact Fourier dualization.
\item V95--V98 proves quadratic \(A/F/I\) endpoint energies and
      distinct- and self-adjacent clean-overlap closure.
\item V99 proves bounded directed words, two-directed/one-paired and
      all-directed compilers, and classifies disconnected
      \(K_{2,2}\) edge sets.
\end{enumerate}

\subsection{The f2 terminal matching theorem}

F2 V100 returns to the two disconnected perfect matchings.  A connected
word refining the associated two-block partition has a unique first
component breaker.  Its complete bank has bounds
\[
 C\min\{1,s_\alpha H_G(A,B)\},\qquad
 C\min\{H_G(A,B),s_\alpha^{-1}\},
\]
at every payer location.  Fully incident matching packets close into
literal trees.  The terminal proof index is
\[
\begin{gathered}
\texttt{lem:V100-no-breaker-zero},\quad
\texttt{thm:V100-breaker-norm},\quad
\texttt{thm:V100-first-breaker-packet},\\
\texttt{thm:V100-matching-closure},\quad
\texttt{prop:V100-missing-tail-reduction}.
\end{gathered}
\]
The last result identifies a missing breaker endpoint but does not
terminate its depleted acquisition.

\section{Current status}\label{sec:V1-status}

\begin{longtable}{>{\raggedright\arraybackslash}p{0.30\textwidth}
                  >{\raggedright\arraybackslash}p{0.14\textwidth}
                  >{\raggedright\arraybackslash}p{0.48\textwidth}}
\toprule
gate & status & exact boundary\\
\midrule
\endfirsthead
\toprule
gate & status & exact boundary\\
\midrule
\endhead
Binary and ternary normalized MTA
& \MGclosed
& Exact Fourier normalization, marks, and one-resolvent estimates are
  proved.\\[0.35em]
Exclusive-pair quartic \([2,2,2]\)
& \MGclosed
& Complete arbitrary-coefficient sector proved in f1 V53.\\[0.35em]
Noncollision quartic \([3,2]\)
& \MGclosed
& Complete assembled sector proved in f1 V56.\\[0.35em]
Clean collision \(2+2\) source
& \MGpartial
& Fixed regrouped monomials, clean overlaps, directed trees, and fully
  incident matching packets are controlled; depleted endpoints and
  nonclean states remain.\\[0.35em]
\([4]\), \([3,3]\), \(3+1\), \(2+1+1\), and discrete sources
& \MGopen
& No complete coefficient-sensitive normalized source theorem covers
  these families.\\[0.35em]
Global quartic and higher MTA
& \MGopen
& Remaining source families and uniform higher-order induction are not
  proved.\\[0.35em]
WUFC, CPM, continuum construction, and physical gap
& \MGopen
& These stages are downstream and have not been reached.\\
\bottomrule
\end{longtable}

\begin{firewall}[Status at the restart]\label{fire:V1-initial-boundary}
Neither archive proves the Yang--Mills existence and mass-gap theorem.
No continuum measure, Osterwalder--Schrader Hilbert space, Hamiltonian,
or uniform positive lower bound on its nonvacuum spectrum has yet been
constructed.  The full Yang--Mills mass gap remains open.
\end{firewall}

\chapter{The archive interface used in the new calculation}
\label{ch:V1-interface}

\section{Four-row response notation}\label{sec:V1-notation}

Let
\[
 \cV=\{p,p^\circ,r,r^\circ\}.
\]
For coordinate \(e\) and row \(x\in\cV\), let \(a_{x,e}\ge0\) be its
input Casimir weight, with \(a_{x,e}=0\) when \(x\) is inactive.  Put
\begin{equation}
 \Xi_x:=\sum_e a_{x,e},
 \qquad
 H_{xy}(G):=\sum_{e\in G}a_{x,e}a_{y,e}.
\label{eq:V1-stocks}
\end{equation}
The clean compact-group calculus has a fixed nontrivial floor
\begin{equation}
 a_{x,e}>0\quad\Longrightarrow\quad a_{x,e}\ge a_*>0.
\label{eq:V1-floor}
\end{equation}
Constants may depend on the fixed group, four-row order, and fixed
decomposition thresholds, but not on support, representations, or
physical multiplicities.

For a bank \(B\), define the oriented response
\begin{equation}
 d_{x\to u}(B):=\frac{H_{xu}(B)}{\Xi_x},\qquad u\ne x.
\label{eq:V1-directed-response}
\end{equation}
It is \(\gamma\)-dominant when
\(d_{x\to u}(B)\ge\gamma a_*\).  F2 V98 proves this whenever a
\(\gamma\)-fraction of row \(x\)'s mass lies in a clean directed cell
whose actual partner is \(u\).

\section{Matching breakers and the inherited packet estimate}

Fix a perfect-matching partition
\[
 \mu=A\mid B,\qquad |A|=|B|=2,
\]
and write \(m(x)\) for the mate of \(x\).  For the complete
first-breaker bank \(G\), put
\begin{equation}
 H_G(A,B)=
 \sum_{\substack{x\in A\\y\in B}}H_{xy}(G).
\label{eq:V1-breaker-stock}
\end{equation}
The symbol \(\Gamma_\otimes(\varepsilon,K)\) is the f2 positive
product-state majorant of the assembled source, including the output
dualizer.  Here \(\varepsilon=0\) is unpaid and \(\varepsilon=1\)
means that the sole final payer has been assigned.

\begin{proposition}[Inherited complete-breaker interface]
\label{prop:V1-inherited-breaker}
There is a complete assembled packet \(K_{\mu,G}^{(\varepsilon)}\)
such that, at every admissible payer location,
\begin{equation}
 \Gamma_\otimes(\varepsilon,K_{\mu,G}^{(\varepsilon)})
 \le C H_G(A,B),
 \qquad \varepsilon\in\{0,1\}.
\label{eq:V1-inherited-breaker}
\end{equation}
\end{proposition}

\begin{proof}
For \(\varepsilon=0\), f2
\(\texttt{thm:V100-first-breaker-packet}\) gives
\[
 C\min\{1,s_\alpha H_G\}\le Cs_0H_G
\]
on \(0<s_\alpha\le s_0\).  For \(\varepsilon=1\), it gives
\[
 C\min\{H_G,s_\alpha^{-1}\}\le CH_G.
\]
That theorem includes payers in the matching word, breaker, finite
separator, and complete suffix, and assembles the operator before
norming it.
\end{proof}

\begin{lemma}[Exact pair allocation of the positive response ledger]
\label{lem:V1-pair-allocation}
For each \(\varepsilon\), there are nonnegative pair-ledger entries
\(Q_{xy}^{(\varepsilon)}\), \(x\in A,y\in B\), such that
\begin{align}
 \Gamma_\otimes(\varepsilon,K_{\mu,G}^{(\varepsilon)})
 &\le\sum_{\substack{x\in A\\y\in B}}Q_{xy}^{(\varepsilon)},
\label{eq:V1-pair-majorant}\\
 Q_{xy}^{(\varepsilon)}&\le C H_{xy}(G).
\label{eq:V1-pair-entry}
\end{align}
The allocation occurs after forming the complete covariance and has no
coordinate-cardinality factor.
\end{lemma}

\begin{proof}
The V100 product telescope first selects \(e\in G\).  Its complete
local two-block covariance is bounded by
\[
 C\sum_{x\in A,y\in B}a_{x,e}a_{y,e}
\]
in the paid response, and by the same expression times \(s_\alpha\)
in the unpaid response.  Retain the finite pair label in this positive
majorant and sum \(e\) afterward.  The \((x,y)\) entry is bounded by
\(CH_{xy}(G)\), and summing the four labels proves
\eqref{eq:V1-pair-majorant}.
\end{proof}

\begin{remark}[Ledger, not operator decomposition]
\label{rem:V1-ledger-not-operator}
The entries \(Q_{xy}^{(\varepsilon)}\) are labelled summands of a
proved positive majorant.  We do not assert that the signed
multicoordinate covariance decomposes into operators localized by
\(H_{xy}\).  That stronger claim is false in the f2 V78 and V84
regressions.
\end{remark}

\section{Literal path currency}

\begin{definition}[Two-internal-vertex path response]
\label{def:V1-path-response}
For distinct breaker endpoints \(x,y\) and partners \(u\ne x\),
\(v\ne y\), define
\begin{equation}
 {\cal T}(u-x-y-v)
 :=d_{x\to u}(B_x)d_{y\to v}(B_y)H_{xy}(G).
\label{eq:V1-path-response}
\end{equation}
\end{definition}

\begin{lemma}[Algebra of the literal path mark]
\label{lem:V1-path-algebra}
If \(u,x,y,v\) are distinct, then
\begin{equation}
 {\cal T}(u-x-y-v)=
 \frac{H_{xu}(B_x)H_{xy}(G)H_{yv}(B_y)}
 {\Xi_x\Xi_y}.
\label{eq:V1-path-algebra}
\end{equation}
This is the f2 V92 aggregate marked-tree response for the literal path
\(u-x-y-v\).
\end{lemma}

\begin{proof}
Substitute \eqref{eq:V1-directed-response} twice.  The denominators are
attached to the two degree-two path vertices.  F2
\(\texttt{prop:V98-same-bank-marking}\) permits coincident coordinate
banks only after their common occurrence has been bounded in the
assembled word, as required in
\Cref{prop:V1-inherited-breaker}.
\end{proof}

\chapter{Pair-localized matching closure and the exact partner table}
\label{ch:V1-partner-table}

\section{The complementary-partner compiler}

\begin{theorem}[Complementary partners close one breaker entry]
\label{thm:V1-complementary-partners}
Fix \(x\in A\), \(y\in B\), and write
\[
 \cV\setminus\{x,y\}=\{z,w\}.
\]
Suppose the assembled word contains dominant responses
\begin{equation}
 d_{x\to u}(B_x)\ge\gamma_xa_*,
 \qquad
 d_{y\to v}(B_y)\ge\gamma_ya_*,
\label{eq:V1-two-endpoint-lower}
\end{equation}
and
\begin{equation}
 \{u,v\}=\{z,w\}.
\label{eq:V1-complementary-partners}
\end{equation}
Then, for \(\varepsilon\in\{0,1\}\),
\begin{equation}
 \boxed{
 Q_{xy}^{(\varepsilon)}
 \le
 \frac{C}{\gamma_x\gamma_ya_*^2}
 d_{x\to u}(B_x)d_{y\to v}(B_y)H_{xy}(G).}
\label{eq:V1-complementary-bound}
\end{equation}
The right side is a legal literal marked-tree response, uniformly in
support and at every payer location admitted by the V100 packet.
\end{theorem}

\begin{proof}
Equation \eqref{eq:V1-two-endpoint-lower} gives
\[
 1\le
 \frac{d_{x\to u}(B_x)d_{y\to v}(B_y)}
 {\gamma_x\gamma_ya_*^2}.
\]
Multiply \eqref{eq:V1-pair-entry} by this inequality.  Condition
\eqref{eq:V1-complementary-partners} says that \(u,x,y,v\) are four
distinct rows.  Lemma~\ref{lem:V1-path-algebra} identifies the product
with the literal path \(u-x-y-v\).  The payer was already allocated
inside \(Q_{xy}^{(\varepsilon)}\), so no second resolvent is introduced.
\end{proof}

\begin{corollary}[The f2 matching-tail theorem is a special case]
\label{cor:V1-V100-special}
If \(u=m(x)\) and \(v=m(y)\), then
\eqref{eq:V1-complementary-partners} holds.  Theorem
\ref{thm:V1-complementary-partners} also permits the crossed choice
\(u=m(y)\), \(v=m(x)\).
\end{corollary}

\begin{proof}
The mates of \(x,y\), which lie in different matching components, are
precisely the two rows outside \(\{x,y\}\).  Interchanging them
preserves their set.
\end{proof}

\section{The complete nine-cell graph table}

\begin{proposition}[Exact endpoint-partner table]
\label{prop:V1-partner-table}
Keep \(\cV\setminus\{x,y\}=\{z,w\}\).  For
\(u\in\{y,z,w\}\), \(v\in\{x,z,w\}\), the graph
\({\cal G}(u,v)=\{xy,xu,yv\}\) has status
\[
\begin{array}{c|ccc}
 &v=x&v=z&v=w\\
\hline
u=y&\text{misses }z,w&\text{misses }w&\text{misses }z\\
u=z&\text{misses }w&\text{misses }w&\text{path}\\
u=w&\text{misses }z&\text{path}&\text{misses }z
\end{array}
\tag{V1.1}\label{eq:V1-partner-table}
\]
Thus the two complementary-partner cells are the only cells giving a
spanning tree with \(x,y\) as its internal vertices.  Every failed
cell names a complementary row absent from the edge set.
\end{proposition}

\begin{proof}
The edge \(xy\) is always present, so its component also contains
\(u,v\).  It spans \(\cV\) iff
\(\{u,v\}\supseteq\{z,w\}\), equivalently
\(\{u,v\}=\{z,w\}\).  This gives the two path cells.

For \((u,v)=(y,x)\), both response edges repeat \(xy\), so \(z,w\)
are absent.  For \(u=v=z\), the graph is the triangle on
\(\{x,y,z\}\), leaving \(w\) absent; \(u=v=w\) is symmetric.  The
remaining four cells contain exactly one complementary row.
\end{proof}

\begin{firewall}[Scalar size cannot alter the graph]
\label{fire:V1-no-relabel}
A large response with partner \(z\) cannot be relabelled as a response
with partner \(w\).  A failed cell of
\eqref{eq:V1-partner-table} is not a marked spanning tree.  This is
the graph-versus-size distinction tested by f2 V83 and V91.
\end{firewall}

\section{Exact localization of the remaining breaker stock}

\begin{definition}[Rows with acquired matching tails]
\label{def:V1-covered-rows}
Let \(C\subseteq\cV\) be the rows having a dominant matching-tail
response \(d_{x\to m(x)}\).  Define
\begin{align}
 H_{\mathrm{good}}
 &:=
 \sum_{\substack{x\in A\cap C\\y\in B\cap C}}H_{xy}(G),
\label{eq:V1-good-stock}\\
 H_{A,\mathrm{miss}}
 &:=
 \sum_{\substack{x\in A\setminus C\\y\in B}}H_{xy}(G),
\label{eq:V1-A-miss}\\
 H_{B,\mathrm{miss}}
 &:=
 \sum_{\substack{x\in A\cap C\\y\in B\setminus C}}H_{xy}(G).
\label{eq:V1-B-miss}
\end{align}
\end{definition}

\begin{lemma}[Disjoint missing-endpoint partition]
\label{lem:V1-missing-partition}
One has
\begin{equation}
 H_G(A,B)=
 H_{\mathrm{good}}+H_{A,\mathrm{miss}}+H_{B,\mathrm{miss}}.
\label{eq:V1-missing-partition}
\end{equation}
Every pair in \(H_{\mathrm{good}}\) closes.  Each other pair is
assigned once to an explicit missing endpoint: first to its
\(A\)-endpoint if missing, otherwise to its missing \(B\)-endpoint.
\end{lemma}

\begin{proof}
Partition \(A\times B\) into the disjoint sets
\[
 (A\cap C)\times(B\cap C),\quad
 (A\setminus C)\times B,\quad
 (A\cap C)\times(B\setminus C).
\]
This proves \eqref{eq:V1-missing-partition}.  In the first set, the
matching mates of \(x,y\) are the two complementary rows, so
\Cref{cor:V1-V100-special} applies.  The order of the other two sets
gives a unique missing-endpoint label.
\end{proof}

\begin{assessment}[Progress over f2 V100]
\label{ass:V1-localization-progress}
The f2 endpoint statement is now an exact nonoverlapping response
ledger.  Moreover, a matching tail is sufficient but not necessary:
crossed complementary partners close the same pair.  The unresolved
directed case is one of seven explicit cells in
\eqref{eq:V1-partner-table}, not an arbitrary four-row graph.
\end{assessment}

\chapter{One-step depleted acquisition and the self-anchored residual}
\label{ch:V1-acquisition}

\section{Anchored used banks}

Fix a missing endpoint \(x\) selected by
\Cref{lem:V1-missing-partition}.  Let
\[
 \cU=\bigcup_{j=1}^m\cB_j
\]
be the union of complete coordinate banks already present in the
assembled word.  The construction records an anchor \(p_j\in\cV\)
for each \(\cB_j\): every coordinate in \(\cB_j\) is active at
\(p_j\), and the bank entered through a certificate rooted there.
The exact four-row source has a universal finite number
\begin{equation}
 m\le M_0
\label{eq:V1-bank-count}
\end{equation}
of such roles.  A physical bank in several roles remains one joint
operator occurrence and is not estimated independently.

Put
\[
 X_x(\cB_j):=\sum_{e\in\cB_j}a_{x,e},
 \qquad
 X_x(\cU):=\sum_{e\in\cU}a_{x,e},
\]
using disjoint bank representatives after common coordinates are
assigned to their earliest role.  Split
\[
 I_{\mathrm{self}}(x)=\{j:p_j=x\},
 \qquad
 I_{\mathrm{ext}}(x)=\{j:p_j\ne x\}.
\]

\begin{lemma}[Anchored-overlap dichotomy]
\label{lem:V1-anchored-overlap}
Assume
\begin{equation}
 X_x(\cU)\ge\delta\Xi_x,\qquad0<\delta<1.
\label{eq:V1-overlap}
\end{equation}
Then either some \(j\in I_{\mathrm{ext}}(x)\) satisfies
\begin{equation}
 d_{x\to p_j}(\cB_j)
 \ge\frac{\delta}{2M_0}a_*,
\label{eq:V1-external-response}
\end{equation}
or the self-anchored banks satisfy
\begin{equation}
 \sum_{j\in I_{\mathrm{self}}(x)}X_x(\cB_j)
 \ge\frac{\delta}{2}\Xi_x.
\label{eq:V1-self-concentration}
\end{equation}
In the first alternative, \(xp_j\) is an actual common-coordinate
edge and \(p_j\ne x\).
\end{lemma}

\begin{proof}
If \eqref{eq:V1-self-concentration} fails, subtract its left side from
\eqref{eq:V1-overlap}:
\[
 \sum_{j\in I_{\mathrm{ext}}(x)}X_x(\cB_j)
 \ge\frac{\delta}{2}\Xi_x.
\]
One of at most \(M_0\) terms is at least
\(\delta\Xi_x/(2M_0)\).  Every coordinate of that bank is active at
\(p_j\), so \eqref{eq:V1-floor} gives
\[
 H_{xp_j}(\cB_j)\ge a_*X_x(\cB_j).
\]
Divide by \(\Xi_x\).  Since \(j\) is external, \(p_j\ne x\).
\end{proof}

\begin{corollary}[One self-anchored bank witnesses the residual]
\label{cor:V1-self-witness}
Under \eqref{eq:V1-self-concentration}, some
\(j\in I_{\mathrm{self}}(x)\) satisfies
\begin{equation}
 X_x(\cB_j)\ge\frac{\delta}{2M_0}\Xi_x.
\label{eq:V1-self-witness}
\end{equation}
\end{corollary}

\begin{proof}
Apply the pigeonhole principle to at most \(M_0\) nonnegative terms.
\end{proof}

\section{Fresh cells versus anchored overlap}

The f2 V84 depleted identity partitions unused mass at \(x\) into
\[
 P_x^{\rm fr},\ U_x^{\rm fr},\ A_x^{\rm fr},\
 F_x^{\rm fr},\ D_x^{\rm fr},\ I_x^{\rm fr}.
\]
For the current used union,
\begin{equation}
 \Xi_x=X_x(\cU)+
 P_x^{\rm fr}+U_x^{\rm fr}+A_x^{\rm fr}+
 F_x^{\rm fr}+D_x^{\rm fr}+I_x^{\rm fr}.
\label{eq:V1-depleted-identity}
\end{equation}

\begin{theorem}[One-step clean acquisition normal form]
\label{thm:V1-one-step-acquisition}
Fix \(0<\delta<1\).  At a clean missing endpoint \(x\), one
application of \eqref{eq:V1-depleted-identity} yields:
\begin{enumerate}[label=\textup{(\alph*)},leftmargin=4em]
\item a fresh \(P\)- or \(U\)-cell of size at least
      \((1-\delta)\Xi_x/6\), entering its complete-moment compiler;
\item a fresh \(A,F,\) or \(I\)-cell of that size, entering the f2
      quadratic endpoint and overlap compilers;
\item a fresh directed response
      \begin{equation}
      d_{x\to q}(\cD)\ge\frac{1-\delta}{18}a_*
      \label{eq:V1-fresh-directed}
      \end{equation}
      for an actual partner \(q\ne x\);
\item the external-anchor response
      \eqref{eq:V1-external-response}; or
\item the one-bank self-anchored witness
      \eqref{eq:V1-self-witness}.
\end{enumerate}
Outside the clean calculus there is separately a nonclean
representation state.  Thus an overlap does not restart an untyped
depleted search: it gives an actual directed edge or one
self-anchored bank.
\end{theorem}

\begin{proof}
If \(X_x(\cU)\ge\delta\Xi_x\), apply
\Cref{lem:V1-anchored-overlap,cor:V1-self-witness}.
Otherwise the six fresh cells in
\eqref{eq:V1-depleted-identity} total more than
\((1-\delta)\Xi_x\), so one has at least one sixth of that mass.
The \(P,U,A,F,I\) routes are the archived compilers.

For a \(D\)-cell, partition by its actual partner
\(q\in\cV\setminus\{x\}\).  Some one of at most three partner cells,
called \(\cD\), carries at least \((1-\delta)\Xi_x/18\).
Every coordinate there is active at \(q\), so
\[
 H_{xq}(\cD)\ge a_*\sum_{e\in\cD}a_{x,e}.
\]
Divide by \(\Xi_x\).  Freshness makes \(\cD\) disjoint from \(\cU\).
\end{proof}

\begin{firewall}[Acquisition is not yet graph closure]
\label{fire:V1-acquisition-not-closure}
A directed edge retains its actual partner and must pass the graph
table.  A self-anchored witness is a repeated complete-bank
occurrence; its two appearances may not be charged by two
first-moment bounds.
\end{firewall}

\section{The finite residual after one-step acquisition}

\begin{theorem}[Pairwise residual normal form]
\label{thm:V1-pairwise-residual}
Fix one positive pair entry \(Q_{xy}^{(\varepsilon)}\).  Apply
\Cref{thm:V1-one-step-acquisition} at each breaker endpoint whose
response is missing.  Within the clean calculus, the pair has one of:
\begin{enumerate}[label=\textup{(\roman*)},leftmargin=4em]
\item a \(P,U,A,F,I\) complete compiler closes its branch;
\item both endpoints have dominant directed responses with
      complementary partners, and
      \Cref{thm:V1-complementary-partners} closes the pair;
\item both endpoints have dominant responses, but their partner cell
      is one of the seven failed entries of
      \eqref{eq:V1-partner-table}, naming a missing row;
\item one endpoint has the one-bank self-anchored witness
      \eqref{eq:V1-self-witness}; or
\item the state is nonclean.
\end{enumerate}
There are four breaker-pair labels and nine partner cells per label.
Thus the directed fibre has \(36\) cells: \(8\) close immediately and
\(28\) have an explicitly named missing row.  This factor is
independent of coordinate support.
\end{theorem}

\begin{proof}
The pair label is fixed first by
\Cref{lem:V1-pair-allocation}.  The endpoint alternatives are
exhaustive by \Cref{thm:V1-one-step-acquisition}.  The established
compilers give \textup{(i)}.  For two directed responses,
\Cref{prop:V1-partner-table} gives exactly \textup{(ii)} or
\textup{(iii)}.  The remaining alternatives are \textup{(iv)--(v)}.
There are \(|A||B|=4\) pair labels; each has nine cells, two paths and
seven failures.  Hence the counts are \(36,8,28\).
\end{proof}

\begin{proposition}[Same-type self-anchored routes]
\label{prop:V1-self-type-routes}
Let \(\cB\) be selected by \Cref{cor:V1-self-witness}.
\begin{enumerate}[label=\textup{(\alph*)},leftmargin=4em]
\item If its two appearances preserve one clean \(A,F,\) or \(I\)
      type, the f2 V97--V98 self-adjacent theorem applies.
\item If they preserve one clean directed \(D\)-type with partner
      \(q\ne x\), the repeated covariance remains in the bounded
      directed word and supplies the literal edge \(xq\) to the table.
\item Private or unbalanced appearances remain inside their archived
      complete-moment compilers.
\end{enumerate}
The new analytic self-anchored residual is therefore a
partition-changing occurrence whose two roles admit no common clean
cell type before the complete bank is assembled.
\end{proposition}

\begin{proof}
Item \textup{(a)} is the scope of
\(\texttt{thm:V98-self-adjacent-sandwich}\) and
\(\texttt{thm:V98-self-adjacent-paid}\).  In \textup{(b)},
\(\texttt{lem:V99-directed-word-norm}\) bounds the complete word before
the response \(H_{xq}(\cB)/\Xi_x\) is inserted; no second first moment
is used.  Item \textup{(c)} is the \(P,U\) incidence route.  When no
common type exists, none of these cited theorems applies, which is the
stated residual.
\end{proof}

\begin{assessment}[The bottleneck has moved upstream]
\label{ass:V1-new-bottleneck}
An externally anchored overlap is a dominant literal edge, while a
self-anchored overlap has one fixed bank witness.  The open object is
the exact V61 partition-state operator containing two different roles
of that bank, together with the \(28\) finite partner-collision cells.
It must be estimated jointly before a positive envelope or separate
first-moment cards are applied.
\end{assessment}

\chapter{Regression checks, result ledger, and V2 mandate}

\begin{proposition}[Archive regressions remain intact]
\label{prop:V1-regressions}
The new calculation satisfies:
\begin{enumerate}[label=\textup{(\roman*)},leftmargin=4em]
\item f2 V15 paths retain actual marks and coincident banks are
      assembled before marking;
\item the V78 Rademacher example is not assigned an operator localized
      by a nonexistent \(H_{ij}\);
\item V83 scalar mass comparisons cannot pass a failed graph cell;
\item the V84 spin-one example receives no false crossing square tail;
\item the V86 binary prefix remains inside the complete word;
\item the V91 path \(\{12,13,34\}\) is never relabelled; and
\item the V100 breaker is summed as \(H_G(A,B)\), not as a coordinate
      count times a cap.
\end{enumerate}
\end{proposition}

\begin{proof}
Items \textup{(i), (v), (vi)} follow from the assembled convention in
\Cref{lem:V1-path-algebra}.  Item \textup{(ii)} is
\Cref{rem:V1-ledger-not-operator}.  Items \textup{(iii)--(iv)} follow
from \Cref{fire:V1-no-relabel} and the absence of any new square-tail
claim.  Item \textup{(vii)} is the exact coordinate sum in
\Cref{lem:V1-pair-allocation}.
\end{proof}

\begin{theorem}[Third-series V1 result ledger]\label{thm:V1-results}
Relative to the frozen archives, this note proves:
\begin{enumerate}[label=\textup{(V1.\arabic*)},leftmargin=5.8em]
\item a standalone archive and claim map;
\item exact four-entry pair allocation of the matching-breaker
      positive response ledger;
\item the complementary-partner compiler
      \eqref{eq:V1-complementary-bound};
\item the exact nine-cell partner table
      \eqref{eq:V1-partner-table};
\item the disjoint breaker-stock partition
      \eqref{eq:V1-missing-partition};
\item the anchored-overlap dichotomy
      \eqref{eq:V1-external-response}--%
      \eqref{eq:V1-self-concentration};
\item the one-step clean acquisition normal form; and
\item reduction to one self-anchored partition-changing bank, finite
      partner-collision cells, or a nonclean state.
\end{enumerate}
\end{theorem}

\begin{proof}
These are
\Cref{rec:V1-archives,
lem:V1-pair-allocation,
thm:V1-complementary-partners,
prop:V1-partner-table,
lem:V1-missing-partition,
lem:V1-anchored-overlap,
thm:V1-one-step-acquisition,
thm:V1-pairwise-residual,
prop:V1-self-type-routes}.
\end{proof}

\begin{programme}[V2: partition-changing self-bank packet]
\label{prog:V1-V2}
\begin{enumerate}[leftmargin=3em]
\item Return to the exact V61 source word for the self-anchored witness
      and write its two partition roles in one common/replicated
      covariance before taking an envelope.
\item Split only by the finite pair label and seven failed partner
      cells; do not split by coordinate or multiplicity copy.
\item Test whether the role difference factors through one additional
      row difference.  If so, retain its actual edge and use
      \Cref{thm:V1-complementary-partners}; otherwise construct the
      smallest \(SU(2)\) counterpacket.
\item If a counterpacket exists, identify whether its obstruction is
      low-output heat, a repeated-bank second moment, or a nonclean
      block, and move upstream to that exact object.
\item After clean matching closure, return to the \(3+1\),
      \(2+1+1\), and discrete first-connectivity families.
\item The next compulsory SM/NS companion audit is active V5.
\end{enumerate}
\end{programme}

\begin{firewall}[Claim boundary after third-series V1]
\label{fire:V1-final-boundary}
This note does not prove the joint partition-changing self-bank
estimate, the \(28\) partner-collision cells, nonclean states, other
quartic source families, higher MTA, WUFC, CPM, cutoff removal,
reflection positivity, Osterwalder--Schrader reconstruction, or a
uniform positive continuum spectral gap.  The full Yang--Mills mass
gap remains open.
\end{firewall}

\paragraph{Restart record.}
This V1 was written after exact V100 was frozen as f2 with the identity
in \Cref{rec:V1-archives}.  It contains no mathematical copy of the
\(49\,581\)-line parent.  The next active V2 appends its chapter to
this complete source.

\clearpage
\chapter{V2: Partition-cover covariances and the anchor-sensitive
role correction}
\label{ch:V2-role-correction}

\section{Why the self-bank role change is revisited}

Third-series V1 reduced the clean matching residual to a finite
partner table and one self-anchored bank whose two source roles need
not have a common \(P,U,A,F,D,I\) type.  The f2 archive already proves
exact first-connectivity, two-block bank telescopes, and
contraction-sandwiched overlap.  What it does not state is the finite
algebra relating two arbitrary replica partitions of the same bank.

This chapter proves that algebra.  A cover in the partition lattice is
exactly one complete composite covariance.  Any two four-row roles
differ by at most three such cover covariances.  Consequently a
partition-changing self-bank is not a new untyped analytic object: it
is a finite signed sum of already admissible two-block covariances.
The remaining issue is whether their genuine bridge stocks touch the
missing anchor row in a dominant way.

\begin{assessment}[Nonduplication boundary]\label{ass:V2-nonduplication}
F2 \(\texttt{thm:V61-first-connectivity}\) splits a connected word at
its first full join.  F2 \(\texttt{thm:V77-bank-telescope}\) treats a
fixed two-block cut against the one-block partition.  V2 proves a
different statement: a saturated-cover telescope between arbitrary,
possibly incomparable, replica partitions, followed by an
anchor-sensitive decomposition.  No local heat estimate is reproved.
\end{assessment}

\section{Partition moments of a complete bank}

Let \(I\subseteq\cV\), and let \(\Pi(I)\) be its partition lattice,
ordered by refinement.  For \(\rho\in\Pi(I)\), recall the exact local
replica moment
\begin{equation}
 M_e(\rho):=\sum_{\pi\le\rho}L_{e,\pi}.
\label{eq:V2-local-partition-moment}
\end{equation}
For a finite coordinate bank \(G\), define
\begin{equation}
 M_G(\rho):=\bigotimes_{e\in G}M_e(\rho).
\label{eq:V2-bank-partition-moment}
\end{equation}
All multiplicity spaces remain inside the complete letters.

\begin{lemma}[Exact block factorization]\label{lem:V2-block-factorization}
For every partition \(\rho\),
\begin{equation}
 M_G(\rho)=\bigotimes_{C\in\rho}M_G(C),
\label{eq:V2-block-factorization}
\end{equation}
where \(M_G(C)\) is the complete common-variable moment of the rows in
\(C\), and different blocks use independent replicas.  In particular,
\(\opnorm{M_G(\rho)}\le1\).
\end{lemma}

\begin{proof}
At one coordinate, a partition letter refining \(\rho\) is uniquely a
choice of a partition on each block \(C\in\rho\).  Hence the finite
moment--cumulant sum in \eqref{eq:V2-local-partition-moment} factors
over those blocks.  Tensoring the identity over \(e\in G\) gives
\eqref{eq:V2-block-factorization}.  Each factor is an expectation of
a product of unitary representation matrices and is a contraction.
\end{proof}

Write \(\rho\lessdot\tau\) when \(\tau\) covers \(\rho\).  Then
\(\tau\) is obtained by merging exactly two blocks \(C,D\in\rho\).
Put
\begin{equation}
 {\mathfrak J}_G^{C|D}
 :=
 M_G(C\cup D)-M_G(C)\otimes M_G(D).
\label{eq:V2-cover-covariance}
\end{equation}

\begin{theorem}[A partition cover is one complete covariance]
\label{thm:V2-cover-covariance}
If \(\rho\lessdot\tau\) merges \(C,D\), then
\begin{equation}
 \boxed{
 M_G(\tau)-M_G(\rho)
 =
 {\mathfrak J}_G^{C|D}
 \otimes
 \bigotimes_{E\in\rho\setminus\{C,D\}}M_G(E).}
\label{eq:V2-cover-factorization}
\end{equation}
With
\begin{equation}
 H_G(C,D):=
 \sum_{e\in G}\sum_{\substack{u\in C\\v\in D}}
 a_{u,e}a_{v,e},
\label{eq:V2-cover-stock}
\end{equation}
the unpaid cover occurrence satisfies
\begin{equation}
 \opnorm{M_G(\tau)-M_G(\rho)}
 \le C\min\{1,s_\alpha H_G(C,D)\}.
\label{eq:V2-cover-unpaid}
\end{equation}
If the sole payer lies in this occurrence, its complete paid envelope
satisfies
\begin{equation}
 \opnorm{(M_G(\tau)-M_G(\rho))^{\rm pay}}
 \le C\min\{H_G(C,D),s_\alpha^{-1}\}.
\label{eq:V2-cover-paid}
\end{equation}
The same bounds hold for product-state capacity.
\end{theorem}

\begin{proof}
Apply \Cref{lem:V2-block-factorization} to \(\rho\) and \(\tau\).
Every block other than \(C,D\) is identical.  Factoring those
spectators gives \eqref{eq:V2-cover-factorization}.

The bracket is precisely the complete composite two-block covariance
of f2 V77, with all coordinates assembled before norming.  Its
ordinary-norm response proof gives
\(C\min\{1,s_\alpha H_G(C,D)\}\); f2 V78 gives the paid minimum
\(C\min\{H_G(C,D),s_\alpha^{-1}\}\).  The remaining complete moments
are contractions.  Product-state capacity is no larger than ordinary
norm.
\end{proof}

\begin{firewall}[A cover stock is a sum of actual bridges]
\label{fire:V2-cover-bridges}
Equation \eqref{eq:V2-cover-stock} is formed inside the complete
covariance estimate.  Only afterward may it be expanded into the
finitely many actual row pairs \(uv\).  No nominal edge is substituted
for those pairs and no uniform coordinate cap is counted over \(G\).
\end{firewall}

\section{Arbitrary and incomparable role changes}

\begin{theorem}[Three-cover role-difference normal form]
\label{thm:V2-three-cover}
For any \(\rho,\sigma\in\Pi(I)\), there exist signs
\(\epsilon_j\in\{-1,1\}\) and cover pairs
\(\alpha_j\lessdot\beta_j\), \(1\le j\le N\), such that
\begin{equation}
 \boxed{
 M_G(\sigma)-M_G(\rho)
 =
 \sum_{j=1}^{N}\epsilon_j
 \bigl(M_G(\beta_j)-M_G(\alpha_j)\bigr),\qquad N\le3.}
\label{eq:V2-three-cover}
\end{equation}
Every summand has the exact covariance form
\eqref{eq:V2-cover-factorization}.  Two transverse two-pair matchings
need only two covers, one from each matching to their one-block join.
\end{theorem}

\begin{proof}
Let \(\eta=\rho\vee\sigma\).  Choose saturated chains
\[
 \rho=\rho_0\lessdot\rho_1\lessdot\cdots\lessdot\rho_a=\eta,
 \qquad
 \sigma=\sigma_0\lessdot\sigma_1\lessdot\cdots\lessdot\sigma_b=\eta.
\]
Telescoping gives
\[
 M_G(\sigma)-M_G(\rho)
 =
 \sum_{j=1}^{a}\bigl(M_G(\rho_j)-M_G(\rho_{j-1})\bigr)
 -
 \sum_{j=1}^{b}\bigl(M_G(\sigma_j)-M_G(\sigma_{j-1})\bigr).
\]
Here
\[
 a+b=|\rho|+|\sigma|-2|\eta|,
\]
where \(|\lambda|\) is the number of blocks.  If the partitions are
comparable, the length is at most \(3\).  Suppose they are
incomparable.  If \(|\eta|\ge2\), the displayed length is at most
\(2\): two distinct three-block partitions are each one merge below
their join, and all other cases are shorter.  If \(|\eta|=1\), two
three-block partitions cannot occur, because each has only one
nontrivial pair and two pair edges cannot connect four vertices.
Thus \(|\rho|+|\sigma|\le5\), and again \(a+b\le3\).
\end{proof}

\begin{corollary}[Support-uniform analytic role bound]
\label{cor:V2-role-bound}
For the cover stocks \(H_j=H_G(C_j,D_j)\) in any normal form
\eqref{eq:V2-three-cover},
\begin{align}
 \opnorm{M_G(\sigma)-M_G(\rho)}
 &\le C\sum_{j=1}^{N}\min\{1,s_\alpha H_j\},
\label{eq:V2-role-unpaid}\\
 \opnorm{(M_G(\sigma)-M_G(\rho))^{\rm pay}}
 &\le C\sum_{j=1}^{N}\min\{H_j,s_\alpha^{-1}\}.
\label{eq:V2-role-paid}
\end{align}
The sums have at most three terms, independently of \(|G|\).
\end{corollary}

\begin{proof}
Apply the triangle inequality to \eqref{eq:V2-three-cover} and then
\Cref{thm:V2-cover-covariance}.  The payer alternatives are allocated
before the finite cover sum, so only one paid occurrence appears in
each source alternative.
\end{proof}

\section{Insertion into a separated self-bank word}

\begin{lemma}[Exact diagonal-plus-role-correction identity]
\label{lem:V2-diagonal-role}
For every bounded separator \(C\),
\begin{equation}
 \boxed{
 M_G(\rho)C M_G(\sigma)
 =
 M_G(\rho)C M_G(\rho)
 +
 M_G(\rho)C\bigl(M_G(\sigma)-M_G(\rho)\bigr).}
\label{eq:V2-diagonal-role}
\end{equation}
Thus the part caused specifically by changing the right replica role
is the second summand; it contains the three-cover normal form.
\end{lemma}

\begin{proof}
Add and subtract \(M_G(\rho)CM_G(\rho)\).  No commutation or trace
cyclicity is used.
\end{proof}

\begin{theorem}[Support-uniform separated role correction]
\label{thm:V2-separated-role}
Let \(A\) be a complete source word with \(\opnorm A\le L_A\), let
\(\opnorm C\le L_C\), and define
\[
 {\cal R}_{\rho,\sigma}:=
 AC\bigl(M_G(\sigma)-M_G(\rho)\bigr).
\]
For any three-cover normal form with stocks \(H_j\),
\begin{equation}
 \opnorm{{\cal R}_{\rho,\sigma}}
 \le
 C L_AL_C\sum_{j=1}^{N}\min\{1,s_\alpha H_j\}.
\label{eq:V2-separated-unpaid}
\end{equation}
If the sole payer lies in the role difference, use
\eqref{eq:V2-role-paid}.  If it lies in \(A\) or \(C\), assume the
corresponding complete paid replacement has norm at most
\(L_{\rm pay}/s_\alpha\).  In all three cases,
\begin{equation}
 \opnorm{{\cal R}_{\rho,\sigma}^{\rm pay}}
 \le
 C L_*\sum_{j=1}^{N}
 \min\{H_j,s_\alpha^{-1}\},
\label{eq:V2-separated-paid}
\end{equation}
where \(L_*\) is a fixed product of the unpaid and paid word constants.
A complete suffix payer has the same bound after the f2 V96 grouped
allocation.
\end{theorem}

\begin{proof}
Insert \eqref{eq:V2-three-cover} and use submultiplicativity together
with \eqref{eq:V2-cover-unpaid}.  This proves
\eqref{eq:V2-separated-unpaid}.

For a payer inside the difference, apply
\eqref{eq:V2-cover-paid}.  For a payer in either bounded companion
factor, multiply the unpaid cover bound by \(C/s_\alpha\) and use
\[
 \frac1{s_\alpha}\min\{1,s_\alpha H_j\}
 =
 \min\{s_\alpha^{-1},H_j\}.
\]
F2 V95 supplies the finite local-separator paid norm and f2 V96 the
grouped suffix norm.  The V2 scalar allocation makes these mutually
exclusive cases, so no product of two paid factors occurs.
\end{proof}

\begin{firewall}[The diagonal term remains a separate source type]
\label{fire:V2-diagonal-separate}
Lemma~\ref{lem:V2-diagonal-role} does not declare
\(M_G(\rho)CM_G(\rho)\) small.  When it has a common clean
\(A,F,I,D,P,\) or \(U\) type, the corresponding f2 theorem applies.
Otherwise it remains a diagonal source problem.  Theorem
\ref{thm:V2-separated-role} proves that changing the partition role
adds no new support-dependent analytic obstruction.
\end{firewall}

\section{When a role change touches the anchor}

For a partition \(\rho\), write \(B_\rho(x)\) for its block containing
\(x\).

\begin{lemma}[An unchanged anchor block factors off exactly]
\label{lem:V2-unchanged-anchor}
Suppose
\[
 B_\rho(x)=B_\sigma(x)=X.
\]
Writing \(\rho',\sigma'\) for the induced partitions of \(I\setminus X\),
\begin{equation}
 \boxed{
 M_G(\sigma)-M_G(\rho)
 =
 M_G(X)\otimes
 \bigl(M_G(\sigma')-M_G(\rho')\bigr).}
\label{eq:V2-anchor-factor}
\end{equation}
The changing factor contains at most three rows and is therefore in
the already proved binary/ternary analytic layer.
\end{lemma}

\begin{proof}
The common block \(X\) occurs in both partitions.  Apply
\eqref{eq:V2-block-factorization} to factor it from both bank moments
and subtract.  Because \(x\in X\), its complement has at most three
rows.  The common factor is a contraction; payer allocation between
the two factors is the inherited lower-order one-payer calculus.
\end{proof}

Now assume \(B_\rho(x)\ne B_\sigma(x)\), and define the changed
partner set
\begin{equation}
 P_x(\rho,\sigma):=
 \bigl(B_\rho(x)\triangle B_\sigma(x)\bigr)\setminus\{x\}.
\label{eq:V2-changed-partners}
\end{equation}
It is nonempty and has at most three elements.  A coordinate
\(e\in G\) is \(x\)-sensitive if \(a_{x,e}>0\) and at least one
\(y\in P_x(\rho,\sigma)\) is active there.  Let \(S_x\) be the
sensitive subbank, \(F_x=G\setminus S_x\), and
\begin{equation}
 X_x^{\rm sens}:=\sum_{e\in S_x}a_{x,e}.
\label{eq:V2-sensitive-mass}
\end{equation}

\begin{lemma}[Changed partners occur in actual cover stocks]
\label{lem:V2-changed-partner-cover}
For every \(y\in P_x(\rho,\sigma)\), some cover in the normal form
\eqref{eq:V2-three-cover} merges a block containing \(x\) with a block
containing \(y\).  Its stock therefore contains \(H_{xy}(G)\).
\end{lemma}

\begin{proof}
Assume first that \(y\in B_\rho(x)\setminus B_\sigma(x)\).
Along the saturated chain from \(\sigma\) to
\(\eta=\rho\vee\sigma\), there is a first cover after which \(x,y\)
lie in one block.  That cover merges their current blocks, and
\eqref{eq:V2-cover-stock} contains \(H_{xy}(G)\).
If \(y\in B_\sigma(x)\setminus B_\rho(x)\), use the chain from
\(\rho\) to \(\eta\).  These are the two signed halves of
\eqref{eq:V2-three-cover}.
\end{proof}

\begin{theorem}[Sensitive-anchor dichotomy]
\label{thm:V2-sensitive-dichotomy}
Suppose the self-anchored witness bank satisfies
\begin{equation}
 X_x(G):=\sum_{e\in G}a_{x,e}\ge\beta\Xi_x.
\label{eq:V2-self-witness}
\end{equation}
Then exactly one of the following threshold alternatives may be
selected:
\begin{enumerate}[label=\textup{(\roman*)},leftmargin=4em]
\item if \(X_x^{\rm sens}\ge\beta\Xi_x/2\), some changed partner
      \(y\in P_x(\rho,\sigma)\) obeys
      \begin{equation}
      \boxed{
      d_{x\to y}(G)=\frac{H_{xy}(G)}{\Xi_x}
      \ge\frac{\beta}{6}a_*;}
      \label{eq:V2-sensitive-response}
      \end{equation}
      the response occurs in an actual cover covariance from
      \eqref{eq:V2-three-cover};
\item if \(X_x^{\rm sens}<\beta\Xi_x/2\), then
      \begin{equation}
      \sum_{e\in F_x}a_{x,e}\ge\frac{\beta}{2}\Xi_x,
      \label{eq:V2-insensitive-large}
      \end{equation}
      and the large part of the witness is locally insensitive to the
      change of \(x\)'s replica block.
\end{enumerate}
\end{theorem}

\begin{proof}
In the first case, every \(e\in S_x\) has an active changed partner.
By \eqref{eq:V1-floor},
\[
 \sum_{y\in P_x(\rho,\sigma)}H_{xy}(G)
 \ge a_*X_x^{\rm sens}.
\]
There are at most three partners, so one \(y\) gives
\[
 H_{xy}(G)\ge\frac{a_*}{3}X_x^{\rm sens}
 \ge\frac{\beta a_*}{6}\Xi_x.
\]
Apply \Cref{lem:V2-changed-partner-cover}.  If the first threshold
fails, subtract \(X_x^{\rm sens}<\beta\Xi_x/2\) from
\eqref{eq:V2-self-witness} to obtain
\eqref{eq:V2-insensitive-large}.
\end{proof}

\begin{lemma}[Exact sensitive/insensitive bank split]
\label{lem:V2-sensitive-split}
Let
\(\Delta_G^{\rho,\sigma}:=M_G(\sigma)-M_G(\rho)\).  Then
\begin{equation}
 \boxed{
 \Delta_G^{\rho,\sigma}
 =
 \Delta_{S_x}^{\rho,\sigma}\otimes M_{F_x}(\sigma)
 +
 M_{S_x}(\rho)\otimes\Delta_{F_x}^{\rho,\sigma}.}
\label{eq:V2-sensitive-split}
\end{equation}
Moreover, there are complete contractions \(A_{F_x}^x\) and
role moments \(R_{F_x}(\rho),R_{F_x}(\sigma)\), whose changing row
sets omit \(x\), such that
\begin{equation}
 \boxed{
 \Delta_{F_x}^{\rho,\sigma}
 =
 A_{F_x}^x\otimes
 \bigl(R_{F_x}(\sigma)-R_{F_x}(\rho)\bigr).}
\label{eq:V2-insensitive-factor}
\end{equation}
\end{lemma}

\begin{proof}
Coordinate factorization gives
\[
 M_G(\lambda)=M_{S_x}(\lambda)\otimes M_{F_x}(\lambda),
 \qquad\lambda\in\{\rho,\sigma\}.
\]
Insert and subtract
\(M_{S_x}(\rho)\otimes M_{F_x}(\sigma)\) to obtain
\eqref{eq:V2-sensitive-split}.

Fix \(e\in F_x\).  If \(x\) is active, no active row belongs to
\(P_x(\rho,\sigma)\).  Hence the block containing \(x\) in the
restrictions of \(\rho\) and \(\sigma\) to the active rows is the same;
call it \(X_e\).  If \(x\) is inactive, take \(X_e=\varnothing\).
By local block factorization,
\[
 M_e(\lambda)=M_e(X_e)\otimes R_e(\lambda),
 \qquad\lambda\in\{\rho,\sigma\},
\]
where the changing factor \(R_e(\lambda)\) contains no row \(x\).
Tensor over \(e\in F_x\), set
\[
 A_{F_x}^x:=\bigotimes_{e\in F_x}M_e(X_e),
 \qquad
 R_{F_x}(\lambda):=\bigotimes_{e\in F_x}R_e(\lambda),
\]
and subtract.  This proves \eqref{eq:V2-insensitive-factor}.
\end{proof}

\begin{corollary}[Response-first anchor dichotomy]
\label{cor:V2-response-first}
Define
\begin{equation}
 r_x(\rho,\sigma;G)
 :=
 \frac1{\Xi_x}
 \sum_{y\in P_x(\rho,\sigma)}H_{xy}(G).
\label{eq:V2-total-changed-response}
\end{equation}
Under \eqref{eq:V2-self-witness}, either some changed partner obeys
\[
 d_{x\to y}(G)\ge\frac{\beta}{6}a_*,
\]
or
\begin{equation}
 r_x(\rho,\sigma;G)<\frac{\beta}{2}a_*,
 \qquad
 X_x^{\rm sens}<\frac{\beta}{2}\Xi_x,
 \qquad
 \sum_{e\in F_x}a_{x,e}\ge\frac{\beta}{2}\Xi_x.
\label{eq:V2-low-response-inert}
\end{equation}
In the second case, the large part of the self-witness occurs only in
the common contraction \(A_{F_x}^x\) of
\eqref{eq:V2-insensitive-factor}; its changing factor omits \(x\).
\end{corollary}

\begin{proof}
If \(r_x\ge\beta a_*/2\), one of at most three summands in
\eqref{eq:V2-total-changed-response} is at least
\(\beta a_*/6\).  Otherwise the floor argument gives
\[
 a_*X_x^{\rm sens}
 \le\sum_{y\in P_x}H_{xy}(G)
 <\frac{\beta a_*}{2}\Xi_x.
\]
Cancel \(a_*\), use \eqref{eq:V2-self-witness}, and apply
\Cref{lem:V2-sensitive-split}.
\end{proof}

\begin{theorem}[Self-anchored role-change reduction]
\label{thm:V2-self-role-reduction}
Let \(G\) be the one-bank self-anchored witness left by V1, with
\(X_x(G)\ge\beta\Xi_x\), and let its replica roles be \(\rho,\sigma\).
\begin{enumerate}[label=\textup{(\roman*)},leftmargin=4em]
\item If \(B_\rho(x)=B_\sigma(x)\), the role difference factors
      through at most three changing rows by
      \eqref{eq:V2-anchor-factor}.
\item If the anchor blocks differ and the first alternative of
      \Cref{cor:V2-response-first} holds, an actual cover covariance
      supplies the dominant edge \(x\to y\), which re-enters the
      finite V1 partner table.
\item Otherwise \eqref{eq:V2-sensitive-split} consists of one
      correction on the \(x\)-sensitive subbank with the strict bound
      \eqref{eq:V2-low-response-inert}, plus an \(x\)-free role
      difference multiplied by a common contraction.
\end{enumerate}
All role corrections have the support-uniform unpaid and once-paid
estimates of \Cref{thm:V2-separated-role}.  Thus the
\emph{partition-changing self-bank} is no longer an untyped analytic
carrier.
\end{theorem}

\begin{proof}
Item \textup{(i)} is \Cref{lem:V2-unchanged-anchor}.  When the blocks
differ, apply \Cref{cor:V2-response-first}.  In its first alternative,
\Cref{lem:V2-changed-partner-cover} places the actual edge in a cover
stock.  In its second alternative,
\Cref{lem:V2-sensitive-split} and
\eqref{eq:V2-low-response-inert} give \textup{(iii)}.
The uniform estimates follow from
\Cref{thm:V2-separated-role} for at most three cover increments.
\end{proof}

\begin{corollary}[Immediate complementary-cell closure]
\label{cor:V2-complementary-reentry}
Fix \(Q_{xy}^{(\varepsilon)}\).  Suppose the self-role reduction at
\(x\) yields \(x\to z\), where
\(\{z,w\}=\cV\setminus\{x,y\}\), and the other endpoint already has
\(y\to w\).  Then
\[
 Q_{xy}^{(\varepsilon)}
 \le C\,d_{x\to z}(G)d_{y\to w}(B_y)H_{xy}(G),
\]
with fixed dominance constants suppressed, and the right side is the
legal path \(z-x-y-w\).  The crossed choice is symmetric.
\end{corollary}

\begin{proof}
Apply \Cref{thm:V1-complementary-partners} with the response supplied
by \Cref{thm:V2-self-role-reduction}.
\end{proof}

\begin{assessment}[Boundary after the upstream move]
\label{ass:V2-new-boundary}
The self-overlap is now one of a lower-row role change, a literal
changed-partner response, or an \(x\)-sensitive correction whose total
changed-partner response is below \(\beta a_*/2\).  Only the last can
still interact with a failed partner cell without immediately
returning an edge at \(x\).  Its cover covariance may instead carry a
bridge between other rows in the block containing \(x\); that bridge
must be routed in the complete packet and may not be renamed.
\end{assessment}

\section{Regression tests and result ledger}

\begin{proposition}[Required regressions]\label{prop:V2-regressions}
\begin{enumerate}[label=\textup{(\roman*)},leftmargin=4em]
\item For the transverse matchings \(12|34\) and \(13|24\), the join
      route has exactly two cover covariances, not four independent
      events.
\item The V78 Rademacher test is respected because
      \eqref{eq:V2-cover-stock} is expanded only after the complete
      covariance estimate; no signed operator is localized by one
      chosen \(H_{uv}\).
\item The V84 spin-one test is respected because the cover estimate is
      the proved composite minimum, not the false crossing square tail.
\item The V91 coherent-tilt carrier is not claimed to cancel between
      partition roles.  V2 supplies an exact finite decomposition and
      retains every bridge stock.
\item The V100 breaker stock remains inside its assembled packet.  A
      response from \eqref{eq:V2-sensitive-response} is inserted only
      through a lower bound and with its actual partner.
\end{enumerate}
\end{proposition}

\begin{proof}
For \textup{(i)}, both matchings are covered by their one-block join,
so \eqref{eq:V2-three-cover} has two terms.  Items
\textup{(ii)--(iii)} follow from
\Cref{fire:V2-cover-bridges,thm:V2-cover-covariance}.  Item
\textup{(iv)} follows because all identities through
\eqref{eq:V2-sensitive-split} precede absolute values and assert no
extra cancellation.  Item \textup{(v)} is
\Cref{lem:V2-changed-partner-cover,cor:V2-complementary-reentry}.
\end{proof}

\begin{theorem}[Third-series V2 result ledger]\label{thm:V2-results}
Relative to V1 and the frozen archives, V2 proves:
\begin{enumerate}[label=\textup{(V2.\arabic*)},leftmargin=5.8em]
\item every partition cover of a complete bank is exactly one
      composite covariance with the unpaid and paid minima
      \eqref{eq:V2-cover-unpaid}--\eqref{eq:V2-cover-paid};
\item arbitrary four-row replica roles differ by at most three cover
      covariances, with no support-dependent fibre;
\item a separated partition-role correction remains support-uniform at
      every sole-payer location;
\item an unchanged anchor block factors from the role difference,
      leaving at most three changing rows;
\item every changed anchor partner occurs in an actual cover stock;
\item a large self-anchored witness either supplies the dominant
      response \eqref{eq:V2-sensitive-response}, or its large part is
      an exact common contraction while the changing factor omits the
      anchor; and
\item complementary responses obtained this way close immediately by
      the V1 path compiler.
\end{enumerate}
\end{theorem}

\begin{proof}
These are
\Cref{thm:V2-cover-covariance,
thm:V2-three-cover,
thm:V2-separated-role,
lem:V2-unchanged-anchor,
lem:V2-changed-partner-cover,
cor:V2-response-first,
thm:V2-self-role-reduction,
cor:V2-complementary-reentry}.
\end{proof}

\begin{programme}[V3: route the low-response sensitive cover]
\label{prog:V2-V3}
\begin{enumerate}[leftmargin=3em]
\item For every anchor-touching cover \(C|D\), split its exact stock
      into the anchor edges
      \[
      \sum_{y\in D}H_{xy}(G)
      \]
      and the shadow bridges
      \[
      \sum_{u\in C\setminus\{x\},\,y\in D}H_{uy}(G).
      \]
      Keep both inside the same complete covariance majorant.
\item Insert each literal bridge into the seven failed partner cells
      of \eqref{eq:V1-partner-table}.  Record which bridge completes a
      path, which produces a star requiring a repeated denominator,
      and which leaves a row absent.
\item In the low-response branch
      \eqref{eq:V2-low-response-inert}, use the cover response as an
      upper currency rather than attempting another lower-response
      insertion.  Determine whether the other endpoint and prefix
      stocks supply the two internal-vertex denominators.
\item Preserve the three-cover signs until the complete V61 source is
      assembled.  If the shadow-bridge terms cannot be routed with a
      bounded fibre, build the smallest physical \(SU(2)\)
      counterpacket and move upstream to its exact partition sum.
\item Treat the diagonal term in
      \eqref{eq:V2-diagonal-role} only by its actual clean type.  Do
      not infer its smallness from the role correction.
\item Recheck the V15, V78, V83, V84, V91, and V100 examples.  The
      next companion audit remains due at active V5.
\end{enumerate}
\end{programme}

\begin{firewall}[Claim boundary after third-series V2]
\label{fire:V2-boundary}
V2 removes the untyped partition-role obstruction but does not close
the low-response sensitive cover, all \(28\) failed partner cells,
the remaining diagonal or nonclean states, the \(3+1\), \(2+1+1\),
and discrete source families, higher MTA, WUFC, CPM, cutoff removal,
reflection positivity, Osterwalder--Schrader reconstruction, or a
uniform positive continuum spectral gap.  The full Yang--Mills mass
gap remains open.
\end{firewall}

\paragraph{Parent-version record.}
V2 was copied byte-for-byte from
\texttt{Renewal\_Yang\_Mills\_Mass\_Gap\_Research\_Notes\_V1.tex}
before this chapter was appended.  The V1 parent had \(901\) source
lines, \(32\,639\) bytes, and SHA--256
\[
 \texttt{a255a0a40967398dc1d5f34bb50ccdda87867d6cb857e9b6dd2e7b5a2267cc68}.
\]
No companion audit is due at V2; the next required audit is V5.

\clearpage
\chapter{V3: Star-denominator closure of the partner-collision cells}
\label{ch:V3-star}

\section{Context and nonduplication}

V1 used the V92 normalization only for a path: the two response
denominators occur at its two degree-two vertices.  V2 showed that a
partition-role change can produce an additional dominant response at
an actual changed partner.  In six of the seven failed V1 partner
cells, that extra response does not replace the other endpoint
response to form a path.  Instead, it shares a tail with an existing
response and forms a three-spoke star.

The star has one vertex of degree three, so its V92 normalization has
the square of one row mass in the denominator.  Two dominant
responses with the same tail supply precisely that square.  This
chapter proves the resulting compiler and the complete augmented
partner table.

\begin{assessment}[Nonduplication boundary]\label{ass:V3-nonduplication}
F2 V99 treats three directed responses with distinct tails as an
arborescence.  V1 treats two distinct endpoint tails plus one raw
breaker as a path.  V3 treats the complementary case: two response
denominators have the same tail and the raw breaker is the third spoke
of a star.  None of those earlier theorems contains this denominator
pattern.
\end{assessment}

\section{The exact denominator-degree criterion}

For a labelled tree \(\tau\) on \(\cV\), with edge banks \(G_{ab}\),
write its aggregate V92 response as
\begin{equation}
 {\cal W}_\tau^{\rm agg}
 :=
 \left(\prod_{v\in\cV}\Xi_v\right)
 \prod_{ab\in E(\tau)}
 \frac{H_{ab}(G_{ab})}{\Xi_a\Xi_b}.
\label{eq:V3-tree-aggregate}
\end{equation}

\begin{lemma}[Degree form of the tree normalization]
\label{lem:V3-degree-form}
For every quartic tree,
\begin{equation}
 \boxed{
 {\cal W}_\tau^{\rm agg}
 =
 \frac{\displaystyle\prod_{ab\in E(\tau)}H_{ab}(G_{ab})}
 {\displaystyle\prod_{v\in\cV}\Xi_v^{\,\deg_\tau(v)-1}}.}
\label{eq:V3-degree-form}
\end{equation}
The total denominator degree is two.
\end{lemma}

\begin{proof}
In \eqref{eq:V3-tree-aggregate}, row \(v\) occurs once in the
numerator factor \(\prod_v\Xi_v\) and once in the denominator for each
incident edge.  Its net denominator exponent is
\(\deg_\tau(v)-1\).  Since a four-vertex tree has three edges,
\[
 \sum_v(\deg_\tau(v)-1)=2|E(\tau)|-|\cV|=6-4=2.
\]
\end{proof}

\begin{proposition}[Two-response denominator classification]
\label{prop:V3-two-response-classification}
Suppose a product of three edge stocks is written as one raw stock and
two oriented responses:
\begin{equation}
 H_{e_0}\,
 \frac{H_{e_1}}{\Xi_{t_1}}\,
 \frac{H_{e_2}}{\Xi_{t_2}},
\label{eq:V3-two-response-product}
\end{equation}
where \(t_i\) is an endpoint of \(e_i\).  Its denominator pattern is
the exact aggregate normalization of the tree
\(\tau=\{e_0,e_1,e_2\}\) precisely in either of the following graph
types:
\begin{enumerate}[label=\textup{(\roman*)},leftmargin=4em]
\item \(\tau\) is a path and \(t_1,t_2\) are its two distinct
      degree-two vertices;
\item \(\tau\) is a star and \(t_1=t_2\) is its degree-three center.
\end{enumerate}
\end{proposition}

\begin{proof}
By \Cref{lem:V3-degree-form}, the formal denominator multiset is
\[
 \{v\text{ repeated }\deg_\tau(v)-1\text{ times}:v\in\cV\}.
\]
A four-vertex tree has degree sequence either \((2,2,1,1)\) or
\((3,1,1,1)\).  In the first case the multiset consists of the two
internal vertices; in the second it consists of two copies of the
center.  These are exactly \textup{(i)--(ii)}.
\end{proof}

\begin{firewall}[Connectivity alone is insufficient]
\label{fire:V3-degree-not-connectivity}
Three edge labels can form a tree while
\eqref{eq:V3-two-response-product} has denominators at its leaves.
Such a product is not the V92 tree response.  V3 checks both the edge
set and the denominator multiset before declaring closure.
\end{firewall}

\section{The breaker-centered star compiler}

Fix a breaker pair \(x,y\), and write
\[
 \cV\setminus\{x,y\}=\{z,w\}.
\]

\begin{theorem}[Two same-tail responses close a breaker star]
\label{thm:V3-star-compiler}
Suppose the assembled matching-breaker word contains
\begin{equation}
 d_{x\to z}(B_z)\ge\gamma_za_*,
 \qquad
 d_{x\to w}(B_w)\ge\gamma_wa_*.
\label{eq:V3-two-spokes}
\end{equation}
Then, for every unpaid or once-paid pair entry,
\begin{equation}
 \boxed{
 Q_{xy}^{(\varepsilon)}
 \le
 \frac{C}{\gamma_z\gamma_wa_*^2}
 d_{x\to z}(B_z)d_{x\to w}(B_w)H_{xy}(G),
 \qquad\varepsilon\in\{0,1\}.}
\label{eq:V3-star-bound}
\end{equation}
The right side is the exact aggregate response of the star with
center \(x\) and spokes \(xy,xz,xw\).  The symmetric statement holds
with center \(y\).
\end{theorem}

\begin{proof}
The two lower bounds in \eqref{eq:V3-two-spokes} give
\[
 1\le
 \frac{d_{x\to z}(B_z)d_{x\to w}(B_w)}
 {\gamma_z\gamma_wa_*^2}.
\]
Multiply the pair majorant
\(Q_{xy}^{(\varepsilon)}\le CH_{xy}(G)\) from
\eqref{eq:V1-pair-entry}.  The resulting scalar is
\[
 \frac{H_{xz}(B_z)H_{xw}(B_w)H_{xy}(G)}{\Xi_x^2},
\]
which is \eqref{eq:V3-degree-form} for the displayed star.  The pair
entry already contains the sole payer, so no new resolvent is used.
If banks coincide, their physical occurrences remain in the one
assembled word before the two scalar lower bounds are inserted, as in
f2 \(\texttt{prop:V98-same-bank-marking}\).
\end{proof}

\begin{remark}[Path and star are the complete two-response atlas]
\label{rem:V3-complete-atlas}
For one raw breaker stock and two response denominators, V1's path
compiler and \Cref{thm:V3-star-compiler} exhaust the two possible
quartic tree degree sequences.  No third normalization pattern exists.
\end{remark}

\section{The augmented seven-cell table}

\begin{theorem}[One extra endpoint response closes six failed cells]
\label{thm:V3-augmented-table}
In the V1 notation, let \(u\) be the existing partner of the response
rooted at \(x\), and \(v\) the existing partner of the response rooted
at \(y\).  Each failed cell of
\eqref{eq:V1-partner-table} has the following exact one-response
repair:
\[
\begin{array}{c|c|c}
(u,v)&\text{one additional response}&\text{tree}\\
\hline
(y,z)&y\to w&\text{star centered at }y\\
(y,w)&y\to z&\text{star centered at }y\\
(z,x)&x\to w&\text{star centered at }x\\
(z,z)&x\to w\ \text{or}\ y\to w&\text{star}\\
(w,x)&x\to z&\text{star centered at }x\\
(w,w)&x\to z\ \text{or}\ y\to z&\text{star}
\end{array}
\tag{V3.1}\label{eq:V3-augmented-table}
\]
The remaining cell \((u,v)=(y,x)\) cannot be made into a correctly
normalized tree by adding only one endpoint response.
\end{theorem}

\begin{proof}
In the first row, the existing response \(y\to z\), the new response
\(y\to w\), and the raw breaker \(xy\) satisfy
\Cref{thm:V3-star-compiler} with center \(y\).  The next five rows are
the same argument with the center and missing complementary spoke
shown in the table.

In the cell \((y,x)\), both existing response edges repeat the raw
breaker \(xy\).  Adding one response can introduce at most one of the
vertices \(z,w\); the other remains absent.  Hence no spanning tree,
regardless of denominator placement, can result.
\end{proof}

\begin{corollary}[V2 changed-partner responses enter the star table]
\label{cor:V3-role-star}
Whenever \Cref{thm:V2-self-role-reduction} supplies the additional
response listed in a row of \eqref{eq:V3-augmented-table}, the
corresponding pair entry closes at every payer location by
\Cref{thm:V3-star-compiler}.  The response is used with its actual
changed partner; no role-cover bridge is relabelled.
\end{corollary}

\begin{proof}
The V2 response has a fixed positive dominance fraction and occurs in
an actual cover stock.  Combine it with the existing same-tail
response specified in the table and apply
\Cref{thm:V3-star-compiler}.
\end{proof}

\section{The complete endpoint-certificate criterion}

For the raw breaker \(xy\), let \({\cal R}_x\) and \({\cal R}_y\)
be the sets of actual partners for which dominant responses rooted at
\(x\) and \(y\), respectively, are present in the assembled word.
Repeated occurrences of one partner do not change these finite sets.

\begin{theorem}[Seeing both complementary rows is necessary and
sufficient]
\label{thm:V3-certificate-complete}
With \(\cV\setminus\{x,y\}=\{z,w\}\), the pair entry
\(Q_{xy}^{(\varepsilon)}\) closes by a V1 path or a V3 star if and
only if
\begin{equation}
 \boxed{
 \bigl({\cal R}_x\cup{\cal R}_y\bigr)\cap\{z,w\}
 =
 \{z,w\}.}
\label{eq:V3-both-complements}
\end{equation}
Equivalently, the complete directed endpoint residual has no hidden
graph types: all its complementary responses point to at most one of
\(z,w\).
\end{theorem}

\begin{proof}
Assume \eqref{eq:V3-both-complements}.  If one of
\({\cal R}_x,{\cal R}_y\) contains both \(z,w\), apply
\Cref{thm:V3-star-compiler} at that root.  Otherwise neither set
contains both.  Because their union contains both, after possibly
interchanging \(z,w\), one has
\[
 z\in{\cal R}_x,\qquad w\in{\cal R}_y.
\]
The two responses and the raw breaker form the path \(z-x-y-w\), so
\Cref{thm:V1-complementary-partners} applies.

Conversely, a star containing the raw spoke \(xy\) must contain the
two other rows as its other spokes, so both \(z,w\) occur as endpoint
partners.  A path normalized by responses rooted at \(x,y\) has those
responses pointing to the two complementary rows by
\Cref{prop:V3-two-response-classification}.  Thus either compiler
implies \eqref{eq:V3-both-complements}.
\end{proof}

\begin{corollary}[The 28 cells collapse to a missing-complement
certificate]
\label{cor:V3-collapse-28}
After all dominant responses already present at the breaker endpoints
are collected, every unclosed one of the \(28\) failed V1
partner-labelled cells has a row
\[
 q\in\{z,w\}
\]
which is not the partner of any dominant response rooted at \(x\) or
\(y\).  If some complementary response exists, \(q\) is unique.  If
none exists, both complementary rows are missing.
\end{corollary}

\begin{proof}
This is the negation of \eqref{eq:V3-both-complements}.  A proper
subset of the two-element set \(\{z,w\}\) is empty or a singleton.
\end{proof}

\begin{corollary}[Exact double-duplicate repair criterion]
\label{cor:V3-double-duplicate}
In the double-duplicate cell
\((u,v)=(y,x)\), let
\[
 {\cal S}\subseteq\{x\to z,x\to w,y\to z,y\to w\}
\]
be the additional dominant endpoint responses.  The cell closes if
and only if \({\cal S}\) contains one of
\[
 \{x\to z,x\to w\},\quad
 \{y\to z,y\to w\},\quad
 \{x\to z,y\to w\},\quad
 \{x\to w,y\to z\}.
\label{eq:V3-double-repairs}
\]
Every maximal nonclosing set is
\(\{x\to z,y\to z\}\) or
\(\{x\to w,y\to w\}\).
\end{corollary}

\begin{proof}
The four displayed pairs are, respectively, the star at \(x\), the
star at \(y\), and the two complementary paths.  By
\Cref{thm:V3-certificate-complete}, closure fails exactly when the
additional responses see at most one complementary row.  The two
two-element maximal sets with that property are those stated.
\end{proof}

\begin{assessment}[New endpoint bottleneck]
\label{ass:V3-new-bottleneck}
The old list of \(28\) failed partner cells was a list of ordered
response choices, not \(28\) genuinely different obstructions.  Once
all responses already in the assembled word are retained, path and
star normalization reduce them to one typed event: a particular
complementary row is not reached from either breaker endpoint.  V2
closes this event whenever its changed-partner cover supplies a
dominant response to that row.  The remaining branch is the
low-response role correction of
\eqref{eq:V2-low-response-inert}, or a depleted acquisition which
continues to avoid the named complementary row.
\end{assessment}

\section{Regression checks and result ledger}

\begin{proposition}[Path/star regression suite]
\label{prop:V3-regressions}
\begin{enumerate}[label=\textup{(\roman*)},leftmargin=4em]
\item In the V15 repeated-coordinate path, the response tails remain
      the two degree-two vertices; V3 does not move them.
\item In the V83 same-side-maximal pattern, two responses with a common
      tail close only if their literal edges are the two distinct
      complementary spokes of the same star.
\item The V84 spin-one family receives no new local square tail; V3
      inserts only already proved dominant lower responses.
\item The V91 path \(\{12,13,34\}\) has denominator vertices \(1,3\)
      and is the path case of
      \Cref{prop:V3-two-response-classification}.
\item A V100 breaker edge is never used simultaneously as two
      different spokes: duplicate response labels are ignored unless
      another literal complementary edge is present.
\end{enumerate}
\end{proposition}

\begin{proof}
Items \textup{(i), (iv)} follow from
\Cref{lem:V3-degree-form}.  Item \textup{(ii)} is the literal-edge
hypothesis of \Cref{thm:V3-star-compiler}.  Item \textup{(iii)}
follows because no analytic upper estimate was changed.  Item
\textup{(v)} follows from the graph proof of
\Cref{thm:V3-augmented-table}.
\end{proof}

\begin{theorem}[Third-series V3 result ledger]\label{thm:V3-results}
Relative to V1--V2 and the frozen archives, V3 proves:
\begin{enumerate}[label=\textup{(V3.\arabic*)},leftmargin=5.8em]
\item the exact denominator-degree form
      \eqref{eq:V3-degree-form} of every quartic tree response;
\item the complete classification of a raw edge plus two response
      denominators into a path or a star;
\item the support-uniform unpaid and once-paid star compiler
      \eqref{eq:V3-star-bound};
\item one-extra-response closure for six of the seven failed cells in
      the V1 table;
\item the necessary-and-sufficient endpoint certificate
      \eqref{eq:V3-both-complements}; and
\item reduction of all \(28\) ordered failed cells to one named
      missing-complement event, with the exact double-duplicate repair
      list \eqref{eq:V3-double-repairs}.
\end{enumerate}
\end{theorem}

\begin{proof}
These are
\Cref{lem:V3-degree-form,
prop:V3-two-response-classification,
thm:V3-star-compiler,
thm:V3-augmented-table,
thm:V3-certificate-complete,
cor:V3-collapse-28,
cor:V3-double-duplicate}.
\end{proof}

\begin{programme}[V4: acquire the named missing complement]
\label{prog:V3-V4}
\begin{enumerate}[leftmargin=3em]
\item Fix a breaker entry \(xy\) and the complementary row \(q\)
      identified by \Cref{cor:V3-collapse-28}.  Return to the two
      exact depleted incidence identities at \(x\) and \(y\), with all
      banks already present in the assembled word removed.
\item Keep incidences whose actual partner is \(q\) separate from
      incidences pointing to \(y\), \(x\), or the other complement.
      Determine whether their combined absence forces a
      \(P,U,A,F,I\) compiler or a role-inert lower-row factor.
\item Insert the low-response role correction
      \eqref{eq:V2-low-response-inert} as an upper cover currency.
      Test its non-anchor shadow bridges against
      \eqref{eq:V3-both-complements}; do not rename a shadow bridge as
      an endpoint response.
\item If both endpoint identities can avoid \(q\) while remaining
      clean and dominant, construct the smallest exact coordinate
      array and its V61 partition word.  Decide whether it is a true
      counterpacket or one of the two matching partitions already
      handled by V100.
\item Preserve payer locations separately and recheck V15, V83, V84,
      V91, and V100.  The next companion audit is active V5.
\end{enumerate}
\end{programme}

\begin{firewall}[Claim boundary after third-series V3]
\label{fire:V3-boundary}
V3 supplies the missing star compiler and collapses the finite graph
residual, but it does not prove acquisition of the named complementary
row, close the low-response sensitive cover, diagonal or nonclean
states, the other quartic source families, higher MTA, WUFC, CPM,
cutoff removal, reflection positivity, Osterwalder--Schrader
reconstruction, or a uniform positive continuum spectral gap.  The
full Yang--Mills mass gap remains open.
\end{firewall}

\paragraph{Parent-version record.}
V3 was copied byte-for-byte from
\texttt{Renewal\_Yang\_Mills\_Mass\_Gap\_Research\_Notes\_V2.tex}
before this chapter was appended.  The V2 parent had \(1\,584\) source
lines, \(57\,401\) bytes, and SHA--256
\[
 \texttt{091214516cc95ef47c9b8bcba00c22c2e99120e3b600323a2ef13cf1169846e4}.
\]
No companion audit is due at V3; the next required audit is V5.

\chapter{V4: Missing-complement acquisition and mass-compatible
response reversal}
\label{ch:V4-missing-complement}

\section{Context and the upstream move}
\label{sec:V4-context}

The V3 endpoint certificate replaced the \(28\) failed ordered cells
by a much smaller statement.  For a breaker edge \(ab\), either the
two complementary rows already occur as partners of endpoint
responses, in which case a path or star closes the entry, or at least
one complementary row is missing.  This chapter fixes such a missing
row and moves the depleted search to that row itself.

This is not the V83 higher-incidence tree exchange in new notation.
That theorem starts with a selected square-tail bank and changes a
nominal \(2+2\) path.  Here the numerator has already been localized
to the raw breaker stock \(H_{ab}\), one endpoint response may already
be present, and the new question is whether a response rooted at the
missing complement can be reversed without corrupting the two
denominators of the final quartic tree.  The only shared ingredient is
the scalar comparison of two row masses.

Fix distinct rows
\[
 \cV=\{a,b,z,q\},
\]
where \(ab\) is the breaker edge.  In the first part of the chapter we
assume that \(z\) is already seen from one breaker endpoint.  After
interchanging \(a,b\) if necessary, write this as
\begin{equation}
 d_{a\to z}(B_z)
 =\frac{H_{az}(B_z)}{\Xi_a}
 \ge\gamma a_*,
 \qquad \gamma>0.
\label{eq:V4-one-complement}
\end{equation}
Thus \(q\) is the unique missing complement in the sense of
\Cref{thm:V3-certificate-complete}.  Every bank below is an actual
bank occurrence in the assembled operator word.  When two displayed
responses use the same complete bank, the word is assembled before
either scalar mark is inserted, as required by the V98 same-bank
marking rule.

\begin{definition}[Mark-admissible simultaneous responses]
\label{def:V4-mark-admissible}
Two response factors are called \emph{mark-admissible} for the fixed
pair entry when their literal edge stocks occur in the already
assembled word and their simultaneous insertion is allowed by the
V15 path or V98 same-bank marking interface.  In particular, this
term never licenses a second operator occurrence, a second first
moment, or a second payer.
\end{definition}

\section{Reversing one actual response}
\label{sec:V4-reversal}

The response \(d_{x\to y}\) is oriented only because its denominator
is \(\Xi_x\); its numerator is symmetric.  A mass comparison is
therefore exactly what is needed to reverse it.

\begin{lemma}[Mass-compatible response reversal]
\label{lem:V4-response-reversal}
Let \(x\ne y\), let \(B\) be an actual bank, and suppose
\begin{equation}
 d_{x\to y}(B)
 =\frac{H_{xy}(B)}{\Xi_x}
 \ge\eta a_*.
\label{eq:V4-forward-response}
\end{equation}
If
\begin{equation}
 \Xi_y\le L\Xi_x,\qquad L\ge1,
\label{eq:V4-mass-compatible}
\end{equation}
then the same literal edge stock gives
\begin{equation}
 \boxed{
 d_{y\to x}(B)
 =\frac{H_{xy}(B)}{\Xi_y}
 \ge\frac{\eta}{L}a_*.}
\label{eq:V4-reversed-response}
\end{equation}
If \eqref{eq:V4-mass-compatible} fails, then
\begin{equation}
 \Xi_y>L\Xi_x,
\label{eq:V4-strict-ascent}
\end{equation}
which is a strict \(L\)-mass ascent from \(x\) to \(y\).
\end{lemma}

\begin{proof}
Symmetry gives \(H_{yx}(B)=H_{xy}(B)\).  Under
\eqref{eq:V4-mass-compatible},
\[
 \frac{H_{xy}(B)}{\Xi_y}
 \ge
 \frac1L\frac{H_{xy}(B)}{\Xi_x}
 \ge\frac{\eta}{L}a_*.
\]
The last assertion is the literal negation of
\eqref{eq:V4-mass-compatible}.
\end{proof}

\begin{remark}[What reversal does not do]
\label{rem:V4-no-relabel}
\Cref{lem:V4-response-reversal} changes only the denominator used to
normalize the same stock \(H_{xy}(B)\).  It does not relabel a
coordinate, replace \(xy\) by a different edge, or localize an
operator norm by \(H_{xy}\).  The complete word estimate must already
have been proved before the scalar comparison is made.
\end{remark}

\section{The relay compiler and the three-partner fan}
\label{sec:V4-fan}

The V1 path compiler covers the case in which the two response tails
are the two internal vertices of a path.  The following displayed
orientation is useful enough in the missing-complement search to
record separately.

\begin{theorem}[One-step relay path]
\label{thm:V4-relay-path}
Assume \eqref{eq:V4-one-complement}, and suppose a mark-admissible
second response satisfies
\begin{equation}
 d_{z\to q}(B_q)\ge\eta a_*.
\label{eq:V4-relay-response}
\end{equation}
Then, for \(\varepsilon\in\{0,1\}\),
\begin{equation}
 \boxed{
 Q_{ab}^{(\varepsilon)}
 \le
 \frac{C}{\gamma\eta a_*^2}\,
 d_{a\to z}(B_z)d_{z\to q}(B_q)H_{ab}(G).}
\label{eq:V4-relay-bound}
\end{equation}
The scalar on the right is the exact V92 aggregate response of the
path \(b-a-z-q\), and the estimate introduces no new payer.
\end{theorem}

\begin{proof}
Multiply the inherited pair estimate
\[
 Q_{ab}^{(\varepsilon)}\le C H_{ab}(G)
\]
by
\[
 1\le
 \frac{d_{a\to z}(B_z)d_{z\to q}(B_q)}
      {\gamma\eta a_*^2}.
\]
After expanding the responses, the inserted scalar is
\[
 \frac{H_{az}(B_z)H_{zq}(B_q)H_{ab}(G)}
      {\Xi_a\Xi_z}.
\]
The three literal edges are precisely the path \(b-a-z-q\), whose
degree-two vertices are \(a,z\).  Hence
\Cref{lem:V3-degree-form} identifies this scalar with the exact
aggregate tree normalization.  The payer was already assigned in
\(Q_{ab}^{(\varepsilon)}\), and only dimensionless response lower
bounds were inserted.
\end{proof}

\begin{theorem}[Missing-complement fan completion]
\label{thm:V4-fan-completion}
Assume \eqref{eq:V4-one-complement}.  Suppose an actual,
mark-admissible response rooted at the missing row obeys
\begin{equation}
 d_{q\to v}(B_q)\ge\eta a_*,
 \qquad v\in\{a,b,z\}.
\label{eq:V4-q-response}
\end{equation}
If
\begin{equation}
 \Xi_v\le L\Xi_q,
\label{eq:V4-q-compatible}
\end{equation}
then \(Q_{ab}^{(\varepsilon)}\) closes by a legal quartic tree, with
\begin{equation}
 \boxed{
 Q_{ab}^{(\varepsilon)}
 \le
 \frac{CL}{\gamma\eta a_*^2}\,
 d_{a\to z}(B_z)d_{v\to q}(B_q)H_{ab}(G).}
\label{eq:V4-fan-bound}
\end{equation}
More precisely:
\begin{center}
\begin{tabular}{c c c}
\toprule
\(v\) & reversed response & exact tree\\
\midrule
\(a\) & \(a\to q\) & star \(b-a-\{z,q\}\)\\
\(b\) & \(b\to q\) & path \(z-a-b-q\)\\
\(z\) & \(z\to q\) & path \(b-a-z-q\)\\
\bottomrule
\end{tabular}
\end{center}
\end{theorem}

\begin{proof}
By \Cref{lem:V4-response-reversal},
\[
 d_{v\to q}(B_q)\ge\frac{\eta}{L}a_*.
\]
Together with \eqref{eq:V4-one-complement}, this gives
\[
 1\le
 \frac{L}{\gamma\eta a_*^2}
 d_{a\to z}(B_z)d_{v\to q}(B_q).
\]
Multiplication by the pair estimate gives
\eqref{eq:V4-fan-bound}.

If \(v=a\), the literal edges are \(ab,az,aq\), and both response
denominators are \(\Xi_a\).  This is the star centered at \(a\), so
\Cref{thm:V3-star-compiler} applies.  If \(v=b\), the edges
\(az,ab,bq\) form the path \(z-a-b-q\), with denominator vertices
\(a,b\); this is the V1 complementary-partner path.  If \(v=z\), the
edges \(ab,az,zq\) form the path \(b-a-z-q\), with denominator
vertices \(a,z\), and \Cref{thm:V4-relay-path} applies.  These are all
three possible partners of \(q\).  Mark-admissibility preserves the
assembled-word and payer hypotheses in each invocation.
\end{proof}

\begin{corollary}[Response or strict ascent]
\label{cor:V4-response-or-ascent}
Under the hypotheses of \Cref{thm:V4-fan-completion}, but without
\eqref{eq:V4-q-compatible}, exactly one of the following two
alternatives holds:
\begin{enumerate}[label=\textup{(\roman*)},leftmargin=4em]
\item \(\Xi_v\le L\Xi_q\), and the pair closes by
      \eqref{eq:V4-fan-bound};
\item \(\Xi_v>L\Xi_q\), and \(q\to v\) is a strict \(L\)-mass ascent.
\end{enumerate}
\end{corollary}

\begin{proof}
The two mass comparisons are complementary.  Apply
\Cref{thm:V4-fan-completion} in the first case and
\Cref{lem:V4-response-reversal} in the second.
\end{proof}

\begin{corollary}[A maximal missing row has no response obstruction]
\label{cor:V4-maximal-q}
If
\[
 \Xi_q=\max_{r\in\cV}\Xi_r,
\]
then every response branch \eqref{eq:V4-q-response} in the
single-missing-complement state closes with \(L=1\).  More generally,
if
\[
 \max_{r\in\cV}\Xi_r\le L\Xi_q,
\]
then every such branch closes with the stated \(L\).
\end{corollary}

\begin{proof}
For every possible partner \(v\), the displayed hypothesis gives
\(\Xi_v\le L\Xi_q\).  Apply
\Cref{thm:V4-fan-completion}.
\end{proof}

\section{Transporting the depleted identity to the missing row}
\label{sec:V4-rowwise-acquisition}

The proof of the V1 acquisition theorem is row-local: neither its
six-cell identity nor its anchor pigeonhole uses the fact that the
chosen row is an endpoint of the original breaker edge.  We record
the resulting transport because the new search is rooted at \(q\).

\begin{lemma}[Rowwise depleted-acquisition transport]
\label{lem:V4-rowwise-transport}
Let \(r\in\cV\), and let \(\cU\) be the disjointly represented union
of at most \(M_0\) complete banks already present in an assembled
four-row word, with an actual anchor recorded for each bank.  Assume
row \(r\) is in the clean calculus.  For every \(0<\delta<1\), the
exact identity
\begin{equation}
 \Xi_r=X_r(\cU)+
 P_r^{\rm fr}+U_r^{\rm fr}+A_r^{\rm fr}+
 F_r^{\rm fr}+D_r^{\rm fr}+I_r^{\rm fr}
\label{eq:V4-rowwise-depleted}
\end{equation}
has one of the following consequences:
\begin{enumerate}[label=\textup{(\alph*)},leftmargin=4em]
\item a fresh \(P,U,A,F,\) or \(I\) cell of size at least
      \((1-\delta)\Xi_r/6\), entering its already proved complete
      compiler;
\item an actual fresh directed response
      \begin{equation}
      d_{r\to v}(D_v)\ge
      \frac{1-\delta}{18}a_*,
      \qquad v\ne r;
      \label{eq:V4-rowwise-D}
      \end{equation}
\item an actual external-anchor response
      \begin{equation}
      d_{r\to v}(B_v)\ge
      \frac{\delta}{2M_0}a_*,
      \qquad v\ne r;
      \label{eq:V4-rowwise-external}
      \end{equation}
\item a single self-anchored bank \(B_r\) with
      \begin{equation}
      X_r(B_r)\ge\frac{\delta}{2M_0}\Xi_r.
      \label{eq:V4-rowwise-self}
      \end{equation}
\end{enumerate}
Outside the hypothesis there is a separately labelled nonclean
state.
\end{lemma}

\begin{proof}
Every step in the proof of
\Cref{lem:V1-anchored-overlap,cor:V1-self-witness,
thm:V1-one-step-acquisition}
uses only the chosen row's weights and the at-most-\(M_0\) bank
anchors.  Replace the symbol \(x\) there by \(r\).

For completeness, if \(X_r(\cU)<\delta\Xi_r\), the six fresh terms in
\eqref{eq:V4-rowwise-depleted} total more than
\((1-\delta)\Xi_r\).  A largest one has at least one sixth of this
mass.  In the \(D\) case, partition by the at most three actual
partners; a largest partner cell has \(r\)-mass at least
\((1-\delta)\Xi_r/18\).  The Casimir floor gives
\eqref{eq:V4-rowwise-D}.  If
\(X_r(\cU)\ge\delta\Xi_r\), the anchor split gives either an external
bank with \(r\)-mass at least \(\delta\Xi_r/(2M_0)\), hence
\eqref{eq:V4-rowwise-external}, or the self-concentration stock.
Pigeonholing the latter over at most \(M_0\) banks gives
\eqref{eq:V4-rowwise-self}.  All selected edges are literal
common-coordinate edges.
\end{proof}

\begin{remark}[Eligibility is an operator statement]
\label{rem:V4-eligibility}
\Cref{lem:V4-rowwise-transport} is a scalar partition only after the
complete banks in \(\cU\) have been identified in the assembled
operator word.  It does not assert that an arbitrary auxiliary copy
of row \(q\) may be introduced.  In applications below, an
\emph{eligible missing row} means precisely that the source word
already supplies this rowwise identity with the stated bank and payer
bookkeeping.
\end{remark}

\section{Single-missing-complement normal form}
\label{sec:V4-single-normal-form}

Set
\begin{equation}
 \eta_{\delta}:=
 \min\left\{\frac{1-\delta}{18},
             \frac{\delta}{2M_0}\right\}.
\label{eq:V4-eta-delta}
\end{equation}
Thus either directed alternative of
\Cref{lem:V4-rowwise-transport} has dominance at least
\(\eta_\delta a_*\).

\begin{theorem}[One-step normal form at the named missing complement]
\label{thm:V4-single-normal-form}
Assume the single-missing-complement state
\eqref{eq:V4-one-complement}, assume \(q\) is eligible in the sense of
\Cref{rem:V4-eligibility}, and apply
\Cref{lem:V4-rowwise-transport} at \(q\).  Every resulting branch is
one of:
\begin{enumerate}[label=\textup{(\roman*)},leftmargin=4em]
\item a \(P,U,A,F,\) or \(I\) complete-compiler branch;
\item a mark-admissible response \(q\to v\), \(v\in\{a,b,z\}\), which
      either closes by \Cref{thm:V4-fan-completion} or satisfies the
      strict ascent \(\Xi_v>L\Xi_q\);
\item a same-type self-anchored branch covered by
      \Cref{prop:V1-self-type-routes};
\item a partition-changing self-bank whose response-first reduction
      gives either a response as in \textup{(ii)}, or the
      low-response \(q\)-sensitive correction and a \(q\)-free role
      difference of \eqref{eq:V2-low-response-inert} and
      \eqref{eq:V2-insensitive-factor};
\item a nonclean state.
\end{enumerate}
No untyped clean branch remains after this one step.
\end{theorem}

\begin{proof}
Alternative \textup{(i)}, the two response sources in
\textup{(ii)}, the self-bank, and the nonclean state are exhaustive
by \Cref{lem:V4-rowwise-transport}.  A response has one of the three
actual partners \(a,b,z\); apply
\Cref{cor:V4-response-or-ascent}.  If the self-bank has a common
clean type, \Cref{prop:V1-self-type-routes} gives \textup{(iii)}.
Otherwise it is the partition-changing bank to which
\Cref{thm:V2-self-role-reduction} applies with anchor \(q\).  Its
response-first branch is again \textup{(ii)}.  Its remaining branch
is exactly the sensitive/insensitive factorization stated in
\textup{(iv)}.  These cases exhaust the rowwise identity without
assigning any shadow bridge the status of an endpoint response.
\end{proof}

\begin{corollary}[What is genuinely left when \(q\) is maximal]
\label{cor:V4-maximal-normal-form}
Under \Cref{thm:V4-single-normal-form}, if \(q\) is a maximal-mass row,
then no directed response branch survives.  The clean residual is
confined to the \(q\)-sensitive low-response correction and the
\(q\)-free lower-row role difference, after the established complete
compilers have been removed.
\end{corollary}

\begin{proof}
\Cref{cor:V4-maximal-q} closes every branch in
\textup{(ii)}, including the response-first half of
\textup{(iv)}.  The same-type and complete-compiler cases are already
assigned to their inherited compilers.  What remains is precisely the
second half of \textup{(iv)}.
\end{proof}

\begin{assessment}[New analytic bottleneck in the singleton state]
\label{ass:V4-single-bottleneck}
The missing-complement response is no longer a graph obstruction at a
maximal row, and at a nonmaximal row its only graph failure is an
explicit strict ascent.  The non-ascent obstruction is now the V2
low-response role correction: a small \(q\)-sensitive change plus a
large common contraction whose changing factor omits \(q\).  Closing
that object requires an upper cover estimate that retains the actual
partner incidences of its lower-row factor; it cannot be obtained by
calling the shadow factor a response.
\end{assessment}

\section{When both complementary rows are initially missing}
\label{sec:V4-double-missing}

The empty endpoint-complement set in
\Cref{cor:V3-collapse-28} is logically different from the singleton
state.  A single response rooted at \(q\) cannot by itself provide the
two response denominators needed beside \(H_{ab}\).  It does,
however, have an exact two-way reduction.

\begin{proposition}[First acquisition in the double-missing state]
\label{prop:V4-double-first}
Suppose no dominant response rooted at \(a\) or \(b\) has partner
\(z\) or \(q\).  Let an eligible acquisition at \(q\) produce
\[
 d_{q\to v}(B_q)\ge\eta a_*,
 \qquad v\in\{a,b,z\}.
\]
If \(\Xi_v\le L\Xi_q\), then:
\begin{enumerate}[label=\textup{(\roman*)},leftmargin=4em]
\item for \(v=a\) or \(v=b\), reversal produces the endpoint response
      \(a\to q\) or \(b\to q\), so the endpoint-complement set becomes
      a singleton;
\item for \(v=z\), the literal bridge \(qz\) is available in both
      orientations with thresholds \(\eta a_*\) and
      \(\eta a_*/L\).  Any later endpoint response from \(a\) or \(b\)
      to either \(q\) or \(z\), mark-admissible with this bridge,
      closes a four-vertex path with the raw edge \(ab\).
\end{enumerate}
If \(\Xi_v>L\Xi_q\), the branch is a strict \(L\)-mass ascent.
\end{proposition}

\begin{proof}
For \(v=a,b\), apply
\Cref{lem:V4-response-reversal}; the reversed response is rooted at a
breaker endpoint and has partner \(q\), proving \textup{(i)}.

For \(v=z\), the original response is \(q\to z\), while
\Cref{lem:V4-response-reversal} supplies \(z\to q\).  Suppose, for
example, the later endpoint response is \(a\to q\).  Use
\(q\to z\); together with \(ab\) the edges form
\[
 b-a-q-z
\]
and the denominators are at its internal vertices \(a,q\).  If the
later response is \(a\to z\), use \(z\to q\), obtaining
\[
 b-a-z-q
\]
with denominator vertices \(a,z\).  Replacing \(a\) by \(b\) gives
the two symmetric paths.  In all four cases
\Cref{lem:V3-degree-form} verifies the exact aggregate
normalization.  The final alternative is the failed mass comparison.
\end{proof}

\begin{corollary}[Double-missing response reduction]
\label{cor:V4-double-reduction}
In the double-missing state, one clean eligible acquisition at a
complementary row has only the following response outcomes:
\begin{enumerate}[label=\textup{(\alph*)},leftmargin=4em]
\item a strict mass ascent;
\item a singleton endpoint-complement state;
\item a bidirectional complementary bridge awaiting one literal
      endpoint-to-complement response.
\end{enumerate}
\end{corollary}

\begin{proof}
This is the three-partner list in
\Cref{prop:V4-double-first}.
\end{proof}

\begin{firewall}[A bridge is not yet a tree]
\label{fire:V4-bridge-not-tree}
The complementary edge \(qz\) and the breaker edge \(ab\) form two
components.  Even when \(qz\) is available in both orientations, V4
does not declare closure until an actual endpoint-to-complement edge
joins those components.  This preserves the V1 no-relabel firewall
and the V3 denominator-degree test.
\end{firewall}

\section{Finite ascent facts and the unresolved continuation}
\label{sec:V4-ascent}

\begin{proposition}[Strict mass ascent cannot cycle]
\label{prop:V4-ascent-acyclic}
Let \(L\ge1\).  Any sequence of rows
\[
 r_0,r_1,\ldots,r_N\in\cV
\]
satisfying
\begin{equation}
 \Xi_{r_{j+1}}>L\Xi_{r_j}
 \quad(0\le j<N)
\label{eq:V4-ascent-chain}
\end{equation}
contains no repeated row.  Consequently \(N\le3\), and
\[
 \Xi_{r_N}>L^N\Xi_{r_0}.
\]
\end{proposition}

\begin{proof}
Iterating \eqref{eq:V4-ascent-chain} gives, for \(i<j\),
\[
 \Xi_{r_j}>L^{j-i}\Xi_{r_i}\ge\Xi_{r_i}.
\]
Thus \(r_j\ne r_i\).  Since \(\cV\) has four elements, at most three
strict transitions are possible.  The displayed quantitative bound
is the same iteration with \(i=0,j=N\).
\end{proof}

\begin{corollary}[Terminal rows are mass-compatible]
\label{cor:V4-terminal-compatible}
A row of maximal mass has no outgoing strict ascent.  Thus, whenever
the acquisition packet can be transported without alteration along
an ascent chain, its response branch becomes mass-compatible after at
most three transitions.
\end{corollary}

\begin{proof}
The first statement is immediate from maximality.  The second is the
conditional combination of
\Cref{prop:V4-ascent-acyclic,cor:V4-response-or-ascent}.
\end{proof}

\begin{firewall}[Acyclicity is not packet transport]
\label{fire:V4-no-ascent-transport}
\Cref{prop:V4-ascent-acyclic} is a finite scalar fact.  It does not
prove that the same localized breaker packet, bank disjointness, role
partition, and payer assignment survive when the depleted search is
restarted at \(r_{j+1}\).  Such a transport theorem is needed before
\Cref{cor:V4-terminal-compatible} can be used as an unconditional
termination argument.
\end{firewall}

\begin{proposition}[V4 regression suite]
\label{prop:V4-regressions}
The new compilers obey the following checks.
\begin{enumerate}[label=\textup{(\roman*)},leftmargin=4em]
\item Every edge in the fan table is a literal stock already present
      in the assembled word.
\item Reversal changes \(H_{qv}/\Xi_q\) to \(H_{qv}/\Xi_v\); it never
      changes the edge \(qv\).
\item The star case has denominator multiset
      \(\{a,a\}\), while the two path cases have respectively
      \(\{a,b\}\) and \(\{a,z\}\).
\item A paid pair remains once-paid because no resolvent or output
      fusion is introduced after the pair allocation.
\item The V83 exchange is cited only as historical precedent; none of
      its square-tail conclusions is used without its selected-bank
      hypotheses.
\item A complementary bridge is recorded as disconnected until a
      literal endpoint bridge occurs.
\end{enumerate}
\end{proposition}

\begin{proof}
Items \textup{(i)--(iii)} follow from
\Cref{lem:V4-response-reversal,thm:V4-fan-completion}.
Item \textup{(iv)} is part of
\Cref{thm:V4-relay-path}.  Item \textup{(v)} follows from the
independent scalar proof of \Cref{lem:V4-response-reversal}.
Item \textup{(vi)} is \Cref{fire:V4-bridge-not-tree}.
\end{proof}

\section{Result ledger and V5 mandate}
\label{sec:V4-ledger}

\begin{theorem}[Third-series V4 result ledger]
\label{thm:V4-results}
Relative to V1--V3 and the two frozen archives, V4 proves:
\begin{enumerate}[label=\textup{(V4.\arabic*)},leftmargin=5.8em]
\item the exact mass-compatible reversal
      \eqref{eq:V4-reversed-response} of an actual directed response;
\item the relay path estimate \eqref{eq:V4-relay-bound};
\item complete closure of all three partner positions for a response
      rooted at the unique missing complement, subject only to the
      explicit mass comparison \eqref{eq:V4-q-compatible};
\item the maximal-row closure
      \Cref{cor:V4-maximal-q};
\item rowwise transport of the exact depleted-acquisition identity;
\item the exhaustive singleton-complement normal form
      \Cref{thm:V4-single-normal-form};
\item reduction of a first response in the double-missing state to a
      strict ascent, a singleton endpoint hit, or a bidirectional
      complementary bridge; and
\item acyclicity and length at most three for strict four-row ascent
      chains.
\end{enumerate}
\end{theorem}

\begin{proof}
The assertions are respectively
\Cref{lem:V4-response-reversal,
thm:V4-relay-path,
thm:V4-fan-completion,
cor:V4-maximal-q,
lem:V4-rowwise-transport,
thm:V4-single-normal-form,
cor:V4-double-reduction,
prop:V4-ascent-acyclic}.
\end{proof}

\begin{programme}[V5: audit, then close packet transport or descend]
\label{prog:V4-V5}
Before adding V5, inspect the current Standard Model and
Navier--Stokes companion notes as required by the five-version audit.
Then:
\begin{enumerate}[leftmargin=3em]
\item compare their latest ancestor-transport or concentration
      mechanisms with the exact Yang--Mills bank and payer data;
\item attempt a packet-preserving ascent lemma carrying the breaker
      localization, mark-admissibility, and sole payer from \(q\) to
      the larger row \(v\);
\item if that transport fails, construct the minimal V61 partition
      word exhibiting the failure and descend to its first role
      change;
\item split the V2 low-response correction at its first changing
      lower row, retaining actual partner incidences rather than
      shadow labels;
\item test whether a complementary bridge from
      \Cref{prop:V4-double-first} supplies the missing connection;
\item re-run the V15, V78, V83, V84, V91, V98, and V100 regression
      examples with payer locations fixed.
\end{enumerate}
\end{programme}

\begin{firewall}[Claim boundary after third-series V4]
\label{fire:V4-boundary}
V4 closes every mass-compatible response branch in the
single-missing-complement state and gives exact finite reductions for
the first double-missing acquisition.  It does not prove
packet-preserving ascent transport, close the low-response sensitive
role correction, manufacture an endpoint edge from a disconnected
complementary bridge, or close diagonal, nonclean, and other quartic
source families.  Higher MTA, WUFC, CPM, cutoff removal, reflection
positivity, Osterwalder--Schrader reconstruction, and a uniform
positive continuum spectral gap also remain open.  The full
Yang--Mills mass gap has not been proved.
\end{firewall}

\paragraph{Parent-version record.}
V4 was copied byte-for-byte from
\texttt{Renewal\_Yang\_Mills\_Mass\_Gap\_Research\_Notes\_V3.tex}
before this chapter was appended.  The V3 parent had \(2\,017\) source
lines, \(73\,742\) bytes, and SHA--256
\[
 \texttt{04ab0b8b5dff89a707de085524b47a47fea055614f5400d1083fecec25686cb2}.
\]
No companion audit is due at V4; the next required audit is V5.

\chapter{V5: Companion audit, reverse-response sharpness, and
lower-row role descent}
\label{ch:V5-audit-descent}

\section{Compulsory companion audit}
\label{sec:V5-audit}

Before choosing the V5 direction, the current Standard Model and
Navier--Stokes notebooks in \texttt{latex\_notes} were read
passively.  The audited files were:
\begin{center}
\small
\begin{tabular}{p{0.48\textwidth} r r}
\toprule
file & lines & bytes\\
\midrule
\texttt{Renewal\_SM\_Einstein\_Continuum\_Research\_Notes\_V83.tex}
& \(33\,149\) & \(1\,199\,038\)\\
\texttt{Renewal\_NS\_Stability\_Research\_Notes\_V31.tex}
& \(23\,052\) & \(887\,537\)\\
\bottomrule
\end{tabular}
\end{center}
Their SHA--256 digests at the audit time were
\[
\begin{split}
&\texttt{4f53891fe399799539d33ab00b4e4d7ab3f26c1fea84d003c92440187adf36db},\\
&\texttt{f1121b62238ad5491bb7cf03635ea9934942edf573ed8faaa9ee79e1d38a6696}.
\end{split}
\]

The SM note's V81--V83 chain transports a complete determinant
connection only along a covariant uniformly gapped local path.  Its
V83 blocking map gives an exact orthogonal decomposition, but the
coarse determinant remains and the missing interpolation gap is
stated explicitly.  The lesson for the present packet is that a
finite scalar chain of row masses cannot transport the breaker
localization, response occurrences, and payer by itself.

The NS note proves an exact dyadic unmarking identity: normalized
probability marks recover the physical tail only after an unbounded
first-moment weight is restored.  A far atom and a genuine smooth
packet saturate that loss.  Its useful warning here is that a
normalized low-response bound does not control the unnormalized
role-changing factor unless the source incidence correlated with the
mark is retained before absolute values.

\begin{record}[V5 audit decision]
\label{rec:V5-audit-decision}
V5 will not promote \Cref{prop:V4-ascent-acyclic} to a termination
theorem.  It will instead:
\begin{enumerate}[label=\textup{(\roman*)},leftmargin=4em]
\item inspect the reverse response itself, rather than infer it only
      from a comparison of total row masses;
\item prove a sharp counterexample to any support-uniform reversal
      without additional structure;
\item show that failure of reverse dominance depletes the actual
      partner-incidence bank at the target row; and
\item keep the V2 insensitive common contraction and its \(q\)-free
      role change in one exact tensor factorization before estimating
      either part.
\end{enumerate}
\end{record}

\section{Direct reversal is the real criterion}
\label{sec:V5-direct-reversal}

Continue with the V4 singleton state
\[
 \cV=\{a,b,z,q\},\qquad
 d_{a\to z}(B_z)\ge\gamma a_*,
\]
where \(q\) is the missing complement.  Suppose an acquired response
is \(q\to v\), \(v\in\{a,b,z\}\).  The total-mass comparison in V4 is
only one sufficient condition for the orientation \(v\to q\); the
tree compiler uses the latter orientation directly.

\begin{theorem}[Direct reverse-response fan criterion]
\label{thm:V5-direct-reverse}
Assume that the same literal \(qv\) bank which produced
\(q\to v\) also satisfies
\begin{equation}
 d_{v\to q}(B_q)
 =\frac{H_{qv}(B_q)}{\Xi_v}
 \ge\theta a_*,
 \qquad\theta>0,
\label{eq:V5-direct-reverse}
\end{equation}
and that it is mark-admissible with the response \(a\to z\).  Then
\begin{equation}
 \boxed{
 Q_{ab}^{(\varepsilon)}
 \le
 \frac{C}{\gamma\theta a_*^2}\,
 d_{a\to z}(B_z)d_{v\to q}(B_q)H_{ab}(G),}
\label{eq:V5-direct-fan}
\end{equation}
and the right side is a legal star for \(v=a\), the path
\(z-a-b-q\) for \(v=b\), or the path \(b-a-z-q\) for \(v=z\).
This conclusion requires no comparison between \(\Xi_q\) and
\(\Xi_v\).
\end{theorem}

\begin{proof}
The lower bounds give
\[
 1\le
 \frac{d_{a\to z}(B_z)d_{v\to q}(B_q)}
      {\gamma\theta a_*^2}.
\]
Multiply the pair-localized estimate
\(Q_{ab}^{(\varepsilon)}\le CH_{ab}(G)\).
The three possible literal-edge and denominator patterns are exactly
those verified in \Cref{thm:V4-fan-completion}; that proof used the
mass comparison only to derive the lower bound
\eqref{eq:V4-reversed-response}.  Here
\eqref{eq:V5-direct-reverse} supplies that bound directly.
\end{proof}

\begin{corollary}[Mass compatibility is sufficient, not necessary]
\label{cor:V5-mass-only-sufficient}
The V4 hypothesis \(\Xi_v\le L\Xi_q\) and forward dominance
\(d_{q\to v}\ge\eta a_*\) imply
\eqref{eq:V5-direct-reverse} with \(\theta=\eta/L\).  Conversely,
\eqref{eq:V5-direct-reverse} may hold with a fixed \(\theta\) even
when \(\Xi_v/\Xi_q\) is arbitrarily large.
\end{corollary}

\begin{proof}
The first statement is
\Cref{lem:V4-response-reversal}.  For the second, the numerator
\(H_{qv}\) may grow proportionally to \(\Xi_v\) on their common
coordinates, independently of the ratio of total row masses.
\end{proof}

\section{Sharpness of reversal from a forward response}
\label{sec:V5-reversal-sharpness}

\begin{proposition}[No support-uniform reverse lower bound]
\label{prop:V5-no-uniform-reverse}
There is no constant \(c>0\), independent of support and row masses,
such that
\[
 d_{q\to v}(B)\ge a_*
 \quad\Longrightarrow\quad
 d_{v\to q}(B)\ge c a_*.
\]
This remains false for clean weights obeying the exact floor
\eqref{eq:V1-floor}.
\end{proposition}

\begin{proof}
For an integer \(N\ge1\), take \(N+1\) coordinates.  On \(e_0\), set
\[
 a_{q,e_0}=a_{v,e_0}=a_*.
\]
On \(e_1,\ldots,e_N\), make only \(v\) active and set
\(a_{v,e_j}=a_*\).  Let \(B=\{e_0\}\).  Then
\[
 \Xi_q=a_*,
 \qquad
 \Xi_v=(N+1)a_*,
 \qquad
 H_{qv}(B)=a_*^2.
\]
Consequently
\[
 d_{q\to v}(B)=a_*,
 \qquad
 d_{v\to q}(B)=\frac{a_*}{N+1}.
\]
Given \(c>0\), choose \(N+1>1/c\).  All nonzero weights equal the
allowed floor, so the example lies in the clean scalar calculus.
\end{proof}

\begin{remark}[The obstruction is off-edge target mass]
\label{rem:V5-off-edge-mass}
The counterexample does not weaken the literal edge \(qv\).  It adds
mass to row \(v\) on coordinates outside that edge.  Thus the correct
response to failed reversal is to deplete the \(qv\)-incidence bank
at \(v\), not to relabel the edge or to sum an unweighted chain of
forward responses.
\end{remark}

\section{Failed reversal forces target-row depletion}
\label{sec:V5-target-depletion}

For rows \(x\ne y\) and a bank \(B\), define its literal incidence
subbank
\begin{equation}
 E_{xy}(B):=
 \{e\in B:a_{x,e}>0\text{ and }a_{y,e}>0\},
\label{eq:V5-incidence-subbank}
\end{equation}
and put
\[
 X_y(E):=\sum_{e\in E}a_{y,e}.
\]
Terms outside \(E_{xy}(B)\) do not contribute to \(H_{xy}(B)\).

\begin{lemma}[Edge-stock domination of incidence mass]
\label{lem:V5-stock-dominates-incidence}
For every clean bank \(B\) and distinct \(x,y\),
\begin{equation}
 H_{xy}(B)
 =H_{xy}\bigl(E_{xy}(B)\bigr)
 \ge a_*X_y\bigl(E_{xy}(B)\bigr),
\label{eq:V5-stock-incidence}
\end{equation}
and symmetrically with \(x,y\) interchanged.
\end{lemma}

\begin{proof}
Outside \(E_{xy}(B)\), at least one factor in
\(a_{x,e}a_{y,e}\) vanishes.  On \(E_{xy}(B)\), the floor gives
\(a_{x,e}\ge a_*\), so
\[
 a_{x,e}a_{y,e}\ge a_*a_{y,e}.
\]
Sum over the incidence subbank.  Symmetry proves the other version.
\end{proof}

\begin{theorem}[Reverse-or-deplete dichotomy]
\label{thm:V5-reverse-or-deplete}
Let \(d_{q\to v}(B_q)\ge\eta a_*\), and fix a target threshold
\(0<\theta<1\).  Exactly one of the following alternatives may be
selected:
\begin{enumerate}[label=\textup{(\roman*)},leftmargin=4em]
\item
      \[
      d_{v\to q}(B_q)\ge\theta a_*;
      \]
      in the singleton missing-complement state,
      \Cref{thm:V5-direct-reverse} closes the pair;
\item
      \begin{equation}
      X_v\bigl(E_{qv}(B_q)\bigr)<\theta\Xi_v,
      \label{eq:V5-target-depleted}
      \end{equation}
      so more than a \((1-\theta)\)-fraction of row \(v\)'s total mass
      lies outside the actual \(qv\)-incidence subbank.
\end{enumerate}
\end{theorem}

\begin{proof}
If \textup{(i)} fails, then
\[
 H_{qv}(B_q)<\theta a_*\Xi_v.
\]
By \Cref{lem:V5-stock-dominates-incidence},
\[
 a_*X_v(E_{qv}(B_q))
 \le H_{qv}(B_q)
 <\theta a_*\Xi_v.
\]
Cancel \(a_*\) to get \eqref{eq:V5-target-depleted}.  Its complement
has \(v\)-mass greater than \((1-\theta)\Xi_v\).  The two alternatives
are exhaustive by the direct reverse threshold.
\end{proof}

\begin{corollary}[The strict-ascent branch has an upstream target]
\label{cor:V5-ascent-target}
Every V4 strict-ascent branch \(q\to v\) admits the sharper
subdivision:
\begin{enumerate}[label=\textup{(\alph*)},leftmargin=4em]
\item the reverse response is already dominant at any chosen fixed
      threshold \(\theta\), and the fan closes without using the
      total-mass ratio; or
\item the actual \(qv\)-incidence bank occupies less than
      \(\theta\Xi_v\), and the next acquisition must examine the
      complementary \(v\)-mass rather than repeat the same edge.
\end{enumerate}
\end{corollary}

\begin{proof}
Apply \Cref{thm:V5-reverse-or-deplete}.  The fact that the total
masses are ascending is not used in the subdivision.
\end{proof}

\begin{firewall}[Scalar depletion does not assemble the new word]
\label{fire:V5-depletion-not-transport}
Equation \eqref{eq:V5-target-depleted} identifies a large stock off
the failed edge, but it does not prove that this stock occurs in a
mark-admissible complete operator bank with the old breaker pair and
sole payer.  This is the exact packet-transport datum missing from
\Cref{fire:V4-no-ascent-transport}.
\end{firewall}

\section{The V2 low-response branch as an exact recursion}
\label{sec:V5-role-recursion}

Return to the partition-changing self-bank at the missing row \(q\).
In the V2 low-response branch,
\[
 \Delta_G^{\rho,\sigma}
 =
 \Delta_{S_q}^{\rho,\sigma}\otimes M_{F_q}(\sigma)
 +
 M_{S_q}(\rho)\otimes
 A_{F_q}^{q}\otimes
 \bigl(R_{F_q}(\sigma)-R_{F_q}(\rho)\bigr).
\]
The last changing factor omits \(q\).  The companion audit requires
that this whole factorization be retained: estimating the common
factor and the lower-row difference separately would discard exactly
the incidence which might reconnect \(q\).

\begin{lemma}[Two-cover bound after deleting one row]
\label{lem:V5-two-cover-lower-row}
Let \(W\subseteq\cV\setminus\{q\}\), so \(|W|\le3\), and let
\(\rho_W,\sigma_W\) be two partitions of \(W\).  There is a signed
path from \(\rho_W\) to \(\sigma_W\) in the cover graph of the
partition lattice using at most two cover increments.  Consequently
the product-moment difference
\[
 R_W(\sigma_W)-R_W(\rho_W)
\]
is a signed sum of at most two exact composite cover covariances,
each having crossing stock supported entirely on pairs in \(W\).
\end{lemma}

\begin{proof}
For \(|W|\le2\), the assertion is immediate.  Let \(|W|=3\) and put
\(\eta_W=\rho_W\vee\sigma_W\).  If the partitions are comparable,
their block counts differ by at most two, so a saturated chain has at
most two covers.

Suppose they are incomparable.  A nontrivial partition of three
points which is incomparable with another one must have two blocks:
it consists of one pair and one singleton.  Two distinct such
partitions have one-block join.  Each is therefore one cover below
\(\eta_W\), giving a two-edge signed path through the join.
Telescoping product moments along this path and applying the exact
cover-covariance identity of \Cref{thm:V2-cover-covariance} proves the
operator statement.  Since row \(q\notin W\), every pair entering a
crossing stock has both endpoints in \(W\).
\end{proof}

\begin{proposition}[Incidence location in the insensitive term]
\label{prop:V5-insensitive-incidence}
In
\begin{equation}
 {\cal I}_q:=
 M_{S_q}(\rho)\otimes A_{F_q}^{q}\otimes
 \bigl(R_{F_q}(\sigma)-R_{F_q}(\rho)\bigr),
\label{eq:V5-insensitive-term}
\end{equation}
all role-changing cover edges lie in
\(\cV\setminus\{q\}\), and at most two cover increments are required.
Therefore every literal edge incident to \(q\) which can reconnect
the missing complement must occur in the common factor
\(M_{S_q}(\rho)\otimes A_{F_q}^{q}\), or in another bank already
present in the surrounding assembled word.
\end{proposition}

\begin{proof}
By \eqref{eq:V2-insensitive-factor}, the changing row sets of
\(R_{F_q}(\lambda)\) omit \(q\).  Apply
\Cref{lem:V5-two-cover-lower-row} to their restrictions.  Its cover
stocks contain no \(q\)-edge.  Tensor multiplication by the two
common factors changes neither those row sets nor their edge labels.
Hence any \(q\)-incident edge in the complete tensor word must come
from a common or surrounding factor.
\end{proof}

\begin{corollary}[Conditional closure from the common factor]
\label{cor:V5-common-factor-closure}
In the singleton missing-complement state, suppose the common factor
in \eqref{eq:V5-insensitive-term} contains a mark-admissible response
\[
 d_{q\to v}(B)\ge\eta a_*,
 \qquad v\in\{a,b,z\}.
\]
Then either its reverse response is dominant and
\Cref{thm:V5-direct-reverse} closes the pair, or its literal
\(qv\)-incidence subbank is depleted at \(v\) as in
\eqref{eq:V5-target-depleted}.
\end{corollary}

\begin{proof}
Apply \Cref{thm:V5-reverse-or-deplete} to that actual response.
The lower-row cover difference remains attached throughout and does
not supply or alter the \(qv\) edge.
\end{proof}

\section{A sharp low-response counterpacket}
\label{sec:V5-low-response-counterpacket}

The next example shows why the normalized low-response inequalities
alone cannot force the common factor to contain a useful
\(q\)-incidence.

\begin{proposition}[Low-response scalar data do not create a
missing-complement edge]
\label{prop:V5-low-response-no-edge}
There are clean four-row incidence data and partitions \(\rho,\sigma\)
for which
\begin{equation}
 X_q(G)=\Xi_q,\qquad
 r_q(\rho,\sigma;G)=0,\qquad
 X_q^{\rm sens}=0,
\label{eq:V5-zero-sensitive-example}
\end{equation}
while
\[
 \Delta_G^{\rho,\sigma}
 =
 A_G^q\otimes\bigl(R_G(\sigma)-R_G(\rho)\bigr)
\]
is a possibly nonzero role difference whose changing factor omits
\(q\).  Adding a breaker edge \(ab\) and one response edge \(az\)
does not produce a connected four-row tree: row \(q\) remains
isolated.
\end{proposition}

\begin{proof}
Take
\[
 \rho=a\mid b\mid z\mid q,
 \qquad
 \sigma=ab\mid z\mid q.
\]
Thus the block of \(q\) is the singleton \(\{q\}\) in both
partitions, and \(P_q(\rho,\sigma)=\varnothing\).

Let the complete bank \(G\) contain a coordinate \(e_q\) at which
only \(q\) is active, with weight \(a_*\), and a disjoint coordinate
\(e_{ab}\) at which \(a,b\) are active and \(q\) is inactive.  Give
row \(q\) no other coordinates.  Then
\[
 \Xi_q=X_q(G)=a_*.
\]
Because the changed-partner set of \(q\) is empty, both the total
changed response and its sensitive mass vanish, proving
\eqref{eq:V5-zero-sensitive-example}.  Coordinate factorization gives
\[
 M_G(\lambda)
 =
 M_{e_q}(\{q\})\otimes M_{e_{ab}}(\lambda),
 \qquad\lambda\in\{\rho,\sigma\},
\]
and subtraction gives
\[
 \Delta_G^{\rho,\sigma}
 =
 M_{e_q}(\{q\})\otimes
 \bigl(M_{e_{ab}}(\sigma)-M_{e_{ab}}(\rho)\bigr).
\]
The second factor contains only \(a,b\), and it is nonzero whenever
the corresponding local \(ab\) covariance is nonzero.

Now add separate literal stocks \(H_{ab}>0\) and \(H_{az}>0\).
Every displayed changing or response edge lies in
\(\{a,b,z\}\).  No edge is incident to \(q\), so the edge graph has
isolated vertex \(q\) and cannot be a spanning tree.  This conclusion
is purely graph-theoretic and does not assert nonvanishing of a
connected cumulant after all source subtractions.
\end{proof}

\begin{firewall}[Scope of the counterpacket]
\label{fire:V5-counterpacket-scope}
\Cref{prop:V5-low-response-no-edge} disproves only an inference from
the V2 scalar low-response data to a \(q\)-edge.  It is not a
counterexample to Yang--Mills clustering or to the full assembled
source estimate.  Additional connected-cumulant subtraction may
annihilate the factorized example; if so, that annihilation is
exactly an operator identity which must be proved and may become the
desired compiler.
\end{firewall}

\section{The corrected residual and result ledger}
\label{sec:V5-results}

\begin{assessment}[Residual after the audit]
\label{ass:V5-residual}
The V4 mass-ascent label contained two different phenomena.  V5
separates them:
\begin{enumerate}[label=\textup{(\roman*)},leftmargin=4em]
\item a large total mass ratio can coexist with a dominant reverse
      response, in which case the fan closes immediately;
\item failure of reverse dominance means the target's actual
      \(qv\)-incidence mass is small, leaving a large off-edge stock
      which must be acquired in the complete packet.
\end{enumerate}
Similarly, the V2 low-response branch is not an amorphous self-bank.
It is a sum of a small \(q\)-sensitive correction and an exact common
factor times a \(q\)-free role difference.  The latter has at most two
cover increments, but those increments cannot reconnect \(q\).
The next decisive object is therefore the first \(q\)-incident edge
in the common factor, together with the off-edge acquisition at the
target row.
\end{assessment}

\begin{theorem}[Third-series V5 result ledger]
\label{thm:V5-results}
Relative to V1--V4 and the frozen archives, V5 proves:
\begin{enumerate}[label=\textup{(V5.\arabic*)},leftmargin=5.8em]
\item the exact direct reverse-response criterion
      \eqref{eq:V5-direct-fan}, which strictly improves the
      mass-comparison formulation;
\item a clean family showing that no support-uniform reverse response
      follows from forward dominance alone;
\item the incidence inequality \eqref{eq:V5-stock-incidence};
\item the reverse-or-deplete dichotomy
      \Cref{thm:V5-reverse-or-deplete};
\item a two-cover theorem for every role change after one row has
      been removed;
\item exact localization of all possible \(q\)-reconnecting
      incidences in the V2 insensitive tensor term;
\item conditional fan closure from any such common-factor response;
      and
\item a factorized clean counterpacket proving that the V2 scalar
      low-response inequalities alone do not create a missing-row
      edge.
\end{enumerate}
\end{theorem}

\begin{proof}
These are
\Cref{thm:V5-direct-reverse,
prop:V5-no-uniform-reverse,
lem:V5-stock-dominates-incidence,
thm:V5-reverse-or-deplete,
lem:V5-two-cover-lower-row,
prop:V5-insensitive-incidence,
cor:V5-common-factor-closure,
prop:V5-low-response-no-edge}.
\end{proof}

\begin{programme}[V6: targeted off-edge acquisition]
\label{prog:V5-V6}
\begin{enumerate}[leftmargin=3em]
\item Fix a failed reverse orientation \(v\to q\) and remove the
      literal incidence subbank \(E_{qv}\) whose \(v\)-mass is small
      by \eqref{eq:V5-target-depleted}.
\item Write the exact depleted identity for the remaining
      \((1-\theta)\Xi_v\) stock, preserving every bank already present
      in the assembled word and the old sole payer.
\item Partition each new directed or anchored outcome by whether its
      actual partner is \(q\).  A \(v\to q\) outcome closes by
      \Cref{thm:V5-direct-reverse}; retain all avoid-\(q\) outcomes as
      a finite labelled table.
\item In the V2 insensitive term, subdivide the common factor
      \(M_{S_q}(\rho)\otimes A_{F_q}^q\) by its first actual
      \(q\)-incidence before applying an operator norm.
\item Test whether the \(q\)-private part of that common factor is
      annihilated by connected-cumulant subtraction.  Prove the
      annihilation from the exact V61 word if true; do not infer it
      from the scalar example.
\item Keep the at-most-two lower-row cover increments attached to the
      common factor and verify every payer location.
\item Re-test the V78 nonlocalization example and the V84 spin-one
      example against any proposed targeted compiler.
\end{enumerate}
\end{programme}

\begin{firewall}[Claim boundary after third-series V5]
\label{fire:V5-boundary}
V5 replaces the coarse mass-ascent branch by an exact
reverse-or-deplete dichotomy and exposes the lower-row recursion in
the V2 insensitive term.  It does not yet assemble the large
off-edge target stock with the original pair and payer, prove
connected-cumulant annihilation of the \(q\)-private common factor, or
close the remaining \(q\)-sensitive, diagonal, nonclean, and other
quartic source families.  Higher MTA, WUFC, CPM, cutoff removal,
reflection positivity, Osterwalder--Schrader reconstruction, and a
uniform positive continuum spectral gap remain open.  The full
Yang--Mills mass gap has not been proved.
\end{firewall}

\paragraph{Parent-version record.}
V5 was copied byte-for-byte from
\texttt{Renewal\_Yang\_Mills\_Mass\_Gap\_Research\_Notes\_V4.tex}
before this chapter was appended.  The V4 parent had \(2\,700\) source
lines, \(98\,637\) bytes, and SHA--256
\[
 \texttt{8163303dd7c0167a0e3dca82a5c7e28abca77123f406926a768e428d400c4901}.
\]
The compulsory V5 companion audit is recorded in
\Cref{sec:V5-audit}; the next required audit is V10.

\chapter{V6: Targeted off-edge acquisition and the common-factor
incidence split}
\label{ch:V6-targeted}

\section{Context and exact target}
\label{sec:V6-context}

V5 replaced the coarse strict-ascent label by the following sharper
event.  A forward response \(q\to v\) has been found at the unique
missing complement \(q\), but its reverse \(v\to q\) is below a fixed
dominance threshold.  The literal incidence subbank \(E_{qv}\) then
carries only a small fraction of row \(v\)'s mass.  The next
acquisition must use the complementary stock, while preserving the
old pair \(ab\), the existing response \(a\to z\), and the sole payer.

This chapter proves two finite facts needed before that operator
assembly can be attempted.  First, it gives the complete graph and
denominator table for a new response rooted at \(v\): among its three
possible partners, only the targeted partner \(q\) closes within the
two-response quartic atlas.  Second, it splits the common contraction
in the V2 insensitive term into a private \(q\)-part and at most three
same-block \(q\)-incidence parts.  A large nonprivate part always
supplies a dominant scalar response candidate rooted at \(q\).

\section{Exact off-edge mass identity}
\label{sec:V6-off-edge}

Let
\[
 E:=E_{qv}(B_q)
\]
be the incidence subbank from \eqref{eq:V5-incidence-subbank}, and
write \(E^c\) for its complement in the complete coordinate support
of row \(v\).

\begin{lemma}[Complementary stock after failed reversal]
\label{lem:V6-complementary-stock}
If
\[
 d_{v\to q}(B_q)<\theta a_*,
 \qquad0<\theta<1,
\]
then
\begin{equation}
 \sum_{e\in E^c}a_{v,e}
 >
 (1-\theta)\Xi_v.
\label{eq:V6-large-complement}
\end{equation}
Moreover, coordinate factorization splits every complete bank
meeting both \(E\) and \(E^c\) into its \(E\)- and \(E^c\)-tensor
factors before any norm or scalar response mark is taken.
\end{lemma}

\begin{proof}
\Cref{thm:V5-reverse-or-deplete} gives
\[
 \sum_{e\in E}a_{v,e}<\theta\Xi_v.
\]
Subtract this strict inequality from the exact definition
\(\Xi_v=\sum_ea_{v,e}\) to obtain
\eqref{eq:V6-large-complement}.  The second statement is the product
definition of a complete coordinate bank:
\[
 M_B(\rho)
 =
 M_{B\cap E}(\rho)\otimes M_{B\cap E^c}(\rho).
\]
This is an identity before an operator norm is applied; it does not
create a second occurrence of \(B\).
\end{proof}

\begin{firewall}[A tensor split is not two capacity payments]
\label{fire:V6-no-double-capacity}
The two factors in \Cref{lem:V6-complementary-stock} are pieces of one
physical bank occurrence.  V6 uses the split only to locate row mass
and actual partner cells.  It does not assign independent first
moments, resolvents, or payers to the two pieces.
\end{firewall}

\section{The targeted partner table}
\label{sec:V6-target-table}

Retain
\[
 S:=a\to z,\qquad F:=q\to v,
\]
where \(S\) is the existing endpoint response and \(F\) the forward
response which failed direct reversal.  Let a new off-edge
acquisition at \(v\) yield
\[
 T:=v\to r,\qquad r\ne v.
\]
If \(r=q\), \(T\) is exactly the desired reverse orientation and V5
closes the fan.  The remaining two choices satisfy
\[
 r\in\cV\setminus\{v,q\}.
\label{eq:V6-avoid-q}
\]

\begin{proposition}[Complete avoid-\(q\) two-response table]
\label{prop:V6-avoid-table}
Assume \eqref{eq:V6-avoid-q}.  Select any two of the three response
factors \(S,F,T\), retain the raw edge \(ab\), and test the product by
\Cref{prop:V3-two-response-classification}.  No selection is an exact
spanning quartic tree response.  The exhaustive table is:
\begin{center}
\small
\begin{tabular}{c c c c c}
\toprule
\(v\) & \(r\) & \(S+F\) & \(S+T\) & \(F+T\)\\
\midrule
\(a\) & \(b\) & star, wrong tails & repeated \(ab\) &
repeated \(ab\)\\
\(a\) & \(z\) & star, wrong tails & repeated \(az\) &
star, wrong tails\\
\(b\) & \(a\) & path, wrong tail & repeated \(ab\) &
repeated \(ab\)\\
\(b\) & \(z\) & path, wrong tail & triangle, \(q\) isolated &
star, wrong tails\\
\(z\) & \(a\) & path, wrong tail & repeated \(az\) &
path, wrong tail\\
\(z\) & \(b\) & path, wrong tail & triangle, \(q\) isolated &
path, wrong tail\\
\bottomrule
\end{tabular}
\end{center}
Here a wrong-tail entry means that the literal edges form a spanning
tree but the two response denominators are not located at the
degree-excess vertices of that tree.
\end{proposition}

\begin{proof}
For \(S+F\), the three edges are
\[
 \{ab,az,qv\}
\]
with response tails \(a,q\).  If \(v=a\), this is the star centered
at \(a\), whose required tail multiset is \(\{a,a\}\).  If \(v=b\),
it is the path \(z-a-b-q\), whose required tails are \(\{a,b\}\).
If \(v=z\), it is the path \(b-a-z-q\), whose required tails are
\(\{a,z\}\).  Thus all three \(S+F\) entries fail exactly at the
forward tail \(q\).

For \(S+T\), the edges are
\[
 \{ab,az,vr\}
\]
with tails \(a,v\).  When \(vr\) equals \(ab\) or \(az\), only two
distinct edges remain.  The only other possibility is \(vr=bz\);
then the three edges form the triangle on \(\{a,b,z\}\), leaving
\(q\) isolated.

For \(F+T\), the edges are
\[
 \{ab,qv,vr\}
\]
with tails \(q,v\).  When \(vr=ab\), there is a repeated raw edge.
The remaining cases are:
\[
\begin{array}{c|c|c}
(v,r)&\text{tree}&\text{required tails}\\
\hline
(a,z)& b-a-\{q,z\}\text{ star}&\{a,a\}\\
(b,z)& a-b-\{q,z\}\text{ star}&\{b,b\}\\
(z,a)& q-z-a-b\text{ path}&\{z,a\}\\
(z,b)& q-z-b-a\text{ path}&\{z,b\}.
\end{array}
\]
Their actual tails are respectively \(\{q,a\}\),
\(\{q,b\}\), \(\{q,z\}\), and \(\{q,z\}\).  Each fails the exact
degree multiset.  This proves every table entry.
\end{proof}

\begin{corollary}[The target \(q\) is necessary in the two-response
atlas]
\label{cor:V6-target-necessary}
For a new response rooted at the failed-reversal target \(v\), the
choice \(v\to q\) closes by
\Cref{thm:V5-direct-reverse}.  Both avoid-\(q\) choices fail every
possible selection of two available responses with the raw pair.
Thus \(q\) is the unique successful partner within the V1/V3
two-response quartic tree atlas.
\end{corollary}

\begin{proof}
The \(v\to q\) statement is
\Cref{thm:V5-direct-reverse}.  The other two partners are
\Cref{prop:V6-avoid-table}.
\end{proof}

\begin{remark}[Four edges require a new analytic numerator]
\label{rem:V6-four-edge}
The assembled word may contain all of \(S,F,T\) and the raw edge.
That four-edge graph can be connected even when every two-response
subselection fails.  But multiplying the pair majorant by three
response lower bounds produces three row denominators, whereas a
quartic tree response has denominator degree two by
\Cref{lem:V3-degree-form}.  Surplus-edge deletion is legal only after
an analytic bound with the corresponding edge-marked numerator has
been proved; graph connectivity alone cannot delete the unwanted
denominator.
\end{remark}

\section{Targeted rowwise acquisition normal form}
\label{sec:V6-targeted-acquisition}

\begin{theorem}[Off-edge targeted acquisition]
\label{thm:V6-targeted-acquisition}
Assume failed reverse dominance at threshold \(\theta\), and suppose
the large complementary stock in \eqref{eq:V6-large-complement} is
eligible for the rowwise clean incidence decomposition at \(v\).
Then one step gives:
\begin{enumerate}[label=\textup{(\roman*)},leftmargin=4em]
\item an inherited \(P,U,A,F,\) or \(I\) compiler;
\item a directed or externally anchored response \(v\to q\), closing
      the pair by \Cref{thm:V5-direct-reverse};
\item a directed or externally anchored response \(v\to r\) with
      \(r\ne v,q\), belonging to one of the six failed cells of
      \Cref{prop:V6-avoid-table};
\item a one-bank self-anchored witness on the off-edge word; or
\item a nonclean state.
\end{enumerate}
All choices are support-uniform, and the partner label in
\textup{(ii)--(iii)} is the actual coordinate partner.
\end{theorem}

\begin{proof}
Apply the proof of \Cref{lem:V4-rowwise-transport} to the coordinate
support \(E^c\), whose total \(v\)-mass is bounded below by
\eqref{eq:V6-large-complement}.  The six incidence types are
exhaustive there.  The complete \(P,U,A,F,I\) types give
\textup{(i)}.  Partition a \(D\)-cell by its at most three actual
partners; likewise retain the actual anchor in an external-overlap
response.  The partner is either \(q\), giving \textup{(ii)}, or one
of the other two rows, giving \textup{(iii)}.  Anchored overlap which
does not give an external response gives one self-bank witness.
Failure of clean typing is \textup{(v)}.  The constants depend only
on \(\theta\), the fixed six-cell split, and \(M_0\), not on support.
\end{proof}

\begin{firewall}[Eligibility remains the packet gate]
\label{fire:V6-target-eligibility}
\Cref{thm:V6-targeted-acquisition} assumes that the \(E^c\) factors
remain in the complete assembled word with the original pair
localization and payer.  The coordinate tensor identity of
\Cref{lem:V6-complementary-stock} does not by itself prove this
operator eligibility.
\end{firewall}

\section{The common-factor \(q\)-incidence partition}
\label{sec:V6-common-factor}

Consider the V2 insensitive factor
\[
 A_{F_q}^{q}
 =
 \bigotimes_{e\in F_q}M_e(X_e),
\]
where \(X_e\) is the block containing \(q\) in the common restriction
of \(\rho,\sigma\) to the active rows at \(e\).  Only coordinates at
which \(q\) is active contribute to the following row-mass ledger.
Fix the global row order and define
\begin{align}
 {\cal P}_q
 &:=
 \{e\in F_q:q\text{ is active and }X_e=\{q\}\},
\label{eq:V6-common-private}\\
 {\cal C}_{q r}
 &:=
 \{e\in F_q:q\text{ is active, }X_e\ne\{q\},\ 
 r=\min(X_e\setminus\{q\})\},
\qquad r\ne q.
\label{eq:V6-common-cells}
\end{align}
These four sets are disjoint and exhaust the \(q\)-active part of
\(F_q\).

\begin{lemma}[Exact common-factor incidence identity]
\label{lem:V6-common-identity}
Put
\[
 P_q^{\rm com}:=\sum_{e\in{\cal P}_q}a_{q,e},
 \qquad
 C_{qr}^{\rm com}:=\sum_{e\in{\cal C}_{qr}}a_{q,e}.
\]
Then
\begin{equation}
 \sum_{e\in F_q}a_{q,e}
 =
 P_q^{\rm com}+\sum_{r\ne q}C_{qr}^{\rm com}.
\label{eq:V6-common-identity}
\end{equation}
For every \(r\ne q\),
\begin{equation}
 H_{qr}({\cal C}_{qr})
 \ge a_*C_{qr}^{\rm com}.
\label{eq:V6-common-response-stock}
\end{equation}
The edge \(qr\) is a literal same-block incidence in the common
factor \(A_{F_q}^{q}\).
\end{lemma}

\begin{proof}
The least-partner rule assigns every nonprivate \(q\)-active
coordinate to exactly one of the three cells
\({\cal C}_{qr}\), proving \eqref{eq:V6-common-identity}.  On such a
cell, row \(r\) is active and the Casimir floor gives
\[
 a_{q,e}a_{r,e}\ge a_*a_{q,e}.
\]
Summation proves \eqref{eq:V6-common-response-stock}.  Since
\(r\in X_e\), both rows occur in the same common moment block, so the
edge is not a shadow edge from the \(q\)-free role difference.
\end{proof}

\begin{theorem}[Private-or-response dichotomy in the insensitive mass]
\label{thm:V6-private-or-response}
Assume the V2 low-response alternative with
\begin{equation}
 \sum_{e\in F_q}a_{q,e}\ge\frac{\beta}{2}\Xi_q.
\label{eq:V6-insensitive-large}
\end{equation}
Then at least one of:
\begin{enumerate}[label=\textup{(\roman*)},leftmargin=4em]
\item the common factor has private concentration
      \begin{equation}
      P_q^{\rm com}\ge\frac{\beta}{4}\Xi_q;
      \label{eq:V6-private-large}
      \end{equation}
\item some \(r\ne q\) gives the dominant same-block response
      candidate
      \begin{equation}
      \boxed{
      d_{q\to r}({\cal C}_{qr})
      =\frac{H_{qr}({\cal C}_{qr})}{\Xi_q}
      \ge\frac{\beta}{12}a_*.}
      \label{eq:V6-common-dominant}
      \end{equation}
\end{enumerate}
\end{theorem}

\begin{proof}
If \eqref{eq:V6-private-large} fails, subtract it from
\eqref{eq:V6-insensitive-large} using
\eqref{eq:V6-common-identity}:
\[
 \sum_{r\ne q}C_{qr}^{\rm com}
 >
 \frac{\beta}{4}\Xi_q.
\]
One of the three partner cells has mass at least
\(\beta\Xi_q/12\).  Apply
\eqref{eq:V6-common-response-stock} and divide by \(\Xi_q\).
\end{proof}

\begin{corollary}[Graph fate of a legal common-factor mark]
\label{cor:V6-common-fate}
In the singleton missing-complement state, suppose the response
candidate in \eqref{eq:V6-common-dominant} is mark-admissible in the
complete V2 word.  For any fixed \(0<\theta<1\), either:
\begin{enumerate}[label=\textup{(\alph*)},leftmargin=4em]
\item its reverse is \(\theta\)-dominant and the pair closes by
      \Cref{thm:V5-direct-reverse}; or
\item the common-factor \(qr\)-incidence bank is
      \(\theta\)-depleted at row \(r\), and the targeted acquisition
      of \Cref{thm:V6-targeted-acquisition} is the next branch.
\end{enumerate}
\end{corollary}

\begin{proof}
Apply \Cref{thm:V5-reverse-or-deplete} with \(v=r\), then use
\Cref{thm:V6-targeted-acquisition} in its second alternative.
\end{proof}

\begin{firewall}[Same-block scalar incidence versus operator marking]
\label{fire:V6-common-marking}
Equation \eqref{eq:V6-common-dominant} proves a literal scalar edge
inside the common contraction.  To call it a marked-tree edge, one
must retain the complete V2 tensor word and verify either the V15
distinct-occurrence marking hypothesis or the specific V98
joint-carrier hypothesis when banks coincide.  The archived V98
proposition permits repeated marks only after its three-occurrence
carrier is bounded; it is not a blanket license for arbitrary
same-bank products.
\end{firewall}

\begin{assessment}[The low-response branch has split]
\label{ass:V6-low-response-split}
The V2 insensitive mass can no longer be treated as a featureless
\(q\)-free correction.  Its common factor carries at least half of
the self-witness \(q\)-mass, and V6 proves that this mass is either
substantially \(q\)-private or produces a dominant same-block
\(q\)-response candidate.  The latter has exactly the V5
reverse-or-deplete fate once its operator marking is certified.  The
private alternative and the marking certification are disjoint
analytic tasks; neither can be replaced by another scalar
pigeonhole.
\end{assessment}

\section{Regression checks and result ledger}
\label{sec:V6-results}

\begin{proposition}[V6 regression suite]
\label{prop:V6-regressions}
\begin{enumerate}[label=\textup{(\roman*)},leftmargin=4em]
\item The failed-edge complement is an exact coordinate tensor split,
      not an independent use of a complete bank.
\item Every failed entry in \Cref{prop:V6-avoid-table} fails either
      literal spanning or the V3 degree-denominator criterion.
\item The forward edge \(qv\) is never silently reoriented.
\item The common-factor edge \(qr\) is taken from \(X_e\), not from
      the \(q\)-free role factor.
\item No V98 repeated-bank conclusion is used without a joint
      assembled carrier.
\item The lower-row difference still has at most two cover
      increments and remains tensor-attached to the common factor.
\end{enumerate}
\end{proposition}

\begin{proof}
Items \textup{(i)--(iii)} are
\Cref{lem:V6-complementary-stock,prop:V6-avoid-table}.
Item \textup{(iv)} is \Cref{lem:V6-common-identity}.  Item
\textup{(v)} is \Cref{fire:V6-common-marking}.  Item \textup{(vi)}
is unchanged from \Cref{lem:V5-two-cover-lower-row,
prop:V5-insensitive-incidence}.
\end{proof}

\begin{theorem}[Third-series V6 result ledger]
\label{thm:V6-results}
Relative to V1--V5 and the frozen archives, V6 proves:
\begin{enumerate}[label=\textup{(V6.\arabic*)},leftmargin=5.8em]
\item the exact large complementary mass identity after reverse
      failure;
\item the complete six-cell avoid-\(q\) response table, including
      every possible pair of the three available responses;
\item uniqueness of \(v\to q\) as the successful targeted response
      within the two-response quartic atlas;
\item the off-edge targeted acquisition normal form;
\item the exact private/three-partner partition of the V2 common
      factor;
\item the quantitative private-or-response dichotomy
      \eqref{eq:V6-private-large}--%
      \eqref{eq:V6-common-dominant}; and
\item the reverse-or-deplete fate of every operator-admissible
      common-factor response.
\end{enumerate}
\end{theorem}

\begin{proof}
These are
\Cref{lem:V6-complementary-stock,
prop:V6-avoid-table,
cor:V6-target-necessary,
thm:V6-targeted-acquisition,
lem:V6-common-identity,
thm:V6-private-or-response,
cor:V6-common-fate}.
\end{proof}

\begin{programme}[V7: certify the common-factor mark or annihilate
the private branch]
\label{prog:V6-V7}
\begin{enumerate}[leftmargin=3em]
\item Write the exact V61 operator word for
      \(M_{S_q}(\rho)\otimes A_{F_q}^q\otimes
      (R_{F_q}(\sigma)-R_{F_q}(\rho))\), including every physical
      occurrence of a cell \({\cal C}_{qr}\).
\item Separate distinct-occurrence cells, where V15 marking applies,
      from coincident cells requiring a genuine V98-style joint
      carrier.  Do not cite V98 outside its three-occurrence
      hypotheses.
\item For the private alternative \eqref{eq:V6-private-large},
      expand the connected source subtraction before norms and test
      whether the one-row common factor cancels.
\item If it does not cancel, insert the archived private complete
      compiler with the full payer ledger and show exactly how it
      remains attached to the two-cover lower-row difference.
\item For off-edge target responses, retain the six failed
      avoid-\(q\) cells rather than iterating an unnamed acquisition.
      Determine whether their four-edge word has a legal surplus-edge
      deletion after a joint numerator estimate.
\item Re-run the V78 and V84 counterexamples on the proposed complete
      word.
\end{enumerate}
\end{programme}

\begin{firewall}[Claim boundary after third-series V6]
\label{fire:V6-boundary}
V6 proves the exact targeted graph table and extracts a dominant
same-block response candidate from every nonprivate V2 insensitive
common factor.  It does not yet prove that this candidate is
mark-admissible in every coincident-bank word, close the substantial
private common factor, or produce a joint four-edge numerator for the
avoid-\(q\) cells.  The sensitive, diagonal, nonclean, and other
quartic source families remain open, as do higher MTA, WUFC, CPM,
cutoff removal, reflection positivity, Osterwalder--Schrader
reconstruction, and a uniform positive continuum spectral gap.  The
full Yang--Mills mass gap has not been proved.
\end{firewall}

\paragraph{Parent-version record.}
V6 was copied byte-for-byte from
\texttt{Renewal\_Yang\_Mills\_Mass\_Gap\_Research\_Notes\_V5.tex}
before this chapter was appended.  The V5 parent had \(3\,281\) source
lines, \(119\,499\) bytes, and SHA--256
\[
 \texttt{2311f497bdfa93df54c24ad1c4116d30bd7b94bbce9aa98d1fecfab1004a18cb}.
\]
The next compulsory companion audit remains V10.

\chapter{V7: First attachment of the missing row in the exact
connected word}
\label{ch:V7-first-attachment}

\section{Why the calculation moves back to V61}
\label{sec:V7-context}

The V6 common-factor split leaves two alternatives.  A nonprivate
common factor contains a dominant same-block response candidate,
whose operator marking still has to be certified.  A private common
factor contains substantial \(q\)-mass but no \(q\)-edge.  Treating
the latter as a terminal obstruction would ignore the exact
connected-source subtraction.

The appropriate upstream object is the V61 linked-word identity.  It
says that a coordinate word occurs in the connected four-row carrier
if and only if the join of all its local row partitions is the
one-block partition.  V7 specializes that identity to one named
missing row.  It proves that a word in which \(q\) remains a singleton
at every coordinate is annihilated exactly, and it disintegrates
every surviving word by the unique first coordinate at which \(q\)
joins another row.  This is an operator classification before any
norm or positive envelope.

\section{Imported linked-word interface}
\label{sec:V7-linked-interface}

Let \(E=\{1<\cdots<n\}\) be the ordered coordinate set and
\(I=\cV=\{a,b,z,q\}\).  For a nonempty \(B\subseteq I\), let
\(k_e(B)\) denote the joint local cumulant carrier at coordinate
\(e\).  For \(\pi\in\Pi(I)\), define
\begin{equation}
 L_{e,\pi}:=\bigotimes_{B\in\pi}k_e(B),
 \qquad
 W(\boldsymbol\pi):=\bigotimes_{e\in E}L_{e,\pi_e}.
\label{eq:V7-local-word}
\end{equation}
Write \(\widehat0,\widehat1\) for the discrete and one-block
partitions.

\begin{proposition}[V61 linked-word identity, specialized to four
rows]
\label{prop:V7-linked-word}
The connected carrier is
\begin{equation}
 \boxed{
 K_I^{E,\mathrm{conn}}
 =
 \sum_{\substack{\boldsymbol\pi\in\Pi(I)^E\\
                  \bigvee_{e\in E}\pi_e=\widehat1}}
 W(\boldsymbol\pi).}
\label{eq:V7-linked-word}
\end{equation}
Every local partition word occurs with coefficient one when its join
is connected and coefficient zero otherwise.
\end{proposition}

\begin{proof}
For a replica partition \(\tau\), the complete moment is
\[
 M_E(\tau)
 =
 \sum_{\pi_e\le\tau\ \forall e}W(\boldsymbol\pi).
\]
M\"obius inversion on \(\Pi(I)\) gives the connected carrier.  If
\(\sigma=\bigvee_e\pi_e\), the coefficient of the fixed word is
\[
 \sum_{\tau:\sigma\le\tau\le\widehat1}
 \mu(\tau,\widehat1)
 =
 \mathbf1_{\{\sigma=\widehat1\}}.
\]
This is exactly the archived V61 finite algebraic identity, included
here only as the interface for the new one-row disintegration.
\end{proof}

\section{Isolation cancellation and first attachment}
\label{sec:V7-first-attachment}

For a partition \(\pi\), write \(B_q(\pi)\) for the block containing
\(q\), and set
\[
 \Pi_q^0:=\{\pi:B_q(\pi)=\{q\}\},
 \qquad
 \Pi_q^+:=\Pi(I)\setminus\Pi_q^0.
\]

\begin{lemma}[Exact annihilation of an everywhere-private row]
\label{lem:V7-private-annihilation}
If \(\pi_e\in\Pi_q^0\) for every \(e\in E\), then
\[
 \bigvee_{e\in E}\pi_e
 \le
 \{q\}\mid(I\setminus\{q\})
 <
 \widehat1.
\]
Consequently that word has coefficient zero in
\(K_I^{E,\mathrm{conn}}\).
\end{lemma}

\begin{proof}
In every equivalence relation \(\pi_e\), the equivalence class of
\(q\) is the singleton.  The equivalence relation generated by their
union therefore contains no pair \(q\sim r\) with \(r\ne q\).  This
generated relation is their join.  It is bounded by the displayed
two-block partition and cannot be \(\widehat1\).  Apply
\Cref{prop:V7-linked-word}.
\end{proof}

\begin{definition}[First \(q\)-attachment data]
\label{def:V7-attachment-data}
For a connected word \(\boldsymbol\pi\), define
\begin{align}
 \tau_q(\boldsymbol\pi)
 &:=
 \min\{e\in E:\pi_e\in\Pi_q^+\},
\label{eq:V7-attachment-time}\\
 k_q(\boldsymbol\pi)
 &:=
 |B_q(\pi_{\tau_q})|,
\label{eq:V7-attachment-size}\\
 r_q(\boldsymbol\pi)
 &:=
 \min\bigl(B_q(\pi_{\tau_q})\setminus\{q\}\bigr),
\label{eq:V7-attachment-partner}
\end{align}
where the minimum uses the fixed global row order.
\end{definition}

\begin{lemma}[The attachment data exist and are unique]
\label{lem:V7-attachment-unique}
Every word in \eqref{eq:V7-linked-word} has unique data
\[
 \tau_q\in E,\qquad
 k_q\in\{2,3,4\},\qquad
 r_q\in I\setminus\{q\}.
\]
\end{lemma}

\begin{proof}
If the defining set in \eqref{eq:V7-attachment-time} were empty,
\Cref{lem:V7-private-annihilation} would contradict connectedness.
Its minimum is unique because \(E\) is ordered.  The block at that
coordinate contains \(q\) and at least one of the three other rows,
so its size is \(2,3,\) or \(4\).  The least other row is unique.
\end{proof}

For \(s\in E\), \(r\ne q\), and \(k\in\{2,3,4\}\), define the exact
subcarrier
\begin{equation}
 {\cal A}_{s,r,k}^{q}
 :=
 \sum_{\substack{\boldsymbol\pi:\ \vee_e\pi_e=\widehat1\\
                  \tau_q(\boldsymbol\pi)=s\\
                  r_q(\boldsymbol\pi)=r\\
                  k_q(\boldsymbol\pi)=k}}
 W(\boldsymbol\pi).
\label{eq:V7-attachment-packet}
\end{equation}

\begin{theorem}[Exact first-\(q\)-attachment disintegration]
\label{thm:V7-attachment-disintegration}
The connected carrier has the disjoint finite decomposition
\begin{equation}
 \boxed{
 K_I^{E,\mathrm{conn}}
 =
 \sum_{s\in E}
 \sum_{r\ne q}
 \sum_{k=2}^4
 {\cal A}_{s,r,k}^{q}.}
\label{eq:V7-attachment-disintegration}
\end{equation}
There is no additional choice of a witnessing coordinate or a
witnessing partner inside any summand.
\end{theorem}

\begin{proof}
By \Cref{lem:V7-attachment-unique}, every connected word has exactly
one triple \((s,r,k)\).  The constraints in
\eqref{eq:V7-attachment-packet} therefore partition the index set of
\eqref{eq:V7-linked-word} into disjoint classes.  Summing those
classes gives \eqref{eq:V7-attachment-disintegration}, with the
coefficient one inherited from the linked-word identity.
\end{proof}

\begin{remark}[No analytic estimate has yet been taken]
\label{rem:V7-no-analytic-sum}
Uniqueness prevents combinatorial overcounting, but it does not by
itself bound the norm of the sum over \(s\).  The source coordinate,
the first attachment coordinate, and the payer may be correlated.
An analytic estimate must retain that correlation, exactly as V61
retains the complete suffix before contraction.
\end{remark}

\section{Exporting attachment from a private common factor}
\label{sec:V7-private-export}

\begin{proposition}[Private-coordinate export]
\label{prop:V7-private-export}
Let \(P\subseteq E\) be a coordinate set such that
\[
 B_q(\pi_e)=\{q\}
 \qquad(e\in P)
\label{eq:V7-private-set}
\]
for every word under consideration.  In every connected word, the
first \(q\)-attachment lies in \(E\setminus P\).  If
\eqref{eq:V7-private-set} also holds on \(E\setminus P\), the full
subcarrier vanishes identically.
\end{proposition}

\begin{proof}
No coordinate in \(P\) belongs to the defining set of
\(\tau_q\), so the minimum lies in its complement.  If the condition
also holds on the complement, \(q\) is singleton at every coordinate
and \Cref{lem:V7-private-annihilation} applies.
\end{proof}

\begin{corollary}[Fate of the V6 private alternative]
\label{cor:V7-private-fate}
In the V6 alternative
\[
 P_q^{\rm com}\ge\frac{\beta}{4}\Xi_q,
\]
the coordinates in \({\cal P}_q\) cannot themselves attach \(q\) in
the V2 common factor.  Every surviving connected word must attach
\(q\) at a coordinate outside \({\cal P}_q\).  If all remaining
factors also keep \(q\) singleton, that word cancels exactly from the
connected carrier.
\end{corollary}

\begin{proof}
For \(e\in{\cal P}_q\), the defining property
\(X_e=\{q\}\) says precisely that the common local partition block of
\(q\) is singleton.  Apply
\Cref{prop:V7-private-export}.  The size of the private mass is not
needed for the algebraic conclusion; it is retained for the later
capacity estimate.
\end{proof}

\begin{firewall}[Export is not contraction]
\label{fire:V7-export-not-contraction}
\Cref{cor:V7-private-fate} does not delete the one-row moment factors
on \({\cal P}_q\).  They remain tensor factors of the word until a
complete-moment or prefix-factor contraction is applied with the
correct payer location.  The result relocates the necessary
attachment; it does not spend the private factor a second time.
\end{firewall}

\section{Binary, ternary, and quartic attachment carriers}
\label{sec:V7-size-trichotomy}

\begin{theorem}[First-attachment size trichotomy]
\label{thm:V7-size-trichotomy}
In a packet \({\cal A}_{s,r,k}^{q}\), let
\[
 B:=B_q(\pi_s)
\]
be the first local block attaching \(q\).  Then:
\begin{enumerate}[label=\textup{(\roman*)},leftmargin=4em]
\item if \(k=2\), the attaching factor is the binary cumulant
      \(k_s(\{q,r\})\), a literal \(qr\) covariance carrier;
\item if \(k=3\), the attaching factor is a ternary connected
      carrier on \(B\); every spanning tree produced by the
      established ternary MTA contains at least one literal edge
      incident to \(q\);
\item if \(k=4\), the attaching factor is the local connected
      quartic carrier \(k_s(I)\), so using the desired quartic MTA to
      estimate it would be circular.
\end{enumerate}
\end{theorem}

\begin{proof}
By \eqref{eq:V7-local-word}, the block \(B\in\pi_s\) contributes
exactly \(k_s(B)\).  If \(|B|=2\), the only other row in \(B\) is the
least partner \(r\), giving \textup{(i)}.  If \(|B|=3\), the archived
ternary MTA expresses the normalized connected carrier as a finite
sum of the three possible two-edge spanning trees on \(B\).  Every
tree on \(B\) contains \(q\) and, being connected, has at least one
edge incident to \(q\).  If \(|B|=4\), then \(B=I\), giving the
quartic local cumulant in \textup{(iii)}.
\end{proof}

\begin{corollary}[Only the size-four attachment is order-circular]
\label{cor:V7-only-quartic-circular}
The \(k=2\) and \(k=3\) first-attachment packets require only binary
and ternary connected estimates before being combined with the
remaining word.  The \(k=4\) packet is the unique attachment-size
class which cannot be bounded by induction on row order.
\end{corollary}

\begin{proof}
Binary and ternary MTA precede quartic MTA in the established
induction.  The final statement is
\Cref{thm:V7-size-trichotomy}\textup{(iii)}.
\end{proof}

\section{Quantitative refinement of the V6 common response}
\label{sec:V7-common-refinement}

In the nonprivate V6 alternative, fix the partner \(r\) for which
\[
 H_{qr}({\cal C}_{qr})
 \ge\frac{\beta}{12}a_*\Xi_q.
\]
Partition its coordinates by the common block size:
\begin{equation}
 {\cal C}_{qr}^{(k)}
 :=
 \{e\in{\cal C}_{qr}:|X_e|=k\},
 \qquad k=2,3,4.
\label{eq:V7-size-cells}
\end{equation}

\begin{lemma}[Dominant size cell]
\label{lem:V7-dominant-size}
For some \(k\in\{2,3,4\}\),
\begin{equation}
 \frac{H_{qr}({\cal C}_{qr}^{(k)})}{\Xi_q}
 \ge\frac{\beta}{36}a_*.
\label{eq:V7-size-dominance}
\end{equation}
\end{lemma}

\begin{proof}
The three cells in \eqref{eq:V7-size-cells} partition
\({\cal C}_{qr}\), and \(H_{qr}\) is additive over disjoint coordinate
sets.  One cell carries at least one third of the total stock.  Divide
the preceding V6 lower bound by three and by \(\Xi_q\).
\end{proof}

\begin{theorem}[Quantitative common-factor attachment trichotomy]
\label{thm:V7-common-trichotomy}
Every nonprivate V2 insensitive common factor has one of the
following dominant subcells:
\begin{enumerate}[label=\textup{(\alph*)},leftmargin=4em]
\item a binary \(qr\) covariance cell satisfying
      \eqref{eq:V7-size-dominance};
\item a ternary same-block cell on a triple containing \(q,r\), with
      the same stock lower bound and a lower-order tree compiler;
\item a local four-row same-block cell satisfying the same lower
      bound, belonging to the \(k=4\) circular attachment class.
\end{enumerate}
Together with \Cref{cor:V7-private-fate}, this is exhaustive for the
large insensitive \(q\)-mass.
\end{theorem}

\begin{proof}
\Cref{thm:V6-private-or-response} gives either the private case or a
partner cell \({\cal C}_{qr}\).  In the latter case apply
\Cref{lem:V7-dominant-size}.  The three possible block sizes have the
operator meanings in \Cref{thm:V7-size-trichotomy}.
\end{proof}

\begin{remark}[Binary stock is still attached to the source packet]
\label{rem:V7-binary-packet}
In alternative \textup{(a)}, the \(qr\) edge is not inferred from a
shadow factor: it is the binary local cumulant at a coordinate where
the common \(q\)-block is exactly \(\{q,r\}\).  Nevertheless its sum
over source and payer positions must be bounded as part of the
complete first-attachment packet.  Only after that joint bound may
the lower response \eqref{eq:V7-size-dominance} enter the V5 fan.
\end{remark}

\section{Exact first-attachment source gate}
\label{sec:V7-source-gate}

\begin{definition}[Normalized first-attachment packet capacity]
\label{def:V7-attachment-capacity}
Let \(\mathfrak A_{r,k}^{q,\varepsilon}\) denote the normalized
positive capacity obtained by summing
\({\cal A}_{s,r,k}^{q}\) over \(s\), after the exact source
decomposition and with \(\varepsilon\in\{0,1\}\) recording whether
the sole payer is absent or assigned at its physical location.
Complete nonpayer suffixes are summed before contraction.
\end{definition}

\begin{proposition}[Sufficient one-row attachment criterion]
\label{prop:V7-attachment-criterion}
Suppose, for each \(r\ne q\) and \(k=2,3\), that
\begin{equation}
 \mathfrak A_{r,k}^{q,\varepsilon}
 \le
 C_{r,k}\sum_{T\in\mathfrak T_4:\,q\in V(T)}
 {\cal W}_T
\label{eq:V7-attachment-criterion}
\end{equation}
uniformly in support and at every payer location, with the attaching
binary or ternary carrier retained until the tree marks are formed.
Then all noncircular first-\(q\)-attachment packets are controlled by
a finite quartic tree sum.  The only first-attachment source left is
\(k=4\).
\end{proposition}

\begin{proof}
There are three choices of \(r\) and two noncircular block sizes.
Sum \eqref{eq:V7-attachment-criterion} over those six types.  The
unique disintegration \eqref{eq:V7-attachment-disintegration}
introduces no further routing multiplicity.  By
\Cref{cor:V7-only-quartic-circular}, the omitted size class is exactly
\(k=4\).
\end{proof}

\begin{firewall}[The criterion is not its proof]
\label{fire:V7-criterion-open}
V7 has not proved \eqref{eq:V7-attachment-criterion}.  In particular,
uniqueness of the first attachment does not bound the sum over its
coordinate, and a suffix payer cannot be replaced by a complete
moment.  The criterion identifies the precise noncircular operator
estimate needed to certify the V6 common-factor mark.
\end{firewall}

\section{Result ledger and V8 mandate}
\label{sec:V7-results}

\begin{theorem}[Third-series V7 result ledger]
\label{thm:V7-results}
Relative to V1--V6 and the frozen archives, V7 proves:
\begin{enumerate}[label=\textup{(V7.\arabic*)},leftmargin=5.8em]
\item exact annihilation of every linked word in which the missing row
      remains private at all coordinates;
\item unique disintegration of every connected word by the first
      \(q\)-attachment coordinate, least partner, and block size;
\item export of the required attachment away from every V6-private
      common-factor coordinate;
\item the binary/ternary/quartic first-attachment trichotomy, with
      only the quartic class circular in row order;
\item a quantitative factor-three refinement of every dominant V6
      common response by its actual block size; and
\item the exact six-type noncircular source criterion
      \eqref{eq:V7-attachment-criterion}.
\end{enumerate}
\end{theorem}

\begin{proof}
These are
\Cref{lem:V7-private-annihilation,
thm:V7-attachment-disintegration,
cor:V7-private-fate,
thm:V7-size-trichotomy,
lem:V7-dominant-size,
prop:V7-attachment-criterion}.
\end{proof}

\begin{programme}[V8: bound the binary/ternary first-attachment
packet]
\label{prog:V7-V8}
\begin{enumerate}[leftmargin=3em]
\item Refine \({\cal A}_{s,r,k}^{q}\) simultaneously by the unique
      first global-connectivity coordinate from V61.  Compare the two
      stopping coordinates without summing either absolutely.
\item If the \(q\)-attachment precedes global connectivity, place it
      in the exact prefix-join factor and use binary/ternary MTA before
      the joining letter is estimated.
\item If it is the first global joining letter, retain its full local
      block and use the two-junction payer allocation.
\item If it lies in the complete suffix, keep the suffix payer case
      separate; nonpayer complete moments may contract only after
      their full partition sums are restored.
\item Prove \eqref{eq:V7-attachment-criterion} first for \(k=2\).
      Then attach the \(k=3\) ternary tree and run the finite surplus
      edge table.
\item Isolate \(k=4\) as a local quartic source family and compare it
      with the already separated diagonal and nonclean gates, avoiding
      any invocation of quartic MTA on itself.
\end{enumerate}
\end{programme}

\begin{firewall}[Claim boundary after third-series V7]
\label{fire:V7-boundary}
V7 resolves the logical fate of an everywhere-private missing row
and reduces every surviving common-factor branch to a unique first
attachment.  It does not prove the uniform capacity estimate for the
sum of binary or ternary attachment positions, nor close the local
quartic attachment, sensitive, diagonal, nonclean, and other quartic
source families.  Higher MTA, WUFC, CPM, cutoff removal, reflection
positivity, Osterwalder--Schrader reconstruction, and a uniform
positive continuum spectral gap remain open.  The full Yang--Mills
mass gap has not been proved.
\end{firewall}

\paragraph{Parent-version record.}
V7 was copied byte-for-byte from
\texttt{Renewal\_Yang\_Mills\_Mass\_Gap\_Research\_Notes\_V6.tex}
before this chapter was appended.  The V6 parent had \(3\,793\) source
lines, \(138\,151\) bytes, and SHA--256
\[
 \texttt{56d1245d4dd7404078d631bb190fa7c0b992ab415509ca70aa3c8e9eb88e0643}.
\]
The next compulsory companion audit remains V10.

\chapter{V8: Ordering first row attachment against first global
connectivity}
\label{ch:V8-two-stops}

\section{Correction to the V7 acquisition order}
\label{sec:V8-correction}

V7 proposed three relative locations for the first \(q\)-attachment:
inside the prefix before global connectivity, at the first global
joining coordinate, or in the complete suffix.  The third case is
impossible.  A globally connected prefix must already have attached
every row, including \(q\).  This simple ordering fact removes an
entire suffix branch and, more importantly, shows that every
pre-source \(q\)-attachment is already contained in a lower-row
connected prefix factor.

There is also a necessary qualification to the V7 circularity
statement.  A binary or ternary \emph{first attachment} is
noncircular as a carrier, but the later first-global-connectivity
letter can still be a four-row local cumulant.  V8 keeps attachment
order and source order separate.

\section{The two stopping coordinates}
\label{sec:V8-stops}

For a connected word
\(\boldsymbol\pi=(\pi_e)_{e=1}^n\), put
\begin{equation}
 j_t(\boldsymbol\pi):=\bigvee_{e\le t}\pi_e,
 \qquad
 g(\boldsymbol\pi):=
 \min\{t:j_t(\boldsymbol\pi)=\widehat1\}.
\label{eq:V8-global-time}
\end{equation}
The time \(\tau_q(\boldsymbol\pi)\) is defined in
\eqref{eq:V7-attachment-time}.

\begin{lemma}[Attachment precedes global connectivity]
\label{lem:V8-order}
Every connected word satisfies
\begin{equation}
 \boxed{\tau_q(\boldsymbol\pi)\le g(\boldsymbol\pi).}
\label{eq:V8-stop-order}
\end{equation}
In particular, the first \(q\)-attachment never lies in the complete
suffix after \(g\).
\end{lemma}

\begin{proof}
At time \(g\), the cumulative join is \(\widehat1\), so its block
containing \(q\) contains all four rows.  If every local partition up
to \(g\) kept \(q\) singleton, their join would keep \(q\) singleton
by the proof of \Cref{lem:V7-private-annihilation}.  Hence some
coordinate \(e\le g\) attaches \(q\), and the first such coordinate
satisfies \(\tau_q\le g\).
\end{proof}

\begin{corollary}[Exact two-case split]
\label{cor:V8-two-cases}
The disintegration \eqref{eq:V7-attachment-disintegration} has only
two timing classes:
\[
 \tau_q<g
 \qquad\text{or}\qquad
 \tau_q=g.
\]
These classes are disjoint and exhaust the connected carrier.
\end{corollary}

\begin{proof}
\Cref{lem:V8-order} rules out \(\tau_q>g\); equality or strict
inequality are the remaining possibilities.
\end{proof}

\section{First-connectivity factorization}
\label{sec:V8-first-connectivity}

For a proper prefix join \(\rho<\widehat1\), define
\begin{align}
 C_{<g}(\rho)
 &:=
 \sum_{\substack{(\pi_e)_{e<g}\\
                  \vee_{e<g}\pi_e=\rho}}
 \bigotimes_{e<g}L_{e,\pi_e},
\label{eq:V8-prefix-carrier}\\
 J_g(\rho)
 &:=
 \sum_{\substack{\pi\in\Pi(I)\\
                  \rho\vee\pi=\widehat1}}
 L_{g,\pi},
\label{eq:V8-joining-letter}\\
 M_{>g}
 &:=
 \bigotimes_{e>g}\sum_{\pi\in\Pi(I)}L_{e,\pi}.
\label{eq:V8-complete-suffix}
\end{align}

\begin{proposition}[V61 first-connectivity interface]
\label{prop:V8-first-connectivity}
The linked carrier has the exact decomposition
\begin{equation}
 K_I^{E,\mathrm{conn}}
 =
 \sum_{g=1}^n\sum_{\rho<\widehat1}
 C_{<g}(\rho)\otimes J_g(\rho)\otimes M_{>g}.
\label{eq:V8-first-connectivity}
\end{equation}
Moreover,
\begin{equation}
 C_{<g}(\rho)
 =
 \bigotimes_{B\in\rho}K_B^{<g,\mathrm{conn}},
\label{eq:V8-prefix-factorization}
\end{equation}
where a singleton block contributes its complete one-row prefix
moment.
\end{proposition}

\begin{proof}
Every word in \eqref{eq:V7-linked-word} has a unique first time \(g\)
at which its cumulative join becomes \(\widehat1\); its preceding join
is a unique \(\rho<\widehat1\).  The suffix is unrestricted after
global connectivity and hence sums to the complete moments in
\eqref{eq:V8-complete-suffix}.  This proves
\eqref{eq:V8-first-connectivity}.

For \eqref{eq:V8-prefix-factorization}, every local partition in a
word with join \(\rho\) refines \(\rho\).  Restrict the word
independently to each block \(B\in\rho\).  Its restriction has
one-block join on every \(B\), and the choices on distinct blocks
factor.  Applying the linked-word identity on each block gives the
displayed tensor product.
\end{proof}

\section{Strict-prefix attachment is already lower order}
\label{sec:V8-prefix-attachment}

\begin{theorem}[Absorption of the strict-prefix attachment time]
\label{thm:V8-prefix-absorption}
Let \(\boldsymbol\pi\) be a connected word with
\(\tau_q<g\), and put
\[
 \rho=j_{g-1}(\boldsymbol\pi).
\]
Then the block \(B_q(\rho)\) containing \(q\) has size \(2\) or \(3\),
and all choices of the internal first-attachment coordinate
\(\tau_q<g\) are already summed inside the single lower-order factor
\begin{equation}
 K_{B_q(\rho)}^{<g,\mathrm{conn}}
\label{eq:V8-q-prefix-factor}
\end{equation}
of \eqref{eq:V8-prefix-factorization}.  No extra sum over a
\(q\)-attachment witness may be added to the source packet.
\end{theorem}

\begin{proof}
Because \(\tau_q<g\), some local partition in the prefix joins \(q\)
to another row, so \(B_q(\rho)\) has at least two elements.  Since
\(g\) is the first global-connectivity time,
\(\rho=j_{g-1}<\widehat1\); hence no block of \(\rho\) can be all four
rows.  Thus \(|B_q(\rho)|\in\{2,3\}\).

The exact prefix factorization
\eqref{eq:V8-prefix-factorization} contains the connected carrier on
that block.  By its linked-word definition, this carrier sums every
prefix word whose restriction joins \(B_q(\rho)\), including every
possible first coordinate at which \(q\) attaches inside that block.
Selecting such a coordinate again would refine an already completed
sum and introduce a false witness multiplicity.
\end{proof}

\begin{corollary}[No new support sum in the prefix branch]
\label{cor:V8-prefix-no-sum}
Assuming the established binary and ternary MTA bounds, the
\(q\)-attachment part of every \(\tau_q<g\) prefix is controlled
support-uniformly before the joining letter is estimated.  The
remaining analytic problem is the ordinary V61 source packet:
combining the fixed prefix join, the joining letter, and the sole
payer without losing the sum over \(g\).
\end{corollary}

\begin{proof}
\Cref{thm:V8-prefix-absorption} makes the \(q\)-block factor binary or
ternary.  Apply the corresponding lower-order MTA to
\eqref{eq:V8-q-prefix-factor}.  The other prefix blocks are also
proper and hence lower order.  The warning about \(g\) is exactly the
source-summation gate left by
\eqref{eq:V8-first-connectivity}.
\end{proof}

\begin{firewall}[A lower-order prefix does not lower the joining
letter]
\label{fire:V8-joining-still-quartic}
Even when \(\tau_q<g\), the joining partition at \(g\) may be
\(\widehat1\), so \(L_{g,\widehat1}=k_g(I)\) is a four-row local
cumulant.  Prefix absorption solves the missing-row attachment inside
the prefix; it does not estimate an order-four joining source by
induction.
\end{firewall}

\section{Synchronous attachment-source pairs}
\label{sec:V8-synchronous}

Suppose now that \(\tau_q=g\).  Then all prefix partitions keep
\(q\) singleton, so
\begin{equation}
 \rho=j_{g-1}
 =
 \{q\}\mid\sigma
\label{eq:V8-q-singleton-prefix}
\end{equation}
for a unique partition \(\sigma\) of the other three rows
\(W=I\setminus\{q\}\).  The joining letter \(\pi=\pi_g\) must satisfy
\(\rho\vee\pi=\widehat1\) and \(B_q(\pi)\ne\{q\}\).

\begin{proposition}[Complete synchronous pair count]
\label{prop:V8-synchronous-count}
There are exactly \(26\) labelled pairs \((\rho,\pi)\) satisfying
\eqref{eq:V8-q-singleton-prefix} and
\(\rho\vee\pi=\widehat1\).  Classified by the number of blocks of
\(\rho\) and the attachment size \(k=|B_q(\pi)|\), their counts are:
\begin{center}
\begin{tabular}{c c c c c}
\toprule
\(|\rho|\) & number of such \(\rho\) & \(k=2\) & \(k=3\) & \(k=4\)\\
\midrule
\(4\) & \(1\) & \(0\) & \(0\) & \(1\)\\
\(3\) & \(3\) & \(6\) & \(6\) & \(3\)\\
\(2\) & \(1\) & \(6\) & \(3\) & \(1\)\\
\midrule
total & \(5\) & \(12\) & \(9\) & \(5\)\\
\bottomrule
\end{tabular}
\end{center}
\end{proposition}

\begin{proof}
There are five partitions \(\sigma\) of the three-row set \(W\).

If \(\sigma\) is discrete, then \(\rho=\widehat0\), so
\(\rho\vee\pi=\pi=\widehat1\).  This gives the single
\((|\rho|,k)=(4,4)\) pair.

If \(\sigma=uv\mid w\), there are three choices of its paired block.
For a fixed choice, a size-two \(q\)-block must be \(qu\) with the
remaining local block \(vw\), or \(qv\) with remaining block \(uw\):
two possibilities.  A size-three \(q\)-block must contain \(q,w\)
and exactly one of \(u,v\), again two possibilities.  The one-block
partition gives one size-four possibility.  Thus each of the three
prefixes contributes \(2,2,1\).

If \(\sigma=W\), the prefix has two blocks \(W\mid q\).  For a
size-two \(q\)-block, choose its other row in three ways and partition
the two unused rows either together or separately, giving \(3\cdot
2=6\).  A size-three \(q\)-block is determined by the omitted row,
giving three choices.  The one-block partition gives one choice.
Adding the three cases yields the displayed table and total
\[
 1+3(2+2+1)+(6+3+1)=26.
\]
\end{proof}

\begin{corollary}[Twenty-one nonquartic synchronous sources]
\label{cor:V8-twenty-one}
Of the \(26\) synchronous pairs, \(12+9=21\) have attachment block
size two or three and hence contain only a binary or ternary
\(q\)-attachment carrier.  The remaining five have
\(\pi=\widehat1\) and contain the local quartic cumulant \(k_g(I)\).
\end{corollary}

\begin{proof}
Read the last row of the table in
\Cref{prop:V8-synchronous-count}.  A partition whose \(q\)-block has
size four is the one-block partition.
\end{proof}

\begin{remark}[The pair count is a source count, not a coordinate
count]
\label{rem:V8-count-scope}
The number \(26\) depends only on the four labelled rows.  Each pair
still carries the exact sum over the first-connectivity coordinate
\(g\) and every possible payer position.  The finite table does not
bound those physical sums.
\end{remark}

\section{Refined source capacities}
\label{sec:V8-capacities}

For each of the \(26\) synchronous partition pairs, let
\(\mathfrak S_{\rho,\pi}^{q,\varepsilon}\) be the normalized capacity
of
\begin{equation}
 \sum_{g=1}^n
 C_{<g}(\rho)\otimes L_{g,\pi}\otimes M_{>g},
\label{eq:V8-fixed-source-packet}
\end{equation}
with the sole payer kept at its exact physical location.

\begin{theorem}[Corrected first-attachment reduction]
\label{thm:V8-corrected-reduction}
To control the missing-row attachment in the connected quartic
carrier, it is sufficient to prove:
\begin{enumerate}[label=\textup{(\roman*)},leftmargin=4em]
\item the ordinary V61 source estimate for strict-prefix attachment,
      with the factor \eqref{eq:V8-q-prefix-factor} bounded by
      lower-order MTA before the joining source is enveloped;
\item a support-uniform marked-tree estimate for the \(21\) fixed
      packets \eqref{eq:V8-fixed-source-packet} having
      \(|B_q(\pi)|=2\) or \(3\); and
\item a separate, noninductive estimate for the five packets with
      \(\pi=\widehat1\).
\end{enumerate}
No suffix attachment packet exists.
\end{theorem}

\begin{proof}
\Cref{cor:V8-two-cases} partitions all connected words.  The strict
case is \Cref{thm:V8-prefix-absorption}.  In the equality case,
\Cref{prop:V8-synchronous-count} exhausts the possible fixed
partition pairs; \Cref{cor:V8-twenty-one} divides them into the
lower-order attachment letters and the five local quartic letters.
\Cref{lem:V8-order} removes the suffix case.
\end{proof}

\begin{assessment}[What V8 actually removes]
\label{ass:V8-removal}
The V7 first-attachment coordinate is not an additional analytic
summation in the strict-prefix branch: it is internal to a binary or
ternary connected prefix carrier.  In the synchronous branch, the
entire missing-row problem is a finite list of \(21\) nonquartic
source types plus five honest local quartic types.  The unresolved
infinite-looking object is only the physical source-coordinate and
payer sum in \eqref{eq:V8-fixed-source-packet}, not a second search
for where \(q\) first appears.
\end{assessment}

\section{Regression checks and V9 mandate}
\label{sec:V8-results}

\begin{proposition}[V8 regression suite]
\label{prop:V8-regressions}
\begin{enumerate}[label=\textup{(\roman*)},leftmargin=4em]
\item No first-\(q\) attachment is placed in a suffix already declared
      globally connected.
\item A singleton \(q\)-prefix is not mistaken for a private
      conclusion; the joining letter may attach it.
\item The internal attachment time of a lower-order prefix factor is
      not counted again.
\item A quartic joining letter remains quartic even when its prefix
      contains only lower-order connected blocks.
\item The \(26\)-pair table is labelled and finite but carries no
      claim of support-uniform source summation.
\item A suffix payer remains a paid local factor and is never replaced
      by a complete moment contraction.
\end{enumerate}
\end{proposition}

\begin{proof}
Items \textup{(i)--(iii)} follow from
\Cref{lem:V8-order,thm:V8-prefix-absorption}.  Item \textup{(iv)} is
\Cref{fire:V8-joining-still-quartic}.  Item \textup{(v)} is
\Cref{rem:V8-count-scope}.  Item \textup{(vi)} is part of the payer
definition in \eqref{eq:V8-fixed-source-packet}.
\end{proof}

\begin{theorem}[Third-series V8 result ledger]
\label{thm:V8-results}
Relative to V1--V7 and the frozen archives, V8 proves:
\begin{enumerate}[label=\textup{(V8.\arabic*)},leftmargin=5.8em]
\item the exact stopping-time order \(\tau_q\le g\);
\item removal of the spurious suffix-attachment branch;
\item absorption of every strict-prefix attachment into a binary or
      ternary connected prefix factor;
\item the complete labelled count of \(26\) synchronous
      attachment/source pairs;
\item separation of those pairs into \(21\) binary/ternary and five
      local-quartic source types; and
\item the corrected three-part source reduction of
      \Cref{thm:V8-corrected-reduction}.
\end{enumerate}
\end{theorem}

\begin{proof}
These are
\Cref{lem:V8-order,
cor:V8-two-cases,
thm:V8-prefix-absorption,
prop:V8-synchronous-count,
cor:V8-twenty-one,
thm:V8-corrected-reduction}.
\end{proof}

\begin{programme}[V9: synchronous binary source before absolute
values]
\label{prog:V8-V9}
\begin{enumerate}[leftmargin=3em]
\item Fix one of the \(12\) synchronous size-two pairs and write its
      joining letter as the binary \(qr\) cumulant times its remaining
      singleton or binary blocks.
\item Keep the exact prefix
      \(C_{<g}(\{q\}\mid\sigma)\) and the joining-coordinate sum in one
      telescoping packet.  Test whether the binary covariance can be
      paid by its actual \(H_{qr}\) stock before a positive envelope.
\item Treat separately a payer in the prefix, at \(g\), and in the
      complete suffix.  Use complete suffix contraction only in the
      first two cases.
\item Insert the V6 common-factor dominance only after the fixed
      source packet is bounded, then apply the V5
      reverse-or-deplete fan.
\item Construct a two-coordinate Rademacher packet for any proposed
      source telescope and check it against the V78 nonlocalization
      regression.
\item Extend the bound to the nine ternary synchronous types only
      after the binary source sum is support-uniform.
\end{enumerate}
\end{programme}

\begin{firewall}[Claim boundary after third-series V8]
\label{fire:V8-boundary}
V8 eliminates a false suffix case and reduces the synchronous
missing-row problem to an exact finite source table.  It does not
prove the source-coordinate/payer estimate for the \(21\)
binary/ternary packets, bound the five local quartic letters, or close
the sensitive, diagonal, nonclean, and other quartic source families.
Higher MTA, WUFC, CPM, cutoff removal, reflection positivity,
Osterwalder--Schrader reconstruction, and a uniform positive
continuum spectral gap remain open.  The full Yang--Mills mass gap
has not been proved.
\end{firewall}

\paragraph{Parent-version record.}
V8 was copied byte-for-byte from
\texttt{Renewal\_Yang\_Mills\_Mass\_Gap\_Research\_Notes\_V7.tex}
before this chapter was appended.  The V7 parent had \(4\,286\) source
lines, \(156\,150\) bytes, and SHA--256
\[
 \texttt{cbb4b39e911e33db828ccd81bc76c19674daca335d3918f68789cc51a87a6714}.
\]
The next compulsory companion audit remains V10.

\chapter{V9: Assembly correction and the pre-envelope
pair/attachment interface}
\label{ch:V9-assembly-boundary}

\section{Internal archive correction}
\label{sec:V9-correction}

The post-V8 archive check changes the proposed V9 calculation.
Compact V62 already enumerated all \(104\) labelled admissible
first-connectivity partition pairs and their twelve simultaneous
row-permutation orbits.  Compact V63 then proved that the individual
joining letters must be reassembled into one cumulant of the prefix
block-products; estimating them separately loses a repeated-row mass.
Compact V65--V67 proved the assembled fully thermal \(2+2\),
\(3+1\), and \(2+1+1\) families.  Compact V68 proved the stronger
packetwise statement that private coordinates cannot be
first-connectivity coordinates.

\begin{record}[Correction of V7--V8 novelty claims]
\label{rec:V9-novelty-correction}
The everywhere-private cancellation in
\Cref{lem:V7-private-annihilation} and its private-coordinate export
are specializations of the archived V68 private-letter theorem, not
new global MTA results.  The V8 count of \(26\) is a named-row
subcensus of the archived V62 table.  The genuinely new information
in V7--V8 is the missing-row bookkeeping:
\begin{enumerate}[label=\textup{(\roman*)},leftmargin=4em]
\item the unique triple consisting of first \(q\)-attachment time,
      least partner, and attachment-block size;
\item the stopping-time order \(\tau_q\le g\);
\item absorption of \(\tau_q<g\) inside the lower-order prefix block;
      and
\item the \(12/9/5\) refinement of the synchronous named-\(q\)
      subcensus.
\end{enumerate}
This record supersedes any broader wording in the V7 and V8 result
ledgers.
\end{record}

\section{The joining letters must be assembled}
\label{sec:V9-assembled-joining}

At coordinate \(g\), write \(X_{i,g}\) for the row variable in the
fixed tensor order.  For \(B\subseteq I\), put
\[
 Z_{B,g}:=\prod_{i\in B}X_{i,g}.
\]
For a proper partition \(\rho\), define the composite-block cumulant
\begin{equation}
 \kappa_g^\rho
 :=
 \kappa\bigl(Z_{B,g}:B\in\rho\bigr).
\label{eq:V9-composite-cumulant}
\end{equation}

\begin{proposition}[Archived assembled joining identity]
\label{prop:V9-assembled-joining}
For every \(\rho<\widehat1\),
\begin{equation}
 \boxed{
 J_g(\rho)
 =
 \sum_{\pi:\rho\vee\pi=\widehat1}L_{g,\pi}
 =
 \kappa_g^\rho.}
\label{eq:V9-assembled-joining}
\end{equation}
The equality holds before a norm, positive envelope, output split, or
payer allocation.
\end{proposition}

\begin{proof}
Apply the product-cumulant identity to the variables
\((Z_{B,g})_{B\in\rho}\).  A partition \(\pi\) of the underlying row
variables contributes precisely when the graph of intersections
between its blocks and the blocks of \(\rho\) is connected.  This is
equivalent to \(\rho\vee\pi=\widehat1\), and its product of local
cumulants is \(L_{g,\pi}\).  This is the V63 joining identity, recalled
here to correct the V8 proposed letterwise estimate.
\end{proof}

\begin{theorem}[The \(26\) synchronous letters collapse to five
assembled sources]
\label{thm:V9-five-sources}
For \(\tau_q=g\), the sum of the \(26\) partition-letter pairs in
\Cref{prop:V8-synchronous-count} is exactly the sum of five fixed
prefix-partition packets:
\begin{equation}
 \boxed{
 \sum_{\substack{\rho=\{q\}\mid\sigma\\ \sigma\in\Pi(W)}}
 \ \sum_{g=1}^n
 C_{<g}(\rho)\otimes\kappa_g^\rho\otimes M_{>g}.}
\label{eq:V9-five-source-sum}
\end{equation}
Up to the labels of \(W=I\setminus\{q\}\), these five packets have
only three analytic types:
\begin{center}
\begin{tabular}{c c c c}
\toprule
\(\sigma\) & number & prefix type & assembled joining carrier\\
\midrule
\(a\mid b\mid z\) & \(1\) & \(1+1+1+1\) &
fourth cumulant\\
\(uv\mid w\) & \(3\) & \(2+1+1\) &
ternary cumulant of block-products\\
\(abz\) & \(1\) & \(3+1\) &
covariance of two block-products\\
\bottomrule
\end{tabular}
\end{center}
\end{theorem}

\begin{proof}
If \(\tau_q=g\), the prefix join has the unique form
\(\rho=\{q\}\mid\sigma\) from
\eqref{eq:V8-q-singleton-prefix}.  For each of the five choices of
\(\sigma\), sum all joining partitions \(\pi\) rather than estimating
their \(12/9/5\) size cells.  By
\Cref{prop:V9-assembled-joining}, their sum is
\(\kappa_g^\rho\), proving \eqref{eq:V9-five-source-sum}.

The partition lattice on three labelled rows has one discrete
partition, three pair--singleton partitions, and one one-block
partition.  Including the singleton block \(\{q\}\) gives the three
prefix types in the table.  The cumulant order is the number of
prefix blocks, yielding order four, three, and two respectively.
\end{proof}

\begin{firewall}[Do not estimate the \(21\) letters separately]
\label{fire:V9-no-letterwise}
The \(12\) binary-attachment and nine ternary-attachment letters in
\Cref{cor:V8-twenty-one} are an exact algebraic expansion of the
composite cumulants in \eqref{eq:V9-five-source-sum}.  They are useful
for graph census only.  Their separate positive envelopes need not
retain the common-row mass or the one-payer cancellation.  V9
therefore withdraws the letterwise V8--V9 programme and keeps each
\(\kappa_g^\rho\) assembled.
\end{firewall}

\section{The positive-envelope boundary}
\label{sec:V9-envelope-boundary}

The current matching calculation begins after V100 has bounded the
exact first-breaker word by a positive product-state majorant and then
allocated its breaker stock among the four pair labels.  The entries
\(Q_{ab}^{(\varepsilon)}\) are ledger summands of that majorant, not an
operator decomposition of the original signed covariance.  Exact
M\"obius cancellation must therefore be used before this step.

\begin{proposition}[Cancellation does not descend through an arbitrary
positive majorant]
\label{prop:V9-cancellation-no-descent}
Let a signed carrier \(P\) satisfy \(P=0\), and suppose
\(|P|\le Q\) for a nonnegative majorant \(Q\).  It does not follow that
\(Q=0\), nor that an arbitrary positive allocation of \(Q\) retains
the cancellation producing \(P=0\).  Consequently the cancellation
in \Cref{lem:V7-private-annihilation} cannot be applied directly to a
post-envelope pair entry \(Q_{ab}^{(\varepsilon)}\).
\end{proposition}

\begin{proof}
The scalar identity
\[
 P=1-1=0
\]
has the valid positive majorant
\[
 |P|\le |1|+|-1|=2.
\]
The two positive terms no longer cancel.  More generally, absolute
value and positive coefficient envelopes are subadditive, not linear,
so they need not preserve a kernel of the signed summation map.  By
\Cref{rem:V1-ledger-not-operator}, \(Q_{ab}^{(\varepsilon)}\) is
exactly such a positive ledger entry rather than a signed localized
operator.  The V7 linked-word cancellation is therefore available
only in the pre-envelope word.
\end{proof}

\begin{corollary}[Legal use of private-row annihilation]
\label{cor:V9-legal-private}
To use the private-row cancellation in the current matching residual,
one must return to the exact V100 carrier
\[
 B_\mu S_G^\mu
\]
and perform the first-\(q\)-attachment classification before the
positive breaker envelope and pair allocation.  If the word is
enveloped first, V7 supplies no cancellation theorem for an
individual \(Q_{ab}^{(\varepsilon)}\).
\end{corollary}

\begin{proof}
The exact linked-word identity applies to the signed partition-state
expansion of \(B_\mu S_G^\mu\).  Apply
\Cref{prop:V9-cancellation-no-descent} after the positive envelope.
\end{proof}

\begin{remark}[Why the V4--V6 graph compilers remain valid]
\label{rem:V9-compilers-valid}
The V4--V6 inequalities do not use signed cancellation inside
\(Q_{ab}^{(\varepsilon)}\).  They multiply an already proved positive
pair bound by actual response lower bounds and then verify the exact
tree denominators.  Those compilers remain valid.  The correction
affects only the attempted V7 use of linked-word annihilation and the
V8 proposal to envelope joining letters separately.
\end{remark}

\section{The missing compatibility interface}
\label{sec:V9-interface}

\begin{definition}[Cancellation-preserving pair/attachment interface]
\label{def:V9-interface}
For the exact first-breaker carrier \(B_\mu S_G^\mu\), a
\emph{pair/attachment interface} consists of pre-envelope packets
\[
 {\cal P}_{xy;\lambda}^{(\varepsilon)},
\]
where \(xy\) is a breaker pair and \(\lambda\) is a first-\(q\)
attachment class, such that:
\begin{enumerate}[label=\textup{(\alph*)},leftmargin=4em]
\item their signed sum is the exact V100 carrier, with every linked
      word occurring once;
\item the no-\(q\)-attachment class is zero by the linked-word
      identity before any absolute value;
\item summing \(\lambda\) for fixed \(xy\) admits the V100 positive
      bound \(CH_{xy}(G)\), with no operator localized by
      \(H_{xy}\) asserted;
\item every surviving response mark is extracted from an actual bank
      occurrence in the same packet; and
\item the sole payer is allocated once after the attachment and
      assembled joining-cumulant sums are formed.
\end{enumerate}
\end{definition}

\begin{proposition}[What such an interface would prove]
\label{prop:V9-interface-suffices}
Assume the interface in \Cref{def:V9-interface} exists with
support-uniform unpaid and once-paid bounds.  Then:
\begin{enumerate}[label=\textup{(\roman*)},leftmargin=4em]
\item every private-\(q\) word with no external attachment disappears
      before pair allocation;
\item every binary common-factor attachment which supplies the two
      required response lower bounds closes by the V4/V5 fan;
\item every failed reverse response has the V5 target-row depletion
      inside the same pre-envelope packet; and
\item the remaining source classes are the assembled ternary and
      local-quartic attachments, plus the previously named sensitive
      and nonclean states.
\end{enumerate}
\end{proposition}

\begin{proof}
Item \textup{(i)} is interface condition \textup{(b)}.  Under
\textup{(c)--(e)}, the hypotheses of
\Cref{thm:V4-fan-completion,thm:V5-direct-reverse} hold in one
assembled word, proving \textup{(ii)}.  Apply
\Cref{thm:V5-reverse-or-deplete} before separating its incidence bank
to get \textup{(iii)}.  The first-attachment size trichotomy and V2
sensitive/insensitive split exhaust the remaining clean roles, giving
\textup{(iv)}.
\end{proof}

\begin{firewall}[The interface is not yet constructed]
\label{fire:V9-interface-open}
V100 proves a support-uniform majorant for the complete breaker and a
four-term scalar allocation of its stock.  V61/V68 prove exact
pre-envelope connectedness and private-factor identities.  Neither
archive theorem proves that these two classifications admit the joint
refinement in \Cref{def:V9-interface}.  Constructing it requires
tracking the first breaker, first \(q\)-attachment, complete composite
joining carrier, and payer in one partition-state word.
\end{firewall}

\section{Corrected residual and V10 mandate}
\label{sec:V9-results}

\begin{assessment}[The actual upstream bottleneck]
\label{ass:V9-bottleneck}
The current obstacle is not a new binary source estimate: those
letters belong to assembled V63--V67 composite cumulants.  Nor may
private-row cancellation be imposed on the already positive
\(Q_{ab}\) ledger.  The smallest honest gate is a compatibility
theorem between two valid but differently ordered constructions:
\[
\begin{array}{c}
\text{exact linked word}
\longrightarrow
\text{first \(q\)-attachment and assembled joining cumulant}
\longrightarrow
\text{positive V100 pair ledger},\\[2mm]
\text{exact linked word}
\longrightarrow
\text{V100 first matching breaker}
\longrightarrow
\text{positive pair ledger and response acquisition}.
\end{array}
\]
V10 must determine whether these routes have a common pre-envelope
refinement without introducing a coordinate or payer multiplicity.
\end{assessment}

\begin{theorem}[Third-series V9 result ledger]
\label{thm:V9-results}
Relative to the corrected archive boundary, V9 proves:
\begin{enumerate}[label=\textup{(V9.\arabic*)},leftmargin=5.8em]
\item the \(26\) synchronous partition letters reassemble exactly
      into five fixed prefix packets and three analytic composite
      cumulant types;
\item the proposed \(21\)-letter separate estimate is invalid as a
      proof strategy because it discards the required assembly;
\item exact linked-word cancellation cannot be inferred from a
      post-envelope positive pair entry;
\item the V4--V6 positive graph compilers are unaffected by that
      correction; and
\item the precise sufficient pair/attachment interface is
      \Cref{def:V9-interface}.
\end{enumerate}
\end{theorem}

\begin{proof}
These are
\Cref{thm:V9-five-sources,
fire:V9-no-letterwise,
prop:V9-cancellation-no-descent,
rem:V9-compilers-valid,
prop:V9-interface-suffices}.
\end{proof}

\begin{programme}[V10: companion audit and common refinement]
\label{prog:V9-V10}
\begin{enumerate}[leftmargin=3em]
\item Perform the compulsory audit of the current SM and NS notes
      before selecting the construction.
\item Return to the exact V100 product telescope for
      \({\mathfrak J}_G^\mu\), before its absolute-value physical
      envelope.
\item On each exact telescope term, retain the cumulative partition
      join and the first \(q\)-attachment class, but sum all joining
      letters with fixed prefix partition into
      \(\kappa_g^\rho\).
\item Determine whether the breaker pair label \(xy\) can be inserted
      as a scalar mark only after the complete covariance is bounded,
      while the attachment class remains an outer exact disjoint
      index.
\item Prove additivity or a common positive pinching sufficient for
      conditions \textup{(a)--(e)} of
      \Cref{def:V9-interface}; if it fails, give the smallest
      two-coordinate matrix counterpacket.
\item Keep all suffix-payer cases inside the complete-product payer
      sum, and re-run V78 and V84 before declaring compatibility.
\end{enumerate}
\end{programme}

\begin{firewall}[Claim boundary after third-series V9]
\label{fire:V9-boundary}
V9 corrects an invalid letterwise direction and locates the exact
pre-envelope interface at which V7 cancellation could legally meet
the V100 pair ledger.  It does not construct that interface, close
the assembled ternary or local-quartic attachment sources in the
current residual, or settle sensitive, diagonal, nonclean, and other
quartic source families.  Higher MTA, WUFC, CPM, cutoff removal,
reflection positivity, Osterwalder--Schrader reconstruction, and a
uniform positive continuum spectral gap remain open.  The full
Yang--Mills mass gap has not been proved.
\end{firewall}

\paragraph{Parent-version record.}
V9 was copied byte-for-byte from
\texttt{Renewal\_Yang\_Mills\_Mass\_Gap\_Research\_Notes\_V8.tex}
before this chapter was appended.  The V8 parent had \(4\,726\) source
lines, \(172\,522\) bytes, and SHA--256
\[
 \texttt{327edf4dce07200aed8b274c968300df8fd6af90463d3f430770d276fd509aa8}.
\]
The compulsory companion audit is due at V10.

\chapter{V10: Companion audit and the assembled early-tail
refinement of the V100 breaker packet}
\label{ch:V10-early-tail}

\section{Compulsory V10 companion audit}
\label{sec:V10-audit}

The current companion notes were read passively before the V10
construction:
\begin{center}
\small
\begin{tabular}{p{0.48\textwidth} r r}
\toprule
file & lines & bytes\\
\midrule
\texttt{Renewal\_SM\_Einstein\_Continuum\_Research\_Notes\_V87.tex}
& \(34\,518\) & \(1\,257\,263\)\\
\texttt{Renewal\_NS\_Stability\_Research\_Notes\_V31.tex}
& \(23\,052\) & \(887\,537\)\\
\bottomrule
\end{tabular}
\end{center}
Their audit-time SHA--256 digests were
\[
\begin{split}
&\texttt{f55322215525cb819de2eb6513541d2b0f2e2fb9a4760aad79e9ce8d15b1bef4},\\
&\texttt{f1121b62238ad5491bb7cf03635ea9934942edf573ed8faaa9ee79e1d38a6696}.
\end{split}
\]

SM V85 explicitly records the Yang--Mills lesson that acyclicity of a
scalar ascent is not packet transport.  SM V86 obtains its own
multiscale transport only after constructing one dimensionless path
template with common gap, locality, and analyticity constants at all
scales.  SM V87 then keeps its global comparison conditional outside
the local gap patch.  No spectral or bordism theorem is imported
here.  The transferable proof-order principle is that all refinements
must be made inside one complete template whose contractive endpoints
are preserved.

NS V31 is unchanged.  Its exact first-moment obstruction and saturating
packets continue to rule out replacing a source-correlated coordinate
sum by a uniform local cap followed by a count.

\begin{record}[V10 audit decision]
\label{rec:V10-audit-decision}
The common Yang--Mills template will be the exact V100 two-block
product telescope.  Before the first breaker, the matching pair
containing the missing endpoint will be split as
\[
 \text{complete joint pair moment}
 =
 \text{independent pair replicas}
 +
 \text{complete pair covariance}.
\]
This sums every early matching-tail attachment into one covariance
and leaves a complete contractive no-early-tail endpoint.  No
restricted first-attachment transfer is propagated.
\end{record}

\section{Return to the exact V100 packet}
\label{sec:V10-V100-interface}

Fix a perfect matching
\[
 \mu=\{x,q\}\mid\{y,z\},
\]
where \(xy\) is a positive breaker-pair label, \(q=m(x)\), and
\(z=m(y)\).  Assume the response \(y\to z\) has already been acquired,
while \(x\to q\) is the missing matching tail.  This is the original
V100 missing-endpoint notation corresponding, after relabelling, to
the singleton missing-complement state of V3--V6.

Let \(t\) denote the first coordinate at which the exact matching word
leaves the two-block partition \(\mu\).  Every local partition before
\(t\) refines \(\mu\).  Hence the prefix factorizes between the two
matching components, and the factor on \(\{x,q\}\) can be treated
without expanding the other component.

At a prefix coordinate \(e<t\), let \(T_e^{xq}\) be the complete joint
two-row moment and let \(R_e^{xq}\) be the moment in which the \(x\)
and \(q\) replicas are independent.  Put
\begin{align}
 T_{<t}^{xq}
 &:=
 \bigotimes_{e<t}T_e^{xq},
\label{eq:V10-prefix-joint}\\
 R_{<t}^{xq}
 &:=
 \bigotimes_{e<t}R_e^{xq},
\label{eq:V10-prefix-independent}\\
 {\mathfrak J}_{<t}^{xq}
 &:=
 T_{<t}^{xq}-R_{<t}^{xq}.
\label{eq:V10-prefix-covariance}
\end{align}

\begin{lemma}[Exact early-tail split]
\label{lem:V10-early-tail-split}
For every first-breaker coordinate \(t\),
\begin{equation}
 \boxed{
 T_{<t}^{xq}
 =
 R_{<t}^{xq}+{\mathfrak J}_{<t}^{xq}.}
\label{eq:V10-early-tail-split}
\end{equation}
In the local partition-word expansion,
\(R_{<t}^{xq}\) contains exactly the words in which \(x,q\) remain
separate throughout the prefix, while
\({\mathfrak J}_{<t}^{xq}\) is the signed assembled sum of all words
in which the matching edge \(xq\) occurs at least once before \(t\).
\end{lemma}

\begin{proof}
The displayed identity is the definition
\eqref{eq:V10-prefix-covariance}.  For its partition interpretation,
the product \(R_{<t}^{xq}\) uses independent replicas at every prefix
coordinate, so every local two-row partition is discrete.  Conversely,
if all local two-row partitions are discrete, their moment word is a
term of that product.  The difference from the complete joint product
therefore contains exactly the complementary linked two-row words,
assembled before any absolute value.  On two rows, linkedness is
equivalent to occurrence of the binary block \(\{x,q\}\) at least
once.
\end{proof}

\begin{corollary}[Two pre-envelope V100 branches]
\label{cor:V10-two-branches}
Insert \eqref{eq:V10-early-tail-split} into the exact matching word
\(B_\mu S_G^\mu\) before its positive envelope.  The result is a
signed sum of two packets:
\begin{enumerate}[label=\textup{(\roman*)},leftmargin=4em]
\item an \emph{early-tail packet} containing the complete covariance
      \({\mathfrak J}_{<t}^{xq}\);
\item a \emph{no-early-tail packet} containing the complete
      independent moment \(R_{<t}^{xq}\).
\end{enumerate}
Every original word occurs once, and no condition is imposed on the
complete suffix after the breaker.
\end{corollary}

\begin{proof}
Tensor distributivity inserts the two terms of
\eqref{eq:V10-early-tail-split} at the single \(\{x,q\}\)-prefix
factor.  There are exactly two summands and their sum is the original
factor.  All other prefix factors, the exact breaker covariance, and
the complete suffix are unchanged.
\end{proof}

\section{The complete early covariance has contractive endpoints}
\label{sec:V10-prefix-bound}

Let
\[
 h_{xq,e}:=a_{x,e}a_{q,e},
 \qquad
 H_{xq}(<t):=\sum_{e<t}h_{xq,e}.
\]

\begin{lemma}[Product telescope for the early tail]
\label{lem:V10-prefix-telescope}
The complete prefix covariance has the exact expansion
\begin{equation}
 {\mathfrak J}_{<t}^{xq}
 =
 \sum_{s<t}
 \left(\bigotimes_{e<s}T_e^{xq}\right)
 \otimes(T_s^{xq}-R_s^{xq})
 \otimes
 \left(\bigotimes_{s<e<t}R_e^{xq}\right).
\label{eq:V10-prefix-telescope}
\end{equation}
Every spectator on both sides of the local covariance is a complete
moment contraction.
\end{lemma}

\begin{proof}
This is the ordered product-difference identity applied to
\(\bigotimes_{e<t}T_e^{xq}-\bigotimes_{e<t}R_e^{xq}\).
Both \(T_e^{xq}\) and \(R_e^{xq}\) are expectations of tensor products
of unitary coefficient matrices, and therefore have operator norm at
most one.
\end{proof}

\begin{theorem}[Support-uniform early-tail response]
\label{thm:V10-prefix-response}
Let \({\cal N}_{<t}^{xq}\) be the complete positive physical envelope
of \({\mathfrak J}_{<t}^{xq}\).  Then
\begin{equation}
 \boxed{
 \|{\cal N}_{<t}^{xq}\|,
 \ \kappa_\otimes({\cal N}_{<t}^{xq})
 \le
 C\min\{1,s_\alpha H_{xq}(<t)\}.}
\label{eq:V10-prefix-unpaid}
\end{equation}
If the sole output numerator and resolvent are assigned to this
prefix covariance, its paid envelope satisfies
\begin{equation}
 \boxed{
 \|{\cal P}_{<t}^{xq}\|,
 \ \kappa_\otimes({\cal P}_{<t}^{xq})
 \le
 C\min\{H_{xq}(<t),s_\alpha^{-1}\}.}
\label{eq:V10-prefix-paid}
\end{equation}
Both estimates are uniform in \(t\) and in the number of prefix
coordinates.
\end{theorem}

\begin{proof}
At the selected coordinate \(s\), the established binary
composite-covariance estimate gives
\[
 \bigl\||T_s^{xq}-R_s^{xq}|_{\rm phys}\bigr\|
 \le Cs_\alpha h_{xq,s}
\]
in the unpaid case and \(Ch_{xq,s}\) when that local factor pays.
Use \Cref{lem:V10-prefix-telescope} and the contractivity of all
spectators, then sum \(s<t\).  This gives the response parts
\(Cs_\alpha H_{xq}(<t)\) and \(CH_{xq}(<t)\).  The two endpoints in
\eqref{eq:V10-prefix-covariance} are complete contractions, giving the
universal unpaid cap; the ordinary one-payer argument gives
\(C/s_\alpha\) in the paid case.  Taking the smaller estimates proves
the two minima.  Capacity is no larger than the same positive envelope
norm.
\end{proof}

\begin{firewall}[No restricted prefix product]
\label{fire:V10-no-restricted-prefix}
V10 never propagates the sum of local letters subject to a
no-attachment-yet restriction.  That restricted product need not be
contractive.  The no-early-tail endpoint is instead the complete
independent-replica moment \(R_{<t}^{xq}\), and all early attachments
are restored before norming as the single difference
\({\mathfrak J}_{<t}^{xq}=T_{<t}^{xq}-R_{<t}^{xq}\).
\end{firewall}

\section{Joining-coordinate sums retain the actual stocks}
\label{sec:V10-ordered-stock}

At breaker coordinate \(t\), define the two-block crossing density
\[
 b_t^\mu
 :=
 \sum_{\substack{u\in\{x,q\}\\v\in\{y,z\}}}
 a_{u,t}a_{v,t}.
\]
For the pair-labelled term \(xy\), write
\[
 b_t^{xy}:=a_{x,t}a_{y,t}.
\]

\begin{lemma}[Early-tail/breaker Cartesian domination]
\label{lem:V10-cartesian}
For every nonnegative matching-tail density \(h_s\) and breaker
density \(b_t\),
\begin{equation}
 \sum_t\left(\sum_{s<t}h_s\right)b_t
 \le
 \left(\sum_sh_s\right)\left(\sum_tb_t\right).
\label{eq:V10-cartesian}
\end{equation}
The same inequality remains true after imposing any additional
ordering condition on one payer coordinate and then dropping that
condition.
\end{lemma}

\begin{proof}
Expand the left side as
\(\sum_{s<t}h_sb_t\).  Its ordered index set is a subset of the full
Cartesian product, and all summands are nonnegative.  Adding a payer
coordinate gives another subset of a Cartesian product; dropping its
ordering constraint can only enlarge the sum.
\end{proof}

\begin{corollary}[No early-tail coordinate multiplicity]
\label{cor:V10-no-coordinate-loss}
After the local covariance at \(s\) and the breaker covariance at
\(t\) have been assembled before their positive envelopes,
\begin{align}
 \sum_t H_{xq}(<t)b_t^\mu
 &\le H_{xq}(E)H_E(\{x,q\},\{y,z\}),
\label{eq:V10-total-stock}\\
 \sum_t H_{xq}(<t)b_t^{xy}
 &\le H_{xq}(E)H_{xy}(E).
\label{eq:V10-pair-stock}
\end{align}
No count of possible early-tail or breaker coordinates occurs.
\end{corollary}

\begin{proof}
Apply \Cref{lem:V10-cartesian} with
\(h_s=h_{xq,s}\), first with \(b_t=b_t^\mu\) and then with
\(b_t=b_t^{xy}\).  The products of total sums are the displayed
stocks.
\end{proof}

\section{Dominant early tail or genuine off-prefix depletion}
\label{sec:V10-dominance}

\begin{theorem}[Pre-breaker tail dichotomy]
\label{thm:V10-tail-dichotomy}
Fix \(0<\gamma<1\).  For every first-breaker coordinate \(t\), exactly
one of:
\begin{enumerate}[label=\textup{(\roman*)},leftmargin=4em]
\item the early matching response is dominant,
      \begin{equation}
      d_{x\to q}(<t)
      :=
      \frac{H_{xq}(<t)}{\Xi_x}
      \ge\gamma a_*;
      \label{eq:V10-early-dominant}
      \end{equation}
\item the \(xq\)-incidence coordinates before \(t\) carry less than a
      \(\gamma\)-fraction of row \(x\)'s total mass,
      \begin{equation}
      \sum_{\substack{e<t\\a_{x,e}a_{q,e}>0}}a_{x,e}
      <\gamma\Xi_x.
      \label{eq:V10-prefix-depleted}
      \end{equation}
\end{enumerate}
In alternative \textup{(i)}, if the other breaker endpoint has its
matching response \(y\to z\), the literal path \(q-x-y-z\) closes
after the early-tail packet and breaker pair have been jointly
enveloped with the sole payer.
\end{theorem}

\begin{proof}
If \eqref{eq:V10-early-dominant} fails, then
\[
 H_{xq}(<t)<\gamma a_*\Xi_x.
\]
On every prefix incidence coordinate, \(a_{q,e}\ge a_*\), so
\[
 a_*\sum_{\substack{e<t\\a_{x,e}a_{q,e}>0}}a_{x,e}
 \le H_{xq}(<t).
\]
Cancel \(a_*\) to obtain \eqref{eq:V10-prefix-depleted}.  The two
threshold alternatives are exhaustive.

In the first alternative, the two response tails are the internal
vertices \(x,y\) of \(q-x-y-z\).  Once the exact early covariance,
breaker, other tail, and payer have been bounded in one assembled
packet, \Cref{thm:V1-complementary-partners} gives the stated path.
\end{proof}

\begin{remark}[Why the final clause remains an interface statement]
\label{rem:V10-joint-envelope}
\Cref{thm:V10-prefix-response} and
\Cref{cor:V10-no-coordinate-loss} provide the prefix and
coordinate-summation estimates required for the early-tail packet.
They do not by themselves prove the full Schatten/product-state
sandwich containing the other matching component and every payer
location.  That joint analytic carrier is the remaining part of the
V9 interface.
\end{remark}

\section{Partial construction of the pair/attachment interface}
\label{sec:V10-interface}

\begin{theorem}[Algebraic and stock parts of the interface]
\label{thm:V10-partial-interface}
The early/no-early split of
\Cref{cor:V10-two-branches} constructs the following parts of
\Cref{def:V9-interface}:
\begin{enumerate}[label=\textup{(\alph*)},leftmargin=4em]
\item it is performed on the exact V100 word before an absolute value;
\item it is a two-term identity, so no word or attachment coordinate
      is counted twice;
\item the no-early branch has the complete contractive endpoint
      \(R_{<t}^{xq}\);
\item the early branch has the support-uniform complete covariance
      bounds \eqref{eq:V10-prefix-unpaid}--%
      \eqref{eq:V10-prefix-paid}; and
\item its joint sum with the breaker coordinate is controlled by the
      actual product of response stocks
      \eqref{eq:V10-total-stock}--%
      \eqref{eq:V10-pair-stock}.
\end{enumerate}
It remains to prove a common product-state sandwich for the early
covariance, the other matching tail, the breaker, and every payer
location.
\end{theorem}

\begin{proof}
Items \textup{(a)--(c)} are
\Cref{lem:V10-early-tail-split,cor:V10-two-branches}.
Item \textup{(d)} is \Cref{thm:V10-prefix-response}, and item
\textup{(e)} is \Cref{cor:V10-no-coordinate-loss}.  None of these
arguments applies a separate capacity to the other matching component,
so the last joint sandwich is not included in the proved statements.
\end{proof}

\begin{assessment}[Residual after the first common refinement]
\label{ass:V10-residual}
The prior statement that a matching tail is missing now has an exact
pre-breaker split.  All early occurrences of that tail are assembled
into one binary covariance with complete contractive endpoints.  If
its stock is dominant, the desired literal path is already visible;
if not, row \(x\) is genuinely depleted on matching incidences before
the breaker.  The no-early packet is not a restricted transfer and
therefore does not recreate the V36/V60 Green problem.

The remaining analytic task is narrower: place the early covariance
inside the V97/V100 one-use Schatten sandwich together with the other
tail and the breaker, and verify the prefix-, breaker-, matching-word-,
separator-, and suffix-payer cases without applying two independent
first moments to the same matching bank.
\end{assessment}

\section{Result ledger and V11 mandate}
\label{sec:V10-results}

\begin{theorem}[Third-series V10 result ledger]
\label{thm:V10-results}
Relative to V1--V9 and the corrected archive boundary, V10 proves:
\begin{enumerate}[label=\textup{(V10.\arabic*)},leftmargin=6.3em]
\item the exact early-tail/no-early-tail decomposition inside the
      pre-envelope V100 word;
\item a complete product telescope for the early matching covariance;
\item support-uniform unpaid and once-paid covariance minima for that
      prefix;
\item Cartesian domination of the joint early-tail/breaker coordinate
      sum by actual stock products;
\item the dominant-early-tail/off-prefix-depletion dichotomy; and
\item the algebraic and stock components of the V9
      pair/attachment interface.
\end{enumerate}
\end{theorem}

\begin{proof}
These are
\Cref{lem:V10-early-tail-split,
lem:V10-prefix-telescope,
thm:V10-prefix-response,
cor:V10-no-coordinate-loss,
thm:V10-tail-dichotomy,
thm:V10-partial-interface}.
\end{proof}

\begin{programme}[V11: the one-use early-tail sandwich]
\label{prog:V10-V11}
\begin{enumerate}[leftmargin=3em]
\item Write the exact occurrence order of
      \({\mathfrak J}_{<t}^{xq}\), the already acquired \(y\to z\)
      matching tail, and the breaker covariance at \(t\) inside the
      V100 matching word.
\item Use one V97 three-factor or Hilbert--Schmidt sandwich, not the
      product of separate prefix and matching-word capacities.
\item Treat the payer in the early covariance, other tail, breaker,
      local separator, and complete suffix as mutually exclusive
      cases.
\item In every case retain the ordered stock sum until
      \Cref{lem:V10-cartesian} can be applied.
\item If the early-tail sandwich closes, apply it to
      \eqref{eq:V10-early-dominant} and continue the no-early branch
      from the exact depletion \eqref{eq:V10-prefix-depleted}.
\item Test coincident matching banks on the V98 joint-carrier
      regression and test the signed source on V78/V84 before taking
      a positive envelope.
\end{enumerate}
\end{programme}

\begin{firewall}[Claim boundary after third-series V10]
\label{fire:V10-boundary}
V10 constructs a genuine pre-envelope early-tail refinement with
complete contractive endpoints and no coordinate multiplicity.  It
does not yet prove the one-use product-state sandwich joining that
prefix covariance to the other tail and breaker at every payer
location, terminate the no-early branch, or close the sensitive,
diagonal, nonclean, and other quartic source families.  Higher MTA,
WUFC, CPM, cutoff removal, reflection positivity,
Osterwalder--Schrader reconstruction, and a uniform positive
continuum spectral gap remain open.  The full Yang--Mills mass gap
has not been proved.
\end{firewall}

\paragraph{Parent-version record.}
V10 was copied byte-for-byte from
\texttt{Renewal\_Yang\_Mills\_Mass\_Gap\_Research\_Notes\_V9.tex}
before this chapter was appended.  The V9 parent had \(5\,102\) source
lines, \(187\,466\) bytes, and SHA--256
\[
 \texttt{b92aa957387ac49420602cf1980ee99e87e8cf98ae1bb5d329b42dd74175e354}.
\]
The compulsory V10 companion audit is recorded in
\Cref{sec:V10-audit}; the next required audit is V15.

\chapter{V11: The one-use early-tail sandwich and its exact row-mass
deficit}
\label{ch:V11-sandwich}

\section{Context}
\label{sec:V11-context}

V10 assembled all matching-tail attachments before the first breaker
into the complete covariance \({\mathfrak J}_{<t}^{xq}\).  The next
question is whether that covariance can be placed in the V97
five-factor sandwich with the breaker, without applying separate
product-state capacities.  The answer has two parts.  The complete
one-use sandwich and every payer location are support-uniform.  But
if it is used merely through its raw stock \(H_{xq}\), the resulting
three-edge numerator is larger than the exact quartic tree response by
one factor \(\Xi_x\).  This deficit is sharp.

Thus the sandwich validates the operator assembly but does not replace
the V10 dominance/depletion dichotomy.  A dominant early response
closes by the existing probability-mark compiler; a nondominant
early covariance must be removed from the target row's mass ledger
rather than treated as a complete tree edge.

\section{Two endpoint energy envelopes}
\label{sec:V11-energies}

For \(h\ge0\), define
\begin{equation}
 U(h):=\min\{1,s_\alpha h\},
 \qquad
 P(h):=\min\{h,s_\alpha^{-1}\}
 =\frac1{s_\alpha}U(h).
\label{eq:V11-UP}
\end{equation}
Set
\[
 h_1:=H_{xq}(<t),
 \qquad
 h_2:=H_{xy}(G_t),
\]
where \(G_t\) denotes the exact breaker occurrence or its
coordinate-resolved bank in the fixed packet.

\begin{lemma}[Quadratic early-tail and breaker energies]
\label{lem:V11-endpoint-energies}
Let \(A_t\) be the self-adjoint complete occurrence representing
\({\mathfrak J}_{<t}^{xq}\) in the V94 physical realization, and let
\(E_t\) be the corresponding signed breaker occurrence.  Then
\begin{align}
 {\cal E}_\otimes(A_t)
 &\le C U(h_1)^2,
 &
 {\cal E}_\otimes(A_t^{\rm pay})
 &\le C P(h_1)^2,
\label{eq:V11-A-energy}\\
 {\cal E}_\otimes(E_t)
 &\le C U(h_2)^2,
 &
 {\cal E}_\otimes(E_t^{\rm pay})
 &\le C P(h_2)^2.
\label{eq:V11-E-energy}
\end{align}
\end{lemma}

\begin{proof}
\Cref{thm:V10-prefix-response} bounds both the norm and product-state
capacity of the complete early covariance by \(CU(h_1)\), and its
paid versions by \(CP(h_1)\).  The V94 one-use energy inequality
\[
 {\cal E}_\otimes(S)\le
 \|S\|\,\kappa_\otimes(S)
\]
gives \eqref{eq:V11-A-energy}.  The same argument applied to the V100
breaker minima gives \eqref{eq:V11-E-energy}.  The standard physical
self-adjoint realization is made before the exact Fourier duality;
no scalar stock is substituted for either endpoint.
\end{proof}

\section{The single Schatten sandwich}
\label{sec:V11-one-use}

Write the exact output-block multiplier of the early-tail packet in
the form
\begin{equation}
 K_{U,t}
 =
 a_{U,t}\,c_{1,U,t}\,b_{U,t}\,
 c_{2,U,t}\,e_{U,t},
\label{eq:V11-five-factor}
\end{equation}
where \(a_{U,t},e_{U,t}\) are the restrictions of \(A_t,E_t\);
\(b_{U,t}\) is the complete intervening matching word, including the
already assembled other component; and the two separators satisfy
uniform bounds
\[
 \sup_{U,t}\|c_{1,U,t}\|\le L_1,
 \qquad
 \sup_{U,t}\|c_{2,U,t}\|\le L_2.
\]
The V100 matching-word theorem supplies
\(\|B_t\|\le C\), or \(C/s_\alpha\) when the sole payer lies in that
middle word or its grouped suffix.

\begin{theorem}[Complete early-tail/breaker payer table]
\label{thm:V11-payer-table}
For every fixed first-breaker packet, the V97 exact Fourier sandwich
gives
\begin{align}
 \Gamma_\otimes(1,K_t)
 &\le C L_1L_2\,U(h_1)U(h_2),
\label{eq:V11-unpaid-sandwich}\\
 \Gamma_\otimes(w,K_t^{\rm pay})
 &\le
 \frac{C L_1L_2}{s_\alpha}\,U(h_1)U(h_2)
\label{eq:V11-paid-sandwich}
\end{align}
at every mutually exclusive payer location:
\begin{enumerate}[label=\textup{(\roman*)},leftmargin=4em]
\item the early prefix covariance;
\item the breaker;
\item the intervening matching word, including the other tail;
\item either bounded local separator; or
\item the complete grouped suffix.
\end{enumerate}
No product of separately optimized endpoint capacities is used.
\end{theorem}

\begin{proof}
In the unpaid case, apply the V97 three-factor endpoint sandwich to
\eqref{eq:V11-five-factor}:
\[
 \Gamma_\otimes(1,K_t)
 \le
 L_1L_2\|B_t\|
 \sqrt{{\cal E}_\otimes(A_t){\cal E}_\otimes(E_t)}.
\]
Use \(\|B_t\|\le C\) and
\Cref{lem:V11-endpoint-energies} to obtain
\eqref{eq:V11-unpaid-sandwich}.

If the early covariance pays, replace \(U(h_1)\) by
\(P(h_1)=U(h_1)/s_\alpha\).  If the breaker pays, make the analogous
replacement at \(h_2\).  If the payer lies in the middle word, a
separator, or the grouped suffix, the V100 complete-word bound changes
\(\|B_t\|\) from \(C\) to \(C/s_\alpha\), while both endpoint energies
remain unpaid.  All cases therefore give the same right side
\eqref{eq:V11-paid-sandwich}.  Since exactly one location pays, none
of these factors is combined with another paid bound.
\end{proof}

\begin{corollary}[Response parts of the sandwich]
\label{cor:V11-response-parts}
The preceding estimates imply
\begin{align}
 \Gamma_\otimes(1,K_t)
 &\le C s_\alpha^2 h_1h_2,
\label{eq:V11-unpaid-response}\\
 \Gamma_\otimes(w,K_t^{\rm pay})
 &\le C s_\alpha h_1h_2.
\label{eq:V11-paid-response}
\end{align}
After summing ordered early-tail and breaker coordinates, the right
sides remain bounded by the corresponding products
\[
 C s_\alpha^2H_{xq}(E)H_{xy}(E),
 \qquad
 C s_\alpha H_{xq}(E)H_{xy}(E),
\]
with no coordinate multiplicity.
\end{corollary}

\begin{proof}
Use \(U(h)\le s_\alpha h\) in
\eqref{eq:V11-unpaid-sandwich}--%
\eqref{eq:V11-paid-sandwich}.  Then apply
\Cref{cor:V10-no-coordinate-loss}.
\end{proof}

\begin{firewall}[The middle tail is not paid twice]
\label{fire:V11-middle-one-use}
The already acquired \(y\to z\) response remains a complete physical
occurrence inside the middle matching word while
\Cref{thm:V11-payer-table} is proved.  Its dominant scalar lower bound
may be inserted only after this joint sandwich.  It is never replaced
first by an independently optimized capacity.
\end{firewall}

\section{The exact missing row denominator}
\label{sec:V11-row-deficit}

Assume
\[
 d_{y\to z}(B_z)
 =
 \frac{H_{yz}(B_z)}{\Xi_y}
 \ge\gamma_y a_*.
\]
Inserting this lower bound into the stock part of the early sandwich
produces
\begin{equation}
 H_{xq}\,H_{xy}\,d_{y\to z}
 =
 \frac{H_{xq}H_{xy}H_{yz}}{\Xi_y}.
\label{eq:V11-stock-three-edge}
\end{equation}
The exact V92 response of the path \(q-x-y-z\) is
\begin{equation}
 {\cal W}_{q-x-y-z}^{\rm agg}
 =
 d_{x\to q}d_{y\to z}H_{xy}
 =
 \frac{H_{xq}H_{xy}H_{yz}}{\Xi_x\Xi_y}.
\label{eq:V11-exact-tree}
\end{equation}

\begin{proposition}[One-row deficit of the stock sandwich]
\label{prop:V11-one-row-deficit}
The stock expression \eqref{eq:V11-stock-three-edge} and the exact
tree response \eqref{eq:V11-exact-tree} satisfy
\begin{equation}
 H_{xq}H_{xy}d_{y\to z}
 =
 \Xi_x\,{\cal W}_{q-x-y-z}^{\rm agg}.
\label{eq:V11-extra-Xi}
\end{equation}
Therefore the one-use sandwich, while operator-valid and
support-uniform, does not by itself supply the missing response
denominator \(\Xi_x^{-1}\).
\end{proposition}

\begin{proof}
Compare \eqref{eq:V11-stock-three-edge} and
\eqref{eq:V11-exact-tree}.  Their numerators and the denominator
\(\Xi_y\) agree; the latter has the additional denominator
\(\Xi_x\).
\end{proof}

\begin{proposition}[Sharp scalar regression for the missing
\(\Xi_x^{-1}\)]
\label{prop:V11-deficit-sharp}
No support-uniform constant \(C\) can satisfy
\begin{equation}
 H_{xq}H_{xy}d_{y\to z}
 \le
 C\,{\cal W}_{q-x-y-z}^{\rm agg}
\label{eq:V11-false-comparison}
\end{equation}
for all clean row weights.
\end{proposition}

\begin{proof}
Choose three disjoint edge coordinates realizing
\[
 H_{xq}=H_{xy}=H_{yz}=a_*^2,
 \qquad
 \Xi_y=2a_*,
\]
with all displayed nonzero edge weights equal to \(a_*\).  Add \(N\)
coordinates private to row \(x\), each of weight \(a_*\).  This leaves
all three edge stocks and \(\Xi_y\) unchanged, while
\[
 \Xi_x=(N+2)a_*.
\]
By \eqref{eq:V11-extra-Xi}, the ratio of the left side of
\eqref{eq:V11-false-comparison} to its right side is \(\Xi_x\), which
is unbounded as \(N\to\infty\).  Every nonzero weight obeys the clean
floor.  The example is a scalar normalization regression, not a
counterexample to an exact Wilson operator inequality.
\end{proof}

\begin{corollary}[Dominance is the legitimate denominator insertion]
\label{cor:V11-dominance-required}
If \(d_{x\to q}\ge\gamma_xa_*\), then
\[
 1\le\frac{d_{x\to q}}{\gamma_xa_*}
\]
legitimately inserts the missing \(\Xi_x^{-1}\), and the V1
complementary-path compiler closes.  If this dominance fails, the raw
early covariance stock cannot replace that probability response by
\Cref{prop:V11-deficit-sharp}.
\end{corollary}

\begin{proof}
The first statement is
\Cref{thm:V1-complementary-partners}.  The second is the sharp failure
of \eqref{eq:V11-false-comparison}.
\end{proof}

\section{Corrected use of the early-tail refinement}
\label{sec:V11-corrected-use}

\begin{theorem}[Early dominance or post-prefix acquisition]
\label{thm:V11-early-or-post}
For a V100 missing-tail endpoint \(x\), the assembled early-tail
refinement has the following rigorous use:
\begin{enumerate}[label=\textup{(\roman*)},leftmargin=4em]
\item if \eqref{eq:V10-early-dominant} holds, the existing
      complementary-path theorem closes the pair;
\item if it fails, the pre-breaker \(xq\)-incidence mass is below
      \(\gamma\Xi_x\) by \eqref{eq:V10-prefix-depleted}, so more than
      \((1-\gamma)\Xi_x\) lies either off the matching edge or at/after
      the breaker coordinate;
\item the complete no-early branch retains the contractive endpoint
      \(R_{<t}^{xq}\), and the failed early covariance is not used as
      a tree mark.
\end{enumerate}
\end{theorem}

\begin{proof}
Item \textup{(i)} is
\Cref{cor:V11-dominance-required}.  Item \textup{(ii)} is the
complement of the mass in \eqref{eq:V10-prefix-depleted}.  Item
\textup{(iii)} follows from
\Cref{cor:V10-two-branches,fire:V10-no-restricted-prefix} and the
deficit in \Cref{prop:V11-one-row-deficit}.
\end{proof}

\begin{assessment}[Bottleneck after the valid sandwich]
\label{ass:V11-bottleneck}
The operator concern raised in V10 is resolved: the early covariance
and breaker do fit in one V97 sandwich with a complete payer table.
The result also proves why this is not the missing MTA estimate.  A
stock-carrying covariance endpoint supplies \(H_{xq}\), whereas a
quartic tree needs the probability response
\(H_{xq}/\Xi_x\).  Private or otherwise off-edge \(x\)-mass makes the
difference arbitrarily large.

The next move is therefore not a stronger Schatten inequality.
It is a post-prefix depleted identity for the large complement in
\Cref{thm:V11-early-or-post}, keeping separate: matching incidences
at the breaker, nonmatching incidences in the breaker bank, later
complete suffix banks, and genuinely private \(x\)-mass.
\end{assessment}

\section{Result ledger and V12 mandate}
\label{sec:V11-results}

\begin{theorem}[Third-series V11 result ledger]
\label{thm:V11-results}
Relative to V1--V10 and the corrected archive boundary, V11 proves:
\begin{enumerate}[label=\textup{(V11.\arabic*)},leftmargin=6.3em]
\item quadratic unpaid and paid energy bounds for the complete early
      covariance and breaker endpoints;
\item one V97 sandwich covering every payer location without a
      product of capacities;
\item support-uniform stock-product response bounds with no ordered
      coordinate multiplicity;
\item the exact extra factor \(\Xi_x\) between that stock product and
      the required quartic path response;
\item a clean scaling family proving that deficit is sharp; and
\item the corrected dominance-or-post-prefix acquisition theorem.
\end{enumerate}
\end{theorem}

\begin{proof}
These are
\Cref{lem:V11-endpoint-energies,
thm:V11-payer-table,
cor:V11-response-parts,
prop:V11-one-row-deficit,
prop:V11-deficit-sharp,
thm:V11-early-or-post}.
\end{proof}

\begin{programme}[V12: exact post-prefix depleted identity]
\label{prog:V11-V12}
\begin{enumerate}[leftmargin=3em]
\item On the failure of \eqref{eq:V10-early-dominant}, define the
      complement of all pre-breaker \(xq\)-incidence coordinates.
\item Partition that complement into the first-breaker bank, other
      already used matching/separator banks, and fresh suffix
      coordinates, before applying the six V79 incidence types.
\item In the breaker bank, distinguish \(xq\) matching incidences from
      \(xy\) and \(xz\) nonmatching incidences.  The latter are already
      component breakers and must not be called matching tails.
\item In later complete banks, use an assembled product difference,
      not a restricted no-tail transfer, to collect all \(xq\)
      attachments.
\item Show that a dominant \(x\to q\) cell closes, while every
      avoid-\(q\) response enters the exact V6 partner table.
\item Keep a private \(x\)-mass branch separate and test its connected
      cancellation before the V100 positive envelope.
\end{enumerate}
\end{programme}

\begin{firewall}[Claim boundary after third-series V11]
\label{fire:V11-boundary}
V11 proves the complete one-use early-tail sandwich and its exact
normalization boundary.  It does not supply the missing
\(\Xi_x^{-1}\) without response dominance, assemble the post-prefix
depleted stock, or close the sensitive, diagonal, nonclean, and other
quartic source families.  Higher MTA, WUFC, CPM, cutoff removal,
reflection positivity, Osterwalder--Schrader reconstruction, and a
uniform positive continuum spectral gap remain open.  The full
Yang--Mills mass gap has not been proved.
\end{firewall}

\paragraph{Parent-version record.}
V11 was copied byte-for-byte from
\texttt{Renewal\_Yang\_Mills\_Mass\_Gap\_Research\_Notes\_V10.tex}
before this chapter was appended.  The V10 parent had \(5\,593\)
source lines, \(205\,286\) bytes, and SHA--256
\[
 \texttt{bf48582bb793648986bf4e20483b1accac090b1f064415c0c580b8efe40be69b}.
\]
The next compulsory companion audit remains V15.

\chapter{V12: Post-prefix depletion and the three-partner
missing-tail table}
\label{ch:V12-post-prefix}

\section{Context and notation}
\label{sec:V12-context}

Continue with the V100 matching
\[
 \mu=\{x,q\}\mid\{y,z\},
\]
the breaker pair \(xy\), the acquired matching response
\[
 d_{y\to z}(B_z)\ge\gamma_y a_*,
\]
and the missing response \(x\to q\).  Let \(t\) be the selected first
breaker coordinate in the exact pre-envelope packet.  V11 proves that
if the pre-breaker response \(d_{x\to q}(<t)\) is dominant, the path
\(q-x-y-z\) closes.  This chapter assumes the complementary branch
\begin{equation}
 H_{xq}(<t)<\gamma a_*\Xi_x,
 \qquad0<\gamma<1.
\label{eq:V12-early-low}
\end{equation}

Define the pre-breaker matching-incidence set
\begin{equation}
 E_-^{xq}(t)
 :=
 \{e<t:a_{x,e}>0,\ a_{q,e}>0\}
\label{eq:V12-early-incidence}
\end{equation}
and its row-\(x\) mass
\[
 A_x^-(t):=\sum_{e\in E_-^{xq}(t)}a_{x,e}.
\]
The Casimir floor and \eqref{eq:V12-early-low} give
\begin{equation}
 A_x^-(t)<\gamma\Xi_x.
\label{eq:V12-early-mass-small}
\end{equation}

Let \(\cU_*(t)\) be the disjointly represented union of all complete
bank roles already present in the V100 packet after the coordinates
in \(E_-^{xq}(t)\) have been assigned to the early covariance.  This
includes the breaker role, matching-word roles, local separators, and
complete suffix roles.  Common coordinates are assigned to their
earliest physical role.  Each nonempty subbank retains its original
anchor, and the number of roles remains bounded by \(M_0\).

\section{The exact post-prefix identity}
\label{sec:V12-identity}

On the complement of
\[
 E_-^{xq}(t)\cup\cU_*(t),
\]
apply the deterministic V79 incidence partition at row \(x\).  Denote
the resulting fresh masses by
\[
 P_x^{\rm fr},\ U_x^{\rm fr},\ A_x^{\rm fr},\
 F_x^{\rm fr},\ D_x^{\rm fr},\ I_x^{\rm fr}.
\]

\begin{lemma}[Post-prefix depleted identity]
\label{lem:V12-post-prefix-identity}
For every fixed first-breaker packet,
\begin{equation}
 \boxed{
 \Xi_x
 =
 A_x^-(t)+X_x(\cU_*(t))
 +P_x^{\rm fr}+U_x^{\rm fr}+A_x^{\rm fr}
 +F_x^{\rm fr}+D_x^{\rm fr}+I_x^{\rm fr}.}
\label{eq:V12-post-prefix-identity}
\end{equation}
All terms are nonnegative and use disjoint coordinate representatives.
No coordinate in the first term can be selected again by any later
alternative.
\end{lemma}

\begin{proof}
First assign every coordinate in
\(E_-^{xq}(t)\) to the complete early covariance.  Among the remaining
coordinates, assign every coordinate occurring in an old complete
bank to its earliest physical role; their union is \(\cU_*(t)\).
Every remaining coordinate is fresh relative to both collections, and
the V79 deterministic labels \(P,U,A,F,D,I\) partition it exactly.
Summing \(a_{x,e}\) over these disjoint sets gives
\eqref{eq:V12-post-prefix-identity}.
\end{proof}

\begin{corollary}[Large post-prefix stock]
\label{cor:V12-large-post-prefix}
Under \eqref{eq:V12-early-low},
\begin{equation}
 X_x(\cU_*(t))
 +P_x^{\rm fr}+U_x^{\rm fr}+A_x^{\rm fr}
 +F_x^{\rm fr}+D_x^{\rm fr}+I_x^{\rm fr}
 >(1-\gamma)\Xi_x.
\label{eq:V12-post-prefix-large}
\end{equation}
\end{corollary}

\begin{proof}
Subtract \eqref{eq:V12-early-mass-small} from
\eqref{eq:V12-post-prefix-identity}.
\end{proof}

\begin{firewall}[The suffix is not refactorized by matching
components]
\label{fire:V12-no-suffix-pair-factor}
Before the breaker, every local partition refines \(\mu\), which
justifies the two-row factorization used in V10.  After the breaker,
the complete suffix contains arbitrary four-row partitions and need
not factor into the two matching components.  V12 therefore treats
suffix coordinates only as anchored complete banks inside
\(\cU_*(t)\); it does not invent a binary suffix covariance
\(T_{>t}^{xq}-R_{>t}^{xq}\).
\end{firewall}

\section{Old-bank versus fresh-stock dichotomy}
\label{sec:V12-old-fresh}

Fix a second threshold \(\delta>0\) satisfying
\begin{equation}
 \gamma+\delta<1.
\label{eq:V12-thresholds}
\end{equation}

\begin{theorem}[One-step post-prefix acquisition]
\label{thm:V12-post-prefix-acquisition}
Under \eqref{eq:V12-early-low}--%
\eqref{eq:V12-thresholds}, one of the following occurs.
\begin{enumerate}[label=\textup{(\roman*)},leftmargin=4em]
\item An old-bank external response
      \begin{equation}
      d_{x\to r}(B_r)
      \ge\frac{\delta}{2M_0}a_*,
      \qquad r\in\{q,y,z\},
      \label{eq:V12-old-response}
      \end{equation}
      is present in a bank of \(\cU_*(t)\).
\item A single old self-anchored bank \(B_x\subseteq\cU_*(t)\) obeys
      \begin{equation}
      X_x(B_x)\ge\frac{\delta}{2M_0}\Xi_x.
      \label{eq:V12-old-self}
      \end{equation}
\item A fresh \(P,U,A,F,\) or \(I\) cell has row-\(x\) mass at least
      \begin{equation}
      \frac{1-\gamma-\delta}{6}\Xi_x
      \label{eq:V12-fresh-complete}
      \end{equation}
      and enters its established complete compiler.
\item A fresh actual-partner response satisfies
      \begin{equation}
      d_{x\to r}(D_r)
      \ge
      \frac{1-\gamma-\delta}{18}a_*,
      \qquad r\in\{q,y,z\}.
      \label{eq:V12-fresh-response}
      \end{equation}
\item The state is nonclean.
\end{enumerate}
The selected old bank in \textup{(i)--(ii)} is disjoint from the weak
early covariance representatives \(E_-^{xq}(t)\).
\end{theorem}

\begin{proof}
Suppose first that
\[
 X_x(\cU_*(t))\ge\delta\Xi_x.
\]
Apply the anchored-overlap dichotomy
\Cref{lem:V1-anchored-overlap,cor:V1-self-witness} to the at most
\(M_0\) remaining bank roles.  An external anchor gives
\eqref{eq:V12-old-response}; otherwise pigeonholing the
self-anchored stock gives \eqref{eq:V12-old-self}.

Suppose instead that
\(X_x(\cU_*(t))<\delta\Xi_x\).  By
\eqref{eq:V12-post-prefix-large}, the six fresh cells total more than
\[
 (1-\gamma-\delta)\Xi_x.
\]
A largest cell has at least the amount in
\eqref{eq:V12-fresh-complete}.  The \(P,U,A,F,I\) cases enter their
established compilers.  In the \(D\) case, partition by the at most
three actual partners \(q,y,z\).  A largest partner cell has
row-\(x\) mass at least
\((1-\gamma-\delta)\Xi_x/18\).  Every coordinate there is active at
that partner, so the Casimir floor gives
\eqref{eq:V12-fresh-response}.  Failure of clean typing is kept as
\textup{(v)}.  Disjointness from the early incidence set is built into
\eqref{eq:V12-post-prefix-identity}.
\end{proof}

\begin{remark}[Fixed constants]
\label{rem:V12-fixed-constants}
For example, taking
\(\gamma=\delta=1/4\) gives fresh complete-cell mass at least
\(\Xi_x/12\), fresh directed dominance at least \(a_*/36\), and
old external dominance at least \(a_*/(8M_0)\).  All constants are
independent of coordinate support, representations, and physical
multiplicity.
\end{remark}

\section{The exact three-partner table}
\label{sec:V12-partner-table}

Any response in
\eqref{eq:V12-old-response} or
\eqref{eq:V12-fresh-response} has one of the three actual partners
\(q,y,z\).

\begin{proposition}[Missing-tail partner table]
\label{prop:V12-partner-table}
Retain the raw breaker edge \(xy\) and the already acquired response
\(y\to z\).  A new response \(x\to r\) has exactly the following
fate:
\begin{center}
\begin{tabular}{c c c}
\toprule
partner \(r\) & literal edges & outcome\\
\midrule
\(q\) & \(xq,xy,yz\) & path \(q-x-y-z\), closes\\
\(y\) & \(xy,xy,yz\) & repeated breaker, \(q\) isolated\\
\(z\) & \(xz,xy,yz\) & triangle on \(x,y,z\), \(q\) isolated\\
\bottomrule
\end{tabular}
\end{center}
Thus \(q\) is the unique successful partner within the
two-response quartic tree atlas.
\end{proposition}

\begin{proof}
For \(r=q\), the edges form the four-vertex path
\(q-x-y-z\), and the response tails \(x,y\) are its two degree-two
vertices.  Hence \Cref{thm:V1-complementary-partners} applies.

For \(r=y\), the new response edge repeats the raw breaker \(xy\), so
only two distinct edges occur and \(q\) is isolated.  For \(r=z\), the
three distinct edges form the triangle \(xy,yz,zx\) on three rows,
again leaving \(q\) isolated.  Neither graph is a spanning tree.
\end{proof}

\begin{corollary}[Post-prefix directed normal form]
\label{cor:V12-directed-normal-form}
Every directed or externally anchored branch of
\Cref{thm:V12-post-prefix-acquisition} either closes through the
desired matching partner \(q\), or is one of two explicitly labelled
avoid-\(q\) cells: the repeated-breaker cell \(r=y\) or the
three-cycle cell \(r=z\).
\end{corollary}

\begin{proof}
Apply \Cref{prop:V12-partner-table}.
\end{proof}

\begin{firewall}[A component breaker is not a matching tail]
\label{fire:V12-breaker-not-tail}
The responses \(x\to y\) and \(x\to z\) are actual
component-breaking incidences, but neither supplies the missing
matching edge \(xq\).  V12 retains their labels and does not use graph
connectivity on the three-row subgraph as a substitute for four-row
spanning.
\end{firewall}

\section{The self-anchored old-bank branch}
\label{sec:V12-self}

\begin{theorem}[Post-prefix self-bank reduction]
\label{thm:V12-self-reduction}
Let \(B_x\) be selected by
\eqref{eq:V12-old-self}.  Within the clean calculus it has one of:
\begin{enumerate}[label=\textup{(\alph*)},leftmargin=4em]
\item a common \(P,U,A,F,I,\) or directed type covered by the
      established same-type self-bank compilers;
\item a partition-changing role whose response-first V2 reduction
      yields \(x\to r\), \(r\in\{q,y,z\}\);
\item a low-response \(x\)-sensitive correction plus an \(x\)-free
      role difference, as in
      \eqref{eq:V2-low-response-inert} and
      \eqref{eq:V2-insensitive-factor}.
\end{enumerate}
In \textup{(b)}, the partner \(q\) closes and the partners \(y,z\)
are exactly the two failed cells of
\Cref{prop:V12-partner-table}.
\end{theorem}

\begin{proof}
Common clean types are
\Cref{prop:V1-self-type-routes}.  If no common type exists, apply
\Cref{thm:V2-self-role-reduction} with anchor \(x\).  Its
response-first alternative gives an actual changed-partner response;
its complement is precisely the sensitive/insensitive factorization.
Apply \Cref{prop:V12-partner-table} to the response partner.
\end{proof}

\begin{assessment}[The depleted search no longer returns to the early
bank]
\label{ass:V12-no-return}
The early \(xq\) covariance has been removed as a disjoint coordinate
representative before the post-prefix identity is applied.  Therefore
an old-bank overlap selected by
\Cref{thm:V12-post-prefix-acquisition} is a genuinely different
physical role or subbank.  A fresh cell is disjoint from both.  The
only continuing clean branches are now:
\[
\begin{array}{c}
\text{two named avoid-\(q\) response cells},\\
\text{one partition-changing self-bank low-response correction},\\
\text{or a complete compiler whose existing theorem must be inserted
with the V100 payer ledger}.
\end{array}
\]
This is a strict finite reduction of the vague instruction to repeat
depleted acquisition.
\end{assessment}

\section{A no-hidden-mass check}
\label{sec:V12-mass-check}

\begin{proposition}[Threshold accounting is sharp at the scalar level]
\label{prop:V12-threshold-sharp}
The lower bound
\((1-\gamma-\delta)\Xi_x\) for the total fresh stock in the small-old-
bank branch cannot be improved using only
\eqref{eq:V12-early-mass-small} and
\(X_x(\cU_*)<\delta\Xi_x\).
\end{proposition}

\begin{proof}
For any \(\varepsilon>0\), choose disjoint coordinate sets carrying
row-\(x\) masses
\[
 (\gamma-\varepsilon)\Xi_x,\qquad
 (\delta-\varepsilon)\Xi_x,\qquad
 (1-\gamma-\delta+2\varepsilon)\Xi_x
\]
in the early, old-bank, and fresh classes respectively.  Choose
weights as integer multiples of the floor after a harmless common
rescaling and approximation.  Both strict hypotheses hold, while the
fresh mass tends to
\((1-\gamma-\delta)\Xi_x\) as \(\varepsilon\downarrow0\).
\end{proof}

\begin{firewall}[Scalar sharpness is not an operator
counterexample]
\label{fire:V12-sharp-scope}
\Cref{prop:V12-threshold-sharp} concerns only the information content
of the two mass inequalities.  It neither asserts that all three
stocks occur in one nonzero Wilson cumulant nor obstructs additional
cancellation in the exact source word.
\end{firewall}

\section{Result ledger and V13 mandate}
\label{sec:V12-results}

\begin{theorem}[Third-series V12 result ledger]
\label{thm:V12-results}
Relative to V1--V11 and the corrected archive boundary, V12 proves:
\begin{enumerate}[label=\textup{(V12.\arabic*)},leftmargin=6.3em]
\item the exact post-prefix depleted identity
      \eqref{eq:V12-post-prefix-identity};
\item a support-uniform old-bank/fresh-stock acquisition theorem which
      cannot reselect the weak early covariance;
\item the complete three-partner table for a missing matching tail;
\item reduction of the self-anchored post-prefix bank to same-type,
      response-first, or low-response role-change alternatives;
\item a precise list of the remaining clean branches; and
\item scalar sharpness of the fresh-stock threshold.
\end{enumerate}
\end{theorem}

\begin{proof}
These are
\Cref{lem:V12-post-prefix-identity,
thm:V12-post-prefix-acquisition,
prop:V12-partner-table,
thm:V12-self-reduction,
ass:V12-no-return,
prop:V12-threshold-sharp}.
\end{proof}

\begin{programme}[V13: assemble the two avoid-\(q\) cells]
\label{prog:V12-V13}
\begin{enumerate}[leftmargin=3em]
\item For the repeated-breaker cell \(x\to y\), keep the response and
      raw breaker as two roles of one complete breaker word.  Test
      whether the V98 joint-carrier method produces a third distinct
      matching edge or only a duplicate mark.
\item For the cycle cell \(x\to z\), retain the triangle
      \(xy,yz,zx\) and return to the exact connected-source subtraction
      before the positive pair envelope.  Determine where the isolated
      row \(q\) first attaches.
\item In both cells, use the V12 post-prefix bank as the endpoint of
      one Schatten sandwich; do not multiply its capacity by the
      V100 breaker capacity.
\item Treat the \(x\)-sensitive low-response correction before the
      \(x\)-free role difference is normed, and test whether its first
      actual \(q\)-incidence supplies the missing tail.
\item Keep complete-compiler and nonclean branches separate from this
      finite directed table.
\end{enumerate}
\end{programme}

\begin{firewall}[Claim boundary after third-series V12]
\label{fire:V12-boundary}
V12 gives an exact post-prefix acquisition which cannot cycle back to
the weak early-tail bank and reduces all directed outcomes to one
closing and two failed partner cells.  It does not close those two
avoid-\(q\) cells, the low-response self-role correction, complete-
compiler payer interfaces, nonclean states, or the remaining quartic
source families.  Higher MTA, WUFC, CPM, cutoff removal, reflection
positivity, Osterwalder--Schrader reconstruction, and a uniform
positive continuum spectral gap remain open.  The full Yang--Mills
mass gap has not been proved.
\end{firewall}

\paragraph{Parent-version record.}
V12 was copied byte-for-byte from
\texttt{Renewal\_Yang\_Mills\_Mass\_Gap\_Research\_Notes\_V11.tex}
before this chapter was appended.  The V11 parent had \(6\,000\)
source lines, \(219\,231\) bytes, and SHA--256
\[
 \texttt{8a83f5bde0b0c2ff47f6041ca6d4707943060b5d353bf33956388eec2f8c1b6a}.
\]
The next compulsory companion audit remains V15.

\chapter{V13: Spectator-stable completion of the two avoid-tail
cells}
\label{ch:V13-spectator-fan}

\section{Context and nonduplication}
\label{sec:V13-context}

V12 leaves two directed partner cells after post-prefix acquisition at
the missing endpoint \(x\): \(x\to y\), which repeats the breaker, and
\(x\to z\), which completes a triangle on \(x,y,z\).  Neither cell
connects the missing mate \(q\).

It would be wasteful and potentially incorrect to build a four-edge
numerator from either cell.  The V4 fan theorem already shows that a
response rooted at \(q\), after direct or mass-compatible reversal,
closes against one existing endpoint response and the raw breaker.
The only new issue is operator bookkeeping: the V12 avoid-\(q\)
occurrence is now present in the word.  V13 proves that it may remain
inside the bounded directed middle word while the closing tree uses
only the other three literal edges.

\section{The isolated-row fan}
\label{sec:V13-fan}

Retain
\[
 \mu=\{x,q\}\mid\{y,z\},
 \qquad
 d_{y\to z}(B_z)\ge\gamma_z a_*,
\]
and the raw breaker stock \(H_{xy}(G)\).  Let \(D_{\rm av}\) denote
zero or one additional complete directed occurrence \(x\to y\) or
\(x\to z\) produced by V12.  The V99 directed-word theorem gives
\begin{equation}
 \|D_{\rm av}\|\le C_{\rm dir},
\label{eq:V13-spectator-norm}
\end{equation}
with the corresponding once-paid bound when its physical occurrence
contains the sole payer.

\begin{theorem}[Spectator-stable isolated-row fan]
\label{thm:V13-spectator-fan}
Suppose an actual response incident to \(q\) is available in the
orientation
\begin{equation}
 d_{r\to q}(B_q)\ge\eta a_*,
 \qquad r\in\{x,y,z\},
\label{eq:V13-reversed-q}
\end{equation}
and is mark-admissible with the existing \(y\to z\) response and the
breaker.  Then every unpaid or once-paid V12 avoid-\(q\) packet obeys
\begin{equation}
 \boxed{
 Q_{xy,{\rm av}}^{(\varepsilon)}
 \le
 \frac{CC_{\rm dir}}{\gamma_z\eta a_*^2}\,
 d_{y\to z}(B_z)d_{r\to q}(B_q)H_{xy}(G),
 \qquad\varepsilon\in\{0,1\}.}
\label{eq:V13-spectator-fan}
\end{equation}
The response \(D_{\rm av}\) is not used as a scalar mark.  The three
displayed edges form:
\begin{center}
\begin{tabular}{c c c}
\toprule
\(r\) & tree & response-tail multiset\\
\midrule
\(x\) & \(q-x-y-z\) & \(\{x,y\}\)\\
\(y\) & star centered at \(y\) & \(\{y,y\}\)\\
\(z\) & \(x-y-z-q\) & \(\{y,z\}\)\\
\bottomrule
\end{tabular}
\end{center}
\end{theorem}

\begin{proof}
Keep \(D_{\rm av}\) in the middle ordinary-norm part of the complete
directed word and use \eqref{eq:V13-spectator-norm}.  The already
proved pair-localized estimate therefore remains
\[
 Q_{xy,{\rm av}}^{(\varepsilon)}
 \le CC_{\rm dir}H_{xy}(G)
\]
at every admitted payer location.  The two dominant lower bounds give
\[
 1\le
 \frac{d_{y\to z}(B_z)d_{r\to q}(B_q)}
      {\gamma_z\eta a_*^2}.
\]
Multiplication proves \eqref{eq:V13-spectator-fan}.

For \(r=x\), the edges \(xq,xy,yz\) form the path \(q-x-y-z\), whose
internal vertices are \(x,y\).  For \(r=y\), the edges
\(yq,yx,yz\) form the star centered at \(y\), and both response
denominators are \(\Xi_y\).  For \(r=z\), the edges \(xy,yz,zq\)
form the path \(x-y-z-q\), whose internal vertices are \(y,z\).
\Cref{lem:V3-degree-form} verifies the exact aggregate normalization
in every row.  The unused V12 occurrence is bounded only once inside
the middle word.
\end{proof}

\begin{corollary}[Forward \(q\)-response fate]
\label{cor:V13-forward-fate}
Suppose instead that acquisition at \(q\) gives
\[
 d_{q\to r}(B_q)\ge\eta_0a_*.
\]
For any \(L\ge1\), either
\(\Xi_r\le L\Xi_q\), in which case
\Cref{thm:V13-spectator-fan} applies with
\(\eta=\eta_0/L\), or \(\Xi_r>L\Xi_q\).  More sharply, for any direct
threshold \(\eta>0\), either \(d_{r\to q}\ge\eta a_*\) and the fan
closes, or
\[
 X_r(E_{qr}(B_q))<\eta\Xi_r.
\]
\end{corollary}

\begin{proof}
The mass-comparison statement is
\Cref{lem:V4-response-reversal}.  The direct statement is
\Cref{thm:V5-reverse-or-deplete}, followed in its first alternative
by \Cref{thm:V13-spectator-fan}.
\end{proof}

\begin{remark}[The V12 response can be ignored, not deleted]
\label{rem:V13-ignore-not-delete}
The extra occurrence \(D_{\rm av}\) remains in the assembled operator
and may contain the payer.  V13 ignores only its scalar lower bound
when choosing the final three tree marks.  It does not remove the
physical factor or claim that a triangle edge has norm at most one
without the V99 complete directed-word estimate.
\end{remark}

\section{Two-sided acquisition normal form}
\label{sec:V13-two-sided}

\begin{theorem}[Endpoint-then-mate acquisition]
\label{thm:V13-two-sided-normal-form}
Apply \Cref{thm:V12-post-prefix-acquisition} first at the missing
breaker endpoint \(x\).  On either directed avoid-\(q\) branch, apply
the rowwise acquisition
\Cref{lem:V4-rowwise-transport} at the mate \(q\), with every existing
bank, including \(D_{\rm av}\), placed in the used union.  Every clean
branch is then one of:
\begin{enumerate}[label=\textup{(\roman*)},leftmargin=4em]
\item a complete \(P,U,A,F,I\) compiler at \(x\) or \(q\);
\item the desired endpoint response \(x\to q\), closing by
      \Cref{prop:V12-partner-table};
\item a \(q\)-rooted response whose reverse is dominant, closing by
      \Cref{thm:V13-spectator-fan};
\item a \(q\)-rooted response whose reverse fails dominance, forcing
      the target-row depletion of
      \Cref{thm:V5-reverse-or-deplete};
\item a same-type self-bank compiler at \(x\) or \(q\);
\item a partition-changing self-bank which gives an actual response
      and re-enters \textup{(ii)--(iv)}, or leaves the V2
      low-response sensitive/insensitive factorization.
\end{enumerate}
Outside the clean calculus there is a labelled nonclean state.
\end{theorem}

\begin{proof}
The \(x\)-side alternatives are
\Cref{thm:V12-post-prefix-acquisition,
thm:V12-self-reduction}.  The partner \(q\) closes, while partners
\(y,z\) produce \(D_{\rm av}\) by
\Cref{cor:V12-directed-normal-form}.

At row \(q\), \Cref{lem:V4-rowwise-transport} is exhaustive.  Its
complete types give \textup{(i)}, and its external or directed
responses have an actual partner \(r\in\{x,y,z\}\).  Apply
\Cref{cor:V13-forward-fate} to obtain \textup{(iii)--(iv)}.
The self-bank alternatives are
\Cref{prop:V1-self-type-routes,thm:V2-self-role-reduction}, giving
\textup{(v)--(vi)}.
\end{proof}

\begin{corollary}[The two V12 graph cells add no new topology]
\label{cor:V13-no-new-topology}
The repeated-breaker and three-cycle cells do not create independent
quartic graph obstructions.  After acquisition is moved to \(q\),
both have exactly the same closing and residual alternatives as a
packet without \(D_{\rm av}\).
\end{corollary}

\begin{proof}
The complete directed occurrence changes only the fixed middle norm
in \Cref{thm:V13-spectator-fan}.  The three closing trees and the
reverse-or-deplete residual do not depend on whether
\(D_{\rm av}=x\to y\), \(x\to z\), or is absent.
\end{proof}

\section{A finite direct-reversal threshold}
\label{sec:V13-threshold}

Fix \(0<\eta<1\).  For a \(q\)-rooted response \(q\to r\), call its
incidence bank \(\eta\)-reversible when
\[
 d_{r\to q}\ge\eta a_*.
\]

\begin{proposition}[Uniform reversible/nonreversible split]
\label{prop:V13-reversible-split}
For each of the at most three partner cells at \(q\), the reversible
part closes with the common constant
\[
 \frac{CC_{\rm dir}}{\gamma_z\eta a_*^2}.
\]
On the nonreversible part,
\[
 X_r(E_{qr})<\eta\Xi_r.
\]
Thus all nonclosed directed branches have a quantitatively named
target row and a quantitatively depleted literal edge.
\end{proposition}

\begin{proof}
Apply \Cref{thm:V13-spectator-fan} on the reversible cell and
\Cref{thm:V5-reverse-or-deplete} on its complement.  There are only
three labelled partner cells, so the same fixed threshold and maximum
constant apply to all.
\end{proof}

\begin{firewall}[Depleted targets are not yet transported packets]
\label{fire:V13-target-transport}
The nonreversible conclusion locates large target-row mass off the
literal \(qr\) edge.  It does not prove that this mass occurs in a
complete bank which can be inserted beside the original V100 breaker
and payer.  The V6 targeted acquisition theorem applies only after
that operator eligibility is verified.
\end{firewall}

\section{Regression checks and result ledger}
\label{sec:V13-results}

\begin{proposition}[V13 regression suite]
\label{prop:V13-regressions}
\begin{enumerate}[label=\textup{(\roman*)},leftmargin=4em]
\item The repeated \(xy\) response is never used as a second copy of
      the raw breaker mark.
\item The triangle \(xy,yz,zx\) is never declared spanning while
      \(q\) is isolated.
\item Only three selected edge stocks enter the final tree; the
      unused directed occurrence remains in ordinary norm.
\item A forward \(q\to r\) edge is reversed only by an actual reverse
      lower bound or the explicit mass comparison.
\item The payer remains in its original occurrence even when that
      occurrence is not selected as a scalar tree mark.
\end{enumerate}
\end{proposition}

\begin{proof}
Items \textup{(i)--(iii)} are built into
\Cref{thm:V13-spectator-fan,rem:V13-ignore-not-delete}.
Item \textup{(iv)} is
\Cref{cor:V13-forward-fate}.  Item \textup{(v)} follows because
\eqref{eq:V13-spectator-norm} uses the paid complete-word estimate
without moving the payer.
\end{proof}

\begin{theorem}[Third-series V13 result ledger]
\label{thm:V13-results}
Relative to V1--V12 and the corrected archive boundary, V13 proves:
\begin{enumerate}[label=\textup{(V13.\arabic*)},leftmargin=6.3em]
\item a spectator-stable path/star compiler for every response
      incident to the isolated row \(q\);
\item closure of its directly or mass-compatibly reversible parts;
\item an exhaustive endpoint-then-mate acquisition normal form;
\item collapse of both V12 directed failures to the already named
      reverse-or-deplete and self-role residuals; and
\item a uniform three-cell reversible/nonreversible partition with a
      named depleted target edge in every failed cell.
\end{enumerate}
\end{theorem}

\begin{proof}
These are
\Cref{thm:V13-spectator-fan,
cor:V13-forward-fate,
thm:V13-two-sided-normal-form,
cor:V13-no-new-topology,
prop:V13-reversible-split}.
\end{proof}

\begin{programme}[V14: packet eligibility at a depleted target]
\label{prog:V13-V14}
\begin{enumerate}[leftmargin=3em]
\item Fix a nonreversible partner cell \(q\to r\) and its small
      incidence subbank \(E_{qr}\).
\item Return to the exact V100 word and factor every physical bank
      across \(E_{qr}\) and its complement before norms, preserving
      one occurrence and one payer.
\item Determine which complement factors lie in the prefix, breaker,
      matching word, separator, or complete suffix roles.
\item Prove an eligibility theorem for the large
      \((1-\eta)\Xi_r\) off-edge mass, or isolate the exact role whose
      tensor factorization fails.
\item Once eligible, apply the V6 targeted acquisition table.  A
      response \(r\to q\) closes; avoid-\(q\) responses remain the six
      explicitly classified cells.
\item Keep the V2 low-response self-bank and nonclean branches outside
      this targeted calculation.
\end{enumerate}
\end{programme}

\begin{firewall}[Claim boundary after third-series V13]
\label{fire:V13-boundary}
V13 proves that the two V12 directed graph failures are harmless
spectators once any reversible \(q\)-incidence is acquired.  It does
not prove operator eligibility of the large off-edge target mass,
close nonreversible target cells, the V2 low-response self-role
factor, complete-compiler payer interfaces, or nonclean states.
Higher MTA, WUFC, CPM, cutoff removal, reflection positivity,
Osterwalder--Schrader reconstruction, and a uniform positive
continuum spectral gap remain open.  The full Yang--Mills mass gap
has not been proved.
\end{firewall}

\paragraph{Parent-version record.}
V13 was copied byte-for-byte from
\texttt{Renewal\_Yang\_Mills\_Mass\_Gap\_Research\_Notes\_V12.tex}
before this chapter was appended.  The V12 parent had \(6\,430\)
source lines, \(234\,499\) bytes, and SHA--256
\[
 \texttt{7ebcf73f0d08a17bdae123cd4256b19975160318f946530f46c5aec8668635d9}.
\]
The next compulsory companion audit remains V15.

% =============================================================================
% BEGIN APPENDED THIRD-SERIES V14 MATERIAL
% =============================================================================

\clearpage
\chapter{V14: Subset-resolved packet eligibility at a depleted target}
\label{ch:V14-eligibility}

\section{Context and the precise upstream correction}
\label{sec:V14-context}

Fix a nonreversible branch from V13.  Thus an actual response
\[
 d_{q\to r}(B_q)\ge\eta_0a_*
\]
has been acquired, but at a chosen direct reversal threshold
\(0<\eta<1\),
\begin{equation}
 d_{r\to q}(B_q)<\eta a_*.
\label{eq:V14-reverse-fails}
\end{equation}
Writing
\[
 E:=E_{qr}(B_q),
\]
\Cref{thm:V5-reverse-or-deplete} gives
\begin{equation}
 X_r(E)<\eta\Xi_r.
\label{eq:V14-E-small}
\end{equation}

V6 observed that complete moment banks factor across \(E\) and
\(E^c\), but correctly left packet eligibility as a firewall.  The
reason is that the V100 word contains complete covariances as well as
moments.  A covariance does not factor into the tensor product of its
two restrictions.  V14 supplies the missing exact two-term
covariance telescope, applies it simultaneously to the finitely many
physical occurrences in the first-breaker packet, and only then uses
the depleted row identity.  This is an upstream operator statement,
not another scalar pigeonhole.

\section{Exact restriction of moments and covariances}
\label{sec:V14-restriction}

Let \(\Omega\) be the finite physical coordinate support of a fixed
V100 first-breaker packet.  For \(S\subseteq\Omega\) and a local
partition state \(\rho\), write
\[
 M_S(\rho):=\bigotimes_{e\in S}M_e(\rho),
 \qquad
 \Delta_S^{\rho,\sigma}:=M_S(\rho)-M_S(\sigma).
\]
Empty tensor products are identities and empty differences vanish.

\begin{lemma}[One-subset covariance telescope]
\label{lem:V14-covariance-telescope}
For every \(E\subseteq\Omega\), every \(S\subseteq\Omega\), and every
pair of partition states \(\rho,\sigma\),
\begin{align}
 M_S(\rho)
 &=
 M_{S\cap E}(\rho)\otimes
 M_{S\cap E^c}(\rho),
\label{eq:V14-moment-split}\\
 \Delta_S^{\rho,\sigma}
 &=
 \Delta_{S\cap E}^{\rho,\sigma}
 \otimes M_{S\cap E^c}(\rho)
 +
 M_{S\cap E}(\sigma)
 \otimes\Delta_{S\cap E^c}^{\rho,\sigma}.
\label{eq:V14-covariance-split}
\end{align}
These are identities before physical absolute values, operator norms,
capacities, response marks, or payer bounds are taken.
\end{lemma}

\begin{proof}
The first line is associativity and the canonical permutation of
coordinate tensor factors.  Put
\[
 A_\rho=M_{S\cap E}(\rho),\qquad
 B_\rho=M_{S\cap E^c}(\rho),
\]
and similarly for \(\sigma\).  Then
\[
 A_\rho\otimes B_\rho-A_\sigma\otimes B_\sigma
 =
 (A_\rho-A_\sigma)\otimes B_\rho
 +A_\sigma\otimes(B_\rho-B_\sigma),
\]
which is \eqref{eq:V14-covariance-split}.
\end{proof}

\begin{firewall}[The false covariance factorization]
\label{fire:V14-no-product-covariance}
In general,
\[
 \Delta_S^{\rho,\sigma}
 \ne
 \Delta_{S\cap E}^{\rho,\sigma}
 \otimes\Delta_{S\cap E^c}^{\rho,\sigma}.
\]
The right side is quadratic in the partition change and omits both
mixed endpoint terms.  Every restriction below uses the exact sum in
\eqref{eq:V14-covariance-split}; no covariance occurrence is replaced
by two independent small factors.
\end{firewall}

\begin{lemma}[Restriction of the one-payer label]
\label{lem:V14-payer-split}
Suppose one complete occurrence supported on \(S\) contains the
unique payer.  Before the payer-coordinate sum is estimated, its
label splits disjointly into
\[
 \{u\in S\cap E\}\ \dot\cup\ \{u\in S\cap E^c\}.
\]
Consequently the restriction expansion has exactly one payer in each
nonzero summand.  The same statement holds for the grouped complete
suffix payer.
\end{lemma}

\begin{proof}
The V2 payer allocation is a sum over one labelled physical
coordinate \(u\in S\).  Insert
\[
 1_{u\in S}=1_{u\in S\cap E}+1_{u\in S\cap E^c}
\]
before taking any norm.  The two indicators are disjoint, and neither
creates a second numerator or resolvent.  For the grouped suffix
allocation, split the barycentric coordinate sum into the same two
disjoint sets.  This is the exact additivity used in the f2 theorem
\texttt{thm:V96-support-uniform-suffix}.
\end{proof}

\begin{theorem}[Subset-resolved V100 packet]
\label{thm:V14-subset-resolved-packet}
Let \(K_{\mu,G}^{(\varepsilon)}\) be a fixed unpaid or once-paid
V100 first-breaker packet, including its matching word, breaker,
finite separators, and complete common suffix.  For every
\(E\subseteq\Omega\), its exact pre-envelope operator is a sum of at
most \(2^{J_*}\) restriction-resolved words, where \(J_*\) is a
universal bound on the number of complete covariance occurrences.
In every summand:
\begin{enumerate}[label=\textup{(\roman*)},leftmargin=4em]
\item each covariance occurrence is localized to exactly one of
      \(E\) and \(E^c\), while an endpoint moment remains on the other
      side;
\item every physical occurrence retains its position in the word;
\item an unpaid word remains unpaid and a paid word contains exactly
      one payer at its original occurrence and side; and
\item the V94 output dualizer and finite regrouping separators remain
      in their original bounded separator positions.
\end{enumerate}
The expansion count is independent of \(|\Omega|\), the input
representations, and physical multiplicities.
\end{theorem}

\begin{proof}
The matching word contains a universally bounded number of directed
covariance occurrences, the breaker contributes one complete
covariance, and the remaining prefix, separator, and suffix factors
are complete moments or the fixed finite regrouping maps described in
f2 V94--V100.  Apply
\eqref{eq:V14-covariance-split} once to each of the \(J\le J_*\)
covariance occurrences and
\eqref{eq:V14-moment-split} to every moment occurrence.  Distribute
the resulting binary sums.  This gives at most \(2^J\) words.

Multiplication by the unchanged finite separators preserves the
operator identity, so they need not be factorized.  The archive
construction absorbs every output dualizer into such a separator
before the V100 norm is taken.  Apply
\Cref{lem:V14-payer-split} at the unique paid occurrence.  Thus each
resulting word has one covariance choice per original occurrence and
at most one payer.  Since \(J_*\) depends only on the fixed four-row
atlas, the number of words is support-uniform.
\end{proof}

\begin{remark}[What the theorem does and does not estimate]
\label{rem:V14-resolution-not-bound}
\Cref{thm:V14-subset-resolved-packet} is an exact finite
normal form.  It does not claim that all resolved covariances are
small.  Their already proved ordinary, paid, energy, or capacity
bounds are applied only after the resolved word has been selected.
This order preserves the V100 breaker response and prevents both
coordinate counting and double payment.
\end{remark}

\section{The exact depleted-target identity inside the packet}
\label{sec:V14-target-identity}

Apply \Cref{thm:V14-subset-resolved-packet} with the set \(E\) from
\eqref{eq:V14-E-small}.  Remove the \(E\)-representatives first.
Among all remaining complete bank roles already present in the
packet, assign a common physical coordinate to its earliest role.
Refine the fixed finite role list by its recorded actual anchor when
necessary.  Including the V13 spectator and the \(q\to r\)
occurrence, the number of nonempty roles is bounded by a universal
constant
\begin{equation}
 m\le M_*.
\label{eq:V14-role-count}
\end{equation}
Write their disjoint union as
\[
 \cU_r^{\rm off}=\dot\bigcup_{j=1}^m B_j^{\rm off},
 \qquad
 B_j^{\rm off}\subseteq E^c.
\]
Every \(B_j^{\rm off}\) retains the one physical occurrence and
actual anchor of its parent role.

On the coordinates outside
\(E\cup\cU_r^{\rm off}\), apply the deterministic V79 partition at
row \(r\), and denote the six masses by
\[
 P_r^{\rm fr},\ U_r^{\rm fr},\ A_r^{\rm fr},\
 F_r^{\rm fr},\ D_r^{\rm fr},\ I_r^{\rm fr}.
\]

\begin{lemma}[Packet-resolved target identity]
\label{lem:V14-target-identity}
For every fixed restriction-resolved packet,
\begin{equation}
 \boxed{
 \Xi_r
 =
 X_r(E)+X_r(\cU_r^{\rm off})
 +P_r^{\rm fr}+U_r^{\rm fr}+A_r^{\rm fr}
 +F_r^{\rm fr}+D_r^{\rm fr}+I_r^{\rm fr}.}
\label{eq:V14-target-identity}
\end{equation}
All displayed coordinate representatives are disjoint.  Each old
term is already an identified factor of the resolved word, and each
fresh term lies in a complete moment factor disjoint from all
previous response marks.
\end{lemma}

\begin{proof}
First assign coordinates in \(E\).  On its complement, assign every
coordinate belonging to a pre-existing complete physical role to the
earliest such role.  The remaining row-\(r\) incidences are
partitioned exactly once by the six V79 types.  These sets are
disjoint and exhaust the support of the weights \(a_{r,e}\); summing
those weights proves \eqref{eq:V14-target-identity}.

Operator eligibility follows from
\Cref{thm:V14-subset-resolved-packet}.  Old representatives occur on
the \(E^c\) side of an identified word factor.  For a fresh cell
\(Z\), the residual complete moment factors further as
\[
 M_F(\rho)=M_Z(\rho)\otimes M_{F\setminus Z}(\rho).
\]
The second factor stays a complete spectator.  Hence the cell is
inserted one-sidedly before norms, exactly as in the f2 fresh-cell
theorem \texttt{thm:V84-fresh-cell}.
\end{proof}

\begin{firewall}[Scalar partition versus packet eligibility]
\label{fire:V14-identity-is-operator}
The equality \eqref{eq:V14-target-identity} alone would only locate
mass.  Eligibility is supplied by the preceding restriction-resolved
operator theorem: every selected old or fresh subbank is a factor of
the same V100 word, with the original breaker pair and sole payer
still present.  No auxiliary copy of row \(r\) is introduced.
\end{firewall}

\section{Eligibility theorem after failed reversal}
\label{sec:V14-eligibility-theorem}

Fix \(0<\delta<1-\eta\).

\begin{theorem}[Depleted-target packet eligibility]
\label{thm:V14-target-eligibility}
Under \eqref{eq:V14-reverse-fails}, every clean
restriction-resolved V100 packet has at least one of the following
operator-eligible alternatives:
\begin{enumerate}[label=\textup{(\roman*)},leftmargin=4em]
\item an old-bank external response
      \begin{equation}
      d_{r\to s}(B_s^{\rm off})
      \ge\frac{\delta}{2M_*}a_*,
      \qquad s\ne r;
      \label{eq:V14-old-response}
      \end{equation}
\item one old self-anchored bank satisfying
      \begin{equation}
      X_r(B_r^{\rm off})
      \ge\frac{\delta}{2M_*}\Xi_r;
      \label{eq:V14-old-self}
      \end{equation}
\item a fresh \(P,U,A,F,\) or \(I\) cell of row-\(r\) mass at least
      \begin{equation}
      \frac{1-\eta-\delta}{6}\Xi_r;
      \label{eq:V14-fresh-complete}
      \end{equation}
\item a fresh actual-partner response
      \begin{equation}
      d_{r\to s}(D_s^{\rm fr})
      \ge
      \frac{1-\eta-\delta}{18}a_*,
      \qquad s\ne r;
      \label{eq:V14-fresh-response}
      \end{equation}
\item a nonclean representation state.
\end{enumerate}
In \textup{(i)--(iv)} the selected occurrence is in the same exact
packet as the original \(q\to r\) response, breaker, spectator word,
and payer.  All routing and analytic constants are support-uniform.
\end{theorem}

\begin{proof}
By \eqref{eq:V14-E-small} and
\eqref{eq:V14-target-identity},
\begin{equation}
 X_r(\cU_r^{\rm off})
 +P_r^{\rm fr}+U_r^{\rm fr}+A_r^{\rm fr}
 +F_r^{\rm fr}+D_r^{\rm fr}+I_r^{\rm fr}
 >(1-\eta)\Xi_r.
\label{eq:V14-large-off-edge}
\end{equation}
If
\[
 X_r(\cU_r^{\rm off})\ge\delta\Xi_r,
\]
apply the anchored-overlap dichotomy
\Cref{lem:V1-anchored-overlap,cor:V1-self-witness} to the at most
\(M_*\) restricted roles.  Its external branch is
\eqref{eq:V14-old-response}; its self branch is
\eqref{eq:V14-old-self}.

Otherwise the six fresh cells in
\eqref{eq:V14-large-off-edge} total more than
\((1-\eta-\delta)\Xi_r\).  A largest cell has at least the amount in
\eqref{eq:V14-fresh-complete}.  The five complete types enter their
established one-sided compilers.  In the \(D\)-case, partition by the
at most three actual partners.  A largest partner cell has
row-\(r\) mass at least
\((1-\eta-\delta)\Xi_r/18\).  Every coordinate in it is active at
that partner, so the Casimir floor gives
\eqref{eq:V14-fresh-response}.  Failure of clean typing gives
\textup{(v)}.

All selections occur after the exact operator resolution in
\Cref{thm:V14-subset-resolved-packet}; hence they preserve the
original word and one-payer ledger.  The only extra routing factor is
bounded by \(2^{J_*}\), and all subsequent partitions have fixed
four-row size.
\end{proof}

\begin{corollary}[The V6 eligibility hypothesis is discharged]
\label{cor:V14-V6-unconditional}
For a clean V100 first-breaker packet, the off-edge acquisition
theorem \Cref{thm:V6-targeted-acquisition} no longer requires an
additional eligibility assumption.  More explicitly, after a failed
reverse \(q\to r\):
\begin{enumerate}[label=\textup{(\alph*)},leftmargin=4em]
\item a selected \(P,U,A,F,I\) cell enters its established complete
      compiler;
\item a selected response \(r\to q\) closes the original pair by
      \Cref{thm:V13-spectator-fan};
\item a selected response \(r\to s\), \(s\ne r,q\), is one of the
      six labelled avoid-\(q\) cells in
      \Cref{prop:V6-avoid-table};
\item an old self-bank enters the established same-type or V2
      self-role reduction; or
\item the state is nonclean.
\end{enumerate}
\end{corollary}

\begin{proof}
Apply \Cref{thm:V14-target-eligibility}.  Its fresh complete cells
give \textup{(a)}.  Both old external and fresh directed responses
have an actual partner.  The partner \(q\) supplies the reverse
orientation required by
\Cref{thm:V13-spectator-fan}; the other two partners are exactly the
exhaustive V6 table.  Apply
\Cref{prop:V1-self-type-routes,thm:V2-self-role-reduction} to
\eqref{eq:V14-old-self}.
\end{proof}

\begin{remark}[A concrete fixed threshold]
\label{rem:V14-fixed-threshold}
Taking \(\eta=\delta=1/4\) gives fresh complete-cell mass at least
\(\Xi_r/12\), fresh response dominance at least \(a_*/36\), and old
external dominance at least \(a_*/(8M_*)\).  The restriction
resolution adds at most the fixed factor \(2^{J_*}\).  None of these
constants depends on the coordinate support.
\end{remark}

\section{No return to the failed incidence bank}
\label{sec:V14-no-return}

\begin{proposition}[Strict no-return property]
\label{prop:V14-no-return}
The old external alternative
\eqref{eq:V14-old-response} cannot select the same
\(q\)-anchored occurrence \(B_q\) with partner \(q\).  Consequently:
\begin{enumerate}[label=\textup{(\roman*)},leftmargin=4em]
\item a response \(r\to q\) selected after the split comes from a
      different physical role and closes; and
\item every continued avoid-\(q\) response either uses a fresh cell
      or a named old role distinct from the depleted incidence
      subbank \(E\).
\end{enumerate}
\end{proposition}

\begin{proof}
Every coordinate of the \(q\)-anchored occurrence \(B_q\) which is
active at both \(q\) and \(r\) belongs, by definition, to
\(E_{qr}(B_q)=E\).  Therefore
\[
 X_r\bigl((B_q\cap E^c)\cap\{a_{q,e}>0\}\bigr)=0.
\]
The \(E\)-representatives were removed before the old-role union was
formed.  Hence \(B_q\cap E^c\) cannot yield an \(r\to q\) response.
Any later such response comes from another recorded role.  Fresh
cells are disjoint from every old role by
\Cref{lem:V14-target-identity}; the remaining statement follows.
\end{proof}

\begin{assessment}[Progress over the V6 and V13 firewalls]
\label{ass:V14-progress}
The large mass in \(E^c\) is no longer merely a scalar reserve.  It is
resolved inside the original operator word into a finite list of
eligible old roles and disjoint fresh cells.  Failed reversal can no
longer send the proof back to the same \(qr\)-incidence occurrence.
The continuing clean response branches are precisely the six
avoid-\(q\) cells already enumerated in V6, together with the named V2
self-role residual.  Thus V14 removes packet transport as an
independent bottleneck without claiming that these remaining
branches are closed.
\end{assessment}

\section{Regression checks and result ledger}
\label{sec:V14-results}

\begin{proposition}[V14 regression suite]
\label{prop:V14-regressions}
\begin{enumerate}[label=\textup{(\roman*)},leftmargin=4em]
\item No covariance is replaced by a product of two restricted
      covariances.
\item Restriction creates algebraic summands, not additional physical
      occurrences.
\item Every paid summand contains exactly one numerator/resolvent
      payer.
\item Coordinates in \(E\) are removed before old-role and fresh-cell
      selection and are never paid or marked a second time.
\item Every selected directed edge records its actual coordinate
      partner.
\item The number of restriction, role, type, and partner choices is
      bounded independently of support.
\end{enumerate}
\end{proposition}

\begin{proof}
Items \textup{(i)--(iii)} are
\Cref{lem:V14-covariance-telescope,
lem:V14-payer-split,
thm:V14-subset-resolved-packet}.  Item \textup{(iv)} is built into
\Cref{lem:V14-target-identity,prop:V14-no-return}.
The V79 actual-partner partition in
\Cref{thm:V14-target-eligibility} proves \textup{(v)}.  The bounds
\(J_*,M_*\), six incidence types, and three possible partners prove
\textup{(vi)}.
\end{proof}

\begin{theorem}[Third-series V14 result ledger]
\label{thm:V14-results}
Relative to V1--V13 and the two frozen archives, V14 proves:
\begin{enumerate}[label=\textup{(V14.\arabic*)},leftmargin=6.3em]
\item the exact two-term restriction formula for every complete
      covariance in the first-breaker packet;
\item a simultaneous support-uniform restriction normal form for the
      whole unpaid or once-paid V100 word;
\item preservation of the unique payer under arbitrary target-edge
      restriction;
\item an exact packet-resolved depleted-target identity;
\item operator eligibility of every old-bank or fresh-cell
      alternative forced by failed reversal; and
\item a strict no-return theorem for the depleted \(qr\)-incidence
      occurrence.
\end{enumerate}
In particular, the eligibility hypothesis in the clean V6 targeted
acquisition theorem is discharged for the V100 first-breaker packet.
\end{theorem}

\begin{proof}
These are
\Cref{lem:V14-covariance-telescope,
thm:V14-subset-resolved-packet,
lem:V14-payer-split,
lem:V14-target-identity,
thm:V14-target-eligibility,
prop:V14-no-return,
cor:V14-V6-unconditional}.
\end{proof}

\begin{programme}[V15: audit, then the six-cell transition graph]
\label{prog:V14-V15}
\begin{enumerate}[leftmargin=3em]
\item Perform the compulsory read-only audit of the newest
      Standard-Model--Einstein and Navier--Stokes notes.
\item For each of the six avoid-\(q\) cells
      \(q\to r\), \(r\to s\) from
      \Cref{prop:V6-avoid-table}, test direct reversal of the newest
      response \(r\to s\).
\item On a reversible cell, enumerate the literal three-edge trees
      obtainable from the breaker and the available orientations,
      retaining unused responses as V99-bounded spectators.
\item On a nonreversible cell, apply the V14 subset-resolved theorem
      at the new target \(s\) and record a finite state transition
      with its depleted edge.
\item Seek a strictly decreasing finite potential built from
      depleted ordered incidences and used physical roles.  If such a
      potential fails, isolate the first recurrent state rather than
      repeat acquisition circularly.
\item Keep the V2 low-response self-role and nonclean states as
      separate analytic branches.
\end{enumerate}
\end{programme}

\begin{firewall}[Claim boundary after third-series V14]
\label{fire:V14-boundary}
V14 proves packet eligibility of the off-edge mass forced by failed
reversal in a clean V100 first-breaker packet.  It does not close the
six avoid-\(q\) target transitions, the V2 low-response self-role
factor, complete-compiler payer interfaces not already covered by the
archive, or nonclean states.  Higher MTA, WUFC, CPM, cutoff removal,
reflection positivity, Osterwalder--Schrader reconstruction, and a
uniform positive continuum spectral gap remain open.  The full
Yang--Mills mass gap has not been proved.
\end{firewall}

\paragraph{Parent-version record.}
V14 was copied byte-for-byte from
\texttt{Renewal\_Yang\_Mills\_Mass\_Gap\_Research\_Notes\_V13.tex}
before this chapter was appended.  The V13 parent had \(6\,760\)
source lines, \(246\,780\) bytes, and SHA--256
\[
 \texttt{e8863c577953b553a619a853748cff5a82eb656a775c7853080ea57e9a6597f6}.
\]
The next iteration is the compulsory V15 companion audit.

% =============================================================================
% BEGIN APPENDED THIRD-SERIES V15 MATERIAL
% =============================================================================

\clearpage
\chapter{V15: Companion audit and latent all-directed closure}
\label{ch:V15-all-directed}

\section{Compulsory companion audit}
\label{sec:V15-audit}

The newest active Standard-Model--Einstein and Navier--Stokes notes
were inspected read-only before choosing the V15 proof direction.

\begin{record}[Companion snapshots used by V15]
\label{rec:V15-companion-snapshots}
The files read at the checkpoint were:
\begin{enumerate}[leftmargin=3em]
\item
\(\texttt{Renewal\_SM\_Einstein\_Continuum\_Research\_Notes\_V92.tex}\),
with \(36\,478\) source lines, \(1\,343\,707\) bytes, and SHA--256
\[
 \texttt{d33db152a4ea86c90b6fa9cf79d64abc9235ed0e4114dc0f19d8c755d2a6a071};
\]
\item
\(\texttt{Renewal\_NS\_Stability\_Research\_Notes\_V31.tex}\),
with \(23\,052\) source lines, \(887\,537\) bytes, and SHA--256
\[
 \texttt{f1121b62238ad5491bb7cf03635ea9934942edf573ed8faaa9ee79e1d38a6696}.
\]
\end{enumerate}
The SM file advanced from V91 to V92 while this audit was being
assembled; the record above is the later snapshot actually used.
No companion file was modified.
\end{record}

\begin{assessment}[Type-correct transfer from the audit]
\label{ass:V15-audit-transfer}
SM V92 repairs a nonglobal local stencil by identifying its invariant
Bott--Thom class before returning to discretization.  Its relevant
proof-order lesson here is to identify the invariant directed-tree
object before proliferating local transition cases.  SM V91 also
retains the complete differentiated supertrace until its cancellations
and first moment are visible.

NS V31 proves that a normalized annular mark removes multiplicity but
cannot be unmarked without the correlated source first moment.  Its
Yang--Mills proof-order analogue is that three already acquired
directed responses must first be tested as one complete marked packet;
they should not be converted into an unweighted scalar transition
walk and reacquired one edge at a time.

No Bott class, index theorem, heat-formation estimate, or
Navier--Stokes cancellation is imported into Yang--Mills.  Only this
order of exact assembly before local continuation is transferred.
\end{assessment}

\begin{assessment}[Internal nonduplication decision]
\label{ass:V15-nonduplication}
The frozen f2 theorem
\texttt{thm:V99-all-directed-tree} already bounds a complete word
containing three distinct directed edge labels with three distinct
tails.  V15 does not reprove that analytic compiler.  It checks the
six response cells produced by V12--V13 against its hypotheses.  This
check was omitted when V13 treated the V12 response only as a bounded
spectator.
\end{assessment}

\section{The six latent arborescences}
\label{sec:V15-six-arborescences}

Retain the V12--V13 notation
\[
 \mu=\{x,q\}\mid\{y,z\}.
\]
The right matching response is
\[
 S:=y\to z,
 \qquad d_S\ge\gamma_y a_*.
\]
On either directed avoid-\(q\) branch of V12, the additional response
is
\[
 D:=x\to u,
 \qquad u\in\{y,z\},
 \qquad d_D\ge\gamma_x a_*.
\]
Acquisition at the mate \(q\) produces
\[
 Q:=q\to r,
 \qquad r\in\{x,y,z\},
 \qquad d_Q\ge\gamma_q a_*.
\]
At this stage no reversal of \(Q\) is required.

\begin{proposition}[Complete six-cell directed graph table]
\label{prop:V15-six-cell-table}
For every \(u\in\{y,z\}\) and \(r\in\{x,y,z\}\), the three undirected
edge labels
\[
 xu,\qquad qr,\qquad yz
\label{eq:V15-three-labels}
\]
are distinct and form a spanning tree.  Their directed tails are
exactly the three distinct vertices \(x,q,y\); the unique root is
\(z\).  The exhaustive table is:
\begin{center}
\small
\begin{tabular}{c c c c}
\toprule
\(D\) & \(Q\) & undirected tree & orientation toward \(z\)\\
\midrule
\(x\to y\) & \(q\to x\) & \(q-x-y-z\) &
 \(q\to x\to y\to z\)\\
\(x\to y\) & \(q\to y\) & star at \(y\) &
 \(x\to y,\ q\to y,\ y\to z\)\\
\(x\to y\) & \(q\to z\) & \(x-y-z-q\) &
 \(x\to y\to z,\ q\to z\)\\
\(x\to z\) & \(q\to x\) & \(q-x-z-y\) &
 \(q\to x\to z,\ y\to z\)\\
\(x\to z\) & \(q\to y\) & \(q-y-z-x\) &
 \(q\to y\to z,\ x\to z\)\\
\(x\to z\) & \(q\to z\) & star at \(z\) &
 \(x\to z,\ q\to z,\ y\to z\)\\
\bottomrule
\end{tabular}
\end{center}
\end{proposition}

\begin{proof}
The edge \(xu\) contains \(x\) but not \(q\), the edge \(qr\)
contains \(q\), and \(yz\) contains neither \(x\) nor \(q\).
Thus \(xu\ne yz\), \(qr\ne yz\), and \(xu\ne qr\).  The table
directly shows connectedness in all six cases.  Three distinct edges
on four vertices which are connected form a tree.

The sources of the three arrows are \(x,q,y\), independently of
\((u,r)\).  Every displayed directed path ends at \(z\), and the star
rows point directly or indirectly to \(z\).  Thus \(z\) is the unique
vertex with no outgoing selected edge.
\end{proof}

\begin{corollary}[Exact aggregate normalization]
\label{cor:V15-aggregate}
For every row of \Cref{prop:V15-six-cell-table},
\begin{equation}
 d_Dd_Qd_S
 =
 \frac{
 H_{xu}(B_D)H_{qr}(B_Q)H_{yz}(B_S)}
 {\Xi_x\Xi_q\Xi_y}
\label{eq:V15-aggregate}
\end{equation}
is the legal oriented marked-tree response of the displayed
arborescence rooted at \(z\).
\end{corollary}

\begin{proof}
An oriented tree response has one row denominator at every nonroot
tail.  By \Cref{prop:V15-six-cell-table}, those tails are
\(x,q,y\), each exactly once.  Substitution of the definitions of the
three directed responses proves \eqref{eq:V15-aggregate}.
\end{proof}

\begin{firewall}[This is not the V6 two-response table]
\label{fire:V15-not-V6-table}
V6 retained the raw breaker and selected only two of
\(S,Q,D\); its six avoid-\(q\) cells correctly fail the corresponding
two-response tree test.  V15 instead uses all three already assembled
directed covariances and leaves the breaker occurrence in the bounded
word.  The scalar numerator and its denominator degree are therefore
different.  Neither table contradicts the other.
\end{firewall}

\section{Analytic insertion and the occurrence gate}
\label{sec:V15-analytic}

\begin{definition}[Joint three-mark admissibility]
\label{def:V15-joint-admissibility}
The triple \((D,Q,S)\) is called jointly mark-admissible when the exact
packet contains a jointly bounded carrier to which the V92
bank-to-mark map assigns the three edge labels in
\eqref{eq:V15-three-labels}.  Distinct physical covariance
occurrences are admissible by f2 V99.  Coincident coordinate supports
are also allowed when they occupy the distinct occurrences of the
V98 same-bank sandwich.  Two desired marks extracted from one
unrepeated physical occurrence are not declared admissible by this
definition.
\end{definition}

\begin{theorem}[Latent all-directed closure]
\label{thm:V15-latent-closure}
Suppose the three responses above are jointly mark-admissible.  Then
the exact V100 packet, including every V12--V14 spectator and the
matching breaker, obeys the inherited all-directed marked-tree bounds
\begin{align}
 \Gamma_\otimes(1,K)
 &\le
 \frac{C_{\rm ad}}
 {\gamma_x\gamma_q\gamma_y a_*^3}\,
 d_Dd_Qd_S,
\label{eq:V15-all-D-unpaid}\\
 \Gamma_\otimes(w,K^{\rm pay})
 &\le
 \frac{C_{\rm ad}}
 {s_\alpha\gamma_x\gamma_q\gamma_y a_*^3}\,
 d_Dd_Qd_S
\label{eq:V15-all-D-paid}
\end{align}
at every admitted payer location.  The right side is the legal tree
response in \eqref{eq:V15-aggregate}.  In particular, none of the six
cells requires reversal of \(q\to r\).
\end{theorem}

\begin{proof}
\Cref{prop:V15-six-cell-table} supplies three distinct undirected
labels and three distinct tails.  The frozen f2 theorem
\texttt{thm:V99-all-directed-tree} and its extra-occurrence corollary
\texttt{cor:V99-more-directed} therefore give
\eqref{eq:V15-all-D-unpaid} and the paid line whenever the payer lies
in a directed occurrence, finite separator, or complete common
suffix.

If the breaker pays, use its universal paid cap
\(\|{\cal P}_G^\mu\|\le C/s_\alpha\) from
\texttt{thm:V100-breaker-norm}, keep the three directed occurrences
in their jointly bounded unpaid carrier, and apply the one-payer
sandwich.  If another complete bounded spectator pays, use its
already assigned V95--V96 paid line in the same way.  Exactly one
factor \(s_\alpha^{-1}\) occurs.  Finally insert the three lower
bounds
\[
 1\le
 \frac{d_Dd_Qd_S}
 {\gamma_x\gamma_q\gamma_y a_*^3}.
\]
\Cref{cor:V15-aggregate} verifies that the resulting scalar product
is a legal directed-tree mark, not merely a connected graph.
\end{proof}

\begin{corollary}[Correction of the V13 response branch]
\label{cor:V15-correct-V13}
On a V12 directed avoid-\(q\) branch, every jointly
mark-admissible \(q\)-rooted response closes immediately by
\Cref{thm:V15-latent-closure}.  Hence alternatives
\textup{(iii)--(iv)} of
\Cref{thm:V13-two-sided-normal-form} are nonminimal in this branch:
the reversible/nonreversible split should not be entered before the
all-directed test.
\end{corollary}

\begin{proof}
The avoid-\(q\) branch contains \(D=x\to y\) or \(x\to z\).
The right matching response \(S=y\to z\) is already present.
Every \(q\)-rooted response has one of the three partners in
\Cref{prop:V15-six-cell-table}.  Apply
\Cref{thm:V15-latent-closure}.
\end{proof}

\begin{corollary}[Exact residual occurrence collision]
\label{cor:V15-collision-residual}
After \Cref{cor:V15-correct-V13}, a clean \(q\)-response branch can
remain open only if at least two of the desired edge marks
\[
 \{xu,\ qr,\ yz\}
\]
are demanded from one physical occurrence for which no joint
multi-mark carrier has been proved.  This is an occurrence-collision
problem, not a six-cell graph problem and not a scalar
reverse-or-deplete problem.
\end{corollary}

\begin{proof}
If the triple is jointly mark-admissible, it closes by
\Cref{thm:V15-latent-closure}.  Failure of joint admissibility is,
by \Cref{def:V15-joint-admissibility}, precisely the asserted
same-occurrence deficit.  The graph certificate remains valid by
\Cref{prop:V15-six-cell-table}; only the analytic simultaneous-mark
carrier is absent.
\end{proof}

\begin{firewall}[Same support is not the same occurrence]
\label{fire:V15-support-occurrence}
The equality of two coordinate supports is harmless when the exact
word contains two corresponding physical occurrences: f2 V98 bounds
their joint sandwich before expanding the two marks.  The unresolved
case in \Cref{cor:V15-collision-residual} is stronger: two edge roles
are requested from a single unrepeated occurrence.  V15 does not turn
one covariance into two occurrences or apply two independent
capacities to it.
\end{firewall}

\section{Effect on V14 and the active proof boundary}
\label{sec:V15-effect}

\begin{proposition}[Correct scope of the V14 eligibility theorem]
\label{prop:V15-V14-scope}
\Cref{thm:V14-target-eligibility} remains a valid packet theorem for
a nonreversible response when no compatible third directed
occurrence is available.  In the main V12 avoid-\(q\) branch,
however, it is downstream of the all-directed test and is invoked only
after the occurrence collision in
\Cref{cor:V15-collision-residual} has been resolved or shown not to
apply.
\end{proposition}

\begin{proof}
V14 uses only the exact subset covariance telescope and the failed
reverse inequality, so its conclusion is unaffected.  But on the
jointly admissible V12 branch, V15 already supplies a marked tree
without any reversal.  Entering V14 there would add a correct but
unnecessary case split.  The only reason the V15 test can fail is the
analytic occurrence gate stated in
\Cref{cor:V15-collision-residual}.
\end{proof}

\begin{assessment}[Directed acquisition after the correction]
\label{ass:V15-directed-status}
The apparent six-state nonreversible walk was an artefact of testing
only two directed responses at a time.  At the graph and denominator
levels, all six states are one arborescence family rooted at \(z\).
At the operator level, every branch with three admissible occurrences
is closed by an archived theorem.  The sole new directed residual is
therefore a finite same-occurrence marking problem.  The V2
low-response sensitive/insensitive role difference and nonclean
states remain separate; V15 makes no claim about them.
\end{assessment}

\section{Regression checks and result ledger}
\label{sec:V15-results}

\begin{proposition}[V15 regression suite]
\label{prop:V15-regressions}
\begin{enumerate}[label=\textup{(\roman*)},leftmargin=4em]
\item The raw breaker is not used as a fourth directed mark.
\item All three selected edge labels are distinct in each of the six
      cells.
\item The three response tails are distinct and equal to
      \(x,q,y\); no response is reversed.
\item The legal denominator is
      \(\Xi_x\Xi_q\Xi_y\), with root \(z\).
\item An unused covariance occurrence remains in the bounded word
      and is not deleted.
\item Coincident supports are accepted only through an existing joint
      repeated-occurrence sandwich.
\item One unrepeated occurrence is never assigned two independent
      capacities or payer debits.
\end{enumerate}
\end{proposition}

\begin{proof}
Items \textup{(i)--(iv)} are
\Cref{prop:V15-six-cell-table,
cor:V15-aggregate,
fire:V15-not-V6-table}.  The V99 extra-occurrence compiler in
\Cref{thm:V15-latent-closure} proves \textup{(v)}.
\Cref{def:V15-joint-admissibility,
fire:V15-support-occurrence} prove \textup{(vi)--(vii)}.
\end{proof}

\begin{theorem}[Third-series V15 result ledger]
\label{thm:V15-results}
Relative to V1--V14 and the frozen V99 interface, V15 proves:
\begin{enumerate}[label=\textup{(V15.\arabic*)},leftmargin=6.3em]
\item the compulsory SM/NS companion audit and its proof-order
      consequence;
\item an exhaustive six-cell table of arborescences rooted at \(z\);
\item exact legal normalization of all six three-directed response
      products;
\item unpaid and once-paid closure of every jointly mark-admissible
      cell by the existing all-directed compiler;
\item correction of the nonminimal V13 reverse-or-deplete branch; and
\item reduction of the surviving directed obstruction to a finite
      same-occurrence multi-mark carrier.
\end{enumerate}
\end{theorem}

\begin{proof}
These are
\Cref{rec:V15-companion-snapshots,
ass:V15-audit-transfer,
prop:V15-six-cell-table,
cor:V15-aggregate,
thm:V15-latent-closure,
cor:V15-correct-V13,
cor:V15-collision-residual}.
\end{proof}

\begin{programme}[V16: exact occurrence-collision partition]
\label{prog:V15-V16}
\begin{enumerate}[leftmargin=3em]
\item For the desired roles
      \(\{D=x\to u,Q=q\to r,S=y\to z\}\), partition the physical
      occurrence map by the five set partitions
      \[
      D|Q|S,\quad DQ|S,\quad DS|Q,\quad QS|D,\quad DQS.
      \]
\item Retain \(D|Q|S\) as the V15 closed case.  In each two-role
      collision, write the exact complete covariance numerator before
      applying any scalar edge marks.
\item Compare the resulting numerator with the valid V83
      higher-incidence square tail and the V84--V85 counterexamples to
      the false crossing-only square tail.  Do not infer a two-mark
      bound from coordinate incidence alone.
\item For the triple collision \(DQS\), descend to the first local
      partition change inside the one occurrence and test whether its
      complete positive envelope supplies one higher-incidence card or
      only one edge stock.
\item Preserve the unique payer inside that occurrence throughout the
      local expansion.
\item After the collision carrier is settled, return to the V2
      low-response sensitive/insensitive factorization as the next
      independent clean branch.
\end{enumerate}
\end{programme}

\begin{firewall}[Claim boundary after third-series V15]
\label{fire:V15-boundary}
V15 closes all six latent directed arborescences when their three
marks have a proved joint carrier.  It does not prove the missing
multi-mark estimate for one unrepeated occurrence, close the V2
low-response self-role factor, complete every remaining payer
interface, or handle nonclean states.  Higher MTA, WUFC, CPM, cutoff
removal, reflection positivity, Osterwalder--Schrader reconstruction,
and a uniform positive continuum spectral gap remain open.  The full
Yang--Mills mass gap has not been proved.
\end{firewall}

\paragraph{Parent-version record.}
V15 was copied byte-for-byte from
\texttt{Renewal\_Yang\_Mills\_Mass\_Gap\_Research\_Notes\_V14.tex}
before this chapter was appended.  The V14 parent had \(7\,305\)
validation records, \(267\,477\) bytes, and SHA--256
\[
 \texttt{1699af0e8aa806c9f052d89b1cee4b1d7801471ba3df0821ecbfe650b4f165d9}.
\]
The next compulsory companion audit is V20.

% =============================================================================
% BEGIN APPENDED THIRD-SERIES V16 MATERIAL
% =============================================================================

\clearpage
\chapter{V16: The one-difference invariant and the occurrence-collision
lattice}
\label{ch:V16-collision}

\section{Context and exact question}
\label{sec:V16-context}

V15 proves the graph and denominator certificate for all six triples
\[
 D=x\to u,\qquad Q=q\to r,\qquad S=y\to z,
\]
where \(u\in\{y,z\}\) and \(r\in\{x,y,z\}\).  If the three desired
edge roles lie in a proved joint three-mark carrier, the archived
all-directed theorem closes the packet.

The remaining issue is not equality of coordinate supports.  V98
already permits equal supports in distinct jointly bounded
occurrences.  The issue is that two or three roles may be extracted
from one unrepeated complete covariance occurrence.  V16 proves that
arbitrary coordinate restriction cannot clone such an occurrence.
It then gives the exhaustive five-state occurrence partition and
identifies the exact higher-incidence locus on which a genuinely new
multi-mark estimate would have to act.

\section{The finite telescope preserves one difference}
\label{sec:V16-one-difference}

Let
\[
 B=B_1\dot\cup\cdots\dot\cup B_m
\]
be an ordered finite partition of one physical bank.  Retain the
notation
\[
 M_A(\rho)=\bigotimes_{e\in A}M_e(\rho),
 \qquad
 \Delta_A^{\rho,\sigma}=M_A(\rho)-M_A(\sigma).
\]

\begin{theorem}[Finite one-difference telescope]
\label{thm:V16-one-difference}
For every \(m\ge1\),
\begin{equation}
 \boxed{
 \Delta_B^{\rho,\sigma}
 =
 \sum_{k=1}^m
 \left(\bigotimes_{j<k}M_{B_j}(\rho)\right)
 \otimes\Delta_{B_k}^{\rho,\sigma}
 \otimes
 \left(\bigotimes_{j>k}M_{B_j}(\sigma)\right).}
\label{eq:V16-one-difference}
\end{equation}
Every summand contains exactly one partition-state difference.
All other subbanks occur as endpoint moments.
\end{theorem}

\begin{proof}
For \(0\le k\le m\), define
\[
 T_k:=
 \left(\bigotimes_{j\le k}M_{B_j}(\rho)\right)
 \otimes
 \left(\bigotimes_{j>k}M_{B_j}(\sigma)\right).
\]
Then \(T_m=M_B(\rho)\), \(T_0=M_B(\sigma)\), and
\[
 T_k-T_{k-1}
 =
 \left(\bigotimes_{j<k}M_{B_j}(\rho)\right)
 \otimes
 \bigl(M_{B_k}(\rho)-M_{B_k}(\sigma)\bigr)
 \otimes
 \left(\bigotimes_{j>k}M_{B_j}(\sigma)\right).
\]
Sum \(T_k-T_{k-1}\) over \(k\).  The left side telescopes to
\(\Delta_B^{\rho,\sigma}\), proving
\eqref{eq:V16-one-difference}.
\end{proof}

\begin{corollary}[No covariance cloning by restriction]
\label{cor:V16-no-cloning}
No sequence of coordinate partitions and exact moment
factorizations turns one occurrence
\(\Delta_B^{\rho,\sigma}\) into a product containing two nontrivial
covariance differences.  Every fully distributed exact term contains
one difference descended from the original occurrence.
\end{corollary}

\begin{proof}
\Cref{thm:V16-one-difference} proves the claim for one finite
partition.  Refining any \(B_k\) applies the same theorem to its sole
difference and leaves all other factors as moments.  Induction over a
finite refinement tree preserves exactly one difference in every
nonzero leaf term.
\end{proof}

\begin{lemma}[One payer under the finite telescope]
\label{lem:V16-one-payer}
If \(\Delta_B^{\rho,\sigma}\) is the occurrence containing the unique
payer, the simultaneous expansion by payer coordinate and by
\eqref{eq:V16-one-difference} has exactly one payer in each nonzero
term.  If another occurrence pays, every term in
\eqref{eq:V16-one-difference} is unpaid.
\end{lemma}

\begin{proof}
Partition the payer-coordinate sum by the disjoint sets
\(B_1,\ldots,B_m\) before estimating it.  In the term indexed by its
unique cell \(B_k\), place the paid local numerator/resolvent in the
single \(k\)-factor and retain endpoint moments on all other cells.
If the payer is external to \(B\), tensoring with
\eqref{eq:V16-one-difference} does not change its location.
\end{proof}

\begin{firewall}[A sum of locations is not a product of marks]
\label{fire:V16-sum-not-product}
Equation \eqref{eq:V16-one-difference} permits the one covariance to
be located on any one subbank in alternative summands.  It does not
permit two of those locations to occur simultaneously in one
summand.  Thus disjoint support for two large response stocks does not
by itself produce their product as two physical covariance marks.
\end{firewall}

\section{The five occurrence partitions}
\label{sec:V16-five-partitions}

Let \({\cal O}\) be the set of physical covariance occurrences in
the exact packet and let
\[
 \pi:\{D,Q,S\}\longrightarrow{\cal O}
\]
send each desired directed edge role to the occurrence from which its
response was acquired.  Equal coordinate support in two distinct
occurrences does not make their \(\pi\)-values equal.

\begin{proposition}[Exhaustive occurrence-collision lattice]
\label{prop:V16-five-partitions}
Exactly one of the following five fibres occurs:
\begin{equation}
 D|Q|S,\qquad
 DQ|S,\qquad
 DS|Q,\qquad
 QS|D,\qquad
 DQS.
\label{eq:V16-five-partitions}
\end{equation}
Here a vertical bar separates distinct physical occurrences and
concatenated letters share one occurrence.  The first state is closed
by V15.  The next three contain one two-role collision.  The last
contains one three-role collision.
\end{proposition}

\begin{proof}
The fibres of a map on a three-element set form a set partition.
There is one partition into three singletons, three choices of a
two-element block and singleton, and one three-element block.  These
are the five displayed states.  The first is jointly admissible by
f2 V99, including equal supports through distinct occurrences, and
hence closes by \Cref{thm:V15-latent-closure}.
\end{proof}

\begin{remark}[The lattice is about occurrences, not banks]
\label{rem:V16-occurrence-not-bank}
Two occurrences may use the same coordinate bank and still belong to
the state \(D|Q|S\); V98 supplies the required joint repeated-bank
sandwich.  Conversely, one occurrence may contain coordinates in
several subbanks, but
\Cref{cor:V16-no-cloning} keeps it in a collision block of
\eqref{eq:V16-five-partitions}.
\end{remark}

\section{Exclusive and genuinely higher-incidence collision support}
\label{sec:V16-higher-incidence}

Fix a two-role collision with distinct undirected edge labels
\[
 e_1=a_1b_1,\qquad e_2=a_2b_2.
\]
Inside its common occurrence \(C\), define the literal incidence
supports
\[
 E_i:=E_{a_ib_i}(C),\qquad
 O_{12}:=E_1\cap E_2.
\]
Then
\[
 C=(E_1\setminus O_{12})
 \dot\cup(E_2\setminus O_{12})
 \dot\cup O_{12}
 \dot\cup\bigl(C\setminus(E_1\cup E_2)\bigr).
\label{eq:V16-four-support-cells}
\]

\begin{lemma}[Incidence order of the overlap cell]
\label{lem:V16-overlap-incidence}
Every coordinate in \(O_{12}\) is active on all vertices of
\(e_1\cup e_2\).  Since the labels are distinct, this means:
\begin{enumerate}[label=\textup{(\roman*)},leftmargin=4em]
\item incidence at least three when \(e_1,e_2\) share one vertex;
\item incidence four when \(e_1,e_2\) are disjoint.
\end{enumerate}
\end{lemma}

\begin{proof}
Membership in \(E_i\) means simultaneous activity at both endpoints
of \(e_i\).  Membership in their intersection therefore activates
the union of the two endpoint sets.  Two distinct edges have a
three-vertex union if they meet and a four-vertex union if they are
disjoint.
\end{proof}

\begin{proposition}[Exact collision-support normal form]
\label{prop:V16-collision-support}
Apply \Cref{thm:V16-one-difference} to the four cells in
\eqref{eq:V16-four-support-cells}.  Every resulting operator summand
is of exactly one of the following forms:
\begin{enumerate}[label=\textup{(\alph*)},leftmargin=4em]
\item the sole covariance difference lies on the
      \(e_1\)-exclusive cell;
\item it lies on the \(e_2\)-exclusive cell;
\item it lies on the higher-incidence overlap \(O_{12}\); or
\item it lies on the remainder, while all three edge-support cells
      are endpoint moments.
\end{enumerate}
Only \textup{(c)} is a candidate for a genuine one-occurrence
two-edge numerator.  Alternatives \textup{(a)}, \textup{(b)}, and
\textup{(d)} cannot be converted into two covariance occurrences by
coordinate factorization.
\end{proposition}

\begin{proof}
The four alternatives are the four terms of
\eqref{eq:V16-one-difference} for the partition
\eqref{eq:V16-four-support-cells}, omitting any empty cell.
\Cref{lem:V16-overlap-incidence} gives the higher-incidence
description of \textup{(c)}.  In the other three terms at most one of
the two desired edge-support cells carries the sole covariance
difference.  The other cells are endpoint moments, so
\Cref{cor:V16-no-cloning} forbids the asserted conversion.
\end{proof}

\begin{corollary}[Triple collisions have one genuine common locus]
\label{cor:V16-triple-locus}
In the state \(DQS\), let \(E_D,E_Q,E_S\) be the three literal edge
supports in the common occurrence.  Refinement by their Boolean
atoms leaves exactly one possible simultaneous three-edge locus,
\[
 O_{DQS}:=E_D\cap E_Q\cap E_S.
\]
Every coordinate of \(O_{DQS}\) is active on all four rows
\(\{x,q,y,z\}\).  All other atoms omit at least one desired edge, and
the exact telescope still places only one covariance difference in
each summand.
\end{corollary}

\begin{proof}
By \Cref{prop:V15-six-cell-table}, the union of the three edge labels
is a spanning tree and hence contains all four vertices.  A coordinate
in all three incidence supports is active at every endpoint in that
union, so it has incidence four.  Every other Boolean atom omits at
least one support by definition.  Apply
\Cref{thm:V16-one-difference} to the Boolean-atom partition.
\end{proof}

\section{Safe fallback routes for two collision classes}
\label{sec:V16-fallback}

The one-difference theorem says where the all-directed proof must
stop, but it does not erase the older one-response route.  Two
collision classes admit that route without double marking.

\begin{proposition}[Two collision classes retain the V13--V14
fallback]
\label{prop:V16-fallback}
In either occurrence state
\[
 DQ|S\qquad\text{or}\qquad DS|Q,
\]
retain \(D\) only as a bounded spectator and use at most one response
mark from each collision occurrence.  Then:
\begin{enumerate}[label=\textup{(\roman*)},leftmargin=4em]
\item if the actual reverse \(r\to q\) of \(Q=q\to r\) is dominant,
      the response \(S=y\to z\), the reverse \(r\to q\), and the raw
      breaker close by \Cref{thm:V13-spectator-fan};
\item if that reverse fails, the packet-resolved target theorem
      \Cref{thm:V14-target-eligibility} applies at \(r\).
\end{enumerate}
Neither alternative requests two scalar marks from the same physical
occurrence.
\end{proposition}

\begin{proof}
In \(DQ|S\), the \(S\)-mark is in the singleton occurrence and the
reverse \(Q\)-mark is the sole selected mark from the collision
occurrence; \(D\) remains unmarked there.  In \(DS|Q\), the
\(S\)-mark is the sole selected mark from its collision occurrence
and the reverse \(Q\)-mark lies in the singleton occurrence.  Thus
the V13 three-edge fan is jointly admissible in either state.

If the reverse fails, V14 uses the forward \(Q\)-response only to
identify its depleted literal incidence set and then resolves the
whole packet before new marks are selected.  The collided \(D\)-role
remains inside its original bounded occurrence.  Hence the
one-difference and one-payer ledgers are preserved.
\end{proof}

\begin{firewall}[The \(QS|D\) and \(DQS\) states are different]
\label{fire:V16-hard-collisions}
The V13 fan uses both \(S\) and the reversed \(Q\).  In the states
\(QS|D\) and \(DQS\), those two desired marks occupy the same
unrepeated occurrence.  Therefore
\Cref{prop:V16-fallback} does not apply to them without an additional
joint higher-incidence estimate.  Reversal of a scalar response does
not split its physical covariance occurrence.
\end{firewall}

\section{What the archived square tail can legitimately test}
\label{sec:V16-V83}

\begin{proposition}[V83 applicability gate]
\label{prop:V16-V83-gate}
Let \(O_{12}\) be the candidate collision locus from
\Cref{prop:V16-collision-support}.  The archived V83
higher-incidence square tail may be invoked only after verifying all
of the following data for one selected row \(k\):
\begin{enumerate}[label=\textup{(\roman*)},leftmargin=4em]
\item \(O_{12}\) is a clean balanced physical bank crossing the
      selected row cut \(k|k^c\);
\item it is retained as the one covariance difference in the exact
      telescope;
\item its row-\(k\) mass is a fixed fraction of \(\Xi_k\);
\item a maximal actual partner \(\ell\) is retained; and
\item the V83 graph and mass-comparison gate converts its
      \(H_{k\ell}/\Xi_k^2\) numerator to the desired tree
      normalization.
\end{enumerate}
Simultaneous activity on the two collision edges proves none of
\textup{(iii)--(v)} by itself.
\end{proposition}

\begin{proof}
Items \textup{(i)--(iv)} are precisely the hypotheses under which the
frozen results
\texttt{thm:V83-selected-bank} and
\texttt{cor:V83-dominant-square} are stated.  The final conversion is
the separate theorem
\texttt{thm:V83-tree-exchange}.  The V84--V85 crossing-square
counterexample shows why an arbitrary cross-stock square cannot
replace these conditions.  Thus incidence order alone is
insufficient.
\end{proof}

\begin{assessment}[Exact boundary after occurrence resolution]
\label{ass:V16-boundary}
The occurrence collision is now a finite, noncircular object.
Disjoint large response supports inside one covariance do not solve
it: the exact telescope assigns the difference to only one support at
a time.  A genuinely new two-mark estimate can live only on a
three- or four-row overlap cell and must pass the V83 row-rooted
square-tail gate.  The \(DQ|S\) and \(DS|Q\) classes retain the safe
V13--V14 one-response fallback.  The hard classes are \(QS|D\) and
\(DQS\), together with any overlap cell which fails V83 dominance or
mass compatibility.
\end{assessment}

\section{Regression checks and result ledger}
\label{sec:V16-results}

\begin{proposition}[V16 regression suite]
\label{prop:V16-regressions}
\begin{enumerate}[label=\textup{(\roman*)},leftmargin=4em]
\item A coordinate restriction of one covariance always has one
      difference per exact summand.
\item The unique payer stays unique under every finite restriction.
\item Equal supports in distinct occurrences are not mislabeled as a
      collision.
\item Exclusive edge supports are not multiplied as independent
      covariance marks.
\item A two-edge common locus is labelled incidence-three or
      incidence-four according to its actual endpoint union.
\item A triple common locus is incidence-four.
\item V83 is used only with its row-cut, dominance, partner, and mass
      hypotheses.
\end{enumerate}
\end{proposition}

\begin{proof}
Items \textup{(i)--(ii)} are
\Cref{thm:V16-one-difference,lem:V16-one-payer}.
\Cref{rem:V16-occurrence-not-bank} proves \textup{(iii)}.
\Cref{prop:V16-collision-support} proves \textup{(iv)}.
\Cref{lem:V16-overlap-incidence,
cor:V16-triple-locus} prove \textup{(v)--(vi)}.
\Cref{prop:V16-V83-gate} proves \textup{(vii)}.
\end{proof}

\begin{theorem}[Third-series V16 result ledger]
\label{thm:V16-results}
Relative to V1--V15, V16 proves:
\begin{enumerate}[label=\textup{(V16.\arabic*)},leftmargin=6.3em]
\item the exact finite telescope for arbitrary restrictions of one
      complete covariance occurrence;
\item invariance of one difference and one payer in every resolved
      summand;
\item the exhaustive five-state occurrence-collision lattice;
\item separation of exclusive response support from the genuine
      incidence-three/four collision locus;
\item identification of the unique incidence-four triple-collision
      locus;
\item safe V13--V14 fallback for the \(DQ|S\) and \(DS|Q\) classes;
      and
\item the exact hypothesis gate for any use of the V83
      higher-incidence square tail.
\end{enumerate}
\end{theorem}

\begin{proof}
These are
\Cref{thm:V16-one-difference,
cor:V16-no-cloning,
lem:V16-one-payer,
prop:V16-five-partitions,
prop:V16-collision-support,
cor:V16-triple-locus,
prop:V16-fallback,
prop:V16-V83-gate}.
\end{proof}

\begin{programme}[V17: numerator test on the hard collision loci]
\label{prog:V16-V17}
\begin{enumerate}[leftmargin=3em]
\item Treat \(QS|D\) first.  Resolve the common \(Q,S\) occurrence
      into the exclusive cells and \(O_{QS}\) by
      \eqref{eq:V16-one-difference}, retaining the payer.
\item On \(O_{QS}\), partition by exact local row partition and a
      V70 maximal partner at each candidate root.  Test the V83
      dominance and mass-comparison hypotheses rather than assuming
      them.
\item If the V83 gate holds, insert the resulting square-tail
      numerator into the raw breaker plus \(D\) packet and check the
      literal tree denominator.
\item If the gate fails, write the exact depleted identity for the
      missing row mass on the exclusive and noncollision atoms; do
      not return to the same common occurrence.
\item Repeat only after \(QS|D\) is resolved for the all-in-one state
      \(DQS\), whose common locus is incidence-four.
\item Test every proposed cross-stock square against the frozen
      V84--V85 counterexample before accepting it.
\item Keep the V2 low-response role difference and nonclean states
      outside this collision calculation.
\end{enumerate}
\end{programme}

\begin{firewall}[Claim boundary after third-series V16]
\label{fire:V16-claim-boundary}
V16 proves that packet restriction cannot clone a covariance and
localizes every multi-mark occurrence collision to exact
higher-incidence atoms.  It does not prove the required multi-mark
numerator on those atoms, close \(QS|D\) or \(DQS\), close every
nonreversible fallback transition, resolve the V2 low-response
self-role factor, or handle nonclean states.  Higher MTA, WUFC, CPM,
cutoff removal, reflection positivity, Osterwalder--Schrader
reconstruction, and a uniform positive continuum spectral gap remain
open.  The full Yang--Mills mass gap has not been proved.
\end{firewall}

\paragraph{Parent-version record.}
V16 was copied byte-for-byte from
\texttt{Renewal\_Yang\_Mills\_Mass\_Gap\_Research\_Notes\_V15.tex}
before this chapter was appended.  The V15 parent had \(7\,743\)
validation records, \(284\,098\) bytes, and SHA--256
\[
 \texttt{5a7d398e9d3fc0c2c62e7cf27bfde1d21527da8ba3dab82b0b6518db16a571d6}.
\]
The next compulsory companion audit remains V20.

% =============================================================================
% BEGIN APPENDED THIRD-SERIES V17 MATERIAL
% =============================================================================

\clearpage
\chapter{V17: Dominant collision roles need not overlap}
\label{ch:V17-no-overlap}

\section{Context and the upstream test}
\label{sec:V17-context}

V16 localizes a two-role occurrence collision to the overlap of the
two literal edge-support sets, but does not show that this overlap is
large.  The hardest two-role class is \(QS|D\): the \(q\)-response
\[
 Q=q\to r
\]
and the right matching response
\[
 S=y\to z
\]
come from one unrepeated covariance \(C\), while
\(D=x\to u\) lies in another occurrence.

It is tempting to argue that dominance of both \(Q\) and \(S\)
forces a dominant incidence-three or incidence-four subbank.  V17
tests and refutes that inference.  It then gives the exact
overlap-or-exclusive mass decomposition and identifies the new
operator needed on the exclusive branch: a mixed
covariance--endpoint-moment compiler, not a second covariance.

\section{Quantitative overlap-or-exclusive decomposition}
\label{sec:V17-decomposition}

Inside the collision occurrence \(C\), put
\[
 E_Q:=E_{qr}(C),\qquad
 E_S:=E_{yz}(C),\qquad
 O:=E_Q\cap E_S,
\]
and
\[
 C_Q:=E_Q\setminus O,\qquad
 C_S:=E_S\setminus O,\qquad
 R:=C\setminus(E_Q\cup E_S).
\]
Thus
\begin{equation}
 C=C_Q\dot\cup O\dot\cup C_S\dot\cup R.
\label{eq:V17-four-cells}
\end{equation}

The V79 acquisition producing a directed response retains a row-mass
certificate before using the Casimir floor.  Write those certificates
as
\begin{equation}
 X_q(E_Q)\ge\kappa_Q\Xi_q,
 \qquad
 X_y(E_S)\ge\kappa_S\Xi_y,
 \qquad
 \kappa_Q,\kappa_S>0.
\label{eq:V17-two-mass-certificates}
\end{equation}
They imply the displayed response lower bounds, but are stronger data
than the scalar \(d_Q,d_S\) alone.

\begin{lemma}[Overlap-or-exclusive alternative]
\label{lem:V17-overlap-exclusive}
Fix thresholds
\[
 0<\theta_Q<\kappa_Q,\qquad
 0<\theta_S<\kappa_S.
\]
Then either
\begin{equation}
 X_q(O)\ge\theta_Q\Xi_q
 \quad\hbox{and}\quad
 X_y(O)\ge\theta_S\Xi_y,
\label{eq:V17-bidominant-overlap}
\end{equation}
or at least one of the following exclusive estimates holds:
\begin{align}
 X_q(C_Q)&\ge(\kappa_Q-\theta_Q)\Xi_q,
\label{eq:V17-Q-exclusive}\\
 X_y(C_S)&\ge(\kappa_S-\theta_S)\Xi_y.
\label{eq:V17-S-exclusive}
\end{align}
\end{lemma}

\begin{proof}
If the first inequality in
\eqref{eq:V17-bidominant-overlap} fails, then
\[
 X_q(C_Q)
 =X_q(E_Q)-X_q(O)
 >(\kappa_Q-\theta_Q)\Xi_q.
\]
If the second fails, the same subtraction gives
\eqref{eq:V17-S-exclusive}.  If neither fails, both inequalities in
\eqref{eq:V17-bidominant-overlap} hold.
\end{proof}

\begin{corollary}[The overlap branch is genuinely higher-incidence]
\label{cor:V17-overlap-higher}
On \eqref{eq:V17-bidominant-overlap}, \(O\) is a fixed-fraction bank
at both directed tails.  It has incidence four when \(r=x\), and
incidence at least three when \(r=y\) or \(r=z\).
\end{corollary}

\begin{proof}
The mass statement is
\eqref{eq:V17-bidominant-overlap}.  The edge labels \(qx\) and \(yz\)
are disjoint, whereas \(qy\) meets \(yz\) at \(y\) and \(qz\) meets
\(yz\) at \(z\).  Apply
\Cref{lem:V16-overlap-incidence}.
\end{proof}

\section{A sharp disjoint-support counterexample}
\label{sec:V17-counterexample}

\begin{proposition}[Two dominant responses do not force overlap]
\label{prop:V17-empty-overlap}
For every \(r\in\{x,y,z\}\), there are clean row-incidence data on a
two-coordinate bank \(C=\{e_Q,e_S\}\) satisfying the Casimir floor,
\[
 d_{q\to r}(C)\ge\frac12a_*,
 \qquad
 d_{y\to z}(C)\ge\frac12a_*,
\label{eq:V17-two-dominant}
\]
but
\[
 E_{qr}(C)\cap E_{yz}(C)=\varnothing.
\label{eq:V17-empty-overlap}
\]
Thus neither directed dominance nor the two retained row-mass
certificates forces a nonempty higher-incidence collision cell.
\end{proposition}

\begin{proof}
Give every active row-coordinate incidence weight \(a_*\), and set
all unlisted weights to zero.

For \(r=x\), activate exactly \(q,x\) at \(e_Q\) and exactly \(y,z\)
at \(e_S\).  Then all four row masses equal \(a_*\), both edge stocks
equal \(a_*^2\), and both responses equal \(a_*\).

For \(r=y\), activate exactly \(q,y\) at \(e_Q\) and exactly \(y,z\)
at \(e_S\).  Then
\[
 \Xi_q=\Xi_z=a_*,
 \qquad
 \Xi_y=2a_*,
 \qquad
 H_{qy}=H_{yz}=a_*^2.
\]
Hence \(d_{q\to y}=a_*\) and
\(d_{y\to z}=a_*/2\).

For \(r=z\), activate exactly \(q,z\) at \(e_Q\) and exactly \(y,z\)
at \(e_S\).  Then
\[
 \Xi_q=\Xi_y=a_*,
 \qquad
 \Xi_z=2a_*,
\]
and both responses with tails \(q,y\) equal \(a_*\).
In all three constructions, \(e_Q\) is not simultaneously active on
\(y,z\), while \(e_S\) is not simultaneously active on \(q,r\).
Therefore the two incidence supports are disjoint.  Every nonzero
weight equals the allowed floor, and each supporting row cell carries
a fixed fraction of its tail mass.
\end{proof}

\begin{firewall}[Scope of the counterexample]
\label{fire:V17-counterexample-scope}
\Cref{prop:V17-empty-overlap} is a clean combinatorial input-bank
test.  It does not assert that every listed incidence assignment
survives every global partition coefficient in the Yang--Mills
source.  It rigorously refutes only an inference based on the row
stocks and Casimir floor: those data alone cannot force a common
higher-incidence coordinate.
\end{firewall}

\begin{corollary}[V83 cannot exhaust the \(QS\) collision]
\label{cor:V17-V83-not-exhaustive}
The V83 higher-incidence square tail cannot be the sole resolution of
the \(QS|D\) class.  Its selected bank \(O\) may be empty while both
directed response thresholds remain uniformly positive.
\end{corollary}

\begin{proof}
Use \Cref{prop:V17-empty-overlap}.  The V83 selected-bank hypotheses
require a nonempty dominant higher-incidence bank.  They are absent
in that example.
\end{proof}

\section{Exact operator normal form on the four support cells}
\label{sec:V17-operator}

Order the four cells in
\eqref{eq:V17-four-cells} as \(C_Q,O,C_S,R\).  The one-difference
telescope gives
\begin{equation}
\begin{split}
 \Delta_C^{\rho,\sigma}
 ={}&
 \Delta_{C_Q}^{\rho,\sigma}
 \otimes M_O(\sigma)\otimes M_{C_S}(\sigma)\otimes M_R(\sigma)\\
 &+
 M_{C_Q}(\rho)\otimes
 \Delta_O^{\rho,\sigma}
 \otimes M_{C_S}(\sigma)\otimes M_R(\sigma)\\
 &+
 M_{C_Q}(\rho)\otimes M_O(\rho)\otimes
 \Delta_{C_S}^{\rho,\sigma}\otimes M_R(\sigma)\\
 &+
 M_{C_Q}(\rho)\otimes M_O(\rho)\otimes M_{C_S}(\rho)
 \otimes\Delta_R^{\rho,\sigma}.
\end{split}
\label{eq:V17-QS-telescope}
\end{equation}

\begin{theorem}[Hard-collision four-term normal form]
\label{thm:V17-four-term-normal-form}
Every unpaid or once-paid \(QS|D\) packet is an exact sum of the four
operator types in \eqref{eq:V17-QS-telescope}, tensored into the
unchanged surrounding word.  They have the following exhaustive
fates:
\begin{enumerate}[label=\textup{(\roman*)},leftmargin=4em]
\item the \(Q\)-exclusive covariance term contains one \(Q\)-edge
      difference and only endpoint moments on the \(S\)-support;
\item the overlap term contains the sole candidate
      higher-incidence \(QS\) difference;
\item the \(S\)-exclusive term contains one \(S\)-edge difference and
      only endpoint moments on the \(Q\)-support;
\item the remainder term contains neither desired edge in its
      covariance difference; both edge supports occur only as
      endpoint moments.
\end{enumerate}
Every term contains exactly one difference and at most one payer.
\end{theorem}

\begin{proof}
Equation \eqref{eq:V17-QS-telescope} is
\Cref{thm:V16-one-difference} for the ordered partition
\eqref{eq:V17-four-cells}.  Tensor it into the exact packet before
norms.  \Cref{lem:V16-one-payer} allocates the unique payer.  The
incidence descriptions follow from the definitions of
\(C_Q,O,C_S,R\).
\end{proof}

\begin{corollary}[Exclusive dominance produces a mixed operator, not
a second covariance]
\label{cor:V17-mixed-operator}
On \eqref{eq:V17-Q-exclusive}, the corresponding exact term has the
schematic form
\[
 \Delta_{C_Q}^{\rho,\sigma}\otimes
 M_{C_S}^{\rm end}\otimes(\text{complete spectators});
\]
on \eqref{eq:V17-S-exclusive}, the symmetric term has
\[
 M_{C_Q}^{\rm end}\otimes
 \Delta_{C_S}^{\rho,\sigma}\otimes(\text{complete spectators}).
\]
The large response stock on the endpoint-moment factor is not a
second covariance occurrence.
\end{corollary}

\begin{proof}
Read the first and third lines of
\eqref{eq:V17-QS-telescope}.  The exclusive mass estimates locate the
large row stock, while
\Cref{cor:V16-no-cloning} determines the operator type.
\end{proof}

\begin{firewall}[The directed-word theorem does not accept an endpoint
moment by relabeling]
\label{fire:V17-no-relabel-moment}
The f2 all-directed compiler assumes three directed covariance
occurrences or a separately proved joint carrier.  In
\Cref{cor:V17-mixed-operator}, one desired role is only an endpoint
moment of the same telescope.  Calling it a second directed
covariance would violate \Cref{cor:V16-no-cloning}.  A new mixed
covariance--moment estimate, or an existing endpoint-energy theorem
whose hypotheses match this exact word, must be inserted explicitly.
\end{firewall}

\section{The exact mixed endpoint gate}
\label{sec:V17-mixed-gate}

\begin{definition}[Collision endpoint-promotion property]
\label{def:V17-endpoint-promotion}
For an exclusive term in
\eqref{eq:V17-QS-telescope}, say that the endpoint role is promoted
when a support-uniform unpaid and once-paid sandwich theorem for that
exact term supplies the endpoint edge stock with its required row
denominator, while:
\begin{enumerate}[label=\textup{(\alph*)},leftmargin=4em]
\item retaining the sole covariance difference on the other cell;
\item keeping the endpoint factor as a moment, not a relabelled
      covariance;
\item preserving one payer; and
\item placing the resulting three edge marks into the V15
      arborescence.
\end{enumerate}
\end{definition}

\begin{proposition}[Finite closure criterion for \(QS|D\)]
\label{prop:V17-QS-criterion}
The \(QS|D\) occurrence class closes if:
\begin{enumerate}[label=\textup{(\roman*)},leftmargin=4em]
\item the overlap term satisfies the V83 applicability gate of
      \Cref{prop:V16-V83-gate} and its resulting numerator matches the
      V15 tree; and
\item both exclusive terms and the remainder term satisfy the
      collision endpoint-promotion property of
      \Cref{def:V17-endpoint-promotion}.
\end{enumerate}
Conversely, the combination of directed dominance, the Casimir floor,
and coordinate factorization alone does not verify these conditions.
\end{proposition}

\begin{proof}
\Cref{thm:V17-four-term-normal-form} is exhaustive.  Under
\textup{(i)}, the overlap line is bounded by a legal tree response.
Under \textup{(ii)}, each of the other three lines is bounded by the
same finite V15 tree atlas with one payer.  Sum the four positive
envelopes; the routing factor is four.

For the converse statement, \Cref{prop:V17-empty-overlap} removes
the V83 line while preserving directed dominance, and
\Cref{cor:V16-no-cloning} shows that factorization alone does not
promote either endpoint moment.  Thus additional operator input is
necessary.
\end{proof}

\begin{assessment}[Why V17 changes the next attack]
\label{ass:V17-direction}
The hard collision cannot be closed by concentrating exclusively on
the higher-incidence overlap.  A dominant overlap is one valid route,
but the empty-overlap example is already compatible with both scalar
response thresholds.  The upstream common problem is therefore the
mixed word
\[
 \text{one complete covariance}\ \otimes\
 \text{one response-carrying endpoint moment}.
\]
The next note should compare this exact word, including its payer,
with the V95--V98 endpoint-energy sandwiches.  Only an exact match of
hypotheses counts as endpoint promotion.
\end{assessment}

\section{Regression checks and result ledger}
\label{sec:V17-results}

\begin{proposition}[V17 regression suite]
\label{prop:V17-regressions}
\begin{enumerate}[label=\textup{(\roman*)},leftmargin=4em]
\item Directed response lower bounds are not used to infer a common
      supporting coordinate.
\item The retained row-mass certificates are distinguished from the
      weaker scalar \(d\)-bounds.
\item Every restricted operator term contains one covariance
      difference.
\item Endpoint moments are not renamed as directed covariances.
\item The V83 square tail is confined to a dominant nonempty overlap
      satisfying its full hypothesis list.
\item The empty-overlap test uses only floor-valued clean incidences
      and fixed response thresholds.
\end{enumerate}
\end{proposition}

\begin{proof}
Items \textup{(i)--(ii)} are
\Cref{lem:V17-overlap-exclusive,
prop:V17-empty-overlap}.  Items \textup{(iii)--(iv)} are
\Cref{thm:V17-four-term-normal-form,
cor:V17-mixed-operator,
fire:V17-no-relabel-moment}.  Item \textup{(v)} is
\Cref{cor:V17-V83-not-exhaustive,
prop:V16-V83-gate}.  The explicit construction in
\Cref{prop:V17-empty-overlap} proves \textup{(vi)}.
\end{proof}

\begin{theorem}[Third-series V17 result ledger]
\label{thm:V17-results}
Relative to V1--V16, V17 proves:
\begin{enumerate}[label=\textup{(V17.\arabic*)},leftmargin=6.3em]
\item a quantitative bidominant-overlap versus exclusive-mass
      alternative for the \(QS\) collision;
\item a sharp clean two-coordinate example with two dominant
      responses and empty overlap;
\item nonexhaustiveness of the V83 higher-incidence route;
\item the exact four-term unpaid/once-paid operator normal form for
      the hard \(QS|D\) collision;
\item identification of the exclusive terms as mixed
      covariance--endpoint-moment words; and
\item a finite necessary operator gate for completing this collision
      class.
\end{enumerate}
\end{theorem}

\begin{proof}
These are
\Cref{lem:V17-overlap-exclusive,
prop:V17-empty-overlap,
cor:V17-V83-not-exhaustive,
thm:V17-four-term-normal-form,
cor:V17-mixed-operator,
prop:V17-QS-criterion}.
\end{proof}

\begin{programme}[V18: audit the endpoint-energy interfaces]
\label{prog:V17-V18}
\begin{enumerate}[leftmargin=3em]
\item Write the first and third lines of
      \eqref{eq:V17-QS-telescope} in the exact V94 Fourier ordering,
      including all local separators and the output dualizer.
\item Compare the response-carrying endpoint moment with each V95
      quadratic \(A/F/I\) energy and the V98 directed-middle endpoint
      hypotheses.  Record the first mismatch if no theorem applies.
\item When an endpoint theorem applies, prove the unpaid line first
      and then enumerate payer in the covariance, endpoint moment,
      separator, and common suffix.
\item Test whether the remainder line in
      \eqref{eq:V17-QS-telescope} cancels with one exclusive line
      before positive envelopes are taken; do not estimate the four
      lines independently if their signs share a common endpoint.
\item Keep the V17 empty-overlap example as a mandatory regression.
\item If endpoint promotion fails, return to the original unsplit
      two-block covariance and seek a single mixed capacity rather
      than another coordinate partition.
\end{enumerate}
\end{programme}

\begin{firewall}[Claim boundary after third-series V17]
\label{fire:V17-claim-boundary}
V17 proves that dominant collision roles need not share any
higher-incidence coordinate and isolates the exact mixed operator
which must replace that failed route.  It does not prove endpoint
promotion, close \(QS|D\) or \(DQS\), close all V13--V14 fallback
transitions, resolve the V2 low-response self-role factor, or handle
nonclean states.  Higher MTA, WUFC, CPM, cutoff removal, reflection
positivity, Osterwalder--Schrader reconstruction, and a uniform
positive continuum spectral gap remain open.  The full Yang--Mills
mass gap has not been proved.
\end{firewall}

\paragraph{Parent-version record.}
V17 was copied byte-for-byte from
\texttt{Renewal\_Yang\_Mills\_Mass\_Gap\_Research\_Notes\_V16.tex}
before this chapter was appended.  The V16 parent had \(8\,225\)
validation records, \(302\,586\) bytes, and SHA--256
\[
 \texttt{13531c2ba1698fc13c6e6429fab181fd33654198417edd50f1377c116255c6b4}.
\]
The next compulsory companion audit remains V20.

% =============================================================================
% BEGIN APPENDED THIRD-SERIES V18 MATERIAL
% =============================================================================

\clearpage
\chapter{V18: Response-carrying moment endpoints and the exact ordering
gate}
\label{ch:V18-endpoint}

\section{Context and archive comparison}
\label{sec:V18-context}

V17 leaves mixed factors of the form
\[
 \Delta_A^{\rho,\sigma}\otimes M_B(\lambda),
\]
where the covariance difference carries one desired response and the
complete endpoint moment carries another.  The moment must not be
renamed as a second covariance.

The V95--V98 audit shows that the appropriate mechanism is the V97
five-factor Schatten sandwich.  A Hilbert--Schmidt endpoint needs a
quadratic one-use energy, but that energy does not intrinsically
require the endpoint to be a covariance.  A complete moment is a
positive contraction.  If its bank carries a fixed dominant response,
the scalar response lower bound converts its universal energy into
the required quadratic response energy.  V18 proves this conversion,
including every payer location, and isolates the remaining datum:
the exact Fourier multiplier must place the moment at an endpoint.

\section{Universal response-to-energy conversion}
\label{sec:V18-energy}

\begin{lemma}[Bounded response occurrence has quadratic energy]
\label{lem:V18-bounded-response-energy}
Let \(A=A^*\) be a complete unpaid occurrence satisfying
\begin{equation}
 \|A\|\le L,\qquad
 \kappa_\otimes(|A|)\le L,
\label{eq:V18-unpaid-caps}
\end{equation}
and suppose its retained physical bank has an actual response
\[
 d_A\ge\gamma_Aa_*.
\]
Then
\begin{equation}
 \boxed{
 {\cal E}_\otimes(A)
 \le
 \frac{L^2}{\gamma_A^2a_*^2}\,d_A^2.}
\label{eq:V18-unpaid-energy}
\end{equation}
If the unique payer is in that occurrence and
\[
 \|A^{\rm pay}\|,
 \ \kappa_\otimes(|A^{\rm pay}|)
 \le\frac{L_{\rm p}}{s_\alpha},
\]
then
\begin{equation}
 \boxed{
 {\cal E}_\otimes(A^{\rm pay})
 \le
 \frac{L_{\rm p}^2}
 {s_\alpha^2\gamma_A^2a_*^2}\,d_A^2.}
\label{eq:V18-paid-energy}
\end{equation}
\end{lemma}

\begin{proof}
By definition of the one-use energy and
\eqref{eq:V18-unpaid-caps},
\[
 {\cal E}_\otimes(A)
 =\|A\|\kappa_\otimes(|A|)
 \le L^2.
\]
The response lower bound gives
\[
 1\le\frac{d_A^2}{\gamma_A^2a_*^2}.
\]
Multiplication proves \eqref{eq:V18-unpaid-energy}.  The paid proof is
identical, starting from
\({\cal E}_\otimes(A^{\rm pay})\le
L_{\rm p}^2/s_\alpha^2\).
\end{proof}

\begin{corollary}[Complete moment endpoint]
\label{cor:V18-moment-endpoint}
Let \(M_B(\lambda)\) be a complete physical endpoint moment.  If its
bank carries an actual response \(d_B\ge\gamma_Ba_*\), then
\begin{align}
 {\cal E}_\otimes(M_B(\lambda))
 &\le
 \frac{d_B^2}{\gamma_B^2a_*^2},
\label{eq:V18-moment-unpaid}\\
 {\cal E}_\otimes(M_B^{\rm pay}(\lambda))
 &\le
 \frac{C}{s_\alpha^2\gamma_B^2a_*^2}\,d_B^2.
\label{eq:V18-moment-paid}
\end{align}
This promotes the moment only as a Hilbert--Schmidt endpoint.  It does
not turn it into a covariance difference.
\end{corollary}

\begin{proof}
A complete heat-kernel moment is a positive contraction, so its norm
and product-state capacity are at most one.  Its paid complete-moment
envelope has both quantities at most \(C/s_\alpha\) by the f2
complete spectator payer cap.  Apply
\Cref{lem:V18-bounded-response-energy}.
\end{proof}

\begin{corollary}[Complete covariance endpoint]
\label{cor:V18-covariance-endpoint}
Let \(\Delta_B^{\rho,\sigma}\) be a complete covariance occurrence
with an actual response \(d_B\ge\gamma_Ba_*\).  Its universal
ordinary norm/capacity and once-paid caps give
\begin{align}
 {\cal E}_\otimes(\Delta_B^{\rho,\sigma})
 &\le
 \frac{C}{\gamma_B^2a_*^2}\,d_B^2,
\label{eq:V18-cov-unpaid}\\
 {\cal E}_\otimes((\Delta_B^{\rho,\sigma})^{\rm pay})
 &\le
 \frac{C}{s_\alpha^2\gamma_B^2a_*^2}\,d_B^2.
\label{eq:V18-cov-paid}
\end{align}
\end{corollary}

\begin{proof}
An unpaid difference of two complete contractions has norm and
capacity at most two.  The f2 universal covariance payer cap bounds
both paid quantities by \(C/s_\alpha\).  Apply
\Cref{lem:V18-bounded-response-energy}.
\end{proof}

\begin{firewall}[A lower response is used only after the complete
energy bound]
\label{fire:V18-lower-after-energy}
The proof first bounds one complete occurrence by one norm and one
capacity.  It then multiplies by a scalar factor already known to be
at least one.  It does not estimate the moment coordinatewise, assign
it a covariance, or create a second payer.
\end{firewall}

\section{A three-response five-factor compiler}
\label{sec:V18-five-factor}

\begin{definition}[Endpoint-separable Fourier ordering]
\label{def:V18-endpoint-separable}
Three response-carrying complete factors \(A,B,E\) are
endpoint-separable when, after the V94 output dualizer has been
absorbed into a separator, every final output multiplier has the
exact equivariant form
\begin{equation}
 K_U=a_Uc_{1,U}b_Uc_{2,U}e_U,
\qquad
 \sup_U\|c_{i,U}\|\le L_i,
\label{eq:V18-endpoint-order}
\end{equation}
where \(a_U,b_U,e_U\) are the restrictions of \(A,B,E\).
The factors \(A,E\) may be complete moments or covariances.  The
middle factor \(B\) is used only through its uniform norm.
\end{definition}

\begin{theorem}[Three-response endpoint compiler]
\label{thm:V18-three-response}
Assume \eqref{eq:V18-endpoint-order}.  Suppose
\[
 d_A\ge\gamma_Aa_*,
 \qquad
 d_B\ge\gamma_Ba_*,
 \qquad
 d_E\ge\gamma_Ea_*,
\]
and suppose \(A,E\) have the universal unpaid and paid caps of
\Cref{cor:V18-moment-endpoint,cor:V18-covariance-endpoint}, while
\[
 \|B\|\le L_B,\qquad
 \|B^{\rm pay}\|\le\frac{L_{B,\rm p}}{s_\alpha}.
\]
Then
\begin{align}
 \Gamma_\otimes(1,K)
 &\le
 \frac{C L_1L_2L_B}
 {\gamma_A\gamma_B\gamma_Ea_*^3}\,
 d_Ad_Bd_E,
\label{eq:V18-three-response-unpaid}\\
 \Gamma_\otimes(w,K^{\rm pay})
 &\le
 \frac{C L_1L_2\max\{L_B,L_{B,\rm p}\}}
 {s_\alpha\gamma_A\gamma_B\gamma_Ea_*^3}\,
 d_Ad_Bd_E
\label{eq:V18-three-response-paid}
\end{align}
for a payer in \(A,B,E\), either finite separator, or the complete
common suffix.
\end{theorem}

\begin{proof}
The exact Fourier sandwich
\(\texttt{thm:V97-five-factor-Fourier}\) gives
\[
 \Gamma_\otimes(1,K)
 \le
 L_1L_2\|B\|
 \sqrt{{\cal E}_\otimes(A){\cal E}_\otimes(E)}.
\]
Insert the two unpaid quadratic energies from
\Cref{lem:V18-bounded-response-energy}; this gives
\[
 \Gamma_\otimes(1,K)
 \le
 \frac{CL_1L_2L_B}
 {\gamma_A\gamma_Ea_*^2}\,d_Ad_E.
\]
Now use
\[
 1\le\frac{d_B}{\gamma_Ba_*}
\]
to prove \eqref{eq:V18-three-response-unpaid}.

If \(A\) or \(E\) pays, its paid quadratic energy contributes exactly
one factor \(s_\alpha^{-1}\) after the square root.  If \(B\) pays,
use \(L_{B,\rm p}/s_\alpha\) for its middle norm and unpaid endpoint
energies.  The finite-separator and complete-suffix payer cases are
the corresponding f2 V95 and V96 one-payer rows.  These locations are
mutually exclusive, so their inverse scales are never multiplied.
Insert the same middle response lower bound to obtain
\eqref{eq:V18-three-response-paid}.
\end{proof}

\begin{corollary}[Legal V15 tree under endpoint ordering]
\label{cor:V18-V15-tree}
If the three edge labels and tails are one row of
\Cref{prop:V15-six-cell-table}, the right sides of
\eqref{eq:V18-three-response-unpaid}--%
\eqref{eq:V18-three-response-paid}
are the legal V15 arborescence response, even when one endpoint factor
is a complete moment.
\end{corollary}

\begin{proof}
The analytic proof of
\Cref{thm:V18-three-response} does not change the three retained
physical edge stocks or their tail denominators.
\Cref{cor:V15-aggregate} identifies their product as the
arborescence rooted at \(z\).
\end{proof}

\begin{firewall}[The theorem is an ordering theorem, not a commuting
rearrangement]
\label{fire:V18-no-reorder}
\Cref{thm:V18-three-response} assumes the exact blockwise ordering
\eqref{eq:V18-endpoint-order}.  It does not commute a moment through
an arbitrary later covariance, output dualizer, or replica
separator.  Coordinate disjointness proves tensor factorization
inside one local occurrence; it does not by itself prove the global
five-factor Fourier order.
\end{firewall}

\section{Application to the V17 exclusive lines}
\label{sec:V18-application}

Consider the first line of
\eqref{eq:V17-QS-telescope}.  It contains
\[
 \Delta_{C_Q}^{\rho,\sigma}
 \otimes M_{C_S}(\sigma)
\]
up to complete moment spectators.  Suppose this line retains
termwise response lower bounds
\begin{equation}
 d_Q(C_Q)\ge\gamma_Q'a_*,
 \qquad
 d_S(C_S)\ge\gamma_S'a_*.
\label{eq:V18-termwise-QS}
\end{equation}
The second estimate follows from
\eqref{eq:V17-S-exclusive} with
\(\gamma_S'=\kappa_S-\theta_S\), by the Casimir floor; the first must
be verified from the actual \(Q\)-stock allocation on \(C_Q\).

\begin{proposition}[Promotion of the \(Q\)-exclusive line]
\label{prop:V18-Q-line}
Assume \eqref{eq:V18-termwise-QS}.  If the exact line is
endpoint-separable with
\[
 A=D,\qquad
 B=\Delta_{C_Q}^{\rho,\sigma},\qquad
 E=M_{C_S}(\sigma),
\]
then it obeys the unpaid and once-paid legal V15 tree bounds with
response product
\[
 d_D\,d_Q(C_Q)\,d_S(C_S).
\]
\end{proposition}

\begin{proof}
The \(D\)-occurrence has the universal complete covariance caps,
the middle \(Q\)-difference has ordinary norm at most two and paid
norm at most \(C/s_\alpha\), and the \(S\)-endpoint is covered by
\Cref{cor:V18-moment-endpoint}.  Apply
\Cref{thm:V18-three-response,cor:V18-V15-tree}.
\end{proof}

\begin{proposition}[Promotion of the \(S\)-exclusive line]
\label{prop:V18-S-line}
Suppose the third line of
\eqref{eq:V17-QS-telescope} retains termwise lower bounds
\[
 d_Q(C_Q)\ge\widetilde\gamma_Qa_*,
 \qquad
 d_S(C_S)\ge\widetilde\gamma_Sa_*.
\]
If it is endpoint-separable with
\[
 A=D,\qquad
 B=\Delta_{C_S}^{\rho,\sigma},\qquad
 E=M_{C_Q}(\rho),
\]
then it obeys the corresponding unpaid and once-paid legal V15 tree
bounds with product
\[
 d_D\,d_S(C_S)\,d_Q(C_Q).
\]
\end{proposition}

\begin{proof}
Interchange \(Q\) and \(S\) in the proof of
\Cref{prop:V18-Q-line}.  The endpoint moment now carries the
\(Q\)-response and the middle difference carries the \(S\)-response.
\end{proof}

\begin{corollary}[The empty-overlap stock example has no energy
obstruction]
\label{cor:V18-empty-overlap-energy}
For the examples in \Cref{prop:V17-empty-overlap}, one has
\[
 O=R=\varnothing,\qquad
 C=C_Q\dot\cup C_S,
\]
and both exclusive subbanks retain fixed termwise response lower
bounds.  Therefore both exact covariance-telescope terms satisfy the
required endpoint energy estimates.  Their only remaining gate is
the endpoint-separable Fourier ordering.
\end{corollary}

\begin{proof}
The computations in \Cref{prop:V17-empty-overlap} give the explicit
lower bounds on \(C_Q=\{e_Q\}\) and \(C_S=\{e_S\}\).  The two-cell
version of \eqref{eq:V17-QS-telescope} has only its first and third
lines.  Apply
\Cref{prop:V18-Q-line,prop:V18-S-line} conditionally on their exact
orders.
\end{proof}

\section{What remains in the overlap and remainder lines}
\label{sec:V18-residual}

\begin{proposition}[Four-line collision status after endpoint
promotion]
\label{prop:V18-four-line-status}
The exact normal form of
\Cref{thm:V17-four-term-normal-form} now has the following status:
\begin{enumerate}[label=\textup{(\roman*)},leftmargin=4em]
\item either exclusive covariance line closes whenever both local
      responses remain dominant and its exact multiplier is
      endpoint-separable;
\item the overlap line remains governed by the full V83 gate of
      \Cref{prop:V16-V83-gate}; and
\item the remainder line still has its covariance difference on
      \(R\), while both \(Q,S\) roles occur only in endpoint moments.
      It is not covered by the one-middle-response theorem
      \Cref{thm:V18-three-response}.
\end{enumerate}
\end{proposition}

\begin{proof}
Item \textup{(i)} is
\Cref{prop:V18-Q-line,prop:V18-S-line}.
Item \textup{(ii)} is the second line of
\eqref{eq:V17-QS-telescope}.  In the remainder line the middle
difference \(\Delta_R\) carries neither the \(Q\)- nor \(S\)-edge, so
there is no third middle response lower bound to insert in
\Cref{thm:V18-three-response}.  This proves \textup{(iii)}.
\end{proof}

\begin{assessment}[The endpoint estimate is no longer the vague gate]
\label{ass:V18-progress}
A response-carrying endpoint moment has the precise unpaid and paid
quadratic energies needed by the V97 Schatten sandwich.  Thus the
exclusive collision terms do not require a new local heat-kernel
estimate.  Their two checkable residual data are:
\[
 \text{termwise response retention}
 \quad\hbox{and}\quad
 \text{endpoint-separable Fourier order}.
\]
The overlap term is a genuine V83 higher-incidence question, while
the remainder term is a distinct two-moment/one-unrelated-difference
problem.  These three operator types must not be conflated.
\end{assessment}

\section{Regression checks and result ledger}
\label{sec:V18-results}

\begin{proposition}[V18 regression suite]
\label{prop:V18-regressions}
\begin{enumerate}[label=\textup{(\roman*)},leftmargin=4em]
\item A complete moment remains a moment throughout endpoint
      promotion.
\item Its response factor is inserted only after a complete energy
      bound.
\item A paid endpoint contributes exactly one
      \(s_\alpha^{-1}\) after the Hilbert--Schmidt square root.
\item A paid middle factor, separator, or suffix is an alternative,
      not an additional debit.
\item No factor is commuted merely from coordinate disjointness.
\item A total response is not assigned to an exclusive telescope term
      without a termwise lower bound.
\item The remainder covariance is not given a \(Q\)- or \(S\)-label.
\end{enumerate}
\end{proposition}

\begin{proof}
Items \textup{(i)--(ii)} are
\Cref{cor:V18-moment-endpoint,
fire:V18-lower-after-energy}.  Items \textup{(iii)--(iv)} are in the
paid proof of \Cref{thm:V18-three-response}.
\Cref{fire:V18-no-reorder} proves \textup{(v)}.
The hypotheses of
\Cref{prop:V18-Q-line,prop:V18-S-line} prove \textup{(vi)}.
\Cref{prop:V18-four-line-status} proves \textup{(vii)}.
\end{proof}

\begin{theorem}[Third-series V18 result ledger]
\label{thm:V18-results}
Relative to V1--V17, V18 proves:
\begin{enumerate}[label=\textup{(V18.\arabic*)},leftmargin=6.3em]
\item universal unpaid and paid quadratic energies for every bounded
      response-carrying complete moment or covariance;
\item a three-response five-factor Fourier compiler with every payer
      location;
\item legal V15 arborescence marking when one Hilbert--Schmidt
      endpoint is a moment;
\item conditional closure of both exclusive lines in the hard
      \(QS|D\) collision;
\item removal of the energy obstruction in the V17 empty-overlap
      tests; and
\item separation of exact ordering, termwise response, overlap, and
      remainder gates.
\end{enumerate}
\end{theorem}

\begin{proof}
These are
\Cref{lem:V18-bounded-response-energy,
cor:V18-moment-endpoint,
cor:V18-covariance-endpoint,
thm:V18-three-response,
cor:V18-V15-tree,
prop:V18-Q-line,
prop:V18-S-line,
cor:V18-empty-overlap-energy,
prop:V18-four-line-status}.
\end{proof}

\begin{programme}[V19: exact Fourier placement of the exclusive
factors]
\label{prog:V18-V19}
\begin{enumerate}[leftmargin=3em]
\item Return to the V61/V94 exact partition-state multiplier before
      its output trace.  Insert the four-cell telescope
      \eqref{eq:V17-QS-telescope} at the original collision
      occurrence.
\item For each exclusive line, follow the cyclic order from the
      separate \(D\)-occurrence to the collision occurrence and list
      every intervening complete separator.
\item Prove the endpoint-separable form
      \eqref{eq:V18-endpoint-order} without commuting through a
      nontrivial shared-coordinate factor, or record the first exact
      obstruction.
\item Track the V2 payer separately in \(D\), the exclusive
      covariance, the endpoint moment, each local separator, and the
      grouped suffix.
\item Partition the total \(Q,S\) row stocks among \(C_Q,O,C_S\) to
      determine which exclusive lines retain both termwise response
      thresholds.
\item Keep the overlap and remainder lines outside the exclusive
      placement proof.
\end{enumerate}
\end{programme}

\begin{firewall}[Claim boundary after third-series V18]
\label{fire:V18-claim-boundary}
V18 proves the mixed endpoint-energy compiler but not its exact
placement in every V100 collision word.  It does not close the
overlap or remainder lines, \(QS|D\), \(DQS\), every V13--V14
fallback transition, the V2 low-response self-role factor, or
nonclean states.  Higher MTA, WUFC, CPM, cutoff removal, reflection
positivity, Osterwalder--Schrader reconstruction, and a uniform
positive continuum spectral gap remain open.  The full Yang--Mills
mass gap has not been proved.
\end{firewall}

\paragraph{Parent-version record.}
V18 was copied byte-for-byte from
\texttt{Renewal\_Yang\_Mills\_Mass\_Gap\_Research\_Notes\_V17.tex}
before this chapter was appended.  The V17 parent had \(8\,685\)
validation records, \(318\,830\) bytes, and SHA--256
\[
 \texttt{15f2ae11ed19dc4e0dd4725d90bb289431d62c51ebac971aa8423b46ac1b74e3}.
\]
The next compulsory companion audit remains V20.

% =============================================================================
% BEGIN APPENDED THIRD-SERIES V19 MATERIAL
% =============================================================================

\clearpage
\chapter{V19: The noncommutative slot ledger for collision endpoints}
\label{ch:V19-slot-ledger}

\section{Context: occurrence data do not determine operator order}
\label{sec:V19-context}

The occurrence map
\[
 \pi:\{D,Q,S\}\longrightarrow{\cal O}
\]
introduced in V16 records which physical covariance occurrence carries
each response role.  It does not record the position of that occurrence
in the exact V61/V94 multiplier.  V18 proves an endpoint-energy compiler
only for the block order
\[
 A\,C_1\,B\,C_2\,E.
\]
Consequently a separate slot datum is required before the mixed
endpoint moment in the V17 telescope may be inserted into the V97
Schatten inequality.

This is not a cosmetic ordering choice.  The product-state density is
fixed on the outside of the multiplier in the V94 dual radius.  Cyclic
trace invariance does not permit an arbitrary factor to cross that
density.  V19 formalizes the missing slot ledger, proves the exact
adjoint reversal which is legal, and gives a finite sufficient
certificate for V18 endpoint promotion.

\section{Why arbitrary cyclic rerooting is illegal}
\label{sec:V19-no-cyclic}

\begin{proposition}[A product-state trace is not freely cyclically
rerooted]
\label{prop:V19-no-cyclic-reroot}
There are a density matrix \(\rho\) and self-adjoint contractions
\(X,Y,Z\) such that
\begin{equation}
 \Tr(\rho XYZ)\ne\Tr(\rho YZX).
\label{eq:V19-noncyclic-example}
\end{equation}
Thus cyclicity of the ordinary trace cannot move \(X\) through the
fixed density \(\rho\) while retaining the form
\(\Tr(\rho\,\cdot)\).
\end{proposition}

\begin{proof}
On \(\mathbb C^2\), let
\[
 \rho=P_0=
 \begin{pmatrix}1&0\\0&0\end{pmatrix},
 \qquad
 X=Y=
 \begin{pmatrix}0&1\\1&0\end{pmatrix},
 \qquad
 Z=P_0.
\]
Then \(XYZ=P_0\), so \(\Tr(\rho XYZ)=1\).  On the other hand,
\[
 YZX=
 \begin{pmatrix}0&1\\1&0\end{pmatrix}
 P_0
 \begin{pmatrix}0&1\\1&0\end{pmatrix}
 =
 \begin{pmatrix}0&0\\0&1\end{pmatrix},
\]
and hence \(\Tr(\rho YZX)=0\).
\end{proof}

\begin{firewall}[What trace cyclicity actually gives]
\label{fire:V19-cyclicity}
One has
\[
 \Tr(\rho XYZ)=\Tr(Z\rho XY),
\]
not \(\Tr(\rho YZX)\).  The latter equality would require an
additional commutation or invariance statement involving \(\rho\).
The supremum over product densities in \(\Gamma_\otimes\) does not
authorize such a pointwise replacement inside the V97 proof.
\end{firewall}

\begin{lemma}[Adjoint reversal is legal]
\label{lem:V19-adjoint-reversal}
For every density matrix \(\rho\) and bounded multiplier \(K\),
\[
 |\Tr(\rho K)|=|\Tr(\rho K^*)|.
\]
If
\[
 K=A\,C_1\,B\,C_2\,E
\]
with \(A,B,E\) self-adjoint, then
\[
 K^*=E\,C_2^*\,B\,C_1^*\,A.
\]
Thus the two Hilbert--Schmidt endpoints may be exchanged by reversing
the entire word.  No interior factor may be cyclically rerooted by
this lemma.
\end{lemma}

\begin{proof}
Complex conjugation and trace cyclicity give
\[
 \overline{\Tr(\rho K)}
 =\Tr(K^*\rho)
 =\Tr(\rho K^*).
\]
Taking absolute values proves the first assertion.  The second is the
adjoint of a product in reverse order.
\end{proof}

\section{The exact slot ledger}
\label{sec:V19-ledger}

For a fixed final output block \(U\), write the exact equivariant
multiplier before the dual contraction as
\begin{equation}
 K_U=T_{0,U}O_{1,U}T_{1,U}\cdots O_{m,U}T_{m,U}.
\label{eq:V19-slot-word}
\end{equation}
Here \(O_j\) are the finitely many distinguished complete physical
occurrences and \(T_j\) are the intervening complete V61 words,
replica regroupings, and local separators.  The complete common
suffix is kept as the separately allocated V96 factor when it is not
inside one \(T_j\).

\begin{definition}[Enriched occurrence certificate]
\label{def:V19-enriched-certificate}
For each response role \(R\in\{D,Q,S\}\), retain:
\begin{enumerate}[label=\textup{(\roman*)},leftmargin=4em]
\item its physical occurrence \(\pi(R)\in{\cal O}\);
\item its slot \(s(R)\in\{1,\ldots,m\}\);
\item its exact local subbank and whether that subbank is a covariance
      difference or an endpoint moment after restriction;
\item its position on the left or right side of every internal
      two-term covariance telescope; and
\item the unique payer flag, if the payer lies in that occurrence.
\end{enumerate}
The collection of these data is the slot ledger
\(\mathfrak L(D,Q,S)\).
\end{definition}

\begin{lemma}[Finite size of the slot ledger]
\label{lem:V19-ledger-finite}
For a four-row V100 first-breaker packet, the number of distinguished
occurrence slots, local telescope sides, and payer alternatives in
\(\mathfrak L(D,Q,S)\) is bounded by a universal constant independent
of coordinate support and representations.
\end{lemma}

\begin{proof}
The matching word has at most four directed covariance occurrences,
the breaker contributes one occurrence, and the V61 local separators
and output dualizer have fixed four-row length.  V12--V14 add only a
fixed number of selected response roles.  Every covariance telescope
has two sides, and the V2 allocation places one payer in one
occurrence, separator, or grouped suffix.  Taking the product of these
finite choices gives a universal bound.
\end{proof}

\begin{remark}[Occurrence partitions are a quotient of the ledger]
\label{rem:V19-ledger-refines}
Forgetting the slot, local side, and payer data in
\(\mathfrak L(D,Q,S)\) leaves the five occurrence partitions of
\Cref{prop:V16-five-partitions}.  Therefore the V16 partition is
necessary for collision detection but is not sufficient for a
noncommutative endpoint estimate.
\end{remark}

\section{Boundary exposure of a moment subfactor}
\label{sec:V19-exposure}

\begin{definition}[Right-exposed collision factor]
\label{def:V19-right-exposed}
Let \(C\) be a collision occurrence and suppose an exact telescope
term restricts on every output block as
\[
 c_U=b_Ue_U,
\]
where \(b_U\) is the sole covariance difference and \(e_U\) is a
self-adjoint complete endpoint moment on a disjoint coordinate
subbank.  The moment is right-exposed relative to a distinct response
occurrence \(A\) when, after removing only complete one-sided
spectators by their established tensor law, the remaining exact
multiplier is
\begin{equation}
 K_U=a_U\ell_Ub_Ue_U,
\qquad
 \sup_U\|\ell_U\|\le L.
\label{eq:V19-right-exposed}
\end{equation}
Left exposure is defined by adjoint reversal.
\end{definition}

\begin{lemma}[Disjoint telescope factors commute internally]
\label{lem:V19-internal-commutation}
In every line of \eqref{eq:V17-QS-telescope}, the covariance
subfactor and all endpoint-moment subfactors belonging to distinct
cells of \eqref{eq:V17-four-cells} commute before the intervening
global word is multiplied.
\end{lemma}

\begin{proof}
The cells are disjoint coordinate sets.  Their operators act on
different coordinate tensor factors, so
\[
 (B\otimes I)(I\otimes E)
 =(I\otimes E)(B\otimes I).
\]
This statement is confined to subfactors of the one collision
occurrence.  It says nothing about a later global occurrence or
separator acting on a shared row representation space.
\end{proof}

\begin{theorem}[Exposed endpoint implies V18 ordering]
\label{thm:V19-exposed-order}
If an exclusive V17 collision line has the right-exposed form
\eqref{eq:V19-right-exposed}, then it is endpoint-separable in the
sense of \Cref{def:V18-endpoint-separable}, with
\[
 A=a,\qquad C_1=\ell,\qquad B=b,\qquad C_2=I,\qquad E=e.
\]
The same conclusion holds for a left-exposed line after applying
\Cref{lem:V19-adjoint-reversal}.  All payer alternatives retain
exactly one paid factor.
\end{theorem}

\begin{proof}
Equation \eqref{eq:V19-right-exposed} is precisely
\eqref{eq:V18-endpoint-order} with \(C_2=I\).
Self-adjointness of \(A,B,E\), equivariance, and the uniform separator
bound are part of the definition.  For left exposure, reverse the
whole multiplier using \Cref{lem:V19-adjoint-reversal}; the norms of
the adjoint separators are unchanged.

If the payer lies in \(A,B,E\), retain the corresponding paid
restriction.  A payer in \(\ell\) uses its paid separator norm, and a
grouped suffix payer remains the V96 alternative removed before
\eqref{eq:V19-right-exposed} is formed.  These cases are disjoint.
\end{proof}

\begin{corollary}[Closure certificate for an exclusive collision
line]
\label{cor:V19-exclusive-closure}
Suppose an exclusive line retains the three termwise response lower
bounds required by
\Cref{prop:V18-Q-line} or
\Cref{prop:V18-S-line}.  If its slot ledger proves right or left
exposure, then that exact line closes at every payer location by the
legal V15 arborescence bound.
\end{corollary}

\begin{proof}
Use \Cref{thm:V19-exposed-order} to verify the ordering hypothesis of
\Cref{prop:V18-Q-line} or
\Cref{prop:V18-S-line}.
\end{proof}

\section{What may be absorbed into the middle separator}
\label{sec:V19-separator}

\begin{proposition}[Uniform interior absorption]
\label{prop:V19-interior-absorption}
In \eqref{eq:V19-slot-word}, any consecutive subword \(W_U\) between
the two proposed endpoints may be absorbed into \(C_1\), \(B\), or
\(C_2\) whenever one has a complete support-uniform bound
\[
 \sup_U\|W_U\|\le L_W.
\]
The constants in the V18 compiler are multiplied by \(L_W\).
If \(W\) pays, it may be absorbed only through its established
one-payer bound \(CL_W/s_\alpha\).
\end{proposition}

\begin{proof}
Group consecutive multiplication factors associatively.  The V97
five-factor theorem uses the separator factors only through their
operator norms and the middle factor only through its norm.  Replacing
one by its product with \(W\) multiplies the corresponding bound by
\(L_W\).  The paid assertion is the same calculation with the unique
paid norm in place of the unpaid norm.
\end{proof}

\begin{firewall}[A stock-carrying occurrence is not a free separator]
\label{fire:V19-not-free-separator}
If a subword has only a stock-dependent capacity or energy estimate
and no support-uniform ordinary norm, it cannot be absorbed by
\Cref{prop:V19-interior-absorption}.  Likewise, a factor outside the
two chosen endpoints cannot be moved across the product density by
calling it a separator.  Its exact slot must first be retained.
\end{firewall}

\begin{proposition}[First exact ordering obstruction]
\label{prop:V19-order-obstruction}
For an exclusive \(QS|D\) line, endpoint exposure can fail only if,
after complete one-sided spectators are removed, at least one of the
following occurs:
\begin{enumerate}[label=\textup{(\roman*)},leftmargin=4em]
\item a nonuniform distinguished occurrence lies outside the cyclic
      interval between \(D\) and the collision occurrence;
\item an intervening subword has no support-uniform ordinary or paid
      norm;
\item the desired endpoint moment is separated from the boundary of
      its collision occurrence by a noncommuting internal factor not
      covered by \Cref{lem:V19-internal-commutation}; or
\item the payer is allocated to a factor for which no corresponding
      paid separator, occurrence, or grouped-suffix line is proved.
\end{enumerate}
\end{proposition}

\begin{proof}
If none of \textup{(i)--(iv)} occurs, choose the word orientation
which places \(D\) first and the collision occurrence last; use
\Cref{lem:V19-adjoint-reversal} if the reverse orientation is needed.
Absorb all interior uniformly bounded subwords by
\Cref{prop:V19-interior-absorption}.  Use
\Cref{lem:V19-internal-commutation} to place the endpoint moment at
the collision boundary.  The remaining word has
\eqref{eq:V19-right-exposed}, including a legal payer row, contrary
to failure of exposure.  Hence failure implies at least one listed
obstruction.
\end{proof}

\section{Audit of the active collision record}
\label{sec:V19-active-audit}

\begin{proposition}[The current record lacks a slot certificate]
\label{prop:V19-missing-slot-data}
The hypotheses recorded in V12--V18 determine:
\begin{enumerate}[label=\textup{(\alph*)},leftmargin=4em]
\item the response tails and actual partner labels;
\item the physical occurrence partition \(\pi\);
\item the coordinate support split
      \(C_Q\dot\cup O\dot\cup C_S\dot\cup R\); and
\item the one-payer restriction telescope.
\end{enumerate}
They do not determine the slot map \(s\), the cyclic interval of the
distinguished occurrences in \eqref{eq:V19-slot-word}, or which side
of every noncommuting separator contains the endpoint moment.
Therefore endpoint exposure is not yet proved for every V17
exclusive line.
\end{proposition}

\begin{proof}
Items \textup{(a)--(d)} are respectively
\Cref{prop:V15-six-cell-table,
prop:V16-five-partitions,
eq:V17-four-cells,
thm:V17-four-term-normal-form}.
Those statements index banks, supports, response labels, and payer
alternatives but contain no ordered factorization of the form
\eqref{eq:V19-slot-word}.  V18 consequently states exposure as an
explicit hypothesis.  Since
\Cref{prop:V19-no-cyclic-reroot} forbids recovering the missing order
by cyclic notation alone, the existing data do not prove the
certificate.
\end{proof}

\begin{assessment}[Upstream correction after V19]
\label{ass:V19-correction}
The exclusive collision is no longer blocked by an unknown endpoint
norm: V18 proves that norm and V19 proves the exact sufficient
placement theorem.  What is absent is a finite combinatorial datum
from the source construction itself.  The next calculation must
return to each V100/V12 packet word, attach a slot number to every
selected response occurrence before taking the positive envelope, and
then apply \Cref{prop:V19-order-obstruction}.  Repeating scalar
acquisition cannot supply this noncommutative information.
\end{assessment}

\section{Regression checks and result ledger}
\label{sec:V19-results}

\begin{proposition}[V19 regression suite]
\label{prop:V19-regressions}
\begin{enumerate}[label=\textup{(\roman*)},leftmargin=4em]
\item No factor is cyclically moved through a product-state density.
\item Reversal is performed only by taking the adjoint of the entire
      word.
\item Internal commutation is used only for disjoint tensor subbanks
      of one collision occurrence.
\item A uniformly bounded interior word is distinguished from a
      stock-carrying endpoint energy.
\item The slot ledger retains occurrence, side, and payer data.
\item Endpoint closure is asserted only after a right- or left-exposure
      certificate is proved.
\end{enumerate}
\end{proposition}

\begin{proof}
Items \textup{(i)--(ii)} are
\Cref{prop:V19-no-cyclic-reroot,
fire:V19-cyclicity,
lem:V19-adjoint-reversal}.  Item \textup{(iii)} is
\Cref{lem:V19-internal-commutation}.  Item \textup{(iv)} is
\Cref{prop:V19-interior-absorption,
fire:V19-not-free-separator}.  Items \textup{(v)--(vi)} are
\Cref{def:V19-enriched-certificate,
thm:V19-exposed-order,
cor:V19-exclusive-closure}.
\end{proof}

\begin{theorem}[Third-series V19 result ledger]
\label{thm:V19-results}
Relative to V1--V18, V19 proves:
\begin{enumerate}[label=\textup{(V19.\arabic*)},leftmargin=6.3em]
\item a concrete noncommutative counterexample to arbitrary cyclic
      rerooting through the product-state density;
\item the exact whole-word adjoint reversal which is permitted;
\item a finite enriched slot ledger refining the five occurrence
      partitions;
\item a right/left endpoint-exposure theorem implying the V18
      five-factor order;
\item a complete list of four possible ordering obstructions; and
\item identification of the missing slot data in the active
      V12--V18 collision record.
\end{enumerate}
\end{theorem}

\begin{proof}
These are
\Cref{prop:V19-no-cyclic-reroot,
lem:V19-adjoint-reversal,
def:V19-enriched-certificate,
lem:V19-ledger-finite,
thm:V19-exposed-order,
prop:V19-order-obstruction,
prop:V19-missing-slot-data}.
\end{proof}

\begin{programme}[V20: compulsory audit and packet-by-packet slot
reconstruction]
\label{prog:V19-V20}
\begin{enumerate}[leftmargin=3em]
\item Perform the compulsory read-only audit of the newest
      Standard-Model--Einstein and Navier--Stokes notes before choosing
      the reconstruction order.
\item Return to the exact V100 first-breaker word
      \(B_\mu S_G^\mu\) and its V61 coordinate order.
\item For the V12 fresh, old-external, and self-response branches,
      record the slot of \(D\); then record the slot of the V13
      \(q\)-response \(Q\) and the inherited right response \(S\).
\item For every \(QS|D\) exclusive telescope line, test the four
      obstructions in
      \Cref{prop:V19-order-obstruction}.
\item Close every exposed line by
      \Cref{cor:V19-exclusive-closure}; retain an exact named word for
      every nonexposed line.
\item Keep the V17 overlap and remainder lines separate from the slot
      reconstruction.
\end{enumerate}
\end{programme}

\begin{firewall}[Claim boundary after third-series V19]
\label{fire:V19-claim-boundary}
V19 proves the sufficient endpoint-order theorem and shows that the
current collision record lacks the slot data needed to invoke it
universally.  It does not reconstruct those slots, close the V17
overlap or remainder lines, \(QS|D\), \(DQS\), every V13--V14
fallback transition, the V2 low-response self-role factor, or
nonclean states.  Higher MTA, WUFC, CPM, cutoff removal, reflection
positivity, Osterwalder--Schrader reconstruction, and a uniform
positive continuum spectral gap remain open.  The full Yang--Mills
mass gap has not been proved.
\end{firewall}

\paragraph{Parent-version record.}
V19 was copied byte-for-byte from
\texttt{Renewal\_Yang\_Mills\_Mass\_Gap\_Research\_Notes\_V18.tex}
before this chapter was appended.  The V18 parent had \(9\,205\)
validation records, \(336\,066\) bytes, and SHA--256
\[
 \texttt{3922dc72b0a74e60316a79dc3fbd40da25d6901bec279e152eb77664498b61bd}.
\]
The next iteration is the compulsory V20 companion audit.

% =============================================================================
% BEGIN APPENDED THIRD-SERIES V20 MATERIAL
% =============================================================================

\clearpage
\chapter{V20: Compulsory audit and ancestry-preserving slot
refinement}
\label{ch:V20-ancestry-slots}

\section{Companion audit and proof-order decision}
\label{sec:V20-audit}

This is the compulsory five-version checkpoint following V15.  The
newest active Standard-Model--Einstein and Navier--Stokes notes were
read passively before the slot reconstruction was chosen.

\begin{record}[Companion snapshots used by V20]
\label{rec:V20-companion-snapshots}
The files actually inspected at the checkpoint were:
\begin{enumerate}[leftmargin=3em]
\item
\(\texttt{Renewal\_SM\_Einstein\_Continuum\_Research\_Notes\_V97.tex}\),
with \(1\,429\,021\) bytes and SHA--256
\[
 \texttt{b300e5aa03b037946ff6e34930acc407c47fcfe0cc631ec879da2e2aedca221c};
\]
\item
\(\texttt{Renewal\_NS\_Stability\_Research\_Notes\_V31.tex}\),
with \(23\,052\) source lines, \(887\,537\) bytes, and SHA--256
\[
 \texttt{f1121b62238ad5491bb7cf03635ea9934942edf573ed8faaa9ee79e1d38a6696}.
\]
\end{enumerate}
The reads were passive.  Neither companion programme nor either
frozen Yang--Mills archive was changed.
\end{record}

\begin{assessment}[Type-correct transfer from the companions]
\label{ass:V20-audit-transfer}
NS V31 proves that a flat probability mark may remove a stopping-time
multiplicity, but exact removal of the mark restores the missing
dyadic first moment.  The Yang--Mills consequence is that the possible
positions of a selected response may not be summed after each
position has been estimated separately.  The response must keep its
parent occurrence and payer until the complete word is bounded.

SM V97 starts from one exact determinant factorization and keeps its
zero-mode line, four local ultraviolet terms, four complementary
infrared terms, and fifth regularized determinant in disjoint
sectors.  Its stopping-ball estimate uses bounded incidence and pays
each curvature-energy occurrence once.  The relevant Yang--Mills
lesson is the same: restriction may refine one occurrence, but it may
not turn that occurrence into an endpoint and an independently
bounded middle factor.

No Navier--Stokes estimate, determinant theorem, Uhlenbeck theorem, or
Standard-Model cancellation is imported here.  Only the exact
one-occurrence/one-payer proof order is transferred.
\end{assessment}

\begin{assessment}[Nonduplication boundary at V20]
\label{ass:V20-nonduplication}
V14 already proves the simultaneous subset restriction of a V100
packet, V16 already proves the arbitrary finite one-difference
telescope, and V19 already defines a slot ledger and an endpoint
exposure certificate.  V20 does not repeat those statements.  It
constructs the missing parent map for every response produced by the
V12--V14 acquisition, proves that the map survives all restriction
telescopes, and partitions the exclusive \(QS|D\) residual by the
actual boundary status of its two parent occurrences.
\end{assessment}

\section{The pre-envelope word and parent occurrences}
\label{sec:V20-parent}

Fix one exact, signed, pre-envelope first-breaker packet.  Before its
V94 output contraction, its connected coordinate carrier has the
V61 factorization
\begin{equation}
 {\cal W}_{g,\rho}
 =
 C_{<g}(\rho)\otimes J_g(\rho)\otimes M_{>g}.
\label{eq:V20-V61-word}
\end{equation}
After the finite replica regroupings and output maps are inserted,
write its exact multiplication word as in V19:
\begin{equation}
 K_U=T_{0,U}O_{1,U}T_{1,U}\cdots O_{m,U}T_{m,U}.
\label{eq:V20-multiplication-word}
\end{equation}
The \(O_j\) are complete physical covariance or moment occurrences,
not scalar stocks.  The \(T_j\) are the unchanged complete V61
intertwiners and finite separators.  A separately grouped common
suffix retains its V96 label.

\begin{definition}[Ancestry tag]
\label{def:V20-ancestry-tag}
Let \(R\in\{D,Q,S\}\) be a response selected during the V12--V14
acquisition.  An ancestry tag for \(R\) is
\begin{equation}
 {\bf a}(R)=
 \bigl(j(R),Z(R),\vartheta(R),\epsilon(R),p(R)\bigr),
\label{eq:V20-ancestry-tag}
\end{equation}
where:
\begin{enumerate}[label=\textup{(\roman*)},leftmargin=4em]
\item \(O_{j(R)}\) is the unique parent occurrence;
\item \(Z(R)\) is its selected coordinate subbank;
\item \(\vartheta(R)\) records whether the selected local factor is a
      covariance difference, an endpoint moment, or a self-role
      difference;
\item \(\epsilon(R)\) records the side of every V14/V16 covariance
      telescope leading to the selected term; and
\item \(p(R)\) is the inherited payer flag, equal to zero when the
      payer is elsewhere.
\end{enumerate}
The slot component of the V19 ledger is
\(s(R):=j(R)\).
\end{definition}

\begin{lemma}[Unique parent for an old or fresh acquisition]
\label{lem:V20-unique-parent}
Every clean response selected by the V12 post-prefix acquisition or
the V14 depleted-target acquisition has a unique ancestry tag before
the positive envelope is taken.
\end{lemma}

\begin{proof}
For an old-bank alternative, V12 and V14 assign every common
coordinate to its earliest physical role.  The resulting old roles
are disjointly represented, and
\Cref{thm:V14-subset-resolved-packet} retains the position of each
parent occurrence in \eqref{eq:V20-multiplication-word}.  Thus the
selected old subbank has one parent \(O_j\).

For a fresh cell \(Z\), the proof of
\Cref{lem:V14-target-identity} factors one residual complete moment
exactly as
\[
 M_F(\rho)=M_Z(\rho)\otimes M_{F\setminus Z}(\rho).
\]
The second tensor factor remains a complete spectator.  Hence \(Z\)
is a child of the unique residual moment occurrence containing
\(M_F(\rho)\); it is not a new occurrence inserted at an arbitrary
point of the multiplication word.

For a self-anchored alternative, the V2 role telescope is applied
inside the selected old parent bank.  Each cover difference and its
complementary endpoint moment therefore inherit that same parent
slot.  Finally, V14 splits a payer label by disjoint coordinate
indicators before norming.  This determines \(\epsilon(R)\) and
\(p(R)\) without copying the payer.  These cases exhaust the clean
acquisition alternatives.
\end{proof}

\begin{firewall}[A coordinate index is not a multiplication slot]
\label{fire:V20-coordinate-not-slot}
The index \(g\) in \eqref{eq:V20-V61-word} orders physical coordinate
tensors.  The index \(j\) in
\eqref{eq:V20-multiplication-word} orders noncommuting multipliers
after replica regrouping and output dualization.  A canonical
permutation of disjoint coordinate tensors does not permit a
permutation of the \(O_j\) through a fixed product-state density.
V20 never identifies \(g\) with \(j\).
\end{firewall}

\begin{lemma}[Ancestry is stable under an arbitrary finite support
telescope]
\label{lem:V20-ancestry-stable}
Apply the V16 one-difference telescope to any finite partition
\(\Omega=\dot\bigcup_{\nu=1}^N B_\nu\) inside one occurrence \(O_j\).
Every nonzero summand has:
\begin{enumerate}[label=\textup{(\alph*)},leftmargin=4em]
\item exactly one covariance child, whose parent slot is \(j\);
\item endpoint-moment children in the same parent slot;
\item the original relative order of all occurrences \(O_i\),
      \(i\ne j\); and
\item exactly the inherited zero-or-one payer flag.
\end{enumerate}
In particular, no refinement changes \(s(R)\).
\end{lemma}

\begin{proof}
The V16 telescope is
\[
 \Delta_\Omega^{\rho,\sigma}
 =
 \sum_{\nu=1}^N
 \left(\bigotimes_{\kappa<\nu}M_{B_\kappa}(\rho)\right)
 \otimes\Delta_{B_\nu}^{\rho,\sigma}
 \otimes
 \left(\bigotimes_{\kappa>\nu}M_{B_\kappa}(\sigma)\right).
\]
All factors displayed on the right replace the single factor
\(\Delta_\Omega^{\rho,\sigma}\) at its original position \(j\).
There is one difference in every summand.  All other global
occurrences are untouched.  The payer-coordinate indicator belongs
to exactly one \(B_\nu\), so the V14 payer split gives at most one
nonzero payer flag in that summand.  This proves all four assertions.
\end{proof}

\section{The branch-by-branch slot table}
\label{sec:V20-slot-table}

Use the V12 labels
\[
 D=x\longrightarrow u,\qquad
 S=y\longrightarrow z,
\]
and let \(Q=q\longrightarrow r\) be the response acquired at the
mate or at a later depleted target.  Denote by \(j_S\) the occurrence
of the inherited matching response inside \(B_\mu\).

\begin{proposition}[Symbolic slot reconstruction for every clean
branch]
\label{prop:V20-slot-table}
Before any positive envelope, the three response slots are determined
by the following exhaustive table:
\begin{center}
\begin{tabular}{c c c}
\toprule
role & acquisition class & slot\\
\midrule
\(S\) & inherited matching response & \(j_S\)\\
\(D\) & old external bank \(B_u\) & \(j(B_u)\)\\
\(D\) & fresh directed cell in \(M_{F_D}\) & \(j(F_D)\)\\
\(D\) & response-first self-role in \(B_x\) & \(j(B_x)\)\\
\(Q\) & old external bank \(B_r\) & \(j(B_r)\)\\
\(Q\) & fresh directed cell in \(M_{F_Q}\) & \(j(F_Q)\)\\
\(Q\) & response-first self-role in \(B_q\) or \(B_r\) &
        \(j(B_q)\) or \(j(B_r)\)\\
\bottomrule
\end{tabular}
\end{center}
The displayed \(j(\cdot)\) is the parent index in the fixed exact word
\eqref{eq:V20-multiplication-word}; it is not chosen after the
response partner is known.  A low-response self-role term retains the
tag of its parent bank but is not promoted to a directed-response row
of this table.
\end{proposition}

\begin{proof}
The \(S\)-response is one of the two matching directed covariances
already assembled in \(B_\mu\), so its parent is fixed.  For \(D\),
the three clean directed possibilities are exactly the old-external,
fresh-directed, and response-first self-role alternatives in
\Cref{thm:V12-post-prefix-acquisition,
thm:V12-self-reduction}.  Apply
\Cref{lem:V20-unique-parent} to each.

In the V13 endpoint-then-mate construction, every existing bank,
including the \(D\)-occurrence, is placed in the used union before
the \(q\)-acquisition.  Therefore a fresh \(Q\) has a residual moment
parent disjoint from the previously selected response subbanks, while
an old or self-role \(Q\) inherits the recorded parent bank and may
collide with \(S\).  The V14 target acquisition has the identical
old/fresh/self trichotomy after its support restriction.  The
low-response alternative is only an upper cover currency by V2 and
has no directed lower bound, which proves the last sentence.
\end{proof}

\begin{corollary}[Collision equalities are exact slot equalities]
\label{cor:V20-collision-slots}
In the slot-refined packet,
\[
 QS|D
 \quad\Longleftrightarrow\quad
 j(Q)=j(S)\ne j(D),
\]
and
\[
 DQS
 \quad\Longleftrightarrow\quad
 j(D)=j(Q)=j(S),
\]
with equality understood at the level of physical parent
occurrences, not merely equal coordinate supports.
\end{corollary}

\begin{proof}
The V16 occurrence partition records equality of physical
occurrences.  By \Cref{lem:V20-unique-parent}, every selected response
has exactly one such parent, and by
\Cref{lem:V20-ancestry-stable} no restriction creates a second one.
The two displayed equivalences are therefore the definitions of the
two occurrence partitions written in parent-slot notation.
\end{proof}

\begin{theorem}[Exact finite slot refinement of the collision
packet]
\label{thm:V20-finite-slot-refinement}
Every fixed unpaid or once-paid V100 packet in the clean
V12--V17 residual is an exact finite signed sum of words on which
\[
 \bigl({\bf a}(D),{\bf a}(Q),{\bf a}(S)\bigr)
\]
is constant.  The number of tag classes is bounded by a universal
constant \(N_{\rm slot}\), independently of coordinate support,
representations, and physical multiplicity.  Every word in the sum
contains exactly the original physical occurrences and at most the
original single payer.
\end{theorem}

\begin{proof}
\Cref{thm:V14-subset-resolved-packet} produces at most \(2^{J_*}\)
restriction words, where \(J_*\) is the universal number of complete
covariance occurrences.  By
\Cref{lem:V19-ledger-finite}, the number of parent slots, telescope
sides, and payer positions is universally bounded.  On each resolved
word, \Cref{lem:V20-unique-parent} assigns one of those tags to each
of the three response roles.  Partitioning by the resulting triple
therefore produces a universal finite number \(N_{\rm slot}\) of
classes.

This is a partition of the already constructed signed expansion; no
absolute value is taken and no sum over a freely chosen coordinate
mark is added.  The occurrence and payer conclusions follow from
\Cref{lem:V20-ancestry-stable}.
\end{proof}

\begin{firewall}[The slot refinement has no hidden first moment]
\label{fire:V20-no-slot-count}
The constant \(N_{\rm slot}\) counts only the fixed four-row role,
telescope-side, and payer alternatives.  It does not count the
coordinates inside \(Z(R)\).  Replacing a complete tag class by a sum
of separately normed coordinate positions would reintroduce the
first-moment loss excluded by V100 and by the NS V31 audit.
\end{firewall}

\section{Boundary status is inherited from the parent}
\label{sec:V20-boundary}

\begin{definition}[Legally reduced multiplication skeleton]
\label{def:V20-reduced-skeleton}
For a fixed tagged word \(K_U\), its legally reduced skeleton is
obtained only by the following operations:
\begin{enumerate}[label=\textup{(\roman*)},leftmargin=4em]
\item remove a positive complete tensor spectator when the exact
      factorization required by the V2 one-sided tensor law has been
      exhibited;
\item keep a grouped common-suffix payer outside the core by the V96
      support-uniform suffix theorem;
\item group consecutive interior multipliers having their established
      support-uniform ordinary or once-paid norm; and
\item reverse the entire word by the V19 adjoint identity.
\end{enumerate}
No general multiplier is moved through the product density.  A parent
slot is a left or right boundary slot when it is the corresponding
outer distinguished occurrence of this reduced skeleton.
\end{definition}

\begin{lemma}[Restriction cannot manufacture boundary exposure]
\label{lem:V20-boundary-inheritance}
Let a response child \(R\) lie in the parent occurrence \(O_j\).
Support restriction and the V14/V16 telescopes preserve the boundary
status of \(j\).  If \(j\) is an interior slot, no child supported on
a subbank of \(O_j\) becomes an outer multiplication endpoint merely
by restriction.  If \(j\) is a boundary slot, a covariance child and
an endpoint-moment child on disjoint subbanks may be placed at that
same boundary in either internal tensor order.
\end{lemma}

\begin{proof}
By \Cref{lem:V20-ancestry-stable}, restriction replaces \(O_j\) at
slot \(j\) and leaves every \(O_i\), \(i\ne j\), in place.  Thus every
global factor preceding or following \(O_j\) remains on the same
side; an interior parent cannot become an outer factor.

For a boundary parent, the covariance and endpoint moment produced by
the support telescope act on disjoint coordinate tensor factors.
\Cref{lem:V19-internal-commutation} permits their internal order to be
chosen without crossing any global separator.  Their parent remains
the same boundary occurrence.  This proves both assertions.
\end{proof}

\begin{definition}[The two exposed \(QS|D\) orientations]
\label{def:V20-exposed-classes}
Fix an exclusive \(QS|D\) telescope line and put
\[
 j_C:=j(Q)=j(S),\qquad j_A:=j(D).
\]
It belongs to \({\mathfrak E}_{A\to C}\) when, after the legal
reductions of \Cref{def:V20-reduced-skeleton}, \(j_A\) is the left
boundary slot, \(j_C\) is the right boundary slot, and the collision
endpoint moment is placed to the right of its covariance child.
It belongs to \({\mathfrak E}_{C\leftarrow A}\) when the adjoint word
has that property.  Let \({\mathfrak I}_{QS|D}\) be the complementary
tagged class.
\end{definition}

\begin{proposition}[The exposed/interior split is exact and
support-uniform]
\label{prop:V20-boundary-split}
For every exclusive \(QS|D\) residual,
\begin{equation}
 {\cal P}_{QS|D}
 =
 {\cal P}_{{\mathfrak E}_{A\to C}}
 +
 {\cal P}_{{\mathfrak E}_{C\leftarrow A}}
 +
 {\cal P}_{{\mathfrak I}_{QS|D}}
\label{eq:V20-boundary-split}
\end{equation}
is an exact disjoint partition of the slot-refined signed packet.
It preserves the termwise response lower bounds, physical occurrence
count, and unique payer.  Its number of classes is at most
\(N_{\rm slot}\).
\end{proposition}

\begin{proof}
On each tag class of
\Cref{thm:V20-finite-slot-refinement}, the ordered parent slots,
telescope side, and payer position are constant.  The boundary test
in \Cref{def:V20-exposed-classes} therefore has one of the three
displayed outcomes.  The classes are mutually exclusive and
exhaustive.  No analytic estimate or new coordinate selection enters
the test, so all lower bounds and the payer label are inherited.
\end{proof}

\begin{theorem}[Closure of both boundary-exposed collision classes]
\label{thm:V20-exposed-closure}
Every clean term of
\({\mathfrak E}_{A\to C}\cup
{\mathfrak E}_{C\leftarrow A}\)
closes at every admitted V100 payer location by the V18
three-response five-factor compiler and the legal V15 directed
arborescence.
\end{theorem}

\begin{proof}
Consider \({\mathfrak E}_{A\to C}\).  By definition and
\Cref{lem:V20-boundary-inheritance}, its reduced exact multiplier is
\[
 a_U\,\ell_U\,b_U\,e_U,
\]
where \(a_U\) is the \(D\)-response endpoint, \(b_U\) is the sole
collision covariance, and \(e_U\) is the disjoint endpoint moment.
The word \(\ell_U\) consists of matching-refining moments, at most
four assembled directed occurrences, finite replica separators, and
possibly the breaker in ordinary norm.  Their uniform bounds are the
V99 directed-word norm and the V100 matching/breaker norms.  If the
one payer lies in this interior word, those same theorems give its
single \(Cs_\alpha^{-1}\) bound; a grouped suffix payer is the
separate V96 alternative.

Thus \Cref{thm:V19-exposed-order} verifies the precise V18 ordering,
and the response lower bounds retained by
\Cref{prop:V20-boundary-split} invoke
\Cref{cor:V19-exclusive-closure}.  For
\({\mathfrak E}_{C\leftarrow A}\), first take the adjoint of the
whole word and apply the identical argument.  No cyclic rerooting is
used.
\end{proof}

\begin{corollary}[The exclusive collision has a finite interior-slot
remainder]
\label{cor:V20-interior-remainder}
The only \(QS|D\) terms not closed by V20 are the explicitly tagged
words in \({\mathfrak I}_{QS|D}\).  Every such word has a
non-removable multiplication prefix or suffix outside the proposed
two endpoints after all proved one-sided tensor reductions have been
performed.  The number of resulting slot classes does not grow with
the coordinate support.
\end{corollary}

\begin{proof}
The exposed terms close by
\Cref{thm:V20-exposed-closure}.  A complementary word fails the
boundary definition, hence at least one proposed parent is not an
outer occurrence of the legally reduced skeleton.  The factors
between the two parents have the uniform norms used in the preceding
proof, and the internal collision subfactors commute on disjoint
coordinate banks.  Therefore the retained obstruction is precisely
an outer multiplication prefix or suffix, with its full tag.  The
finite assertion is \Cref{thm:V20-finite-slot-refinement}.
\end{proof}

\begin{firewall}[Uniform norm does not erase an outer multiplier]
\label{fire:V20-outer-factor}
An outer factor in
\({\mathfrak I}_{QS|D}\) is not moved across the product density and
is not absorbed into an endpoint merely because its operator norm is
bounded.  Such an absorption would require a product-state energy
inequality for the dressed endpoint, which has not been proved.  V20
keeps the entire outer word.
\end{firewall}

\section{What the compressed V100 interface does not determine}
\label{sec:V20-compressed-interface}

\begin{proposition}[The aggregate \(B_\mu S_G^\mu\) theorem loses the
boundary flags]
\label{prop:V20-aggregate-loses-boundary}
The archived statements
\[
 K_{\mu,G}=B_\mu S_G^\mu,
 \qquad
 \|B_\mu\|\le C
\]
determine that the complete breaker occurrence is terminal relative
to the aggregated matching word.  They do not determine which
matching directed occurrence is the first or last non-removable
multiplier inside \(B_\mu\), nor the boundary status of an old or
self-role child of such an occurrence.  Hence the aggregate V100
interface alone cannot decide membership in the two exposed classes
of \Cref{def:V20-exposed-classes}.
\end{proposition}

\begin{proof}
The definition of \(B_\mu\) lists matching-refining prefix carriers,
two matching directed covariances, replica regroupings, and bounded
separators, but its displayed norm theorem replaces the assembled
internal word by one number.  Neither that definition nor the bound
records the ordered factorization
\eqref{eq:V20-multiplication-word}.  In particular, two internal
factors may have the same individual norm bounds and occurrence
labels while occupying different sides of a separator.

The terminal position of \(S_G^\mu\) follows from the displayed exact
product \(B_\mu S_G^\mu\).  No analogous displayed factorization is
given inside \(B_\mu\).  Since
\Cref{prop:V19-no-cyclic-reroot} forbids recovering an internal
boundary flag by cyclic trace notation, the missing flags are not a
consequence of the aggregate interface.
\end{proof}

\begin{assessment}[Upstream location after V20]
\label{ass:V20-upstream-location}
The response-to-occurrence part of the V19 slot ledger is now
complete: old, fresh, and response-first self-role branches all have
unique parent slots and retain one payer through every support
telescope.  The exposed portion of \(QS|D\) is closed.  The remaining
problem is smaller and more upstream: expand the finitely many actual
V94/V99 multipliers hidden inside \(B_\mu\), record their two outer
non-tensor factors, and either expose the two parents or prove a
dressed-endpoint inequality for the retained outer word.
\end{assessment}

\section{Regression checks and result ledger}
\label{sec:V20-results}

\begin{proposition}[V20 regression suite]
\label{prop:V20-regressions}
\begin{enumerate}[label=\textup{(\roman*)},leftmargin=4em]
\item No coordinate order is identified with a noncommutative
      multiplication order.
\item Every old, fresh, or self-role response has one parent
      occurrence.
\item Restriction changes an internal tensor factor, not its global
      slot.
\item No slot class is summed coordinate by coordinate after norming.
\item Whole-word adjunction is the only reversal used.
\item An outer bounded multiplier is retained unless an exact
      one-sided tensor factorization removes it.
\item The low-response self-role term is not relabelled as a directed
      response.
\item The \(DQS\), overlap, and remainder collision classes are not
      included in the \(QS|D\) endpoint closure.
\end{enumerate}
\end{proposition}

\begin{proof}
Items \textup{(i)--(iii)} are
\Cref{fire:V20-coordinate-not-slot,
lem:V20-unique-parent,
lem:V20-ancestry-stable,
lem:V20-boundary-inheritance}.  Item \textup{(iv)} is
\Cref{fire:V20-no-slot-count}.  Items \textup{(v)--(vi)} are built
into \Cref{def:V20-reduced-skeleton,fire:V20-outer-factor}.  Item
\textup{(vii)} is explicit in
\Cref{prop:V20-slot-table}.  Item \textup{(viii)} follows from the
exclusive hypothesis in
\Cref{def:V20-exposed-classes}.
\end{proof}

\begin{theorem}[Third-series V20 result ledger]
\label{thm:V20-results}
Relative to V1--V19, V20 proves:
\begin{enumerate}[label=\textup{(V20.\arabic*)},leftmargin=6.3em]
\item the compulsory SM V97/NS V31 audit and its type-correct
      one-occurrence proof-order transfer;
\item a unique ancestry tag for every clean old, fresh, and
      response-first self-role acquisition;
\item invariance of that tag and the sole payer under arbitrary
      finite support telescopes;
\item the complete symbolic slot table for \(D,Q,S\);
\item an exact, support-uniform finite refinement of every collision
      packet by three ancestry tags;
\item inheritance of global boundary status from the parent
      occurrence;
\item closure of both boundary-exposed \(QS|D\) orientations at every
      payer location; and
\item reduction of the unclosed exclusive collision to a finite list
      of fully tagged interior-slot words.
\end{enumerate}
\end{theorem}

\begin{proof}
These are
\Cref{rec:V20-companion-snapshots,
ass:V20-audit-transfer,
lem:V20-unique-parent,
lem:V20-ancestry-stable,
prop:V20-slot-table,
thm:V20-finite-slot-refinement,
lem:V20-boundary-inheritance,
thm:V20-exposed-closure,
cor:V20-interior-remainder}.
\end{proof}

\begin{programme}[V21: expand the interior-slot multipliers]
\label{prog:V20-V21}
\begin{enumerate}[leftmargin=3em]
\item Return to the exact V94 Fourier multiplier for each of the two
      matching directed occurrences inside \(B_\mu\); do not use only
      the aggregate norm \(\|B_\mu\|\le C\).
\item For every tagged word in
      \({\mathfrak I}_{QS|D}\), record the complete prefix before the
      earlier parent and the complete suffix after the later parent.
\item Test each outer factor against the exact V2 positive
      one-sided tensor hypothesis.  Remove it only if the required
      coordinate factorization is present.
\item If an outer factor shares the endpoint row spaces, test a
      dressed quadratic energy
      \[
       \sup_{\rho\ {\rm product}}
       \Tr\!\left(\rho\,(PA)^*(PA)\right)
      \]
      at the response-square scale.  If it fails, retain the smallest
      matrix counterword and return to the pre-dual partition sum.
\item Close newly exposed \(QS|D\) words by
      \Cref{thm:V20-exposed-closure}; keep the \(DQS\), V17 overlap
      and remainder, V2 low-response self-role, and nonclean branches
      separate.
\item Preserve the sole payer and re-run the V15, V83, V84, V91,
      V98, and V100 counterpacket tests.
\end{enumerate}
\end{programme}

\begin{firewall}[Claim boundary after third-series V20]
\label{fire:V20-claim-boundary}
V20 reconstructs response ancestry and closes the boundary-exposed
part of the exclusive \(QS|D\) collision.  It does not prove a
dressed-endpoint estimate for the remaining interior-slot words,
close the \(DQS\), V17 overlap or remainder lines, the V2 low-response
self-role factor, or nonclean representation states.  Higher MTA,
WUFC, CPM, cutoff removal, reflection positivity,
Osterwalder--Schrader reconstruction, construction of the
four-dimensional continuum Yang--Mills measure, and a uniform
positive continuum spectral gap remain open.  The full Yang--Mills
mass gap has not been proved.
\end{firewall}

\paragraph{Parent-version record.}
V20 was copied byte-for-byte from
\texttt{Renewal\_Yang\_Mills\_Mass\_Gap\_Research\_Notes\_V19.tex}
before this chapter was appended.  The V19 parent had \(9\,684\)
source lines, \(353\,984\) bytes, and SHA--256
\[
 \texttt{5690e036f5c6a0706ac77f24d24e0b99639987f2e9a66da5a383fae3af57d5ca}.
\]
The next compulsory companion audit is V25.

% =============================================================================
% BEGIN APPENDED THIRD-SERIES V21 MATERIAL
% =============================================================================

\clearpage
\chapter{V21: Dressed endpoint energies and the matching-prefix
atlas}
\label{ch:V21-dressed-endpoints}

\section{Context: the exact outer-word problem}
\label{sec:V21-context}

V20 reduces the exclusive \(QS|D\) collision to a finite list of
tagged words for which one proposed Hilbert--Schmidt endpoint has a
non-removable multiplication prefix or suffix.  The archived V94
dual theorem explains why the ordinary estimate
\(\|P\|\le L\) does not by itself remove a prefix \(P\): the
product-state density remains outside \(P\), and conditioning through
a row-coupled operator need not preserve a product state.

This chapter writes the exact replacement for the V97 endpoint
sandwich.  The left endpoint is measured by \(P A^2P^*\), and the
right endpoint by \(Q^*E^2Q\).  Tensor-separated and row-product
outer words preserve the old response-square energy.  A
dimension-free norm-only replacement is disproved by an explicit
rank-one example.  Finally the V61 matching-prefix factorization
reduces every genuinely coupled prefix to one or two connected
two-row blocks.

\begin{assessment}[Nonduplication boundary at V21]
\label{ass:V21-nonduplication}
The f2 archive proves the undressed V94/V97 three- and five-factor
sandwiches and the V2 one-sided tensor law.  V18 converts a response
lower bound into an undressed endpoint energy.  V21 proves a different
inequality with arbitrary outer words \(P,Q\), identifies two exact
dressing classes on which the V18 energy survives, and proves a
no-go theorem outside those classes.
\end{assessment}

\section{The two-sided dressed Fourier sandwich}
\label{sec:V21-dressed-sandwich}

\begin{definition}[Left and right dressed quadratic radii]
\label{def:V21-dressed-radii}
Let \(A=A^*\), \(E=E^*\), and let \(P,Q\) be bounded complete
equivariant operators on the four-row Hilbert space.  Put
\begin{align}
 {\cal D}_{\otimes}^{L}(P;A)
 &:=
 \sup_{\rho=\rho_1\otimes\cdots\otimes\rho_4}
 \Tr\!\left(\rho\,P A^2P^*\right),
\label{eq:V21-left-radius}\\
 {\cal D}_{\otimes}^{R}(E;Q)
 &:=
 \sup_{\rho=\rho_1\otimes\cdots\otimes\rho_4}
 \Tr\!\left(\rho\,Q^*E^2Q\right),
\label{eq:V21-right-radius}
\end{align}
where the \(\rho_i\) are density matrices.  Paid radii are defined by
replacing the unique paid factor at its actual location before the
square is formed.
\end{definition}

\begin{theorem}[Two-sided dressed endpoint sandwich]
\label{thm:V21-dressed-sandwich}
Suppose the exact V94 dual multiplier, after its output contraction
has been absorbed into the first separator, has the complete form
\begin{equation}
 {\cal T}_{1,K}(\boldsymbol D)
 =
 P\,A\,C_{1,\boldsymbol D}\,B\,C_2\,E\,Q,
\label{eq:V21-dressed-word}
\end{equation}
where \(A,B,E\) are self-adjoint,
\[
 \sup_{\boldsymbol D}\|C_{1,\boldsymbol D}\|\le L_1,
 \qquad
 \|C_2\|\le L_2.
\]
Then
\begin{equation}
 \boxed{
 \Gamma_\otimes(1,K)
 \le
 L_1L_2\|B\|
 \sqrt{
 {\cal D}_{\otimes}^{L}(P;A)
 {\cal D}_{\otimes}^{R}(E;Q)}.}
\label{eq:V21-dressed-bound}
\end{equation}
The same estimate holds in every once-paid row after the paid factor
is placed inside the corresponding dressed radius, middle norm, or
separator norm.  Exactly one inverse-\(s_\alpha\) debit is used.
\end{theorem}

\begin{proof}
Fix a product density \(\rho\) and a dual family
\(\boldsymbol D\).  Schatten H\"older gives
\begin{align*}
 &\left\|
 \rho^{1/2}P A C_{1,\boldsymbol D}B C_2E Q\rho^{1/2}
 \right\|_1\\
 &\qquad\le
 L_1L_2\|B\|\,
 \|\rho^{1/2}PA\|_2\,
 \|EQ\rho^{1/2}\|_2.
\end{align*}
Cyclicity of the ordinary trace, without moving a multiplier through
\(\rho\), gives
\begin{align*}
 \|\rho^{1/2}PA\|_2^2
 &=\Tr\!\left(A P^*\rho P A\right)
 =\Tr\!\left(\rho P A^2P^*\right),\\
 \|EQ\rho^{1/2}\|_2^2
 &=\Tr\!\left(\rho Q^*E^2Q\right).
\end{align*}
Bound the two factors by
\eqref{eq:V21-left-radius}--%
\eqref{eq:V21-right-radius}, then take the suprema over
\(\rho\) and \(\boldsymbol D\).

If an endpoint or an outer word pays, form the corresponding square
with that paid factor.  If \(B\) or a separator pays, use its paid
norm.  A grouped common-suffix payer is first separated by the f2 V96
suffix theorem.  The V2 allocation makes these alternatives disjoint,
so only one paid substitution occurs.
\end{proof}

\begin{remark}[Recovery of the old endpoint theorem]
\label{rem:V21-undressed-case}
When \(P=Q=I\),
\[
 {\cal D}_{\otimes}^{L}(I;A)
 =\sup_{\rho\ {\rm product}}\Tr(\rho A^2)
 \le\|A\|\kappa_\otimes(|A|)
 ={\cal E}_\otimes(A),
\]
and similarly for \(E\).  Thus
\Cref{thm:V21-dressed-sandwich} recovers the exact analytic core of
the V97/V18 compiler.
\end{remark}

\begin{corollary}[Dressed three-response compiler]
\label{cor:V21-dressed-response}
Assume the three response lower bounds
\[
 d_A\ge\gamma_Aa_*,
 \qquad
 d_B\ge\gamma_Ba_*,
 \qquad
 d_E\ge\gamma_Ea_*,
\]
and suppose
\begin{align}
 {\cal D}_{\otimes}^{L}(P;A)
 &\le
 \frac{C_L}{\gamma_A^2a_*^2}\,d_A^2,
\label{eq:V21-left-response-energy}\\
 {\cal D}_{\otimes}^{R}(E;Q)
 &\le
 \frac{C_R}{\gamma_E^2a_*^2}\,d_E^2.
\label{eq:V21-right-response-energy}
\end{align}
Then
\begin{equation}
 \Gamma_\otimes(1,K)
 \le
 \frac{L_1L_2\|B\|\sqrt{C_LC_R}}
 {\gamma_A\gamma_B\gamma_Ea_*^3}\,
 d_Ad_Bd_E.
\label{eq:V21-dressed-response-bound}
\end{equation}
If the paid versions of
\eqref{eq:V21-left-response-energy} or
\eqref{eq:V21-right-response-energy} contain
\(s_\alpha^{-2}\), their square root gives exactly the once-paid
analogue of \eqref{eq:V21-dressed-response-bound}.
\end{corollary}

\begin{proof}
Insert
\eqref{eq:V21-left-response-energy}--%
\eqref{eq:V21-right-response-energy}
into \eqref{eq:V21-dressed-bound}.  Then multiply by
\[
 1\le\frac{d_B}{\gamma_Ba_*}.
\]
The paid statement is the identical calculation with one paid
alternative from \Cref{thm:V21-dressed-sandwich}.
\end{proof}

\section{Two dressing classes which preserve endpoint energy}
\label{sec:V21-safe-dressing}

\begin{lemma}[Tensor-separated dressing]
\label{lem:V21-tensor-dressing}
For each row let
\(\mathcal H_i=\mathcal C_i\otimes\mathcal S_i\).
Suppose \(A\) acts on \(\bigotimes_i\mathcal C_i\), while \(P\) acts
on \(\bigotimes_i\mathcal S_i\).  Then
\begin{equation}
 {\cal D}_{\otimes}^{L}(I_{\mathcal C}\otimes P;
                       A\otimes I_{\mathcal S})
 =
 \kappa_\otimes(A^2\otimes PP^*)
 \le
 \|P\|^2\kappa_\otimes(A^2).
\label{eq:V21-tensor-left}
\end{equation}
Similarly, if \(E\) and \(Q\) are tensor-separated,
\begin{equation}
 {\cal D}_{\otimes}^{R}(E\otimes I_{\mathcal S};
                       I_{\mathcal C}\otimes Q)
 \le
 \|Q\|^2\kappa_\otimes(E^2).
\label{eq:V21-tensor-right}
\end{equation}
The assertions remain true with one paid tensor factor.
\end{lemma}

\begin{proof}
The squared left endpoint is exactly the positive tensor product
\[
 (I\otimes P)(A^2\otimes I)(I\otimes P^*)
 =A^2\otimes PP^*.
\]
Apply the second orientation of the f2 V2 one-sided tensor law to
obtain
\[
 \kappa_\otimes(A^2\otimes PP^*)
 \le
 \kappa_\otimes(A^2)\|PP^*\|
 =
 \|P\|^2\kappa_\otimes(A^2).
\]
The right calculation uses \(Q^*Q\), and the paid assertion merely
replaces one of these positive factors before applying the same
one-sided inequality.
\end{proof}

\begin{lemma}[Row-product dressing]
\label{lem:V21-row-product-dressing}
If \(P=P_1\otimes\cdots\otimes P_4\), then for every self-adjoint
\(A\),
\begin{equation}
 {\cal D}_{\otimes}^{L}(P;A)
 \le\|P\|^2\kappa_\otimes(A^2).
\label{eq:V21-row-product-left}
\end{equation}
If \(Q=Q_1\otimes\cdots\otimes Q_4\), then
\begin{equation}
 {\cal D}_{\otimes}^{R}(E;Q)
 \le\|Q\|^2\kappa_\otimes(E^2).
\label{eq:V21-row-product-right}
\end{equation}
\end{lemma}

\begin{proof}
For \(\rho=\bigotimes_i\rho_i\), put
\[
 \sigma_i=P_i^*\rho_iP_i,\qquad t_i=\Tr\sigma_i.
\]
If some \(t_i=0\), the left expectation vanishes.  Otherwise
\(\widehat\sigma_i=\sigma_i/t_i\) are density matrices and
\[
 \Tr(\rho P A^2P^*)
 =
 \Tr\!\left[\left(\bigotimes_i\sigma_i\right)A^2\right]
 \le
 \left(\prod_it_i\right)\kappa_\otimes(A^2).
\]
Since \(t_i\le\|P_i\|^2\) and
\(\prod_i\|P_i\|^2=\|P\|^2\), this proves
\eqref{eq:V21-row-product-left}.  The right statement follows by
applying the same argument to \(Q^*E^2Q\), with
\(\sigma_i=Q_i\rho_iQ_i^*\).
\end{proof}

\begin{corollary}[Safe outer words close the V20 interior term]
\label{cor:V21-safe-outer-closure}
Let a tagged word of \({\mathfrak I}_{QS|D}\) have outer prefix
\(P\) and outer suffix \(Q\) after all uniformly bounded middle
factors have been grouped.  If each outer word is either
tensor-separated from its adjacent response endpoint as in
\Cref{lem:V21-tensor-dressing} or is a row product as in
\Cref{lem:V21-row-product-dressing}, then the word closes at every
admitted payer location.
\end{corollary}

\begin{proof}
V18 gives
\[
 \kappa_\otimes(A^2)
 \le {\cal E}_\otimes(A)
 \le
 \frac{C}{\gamma_A^2a_*^2}d_A^2
\]
and the analogous right-endpoint estimate, including the
\(s_\alpha^{-2}\) paid row.  Apply
\Cref{lem:V21-tensor-dressing} or
\Cref{lem:V21-row-product-dressing} to obtain
\eqref{eq:V21-left-response-energy}--%
\eqref{eq:V21-right-response-energy}.  Now use
\Cref{cor:V21-dressed-response} and the V15 arborescence bound.
\end{proof}

\section{Why a norm-only outer word is insufficient}
\label{sec:V21-no-go}

\begin{proposition}[Unbounded product-energy amplification by a
norm-one outer word]
\label{prop:V21-norm-no-go}
For every \(d\ge2\), there are two row spaces
\(\mathbb C^d,\mathbb C^d\), a self-adjoint rank-one projection
\(A_d\), and a unitary \(P_d\) such that
\begin{equation}
 \kappa_\otimes(A_d^2)=\frac1d,
 \qquad
 \|P_d\|=1,
 \qquad
 \kappa_\otimes(P_dA_d^2P_d^*)=1.
\label{eq:V21-amplification}
\end{equation}
Consequently no dimension-independent constant \(C\) can make
\[
 {\cal D}_{\otimes}^{L}(P;A)
 \le C\|P\|^2\kappa_\otimes(A^2)
\]
valid for arbitrary row-coupled \(P\).
\end{proposition}

\begin{proof}
Let
\[
 \Omega_d=d^{-1/2}\sum_{k=1}^de_k\otimes e_k,
 \qquad
 A_d=|\Omega_d\rangle\langle\Omega_d|.
\]
The largest squared overlap of \(\Omega_d\) with a product unit
vector is its largest Schmidt coefficient squared, namely \(1/d\).
Thus \(\kappa_\otimes(A_d)=1/d\), and \(A_d^2=A_d\).

Choose a unitary \(P_d\) sending \(\Omega_d\) to
\(e_1\otimes e_1\).  Then
\[
 P_dA_dP_d^*
 =
 |e_1\otimes e_1\rangle\langle e_1\otimes e_1|,
\]
whose expectation in the corresponding product state is one.
This proves \eqref{eq:V21-amplification}.  The ratio of the last and
first quantities is \(d\), proving the no-go statement.
\end{proof}

\begin{firewall}[The no-go theorem is not a Yang--Mills
counterpacket]
\label{fire:V21-no-go-scope}
\Cref{prop:V21-norm-no-go} disproves an inequality based only on the
operator norm of an arbitrary outer word.  It does not show that a
physical V61 matching-prefix carrier realizes the unitaries \(P_d\).
Such carriers have additional partition, covariance, heat-kernel,
and response structure.  That structure, rather than a bare norm,
must be used next.
\end{firewall}

\section{The exact matching-prefix atlas}
\label{sec:V21-prefix-atlas}

Fix the V100 matching
\[
 \mu=A\mid B,\qquad |A|=|B|=2.
\]
Before the first \(\mu\)-breaking coordinate, every local partition
refines \(\mu\).  Let \(\rho\le\mu\) be the cumulative prefix join.

\begin{proposition}[Four matching-prefix join types]
\label{prop:V21-four-prefix-types}
There are exactly four possible prefix joins:
\begin{enumerate}[label=\textup{(\roman*)},leftmargin=4em]
\item all four rows are singleton;
\item the two rows of \(A\) form one block and the rows of \(B\) are
      singleton;
\item the two rows of \(B\) form one block and the rows of \(A\) are
      singleton;
\item the two blocks are exactly \(A\) and \(B\), so
      \(\rho=\mu\).
\end{enumerate}
For every type, the V61 factorization is
\begin{equation}
 C_{<g}(\rho)
 =
 \bigotimes_{C\in\rho}K_C^{<g,\mathrm{conn}}.
\label{eq:V21-prefix-factorization}
\end{equation}
Type \textup{(i)} is a row product.  Types
\textup{(ii)--(iv)} contain respectively one, one, or two connected
two-row factors and no connected three- or four-row prefix factor.
\end{proposition}

\begin{proof}
A partition below \(\mu\) is obtained by deciding independently
whether the two elements of \(A\) are joined and whether the two
elements of \(B\) are joined.  This gives the four listed
possibilities and no others.  Equation
\eqref{eq:V21-prefix-factorization} is
\Cref{prop:V8-first-connectivity}.  A singleton block contributes its
complete one-row moment, while a joined two-element block contributes
one complete connected two-row carrier.  The final assertions follow
immediately.
\end{proof}

\begin{corollary}[Discrete prefixes and complete coordinate suffixes
are safe dressings]
\label{cor:V21-discrete-prefix-safe}
If an outer V61 prefix has type \textup{(i)} in
\Cref{prop:V21-four-prefix-types}, it obeys the row-product dressing
bound.  A nonpayer complete coordinate suffix which is exactly
tensor-separated from the core obeys the tensor-dressing bound; its
paid alternative is the already grouped V96 suffix row.  Neither
produces an interior-slot obstruction.
\end{corollary}

\begin{proof}
For type \textup{(i)},
\eqref{eq:V21-prefix-factorization} is a tensor product of four
one-row complete moments, so
\Cref{lem:V21-row-product-dressing} applies.  The suffix occupies
coordinate tensor factors disjoint from the retained core, so
\Cref{lem:V21-tensor-dressing} applies before the output norm.  The
paid suffix statement is precisely why that row was kept separate in
the V20 reduced skeleton.
\end{proof}

\begin{theorem}[Pair-prefix reduction of the V20 interior remainder]
\label{thm:V21-pair-prefix-reduction}
After the closures in
\Cref{thm:V20-exposed-closure,cor:V21-safe-outer-closure},
every remaining clean \(QS|D\) word has, on at least one obstructed
side, one of:
\begin{enumerate}[label=\textup{(\alph*)},leftmargin=4em]
\item a connected two-row V61 prefix carrier on the matching
      component \(A\);
\item a connected two-row V61 prefix carrier on the matching
      component \(B\);
\item one carrier of each kind; or
\item a fixed finite replica/regrouping multiplier which shares the
      response endpoint row spaces and has not been proved
      row-product or tensor-separated.
\end{enumerate}
The four cases form an exact finite refinement.  Their number is
independent of the prefix length and coordinate support.
\end{theorem}

\begin{proof}
Use the exact V20 boundary split.  Every outer word already proved
tensor-separated or row-product closes by
\Cref{cor:V21-safe-outer-closure}.  For a remaining V61
matching-refining prefix, apply
\Cref{prop:V21-four-prefix-types}.  The discrete type has just closed;
the other three types are exactly
\textup{(a)--(c)}.  Factors not belonging to the coordinate prefix or
complete suffix are among the universally finite V94 replica
regroupings and local separators.  A factor with an established
middle norm may be grouped between the endpoints; if it remains
outside and shares their row spaces, it is \textup{(d)}.
The partition atlas has four members and the separator list is fixed,
so the refinement is support-uniform.
\end{proof}

\begin{assessment}[The new analytic bottleneck]
\label{ass:V21-new-bottleneck}
The V20 outer-word remainder is no longer arbitrary.  Complete
coordinate suffixes and discrete prefixes are discharged.  A genuine
V61 obstruction is a connected two-row matching-prefix carrier
dressing a response endpoint.  The norm-only estimate is false by
\Cref{prop:V21-norm-no-go}; the next theorem must use that carrier's
exact two-row covariance/moment decomposition and its physical
matching response.
\end{assessment}

\section{Regression checks and result ledger}
\label{sec:V21-results}

\begin{proposition}[V21 regression suite]
\label{prop:V21-regressions}
\begin{enumerate}[label=\textup{(\roman*)},leftmargin=4em]
\item No outer multiplier is moved through a product density.
\item The V94 output contraction is absorbed only into its established
      separator.
\item The left and right dressed radii keep the actual order
      \(P A\) and \(E Q\).
\item A norm-only dressing estimate is not assumed.
\item Tensor separation is used only after an exact coordinate
      factorization.
\item A discrete prefix is distinguished from a connected two-row
      prefix.
\item The suffix payer remains the separate V96 alternative.
\item No matrix no-go example is promoted to a physical Yang--Mills
      counterpacket.
\end{enumerate}
\end{proposition}

\begin{proof}
Items \textup{(i)--(iii)} are built into
\Cref{thm:V21-dressed-sandwich}.  Item \textup{(iv)} is
\Cref{prop:V21-norm-no-go}.  Items \textup{(v)--(vii)} are
\Cref{lem:V21-tensor-dressing,
prop:V21-four-prefix-types,
cor:V21-discrete-prefix-safe}.  Item \textup{(viii)} is
\Cref{fire:V21-no-go-scope}.
\end{proof}

\begin{theorem}[Third-series V21 result ledger]
\label{thm:V21-results}
Relative to V1--V20, V21 proves:
\begin{enumerate}[label=\textup{(V21.\arabic*)},leftmargin=6.3em]
\item the exact two-sided dressed endpoint sandwich
      \eqref{eq:V21-dressed-bound};
\item its three-response unpaid and once-paid compilers;
\item preservation of endpoint energy under tensor-separated
      dressing;
\item preservation under arbitrary row-product dressing;
\item a dimension-growing counterexample to every general
      norm-only dressing bound;
\item the four-type matching-prefix join atlas;
\item closure of discrete-prefix and complete-suffix outer words; and
\item reduction of the remaining exclusive collision to connected
      two-row matching prefixes or a fixed finite coupled separator.
\end{enumerate}
\end{theorem}

\begin{proof}
These are
\Cref{thm:V21-dressed-sandwich,
cor:V21-dressed-response,
lem:V21-tensor-dressing,
lem:V21-row-product-dressing,
prop:V21-norm-no-go,
prop:V21-four-prefix-types,
cor:V21-discrete-prefix-safe,
thm:V21-pair-prefix-reduction}.
\end{proof}

\begin{programme}[V22: connected pair-prefix dressing]
\label{prog:V21-V22}
\begin{enumerate}[leftmargin=3em]
\item Write the exact two-row connected prefix carrier
      \(K_{\{i,j\}}^{<g,\mathrm{conn}}\) as its V77 complete
      covariance before its positive envelope.
\item For a response endpoint \(A\), estimate the actual dressed
      positive operator
      \[
       K_{\{i,j\}}^{<g,\mathrm{conn}}\,
       A^2\,
       \bigl(K_{\{i,j\}}^{<g,\mathrm{conn}}\bigr)^*
      \]
      against product states.  Keep its matching-edge stock and the
      endpoint response in one occurrence ledger.
\item Test whether the pair prefix is disjoint from the endpoint
      subbank.  Close the disjoint part by the V2 one-sided tensor
      law and retain the shared part as one joint two-occurrence word.
\item On the shared part, attempt the exact V94 two-endpoint sandwich
      with the prefix covariance as one endpoint rather than as a
      normed outer factor.
\item Test every proposed estimate on
      \Cref{prop:V21-norm-no-go} and on the archived V2
      entanglement-swapping family.  Any improvement must use physical
      covariance structure absent from those examples.
\item Keep case \textup{(d)} of
      \Cref{thm:V21-pair-prefix-reduction}, the \(DQS\) collision,
      the V17 overlap/remainder, low-response self-role, and nonclean
      states separate.
\end{enumerate}
\end{programme}

\begin{firewall}[Claim boundary after third-series V21]
\label{fire:V21-claim-boundary}
V21 proves the exact dressed endpoint inequality and closes all
V20 outer words which are tensor-separated or row-product.  It does
not prove the needed dressed energy for a connected two-row prefix,
close a row-coupled finite outer separator, the remaining \(QS|D\)
words, \(DQS\), V17 overlap or remainder lines, the V2 low-response
self-role factor, or nonclean states.  Higher MTA, WUFC, CPM, cutoff
removal, reflection positivity, Osterwalder--Schrader reconstruction,
construction of the continuum Yang--Mills measure, and a uniform
positive continuum spectral gap remain open.  The full Yang--Mills
mass gap has not been proved.
\end{firewall}

\paragraph{Parent-version record.}
V21 was copied byte-for-byte from
\texttt{Renewal\_Yang\_Mills\_Mass\_Gap\_Research\_Notes\_V20.tex}
before this chapter was appended.  The V20 parent had \(10\,349\)
source lines, \(381\,046\) bytes, and SHA--256
\[
 \texttt{3e374d83d098433afc244ed49bfddce6ef59f42717041464301e0a395a2fa3c8}.
\]
The next compulsory companion audit remains V25.

% =============================================================================
% BEGIN APPENDED THIRD-SERIES V22 MATERIAL
% =============================================================================

\clearpage
\chapter{V22: Commutator localization of connected pair-prefix
dressing}
\label{ch:V22-commutator}

\section{Context: isolate the cost of row coupling}
\label{sec:V22-context}

V21 shows that a connected two-row matching prefix cannot be removed
from an endpoint by its operator norm alone.  It also identifies the
exact quantity to estimate:
\[
 {\cal D}_{\otimes}^{L}(P;A)
 =
 \kappa_\otimes(PA^2P^*).
\]
The elementary identity
\[
 PA=AP+[P,A]
\]
separates the part controlled by the old endpoint energy from the
part measuring genuine noncommutative row coupling.  V22 carries this
identity through the product-state capacity, then applies the V14
support telescope to prove that the commutator is supported only on
the physical coordinate overlap of the prefix and endpoint banks.

\begin{assessment}[Nonduplication boundary at V22]
\label{ass:V22-nonduplication}
V14 localizes one covariance occurrence across a support subset, and
V21 defines dressed endpoint radii.  V22 does not repeat either
result.  It proves a new positive commutator-square bound for those
radii, an exact two-cell commutator formula for a restricted complete
covariance, and a necessary local target for the remaining pair-prefix
words.
\end{assessment}

\section{A universal commutator dressing inequality}
\label{sec:V22-universal}

\begin{lemma}[Positive two-summand square]
\label{lem:V22-two-summand-square}
For bounded operators \(X,Y\),
\begin{equation}
 (X+Y)(X+Y)^*
 \preccurlyeq
 2XX^*+2YY^*.
\label{eq:V22-two-summand-square}
\end{equation}
The adjoint-oriented analogue
\[
 (X+Y)^*(X+Y)\preccurlyeq2X^*X+2Y^*Y
\]
also holds.
\end{lemma}

\begin{proof}
The difference between the right and left sides of
\eqref{eq:V22-two-summand-square} is
\[
 XX^*+YY^*-XY^*-YX^*
 =(X-Y)(X-Y)^*\succeq0.
\]
Apply the same argument to \(X^*,Y^*\) for the second assertion.
\end{proof}

\begin{theorem}[Commutator control of both dressed radii]
\label{thm:V22-commutator-dressing}
For bounded \(P,Q\) and self-adjoint \(A,E\),
\begin{align}
 {\cal D}_{\otimes}^{L}(P;A)
 &\le
 2\|P\|^2\kappa_\otimes(A^2)
 +2\kappa_\otimes\!\left(
 [P,A][P,A]^*\right),
\label{eq:V22-left-commutator}\\
 {\cal D}_{\otimes}^{R}(E;Q)
 &\le
 2\|Q\|^2\kappa_\otimes(E^2)
 +2\kappa_\otimes\!\left(
 [E,Q]^*[E,Q]\right).
\label{eq:V22-right-commutator}
\end{align}
The inequalities are valid before or after the unique payer is
inserted at its actual location.
\end{theorem}

\begin{proof}
Write
\[
 PA=AP+[P,A].
\]
Since \(PA^2P^*=(PA)(PA)^*\),
\Cref{lem:V22-two-summand-square} gives
\[
 PA^2P^*
 \preccurlyeq
 2APP^*A+2[P,A][P,A]^*
 \preccurlyeq
 2\|P\|^2A^2+2[P,A][P,A]^*.
\]
Take the expectation in an arbitrary product density and then the
supremum to obtain \eqref{eq:V22-left-commutator}.

Similarly,
\[
 EQ=QE+[E,Q],
\qquad
 Q^*E^2Q=(EQ)^*(EQ).
\]
The adjoint-oriented part of
\Cref{lem:V22-two-summand-square} yields
\[
 Q^*E^2Q
 \preccurlyeq
 2EQ^*QE+2[E,Q]^*[E,Q]
 \preccurlyeq
 2\|Q\|^2E^2+2[E,Q]^*[E,Q].
\]
This proves \eqref{eq:V22-right-commutator}.  A paid occurrence is
simply part of \(P,A,Q,\) or \(E\) before the same algebra is applied;
no second payer is introduced.
\end{proof}

\begin{corollary}[Commutator-square response criterion]
\label{cor:V22-commutator-response}
Suppose \(A\) carries
\(d_A\ge\gamma_Aa_*\), \(\|P\|\le L_P\), and
\begin{equation}
 \kappa_\otimes([P,A][P,A]^*)
 \le
 \frac{C_{\rm com}}{\gamma_A^2a_*^2}\,d_A^2.
\label{eq:V22-left-commutator-target}
\end{equation}
Then
\[
 {\cal D}_{\otimes}^{L}(P;A)
 \le
 \frac{C( L_P^2+C_{\rm com})}
 {\gamma_A^2a_*^2}\,d_A^2.
\]
There is a symmetric right-endpoint statement with
\([E,Q]^*[E,Q]\).  If the paid commutator-square bound has one
\(s_\alpha^{-2}\), the V21 dressed compiler produces exactly one
\(s_\alpha^{-1}\).
\end{corollary}

\begin{proof}
By V18,
\[
 \kappa_\otimes(A^2)
 \le {\cal E}_\otimes(A)
 \le
 \frac{C}{\gamma_A^2a_*^2}\,d_A^2.
\]
Insert this and
\eqref{eq:V22-left-commutator-target} into
\eqref{eq:V22-left-commutator}.  The right and paid statements are
identical.
\end{proof}

\begin{corollary}[Commuting outer words are harmless]
\label{cor:V22-commuting-safe}
If \([P,A]=0\), then
\[
 {\cal D}_{\otimes}^{L}(P;A)
 \le
 \|P\|^2\kappa_\otimes(A^2).
\]
If \([E,Q]=0\), the analogous right bound holds.
\end{corollary}

\begin{proof}
When \(P\) commutes with \(A\),
\[
 PA^2P^*=A PP^*A\preccurlyeq\|P\|^2A^2.
\]
The right assertion is the adjoint analogue.
\end{proof}

\section{Exact localization to the shared coordinate bank}
\label{sec:V22-localization}

Let a complete two-row prefix covariance have support \(F\):
\[
 \Delta_F^{\rho,\sigma}
 =
 M_F(\rho)-M_F(\sigma).
\]
Work on one fixed restriction-resolved endpoint term with physical
support \(Z\).  Put
\[
 O=F\cap Z,\qquad R=F\setminus Z,\qquad W=Z\setminus F.
\]
After applying the V16 telescope inside the endpoint if necessary,
write
\[
 A=A_O\otimes I_R\otimes A_W,
 \qquad \|A_W\|\le L_A,
\]
where \(L_A\) is the universal complete moment/covariance norm.

\begin{proposition}[Two-cell commutator identity]
\label{prop:V22-two-cell-commutator}
The exact V14 covariance telescope gives
\begin{align}
 \Delta_F^{\rho,\sigma}
 &=
 \Delta_O^{\rho,\sigma}\otimes M_R(\rho)
 +
 M_O(\sigma)\otimes\Delta_R^{\rho,\sigma},
\label{eq:V22-covariance-split}\\
 [\Delta_F^{\rho,\sigma},A]
 &=
 [\Delta_O^{\rho,\sigma},A_O]\otimes M_R(\rho)
 \otimes A_W
 +
 [M_O(\sigma),A_O]\otimes\Delta_R^{\rho,\sigma}
 \otimes A_W,
\label{eq:V22-commutator-split}
\end{align}
where identity factors on \(W\) are suppressed in the first line.
Thus every nontrivial commutator letter is supported on \(O\);
the \(W\)-factor is retained as a complete endpoint spectator.
\end{proposition}

\begin{proof}
Equation \eqref{eq:V22-covariance-split} is the V14 two-term
restriction formula.  The endpoint acts trivially on \(R\), so
\[
 [X_O\otimes Y_R\otimes I_W,
  A_O\otimes I_R\otimes A_W]
 =
 [X_O,A_O]\otimes Y_R\otimes A_W.
\]
Apply this identity to both terms of
\eqref{eq:V22-covariance-split}.  Factors on
\(W=Z\setminus F\) commute with the prefix because that set is
disjoint from \(F\), and remain in the endpoint occurrence.  This proves
\eqref{eq:V22-commutator-split} and its support statement.
\end{proof}

\begin{corollary}[Quantitative localization of the commutator square]
\label{cor:V22-local-commutator-bound}
With the notation above,
\begin{align}
 &\kappa_\otimes\!\left(
 [\Delta_F,A][\Delta_F,A]^*\right)
\nonumber\\
 &\quad\le
 2L_A^2\,
 \kappa_\otimes\!\left(
 [\Delta_O,A_O][\Delta_O,A_O]^*\right)
 +
 8L_A^2\,
 \kappa_\otimes\!\left(
 [M_O(\sigma),A_O][M_O(\sigma),A_O]^*\right).
\label{eq:V22-local-commutator-bound}
\end{align}
The constants are independent of \(|F|\) and \(|Z|\).
\end{corollary}

\begin{proof}
Call the two terms in
\eqref{eq:V22-commutator-split} \(X\) and \(Y\).
By \Cref{lem:V22-two-summand-square},
\[
 [\Delta_F,A][\Delta_F,A]^*
 \preccurlyeq2XX^*+2YY^*.
\]
The moment \(M_R(\rho)\) is a contraction.  The f2 V2 one-sided
tensor law gives
\[
 \kappa_\otimes(XX^*)
 \le
 L_A^2\kappa_\otimes(
 [\Delta_O,A_O][\Delta_O,A_O]^*).
\]
Likewise \(\|\Delta_R^{\rho,\sigma}\|\le2\), so
\[
 \kappa_\otimes(YY^*)
 \le
 4L_A^2\,
 \kappa_\otimes(
 [M_O(\sigma),A_O][M_O(\sigma),A_O]^*).
\]
Combine the two inequalities.  No coordinatewise sum has been taken.
\end{proof}

\begin{corollary}[Empty overlap closes exactly]
\label{cor:V22-empty-overlap-closes}
If \(F\cap Z=\varnothing\), then
\[
 [\Delta_F^{\rho,\sigma},A]=0,
\]
and the connected pair prefix is a safe dressing of \(A\).
\end{corollary}

\begin{proof}
For \(O=\varnothing\), the two local commutators in
\eqref{eq:V22-commutator-split} vanish.  Apply
\Cref{cor:V22-commuting-safe}.
\end{proof}

\begin{firewall}[A shared support does not imply a large
commutator]
\label{fire:V22-overlap-not-commutator}
The inclusion \(O\ne\varnothing\) only locates the possible
commutator.  It supplies neither a lower nor an upper response-scale
bound for the two local positive squares in
\eqref{eq:V22-local-commutator-bound}.  Conversely, equal physical
support does not forbid exact commutation.  The local operators, not
incidence alone, must be evaluated.
\end{firewall}

\section{The commutator term is genuinely necessary}
\label{sec:V22-necessity}

\begin{proposition}[A self-adjoint norm-one dressing can have an
order-one commutator defect]
\label{prop:V22-commutator-necessary}
For every \(d\ge3\), the example of
\Cref{prop:V21-norm-no-go} may be chosen with \(P_d=P_d^*=P_d^{-1}\).
For that choice,
\begin{equation}
 \kappa_\otimes([P_d,A_d][P_d,A_d]^*)
 \ge\frac12-\frac1d.
\label{eq:V22-commutator-lower}
\end{equation}
Thus the commutator-square term in
\eqref{eq:V22-left-commutator} cannot be discarded under a
self-adjoint norm-one hypothesis.
\end{proposition}

\begin{proof}
The vectors
\(\Omega_d\) and \(e_1\otimes e_1\) have real positive inner
product.  The Householder reflection across their angle bisector is a
self-adjoint unitary \(P_d\) sending the first vector to the second.
Hence the three equalities in
\eqref{eq:V21-amplification} still hold.

Apply \eqref{eq:V22-left-commutator}:
\[
 1
 =
 {\cal D}_{\otimes}^{L}(P_d;A_d)
 \le
 \frac2d+
 2\kappa_\otimes([P_d,A_d][P_d,A_d]^*).
\]
Rearranging gives \eqref{eq:V22-commutator-lower}.
\end{proof}

\begin{remark}[Consistency with the V2 swapping regression]
\label{rem:V22-swapping}
The V2 entanglement-swapping example and
\Cref{prop:V22-commutator-necessary} diagnose the same missing step
in different forms.  Conditioning through a shared row-coupled
factor can alter the class of product states; the commutator square
measures the failure of moving that factor next to the bounded middle
word without changing the endpoint.
\end{remark}

\section{Application to the matching-prefix remainder}
\label{sec:V22-application}

\begin{definition}[The two local pair-prefix defects]
\label{def:V22-local-defects}
For a connected pair prefix \(\Delta_F^{\rho,\sigma}\) dressing the
endpoint \(A_Z\), define
\begin{align}
 {\cal C}_{\Delta}(O;A)
 &:=
 [\Delta_O^{\rho,\sigma},A_O]
 [\Delta_O^{\rho,\sigma},A_O]^*,
\label{eq:V22-cov-defect}\\
 {\cal C}_{M}(O;A)
 &:=
 [M_O(\sigma),A_O]
 [M_O(\sigma),A_O]^*.
\label{eq:V22-moment-defect}
\end{align}
They are kept in the same parent-slot and one-payer ledger as the
original prefix and endpoint occurrences.
\end{definition}

\begin{theorem}[Pair-prefix commutator reduction]
\label{thm:V22-pair-prefix-reduction}
Every connected pair-prefix word left by
\Cref{thm:V21-pair-prefix-reduction} has an exact support-uniform
split into:
\begin{enumerate}[label=\textup{(\roman*)},leftmargin=4em]
\item an empty-overlap part, closed by
      \Cref{cor:V22-empty-overlap-closes}; and
\item a shared part for which the only additional analytic quantities
      required by the V21 dressed compiler are
      \[
       \kappa_\otimes({\cal C}_{\Delta}(O;A)),
       \qquad
       \kappa_\otimes({\cal C}_{M}(O;A)).
      \]
\end{enumerate}
If both quantities on the shared part satisfy the response-square
bound
\begin{equation}
 \kappa_\otimes({\cal C}_{\Delta}(O;A))
 +
 \kappa_\otimes({\cal C}_{M}(O;A))
 \le
 \frac{C}{\gamma_A^2a_*^2}\,d_A^2,
\label{eq:V22-local-target}
\end{equation}
and the analogous paid line has one
\(s_\alpha^{-2}\), then that entire pair-prefix word closes at every
payer location.
\end{theorem}

\begin{proof}
Split the prefix at \(O=F\cap Z\) before any positive envelope by
\Cref{prop:V22-two-cell-commutator}.  The two overlap alternatives
are deterministic and use no coordinate choice.  The empty cell
closes by \Cref{cor:V22-empty-overlap-closes}.  On the shared cell,
\Cref{cor:V22-local-commutator-bound} bounds the full commutator
square by the two displayed local defects.

Under \eqref{eq:V22-local-target},
\Cref{cor:V22-commutator-response} supplies the left dressed
response energy required by
\Cref{cor:V21-dressed-response}.  The right-dressed case is obtained
from the adjoint word.  The payer split of V14 assigns the payer to
exactly one local or exterior tensor factor, so the paid square
contributes one inverse \(s_\alpha\) after the Hilbert--Schmidt square
root.
\end{proof}

\begin{corollary}[At most two local defects per obstructed side]
\label{cor:V22-two-defects}
One connected pair-prefix factor creates at most the covariance
defect \({\cal C}_{\Delta}\) and the moment defect
\({\cal C}_M\).  If both sides of a V21 dressed word contain a
connected pair prefix, the resulting left/right defect table has at
most four entries.  This count is independent of physical support.
\end{corollary}

\begin{proof}
There are exactly two terms in
\eqref{eq:V22-commutator-split}.  Apply the same identity separately
on the two sides.  The Cartesian product has at most four entries.
\end{proof}

\begin{assessment}[Upstream location after V22]
\label{ass:V22-upstream}
The general row-coupled outer-word problem has disappeared from the
connected V61 prefix branch.  It is now the local capacity of two
explicit positive commutator squares on the shared coordinate bank.
The next calculation must open the one-link heat/Wilson moment inside
\(M_O\) and \(\Delta_O\).  Repeating a global covariance norm or
response acquisition cannot estimate these commutators.
\end{assessment}

\section{Regression checks and result ledger}
\label{sec:V22-results}

\begin{proposition}[V22 regression suite]
\label{prop:V22-regressions}
\begin{enumerate}[label=\textup{(\roman*)},leftmargin=4em]
\item The outer word remains on its original side of the endpoint.
\item The product density is never conditioned and then declared
      product.
\item The support telescope has one covariance child in each term.
\item The moment child is not renamed as a covariance.
\item Empty overlap, nonempty overlap, and exact commutation are
      distinguished.
\item The payer is inserted before the commutator square is formed.
\item The no-go example is compatible with self-adjoint outer words.
\item No commutator capacity bound is inferred from incidence alone.
\end{enumerate}
\end{proposition}

\begin{proof}
Items \textup{(i)--(ii)} follow from
\Cref{thm:V22-commutator-dressing}.  Items
\textup{(iii)--(v)} are
\Cref{prop:V22-two-cell-commutator,
cor:V22-empty-overlap-closes,
fire:V22-overlap-not-commutator}.  Item \textup{(vi)} is the paid
statement of
\Cref{thm:V22-commutator-dressing}.  Item \textup{(vii)} is
\Cref{prop:V22-commutator-necessary}; item \textup{(viii)} is the
firewall just cited.
\end{proof}

\begin{theorem}[Third-series V22 result ledger]
\label{thm:V22-results}
Relative to V1--V21, V22 proves:
\begin{enumerate}[label=\textup{(V22.\arabic*)},leftmargin=6.3em]
\item universal left and right commutator bounds for dressed endpoint
      radii;
\item the sufficient response-square commutator criterion;
\item exact localization of a connected pair-prefix commutator to
      the shared coordinate bank;
\item a support-uniform two-term bound by covariance and moment
      commutator squares;
\item closure of every empty-overlap or exactly commuting prefix;
\item a self-adjoint norm-one example proving the commutator term
      necessary;
\item a two-local-defect normal form for every shared pair prefix; and
\item reduction of both-sided prefix dressing to at most four local
      commutator defects.
\end{enumerate}
\end{theorem}

\begin{proof}
These are
\Cref{thm:V22-commutator-dressing,
cor:V22-commutator-response,
prop:V22-two-cell-commutator,
cor:V22-local-commutator-bound,
cor:V22-empty-overlap-closes,
prop:V22-commutator-necessary,
thm:V22-pair-prefix-reduction,
cor:V22-two-defects}.
\end{proof}

\begin{programme}[V23: one-link commutator capacities]
\label{prog:V22-V23}
\begin{enumerate}[leftmargin=3em]
\item Write the exact one-link Fourier multipliers which generate
      \(M_O(\sigma)\), \(\Delta_O^{\rho,\sigma}\), and the local
      response endpoint \(A_O\).
\item Determine which pairs are functions of commuting Casimirs and
      therefore have zero defect before applying any norm.
\item For a noncommuting local pair, derive its heat-semigroup
      commutator Duhamel formula and retain the exact shared
      representation labels.
\item Estimate the positive squares
      \({\cal C}_M\) and \({\cal C}_\Delta\) at the target scale
      \eqref{eq:V22-local-target}, first unpaid and then at each
      one-payer position.
\item Test the proposed estimate on the V21 Householder no-go family,
      the archived V2 swapping example, and the V84 spin-one
      intermediate-output family.
\item Keep the finite coupled-separator row of V21, \(DQS\), V17
      overlap/remainder, the low-response self-role factor, and
      nonclean states separate.
\end{enumerate}
\end{programme}

\begin{firewall}[Claim boundary after third-series V22]
\label{fire:V22-claim-boundary}
V22 reduces connected pair-prefix dressing to two local positive
commutator squares and closes the empty-overlap and exactly commuting
parts.  It does not prove the response-square capacity
\eqref{eq:V22-local-target} on a shared noncommuting bank, close the
finite coupled-separator row, the remaining \(QS|D\) words, \(DQS\),
V17 overlap or remainder lines, the V2 low-response self-role factor,
or nonclean states.  Higher MTA, WUFC, CPM, cutoff removal, reflection
positivity, Osterwalder--Schrader reconstruction, construction of the
continuum Yang--Mills measure, and a uniform positive continuum
spectral gap remain open.  The full Yang--Mills mass gap has not been
proved.
\end{firewall}

\paragraph{Parent-version record.}
V22 was copied byte-for-byte from
\texttt{Renewal\_Yang\_Mills\_Mass\_Gap\_Research\_Notes\_V21.tex}
before this chapter was appended.  The V21 parent had \(10\,948\)
source lines, \(401\,672\) bytes, and SHA--256
\[
 \texttt{6eec149ba6ffee1dcc9b7ecf0bb56cd6d4d2cca327e0b08cbdae87636c7cdba8}.
\]
The next compulsory companion audit remains V25.

% =============================================================================
% BEGIN APPENDED THIRD-SERIES V23 MATERIAL
% =============================================================================

\clearpage
\chapter{V23: Universal bounded dressing and complete exclusive
collision closure}
\label{ch:V23-bounded-dressing}

\section{Context: the commutator estimate is stronger than needed}
\label{sec:V23-context}

V21--V22 ask whether an outer word preserves the undressed
product-state capacity of an endpoint.  The V21 example correctly
shows that the relative estimate
\[
 \kappa_\otimes(PA^2P^*)
 \lesssim
 \|P\|^2\kappa_\otimes(A^2)
\]
is false for arbitrary row-coupled \(P\).  V22 therefore isolates a
commutator-square correction.

The active \(QS|D\) compiler needs less.  Its endpoint already carries
a fixed dominant response
\[
 d_A\ge\gamma_Aa_*.
\]
Exactly as in V18, a universal operator bound can first be proved and
then converted to response-square scale by multiplying by a scalar
known to be at least one.  Since every complete V100 outer word has a
support-uniform ordinary norm, and a support-uniform
\(s_\alpha^{-1}\) norm when it pays, this elementary route handles
arbitrary noncommuting dressing.  No commutator estimate is needed for
the exclusive collision.

\begin{assessment}[Upstream correction to the V21--V22 mandate]
\label{ass:V23-upstream-correction}
The V21 no-go theorem remains a valid sharpness result for relative
product capacity, and the V22 commutator localization remains a valid
finer identity.  Neither is the minimal hypothesis of the active
three-response compiler.  V23 replaces the unnecessary relative
estimate by an absolute dressed norm bound followed by the already
retained response lower bound.
\end{assessment}

\section{Absolute dressed response energies}
\label{sec:V23-absolute-dressing}

\begin{lemma}[Universal bounded left dressing]
\label{lem:V23-left-dressing}
Let \(A=A^*\), let \(P\) be arbitrary, and suppose
\[
 \|A\|\le L_A,\qquad
 \|P\|\le L_P,\qquad
 d_A\ge\gamma_Aa_*.
\]
Then
\begin{equation}
 \boxed{
 {\cal D}_{\otimes}^{L}(P;A)
 \le
 \frac{L_P^2L_A^2}{\gamma_A^2a_*^2}\,d_A^2.}
\label{eq:V23-left-dressing}
\end{equation}
If the unique payer lies in \(P\) and
\(\|P^{\rm pay}\|\le L_{P,\rm p}/s_\alpha\), then
\begin{equation}
 {\cal D}_{\otimes}^{L}(P^{\rm pay};A)
 \le
 \frac{L_{P,\rm p}^2L_A^2}
 {s_\alpha^2\gamma_A^2a_*^2}\,d_A^2.
\label{eq:V23-left-prefix-paid}
\end{equation}
If it lies in \(A\) and
\(\|A^{\rm pay}\|\le L_{A,\rm p}/s_\alpha\), then
\begin{equation}
 {\cal D}_{\otimes}^{L}(P;A^{\rm pay})
 \le
 \frac{L_P^2L_{A,\rm p}^2}
 {s_\alpha^2\gamma_A^2a_*^2}\,d_A^2.
\label{eq:V23-left-endpoint-paid}
\end{equation}
\end{lemma}

\begin{proof}
For every density matrix, product or otherwise,
\[
 \Tr(\rho PA^2P^*)
 \le
 \|PA^2P^*\|
 \le L_P^2L_A^2.
\]
The response lower bound gives
\[
 1\le\frac{d_A^2}{\gamma_A^2a_*^2}.
\]
Multiplication proves \eqref{eq:V23-left-dressing}.  Replacing the
norm of the unique paid factor before the same calculation proves
\eqref{eq:V23-left-prefix-paid} and
\eqref{eq:V23-left-endpoint-paid}.
\end{proof}

\begin{lemma}[Universal bounded right dressing]
\label{lem:V23-right-dressing}
Let \(E=E^*\), let \(Q\) be arbitrary, and suppose
\[
 \|E\|\le L_E,\qquad
 \|Q\|\le L_Q,\qquad
 d_E\ge\gamma_Ea_*.
\]
Then
\begin{equation}
 \boxed{
 {\cal D}_{\otimes}^{R}(E;Q)
 \le
 \frac{L_Q^2L_E^2}{\gamma_E^2a_*^2}\,d_E^2.}
\label{eq:V23-right-dressing}
\end{equation}
Payer in \(Q\) or in \(E\) gives the corresponding bound with exactly
one factor \(s_\alpha^{-2}\).
\end{lemma}

\begin{proof}
Use
\[
 \Tr(\rho Q^*E^2Q)
 \le\|Q^*E^2Q\|
 \le L_Q^2L_E^2
\]
and the inequality
\(1\le d_E^2/(\gamma_E^2a_*^2)\).
The paid calculation is identical.
\end{proof}

\begin{firewall}[What the absolute bound does not claim]
\label{fire:V23-not-relative}
\Cref{lem:V23-left-dressing,lem:V23-right-dressing} do not bound a
dressed product capacity by the undressed product capacity.  They use
the full operator norm and then the physical response lower bound.
Therefore they do not contradict
\Cref{prop:V21-norm-no-go} and do not condition a product density
through an entangling outer word.
\end{firewall}

\begin{corollary}[Arbitrary uniformly bounded outer words are
response-safe]
\label{cor:V23-arbitrary-outer-safe}
In the V21 dressed word
\[
 P\,A\,C_1\,B\,C_2\,E\,Q,
\]
assume \(A,E\) carry their retained dominant responses, \(P,Q\) have
support-uniform ordinary norms, and every once-paid alternative has
its established \(Cs_\alpha^{-1}\) assembled norm.  Then the dressed
endpoint hypotheses
\eqref{eq:V21-left-response-energy}--%
\eqref{eq:V21-right-response-energy}
hold without tensor separation, row-product structure, or a
commutator estimate.
\end{corollary}

\begin{proof}
Apply
\Cref{lem:V23-left-dressing,lem:V23-right-dressing}.  Their paid
lines give the exact squared scale required before the
Hilbert--Schmidt square root.
\end{proof}

\section{Every exclusive collision has a dressed five-factor order}
\label{sec:V23-order}

\begin{lemma}[Dressed normal form for a \(QS|D\) line]
\label{lem:V23-dressed-normal-form}
Fix one ancestry-tag class of an exclusive \(QS|D\) telescope line.
After at most one whole-word adjoint reversal, and after commuting
only disjoint tensor subfactors inside their own parent occurrence,
its exact V94 dual multiplier has the form
\begin{equation}
 P\,A\,C_{1,\boldsymbol D}\,B\,C_2\,E\,Q.
\label{eq:V23-exclusive-word}
\end{equation}
Here:
\begin{enumerate}[label=\textup{(\roman*)},leftmargin=4em]
\item \(A\) is the response-carrying \(D\)-parent subfactor;
\item \(B\) is the sole covariance child in the \(QS\)-collision
      occurrence;
\item \(E\) is its response-carrying endpoint-moment child;
\item \(P,Q\) are the complete multiplication prefix and suffix
      outside the two parent occurrences; and
\item \(C_1,C_2\) contain all intervening parent remainders,
      complete V61 words, output dualizer, and finite separators.
\end{enumerate}
All factors retain their V20 ancestry and payer tags.
\end{lemma}

\begin{proof}
The exclusive occurrence identity is
\[
 j(Q)=j(S)\ne j(D)
\]
by \Cref{cor:V20-collision-slots}.  Hence the two parent occurrences
have a definite linear order in the exact word.  If the \(D\)-parent
comes second, take the adjoint of the entire word using the V19
identity.  This reverses all factors without cyclically moving one
through the product density.

Inside the \(D\)-parent, move only its disjoint coordinate-tensor
spectators to the side of \(A\) belonging to \(C_1\) or \(P\).
Inside the collision parent, the covariance child and endpoint moment
act on disjoint cells of the V17 telescope and commute by
\Cref{lem:V19-internal-commutation}; put the covariance before the
moment.  The factors strictly outside the two parents form \(P,Q\);
the factors between them form \(C_1,C_2\).  V20 ancestry stability
shows that none of these operations changes a parent slot, telescope
side, or payer flag.
\end{proof}

\begin{lemma}[Uniform assembled norms in the dressed normal form]
\label{lem:V23-assembled-norms}
In \eqref{eq:V23-exclusive-word}, the unpaid factors satisfy
\[
 \|P\|,\|Q\|,\|C_1\|,\|C_2\|,\|B\|\le C_*.
\]
If the unique payer lies in any one of them, its complete assembled
paid norm is at most \(C_*/s_\alpha\).  If it lies in \(A\) or \(E\),
their paid complete-occurrence norms have the same scale.  The
constant \(C_*\) is independent of coordinate support,
representations, and physical multiplicity.
\end{lemma}

\begin{proof}
The exact word contains at most four complete directed covariances,
the complete matching breaker, complete matching-refining moments,
finite V61/V94 separators, and the common suffix.  Complete moments
are contractions.  The archived V99 directed-word theorem bounds any
assembled subword of directed covariances and finite separators.
The V100 breaker has a uniform ordinary norm, and its paid minimum is
at most \(C/s_\alpha\).  Local separator payment is the archived V95
row.  The complete common suffix is contracted when unpaid and is the
separate grouped V96 row when paid.

Multiplying only the fixed number of these bounds gives \(C_*\).
The V2 one-payer allocation makes the paid alternatives disjoint, so
no product contains two inverse-\(s_\alpha\) factors.  Complete moment
and covariance endpoints have the V18 ordinary and paid caps.
\end{proof}

\section{Closure of the full exclusive \(QS|D\) class}
\label{sec:V23-exclusive-closure}

\begin{theorem}[Complete exclusive-collision closure]
\label{thm:V23-complete-QS-D}
Every clean, termwise response-admissible \(QS|D\) line in the V17
four-term covariance telescope satisfies the legal unpaid and
once-paid V18 three-response marked-tree bounds.  This includes both
V20 boundary-exposed orientations and every word formerly retained in
\({\mathfrak I}_{QS|D}\).
\end{theorem}

\begin{proof}
Fix one exact ancestry-tag class.  Use
\Cref{lem:V23-dressed-normal-form} and
\Cref{lem:V23-assembled-norms}.  The \(D\)-endpoint \(A\) and
collision moment \(E\) retain the response lower bounds proved
termwise in V17--V18.  Therefore
\Cref{lem:V23-left-dressing,lem:V23-right-dressing}
give the two dressed response-square estimates.

The collision covariance \(B\) has a uniform ordinary norm and keeps
its third response lower bound.  Apply
\Cref{cor:V21-dressed-response}.  If the payer is in \(P,A\), use the
paid left radius; if it is in \(E,Q\), use the paid right radius; if
it is in \(C_1,B,C_2\), use the corresponding single paid middle
norm.  The grouped-suffix alternative is the V96 row.  In every case
the Hilbert--Schmidt square root produces exactly one
\(s_\alpha^{-1}\).

Finally use the response assignment of
\Cref{prop:V18-Q-line,prop:V18-S-line} and the corrected V15
six-arborescence table.  The three retained scalar responses mark the
legal tree, while \(P,Q,C_1,C_2\) are bounded only as parts of the
same complete word.  Summing the universally finite tag classes
preserves the stated scale.
\end{proof}

\begin{corollary}[The V20 interior-slot remainder is discharged]
\label{cor:V23-V20-discharged}
The packet
\({\cal P}_{{\mathfrak I}_{QS|D}}\) in
\eqref{eq:V20-boundary-split} is controlled by the same support-uniform
unpaid and once-paid marked-tree estimate as its two exposed
companions.  Endpoint boundary status is no longer an analytic gate
for the exclusive collision.
\end{corollary}

\begin{proof}
Apply \Cref{thm:V23-complete-QS-D} to every interior tag class and use
the finite sum in
\Cref{thm:V20-finite-slot-refinement}.
\end{proof}

\begin{corollary}[Pair-prefix commutator capacities are not needed
for \(QS|D\)]
\label{cor:V23-commutator-not-needed}
The target \eqref{eq:V22-local-target} may be omitted from the proof
of the exclusive \(QS|D\) line.  The V22 defect identities remain
available for sharper estimates or other branches, but their failure
cannot obstruct \Cref{thm:V23-complete-QS-D}.
\end{corollary}

\begin{proof}
The proof of
\Cref{thm:V23-complete-QS-D} uses the absolute dressed bounds
\eqref{eq:V23-left-dressing}--%
\eqref{eq:V23-right-dressing} and never invokes a commutator
capacity.
\end{proof}

\begin{firewall}[Why an outer word is not counted as another mark]
\label{fire:V23-outer-no-mark}
The outer factors \(P,Q\) enter only through their assembled ordinary
or once-paid norms.  V23 does not attach their coordinate stocks to
the marked tree, does not infer a response from their norm, and does
not create a fourth scalar mark.  The only response conversion in the
dressed radii uses the already selected endpoint response.
\end{firewall}

\section{Updated collision ledger}
\label{sec:V23-collision-ledger}

\begin{proposition}[Status of the five occurrence partitions]
\label{prop:V23-five-partition-status}
For the five V16 occurrence partitions:
\begin{enumerate}[label=\textup{(\roman*)},leftmargin=4em]
\item \(D|Q|S\) is covered by the corrected V15 all-directed
      compiler whenever the three marks are jointly admissible;
\item \(DQ|S\) and \(DS|Q\) retain the V13--V14
      reverse-or-deplete and self-role fallback, without a new
      operator-order obstruction;
\item \(QS|D\) is closed analytically by
      \Cref{thm:V23-complete-QS-D}; and
\item \(DQS\) remains the genuine one-unrepeated-occurrence
      multi-mark collision.
\end{enumerate}
The remaining overlap and remainder lines of the V17 support
telescope are kept separate from this occurrence-partition ledger.
\end{proposition}

\begin{proof}
Items \textup{(i)--(ii)} are the V15--V16 classification.
Item \textup{(iii)} is the theorem just proved.  In \(DQS\), all
three desired roles have one parent occurrence by
\Cref{cor:V20-collision-slots}; the one-difference telescope cannot
turn that parent into three independent physical occurrences.
The last sentence is the V17 four-cell decomposition.
\end{proof}

\begin{assessment}[Bottleneck after V23]
\label{ass:V23-bottleneck}
The noncommutative endpoint-order problem for the exclusive collision
is complete.  The first remaining collision-specific gate is no
longer analytic dressing; it is the marking theorem for three desired
edge roles carried by one unrepeated physical occurrence.  That is
the \(DQS\) class.  The next attack must return to the exact local
partition-state sum of that occurrence before any three scalar marks
are assigned.
\end{assessment}

\section{Regression checks and result ledger}
\label{sec:V23-results}

\begin{proposition}[V23 regression suite]
\label{prop:V23-regressions}
\begin{enumerate}[label=\textup{(\roman*)},leftmargin=4em]
\item The V21 dimension-\(d\) no-go is not contradicted.
\item No cyclic rerooting through a product density is used.
\item The outer prefix and suffix remain inside the exact multiplier.
\item Response conversion occurs only after a universal dressed norm
      bound.
\item Every payer alternative produces one inverse
      \(s_\alpha\), never two.
\item No outer factor is used as an additional scalar edge mark.
\item Only the exclusive \(QS|D\) telescope lines are closed.
\item The one-occurrence \(DQS\) marking problem remains explicit.
\end{enumerate}
\end{proposition}

\begin{proof}
Item \textup{(i)} is
\Cref{fire:V23-not-relative}.  Items \textup{(ii)--(iii)} are
\Cref{lem:V23-dressed-normal-form}.  Item \textup{(iv)} is
\Cref{lem:V23-left-dressing,lem:V23-right-dressing}.  Item
\textup{(v)} is
\Cref{lem:V23-assembled-norms,thm:V23-complete-QS-D}.
Item \textup{(vi)} is \Cref{fire:V23-outer-no-mark}.  Items
\textup{(vii)--(viii)} are
\Cref{prop:V23-five-partition-status}.
\end{proof}

\begin{theorem}[Third-series V23 result ledger]
\label{thm:V23-results}
Relative to V1--V22, V23 proves:
\begin{enumerate}[label=\textup{(V23.\arabic*)},leftmargin=6.3em]
\item universal unpaid and once-paid response-square bounds for
      arbitrary bounded left and right endpoint dressings;
\item an exact dressed normal form for every exclusive collision tag
      class;
\item support-uniform ordinary and once-paid bounds for every outer
      word in that normal form;
\item closure of every clean, termwise admissible \(QS|D\) line;
\item discharge of the complete V20 interior-slot remainder;
\item removal of the V22 commutator target from this branch; and
\item identification of \(DQS\) as the first remaining
      collision-specific one-occurrence marking gate.
\end{enumerate}
\end{theorem}

\begin{proof}
These are
\Cref{lem:V23-left-dressing,
lem:V23-right-dressing,
lem:V23-dressed-normal-form,
lem:V23-assembled-norms,
thm:V23-complete-QS-D,
cor:V23-V20-discharged,
cor:V23-commutator-not-needed,
prop:V23-five-partition-status}.
\end{proof}

\begin{programme}[V24: the \(DQS\) one-occurrence marking theorem]
\label{prog:V23-V24}
\begin{enumerate}[leftmargin=3em]
\item Fix one \(DQS\) parent occurrence and retain its exact V61 local
      partition-state operator before the V14/V16 support telescope.
\item Partition its coordinates simultaneously by the three actual
      directed response labels, including all pairwise and triple
      overlaps.  Keep this as one occurrence partition, not three
      copied banks.
\item Determine which cells admit the archived V98
      same-bank joint-marking theorem and which cells contain only one
      unrepeated covariance factor.
\item On a genuine triple-overlap cell, seek one complete
      three-response Fourier estimate before expanding scalar marks.
\item On pairwise-exclusive cells, use the exact finite
      one-difference telescope and test whether two response roles can
      be placed in endpoint moments while the sole covariance remains
      in the middle.
\item Retain every payer side and test the V15, V83, V84, V91, V98,
      and V21 no-go examples.
\item Keep V17 overlap/remainder lines, V2 low-response self-role, and
      nonclean states separate.  The next compulsory companion audit
      is V25.
\end{enumerate}
\end{programme}

\begin{firewall}[Claim boundary after third-series V23]
\label{fire:V23-claim-boundary}
V23 closes the complete clean, termwise response-admissible exclusive
\(QS|D\) collision class.  It does not prove joint marking of the
one-occurrence \(DQS\) class, close the V17 overlap or remainder lines,
all V13--V14 self-role fallbacks, the V2 low-response self-role factor,
or nonclean states.  Higher MTA, WUFC, CPM, cutoff removal, reflection
positivity, Osterwalder--Schrader reconstruction, construction of the
continuum Yang--Mills measure, and a uniform positive continuum
spectral gap remain open.  The full Yang--Mills mass gap has not been
proved.
\end{firewall}

\paragraph{Parent-version record.}
V23 was copied byte-for-byte from
\texttt{Renewal\_Yang\_Mills\_Mass\_Gap\_Research\_Notes\_V22.tex}
before this chapter was appended.  The V22 parent had \(11\,525\)
source lines, \(419\,747\) bytes, and SHA--256
\[
 \texttt{76ffac99104f6efb050db85861316cfdd0c4098fd532b4e00f3a9727f0f4ea05}.
\]
The next compulsory companion audit remains V25.

% =============================================================================
% BEGIN APPENDED THIRD-SERIES V24 MATERIAL
% =============================================================================

\clearpage
\chapter{V24: Absolute attachment and the one-occurrence triple mark}
\label{ch:V24-absolute-triple}

\section{Context: separate combinatorial repetition from analytic use}
\label{sec:V24-context}

V23 showed that a dominant response may be attached to a uniformly
bounded dressed endpoint after the complete operator estimate has
already been proved.  The remaining \(DQS\) class asks for three
directed responses whose ancestry tags point to one unrepeated
physical covariance occurrence.  The archived V92 marking firewall
forbids treating that occurrence as three independent first-moment
cards.  It does not forbid the following proof order:
\begin{enumerate}[label=\textup{(\roman*)},leftmargin=4em]
\item assemble and bound the complete \(DQS\) word once;
\item use the three already proved scalar response lower bounds only
      after that operator bound;
\item expand their product by the exact bank-to-mark identity.
\end{enumerate}
This chapter proves that this order is sufficient.  It also
disintegrates the scalar product into its five coordinate-equality
strata, so that a repeated coordinate is visible and is never hidden
as an independent analytic use.

\begin{assessment}[Upstream correction to the V24 programme]
\label{ass:V24-upstream-correction}
The proposed new three-response Fourier estimate on a triple-overlap
cell is stronger than the clean, termwise response-admissible \(DQS\)
compiler requires.  A support-uniform absolute norm of the complete
single occurrence, followed by aggregate response attachment, closes
that compiler without splitting or cloning the occurrence.
\end{assessment}

\section{The seven signature atoms}
\label{sec:V24-signatures}

Retain one row of the V15 table.  Write
\[
 {\cal R}:=\{D,Q,S\},\qquad
 g_R=t_Rh_R\quad(R\in{\cal R}),
\]
where \(t_D=x,t_Q=q,t_S=y\), and where the three undirected labels
\(g_D,g_Q,g_S\) are distinct and form the displayed tree rooted at
\(z\).  Let \(C\) be the complete coordinate bank of the unique
\(DQS\) parent occurrence.  For \(e\in C\), put
\begin{equation}
 b_R(e):=a_{t_R,e}a_{h_R,e},
 \qquad
 E_R:=\{e\in C:b_R(e)>0\}.
\label{eq:V24-coordinate-stock}
\end{equation}
For \(B\subseteq C\), define
\begin{equation}
 H_R(B):=\sum_{e\in B}b_R(e),
 \qquad
 d_R(B):=\frac{H_R(B)}{\Xi_{t_R}}.
\label{eq:V24-restricted-response}
\end{equation}

\begin{definition}[Signature atoms]
\label{def:V24-signature-atoms}
For \(J\subseteq{\cal R}\), let
\begin{equation}
 C_J:=
 \left(\bigcap_{R\in J}E_R\right)
 \cap
 \left(\bigcap_{R\in{\cal R}\setminus J}(C\setminus E_R)\right).
\label{eq:V24-signature-atom}
\end{equation}
Thus \(C_J\) consists of coordinates which support exactly the edge
roles in \(J\).  The seven cells with
\(\varnothing\ne J\subseteq{\cal R}\) are the active signature
atoms; \(C_\varnothing\) is inactive for all three desired edges.
\end{definition}

\begin{lemma}[Exact signature disintegration]
\label{lem:V24-signature-disintegration}
The eight sets \(C_J\), \(J\subseteq{\cal R}\), are pairwise disjoint
and partition \(C\).  For every \(R\in{\cal R}\),
\begin{equation}
 d_R(C)
 =
 \sum_{\substack{J\subseteq{\cal R}\\R\in J}}d_R(C_J).
\label{eq:V24-response-four-cells}
\end{equation}
The sum has exactly four terms, some of which may vanish.
\end{lemma}

\begin{proof}
Every coordinate has a unique membership signature
\(\{R:e\in E_R\}\subseteq{\cal R}\), proving the partition claim.
Moreover
\[
 E_R=\bigsqcup_{J\ni R}C_J.
\]
The edge stock is a finite nonnegative coordinate sum, so it is
additive on this disjoint union.  Division by the fixed positive
\(\Xi_{t_R}\) proves \eqref{eq:V24-response-four-cells}.  There are
\(2^2=4\) subsets of \({\cal R}\) containing a fixed role \(R\).
\end{proof}

\begin{corollary}[Deterministic heavy signature cell]
\label{cor:V24-heavy-signature}
Fix once and for all a total order on the seven nonempty subsets of
\({\cal R}\).  For each \(R\), let \(J_R^\star\) be the first
maximizer of \(d_R(C_J)\) among \(J\ni R\), and put
\(B_R:=C_{J_R^\star}\).  Then
\begin{equation}
 d_R(B_R)\ge\frac14d_R(C).
\label{eq:V24-heavy-one}
\end{equation}
Consequently, if
\(d_R(C)\ge\gamma_Ra_*\) for all three roles, then
\begin{equation}
 \prod_{R\in{\cal R}}d_R(B_R)
 \ge
 \frac{\gamma_D\gamma_Q\gamma_S}{64}\,a_*^3.
\label{eq:V24-heavy-product}
\end{equation}
\end{corollary}

\begin{proof}
The four nonnegative terms in
\eqref{eq:V24-response-four-cells} have maximum at least one quarter
of their sum.  Multiply the resulting three inequalities and insert
the response lower bounds.
\end{proof}

\begin{proposition}[Five selected-cell incidence types]
\label{prop:V24-five-cell-types}
Because signature atoms are disjoint, the equality relation among
\((B_D,B_Q,B_S)\) is exactly one of
\begin{equation}
 D|Q|S,\qquad DQ|S,\qquad DS|Q,\qquad QS|D,\qquad DQS.
\label{eq:V24-five-cell-types}
\end{equation}
If two roles share a selected atom, its signature contains both
roles.  If all three share a selected atom, that atom is necessarily
\(C_{\{D,Q,S\}}\).
\end{proposition}

\begin{proof}
Equality of three objects has the five Bell-partition types in
\eqref{eq:V24-five-cell-types}.  If \(B_R=B_{R'}=C_J\), then the
definitions of \(J_R^\star\) and \(J_{R'}^\star\) require
\(R,R'\in J\).  If the same holds for all three roles, then
\({\cal R}\subseteq J\subseteq{\cal R}\), hence \(J={\cal R}\).
\end{proof}

\begin{firewall}[A signature atom is not a new occurrence]
\label{fire:V24-atom-not-occurrence}
The sets \(C_J\) are a coordinate partition of the one \(DQS\)
parent.  They are not independent covariance occurrences.  Applying
the V16 restriction telescope to them produces one covariance
difference per summand, not seven differences.
\end{firewall}

\section{Exact coordinate-equality strata}
\label{sec:V24-diagonals}

Let \(B_D,B_Q,B_S\subseteq C\) be arbitrary edge banks, not
necessarily the heavy banks above.  Put
\[
 \Omega:=B_D\times B_Q\times B_S
\]
and, for \(\mathbf e=(e_D,e_Q,e_S)\in\Omega\), set
\begin{equation}
 \omega(\mathbf e):=
 \frac{b_D(e_D)b_Q(e_Q)b_S(e_S)}
 {\Xi_x\Xi_q\Xi_y}.
\label{eq:V24-triple-coordinate-weight}
\end{equation}
Coordinates for which one displayed \(b_R\) vanishes contribute
zero and may be retained harmlessly.

\begin{definition}[The five diagonal strata]
\label{def:V24-diagonal-strata}
Partition \(\Omega\) into
\begin{align}
 \Delta_{D|Q|S}
 &:=
 \{\mathbf e:e_D,e_Q,e_S\text{ are pairwise distinct}\},
\notag\\
 \Delta_{DQ|S}
 &:=
 \{\mathbf e:e_D=e_Q\ne e_S\},
\notag\\
 \Delta_{DS|Q}
 &:=
 \{\mathbf e:e_D=e_S\ne e_Q\},
\notag\\
 \Delta_{QS|D}
 &:=
 \{\mathbf e:e_Q=e_S\ne e_D\},
\notag\\
 \Delta_{DQS}
 &:=
 \{\mathbf e:e_D=e_Q=e_S\}.
\label{eq:V24-five-diagonals}
\end{align}
For a type \(\pi\) in this list, define
\begin{equation}
 {\mathfrak M}_\pi(B_D,B_Q,B_S)
 :=
 \sum_{\mathbf e\in\Delta_\pi}\omega(\mathbf e).
\label{eq:V24-diagonal-mass}
\end{equation}
\end{definition}

\begin{theorem}[Exact five-stratum mark identity]
\label{thm:V24-five-stratum-identity}
For arbitrary finite edge banks,
\begin{equation}
 \boxed{
 d_D(B_D)d_Q(B_Q)d_S(B_S)
 =
 \sum_{\pi\in{\mathfrak P}_3}
 {\mathfrak M}_\pi(B_D,B_Q,B_S),}
\label{eq:V24-five-stratum-identity}
\end{equation}
where
\[
 {\mathfrak P}_3
 =
 \{D|Q|S,DQ|S,DS|Q,QS|D,DQS\}.
\]
Every term is nonnegative, the decomposition has no support-size
factor, and each coordinate triple occurs exactly once.
\end{theorem}

\begin{proof}
Expanding the three finite sums gives
\[
 d_D(B_D)d_Q(B_Q)d_S(B_S)
 =
 \sum_{(e_D,e_Q,e_S)\in\Omega}
 \frac{b_D(e_D)b_Q(e_Q)b_S(e_S)}
 {\Xi_x\Xi_q\Xi_y}.
\]
Equality of the three coordinate values has exactly the five
patterns in \eqref{eq:V24-five-diagonals}.  Those sets are pairwise
disjoint and exhaust \(\Omega\).  Splitting the finite sum proves
\eqref{eq:V24-five-stratum-identity}.  No estimate or
pigeonhole argument is used, so no multiplicity factor appears.
\end{proof}

\begin{corollary}[Compatibility of atom and coordinate collisions]
\label{cor:V24-atom-coordinate-compatibility}
For the heavy signature banks of
\Cref{cor:V24-heavy-signature}:
\begin{enumerate}[label=\textup{(\roman*)},leftmargin=4em]
\item if \(B_D,B_Q,B_S\) are pairwise distinct, only
      \({\mathfrak M}_{D|Q|S}\) can be nonzero;
\item if exactly \(B_R=B_{R'}\ne B_T\), only
      \({\mathfrak M}_{D|Q|S}\) and
      \({\mathfrak M}_{RR'|T}\) can be nonzero; and
\item if \(B_D=B_Q=B_S=C_{\{D,Q,S\}}\), all five coordinate
      strata are possible.
\end{enumerate}
\end{corollary}

\begin{proof}
Distinct signature atoms are disjoint.  Thus coordinates selected
from distinct atoms cannot be equal.  This excludes every equality
which crosses two distinct selected atoms and proves
\textup{(i)--(ii)}.  In \textup{(iii)} all three coordinate choices
belong to the same atom, so none of the five equality patterns is
excluded set-theoretically.
\end{proof}

\begin{remark}[Explicit diagonal formulas]
\label{rem:V24-diagonal-formulas}
For example,
\begin{align*}
 {\mathfrak M}_{DQ|S}
 &=
 \frac1{\Xi_x\Xi_q\Xi_y}
 \sum_{\substack{e\in B_D\cap B_Q\\f\in B_S,\ f\ne e}}
 b_D(e)b_Q(e)b_S(f),\\
 {\mathfrak M}_{DQS}
 &=
 \frac1{\Xi_x\Xi_q\Xi_y}
 \sum_{e\in B_D\cap B_Q\cap B_S}
 b_D(e)b_Q(e)b_S(e).
\end{align*}
Thus pair and triple coordinate reuse is displayed as a joint local
weight; it is not silently replaced by two or three analytic
covariance estimates.
\end{remark}

\section{Absolute attachment of several responses}
\label{sec:V24-attachment}

\begin{lemma}[Absolute aggregate-response attachment]
\label{lem:V24-absolute-attachment}
Let \(K\) be a complete assembled multiplier and let
\(d_1,\ldots,d_m\) be nonnegative responses satisfying
\[
 d_i\ge\gamma_i a_*,
 \qquad \gamma_i>0,
 \qquad 1\le i\le m.
\]
Suppose
\begin{equation}
 \Gamma_\otimes(1,K)\le C_0.
\label{eq:V24-unpaid-absolute}
\end{equation}
Then
\begin{equation}
 \boxed{
 \Gamma_\otimes(1,K)
 \le
 \frac{C_0}{a_*^m\prod_{i=1}^m\gamma_i}
 \prod_{i=1}^m d_i.}
\label{eq:V24-unpaid-attachment}
\end{equation}
If the unique payer has already been assigned inside the complete
multiplier and
\begin{equation}
 \Gamma_\otimes(w,K^{\rm pay})\le\frac{C_{\rm p}}{s_\alpha},
\label{eq:V24-paid-absolute}
\end{equation}
then
\begin{equation}
 \boxed{
 \Gamma_\otimes(w,K^{\rm pay})
 \le
 \frac{C_{\rm p}}
 {s_\alpha a_*^m\prod_{i=1}^m\gamma_i}
 \prod_{i=1}^m d_i.}
\label{eq:V24-paid-attachment}
\end{equation}
The lemma does not require the responses to arise from distinct
physical occurrences or distinct coordinate banks.
\end{lemma}

\begin{proof}
The hypotheses imply
\[
 1\le
 \prod_{i=1}^m\frac{d_i}{\gamma_i a_*}.
\]
Multiply \eqref{eq:V24-unpaid-absolute} or
\eqref{eq:V24-paid-absolute} by this scalar.  The operation is
performed after the complete analytic bound and neither factors nor
duplicates \(K\).
\end{proof}

\begin{firewall}[Why this does not violate the V92 warning]
\label{fire:V24-V92-compatible}
The forbidden argument would apply separate first-moment estimates
to several coordinate marks inside one unrepeated covariance
occurrence.  \Cref{lem:V24-absolute-attachment} applies exactly one
absolute estimate to the complete multiplier and only then multiplies
by a scalar at least one.  The subsequent V92 expansion is an
identity for that scalar majorant.  It makes no claim that the source
operator analytically supplied one covariance factor for each marked
edge.
\end{firewall}

\begin{proposition}[The attached triple is a legal tree mark]
\label{prop:V24-legal-triple-mark}
For the three V15 responses,
\begin{equation}
 d_D(B_D)d_Q(B_Q)d_S(B_S)
 =
 \frac{
 H_{g_D}(B_D)H_{g_Q}(B_Q)H_{g_S}(B_S)}
 {\Xi_x\Xi_q\Xi_y}
\label{eq:V24-oriented-product}
\end{equation}
is the legal oriented marked-tree response of the corresponding V15
arborescence.  This remains true when the banks overlap or coincide.
Its coordinate expansion is the disjoint sum in
\eqref{eq:V24-five-stratum-identity}.
\end{proposition}

\begin{proof}
The six-row table of
\Cref{prop:V15-six-cell-table} proves that the undirected labels are
the three edges of a tree and that their distinct directed tails are
\(x,q,y\), with root \(z\).  Substitution of
\eqref{eq:V24-restricted-response} proves
\eqref{eq:V24-oriented-product}; this is exactly the normalization
already verified in \Cref{cor:V15-aggregate}.  The archived V92
bank-to-mark bijection allows the same coordinate on more than one
edge and introduces no support-cardinality factor.  Finally apply
\Cref{thm:V24-five-stratum-identity}.
\end{proof}

\section{Closure of the clean termwise \(DQS\) class}
\label{sec:V24-DQS-closure}

\begin{lemma}[A \(DQS\) word is bounded before marking]
\label{lem:V24-DQS-absolute-word}
Fix one exact V20 ancestry-tag class whose occurrence partition is
\(DQS\).  Let \(K_{DQS}\) be its complete V94 dual multiplier,
without splitting its unique response parent.  There are constants
\(C_{\rm abs},C_{\rm pay}\), independent of coordinate support,
representations, and physical multiplicity, such that
\begin{align}
 \Gamma_\otimes(1,K_{DQS})
 &\le C_{\rm abs},
\label{eq:V24-DQS-absolute}\\
 \Gamma_\otimes(w,K_{DQS}^{\rm pay})
 &\le \frac{C_{\rm pay}}{s_\alpha}
\label{eq:V24-DQS-paid}
\end{align}
at every admitted unique-payer location.
\end{lemma}

\begin{proof}
Keep the \(DQS\) parent as one complete occurrence.  Every other
factor is a complete directed covariance, the complete matching
breaker, a complete matching-refining moment, a finite V61/V94
separator, or the complete common suffix.  Complete moments are
contractions.  The archived V99 directed-word theorem, the V100
breaker bound, and the V95--V96 separator and suffix rows give
support-uniform ordinary norms.  Their paid alternatives give one
factor \(C/s_\alpha\) at the assigned payer.

The complete \(DQS\) parent itself has the V18 complete-occurrence
ordinary and paid caps; no coordinate restriction or first-moment
extraction is performed.  Multiplying the fixed finite number of
ordinary bounds proves \eqref{eq:V24-DQS-absolute}.  The V2
one-payer allocation puts exactly one factor on a paid line and all
others on ordinary lines, proving \eqref{eq:V24-DQS-paid}.
Taking the product-state supremum cannot exceed the resulting
complete operator-norm bound.
\end{proof}

\begin{theorem}[Complete one-occurrence triple-collision closure]
\label{thm:V24-complete-DQS}
Every clean, termwise response-admissible \(DQS\) ancestry-tag class
satisfies the legal unpaid and once-paid three-response marked-tree
bounds
\begin{align}
 \Gamma_\otimes(1,K_{DQS})
 &\le
 \frac{C_{\rm abs}}
 {\gamma_D\gamma_Q\gamma_Sa_*^3}\,
 d_D(C)d_Q(C)d_S(C),
\label{eq:V24-DQS-unpaid-tree}\\
 \Gamma_\otimes(w,K_{DQS}^{\rm pay})
 &\le
 \frac{C_{\rm pay}}
 {s_\alpha\gamma_D\gamma_Q\gamma_Sa_*^3}\,
 d_D(C)d_Q(C)d_S(C).
\label{eq:V24-DQS-paid-tree}
\end{align}
The response product is the V15 arborescence mark.  The theorem
allows arbitrary pairwise or triple overlap of the three coordinate
supports.
\end{theorem}

\begin{proof}
Termwise response admissibility gives
\[
 d_D(C)\ge\gamma_Da_*,
 \qquad
 d_Q(C)\ge\gamma_Qa_*,
 \qquad
 d_S(C)\ge\gamma_Sa_*.
\]
Apply \Cref{lem:V24-absolute-attachment} with \(m=3\) to
\Cref{lem:V24-DQS-absolute-word}.  This proves
\eqref{eq:V24-DQS-unpaid-tree}--%
\eqref{eq:V24-DQS-paid-tree}.  By
\Cref{prop:V24-legal-triple-mark}, the scalar product is the legal
directed-tree response and has the exact five-stratum coordinate
expansion.  The analytic multiplier was bounded once before that
expansion, so a diagonal marked coordinate is not a second use of
the physical occurrence.
\end{proof}

\begin{corollary}[No local triple-overlap capacity is needed for
\(DQS\)]
\label{cor:V24-no-triple-capacity}
The proposed V23--V24 target of a separate three-response Fourier
capacity on \(C_{\{D,Q,S\}}\) may be removed from the proof of the
clean, termwise response-admissible \(DQS\) class.
\end{corollary}

\begin{proof}
\Cref{thm:V24-complete-DQS} covers coincident banks and its proof
uses no local overlap capacity.
\end{proof}

\begin{corollary}[Occurrence collision is not the next active gate]
\label{cor:V24-occurrence-ledger}
Among the five occurrence partitions of
\Cref{prop:V23-five-partition-status}, the previously unresolved
\(DQS\) partition is now closed whenever its three responses are
termwise admissible.  Thus no clean termwise collision remains whose
only obstruction is that several desired roles have the same parent
occurrence.
\end{corollary}

\begin{proof}
The exclusive \(QS|D\) partition is V23, and
\Cref{thm:V24-complete-DQS} is the remaining one-parent case.
The other partitions retain the V15 compiler or V13--V14 fallback
recorded in \Cref{prop:V23-five-partition-status}.
\end{proof}

\section{Regression checks and next upstream target}
\label{sec:V24-results}

\begin{proposition}[V24 regression suite]
\label{prop:V24-regressions}
\begin{enumerate}[label=\textup{(\roman*)},leftmargin=4em]
\item The unique \(DQS\) parent is bounded exactly once.
\item Signature atoms are not renamed as physical occurrences.
\item Coordinate triples are partitioned by equality without loss or
      duplication.
\item A repeated coordinate remains visible in a pair- or
      triple-diagonal joint weight.
\item No independent first-moment card is applied to such a repeated
      coordinate.
\item Response insertion occurs only after the complete absolute
      multiplier bound.
\item The scalar response product has the exact V15 tree labels and
      the three distinct tails \(x,q,y\).
\item Every paid alternative has one factor
      \(s_\alpha^{-1}\), never two.
\item Only clean, termwise response-admissible \(DQS\) lines are
      closed.
\end{enumerate}
\end{proposition}

\begin{proof}
Items \textup{(i)}, \textup{(v)}, \textup{(vi)}, and
\textup{(viii)} are
\Cref{lem:V24-DQS-absolute-word,
lem:V24-absolute-attachment,fire:V24-V92-compatible}.
Items \textup{(ii)--(iv)} are
\Cref{lem:V24-signature-disintegration,
thm:V24-five-stratum-identity,
cor:V24-atom-coordinate-compatibility}.
Item \textup{(vii)} is
\Cref{prop:V24-legal-triple-mark}.  Item \textup{(ix)} is a
hypothesis of \Cref{thm:V24-complete-DQS}.
\end{proof}

\begin{theorem}[Third-series V24 result ledger]
\label{thm:V24-results}
Relative to V1--V23, V24 proves:
\begin{enumerate}[label=\textup{(V24.\arabic*)},leftmargin=6.3em]
\item an exact seven-atom signature partition for the three edge
      roles of one \(DQS\) occurrence;
\item a support-uniform heavy-cell selector losing at most \(4^3\);
\item the exhaustive five selected-cell incidence types;
\item an exact five-stratum decomposition of the triple mark by
      physical-coordinate equality;
\item a general absolute attachment lemma for any finite number of
      already dominant scalar responses;
\item compatibility of that lemma with the archived V92 repeated-mark
      firewall; and
\item complete closure of every clean, termwise
      response-admissible \(DQS\) ancestry-tag class.
\end{enumerate}
\end{theorem}

\begin{proof}
These are
\Cref{lem:V24-signature-disintegration,
cor:V24-heavy-signature,
prop:V24-five-cell-types,
thm:V24-five-stratum-identity,
lem:V24-absolute-attachment,
fire:V24-V92-compatible,
thm:V24-complete-DQS}.
\end{proof}

\begin{assessment}[Bottleneck after V24]
\label{ass:V24-bottleneck}
The one-unrepeated-occurrence collision is no longer the first
obstruction.  The next upstream target is response admissibility
itself on the V17 overlap and remainder lines and on the unresolved
V13--V14 self-role fallbacks.  Those lines may fail to supply one of
the three dominant responses, so the absolute attachment lemma cannot
manufacture the missing lower bound.  V25 must first perform the
compulsory companion audit and then choose between these depleted
response gates.
\end{assessment}

\begin{programme}[V25: audit and depleted-response normal form]
\label{prog:V24-V25}
\begin{enumerate}[leftmargin=3em]
\item Inspect the newest active Standard-Model--Einstein and
      Navier--Stokes notes before choosing the branch.
\item Return to the exact V17 four-cell support telescope and record
      which response lower bound fails on each overlap or remainder
      line.
\item Merge that ledger with the V13--V14 reverse-or-deplete and
      low-response self-role alternatives.
\item Seek a finite response-or-depletion theorem before introducing
      any new operator capacity.
\item Keep nonclean states and the higher-order source families
      separate.
\end{enumerate}
\end{programme}

\begin{firewall}[Claim boundary after third-series V24]
\label{fire:V24-claim-boundary}
V24 closes the clean, termwise response-admissible one-occurrence
\(DQS\) collision.  It does not prove the missing response lower
bounds on V17 overlap or remainder lines, close every V13--V14
self-role fallback, the V2 low-response self-role factor, or
nonclean states.  Higher MTA, WUFC, CPM, cutoff removal, reflection
positivity, Osterwalder--Schrader reconstruction, construction of the
continuum Yang--Mills measure, and a uniform positive continuum
spectral gap remain open.  The full Yang--Mills mass gap has not been
proved.
\end{firewall}

\paragraph{Parent-version record.}
V24 was copied byte-for-byte from
\texttt{Renewal\_Yang\_Mills\_Mass\_Gap\_Research\_Notes\_V23.tex}
before this chapter was appended.  The V23 parent had \(12\,004\)
source lines, \(437\,747\) bytes, and SHA--256
\[
 \texttt{b3d3d309cf0fd2f95a0e0e95c1f936495cad6e5ce059afcd547e3b8ae75e4e2b}.
\]
The next iteration is the compulsory V25 companion audit.

% =============================================================================
% BEGIN APPENDED THIRD-SERIES V25 MATERIAL
% =============================================================================

\clearpage
\chapter{V25: Companion audit and parent-response stability under exact
telescopes}
\label{ch:V25-parent-response}

\section{Compulsory companion audit}
\label{sec:V25-audit}

The newest active Standard-Model--Einstein and Navier--Stokes notes
were inspected read-only before the V25 direction was chosen.

\begin{record}[Companion snapshots used by V25]
\label{rec:V25-companion-snapshots}
The files actually read at this checkpoint were:
\begin{enumerate}[leftmargin=3em]
\item
\(\texttt{Renewal\_SM\_Einstein\_Continuum\_Research\_Notes\_V4.tex}\),
with \(1\,650\) source lines, \(66\,347\) bytes, and SHA--256
\[
 \texttt{c1e312d222369be51f8e458001ef662e196837e7fda3a4526ca84fe66fdb2d4e};
\]
\item
\(\texttt{Renewal\_NS\_Stability\_Research\_Notes\_V31.tex}\),
with \(23\,052\) source lines, \(887\,537\) bytes, and SHA--256
\[
 \texttt{f1121b62238ad5491bb7cf03635ea9934942edf573ed8faaa9ee79e1d38a6696}.
\]
\end{enumerate}
The active SM file advanced while the audit was being started; V4 is
the later snapshot actually used.  No companion file was modified.
\end{record}

\begin{assessment}[Type-correct transfer from the audit]
\label{ass:V25-audit-transfer}
SM V4 defines one exact connected shell map and only then projects
its output into marginal and irrelevant coordinates.  Its Taylor
allocation assigns every connected occurrence to one projection; it
does not require each projected child to contain a second copy of the
parent interaction.  It also warns that a contraction of the
irrelevant complement is conditional on retaining the complete local
projection before taking norms.

NS V31 proves that probability marking removes stopping-time
multiplicity, while unmarking costs an actual first moment.  Its
formation bound is applied to the complete Duhamel source before the
dyadic mark is expanded, and its firewall forbids replacing a
restricted correlated source by a continuation-class quantity and
calling the result a priori.

The Yang--Mills transfer is only proof order.  A support-telescope
child will be bounded as one complete multiplier; the response
product already proved for its parent may then be attached to the
resulting scalar majorant.  The child is not said to own those
responses.  This can preserve an existing first moment across a
finite exact decomposition, but it cannot manufacture a parent
response which was never acquired.  No RG, entropy, Duhamel, or
Navier--Stokes estimate is imported.
\end{assessment}

\section{Parent responses survive finite analytic decomposition}
\label{sec:V25-abstract}

\begin{definition}[Parent-admissible response packet]
\label{def:V25-parent-admissible}
Let \(K\) be a complete assembled packet and suppose its construction
has already produced three parent responses
\begin{equation}
 d_D\ge\gamma_Da_*,
 \qquad
 d_Q\ge\gamma_Qa_*,
 \qquad
 d_S\ge\gamma_Sa_*,
\label{eq:V25-parent-responses}
\end{equation}
whose physical edge labels and tails are one row of the V15
six-arborescence table.  The packet is called
\emph{parent-admissible}.  This term makes no assertion that every
child in a later exact decomposition retains the three lower bounds
on its restricted coordinate support.
\end{definition}

\begin{lemma}[Finite parent-response attachment]
\label{lem:V25-finite-parent-attachment}
Suppose an exact operator identity has the form
\begin{equation}
 K=\sum_{\lambda\in\Lambda}\epsilon_\lambda K_\lambda,
 \qquad
 \epsilon_\lambda\in\{-1,1\},
 \qquad
 |\Lambda|\le N_*,
\label{eq:V25-finite-decomposition}
\end{equation}
where \(N_*\) is support-uniform.  Assume the parent packet is
parent-admissible and each complete child has the absolute bounds
\begin{align}
 \Gamma_\otimes(1,K_\lambda)
 &\le C_\lambda,
\label{eq:V25-child-unpaid}\\
 \Gamma_\otimes(w,K_\lambda^{\rm pay})
 &\le \frac{C_{\lambda,\rm p}}{s_\alpha}
\label{eq:V25-child-paid}
\end{align}
at every disjoint assignment of the unique payer.  Then
\begin{align}
 \Gamma_\otimes(1,K)
 &\le
 \frac{\sum_{\lambda}C_\lambda}
 {\gamma_D\gamma_Q\gamma_Sa_*^3}\,
 d_Dd_Qd_S,
\label{eq:V25-parent-unpaid}\\
 \Gamma_\otimes(w,K^{\rm pay})
 &\le
 \frac{\sum_{\lambda}C_{\lambda,\rm p}}
 {s_\alpha\gamma_D\gamma_Q\gamma_Sa_*^3}\,
 d_Dd_Qd_S.
\label{eq:V25-parent-paid}
\end{align}
No child-response lower bound is required.
\end{lemma}

\begin{proof}
Subadditivity of the product-state majorant and
\eqref{eq:V25-finite-decomposition} give
\[
 \Gamma_\otimes(1,K)
 \le\sum_{\lambda\in\Lambda}
 \Gamma_\otimes(1,K_\lambda)
 \le\sum_\lambda C_\lambda.
\]
The parent response lower bounds imply
\[
 1\le
 \frac{d_Dd_Qd_S}
 {\gamma_D\gamma_Q\gamma_Sa_*^3}.
\]
Multiplying the preceding complete bound by this scalar proves
\eqref{eq:V25-parent-unpaid}.

For a paid packet, the exact payer allocation is a disjoint sum of
children and payer locations.  Apply
\eqref{eq:V25-child-paid} once in each such line and use ordinary
bounds for every other factor.  Subadditivity gives the numerator
\(\sum_\lambda C_{\lambda,\rm p}/s_\alpha\), with only one inverse
scale.  Attach the same parent scalar to prove
\eqref{eq:V25-parent-paid}.
\end{proof}

\begin{corollary}[Signs and Möbius coefficients cause no new gate]
\label{cor:V25-signed-decomposition}
The conclusion of
\Cref{lem:V25-finite-parent-attachment} remains valid for a fixed
finite coefficient family
\[
 K=\sum_{\lambda\in\Lambda}c_\lambda K_\lambda
\]
after replacing \(C_\lambda\) by
\(|c_\lambda|C_\lambda\) and similarly on the paid line.  Hence a
fixed partition-lattice Möbius expansion does not require the parent
responses to be reassigned to its signed children.
\end{corollary}

\begin{proof}
Use absolute homogeneity and subadditivity of
\(\Gamma_\otimes\), then repeat the proof of
\Cref{lem:V25-finite-parent-attachment}.
\end{proof}

\begin{firewall}[The parent response is external currency, not child
ancestry]
\label{fire:V25-no-false-inheritance}
Equations \eqref{eq:V25-parent-unpaid}--%
\eqref{eq:V25-parent-paid} do not assert
\[
 d_R(K_\lambda)\ge\gamma_Ra_*.
\]
They use the one parent value \(d_R\) after the sum of complete child
bounds has been formed.  Thus a response is neither copied among
children nor localized to a child whose covariance lies on another
support cell.
\end{firewall}

\begin{proposition}[The attached parent product is still the same
marked tree]
\label{prop:V25-parent-tree}
Under \eqref{eq:V25-parent-responses}, the product
\(d_Dd_Qd_S\) on the right of
\eqref{eq:V25-parent-unpaid}--%
\eqref{eq:V25-parent-paid} is the legal V15 arborescence response.
If any of its edge banks overlap, its V92 coordinate expansion
retains every diagonal marking exactly once.
\end{proposition}

\begin{proof}
The exact decomposition changes the operator \(K\), not the three
physical edge banks and their tail normalizations.  Apply
\Cref{cor:V15-aggregate} to identify the directed tree.  The
archived V92 bank-to-mark identity and
\Cref{thm:V24-five-stratum-identity} give the overlap expansion
without a support-cardinality factor.
\end{proof}

\begin{firewall}[The attachment lemma cannot create a first moment]
\label{fire:V25-no-created-response}
If one of the inequalities in
\eqref{eq:V25-parent-responses} is absent, the factor at least one
used in the proof is absent.  An ordinary norm bound on \(K\), a
small child response, or a large unnormalized row mass cannot replace
it.  This is the exact analogue of the first-moment boundary observed
in the NS audit.
\end{firewall}

\section{Application to the V17 four-cell telescope}
\label{sec:V25-V17}

\begin{lemma}[Uniform absolute bounds for all four V17 lines]
\label{lem:V25-four-line-absolute}
Fix a clean \(QS|D\) parent packet whose total parent responses
\(d_D,d_Q,d_S\) satisfy
\eqref{eq:V25-parent-responses}.  Each of the four exact children in
\eqref{eq:V17-QS-telescope}, including the overlap and remainder
children, has a support-uniform ordinary bound and a
support-uniform once-paid \(C/s_\alpha\) bound inside the unchanged
surrounding word.
\end{lemma}

\begin{proof}
Each child contains exactly one restricted covariance difference,
whose ordinary norm is at most two and whose complete once-paid norm
is at most \(C/s_\alpha\).  Every other restriction of the parent
occurrence is an endpoint moment and hence a contraction.  Factors
outside that occurrence are the same complete directed words,
matching breaker, finite separators, and grouped suffix bounded in
\Cref{lem:V23-assembled-norms}.

The V14/V16 payer split assigns the sole payer to exactly one
restricted occurrence or one unchanged outside factor.  Put that
factor on its established paid line and all other factors on ordinary
lines.  There are four children and a fixed number of outer factors,
so the resulting constants are support-uniform.
\end{proof}

\begin{theorem}[Unconditional parent-admissible \(QS|D\) closure]
\label{thm:V25-QS-D-parent-closure}
Every clean \(QS|D\) packet whose three \emph{parent} responses obey
\eqref{eq:V25-parent-responses} satisfies the legal unpaid and
once-paid V15 marked-tree bounds.  No restricted child in the V17
overlap, \(Q\)-exclusive, \(S\)-exclusive, or remainder line is
required to retain all three response lower bounds.
\end{theorem}

\begin{proof}
The exact four-term identity
\eqref{eq:V17-QS-telescope} is
\eqref{eq:V25-finite-decomposition} with \(N_*=4\).
\Cref{lem:V25-four-line-absolute} supplies its child bounds.
Apply
\Cref{lem:V25-finite-parent-attachment,
prop:V25-parent-tree}.
\end{proof}

\begin{corollary}[The V17 overlap and remainder gates are removed]
\label{cor:V25-V17-gates-removed}
For a parent-admissible clean \(QS|D\) packet:
\begin{enumerate}[label=\textup{(\roman*)},leftmargin=4em]
\item the overlap line does not require the V83 dominant-square gate;
\item the remainder line does not require its covariance difference
      to carry either the \(Q\)- or \(S\)-response; and
\item the exclusive lines do not require termwise response retention
      on their restricted cells.
\end{enumerate}
The sharper V83 and endpoint-energy theorems remain valid optional
estimates but are not gates for this aggregate compiler.
\end{corollary}

\begin{proof}
All four lines are included in
\Cref{thm:V25-QS-D-parent-closure}, whose proof uses only their
complete absolute bounds and the total parent responses.
\end{proof}

\begin{corollary}[Strengthening of V23]
\label{cor:V25-strengthen-V23}
The phrase \emph{termwise response-admissible} in
\Cref{thm:V23-complete-QS-D} may be replaced by
\emph{parent-admissible}.  In particular, the termwise response
caveat in the V23 claim boundary is discharged for clean \(QS|D\)
packets.
\end{corollary}

\begin{proof}
\Cref{thm:V25-QS-D-parent-closure} has exactly the stronger
hypothesis and conclusion.
\end{proof}

\section{The collision-erasure theorem}
\label{sec:V25-collision-erasure}

\begin{theorem}[Finite support and occurrence telescopes do not
obstruct an acquired directed triple]
\label{thm:V25-collision-erasure}
Let a clean complete first-breaker packet have the three acquired
responses \(D,Q,S\) of a V15 arborescence.  Apply any fixed finite
sequence of:
\begin{enumerate}[label=\textup{(\alph*)},leftmargin=4em]
\item V14/V16 coordinate-support telescopes;
\item V20 ancestry and slot refinements;
\item fixed partition-lattice cover or Möbius expansions; and
\item disjoint unique-payer allocations.
\end{enumerate}
If every resulting complete word uses only the established complete
occurrence, breaker, separator, and suffix factors, then the whole
expanded packet satisfies the legal unpaid and once-paid
three-response marked-tree bounds.  The conclusion is independent of
which of
\[
 D|Q|S,\quad DQ|S,\quad DS|Q,\quad QS|D,\quad DQS
\]
is the parent occurrence partition.
\end{theorem}

\begin{proof}
Each listed operation has a universally finite branching number and
is an exact identity before norms.  V14/V16 preserve one covariance
difference per parent occurrence and one payer; V20 preserves
ancestry and records a finite number of slots.  Fixed lattice
expansions have fixed coefficients.  Therefore the final expansion
has the form of
\Cref{lem:V25-finite-parent-attachment} or
\Cref{cor:V25-signed-decomposition}.

The hypothesis on complete word types and the proof of
\Cref{lem:V23-assembled-norms} give uniform ordinary and one-payer
bounds for every child.  Apply finite parent-response attachment.
\Cref{prop:V25-parent-tree} identifies the common scalar product as
the V15 tree.  Nothing in this argument uses the equality pattern of
the three parent slots, proving the last assertion.
\end{proof}

\begin{corollary}[Complete clean acquired-triple collision module]
\label{cor:V25-complete-collision-module}
Once the three actual directed responses \(D,Q,S\) have been
acquired in a clean first-breaker packet, occurrence collision,
coordinate overlap, response loss under a support telescope, and
noncommutative child placement are not further obstructions to the
aggregate quartic tree majorant.
\end{corollary}

\begin{proof}
Occurrence collision and coordinate overlap are allowed in
\Cref{thm:V25-collision-erasure}; child response loss is handled by
\Cref{fire:V25-no-false-inheritance}; noncommutative placement is
irrelevant because each child is bounded as a complete word before
response attachment.
\end{proof}

\begin{firewall}[Scope of collision erasure]
\label{fire:V25-collision-scope}
\Cref{thm:V25-collision-erasure} is conditional on an already
acquired V15 directed triple and on the established uniform absolute
bounds for the complete word types.  It does not close a branch whose
third response is missing, a low-response self-role correction, a
nonclean representation state, or a source family not contained in
the clean first-breaker packet.
\end{firewall}

\section{Updated gate ledger}
\label{sec:V25-ledger}

\begin{proposition}[What remains of the V13--V17 decision tree]
\label{prop:V25-updated-tree}
In the clean first-breaker analysis:
\begin{enumerate}[label=\textup{(\roman*)},leftmargin=4em]
\item every branch which reaches the V15 triple \(D,Q,S\) closes,
      without any occurrence-partition subcase;
\item direct reverse and complete \(P,U,A,F,I\) branches retain their
      earlier closures;
\item the remaining directed search is needed only where one of the
      three V15 roles has not yet been acquired;
\item the partition-changing self-bank residual is reduced to the
      V2 low-response sensitive correction plus an anchor-free role
      difference; and
\item nonclean states remain outside the clean calculus.
\end{enumerate}
\end{proposition}

\begin{proof}
Item \textup{(i)} is
\Cref{cor:V25-complete-collision-module}.  Item \textup{(ii)} is
\Cref{thm:V13-spectator-fan,
thm:V14-target-eligibility}.  The collision module removes every
post-acquisition obstruction, proving \textup{(iii)}.  Item
\textup{(iv)} is
\Cref{thm:V2-self-role-reduction}; item \textup{(v)} is the clean
typing boundary retained throughout V1--V24.
\end{proof}

\begin{assessment}[Bottleneck after V25]
\label{ass:V25-bottleneck}
The active quartic obstruction has moved upstream from collision
analysis to response acquisition.  The sharp residual is no longer
the allocation of a response to a telescope child, but whether the
exact connected packet acquires the missing V15 response before it
reaches the V2 low-response anchor-free factor.  The V14 no-return
theorem prevents immediate return to one depleted incidence bank, but
does not yet supply a global decreasing quantity for repeated
avoid-partner/self-role transitions.
\end{assessment}

\section{Regression checks and result ledger}
\label{sec:V25-results}

\begin{proposition}[V25 regression suite]
\label{prop:V25-regressions}
\begin{enumerate}[label=\textup{(\roman*)},leftmargin=4em]
\item Companion results are transferred only as proof-order
      constraints.
\item Every exact child is bounded before the parent response is
      attached.
\item A child is never said to carry a response absent from its
      support.
\item The parent response product is expanded only after the complete
      analytic bound.
\item No support cardinality enters a finite telescope.
\item Signs are paid by a fixed absolute coefficient sum.
\item Every paid line contains one \(s_\alpha^{-1}\).
\item All five parent occurrence partitions are covered only after
      the V15 triple has been acquired.
\item Missing responses, low-response self-role, and nonclean states
      remain explicit.
\end{enumerate}
\end{proposition}

\begin{proof}
Item \textup{(i)} is
\Cref{ass:V25-audit-transfer}.  Items
\textup{(ii)--(iv)} are
\Cref{lem:V25-finite-parent-attachment,
fire:V25-no-false-inheritance,
prop:V25-parent-tree}.  Items \textup{(v)--(vii)} follow from the
fixed finite expansion and disjoint payer proof in
\Cref{lem:V25-finite-parent-attachment,
cor:V25-signed-decomposition}.  Items
\textup{(viii)--(ix)} are
\Cref{thm:V25-collision-erasure,fire:V25-collision-scope}.
\end{proof}

\begin{theorem}[Third-series V25 result ledger]
\label{thm:V25-results}
Relative to V1--V24, V25 proves:
\begin{enumerate}[label=\textup{(V25.\arabic*)},leftmargin=6.3em]
\item a compulsory type-correct audit of active SM V4 and NS V31;
\item finite parent-response attachment across arbitrary exact signed
      child decompositions;
\item preservation of the exact V15/V92 marked-tree product under
      that attachment;
\item uniform absolute bounds for all four V17 support-telescope
      lines;
\item unconditional parent-admissible closure of the clean
      \(QS|D\) overlap, exclusive, and remainder lines;
\item removal of termwise response retention and the V83 overlap
      gate from that aggregate compiler; and
\item a collision-erasure theorem covering all five occurrence
      partitions after the directed triple has been acquired.
\end{enumerate}
\end{theorem}

\begin{proof}
These are
\Cref{rec:V25-companion-snapshots,
ass:V25-audit-transfer,
lem:V25-finite-parent-attachment,
cor:V25-signed-decomposition,
prop:V25-parent-tree,
lem:V25-four-line-absolute,
thm:V25-QS-D-parent-closure,
cor:V25-V17-gates-removed,
thm:V25-collision-erasure}.
\end{proof}

\begin{programme}[V26: a well-founded acquisition ledger]
\label{prog:V25-V26}
\begin{enumerate}[leftmargin=3em]
\item Encode a clean response-search state by its missing V15 tail,
      depleted literal incidence sets, used physical roles, and
      unresolved self-role factors.
\item Use the V14 strict no-return property to determine which
      transitions strictly enlarge the used-role set.
\item Quotient transitions which differ only by a closed
      parent-admissible collision, now harmless by V25.
\item Test whether every avoid-partner cycle enters either a complete
      compiler, an acquired \(D,Q,S\) triple, or the V2
      low-response anchor-free factor in a uniformly bounded number
      of steps.
\item If a cycle survives, isolate its first repeated state and move
      upstream to the exact connected partition word rather than
      iterating scalar acquisition again.
\item Keep nonclean states and higher-order source families separate.
      The next compulsory companion audit is V30.
\end{enumerate}
\end{programme}

\begin{firewall}[Claim boundary after third-series V25]
\label{fire:V25-claim-boundary}
V25 closes all clean collision/support-telescope obstructions after
the V15 directed triple is acquired.  It does not prove termination
of the remaining response-acquisition search, close the V2
low-response anchor-free role factor, or handle nonclean states.
Higher MTA, WUFC, CPM, cutoff removal, reflection positivity,
Osterwalder--Schrader reconstruction, construction of the continuum
Yang--Mills measure, and a uniform positive continuum spectral gap
remain open.  The full Yang--Mills mass gap has not been proved.
\end{firewall}

\paragraph{Parent-version record.}
V25 was copied byte-for-byte from
\texttt{Renewal\_Yang\_Mills\_Mass\_Gap\_Research\_Notes\_V24.tex}
before this chapter was appended.  The V24 parent had \(12\,634\)
source lines, \(459\,461\) bytes, and SHA--256
\[
 \texttt{7acf32a855b728fb2c3b57c257b38fcddbbba88844fe3455a6f056600d617d5f}.
\]
The compulsory V25 companion audit is recorded in
\Cref{sec:V25-audit}; the next required audit is V30.

% =============================================================================
% BEGIN APPENDED THIRD-SERIES V26 MATERIAL
% =============================================================================

\clearpage
\chapter{V26: Two-stage acquisition terminates at the private-anchor
residual}
\label{ch:V26-two-stage}

\section{Context: reversal is no longer part of the active route}
\label{sec:V26-context}

The V13 reverse-or-deplete analysis was necessary when the proof
used only the raw breaker, the right matching response, and a
reversed \(q\)-edge.  V15 observed that, after an avoid-\(q\)
response \(D\) has been acquired at \(x\), every forward
\(q\)-rooted response \(Q\) already completes the directed triple
\(D,Q,S\).  At that time occurrence collisions still prevented a
uniform conclusion.  V25 has now removed that last condition.

Consequently the active first-breaker route has only two acquisition
stages.  The \(x\)-stage either closes, enters a complete compiler,
produces \(D\), or enters its V2 low-response self-role residual.
After \(D\), the first \(q\)-stage either closes, enters a complete
compiler, or enters the analogous V2 residual.  A forward response at
\(q\) is terminal; it must not be reversed and sent through V14.

\begin{assessment}[Correction of the V25 acquisition programme]
\label{ass:V26-programme-correction}
There is no avoid-partner cycle to which a global decreasing
used-role quantity must be attached in the active clean packet.
The V12--V13 state tree has depth two once
\Cref{thm:V25-collision-erasure} is inserted at the first
\(q\)-rooted response.  The genuine residual is the response-poor
partition-changing self-role factor, not repeated directed
acquisition.
\end{assessment}

\section{The exact two-stage decision tree}
\label{sec:V26-decision-tree}

\begin{definition}[Terminal and residual leaves]
\label{def:V26-leaves}
A clean branch is called:
\begin{enumerate}[label=\textup{(\roman*)},leftmargin=4em]
\item \emph{tree-terminal} if it is bounded by a legal quartic
      marked-tree response;
\item \emph{compiler-terminal} if it enters one of the established
      complete \(P,U,A,F,I\) or same-type self-bank compilers; or
\item \emph{private-anchor residual} if a V2
      partition-changing self-bank reaches the low-response
      factorization and every response candidate in its common
      anchor factor is below the fixed threshold introduced below.
\end{enumerate}
\end{definition}

\begin{lemma}[Every \(x\)-response has a terminal or \(D\)-fate]
\label{lem:V26-x-response-fate}
Retain the right matching response \(S=y\to z\) and the raw breaker
\(xy\).  Every actual response \(x\to r\) produced by the V12
post-prefix acquisition or its response-first self-role branch has
exactly one of:
\begin{enumerate}[label=\textup{(\alph*)},leftmargin=4em]
\item \(r=q\), in which case the path \(q-x-y-z\) is
      tree-terminal; or
\item \(r\in\{y,z\}\), in which case it is the response
      \(D=x\to r\) and the next and final acquisition row is \(q\).
\end{enumerate}
\end{lemma}

\begin{proof}
This is the exhaustive partner table in
\Cref{prop:V12-partner-table}.  Its \(r=q\) row is the earlier
complementary-partner closure.  Its other two rows are exactly the
two directed avoid-\(q\) cells used to define \(D\) in V15.
\end{proof}

\begin{lemma}[The first \(q\)-response after \(D\) is terminal]
\label{lem:V26-q-response-terminal}
Suppose \(D=x\to u\), \(u\in\{y,z\}\), has been acquired.  Every
actual response
\[
 Q=q\to r,\qquad r\in\{x,y,z\},
\]
produced by rowwise acquisition or a response-first self-role
reduction at \(q\) is tree-terminal, regardless of its physical
occurrence, coordinate overlap, or reverse response.
\end{lemma}

\begin{proof}
Together \(D,Q,S\) are one of the six V15 arborescences.
\Cref{thm:V25-collision-erasure} gives the complete unpaid and
once-paid marked-tree bound for every occurrence partition and every
finite support telescope.  No reverse \(r\to q\) is a hypothesis of
that theorem.
\end{proof}

\begin{theorem}[Two-stage clean acquisition theorem]
\label{thm:V26-two-stage}
For every clean V100 first-breaker packet in the V12 missing-tail
state, the exact V12--V13 acquisition tree has depth at most two and
ends in one of:
\begin{enumerate}[label=\textup{(\roman*)},leftmargin=4em]
\item a tree-terminal branch;
\item a compiler-terminal branch;
\item a V2 low-response partition-changing self-role factor at
      \(x\) or, after one \(D\)-response, at \(q\); or
\item a nonclean representation state.
\end{enumerate}
In particular, no clean forward \(q\)-response enters a
reverse-or-deplete continuation in the minimal proof.
\end{theorem}

\begin{proof}
If the pre-breaker \(x\to q\) response is dominant, V11 is
tree-terminal.  Otherwise apply
\Cref{thm:V12-post-prefix-acquisition}.  Its fresh
\(P,U,A,F,I\) cells and common same-type self-bank cells are
compiler-terminal.  Every old or fresh external response, and every
response-first partition-changing self-bank, is covered by
\Cref{lem:V26-x-response-fate}.  The \(x\to q\) partner is terminal.
The two avoid-\(q\) partners produce \(D\).  The only other clean
\(x\)-leaf is the V2 low-response factor in
\Cref{thm:V12-self-reduction}.

After \(D\), apply the rowwise acquisition at \(q\) exactly once, as
in \Cref{thm:V13-two-sided-normal-form}.  Complete types and
same-type self-banks are compiler-terminal.  Every external,
directed, or response-first self-role response is
\Cref{lem:V26-q-response-terminal}.  The remaining clean leaf is the
V2 low-response factor at \(q\).  Failure of clean typing gives
\textup{(iv)}.  This list is exhaustive and has at most the two
stages just described.
\end{proof}

\begin{corollary}[The V13--V14 reversal branch is nonminimal here]
\label{cor:V26-no-reversal}
Alternatives \textup{(iii)--(iv)} of
\Cref{thm:V13-two-sided-normal-form} both close at their common
forward \(q\)-response by
\Cref{lem:V26-q-response-terminal}.  The V13 direct-reversal test and
V14 depleted-target acquisition remain valid theorems, but they are
not entered by the minimal clean first-breaker proof after V25.
\end{corollary}

\begin{proof}
Both alternatives begin with an already dominant
\(q\to r\).  Their split depends only on the size of the reverse
\(r\to q\), which
\Cref{lem:V26-q-response-terminal} does not use.
\end{proof}

\begin{firewall}[Depth two does not include self-role analysis]
\label{fire:V26-depth-scope}
The depth count concerns rowwise external/directed acquisition.
It does not declare the V2 low-response self-role factor zero or
compiler-terminal.  That factor is the exceptional leaf retained in
\Cref{thm:V26-two-stage}.
\end{firewall}

\section{A response inside the common self-role moment is now legal}
\label{sec:V26-common-response}

Fix a V2 low-response self-bank with anchor \(v\in\{x,q\}\) and
parameter \(\beta>0\).  Retain the exact factorization
\begin{equation}
 \Delta_G^{\rho,\sigma}
 =
 \Delta_{S_v}^{\rho,\sigma}\otimes M_{F_v}(\sigma)
 +
 M_{S_v}(\rho)\otimes A_{F_v}^v\otimes
 \bigl(R_{F_v}(\sigma)-R_{F_v}(\rho)\bigr)
\label{eq:V26-low-response-factorization}
\end{equation}
from V2.  The low-response data are
\begin{equation}
 r_v(\rho,\sigma;G)<\frac{\beta}{2}a_*,
 \qquad
 \sum_{e\in F_v}a_{v,e}\ge\frac{\beta}{2}\Xi_v.
\label{eq:V26-low-response-data}
\end{equation}

Use the V6 least-partner partition of the common factor:
\[
 F_v^{(v)}
 =
 {\cal P}_v\dot\cup
 \bigdotcup_{r\ne v}{\cal C}_{vr},
\]
where \(F_v^{(v)}\) denotes its \(v\)-active part,
\({\cal P}_v\) is \(v\)-private, and
\({\cal C}_{vr}\) has \(r\) in the common local block.  Put
\[
 P_v^{\rm com}:=\sum_{e\in{\cal P}_v}a_{v,e},
 \qquad
 C_{vr}^{\rm com}:=\sum_{e\in{\cal C}_{vr}}a_{v,e}.
\]

\begin{lemma}[Exclusive common-response/private alternative]
\label{lem:V26-common-exclusive}
Fix the response threshold
\[
 \eta_v:=\frac{\beta}{12}.
\]
Exactly one of the following alternatives may be selected:
\begin{enumerate}[label=\textup{(\alph*)},leftmargin=4em]
\item some \(r\ne v\) obeys
      \begin{equation}
      d_{v\to r}({\cal C}_{vr})
      \ge\eta_va_*;
      \label{eq:V26-common-response}
      \end{equation}
\item every common response is below that threshold, and
      \begin{align}
      \sum_{r\ne v}C_{vr}^{\rm com}
      &<\frac{\beta}{4}\Xi_v,
      \label{eq:V26-common-small}\\
      P_v^{\rm com}
      &>\frac{\beta}{4}\Xi_v.
      \label{eq:V26-private-large}
      \end{align}
\end{enumerate}
\end{lemma}

\begin{proof}
Select \textup{(a)} if any displayed response reaches its threshold.
Otherwise, for every \(r\ne v\), the Casimir-floor inequality
\[
 a_*C_{vr}^{\rm com}
 \le H_{vr}({\cal C}_{vr})
 <\frac{\beta}{12}a_*\Xi_v
\]
gives
\(C_{vr}^{\rm com}<\beta\Xi_v/12\).
There are at most three partners, proving
\eqref{eq:V26-common-small}.  The \(v\)-active common-factor identity
and \eqref{eq:V26-low-response-data} then give
\[
 P_v^{\rm com}
 =
 \sum_{e\in F_v}a_{v,e}
 -
 \sum_{r\ne v}C_{vr}^{\rm com}
 >
 \frac{\beta}{4}\Xi_v,
\]
which is \eqref{eq:V26-private-large}.  The threshold convention
makes the two selected alternatives disjoint.
\end{proof}

\begin{lemma}[Common-moment response attachment]
\label{lem:V26-common-moment-attachment}
In alternative \eqref{eq:V26-common-response}, the edge
\(v\to r\) is a legal acquired response for the two-stage decision
tree.  No separate first-moment estimate or covariance relabelling of
the common moment \(A_{F_v}^v\) is required.
\end{lemma}

\begin{proof}
The V6 construction gives the literal physical stock
\(H_{vr}({\cal C}_{vr})\) and its correct tail denominator
\(\Xi_v\).  Keep the entire low-response factor
\eqref{eq:V26-low-response-factorization} in the complete assembled
word and apply its established support-uniform ordinary or
once-paid bound once.  Then attach the scalar response by
\Cref{lem:V24-absolute-attachment}.  This is the proof order of
\Cref{fire:V25-no-false-inheritance}: the common factor remains a
moment and is not used as another analytic covariance occurrence.
\end{proof}

\begin{corollary}[Every nonprivate low-response leaf re-enters and
terminates]
\label{cor:V26-nonprivate-terminates}
If \eqref{eq:V26-common-response} occurs at the \(x\)-stage, it either
closes directly when \(r=q\), or supplies \(D\) and terminates after
one \(q\)-stage.  If it occurs at the \(q\)-stage, it terminates
immediately as the response \(Q\).
\end{corollary}

\begin{proof}
By \Cref{lem:V26-common-moment-attachment}, the common-moment stock
is an actual response in the decision tree.
Apply \Cref{lem:V26-x-response-fate} when \(v=x\), and
\Cref{lem:V26-q-response-terminal} when \(v=q\).
\end{proof}

\section{Exact private-anchor terminal normal form}
\label{sec:V26-private}

\begin{definition}[Private-anchor residual]
\label{def:V26-private-residual}
The residual \({\mathfrak R}_{\rm priv}(v)\) consists of the exact
two terms in \eqref{eq:V26-low-response-factorization} under:
\begin{align}
 \sum_{r\in P_v(\rho,\sigma)}
 d_{v\to r}(S_v)
 &<\frac{\beta}{2}a_*,
\label{eq:V26-sensitive-small}\\
 d_{v\to r}({\cal C}_{vr})
 &<\frac{\beta}{12}a_*
 \quad(r\ne v),
\label{eq:V26-common-responses-small}\\
 P_v^{\rm com}
 &>\frac{\beta}{4}\Xi_v.
\label{eq:V26-private-residual-mass}
\end{align}
Its role difference
\[
 R_{F_v}(\sigma)-R_{F_v}(\rho)
\]
acts on at most three rows and omits \(v\), while the large
\(v\)-mass lies in the common one-row factor
\[
 A_{{\cal P}_v}^v
 =
 \bigotimes_{e\in{\cal P}_v}M_e(\{v\}).
\]
\end{definition}

\begin{proposition}[Complete reduction of clean acquisition]
\label{prop:V26-complete-reduction}
Every clean first-breaker packet is a support-uniform finite sum of:
\begin{enumerate}[label=\textup{(\roman*)},leftmargin=4em]
\item tree-terminal packets;
\item compiler-terminal packets; and
\item at most two private-anchor residual types,
      \({\mathfrak R}_{\rm priv}(x)\) and, conditional on a prior
      \(D\)-response, \({\mathfrak R}_{\rm priv}(q)\).
\end{enumerate}
The remaining alternative is a separately labelled nonclean state.
\end{proposition}

\begin{proof}
Apply \Cref{thm:V26-two-stage}.  At each of its V2 low-response
leaves, use \Cref{lem:V26-common-exclusive}.  Its response
alternative terminates by
\Cref{cor:V26-nonprivate-terminates}; its complementary alternative
is exactly \Cref{def:V26-private-residual}.  The V12 and V13
acquisition theorems have fixed finite type, role, and partner
partitions, so the total number of summands is support-uniform.
\end{proof}

\begin{firewall}[Private mass is not a missing edge]
\label{fire:V26-private-no-edge}
The lower bound \eqref{eq:V26-private-residual-mass} is a one-row
mass statement.  It supplies no edge stock and no tree mark.  The
small-response inequalities do not make the sensitive operator term
small in norm.  Thus \({\mathfrak R}_{\rm priv}(v)\) cannot be
discarded, marked, or absorbed using only the scalar ledger.
\end{firewall}

\begin{assessment}[Bottleneck after V26]
\label{ass:V26-bottleneck}
The clean first-breaker response search is now finite and complete.
Its sole unclosed clean leaf has a sharply constrained geometry: a
large anchor-private common contraction multiplies a lower-row role
difference, while both the changed-partner and common-block response
families at the anchor are subthreshold.  The next attack must return
to the exact connected V61 partition sum and decide whether its
four-row Möbius subtraction annihilates this private-anchor leaf or
forces an external occurrence to attach the anchor before the
absolute norm is taken.
\end{assessment}

\section{Regression checks and result ledger}
\label{sec:V26-results}

\begin{proposition}[V26 regression suite]
\label{prop:V26-regressions}
\begin{enumerate}[label=\textup{(\roman*)},leftmargin=4em]
\item A forward \(q\)-response is never reversed in the minimal
      active route.
\item The two-stage proof uses the exhaustive V12--V13 alternatives.
\item A response in a common moment is attached after one complete
      operator bound and is not renamed a covariance.
\item All occurrence and coordinate collisions are delegated only to
      the proved V25 collision-erasure theorem.
\item The private alternative has a fixed quantitative lower bound.
\item Every anchor-partner response family in the residual has a
      fixed quantitative upper bound.
\item The anchor-free role difference is retained inside the same
      exact multiplier.
\item No small scalar response is converted into a small operator
      norm.
\item Nonclean states remain separate.
\end{enumerate}
\end{proposition}

\begin{proof}
Items \textup{(i)--(ii)} are
\Cref{thm:V26-two-stage,cor:V26-no-reversal}.
Items \textup{(iii)--(iv)} are
\Cref{lem:V26-common-moment-attachment,
lem:V26-q-response-terminal}.  Items
\textup{(v)--(vii)} are
\Cref{lem:V26-common-exclusive,
def:V26-private-residual}.  Item \textup{(viii)} is
\Cref{fire:V26-private-no-edge}; item \textup{(ix)} is retained in
\Cref{prop:V26-complete-reduction}.
\end{proof}

\begin{theorem}[Third-series V26 result ledger]
\label{thm:V26-results}
Relative to V1--V25, V26 proves:
\begin{enumerate}[label=\textup{(V26.\arabic*)},leftmargin=6.3em]
\item exact depth-two termination of clean external/directed
      acquisition in the V100 first-breaker packet;
\item removal of the V13--V14 reverse-or-deplete branch from the
      minimal post-\(D\) proof;
\item an exclusive common-response versus private-anchor alternative
      inside the V2 low-response factor;
\item legal absolute attachment of a response carried by the common
      self-role moment;
\item termination of every nonprivate low-response branch; and
\item reduction of the entire clean acquisition tree to terminal
      packets plus the two explicitly quantified private-anchor
      residual types.
\end{enumerate}
\end{theorem}

\begin{proof}
These are
\Cref{thm:V26-two-stage,
cor:V26-no-reversal,
lem:V26-common-exclusive,
lem:V26-common-moment-attachment,
cor:V26-nonprivate-terminates,
prop:V26-complete-reduction}.
\end{proof}

\begin{programme}[V27: connected subtraction of the private anchor]
\label{prog:V26-V27}
\begin{enumerate}[leftmargin=3em]
\item Insert \({\mathfrak R}_{\rm priv}(v)\) into the exact V61
      connected partition-state sum before any norm.
\item Separate the terms in which all other occurrences refine
      \(\{v\}\mid(\cV\setminus\{v\})\) from those containing the
      first external attachment to \(v\).
\item Prove exact Möbius cancellation of the permanently isolated
      terms.
\item In the complementary terms, expose the unique first
      \(v\)-attaching occurrence and record its actual partner.
\item Test whether its complete bank supplies a dominant response or
      a depleted target identity without charging the private bank
      again.
\item Preserve the sensitive correction, lower-row role difference,
      and unique payer as parts of the same word.
\item Keep nonclean states and higher-order source families separate.
      The next compulsory companion audit is V30.
\end{enumerate}
\end{programme}

\begin{firewall}[Claim boundary after third-series V26]
\label{fire:V26-claim-boundary}
V26 proves finite termination of the clean response-search tree and
localizes its only unresolved clean leaf to the quantified
private-anchor residual.  It does not prove connected-cumulant
annihilation or external attachment for that residual and does not
handle nonclean states.  Higher MTA, WUFC, CPM, cutoff removal,
reflection positivity, Osterwalder--Schrader reconstruction,
construction of the continuum Yang--Mills measure, and a uniform
positive continuum spectral gap remain open.  The full Yang--Mills
mass gap has not been proved.
\end{firewall}

\paragraph{Parent-version record.}
V26 was copied byte-for-byte from
\texttt{Renewal\_Yang\_Mills\_Mass\_Gap\_Research\_Notes\_V25.tex}
before this chapter was appended.  The V25 parent had \(13\,162\)
source lines, \(479\,892\) bytes, and SHA--256
\[
 \texttt{2ca2ae97e55d9f2669e54326271b49b7ec8eba12d3cb267435033f68e3cde05a}.
\]
The next compulsory companion audit remains V30.

% =============================================================================
% BEGIN APPENDED THIRD-SERIES V27 MATERIAL
% =============================================================================

\clearpage
\chapter{V27: Physical-private refinement and the first anchor
attachment}
\label{ch:V27-private-attachment}

\section{Context: role-private is not physically private}
\label{sec:V27-context}

The V26 terminal cell \({\cal P}_v\) was called private because the
anchor block of the common self-role factor is the singleton
\(\{v\}\).  This is replica-role privacy, not necessarily physical
coordinate privacy.  Other rows may be active at the same coordinate
inside the complementary factor
\(R_{F_v}(\lambda)\).  The archived V68 private multiplier applies
only when no other row is physically active.

This distinction gives a new dichotomy.  If a fixed fraction of the
V26 private-anchor mass is physically shared, the Casimir floor
produces an actual response and V26 terminates.  If not, a fixed
fraction is genuinely physical-private.  That bank factors as one
exact Wilson multiplier from the complete V61 first-connectivity
family.  Its first moment pays the missing anchor denominator on the
first shared coordinate which attaches \(v\) to the other rows.

\section{Physical refinement of the role-private cell}
\label{sec:V27-refinement}

For a physical coordinate \(e\), write
\[
 A_e:=\{w\in\cV:a_{w,e}>0\}
\]
for its active-row set.  Inside the role-private cell
\({\cal P}_v\) of \Cref{def:V26-private-residual}, define
\begin{align}
 {\cal P}_v^{\rm phys}
 &:=
 \{e\in{\cal P}_v:A_e=\{v\}\},
\label{eq:V27-physical-private}\\
 {\cal S}_{vr}^{\rm phys}
 &:=
 \{e\in{\cal P}_v:|A_e|\ge2,\
 r=\min(A_e\setminus\{v\})\},
 \qquad r\ne v,
\label{eq:V27-physical-shared}
\end{align}
where the fixed global row order selects the least physical partner.
Put
\begin{align}
 p_v^{\rm phys}
 &:=
 \sum_{e\in{\cal P}_v^{\rm phys}}a_{v,e},
\label{eq:V27-private-mass}\\
 s_{vr}^{\rm phys}
 &:=
 \sum_{e\in{\cal S}_{vr}^{\rm phys}}a_{v,e}.
\label{eq:V27-shared-mass}
\end{align}

\begin{lemma}[Exact physical-private/shared identity]
\label{lem:V27-physical-split}
The four sets in
\eqref{eq:V27-physical-private}--%
\eqref{eq:V27-physical-shared} are pairwise disjoint and exhaust
\({\cal P}_v\).  Hence
\begin{equation}
 P_v^{\rm com}
 =
 p_v^{\rm phys}+\sum_{r\ne v}s_{vr}^{\rm phys}.
\label{eq:V27-physical-identity}
\end{equation}
Moreover,
\begin{equation}
 H_{vr}({\cal S}_{vr}^{\rm phys})
 \ge a_*s_{vr}^{\rm phys}.
\label{eq:V27-physical-floor}
\end{equation}
\end{lemma}

\begin{proof}
Every coordinate in \({\cal P}_v\) either has no other physically
active row or has a unique least active partner.  This proves the
partition and \eqref{eq:V27-physical-identity}.  On
\({\cal S}_{vr}^{\rm phys}\), row \(r\) is active, so
\(a_{r,e}\ge a_*\).  Thus
\(a_{v,e}a_{r,e}\ge a_*a_{v,e}\); summation gives
\eqref{eq:V27-physical-floor}.
\end{proof}

\begin{proposition}[Shared response or genuine private dominance]
\label{prop:V27-shared-or-private}
In the V26 private-anchor residual, exactly one of the following
threshold alternatives may be selected:
\begin{enumerate}[label=\textup{(\alph*)},leftmargin=4em]
\item for some \(r\ne v\),
      \begin{equation}
      d_{v\to r}({\cal S}_{vr}^{\rm phys})
      \ge\frac{\beta}{24}a_*;
      \label{eq:V27-shared-response}
      \end{equation}
\item the total physically shared \(v\)-mass in \({\cal P}_v\) is
      below \(\beta\Xi_v/8\), and
      \begin{equation}
      p_v^{\rm phys}>
      \frac{\beta}{8}\Xi_v.
      \label{eq:V27-physical-private-dominant}
      \end{equation}
\end{enumerate}
\end{proposition}

\begin{proof}
If
\(\sum_{r\ne v}s_{vr}^{\rm phys}\ge\beta\Xi_v/8\),
one of at most three partner cells has mass at least
\(\beta\Xi_v/24\).  Divide
\eqref{eq:V27-physical-floor} by \(\Xi_v\) to obtain
\eqref{eq:V27-shared-response}.  Select \textup{(a)} in this case.
Otherwise, use
\eqref{eq:V26-private-residual-mass} and
\eqref{eq:V27-physical-identity}:
\[
 p_v^{\rm phys}
 >
 \frac{\beta}{4}\Xi_v-\frac{\beta}{8}\Xi_v
 =
 \frac{\beta}{8}\Xi_v.
\]
This is \textup{(b)}, and the threshold convention makes the
selection disjoint.
\end{proof}

\begin{corollary}[The physically shared branch terminates]
\label{cor:V27-shared-terminates}
Alternative \eqref{eq:V27-shared-response} is a legal actual
response in the V26 two-stage tree even though the self-role
partition places \(v\) in a singleton block at that coordinate.
It therefore terminates as in
\Cref{cor:V26-nonprivate-terminates}.
\end{corollary}

\begin{proof}
The scalar stock in \eqref{eq:V27-shared-response} is a literal
physical overlap with the correct tail denominator.  Bound the
complete self-role packet once and attach this lower response using
\Cref{lem:V24-absolute-attachment}; no local replica block is
renamed as a covariance.  Apply
\Cref{lem:V26-x-response-fate} or
\Cref{lem:V26-q-response-terminal}.
\end{proof}

\begin{firewall}[Replica isolation does not erase physical overlap]
\label{fire:V27-role-physical}
The condition \(X_e=\{v\}\) says that the selected partition letter
does not join \(v\) to another row.  It does not say
\(a_{r,e}=0\) for every \(r\ne v\).  Conversely, only
\({\cal P}_v^{\rm phys}\) has the scalar packetwise multiplier of
the archived physical-private theorem.
\end{firewall}

\section{Exact multiplier and private moment}
\label{sec:V27-private-multiplier}

For \(e\in{\cal P}_v^{\rm phys}\), let \(q_{v,e}\) be the normalized
one-link Wilson expectation in the fixed physical input block, and
put
\begin{equation}
 q_{P_v}:=\prod_{e\in{\cal P}_v^{\rm phys}}q_{v,e}.
\label{eq:V27-private-product}
\end{equation}

\begin{lemma}[Exact physical-private multiplier on the regrouped
family]
\label{lem:V27-exact-private-factor}
Return from the selected self-role child to its complete V61
first-connectivity family before taking a positive envelope.  In
every fixed physical output block, that family has the exact factor
\begin{equation}
 {\cal S}=q_{P_v}{\cal S}^{\rm del},
\label{eq:V27-exact-private-factor}
\end{equation}
where \({\cal S}^{\rm del}\) is obtained by deleting the coordinates
in \({\cal P}_v^{\rm phys}\) and retaining the relative order of all
others.  No coordinate in this bank can be a first-connectivity or
first-\(v\)-attachment coordinate.
\end{lemma}

\begin{proof}
This is the archived packetwise physical-private theorem
\(\texttt{thm:V68-packetwise-private-factor}\), restricted to the
selected finite physical-private bank.  Its proof applies because at
each such coordinate only row \(v\) is nontrivial: the local letter
is the discrete partition letter and its normalized expectation is
the scalar \(q_{v,e}\), independently of the prefix state.  Hence
every word in the complete first-connectivity family contains the
same product once.  Factoring is performed on the whole family, as
required by the archived V68 orbit firewall, not on one separately
enveloped partition letter.
\end{proof}

\begin{lemma}[Private denominator compiler]
\label{lem:V27-private-denominator}
On \eqref{eq:V27-physical-private-dominant}, there is a constant
\(C_\beta\), independent of support and representations, such that
for every nonnegative cross-edge stock \(H\),
\begin{align}
 s_\alpha q_{P_v}H
 &\le C_\beta\frac{H}{\Xi_v},
\label{eq:V27-private-unpaid-edge}\\
 q_{P_v}H
 &\le \frac{C_\beta}{s_\alpha}\frac{H}{\Xi_v}.
\label{eq:V27-private-paid-edge}
\end{align}
More generally, for \(m=1,2,3\),
\begin{equation}
 (s_\alpha\Xi_v)^m q_{P_v}
 \le C_{\beta,m}.
\label{eq:V27-private-three-moments}
\end{equation}
\end{lemma}

\begin{proof}
The archived fixed-order private-product estimate gives
\[
 (s_\alpha p_v^{\rm phys})^m q_{P_v}\le C_m,
 \qquad m=1,2,3.
\]
By \eqref{eq:V27-physical-private-dominant},
\(\Xi_v<(8/\beta)p_v^{\rm phys}\), proving
\eqref{eq:V27-private-three-moments}.  Its \(m=1\) case is
\[
 s_\alpha\Xi_v q_{P_v}\le C_{\beta,1}.
\]
Multiply by \(H/\Xi_v\) to prove
\eqref{eq:V27-private-unpaid-edge}; divide it by \(s_\alpha\)
to prove \eqref{eq:V27-private-paid-edge}.
\end{proof}

\begin{firewall}[Damping is used before deleting the private bank]
\label{fire:V27-keep-private}
Although \({\cal S}^{\rm del}\) has no physical-private coordinates,
the scalar \(q_{P_v}\) in
\eqref{eq:V27-exact-private-factor} must remain attached until
\Cref{lem:V27-private-denominator} pays the missing row
normalization.  Replacing it by the bound \(q_{P_v}\le1\) would
destroy the only estimate available on the private-dominant branch.
\end{firewall}

\section{First attachment of a physically private-dominant row}
\label{sec:V27-first-attachment}

Let
\[
 \rho_v:=\{v\}\mid(\cV\setminus\{v\})
\]
and, at a nonprivate coordinate \(e\), define
\begin{equation}
 I_e^v:=M_e(\rho_v),
 \qquad
 J_e^v:=M_e-I_e^v.
\label{eq:V27-cut-letters}
\end{equation}
Thus \(I_e^v\) is the complete local moment with the two sides of the
cut independent, while \(J_e^v\) is their complete composite
covariance.

\begin{lemma}[Exact first-\(v\)-attachment telescope]
\label{lem:V27-first-attachment}
On any finite ordered nonprivate coordinate set \(E_{\rm sh}\),
\begin{equation}
 \bigotimes_{e\in E_{\rm sh}}M_e
 -
 \bigotimes_{e\in E_{\rm sh}}I_e^v
 =
 \sum_{f\in E_{\rm sh}}
 \left(\bigotimes_{e<f}I_e^v\right)
 \otimes J_f^v\otimes
 \left(\bigotimes_{e>f}M_e\right).
\label{eq:V27-first-attachment}
\end{equation}
After applying the connected four-row Möbius functional, the second
product on the left vanishes.  Hence every connected word is assigned
once to its first coordinate whose local partition joins \(v\) to
another row.
\end{lemma}

\begin{proof}
The displayed identity is the standard one-difference product
telescope, obtained by replacing the factors in increasing order.
The product of the \(I_e^v\)'s factors across the nontrivial cut
\(\rho_v\).  Its connected four-row cumulant is therefore zero by the
Möbius identity
\[
 \sum_{\tau:\,\sigma\le\tau\le\widehat1}
 \mu(\tau,\widehat1)=0
 \qquad(\sigma<\widehat1).
\]
Equivalently, in the V61 linked-word formula all of its local
partitions refine \(\rho_v\), so their join cannot be
\(\widehat1\).  In a nonzero word on the right, the prefix still
refines \(\rho_v\) and \(f\) is its unique first crossing.
\end{proof}

\begin{lemma}[Cross-edge envelope at the first attachment]
\label{lem:V27-first-attachment-bound}
Put
\begin{equation}
 H_{v|\bar v}(B)
 :=
 \sum_{r\ne v}H_{vr}(B).
\label{eq:V27-cut-stock}
\end{equation}
The complete first-attachment covariance on a bank \(B\) has the
inherited bounds
\begin{align}
 \|J_B^v\|_{\rm env}
 &\le Cs_\alpha H_{v|\bar v}(B),
\label{eq:V27-joining-unpaid}\\
 \|(J_B^v)^{\rm pay}\|_{\rm env}
 &\le C H_{v|\bar v}(B),
\label{eq:V27-joining-paid}
\end{align}
while every complete prefix and suffix factor in
\eqref{eq:V27-first-attachment} has a support-uniform ordinary bound.
If the payer lies outside \(J_B^v\), that one outside factor has its
usual \(C/s_\alpha\) bound.
\end{lemma}

\begin{proof}
\(J_B^v=M_B-M_B(\rho_v)\) is the complete covariance across the
two-block cut \(v|\bar v\).  The archived V65 complete composite
covariance theorem gives the unpaid minimum
\(C\min\{1,s_\alpha H_{v|\bar v}(B)\}\) and the paid minimum
\(C\min\{H_{v|\bar v}(B),s_\alpha^{-1}\}\), which imply the two
displayed bounds.  Complete \(I^v\)-prefix moments and complete
suffix moments are contractions.  The fixed connected Möbius
functional and the finite V61 separators change only the universal
constant.  The final statement is the established one-payer
allocation.
\end{proof}

\begin{theorem}[Private-normalized first-attachment response]
\label{thm:V27-private-first-response}
On the physical-private alternative
\eqref{eq:V27-physical-private-dominant}, the complete regrouped
first-attachment packet satisfies
\begin{align}
 \Gamma_\otimes(1,{\cal S})
 &\le
 C_\beta\sum_{r\ne v}
 d_{v\to r}(B_{\rm att}),
\label{eq:V27-first-response-unpaid}\\
 \Gamma_\otimes(w,{\cal S}^{\rm pay})
 &\le
 \frac{C_\beta}{s_\alpha}
 \sum_{r\ne v}
 d_{v\to r}(B_{\rm att})
\label{eq:V27-first-response-paid}
\end{align}
up to the already retained scalar factors of the surrounding
first-breaker packet.  Here \(B_{\rm att}\) is the complete
first-attachment bank, and the sum over its three actual partners has
no support-cardinality factor.  No dominance lower bound on these
responses is assumed.
\end{theorem}

\begin{proof}
Factor \(q_{P_v}\) from the complete family by
\Cref{lem:V27-exact-private-factor}, then apply the exact telescope
\eqref{eq:V27-first-attachment}.  If the joining covariance is
unpaid, combine \eqref{eq:V27-joining-unpaid} with
\eqref{eq:V27-private-unpaid-edge}.  If it pays, combine
\eqref{eq:V27-joining-paid} with
\eqref{eq:V27-private-paid-edge}.  If the payer is outside, its
\(s_\alpha^{-1}\) bound multiplies the unpaid joining line and again
gives \eqref{eq:V27-first-response-paid}.

Finally use
\[
 \frac{H_{v|\bar v}(B_{\rm att})}{\Xi_v}
 =
 \sum_{r\ne v}
 \frac{H_{vr}(B_{\rm att})}{\Xi_v}.
\]
The partner split is a three-term row identity, not a coordinate
selection, so its constant is support-uniform.
\end{proof}

\begin{firewall}[An upper response needs no dominance]
\label{fire:V27-upper-response}
The responses on the right of
\eqref{eq:V27-first-response-unpaid} are supplied as upper
numerator factors by the joining covariance and private damping.
They are not inserted by multiplying with a scalar at least one.
Therefore they may be arbitrarily small.  Only any additional
responses attached from other parent roles must satisfy their stated
dominance lower bounds.
\end{firewall}

\section{Closure of all single-private leaves}
\label{sec:V27-single-private}

\begin{lemma}[Mixed upper/lower tree attachment]
\label{lem:V27-mixed-attachment}
Suppose a complete packet has already been bounded by
\[
 C\,d_R
\]
or \(Cs_\alpha^{-1}d_R\), where \(d_R\) is one correctly normalized
actual directed edge.  If the remaining one or two edge roles of a
legal V15 tree have dominant lower bounds, they may be attached after
this estimate by the corresponding factors at least one.  The result
is the legal product of all tree responses with the same payer
count.
\end{lemma}

\begin{proof}
Multiply the established upper bound by
\(d_T/(\gamma_Ta_*)\ge1\) for each remaining role \(T\).
The operation is scalar and occurs after the analytic estimate.
The V15 table verifies the edge labels and tail denominators.
\end{proof}

\begin{proposition}[Single-private terminal cases]
\label{prop:V27-single-private-cases}
The following physical-private residuals close:
\begin{enumerate}[label=\textup{(\roman*)},leftmargin=4em]
\item a \(q\)-residual reached after a dominant \(D\)-response;
\item an \(x\)-residual whose first attached partner is \(q\);
\item an \(x\)-residual whose first attached partner is \(y\) or
      \(z\), provided the following \(q\)-stage produces a dominant
      response or a compiler-terminal branch.
\end{enumerate}
\end{proposition}

\begin{proof}
For \textup{(i)}, apply
\Cref{thm:V27-private-first-response} with \(v=q\), split the three
partners, and attach the already dominant \(D,S\) responses using
\Cref{lem:V27-mixed-attachment}.  Every resulting product is a V15
tree.

For \textup{(ii)}, the upper response is \(x\to q\).  Retain the
V100 raw breaker estimate and attach the dominant
\(S=y\to z\) response.  The three edges are the path
\(q-x-y-z\), with the correct internal tail denominators.

For \textup{(iii)}, the upper response is
\(D=x\to y\) or \(x\to z\).  A dominant \(q\)-response is \(Q\);
attach it and \(S\) by
\Cref{lem:V27-mixed-attachment}.  The complete-compiler alternatives
retain their established terminal bounds.
\end{proof}

\begin{proposition}[Distinct double-private attachments close]
\label{prop:V27-distinct-double-private}
Suppose an \(x\)-private first attachment supplies an upper
\(D\)-response and the subsequent \(q\)-stage is also
physical-private.  If the two first-attachment covariances occupy
distinct physical occurrences in the exact word, the packet satisfies
the legal V15 tree bound and closes.
\end{proposition}

\begin{proof}
Keep the two physical-private multipliers attached to their disjoint
row-private banks.  Apply the two distinct complete covariance
estimates before taking the final product-state envelope.  Unpaid,
their scalar factor is
\[
 (s_\alpha q_{P_x}H_D)
 (s_\alpha q_{P_q}H_Q)
 \le C_{\beta_x,\beta_q}d_Dd_Q
\]
by \Cref{lem:V27-private-denominator} on the two rows.  If either
covariance pays, replace its first parenthesis by the paid line; if an
outside factor pays, multiply the unpaid product by its single
\(s_\alpha^{-1}\) bound.  In every case the paid scale is
\(C s_\alpha^{-1}d_Dd_Q\).

Attach the dominant \(S\)-response by
\Cref{lem:V27-mixed-attachment}.  The V15 table gives the legal
tree.  Distinctness of the occurrences is used only to apply the two
analytic covariance bounds; no coordinate disjointness is assumed.
\end{proof}

\begin{firewall}[The same-occurrence double upper mark remains]
\label{fire:V27-double-private-collision}
If the two first-attachment roles \(D\) and \(Q\) lie in one
unrepeated physical covariance occurrence, the V16 telescope supplies
only one covariance difference per child.  Unlike V25, neither
\(d_D\) nor \(d_Q\) is known to be bounded below, so absolute
parent-response attachment cannot create the missing upper factor.
\Cref{prop:V27-distinct-double-private} therefore does not cover this
same-occurrence double-private carrier.
\end{firewall}

\begin{theorem}[Exact clean residual after V27]
\label{thm:V27-exact-residual}
Every clean first-breaker packet is closed except possibly a
double-private attachment packet with the following simultaneous
properties:
\begin{enumerate}[label=\textup{(\alph*)},leftmargin=4em]
\item the V26 low-response residual occurs first at \(x\) and then at
      \(q\);
\item both role-private cells are physical-private-dominant;
\item the first \(x\)-attachment has partner \(y\) or \(z\);
\item the two resulting upper response roles \(D,Q\) have the same
      unrepeated physical covariance parent; and
\item the dominant \(S=y\to z\) response and unique payer remain in
      the complete word.
\end{enumerate}
Nonclean states remain a separate residual.
\end{theorem}

\begin{proof}
Apply \Cref{prop:V27-shared-or-private} to each V26
private-anchor leaf.  The physically shared alternatives terminate
by \Cref{cor:V27-shared-terminates}.  For a physical-private
alternative, apply
\Cref{prop:V27-single-private-cases}.  Only an \(x\)-attachment with
avoid-\(q\) partner can reach a second private leaf at \(q\).
\Cref{prop:V27-distinct-double-private} closes it when the attachment
occurrences differ.  The complement is exactly
\textup{(a)--(e)}.  The V26 nonclean alternative was never included
in this clean analysis.
\end{proof}

\begin{assessment}[Bottleneck after V27]
\label{ass:V27-bottleneck}
The former large private-anchor family has collapsed to one joint
carrier: two non-dominant response numerators, each normalized by a
different physical-private multiplier, meet inside one unrepeated
covariance occurrence.  The correct next object is not another
response search.  It is a two-row-private, one-covariance
first-attachment estimate formed before the V16 support telescope.
Its target is an upper product \(d_Dd_Q\), with the sole payer on
either private bank, the common covariance, or the surrounding word.
\end{assessment}

\section{Regression checks and result ledger}
\label{sec:V27-results}

\begin{proposition}[V27 regression suite]
\label{prop:V27-regressions}
\begin{enumerate}[label=\textup{(\roman*)},leftmargin=4em]
\item Role-private and physical-private coordinates are distinguished.
\item The physical shared branch uses actual row activity and the
      Casimir floor.
\item A physical-private multiplier is factored from the complete
      V61 family, not a single orbit envelope.
\item Its damping is retained until the missing row denominator is
      paid.
\item Permanently \(v\)-disconnected words vanish by the exact
      connected Möbius identity.
\item First-attachment grouping has one joining covariance and a
      complete suffix.
\item An upper response is never assumed dominant.
\item Two covariance estimates are multiplied only for distinct
      physical occurrences.
\item The same-occurrence double-private carrier remains explicit.
\end{enumerate}
\end{proposition}

\begin{proof}
Items \textup{(i)--(ii)} are
\Cref{lem:V27-physical-split,
prop:V27-shared-or-private,
fire:V27-role-physical}.  Items \textup{(iii)--(iv)} are
\Cref{lem:V27-exact-private-factor,
lem:V27-private-denominator,
fire:V27-keep-private}.  Items \textup{(v)--(vi)} are
\Cref{lem:V27-first-attachment,
lem:V27-first-attachment-bound}.  Item \textup{(vii)} is
\Cref{fire:V27-upper-response}.  Items
\textup{(viii)--(ix)} are
\Cref{prop:V27-distinct-double-private,
fire:V27-double-private-collision}.
\end{proof}

\begin{theorem}[Third-series V27 result ledger]
\label{thm:V27-results}
Relative to V1--V26, V27 proves:
\begin{enumerate}[label=\textup{(V27.\arabic*)},leftmargin=6.3em]
\item an exact physical refinement of the V26 role-private cell;
\item a dominant shared response or genuine
      physical-private-dominance alternative;
\item exact factoring of the archived V68 private multiplier from the
      complete regrouped family;
\item private moment payment of a missing anchor denominator;
\item an exact first-anchor-attachment telescope and its composite
      covariance envelope;
\item a normalized upper response estimate requiring no dominance;
\item closure of every single-private and distinct double-private
      leaf; and
\item reduction of the remaining clean sector to one
      same-occurrence double-private carrier.
\end{enumerate}
\end{theorem}

\begin{proof}
These are
\Cref{lem:V27-physical-split,
prop:V27-shared-or-private,
lem:V27-exact-private-factor,
lem:V27-private-denominator,
lem:V27-first-attachment,
thm:V27-private-first-response,
prop:V27-single-private-cases,
prop:V27-distinct-double-private,
thm:V27-exact-residual}.
\end{proof}

\begin{programme}[V28: the common-covariance two-private sandwich]
\label{prog:V27-V28}
\begin{enumerate}[leftmargin=3em]
\item Keep the two physical-private multipliers
      \(q_{P_x},q_{P_q}\) outside the exact common covariance before
      any support telescope.
\item Partition the common covariance coordinates by the two actual
      response edge labels \(D\) and \(Q\), retaining overlap and
      exclusive cells in one complete two-cut covariance.
\item Seek the joint upper estimate
      \[
      q_{P_x}q_{P_q}\,
      \Gamma_\otimes(\Delta_C)
      \lesssim d_Dd_Q
      \]
      unpaid and \(s_\alpha^{-1}d_Dd_Q\) once-paid.
\item Test the empty-overlap and one-coordinate diagonal cases
      separately; do not infer a common coordinate from the two
      attachment roles.
\item Use the two private moment inequalities only after the complete
      two-cut carrier is bounded.
\item If the product estimate fails, identify whether one covariance
      supplies only a sum \(d_D+d_Q\); then return to its exact local
      block-product covariance rather than cloning it.
\item Keep the dominant \(S\)-response, outer word, and unique payer
      attached.  The next compulsory companion audit is V30.
\end{enumerate}
\end{programme}

\begin{firewall}[Claim boundary after third-series V27]
\label{fire:V27-claim-boundary}
V27 closes every clean first-breaker branch except the explicitly
identified same-occurrence double-private carrier.  It does not prove
the required joint \(d_Dd_Q\) upper estimate for that carrier and does
not handle nonclean states.  Higher MTA, WUFC, CPM, cutoff removal,
reflection positivity, Osterwalder--Schrader reconstruction,
construction of the continuum Yang--Mills measure, and a uniform
positive continuum spectral gap remain open.  The full Yang--Mills
mass gap has not been proved.
\end{firewall}

\paragraph{Parent-version record.}
V27 was copied byte-for-byte from
\texttt{Renewal\_Yang\_Mills\_Mass\_Gap\_Research\_Notes\_V26.tex}
before this chapter was appended.  The V26 parent had \(13\,637\)
source lines, \(497\,989\) bytes, and SHA--256
\[
 \texttt{2f6f7a9e0c4d4cc4abd6417dbfb124d358c26e5060db9867aaafb1faa7e1142e}.
\]
The next compulsory companion audit remains V30.

% =============================================================================
% BEGIN APPENDED THIRD-SERIES V28 MATERIAL
% =============================================================================

\clearpage
\chapter{V28: Rank-two quotient joining in the double-private
collision}
\label{ch:V28-rank-two}

\section{Context: one linear covariance cannot pay two upper marks}
\label{sec:V28-context}

V27 leaves one clean carrier in which physical-private damping at
rows \(x\) and \(q\) must normalize two non-dominant upper responses
\(D\) and \(Q\), while their ancestry points to the same unrepeated
physical parent.  The proposed bound
\[
 q_{P_x}q_{P_q}\,
 \Gamma_\otimes(\Delta_C)
 \lesssim d_Dd_Q
\]
cannot follow from the ordinary two-block covariance estimate.  That
estimate is linear in the total crossing stock, whereas the target is
quadratic in two independently variable edge stocks.

The correct upstream object is selected before the parent is
decomposed into cover covariances.  Immediately before a local letter
first attaches both private-dominant rows, \(x\) and \(q\) are
singleton quotient blocks.  A letter carrying the two distinct
response roles reduces the number of quotient components by at least
two.  It therefore belongs to a ternary block-product cumulant or to
a rank-two part of the four-block joining carrier, not to one
rank-one cover covariance.

\section{A sharp abstract no-go for the linear route}
\label{sec:V28-no-go}

\begin{proposition}[A one-covariance bound does not imply a response
product]
\label{prop:V28-linear-no-go}
There is no universal constant \(C\) for which the data
\begin{equation}
 0\le G\le H_D+H_Q,
 \qquad
 0<d_D=H_D,\quad 0<d_Q=H_Q
\label{eq:V28-linear-data}
\end{equation}
alone imply
\begin{equation}
 G\le C d_Dd_Q.
\label{eq:V28-false-product}
\end{equation}
The failure persists after multiplication by two uniformly bounded
private factors satisfying every fixed first-moment upper bound.
\end{proposition}

\begin{proof}
For \(0<\varepsilon<1\), take
\[
 H_D=H_Q=\varepsilon,
 \qquad
 d_D=d_Q=\varepsilon,
 \qquad
 G=\varepsilon.
\]
Then \eqref{eq:V28-linear-data} holds, while
\eqref{eq:V28-false-product} would require
\(\varepsilon\le C\varepsilon^2\), which fails as
\(\varepsilon\downarrow0\).

Set \(s_\alpha=\Xi_x=\Xi_q=1\) and take both abstract private
multipliers equal to one.  Then
\((s_\alpha\Xi_v)^m q_{P_v}=1\) for every fixed \(m\), so all private
moment upper bounds hold.  Multiplying \(G\) and the target by these
factors does not change the contradiction.
\end{proof}

\begin{firewall}[Scope of the no-go]
\label{fire:V28-no-go-scope}
\Cref{prop:V28-linear-no-go} is not a Wilson-law counterexample and
does not assert that the complete connected double-private carrier is
linear.  It proves that the V65 two-block covariance bound, private
moment upper bounds, and positivity do not by themselves contain the
needed second crossing stock.  A valid proof must retain a
rank-two connected joining carrier or another genuinely quadratic
Wilson identity.
\end{firewall}

\begin{corollary}[The V16 one-difference telescope is too fine here]
\label{cor:V28-V16-too-fine}
Expanding the common parent by the V16 support telescope and bounding
each child only by its one covariance difference cannot prove the
double-private \(d_Dd_Q\) target.  The children must first be
reassembled into the connected rank-two quotient carrier.
\end{corollary}

\begin{proof}
Each V16 child has one covariance difference and hence only a linear
crossing-stock estimate of the form tested in
\Cref{prop:V28-linear-no-go}.  A triangle inequality before
reassembly retains only a sum of such linear bounds.
\end{proof}

\section{The prefix quotient has only two types}
\label{sec:V28-prefix}

Let \(c\) denote the common physical parent occurrence from
\Cref{thm:V27-exact-residual}.  Return to its exact V61 partition
word and let
\[
 \rho^-:=\bigvee_{e<c}\pi_e
\]
be the cumulative row partition immediately before the first local
letter in that occurrence which attaches either private-dominant
anchor.  The ordering is the inherited physical coordinate order
inside the complete first-connectivity family.

\begin{lemma}[Two possible double-private prefix joins]
\label{lem:V28-two-prefixes}
If the same occurrence is the first physical attachment of both
\(x\) and \(q\), then
\begin{equation}
 \rho^-\in
 \left\{
 \rho_4:=x|q|y|z,\quad
 \rho_3:=x|q|yz
 \right\}.
\label{eq:V28-two-prefixes}
\end{equation}
\end{lemma}

\begin{proof}
Before their first attachments, \(x\) and \(q\) are singleton blocks
of the cumulative join.  The remaining set is \(\{y,z\}\), whose
induced partition is either \(y|z\) or \(yz\).  These are exactly the
two partitions displayed in \eqref{eq:V28-two-prefixes}.
\end{proof}

\begin{definition}[Quotient rank]
\label{def:V28-quotient-rank}
For \(\rho,\pi\in\Pi(\cV)\), define the rank gained by adding the
local letter \(\pi\) to the prefix \(\rho\) by
\begin{equation}
 {\rm rk}_\rho(\pi):=
 b(\rho)-b(\rho\vee\pi),
\label{eq:V28-rank}
\end{equation}
where \(b(\sigma)\) is the number of blocks of a partition
\(\sigma\).
\end{definition}

\begin{lemma}[Two distinct anchor attachments have rank at least two]
\label{lem:V28-rank-two}
Let \(\rho\) be either partition in
\eqref{eq:V28-two-prefixes}.  Suppose a local partition letter
\pi supplies:
\begin{enumerate}[label=\textup{(\roman*)},leftmargin=4em]
\item an actual \(D\)-edge joining \(x\) to \(y\) or \(z\); and
\item a distinct actual \(Q\)-edge joining \(q\) to one of
      \(x,y,z\).
\end{enumerate}
Then
\begin{equation}
 {\rm rk}_\rho(\pi)\ge2.
\label{eq:V28-rank-two}
\end{equation}
\end{lemma}

\begin{proof}
Pass to the quotient graph whose vertices are the blocks of \(\rho\)
and whose edges join quotient blocks met by a block of \(\pi\).
The \(D\)-edge joins the singleton \(x\) to a block containing
\(y\) or \(z\).  The distinct \(Q\)-edge is incident to the
singleton \(q\).  If it joins \(q\) to the same non-\(x\) quotient
component, the three vertices \(x,q,\bar r\) merge and the block count
drops by two.  If it joins \(q\) to \(x\), the \(D\)-edge then merges
the resulting \(xq\) block with the \(y/z\) block, again dropping the
count by two.  If its other endpoint is a different quotient block,
two disjoint merges occur.  In every case the join loses at least two
blocks, proving \eqref{eq:V28-rank-two}.
\end{proof}

\begin{corollary}[A cover increment cannot be the double attachment]
\label{cor:V28-no-cover}
A partition cover of either prefix in
\eqref{eq:V28-two-prefixes} has quotient rank one and therefore
cannot carry both distinct \(D\) and \(Q\) roles of the V27 residual.
\end{corollary}

\begin{proof}
A cover merges exactly two blocks, so its rank is one.  Apply
\Cref{lem:V28-rank-two}.
\end{proof}

\begin{firewall}[Parent occurrence versus cover child]
\label{fire:V28-parent-cover}
The V20 ancestry tag may assign \(D,Q\) to the same physical parent
even though a cover decomposition writes that parent as a signed sum
of rank-one covariance children.  No one child then contains the
double attachment.  The quadratic contribution, if present, lives in
the connected reassembly of those children; it must not be sought by
duplicating one child.
\end{firewall}

\section{The exact rank-two carriers}
\label{sec:V28-carriers}

\begin{theorem}[Three-block double attachment is an assembled ternary
cumulant]
\label{thm:V28-three-block-carrier}
For the prefix
\[
 \rho_3=x|q|yz,
\]
let \(X_x,X_q,Z_{yz}\) denote the local variables of its three
quotient blocks.  The complete local carrier which makes the
cumulative join connected is
\begin{equation}
 \boxed{
 J_c(\rho_3)
 =
 \sum_{\pi:\,\rho_3\vee\pi=\widehat1}L_{c,\pi}
 =
 \kappa_3(X_x,X_q,Z_{yz}).}
\label{eq:V28-ternary-carrier}
\end{equation}
Every contributing partition has quotient rank two.  The sum must be
formed before a positive envelope.
\end{theorem}

\begin{proof}
The first equality is the definition of the V61 joining letter.
The archived assembled joining-letter identity
\(\texttt{thm:V63-joining-cumulant}\) identifies it with the
connected cumulant of the prefix block-products, giving the second
equality.  Since \(b(\rho_3)=3\) and
\(\rho_3\vee\pi=\widehat1\), every contributing term has
\({\rm rk}_{\rho_3}(\pi)=2\).  The cumulant identity itself contains
the cancellation among its moment partitions, so a triangle
inequality may be applied only after the displayed sum is formed.
\end{proof}

For the discrete prefix \(\rho_4=x|q|y|z\), let
\({\cal A}_{xq}\subseteq\Pi(\cV)\) be the set of partitions
\(\pi\) satisfying:
\begin{enumerate}[label=\textup{(\alph*)},leftmargin=4em]
\item the block of \(x\) meets \(\{y,z\}\);
\item \(q\) is not a singleton block; and
\item \({\rm rk}_{\rho_4}(\pi)\ge2\).
\end{enumerate}
Define the assembled local double-attachment letter
\begin{equation}
 J_c^{xq,(2+)}
 :=
 \sum_{\pi\in{\cal A}_{xq}}L_{c,\pi}.
\label{eq:V28-four-block-carrier}
\end{equation}

\begin{proposition}[Exhaustive four-block rank-two carrier]
\label{prop:V28-four-block-carrier}
Every local partition letter at the common occurrence which supplies
the distinct roles \(D=x\to y\) or \(x\to z\) and
\(Q=q\to r\) belongs to \({\cal A}_{xq}\).
Conversely, every letter in \({\cal A}_{xq}\) has a quotient graph
with an \(x\)-to-\(\{y,z\}\) edge and an additional edge incident to
\(q\); after choosing actual partners it supplies one of the six V15
\((D,Q)\) pairs.  The number of partner routings is universal.
\end{proposition}

\begin{proof}
The first direction is
\Cref{lem:V28-rank-two}: the \(D\)-edge gives condition
\textup{(a)}, the \(Q\)-edge gives \textup{(b)}, and their
distinctness gives \textup{(c)}.  Conversely, condition
\textup{(a)} supplies a block path from \(x\) to \(y\) or \(z\).
Choose its first row edge as \(D\).  Condition \textup{(b)} supplies
an edge from \(q\) to another row.  If the first such edge coincides
with a previously selected quotient connection, choose the next edge
in a spanning forest of the block; the rank condition guarantees two
independent quotient mergers.  The resulting undirected labels are
distinct.  There are four rows and fifteen partitions, so the choices
have a fixed finite routing multiplicity.
\end{proof}

\begin{remark}[The familiar local types]
\label{rem:V28-four-local-types}
The rank-two part contains the \(2+2\) letters such as
\(xy|qz\), the \(3+1\) letters in which the triple contains
appropriate \(D,Q\) paths, and the one-block fourth cumulant letter.
Rank-one \(2+1+1\) covers are excluded.  Keeping their permitted sum
in \eqref{eq:V28-four-block-carrier} avoids assigning an independent
positive norm to each internal moment partition.
\end{remark}

\begin{proposition}[Exact residual carrier dichotomy]
\label{prop:V28-carrier-dichotomy}
The same-occurrence double-private residual of
\Cref{thm:V27-exact-residual}, after undoing cover expansion and
before positive envelopes, is a finite sum of:
\begin{enumerate}[label=\textup{(\roman*)},leftmargin=4em]
\item a three-block packet with the complete local carrier
      \(J_c(\rho_3)=\kappa_3(X_x,X_q,Z_{yz})\); and
\item a discrete-prefix packet with the assembled local letter
      \(J_c^{xq,(2+)}\), followed when necessary by the first suffix
      occurrence which completes the \(y,z\) connection.
\end{enumerate}
Both packets retain the two physical-private multipliers, the
dominant \(S\)-response, and the unique payer.
\end{proposition}

\begin{proof}
\Cref{lem:V28-two-prefixes} is exhaustive.  For \(\rho_3\), the
common occurrence makes the join connected and
\Cref{thm:V28-three-block-carrier} gives \textup{(i)}.  For
\(\rho_4\), collect exactly the rank-two letters identified in
\Cref{prop:V28-four-block-carrier}.  If such a letter does not make
the join one-block, the V61 first-connectivity formula retains the
complete suffix until its unique first completing occurrence; group
by that occurrence.  This is \textup{(ii)}.  The two prefix joins
are disjoint.  The private factors were extracted packetwise before
this split, and the V2 payer label stays in its original factor.
\end{proof}

\section{The correct two-edge quotient estimate}
\label{sec:V28-target}

For a common bank \(C\), define the legal two-response numerator
\begin{equation}
 {\cal H}_{DQ}(C)
 :=
 \sum_{\substack{u\in\{y,z\}\\r\in\{x,y,z\}}}
 H_{xu}(C)H_{qr}(C),
\label{eq:V28-DQ-stock}
\end{equation}
with any terms absent from the exact partner routing omitted.

\begin{definition}[Two-edge quotient-joining estimate]
\label{def:V28-QJ2}
The \(QJ_2\) estimate for either carrier in
\Cref{prop:V28-carrier-dichotomy} is
\begin{align}
 \Gamma_\otimes(1,J_C^{DQ})
 &\le C s_\alpha^2{\cal H}_{DQ}(C),
\label{eq:V28-QJ2-unpaid}\\
 \Gamma_\otimes(w,(J_C^{DQ})^{\rm pay})
 &\le C s_\alpha{\cal H}_{DQ}(C)
\label{eq:V28-QJ2-paid}
\end{align}
at the local payer, together with the same second line when the sole
payer lies in a fixed surrounding factor.  The carrier
\(J_C^{DQ}\) means the complete ternary or rank-two four-block packet,
not one cover covariance.
\end{definition}

\begin{theorem}[\(QJ_2\) closes the same-occurrence double-private
residual]
\label{thm:V28-QJ2-suffices}
If \eqref{eq:V28-QJ2-unpaid}--%
\eqref{eq:V28-QJ2-paid} hold uniformly for both carriers in
\Cref{prop:V28-carrier-dichotomy}, then the clean residual of
\Cref{thm:V27-exact-residual} satisfies the legal unpaid and
once-paid V15 tree bound and closes.
\end{theorem}

\begin{proof}
Retain the exact multipliers \(q_{P_x}q_{P_q}\).  A summand of
\eqref{eq:V28-DQ-stock} obeys
\begin{align*}
 q_{P_x}q_{P_q}s_\alpha^2
 H_{xu}(C)H_{qr}(C)
 &=
 \bigl(s_\alpha q_{P_x}H_{xu}(C)\bigr)
 \bigl(s_\alpha q_{P_q}H_{qr}(C)\bigr)\\
 &\le C\,d_{x\to u}(C)d_{q\to r}(C)
\end{align*}
by \Cref{lem:V27-private-denominator}.  On the paid line,
write the factor \(s_\alpha H_{xu}H_{qr}\) as one paid private-edge
line times one unpaid line.  This gives
\[
 q_{P_x}q_{P_q}s_\alpha
 H_{xu}H_{qr}
 \le\frac{C}{s_\alpha}
 d_{x\to u}d_{q\to r}.
\]
An outside payer gives the same scale by multiplying the unpaid line
by its unique \(s_\alpha^{-1}\) factor.

Finally attach the dominant \(S=y\to z\) response by
\Cref{lem:V27-mixed-attachment}.  Every routed pair \((u,r)\) is one
row of the V15 table, so the resulting product is a legal tree.
There are only six routings and two prefix types.
\end{proof}

\begin{firewall}[Why a product of total cut stocks is too coarse]
\label{fire:V28-cut-product}
The expression
\[
 H_{x|\bar x}(C)H_{q|\bar q}(C)
\]
contains repeated or incompatible edge choices, including choices
which do not match the actual local partition path.  The target
\eqref{eq:V28-DQ-stock} is formed only after the exact carrier supplies
actual partners.  A proof of \(QJ_2\) must preserve that routing and
may dominate it by the total cut product only at the final
nonnegative scalar step.
\end{firewall}

\section{Regression checks and result ledger}
\label{sec:V28-results}

\begin{proposition}[V28 regression suite]
\label{prop:V28-regressions}
\begin{enumerate}[label=\textup{(\roman*)},leftmargin=4em]
\item A linear covariance estimate is not promoted to a quadratic
      response product.
\item Cover children are reassembled before a positive envelope.
\item The exact prefix has singleton \(x,q\) blocks.
\item Two distinct response roles require quotient rank at least two.
\item The three-block carrier is the complete block-product
      cumulant.
\item The four-block carrier excludes rank-one covers and preserves
      the finite permitted sum.
\item Both physical-private multipliers and the unique payer remain
      attached.
\item \(QJ_2\) uses two crossing-stock powers and exactly one fewer
      \(s_\alpha\) on a paid line.
\item The full Yang--Mills conclusion is not inferred from the
      conditional \(QJ_2\) criterion.
\end{enumerate}
\end{proposition}

\begin{proof}
Items \textup{(i)--(ii)} are
\Cref{prop:V28-linear-no-go,
cor:V28-V16-too-fine,
fire:V28-parent-cover}.  Items \textup{(iii)--(iv)} are
\Cref{lem:V28-two-prefixes,lem:V28-rank-two}.  Items
\textup{(v)--(vi)} are
\Cref{thm:V28-three-block-carrier,
prop:V28-four-block-carrier}.  Item \textup{(vii)} is
\Cref{prop:V28-carrier-dichotomy}.  Item \textup{(viii)} is
\Cref{def:V28-QJ2,thm:V28-QJ2-suffices}; item
\textup{(ix)} is the claim boundary below.
\end{proof}

\begin{theorem}[Third-series V28 result ledger]
\label{thm:V28-results}
Relative to V1--V27, V28 proves:
\begin{enumerate}[label=\textup{(V28.\arabic*)},leftmargin=6.3em]
\item a sharp abstract no-go for deriving two upper response marks
      from one linear covariance bound;
\item the exhaustive two-type prefix classification of the
      double-private same-occurrence residual;
\item a quotient-rank-two necessity theorem;
\item identification of the three-block carrier with one assembled
      ternary block-product cumulant;
\item an exact assembled rank-two carrier for the discrete prefix;
\item reassembly of the V27 residual into those two carriers; and
\item a precise unpaid/once-paid \(QJ_2\) estimate sufficient to
      close the last clean first-breaker packet.
\end{enumerate}
\end{theorem}

\begin{proof}
These are
\Cref{prop:V28-linear-no-go,
lem:V28-two-prefixes,
lem:V28-rank-two,
thm:V28-three-block-carrier,
prop:V28-four-block-carrier,
prop:V28-carrier-dichotomy,
def:V28-QJ2,
thm:V28-QJ2-suffices}.
\end{proof}

\begin{assessment}[Bottleneck after V28]
\label{ass:V28-bottleneck}
The last clean first-breaker obstruction is no longer described as
two marks from one covariance.  It is the Wilson-specific
rank-two quotient estimate \(QJ_2\).  The smallest case is the
three-block carrier
\(\kappa_3(X_x,X_q,Z_{yz})\), for which the desired two crossing
stocks are exactly the two edges of a tree on three quotient blocks.
V29 should compare this carrier directly with the archived normalized
ternary marked-tree theorem while retaining the composite \(yz\)
block and both private multipliers.
\end{assessment}

\begin{programme}[V29: the ternary quotient-tree interface]
\label{prog:V28-V29}
\begin{enumerate}[leftmargin=3em]
\item Fix \(\rho_3=x|q|yz\) and keep
      \(\kappa_3(X_x,X_q,Z_{yz})\) assembled.
\item Treat \(Z_{yz}\) as one quotient variable until the ternary
      capacity and payer have been allocated.
\item Expand each quotient edge only afterward into its actual row
      partner and verify the six V15 routes.
\item Prove the two-crossing-stock powers
      \(s_\alpha^2{\cal H}_{DQ}\) unpaid and
      \(s_\alpha{\cal H}_{DQ}\) once-paid, or identify the exact
      duplicated row/output mass.
\item Test coincident physical coordinates without applying two
      independent first-moment cards.
\item Then attack the discrete-prefix rank-two carrier with the
      successful interface or record why its fourth-order local
      cancellation is different.
\item Keep nonclean states separate.  The next compulsory companion
      audit is V30.
\end{enumerate}
\end{programme}

\begin{firewall}[Claim boundary after third-series V28]
\label{fire:V28-claim-boundary}
V28 proves the exact rank-two carrier reduction and the sufficiency
of \(QJ_2\), but it does not prove either \(QJ_2\) estimate.  Hence
the same-occurrence double-private clean packet and nonclean states
remain open.  Higher MTA, WUFC, CPM, cutoff removal, reflection
positivity, Osterwalder--Schrader reconstruction, construction of the
continuum Yang--Mills measure, and a uniform positive continuum
spectral gap remain open.  The full Yang--Mills mass gap has not been
proved.
\end{firewall}

\paragraph{Parent-version record.}
V28 was copied byte-for-byte from
\texttt{Renewal\_Yang\_Mills\_Mass\_Gap\_Research\_Notes\_V27.tex}
before this chapter was appended.  The V27 parent had \(14\,302\)
source lines, \(522\,415\) bytes, and SHA--256
\[
 \texttt{c52555aa304c2587c4f72f8bdf8a7bf8a1dfc7381e127a13e0281b6f0818aefd}.
\]
The next compulsory companion audit remains V30.

% =============================================================================
% BEGIN APPENDED THIRD-SERIES V29 MATERIAL
% =============================================================================

\clearpage
\chapter{V29: The composite-ternary quotient tree and paid private
moments}
\label{ch:V29-ternary-QJ2}

\section{Context and archive interface}
\label{sec:V29-context}

The three-block carrier isolated in V28 is
\[
 J_c(x|q|yz)=\kappa_3(X_x,X_q,Z_{yz}).
\]
This is precisely the local composite-ternary joining carrier treated
in the f2 archive at V67.  That theorem was proved after the V59
regression and keeps the full block-product cumulant assembled.  Its
local numerator is the three-term tree polynomial on the quotient
blocks and has the required powers \(s_\alpha^2\) unpaid and
\(s_\alpha\) when the joining carrier pays.

There is one payer interface not explicit in V27.  If the sole payer
lies in a physical-private bank, replacing that bank by a coarse
\(C/s_\alpha\) bound would discard the private damping needed to pay
its row denominator.  This chapter proves a paid private-moment
estimate which retains both effects.  It then verifies every payer
location and closes the complete three-block part of the V28
residual.

\section{Paid moments of a complete physical-private bank}
\label{sec:V29-paid-private}

Let \(P\) be a nonempty physical-private bank on one row.  Put
\[
 Y_P:=\sum_{e\in P}a_e,
 \qquad
 q_P:=\prod_{e\in P}\eta_e,
\]
where \(0\le\eta_e\le1\) are the normalized one-link Wilson
expectations.  Its complete paid scalar is
\begin{equation}
 {\mathfrak p}_P:=
 \frac{Y_Pq_P}{1-q_P}.
\label{eq:V29-paid-private}
\end{equation}
The zero-denominator case is interpreted by continuity; for a
nontrivial bank in the standing heat-kernel range one has \(q_P<1\).

\begin{lemma}[All fixed paid private moments]
\label{lem:V29-paid-private-moments}
For every fixed integer \(m\ge0\), there is \(C_m^{\rm pp}<\infty\)
such that
\begin{equation}
 \boxed{
 (s_\alpha Y_P)^m{\mathfrak p}_P
 \le\frac{C_m^{\rm pp}}{s_\alpha}.}
\label{eq:V29-paid-private-moments}
\end{equation}
The constant is uniform in the size of \(P\), its representations,
and physical multiplicities.
\end{lemma}

\begin{proof}
The archived one-link failure debit and its Bernoulli product
compiler give
\begin{equation}
 s_\alpha Y_Pq_P\le D(1-q_P).
\label{eq:V29-product-failure}
\end{equation}
The archived fixed-order private-product theorem gives, for every
fixed \(k\),
\begin{equation}
 (s_\alpha Y_P)^kq_P\le C_k.
\label{eq:V29-private-moments-inherited}
\end{equation}

Write \(t=s_\alpha Y_P\).  If \(t\le1\), multiply
\eqref{eq:V29-product-failure} by \(t^m/(1-q_P)\) to get
\[
 t^{m+1}\frac{q_P}{1-q_P}\le Dt^m\le D.
\]
If \(t>1\), \eqref{eq:V29-product-failure} implies
\[
 tq_P\le D(1-q_P),
 \qquad
 1-q_P\ge\frac{t}{D+t}\ge\frac1{D+1}.
\]
Therefore
\[
 t^{m+1}\frac{q_P}{1-q_P}
 \le(D+1)t^{m+1}q_P
 \le(D+1)C_{m+1}.
\]
In both cases,
\[
 t^m{\mathfrak p}_P
 =
 \frac1{s_\alpha}
 t^{m+1}\frac{q_P}{1-q_P}
\]
has the asserted bound.
\end{proof}

\begin{corollary}[A paid private bank still pays its row denominator]
\label{cor:V29-paid-private-denominator}
If \(Y_P\ge c\Xi_v\), \(c>0\), then for every edge stock \(H\),
\begin{equation}
 s_\alpha{\mathfrak p}_P H
 \le
 \frac{C_c}{s_\alpha}\frac{H}{\Xi_v}.
\label{eq:V29-paid-edge}
\end{equation}
For \(m=2\), the same bank can pay two powers of \(\Xi_v\) while
retaining the unique \(s_\alpha^{-1}\) paid scale.
\end{corollary}

\begin{proof}
By \(Y_P\ge c\Xi_v\) and
\Cref{lem:V29-paid-private-moments} with \(m=1\),
\[
 s_\alpha\Xi_v{\mathfrak p}_P
 \le c^{-1}s_\alpha Y_P{\mathfrak p}_P
 \le\frac{C_c}{s_\alpha}.
\]
Multiply by \(H/\Xi_v\).  The \(m=2\) statement is identical with
the square of \(s_\alpha\Xi_v\).
\end{proof}

\begin{firewall}[A private payer is not a coarse outside payer]
\label{fire:V29-private-payer}
When a private bank pays, its factor is
\({\mathfrak p}_P\), not an independent \(C/s_\alpha\) spectator
multiplied by a second copy of \(q_P\).  Equation
\eqref{eq:V29-paid-private-moments} is one estimate on that complete
paid bank and uses the final denominator once.
\end{firewall}

\begin{corollary}[Completion of the V27 payer ledger]
\label{cor:V29-correct-V27-payer}
The paid claims in
\Cref{thm:V27-private-first-response,
prop:V27-distinct-double-private}
remain valid when the sole payer lies in either physical-private
bank: use \Cref{cor:V29-paid-private-denominator} at that bank,
the unpaid private denominator estimate at every other private bank,
and ordinary bounds elsewhere.
\end{corollary}

\begin{proof}
The joining numerator supplies one factor
\(s_\alpha H\) for the row whose private bank pays.  Combine it with
\eqref{eq:V29-paid-edge}; this gives
\(C s_\alpha^{-1}H/\Xi_v\).  Every other joining edge is treated by
\eqref{eq:V27-private-unpaid-edge}.  No second payer is introduced.
\end{proof}

\section{The exact local quotient-tree polynomial}
\label{sec:V29-local-tree}

Regard the three quotient blocks as
\[
 A=\{y,z\},\qquad B=\{x\},\qquad C=\{q\}.
\]
At a local joining coordinate \(c\), define
\begin{align}
 b_{AB,c}&:=h_{xy,c}+h_{xz,c},
\notag\\
 b_{AC,c}&:=h_{qy,c}+h_{qz,c},
\notag\\
 b_{BC,c}&:=h_{xq,c}.
\label{eq:V29-cluster-stocks}
\end{align}
Put
\begin{equation}
 {\cal Q}_c^{x|q|yz}
 :=
 b_{AB,c}b_{AC,c}
 +b_{AB,c}b_{BC,c}
 +b_{AC,c}b_{BC,c}.
\label{eq:V29-local-Q}
\end{equation}

\begin{theorem}[Inherited assembled composite-ternary estimate]
\label{thm:V29-local-ternary}
Let \({\cal N}_c\) be the complete positive envelope of
\(\kappa_3(X_x,X_q,Z_{yz})\) before payment, and let
\({\cal P}_c\) be its envelope when that joining carrier contains the
sole payer.  Then
\begin{align}
 \|{\cal N}_c\|
 &\le Cs_\alpha^2{\cal Q}_c^{x|q|yz},
\label{eq:V29-local-unpaid}\\
 \|{\cal P}_c\|
 &\le Cs_\alpha{\cal Q}_c^{x|q|yz}.
\label{eq:V29-local-paid}
\end{align}
The sum in \eqref{eq:V29-local-Q} is formed before the positive
envelope; its factors are actual row-pair influences.
\end{theorem}

\begin{proof}
Apply the archived theorem
\(\texttt{lem:V67-local-ternary}\) after the row relabelling
\((1,2,3,4)=(y,z,x,q)\).  It uses the operator-martingale ternary
identity on the complete block-product cumulant.  Replacing the
composite variable \(Z_{yz}\) telescopes in a fixed order, producing
the two sums in
\eqref{eq:V29-cluster-stocks}; the singleton--singleton influence is
\(h_{xq,c}\).  The ternary replicated-tree formula gives exactly the
three products in \eqref{eq:V29-local-Q}.  Its one-resolvent
allocation removes one power of \(s_\alpha\) only after the three
tree terms have been assembled, proving the paid line.
\end{proof}

\begin{firewall}[This is not the withdrawn orbit expansion]
\label{fire:V29-not-orbits}
\Cref{thm:V29-local-ternary} is an estimate on one complete
third cumulant of quotient-block products.  It is not obtained by
separately norming its \(2+1+1\) partition letters.  Hence the V59
graphwise regression and the V67 no-orbit firewall are respected.
\end{firewall}

\begin{lemma}[The three quotient trees lift to legal directed
four-row trees]
\label{lem:V29-tree-lift}
Multiply any term of \({\cal Q}_c^{x|q|yz}\) by the already retained
directed response \(S=y\to z\).  After expanding the two cluster
sums, orient the two quotient edges with tails \(x,q\) as follows:
\begin{enumerate}[label=\textup{(\roman*)},leftmargin=4em]
\item \(b_{AB}b_{AC}\) gives
      \(x\to u,\ q\to r,\ y\to z\), with
      \(u,r\in\{y,z\}\);
\item \(b_{AB}b_{BC}\) gives
      \(x\to u,\ q\to x,\ y\to z\), with
      \(u\in\{y,z\}\);
\item \(b_{AC}b_{BC}\) gives
      \(x\to q,\ q\to r,\ y\to z\), with
      \(r\in\{y,z\}\).
\end{enumerate}
Every displayed graph is a spanning tree with the three distinct
tails \(x,q,y\) and root \(z\).
\end{lemma}

\begin{proof}
In \textup{(i)}, both singleton quotient vertices attach to the
connected block \(yz\); adding its internal edge \(yz\) gives a
tree.  These are four rows of the V15 table.  In \textup{(ii)}, the
edges are \(xu,xq,yz\) with orientations
\(x\to u,q\to x,y\to z\); these are the two V15 rows with
\(Q=q\to x\).  In \textup{(iii)}, the edges are
\(xq,qr,yz\).  For \(r=y\) they form the path
\(x-q-y-z\); for \(r=z\) they form the tree with edges
\(xq,qz,yz\).  The displayed orientations lead to \(z\) in each
case.  No edge is repeated.
\end{proof}

\begin{remark}[The third quotient tree is an alternate closing route]
\label{rem:V29-third-route}
The term \(b_{AC}b_{BC}\) does not use the V12 label
\(D=x\to y\) or \(x\to z\).  It gives the stronger alternate
response \(x\to q\), followed by \(q\to r\).  It must be retained as
a legal tree rather than forced into the narrower product
\eqref{eq:V28-DQ-stock}.
\end{remark}

\section{Private normalization of the ternary joining carrier}
\label{sec:V29-normalization}

\begin{lemma}[Unpaid private normalization of all three quotient
trees]
\label{lem:V29-unpaid-normalization}
On the double physical-private branch of V27,
\begin{equation}
 q_{P_x}q_{P_q}s_\alpha^2
 {\cal Q}_c^{x|q|yz}
 \le
 C\sum_{\tau\in{\mathfrak T}_{xq|yz}}
 d_{\tau,x}(c)d_{\tau,q}(c),
\label{eq:V29-unpaid-normalization}
\end{equation}
where \({\mathfrak T}_{xq|yz}\) is the finite set of directed
two-edge quotient trees in
\Cref{lem:V29-tree-lift}, and \(d_{\tau,x},d_{\tau,q}\) are their
actual edge stocks divided respectively by \(\Xi_x,\Xi_q\).
\end{lemma}

\begin{proof}
Expand \eqref{eq:V29-local-Q} into actual row pairs.  In every
monomial assign the edge with tail \(x\) to \(q_{P_x}\) and the edge
with tail \(q\) to \(q_{P_q}\), exactly as listed in
\Cref{lem:V29-tree-lift}.  Then
\[
 (s_\alpha q_{P_x}H_{x u})
 (s_\alpha q_{P_q}H_{q r})
 \le C
 \frac{H_{xu}}{\Xi_x}\frac{H_{qr}}{\Xi_q}
\]
by \Cref{lem:V27-private-denominator}.  The two cases using \(xq\)
assign its denominator according to the stated tail; symmetry of
\(H_{xq}\) makes both assignments legitimate.  Sum the fixed finite
row-pair choices.
\end{proof}

\begin{lemma}[Every payer has the same normalized scale]
\label{lem:V29-payer-ledger}
For the complete three-block packet, every admitted location of the
unique payer yields
\begin{equation}
 \frac{C}{s_\alpha}
 \sum_{\tau\in{\mathfrak T}_{xq|yz}}
 d_{\tau,x}(c)d_{\tau,q}(c)
\label{eq:V29-paid-normalization}
\end{equation}
after the two physical-private banks are retained.
\end{lemma}

\begin{proof}
If the joining cumulant pays, use
\eqref{eq:V29-local-paid}; it has one fewer power of
\(s_\alpha\) than \eqref{eq:V29-local-unpaid}.  Repeat the
factorization in
\Cref{lem:V29-unpaid-normalization}; equivalently factor out
\(s_\alpha^{-1}\).

If a nonprivate prefix, separator, or suffix factor pays, use its
complete \(C/s_\alpha\) bound and the unpaid local line.  If the
private bank at \(x\) or \(q\) pays, use
\Cref{cor:V29-paid-private-denominator} on its assigned quotient
edge and \Cref{lem:V27-private-denominator} on the other.  These
payer locations are disjoint.  Thus exactly one
\(s_\alpha^{-1}\) occurs in every case.
\end{proof}

\begin{theorem}[Complete three-block \(QJ_2\) closure]
\label{thm:V29-three-block-close}
The \(\rho_3=x|q|yz\) carrier in
\Cref{prop:V28-carrier-dichotomy} satisfies the legal unpaid and
once-paid quartic marked-tree bounds.  Hence it is removed from the
V28 residual.
\end{theorem}

\begin{proof}
Use the exact private multiplier
\eqref{eq:V27-exact-private-factor} for both physical-private banks
and the assembled local bounds in
\Cref{thm:V29-local-ternary}.  Apply
\Cref{lem:V29-unpaid-normalization} or
\Cref{lem:V29-payer-ledger}.  The remaining response
\(S=y\to z\) is already dominant; attach it after this complete
estimate by \Cref{lem:V27-mixed-attachment}.  Every resulting
three-edge product is a legal tree by
\Cref{lem:V29-tree-lift}.  The complete prefix and suffix stay
assembled, and the number of row and payer routings is fixed.
\end{proof}

\begin{corollary}[Three-block \(QJ_2\) with an alternate tree term]
\label{cor:V29-QJ2-status}
The first two terms of
\eqref{eq:V29-local-Q} verify the V28 \(QJ_2\) target for the original
\((D,Q)\) labels.  The third term closes by the alternate
\(x\to q,\ q\to r\) tree.  Thus the complete carrier closes even
though the narrower scalar polynomial
\eqref{eq:V28-DQ-stock} does not list every useful quotient tree.
\end{corollary}

\begin{proof}
Read the three cases of
\Cref{lem:V29-tree-lift} and apply
\Cref{thm:V29-three-block-close}.
\end{proof}

\section{The remaining discrete-prefix carrier}
\label{sec:V29-residual}

\begin{theorem}[Exact clean residual after V29]
\label{thm:V29-exact-residual}
Every clean first-breaker packet is closed except possibly the
discrete-prefix branch of
\Cref{prop:V28-carrier-dichotomy}.  Its exact local carrier is
\[
 J_c^{xq,(2+)}
 =
 \sum_{\pi\in{\cal A}_{xq}}L_{c,\pi},
\]
with both physical-private multipliers, the complete suffix
containing the \(S\)-role when needed, and the unique payer retained.
Nonclean states remain separate.
\end{theorem}

\begin{proof}
\Cref{thm:V27-exact-residual} reduced the clean sector to one
same-occurrence double-private packet.
\Cref{prop:V28-carrier-dichotomy} split it into the three-block and
discrete-prefix carriers.  The former is
\Cref{thm:V29-three-block-close}; the latter is exactly the displayed
packet and has not been estimated here.
\end{proof}

\begin{assessment}[Bottleneck after V29]
\label{ass:V29-bottleneck}
The smallest quotient-cumulant interface is complete, including the
previously implicit private-bank payer.  The sole remaining clean
first-breaker carrier begins from four singleton prefix blocks.  Its
rank-two local sum contains \(2+2\), \(3+1\), and one-block letters.
The next proof must determine whether these are the influence-tree
parts of the complete local fourth cumulant and its exact lower-order
corrections.  They must be assembled before capacity; applying the
V67 ternary estimate separately to a \(3+1\) letter would repeat the
withdrawn orbit error.
\end{assessment}

\section{Regression checks and result ledger}
\label{sec:V29-results}

\begin{proposition}[V29 regression suite]
\label{prop:V29-regressions}
\begin{enumerate}[label=\textup{(\roman*)},leftmargin=4em]
\item The private-bank payer is bounded as one complete bank.
\item Its paid estimate retains the moments needed for a row
      denominator.
\item The joining carrier remains a complete block-product third
      cumulant.
\item Cluster edges are expanded into actual row pairs only after the
      ternary tree formula.
\item Each private multiplier is assigned to the edge with its own
      tail.
\item The \(xq\) edge is oriented according to that assignment, not
      counted twice.
\item The third quotient tree is retained as an alternate legal route.
\item Every payer line has exactly one \(s_\alpha^{-1}\).
\item The discrete-prefix fourth-order carrier remains explicit.
\end{enumerate}
\end{proposition}

\begin{proof}
Items \textup{(i)--(ii)} are
\Cref{lem:V29-paid-private-moments,
cor:V29-paid-private-denominator,
fire:V29-private-payer}.  Items \textup{(iii)--(iv)} are
\Cref{thm:V29-local-ternary,fire:V29-not-orbits}.
Items \textup{(v)--(vii)} are
\Cref{lem:V29-tree-lift,
lem:V29-unpaid-normalization,
rem:V29-third-route}.  Item \textup{(viii)} is
\Cref{lem:V29-payer-ledger}; item \textup{(ix)} is
\Cref{thm:V29-exact-residual}.
\end{proof}

\begin{theorem}[Third-series V29 result ledger]
\label{thm:V29-results}
Relative to V1--V28, V29 proves:
\begin{enumerate}[label=\textup{(V29.\arabic*)},leftmargin=6.3em]
\item all fixed moments of a complete paid physical-private bank;
\item completion of every private-bank payer case left implicit in
      V27;
\item the exact V67 composite-ternary quotient-tree interface for
      \(x|q|yz\);
\item legal lifting of all three quotient trees with tails \(x,q,y\);
\item unpaid and every-payer private normalization of the two local
      quotient edges;
\item complete closure of the three-block \(QJ_2\) carrier; and
\item reduction of the remaining clean first-breaker sector to the
      discrete-prefix rank-two local carrier.
\end{enumerate}
\end{theorem}

\begin{proof}
These are
\Cref{lem:V29-paid-private-moments,
cor:V29-correct-V27-payer,
thm:V29-local-ternary,
lem:V29-tree-lift,
lem:V29-unpaid-normalization,
lem:V29-payer-ledger,
thm:V29-three-block-close,
thm:V29-exact-residual}.
\end{proof}

\begin{programme}[V30: companion audit and discrete-prefix
rank-two assembly]
\label{prog:V29-V30}
\begin{enumerate}[leftmargin=3em]
\item Perform the compulsory audit of the newest active SM and NS
      notes before choosing the exact fourth-order decomposition.
\item Express \(J_c^{xq,(2+)}\) as a finite combination of the
      complete local fourth cumulant, composite ternary corrections,
      and products of two binary covariances.
\item Keep each connected carrier assembled until its influence tree
      is formed.
\item Apply the two paid private moments at \(x,q\) only after the two
      quotient edge stocks have appeared.
\item Retain the suffix \(S\)-role and test all payer locations,
      including both private banks.
\item If a rank-one or disconnected correction survives, prove its
      Möbius cancellation from the exact discrete prefix rather than
      assigning it a positive orbit norm.
\item Keep nonclean states separate.  The V30 audit will set the next
      cross-program checkpoint at V35.
\end{enumerate}
\end{programme}

\begin{firewall}[Claim boundary after third-series V29]
\label{fire:V29-claim-boundary}
V29 closes the complete three-block double-private carrier but not
the discrete-prefix rank-two carrier or nonclean states.  Higher MTA,
WUFC, CPM, cutoff removal, reflection positivity,
Osterwalder--Schrader reconstruction, construction of the continuum
Yang--Mills measure, and a uniform positive continuum spectral gap
remain open.  The full Yang--Mills mass gap has not been proved.
\end{firewall}

\paragraph{Parent-version record.}
V29 was copied byte-for-byte from
\texttt{Renewal\_Yang\_Mills\_Mass\_Gap\_Research\_Notes\_V28.tex}
before this chapter was appended.  The V28 parent had \(14\,835\)
source lines, \(542\,375\) bytes, and SHA--256
\[
 \texttt{44a7ff620d87ac2855afcd996e6acee84e56e85ff349d9aa6ca7f2791915356c}.
\]
The next iteration is the compulsory V30 companion audit.

% =============================================================================
% V30 APPENDIX CHAPTER
% =============================================================================

\clearpage
\chapter{V30: Exact assembly of the discrete-prefix carrier}
\label{ch:V30-discrete-assembly}

\section{Compulsory companion audit and scope}
\label{sec:V30-audit}

Before choosing a decomposition of the last clean first-breaker
carrier, the newest active companion notes were inspected as required.
The audit inputs were
\begin{center}
\begin{tabular}{@{}p{0.21\linewidth}p{0.68\linewidth}@{}}
\toprule
Programme & Audited file and fingerprint\\
\midrule
SM--Einstein&
\texttt{Renewal\_SM\_Einstein\_Continuum\_Research\_Notes\_V6.tex},
\(2\,505\) source lines, \(100\,591\) bytes,
SHA--256
\texttt{19e6fb21dffa9cd36b8350f7f9762b695e4b7118822a10eec98241ebf15cb1b4};
\\[2mm]
Navier--Stokes&
\texttt{Renewal\_NS\_Stability\_Research\_Notes\_V31.tex},
\(23\,052\) source lines, \(887\,537\) bytes,
SHA--256
\texttt{f1121b62238ad5491bb7cf03635ea9934942edf573ed8faaa9ee79e1d38a6696}.
\\
\bottomrule
\end{tabular}
\end{center}

The SM V5--V6 shell-Hessian calculation gives one exact polarized
packet: the two response contributions inside it do not possess
separate signs and are not to be estimated independently.  The NS
V31 audit makes the complementary point that normalized marks do not
manufacture a missing first moment: the actual correlated source must
remain present.  The YM consequence is a proof-order rule, not an
imported estimate.  Namely, the five local partition letters left in
\(J_c^{xq,(2+)}\) must be assembled before a positive envelope is
taken; the two private moments and the suffix \(S\)-response must be
attached only to that assembled object.

\begin{firewall}[Type-correct companion transfer]
\label{fire:V30-audit-types}
No SM spectral estimate, gravitational resolvent bound, NS formation
estimate, or fluid mark is used below.  The companion audit transfers
only the algebraic discipline of retaining a complete connected
packet through the singular estimate.  The identity, Wilson influence
bound, private-moment payment, and tree normalization used in this
chapter are all statements already established inside the YM
programme.
\end{firewall}

\section{Exact enumeration of the rank-two set}
\label{sec:V30-enumeration}

Write a partition by placing elements in the same block next to one
another and separating blocks by vertical bars.  Thus, for example,
\(xqy|z\) denotes
\(\{\{x,q,y\},\{z\}\}\).  Recall from
\eqref{eq:V28-four-block-carrier} that
\({\cal A}_{xq}\) consists of the partitions \(\pi\) of
\(\{x,q,y,z\}\) for which
\begin{enumerate}[label=\textup{(\alph*)},leftmargin=4em]
\item the block of \(x\) meets \(\{y,z\}\);
\item \(q\) is not a singleton; and
\item the rank relative to the discrete prefix
      \(x|q|y|z\) is at least two.
\end{enumerate}

\begin{lemma}[The five and only five admissible partitions]
\label{lem:V30-five-partitions}
One has the exact set identity
\begin{equation}
 \boxed{
 {\cal A}_{xq}
 =
 \{\,xqyz,\ xqy|z,\ xqz|y,\ xy|qz,\ xz|qy\,\}.}
\label{eq:V30-five-partitions}
\end{equation}
\end{lemma}

\begin{proof}
For the discrete prefix, the quotient rank of \(\pi\) is
\(4-|\pi|\).  Condition \textup{(c)} therefore leaves only partitions
with one or two blocks.

There is one one-block partition, \(xqyz\), and it satisfies
\textup{(a)--(c)}.  The two-block partitions have type \(3+1\) or
\(2+2\).  The four partitions of type \(3+1\) are
\[
 xqy|z,\qquad xqz|y,\qquad xyz|q,\qquad qyz|x.
\]
The first two are admissible.  The third violates \textup{(b)}, and
the fourth violates \textup{(a)}.  The three partitions of type
\(2+2\) are
\[
 xy|qz,\qquad xz|qy,\qquad xq|yz.
\]
The first two are admissible and the last violates \textup{(a)}.
This exhausts all set partitions with at most two blocks and proves
\eqref{eq:V30-five-partitions}.
\end{proof}

\begin{remark}[Correction of any abstract routing ambiguity]
\label{rem:V30-routing-correction}
The explicit list \eqref{eq:V30-five-partitions}, rather than a
choice of graphical partners inside a nonsingleton block, is the
definitive description of the V28 carrier.  In particular it excludes
\(xq|yz\), \(xyz|q\), and \(qyz|x\) before any positive estimate.
\end{remark}

\section{The five letters are one composite ternary cumulant}
\label{sec:V30-identity}

At the common occurrence \(c\), abbreviate
\[
 F_i:=X_{i,c},\qquad
 Z_{yz}:=F_yF_z.
\]
Products retain the fixed tensor order used in the archived
joining-letter identity.  If \(B\) is a block, write
\(\kappa_c(F_i:i\in B)\) for its local cumulant.  The partition letter
is
\begin{equation}
 L_{c,\pi}
 =
 \prod_{B\in\pi}\kappa_c(F_i:i\in B),
\label{eq:V30-partition-letter}
\end{equation}
with the same fixed tensor order.

\begin{theorem}[Exact discrete-to-composite identity]
\label{thm:V30-discrete-composite}
The residual discrete-prefix carrier is exactly the composite ternary
cumulant:
\begin{equation}
 \boxed{
 J_c^{xq,(2+)}
 =
 \kappa_{3,c}(Z_{yz},F_x,F_q).}
\label{eq:V30-discrete-composite}
\end{equation}
More explicitly,
\begin{align}
 J_c^{xq,(2+)}
 ={}&
 \kappa_{4,c}(F_x,F_q,F_y,F_z)
\mathbb E_cF_y\,\kappa_{3,c}(F_z,F_x,F_q)
\mathbb E_cF_z\,\kappa_{3,c}(F_y,F_x,F_q)
\notag\\
&+\operatorname{Cov}_c(F_y,F_x)
       \operatorname{Cov}_c(F_z,F_q)
\operatorname{Cov}_c(F_y,F_q)
       \operatorname{Cov}_c(F_z,F_x).
\label{eq:V30-five-term-identity}
\end{align}
\end{theorem}

\begin{proof}
Apply the Leonov--Shiryaev product--cumulant formula to the three
products
\[
 \prod_{i\in\{y,z\}}F_i,\qquad F_x,\qquad F_q.
\]
A partition of the underlying four variables contributes precisely
when its join with \(yz|x|q\) is the one-block partition.  Enumerating
those partitions gives, in the order displayed in
\eqref{eq:V30-five-term-identity},
\[
 xqyz,\quad xqz|y,\quad xqy|z,\quad xy|qz,\quad xz|qy.
\]
Using \eqref{eq:V30-partition-letter}, their sum is
\(\kappa_{3,c}(Z_{yz},F_x,F_q)\).  By
\Cref{lem:V30-five-partitions} the same five partitions, and no
others, comprise \({\cal A}_{xq}\).  The definition
\eqref{eq:V28-four-block-carrier} of \(J_c^{xq,(2+)}\) therefore gives
\eqref{eq:V30-discrete-composite}.
\end{proof}

\begin{corollary}[No rank-one correction survives]
\label{cor:V30-no-rank-one}
There is no additional rank-one or disconnected correction outside
the right-hand side of \eqref{eq:V30-discrete-composite}.  In
particular, applying a positive envelope separately to the five terms
of \eqref{eq:V30-five-term-identity} is unnecessary and is not used.
\end{corollary}

\begin{proof}
The definition of the residual carrier is the sum over
\({\cal A}_{xq}\), and \Cref{lem:V30-five-partitions} exhausts that
set.  \Cref{thm:V30-discrete-composite} identifies its complete sum
with one cumulant before an envelope is taken.
\end{proof}

\section{Influence estimate with the suffix response retained}
\label{sec:V30-estimate}

Use the quotient blocks
\[
 A=\{y,z\},\qquad B=\{x\},\qquad C=\{q\},
\]
and the local influence stocks from V29,
\begin{equation}
 b_{AB,c}=h_{xy,c}+h_{xz,c},\qquad
 b_{AC,c}=h_{qy,c}+h_{qz,c},\qquad
 b_{BC,c}=h_{xq,c}.
\label{eq:V30-block-stocks}
\end{equation}
Set
\begin{equation}
 {\cal Q}_c
 :=
 b_{AB,c}b_{AC,c}
 +b_{AB,c}b_{BC,c}
 +b_{AC,c}b_{BC,c}.
\label{eq:V30-Q}
\end{equation}
The occurrence carrying the already selected \(S\)-role is retained
in the word.  At the influence level it supplies an actual
\(y\)-to-\(z\) stock \(h_{yz,d}\), where \(d\ge c\) is its occurrence
coordinate; when it occurs at a composite suffix, V25 attaches the
complete lower parent response only after the local carrier has been
bounded.

\begin{lemma}[Every quotient term lifts with the \(S\)-edge]
\label{lem:V30-three-tree-lifts}
After expanding one summand from each composite stock in
\eqref{eq:V30-block-stocks}, every monomial in
\[
 {\cal Q}_c h_{yz,d}
\]
contains a spanning tree on the four labelled rows
\(\{x,q,y,z\}\).  The three quotient-tree cases lift as follows:
\begin{align*}
 b_{AB,c}b_{AC,c}h_{yz,d}
 &\leadsto
 h_{xr,c}h_{qr',c}h_{yz,d},
 &&r,r'\in\{y,z\},\\
 b_{AB,c}b_{BC,c}h_{yz,d}
 &\leadsto
 h_{xr,c}h_{xq,c}h_{yz,d},
 &&r\in\{y,z\},\\
 b_{AC,c}b_{BC,c}h_{yz,d}
 &\leadsto
 h_{qr,c}h_{xq,c}h_{yz,d},
 &&r\in\{y,z\}.
\end{align*}
\end{lemma}

\begin{proof}
In the first line, the two local edges attach both singleton vertices
\(x,q\) to the set \(\{y,z\}\), and the \(S\)-edge joins \(y\) to
\(z\).  Hence the resulting graph is connected and has three edges
on four vertices, so it is a tree.  In the second line the edge
\(xq\) joins \(q\) to \(x\), the first edge joins \(x\) to one of
\(y,z\), and the \(S\)-edge joins the remaining two physical rows;
again the graph is a tree.  The third line is identical with \(x\)
and \(q\) interchanged.  This proves every choice in the displayed
expansions, including \(r=r'\) in the first line.
\end{proof}

\begin{lemma}[Assembled local estimate]
\label{lem:V30-local-estimate}
Let \({\cal N}^{\mathrm{disc}}_c\) be the complete positive envelope
of the discrete-prefix carrier after
\eqref{eq:V30-discrete-composite}, and let
\({\cal P}^{\mathrm{disc}}_c\) denote the corresponding envelope when
the common joining carrier contains the sole payer.  Then
\begin{align}
 \|{\cal N}^{\mathrm{disc}}_c\|
 &\le C s_\alpha^2{\cal Q}_c,
\label{eq:V30-local-unpaid}\\
 \|{\cal P}^{\mathrm{disc}}_c\|
 &\le C s_\alpha{\cal Q}_c.
\label{eq:V30-local-paid}
\end{align}
\end{lemma}

\begin{proof}
By \Cref{thm:V30-discrete-composite}, the object to be bounded is
exactly
\(\kappa_{3,c}(Z_{yz},F_x,F_q)\), up to the symmetry of the three
cumulant arguments.  This is the assembled composite-ternary carrier
of \Cref{thm:V29-local-ternary}, with the same three block stocks
\eqref{eq:V30-block-stocks}.  Equations
\eqref{eq:V29-local-unpaid} and \eqref{eq:V29-local-paid} therefore
give the two asserted bounds.  No term of
\eqref{eq:V30-five-term-identity} is separately normed.
\end{proof}

\begin{proposition}[Payer ledger for the discrete branch]
\label{prop:V30-payer-ledger}
For the complete discrete-prefix packet, the unique payer may lie:
\begin{enumerate}[label=\textup{(\roman*)},leftmargin=4em]
\item in the assembled joining carrier at \(c\);
\item in the retained \(S\)-occurrence or its complete lower parent;
\item in the physical-private bank at \(x\); or
\item in the physical-private bank at \(q\).
\end{enumerate}
In every case the normalized packet has precisely one factor
\(s_\alpha^{-1}\), retains both physical-private dampings, and is
summable by a three-edge directed tree from
\Cref{lem:V30-three-tree-lifts}.
\end{proposition}

\begin{proof}
In case \textup{(i)}, use the paid line
\eqref{eq:V30-local-paid}; the two private banks use the unpaid
fixed-moment line of \Cref{lem:V29-paid-private-moments}.  In case
\textup{(ii)}, use \eqref{eq:V30-local-unpaid}, attach the complete
\(S\)-parent response by the parent-response attachment lemma of V25,
and allocate its one response denominator only after multiplication.
In cases \textup{(iii)--(iv)}, use
\eqref{eq:V30-local-unpaid} and the paid private-moment estimate
\[
 (s_\alpha Y)^m\,{Yq\over1-q}\le {C_m\over s_\alpha}
\]
from \Cref{lem:V29-paid-private-moments} at the payer bank, leaving
the other private bank unpaid.  Thus exactly one of the local or bank
factors loses one power of \(s_\alpha\).

In all four cases the first-v-attachment telescope of V27 and the
private normalization of
\Cref{lem:V29-unpaid-normalization,lem:V29-payer-ledger} apply without
change.  Finally \Cref{lem:V30-three-tree-lifts} converts every
quotient-tree monomial together with the retained \(S\)-stock into a
legal three-edge row tree.  The directed-tree summation theorem then
gives the required row/output normalization.
\end{proof}

\begin{firewall}[Why the suffix may not be discarded]
\label{fire:V30-suffix}
For the discrete prefix, \(y\) and \(z\) were not connected in the
earlier prefix.  The composite notation \(Z_{yz}\) is an algebraic
assembly device and does not by itself create a physical \(yz\)
edge.  The actual third tree edge is the retained \(S\)-response in
\Cref{lem:V30-three-tree-lifts}.  Removing it before normalization
would turn the proof into the false rank-one attachment argument
excluded in V28.
\end{firewall}

\begin{theorem}[Closure of the discrete-prefix \(QJ_2\) carrier]
\label{thm:V30-discrete-close}
Every clean first-breaker packet in the discrete-prefix branch of
\Cref{prop:V28-carrier-dichotomy} satisfies the required once-paid
tree majorant, uniformly in the finite volume, cutoff, occurrence
location, and all equality patterns among the physical coordinates.
\end{theorem}

\begin{proof}
Keep the five local letters assembled and apply
\Cref{thm:V30-discrete-composite}.  The local unpaid and paid bounds
are \Cref{lem:V30-local-estimate}.  Preserve the two physical-private
multipliers and the full suffix carrying \(S\); all payer locations
are covered by \Cref{prop:V30-payer-ledger}.  Expand block stocks only
after these allocations.  \Cref{lem:V30-three-tree-lifts} supplies a
legal row tree term by term.

If physical coordinates coincide, the exact signature-stratum
decomposition of V24 is used before this argument.  It changes neither
the labelled row tree nor the two separate physical-private
multipliers.  The collision-erasure and parent-response results of
V25 therefore apply.  This proves the bound on every equality stratum
with the same constant, hence on their finite sum.
\end{proof}

\section{Completion of the clean first-breaker sector}
\label{sec:V30-clean-completion}

\begin{theorem}[Clean first-breaker closure]
\label{thm:V30-clean-first-breaker}
Under the standing V100 clean matching hypotheses, every packet in
the \(D,Q,S\) first-breaker sector is once-paid and obeys the required
three-edge directed-tree majorant.
\end{theorem}

\begin{proof}
V26 proves that clean external/directed acquisition has depth at most
two.  V27 closes the zero-private, single-private, and distinct
double-private cases and reduces the same-occurrence double-private
case to the two carriers of
\Cref{prop:V28-carrier-dichotomy}.  V29 closes its three-block carrier.
\Cref{thm:V30-discrete-close} closes its discrete-prefix carrier.
These branches are disjoint and exhaustive by the exact V27--V28
reduction.  Hence every clean first-breaker packet is closed.
\end{proof}

\begin{corollary}[The \(QJ_2\) target is discharged]
\label{cor:V30-QJ2}
The Wilson-specific rank-two quotient estimate \(QJ_2\) formulated in
V28 holds for both possible prefix shapes
\[
 x|q|yz,\qquad x|q|y|z,
\]
with the unique payer and dominant \(S\)-response retained.
\end{corollary}

\begin{proof}
The three-block statement is
\Cref{thm:V29-three-block-close}; the discrete statement is
\Cref{thm:V30-discrete-close}.
\end{proof}

\begin{assessment}[What is and is not completed]
\label{ass:V30-status}
The clean \(D,Q,S\) first-breaker obstruction isolated at the frozen
V100 endpoint is now closed.  It does not yet follow merely from this
local statement that the complete quartic master tree allocation has
been restored: the frozen V59--V100 archive contains regressions,
source-family reductions, and clean/nonclean qualifications whose
logical dependencies must be reconciled explicitly.  V31 will audit
that terminal ledger and determine whether any nonclean or other
quartic source family remains outside the hypotheses of
\Cref{thm:V30-clean-first-breaker}.
\end{assessment}

\begin{theorem}[Third-series V30 result ledger]
\label{thm:V30-results}
Relative to V1--V29, V30 proves:
\begin{enumerate}[label=\textup{(V30.\arabic*)},leftmargin=6.3em]
\item the exact five-element census of
      \({\cal A}_{xq}\);
\item the identity
      \(J_c^{xq,(2+)}=\kappa_3(Z_{yz},F_x,F_q)\);
\item absence of a residual rank-one correction;
\item applicability of the V29 unpaid and once-paid composite-ternary
      estimates to the discrete prefix;
\item legal lifting of all three quotient-tree monomials using the
      retained \(S\)-edge;
\item closure for every payer location, including both
      physical-private banks;
\item closure of the discrete-prefix \(QJ_2\) carrier; and
\item completion of the clean \(D,Q,S\) first-breaker sector.
\end{enumerate}
\end{theorem}

\begin{proof}
These statements are respectively
\Cref{lem:V30-five-partitions,
thm:V30-discrete-composite,
cor:V30-no-rank-one,
lem:V30-local-estimate,
lem:V30-three-tree-lifts,
prop:V30-payer-ledger,
thm:V30-discrete-close,
thm:V30-clean-first-breaker}.
\end{proof}

\begin{programme}[V31: terminal-ledger reconciliation]
\label{prog:V30-V31}
\begin{enumerate}[leftmargin=3em]
\item Reconstruct the exact dependency chain from the frozen V59
      regression through the V100 terminal theorem.
\item Identify the formal clean-state definition and list every
      excluded state without paraphrase.
\item Match the V30 theorem hypothesis-by-hypothesis against the
      terminal missing-tail statement.
\item Determine whether the restored statement is the full quartic
      MTA or only one matching/source subfamily.
\item For every remaining family, record an exact finite residual
      carrier rather than a prose bottleneck.
\item If the terminal implication is circular, move upstream to the
      first invalid dependency and replace it.
\item Preserve the next compulsory companion audit for V35.
\end{enumerate}
\end{programme}

\begin{firewall}[Claim boundary after third-series V30]
\label{fire:V30-claim-boundary}
V30 completes the clean first-breaker sector under the standing V100
matching hypotheses.  It has not yet established the full quartic MTA
outside those hypotheses.  Higher MTA, WUFC, CPM, cutoff removal,
reflection positivity, Osterwalder--Schrader reconstruction,
construction of the continuum Yang--Mills measure, and a uniform
positive continuum spectral gap remain open.  The full Yang--Mills
mass gap has not been proved.
\end{firewall}

\paragraph{Parent-version record.}
V30 was copied byte-for-byte from
\texttt{Renewal\_Yang\_Mills\_Mass\_Gap\_Research\_Notes\_V29.tex}
before this chapter was appended.  The V29 parent had \(15\,358\)
source lines, \(560\,716\) bytes, and SHA--256
\[
 \texttt{4075bb65ad0c5653f816e82bab96c08e2db09015698fd1e82df673dcbdc3c871}.
\]
The next compulsory companion audit is V35.

% =============================================================================
% V31 APPENDIX CHAPTER
% =============================================================================

\clearpage
\chapter{V31: Terminal-ledger reconciliation and the exact scope of
clean \(2+2\) closure}
\label{ch:V31-reconciliation}

\section{Purpose}
\label{sec:V31-purpose}

V30 closes the last local carrier produced by the active
missing-endpoint search.  This chapter checks the logical consequence
against the frozen f2 terminal ledger rather than declaring quartic
MTA by proximity.  Three questions are separated:
\begin{enumerate}[label=\textup{(\roman*)},leftmargin=4em]
\item Does the active V1--V30 chain discharge the clean residual
      explicitly left by f2 V100?
\item Is the invoked V67 composite-ternary estimate independent of
      the quartic statement being sought?
\item Which source families and representation states remain outside
      that implication?
\end{enumerate}

\section{Exact frozen-to-active dependency chain}
\label{sec:V31-chain}

The frozen terminal classification is
\(\texttt{thm:V99-endpoint-reduction}\).  Inside the clean
\(SU(2)\) \(S_{22}\) calculus its alternatives are:
\begin{enumerate}[label=\textup{(\alph*)},leftmargin=4em]
\item a private or unbalanced complete-moment compiler;
\item a fresh \(A,F,I\) endpoint;
\item one missing endpoint followed by a directed cell;
\item an all-directed word whose edge set contains a spanning tree;
\item an identified overlap with an \(A,F,I\) or directed bank;
\item a disconnected perfect matching or a row with no acquired
      incidence; and
\item a local representation state outside the standing compiler.
\end{enumerate}
The theorem states that \textup{(a)--(e)} were already controlled and
that only \textup{(f)--(g)} were new residuals.

\begin{lemma}[The f2 V100 residual is exactly the active starting
point]
\label{lem:V31-exact-handoff}
Within alternative \textup{(f)}, the f2 V100 ledger has precisely the
following exhaustive refinement:
\begin{enumerate}[label=\textup{(\roman*)},leftmargin=4em]
\item a matching-refining word with no component breaker, whose
      connected contribution is zero;
\item a matching word for which all breaker-visible endpoints possess
      fixed-fraction matching-tail responses, closed by the V100
      first-breaker theorem; or
\item an explicit positive breaker term with an explicit endpoint
      lacking its matching-tail response, reduced to one row-local
      depleted acquisition.
\end{enumerate}
The third-series V1 incidence identity starts from precisely
\textup{(iii)}.
\end{lemma}

\begin{proof}
The no-breaker statement is
\(\texttt{lem:V100-no-breaker-zero}\).  The fully incident statement
is
\(\texttt{thm:V100-matching-closure}\), with
\(\texttt{cor:V100-two-matchings-closed}\) covering both perfect
matchings.  Negating its endpoint hypothesis and discarding zero
breaker-stock terms is exactly
\(\texttt{prop:V100-missing-tail-reduction}\).  That proposition
applies the depleted identity at the named endpoint after every bank
already used in the matching word has been removed.

Third-series V1 begins with the same used-bank union, breaker pair,
and depleted row identity
\eqref{eq:V1-depleted-identity}; its one-step normal form
\Cref{thm:V1-one-step-acquisition} refines the fresh/overlap
alternatives left by f2 V100.  Thus no additional state is inserted
or silently omitted at the archive handoff.
\end{proof}

\begin{theorem}[Acyclic clean \(S_{22}\) closure]
\label{thm:V31-clean-S22}
Every packet in the clean \(SU(2)\) \(S_{22}\) calculus appearing in
\(\texttt{thm:V99-endpoint-reduction}\) satisfies the unpaid and
once-paid literal quartic marked-tree bound.
\end{theorem}

\begin{proof}
Alternatives \textup{(a)--(e)} of the V99 classification are closed
by that theorem's cited V92--V99 compilers.  Consider
alternative \textup{(f)}.  Apply
\Cref{lem:V31-exact-handoff}.  Its first two subcases are respectively
the f2 V100 zero and fully incident theorems.  In the third subcase,
the active V1--V26 acquisition chain terminates every nonprivate clean
search and reduces every private branch to the two exact V28
rank-two carriers.  V29 closes the three-block carrier and V30 closes
the discrete-prefix carrier.  This is
\Cref{thm:V30-clean-first-breaker}.  Hence every clean branch of
\textup{(f)} closes.  Alternative \textup{(g)} is excluded by the
hypothesis of the present theorem, so the clean classification is
exhausted.
\end{proof}

\begin{lemma}[No use of quartic MTA inside the V29--V30 local
estimate]
\label{lem:V31-no-circularity}
The composite-ternary estimates
\eqref{eq:V29-local-unpaid}--\eqref{eq:V29-local-paid}, and hence
\Cref{thm:V30-discrete-close}, do not assume the quartic MTA conclusion
of \Cref{thm:V31-clean-S22}.
\end{lemma}

\begin{proof}
The only inherited analytic input in
\Cref{thm:V29-local-ternary} is the archived local result
\(\texttt{lem:V67-local-ternary}\).  Its proof invokes the
operator-martingale ternary identity from V38, a fixed-order
replacement telescope for the composite variable, the ternary
replicated-tree formula, one-link Wilson heat estimates, and the
scalar one-resolvent allocation.  It is a one-coordinate third
cumulant estimate.  It neither cites global quartic MTA nor uses an
\(S_{22}\) source conclusion.  V30 invokes that estimate only after
the exact algebraic identity
\eqref{eq:V30-discrete-composite}; its third physical edge is supplied
by the retained \(S\)-occurrence, not by a presumed quartic tree
bound.  Therefore the dependency chain is acyclic.
\end{proof}

\section{Why this is not full quartic MTA}
\label{sec:V31-not-global}

\begin{lemma}[Exact census of proper quartic prefix types]
\label{lem:V31-prefix-census}
The proper partitions of four labelled rows have the following type
census:
\begin{center}
\begin{tabular}{@{}c c c@{}}
\toprule
type & number & block counts\\
\midrule
\(3+1\)&\(4\)&two\\
\(2+2\)&\(3\)&two\\
\(2+1+1\)&\(6\)&three\\
\(1+1+1+1\)&\(1\)&four\\
\bottomrule
\end{tabular}
\end{center}
Their total is \(14\), the number of proper elements of
\(\Pi([4])\).
\end{lemma}

\begin{proof}
For \(3+1\), choose the singleton in four ways.  For \(2+2\), divide
the four labels into two unordered pairs, giving three choices.  For
\(2+1+1\), choose the unique pair, giving
\(\binom42=6\) choices.  The discrete partition is unique.  The sum
is \(4+3+6+1=14=B_4-1\), since the Bell number \(B_4=15\) and the
one-block partition is excluded.
\end{proof}

\begin{proposition}[Source-family boundary]
\label{prop:V31-source-boundary}
\Cref{thm:V31-clean-S22} treats the \(2+2\) source family only, and
within it only the clean \(SU(2)\) state space used by the frozen
V99 theorem.  It supplies no estimate for the \(3+1\),
\(2+1+1\), or discrete prefix families of
\Cref{lem:V31-prefix-census}.  It also supplies no new theorem for any
separately retained local \([4]\) or \([3,3]\) source class.
\end{proposition}

\begin{proof}
The antecedent of \Cref{thm:V31-clean-S22} is the V99
\(S_{22}\) classification, whose fixed prefix is a two-block perfect
matching.  None of the other three prefix types occurs in that
antecedent.  The active V1--V30 proof refines only the missing endpoint
of the same matching word and therefore cannot enlarge its source
domain.  The final two local source labels are likewise absent from
the theorem's hypotheses.  A conclusion cannot cover a disjoint
source class that never appears in its exact decomposition.
\end{proof}

\begin{firewall}[The V30 discrete prefix is not the global discrete
source]
\label{fire:V31-two-discretes}
The expression \(x|q|y|z\) in V28--V30 is the local prefix of one
rank-two carrier occurring inside the \(S_{22}\) missing-tail
decomposition.  The discrete first-connectivity family in
\Cref{lem:V31-prefix-census} is the global V61 source family whose
entire earlier prefix has four singleton blocks.  The former identity
does not close the latter family.
\end{firewall}

\section{The nonclean predicate is not yet a theorem}
\label{sec:V31-nonclean}

\begin{proposition}[Document-level predicate gap]
\label{prop:V31-predicate-gap}
In the frozen f2 chain, the phrase nonclean local representation
state is not introduced by a definition or by a quantitative
necessary-and-sufficient condition.  Its sharpest occurrence is item
\textup{(v)} of \(\texttt{thm:V83-residual}\): a local representation
pattern outside the \(SU(2)\) protected/gapped calculus of V70 and
V83.  The same placeholder is then propagated through V92, V99, and
V100 without an estimate.
\end{proposition}

\begin{proof}
The exact occurrences are:
\begin{enumerate}[label=\textup{(\roman*)},leftmargin=4em]
\item \(\texttt{thm:V83-residual}\), item \textup{(v)}, which records
      the complement of the \(SU(2)\) protected/gapped input;
\item the V92 clean reduction, which retains the same complement;
\item \(\texttt{thm:V99-endpoint-reduction}\), item \textup{(g)},
      which calls it outside the standing compiler; and
\item \(\texttt{prop:V100-missing-tail-reduction}\), item
      \textup{(iii)}, which asserts no conclusion for it.
\end{enumerate}
None of those statements supplies a formula for the complement or a
bound on its carrier.  The explicit clean definitions that do occur
in f2 concern selected banks or chains: disjointness from marks and
previous banks, balance, a protected/gapped sector, and a terminal
moment certificate.  They do not define a global partition of all
local representation assignments.  Thus the quoted phrase is a
scope restriction, not a proved residual estimate.
\end{proof}

\begin{corollary}[No unconditional \(2+2\) theorem yet]
\label{cor:V31-no-unconditional-22}
The present notes may record clean \(SU(2)\) \(S_{22}\) closure, but
may not record an unconditional \(2+2\) source theorem until the
predicate gap in \Cref{prop:V31-predicate-gap} is removed.
\end{corollary}

\begin{proof}
The alternative labelled nonclean is an explicit branch of the
exhaustive V99 and V100 reductions and is not in the antecedent of
\Cref{thm:V31-clean-S22}.  No estimate in the current dependency chain
controls it.
\end{proof}

\section{Upstream replacement target}
\label{sec:V31-upstream}

Let \(G\) be the fixed compact gauge group and
\(\widehat G\) its irreducible dual.  For a nontrivial
\(\lambda\in\widehat G\), write
\[
 a_\lambda:=1+C_2(\lambda),\qquad d_\lambda:=\dim V_\lambda .
\]
At one balanced coordinate \(u\), let \(\mathsf T_u\) be the complete
local moment and let \(\mathsf P_u\) be its all-trivial-output
compression, with
\(\mathsf G_u=\mathsf T_u-\mathsf P_u\).

\begin{definition}[Group-uniform local compiler \(G\)-PG]
\label{def:V31-GPG}
The \(G\)-PG target consists of constants
\(c_G>0\), \(0<\rho_G<1\), and \(C_G<\infty\) such that:
\begin{align}
 \|\mathsf G_u\|
 &\le 1-c_Gs_\alpha,
\label{eq:V31-GPG-gap}\\
 \kappa_\otimes^{\,k|k^c}(\mathsf P_u)
 &\le\rho_G,
\label{eq:V31-GPG-capacity}\\
 a_{\lambda_{k,u}}^2
 &\le
 C_G a_{\lambda_{k,u}}a_{\lambda_{\ell,u}}
\label{eq:V31-GPG-partner}
\end{align}
for at least one other active row \(\ell\), whenever the local
invariant block containing row \(k\) is balanced.  Marked occurrences
are compressed only after their complete covariance difference is
formed.
\end{definition}

\begin{proposition}[Why \(G\)-PG is the first missing input]
\label{prop:V31-GPG-suffices-locally}
If \eqref{eq:V31-GPG-gap}--\eqref{eq:V31-GPG-partner} hold, the V70
row-rooted power estimate, protected/gapped bank currencies, and V83
powered one-failure resolver extend with group-dependent constants.
Consequently the representation-state exclusion named in
\Cref{prop:V31-predicate-gap} is removed from the local
protected/gapped step.
\end{proposition}

\begin{proof}
Equation \eqref{eq:V31-GPG-partner} replaces the only spin-order step
in \(\texttt{lem:V70-row-power}\).  Cauchy--Schwarz over a partner
bank then gives
\[
 S_{k\ell}(D)^2
 \le C_G |D_{k\ell}|H_{k\ell}^{(k)}(D).
\]
Tensoring \eqref{eq:V31-GPG-capacity} over the all-protected sector
gives \(\rho_G^N\); since
\(\sup_{N\ge1}N^2\rho_G^N<\infty\), the V70 protected currency follows.
Equation \eqref{eq:V31-GPG-gap} gives the complementary exponential
currency.  The V83 powered decomposition repeats with
\(\rho_G\) and \(c_G\), because it uses only orthogonality of the
protected/gapped sectors, their two one-coordinate bounds, and the
row-rooted power inequality.  These are exactly the three displayed
inputs.
\end{proof}

\begin{assessment}[Chosen V32 direction]
\label{ass:V31-direction}
The predicate audit prevents an unconditional source theorem, but it
also localizes the first upstream repair.  V32 should prove \(G\)-PG
from compact-group representation theory: a Casimir spectral floor
for nontrivial outputs, a Schmidt bound for the invariant projection
across one row versus the rest, and a highest-weight triangle
inequality yielding a comparable partner.  This is preferable to
continuing downstream with an undefined nonclean branch.
\end{assessment}

\begin{theorem}[Third-series V31 result ledger]
\label{thm:V31-results}
Relative to V1--V30, V31 proves:
\begin{enumerate}[label=\textup{(V31.\arabic*)},leftmargin=6.3em]
\item exact identification of the f2 V100-to-active-V1 handoff;
\item clean \(SU(2)\) \(S_{22}\) closure;
\item acyclicity of the V67/V29/V30 composite-ternary dependency;
\item the fourteen-element census of the four proper quartic prefix
      families;
\item separation of the local V30 discrete carrier from the global
      discrete source;
\item identification of an undefined nonclean-state predicate in the
      archived terminal ledger; and
\item the three-line \(G\)-PG representation-theoretic target
      sufficient to replace that local exclusion.
\end{enumerate}
\end{theorem}

\begin{proof}
These are
\Cref{lem:V31-exact-handoff,
thm:V31-clean-S22,
lem:V31-no-circularity,
lem:V31-prefix-census,
fire:V31-two-discretes,
prop:V31-predicate-gap,
def:V31-GPG,
prop:V31-GPG-suffices-locally}.
\end{proof}

\begin{programme}[V32: compact-group protected/gapped compiler]
\label{prog:V31-V32}
\begin{enumerate}[leftmargin=3em]
\item State the fixed compact connected semisimple group hypotheses
      and normalize the Casimir and highest-weight norm.
\item Prove uniform comparison
      \(1+C_2(\lambda)\asymp_G1+\|\lambda\|^2\).
\item If an invariant occurs in
      \(V_{\lambda_1}\otimes\cdots\otimes V_{\lambda_m}\), prove that
      every \(\|\lambda_k\|\) is bounded by a group constant times the
      sum of the other highest-weight norms.
\item Select one of the at most three other rows with comparable
      Casimir and derive \eqref{eq:V31-GPG-partner}.
\item Prove the exact invariant-projection product-state capacity
      across \(V_{\lambda_k}\mid W\) is at most
      \(1/d_{\lambda_k}\), including arbitrary multiplicity.
\item Combine the nontrivial Casimir floor and heat multiplier to
      prove \eqref{eq:V31-GPG-gap}.
\item Re-run the V70 and V83 regression hypotheses without asserting
      closure of marked or overlapping branches not actually included.
\end{enumerate}
\end{programme}

\begin{firewall}[Claim boundary after third-series V31]
\label{fire:V31-claim-boundary}
V31 proves the clean \(SU(2)\) \(S_{22}\) source implication and
identifies its acyclic dependencies.  It does not close the undefined
nonclean complement, the \(3+1\), \(2+1+1\), discrete, local
\([4]\), or \([3,3]\) source families.  Higher MTA, WUFC, CPM, cutoff
removal, reflection positivity, Osterwalder--Schrader reconstruction,
construction of the continuum Yang--Mills measure, and a uniform
positive continuum spectral gap remain open.  The full Yang--Mills
mass gap has not been proved.
\end{firewall}

\paragraph{Parent-version record.}
V31 was copied byte-for-byte from
\texttt{Renewal\_Yang\_Mills\_Mass\_Gap\_Research\_Notes\_V30.tex}
before this chapter was appended.  The V30 parent had \(15\,851\)
source lines, \(579\,070\) bytes, and SHA--256
\[
 \texttt{f2874016202d63eebf987c6e878598a495057ee3938ebabf83a0c82a664f25cb}.
\]
The next compulsory companion audit is V35.

% =============================================================================
% V32 APPENDIX CHAPTER
% =============================================================================

\clearpage
\chapter{V32: Compact-group protected/gapped representation compiler}
\label{ch:V32-compact-group-PG}

\section{Fixed group and normalization}
\label{sec:V32-setup}

Let \(G\) be a fixed compact connected semisimple Lie group.  Choose
an \(\operatorname{Ad}G\)-invariant inner product and a positive Weyl
chamber.  An irreducible unitary representation is denoted
\(V_\lambda\), where \(\lambda\) is its dominant highest weight, and
\[
 a_\lambda:=1+C_2(\lambda).
\]
With the compatible Casimir normalization,
\begin{equation}
 C_2(\lambda)=\langle\lambda,\lambda+2\rho\rangle ,
\label{eq:V32-Casimir-formula}
\end{equation}
where \(\rho\) is the half-sum of the positive roots, including the
direct-sum interpretation when \(G\) has several semisimple factors.
All constants below may depend on this fixed normalization and on
\(G\), but not on \(\lambda\), coordinate support, or physical
multiplicity.

\begin{lemma}[Uniform Casimir--weight comparison]
\label{lem:V32-Casimir-weight}
There are constants \(0<c_G\le C_G<\infty\) such that for every
dominant integral \(\lambda\),
\begin{equation}
 c_G(1+\|\lambda\|^2)
 \le
 1+C_2(\lambda)
 \le
 C_G(1+\|\lambda\|^2).
\label{eq:V32-Casimir-comparison}
\end{equation}
Moreover
\begin{equation}
 c_{\mathrm{gap},G}
 :=
 \min_{\lambda\ne0}C_2(\lambda)>0.
\label{eq:V32-Casimir-floor}
\end{equation}
\end{lemma}

\begin{proof}
For dominant \(\lambda\),
\(\langle\lambda,\rho\rangle\ge0\), so
\[
 C_2(\lambda)\ge\|\lambda\|^2.
\]
Conversely, Cauchy--Schwarz and
\(2ab\le a^2+b^2\) give
\[
 C_2(\lambda)
 \le
 \|\lambda\|^2+2\|\rho\|\,\|\lambda\|
 \le
 2\|\lambda\|^2+\|\rho\|^2.
\]
This proves \eqref{eq:V32-Casimir-comparison}, after changing fixed
constants.  The dominant integral weights form a discrete lattice
cone and \eqref{eq:V32-Casimir-formula} is strictly positive away from
\(\lambda=0\).  Its minimum on the nonzero lattice is therefore
positive, proving \eqref{eq:V32-Casimir-floor}.
\end{proof}

\section{Tensor-product triangle inequalities}
\label{sec:V32-triangle}

\begin{lemma}[Every weight lies in the highest-weight ball]
\label{lem:V32-weight-ball}
If \(\mu\) is a weight of \(V_\lambda\), then
\[
 \|\mu\|\le\|\lambda\|.
\label{eq:V32-weight-ball}
\]
\end{lemma}

\begin{proof}
The weights of \(V_\lambda\) lie in the convex hull of the Weyl orbit
\(W\lambda\).  The chosen norm is Weyl invariant.  The norm is convex,
so its maximum on that convex hull is attained at an extreme point
and equals \(\|\lambda\|\).
\end{proof}

\begin{theorem}[Two-sided tensor constituent triangle]
\label{thm:V32-tensor-triangle}
Suppose
\[
 V_\nu
 \subseteq
 V_{\lambda_1}\otimes\cdots\otimes V_{\lambda_m}.
\]
Then
\begin{align}
 \|\nu\|
 &\le\sum_{i=1}^m\|\lambda_i\|,
\label{eq:V32-forward-triangle}\\
 \|\lambda_k\|
 &\le
 \|\nu\|+\sum_{i\ne k}\|\lambda_i\|,
 \qquad 1\le k\le m.
\label{eq:V32-reverse-triangle}
\end{align}
\end{theorem}

\begin{proof}
The highest weight \(\nu\) of the constituent is a weight of the
tensor product.  Hence
\(\nu=\sum_i\mu_i\), where \(\mu_i\) is a weight of
\(V_{\lambda_i}\).  Equation \eqref{eq:V32-forward-triangle} follows
from the ordinary triangle inequality and
\Cref{lem:V32-weight-ball}.

By Frobenius reciprocity,
\[
 V_{\lambda_k}^*
 \subseteq
 V_\nu^*\otimes\bigotimes_{i\ne k}V_{\lambda_i}^*.
\]
Apply \eqref{eq:V32-forward-triangle} to this inclusion.  Duality
replaces a highest weight by \(-w_0\lambda\), which has the same norm.
This gives \eqref{eq:V32-reverse-triangle}.
\end{proof}

\begin{corollary}[Invariant blocks force a physical partner]
\label{cor:V32-invariant-partner}
If
\[
 \left(
 V_{\lambda_1}\otimes\cdots\otimes V_{\lambda_m}
 \right)^G\ne0
\]
and \(\lambda_k\ne0\), then
\begin{equation}
 \|\lambda_k\|
 \le
 \sum_{i\ne k}\|\lambda_i\|.
\label{eq:V32-invariant-triangle}
\end{equation}
Consequently some \(i\ne k\) is nontrivial and satisfies
\[
 \|\lambda_i\|
 \ge\frac{\|\lambda_k\|}{m-1}.
\label{eq:V32-invariant-max-partner}
\]
\end{corollary}

\begin{proof}
Take \(\nu=0\) in
\eqref{eq:V32-reverse-triangle}, and then apply the pigeonhole
principle to its \(m-1\) summands.
\end{proof}

\section{Balanced and unbalanced row compilers}
\label{sec:V32-balance}

At a coordinate involving \(m\le4\) active rows, call the assignment
unbalanced at row \(k\) if
\begin{equation}
 \sum_{i\ne k}\|\lambda_i\|
 \le\frac12\|\lambda_k\|.
\label{eq:V32-unbalanced}
\end{equation}
Call it balanced if no row satisfies
\eqref{eq:V32-unbalanced}.  This is the group-theoretic replacement
of the archived spin definition
\(\texttt{eq:V69-unbalanced}\).

\begin{lemma}[Balanced maximal partner and Casimir square]
\label{lem:V32-balanced-partner}
Fix an active row \(k\) at a balanced coordinate.  Some other active
row \(\ell\) satisfies
\begin{align}
 \|\lambda_\ell\|
 &>
 \frac{\|\lambda_k\|}{2(m-1)}
 \ge\frac{\|\lambda_k\|}{6},
\label{eq:V32-balanced-norm-partner}\\
 a_{\lambda_k}
 &\le C_Ga_{\lambda_\ell},
\label{eq:V32-balanced-Casimir-partner}\\
 a_{\lambda_k}^2
 &\le C_Ga_{\lambda_k}a_{\lambda_\ell}.
\label{eq:V32-balanced-power}
\end{align}
\end{lemma}

\begin{proof}
The negation of \eqref{eq:V32-unbalanced} gives
\[
 \sum_{i\ne k}\|\lambda_i\|>\frac12\|\lambda_k\|.
\]
One of at most \(m-1\le3\) summands gives
\eqref{eq:V32-balanced-norm-partner}.  Use
\Cref{lem:V32-Casimir-weight} and
\(\|\lambda_k\|\le6\|\lambda_\ell\|\) to obtain
\[
 a_{\lambda_k}
 \le C_G(1+\|\lambda_k\|^2)
 \le C_G(1+\|\lambda_\ell\|^2)
 \le C_Ga_{\lambda_\ell}.
\]
Multiplication by \(a_{\lambda_k}\) proves
\eqref{eq:V32-balanced-power}.
\end{proof}

\begin{corollary}[Row-rooted bank power inequality]
\label{cor:V32-bank-power}
Assign each balanced coordinate active at row \(k\) to one maximal
partner \(\ell\) satisfying
\Cref{lem:V32-balanced-partner}.  For a bank \(D_{k\ell}\), put
\[
 S_{k\ell}:=\sum_{u\in D_{k\ell}}a_{k,u},
 \qquad
 H_{k\ell}^{(k)}
 :=
 \sum_{u\in D_{k\ell}}a_{k,u}a_{\ell,u}.
\]
Then
\begin{equation}
 \boxed{
 S_{k\ell}^2
 \le
 C_G|D_{k\ell}|H_{k\ell}^{(k)}.}
\label{eq:V32-bank-power}
\end{equation}
\end{corollary}

\begin{proof}
Cauchy--Schwarz and
\eqref{eq:V32-balanced-power} give
\[
 S_{k\ell}^2
 \le
 |D_{k\ell}|\sum_{u\in D_{k\ell}}a_{k,u}^2
 \le
 C_G|D_{k\ell}|H_{k\ell}^{(k)}.
\]
\end{proof}

\begin{lemma}[Unbalanced output retains the dominant Casimir]
\label{lem:V32-unbalanced-output}
Suppose \eqref{eq:V32-unbalanced} holds and
\(V_\nu\) is any irreducible constituent of the local tensor product.
Then
\begin{align}
 \|\nu\|
 &\ge\frac12\|\lambda_k\|,
\label{eq:V32-unbalanced-output-norm}\\
 1+C_2(\nu)
 &\ge c_Ga_{\lambda_k}.
\label{eq:V32-unbalanced-output-Casimir}
\end{align}
In particular, an unbalanced local tensor product contains no
invariant vector.
\end{lemma}

\begin{proof}
Rearrange \eqref{eq:V32-reverse-triangle}:
\[
 \|\nu\|
 \ge
 \|\lambda_k\|-\sum_{i\ne k}\|\lambda_i\|
 \ge\frac12\|\lambda_k\|.
\]
Casimir--weight comparison then proves
\eqref{eq:V32-unbalanced-output-Casimir}.  If \(\nu=0\), the first
inequality contradicts the activity of row \(k\).
\end{proof}

\section{Invariant projection capacity with arbitrary multiplicity}
\label{sec:V32-capacity}

\begin{theorem}[Exact one-row invariant capacity]
\label{thm:V32-invariant-capacity}
Let \(V\) be a nontrivial irreducible unitary \(G\)-module of
dimension \(d\), let \(W\) be any finite-dimensional unitary
\(G\)-module, and let \(P_{\mathrm{inv}}\) be the orthogonal
projection onto \((V\otimes W)^G\).  For all density matrices
\(\rho_V,\rho_W\),
\begin{equation}
 \boxed{
 \operatorname{Tr}\!\left[
 P_{\mathrm{inv}}(\rho_V\otimes\rho_W)
 \right]
 \le\frac1d.}
\label{eq:V32-invariant-capacity}
\end{equation}
The bound is independent of the multiplicity of \(V^*\) in \(W\).
\end{theorem}

\begin{proof}
Only the \(V^*\)-isotypic summand of \(W\) contributes.  Write it
unitarily as
\[
 W[V^*]\simeq V^*\otimes M,
\]
where \(M\) is the multiplicity space.  Choose orthonormal bases
\((e_i)_{i=1}^d\) of \(V\), its dual basis in \(V^*\), and
\((f_a)\) of \(M\).  The invariant subspace has orthonormal basis
\[
 \Omega_a
 =
 \frac1{\sqrt d}
 \sum_{i=1}^d e_i\otimes e_i^*\otimes f_a .
\]
For unit vectors \(v\in V\) and \(w\in W\), write the isotypic
projection of \(w\) as
\(\sum_i e_i^*\otimes m_i\).  Then
\begin{align*}
 \langle v\otimes w,P_{\mathrm{inv}}v\otimes w\rangle
 &=
 \sum_a|\langle\Omega_a,v\otimes w\rangle|^2\\
 &=
 \frac1d
 \left\|
 \sum_i\langle e_i,v\rangle m_i
 \right\|^2\\
 &\le
 \frac1d\sum_i\|m_i\|^2
 \le\frac1d.
\end{align*}
The first inequality is Cauchy--Schwarz.  Every product of mixed
states is a convex combination of pure products, so the same bound
holds for density matrices.
\end{proof}

\begin{corollary}[All-trivial local sector]
\label{cor:V32-protected-capacity}
Fix one local moment-partition state and a nontrivial active row \(k\).
Let \(\mathsf P_u\) be the orthogonal sector in which every output
block is trivial.  Then
\begin{equation}
 \kappa_\otimes^{\,k|k^c}(\mathsf P_u)
 \le
 \frac1{d_{\lambda_{k,u}}}
 \le
 \rho_G,
\qquad
 \rho_G:=
 \frac1{d_{\min,G}}<1,
\label{eq:V32-protected-capacity}
\end{equation}
where
\[
 d_{\min,G}:=
 \min_{\lambda\ne0}\dim V_\lambda\ge2.
\]
\end{corollary}

\begin{proof}
In the fixed partition state, row \(k\) belongs to one invariant
output block.  Regroup that block as
\(V_{\lambda_{k,u}}\otimes W\) and apply
\Cref{thm:V32-invariant-capacity}.  All other invariant blocks are
orthogonal contractions.  A nontrivial irreducible representation of
a connected group cannot be one-dimensional when the group is
semisimple, so \(d_{\min,G}\ge2\).  Discreteness and Weyl's dimension
formula ensure that the displayed minimum exists.
\end{proof}

\begin{firewall}[Multiplicity and partition sums]
\label{fire:V32-capacity-scope}
The multiplicity space is already included in
\Cref{thm:V32-invariant-capacity}; no multiplicity factor may be
inserted afterward.  On the other hand,
\eqref{eq:V32-protected-capacity} is statewise for one complete local
moment-partition state.  A finite signed covariance combination must
be formed first, and a global finite partition sum is taken only after
the bank estimate.  The theorem does not claim that capacities of
overlapping signed summands multiply.
\end{firewall}

\section{The gapped output sector}
\label{sec:V32-gapped}

Assume the normalized one-link Wilson heat multiplier on an output
\(\nu\) is bounded by
\[
 e^{-s_\alpha C_2(\nu)}
\]
relative to the trivial output, as in the standing local heat
estimate.  Let \(\mathsf G_u\) be the orthogonal complement of the
all-trivial sector inside a complete local moment.

\begin{lemma}[Uniform compact-group output gap]
\label{lem:V32-output-gap}
For \(0<s_\alpha\le s_0\),
\begin{equation}
 \boxed{
 \|\mathsf G_u\|
 \le
 e^{-c_{\mathrm{gap},G}s_\alpha}
 \le
 1-\gamma_Gs_\alpha,}
\qquad
\gamma_G:=
\frac{1-e^{-c_{\mathrm{gap},G}s_0}}{s_0}>0.
\label{eq:V32-output-gap}
\end{equation}
For every integer \(t\ge1\), the \(t\)-powered heat block satisfies
\[
 \|\mathsf G_u^{(t)}\|
 \le e^{-t c_{\mathrm{gap},G}s_\alpha}.
\label{eq:V32-powered-gap}
\]
\end{lemma}

\begin{proof}
Every summand in \(\mathsf G_u\) has at least one nontrivial output
block.  Its Casimir is at least
\(c_{\mathrm{gap},G}\) by
\eqref{eq:V32-Casimir-floor}; all other normalized moment factors are
contractions.  Orthogonality of the output sectors gives the first
bound.  The function
\[
 s\longmapsto\frac{1-e^{-c_{\mathrm{gap},G}s}}s
\]
is decreasing on \((0,\infty)\), so its value on
\((0,s_0]\) is at least its value at \(s_0\).  This gives the affine
bound.  Replacing the heat time by \(ts_\alpha\) proves the powered
line.
\end{proof}

\begin{remark}[Marked local factors are finite separators]
\label{rem:V32-marked-separators}
At a marked joining or prefix coordinate, one first forms the complete
signed covariance and then compresses it into protected and gapped
output sectors.  A fixed local constant may result.  Such marked
coordinates are among the at most three tree marks and are treated as
finite separators; they are not tensorized over an arbitrarily long
bank.  Thus the strict factors
\(\rho_G<1\) and \(1-\gamma_Gs_\alpha\) are used only on unmarked
complete bank coordinates, where
\Cref{cor:V32-protected-capacity,lem:V32-output-gap} apply exactly.
\end{remark}

\section{Completion of the \(G\)-PG target}
\label{sec:V32-GPG}

\begin{theorem}[Compact-group \(G\)-PG compiler]
\label{thm:V32-GPG}
For every fixed compact connected semisimple \(G\), the three
representation-theoretic targets in
\Cref{def:V31-GPG} hold on unmarked complete local moment states, with
constants depending only on \(G\).  More precisely:
\begin{enumerate}[label=\textup{(\roman*)},leftmargin=4em]
\item the gapped sector obeys \eqref{eq:V32-output-gap};
\item the protected sector obeys
      \eqref{eq:V32-protected-capacity}; and
\item every balanced active row has a maximal-partner assignment
      obeying \eqref{eq:V32-balanced-power}.
\end{enumerate}
Unbalanced states obey the retained-output estimate
\eqref{eq:V32-unbalanced-output-Casimir}.
\end{theorem}

\begin{proof}
Items \textup{(i)--(iii)} are respectively
\Cref{lem:V32-output-gap,
cor:V32-protected-capacity,
lem:V32-balanced-partner}.
The final assertion is \Cref{lem:V32-unbalanced-output}.
\end{proof}

\begin{corollary}[Group-general V70 local bank currencies]
\label{cor:V32-V70-currencies}
For an unmarked row-rooted balanced bank of size \(N\), the V70
protected and gapped estimates hold with group-dependent constants:
\begin{align}
 \kappa_\otimes(\mathcal P_{k\ell}(D))
 &\le
 \frac{C_G}{N}
 \frac{H_{k\ell}^{(k)}(D)}{S_{k\ell}(D)^2},
\label{eq:V32-protected-currency}\\
 \kappa_\otimes(\mathcal G_{k\ell}(D))
 &\le
 C_G\min\left\{
 e^{-\gamma_Gs_\alpha N},
 \frac{H_{k\ell}^{(k)}(D)}
 {s_\alpha S_{k\ell}(D)^2}
 \right\}.
\label{eq:V32-gapped-currency}
\end{align}
\end{corollary}

\begin{proof}
Tensoring \eqref{eq:V32-protected-capacity} gives
\(\rho_G^N\).  By \eqref{eq:V32-bank-power},
\[
 \frac{H_{k\ell}^{(k)}(D)}{S_{k\ell}(D)^2}
 \ge\frac1{C_GN}.
\]
Since \(\sup_{N\ge1}N^2\rho_G^N<\infty\), the first line follows.
Tensoring \eqref{eq:V32-output-gap} gives the first entry in the
second line.  The affine replacement telescope used in V70, together
with \eqref{eq:V32-bank-power}, gives its second entry; that argument
uses no rank-one spin ordering once the row power inequality is
available.
\end{proof}

\begin{firewall}[What has actually been generalized]
\label{fire:V32-generalization-boundary}
\Cref{thm:V32-GPG,cor:V32-V70-currencies} remove the
\(SU(2)\)-specific local highest-spin and invariant-capacity inputs
from V70.  They do not by themselves replay every V71--V99 selected
bank, marked-overlap, payer, and one-failure summation for general
\(G\).  In particular, the phrase nonclean may be deleted from the
local representation decomposition only after those downstream
hypotheses are checked line by line against the group-general
currencies.
\end{firewall}

\begin{assessment}[Residual after the local repair]
\label{ass:V32-residual}
The upstream representation obstruction is now split cleanly.
There is no exceptional local highest-weight pattern for a fixed
compact connected semisimple group: every coordinate is balanced or
unbalanced, every balanced row has a comparable physical partner,
and every output state is protected or gapped.  What remains is a
proof-interface audit: the V71--V99 routing chain may also contain
rank-one spin sums, equal-spin channel formulae, or marked covariance
constants not covered by the unmarked \(G\)-PG theorem.  V33 will
enumerate those dependencies and isolate the first genuinely
\(SU(2)\)-specific downstream lemma.
\end{assessment}

\begin{theorem}[Third-series V32 result ledger]
\label{thm:V32-results}
Relative to V1--V31, V32 proves:
\begin{enumerate}[label=\textup{(V32.\arabic*)},leftmargin=6.3em]
\item uniform Casimir--highest-weight comparison and a positive
      nontrivial Casimir floor;
\item forward and reverse tensor-product highest-weight triangle
      inequalities;
\item the group-general balanced maximal-partner and bank power
      estimates;
\item retained output Casimir on every unbalanced coordinate;
\item the exact \(1/\dim V_\lambda\) invariant-projector capacity
      across one row, independent of multiplicity;
\item a uniform protected/gapped local spectral alternative;
\item group-general V70 protected and gapped bank currencies; and
\item localization of the remaining group-extension problem to the
      downstream routing interfaces.
\end{enumerate}
\end{theorem}

\begin{proof}
These are
\Cref{lem:V32-Casimir-weight,
thm:V32-tensor-triangle,
lem:V32-balanced-partner,
cor:V32-bank-power,
lem:V32-unbalanced-output,
thm:V32-invariant-capacity,
cor:V32-protected-capacity,
lem:V32-output-gap,
thm:V32-GPG,
cor:V32-V70-currencies,
fire:V32-generalization-boundary}.
\end{proof}

\begin{programme}[V33: downstream group-interface audit]
\label{prog:V32-V33}
\begin{enumerate}[leftmargin=3em]
\item List every V71--V99 lemma used by
      \(\texttt{thm:V99-endpoint-reduction}\) and classify its local
      analytic input as group-neutral, supplied by V32, or
      genuinely \(SU(2)\)-specific.
\item Check the V23/V83 powered one-failure sums and determine whether
      they sum replica powers or representation labels.
\item Separate equal-spin Clebsch--Gordan regression tests from
      hypotheses used in the final positive theorem.
\item Replace the first rank-one spin sum by a highest-weight-lattice
      estimate using the Weyl dimension and Casimir formulae, if one
      is actually required.
\item Retain marked joining coordinates as finite separators and do
      not tensorize their signed local constants.
\item State an exact updated clean/nonclean predicate only after this
      audit is complete.
\end{enumerate}
\end{programme}

\begin{firewall}[Claim boundary after third-series V32]
\label{fire:V32-claim-boundary}
V32 proves the compact-group local protected/gapped compiler and its
unmarked bank currencies.  It does not yet extend the complete
V71--V99 \(S_{22}\) routing theorem beyond its stated \(SU(2)\)
hypotheses, nor close the \(3+1\), \(2+1+1\), discrete, local
\([4]\), or \([3,3]\) sources.  Higher MTA, WUFC, CPM, cutoff removal,
reflection positivity, Osterwalder--Schrader reconstruction,
construction of the continuum Yang--Mills measure, and a uniform
positive continuum spectral gap remain open.  The full Yang--Mills
mass gap has not been proved.
\end{firewall}

\paragraph{Parent-version record.}
V32 was copied byte-for-byte from
\texttt{Renewal\_Yang\_Mills\_Mass\_Gap\_Research\_Notes\_V31.tex}
before this chapter was appended.  The V31 parent had \(16\,245\)
source lines, \(595\,190\) bytes, and SHA--256
\[
 \texttt{ee08e95db5078cf765758ba584e79256119bfaa5ec8188fec89ce1409e0b7ab8}.
\]
The next compulsory companion audit is V35.

% =============================================================================
% V33 APPENDIX CHAPTER
% =============================================================================

\clearpage
\chapter{V33: Group-general one-failure resolver and the corrected
output-mass interface}
\label{ch:V33-one-failure}

\section{What the V83 audit actually finds}
\label{sec:V33-audit}

The downstream audit begins at
\(\texttt{lem:V83-one-failure-lift}\), the first place where the
frozen text explicitly says that an \(SU(2)\) analytic input is being
used.  Its so-called integer and half-integer sums are not sums over
integer and half-integer spins.  They are the two geometric
resolvent-power series
\[
 \sum_{n\ge1}\Phi_N(r^n),
 \qquad
 \sum_{n\ge0}\Phi_N(r^{n+1/2}).
\]
Their variables are the powers \(n\) and \(n+\tfrac12\), and their
proof uses only a number \(0<\rho<1\) and an affine heat gap.  Thus
V32 supplies all group-dependent local input to that step.

The next line of V83 is
\[
 \Lambda_{\boldsymbol U}\le C(N+Z_{\boldsymbol U}).
\]
It is essential to type this line correctly.  Here
\(\boldsymbol U\) is a tuple of coordinatewise output blocks and
\(\Lambda_{\boldsymbol U}\) is the additive part of the final output
mass allocated to this coordinate bank.  It is not the Casimir of an
arbitrary irreducible constituent obtained by fusing all \(N\)
coordinates into one representation.  Such a different assertion
would in general lose a factor \(N\).  This chapter proves the
correct, coordinate-additive statement.

\section{An abstract protected/gapped product}
\label{sec:V33-abstract-product}

Let \(E\) be a finite set of size \(N\).  For each \(e\in E\), let
\(P_e\) be an orthogonal projection and let \(M_e(t)\) be a positive
contraction satisfying
\begin{equation}
 0\preccurlyeq M_e(t)
 \preccurlyeq
 P_e+r_t(I-P_e),
 \qquad
 r_t:=(1-\delta)^t,
\label{eq:V33-local-majorant}
\end{equation}
where \(0<\delta\le\tfrac12\).  Across one fixed tensor cut, suppose
\begin{equation}
 \kappa_\otimes(P_e)\le\rho<1.
\label{eq:V33-local-capacity}
\end{equation}
Set
\[
 M_E(t):=\bigotimes_{e\in E}M_e(t),
\qquad
 P_*:=\bigotimes_{e\in E}P_e
\]
when every \(P_e\ne0\), and put \(P_*=0\) otherwise.  Define
\begin{align}
 A_\rho(r)&:=\rho+(1-\rho)r,&
 B_\rho(r)&:=\rho(1-r),
\label{eq:V33-AB}\\
 \Phi_{N,\rho}^{(0)}(r)
 &:=
 A_\rho(r)^N-B_\rho(r)^N,&
 \Phi_{N,\rho}^{(1)}(r)
 &:=
 rA_\rho(r)^{N-1}.
\label{eq:V33-Phi}
\end{align}

\begin{lemma}[Abstract one-failure product bound]
\label{lem:V33-abstract-one-failure}
Under \eqref{eq:V33-local-majorant}--%
\eqref{eq:V33-local-capacity},
\begin{equation}
 \kappa_\otimes(M_E(t)-P_*)
 \le
 \begin{cases}
 \Phi_{N,\rho}^{(0)}(r_t),&P_*\ne0,\\
 \Phi_{N,\rho}^{(1)}(r_t),&P_*=0.
 \end{cases}
\label{eq:V33-one-failure}
\end{equation}
\end{lemma}

\begin{proof}
Expand the tensor product of the positive upper bounds in
\eqref{eq:V33-local-majorant}.  If \(A\subseteq E\) is the set on
which \(P_e\) is selected, tensor submultiplicativity of product-state
capacity across the one fixed cut gives a factor at most
\(\rho^{|A|}\); every complementary factor contributes \(r_t\).
The binomial sum is \(A_\rho(r_t)^N\).

When all \(P_e\ne0\), subtracting \(P_*\) removes the all-protected
term.  Re-expanding that removal in the same positive difference
majorant gives \(B_\rho(r_t)^N\), hence
\(\Phi^{(0)}\).  If some \(P_{e_0}=0\), at least one coordinate is
forced to contribute \(r_t\); bounding the remaining \(N-1\)
coordinates by the binomial sum gives
\(r_tA_\rho(r_t)^{N-1}=\Phi^{(1)}\).
\end{proof}

\begin{lemma}[The resolvent-power sums are group-neutral]
\label{lem:V33-power-sums}
For \(\nu\in\{0,1\}\) and \(0<\delta\le\tfrac12\),
\begin{align}
 \sum_{n\ge1}
 \Phi_{N,\rho}^{(\nu)}((1-\delta)^n)
 &\le
 C_\rho e^{-c_\rho\delta N}
 \left(1+\frac1{\delta N}\right),
\label{eq:V33-integer-powers}\\
 \sum_{n\ge0}
 \Phi_{N,\rho}^{(\nu)}((1-\delta)^{n+1/2})
 &\le
 C_\rho e^{-c_\rho\delta N}
 \left(1+\frac1{\delta N}\right).
\label{eq:V33-half-powers}
\end{align}
In particular either sum is at most
\(C_\rho/(\delta N)\).
\end{lemma}

\begin{proof}
For \(\nu=1\),
\[
 \int_0^1\frac{\Phi_{N,\rho}^{(1)}(r)}r\,dr
 =
 \int_0^1A_\rho(r)^{N-1}\,dr
 \le\frac1{(1-\rho)N}.
\]
For \(\nu=0\), observe that
\(A_\rho(r)-B_\rho(r)=r\).  On \(0<r\le\tfrac12\), the mean-value
theorem gives
\[
 \frac{A_\rho(r)^N-B_\rho(r)^N}{r}
 \le
 N A_\rho(\tfrac12)^{N-1}.
\]
On \([\tfrac12,1]\), direct integration of
\(2A_\rho(r)^N\) is \(O_\rho(N^{-1})\).  Thus
\begin{equation}
 \int_0^1\frac{\Phi_{N,\rho}^{(\nu)}(r)}r\,dr
 \le\frac{C_\rho}{N}.
\label{eq:V33-Phi-integral}
\end{equation}

Put \(\lambda=-\log(1-\delta)\), so
\(\delta\le\lambda\le2\delta\).
Monotonicity and geometric interval comparison give
\[
 \sum_{n\ge1}\Phi(e^{-\lambda n})
 \le
 \Phi(e^{-\lambda})
 \frac1\lambda
 \int_0^{e^{-\lambda}}\frac{\Phi(r)}r\,dr.
\]
Because
\[
 A_\rho(e^{-\lambda})
 \le
 \exp[-(1-\rho)(1-e^{-\lambda})],
\]
the endpoint and the truncated version of
\eqref{eq:V33-Phi-integral} acquire
\(e^{-c_\rho\delta N}\).  This proves
\eqref{eq:V33-integer-powers}.  Starting the same comparison at
\(e^{-\lambda/2}\) proves
\eqref{eq:V33-half-powers}.  Finally
\(e^{-cx}(1+x^{-1})\le C/x\).
\end{proof}

\section{Casimir-weighted resolver without a rank-one sum}
\label{sec:V33-Casimir-resolver}

\begin{lemma}[Scalar half-heat domination]
\label{lem:V33-half-heat}
For every \(x>0\),
\begin{equation}
 \frac{xe^{-x}}{1-e^{-x}}
 \le C e^{-x/2}.
\label{eq:V33-half-heat}
\end{equation}
\end{lemma}

\begin{proof}
On \(0<x\le1\), the left side is bounded and
\(e^{-x/2}\ge e^{-1/2}\).  On \(x\ge1\),
\((1-e^{-x})^{-1}\le(1-e^{-1})^{-1}\) and
\(xe^{-x/2}\) is bounded.  Multiplying the latter bound by
\(e^{-x/2}\) proves the claim.
\end{proof}

For an output tuple
\(\boldsymbol U=(U_e)_{e\in E}\), put
\[
 Z_{\boldsymbol U}:=\sum_{e\in E}C_2(U_e),
 \qquad
 Q_{\boldsymbol U}:=
 e^{-s_\alpha Z_{\boldsymbol U}}.
\]
If a coordinate has several output blocks, their Casimirs are all
included in the corresponding \(e\)-summand; the number of such
blocks is at most four and changes fixed constants only.

\begin{proposition}[Group-neutral geometric and Casimir resolvers]
\label{prop:V33-two-resolvers}
Assume \(\delta=\gamma_Gs_\alpha\) in
\eqref{eq:V33-local-majorant}, with \(s_\alpha\le s_0\) chosen so
that \(\delta\le\tfrac12\).  On the complement of \(P_*\), define
\begin{align*}
 \mathscr G_E
 &:=
 \sum_{\boldsymbol U}
 \frac{Q_{\boldsymbol U}}{1-Q_{\boldsymbol U}}
 P_{\boldsymbol U},\\
 \mathscr C_E
 &:=
 \sum_{\boldsymbol U}
 \frac{Z_{\boldsymbol U}Q_{\boldsymbol U}}
 {1-Q_{\boldsymbol U}}
 P_{\boldsymbol U}.
\end{align*}
Then
\begin{align}
 \kappa_\otimes(\mathscr G_E)
 &\le
 C_Ge^{-c_Gs_\alpha N}
 \left(1+\frac1{s_\alpha N}\right),
\label{eq:V33-geometric-resolver}\\
 \kappa_\otimes(\mathscr C_E)
 &\le
 \frac{C_G}{s_\alpha}
 e^{-c_Gs_\alpha N}.
\label{eq:V33-Casimir-resolver}
\end{align}
\end{proposition}

\begin{proof}
The geometric identity
\[
 \frac Q{1-Q}=\sum_{n\ge1}Q^n
\]
and \Cref{lem:V33-abstract-one-failure} identify the capacity of the
\(n\)-th term with the corresponding
\(\Phi_{N,\rho_G}((1-\gamma_Gs_\alpha)^n)\) upper bound.
Subadditivity and \eqref{eq:V33-integer-powers} prove
\eqref{eq:V33-geometric-resolver}.

For the Casimir line, put \(x=s_\alpha Z_{\boldsymbol U}\).
\Cref{lem:V33-half-heat} gives the blockwise inequality
\[
 \frac{Z_{\boldsymbol U}Q_{\boldsymbol U}}
 {1-Q_{\boldsymbol U}}
 \le
 \frac{C}{s_\alpha}Q_{\boldsymbol U}^{1/2}.
\]
The all-protected tuple has \(Z_{\boldsymbol U}=0\) and is absent.
Apply \Cref{lem:V33-abstract-one-failure} at \(t=\tfrac12\).
Since
\[
 A_{\rho_G}((1-\gamma_Gs_\alpha)^{1/2})^N
 \le e^{-c_Gs_\alpha N},
\]
and the same bound holds for both \(\Phi\)-types after changing
constants, \eqref{eq:V33-Casimir-resolver} follows.
\end{proof}

\section{The coordinate-additive output mass}
\label{sec:V33-output-mass}

\begin{lemma}[Fixed-arity compact-group fusion]
\label{lem:V33-fixed-fusion}
If \(m\le4\) and
\[
 V_\nu\subseteq
 V_{\lambda_1}\otimes\cdots\otimes V_{\lambda_m},
\]
then
\begin{equation}
 a_\nu
 \le
 C_G\sum_{i=1}^m a_{\lambda_i}.
\label{eq:V33-fixed-fusion}
\end{equation}
\end{lemma}

\begin{proof}
By \eqref{eq:V32-forward-triangle},
\[
 \|\nu\|^2
 \le
 \left(\sum_i\|\lambda_i\|\right)^2
 \le
 m\sum_i\|\lambda_i\|^2.
\]
Use \(m\le4\) and
\Cref{lem:V32-Casimir-weight} on both sides.
\end{proof}

\begin{lemma}[Exact meaning of the V83 bank payer mass]
\label{lem:V33-additive-payer}
Let a coordinate bank \(E\) have \(|E|=N\).  Resolve, separately at
each \(e\), the fixed-arity output blocks \(U_{e,b}\) of one complete
local partition state.  The portion of the product-group output mass
allocated to this bank satisfies
\begin{equation}
 \boxed{
 \Lambda_{\boldsymbol U}
 \le
 C_G\left(
 N+\sum_{e,b}C_2(U_{e,b})
 \right)
 =
 C_G(N+Z_{\boldsymbol U}).}
\label{eq:V33-additive-payer}
\end{equation}
\end{lemma}

\begin{proof}
The Fourier output over a coordinate set is a representation of the
product group \(G^E\).  Its Laplace/Casimir mass is additive over
coordinates.  At a fixed coordinate, the number of partition blocks
is at most four.  Apply
\Cref{lem:V33-fixed-fusion} inside each fixed-arity local block and
then the exact scalar one-payer allocation among those blocks.
Summing the allocated local masses over \(e\in E\) contributes at most
a fixed multiple of \(N+\sum_{e,b}C_2(U_{e,b})\).  There is no fusion
of the \(N\) coordinate outputs into one representation of \(G\), so
no support factor is generated.
\end{proof}

\begin{firewall}[False alternative explicitly excluded]
\label{fire:V33-no-global-fusion}
If one instead asked for the Casimir of a highest constituent of
\(\bigotimes_{e\in E}U_e\) as a representation of a single diagonal
copy of \(G\), the general bound furnished by
\Cref{lem:V33-fixed-fusion} would be of order
\(N\sum_e a_{U_e}\), and the factor \(N\) can be sharp.  That is not
the output mass in \Cref{lem:V33-additive-payer}.  The product-group
coordinate structure is indispensable.
\end{firewall}

\section{Group-general higher-incidence square tail}
\label{sec:V33-square-tail}

Use the V32 balanced partner bank stocks
\[
 S=\sum_{e\in E}a_{k,e},\qquad
 H=\sum_{e\in E}a_{k,e}a_{\ell,e},\qquad
 N=|E|.
\]
Then
\begin{equation}
 S^2\le C_GNH
\label{eq:V33-row-power}
\end{equation}
by \eqref{eq:V32-bank-power}.

\begin{theorem}[Compact-group V83 paid square tail]
\label{thm:V33-paid-square-tail}
Let \(E\) be an unmarked complete balanced bank, disjoint from
previously selected transfer and payer banks, in a fixed global
partition state.  Form its complete physical endpoint moment before
resolving outputs.  On the non-all-protected sector, let
\(\mathcal P_E\) be the positive envelope carrying the unique
allocated output numerator and resolvent.  Then
\begin{align}
 \kappa_\otimes^{\,k|k^c}(\mathcal P_E)
 &\le\frac{C_G}{s_\alpha},
\label{eq:V33-paid-universal}\\
 \boxed{
 s_\alpha^2S^2
 \kappa_\otimes^{\,k|k^c}(\mathcal P_E)
 \le C_GH.}
\label{eq:V33-paid-square-tail}
\end{align}
Both constants are independent of \(E\), its representations, and
their multiplicities.
\end{theorem}

\begin{proof}
By \Cref{lem:V33-additive-payer},
\[
 \mathcal P_E
 \preccurlyeq
 C_G(N\mathscr G_E+\mathscr C_E).
\]
Equations
\eqref{eq:V33-geometric-resolver}--%
\eqref{eq:V33-Casimir-resolver} immediately give the universal
\(C_G/s_\alpha\) bound.

For the marked line put \(x=s_\alpha N\).  If \(x\le1\), use
\(\kappa_\otimes(\mathscr G_E)\le C_G/(s_\alpha N)\),
\(\kappa_\otimes(\mathscr C_E)\le C_G/s_\alpha\), and
\eqref{eq:V33-row-power}:
\begin{align*}
 s_\alpha^2S^2N\kappa_\otimes(\mathscr G_E)
 &\le C_Gs_\alpha S^2
 \le C_GxH,\\
 s_\alpha^2S^2\kappa_\otimes(\mathscr C_E)
 &\le C_Gs_\alpha S^2
 \le C_GxH.
\end{align*}
Since \(x\le1\), both are at most \(C_GH\).

If \(x>1\), use the exponentially refined bounds.  The first term is
at most
\[
 C_Gx^2e^{-c_Gx}(1+x^{-1})H,
\]
and the second at most
\[
 C_Gxe^{-c_Gx}H.
\]
Both scalar functions are uniformly bounded on \([1,\infty)\).
This proves \eqref{eq:V33-paid-square-tail}.
\end{proof}

\begin{corollary}[Unpaid selected-bank certificate]
\label{cor:V33-unpaid-square-tail}
For the same bank, its common endpoint, cut-replicated endpoint, and
their assembled covariance satisfy
\begin{equation}
 \kappa_\otimes^{\,k|k^c}(\mathcal N_E)
 \le
 C_G\frac{H}{s_\alpha S^2}.
\label{eq:V33-unpaid-square-tail}
\end{equation}
With the payer they satisfy
\[
 \kappa_\otimes^{\,k|k^c}(\mathcal P_E)
 \le
 C_G\frac{H}{s_\alpha^2S^2}.
\label{eq:V33-paid-selected}
\]
\end{corollary}

\begin{proof}
For the unpaid line, split the complete bank into its all-protected
and complementary sectors.  The protected capacity is
\(\rho_G^N\); the complement is bounded by the \(t=1\) instance of
\Cref{lem:V33-abstract-one-failure}.  Combining those bounds with
\eqref{eq:V33-row-power}, separately in
\(s_\alpha N\le1\) and \(s_\alpha N>1\), is the proof of the V70
balanced tail with \(G\)-constants.  Subadditivity treats the
difference of the common and replicated endpoints only after both
complete endpoint moments have been formed.

The paid line is
\eqref{eq:V33-paid-square-tail} divided by
\(s_\alpha^2S^2\).
\end{proof}

\begin{theorem}[Removal of the V83 local representation exception]
\label{thm:V33-remove-V83-exception}
For a fixed compact connected semisimple gauge group, item
\textup{(v)} of the archived
\(\texttt{thm:V83-residual}\), namely a local representation pattern
outside the \(SU(2)\) protected/gapped calculus, is unnecessary.
Every local representation assignment is covered by exactly one of:
\begin{enumerate}[label=\textup{(\roman*)},leftmargin=4em]
\item the V32 unbalanced retained-output branch; or
\item the V32 balanced maximal-partner branch, whose unmarked complete
      bank has the unpaid and paid certificates
      \eqref{eq:V33-unpaid-square-tail} and
      \eqref{eq:V33-paid-selected}.
\end{enumerate}
\end{theorem}

\begin{proof}
The definition \eqref{eq:V32-unbalanced} and its negation partition
all local highest-weight assignments.  The unbalanced alternative is
\Cref{lem:V32-unbalanced-output}.  In the balanced alternative,
\Cref{lem:V32-balanced-partner} selects one of at most three actual
partners, and
\Cref{thm:V32-GPG} gives the protected/gapped decomposition.
\Cref{thm:V33-paid-square-tail,cor:V33-unpaid-square-tail} supply the
exact V83 endpoint estimates.  No further representation type
remains.
\end{proof}

\begin{firewall}[Clean-bank hypotheses still remain]
\label{fire:V33-clean-bank-boundary}
\Cref{thm:V33-remove-V83-exception} removes only the local
representation-pattern exception.  It does not remove the other V83
residual items: same-side maximal partners, strict mass ascent, absence
of a dominant bank after used coordinates are removed, or an
output-ancestry payer on an overlapping coordinate.  Those branches
were handled later in the frozen \(SU(2)\) chain, but their
group-general dependency must still be audited before the word
nonclean is deleted from V99 and V100.
\end{firewall}

\section{V33 status and next interface}
\label{sec:V33-status}

\begin{theorem}[Third-series V33 result ledger]
\label{thm:V33-results}
Relative to V1--V32, V33 proves:
\begin{enumerate}[label=\textup{(V33.\arabic*)},leftmargin=6.3em]
\item the abstract protected/gapped one-failure product bound;
\item support-uniform integer and half-integer resolvent-power sums
      for every fixed \(\rho<1\);
\item a rank-independent Casimir-resolver bound by half-heat
      domination;
\item the fixed-arity compact-group fusion inequality;
\item the coordinate-additive meaning and proof of the V83
      \(N+Z_{\boldsymbol U}\) payer mass;
\item the compact-group unpaid and paid higher-incidence square tails;
      and
\item removal of the local representation-pattern exception from the
      V83 residual theorem.
\end{enumerate}
\end{theorem}

\begin{proof}
These are
\Cref{lem:V33-abstract-one-failure,
lem:V33-power-sums,
lem:V33-half-heat,
prop:V33-two-resolvers,
lem:V33-fixed-fusion,
lem:V33-additive-payer,
thm:V33-paid-square-tail,
cor:V33-unpaid-square-tail,
thm:V33-remove-V83-exception}.
\end{proof}

\begin{assessment}[Next bottleneck]
\label{ass:V33-next}
The first purportedly \(SU(2)\)-specific downstream lemma was in fact
abstract, once the local invariant capacity and Casimir floor were
proved.  The audit should now continue through V84--V99.  The
highest-risk interfaces are not the \(SU(2)\) counterexamples, which
serve only as regressions, but the positive theorems using
equal-spin binary-prefix heat cards, exact Clebsch--Gordan adjacency,
or a claim that a marked overlap is a contraction.  V34 will trace
those dependencies and isolate the first input not already supplied
by V32--V33.
\end{assessment}

\begin{programme}[V34: V84--V99 group-dependency audit]
\label{prog:V33-V34}
\begin{enumerate}[leftmargin=3em]
\item Trace the positive dependencies of
      \(\texttt{thm:V99-endpoint-reduction}\) backward through
      V92--V98, ignoring counterexample-only hypotheses.
\item Separate group-neutral algebra, graph routing, and scalar
      one-payer allocation from representation-specific heat cards.
\item Check whether the V86--V91 equal-spin calculations are used in
      any terminal positive estimate or only falsify a discarded
      crossing-square route.
\item At the first genuinely group-specific positive lemma, state a
      compact-group replacement with exact constants and no implicit
      multiplicity sum.
\item Preserve complete marked covariance packets and keep the
      coordinate-additive payer convention of
      \Cref{lem:V33-additive-payer}.
\item The next compulsory SM/NS audit remains V35.
\end{enumerate}
\end{programme}

\begin{firewall}[Claim boundary after third-series V33]
\label{fire:V33-claim-boundary}
V33 removes the local representation-pattern exception from the V83
higher-incidence theorem for fixed compact connected semisimple
groups.  It does not yet prove that every later V84--V99 clean routing
compiler is group-general, nor close the \(3+1\), \(2+1+1\),
discrete, local \([4]\), or \([3,3]\) source families.  Higher MTA,
WUFC, CPM, cutoff removal, reflection positivity,
Osterwalder--Schrader reconstruction, construction of the continuum
Yang--Mills measure, and a uniform positive continuum spectral gap
remain open.  The full Yang--Mills mass gap has not been proved.
\end{firewall}

\paragraph{Parent-version record.}
V33 was copied byte-for-byte from
\texttt{Renewal\_Yang\_Mills\_Mass\_Gap\_Research\_Notes\_V32.tex}
before this chapter was appended.  The V32 parent had \(16\,858\)
source lines, \(614\,276\) bytes, and SHA--256
\[
 \texttt{f2c78d434e74548cabb418d2c22d5135c5bc1072f1947be0db838788ee3936da}.
\]
The next compulsory companion audit is V35.

% =============================================================================
% V34 APPENDIX CHAPTER
% =============================================================================

\clearpage
\chapter{V34: Heat-witness-free endpoint classification and
compact-group \(2+2\) closure}
\label{ch:V34-G-S22}

\section{The first positive dependency after V83}
\label{sec:V34-first-dependency}

The V84--V91 \(SU(2)\) calculations perform two indispensable
regression tasks: they disprove a separately enveloped crossing-square
estimate and forbid deletion or relabelling of the actual prefix and
joining edges.  One positive shortcut also appears: an equal-spin
superhot binary-prefix heat card can certify an endpoint even when its
ordinary overlap stock is small.

That shortcut is not needed for exhaustion.  If the ordinary
pair-overlap stock is small compared with the endpoint row mass, then
the exact incidence identity already forces the missing mass outside
the binary prefix.  The outside-prefix routing of V92--V99 was built
precisely for that alternative.  Thus a two-branch, group-neutral
classification can bypass the equal-spin heat card while retaining
all of its regression restrictions.

\section{A group-neutral binary-prefix dichotomy}
\label{sec:V34-prefix-dichotomy}

Let \(D_{uv}\) be a complete pair-prefix bank.  Put
\begin{align}
 A_u(D_{uv})
 &:=
 \sum_{e\in D_{uv}}a_{u,e},
 &
 H_{uv}(D_{uv})
 &:=
 \sum_{e\in D_{uv}}a_{u,e}a_{v,e},
\label{eq:V34-prefix-stocks}\\
 \mathfrak b_s(D_{uv})
 &:=
 \min\{1,s_\alpha H_{uv}(D_{uv})\},
 &
 \gamma_s(D_{uv})
 &:=
 \frac{\mathfrak b_s(D_{uv})}{H_{uv}(D_{uv})}
\label{eq:V34-prefix-response}
\end{align}
when \(H_{uv}>0\).  Let \(\Xi_u\) be the complete row mass.  Every
active coordinate has \(a_{v,e}\ge a_*>0\), so
\begin{equation}
 H_{uv}(D_{uv})\ge a_*A_u(D_{uv}).
\label{eq:V34-H-dominates-A}
\end{equation}

\begin{lemma}[Two immediate endpoint certificates]
\label{lem:V34-two-certificates}
Fix \(\Lambda\ge1\) and \(0<\eta<1\).  Each of
\begin{align}
 s_\alpha\Xi_u&\le\Lambda,
\label{eq:V34-row-cold}\\
 H_{uv}(D_{uv})&\ge\eta\Xi_u
\label{eq:V34-overlap-large}
\end{align}
implies
\begin{equation}
 \gamma_s(D_{uv})
 \le
 \frac{C_{\Lambda,\eta}}{\Xi_u}.
\label{eq:V34-endpoint-card}
\end{equation}
\end{lemma}

\begin{proof}
If \eqref{eq:V34-row-cold} holds, then
\[
 \gamma_s(D_{uv})
 =
 \min\{H_{uv}^{-1},s_\alpha\}
 \le s_\alpha
 \le\frac{\Lambda}{\Xi_u}.
\]
If \eqref{eq:V34-overlap-large} holds, then
\[
 \gamma_s(D_{uv})
 \le H_{uv}^{-1}
 \le\frac1{\eta\Xi_u}.
\]
\end{proof}

\begin{theorem}[Heat-witness-free endpoint alternative]
\label{thm:V34-prefix-or-outside}
Fix \(0<\eta<a_*/2\).  At every row \(u\) incident to a binary-prefix
bank, exactly the following exhaustive implication is available:
\begin{enumerate}[label=\textup{(\roman*)},leftmargin=4em]
\item one of \eqref{eq:V34-row-cold} or
      \eqref{eq:V34-overlap-large} holds, and the prefix supplies the
      endpoint card \eqref{eq:V34-endpoint-card}; or
\item
      \[
      s_\alpha\Xi_u>\Lambda,\qquad
      A_u(D_{uv})<\frac{\eta}{a_*}\Xi_u,
      \]
      and the other five cells in the exact V79 row-incidence
      identity carry more than
      \[
      \left(1-\frac{\eta}{a_*}\right)\Xi_u.
      \tag{V34.1}\label{eq:V34-outside-large}
      \]
\end{enumerate}
No equality of representations and no Clebsch--Gordan formula enters
this alternative.
\end{theorem}

\begin{proof}
If either displayed hypothesis in
\Cref{lem:V34-two-certificates} holds, use item \textup{(i)}.
Otherwise their negations give
\[
 s_\alpha\Xi_u>\Lambda,
 \qquad
 H_{uv}(D_{uv})<\eta\Xi_u.
\]
Equation \eqref{eq:V34-H-dominates-A} then gives
\[
 A_u(D_{uv})<\frac{\eta}{a_*}\Xi_u.
\]
The V79 identity is a disjoint sum of this own-prefix cell and the
private, unbalanced, foreign, directed, and internal/output-ancestry
cells.  Subtracting the own-prefix bound proves
\eqref{eq:V34-outside-large}.  The two alternatives are exhaustive
because they split by the truth values of
\eqref{eq:V34-row-cold} and
\eqref{eq:V34-overlap-large}.
\end{proof}

\begin{corollary}[The superhot certificate is optional]
\label{cor:V34-superhot-optional}
The V86--V87 equal-spin superhot heat card closes an additional subset
of item \textup{(ii)} of
\Cref{thm:V34-prefix-or-outside}, but is not needed to prove the
outside-prefix lower bound or to exhaust endpoint states.
\end{corollary}

\begin{proof}
The proof of \Cref{thm:V34-prefix-or-outside} uses only the ordinary
response \(\min\{1,s_\alpha H\}\), the positive floor \(a_*\), and the
exact row-incidence identity.  A superhot witness can only reassign a
state already in the second branch to a proved prefix-card branch; it
cannot create a third state.
\end{proof}

\section{The group-general unbalanced complete bank}
\label{sec:V34-unbalanced}

Let \(U\) be a bank of coordinates which are unbalanced in the sense
of \eqref{eq:V32-unbalanced}.  At each \(e\in U\), choose a dominant
row \(k(e)\) and put
\[
 w_e:=a_{\lambda_{k(e),e}},
 \qquad
 W:=\sum_{e\in U}w_e.
\]

\begin{lemma}[Two-sided output comparison on an unbalanced
coordinate]
\label{lem:V34-unbalanced-output}
For every allowed local output constituent \(\nu_e\),
\begin{equation}
 c_Gw_e
 \le
 1+C_2(\nu_e)
 \le
 C_Gw_e.
\label{eq:V34-unbalanced-output}
\end{equation}
\end{lemma}

\begin{proof}
The lower bound is
\eqref{eq:V32-unbalanced-output-Casimir}.  For the upper bound use
\Cref{lem:V33-fixed-fusion}.  The unbalanced condition gives
\[
 \sum_{i\ne k(e)}\|\lambda_{i,e}\|
 \le\frac12\|\lambda_{k(e),e}\|.
\]
Casimir--weight comparison therefore bounds the sum of all active
input masses, of which there are at most four, by \(C_Gw_e\).
Equation \eqref{eq:V33-fixed-fusion} proves the upper line.
\end{proof}

\begin{theorem}[Unbalanced bank moment and payer compiler]
\label{thm:V34-unbalanced-bank}
Let \(M_U\) be the complete Wilson moment on \(U\), with output
tuples resolved only after the complete product has been formed.
For each fixed integer \(m\ge0\),
\begin{equation}
 (s_\alpha W)^m\|M_U\|\le C_{G,m}.
\label{eq:V34-unbalanced-moments}
\end{equation}
If this bank contains the unique output numerator and resolvent, its
complete paid envelope satisfies
\begin{equation}
 \|\mathcal P_U\|\le\frac{C_G}{s_\alpha}.
\label{eq:V34-unbalanced-paid}
\end{equation}
\end{theorem}

\begin{proof}
By \Cref{lem:V34-unbalanced-output}, every resolved tuple has heat
multiplier
\[
 Q_{\boldsymbol\nu}
 \le
 \exp[-c_Gs_\alpha W].
\]
All recoupling maps and normalized nonpayer blocks are contractions.
Thus
\[
 \|M_U\|\le e^{-c_Gs_\alpha W},
\]
and \eqref{eq:V34-unbalanced-moments} follows from boundedness of
\(x^me^{-c_Gx}\).

For the payer, the coordinate-additive fusion theorem
\Cref{lem:V33-additive-payer} and
\eqref{eq:V34-unbalanced-output} give
\[
 \Lambda_{\boldsymbol\nu}\le C_GW.
\]
Therefore each paid block is bounded by
\[
 C_G
 \frac{W e^{-c_Gs_\alpha W}}
 {1-e^{-c_Gs_\alpha W}}
 \le\frac{C_G}{s_\alpha},
\]
using \(x e^{-x}/(1-e^{-x})\le1\) after rescaling \(x\) by the fixed
constant.  Orthogonality of the complete physical blocks proves the
operator-norm statement without a representation or multiplicity
sum.
\end{proof}

\section{Audit of the positive V84--V99 chain}
\label{sec:V34-dependency-audit}

\begin{proposition}[Representation-specific calculations are
regressions or optional shortcuts]
\label{prop:V34-SU2-optional}
In the positive dependency chain leading to
\(\texttt{thm:V99-endpoint-reduction}\), the explicit \(SU(2)\)
Clebsch--Gordan calculations in V84--V91 are used in exactly two ways:
\begin{enumerate}[label=\textup{(\roman*)},leftmargin=4em]
\item to refute the isolated crossing-square and separately enveloped
      joining estimates; and
\item to provide the optional superhot binary-prefix certificate
      discussed in \Cref{cor:V34-superhot-optional}.
\end{enumerate}
They are not needed for the remaining positive endpoint
classification once \Cref{thm:V34-prefix-or-outside} is used.
\end{proposition}

\begin{proof}
The V84 spin-one and V85--V91 tilted/coherent families occur in the
terminal chain as regression tests forbidding the discarded
crossing-square proof order.  The one positive formula
\(\texttt{thm:V86-binary-heat-capacity}\) feeds the V87 superhot
witness.  V88 itself proves that failure of the row-cold, ordinary,
and superhot cards forces outside-prefix mass, but its proof of the
outside mass uses only failure of the ordinary card:
\[
 H_{uv}<\eta\Xi_u
 \quad\Longrightarrow\quad
 A_u(D_{uv})<(\eta/a_*)\Xi_u.
\]
This is exactly
\Cref{thm:V34-prefix-or-outside}.  Hence omitting the superhot shortcut
merely sends more states, with the same quantitative lower bound, to
the already constructed outside-prefix routing.  All counterexample
firewalls remain valid.
\end{proof}

\begin{proposition}[Group-neutral terminal interfaces]
\label{prop:V34-terminal-interfaces}
After the replacements
\Cref{thm:V32-GPG,
thm:V33-paid-square-tail,
thm:V34-prefix-or-outside,
thm:V34-unbalanced-bank},
the positive V92--V100 terminal interfaces use no rank-one
representation formula.  Their inputs are:
\begin{center}
\small
\begin{tabular}{@{}p{0.23\linewidth}p{0.68\linewidth}@{}}
\toprule
archive interface & group-neutral content\\
\midrule
V92&
exact expansion of products of nonnegative edge-bank stocks into
literal coordinate markings; finite partition/payer fibres;\\
V93--V94&
coarse block pinching, trace-radius inequalities, and exact Fourier
duality without multiplicity resolution;\\
V95&
quadratic endpoint energy from a complete occurrence and a one-use
output dualizer;\\
V96&
three-factor sandwiching and support-uniform grouping of a common
suffix;\\
V97--V98&
finite five-factor Schatten sandwiches, fresh/self-adjacent bank
splitting, and bounded assembled covariance words;\\
V99&
insertion of fixed lower response bounds and finite graph
classification in \(K_{2,2}\);\\
V100&
the exact first component breaker, two-block covariance telescope,
and matching-tail acquisition.
\\
\bottomrule
\end{tabular}
\end{center}
All group dependence is confined to constants in the local heat,
Casimir, invariant-capacity, and complete-bank estimates already
supplied by V32--V34.
\end{proposition}

\begin{proof}
The displayed interfaces are statements about finite-dimensional
equivariant blocks, Schatten norms, scalar response stocks, exact
coordinate partitions, or four-vertex graphs.  Their proofs use only:
\begin{enumerate}[label=\textup{(\alph*)},leftmargin=4em]
\item contraction of normalized Wilson moments;
\item submultiplicativity and Cauchy--Schwarz for the one fixed
      product-state cut;
\item the unique scalar resolvent allocation;
\item the complete local payer cap and coordinate-additive fusion;
\item exact bank disjointness or an explicitly assembled overlap word;
      and
\item the positive nontrivial mass floor \(a_*\).
\end{enumerate}
Items \textup{(a)--(d)} are group-neutral once V32--V34 replace their
local constants.  Items \textup{(e)--(f)} are combinatorial.  The
places where V86 and V91 are cited in V97--V99 are labelled
regression tests: they verify that the abstract sandwiches do not
recreate the false edge relabelling.  They are not hypotheses in the
proofs of the terminal positive theorems.  Finally
\Cref{thm:V34-prefix-or-outside} replaces the only optional V86--V87
positive shortcut.  This proves the assertion.
\end{proof}

\begin{firewall}[No deletion of the \(SU(2)\) counterexamples]
\label{fire:V34-keep-regressions}
The V84--V91 examples remain mandatory regression tests.  A
group-general theorem must still retain the actual prefix edges, keep
the joining covariance assembled until its surrounding word is
bounded, and never substitute a same-side maximal partner for a
crossing edge.  Bypassing their optional heat card does not bypass
these negative conclusions.
\end{firewall}

\section{Removal of the terminal nonclean branch}
\label{sec:V34-remove-nonclean}

\begin{definition}[Admissible compact-group \(S_{22}\) packet]
\label{def:V34-admissible-S22}
An admissible packet is an exact V61 first-connectivity summand with
global prefix type \(2+2\), for a fixed compact connected semisimple
\(G\), in which:
\begin{enumerate}[label=\textup{(\roman*)},leftmargin=4em]
\item complete local partition and output sums are formed before
      positive envelopes;
\item the unique output numerator/resolvent is allocated once;
\item selected coordinate banks are removed before a fresh-cell
      identity is applied; and
\item an overlap is retained as one of the assembled V95--V98 words,
      rather than treated as two independent events.
\end{enumerate}
These are proof-order requirements, not restrictions on the local
highest weights.
\end{definition}

\begin{lemma}[Exhaustive local state split]
\label{lem:V34-exhaustive-local}
Every local representation assignment in an admissible packet lies
in exactly one balanced/unbalanced class.  Every balanced complete
bank lies in the protected/gapped decomposition of V32--V33, and every
unbalanced bank lies in
\Cref{thm:V34-unbalanced-bank}.  Hence no residual local
representation class remains outside the terminal compiler.
\end{lemma}

\begin{proof}
The inequality \eqref{eq:V32-unbalanced} and its negation are
exhaustive.  V32 proves the maximal-partner and protected/gapped
statements on the balanced side.  V33 proves their support-uniform
unpaid and paid bank resolvers.  V34 proves the complete unbalanced
bank bounds.  All constants are uniform over highest weights and
multiplicities for fixed \(G\).  Thus the complement previously named
nonclean is empty under
\Cref{def:V34-admissible-S22}.
\end{proof}

\begin{theorem}[Compact-group replacement of the V99 endpoint
reduction]
\label{thm:V34-G-endpoint-reduction}
For every admissible compact-group \(S_{22}\) packet, the endpoint
classification of the archived
\(\texttt{thm:V99-endpoint-reduction}\) holds with alternatives
\textup{(a)--(f)} only; its alternative \textup{(g)} is absent.
\end{theorem}

\begin{proof}
Use \Cref{thm:V34-prefix-or-outside} at a prefix endpoint.  The prefix
card is the archived paired branch.  In the outside branch, apply the
exact V79/V84 depleted identity after removing used banks.  Private,
unbalanced, \(A,F,I\), directed, and overlap cells are routed through
the group-neutral interfaces of
\Cref{prop:V34-terminal-interfaces}; the unbalanced input is
\Cref{thm:V34-unbalanced-bank}.  Balanced higher-incidence directed
cells use V32--V33.  The graph alternatives are the V99 spanning-tree,
missing-incidence, and disconnected-matching cases, which are
exhaustive on four rows.  By
\Cref{lem:V34-exhaustive-local}, no local representation complement
remains.  Therefore the former item \textup{(g)} is removed.
\end{proof}

\section{The complete \(2+2\) source theorem}
\label{sec:V34-complete-S22}

\begin{theorem}[Compact-group \(2+2\) marked-tree allocation]
\label{thm:V34-complete-S22}
Let \(G\) be a fixed compact connected semisimple gauge group.
Every admissible quartic first-connectivity source with global prefix
type \(2+2\) satisfies the support-uniform unpaid and once-paid literal
marked-tree bound, with a constant depending only on \(G\) and the
fixed four-row thresholds.
\end{theorem}

\begin{proof}
Apply \Cref{thm:V34-G-endpoint-reduction}.  Its alternatives
\textup{(a)--(e)} are closed by the positive interfaces listed in
\Cref{prop:V34-terminal-interfaces}.  For alternative \textup{(f)},
the frozen V100 first-breaker decomposition gives three cases:
no breaker, a fully incident matching, or a depleted explicit breaker
endpoint.  The first is zero, and the second is the V100 matching
theorem.  The third is exactly the active V1 starting state by
\Cref{lem:V31-exact-handoff}; its clean acquisition is completely
closed by \Cref{thm:V30-clean-first-breaker}.  The local
representation qualifier in that active theorem is supplied for
general \(G\) by V32--V34.  Thus every branch closes with its unique
payer and actual coordinate marks retained.
\end{proof}

\begin{corollary}[The global \(2+2\) family is no longer a quartic
source gap]
\label{cor:V34-S22-no-gap}
In the proper-prefix census of
\Cref{lem:V31-prefix-census}, all three labelled \(2+2\) partitions
obey the quartic MTA source estimate for fixed compact connected
semisimple \(G\).
\end{corollary}

\begin{proof}
\Cref{thm:V34-complete-S22} is invariant under permutation of the four
row labels.  The three \(2+2\) partitions are the three perfect
matchings of the four labels.
\end{proof}

\begin{firewall}[Why quartic MTA is still open]
\label{fire:V34-not-quartic-MTA}
The theorem closes one of the four global proper-prefix families.  It
does not estimate the four \(3+1\), six \(2+1+1\), or the one global
discrete source, and it does not close the separately retained local
\([4]\) or \([3,3]\) carriers.  The word discrete in V30 referred to
a local carrier inside \(S_{22}\), as explained in
\Cref{fire:V31-two-discretes}; it cannot be counted again here.
\end{firewall}

\section{V34 ledger and the next source}
\label{sec:V34-ledger}

\begin{theorem}[Third-series V34 result ledger]
\label{thm:V34-results}
Relative to V1--V33, V34 proves:
\begin{enumerate}[label=\textup{(V34.\arabic*)},leftmargin=6.3em]
\item a group-neutral prefix-card versus outside-prefix dichotomy;
\item proof that the equal-spin superhot card is optional for
      exhaustion;
\item compact-group unbalanced output comparison, bank moments, and
      complete payer control;
\item separation of the V84--V91 regression calculations from the
      positive terminal dependencies;
\item a group-neutral audit of the V92--V100 positive interfaces;
\item deletion of the undefined terminal nonclean representation
      branch under an exact admissible-packet definition; and
\item the full \(2+2\) proper-prefix source theorem for fixed compact
      connected semisimple \(G\).
\end{enumerate}
\end{theorem}

\begin{proof}
These are
\Cref{lem:V34-two-certificates,
thm:V34-prefix-or-outside,
cor:V34-superhot-optional,
lem:V34-unbalanced-output,
thm:V34-unbalanced-bank,
prop:V34-SU2-optional,
prop:V34-terminal-interfaces,
def:V34-admissible-S22,
lem:V34-exhaustive-local,
thm:V34-G-endpoint-reduction,
thm:V34-complete-S22}.
\end{proof}

\begin{assessment}[Choice before the V35 audit]
\label{ass:V34-next}
The first large quartic source family is now closed without the
nonclean placeholder.  The next source should not be chosen by visual
similarity.  V35 is a compulsory SM/NS audit; after it, the exact V63
joining-cumulant formula should be used to compare the remaining
\(3+1\), \(2+1+1\), and global discrete prefixes.  The most economical
candidate is \(2+1+1\), because its local composite-ternary carrier is
already the V67/V29 theorem, but its hot-row endpoint normalization
must be audited separately from the local V30 discrete branch.
\end{assessment}

\begin{programme}[V35: companion audit and \(2+1+1\) source
re-entry]
\label{prog:V34-V35}
\begin{enumerate}[leftmargin=3em]
\item Inspect the newest active SM and NS notes and record exact
      filenames, line counts, bytes, and SHA--256 fingerprints.
\item Transfer only proof-order information relevant to assembling a
      complete source before a singular estimate.
\item Write the exact V63 \(2+1+1\) first-connectivity word, including
      its binary prefix, composite ternary joining carrier, complete
      suffix, and all payer locations.
\item Reconcile the V67 fully thermal theorem with its hot-row
      reduction; do not redo the local ternary estimate already used
      in V29--V30.
\item Classify the hot endpoint by the group-general V32--V34
      balanced/unbalanced and prefix/outside alternatives.
\item Isolate the first global tree-normalization branch not already
      supplied by the \(2+2\) compiler.
\end{enumerate}
\end{programme}

\begin{firewall}[Claim boundary after third-series V34]
\label{fire:V34-claim-boundary}
V34 proves the complete global \(2+2\) proper-prefix source estimate
for fixed compact connected semisimple \(G\).  It does not close the
\(3+1\), \(2+1+1\), global discrete, local \([4]\), or \([3,3]\)
source families.  Higher MTA, WUFC, CPM, cutoff removal, reflection
positivity, Osterwalder--Schrader reconstruction, construction of the
continuum Yang--Mills measure, and a uniform positive continuum
spectral gap remain open.  The full Yang--Mills mass gap has not been
proved.
\end{firewall}

\paragraph{Parent-version record.}
V34 was copied byte-for-byte from
\texttt{Renewal\_Yang\_Mills\_Mass\_Gap\_Research\_Notes\_V33.tex}
before this chapter was appended.  The V33 parent had \(17\,454\)
source lines, \(633\,137\) bytes, and SHA--256
\[
 \texttt{70866f4a779b0b365f109c5ae842b20d19a29f7416bc5b4033a0428645a30014}.
\]
The next iteration is the compulsory V35 companion audit.

% =============================================================================
% V35 APPENDIX CHAPTER
% =============================================================================

\clearpage
\chapter{V35: Companion audit and exact hot-demand reduction for the
\(2+1+1\) source}
\label{ch:V35-hot-211}

\section{Compulsory companion audit}
\label{sec:V35-audit}

The newest active companion files in
\texttt{latex\_notes} were inspected read-only before choosing this
direction:
\begin{center}
\begin{tabular}{@{}p{0.20\linewidth}p{0.69\linewidth}@{}}
\toprule
programme & audited file and fingerprint\\
\midrule
SM--Einstein&
\texttt{Renewal\_SM\_Einstein\_Continuum\_Research\_Notes\_V13.tex},
\(4\,623\) source lines, \(186\,516\) bytes,
SHA--256
\texttt{a98b100f506d76c406928b7dc96c520f756637fa9cdd96c32403a2d1ebc3a5b1};
\\[2mm]
Navier--Stokes&
\texttt{Renewal\_NS\_Stability\_Research\_Notes\_V31.tex},
\(23\,052\) source lines, \(887\,537\) bytes,
SHA--256
\texttt{f1121b62238ad5491bb7cf03635ea9934942edf573ed8faaa9ee79e1d38a6696}.
\\
\bottomrule
\end{tabular}
\end{center}

SM V13 disintegrates a finite-cutoff measure near regular flat-toron
strata.  At a centralizer wall it does not invert a vanishing hard
determinant; it transfers the entire soft eigenspace to the parent
variable and keeps the remaining hard determinant uniformly
separated.  The type-correct YM lesson is that a failed endpoint
response is not to be inverted: the complete row mass must be
transferred to its exact incidence cell before a new estimate is
selected.

NS V31 is unchanged since the V30 audit.  It proves that flat dyadic
marks do not create the missing critical first moment; removing a
mark costs its actual radius weight, and a persistent interval cannot
be charged again.  The type-correct YM lesson is that every added tree
edge must be an actual coordinate stock and a bank already used by
the source or payer may not be debited again.

\begin{firewall}[Companion type boundary]
\label{fire:V35-audit-types}
No toron determinant, Einstein or Standard-Model shell estimate,
Navier--Stokes coarea bound, or formation estimate is used in this
chapter.  Only proof order transfers: move a degenerating object into
the retained parent packet rather than invert it, and remove a
normalized mark only with its actual first moment.  All mathematical
estimates below are YM identities or consequences of V32--V34.
\end{firewall}

\section{The exact archived \(2+1+1\) word}
\label{sec:V35-exact-word}

For the prefix \(12|3|4\), the V63 first-connectivity identity gives
\begin{equation}
 \mathcal S_{211}
 =
 \sum_{r=1}^n
 K_{12}^{<r,\mathrm{conn}}
 \otimes M_{\{3\}}^{<r}
 \otimes M_{\{4\}}^{<r}
 \otimes J_r(12|3|4)
 \otimes M_{>r}.
\label{eq:V35-S211}
\end{equation}
The joining carrier is the assembled composite third cumulant
\[
 J_r(12|3|4)
 =
 \kappa_3(Z_{12,r},X_{3,r},X_{4,r}).
\]
Its three quotient-tree stocks are
\begin{align}
 B_{AB}&:=H_{13}+H_{23},&
 B_{AC}&:=H_{14}+H_{24},&
 B_{BC}&:=H_{34}.
\label{eq:V35-quotient-stocks}
\end{align}
The binary prefix contributes \(H_{12}\).  The V67 payer ledger,
which includes a payer in the binary prefix, joining carrier, either
singleton bank, or the complete suffix, gives
\begin{equation}
 \mathfrak C_{211}^{\mathrm{pay}}
 \le
 Cs_\alpha^2H_{12}
 \left(
 B_{AB}B_{AC}
 +B_{AB}B_{BC}
 +B_{AC}B_{BC}
 \right).
\label{eq:V35-V67-carrier}
\end{equation}
This local and bankwise theorem is already proved and is not redone.

\begin{firewall}[What V67 actually closed]
\label{fire:V35-V67-scope}
The archived theorem
\(\texttt{thm:V67-thermal-211}\) turns
\eqref{eq:V35-V67-carrier} into a normalized tree bound only under
\(s_\alpha\Xi_i\le M\) for all four rows.  Its corollary explicitly
leaves a hot input row.  Later f2 chapters chose the \(2+2\) source
first; V34 closes that family but does not retrospectively close the
hot \(2+1+1\) normalization.
\end{firewall}

\section{The eight literal tree terms}
\label{sec:V35-tree-census}

For \(i,j\in\{1,2\}\), define
\begin{align}
 T_{ij}^{(0)}
 &:=
 \{12,i3,j4\},
\label{eq:V35-tree-zero}\\
 T_i^{(3)}
 &:=
 \{12,i3,34\},
\label{eq:V35-tree-three}\\
 T_j^{(4)}
 &:=
 \{12,j4,34\}.
\label{eq:V35-tree-four}
\end{align}

\begin{lemma}[Exact tree expansion]
\label{lem:V35-eight-trees}
Expanding the three terms of
\eqref{eq:V35-V67-carrier} gives precisely the following eight
labelled tree monomials:
\begin{equation}
 \{T_{ij}^{(0)}:i,j\in\{1,2\}\}
 \cup
 \{T_i^{(3)}:i\in\{1,2\}\}
 \cup
 \{T_j^{(4)}:j\in\{1,2\}\}.
\label{eq:V35-eight-trees}
\end{equation}
Every listed edge set is a spanning tree on \([4]\).
\end{lemma}

\begin{proof}
The product \(H_{12}B_{AB}B_{AC}\) has the four terms
\[
 H_{12}H_{i3}H_{j4},
 \qquad i,j\in\{1,2\},
\]
giving \(T_{ij}^{(0)}\).  The products
\(H_{12}B_{AB}B_{BC}\) and
\(H_{12}B_{AC}B_{BC}\) give respectively the two terms
\[
 H_{12}H_{i3}H_{34},
 \qquad
 H_{12}H_{j4}H_{34}.
\]
The three families are disjoint as labelled edge sets.  Each has four
vertices and three edges.  The first attaches rows \(3,4\) to the
connected edge \(12\); the second attaches \(4\) to \(3\) and \(3\)
to the connected edge \(12\); the third is symmetric.  Hence every
edge set is connected and therefore a tree.
\end{proof}

For a tree \(T\), define its excess-demand vector
\begin{equation}
 \delta_T(k):=\deg_T(k)-1,
 \qquad
 {\cal D}_T:=
 \prod_{k=1}^4(s_\alpha\Xi_k)^{\delta_T(k)}.
\label{eq:V35-demand}
\end{equation}

\begin{lemma}[Exact demand census]
\label{lem:V35-demand-census}
Every tree in \eqref{eq:V35-eight-trees} has
\[
 \sum_{k=1}^4\delta_T(k)=2.
\label{eq:V35-two-demands}
\]
More precisely:
\begin{enumerate}[label=\textup{(\roman*)},leftmargin=4em]
\item \(T_{11}^{(0)}\) has demand \(2e_1\), and
      \(T_{22}^{(0)}\) has demand \(2e_2\);
\item \(T_{12}^{(0)}\) and \(T_{21}^{(0)}\) have demand
      \(e_1+e_2\);
\item each \(T_i^{(3)}\) and \(T_j^{(4)}\) is a path and has demand
      one at each of its two internal vertices.
\end{enumerate}
Thus the only double demand at one row occurs at the centers of the
two stars \(T_{11}^{(0)}\) and \(T_{22}^{(0)}\).
\end{lemma}

\begin{proof}
For any four-vertex tree,
\[
 \sum_k(\deg_T(k)-1)=2|E(T)|-4=6-4=2.
\]
The two diagonal \(T_{ii}^{(0)}\) have center \(i\) of degree three
and all other vertices of degree one.  The two off-diagonal terms are
paths with internal vertices \(1,2\).  Direct inspection of
\eqref{eq:V35-tree-three}--\eqref{eq:V35-tree-four} gives paths and
their two degree-two vertices.
\end{proof}

\begin{lemma}[Normalization identity]
\label{lem:V35-normalization}
For every tree monomial in
\eqref{eq:V35-V67-carrier},
\begin{equation}
 \frac{
 s_\alpha^2\prod_{uv\in E(T)}H_{uv}}
 {\prod_{k=1}^4\Xi_k}
 =
 {\cal D}_T
 \prod_{uv\in E(T)}\theta_{uv},
\qquad
\theta_{uv}:=\frac{H_{uv}}{\Xi_u\Xi_v}.
\label{eq:V35-normalization}
\end{equation}
Consequently only rows with \(\delta_T(k)>0\) can obstruct the V67
thermal proof.
\end{lemma}

\begin{proof}
Substitute
\(H_{uv}=\theta_{uv}\Xi_u\Xi_v\).
The power of \(\Xi_k\) in the product is \(\deg_T(k)\); division by
\(\prod_k\Xi_k\) leaves
\(\Xi_k^{\deg_T(k)-1}\).  Equation
\eqref{eq:V35-two-demands} distributes the two factors of
\(s_\alpha\) exactly as in
\eqref{eq:V35-demand}.  A leaf has exponent zero, proving the final
claim.
\end{proof}

\begin{corollary}[Hot leaves are harmless]
\label{cor:V35-hot-leaves}
A row with \(s_\alpha\Xi_k>M\) but
\(\delta_T(k)=0\) requires no new moment or response estimate in the
tree term \(T\).  The hot residual is supported only on the one or two
internal vertices named by
\Cref{lem:V35-demand-census}.
\end{corollary}

\begin{proof}
Its row mass does not occur in
\({\cal D}_T\) in \eqref{eq:V35-normalization}.
\end{proof}

\section{Exact row-demand allocation}
\label{sec:V35-row-allocation}

Fix a tree \(T\) and an internal row \(k\).  Let
\(Q_T\) be the set of the at most three physical coordinates used by
the chosen coordinate marking of \(T\), together with the one
coordinate carrying a local payer if it is different.  Let
\[
 X_k(Q_T):=\sum_{e\in Q_T}a_{k,e}.
\]
Remove \(Q_T\) and every bank already used by the assembled
\(2+1+1\) word before classifying the remaining row incidences.

\begin{lemma}[Marked-heavy or fresh-cell demand split]
\label{lem:V35-heavy-or-fresh}
Fix \(0<\varepsilon<1\).  At a demanded row \(k\), one of the
following holds:
\begin{enumerate}[label=\textup{(\roman*)},leftmargin=4em]
\item some \(e\in Q_T\) satisfies
      \[
      a_{k,e}
      \ge
      \frac{\varepsilon}{|Q_T|}\Xi_k;
      \tag{heavy mark}\label{eq:V35-heavy-mark}
      \]
\item one identified previously used bank outside \(Q_T\) carries a
      fixed fraction of \(\Xi_k\);
\item a fresh \(P,U,A,F,D,\) or \(I\) cell carries at least
      \[
      \frac{1-\varepsilon-\varepsilon'}6\Xi_k
      \tag{fresh cell}\label{eq:V35-fresh-cell}
      \]
      for a fixed used-bank threshold \(\varepsilon'\); or
\item the already selected prefix bank supplies the endpoint card of
      \Cref{lem:V34-two-certificates}.
\end{enumerate}
The alternatives use actual disjoint coordinate stocks and have a
support-independent finite fibre.
\end{lemma}

\begin{proof}
First apply \Cref{thm:V34-prefix-or-outside}.  If its prefix-card
branch holds, this is item \textup{(iv)}.  Otherwise the complementary
row mass lies in the other incidence cells.  If
\(X_k(Q_T)\ge\varepsilon\Xi_k\), the at most \(|Q_T|\) marked
coordinates give \eqref{eq:V35-heavy-mark} by pigeonhole.

Suppose instead that the marked mass is smaller.  Split the remaining
row mass between banks already used by the word and their coordinate
complement.  If the former is at least
\(\varepsilon'\Xi_k\), record the actual bank attaining a fixed
fraction; there are only finitely many word roles, giving item
\textup{(ii)}.  Otherwise the six exact V79 cells on the fresh
complement total more than
\((1-\varepsilon-\varepsilon')\Xi_k\).  One is at least one sixth of
that amount, proving item \textup{(iii)}.  All splittings are by
disjoint coordinate sets, so no row incidence is copied.
\end{proof}

\begin{proposition}[Immediate demand compilers]
\label{prop:V35-immediate-compilers}
For either one demand token or the two tokens at a star center:
\begin{enumerate}[label=\textup{(\roman*)},leftmargin=4em]
\item a fresh or used physical-private \(P\)-cell is controlled by the
      all-order private product moments;
\item an unbalanced \(U\)-cell is controlled to every fixed order by
      \eqref{eq:V34-unbalanced-moments}, including its payer case;
\item a complete prefix/foreign/internal \(A,F,I\) cell with an actual
      response certificate pays one token by the quadratic endpoint
      energy interface; and
\item a directed \(D\)-cell supplies an actual partner edge and pays
      one token whenever its partner mass is compatible; otherwise it
      produces the strict finite mass ascent.
\end{enumerate}
\end{proposition}

\begin{proof}
Items \textup{(i)--(ii)} are respectively the archived all-order
private-bank theorem and
\Cref{thm:V34-unbalanced-bank}.  The quadratic energy of an
\(A,F,I\) occurrence is the V94--V95 endpoint theorem used
group-neutrally in
\Cref{prop:V34-terminal-interfaces}; its actual edge stock is retained.
The directed statement is the V83 partner comparison after
\Cref{thm:V33-paid-square-tail}: the bank numerator is the actual
\(H_{k\ell}\), and comparison with
\(\Xi_k\Xi_\ell\) is legal when
\(\Xi_\ell\le L\Xi_k\).  Failure is the strict ascent
\(\Xi_\ell>L\Xi_k\).  For a two-token private or unbalanced cell, use
the corresponding fixed second moment once, not two independent
first moments.
\end{proof}

\begin{firewall}[One-token response is not a two-token star compiler]
\label{fire:V35-one-not-two}
At \(T_{11}^{(0)}\) or \(T_{22}^{(0)}\), the center has demand two.
A single \(A,F,I,D\) response certificate pays at most one token.
It cannot be squared unless two distinct response occurrences or a
proved same-bank quadratic carrier are retained.  Private and
unbalanced banks are different: their complete second-moment theorems
pay the two-token demand without duplication.
\end{firewall}

\section{The exact remaining hot carriers}
\label{sec:V35-residual}

\begin{definition}[Single- and double-demand residuals]
\label{def:V35-residuals}
After the immediate compilers of
\Cref{prop:V35-immediate-compilers}, define:
\begin{enumerate}[label=\textup{(\roman*)},leftmargin=4em]
\item \(R_1(k;T)\): a one-token internal row whose mass is concentrated
      in a marked or previously used bank and for which no unused
      actual response has yet been exposed;
\item \(R_2(k;T)\): a star center whose two-token mass is carried by
      one or two \(A,F,I,D\)-type banks for which only linear response
      estimates have been assigned.
\end{enumerate}
Both residuals retain the complete source word, unique payer, actual
tree marking, and all bank intersections.
\end{definition}

\begin{theorem}[Hot \(2+1+1\) reduction]
\label{thm:V35-hot-reduction}
Every \(2+1+1\) source term is closed by the V67 thermal theorem or by
the group-general V32--V34 compilers, except possibly a finite sum of
the residual carriers \(R_1(k;T)\) and \(R_2(k;T)\) from
\Cref{def:V35-residuals}.  More precisely:
\begin{enumerate}[label=\textup{(\alph*)},leftmargin=4em]
\item \(R_2\) occurs only for the star terms
      \(T_{11}^{(0)}\) and \(T_{22}^{(0)}\);
\item every path term has at most two \(R_1\) carriers, one at each
      internal vertex;
\item a strict mass-ascent chain has length at most three and at each
      new row the fresh/overlap split
      \Cref{lem:V35-heavy-or-fresh} is applied only after the new cell
      is acquired; and
\item no hot leaf appears in the residual.
\end{enumerate}
\end{theorem}

\begin{proof}
Expand into the eight trees by
\Cref{lem:V35-eight-trees}.  If every demanded row is thermal, V67
applies termwise through
\eqref{eq:V35-normalization}.  At a hot demanded row apply
\Cref{lem:V35-heavy-or-fresh}.  Private and unbalanced cases close to
the required demand order; certified \(A,F,I,D\) cases close one token
by \Cref{prop:V35-immediate-compilers}.  A failed compatible-partner
comparison gives a strict ascent, which cannot revisit a row and hence
has length at most three on four vertices.

Any branch not closed in this process is, by construction, \(R_1\) or
\(R_2\).  \Cref{lem:V35-demand-census} proves
\textup{(a)--(b)}, and
\Cref{cor:V35-hot-leaves} proves \textup{(d)}.  The construction
applies the exact fresh-cell identity after each new acquired cell,
proving \textup{(c)} and excluding circular use of mass ascent.
\end{proof}

\begin{proposition}[The smallest unresolved carrier is the star
double demand]
\label{prop:V35-smallest-carrier}
Up to the interchange \(1\leftrightarrow2\), the minimal genuinely
new residual is
\[
 T_{11}^{(0)}=\{12,13,14\}
\]
with two endpoint tokens at row \(1\), one unique payer, and a
previously used or marked \(A,F,I,D\) bank carrying a fixed fraction
of \(\Xi_1\).  The required next estimate is a joint second-order
row-response bound on the complete multiplier; two copies of a
linear endpoint card are not sufficient.
\end{proposition}

\begin{proof}
Every path has only one token at a given internal row and hence first
reduces to \(R_1\).  A two-token residual can occur only in the two
stars by \Cref{lem:V35-demand-census}.  Relabeling the binary prefix
interchanges their centers.  The private and unbalanced second moments
already close, leaving the stated bank types.  The impossibility of
squaring one linear response is
\Cref{fire:V35-one-not-two}.
\end{proof}

\begin{assessment}[Upstream direction after V35]
\label{ass:V35-direction}
The companion audits point to the same proof order.  The star-center
bank must be transferred into the parent source word, not inverted as
a stand-alone response, and its two demand tokens must be paid by an
actual second moment, not by two formal marks.  V36 should form the
exact two-occurrence star-center multiplier and test whether the
five-factor Schatten sandwich already yields a quadratic response
energy.  If it gives only a linear stock, the correct next object is
the local composite covariance between the two star arms.
\end{assessment}

\begin{theorem}[Third-series V35 result ledger]
\label{thm:V35-results}
Relative to V1--V34, V35 proves:
\begin{enumerate}[label=\textup{(V35.\arabic*)},leftmargin=6.3em]
\item the compulsory SM V13/NS V31 audit with exact fingerprints;
\item the exact eight-tree census of the assembled \(2+1+1\) source;
\item the excess-demand vector and normalization identity for every
      tree;
\item proof that hot leaves require no new estimate;
\item the marked-heavy/used-bank/fresh-cell/prefix-card split at every
      demanded row;
\item closure of every private, unbalanced, and one-response
      compatible demand branch;
\item reduction of the entire hot \(2+1+1\) source to the finite
      carriers \(R_1,R_2\); and
\item localization of the first new obstruction to one two-token star
      center.
\end{enumerate}
\end{theorem}

\begin{proof}
These are
\Cref{fire:V35-audit-types,
lem:V35-eight-trees,
lem:V35-demand-census,
lem:V35-normalization,
cor:V35-hot-leaves,
lem:V35-heavy-or-fresh,
prop:V35-immediate-compilers,
thm:V35-hot-reduction,
prop:V35-smallest-carrier}.
\end{proof}

\begin{programme}[V36: joint star-center response energy]
\label{prog:V35-V36}
\begin{enumerate}[leftmargin=3em]
\item Fix \(T_{11}^{(0)}=\{12,13,14\}\) and retain the exact binary
      prefix, composite ternary joining carrier, singleton banks, and
      suffix.
\item For a center bank occurring in two source roles, form its
      complete two-occurrence multiplier before a positive envelope.
\item Apply the V94 quadratic product-state energy and V97
      three-factor sandwich with both star arms still present.
\item Prove a bound carrying two genuine center denominators, or show
      by a fixed \(SU(2)\) counterpacket that only one survives.
\item Treat a marked-heavy center coordinate by its complete local
      fourth-order carrier, never by two independent first-moment
      estimates.
\item Preserve all four payer locations from the V67 ledger.
\item If the joint energy is insufficient, replace it upstream by the
      assembled composite covariance of the two star arms.
\end{enumerate}
\end{programme}

\begin{firewall}[Claim boundary after third-series V35]
\label{fire:V35-claim-boundary}
V35 gives an exact hot-demand reduction for the \(2+1+1\) source but
does not close the \(R_1\) and \(R_2\) used/marked-bank residuals.
The \(3+1\), global discrete, local \([4]\), and \([3,3]\) source
families also remain open.  Higher MTA, WUFC, CPM, cutoff removal,
reflection positivity, Osterwalder--Schrader reconstruction,
construction of the continuum Yang--Mills measure, and a uniform
positive continuum spectral gap remain open.  The full Yang--Mills
mass gap has not been proved.
\end{firewall}

\paragraph{Parent-version record.}
V35 was copied byte-for-byte from
\texttt{Renewal\_Yang\_Mills\_Mass\_Gap\_Research\_Notes\_V34.tex}
before this chapter was appended.  The V34 parent had \(18\,026\)
source lines, \(653\,989\) bytes, and SHA--256
\[
 \texttt{bc4595361a970bfc428e23b93202d3b4874936d9b421af970883843b2aac00ee}.
\]
The next compulsory companion audit is V40.

% =============================================================================
% V36 APPENDIX CHAPTER
% =============================================================================

\clearpage
\chapter{V36: Exact occurrence audit and the balanced marked
two-debit carrier}
\label{ch:V36-star-audit}

\section{Why the proposed five-factor insertion must be corrected}
\label{sec:V36-correction}

V35 isolated the two center demands of the star
\[
 T_{11}^{(0)}=\{12,13,14\}
\]
and proposed placing its two arms at the endpoints of the archived
V97 five-factor sandwich.  That proposal is not justified by the
exact V67 source word.  The \(13\) and \(14\) factors are two scalar
influence marks obtained only after the single local operator
\[
 J_r(12|3|4)=\kappa_3(Z_{12,r},X_{3,r},X_{4,r})
\]
has been bounded.  They are not two complete self-adjoint operator
occurrences in the unsplit Fourier multiplier.

This chapter makes that distinction exact.  It proves that V97 does
close a two-demand branch whenever two genuine complete response
occurrences survive, but it cannot be invoked merely because the
positive influence polynomial contains a square or a product.  For a
marked-heavy star it then identifies the precise missing estimate:
two inverse powers of the marked center mass on one complete balanced
local carrier.

\begin{assessment}[Nonduplication boundary]
\label{ass:V36-nonduplication}
Archived V97--V98 proves a joint Schatten estimate for an already
factorized word
\[
 A_1C_1A_2C_2A_3.
\]
V29 proves the local composite-ternary influence polynomial, and V30
identifies a different discrete-prefix partition sum with a third
cumulant.  Neither result proves that the two scalar arms of the V67
polynomial are separate complete occurrences.  V36 audits this
missing factorization and does not repeat any of those theorems.
\end{assessment}

\section{The ordered star stock}
\label{sec:V36-ordered-stock}

Retain the notation
\[
 h_{ij,r}=a_{i,r}a_{j,r},
 \qquad
 H_{12}^{<r}=\sum_{e<r}h_{12,e}.
\]
Before the Cartesian enlargement in the proof of the archived V67
payer ledger, the \(T_{11}^{(0)}\) stock is
\begin{equation}
 \Omega_{11}
 :=
 \sum_r H_{12}^{<r}h_{13,r}h_{14,r}.
\label{eq:V36-ordered-star}
\end{equation}

\begin{lemma}[Exact marked-coordinate expansion]
\label{lem:V36-ordered-expansion}
One has
\begin{equation}
 \Omega_{11}
 =
 \sum_{e<r}
 a_{1,e}a_{2,e}\,
 a_{1,r}^{\,2}a_{3,r}a_{4,r}.
\label{eq:V36-star-expanded}
\end{equation}
In particular, the binary-prefix mark and joining mark are distinct,
whereas both quotient arms use the same joining coordinate \(r\).
Moreover,
\begin{equation}
 \Omega_{11}\le H_{12}H_{13}H_{14}.
\label{eq:V36-Cartesian-enlargement}
\end{equation}
\end{lemma}

\begin{proof}
Insert
\(H_{12}^{<r}=\sum_{e<r}a_{1,e}a_{2,e}\) and
\[
 h_{13,r}h_{14,r}
 =
 a_{1,r}a_{3,r}a_{1,r}a_{4,r}.
\]
This gives \eqref{eq:V36-star-expanded}.  The inequality \(e<r\)
proves distinctness.  Finally enlarge the diagonal pair of joining
marks \((r,r)\) to independent coordinates in the nonnegative sum:
\[
 \sum_{e<r}h_{12,e}h_{13,r}h_{14,r}
 \le
 \sum_{e,u,v}h_{12,e}h_{13,u}h_{14,v}.
\]
The latter is the right side of
\eqref{eq:V36-Cartesian-enlargement}.
\end{proof}

\begin{corollary}[The second center incidence is local]
\label{cor:V36-second-incidence-local}
For every monomial of \(\Omega_{11}\), the two \(1\)-incidences
belonging to the star arms are the two powers of \(a_{1,r}\) inside
one local joining carrier.  They do not define two coordinate banks
and cannot be optimized against two independently chosen product
states.
\end{corollary}

\begin{proof}
This is the factor
\(a_{1,r}^{\,2}a_{3,r}a_{4,r}\) in
\eqref{eq:V36-star-expanded}.  Both copies have the same coordinate
label \(r\) and arise from the same norm estimate on
\(J_r(12|3|4)\).
\end{proof}

The \(T_{22}^{(0)}\) formula is obtained throughout by interchanging
rows \(1\) and \(2\).

\section{What a quadratic endpoint energy does and does not buy}
\label{sec:V36-energy-audit}

\begin{definition}[Complete response occurrence]
\label{def:V36-complete-occurrence}
A response occurrence is called complete if, before an absolute
value, product-state supremum, positive envelope, or output-block
sum, it is a self-adjoint equivariant factor \(S\) in the exact
multiplier.  Its one-use energy is the archived quantity
\[
 {\cal E}_\otimes(S)
 =
 \|S\|\kappa_\otimes(|S|).
\]
A scalar influence factor appearing only in a positive bound for
\(\|S\|\) is not a second occurrence of \(S\).
\end{definition}

\begin{lemma}[Square-root accounting]
\label{lem:V36-square-root-accounting}
Suppose the hypotheses of the archived V97 five-factor theorem hold
for
\[
 A\,C_1\,B\,C_2\,E
\]
with \(\|B\|,\|C_1\|,\|C_2\|\le L\).  If
\[
 {\cal E}_\otimes(A)\le C\mathfrak h_A^2,
 \qquad
 {\cal E}_\otimes(E)\le C\mathfrak h_E^2,
\]
then
\begin{equation}
 \Gamma_\otimes(1,K)
 \le C_L\mathfrak h_A\mathfrak h_E.
\label{eq:V36-two-endpoint-product}
\end{equation}
If instead only one stock-carrying occurrence is available and the
other endpoint is the identity, the same argument yields only
\begin{equation}
 \Gamma_\otimes(1,K)
 \le C_L\mathfrak h_A.
\label{eq:V36-one-endpoint-linear}
\end{equation}
\end{lemma}

\begin{proof}
The V97 inequality is
\[
 \Gamma_\otimes(1,K)
 \le
 L^3\sqrt{
 {\cal E}_\otimes(A){\cal E}_\otimes(E)}.
\]
Insert the two quadratic energies to obtain
\eqref{eq:V36-two-endpoint-product}.  For the second statement put
\(E=I\) and use
\({\cal E}_\otimes(I)=1\).
\end{proof}

\begin{firewall}[Quadratic energy is still one use]
\label{fire:V36-quadratic-one-use}
The adjective quadratic refers to the product of a norm and a
capacity inside \({\cal E}_\otimes(S)\).  The square root in the
Schatten sandwich returns one copy of the stock of \(S\).  A square
of a stock in the final estimate requires two complete endpoints
with quadratic energies, or a different proved second-order carrier.
\end{firewall}

\begin{proposition}[Exact failure of the proposed star-arm insertion]
\label{prop:V36-no-star-arm-factorization}
The source identity \eqref{eq:V35-S211} and the local theorem used in
\eqref{eq:V35-V67-carrier} do not supply a factorization in which
\(h_{13,r}\) and \(h_{14,r}\) are the stocks of two complete endpoint
occurrences.  Consequently the V97--V98 three-occurrence theorem
cannot, on the presently proved identities, be applied with the
\(13\) and \(14\) arms as its two endpoints.
\end{proposition}

\begin{proof}
Before positive domination, the local factor in
\eqref{eq:V35-S211} is the single operator
\(J_r(12|3|4)\).  The archived local ternary theorem first applies
the operator-martingale identity to that whole third cumulant and
then bounds its norm by
\[
 Cs_\alpha^2
 \bigl(
 b_{AB,r}b_{AC,r}
 +b_{AB,r}b_{BC,r}
 +b_{AC,r}b_{BC,r}
 \bigr).
\]
Thus \(h_{13,r}\) and \(h_{14,r}\) occur on the right side of an
operator-norm inequality, not as factors of an operator equality.
Definition \ref{def:V36-complete-occurrence} therefore excludes both
as separately selectable endpoints.  The hypotheses of the
five-factor theorem have not been established for this choice.
\end{proof}

\begin{remark}[This is not a no-go theorem for the local cumulant]
\label{rem:V36-not-physical-no-go}
\Cref{prop:V36-no-star-arm-factorization} proves non-applicability of
one compiler.  It does not prove that the required local estimate is
false.  Cancellation inside the complete signed third cumulant may
supply the two-debit estimate, but that cancellation is absent from
the positive V67 influence polynomial.
\end{remark}

\section{A scalar regression for stock-only normalization}
\label{sec:V36-stock-regression}

For a nonzero \(T_{11}^{(0)}\) monomial, the V67 Cartesian majorant
and the desired normalized tree monomial differ by the exact factor
\begin{equation}
 \frac{
 s_\alpha^2H_{12}H_{13}H_{14}/
 (\Xi_1\Xi_2\Xi_3\Xi_4)}
 {H_{12}H_{13}H_{14}/
 (\Xi_1^3\Xi_2\Xi_3\Xi_4)}
 =
 (s_\alpha\Xi_1)^2.
\label{eq:V36-exact-scalar-ratio}
\end{equation}

\begin{proposition}[No stock-only hot-star bound]
\label{prop:V36-stock-only-no-go}
There is no support- and representation-uniform constant \(C\) for
which the V67 positive stock majorant of every hot star is at most
\(C\) times its normalized tree monomial.  Any successful estimate
must retain additional Wilson damping, output capacity, or signed
cumulant cancellation.
\end{proposition}

\begin{proof}
Equation \eqref{eq:V36-exact-scalar-ratio} would force
\((s_\alpha\Xi_1)^2\le C\).  Fix \(s_\alpha>0\) and take two active
coordinates with row-\(1\) Casimir weights tending to infinity, while
choosing fixed nontrivial row-\(2\), row-\(3\), and row-\(4\) weights
so that the three displayed stocks are nonzero.  Then
\(\Xi_1\to\infty\), contradicting the asserted uniform bound.
\end{proof}

\begin{firewall}[Scope of the scalar regression]
\label{fire:V36-scalar-scope}
The preceding construction is a counterexample to a scalar
stock-comparison step, not a physical Yang--Mills counterpacket.
Heat-kernel multipliers and the signs of the complete cumulant were
deliberately discarded in the V67 positive stock bound.  They are
exactly the structures which a sharper theorem may exploit.
\end{firewall}

\section{The exact marked two-debit target}
\label{sec:V36-two-debit-target}

For a monomial indexed by \(e<r\), let \(q\) denote the physical
coordinate of the sole payer when it is not already \(e\) or \(r\),
and put
\begin{equation}
 Q_{e,r,q}:=\{e,r\}\cup\{q\},
 \qquad
 A_1(Q_{e,r,q})
 :=
 \sum_{u\in Q_{e,r,q}}a_{1,u}.
\label{eq:V36-marked-mass}
\end{equation}
The set notation removes repetitions.  Let
\(\mathcal W_{e,r,q}^{\rm pay}\) denote the corresponding complete
signed source summand after the unique payer has been allocated but
before its local partition letters or Fourier output blocks are
separated.

\begin{definition}[Balanced star-collision two-debit card]
\label{def:V36-BSC2}
The \(BSC_2\) estimate at row \(1\) is
\begin{equation}
 \boxed{
 \kappa_\otimes(
 \mathcal W_{e,r,q}^{\rm pay})
 \le
 C
 \min\left\{
 s_\alpha^2,
 \frac1{A_1(Q_{e,r,q})^2}
 \right\}
 h_{12,e}h_{13,r}h_{14,r}.}
\label{eq:V36-BSC2}
\end{equation}
It is required only on the balanced protected local sector not
already covered by private or unbalanced moment compilers.  The
definition includes every payer location from the V67 ledger.
\end{definition}

\begin{lemma}[\(BSC_2\) is the exact marked-heavy interface]
\label{lem:V36-BSC2-suffices}
Suppose
\[
 A_1(Q_{e,r,q})\ge\varepsilon\Xi_1.
\]
Then \eqref{eq:V36-BSC2}, after division by the four input masses,
implies
\begin{equation}
 \frac{
 \kappa_\otimes(
 \mathcal W_{e,r,q}^{\rm pay})}
 {\Xi_1\Xi_2\Xi_3\Xi_4}
 \le
 \frac{C}{\varepsilon^2}
 \frac{
 h_{12,e}h_{13,r}h_{14,r}}
 {\Xi_1^3\Xi_2\Xi_3\Xi_4}.
\label{eq:V36-marked-tree-bound}
\end{equation}
\end{lemma}

\begin{proof}
Use the second entry of the minimum and
\[
 A_1(Q_{e,r,q})^{-2}
 \le
 \varepsilon^{-2}\Xi_1^{-2}.
\]
The remaining factor \(\prod_i\Xi_i^{-1}\) is the source
normalization.  The powers on the right are exactly the vertex
degrees of \(\{12,13,14\}\).
\end{proof}

\begin{lemma}[Exponential damping supplies two debits]
\label{lem:V36-exponential-two-debit}
If a complete local envelope satisfies
\begin{equation}
 \kappa_\otimes(
 \mathcal W_{e,r,q}^{\rm pay})
 \le
 Cs_\alpha^2
 e^{-c s_\alpha A_1(Q_{e,r,q})}
 h_{12,e}h_{13,r}h_{14,r},
\label{eq:V36-exponential-card}
\end{equation}
then it satisfies \(BSC_2\).
\end{lemma}

\begin{proof}
The V67 estimate supplies the first entry of the minimum.  For the
second, set \(t=s_\alpha A_1(Q_{e,r,q})\) and use
\[
 \sup_{t\ge0}t^2e^{-ct}<\infty.
\]
Thus
\[
 s_\alpha^2e^{-c s_\alpha A_1(Q)}
 \le
 \frac{C_c}{A_1(Q)^2}.
\]
\end{proof}

\begin{corollary}[The marked unbalanced sector is closed]
\label{cor:V36-unbalanced-marked-close}
Whenever the marked packet has the complete exponential output
separation of the V34 unbalanced-bank theorem, its two star-center
demands obey \eqref{eq:V36-marked-tree-bound}.
\end{corollary}

\begin{proof}
The V34 all-order unbalanced estimate retains exponential damping
before the fixed second moment is taken, giving
\eqref{eq:V36-exponential-card}.  Apply
\Cref{lem:V36-exponential-two-debit,
lem:V36-BSC2-suffices}.
\end{proof}

\section{Why the generic protected capacity is not enough}
\label{sec:V36-protected-no-go}

\begin{proposition}[\(SU(2)\) scaling obstruction for the generic
projector card]
\label{prop:V36-projector-insufficient}
The representation-independent invariant-projector capacity bound
proved in V32 cannot by itself imply \(BSC_2\).  Already for
\(G=SU(2)\), a spin-\(j\) row has
\[
 a_j=1+j(j+1),
 \qquad
 d_j=2j+1,
\]
whereas the generic product-state invariant-projector card is only
\(d_j^{-1}\).  There is no \(C\) such that
\begin{equation}
 \frac1{d_j}\le\frac C{a_j^2}
\label{eq:V36-projector-false-scale}
\end{equation}
for all \(j\).
\end{proposition}

\begin{proof}
One has
\[
 \frac{a_j^2}{d_j}
 =
 \frac{(1+j(j+1))^2}{2j+1}
 \sim \frac12j^3
 \longrightarrow\infty.
\]
This contradicts \eqref{eq:V36-projector-false-scale}.
\end{proof}

\begin{firewall}[The obstruction is to the interface, not the theorem]
\label{fire:V36-projector-scope}
V32 gives a sharp multiplicity-free generic projector estimate, but
it does not claim two inverse Casimir powers.  The failure of
\eqref{eq:V36-projector-false-scale} says that \(BSC_2\) must use
more than generic invariant projection: either cancellation in the
third cumulant, a stronger local character estimate, or a second
genuine response occurrence.
\end{firewall}

\section{The genuine two-occurrence branch}
\label{sec:V36-two-occurrence}

\begin{definition}[Occurrence number]
\label{def:V36-occurrence-number}
For a physical bank \(\mathcal B\) inside a fixed complete source
word \(W\), let \(\nu_{\mathcal B}(W)\) be the number of
stock-carrying complete response occurrences of \(\mathcal B\) in
the sense of \Cref{def:V36-complete-occurrence}.  Scalar marks in a
positive envelope do not increase \(\nu_{\mathcal B}(W)\).
\end{definition}

\begin{proposition}[Two genuine endpoints close the analytic
double-use gate]
\label{prop:V36-two-genuine-endpoints}
Suppose a used \(A,F,I\) bank in an \(R_2\) carrier has
\(\nu_{\mathcal B}(W)\ge2\), and the exact multiplier can be ordered
as
\[
 A_1C_1BC_2A_3
\]
with the two occurrences at \(A_1,A_3\), uniformly bounded middle
word, and row-response stocks
\(\mathfrak h_1,\mathfrak h_3\).  Then
\begin{align}
 \Gamma_\otimes(1,K)
 &\le C\mathfrak h_1\mathfrak h_3,
\label{eq:V36-genuine-unpaid}\\
 \Gamma_\otimes(w,K^{\rm pay})
 &\le
 \frac C{s_\alpha}\mathfrak h_1\mathfrak h_3
\label{eq:V36-genuine-paid}
\end{align}
for every payer position allowed by the complete word.
Consequently, whenever the acquisition ledger assigns the two
center demands to these two actual response stocks, that \(R_2\)
branch closes.
\end{proposition}

\begin{proof}
The archived V95 quadratic energy table gives
\[
 {\cal E}_\otimes(A_i)
 \le C\mathfrak h_i^2,
 \qquad i=1,3.
\]
Apply \Cref{lem:V36-square-root-accounting} for the unpaid line.
For a payer in an endpoint use its paid quadratic energy; for a payer
in the middle word or suffix use its complete \(C/s_\alpha\) norm.
This is the archived V97--V98 payer ledger and gives
\eqref{eq:V36-genuine-paid}.  The final assertion uses two distinct
occurrence slots and hence does not duplicate a capacity.
\end{proof}

\begin{firewall}[Directed covariances are not silently retyped]
\label{fire:V36-directed-separate}
The preceding proposition states \(A,F,I\), whose quadratic energy
table is proved.  A directed \(D\)-cell may be used as a uniformly
bounded middle covariance or through its proved one-token partner
edge.  It is not declared to have an \(A,F,I\) quadratic endpoint
energy without an additional theorem.
\end{firewall}

\section{The corrected residual after V36}
\label{sec:V36-residual}

\begin{theorem}[Refined star-center residual]
\label{thm:V36-refined-residual}
For the two star terms in the V35 \(2+1+1\) reduction, the
double-demand residual \(R_2\) is closed in each of the following
subcases:
\begin{enumerate}[label=\textup{(\roman*)},leftmargin=4em]
\item the center mass lies in a private or unbalanced complete bank;
\item the marked packet has the unbalanced exponential separation of
      \Cref{cor:V36-unbalanced-marked-close}; or
\item two genuine complete \(A,F,I\) occurrences satisfying
      \Cref{prop:V36-two-genuine-endpoints} are present and carry the
      two acquired response roles.
\end{enumerate}
Every remaining \(R_2\) carrier belongs to at least one of:
\begin{enumerate}[label=\textup{(\alph*)},leftmargin=4em]
\item a balanced marked packet requiring \(BSC_2\);
\item a used \(A,F,I\) bank with at most one genuine stock-carrying
      occurrence in the complete word; or
\item a used directed or coupled bank for which only the one-token
      response or middle-factor norm has been proved.
\end{enumerate}
The one-token residual \(R_1\) is unchanged.
\end{theorem}

\begin{proof}
Private and unbalanced all-order moments are
\Cref{prop:V35-immediate-compilers}.  The marked unbalanced case is
\Cref{cor:V36-unbalanced-marked-close}, and the two-occurrence case is
\Cref{prop:V36-two-genuine-endpoints}.

Apply the V35 marked-heavy/used-bank/fresh-cell/prefix split to any
unclosed \(R_2\).  The fresh private and unbalanced alternatives have
just closed.  A marked-heavy alternative is either unbalanced, hence
closed, or balanced and is exactly item \textup{(a)} by
\Cref{lem:V36-BSC2-suffices}.  In a used \(A,F,I\) bank, two proved
complete occurrences close; failure of that hypothesis is
\textup{(b)}.  The remaining admitted used-bank type is directed or
coupled and gives \textup{(c)}.  The list is exhaustive by
\Cref{lem:V35-heavy-or-fresh}.  No argument in this chapter adds a
second response to a one-demand path row, so \(R_1\) remains as
stated.
\end{proof}

\begin{assessment}[Upstream object selected for V37]
\label{ass:V36-upstream}
The smallest residual is no longer an abstract need for a
second-order response.  It is the balanced protected part of the
complete signed two-site word
\[
 K_{12}^{<r,\mathrm{conn}}
 \otimes
 \kappa_3(Z_{12,r},X_{3,r},X_{4,r}),
\]
with \(e<r\), the unique payer retained, and the \(r\)-coordinate
carrying both star arms.  Generic heat and invariant-projector bounds
are insufficient.  V37 must expand only the balanced protected output
of this signed word and test whether its cumulant M\"obius
cancellation removes the leading invariant channel.  If it does not,
the correct response is a physical \(SU(2)\) counterpacket to
\eqref{eq:V36-BSC2}, not another positive-envelope decomposition.
\end{assessment}

\begin{theorem}[Third-series V36 result ledger]
\label{thm:V36-results}
Relative to V1--V35, V36 proves:
\begin{enumerate}[label=\textup{(V36.\arabic*)},leftmargin=6.3em]
\item the exact ordered-coordinate stock for both \(2+1+1\) star
      terms;
\item proof that the two quotient arms share one joining coordinate;
\item the one-use square-root accounting for quadratic endpoint
      energies;
\item non-applicability of the V97 endpoint sandwich to scalar arms
      of the already-enveloped local cumulant;
\item a scalar no-go theorem for stock-only hot-star normalization;
\item the exact \(BSC_2\) estimate sufficient for a marked-heavy
      center;
\item exponential closure of its unbalanced sector;
\item proof that generic invariant-projector capacity cannot supply
      the two required Casimir debits;
\item closure of every branch with two genuine complete \(A,F,I\)
      endpoints; and
\item reduction of the remaining double-demand sector to three
      precisely typed carriers.
\end{enumerate}
\end{theorem}

\begin{proof}
These are
\Cref{lem:V36-ordered-expansion,
cor:V36-second-incidence-local,
lem:V36-square-root-accounting,
prop:V36-no-star-arm-factorization,
prop:V36-stock-only-no-go,
def:V36-BSC2,
lem:V36-BSC2-suffices,
lem:V36-exponential-two-debit,
cor:V36-unbalanced-marked-close,
prop:V36-projector-insufficient,
prop:V36-two-genuine-endpoints,
thm:V36-refined-residual}.
\end{proof}

\begin{programme}[V37: balanced protected star collision]
\label{prog:V36-V37}
\begin{enumerate}[leftmargin=3em]
\item Freeze one ordered pair \(e<r\) and one payer location; keep the
      binary prefix and local third cumulant signed and complete.
\item Resolve only the balanced protected output at coordinate \(r\).
      Do not expand the gapped or unbalanced parts again.
\item Write the exact M\"obius expansion of
      \(\kappa_3(Z_{12,r},X_{3,r},X_{4,r})\) on that protected output
      and identify its leading invariant channel.
\item Test whether the constant-output component cancels before
      absolute values.  If so, quantify the first surviving Casimir
      or character difference to order two.
\item If it does not cancel, construct a physical \(SU(2)\) family
      with the binary prefix, joining coordinate, and payer all
      realized, and compare it directly with
      \eqref{eq:V36-BSC2}.
\item Keep the single-occurrence and directed used-bank residuals
      separate until the local balanced decision is known.
\item No companion audit is due before V40.
\end{enumerate}
\end{programme}

\begin{firewall}[Claim boundary after third-series V36]
\label{fire:V36-claim-boundary}
V36 corrects the proposed star-arm factorization and closes several
genuine two-debit subcases.  It does not prove the balanced
\(BSC_2\) card, the single-occurrence or directed \(R_2\) branches,
or the \(R_1\) carriers.  Hence the global \(2+1+1\) source remains
open.  The \(3+1\), global discrete, local \([4]\), and \([3,3]\)
families, higher MTA, WUFC, CPM, cutoff removal, reflection
positivity, Osterwalder--Schrader reconstruction, construction of the
continuum Yang--Mills measure, and a uniform positive continuum
spectral gap also remain open.  The full Yang--Mills mass gap has not
been proved.
\end{firewall}

\paragraph{Parent-version record.}
V36 was copied byte-for-byte from
\texttt{Renewal\_Yang\_Mills\_Mass\_Gap\_Research\_Notes\_V35.tex}
before this chapter was appended.  The V35 parent had \(18\,558\)
source lines, \(673\,335\) bytes, and SHA--256
\[
 \texttt{8b375ebc8a980641df66eecf2656b33ce5ed45908ebf30ff8d0b9090ad478c88}.
\]
The next compulsory companion audit is V40.

% =============================================================================
% V37 APPENDIX CHAPTER
% =============================================================================

\clearpage
\chapter{V37: Five-letter reassembly and the binary--ternary
star endpoint}
\label{ch:V37-five-letter}

\section{Archive reconciliation before another estimate}
\label{sec:V37-archive-reconciliation}

The upstream search required by V36 reaches two archived statements
which must be read together.  Frozen f2 V34 proves the complete
one-selected incidence-four junction carrier for \(SU(2)\).  Frozen
f2 V36 later withdraws only the global conclusion obtained by summing
possible four-block locations: the restricted prefix
\(\bigotimes(M-K_4)\) is not a contraction.  The local selected
carrier, invariant-space normal form, and mixed-incidence bank
compilers remain valid.

This distinction matters here.  The coordinate \(r\) in
\eqref{eq:V35-S211} is selected uniquely by first connectivity, so
its fourth-cumulant letter does not incur the global overcount.
However, the local joining carrier is not just that letter.  It is a
composite third cumulant containing five exact connected
moment-partition terms.  The selected \([4]\) theorem can be applied
to the fourth-cumulant term, but not to the other four by relabelling.

\begin{record}[Exact archived scope]
\label{rec:V37-archived-scope}
The following statements are retained:
\begin{enumerate}[label=\textup{(\roman*)},leftmargin=4em]
\item the \(SU(2)\) one-selected four-row junction theorem;
\item its full invariant-isotypic projector, with no intermediate-spin
      multiplicity loss; and
\item the correction that no global multiple-selection theorem
      follows from it without controlling the restricted
      \(M-K_4\) prefix.
\end{enumerate}
No global embedded-\([4]\) theorem is used below.
\end{record}

\begin{firewall}[Relation to active V32]
\label{fire:V37-G-boundary}
Active V32 replaces the one-row invariant-projector capacity,
balanced partner, and unbalanced output lemmas by theorems for every
fixed compact connected semisimple \(G\).  It explicitly treats
marked coordinates only as finite separators and does not replay the
full selected-junction payer proof.  Therefore
\Cref{rec:V37-archived-scope} closes the selected fourth-cumulant
component for \(SU(2)\); its extension to general \(G\) is retained
as a separate typed interface rather than inferred silently.
\end{firewall}

\section{The five exact local roles}
\label{sec:V37-five-roles}

At the joining coordinate suppress \(r\) and put \(X_i=X_{i,r}\).
The active V30 Leonov--Shiryaev identity, relabelled by
\((y,z,x,q)=(1,2,3,4)\), is
\begin{align}
 \kappa_3(X_1X_2,X_3,X_4)
 ={}&
 \kappa_4(X_1,X_2,X_3,X_4)
\notag\\
 &+\mathbb E X_1\,\kappa_3(X_2,X_3,X_4)
  +\mathbb E X_2\,\kappa_3(X_1,X_3,X_4)
\notag\\
 &+\operatorname{Cov}(X_1,X_3)
       \operatorname{Cov}(X_2,X_4)
\notag\\
 &+\operatorname{Cov}(X_1,X_4)
       \operatorname{Cov}(X_2,X_3).
\label{eq:V37-five-role-identity}
\end{align}
The identity is formed before a norm, positive envelope, output
resolution, or payer allocation.

\begin{assessment}[Nonduplication]
\label{ass:V37-nonduplication}
Equation \eqref{eq:V37-five-role-identity} is not a new
product--cumulant formula: it is the specialization of
\Cref{thm:V30-discrete-composite}.  What is new is its use to audit
the occurrence structure of the global \(12|3|4\) source and to
separate the selected-four-row interface from the binary--ternary
overlap.
\end{assessment}

Let
\begin{align*}
 C_4&:=\kappa_4(X_1,X_2,X_3,X_4),\\
 C_{31}^{(1)}
 &:=\mathbb E X_1\,\kappa_3(X_2,X_3,X_4),\\
 C_{31}^{(2)}
 &:=\mathbb E X_2\,\kappa_3(X_1,X_3,X_4),\\
 C_{22}^{(a)}
 &:=\operatorname{Cov}(X_1,X_3)
       \operatorname{Cov}(X_2,X_4),\\
 C_{22}^{(b)}
 &:=\operatorname{Cov}(X_1,X_4)
       \operatorname{Cov}(X_2,X_3).
\end{align*}

\begin{lemma}[Exact lower-role tree table]
\label{lem:V37-lower-role-table}
Adjoin the already connected prefix edge \(12\).  Expanding only the
lower-order connected factors in
\eqref{eq:V37-five-role-identity} gives exactly the eight trees of
\eqref{eq:V35-eight-trees}:
\begin{center}
\begin{tabular}{@{}p{0.21\linewidth}p{0.66\linewidth}@{}}
\toprule
local role & resulting row trees\\
\midrule
\(C_{31}^{(1)}\)&
\(\{12,23,24\}=T_{22}^{(0)}\),
\(\{12,23,34\}=T_2^{(3)}\),
\(\{12,24,34\}=T_2^{(4)}\);\\
\(C_{31}^{(2)}\)&
\(\{12,13,14\}=T_{11}^{(0)}\),
\(\{12,13,34\}=T_1^{(3)}\),
\(\{12,14,34\}=T_1^{(4)}\);\\
\(C_{22}^{(a)}\)&
\(\{12,13,24\}=T_{12}^{(0)}\);\\
\(C_{22}^{(b)}\)&
\(\{12,14,23\}=T_{21}^{(0)}\).\\
\bottomrule
\end{tabular}
\end{center}
Every star double demand in this lower-role table occurs in exactly
one binary--ternary product.
\end{lemma}

\begin{proof}
The ternary connected factor on rows \(\{2,3,4\}\) has the three
trees
\[
 \{23,24\},\qquad\{23,34\},\qquad\{24,34\}.
\]
Adjoining \(12\) gives the first table row.  The corresponding three
trees on \(\{1,3,4\}\) give the second row.  Each covariance product
already supplies its two displayed disjoint edges, giving the last
two rows after \(12\) is adjoined.  Comparison with
\eqref{eq:V35-tree-zero}--\eqref{eq:V35-tree-four} proves that every
one of the eight trees occurs once.  Only
\(\{12,23,24\}\) and \(\{12,13,14\}\) have a vertex of degree three.
\end{proof}

\begin{corollary}[The star is not a three-copy same-bank word]
\label{cor:V37-star-two-occurrences}
For the star centered at row \(1\), the exact lower-order word is
\begin{equation}
 K_{12}^{<r,\mathrm{conn}}
 \otimes
 \mathbb E_rX_{2,r}
 \otimes
 \kappa_{3,r}(X_{1,r},X_{3,r},X_{4,r}),
\label{eq:V37-binary-ternary-word}
\end{equation}
up to the complete prefix singleton and suffix factors.  Its two
center roles are one binary connected occurrence and one ternary
connected occurrence.  They are not separate \(13\) and \(14\)
occurrences, and they are not three occurrences of a common bank.
\end{corollary}

\begin{proof}
The only lower-order row producing \(T_{11}^{(0)}\) in
\Cref{lem:V37-lower-role-table} is \(C_{31}^{(2)}\).  Restoring its
binary prefix gives \eqref{eq:V37-binary-ternary-word}.  The ternary
factor is complete and contains both local arms; V36 proves that its
two scalar influence marks cannot be counted as two operator
occurrences.
\end{proof}

\section{The selected fourth-order component}
\label{sec:V37-selected-four}

\begin{definition}[General-\(G\) selected-junction interface]
\label{def:V37-GSEL4}
\(GSEL_4\) is the assertion that a complete quartic word with one
first-connectivity-selected four-row local cumulant, arbitrary
complete lower-order prefix factors, complete positive moment
spectators, and one allocated final payer obeys the literal
coefficient-sensitive marked-tree bound, uniformly in support and
representations for the fixed group \(G\).
\end{definition}

\begin{proposition}[Exact status of the \(C_4\) component]
\label{prop:V37-C4-status}
The \(C_4\) contribution to
\eqref{eq:V37-five-role-identity} satisfies \(GSEL_4\) for
\(G=SU(2)\).  For a general fixed compact connected semisimple group,
active V32 supplies the invariant-capacity, balanced-partner, and
unbalanced-output inputs, but the marked selected-junction payer
argument has not yet been proved in the active series.
\end{proposition}

\begin{proof}
The source coordinate \(r\) in
\eqref{eq:V35-S211} is the unique first coordinate at which the
prefix partition \(12|3|4\) joins.  Selecting the \(C_4\) letter at
that coordinate is therefore disjoint; there is no sum over
alternative selected four-block letters.  It lies exactly in the
one-selected carrier retained by frozen f2 V36 from the f2 V34
theorem, proving the \(SU(2)\) assertion.  The general-\(G\) inputs
listed in the second assertion are
\Cref{lem:V32-balanced-partner,
lem:V32-unbalanced-output,
thm:V32-invariant-capacity}.  The marked-coordinate limitation is
\Cref{rem:V32-marked-separators}; hence it would be circular to claim
the full payer conclusion from those inputs alone.
\end{proof}

\begin{firewall}[The withdrawn theorem is not restored]
\label{fire:V37-no-global-four}
The proof of \Cref{prop:V37-C4-status} uses the already selected
first-connectivity coordinate \(r\).  It does not sum the first
four-block coordinate through a restricted product
\(\bigotimes(M-K_4)\), and it proves no global \([4]\) aggregation
statement.
\end{firewall}

\section{The protected ternary endpoint does not vanish}
\label{sec:V37-protected-test}

The exact word \eqref{eq:V37-binary-ternary-word} suggests trying to
annihilate its protected local output by cumulant cancellation.  The
following physical harmonic test rules out vanishing as a general
identity.

\begin{proposition}[Nonzero invariant channel for normalized
\(SU(2)\) characters]
\label{prop:V37-character-noncancellation}
Let \(\chi_j\) be the spin-\(j\) character of \(SU(2)\),
\(d_j=2j+1\), and \(x_j=\chi_j/d_j\).  With Haar expectation,
\begin{equation}
 \boxed{
 \kappa_3(x_j^2,x_j,x_j)
 =
 \frac{d_j-1}{d_j^4}>0
 \qquad(j>0).}
\label{eq:V37-character-cumulant}
\end{equation}
Thus the invariant component of the composite ternary correction is
not killed identically by centering.
\end{proposition}

\begin{proof}
Orthogonality of irreducible characters gives
\[
 \mathbb E x_j=0,
 \qquad
 \mathbb E x_j^2=d_j^{-2}.
\]
The multiplicity-free Clebsch--Gordan rule
\[
 V_j\otimes V_j
 =
 \bigoplus_{\ell=0}^{2j}V_\ell
\]
contains \(d_j=2j+1\) summands.  Hence
\[
 \mathbb E\chi_j^4
 =
 \sum_{\ell=0}^{2j}1=d_j,
 \qquad
 \mathbb E x_j^4=d_j^{-3}.
\]
Since the last two arguments have mean zero,
\[
 \kappa_3(x_j^2,x_j,x_j)
 =
 \mathbb E x_j^4
 -
 (\mathbb E x_j^2)^2
 =
 d_j^{-3}-d_j^{-4},
\]
which is \eqref{eq:V37-character-cumulant}.
\end{proof}

\begin{firewall}[What the character test proves]
\label{fire:V37-character-scope}
The Haar calculation is not a counterexample to the desired
heat-kernel estimate and contains no global prefix or payer.  It
proves only that the next argument cannot delete the protected
binary--ternary word by asserting that its local third cumulant
vanishes on invariant output.
\end{firewall}

\section{The exact binary--ternary endpoint target}
\label{sec:V37-BTE}

For the row-\(1\) star define its literal normalized stock
\begin{equation}
 {\mathfrak t}_{1;r}
 :=
 \frac{
 H_{12}^{<r}h_{13,r}h_{14,r}}
 {\Xi_1^3\Xi_2\Xi_3\Xi_4}.
\label{eq:V37-star-stock}
\end{equation}
Let \(\mathcal W_{BT,r}^{\rm pay}\) be the complete word
\eqref{eq:V37-binary-ternary-word}, with its prefix singleton
moments, suffix, physical output compression, and exactly one payer
retained.

\begin{definition}[Binary--ternary endpoint estimate]
\label{def:V37-BTE2}
The \(BTE_2\) estimate is
\begin{equation}
 \boxed{
 \frac{
 \kappa_\otimes(\mathcal W_{BT,r}^{\rm pay})}
 {\Xi_1\Xi_2\Xi_3\Xi_4}
 \le
 C{\mathfrak t}_{1;r}.}
\label{eq:V37-BTE2}
\end{equation}
The symmetric estimate interchanges rows \(1,2\).
\end{definition}

\begin{lemma}[A normalized two-endpoint energy is sufficient]
\label{lem:V37-BTE-energy-suffices}
Suppose the exact normalized multiplier of
\(\mathcal W_{BT,r}^{\rm pay}\) admits the V97 order
\[
 A_{12}\,C_1\,B\,C_2\,E_{134},
\]
where \(A_{12}\) is the complete binary-prefix occurrence,
\(E_{134}\) is the complete ternary occurrence, and the other factors
have uniformly bounded complete norms.  Suppose further that their
coefficient-normalized one-use energies satisfy
\begin{align}
 {\cal E}_\otimes(A_{12})
 &\le
 C\left(
 \frac{H_{12}^{<r}}{\Xi_1\Xi_2}
 \right)^2,
\label{eq:V37-binary-energy}\\
 {\cal E}_\otimes(E_{134})
 &\le
 C\left(
 \frac{h_{13,r}h_{14,r}}
 {\Xi_1^2\Xi_3\Xi_4}
 \right)^2,
\label{eq:V37-ternary-star-energy}
\end{align}
with the corresponding paid estimate at whichever factor contains
the unique payer.  Then \(BTE_2\) holds.
\end{lemma}

\begin{proof}
Apply the V97 five-factor sandwich to the two complete endpoints.
Its geometric mean, using
\eqref{eq:V37-binary-energy}--%
\eqref{eq:V37-ternary-star-energy}, is
\[
 C
 \frac{H_{12}^{<r}}{\Xi_1\Xi_2}
 \frac{h_{13,r}h_{14,r}}
 {\Xi_1^2\Xi_3\Xi_4}
 =
 C{\mathfrak t}_{1;r}.
\]
The paid endpoint or paid middle-word variant contributes the single
already allocated payer factor and no second resolvent, exactly as in
\Cref{prop:V36-two-genuine-endpoints}.  This proves
\eqref{eq:V37-BTE2}.
\end{proof}

\begin{proposition}[The missing energy is ternary, not binary]
\label{prop:V37-missing-ternary-energy}
The active and frozen binary \(A,F,I\) energy tables supply the
binary endpoint form \eqref{eq:V37-binary-energy} whenever its
certified prefix branch is selected.  No theorem used in V1--V36
supplies the coefficient-normalized quadratic ternary endpoint
energy \eqref{eq:V37-ternary-star-energy} with all local output
blocks and payer positions retained.
\end{proposition}

\begin{proof}
The binary statement is the quadratic endpoint interface audited in
\Cref{prop:V34-terminal-interfaces}.  The local composite-ternary
theorem used in V29, V30, and V35 is an ordinary-norm influence-tree
bound with scales \(s_\alpha^2{\cal Q}_r\) and
\(s_\alpha{\cal Q}_r\).  It does not state a product-state energy for
the complete ternary multiplier after coefficient normalization.
The archived normalized three-cluster response already assigns the
final resolvent to the ternary carrier; using it as an endpoint while
also allowing the prefix to pay would not prove the one-payer
two-occurrence statement.  Therefore
\eqref{eq:V37-ternary-star-energy} is a genuinely missing interface.
\end{proof}

\section{Corrected residual after five-letter reassembly}
\label{sec:V37-residual}

\begin{theorem}[Exact double-demand reduction after V37]
\label{thm:V37-double-demand-reduction}
After applying the exact identity
\eqref{eq:V37-five-role-identity}, every double-demand contribution
to the two \(2+1+1\) star sources belongs to one of:
\begin{enumerate}[label=\textup{(\roman*)},leftmargin=4em]
\item the selected fourth-cumulant interface \(GSEL_4\);
\item the binary--ternary endpoint \(BTE_2\) in
      \eqref{eq:V37-binary-ternary-word}; or
\item the private, unbalanced, or genuine two-\(A,F,I\)-endpoint
      branches already closed in V35--V36.
\end{enumerate}
The covariance-product corrections and the two nonstar trees of each
triple correction have only one demand at each internal row and hence
belong to the \(R_1\), not \(R_2\), ledger.
\end{theorem}

\begin{proof}
The five local roles are exhaustive by
\eqref{eq:V37-five-role-identity}.  The fourth-cumulant role is
\Cref{def:V37-GSEL4}.  By
\Cref{lem:V37-lower-role-table}, the only lower-order stars are the
first tree of \(C_{31}^{(1)}\) and the first tree of
\(C_{31}^{(2)}\); their exact words are the two row permutations of
\eqref{eq:V37-binary-ternary-word}, hence are \(BTE_2\).
All other lower-role trees are paths and have demand vector supported
once on each of two internal rows by
\Cref{lem:V35-demand-census}.  The already closed bank alternatives
are \Cref{thm:V36-refined-residual}.  No other local role or demand
pattern remains.
\end{proof}

\begin{assessment}[Upstream direction after V37]
\label{ass:V37-upstream}
The protected local third cumulant cannot be discarded, but it need
not be forced into two fictitious arm occurrences.  The exact object
is one complete ternary endpoint paired with the binary prefix.
V38 should prove or disprove
\eqref{eq:V37-ternary-star-energy} by applying the product-state
Hilbert--Schmidt square before the ternary martingale tree is
positively expanded.  In parallel it should replay the selected
four-row payer proof with the active general-\(G\) inputs.  If the
ternary energy fails, the next object is its squared
Leonov--Shiryaev multiplier, not another row-cell decomposition.
\end{assessment}

\begin{theorem}[Third-series V37 result ledger]
\label{thm:V37-results}
Relative to V1--V36, V37 proves:
\begin{enumerate}[label=\textup{(V37.\arabic*)},leftmargin=6.3em]
\item reconciliation of the retained one-selected incidence-four
      theorem with the withdrawn global \([4]\) claim;
\item the exact five-role decomposition of the \(2+1+1\) joining
      carrier;
\item a bijection from its four lower-order roles to the eight V35
      trees;
\item localization of both lower-order stars to binary--ternary
      products;
\item nonvanishing of their invariant local channel on an exact
      normalized-character family;
\item the precise \(BTE_2\) marked-tree target;
\item a sufficient two-endpoint energy theorem for \(BTE_2\);
\item identification of the missing input as a normalized ternary
      quadratic energy; and
\item the exhaustive three-class double-demand residual.
\end{enumerate}
\end{theorem}

\begin{proof}
These are
\Cref{rec:V37-archived-scope,
lem:V37-lower-role-table,
cor:V37-star-two-occurrences,
prop:V37-C4-status,
prop:V37-character-noncancellation,
def:V37-BTE2,
lem:V37-BTE-energy-suffices,
prop:V37-missing-ternary-energy,
thm:V37-double-demand-reduction}.
\end{proof}

\begin{programme}[V38: ternary endpoint energy and general-\(G\)
selected junction]
\label{prog:V37-V38}
\begin{enumerate}[leftmargin=3em]
\item Keep the exact complete multiplier of
      \(\kappa_3(X_1,X_3,X_4)\) before its martingale tree envelope.
\item Compute its product-state energy
      \(\|E_{134}\|\kappa_\otimes(|E_{134}|)\) with the coefficient
      normalization already inserted.
\item Split protected and gapped outputs only inside that single
      energy.  Test the protected term on
      \eqref{eq:V37-character-cumulant}.
\item Prove \eqref{eq:V37-ternary-star-energy}, or produce a physical
      representation family for which its two sides have incompatible
      scaling.
\item Separately audit the f2 V26 local output moments and f2 V34
      selected-junction payer proof against
      \Cref{thm:V32-GPG}; state the first marked step not covered for
      general \(G\).
\item Do not revisit covariance-product path corrections until the
      \(R_2\) decision is complete.
\item The next companion audit remains V40.
\end{enumerate}
\end{programme}

\begin{firewall}[Claim boundary after third-series V37]
\label{fire:V37-claim-boundary}
V37 gives an exact algebraic and occurrence-level reduction of the
\(2+1+1\) double-demand sector.  It does not prove the ternary energy
\eqref{eq:V37-ternary-star-energy}, general-\(G\) \(GSEL_4\), the
remaining \(R_1\) paths, or therefore the full \(2+1+1\) source.
The \(3+1\), global discrete, globally aggregated \([4]\), and
\([3,3]\) families, higher MTA, WUFC, CPM, cutoff removal, reflection
positivity, Osterwalder--Schrader reconstruction, construction of the
continuum Yang--Mills measure, and a uniform positive continuum
spectral gap remain open.  The full Yang--Mills mass gap has not been
proved.
\end{firewall}

\paragraph{Parent-version record.}
V37 was copied byte-for-byte from
\texttt{Renewal\_Yang\_Mills\_Mass\_Gap\_Research\_Notes\_V36.tex}
before this chapter was appended.  The V36 parent had \(19\,215\)
source lines, \(695\,714\) bytes, and SHA--256
\[
 \texttt{585e11674a8b44b9ea6e8459f242673e4f460949649c15287891d3e3e5279fea}.
\]
The next compulsory companion audit is V40.

% =============================================================================
% V38 APPENDIX CHAPTER
% =============================================================================

\clearpage
\chapter{V38: Compact-group selected-junction closure and an energy
interface no-go}
\label{ch:V38-G-selected}

\section{The two independent tasks left by V37}
\label{sec:V38-purpose}

V37 separated the star double-demand sector into a selected local
fourth cumulant and a binary--ternary endpoint word.  Their analytic
issues are different.  The selected fourth cumulant requires a
compact-group replay of an archived \(SU(2)\) theorem.  The ternary
endpoint requires a new coefficient-normalized energy, not another
representation classification.

The first task is completed in this chapter.  The active V32
highest-weight and invariant-capacity theorems replace every
\(SU(2)\)-specific representation input in the retained
one-selected-junction proof.  For the second task, a maximally
entangled projector gives a sharp logical obstruction: the known
product-state capacity of an endpoint cannot simply be squared to
obtain its one-use energy.  Thus the remaining proof must retain the
signed ternary multiplier in its ordinary norm factor.

\begin{assessment}[Nonduplication boundary]
\label{ass:V38-nonduplication}
This chapter does not reprove the frozen one-selected-junction
case analysis.  It proves the only local output-moment input whose
group dependence had not been audited, lists every
representation-specific premise of that case analysis, and replaces
each by an active compact-group theorem.  The global multiple-junction
withdrawal remains in force.
\end{assessment}

\section{Compact-group local output moments}
\label{sec:V38-local-moments}

Let \(G\) be the fixed compact connected semisimple group of V32.
At one four-active coordinate, let
\[
 K_4=\kappa_4(X_1,X_2,X_3,X_4)
\]
be the complete signed local fourth cumulant.  Resolve one complete
physical output block \(U\), including its full multiplicity space,
and denote its compression by \(K_{4,U}\).  Put
\[
 a_U:=1+C_2(U).
\]

\begin{lemma}[Fixed-arity fusion in Casimir mass]
\label{lem:V38-Casimir-fusion}
If \(U\) is an irreducible constituent obtained from at most four
intermediate block outputs \(S_1,\ldots,S_b\), then
\begin{equation}
 a_U\le C_G\sum_{\nu=1}^b a_{S_\nu}.
\label{eq:V38-Casimir-fusion}
\end{equation}
\end{lemma}

\begin{proof}
By \Cref{thm:V32-tensor-triangle},
\[
 \|\lambda_U\|
 \le
 \sum_{\nu=1}^b\|\lambda_{S_\nu}\|.
\]
Since \(b\le4\), Cauchy--Schwarz gives
\[
 1+\|\lambda_U\|^2
 \le
 1+4\sum_\nu\|\lambda_{S_\nu}\|^2
 \le
 4\sum_\nu(1+\|\lambda_{S_\nu}\|^2).
\]
Apply the two sides of
\eqref{eq:V32-Casimir-comparison}.
\end{proof}

\begin{lemma}[One-link debits through cubic order]
\label{lem:V38-one-link-debits}
For \(k=0,1,2,3\), every normalized heat multiplier
\(q_{\alpha,S}\le e^{-s_\alpha C_2(S)}\) obeys
\begin{equation}
 (s_\alpha a_S)^kq_{\alpha,S}\le C_{k,G}
\label{eq:V38-one-link-debits}
\end{equation}
uniformly for \(0<s_\alpha\le s_0\) and every irreducible \(S\).
\end{lemma}

\begin{proof}
If \(S\) is trivial, the left side is \(s_\alpha^k\le
\max\{1,s_0^3\}\).  Otherwise
\(a_S=1+C_2(S)\le C_{G,s_0}C_2(S)\), using the positive Casimir floor
\eqref{eq:V32-Casimir-floor}.  With
\(t=s_\alpha C_2(S)\), the assertion follows from
\(\sup_{t\ge0}t^ke^{-t}<\infty\).
\end{proof}

\begin{theorem}[Compact-group complete local output moments]
\label{thm:V38-G-local-output-moments}
For \(m=1,2,3\),
\begin{equation}
 \boxed{
 (s_\alpha a_U)^m\|K_{4,U}\|
 \le C_{m,G}.}
\label{eq:V38-G-local-output-moments}
\end{equation}
The constants are independent of the four input representations,
the resolved output, and every intermediate or final multiplicity.
\end{theorem}

\begin{proof}
Use the finite moment--cumulant identity
\[
 K_{4,U}
 =
 \sum_{\pi\in\Pi_4}
 (|\pi|-1)!(-1)^{|\pi|-1}M_{\pi,U}.
\]
Fix \(\pi\) and a complete resolved intermediate path.  If \(S_B\)
is the output of block \(B\in\pi\), all recouplings and compressions
are contractions, so
\[
 \|M_{\pi,U}\|
 \le
 \prod_{B\in\pi}q_{\alpha,S_B}.
\]
By \Cref{lem:V38-Casimir-fusion},
\[
 (s_\alpha a_U)^m
 \le
 C_{m,G}
 \left(\sum_{B\in\pi}s_\alpha a_{S_B}\right)^m.
\]
Expand the last power as ordered \(m\)-tuples of blocks.  If a block
\(B\) appears \(k_B\le m\) times, keep all its powers with its single
heat multiplier and apply
\eqref{eq:V38-one-link-debits}.  Unselected block multipliers are at
most one.  There are at most four blocks and \(m\le3\), so the number
of assignments is fixed.  Finally sum the fifteen partitions with
their fixed M\"obius coefficients.

The argument bounds a complete multiplicity block at every step.
Orthogonal physical compression therefore introduces neither an
intermediate-channel nor a final-output count.
\end{proof}

\begin{corollary}[Hot selected output decay]
\label{cor:V38-hot-selected-output}
On \(s_\alpha a_U>1\),
\begin{equation}
 \frac{a_U}{1-q_{\alpha,U}}\|K_{4,U}\|
 \le
 \frac{C_G}{s_\alpha^3a_U^2}.
\label{eq:V38-hot-selected-output}
\end{equation}
\end{corollary}

\begin{proof}
The hot gap gives
\((1-q_{\alpha,U})^{-1}\le C_G\).  Apply
\Cref{thm:V38-G-local-output-moments} with \(m=3\) and multiply by
\(a_U\).
\end{proof}

\section{General-\(G\) one-selected junction}
\label{sec:V38-GSEL4}

\begin{definition}[Admissible selected-junction packet]
\label{def:V38-admissible-selected}
An admissible \(GSEL_4\) packet consists of:
\begin{enumerate}[label=\textup{(\roman*)},leftmargin=4em]
\item one coordinate \(r\) selected by an exact first-connectivity
      identity and carrying the complete signed four-row cumulant
      \(K_{4,r}\);
\item complete lower-order connected prefix carriers and complete
      positive moment spectators, with their coordinate order intact;
\item unmarked shared coordinates split by their actual incidence
      \(2,3,\) or \(4\), then by the V32 balanced/unbalanced
      alternative;
\item exact physical output compression with full multiplicity
      spaces; and
\item exactly one final output mass and resolvent, allocated to the
      selected junction, one complete shared or private bank, a
      finite separator, or the complete suffix.
\end{enumerate}
No selection or sum over a possible second four-block letter is
included.
\end{definition}

\begin{lemma}[Representation-input replacement table]
\label{lem:V38-replacement-table}
Every representation-dependent input of the retained frozen
one-selected \(SU(2)\) junction proof has the following
compact-group replacement:
\begin{center}
\begin{tabular}{@{}p{0.42\linewidth}p{0.47\linewidth}@{}}
\toprule
required input & compact-group theorem\\
\midrule
Casimir fusion through fixed arity&
\Cref{lem:V38-Casimir-fusion};\\
comparable partner and squared dominant-row mass&
\Cref{lem:V32-balanced-partner,cor:V32-bank-power};\\
absence of invariant output and retained dominant Casimir in the
unbalanced case&
\Cref{lem:V32-unbalanced-output};\\
full invariant-space product capacity without a multiplicity count&
\Cref{thm:V32-invariant-capacity,
cor:V32-protected-capacity};\\
unmarked protected/gapped complete-bank currencies&
\Cref{thm:V32-GPG,cor:V32-V70-currencies};\\
selected local output moments through order three&
\Cref{thm:V38-G-local-output-moments,
cor:V38-hot-selected-output}.\\
\bottomrule
\end{tabular}
\end{center}
\end{lemma}

\begin{proof}
The entries are direct citations.  They are exhaustive because the
retained proof uses representation theory only to: compare an output
mass to finitely many inputs; select a comparable input at a balanced
coordinate; force output mass at an unbalanced coordinate; control
the all-invariant projection; sum unmarked protected/gapped banks; and
control the resolved selected-junction payer.  Its remaining steps are
the finite tree census, Cauchy--Schwarz, Bernoulli failure debits,
orthogonality of output sectors, and the one-payer allocation.
\end{proof}

\begin{theorem}[Compact-group one-selected four-row junction]
\label{thm:V38-GSEL4}
For every fixed compact connected semisimple \(G\), every admissible
selected-junction packet of
\Cref{def:V38-admissible-selected} obeys the coefficient-sensitive
once-paid quartic marked-tree estimate, with a constant depending
only on \(G\) and the fixed decomposition thresholds.  Equivalently,
\(GSEL_4\) of \Cref{def:V37-GSEL4} holds.
\end{theorem}

\begin{proof}
Replay the retained one-selected-junction proof, keeping its proof
order.  On unmarked balanced banks, assign each coordinate to the
comparable physical partner supplied by
\Cref{lem:V32-balanced-partner}.  Then
\eqref{eq:V32-bank-power} replaces the \(SU(2)\) spin-square
inequality, while
\Cref{thm:V32-invariant-capacity,cor:V32-V70-currencies} controls the
protected and gapped output sums without a multiplicity factor.

On an unmarked unbalanced bank,
\eqref{eq:V32-unbalanced-output-Casimir} retains the dominant row
mass in every output.  The fixed-order heat moments therefore give
the quadratic and cubic bank debits used by the cold and hot
selected-output branches.

At the marked coordinate \(r\), do not tensorize a local projector
estimate.  Keep the complete signed \(K_{4,r}\).  Low selected output
uses the finite tree numerator and the scalar low-output gap; high
selected output uses
\eqref{eq:V38-hot-selected-output}.  This is precisely the local
good/bad split, now justified for \(G\) by
\Cref{thm:V38-G-local-output-moments}.  If one input at \(r\) is a
fixed fraction of its global row mass, the elementary star
comparison
\[
 \frac{a_{c,r}^3\prod_{i\ne c}a_{i,r}}{\Xi_c^2}
 \ge
 c^2\prod_{i=1}^4a_{i,r}
\]
closes the local-dominant branch.  Otherwise the unmarked bank
alternatives carry the missing row masses.

Allocate the sole payer before these estimates.  A selected-junction
payer uses the preceding low/high split; a complete shared or private
bank uses its paid moment compiler; a finite separator or suffix uses
its established single-debit norm.  All remaining factors are
contractions or the exact lower-order connected prefix factors in the
first-connectivity word.  The same finite tree routing as in the
retained proof then gives a literal marked tree.  Every
representation-dependent line has been replaced in
\Cref{lem:V38-replacement-table}, so the constant is uniform in
representations, support, and multiplicity.
\end{proof}

\begin{corollary}[The selected \(C_4\) residual is closed for \(G\)]
\label{cor:V38-C4-closed}
The \(C_4\) alternative in
\Cref{thm:V37-double-demand-reduction} is removed for every fixed
compact connected semisimple \(G\).
\end{corollary}

\begin{proof}
At the first-connectivity coordinate \(r\), the \(C_4\) term of
\eqref{eq:V37-five-role-identity} is an admissible packet in the
sense of \Cref{def:V38-admissible-selected}.  Apply
\Cref{thm:V38-GSEL4}.
\end{proof}

\begin{firewall}[No multiple-selection conclusion]
\label{fire:V38-no-global-selection}
The theorem begins after \(r\) has been selected by the disjoint
first-connectivity identity.  It never tensors a restricted
\(M-K_4\) prefix and never sums possible first four-block locations.
The frozen withdrawal of global embedded-\([4]\) aggregation is
unchanged.
\end{firewall}

\section{Capacity does not square into energy}
\label{sec:V38-energy-no-go}

\begin{proposition}[Sharp abstract endpoint obstruction]
\label{prop:V38-capacity-not-energy}
There is no universal implication
\begin{equation}
 \|S\|\le1,\qquad
 \kappa_\otimes(S)\le\mathfrak h
 \quad\Longrightarrow\quad
 {\cal E}_\otimes(S)\le C\mathfrak h^2
\label{eq:V38-false-energy-implication}
\end{equation}
for positive bipartite operators.
\end{proposition}

\begin{proof}
On \(\mathbb C^d\otimes\mathbb C^d\), let
\[
 \Omega_d=d^{-1/2}\sum_{k=1}^de_k\otimes e_k,
 \qquad
 S_d=|\Omega_d\rangle\langle\Omega_d|.
\]
Then \(\|S_d\|=1\).  Its maximum expectation on a product vector is
\(d^{-1}\), so
\[
 \kappa_\otimes(S_d)=d^{-1},
 \qquad
 {\cal E}_\otimes(S_d)
 =
 \|S_d\|\kappa_\otimes(S_d)=d^{-1}.
\]
Taking \(\mathfrak h=d^{-1}\) in
\eqref{eq:V38-false-energy-implication} would require
\(d^{-1}\le Cd^{-2}\), which fails as \(d\to\infty\).
\end{proof}

\begin{corollary}[The archived ternary response is not the missing
energy]
\label{cor:V38-ternary-response-not-energy}
A coefficient-normalized product-state response bound for the
complete ternary carrier, even together with a unit ordinary-norm
cap, does not imply
\eqref{eq:V37-ternary-star-energy}.  A proof of that energy must show
that the ordinary norm of the same signed ternary occurrence carries
an additional copy of its normalized star stock, or must replace the
energy route.
\end{corollary}

\begin{proof}
One-use energy is the product of ordinary norm and absolute
product-state capacity.  \Cref{prop:V38-capacity-not-energy} shows
that a small capacity can be caused entirely by entanglement while
the norm stays one.  In that case the energy retains only one power
of the small parameter.
\end{proof}

\begin{firewall}[The abstract projector is a diagnostic]
\label{fire:V38-projector-scope}
The family \(S_d\) is not asserted to equal the physical ternary
Wilson multiplier.  It proves that the desired quadratic energy
cannot be cited from the known response capacity without using
additional signed or heat-kernel structure of that multiplier.
\end{firewall}

\section{The new exact residual}
\label{sec:V38-residual}

\begin{theorem}[Double-demand residual after selected-junction
closure]
\label{thm:V38-residual}
For every fixed compact connected semisimple \(G\), the selected
fourth-cumulant component of each \(2+1+1\) star is closed.  The
remaining lower-order double-demand component is exactly the two
row-permutations of the binary--ternary word
\eqref{eq:V37-binary-ternary-word}.  The five-factor route closes it
if \eqref{eq:V37-ternary-star-energy} is proved, but that energy does
not follow from the existing ternary response-capacity theorem.
\end{theorem}

\begin{proof}
The selected component is
\Cref{cor:V38-C4-closed}.  The exhaustive five-role classification
and localization of the lower-order stars are
\Cref{lem:V37-lower-role-table,
thm:V37-double-demand-reduction}.  Sufficiency of the ternary energy
is \Cref{lem:V37-BTE-energy-suffices}, while its non-implication from
capacity is \Cref{cor:V38-ternary-response-not-energy}.
\end{proof}

\begin{assessment}[Upstream decision after V38]
\label{ass:V38-upstream}
The \(R_2\) search should now remain inside one complete ternary
occurrence.  The first calculation is the square
\(E_{134}^*E_{134}\) before the martingale difference is positively
enveloped.  If its cross terms reconstruct a complete
four-replica covariance, that joint carrier may provide the second
stock.  If triangle inequality removes those cross terms, the
maximally entangled diagnostic shows why the energy route stalls and
the proof must use a two-occurrence sandwich with a different
complete endpoint.
\end{assessment}

\begin{theorem}[Third-series V38 result ledger]
\label{thm:V38-results}
Relative to V1--V37, V38 proves:
\begin{enumerate}[label=\textup{(V38.\arabic*)},leftmargin=6.3em]
\item compact-group fixed-arity Casimir fusion;
\item one-link heat debits through cubic order;
\item compact-group complete local fourth-cumulant output moments
      through cubic order, without multiplicity loss;
\item an exhaustive replacement table for the retained
      one-selected-junction proof;
\item \(GSEL_4\) for every fixed compact connected semisimple \(G\);
\item closure of the selected \(C_4\) component of the \(2+1+1\)
      star;
\item a sharp abstract counterexample to squaring product-state
      capacity into one-use energy; and
\item reduction of the double-demand residual to the
      binary--ternary word and its missing signed endpoint energy.
\end{enumerate}
\end{theorem}

\begin{proof}
These are
\Cref{lem:V38-Casimir-fusion,
lem:V38-one-link-debits,
thm:V38-G-local-output-moments,
lem:V38-replacement-table,
thm:V38-GSEL4,
cor:V38-C4-closed,
prop:V38-capacity-not-energy,
thm:V38-residual}.
\end{proof}

\begin{programme}[V39: the squared signed ternary occurrence]
\label{prog:V38-V39}
\begin{enumerate}[leftmargin=3em]
\item Write the exact martingale-difference representation of
      \(E_{134}=\kappa_3(X_1,X_3,X_4)\) before its influence-tree
      envelope.
\item Expand \(E_{134}^*E_{134}\) in two replicas, preserving diagonal
      and cross terms as one positive complete carrier.
\item Determine whether the product-state capacity of that carrier
      yields the square of the normalized star stock in
      \eqref{eq:V37-ternary-star-energy}.
\item Test every proposed inequality on the maximally entangled
      family of \Cref{prop:V38-capacity-not-energy} and on the
      normalized \(SU(2)\) character calculation
      \eqref{eq:V37-character-cumulant}.
\item If the cross terms have no fixed sign, keep the four-replica
      carrier assembled and identify its first-connectivity partition
      rather than applying triangle inequality.
\item Leave the \(R_1\) path ledger untouched until this \(R_2\)
      decision is made.
\item V40 remains the next compulsory companion audit.
\end{enumerate}
\end{programme}

\begin{firewall}[Claim boundary after third-series V38]
\label{fire:V38-claim-boundary}
V38 closes the selected fourth-cumulant part of the \(2+1+1\) star
for fixed compact connected semisimple \(G\).  It does not prove the
binary--ternary endpoint energy, the remaining double-demand word,
the \(R_1\) paths, or the full \(2+1+1\) source.  The \(3+1\), global
discrete, globally aggregated \([4]\), and \([3,3]\) families, higher
MTA, WUFC, CPM, cutoff removal, reflection positivity,
Osterwalder--Schrader reconstruction, construction of the continuum
Yang--Mills measure, and a uniform positive continuum spectral gap
remain open.  The full Yang--Mills mass gap has not been proved.
\end{firewall}

\paragraph{Parent-version record.}
V38 was copied byte-for-byte from
\texttt{Renewal\_Yang\_Mills\_Mass\_Gap\_Research\_Notes\_V37.tex}
before this chapter was appended.  The V37 parent had \(19\,744\)
source lines, \(714\,855\) bytes, and SHA--256
\[
 \texttt{b0bbe787b0749a0a245b91c3cd16958c745970a36a43e71bdbd7a0984807c336}.
\]
The next compulsory companion audit is V40.

% =============================================================================
% V39 APPENDIX CHAPTER
% =============================================================================

\clearpage
\chapter{V39: The exact two-replica square of the ternary endpoint}
\label{ch:V39-ternary-square}

\section{Why the square must be formed before the tree envelope}
\label{sec:V39-purpose}

The missing V37 endpoint estimate concerns one complete signed
ternary occurrence.  V38 showed that its known product-state response
capacity cannot be squared abstractly: a small capacity may come only
from entanglement while the ordinary norm remains one.  The only
remaining energy route is therefore to compute the norm and capacity
of the same signed occurrence jointly.

This chapter constructs the exact positive object on which that
question lives.  The third cumulant is first centered, then squared
with two complete replicas.  The resulting double integral is kept as
one positive operator.  A general energy-lift lemma proves that two
quadratic bounds on this carrier are sufficient for the required
one-use energy.  It also proves that applying the V67 positive
influence envelope before the square loses the cross terms and cannot
establish those bounds.

\begin{assessment}[Nonduplication boundary]
\label{ass:V39-nonduplication}
V94 defines one-use energy and V97 uses it in a five-factor
sandwich.  V29/V67 bounds a ternary cumulant after a martingale-tree
envelope.  Neither archive writes the positive two-replica carrier of
the complete signed ternary occurrence.  That carrier, and the exact
reduction of the missing energy to it, are the new results below.
\end{assessment}

\section{Centered form of the local third cumulant}
\label{sec:V39-centered}

Let \((\Omega,\mu)\) be the one-coordinate Wilson probability space.
Let
\[
 A:\Omega\to{\cal B}(H_1),\qquad
 B:\Omega\to{\cal B}(H_3),\qquad
 C:\Omega\to{\cal B}(H_4)
\]
be bounded measurable row operators.  They are embedded in
\({\cal B}(H_1\otimes H_3\otimes H_4)\) on their respective tensor
legs and therefore commute pointwise across distinct letters.  Put
\[
 A^\circ=A-\mathbb EA,\qquad
 B^\circ=B-\mathbb EB,\qquad
 C^\circ=C-\mathbb EC.
\]

\begin{lemma}[Exact centered ternary identity]
\label{lem:V39-centered-identity}
The complete local third cumulant is
\begin{equation}
 S:=\kappa_3(A,B,C)
 =
 \int_\Omega
 A^\circ(g)\otimes B^\circ(g)\otimes C^\circ(g)\,d\mu(g).
\label{eq:V39-centered-identity}
\end{equation}
This identity is valid before physical output compression.
\end{lemma}

\begin{proof}
Expand the right side.  The uncentered term is
\(\mathbb E[ABC]\).  The three terms with exactly one expectation are
\[
 -(\mathbb EA)\mathbb E[BC]
 -(\mathbb EB)\mathbb E[AC]
 -(\mathbb EC)\mathbb E[AB].
\]
The three terms with two expectations contribute
\(3(\mathbb EA)(\mathbb EB)(\mathbb EC)\), while the term with three
centered subtractions contributes
\(-(\mathbb EA)(\mathbb EB)(\mathbb EC)\).  Their sum has coefficient
two, giving the moment--cumulant formula for \(\kappa_3\).
All products have their fixed tensor order, and operators belonging to
different row legs commute, so the expansion is legitimate.
\end{proof}

\begin{firewall}[Composite blocks]
\label{fire:V39-composite-block}
For the V67 joining carrier, take
\((A,B,C)=(X_1,X_3,X_4)\) only after the exact
Leonov--Shiryaev split \eqref{eq:V37-five-role-identity} has selected
the correction
\(\mathbb EX_2\,\kappa_3(X_1,X_3,X_4)\).
Equation \eqref{eq:V39-centered-identity} is not applied separately
to the five letters of the original composite cumulant.
\end{firewall}

\section{The exact two-replica carrier}
\label{sec:V39-two-replica}

Let \(g,h\) be independent coordinates with law \(\mu\).  Define
\begin{align}
 {\cal K}_{g,h}
 &:=
 A^\circ(g)^*A^\circ(h)
 \otimes
 B^\circ(g)^*B^\circ(h)
 \otimes
 C^\circ(g)^*C^\circ(h),
\label{eq:V39-two-replica-kernel}\\
 {\cal Q}_R(S)
 &:=
 \iint_{\Omega^2}{\cal K}_{g,h}\,d\mu(g)d\mu(h).
\label{eq:V39-right-square}
\end{align}
Define the left square analogously by reversing the factors,
\begin{equation}
 {\cal Q}_L(S)
 :=
 \iint_{\Omega^2}
 A^\circ(g)A^\circ(h)^*
 \otimes
 B^\circ(g)B^\circ(h)^*
 \otimes
 C^\circ(g)C^\circ(h)^*
 \,d\mu(g)d\mu(h).
\label{eq:V39-left-square}
\end{equation}

\begin{theorem}[Exact square identity]
\label{thm:V39-exact-square}
One has
\begin{equation}
 \boxed{
 {\cal Q}_R(S)=S^*S,\qquad
 {\cal Q}_L(S)=SS^*.}
\label{eq:V39-exact-square}
\end{equation}
In particular both double integrals are positive even though their
pointwise kernels need not be positive.
\end{theorem}

\begin{proof}
Taking the adjoint of
\eqref{eq:V39-centered-identity} and multiplying gives
\[
 S^*S
 =
 \left(\int
 A^\circ(g)^*\otimes B^\circ(g)^*\otimes C^\circ(g)^*
 \,d\mu(g)\right)
 \left(\int
 A^\circ(h)\otimes B^\circ(h)\otimes C^\circ(h)
 \,d\mu(h)\right).
\]
Fubini's theorem applies because all factors are bounded.  Multiplying
the tensor legs gives \eqref{eq:V39-right-square}.  The proof of the
left identity is identical.  Positivity follows from the two operator
squares, not from the individual kernels.
\end{proof}

\begin{firewall}[Do not take absolute values inside the replica
integral]
\label{fire:V39-no-kernel-envelope}
Replacing \({\cal K}_{g,h}\) by a pointwise positive envelope before
the double integration discards the cross-replica cancellation which
makes \({\cal Q}_R(S)=S^*S\).  It can also create two independently
optimized joining marks.  The square identity authorizes only an
estimate of the assembled positive double integral.
\end{firewall}

\section{Hermitianization and the exact energy lift}
\label{sec:V39-energy-lift}

The ternary occurrence need not be self-adjoint for a general compact
group.  On \(\mathbb C^2\otimes H\), define its Hermitianization
\begin{equation}
 {\mathscr H}(S)
 :=
 \begin{pmatrix}
 0&S\\ S^*&0
 \end{pmatrix}.
\label{eq:V39-Hermitianization}
\end{equation}
The auxiliary two-dimensional factor is finite and carries no row
mass.

\begin{lemma}[Hermitianized absolute value]
\label{lem:V39-Hermitianized-absolute}
One has
\begin{align}
 \|{\mathscr H}(S)\|&=\|S\|,
\label{eq:V39-Hermitianized-norm}\\
 |{\mathscr H}(S)|
 &=
 \begin{pmatrix}
 (SS^*)^{1/2}&0\\
 0&(S^*S)^{1/2}
 \end{pmatrix}.
\label{eq:V39-Hermitianized-modulus}
\end{align}
Consequently
\begin{equation}
 \kappa_\otimes(|{\mathscr H}(S)|)
 =
 \max\left\{
 \kappa_\otimes((SS^*)^{1/2}),
 \kappa_\otimes((S^*S)^{1/2})
 \right\},
\label{eq:V39-Hermitianized-capacity}
\end{equation}
where the supremum also ranges over the auxiliary density matrix.
\end{lemma}

\begin{proof}
Squaring \eqref{eq:V39-Hermitianization} gives the diagonal operator
\(\operatorname{diag}(SS^*,S^*S)\), proving the first two identities
by functional calculus.  The expectation of the positive diagonal
operator in an auxiliary density matrix is a convex combination of
the two diagonal expectations.  Its supremum is their maximum, and
either value is attained by an auxiliary basis state.
\end{proof}

\begin{lemma}[Product-state Cauchy--Schwarz for a square root]
\label{lem:V39-product-CS}
For every bounded \(S\),
\begin{align}
 \kappa_\otimes((S^*S)^{1/2})^2
 &\le\kappa_\otimes(S^*S),
\label{eq:V39-right-CS}\\
 \kappa_\otimes((SS^*)^{1/2})^2
 &\le\kappa_\otimes(SS^*).
\label{eq:V39-left-CS}
\end{align}
\end{lemma}

\begin{proof}
For a product density \(\rho\), Cauchy--Schwarz in the positive
functional \(X\mapsto\operatorname{Tr}(\rho X)\) gives
\[
 \operatorname{Tr}(\rho (S^*S)^{1/2})^2
 \le
 \operatorname{Tr}(\rho)\operatorname{Tr}(\rho S^*S)
 =
 \operatorname{Tr}(\rho S^*S).
\]
Take the supremum over product densities.  The left-square statement
is identical.
\end{proof}

\begin{theorem}[Two-replica energy lift]
\label{thm:V39-energy-lift}
For every bounded \(S\),
\begin{equation}
 \boxed{
 {\cal E}_\otimes({\mathscr H}(S))
 \le
 \sqrt{
 \max\{\|S^*S\|,\|SS^*\|\}
 \max\{
 \kappa_\otimes(S^*S),
 \kappa_\otimes(SS^*)\}}.}
\label{eq:V39-energy-lift}
\end{equation}
Hence, if for a scalar stock \(\mathfrak h\)
\begin{align}
 \|{\cal Q}_R(S)\|,\|{\cal Q}_L(S)\|
 &\le C\mathfrak h^2,
\label{eq:V39-square-norm-target}\\
 \kappa_\otimes({\cal Q}_R(S)),
 \kappa_\otimes({\cal Q}_L(S))
 &\le C\mathfrak h^2,
\label{eq:V39-square-capacity-target}
\end{align}
then
\begin{equation}
 {\cal E}_\otimes({\mathscr H}(S))
 \le C\mathfrak h^2.
\label{eq:V39-energy-conclusion}
\end{equation}
\end{theorem}

\begin{proof}
By \Cref{lem:V39-Hermitianized-absolute},
\[
 {\cal E}_\otimes({\mathscr H}(S))
 =
 \|S\|
 \max\left\{
 \kappa_\otimes((SS^*)^{1/2}),
 \kappa_\otimes((S^*S)^{1/2})
 \right\}.
\]
Use
\(\|S\|^2=\|S^*S\|=\|SS^*\|\) and
\Cref{lem:V39-product-CS}.  This proves
\eqref{eq:V39-energy-lift}.  Insert
\eqref{eq:V39-square-norm-target}--%
\eqref{eq:V39-square-capacity-target} to get
\eqref{eq:V39-energy-conclusion}.
\end{proof}

\section{The normalized ternary-square target}
\label{sec:V39-TQ2}

For the star endpoint at coordinate \(r\), put
\begin{equation}
 \mathfrak h_{134,r}
 :=
 \frac{h_{13,r}h_{14,r}}
 {\Xi_1^2\Xi_3\Xi_4}.
\label{eq:V39-ternary-stock}
\end{equation}
Let \(\widetilde S_{134,r}\) denote the complete coefficient-normalized
physical multiplier of
\(\kappa_{3,r}(X_{1,r},X_{3,r},X_{4,r})\), with full output
multiplicity and without the final payer.  All scalar coefficient
normalizations are included before its absolute value is taken.

\begin{definition}[Ternary quadratic two-replica target]
\label{def:V39-TQ2}
\(TQ_2\) is the pair of estimates
\begin{align}
 \|
 \widetilde S_{134,r}^*\widetilde S_{134,r}
 \|,
 \|
 \widetilde S_{134,r}\widetilde S_{134,r}^*
 \|
 &\le C\mathfrak h_{134,r}^2,
\label{eq:V39-TQ2-norm}\\
 \kappa_\otimes(
 \widetilde S_{134,r}^*\widetilde S_{134,r}),
 \kappa_\otimes(
 \widetilde S_{134,r}\widetilde S_{134,r}^*)
 &\le C\mathfrak h_{134,r}^2,
\label{eq:V39-TQ2-capacity}
\end{align}
together with the corresponding once-paid versions at every local
or adjacent payer position.
\end{definition}

\begin{corollary}[\(TQ_2\) closes the double-demand endpoint]
\label{cor:V39-TQ2-suffices}
If \(TQ_2\) holds, then the missing ternary energy
\eqref{eq:V37-ternary-star-energy}, and hence \(BTE_2\), holds.
\end{corollary}

\begin{proof}
Apply \Cref{thm:V39-energy-lift} with
\(S=\widetilde S_{134,r}\) and
\(\mathfrak h=\mathfrak h_{134,r}\).  Hermitianization adds only a
fixed auxiliary factor and can be propagated through the two bounded
separators in the V97 word by the standard off-diagonal block
doubling.  Thus
\eqref{eq:V39-energy-conclusion} is exactly the endpoint hypothesis
\eqref{eq:V37-ternary-star-energy}.  Apply
\Cref{lem:V37-BTE-energy-suffices}.
\end{proof}

\begin{firewall}[The target is not two separately paid replicas]
\label{fire:V39-one-payer}
The paid form of \(TQ_2\) contains the original sole payer in one
complete doubled carrier.  It does not assign a resolvent to both
replicas.  Doing so would prove a different, two-payer inequality and
would lose the required power of \(s_\alpha\).
\end{firewall}

\section{Why the V67 envelope cannot prove \(TQ_2\)}
\label{sec:V39-envelope-no-go}

\begin{proposition}[Envelope-first squaring loses the required
carrier]
\label{prop:V39-envelope-first-loss}
The ordinary-norm inequality
\[
 \|S\|
 \le
 Cs_\alpha^2
 \bigl(h_{13,r}h_{14,r}
       +h_{13,r}h_{34,r}
       +h_{14,r}h_{34,r}\bigr)
\]
does not imply either line of \(TQ_2\).  Squaring its right side is
only a scalar bound on \(\|S\|^2\); it supplies neither the two
additional row-\(1\) coefficient denominators nor the product-state
capacity of the assembled operators \(S^*S,SS^*\).
\end{proposition}

\begin{proof}
The first assertion follows directly from the absence of every
\(\Xi_i\) on the right side of the displayed local estimate.  More
structurally, the inequality was obtained after the signed martingale
differences were replaced by positive influence factors.  It contains
no operator whose product-state capacity equals that of
\eqref{eq:V39-right-square} or \eqref{eq:V39-left-square}.

Finally, \Cref{prop:V38-capacity-not-energy} gives positive operators
with the same logical data---unit norm and small product capacity---for
which the required squared capacity scale fails.  Therefore no
abstract combination of the envelope norm and the known one-use
capacity can establish \(TQ_2\).
\end{proof}

\begin{remark}[What remains available]
\label{rem:V39-cross-terms}
The exact double integrals
\eqref{eq:V39-right-square}--\eqref{eq:V39-left-square} retain
the cross terms \(g\ne h\).  They may be reorganized by
first connectivity on the doubled replica alphabet.  This is the
only remaining information in the energy route that is absent from
the V67 envelope.
\end{remark}

\section{Regression on normalized characters}
\label{sec:V39-character-regression}

\begin{proposition}[The character family is compatible with the
square construction]
\label{prop:V39-character-square}
In the normalized \(SU(2)\) character family of
\Cref{prop:V37-character-noncancellation}, the complete ternary
endpoint is the scalar
\[
 S_j=\frac{d_j-1}{d_j^4},
\]
and hence
\begin{equation}
 S_j^*S_j=S_jS_j^*
 =
 \frac{(d_j-1)^2}{d_j^8}
 =
 O(d_j^{-6}).
\label{eq:V39-character-square}
\end{equation}
Thus this protected invariant family does not by itself contradict a
quadratic two-replica estimate.
\end{proposition}

\begin{proof}
The value of \(S_j\) is
\eqref{eq:V37-character-cumulant}.  It is a nonnegative scalar.
Squaring gives the displayed identity, and
\((d_j-1)^2d_j^{-8}\le d_j^{-6}\).
\end{proof}

\begin{assessment}[Interpretation of the two regressions]
\label{ass:V39-regressions}
The maximally entangled projector of V38 forbids deriving \(TQ_2\)
from capacity alone.  The character calculation shows that the most
obvious invariant scalar channel has additional cumulant decay and
does not refute \(TQ_2\).  The decisive case is therefore a
coefficient block with large invariant multiplicity and a signed
ternary multiplier whose norm remains large.  It must be tested inside
the exact two-replica carrier, not after character tracing.
\end{assessment}

\section{V39 status}
\label{sec:V39-status}

\begin{theorem}[Exact residual after the energy lift]
\label{thm:V39-residual}
The double-demand \(2+1+1\) residual is reduced to the two-replica
estimates \(TQ_2\) for the complete ternary endpoint and their
once-paid variants.  No additional local representation case remains:
the selected fourth-cumulant component is closed by V38, and the
binary endpoint is already covered by the quadratic \(A,F,I\) table.
\end{theorem}

\begin{proof}
The V38 residual theorem leaves only \(BTE_2\).
\Cref{cor:V39-TQ2-suffices} reduces \(BTE_2\) to \(TQ_2\).
The other two assertions are
\Cref{cor:V38-C4-closed,
prop:V37-missing-ternary-energy}.
\end{proof}

\begin{theorem}[Third-series V39 result ledger]
\label{thm:V39-results}
Relative to V1--V38, V39 proves:
\begin{enumerate}[label=\textup{(V39.\arabic*)},leftmargin=6.3em]
\item the exact centered representation of the complete ternary
      occurrence;
\item its exact positive left and right two-replica squares;
\item a non-self-adjoint Hermitianization compatible with
      product-state capacity;
\item product-state Cauchy--Schwarz for the two square roots;
\item the two-replica energy-lift theorem;
\item the precise \(TQ_2\) norm, capacity, and one-payer targets;
\item proof that \(TQ_2\) implies the remaining \(BTE_2\) estimate;
\item proof that envelope-first squaring cannot establish \(TQ_2\);
      and
\item compatibility of the normalized-character regression with the
      exact square.
\end{enumerate}
\end{theorem}

\begin{proof}
These are
\Cref{lem:V39-centered-identity,
thm:V39-exact-square,
lem:V39-Hermitianized-absolute,
lem:V39-product-CS,
thm:V39-energy-lift,
def:V39-TQ2,
cor:V39-TQ2-suffices,
prop:V39-envelope-first-loss,
prop:V39-character-square,
thm:V39-residual}.
\end{proof}

\begin{programme}[V40: audit and doubled first connectivity]
\label{prog:V39-V40}
\begin{enumerate}[leftmargin=3em]
\item Perform the compulsory read-only audit of the newest active SM
      and NS notes before choosing a doubled-replica decomposition.
\item Expand the centered factors in
      \eqref{eq:V39-two-replica-kernel} only after the \(g,h\)
      integrations have been assembled.
\item Classify the equality stratum \(g=h\) and the ordered strata
      \(g<h,h<g\) by first connectivity on the six row-replica slots.
\item Determine whether the off-diagonal strata tensor into one
      ternary occurrence followed by a complete moment, or whether
      they form a genuinely new six-slot cumulant.
\item Prove the \(TQ_2\) norm line first.  Attempt the product-state
      capacity line only on the same assembled carrier.
\item Preserve the one-payer label through both replicas.
\item If the doubled carrier recreates an uncontrolled global
      \([3,3]\) transfer, go upstream to its exact restricted factor
      rather than multiplying local bounds.
\end{enumerate}
\end{programme}

\begin{firewall}[Claim boundary after third-series V39]
\label{fire:V39-claim-boundary}
V39 constructs the exact positive carrier sufficient for the missing
ternary energy but does not prove its \(TQ_2\) bounds.  Therefore the
binary--ternary double-demand word, the \(R_1\) paths, and the full
\(2+1+1\) source remain open.  The \(3+1\), global discrete,
globally aggregated \([4]\), and \([3,3]\) families, higher MTA,
WUFC, CPM, cutoff removal, reflection positivity,
Osterwalder--Schrader reconstruction, construction of the continuum
Yang--Mills measure, and a uniform positive continuum spectral gap
also remain open.  The full Yang--Mills mass gap has not been proved.
\end{firewall}

\paragraph{Parent-version record.}
V39 was copied byte-for-byte from
\texttt{Renewal\_Yang\_Mills\_Mass\_Gap\_Research\_Notes\_V38.tex}
before this chapter was appended.  The V38 parent had \(20\,239\)
source lines, \(733\,161\) bytes, and SHA--256
\[
 \texttt{caf65ef723d8571e4f4f32922c4c5afebae19cab13a458deacf7cf638cdcce8a}.
\]
The next iteration is the compulsory V40 companion audit.

% =============================================================================
% V40 APPENDIX CHAPTER
% =============================================================================

\clearpage
\chapter{V40: Companion audit, replica correction, and the exact
four-replica polarization}
\label{ch:V40-audit-polarization}

\section{Compulsory companion audit}
\label{sec:V40-audit}

The newest active companion notes were inspected read-only before the
V39 doubled carrier was decomposed:
\begin{center}
\begin{tabular}{@{}p{0.20\linewidth}p{0.69\linewidth}@{}}
\toprule
programme & audited file and fingerprint\\
\midrule
SM--Einstein&
\texttt{Renewal\_SM\_Einstein\_Continuum\_Research\_Notes\_V17.tex},
\(6\,022\) source lines, \(240\,589\) bytes,
SHA--256
\texttt{fc378f347a3b052cdb99b15526fa571f0b03476ac5e6f44052d5d8586d3f02b8};
\\[2mm]
Navier--Stokes&
\texttt{Renewal\_NS\_Stability\_Research\_Notes\_V31.tex},
\(23\,052\) source lines, \(887\,537\) bytes,
SHA--256
\texttt{f1121b62238ad5491bb7cf03635ea9934942edf573ed8faaa9ee79e1d38a6696}.
\\
\bottomrule
\end{tabular}
\end{center}

SM V15 proves an exact two-mark product rule.  Its two disjoint cases
are one genuine second derivative of a single activity and two
genuine first derivatives on two distinct activities of the same
connected cluster.  SM V16 retains both terms in
\[
 D^2\operatorname{Tr}\log(1+K)
 =
 \operatorname{Tr}\bigl(
 QD^2K-QDK\,QDK\bigr),
\]
and SM V17 keeps the quartic seagull, the twice-marked cubic
occurrence, and the two separately marked cubic occurrences in one
cumulant packet.  None of these results infers a second response by
squaring a first-response upper bound.

NS V31 is byte-for-byte unchanged from the V35 audit.  Its flat
annular mark reconstructs an unweighted integral only after paying an
actual first radial moment; heat smoothing alone does not create it.
The corresponding YM rule remains that two endpoint debits require
two actual replica occurrences in one complete carrier, and an
absolute envelope cannot manufacture either debit.

\begin{firewall}[Companion type boundary]
\label{fire:V40-audit-types}
No SM cluster smallness, trace-log resolvent estimate, seagull
identity, NS heat-formation bound, or fluid coarea estimate is used as
a Yang--Mills inequality.  The audit transfers only proof order:
differentiate or replicate the complete connected object first,
separate exhaustive actual mark locations, and take absolute bounds
afterward.
\end{firewall}

\section{Correction of the V39 replica programme}
\label{sec:V40-correction}

In V39, \(g,h\in\Omega\) denote two independent samples of the
one-coordinate Wilson probability space.  They are not physical
coordinate positions in the ordered source word.

\begin{theorem}[Withdrawal of ordered replica strata]
\label{thm:V40-withdraw-order}
The proposed V40 split into \(g=h\), \(g<h\), and \(h<g\) strata in
\Cref{prog:V39-V40} is invalid and is withdrawn.  There is no
canonical order on two probability-space samples, and the diagonal
need not be a distinguished positive component of the product
measure.
\end{theorem}

\begin{proof}
The variables \(g,h\) enter through Fubini's identity
\[
 S^*S
 =
 \iint_{\Omega^2}{\cal K}_{g,h}\,d\mu(g)d\mu(h).
\]
Their labels record independent replicas only.  An order relation on
the physical coordinate set does not induce an order on \(\Omega\).
Moreover, for an atomless measure the diagonal has product measure
zero, whereas for a measure with atoms its mass depends on the atomic
law; neither case gives a universal first-connectivity component.
\end{proof}

\begin{corollary}[No global \([3,3]\) transfer is created by the
replica labels]
\label{cor:V40-no-coordinate-transfer}
The two-replica square is a local one-coordinate operator identity.
It cannot be classified by the global coordinate automaton used for
the \([3,3]\) accepted-state transfer unless an additional physical
coordinate telescope is first introduced.
\end{corollary}

\begin{proof}
Replica labels belong to copies of the local probability space,
whereas the accepted-state automaton is indexed by ordered physical
coordinates.  Identifying them would conflate two different index
sets.
\end{proof}

\begin{assessment}[Correction record]
\label{ass:V40-correction-record}
The exact identities and energy lift proved in V39 remain valid.
Only the proposed decomposition in its programme is withdrawn.  The
replacement is an independent-replica polarization of the local
third cumulant.
\end{assessment}

\section{Four-replica polarization of a third cumulant}
\label{sec:V40-four-replica}

Let \(\xi_0,\xi_1,\xi_2,\xi_3\) be independent samples with law
\(\mu\).  For the three row functions \(A,B,C\), define
\begin{equation}
 \Delta_1A:=A(\xi_0)-A(\xi_1),\quad
 \Delta_2B:=B(\xi_0)-B(\xi_2),\quad
 \Delta_3C:=C(\xi_0)-C(\xi_3).
\label{eq:V40-three-differences}
\end{equation}
All three differences share the base replica \(\xi_0\), while their
replacement replicas are distinct.

\begin{theorem}[Exact four-replica polarization]
\label{thm:V40-four-replica}
\begin{equation}
 \boxed{
 \kappa_3(A,B,C)
 =
 \mathbb E_{\xi_0,\xi_1,\xi_2,\xi_3}
 \left[
 \Delta_1A\otimes\Delta_2B\otimes\Delta_3C
 \right].}
\label{eq:V40-four-replica}
\end{equation}
The formula is valid for bounded operators on separate row tensor
legs and preserves their fixed order.
\end{theorem}

\begin{proof}
Expand the product of the three differences.  The term using
\(\xi_0\) in all slots is \(\mathbb E[ABC]\).  The three terms with
one replacement replica are
\[
 -(\mathbb EA)\mathbb E[BC],
 \quad
 -(\mathbb EB)\mathbb E[AC],
 \quad
 -(\mathbb EC)\mathbb E[AB].
\]
Each of the three terms with two replacement replicas factors as
\((\mathbb EA)(\mathbb EB)(\mathbb EC)\), and the term with all three
replacements has the negative of that value.  The net coefficient of
the product of means is \(3-1=2\).  These are exactly the five terms
of the third moment--cumulant identity.  Independence justifies every
factorization.
\end{proof}

\begin{remark}[Why one common base is necessary]
\label{rem:V40-common-base}
Replacing \(A,B,C\) by differences between the same two complete
replicas would give zero: the all-base and all-replacement third
moments cancel.  In \eqref{eq:V40-four-replica}, the common base
retains the connected third moment while three independent
replacement samples perform the centering.
\end{remark}

\begin{firewall}[The differences are actual marks]
\label{fire:V40-actual-differences}
Each factor in \eqref{eq:V40-three-differences} is one exact
replacement difference.  It is not a second copy of a bound for a
centered variable.  The three replacements are formed in the signed
operator before any influence norm is used.
\end{firewall}

\section{The eight-replica positive square}
\label{sec:V40-eight-replica}

Let
\(\eta_0,\eta_1,\eta_2,\eta_3\) be an independent second copy of the
four replicas and put
\begin{align}
 D(\boldsymbol\xi)
 &:=
 \Delta_1A(\boldsymbol\xi)
 \otimes\Delta_2B(\boldsymbol\xi)
 \otimes\Delta_3C(\boldsymbol\xi),
\label{eq:V40-D-xi}\\
 D(\boldsymbol\eta)
 &:=
 \Delta_1A(\boldsymbol\eta)
 \otimes\Delta_2B(\boldsymbol\eta)
 \otimes\Delta_3C(\boldsymbol\eta).
\label{eq:V40-D-eta}
\end{align}

\begin{theorem}[Exact polarized square]
\label{thm:V40-polarized-square}
For \(S=\kappa_3(A,B,C)\),
\begin{align}
 S^*S
 &=
 \mathbb E_{\boldsymbol\xi,\boldsymbol\eta}
 \bigl[
 D(\boldsymbol\xi)^*D(\boldsymbol\eta)
 \bigr],
\label{eq:V40-polarized-right}\\
 SS^*
 &=
 \mathbb E_{\boldsymbol\xi,\boldsymbol\eta}
 \bigl[
 D(\boldsymbol\xi)D(\boldsymbol\eta)^*
 \bigr].
\label{eq:V40-polarized-left}
\end{align}
Each square contains two actual \(A\)-differences, one in each
four-replica packet.
\end{theorem}

\begin{proof}
By \Cref{thm:V40-four-replica},
\[
 S=\mathbb E_{\boldsymbol\xi}D(\boldsymbol\xi).
\]
Multiply this identity by its adjoint and use independence and
Fubini.  The two \(A\)-differences are visibly
\(\Delta_1A(\boldsymbol\xi)\) and
\(\Delta_1A(\boldsymbol\eta)\).
\end{proof}

\begin{corollary}[Exact two-mark interpretation]
\label{cor:V40-two-mark}
The center-row square in
\eqref{eq:V40-polarized-right}--%
\eqref{eq:V40-polarized-left} is a genuine two-mark object in the
sense of the V40 companion audit: the two marks act on two actual
occurrences inside one assembled positive square.
\end{corollary}

\begin{proof}
There is one center-row replacement in each of the two complete
factors whose product is \(S^*S\) or \(SS^*\).  The two placements
are exhaustive and are created by the exact square, not by duplicating
an upper bound.
\end{proof}

\begin{firewall}[Positivity remains global]
\label{fire:V40-global-positivity}
The expectations in
\eqref{eq:V40-polarized-right}--%
\eqref{eq:V40-polarized-left} are positive because they equal
operator squares.  A fixed kernel
\(D(\boldsymbol\xi)^*D(\boldsymbol\eta)\) need not be positive.
Triangle inequality may therefore be applied only after deciding
whether its cross-replica terms are needed for \(TQ_2\).
\end{firewall}

\section{What the square can and cannot improve}
\label{sec:V40-square-limit}

\begin{lemma}[The norm line of \(TQ_2\) is equivalent to a signed
ternary norm]
\label{lem:V40-norm-equivalence}
For the normalized ternary occurrence
\(\widetilde S=\widetilde S_{134,r}\),
\begin{equation}
 \|\widetilde S^*\widetilde S\|
 =
 \|\widetilde S\widetilde S^*\|
 =
 \|\widetilde S\|^2.
\label{eq:V40-norm-equivalence}
\end{equation}
Consequently \eqref{eq:V39-TQ2-norm} holds if and only if
\begin{equation}
 \boxed{
 \|\widetilde S_{134,r}\|
 \le C\mathfrak h_{134,r}.}
\label{eq:V40-normalized-ternary-norm}
\end{equation}
\end{lemma}

\begin{proof}
The \(C^*\)-identity gives
\(\|S^*S\|=\|S\|^2=\|SS^*\|\).  Take square roots in
\eqref{eq:V39-TQ2-norm}, absorbing the square root of its constant.
The converse follows by squaring
\eqref{eq:V40-normalized-ternary-norm}.
\end{proof}

\begin{assessment}[The positive square is not a free norm gain]
\label{ass:V40-no-free-norm}
The eight-replica formula is essential for the product-state capacity
line of \(TQ_2\), but it cannot improve the ordinary norm line:
that line is exactly the missing normalized signed ternary norm
\eqref{eq:V40-normalized-ternary-norm}.  Thus the next proof must
either establish that norm before the positive square is used, or
replace the V97 energy route by a sandwich which does not require it.
\end{assessment}

\begin{proposition}[The universal difference bound is insufficient]
\label{prop:V40-universal-difference}
If \(\|A\|_\infty,\|B\|_\infty,\|C\|_\infty\le1\), then
\begin{equation}
 \|S\|\le8,\qquad
 \|S^*S\|,\|SS^*\|\le64.
\label{eq:V40-universal-difference}
\end{equation}
These support-uniform bounds do not imply
\eqref{eq:V40-normalized-ternary-norm} when
\(\mathfrak h_{134,r}\to0\).
\end{proposition}

\begin{proof}
Every difference in
\eqref{eq:V40-three-differences} has norm at most two.  Hence
\(\|D(\boldsymbol\xi)\|\le8\), and
\Cref{thm:V40-four-replica} gives \(\|S\|\le8\).
The square bounds follow from the \(C^*\)-identity.  A fixed constant
cannot be bounded by \(C\mathfrak h\) uniformly as a positive scalar
\(\mathfrak h\) tends to zero.
\end{proof}

\section{Mass locality of the ternary endpoint}
\label{sec:V40-mass-locality}

The normalized stock
\[
 \mathfrak h_{134,r}
 =
 \frac{h_{13,r}h_{14,r}}
 {\Xi_1^2\Xi_3\Xi_4}
\]
uses global row masses.  A local occurrence at \(r\) cannot pay mass
which is concentrated in unrelated spectator coordinates unless
those spectators remain in the same complete endpoint.

\begin{lemma}[Local occurrence versus external mass]
\label{lem:V40-local-external}
Fix the local representations at \(r\) and enlarge only the
off-\(r\) row-\(1\) mass.  Then
\(\mathfrak h_{134,r}\to0\).  Therefore a uniform proof of
\eqref{eq:V40-normalized-ternary-norm} must use one of:
\begin{enumerate}[label=\textup{(\roman*)},leftmargin=4em]
\item a marked-dominance condition
      \(a_{1,r}\ge c\Xi_1\);
\item a complete off-\(r\) bank multiplier inside
      \(\widetilde S_{134,r}\); or
\item cancellation coupling the local occurrence to the external
      source word.
\end{enumerate}
\end{lemma}

\begin{proof}
The numerator \(h_{13,r}h_{14,r}\) is fixed, while the denominator
contains \(\Xi_1^2\).  Hence the stock tends to zero.  If the local
operator and all of its multipliers were unchanged and nonzero, its
norm could not obey a bound tending to zero.  Thus a valid uniform
proof must make the local representation comparable to the global
mass, retain a multiplier which changes with the added mass, or use a
signed identity involving the external word.  These are the three
listed possibilities.
\end{proof}

\begin{corollary}[Correction of the isolated endpoint target]
\label{cor:V40-dressed-target}
The ternary energy target in
\eqref{eq:V37-ternary-star-energy} is a valid isolated endpoint
target only on the marked-dominant subcell.  On a used-bank subcell,
the endpoint must be enlarged to include the actual bank which carries
the fixed fraction of \(\Xi_1\).  Its energy is a dressed
binary--bank--ternary problem, not a local ternary estimate.
\end{corollary}

\begin{proof}
Apply \Cref{lem:V40-local-external}.  The V35 row-demand split
distinguishes marked mass from a previously used or fresh external
bank precisely before an endpoint estimate is selected.
\end{proof}

\begin{firewall}[No spectator duplication]
\label{fire:V40-no-spectator-duplication}
The bank in \Cref{cor:V40-dressed-target} must be inserted as its
single complete occurrence in the source word.  It cannot be attached
once to the norm and independently again to the capacity.  Such a
step would repeat both the V21 entanglement-swapping error and the NS
V31 first-moment error.
\end{firewall}

\section{The corrected V40 residual}
\label{sec:V40-residual}

\begin{definition}[Marked and dressed ternary-square targets]
\label{def:V40-two-targets}
The remaining \(TQ_2\) problem is split as follows.
\begin{enumerate}[label=\textup{(\roman*)},leftmargin=4em]
\item \(TQ_2^{\rm mark}\) is
      \eqref{eq:V39-TQ2-norm}--%
      \eqref{eq:V39-TQ2-capacity} under
      \(a_{1,r}\ge c\Xi_1\), or the analogous dominance of another
      actual marked coordinate in the complete endpoint.
\item \(TQ_2^{\rm dress}\) is the same two-replica energy statement
      after the unique complete used bank carrying
      \(c\Xi_1\) has been inserted in its exact position between the
      binary prefix and ternary occurrence.
\end{enumerate}
Private, unbalanced, and two-genuine-\(A,F,I\) banks remain excluded
because they were closed in V35--V36.
\end{definition}

\begin{theorem}[Exact residual after the checkpoint correction]
\label{thm:V40-residual}
The double-demand \(2+1+1\) star sector is reduced to
\(TQ_2^{\rm mark}\) and \(TQ_2^{\rm dress}\) from
\Cref{def:V40-two-targets}.  The former requires the signed normalized
ternary norm \eqref{eq:V40-normalized-ternary-norm} and the capacity
of its exact eight-replica square.  The latter requires the same two
bounds on one bank-dressed complete occurrence.  No ordering of
probability replicas is part of either target.
\end{theorem}

\begin{proof}
V39 reduces the binary--ternary endpoint to \(TQ_2\).
\Cref{cor:V40-dressed-target} separates the marked-dominant and
external-mass alternatives.  In the first,
\Cref{lem:V40-norm-equivalence} gives the exact norm target and
\Cref{thm:V40-polarized-square} gives its positive capacity carrier.
The same statements apply after inserting the complete external bank
in the second alternative.  The excluded bank types are those in
\Cref{thm:V36-refined-residual}.  Finally
\Cref{thm:V40-withdraw-order} removes the invalid replica-order
classification.
\end{proof}

\begin{assessment}[Direction selected by the V40 audit]
\label{ass:V40-direction}
The SM companion result and the corrected replica algebra agree on
the next object.  For \(TQ_2^{\rm mark}\), keep the two actual
center differences in
\eqref{eq:V40-polarized-right} and seek a one-link
second-difference estimate on the complete ternary multiplier.  For
\(TQ_2^{\rm dress}\), form the actual three-factor word consisting of
binary prefix, mass-carrying bank, and ternary occurrence before any
norm.  A local ternary estimate with global masses inserted
afterward is not a valid substitute.
\end{assessment}

\begin{theorem}[Third-series V40 result ledger]
\label{thm:V40-results}
Relative to V1--V39, V40 proves:
\begin{enumerate}[label=\textup{(V40.\arabic*)},leftmargin=6.3em]
\item the compulsory SM V17/NS V31 audit with exact fingerprints;
\item withdrawal of the invalid ordered-replica programme;
\item proof that local probability replicas do not define a global
      \([3,3]\) coordinate transfer;
\item an exact four-replica polarization of the third cumulant;
\item exact eight-replica formulae for its positive left and right
      squares;
\item identification of two actual center marks inside those squares;
\item equivalence of the \(TQ_2\) norm target to a normalized signed
      ternary norm;
\item proof that the universal difference bound gives no normalized
      stock gain;
\item the local-versus-external mass trichotomy; and
\item the corrected split into marked and bank-dressed
      two-replica targets.
\end{enumerate}
\end{theorem}

\begin{proof}
These are
\Cref{fire:V40-audit-types,
thm:V40-withdraw-order,
cor:V40-no-coordinate-transfer,
thm:V40-four-replica,
thm:V40-polarized-square,
cor:V40-two-mark,
lem:V40-norm-equivalence,
prop:V40-universal-difference,
lem:V40-local-external,
cor:V40-dressed-target,
thm:V40-residual}.
\end{proof}

\begin{programme}[V41: genuine second difference at a marked ternary
endpoint]
\label{prog:V40-V41}
\begin{enumerate}[leftmargin=3em]
\item Work first on \(TQ_2^{\rm mark}\) with
      \(a_{1,r}\ge c\Xi_1\); keep the complete local output and payer.
\item Express the two center differences in
      \eqref{eq:V40-polarized-right} through the local Wilson
      semigroup carré du champ, without taking separate suprema over
      the two replicas.
\item Prove the signed ordinary-norm estimate
      \eqref{eq:V40-normalized-ternary-norm}, or construct a physical
      compact-group block for which it fails.
\item Only if the norm estimate holds, prove the product-state
      capacity line on the same eight-replica carrier.
\item Then insert one used bank in its exact source position and test
      \(TQ_2^{\rm dress}\); do not reuse the bank in both factors.
\item If the marked norm fails on a protected output, abandon the
      V97 energy route and return to the complete binary--ternary
      product before endpoint separation.
\item The next compulsory companion audit is V45.
\end{enumerate}
\end{programme}

\begin{firewall}[Claim boundary after third-series V40]
\label{fire:V40-claim-boundary}
V40 corrects the replica decomposition and constructs the exact
polarized square, but it does not prove
\(TQ_2^{\rm mark}\), \(TQ_2^{\rm dress}\), the remaining
double-demand word, the \(R_1\) paths, or the full \(2+1+1\) source.
The \(3+1\), global discrete, globally aggregated \([4]\), and
\([3,3]\) families, higher MTA, WUFC, CPM, cutoff removal, reflection
positivity, Osterwalder--Schrader reconstruction, construction of the
continuum Yang--Mills measure, and a uniform positive continuum
spectral gap remain open.  The full Yang--Mills mass gap has not been
proved.
\end{firewall}

\paragraph{Parent-version record.}
V40 was copied byte-for-byte from
\texttt{Renewal\_Yang\_Mills\_Mass\_Gap\_Research\_Notes\_V39.tex}
before this chapter was appended.  The V39 parent had \(20\,794\)
source lines, \(751\,273\) bytes, and SHA--256
\[
 \texttt{27b5e499229badebee1eb903a67db93146d1e33df6030c48ba235901d3d013c8}.
\]
The next compulsory companion audit is V45.

% =====================================================================
% V41
% =====================================================================

\clearpage
\chapter{V41: Norm dominance, all-row mass locality, and the finite
external-demand ledger}
\label{ch:V41-norm-locality}

\section{The capacity half of \(TQ_2\) is redundant}
\label{sec:V41-capacity-redundant}

V39 introduced norm and product-capacity estimates for both positive
squares of the signed ternary occurrence.  There is a useful
one-way order which was not used there: on a positive operator,
product-state capacity is bounded by the ordinary operator norm.
Consequently the genuinely missing analytic statement is a single
signed norm estimate, together with the same estimate for each
allowed once-paid occurrence.

\begin{lemma}[Positive capacity is norm dominated]
\label{lem:V41-capacity-norm}
Let \(X\ge0\) act on a finite tensor product.  Then
\begin{equation}
 0\le \kappa_\otimes(X)\le \|X\|.
\label{eq:V41-capacity-norm}
\end{equation}
The same conclusion holds after adjoining a fixed finite-dimensional
auxiliary factor.
\end{lemma}

\begin{proof}
For every product density matrix
\(\rho=\rho_1\otimes\cdots\otimes\rho_m\), positivity and the
operator inequality \(0\le X\le\|X\|I\) give
\[
 0\le\operatorname{Tr}(\rho X)
 \le \|X\|\operatorname{Tr}\rho=\|X\|.
\]
Taking the supremum over product densities proves
\eqref{eq:V41-capacity-norm}.  Nothing in the argument depends on the
number or dimensions of the tensor factors.
\end{proof}

\begin{theorem}[One-line form of the ternary-square target]
\label{thm:V41-TQ2-collapse}
Let \(S\) be any bounded complete ternary occurrence and let
\(\mathfrak h\ge0\).  Up to changing constants, the following are
equivalent:
\begin{enumerate}[label=\textup{(\roman*)},leftmargin=4em]
\item \(\|S\|\le C\mathfrak h\);
\item
\[
 \|S^*S\|,\ \|SS^*\|\le C\mathfrak h^2;
\]
\item the two norm bounds in \textup{(ii)} and the two capacity
      bounds
\[
 \kappa_\otimes(S^*S),\ \kappa_\otimes(SS^*)
 \le C\mathfrak h^2
\]
hold simultaneously.
\end{enumerate}
Moreover
\begin{equation}
 {\cal E}_\otimes({\mathscr H}(S))\le \|S\|^2.
\label{eq:V41-energy-by-norm}
\end{equation}
The assertions remain true occurrence by occurrence for every
once-paid version \(S^{\rm pay}\).
\end{theorem}

\begin{proof}
The \(C^*\)-identity gives
\[
 \|S^*S\|=\|SS^*\|=\|S\|^2.
\]
Thus \textup{(i)} and \textup{(ii)} are equivalent after taking a
square or a square root.  By
\Cref{lem:V41-capacity-norm}, \textup{(ii)} implies all the capacity
bounds in \textup{(iii)}, while \textup{(iii)} contains
\textup{(ii)}.

For the energy statement, use
\Cref{lem:V39-Hermitianized-absolute} and then
\Cref{lem:V41-capacity-norm}:
\[
 {\cal E}_\otimes({\mathscr H}(S))
 =\|S\|\,\kappa_\otimes(|{\mathscr H}(S)|)
 \le \|S\|\,\|{\mathscr H}(S)\|
 =\|S\|^2.
\]
The proof is purely operator-theoretic and can therefore be repeated
for each paid occurrence without moving or duplicating its payer.
\end{proof}

\begin{corollary}[Corrected analytic content of \(TQ_2\)]
\label{cor:V41-TQ1}
The unpaid part of \(TQ_2\) in \Cref{def:V39-TQ2} is equivalent to
the single estimate
\begin{equation}
 \boxed{
 \|\widetilde S_{134,r}\|
 \le C\mathfrak h_{134,r}.}
\label{eq:V41-TQ1}
\end{equation}
The complete target, including the one-payer ledger, is equivalent to
\eqref{eq:V41-TQ1} and the analogous signed norm estimate for every
allowed paid occurrence.  No independent proof of
\eqref{eq:V39-TQ2-capacity} is required after those norm estimates
are known.
\end{corollary}

\begin{proof}
Apply \Cref{thm:V41-TQ2-collapse} to the unpaid occurrence and then
separately to every occurrence in the finite payer ledger of
\Cref{def:V39-TQ2}.  The equivalence
\eqref{eq:V40-normalized-ternary-norm} is recovered as the unpaid
case.
\end{proof}

\begin{firewall}[Role of the eight-replica identity]
\label{fire:V41-eight-replica-role}
The exact polarization in \Cref{thm:V40-polarized-square} remains a
valid possible device for proving a signed norm bound.  It is not a
second independent endpoint obligation: once the ordinary norm of
the signed occurrence has the required stock, positivity and
\Cref{lem:V41-capacity-norm} supply its product-state capacity for
free.  In particular, a separate semigroup argument aimed only at
capacity cannot close the residual.
\end{firewall}

\section{Exact factorization by local row fractions}
\label{sec:V41-row-fractions}

Assume the three relevant row masses are positive; the zero-mass
cases have vanishing source and are removed before normalization.
Put
\begin{equation}
 \lambda_{i,r}:=\frac{a_{i,r}}{\Xi_i},
 \qquad i\in\{1,3,4\}.
\label{eq:V41-local-fractions}
\end{equation}
Since \(\Xi_i\) is the complete row mass, one has
\(0\le\lambda_{i,r}\le1\).

\begin{lemma}[Exact three-row factorization]
\label{lem:V41-stock-factorization}
The normalized ternary star stock is
\begin{equation}
 \boxed{
 \mathfrak h_{134,r}
 =\lambda_{1,r}^{2}\lambda_{3,r}\lambda_{4,r}.}
\label{eq:V41-stock-factorization}
\end{equation}
Thus its row-demand multiset is
\begin{equation}
 \nu_{134}=\{1,1,3,4\},
 \qquad (\nu_1,\nu_3,\nu_4)=(2,1,1).
\label{eq:V41-demand-vector}
\end{equation}
\end{lemma}

\begin{proof}
By definition \(h_{ij,r}=a_{i,r}a_{j,r}\).  Hence
\[
 \mathfrak h_{134,r}
 =\frac{(a_{1,r}a_{3,r})(a_{1,r}a_{4,r})}
 {\Xi_1^2\Xi_3\Xi_4}
 =\left(\frac{a_{1,r}}{\Xi_1}\right)^2
  \frac{a_{3,r}}{\Xi_3}\frac{a_{4,r}}{\Xi_4}.
\]
The exponents give \eqref{eq:V41-demand-vector}.
\end{proof}

\begin{proposition}[All-row mass-locality obstruction]
\label{prop:V41-all-row-locality}
Fix the local data at \(r\), and suppose the corresponding bare local
multiplier \(\widetilde S_{134,r}\) is nonzero and unchanged when
spectator coordinates are adjoined.  If the off-\(r\) mass of any
one row \(i\in\{1,3,4\}\) tends to infinity while the other data are
fixed, then
\[
 \mathfrak h_{134,r}\longrightarrow0
 \quad\text{but}\quad
 \|\widetilde S_{134,r}\|\not\longrightarrow0.
\]
Consequently a uniform proof of the bare estimate
\eqref{eq:V41-TQ1} must, for every row whose local fraction is not
bounded below, retain a complete external carrier or a signed
cancellation involving that row.  Dominance of row \(1\) alone is
not sufficient when either leaf fraction tends to zero.
\end{proposition}

\begin{proof}
The factorization \eqref{eq:V41-stock-factorization} contains a
positive power of each of \(\Xi_1^{-1},\Xi_3^{-1},\Xi_4^{-1}\).
Increasing off-\(r\) mass in row \(i\) keeps the numerator and local
operator fixed and sends the corresponding fraction to zero.  The
right side of \eqref{eq:V41-TQ1} therefore tends to zero, whereas its
left side is a fixed positive number.  The only ways to avoid this
contradiction are for the relevant local fraction to stay bounded
below, for the complete operator word itself to contain an external
factor which changes along the family, or for the local term to be
part of a larger signed cancellation.  Taking \(i=3\) or \(i=4\)
proves the final assertion even under a lower bound on
\(\lambda_{1,r}\).
\end{proof}

\begin{corollary}[Second correction of the bare marked target]
\label{cor:V41-correct-mark-target}
In the absence of external dressings or source-level cancellation,
the bare local part of \(TQ_2^{\rm mark}\) from
\Cref{def:V40-two-targets} must be restricted to
\begin{equation}
 \lambda_{1,r}\ge c_1,\qquad
 \lambda_{3,r}\ge c_3,\qquad
 \lambda_{4,r}\ge c_4
\label{eq:V41-fully-local-cell}
\end{equation}
for fixed positive thresholds.  A lower bound on only one marked row
does not isolate a valid bare endpoint target.
\end{corollary}

\begin{proof}
Apply \Cref{prop:V41-all-row-locality} to each row not covered by a
lower bound.  This refines, rather than reverses, the trichotomy in
\Cref{lem:V40-local-external}: that lemma tested row \(1\), whereas
the exact factorization shows that the same test is compulsory for
both leaves.
\end{proof}

\section{Closure of the fully local-dominant cell}
\label{sec:V41-local-close}

The preceding correction also removes one genuine subproblem.  When
all three local fractions are bounded below, the normalized stock is
itself bounded below, and the standing contraction and payer caps
are sufficient.  No carré-du-champ estimate is needed on this cell.

\begin{theorem}[Fully local ternary endpoint]
\label{thm:V41-fully-local-TQ2}
Fix \(c_1,c_3,c_4>0\).  On the cell
\eqref{eq:V41-fully-local-cell}, the unpaid ternary occurrence obeys
\begin{equation}
 \|\widetilde S_{134,r}\|
 \le
 \frac{8}{c_1^2c_3c_4}\,\mathfrak h_{134,r}.
\label{eq:V41-local-signed-norm}
\end{equation}
Hence both square-norm bounds, both product-capacity bounds, and the
quadratic energy bound required by \(TQ_2\) hold on this cell.

The same conclusion holds for every local or adjacent once-paid
occurrence to which the standing complete payer cap applies, with
the constant multiplied by that cap.  In particular, under the
normalization and payer hypotheses recorded in
\Cref{prop:V34-terminal-interfaces}, the complete fully local
\(TQ_2\) payer ledger closes.
\end{theorem}

\begin{proof}
Normalized Wilson row functions and the physical output compression
are contractions.  Thus
\Cref{prop:V40-universal-difference} gives
\(\|\widetilde S_{134,r}\|\le8\).  On
\eqref{eq:V41-fully-local-cell},
\[
 \mathfrak h_{134,r}
 =\lambda_{1,r}^2\lambda_{3,r}\lambda_{4,r}
 \ge c_1^2c_3c_4.
\]
Multiplying this inequality by \(8/(c_1^2c_3c_4)\) proves
\eqref{eq:V41-local-signed-norm}.  The square, capacity, and energy
conclusions follow from \Cref{thm:V41-TQ2-collapse}.

For a paid occurrence, keep the payer in its original position.  The
complete payer cap bounds the norm of that exact composition by a
fixed constant \(C_{\rm pay}\); all remaining normalized moment,
separator, and compression maps are contractions.  Replacing \(8\)
by \(C_{\rm pay}\) in the preceding argument proves the signed norm
bound for that occurrence.  Apply
\Cref{thm:V41-TQ2-collapse} occurrence by occurrence.  No copy of the
payer is formed.
\end{proof}

\begin{firewall}[What closed the local cell]
\label{fire:V41-local-scope}
The proof of \Cref{thm:V41-fully-local-TQ2} uses a lower bound on all
four factors in the demand multiset \(\{1,1,3,4\}\).  It does not
show heat-kernel smallness of the ternary cumulant and cannot be
continued to a cell on which one of the three row fractions tends to
zero.  Its purpose is to remove the compact fully local cell before
the external banks are analyzed.
\end{firewall}

\section{The seven external-demand cells}
\label{sec:V41-seven-cells}

Choose fixed thresholds \(0<c_i<1\) for
\(i\in\{1,3,4\}\).  Declare row \(i\) local when
\(\lambda_{i,r}\ge c_i\) and external otherwise.  For an external
row,
\begin{equation}
 \sum_{e\ne r}a_{i,e}
 =\Xi_i-a_{i,r}
 >(1-c_i)\Xi_i.
\label{eq:V41-external-mass}
\end{equation}
Equation \eqref{eq:V41-external-mass} is only a scalar localization
statement; it does not yet select a legal operator bank.

\begin{proposition}[Exact finite external-demand ledger]
\label{prop:V41-seven-cell-ledger}
Apart from the fully local cell closed in
\Cref{thm:V41-fully-local-TQ2}, the ternary endpoint has exactly the
following seven row cells and missing demand multisets:
\begin{center}
\begin{tabular}{@{}ccl@{}}
\toprule
external rows & missing demand & required source content\\
\midrule
\(\{3\}\) & \(\{3\}\) & one genuine row-\(3\) carrier\\
\(\{4\}\) & \(\{4\}\) & one genuine row-\(4\) carrier\\
\(\{3,4\}\) & \(\{3,4\}\) & both leaf demands\\
\(\{1\}\) & \(\{1,1\}\) & a genuine quadratic row-\(1\) carrier\\
\(\{1,3\}\) & \(\{1,1,3\}\) & quadratic center plus leaf \(3\)\\
\(\{1,4\}\) & \(\{1,1,4\}\) & quadratic center plus leaf \(4\)\\
\(\{1,3,4\}\) & \(\{1,1,3,4\}\) & all external demands\\
\bottomrule
\end{tabular}
\end{center}
Here ``required source content'' means a complete occurrence in the
exact binary--ternary word, or an exact signed cancellation involving
that occurrence.  It is not permission to attach an abstract scalar
mass after taking norms.
\end{proposition}

\begin{proof}
There are three binary local/external decisions, hence eight cells.
The empty external set is the fully local cell.  For every nonempty
set \(E\subset\{1,3,4\}\), restrict the demand multiset
\(\{1,1,3,4\}\) from \eqref{eq:V41-demand-vector} to the labels in
\(E\).  This gives the seven displayed rows and proves exhaustion and
mutual exclusivity.  The multiplicity two for row \(1\) follows from
the square \(\lambda_{1,r}^2\), not from two names for one bank.
Finally \Cref{prop:V41-all-row-locality} shows why every missing
demand must remain coupled to changing external data in the complete
word.
\end{proof}

\begin{firewall}[A double center demand is not a duplicated bank]
\label{fire:V41-no-center-duplication}
In the four cells containing external row \(1\), one complete linear
row-\(1\) bank cannot be used twice.  The demand \(\{1,1\}\) may be
discharged only by two actual source occurrences, by a quadratic or
cross term of one complete bank retained before the positive
envelope, or by a signed identity which produces the same two marks.
This is the exact analogue of the two-mark rule in the SM audit and
does not follow merely from the scalar inequality
\eqref{eq:V41-external-mass}.
\end{firewall}

\begin{corollary}[Corrected residual after V41]
\label{cor:V41-residual}
For the double-demand \(2+1+1\) star, the fully local component of
\(TQ_2\) is closed and the separate capacity obligation is removed.
The remaining ternary problem consists exactly of signed ordinary
norm estimates for the seven external-demand cells of
\Cref{prop:V41-seven-cell-ledger}, including their exact once-paid
versions.  The cells with external center row require two genuine
center marks; the other three nonempty cells require only the listed
leaf marks in addition to whatever local rows are already present.
\end{corollary}

\begin{proof}
Use \Cref{thm:V41-fully-local-TQ2} on the empty external set,
\Cref{cor:V41-TQ1} to remove the redundant capacity obligations, and
\Cref{prop:V41-seven-cell-ledger} for the exhaustive complement.
The center multiplicity is protected by
\Cref{fire:V41-no-center-duplication}.
\end{proof}

\begin{assessment}[Why the V40 semigroup task moves upstream]
\label{ass:V41-upstream}
The proposed one-link carré-du-champ calculation is unnecessary on
the fully local cell and insufficient on an external-row cell unless
the changing external carrier is present in the same operator word.
The next calculation should therefore start with the two singly
external leaf cells \(\{3\}\) and \(\{4\}\).  Their exact prefix
singleton/suffix positions must be restored before any norm is
taken.  Only after those one-demand words are settled should the
quadratic center cells be attacked.
\end{assessment}

\begin{theorem}[Third-series V41 result ledger]
\label{thm:V41-results}
Relative to V1--V40, V41 proves:
\begin{enumerate}[label=\textup{(V41.\arabic*)},leftmargin=6.3em]
\item positive product-state capacity is norm dominated;
\item the four norm/capacity lines of \(TQ_2\) reduce to one signed
      norm line, occurrence by occurrence in the payer ledger;
\item the endpoint energy is bounded directly by the square of that
      signed norm;
\item the exact local-fraction factorization
      \(\mathfrak h=\lambda_1^2\lambda_3\lambda_4\);
\item the all-row mass-locality obstruction and correction of the
      one-row marked target;
\item closure of the fully local-dominant unpaid and capped-payer
      cells; and
\item the exhaustive seven-cell external-demand ledger with center
      multiplicity two protected against bank duplication.
\end{enumerate}
\end{theorem}

\begin{proof}
These are
\Cref{lem:V41-capacity-norm,
thm:V41-TQ2-collapse,
cor:V41-TQ1,
lem:V41-stock-factorization,
prop:V41-all-row-locality,
cor:V41-correct-mark-target,
thm:V41-fully-local-TQ2,
prop:V41-seven-cell-ledger,
cor:V41-residual}.
\end{proof}

\begin{programme}[V42: exact singly external leaf word]
\label{prog:V41-V42}
\begin{enumerate}[leftmargin=3em]
\item Restore the complete prefix-singleton and suffix factors
      suppressed in \eqref{eq:V37-binary-ternary-word}.
\item On the \(\{3\}\) cell, split the off-\(r\) row-\(3\) mass by
      its actual position before or after the joining coordinate.
\item Write a disjoint first-hit telescope selecting one genuine
      row-\(3\) occurrence; do not infer an occurrence from
      \eqref{eq:V41-external-mass} alone.
\item Prove its signed norm bound with the ternary factor kept
      assembled, and repeat by symmetry for the \(\{4\}\) cell.
\item Check every local, bank, separator, and suffix payer position
      before merging the two leaf cases.
\item If no complete leaf occurrence exists in one positional cell,
      return to the exact Leonov--Shiryaev source identity rather than
      inventing a bank.
\item Leave the quadratic center cells untouched until the
      one-demand leaf word is certified.  The next compulsory
      companion audit remains V45.
\end{enumerate}
\end{programme}

\begin{firewall}[Claim boundary after third-series V41]
\label{fire:V41-claim-boundary}
V41 closes only the fully local-dominant part of the ternary endpoint
and simplifies the remaining analytic target.  It does not prove the
seven external-demand cells, the full binary--ternary endpoint, the
remaining \(R_1\) paths, or the full \(2+1+1\) source.  The \(3+1\),
global discrete, globally aggregated \([4]\), and \([3,3]\) families,
higher MTA, WUFC, CPM, cutoff removal, reflection positivity,
Osterwalder--Schrader reconstruction, construction of the continuum
Yang--Mills measure, and a uniform positive continuum spectral gap
remain open.  The full Yang--Mills mass gap has not been proved.
\end{firewall}

\paragraph{Parent-version record.}
V41 was copied byte-for-byte from
\texttt{Renewal\_Yang\_Mills\_Mass\_Gap\_Research\_Notes\_V40.tex}
before this chapter was appended.  The V40 parent had \(21\,339\)
source lines, \(771\,054\) bytes, and SHA--256
\[
 \texttt{b99c89e805430fd4a6c01b7ab5dd2a87a67ff4878c3fb8e3d198585952bae59d}.
\]
The next compulsory companion audit is V45.

% =====================================================================
% V42
% =====================================================================

\clearpage
\chapter{V42: Undoing the overstrong endpoint split and restoring the
exact center demand}
\label{ch:V42-center-demand}

\section{Clear the source normalization before splitting endpoints}
\label{sec:V42-clear-normalization}

V41 correctly simplified the operator-theoretic content of the
isolated target \(TQ_2\), but its seven-cell residual treated every
denominator inside the isolated ternary stock as an excess demand of
the original source.  The exact \(2+1+1\) normalization in V35 shows
that this is false: rows \(2,3,4\) are leaves of the star and their
normalizations cancel when the actual binary--ternary target is
cleared.

\begin{theorem}[Exact cleared form of \(BTE_2\)]
\label{thm:V42-cleared-BTE2}
Assume \(\Xi_i>0\) for \(1\le i\le4\).  The binary--ternary endpoint
estimate \eqref{eq:V37-BTE2} is equivalent to
\begin{equation}
 \boxed{
 \kappa_\otimes(\mathcal W_{BT,r}^{\rm pay})
 \le
 C\frac{H_{12}^{<r}h_{13,r}h_{14,r}}{\Xi_1^2}.}
\label{eq:V42-cleared-BTE2}
\end{equation}
The row-interchanged star has the identical formula with
\(1\leftrightarrow2\).
\end{theorem}

\begin{proof}
Insert \eqref{eq:V37-star-stock} into
\eqref{eq:V37-BTE2}:
\[
 \frac{\kappa_\otimes(\mathcal W_{BT,r}^{\rm pay})}
 {\Xi_1\Xi_2\Xi_3\Xi_4}
 \le
 C\frac{H_{12}^{<r}h_{13,r}h_{14,r}}
 {\Xi_1^3\Xi_2\Xi_3\Xi_4}.
\]
Multiplication by the positive product
\(\Xi_1\Xi_2\Xi_3\Xi_4\) gives
\eqref{eq:V42-cleared-BTE2}; division by the same product proves the
converse.  Relabelling the center proves the symmetric statement.
\end{proof}

\begin{corollary}[Only the star center has excess mass demand]
\label{cor:V42-only-center}
The denominator-demand vector of the cleared row-\(1\) star target is
\begin{equation}
 (2,0,0,0).
\label{eq:V42-center-demand-vector}
\end{equation}
In particular no off-\(r\) row-\(3\) or row-\(4\) bank is required
merely because \(a_{3,r}/\Xi_3\) or
\(a_{4,r}/\Xi_4\) is small.
\end{corollary}

\begin{proof}
The only global mass remaining in
\eqref{eq:V42-cleared-BTE2} is \(\Xi_1^{-2}\).  This is also the
excess-demand vector \(2e_1\) of \(T_{11}^{(0)}\) in
\Cref{lem:V35-demand-census}.  Rows \(3,4\) have degree one in the
tree, and \Cref{lem:V35-normalization,cor:V35-hot-leaves} show
directly that leaf rows have zero excess demand.
\end{proof}

\begin{proposition}[Reconciliation of V35 and V41]
\label{prop:V42-reconcile-V41}
The identity
\(\mathfrak h_{134,r}=\lambda_{1,r}^2
\lambda_{3,r}\lambda_{4,r}\) in
\Cref{lem:V41-stock-factorization} remains correct for the isolated
ternary endpoint stock.  The all-row obstruction in
\Cref{prop:V41-all-row-locality} remains correct under its explicit
hypothesis that the bare local multiplier is unchanged as spectator
mass is added.  However the inference in
\Cref{cor:V41-residual} from this isolated target to seven residual
cells of the complete source is withdrawn.  The actual source has
only the center demand \eqref{eq:V42-center-demand-vector}.
\end{proposition}

\begin{proof}
The first two assertions are the algebra and fixed-local-operator
argument already proved in V41.  They concern the sufficient
endpoint target \(TQ_2\).  The complete source target is instead
\eqref{eq:V42-cleared-BTE2}, in which the leaf masses have canceled.
Therefore the leaf cells in the V41 table measure the extra strength
of isolating \(TQ_2\); they are not residual source sectors.  The
center vector follows from \Cref{cor:V42-only-center}.
\end{proof}

\begin{firewall}[Hot leaves remain harmless]
\label{fire:V42-hot-leaves}
No subsequent proof may introduce a row-\(3\) or row-\(4\) external
bank solely to repair the isolated denominator in
\(\mathfrak h_{134,r}\).  Such a bank is legal only if it already
occurs in the exact source word for an independent reason.  This
firewall restores the source-level conclusion of
\Cref{cor:V35-hot-leaves}.
\end{firewall}

\section{Why \(TQ_2\) is sufficient but need not be necessary}
\label{sec:V42-TQ2-overstrong}

\begin{proposition}[A product estimate does not imply an endpoint
estimate]
\label{prop:V42-no-reverse-endpoint}
There is no universal implication of the form
\begin{equation}
 \|AS\|\le Cpq
 \quad\Longrightarrow\quad
 \|S\|\le C' q
\label{eq:V42-false-reverse}
\end{equation}
for bounded operators \(A,S\) and positive scalar stocks \(p,q\),
without a quantitative lower bound or invertibility statement for
\(A\).  The conclusion remains false if the left norm is replaced by
product-state capacity.
\end{proposition}

\begin{proof}
On any nonzero finite-dimensional Hilbert space take, for
\(0<\varepsilon<1\),
\[
 A=\varepsilon^2I,qquad S=I,qquad p=q=\varepsilon.
\]
Then \(\|AS\|=\varepsilon^2=pq\), but
\(\|S\|=1\) cannot be bounded by \(C'\varepsilon\) uniformly as
\(\varepsilon\downarrow0\).  These scalar multiples of the identity
have the same norm and product-state capacity, proving the final
assertion.
\end{proof}

\begin{corollary}[The endpoint-energy route may be abandoned]
\label{cor:V42-abandon-TQ2}
The implication \(TQ_2\Rightarrow BTE_2\) proved in
\Cref{cor:V39-TQ2-suffices} is valid, but no converse has been proved
or is available from abstract operator inequalities.  Failure to
prove the isolated signed norm \eqref{eq:V41-TQ1} therefore does not
imply failure of \(BTE_2\).  The legitimate weaker target is the
complete-word estimate \eqref{eq:V42-cleared-BTE2}.
\end{corollary}

\begin{proof}
The forward implication is the cited result.  The abstract
non-implication is \Cref{prop:V42-no-reverse-endpoint}; the physical
binary prefix can carry damping or cancellation which disappears if
the ternary endpoint is estimated separately.  The exact weaker
target is equivalent to \(BTE_2\) by
\Cref{thm:V42-cleared-BTE2}.
\end{proof}

\section{Return to the complete two-center word}
\label{sec:V42-return-BSC2}

For a marked coordinate tuple \((e,r,q)\), V36 already defined the
complete signed summand \(\mathcal W_{e,r,q}^{\rm pay}\), its actual
marked mass \(A_1(Q_{e,r,q})\), and the balanced star-collision card
\(BSC_2\).  These objects occur before the later binary--ternary
endpoint separation.

\begin{theorem}[The exact marked-heavy gate is still \(BSC_2\)]
\label{thm:V42-BSC2-restored}
On a balanced marked packet satisfying
\[
 A_1(Q_{e,r,q})\ge\varepsilon\Xi_1,
\]
the required two center debits follow from the single complete-word
estimate
\begin{equation}
 \kappa_\otimes(\mathcal W_{e,r,q}^{\rm pay})
 \le C
 \min\left\{s_\alpha^2,
 A_1(Q_{e,r,q})^{-2}\right\}
 h_{12,e}h_{13,r}h_{14,r}.
\label{eq:V42-restored-BSC2}
\end{equation}
After summing the disjoint prefix marks \(e<r\), its inverse-mass
branch gives precisely the cleared target
\eqref{eq:V42-cleared-BTE2}.  It creates no leaf demand and contains
the payer exactly once.
\end{theorem}

\begin{proof}
Equation \eqref{eq:V42-restored-BSC2} is the card
\eqref{eq:V36-BSC2} with its complete source occurrence retained.
By marked heaviness,
\[
 A_1(Q_{e,r,q})^{-2}
 \le\varepsilon^{-2}\Xi_1^{-2}.
\]
Therefore the second branch of the minimum gives
\[
 \kappa_\otimes(\mathcal W_{e,r,q}^{\rm pay})
 \le
 C_\varepsilon\Xi_1^{-2}
 h_{12,e}h_{13,r}h_{14,r}.
\]
The first-connectivity prefix marks are disjoint and
\(\sum_{e<r}h_{12,e}=H_{12}^{<r}\).  Summing gives
\eqref{eq:V42-cleared-BTE2}.  This is the proof of
\Cref{lem:V36-BSC2-suffices}, now compared explicitly with the
cleared \(BTE_2\) normalization.  Neither the summation nor the
inequality changes the unique payer position.
\end{proof}

\begin{firewall}[Do not square the center mass after separation]
\label{fire:V42-no-separated-center-square}
The factor \(A_1(Q_{e,r,q})^{-2}\) in
\eqref{eq:V42-restored-BSC2} must be obtained from one estimate on the
complete signed summand.  It may not be reconstructed by giving one
copy of \(A_1(Q)^{-1}\) to the binary prefix and another to the
ternary endpoint unless the exact word contains the two corresponding
complete response occurrences.  This is the original V36 occurrence
rule, unaffected by the V39 square identity.
\end{firewall}

\begin{theorem}[Recovered exact double-demand residual]
\label{thm:V42-recovered-residual}
For the two star terms in the hot \(2+1+1\) reduction, all private,
unbalanced, and genuine-two-\(A,F,I\)-occurrence cases remain closed.
The unresolved double-center sector is exactly the three classes in
\Cref{thm:V36-refined-residual}:
\begin{enumerate}[label=\textup{(\alph*)},leftmargin=4em]
\item a balanced marked packet requiring the complete \(BSC_2\) card;
\item a used \(A,F,I\) bank with at most one genuine complete
      stock-carrying occurrence; or
\item a used directed or coupled bank for which only a one-token
      response or a middle-factor norm is available.
\end{enumerate}
The isolated targets \(TQ_2\) may still close subcells, including the
fully local-dominant subcell of
\Cref{thm:V41-fully-local-TQ2}, but they are no longer imposed on the
whole residual.  The one-token path residual \(R_1\) remains
unchanged.
\end{theorem}

\begin{proof}
The exact five-role decomposition in V37 changes the presentation of
the two lower-order stars but not their source word or their center
demand.  The closures and exhaustive residual classification are
\Cref{thm:V36-refined-residual}.  V39 proves only the sufficient
implication \(TQ_2\Rightarrow BTE_2\), and
\Cref{cor:V42-abandon-TQ2} shows that it is not a necessary interface.
Thus removing the overstrong endpoint obligation returns exactly to
the V36 list, with no new cell.  The V41 local theorem remains a valid
sufficient closure on its stated subcell.  No later chapter changes
the independent \(R_1\) ledger.
\end{proof}

\begin{assessment}[The next analytic object]
\label{ass:V42-next-object}
The smallest unresolved object is again the balanced protected
complete summand \(\mathcal W_{e,r,q}^{\rm pay}\) in
\eqref{eq:V42-restored-BSC2}.  The V37 Leonov--Shiryaev identity and
the V40 replica polarization may be used inside it, but the binary
prefix, complete ternary correction, separators, output compression,
and sole payer must be assembled before a positive norm or capacity
is taken.  The goal is the second branch
\(A_1(Q)^{-2}\), not a separately normalized leaf endpoint.
\end{assessment}

\begin{theorem}[Third-series V42 result ledger]
\label{thm:V42-results}
Relative to V1--V41, V42 proves:
\begin{enumerate}[label=\textup{(V42.\arabic*)},leftmargin=6.3em]
\item the exact cleared form
      \eqref{eq:V42-cleared-BTE2} of the binary--ternary target;
\item cancellation of all leaf-row masses and recovery of the sole
      excess-demand vector \((2,0,0,0)\);
\item reconciliation with the V35 hot-leaf theorem and withdrawal of
      the V41 seven-cell table as a source-residual classification;
\item an abstract proof that a complete product bound need not imply
      the isolated endpoint bound used by \(TQ_2\);
\item proof that \(TQ_2\) is sufficient but not a necessary declared
      interface for \(BTE_2\);
\item restoration of the complete \(BSC_2\) marked-heavy gate; and
\item recovery of the exact three-class V36 double-center residual.
\end{enumerate}
\end{theorem}

\begin{proof}
These are
\Cref{thm:V42-cleared-BTE2,
cor:V42-only-center,
prop:V42-reconcile-V41,
fire:V42-hot-leaves,
prop:V42-no-reverse-endpoint,
cor:V42-abandon-TQ2,
thm:V42-BSC2-restored,
thm:V42-recovered-residual}.
\end{proof}

\begin{programme}[V43: balanced protected complete-word square]
\label{prog:V42-V43}
\begin{enumerate}[leftmargin=3em]
\item Fix one marked-heavy balanced packet
      \(\mathcal W_{e,r,q}^{\rm pay}\) and one allowed payer
      position; do not sum payer positions.
\item Insert the exact V37 five-role identity at the joining
      coordinate while retaining the binary prefix and all complete
      spectators.
\item Apply the V40 polarization only after multiplication by the
      actual binary-prefix occurrence.  Square the complete word, not
      the isolated ternary endpoint.
\item Identify the two genuine row-\(1\) marks in that square and
      compare them with \(A_1(Q_{e,r,q})^2\) before any positive
      envelope.
\item Test the proposed inequality on the normalized \(SU(2)\)
      character channel and the maximally entangled projector
      diagnostic from V37--V38.
\item If the complete-word square still yields only one center mark,
      construct the smallest physical counterpacket to \(BSC_2\) and
      return to the used-bank classification instead of adding a
      fictitious occurrence.
\item Leave the one-occurrence \(A,F,I\), directed/coupled, and
      \(R_1\) sectors unchanged.  The next compulsory companion audit
      remains V45.
\end{enumerate}
\end{programme}

\begin{firewall}[Claim boundary after third-series V42]
\label{fire:V42-claim-boundary}
V42 corrects the endpoint overfactorization and restores the exact
complete-word residual, but it does not prove the balanced protected
\(BSC_2\) card, the one-occurrence or directed/coupled center cases,
the \(R_1\) paths, or the full \(2+1+1\) source.  The \(3+1\), global
discrete, globally aggregated \([4]\), and \([3,3]\) families,
higher MTA, WUFC, CPM, cutoff removal, reflection positivity,
Osterwalder--Schrader reconstruction, construction of the continuum
Yang--Mills measure, and a uniform positive continuum spectral gap
remain open.  The full Yang--Mills mass gap has not been proved.
\end{firewall}

\paragraph{Parent-version record.}
V42 was copied byte-for-byte from
\texttt{Renewal\_Yang\_Mills\_Mass\_Gap\_Research\_Notes\_V41.tex}
before this chapter was appended.  The V41 parent had \(21\,813\)
source lines, \(789\,225\) bytes, and SHA--256
\[
 \texttt{607e91a2ed04bda8aba3095b7ddd741c9887220cb0ee2d372f8cd02d9c2b4bb3}.
\]
The next compulsory companion audit is V45.

% =====================================================================
% V43
% =====================================================================

\clearpage
\chapter{V43: The complete binary--ternary square and its exact
replica carrier}
\label{ch:V43-complete-square}

\section{Cauchy--Schwarz after, not before, word assembly}
\label{sec:V43-word-CS}

For clarity, write
\[
 \kappa_\otimes^{\rm abs}(W):=\kappa_\otimes(|W|).
\]
This is the product-state quantity used when a signed physical word
is passed to a positive majorant.  The following elementary estimate
allows the complete word to be squared without first estimating its
ternary endpoint in ordinary norm.

\begin{lemma}[Complete-word product-state Cauchy--Schwarz]
\label{lem:V43-complete-word-CS}
For every bounded operator \(W\),
\begin{align}
 \kappa_\otimes^{\rm abs}(W)^2
 &\le \kappa_\otimes(W^*W),
\label{eq:V43-right-CS}\\
 \kappa_\otimes^{\rm abs}(W^*)^2
 &\le \kappa_\otimes(WW^*).
\label{eq:V43-left-CS}
\end{align}
\end{lemma}

\begin{proof}
For a product density matrix \(\rho\), Cauchy--Schwarz in the
positive functional \(X\mapsto\operatorname{Tr}(\rho X)\) gives
\[
 \operatorname{Tr}(\rho|W|)^2
 \le \operatorname{Tr}(\rho)\operatorname{Tr}(\rho|W|^2)
 =\operatorname{Tr}(\rho W^*W).
\]
Take the supremum over product densities.  Since the square function
is increasing on nonnegative scalars, the supremum of the left side
is \(\kappa_\otimes^{\rm abs}(W)^2\).  Replacing \(W\) by \(W^*\)
proves \eqref{eq:V43-left-CS}.
\end{proof}

For a fixed marked packet put
\begin{equation}
 b_{e,r,q}
 :=
 \min\left\{s_\alpha^2,
 A_1(Q_{e,r,q})^{-2}\right\}
 h_{12,e}h_{13,r}h_{14,r}.
\label{eq:V43-BSC-stock}
\end{equation}

\begin{definition}[Complete-word quadratic \(BSC_2\) target]
\label{def:V43-CWQ2}
For one fixed payer placement \(\pi\), let
\(W_\pi=\mathcal W_{e,r,q}^{\rm pay,\pi}\) be the exact complete
signed word.  The right \(CWQ_2(\pi)\) target is
\begin{equation}
 \boxed{
 \kappa_\otimes(W_\pi^*W_\pi)
 \le C b_{e,r,q}^{2}.}
\label{eq:V43-CWQ2-right}
\end{equation}
The left target replaces \(W_\pi^*W_\pi\) by
\(W_\pi W_\pi^*\).  Only the orientation required by the fixed
physical output duality is imposed; no sum over payer placements is
part of the definition.
\end{definition}

\begin{theorem}[The complete-word square is sufficient for
\(BSC_2\)]
\label{thm:V43-CWQ2-suffices}
If the appropriate orientation of \(CWQ_2(\pi)\) holds for each
allowed payer placement \(\pi\), then the balanced protected packet
satisfies \eqref{eq:V42-restored-BSC2} and hence the marked-heavy
part of \(BTE_2\).
\end{theorem}

\begin{proof}
Use \eqref{eq:V43-right-CS} or
\eqref{eq:V43-left-CS}, according to the output orientation, and then
take the positive square root of
\eqref{eq:V43-CWQ2-right}.  This gives
\[
 \kappa_\otimes^{\rm abs}(W_\pi)
 \le C b_{e,r,q},
\]
which is exactly \eqref{eq:V42-restored-BSC2} for that placement.
The finite payer ledger is checked placement by placement, not summed.
Apply \Cref{thm:V42-BSC2-restored}.
\end{proof}

\begin{firewall}[Squaring a fixed paid word does not allocate a
second payer]
\label{fire:V43-square-one-payer}
The two appearances of \(W_\pi\) in \(W_\pi^*W_\pi\) are forced by
the exact Cauchy--Schwarz identity for the one already chosen paid
operator.  They are not two independent source terms and do not
authorize choosing two payer positions.  Any proof of
\eqref{eq:V43-CWQ2-right} must keep the same label \(\pi\) in both
factors and recover the square of the one-payer scale before taking
the final square root.
\end{firewall}

\section{Exact dressed replica formula}
\label{sec:V43-dressed-replica}

Fix first a payer placement outside the local ternary cumulant, so
that the exact word can be written
\begin{equation}
 W_\pi=L_\pi S R_\pi,
 \qquad S=\kappa_3(A,B,C),
\label{eq:V43-LSR-word}
\end{equation}
where \(L_\pi\) contains the binary prefix and all factors to its
side, while \(R_\pi\) contains the remaining complete suffix.  No
norm has yet been taken.  Let \(D(\boldsymbol\xi)\) be the polarized
four-replica difference tensor in
\eqref{eq:V40-D-xi}.

\begin{theorem}[Exact complete-word replica squares]
\label{thm:V43-dressed-square}
With independent four-replica packets \(\boldsymbol\xi\) and
\(\boldsymbol\eta\),
\begin{align}
 W_\pi^*W_\pi
 &={\mathbb E}_{\boldsymbol\eta,\boldsymbol\xi}
 \left[
 R_\pi^*D(\boldsymbol\eta)^*L_\pi^*L_\pi
 D(\boldsymbol\xi)R_\pi
 \right],
\label{eq:V43-dressed-right}\\
 W_\pi W_\pi^*
 &={\mathbb E}_{\boldsymbol\xi,\boldsymbol\eta}
 \left[
 L_\pi D(\boldsymbol\xi)R_\pi R_\pi^*
 D(\boldsymbol\eta)^*L_\pi^*
 \right].
\label{eq:V43-dressed-left}
\end{align}
Both expectations are positive after integration, although their
pointwise integrands need not be positive.  The binary prefix remains
inside \(L_\pi^*L_\pi\) or its left-square analogue.
\end{theorem}

\begin{proof}
By \Cref{thm:V40-four-replica},
\(S={\mathbb E}_{\boldsymbol\xi}D(\boldsymbol\xi)\).  Hence
\[
 W_\pi^*W_\pi
 =R_\pi^*S^*L_\pi^*L_\pi SR_\pi.
\]
Insert independent integral representations for \(S^*\) and \(S\)
and apply Fubini; all factors are bounded.  This proves
\eqref{eq:V43-dressed-right}.  The calculation
\(W_\pi W_\pi^*=L_\pi SR_\pi R_\pi^*S^*L_\pi^*\) proves
\eqref{eq:V43-dressed-left}.  Positivity follows from the integrated
operator squares themselves, not from pointwise kernels.  The final
statement is visible in the displayed formulas.
\end{proof}

\begin{corollary}[Actual center-bearing factors in the square]
\label{cor:V43-center-factors}
Each integrand of \eqref{eq:V43-dressed-right} contains two local
row-\(1\) difference factors, one in
\(D(\boldsymbol\eta)^*\) and one in
\(D(\boldsymbol\xi)\), together with the two conjugate occurrences of
the row-\(1\) binary prefix inside \(L_\pi^*L_\pi\).  These are actual
operator occurrences in the quadratic carrier.  They do not yet
constitute a proof of two inverse powers of
\(A_1(Q_{e,r,q})\).
\end{corollary}

\begin{proof}
The first assertion follows from the definition of
\(D\) in \eqref{eq:V40-D-xi}: each replica packet contains exactly
one \(A\)-difference.  The factor \(L_\pi\) contains the complete
binary prefix by \eqref{eq:V43-LSR-word}; hence
\(L_\pi^*L_\pi\) contains its adjoint and original occurrences.
The last sentence is necessary because occurrence counting alone
does not supply a heat, response, or inverse-mass estimate.
\end{proof}

\begin{firewall}[Replica samples are not new lattice coordinates]
\label{fire:V43-replica-coordinate}
The two copies of a local difference in
\eqref{eq:V43-dressed-right} arise from the exact positive square.
They may support a quadratic estimate on that square, but their
replica labels may not be ordered or retyped as two global coordinate
banks.  The V40 correction remains in force.
\end{firewall}

\section{The forbidden early estimate}
\label{sec:V43-forbidden-early}

\begin{proposition}[Norming the dressing first returns to the
overstrong endpoint target]
\label{prop:V43-norm-first-loss}
For the factorization \eqref{eq:V43-LSR-word},
\begin{equation}
 \kappa_\otimes(W_\pi^*W_\pi)
 \le
 \|L_\pi\|^2\|R_\pi\|^2\|S\|^2.
\label{eq:V43-coarse-square}
\end{equation}
If \(L_\pi,R_\pi\) are replaced only by support-uniform constants,
then proving \(CWQ_2(\pi)\) from
\eqref{eq:V43-coarse-square} requires the same isolated signed norm
smallness that made \(TQ_2\) overstrong.  Thus the binary prefix must
remain inside the expectation in
\eqref{eq:V43-dressed-right}.
\end{proposition}

\begin{proof}
By \Cref{lem:V41-capacity-norm} and submultiplicativity,
\[
 \kappa_\otimes(W_\pi^*W_\pi)
 \le\|W_\pi^*W_\pi\|
 =\|W_\pi\|^2
 \le\|L_\pi\|^2\|S\|^2\|R_\pi\|^2.
\]
After replacing the outside factors by constants, the only remaining
source of the small scale \(b_{e,r,q}^2\) is \(\|S\|^2\).  This is
precisely the isolated ordinary-norm obligation avoided in V42.
Keeping \(L_\pi^*L_\pi\) between the two replica differences is
therefore necessary for this route.
\end{proof}

\begin{firewall}[No pointwise positive envelope]
\label{fire:V43-no-pointwise-envelope}
The integrands in \eqref{eq:V43-dressed-right}--%
\eqref{eq:V43-dressed-left} need not be positive.  Replacing them by
pointwise absolute values before integration destroys both the
cross-replica cancellation and the interaction with
\(L_\pi^*L_\pi\).  Positivity belongs to the assembled expectation,
not to its individual replica kernels.
\end{firewall}

\section{Local-payer boundary and the V43 residual}
\label{sec:V43-residual}

\begin{proposition}[What remains when the payer lies in the ternary
factor]
\label{prop:V43-local-payer-boundary}
If the unique payer is inside the local ternary occurrence, write the
exact word as \(W_\pi=L S_\pi R\).  The algebraic square identities
\[
 W_\pi^*W_\pi=R^*S_\pi^*L^*LS_\pi R,
 \qquad
 W_\pi W_\pi^*=LS_\pi RR^*S_\pi^*L^*
\]
remain exact.  Equations \eqref{eq:V43-dressed-right}--%
\eqref{eq:V43-dressed-left} may be used only after an exact
polarization of the paid occurrence \(S_\pi\) is written; the unpaid
polarization of \(S\) cannot simply be multiplied by a second payer.
\end{proposition}

\begin{proof}
The two displayed identities are direct multiplication.  A payer
inside \(S_\pi\) changes that operator before the square is formed.
Replacing it by the unpaid \(S\) and then inserting a payer into each
replica factor would change the occurrence and would not follow from
\Cref{thm:V40-four-replica}.  Hence a paid polarization is a separate
algebraic prerequisite.
\end{proof}

\begin{theorem}[Exact V43 reduction of the balanced marked sector]
\label{thm:V43-residual}
The balanced marked-heavy \(BSC_2\) sector is reduced to:
\begin{enumerate}[label=\textup{(\roman*)},leftmargin=4em]
\item the oriented complete-word quadratic estimates
      \(CWQ_2(\pi)\) for payer placements outside the ternary
      occurrence, with exact carriers
      \eqref{eq:V43-dressed-right}--%
      \eqref{eq:V43-dressed-left}; and
\item the same positive-square estimates for local-payer words,
      preceded by an exact paid polarization as required in
      \Cref{prop:V43-local-payer-boundary}.
\end{enumerate}
Either collection is sufficient by
\Cref{thm:V43-CWQ2-suffices}.  Neither collection is proved in V43.
The used one-occurrence \(A,F,I\), directed/coupled center, and
one-token \(R_1\) residuals are unchanged.
\end{theorem}

\begin{proof}
Fix a payer placement.  If it is outside the ternary occurrence,
\Cref{thm:V43-dressed-square} supplies the exact positive square.  If
it is inside, use \Cref{prop:V43-local-payer-boundary}.  In both
cases \Cref{def:V43-CWQ2} is the sufficient quantitative target and
\Cref{thm:V43-CWQ2-suffices} gives \(BSC_2\).  The final residual
classes are those preserved in \Cref{thm:V42-recovered-residual}.
\end{proof}

\begin{assessment}[The remaining calculation is now positive but
not pointwise]
\label{ass:V43-next}
The complete-word route no longer requires ordinary-norm decay of the
isolated ternary occurrence.  Its missing estimate is instead the
product-state capacity of one assembled positive square.  The next
step is to expand \(L_\pi^*L_\pi\) by the exact binary-prefix
covariance telescope and determine whether its cross terms supply
the two center debits without pointwise absolute values.
\end{assessment}

\begin{theorem}[Third-series V43 result ledger]
\label{thm:V43-results}
Relative to V1--V42, V43 proves:
\begin{enumerate}[label=\textup{(V43.\arabic*)},leftmargin=6.3em]
\item product-state Cauchy--Schwarz for a complete signed word;
\item the precise oriented \(CWQ_2\) positive-square target
      sufficient for \(BSC_2\);
\item exact right and left eight-replica formulas with the binary
      prefix retained inside the square;
\item the inventory of actual center-bearing factors in that carrier;
\item proof that norming the dressing first returns to the overstrong
      isolated endpoint target;
\item the no-pointwise-envelope and one-payer firewalls; and
\item separation of outside-ternary and local-ternary payer
      prerequisites.
\end{enumerate}
\end{theorem}

\begin{proof}
These are
\Cref{lem:V43-complete-word-CS,
def:V43-CWQ2,
thm:V43-CWQ2-suffices,
thm:V43-dressed-square,
cor:V43-center-factors,
prop:V43-norm-first-loss,
fire:V43-no-pointwise-envelope,
prop:V43-local-payer-boundary,
thm:V43-residual}.
\end{proof}

\begin{programme}[V44: doubled binary-prefix first connectivity]
\label{prog:V43-V44}
\begin{enumerate}[leftmargin=3em]
\item Fix an outside-ternary payer placement and the right-square
      orientation \eqref{eq:V43-dressed-right}.
\item Substitute the exact binary covariance telescope into both
      occurrences inside \(L_\pi^*L_\pi\), keeping their coordinate
      order and cross terms.
\item Partition the doubled prefix by the first coordinate at which
      the two copies jointly connect the center output; do not order
      the probability replicas.
\item Identify diagonal, distinct-coordinate, and shared-bank terms
      as operator carriers before applying a positive majorant.
\item Prove a square-stock bound at scale \(b_{e,r,q}^2\), or isolate
      the first carrier on which only one factor of
      \(A_1(Q)^{-1}\) is obtained.
\item Repeat only after success for the left orientation and for the
      finite outside-ternary payer placements.
\item Leave the local-ternary payer for its separate paid
      polarization.  The next compulsory companion audit remains
      V45.
\end{enumerate}
\end{programme}

\begin{firewall}[Claim boundary after third-series V43]
\label{fire:V43-claim-boundary}
V43 constructs a correct complete-word positive-square target and
its exact replica carrier, but it does not prove \(CWQ_2\),
\(BSC_2\), the one-occurrence or directed/coupled center cases, the
\(R_1\) paths, or the full \(2+1+1\) source.  The \(3+1\), global
discrete, globally aggregated \([4]\), and \([3,3]\) families,
higher MTA, WUFC, CPM, cutoff removal, reflection positivity,
Osterwalder--Schrader reconstruction, construction of the continuum
Yang--Mills measure, and a uniform positive continuum spectral gap
remain open.  The full Yang--Mills mass gap has not been proved.
\end{firewall}

\paragraph{Parent-version record.}
V43 was copied byte-for-byte from
\texttt{Renewal\_Yang\_Mills\_Mass\_Gap\_Research\_Notes\_V42.tex}
before this chapter was appended.  The V42 parent had \(22\,163\)
source lines, \(803\,081\) bytes, and SHA--256
\[
 \texttt{14326aa33a7b6917e1a1bc8df2c9366da5844352be56decf41985a9489554fee}.
\]
The next compulsory companion audit is V45.

% =====================================================================
% V44 APPENDIX: DOUBLED BINARY-PREFIX TELESCOPE AND ONE-DEBIT CEILING
% =====================================================================

\chapter{V44: Doubled Binary-Prefix Telescope and the One-Debit Ceiling}
\label{ch:V44-doubled-prefix}

\section{Purpose and inheritance}
\label{sec:V44-purpose}

V43 reduced the balanced marked-heavy cell to a positive quadratic
estimate on the \emph{complete} word.  In the right-square
orientation, the exact carrier is
\[
 {\mathbb E}_{\boldsymbol\eta,\boldsymbol\xi}
 \left[
 R_\pi^*D(\boldsymbol\eta)^*L_\pi^*L_\pi
 D(\boldsymbol\xi)R_\pi
 \right].
\]
The binary covariance prefix is contained in both occurrences of
\(L_\pi\).  The next proposed move was therefore to telescope those
two copies simultaneously.  This chapter performs the resulting
finite algebra exactly and answers a narrower question:
\begin{quote}
Does estimating the doubled binary prefix, before exploiting its
interaction with the ternary replicas, supply the two center debits
required by \(CWQ_2\)?
\end{quote}
The answer is no.  The square is useful and uniformly bounded, but
the product-state Cauchy--Schwarz step takes its square root and
returns only the single response stock already present in V10.  This
identifies the first operation at which a prefix-only proof loses the
second center debit.

Throughout this chapter, \(xq\) is the binary response pair, \(t\) is
the first breaker coordinate, and
\[
 h_s:=h_{xq,s}=a_{x,s}a_{q,s},
 \qquad
 H:=H_{xq}(<t)=\sum_{s<t}h_s.
\]
All sums below are finite.  No order is imposed on probability
replicas; the order \(s<u\) is only the pre-existing order of lattice
coordinates.

\section{Exact doubled telescope}
\label{sec:V44-exact-square}

For \(s<t\), define
\begin{equation}
 U_s:=
 \left(\bigotimes_{e<s}T_e^{xq}\right)
 \otimes\Delta_s\otimes
 \left(\bigotimes_{s<e<t}R_e^{xq}\right),
 \qquad
 \Delta_s:=T_s^{xq}-R_s^{xq}.
\label{eq:V44-prefix-U}
\end{equation}
By \Cref{lem:V10-prefix-telescope},
\begin{equation}
 {\mathfrak J}_{<t}^{xq}=\sum_{s<t}U_s.
\label{eq:V44-J-sum}
\end{equation}

\begin{theorem}[Exact doubled binary-prefix telescope]
\label{thm:V44-double-telescope}
The right and left squares of the complete prefix covariance are
\begin{align}
 ({\mathfrak J}_{<t}^{xq})^*{\mathfrak J}_{<t}^{xq}
 &=
 \sum_{s<t}U_s^*U_s
 +\sum_{s<u<t}\bigl(U_s^*U_u+U_u^*U_s\bigr),
\label{eq:V44-double-telescope-right}\\
 {\mathfrak J}_{<t}^{xq}({\mathfrak J}_{<t}^{xq})^*
 &=
 \sum_{s<t}U_sU_s^*
 +\sum_{s<u<t}\bigl(U_sU_u^*+U_uU_s^*\bigr).
\label{eq:V44-double-telescope-left}
\end{align}
The whole left side of each identity is positive.  Each diagonal
summand is positive, while an individual paired off-diagonal
summand is self-adjoint but need not be positive.
\end{theorem}

\begin{proof}
Insert \eqref{eq:V44-J-sum} and distribute:
\[
 \left(\sum_sU_s\right)^*\left(\sum_uU_u\right)
 =\sum_{s,u}U_s^*U_u.
\]
Partition the ordered pairs into \(s=u\), \(s<u\), and \(u<s\).
In the last class interchange the dummy indices.  This gives
\eqref{eq:V44-double-telescope-right}.  The left-square identity is
identical.  An operator square is positive, as are \(U_s^*U_s\) and
\(U_sU_s^*\).  The adjoint of
\(U_s^*U_u+U_u^*U_s\) is itself; self-adjointness alone does not imply
positivity.
\end{proof}

\begin{corollary}[The coordinate ordering is legal]
\label{cor:V44-coordinate-order}
The decomposition in \Cref{thm:V44-double-telescope} does not use an
ordering of the independent replica variables
\(\boldsymbol\xi,\boldsymbol\eta\).  It only uses the fixed order on
the finite coordinate set \(\{e:e<t\}\).  Consequently it does not
invoke the invalid replica-ordering step excluded in V40.
\end{corollary}

\begin{proof}
The indices \(s,u\) label the local factors in the tensor products
\eqref{eq:V44-prefix-U}; the replica variables do not occur in their
definition.  The inequality \(s<u\) merely selects one representative
from each off-diagonal pair \(\{(s,u),(u,s)\}\).
\end{proof}

\begin{firewall}[Off-diagonal terms may not be discarded]
\label{fire:V44-offdiagonal}
Although the assembled expressions in
\eqref{eq:V44-double-telescope-right}--%
\eqref{eq:V44-double-telescope-left} are positive, the off-diagonal
pairs do not have a fixed sign.  Therefore positivity of the whole
square does not permit one to retain only \(\sum_sU_s^*U_s\), nor to
replace the cross terms by a favorable signed contribution.
\end{firewall}

\section{Support-uniform square bounds}
\label{sec:V44-square-bounds}

\begin{lemma}[Local telescope summands]
\label{lem:V44-local-U}
There is a constant \(C\), independent of \(t\) and of the number of
prefix coordinates, such that
\begin{equation}
 \|U_s\|\le C s_\alpha h_s
 \qquad(s<t).
\label{eq:V44-U-bound}
\end{equation}
If the sole output numerator and resolvent are assigned to the local
covariance at \(s\), the corresponding paid summand satisfies
\begin{equation}
 \|U_s^{\rm pay}\|\le C h_s.
\label{eq:V44-U-paid-bound}
\end{equation}
\end{lemma}

\begin{proof}
Every factor strictly to the left or right of \(\Delta_s\) in
\eqref{eq:V44-prefix-U} is a complete moment contraction by
\Cref{lem:V10-prefix-telescope}.  The local binary
composite-covariance estimate used in the proof of
\Cref{thm:V10-prefix-response} gives
\(\||\Delta_s|_{\rm phys}\|\le Cs_\alpha h_s\), and gives \(Ch_s\)
when that local occurrence pays.  Submultiplicativity proves both
claims.
\end{proof}

\begin{theorem}[The doubled prefix is support-uniform]
\label{thm:V44-prefix-square-bound}
Let \({\mathfrak J}:={\mathfrak J}_{<t}^{xq}\).  Then
\begin{align}
 \|{\mathfrak J}^*{\mathfrak J}\|,
 \ \kappa_\otimes({\mathfrak J}^*{\mathfrak J}),
 \ \|{\mathfrak J}{\mathfrak J}^*\|,
 \ \kappa_\otimes({\mathfrak J}{\mathfrak J}^*)
 &\le C\min\{1,s_\alpha H\}^2.
\label{eq:V44-prefix-square-unpaid}
\end{align}
For an exact paid prefix occurrence \({\mathfrak J}^{\rm pay}\) of
the kind covered by \eqref{eq:V10-prefix-paid},
\begin{align}
 \|({\mathfrak J}^{\rm pay})^*{\mathfrak J}^{\rm pay}\|,
 \ \kappa_\otimes(( {\mathfrak J}^{\rm pay})^*
                         {\mathfrak J}^{\rm pay}),
 \ \|{\mathfrak J}^{\rm pay}({\mathfrak J}^{\rm pay})^*\|,
 \ \kappa_\otimes({\mathfrak J}^{\rm pay}
                         ({\mathfrak J}^{\rm pay})^*)
 &\le C\min\{H,s_\alpha^{-1}\}^2.
\label{eq:V44-prefix-square-paid}
\end{align}
These bounds have no support-cardinality loss.
\end{theorem}

\begin{proof}
First, \Cref{lem:V44-local-U} and the triangle inequality give
\[
 \|{\mathfrak J}\|
 \le\sum_{s<t}\|U_s\|
 \le Cs_\alpha\sum_{s<t}h_s
 =Cs_\alpha H.
\]
The complete endpoints in
\({\mathfrak J}=T_{<t}^{xq}-R_{<t}^{xq}\) are contractions, so
\(\|{\mathfrak J}\|\le2\).  Absorbing fixed constants yields
\[
 \|{\mathfrak J}\|\le C\min\{1,s_\alpha H\}.
\]
The \(C^*\)-identity gives
\[
 \|{\mathfrak J}^*{\mathfrak J}\|
 =\|{\mathfrak J}{\mathfrak J}^*\|
 =\|{\mathfrak J}\|^2.
\]
Both squares are positive, and
\Cref{lem:V41-capacity-norm} bounds their product-state capacities by
their norms.  This proves \eqref{eq:V44-prefix-square-unpaid}.

For completeness, estimating the exact expansion directly gives
\[
 \sum_s\|U_s\|^2
 +2\sum_{s<u}\|U_s\|\|U_u\|
 =\left(\sum_s\|U_s\|\right)^2
 \le C s_\alpha^2H^2,
\]
so the doubled telescope itself introduces no coordinate count.
The paid assertion follows in the same way from the already proved
support-uniform estimate \eqref{eq:V10-prefix-paid} and the
\(C^*\)-identity.  Notice that this invokes one exact paid occurrence;
it does not duplicate the sole payer.
\end{proof}

\begin{corollary}[Positive envelopes obey the same scale]
\label{cor:V44-envelope-square}
If \({\cal N}_{<t}^{xq}\) and \({\cal P}_{<t}^{xq}\) are the positive
physical envelopes of \Cref{thm:V10-prefix-response}, then
\begin{align*}
 \|({\cal N}_{<t}^{xq})^2\|,
 \ \kappa_\otimes(( {\cal N}_{<t}^{xq})^2)
 &\le C\min\{1,s_\alpha H\}^2,\\
 \|({\cal P}_{<t}^{xq})^2\|,
 \ \kappa_\otimes(( {\cal P}_{<t}^{xq})^2)
 &\le C\min\{H,s_\alpha^{-1}\}^2.
\end{align*}
\end{corollary}

\begin{proof}
The envelopes are positive.  Square the norm bounds in
\eqref{eq:V10-prefix-unpaid}--\eqref{eq:V10-prefix-paid}, apply the
\(C^*\)-identity, and use
\Cref{lem:V41-capacity-norm}.
\end{proof}

\section{The one-debit ceiling}
\label{sec:V44-one-debit}

The square estimate above is not yet the desired complete-word
quadratic estimate.  The distinction is decisive: in the dressed
word the two prefix copies surround the ternary replica carriers,
whereas \Cref{thm:V44-prefix-square-bound} estimates the prefix block
before that interaction is used.

\begin{proposition}[Square-root recovery law]
\label{prop:V44-square-root-law}
Let \(X\) be any bounded operator and let \(p\ge0\).  If the only
quantitative input on \(X\) is
\begin{equation}
 \kappa_\otimes(X^*X)\le Cp^2,
\label{eq:V44-abstract-square}
\end{equation}
then product-state Cauchy--Schwarz yields at most
\begin{equation}
 \kappa_\otimes^{\rm abs}(X)\le C^{1/2}p.
\label{eq:V44-abstract-root}
\end{equation}
Thus forming and estimating a square doubles a response stock only
before the final square root; it does not create a second copy of
that stock in the resulting linear estimate.
\end{proposition}

\begin{proof}
Apply \Cref{lem:V43-complete-word-CS} and then
\eqref{eq:V44-abstract-square}:
\[
 \bigl(\kappa_\otimes^{\rm abs}(X)\bigr)^2
 \le\kappa_\otimes(X^*X)
 \le Cp^2.
\]
Both sides are nonnegative, so taking square roots gives
\eqref{eq:V44-abstract-root}.  No stronger power of \(p\) follows
from the displayed hypotheses: on the one-dimensional algebra with
its unique state, the scalar operator \(X=p\) has
\(\kappa_\otimes(X^*X)=p^2\) and
\(\kappa_\otimes^{\rm abs}(X)=p\), giving equality.
\end{proof}

\begin{theorem}[One-debit ceiling for an early prefix-only estimate]
\label{thm:V44-one-debit}
Suppose the two occurrences of the binary prefix in
\eqref{eq:V43-dressed-right} are bounded before any signed
interaction with
\(D(\boldsymbol\eta)^*\) and \(D(\boldsymbol\xi)\) is used.  Then the
unpaid prefix contributes, after the complete-word
Cauchy--Schwarz square root, at most the linear scale
\begin{equation}
 C\min\{1,s_\alpha H\},
\label{eq:V44-one-unpaid-debit}
\end{equation}
and a singly paid prefix contributes at most
\begin{equation}
 C\min\{H,s_\alpha^{-1}\}.
\label{eq:V44-one-paid-debit}
\end{equation}
In particular, the doubled prefix estimate alone cannot establish
the \(A_1(Q_{e,r,q})^{-2}\) branch of the linear \(BSC_2\) target, or
equivalently its fourth-power branch inside \(CWQ_2\), unless a
second inverse center mass has already been obtained from another
factor.
\end{theorem}

\begin{proof}
By \Cref{thm:V44-prefix-square-bound}, the separated prefix square
has precisely the squares of the scales in
\eqref{eq:V44-one-unpaid-debit}--%
\eqref{eq:V44-one-paid-debit}.  Applying
\Cref{prop:V44-square-root-law} recovers the corresponding linear
scale and no additional power.  The \(BSC_2\) center branch defined
in \Cref{def:V43-CWQ2} requires two inverse powers of
\(A_1(Q_{e,r,q})\) in the linear bound; its quadratic target therefore
requires four.  A prefix-only bound supplies only the stock encoded
by its one binary response occurrence.  Without an independently
proved second inverse-mass factor, the required exponent cannot be
deduced.
\end{proof}

\begin{firewall}[No debit from multiplicity alone]
\label{fire:V44-no-multiplicity-debit}
The occurrence of \(L_\pi\) twice in \(L_\pi^*L_\pi\) is forced by
the positive square.  It is not a license to charge the single
binary response twice.  Any proof of two center debits must display
two quantitative estimates before the final square root, or one
genuinely quadratic estimate whose scale already contains the fourth
power required by \(CWQ_2\).
\end{firewall}

\begin{firewall}[The paid square still has one payer]
\label{fire:V44-paid-square}
Equation \eqref{eq:V44-prefix-square-paid} concerns the square of one
already-paid operator occurrence.  The adjoint copy in that square
does not constitute a second output numerator or a second resolvent.
Interpreting it as a new payer would violate the one-payer ledger.
\end{firewall}

\section{Residual and next route}
\label{sec:V44-residual}

\begin{theorem}[First unresolved carrier after the prefix audit]
\label{thm:V44-residual}
The exact doubled-prefix expansion does not close \(CWQ_2\).  After
all results of V43 and this chapter, the first unresolved operation
is a \emph{joint} estimate on
\begin{equation}
 {\mathbb E}_{\boldsymbol\eta,\boldsymbol\xi}
 \left[
 R_\pi^*D(\boldsymbol\eta)^*L_\pi^*L_\pi
 D(\boldsymbol\xi)R_\pi
 \right]
\label{eq:V44-joint-carrier}
\end{equation}
before the binary prefix, the two ternary replica packets, and their
signed cross terms are separated.  A successful continuation must
obtain the second center debit from at least one of:
\begin{enumerate}[label=\textup{(\roman*)},leftmargin=4em]
\item a binary--ternary cross estimate inside
      \eqref{eq:V44-joint-carrier};
\item a complete-bank quadratic moment carrying two center response
      factors before Cauchy--Schwarz; or
\item an exact signed cancellation in the assembled square.
\end{enumerate}
No one of these three mechanisms is proved here.
\end{theorem}

\begin{proof}
The sufficiency of \(CWQ_2\) is
\Cref{thm:V43-CWQ2-suffices}.  The exact carrier is
\Cref{thm:V43-dressed-square}.  By
\Cref{thm:V44-one-debit}, separating and estimating the prefix first
cannot provide both required center debits.  Hence any proof that
continues through this carrier must retain a joint interaction or
produce the missing estimate elsewhere.  The three listed mechanisms
exhaust those possibilities at the level of the current factorization:
the second factor is obtained from a cross interaction, is already
present in a stronger quadratic bank estimate, or arises by signed
cancellation before absolute values.  This is a classification of
routes, not a proof that any route succeeds.
\end{proof}

\begin{assessment}[The useful outcome of the failed route]
\label{ass:V44-outcome}
The doubled telescope is not wasted: it proves that coordinate
cardinality is not the obstruction and that all off-diagonal terms
can be retained in a legal coordinate ordering.  The obstruction is
instead the exponent after the inevitable square root.  Consequently
V45 should not repeat a prefix norm or capacity estimate.  It should
audit the companion Standard-Model and Navier--Stokes notes as
required, then seek an upstream two-mark or joint-quadratic identity
that applies to \eqref{eq:V44-joint-carrier} without duplicating the
single payer.
\end{assessment}

\begin{theorem}[Third-series V44 result ledger]
\label{thm:V44-results}
Relative to V1--V43, V44 proves:
\begin{enumerate}[label=\textup{(V44.\arabic*)},leftmargin=6.3em]
\item the exact right and left doubled binary-prefix telescopes;
\item the legality of ordering coordinate labels while leaving
      probability replicas unordered;
\item support-uniform unpaid and paid prefix-square bounds;
\item the abstract square-root recovery law, with a sharp scalar
      example;
\item the one-debit ceiling for every prefix-first implementation of
      the V43 complete-word route;
\item the no-off-diagonal-discard, no-multiplicity-debit, and
      one-payer firewalls; and
\item reduction of the next step to a genuinely joint
      binary--ternary quadratic carrier.
\end{enumerate}
\end{theorem}

\begin{proof}
These statements are
\Cref{thm:V44-double-telescope,
cor:V44-coordinate-order,
thm:V44-prefix-square-bound,
prop:V44-square-root-law,
thm:V44-one-debit,
fire:V44-offdiagonal,
fire:V44-no-multiplicity-debit,
fire:V44-paid-square,
thm:V44-residual}.
\end{proof}

\begin{programme}[V45: compulsory companion audit and joint carrier]
\label{prog:V44-V45}
\begin{enumerate}[leftmargin=3em]
\item Inspect the newest Standard-Model and Navier--Stokes notes in
      \texttt{latex\_notes}, recording exact filenames, hashes, and
      the specific newest results examined.
\item Apply the type firewall: import only an algebraic or analytic
      lemma whose hypotheses and object types are verified in the
      Yang--Mills carrier.
\item Compare any companion two-mark or quadratic-response mechanism
      with \eqref{eq:V44-joint-carrier}.
\item If no valid transfer exists, state that result and work directly
      upstream from the earliest joint binary--ternary cross term.
\item Do not estimate the binary prefix separately again; V44 proves
      that such a route has a one-debit ceiling.
\item Preserve the local-payer polarization boundary and the V40
      replica-ordering correction.
\end{enumerate}
\end{programme}

\begin{firewall}[Claim boundary after third-series V44]
\label{fire:V44-claim-boundary}
V44 proves exact doubled-prefix identities and a no-go theorem for a
prefix-first estimate.  It does not prove the joint carrier estimate
\eqref{eq:V44-joint-carrier}, \(CWQ_2\), \(BSC_2\), the
one-occurrence or directed/coupled center cases, the \(R_1\) paths,
or the full \(2+1+1\) source.  The \(3+1\), global discrete,
globally aggregated \([4]\), and \([3,3]\) families, higher MTA,
WUFC, CPM, cutoff removal, reflection positivity,
Osterwalder--Schrader reconstruction, construction of the continuum
Yang--Mills measure, and a uniform positive continuum spectral gap
remain open.  The full Yang--Mills mass gap has not been proved.
\end{firewall}

\paragraph{Parent-version record.}
V44 was copied byte-for-byte from
\texttt{Renewal\_Yang\_Mills\_Mass\_Gap\_Research\_Notes\_V43.tex}
before this chapter was appended.  The V43 parent had \(22\,552\)
source lines, \(817\,634\) bytes, and SHA--256
\[
 \texttt{53e84071b073870f26433f45918ebf6ec2b4832514c1d8fd3e73c3b4ccdeefc1}.
\]
The next compulsory companion audit is V45.

% =====================================================================
% V45 APPENDIX: COMPANION AUDIT AND AGGREGATED PREFIX TARGET
% =====================================================================

\chapter{V45: Companion Audit and the Aggregated Marked-Heavy Target}
\label{ch:V45-aggregate}

\section{Compulsory companion audit}
\label{sec:V45-audit}

This is the compulsory five-version audit requested for the parallel
programmes.  At the time of inspection, the newest non-frozen
companion files in \texttt{latex\_notes} were
\begin{align*}
 &\texttt{Renewal\_SM\_Einstein\_Continuum\_Research\_Notes\_V23.tex},\\
 &\texttt{Renewal\_NS\_Stability\_Research\_Notes\_V31.tex}.
\end{align*}
The SM file had \(8\,238\) source lines, \(321\,056\) bytes, and
SHA--256
\[
 \texttt{1e761d967ca725c164cd2808d30fd2995b883d169eac75aa521eb03f55b10ac5}.
\]
The NS file had \(23\,052\) source lines, \(887\,537\) bytes, and
SHA--256
\[
 \texttt{f1121b62238ad5491bb7cf03635ea9934942edf573ed8faaa9ee79e1d38a6696}.
\]

The SM programme has advanced from the V17 file inspected at V40 to
V23.  The pieces examined here were its V20 two-background-mark
theorem, V22 first-failure ownership, and V23 relative-Hodge,
Schur-harmonic-extension, and nonlinear plaquette-chart results.  The
NS file is unchanged from the V40 audit; its newest V31 results were
rechecked, in particular the exact heat-formation ladder, its
forced-mode sharpness test, the high-parent circularity firewall, and
the balanced-sector nullspace of the signed hinge.

\begin{theorem}[V45 companion transfer decision]
\label{thm:V45-audit-decision}
Neither active companion file supplies a Yang--Mills estimate that
can be imported into \eqref{eq:V44-joint-carrier}.  They do, however,
support the following proof-order decision:
\begin{enumerate}[label=\textup{(\roman*)},leftmargin=4em]
\item two debits must be attached to two actual occurrences in one
      assembled object, as in the SM V20 two-mark theorem, and not to
      two copies of a one-mark upper bound;
\item a scale or mass factor lost by an early positive estimate is not
      restored by the subsequent linear propagator or square root, as
      shown sharply by the NS V31 formation tests; and
\item signed cancellation may be used only on the branch on which its
      nullspace has been excluded, as the NS balanced-hinge example
      demonstrates.
\end{enumerate}
For the present YM sector, these principles require retaining the
complete sum over actual prefix marks before taking product-state
capacity.
\end{theorem}

\begin{proof}
The SM V20 estimate assumes differentiated polymer factors bounded by
one common cutoff-transition event and uses an independent subset of
those events.  The YM Wilson carrier has neither those scalar events
nor the finite-dependence graph in the hypotheses, so that estimate
does not transfer.  The SM V23 Hodge result is a finite-dimensional
cochain estimate under a complete obstruction map and an elliptic
comparison; it provides no product-state operator-capacity estimate.

The NS V31 heat theorem acts on Fourier spectral shells and its sharp
examples concern Duhamel smoothing.  Its hinge theorem concerns a
helical sector difference and has an explicit positive-total-sector
nullspace.  These object types do not match Wilson tensor moments.
Thus none of the companion inequalities applies.  The three stated
proof-order rules are precisely the logically type-free content of
their actual-occurrence, sharpness, and nullspace firewalls.
\end{proof}

\begin{firewall}[Companion type boundary at V45]
\label{fire:V45-companion-types}
No SM KP estimate, Gaussian tail, Schur-complement estimate, relative
Hodge constant, or BCH bound is used as a Wilson estimate.  No NS
heat-formation, Leray, coarea, spectral-tail, or hinge estimate is
used as a Wilson estimate.  Only the verified proof-order lessons in
\Cref{thm:V45-audit-decision} affect the direction of the YM
programme.
\end{firewall}

\section{Scope correction: fixed marks versus the assembled prefix}
\label{sec:V45-scope}

The V36 definition of \(BSC_2\) freezes one ordered pair \(e<r\).
Thus \(\mathcal W_{e,r,q}^{\rm pay}\) contains the selected
first-connectivity occurrence at \(e\), with complete spectators, but
does not contain the sum over all possible selected coordinates.
By contrast, the exact V10 prefix covariance and its V44 square are
assembled sums over those coordinates.

\begin{proposition}[Precise scope of the V44 doubled telescope]
\label{prop:V45-V44-scope}
Let \(U_e\) be the V44 telescope summand.  Then:
\begin{enumerate}[label=\textup{(\alph*)},leftmargin=4em]
\item the diagonal estimate for \(U_e^*U_e\), followed by a square
      root, has the one-debit scale \(Cs_\alpha h_{12,e}\);
\item the off-diagonal terms \(U_e^*U_f\), \(e\ne f\), belong to the
      square of the assembled prefix \(\sum_eU_e\); and
\item those off-diagonal terms cannot be used to prove the stronger
      fixed-\(e\) card \(BSC_2\) without first proving a separate
      identity that puts them into that fixed-\(e\) word.
\end{enumerate}
The V44 one-debit conclusion remains valid.  Its doubled-coordinate
cross terms must, however, be assigned to an aggregate target rather
than to an individual \(BSC_2\) summand.
\end{proposition}

\begin{proof}
Item \textup{(a)} follows from
\eqref{eq:V44-U-bound}, the \(C^*\)-identity, and
\Cref{prop:V44-square-root-law}.  Item \textup{(b)} is the exact
identity \eqref{eq:V44-double-telescope-right}.  For item
\textup{(c)}, the definition
\eqref{eq:V36-marked-mass}--\eqref{eq:V36-BSC2} fixes \(e\), while
\(U_e^*U_f\) contains the distinct coordinate label \(f\).  It is
therefore absent from \(U_e^*U_e\).  Adding it would change the
operator whose capacity is being estimated.  The final statement
only changes the scope assigned to the V44 cross terms; its diagonal
one-debit calculation is unchanged.
\end{proof}

\begin{lemma}[Cancellation of an aggregate does not prove its
summand bounds]
\label{lem:V45-aggregate-not-individual}
For a family of signed bounded operators \(X_e\), a bound on
\(\kappa_\otimes^{\rm abs}(\sum_eX_e)\) does not imply corresponding
bounds on the individual \(\kappa_\otimes^{\rm abs}(X_e)\).
Consequently cross-coordinate cancellation is available only after
the source target itself has been formulated for the assembled sum.
\end{lemma}

\begin{proof}
On the scalar algebra take \(X_1=M\), \(X_2=-M\).  The aggregate is
zero for every \(M\), whereas each absolute capacity equals \(|M|\).
Thus no individual estimate with a constant controlled by the
aggregate can follow.  This scalar example proves the logical claim
without making any assertion that the physical YM words cancel in
that particular way.
\end{proof}

\begin{firewall}[No retrospective use of cross-prefix cancellation]
\label{fire:V45-no-retrospective-cross}
The off-diagonal V44 terms may not be inserted retrospectively into
the fixed-\(e\) definition of \(CWQ_2(\pi)\).  They become legitimate
only after replacing the sufficient interface by an aggregate source
estimate whose left side really contains
\(\sum_e\mathcal W_{e,r,q}^{\rm pay,\pi}\).
\end{firewall}

\section{An aggregate interface sufficient for the source}
\label{sec:V45-aggregate-interface}

Fix the joining coordinate \(r\), the remaining external labels
(including \(q\) when it is distinct), and one payer placement
\(\pi\).  Let \(I_{r,q,\pi}^{\rm mh}\) be the set of prefix marks
\(e<r\) for which the packet is balanced, protected, and
marked-heavy:
\[
 A_1(Q_{e,r,q})\ge\varepsilon\Xi_1.
\]
Put
\begin{align}
 H_{12}(I_{r,q,\pi}^{\rm mh})
 &:=\sum_{e\in I_{r,q,\pi}^{\rm mh}}h_{12,e},
\label{eq:V45-HI}\\
 \mathbf W_{r,q,\pi}^{\rm mh}
 &:=\sum_{e\in I_{r,q,\pi}^{\rm mh}}
       \mathcal W_{e,r,q}^{\rm pay,\pi}.
\label{eq:V45-Wagg}
\end{align}
The sum is taken in the original signed source before a positive
envelope.

\begin{definition}[Aggregated balanced star-collision card]
\label{def:V45-ABSC2}
The aggregate card \(ABSC_2(r,q,\pi)\) is
\begin{equation}
 \boxed{
 \kappa_\otimes^{\rm abs}(\mathbf W_{r,q,\pi}^{\rm mh})
 \le C_\varepsilon
 \min\{s_\alpha^2,\Xi_1^{-2}\}
 H_{12}(I_{r,q,\pi}^{\rm mh})
 h_{13,r}h_{14,r}.}
\label{eq:V45-ABSC2}
\end{equation}
No individual prefix-mark estimate is included in this definition.
\end{definition}

\begin{theorem}[The aggregate card is sufficient at the prefix-sum
interface]
\label{thm:V45-ABSC2-suffices}
For fixed \((r,q,\pi)\), \(ABSC_2(r,q,\pi)\) implies
\begin{equation}
 \frac{
 \kappa_\otimes^{\rm abs}(\mathbf W_{r,q,\pi}^{\rm mh})}
 {\Xi_1\Xi_2\Xi_3\Xi_4}
 \le
 C_\varepsilon
 \frac{
 H_{12}(I_{r,q,\pi}^{\rm mh})h_{13,r}h_{14,r}}
 {\Xi_1^3\Xi_2\Xi_3\Xi_4}.
\label{eq:V45-aggregate-tree}
\end{equation}
This is exactly the row-\(1\) double-center tree normalization for
the marked-heavy prefix sum.  Hence, after the already established
summations of the remaining external labels and finite payer types,
the aggregate card is sufficient wherever the earlier proof used the
sum of the fixed-\(e\) \(BSC_2\) cards.
\end{theorem}

\begin{proof}
Use the inverse-mass branch of \eqref{eq:V45-ABSC2} and divide by the
four input masses.  This gives \eqref{eq:V45-aggregate-tree}.  The
power \(\Xi_1^{-3}\) consists of the input normalization and the two
center debits; the other row powers are unchanged.  By
\eqref{eq:V45-HI}, the numerator is the exact sum of the binary tree
stocks over the marked-heavy prefix marks.  The prior source
summation is linear before capacity is taken, and the remaining
external-label and payer ledgers do not require a separate absolute
value for each \(e\).  Thus the aggregate estimate can replace the
stronger sequence of proving individual \(BSC_2\) cards and then
using the triangle inequality at precisely this interface.
\end{proof}

\begin{proposition}[The fixed-mark absolute card implies the
aggregate card]
\label{prop:V45-fixed-implies-aggregate}
If the positive-majorant form of \eqref{eq:V36-BSC2} used in V43,
with
\(\kappa_\otimes^{\rm abs}
(\mathcal W_{e,r,q}^{\rm pay,\pi})\) on the left, holds uniformly for
every \(e\in I_{r,q,\pi}^{\rm mh}\), then
\eqref{eq:V45-ABSC2} holds.  The converse is false as a matter of
operator inequalities.  Thus \(ABSC_2\) is a genuine relaxation of
the old sufficient interface.
\end{proposition}

\begin{proof}
The triangle inequality and marked heaviness give
\begin{align*}
 \kappa_\otimes^{\rm abs}(\mathbf W_{r,q,\pi}^{\rm mh})
 &\le\sum_{e\in I_{r,q,\pi}^{\rm mh}}
       \kappa_\otimes^{\rm abs}
       (\mathcal W_{e,r,q}^{\rm pay,\pi})\\
 &\le C\sum_e
 \min\{s_\alpha^2,A_1(Q_{e,r,q})^{-2}\}
 h_{12,e}h_{13,r}h_{14,r}\\
 &\le C_\varepsilon
 \min\{s_\alpha^2,\Xi_1^{-2}\}
 H_{12}(I_{r,q,\pi}^{\rm mh})h_{13,r}h_{14,r}.
\end{align*}
For the last inequality, use
\(A_1(Q_{e,r,q})^{-2}\le\varepsilon^{-2}\Xi_1^{-2}\)
and absorb \(\varepsilon^{-2}\).  Failure of the converse follows
from \Cref{lem:V45-aggregate-not-individual}, with equal positive
weights assigned to the two scalar summands.
\end{proof}

\begin{firewall}[Aggregation does not change payer ownership]
\label{fire:V45-aggregate-payer}
The aggregate in \eqref{eq:V45-Wagg} is formed only after fixing the
payer placement type \(\pi\) and all non-prefix labels.  It neither
sums two payer choices inside one word nor promotes the adjoint copy
in a square to a second payer.  Local-ternary payer words still
require the separate paid polarization identified in
\Cref{prop:V43-local-payer-boundary}.
\end{firewall}

\section{The exact aggregate quadratic carrier}
\label{sec:V45-aggregate-square}

Consider first an outside-ternary payer placement.  For each
\(e\in I_{r,q,\pi}^{\rm mh}\), preserve the exact factorization
\begin{equation}
 \mathcal W_{e,r,q}^{\rm pay,\pi}=L_{e,\pi}S_rR_{e,\pi},
 \qquad
 S_r=\kappa_3(A_r,B_r,C_r).
\label{eq:V45-mark-factor}
\end{equation}
The factors are allowed to depend on \(e\); in particular, the
selected prefix occurrence and its complete spectators are part of
\(L_{e,\pi}\).

\begin{theorem}[Exact aggregate marked-prefix square]
\label{thm:V45-aggregate-square}
Let \(D_r(\boldsymbol\xi)\) be the V40 four-replica polarization of
\(S_r\), and abbreviate \(I=I_{r,q,\pi}^{\rm mh}\).  Then
\begin{align}
 (\mathbf W_{r,q,\pi}^{\rm mh})^*
 \mathbf W_{r,q,\pi}^{\rm mh}
 ={}&
 \sum_{e,f\in I}
 {\mathbb E}_{\boldsymbol\eta,\boldsymbol\xi}
 \left[
 R_{e,\pi}^*D_r(\boldsymbol\eta)^*
 L_{e,\pi}^*L_{f,\pi}
 D_r(\boldsymbol\xi)R_{f,\pi}
 \right].
\label{eq:V45-aggregate-square}
\end{align}
Equivalently, it is the sum of the diagonal \(e=f\) carriers and the
paired cross-coordinate carriers with \(e<f\).  The complete sum is
positive.  A diagonal carrier or a paired cross carrier need not be
positive after the replica expectation is expanded.
\end{theorem}

\begin{proof}
By \eqref{eq:V45-Wagg} and \eqref{eq:V45-mark-factor},
\[
 \mathbf W_{r,q,\pi}^{\rm mh}
 =\sum_{e\in I}L_{e,\pi}S_rR_{e,\pi}.
\]
Multiply the finite sum by its adjoint.  In each \((e,f)\) term,
insert
\(S_r={\mathbb E}_{\boldsymbol\xi}D_r(\boldsymbol\xi)\) and its
independent adjoint copy, then apply Fubini.  This proves
\eqref{eq:V45-aggregate-square}.  Splitting the finite ordered-pair
sum into \(e=f\), \(e<f\), and \(f<e\) gives the equivalent paired
form.  Positivity belongs to the left side, which is an operator
square; the individual kernels have no asserted sign.
\end{proof}

\begin{definition}[Aggregate complete-word quadratic target]
\label{def:V45-ACWQ2}
The right-oriented \(ACWQ_2(r,q,\pi)\) target is
\begin{equation}
 \boxed{
 \kappa_\otimes\left(
 (\mathbf W_{r,q,\pi}^{\rm mh})^*
 \mathbf W_{r,q,\pi}^{\rm mh}
 \right)
 \le C_\varepsilon
 \left[
 \min\{s_\alpha^2,\Xi_1^{-2}\}
 H_{12}(I_{r,q,\pi}^{\rm mh})
 h_{13,r}h_{14,r}
 \right]^2.}
\label{eq:V45-ACWQ2}
\end{equation}
The left-oriented target uses the opposite square.
\end{definition}

\begin{corollary}[The aggregate quadratic target is sufficient]
\label{cor:V45-ACWQ2-suffices}
The output-oriented version of \(ACWQ_2(r,q,\pi)\) implies
\(ABSC_2(r,q,\pi)\), and hence the marked-heavy prefix-sum estimate
\eqref{eq:V45-aggregate-tree}.
\end{corollary}

\begin{proof}
Apply the complete-word product-state Cauchy--Schwarz inequality
\eqref{eq:V43-right-CS} or \eqref{eq:V43-left-CS} to the assembled
operator \(\mathbf W_{r,q,\pi}^{\rm mh}\), then take the nonnegative
square root of \eqref{eq:V45-ACWQ2}.  This is
\eqref{eq:V45-ABSC2}; apply
\Cref{thm:V45-ABSC2-suffices}.
\end{proof}

\begin{proposition}[Why the aggregate carrier is the first usable
doubled-prefix object]
\label{prop:V45-first-usable-cross}
The cross terms with \(e\ne f\) in
\eqref{eq:V45-aggregate-square} are actual occurrences in the square
of the exact signed source aggregate.  They are therefore available
for cancellation in \(ACWQ_2\).  They were unavailable in the
fixed-mark \(CWQ_2(e,r,q,\pi)\) target.  Estimating them term by term
in absolute capacity reduces \(ACWQ_2\) to the same one-debit
prefix-first route excluded by V44.
\end{proposition}

\begin{proof}
The first statement is the algebraic identity in
\Cref{thm:V45-aggregate-square}; no term has been added to the
operator square.  The second is
\Cref{prop:V45-V44-scope}.  If every \((e,f)\) kernel is replaced by
an absolute upper bound, the signs of the paired cross terms are
lost.  The resulting double sum is bounded by the product of the two
separate prefix \(\ell^1\) response sums, exactly the estimate in the
proof of \Cref{thm:V44-prefix-square-bound}.  Its final square root
is limited by \Cref{thm:V44-one-debit}.
\end{proof}

\begin{firewall}[The aggregate target is sufficient, not proved]
\label{fire:V45-aggregate-open}
Equations \eqref{eq:V45-aggregate-square} and
\eqref{eq:V45-ACWQ2} identify a legitimate place for cross-prefix
cancellation.  They do not prove that this cancellation supplies the
second inverse power of \(\Xi_1\).  In particular, the positivity of
the whole square cannot be used to choose favorable signs for its
off-diagonal kernels.
\end{firewall}

\section{V45 residual and acquisition order}
\label{sec:V45-residual}

\begin{theorem}[Exact V45 reduction of the balanced marked-heavy
prefix sector]
\label{thm:V45-residual}
For each fixed set of non-prefix labels and each outside-ternary payer
placement, the old family of per-mark \(CWQ_2\) estimates may be
replaced by the single weaker target \(ACWQ_2\) on the exact aggregate
carrier \eqref{eq:V45-aggregate-square}.  A proof of that target
closes the corresponding marked-heavy prefix sum.  The unresolved
analytic question is now whether the assembled cross-coordinate
matrix
\[
 \left(
 {\mathbb E}
 [R_{e,\pi}^*D_r(\boldsymbol\eta)^*
  L_{e,\pi}^*L_{f,\pi}D_r(\boldsymbol\xi)R_{f,\pi}]
 \right)_{e,f\in I}
\]
has a signed two-center estimate stronger than its entrywise absolute
majorant.  No such estimate is proved in V45.
\end{theorem}

\begin{proof}
The replacement is sufficient by
\Cref{cor:V45-ACWQ2-suffices,thm:V45-ABSC2-suffices} and is weaker by
\Cref{prop:V45-fixed-implies-aggregate}.  The exact matrix entries
are those of \eqref{eq:V45-aggregate-square}.  By
\Cref{prop:V45-first-usable-cross}, entrywise absolute bounds lose the
only newly available structure.  Hence the stated signed matrix
estimate is precisely the remaining analytic question at this
interface.
\end{proof}

\begin{theorem}[Third-series V45 result ledger]
\label{thm:V45-results}
Relative to V1--V44, V45 proves:
\begin{enumerate}[label=\textup{(V45.\arabic*)},leftmargin=6.3em]
\item the compulsory SM V23 and NS V31 audit with an explicit
      no-transfer decision;
\item the fixed-mark/assembled-prefix scope correction for the V44
      doubled telescope;
\item a scalar proof that aggregate cancellation cannot establish
      individual marked-word bounds;
\item the weaker aggregate \(ABSC_2\) interface and its sufficiency
      for the normalized marked-heavy prefix sum;
\item proof that the original fixed-\(e\) \(BSC_2\) family implies,
      but is not implied by, the aggregate interface;
\item the exact aggregate eight-replica square, including all legal
      cross-coordinate terms;
\item the sufficient aggregate quadratic target \(ACWQ_2\); and
\item reduction of the next analytic step to a signed
      cross-coordinate operator-matrix estimate.
\end{enumerate}
\end{theorem}

\begin{proof}
These are
\Cref{thm:V45-audit-decision,
prop:V45-V44-scope,
lem:V45-aggregate-not-individual,
def:V45-ABSC2,
thm:V45-ABSC2-suffices,
prop:V45-fixed-implies-aggregate,
thm:V45-aggregate-square,
def:V45-ACWQ2,
cor:V45-ACWQ2-suffices,
thm:V45-residual}.
\end{proof}

\begin{programme}[V46: signed prefix-matrix structure]
\label{prog:V45-V46}
\begin{enumerate}[leftmargin=3em]
\item Derive the exact dependence of \(L_{e,\pi}^*L_{f,\pi}\) on the
      first-connectivity positions for \(e<f\); do not replace either
      side by its norm.
\item Partition only by the legal coordinate order \(e<f\), retaining
      the common ternary replicas and all complete spectators.
\item Test whether the sum of diagonal and paired off-diagonal
      kernels is a discrete second difference or a Gram form of a
      centered prefix martingale.
\item If such a representation exists, prove its exact boundary
      terms and determine whether its quadratic variation contains
      two actual row-\(1\) response occurrences.
\item If it does not exist, construct a scalar or \(SU(2)\) packet
      showing that \(ACWQ_2\) cannot follow from the current local
      inputs, and return upstream to the complete bank rather than
      reimposing per-mark \(BSC_2\).
\item Keep local-ternary payer placements separate until their paid
      polarization is written.
\item Do not perform another companion audit before V50.
\end{enumerate}
\end{programme}

\begin{firewall}[Claim boundary after third-series V45]
\label{fire:V45-claim-boundary}
V45 replaces an unnecessarily strong per-prefix interface by an
exact aggregate target, but it does not prove \(ACWQ_2\),
\(ABSC_2\), the complete balanced marked-heavy contribution, the
local-payer polarization, the one-occurrence or directed/coupled
center cases, the \(R_1\) paths, or the full \(2+1+1\) source.  The
\(3+1\), global discrete, globally aggregated \([4]\), and
\([3,3]\) families, higher MTA, WUFC, CPM, cutoff removal,
reflection positivity, Osterwalder--Schrader reconstruction,
construction of the continuum Yang--Mills measure, and a uniform
positive continuum spectral gap remain open.  The full Yang--Mills
mass gap has not been proved.
\end{firewall}

\paragraph{Parent-version record.}
V45 was copied byte-for-byte from
\texttt{Renewal\_Yang\_Mills\_Mass\_Gap\_Research\_Notes\_V44.tex}
before this chapter was appended.  The corrected V44 parent had
\(23\,026\) source lines, \(835\,475\) bytes, and SHA--256
\[
 \texttt{8fe21478a9d51095b034eeae6512cc434e707b60a91ca586b15b3570eb9c97c0}.
\]
The next compulsory companion audit is V50.

% =====================================================================
% V46 APPENDIX: PREFIX PATH, CROSS TERMS, AND ABSTRACT OBSTRUCTION
% =====================================================================

\chapter{V46: The Prefix Path and the Failure of Automatic
Cross-Coordinate Cancellation}
\label{ch:V46-prefix-path}

\section{The exact interpolation path}
\label{sec:V46-path}

Fix a finite ordered prefix coordinate set
\[
 E=\{e_1<e_2<\cdots<e_n\}.
\]
For brevity write \(T_j=T_{e_j}^{xq}\),
\(R_j=R_{e_j}^{xq}\), and
\(\Delta_j=T_j-R_j\).  Define the interpolation path
\begin{equation}
 P_k:=
 \left(\bigotimes_{j\le k}T_j\right)
 \otimes
 \left(\bigotimes_{j>k}R_j\right),
 \qquad 0\le k\le n,
\label{eq:V46-Pk}
\end{equation}
with empty tensor products interpreted as identities.  Thus
\(P_0=\bigotimes_jR_j\) and \(P_n=\bigotimes_jT_j\).

\begin{lemma}[The V10 telescope summands are path increments]
\label{lem:V46-increments}
For \(1\le k\le n\),
\begin{equation}
 U_k:=P_k-P_{k-1}
 =
 \left(\bigotimes_{j<k}T_j\right)
 \otimes\Delta_k\otimes
 \left(\bigotimes_{j>k}R_j\right).
\label{eq:V46-Uk}
\end{equation}
Consequently
\begin{equation}
 {\mathfrak J}_{<t}^{xq}=P_n-P_0=\sum_{k=1}^nU_k.
\label{eq:V46-J-path}
\end{equation}
\end{lemma}

\begin{proof}
The two products \(P_k\) and \(P_{k-1}\) agree at every coordinate
except \(e_k\).  At that coordinate their difference is
\(T_k-R_k=\Delta_k\), proving \eqref{eq:V46-Uk}.  Summing the
increments telescopes to \(P_n-P_0\), which is the V10 complete
prefix covariance.
\end{proof}

\begin{theorem}[Coordinatewise normal form of every right cross term]
\label{thm:V46-cross-normal-form}
For \(k<\ell\), the right cross term in the doubled prefix is exactly
\begin{align}
 U_k^*U_\ell
 ={}&
 \left(\bigotimes_{j<k}T_j^*T_j\right)
 \otimes(\Delta_k^*T_k)
 \otimes
 \left(\bigotimes_{k<j<\ell}R_j^*T_j\right)
 \notag\\
 &\otimes(R_\ell^*\Delta_\ell)
 \otimes
 \left(\bigotimes_{j>\ell}R_j^*R_j\right).
\label{eq:V46-right-cross}
\end{align}
The diagonal is
\begin{equation}
 U_k^*U_k
 =
 \left(\bigotimes_{j<k}T_j^*T_j\right)
 \otimes(\Delta_k^*\Delta_k)
 \otimes
 \left(\bigotimes_{j>k}R_j^*R_j\right).
\label{eq:V46-right-diagonal}
\end{equation}
The left cross term has the exact form
\begin{align}
 U_kU_\ell^*
 ={}&
 \left(\bigotimes_{j<k}T_jT_j^*\right)
 \otimes(\Delta_kT_k^*)
 \otimes
 \left(\bigotimes_{k<j<\ell}R_jT_j^*\right)
 \notag\\
 &\otimes(R_\ell\Delta_\ell^*)
 \otimes
 \left(\bigotimes_{j>\ell}R_jR_j^*\right).
\label{eq:V46-left-cross}
\end{align}
In particular, an off-diagonal term has two actual local covariance
occurrences, one at \(e_k\) and one at \(e_\ell\), separated by
complete mixed moments.
\end{theorem}

\begin{proof}
At a coordinate \(j<k\), both \(U_k\) and \(U_\ell\) contain \(T_j\);
at \(j=k\), they contain \(\Delta_k\) and \(T_k\); between \(k\) and
\(\ell\), they contain \(R_j\) and \(T_j\); at \(j=\ell\), they
contain \(R_\ell\) and \(\Delta_\ell\); and after \(\ell\), both
contain \(R_j\).  Multiplying the first word adjoint-first gives
\eqref{eq:V46-right-cross}.  The same coordinatewise multiplication
with \(k=\ell\) gives \eqref{eq:V46-right-diagonal}.  Reversing the
adjoint orientation gives \eqref{eq:V46-left-cross}.  The two
difference factors in \eqref{eq:V46-right-cross} are visibly
\(\Delta_k^*\) and \(\Delta_\ell\).
\end{proof}

\begin{corollary}[There is no hidden coordinate multiplicity]
\label{cor:V46-no-coordinate-multiplicity}
Using contractivity of \(T_j,R_j\) and the local response bounds,
\begin{align}
 \|U_k^*U_\ell\|
 &\le C s_\alpha^2h_{12,e_k}h_{12,e_\ell}
 \qquad(k\ne\ell),
\label{eq:V46-cross-bound}\\
 \|U_k^*U_k\|
 &\le C s_\alpha^2h_{12,e_k}^2.
\label{eq:V46-diagonal-bound}
\end{align}
Summing these entrywise bounds gives \(Cs_\alpha^2H_{12}(E)^2\),
with no additional factor depending on \(|E|\).
\end{corollary}

\begin{proof}
All factors in
\eqref{eq:V46-right-cross}--\eqref{eq:V46-right-diagonal} other than
the displayed local differences have norm at most one.  Apply the
local bounds
\(\|\Delta_k\|\le Cs_\alpha h_{12,e_k}\) and
\(\|\Delta_\ell\|\le Cs_\alpha h_{12,e_\ell}\).
Finally,
\[
 \sum_k h_{12,e_k}^2+
 2\sum_{k<\ell}h_{12,e_k}h_{12,e_\ell}
 =
 H_{12}(E)^2.
\]
\end{proof}

\section{The increments are not martingale differences}
\label{sec:V46-no-martingale}

Put \(Q_k:=P_k-P_0=\sum_{j\le k}U_j\), so \(Q_0=0\) and
\(Q_n={\mathfrak J}_{<t}^{xq}\).

\begin{lemma}[Exact path-energy identity]
\label{lem:V46-path-energy}
\begin{equation}
 ({\mathfrak J}_{<t}^{xq})^*{\mathfrak J}_{<t}^{xq}
 =
 \sum_{k=1}^n
 \left(
 U_k^*U_k+Q_{k-1}^*U_k+U_k^*Q_{k-1}
 \right).
\label{eq:V46-path-energy}
\end{equation}
One sufficient condition for the off-diagonal terms to vanish
term by term under a positive functional \(\varphi\) is the additional
orthogonality family
\begin{equation}
 \operatorname{Re}\varphi(Q_{k-1}^*U_k)=0
 \qquad(1\le k\le n).
\label{eq:V46-orthogonality-needed}
\end{equation}
These conditions do not follow from contractivity or from the local
response estimates.
\end{lemma}

\begin{proof}
Since \(Q_k=Q_{k-1}+U_k\),
\[
 Q_k^*Q_k-Q_{k-1}^*Q_{k-1}
 =
 U_k^*U_k+Q_{k-1}^*U_k+U_k^*Q_{k-1}.
\]
Sum over \(k\).  The left side telescopes from zero to
\(Q_n^*Q_n\), proving \eqref{eq:V46-path-energy}.  Applying a positive
functional and pairing the last two terms gives
\(2\operatorname{Re}\varphi(Q_{k-1}^*U_k)\).  Conditions
\eqref{eq:V46-orthogonality-needed} make each such paired contribution
zero.  Norm bounds alone impose no such equality, as the scalar
example below shows.  Cancellation among different \(k\) could also
make their total zero, but that too would be an additional signed
identity.
\end{proof}

\begin{firewall}[A deterministic interpolation is not a martingale]
\label{fire:V46-not-martingale}
The index \(k\) in \eqref{eq:V46-Pk} changes which complete local
moment is used.  No conditional expectation, filtration, or
mean-zero property has been constructed.  Calling \(U_k\) a
martingale difference would therefore add an unproved hypothesis.
\end{firewall}

\begin{proposition}[Sharp scalar saturation of the cross terms]
\label{prop:V46-scalar-prefix}
Let \(0\le\delta_j\le1\), take the scalar contractions
\begin{equation}
 T_j=1,\qquad R_j=1-\delta_j,
\label{eq:V46-scalar-TR}
\end{equation}
and put \(D:=\sum_j\delta_j\).  Then
\begin{align}
 U_k&=\delta_k\prod_{j>k}(1-\delta_j)\ge0,
\label{eq:V46-scalar-U}\\
 {\mathfrak J}_{<t}^{xq}
 &=1-\prod_{j=1}^n(1-\delta_j).
\label{eq:V46-scalar-J}
\end{align}
Every off-diagonal product \(U_kU_\ell\) is nonnegative.  Moreover,
if \(D\le1\), then
\begin{equation}
 \frac D2
 \le {\mathfrak J}_{<t}^{xq}
 \le D.
\label{eq:V46-scalar-sharp}
\end{equation}
Hence the available contraction and local-difference hypotheses
permit fully constructive cross terms and a prefix square of order
\(D^2\).
\end{proposition}

\begin{proof}
Equations \eqref{eq:V46-scalar-U}--\eqref{eq:V46-scalar-J} follow
directly from \eqref{eq:V46-Pk}--\eqref{eq:V46-J-path}.  The upper
bound in \eqref{eq:V46-scalar-sharp} is the elementary union bound
\(1-\prod_j(1-\delta_j)\le\sum_j\delta_j\).  Since
\(1-\delta_j\le e^{-\delta_j}\),
\[
 1-\prod_j(1-\delta_j)\ge1-e^{-D}.
\]
For \(0\le D\le1\), concavity or the Taylor bound gives
\(1-e^{-D}\ge D/2\).  Positivity of all \(U_k\) gives the sign claim.
\end{proof}

\begin{corollary}[No automatic signed gain in the V45 prefix matrix]
\label{cor:V46-no-automatic-gain}
The exact coordinate order and the two actual difference occurrences
in \eqref{eq:V46-right-cross} do not, by themselves, imply a
cancellation or an \(\ell^2\) improvement in the aggregate prefix
matrix.  Any such gain must use a further identity special to the
Wilson moments, the ternary carrier, the payer dressing, or the
product-state output.
\end{corollary}

\begin{proof}
The scalar family in \Cref{prop:V46-scalar-prefix} satisfies
contractivity and has
\(\|\Delta_j\|=\delta_j\), so it satisfies the same local response
form after setting \(\delta_j=s_\alpha h_{12,e_j}\).  Nevertheless
all cross terms have the same sign and the aggregate is comparable
to the \(\ell^1\) response sum.  Therefore those hypotheses cannot
imply cancellation or square-function improvement.
\end{proof}

\section{An abstract obstruction to the inverse-mass branch}
\label{sec:V46-abstract-obstruction}

\begin{theorem}[Current abstract inputs do not imply \(ABSC_2\)]
\label{thm:V46-abstract-no-ABSC2}
There is no deduction of \eqref{eq:V45-ABSC2} from only the following
four inputs:
\begin{enumerate}[label=\textup{(\roman*)},leftmargin=4em]
\item \(T_j\) and \(R_j\) are contractions;
\item \(\|T_j-R_j\|\le Cs_\alpha h_{12,e_j}\);
\item the ternary factor is bounded and has the exact V40
      four-replica polarization; and
\item the exterior dressing factors are contractions.
\end{enumerate}
More precisely, scalar data satisfying all four properties can have a
nonzero aggregate word independent of a marked center mass
\(\Xi_1=M\), while the inverse-mass branch on the right of
\eqref{eq:V45-ABSC2} tends to zero as \(M\to\infty\).
\end{theorem}

\begin{proof}
Take one prefix coordinate, \(s_\alpha=1\), \(T_1=1\),
\(R_1=1-\delta\) with fixed \(0<\delta<1\), and
\(h_{12,e_1}=\delta\).  All exterior dressings are the scalar one.
Let \(X\) be Bernoulli with parameter \(p=1/4\) and put
\(A=B=C=X\).  Its third cumulant is
\[
 S=\kappa_3(X,X,X)
 =p(1-p)(1-2p)=\frac3{32}.
\]
By \Cref{thm:V40-four-replica}, this scalar \(S\) has exactly the
required four-replica polarization.  The aggregate word is therefore
\[
 \mathbf W=(T_1-R_1)S=\frac{3\delta}{32}\ne0.
\]

Assign fixed positive values to \(h_{13,r}\) and \(h_{14,r}\).  Let
the marked set contain a third coordinate \(q\) with row-\(1\) mass
\(a_{1,q}=M\), and take \(\Xi_1=M+O(1)\).  Then
\(A_1(Q_{e_1,r,q})\ge M\), so the packet is marked-heavy for every
fixed \(\varepsilon<1\) when \(M\) is large.  The inverse-mass side of
\eqref{eq:V45-ABSC2} is
\[
 C_\varepsilon\Xi_1^{-2}
 \delta h_{13,r}h_{14,r}\longrightarrow0,
\]
whereas \(\kappa_\otimes^{\rm abs}(\mathbf W)=3\delta/32\).
Thus the four listed inputs do not imply the aggregate card.
\end{proof}

\begin{firewall}[Scope of the scalar obstruction]
\label{fire:V46-scalar-scope}
The scalar construction in
\Cref{thm:V46-abstract-no-ABSC2} is not asserted to be a physical
Yang--Mills packet.  It deliberately omits every Wilson-specific
relation between the marked Casimir mass and the heat, character,
resolvent, and output factors.  It proves exactly that such a relation
is a necessary new input; it does not disprove \(ABSC_2\) for the
physical carrier.
\end{firewall}

\begin{corollary}[Mass visibility is the next mandatory input]
\label{cor:V46-mass-visibility}
A proof of \(ABSC_2\) or \(ACWQ_2\) must use a quantitative hypothesis
which fails on the scalar family above and therefore makes the marked
mass \(A_1(Q_{e,r,q})\) visible inside the complete operator.  Since
\[
 A_1(Q_{e,r,q})
 =
 \sum_{u\in\{e,r,q\}}a_{1,u},
\]
where the braces delete repeated coordinates, the new estimate must retain
mass-dependent damping or cancellation in at least one of the prefix
coordinate \(e\), the ternary coordinate \(r\), or the payer/exterior
coordinate \(q\).
\end{corollary}

\begin{proof}
If all estimates used only the four mass-blind inputs of
\Cref{thm:V46-abstract-no-ABSC2}, that theorem would contradict the
claimed conclusion.  Hence some additional estimate must depend on
the marked mass.  By definition \eqref{eq:V36-marked-mass}, that mass
is the sum of the row-\(1\) weights at the distinct coordinates among
\(e,r,q\), which gives the stated exhaustive list.
\end{proof}

\section{V46 conclusion and the upstream split}
\label{sec:V46-conclusion}

\begin{theorem}[Exact V46 status of the aggregate route]
\label{thm:V46-status}
The V45 aggregate reformulation is algebraically valid, but its cross
terms have no automatic favorable sign, martingale orthogonality, or
mass decay.  The first unresolved input is now a Wilson-specific
mass-visibility estimate on the complete marked word.  The appropriate
upstream division is the exhaustive center-mass split
\begin{equation}
 \max\{a_{1,e},a_{1,r},a_{1,q}\}
 \ge\frac13A_1(Q_{e,r,q}),
\label{eq:V46-three-way-mass}
\end{equation}
with repeated coordinates deleted.  The three cases must be tested
against, respectively, the exact binary covariance, the protected
ternary character channel, and the paid exterior/resolvent carrier.
\end{theorem}

\begin{proof}
The cross-term assertions are
\Cref{lem:V46-path-energy,
prop:V46-scalar-prefix,
cor:V46-no-automatic-gain}.  The need for mass visibility is
\Cref{thm:V46-abstract-no-ABSC2,cor:V46-mass-visibility}.  Finally,
\eqref{eq:V46-three-way-mass} is the pigeonhole principle applied to
the at most three nonnegative summands in
\(A_1(Q_{e,r,q})\).  This division concerns only the center row and
therefore does not reinstate the erroneous all-row locality
classification withdrawn in V42.
\end{proof}

\begin{theorem}[Third-series V46 result ledger]
\label{thm:V46-results}
Relative to V1--V45, V46 proves:
\begin{enumerate}[label=\textup{(V46.\arabic*)},leftmargin=6.3em]
\item an exact interpolation path whose increments are the V10
      binary-prefix telescope summands;
\item coordinatewise normal forms for every diagonal and
      off-diagonal prefix-square term;
\item an exact path-energy identity identifying the missing
      orthogonality conditions;
\item proof that the deterministic interpolation is not a
      martingale under the presently available hypotheses;
\item a sharp scalar family in which all cross terms are constructive
      and the prefix response has its full \(\ell^1\) size;
\item a scalar no-go theorem showing that the current mass-blind
      abstract inputs cannot imply \(ABSC_2\);
\item the consequent mass-visibility requirement; and
\item the exhaustive three-coordinate center-mass split which the
      next physical Wilson estimate must address.
\end{enumerate}
\end{theorem}

\begin{proof}
These are
\Cref{lem:V46-increments,
thm:V46-cross-normal-form,
cor:V46-no-coordinate-multiplicity,
lem:V46-path-energy,
fire:V46-not-martingale,
prop:V46-scalar-prefix,
thm:V46-abstract-no-ABSC2,
cor:V46-mass-visibility,
thm:V46-status}.
\end{proof}

\begin{programme}[V47: Wilson mass-visibility trichotomy]
\label{prog:V46-V47}
\begin{enumerate}[leftmargin=3em]
\item Freeze one marked-heavy word and split it by
      \eqref{eq:V46-three-way-mass}; do not impose dominance on the
      leaf rows.
\item In the \(e\)-dominant case, retain the exact local heat and
      difference formula for \(T_e-R_e\) and determine whether one
      local binary covariance carries two center derivatives or only
      one.
\item In the \(r\)-dominant case, reuse only the valid local portion
      of V41 and test the protected balanced character channel,
      including the \(SU(2)\) obstruction from V36.
\item In the \(q\)-dominant case, write the payer numerator and
      resolvent in their exact spectral order and determine whether
      the heat factor supplies a second inverse Casimir without
      allocating a second payer.
\item If a case has only one visible center response, retain the
      aggregate cross matrix and identify the exact mixed factor that
      must supply the other.
\item Do not replace any of the three physical carriers by the scalar
      diagnostic; it established insufficiency only.
\item No companion audit is due before V50.
\end{enumerate}
\end{programme}

\begin{firewall}[Claim boundary after third-series V46]
\label{fire:V46-claim-boundary}
V46 proves the exact prefix cross-term structure and an abstract
no-go theorem, but it does not prove a physical mass-visibility
estimate, \(ACWQ_2\), \(ABSC_2\), the local-payer polarization, the
one-occurrence or directed/coupled center cases, the \(R_1\) paths,
or the full \(2+1+1\) source.  The \(3+1\), global discrete, globally
aggregated \([4]\), and \([3,3]\) families, higher MTA, WUFC, CPM,
cutoff removal, reflection positivity, Osterwalder--Schrader
reconstruction, construction of the continuum Yang--Mills measure,
and a uniform positive continuum spectral gap remain open.  The full
Yang--Mills mass gap has not been proved.
\end{firewall}

\paragraph{Parent-version record.}
V46 was copied byte-for-byte from
\texttt{Renewal\_Yang\_Mills\_Mass\_Gap\_Research\_Notes\_V45.tex}
before this chapter was appended.  The V45 parent had \(23\,544\)
source lines, \(856\,405\) bytes, and SHA--256
\[
 \texttt{c5116d0a54f9d4b5b42c00ea19fabf1deffb7e5cfc509f9596e58e942e2d7894}.
\]
The next compulsory companion audit is V50.

% =====================================================================
% V47 APPENDIX: PHYSICAL MASS-VISIBILITY TRICHOTOMY
% =====================================================================

\chapter{V47: Physical Center-Mass Visibility and the Balanced
Low-Output Residual}
\label{ch:V47-mass-trichotomy}

\section{A disjoint dominant-coordinate partition}
\label{sec:V47-dominant-partition}

Fix a marked-heavy word
\(\mathcal W_{e,r,q}^{\rm pay,\pi}\).  The set
\(Q_{e,r,q}=\{e,r\}\cup\{q\}\) contains at most three distinct
coordinates.  Choose a deterministic order on the three role names,
and let
\[
 d=d(e,r,q)\in Q_{e,r,q}
\]
be the first coordinate in that order at which
\(\max_{u\in Q_{e,r,q}}a_{1,u}\) is attained.  This tie rule partitions
the marked-heavy family into \(e\)-, \(r\)-, and \(q\)-dominant
classes without duplicating a word; when role names denote the same
coordinate, the first such role owns it.

\begin{lemma}[Dominant marked-coordinate inequality]
\label{lem:V47-dominant-coordinate}
For every marked-heavy word,
\begin{equation}
 a_{1,d}
 \ge\frac13A_1(Q_{e,r,q})
 \ge\frac{\varepsilon}{3}\Xi_1.
\label{eq:V47-dominant-mass}
\end{equation}
Consequently
\begin{equation}
 A_1(Q_{e,r,q})^{-2}\le a_{1,d}^{-2},
 \qquad
 a_{1,d}^{-2}\le9A_1(Q_{e,r,q})^{-2},
\label{eq:V47-mass-equivalence}
\end{equation}
so the local dominant mass and the marked mass are comparable up to
the fixed factor nine.
\end{lemma}

\begin{proof}
The marked mass is a sum of at most three nonnegative terms.  Its
largest term is therefore at least one third of the sum.  Marked
heaviness gives the second inequality in
\eqref{eq:V47-dominant-mass}.  Since
\(a_{1,d}\le A_1(Q_{e,r,q})\le3a_{1,d}\), taking inverse squares gives
\eqref{eq:V47-mass-equivalence}.
\end{proof}

\begin{corollary}[Three local two-debit targets suffice]
\label{cor:V47-local-targets}
On the class owned by \(d\in\{e,r,q\}\), the fixed-word estimate
\begin{equation}
 \kappa_\otimes^{\rm abs}
 (\mathcal W_{e,r,q}^{\rm pay,\pi})
 \le C
 \min\{s_\alpha^2,a_{1,d}^{-2}\}
 h_{12,e}h_{13,r}h_{14,r}
\label{eq:V47-local-two-debit}
\end{equation}
implies \(BSC_2\) for that word, with a changed absolute constant.
The same conclusion holds for an aggregate estimate over a
dominance class if its right side uses the corresponding signed
source sum and the common bound
\(a_{1,d}\ge(\varepsilon/3)\Xi_1\).
\end{corollary}

\begin{proof}
The first assertion follows from
\eqref{eq:V47-mass-equivalence}.  For the aggregate assertion use
\(a_{1,d}^{-2}\le9\varepsilon^{-2}\Xi_1^{-2}\) termwise and repeat
the normalization in \Cref{thm:V45-ABSC2-suffices}.
\end{proof}

\begin{firewall}[Center dominance only]
\label{fire:V47-center-only}
The partition above uses only the three row-\(1\) masses in the marked
set.  It makes no locality or dominance assertion about rows
\(2,3,4\), and therefore does not restore the seven all-row cells
withdrawn in V42.
\end{firewall}

\section{The already closed unbalanced branch}
\label{sec:V47-unbalanced}

At a coordinate \(u\), call a resolved local representation tuple
unbalanced when it satisfies the V32 highest-weight inequality
\eqref{eq:V32-unbalanced}; otherwise it is balanced.  The word
complete below means that all local output blocks are assembled before
the positive envelope, as required in V34.

\begin{theorem}[Dominant mass in a compiled unbalanced bank closes]
\label{thm:V47-unbalanced-closes}
Suppose the dominant coordinate \(d\) of
\Cref{lem:V47-dominant-coordinate} belongs to a complete unbalanced
bank to which \Cref{thm:V34-unbalanced-bank} applies, with the payer
retained at its actual location.  Then the corresponding
marked-heavy word satisfies \(BSC_2\).
\end{theorem}

\begin{proof}
By \Cref{lem:V34-unbalanced-output}, the resolved nontrivial output
mass at \(d\) is comparable to the dominant input Casimir mass.
\Cref{thm:V34-unbalanced-bank} retains the resulting exponential heat
damping and controls the unique numerator/resolvent when the bank
pays.  Thus the complete word has the exponential marked form
\[
 C s_\alpha^2
 e^{-c s_\alpha a_{1,d}}
 h_{12,e}h_{13,r}h_{14,r},
\]
after harmless fixed-group constants and the already compiled
spectators are absorbed.  By
\eqref{eq:V47-dominant-mass},
\(a_{1,d}\ge A_1(Q_{e,r,q})/3\).  Apply
\Cref{lem:V36-exponential-two-debit} with the exponent constant
reduced by three.  This proves \(BSC_2\) without assigning a second
payer.
\end{proof}

\begin{corollary}[The mass-visibility residual is balanced]
\label{cor:V47-balanced-residual}
After applying all existing private, unbalanced, and complete-bank
compilers, every still unresolved dominant-coordinate class is a
balanced protected or balanced low-output class.  Exponential
unbalanced output separation need not be reconsidered there.
\end{corollary}

\begin{proof}
The unbalanced alternative is
\Cref{thm:V47-unbalanced-closes}; the private and earlier compiled
alternatives are part of
\Cref{thm:V36-refined-residual}.  Their complement is the balanced
protected/low-output sector by the exhaustive V32
balanced--unbalanced split.
\end{proof}

\section{What a dominant prefix coordinate actually supplies}
\label{sec:V47-e-dominant}

Assume now that \(d=e\).  Since every active row mass has the fixed
floor \(a_*>0\),
\begin{equation}
 h_{12,e}=a_{1,e}a_{2,e}\ge a_*a_{1,e}.
\label{eq:V47-he-lower}
\end{equation}
For the singleton binary bank \(\{e\}\), define its ordinary response
ratio
\begin{equation}
 \gamma_e:=
 \frac{\min\{1,s_\alpha h_{12,e}\}}{h_{12,e}}
 =
 \min\{h_{12,e}^{-1},s_\alpha\}.
\label{eq:V47-gamma-e}
\end{equation}

\begin{proposition}[The \(e\)-dominant prefix gives one center debit]
\label{prop:V47-e-one-debit}
On the \(e\)-dominant marked-heavy class,
\begin{equation}
 \gamma_e
 \le
 \frac{3}{a_*\varepsilon\,\Xi_1}.
\label{eq:V47-e-one-debit}
\end{equation}
This is one genuine center response.  The presently available
ordinary binary-prefix estimate does not provide the second factor
\(\Xi_1^{-1}\).
\end{proposition}

\begin{proof}
Equations \eqref{eq:V47-he-lower}--\eqref{eq:V47-gamma-e} give
\(\gamma_e\le(a_*a_{1,e})^{-1}\).  Apply
\eqref{eq:V47-dominant-mass}.  The final statement is exactly the
one-debit ceiling proved in
\Cref{thm:V44-one-debit}; V46 shows that cross-coordinate
contractivity alone does not improve it.
\end{proof}

\begin{firewall}[A large center does not make the named partner
large]
\label{fire:V47-e-partner}
The inequality \(a_{1,e}\gtrsim\Xi_1\) does not imply
\(a_{2,e}\gtrsim\Xi_2\) or \(a_{2,e}\gtrsim a_{1,e}\).  Hence one may
not replace \eqref{eq:V47-e-one-debit} by a squared binary response
using \(h_{12,e}\) twice.  A second debit must come from an actual
Wilson factor in the complete word.
\end{firewall}

\section{The \(r\)-dominant local target}
\label{sec:V47-r-dominant}

\begin{proposition}[Exact cancellation of the two center coefficients]
\label{prop:V47-r-stock}
On the \(r\)-dominant class, the inverse local-mass branch of
\eqref{eq:V47-local-two-debit} has scalar stock
\begin{equation}
 a_{1,r}^{-2}
 h_{12,e}h_{13,r}h_{14,r}
 =
 h_{12,e}a_{3,r}a_{4,r}.
\label{eq:V47-r-cancel}
\end{equation}
Thus the required two center powers are already represented by the
two actual row-\(1\) coefficients at the ternary coordinate.  The
remaining analytic problem is to prove that the balanced protected
operator carrier does not lose those two powers before output
capacity is taken.
\end{proposition}

\begin{proof}
Use
\(h_{13,r}=a_{1,r}a_{3,r}\) and
\(h_{14,r}=a_{1,r}a_{4,r}\), then cancel
\(a_{1,r}^2\).  The two coefficient occurrences correspond to the two
star arms \(13\) and \(14\); this is an algebraic inventory, not an
operator estimate.
\end{proof}

\begin{corollary}[The valid V41 local subcell remains closed]
\label{cor:V47-r-V41}
If, in addition to \(r\)-dominance of the center, the leaf locality
hypotheses of \eqref{eq:V41-fully-local-cell} hold, then the V41
ternary theorem closes this subcell.  Without those extra leaf
hypotheses, V41 supplies no estimate for the remaining balanced
protected \(r\)-dominant carrier.
\end{corollary}

\begin{proof}
The first assertion is
\Cref{thm:V41-fully-local-TQ2}.  Its proof explicitly uses lower
bounds for the normalized row-\(3\) and row-\(4\) local masses, so it
cannot be invoked after those hypotheses are removed.
\end{proof}

\section{Why a dominant payer coordinate need not be heat-visible}
\label{sec:V47-q-dominant}

The unique payer acts on a resolved output Casimir, not directly on
an arbitrarily chosen input Casimir.  In an unbalanced channel V34
compares the two, which is why
\Cref{thm:V47-unbalanced-closes} applies.  Balanced tensor products
have low-output channels for which this comparison fails.

\begin{proposition}[\(SU(2)\) balanced low-output obstruction]
\label{prop:V47-SU2-low-output}
For every spin \(j\ge1\),
\begin{equation}
 V_j\otimes V_j
 \cong
 \bigoplus_{\ell=0}^{2j}V_\ell.
\label{eq:V47-CG}
\end{equation}
Hence this balanced tensor product contains both the trivial output
\(V_0\) and the fixed nontrivial output \(V_1\), while
\[
 a_j=1+j(j+1)\longrightarrow\infty.
\]
There is therefore no constant \(c>0\) such that every nonzero
resolved output mass \(\Lambda_{\rm out}\) in a balanced coordinate
satisfies
\begin{equation}
 \Lambda_{\rm out}\ge c\,a_j.
\label{eq:V47-false-output-lower}
\end{equation}
\end{proposition}

\begin{proof}
Equation \eqref{eq:V47-CG} is the Clebsch--Gordan decomposition for
two equal \(SU(2)\) spins.  The spin-one output has fixed Casimir
\(C_2(1)=2\), while \(a_j\to\infty\).  Thus the ratio of its output
mass to \(a_j\) tends to zero, contradicting
\eqref{eq:V47-false-output-lower}.  The spin-zero summand gives the
even stronger zero-output instance, although an output-Casimir
numerator may annihilate that particular block.
\end{proof}

\begin{corollary}[The payer heat is not a generic second input debit]
\label{cor:V47-payer-not-input-debit}
On the balanced \(q\)-dominant class, the scalar
numerator/resolvent and its output heat multiplier cannot by
themselves be converted into a factor \(a_{1,q}^{-2}\) by a uniform
lower comparison between output and input Casimirs.  Any closure of
that class must additionally use product-state suppression of the
low-output projection or signed coupling to the prefix--ternary
carrier.
\end{corollary}

\begin{proof}
The proposed conversion would require a comparison of the form
\eqref{eq:V47-false-output-lower} in order for output heat to see the
large input mass.  \Cref{prop:V47-SU2-low-output} disproves such a
comparison already for a fixed compact group.  The listed
alternatives are the remaining structures retained in the complete
balanced carrier.
\end{proof}

\begin{firewall}[The zero-output block and the whole protected sector]
\label{fire:V47-zero-output}
If the payer numerator is the output Casimir, its trivial-output block
vanishes.  This does not close the entire protected sector: the
spin-one channel in \Cref{prop:V47-SU2-low-output} has nonzero but
uniformly bounded output Casimir while the input mass diverges.
Deleting only the invariant block therefore does not produce two
inverse input-Casimir powers.
\end{firewall}

\begin{proposition}[Conditional closure when the output really sees
the payer mass]
\label{prop:V47-output-visible-closes}
Suppose a \(q\)-dominant complete word has an output block satisfying
\(\Lambda_{\rm out}\ge c\,a_{1,q}\), and suppose its already-paid
complete envelope retains the heat bound
\begin{equation}
 C s_\alpha^2e^{-c's_\alpha\Lambda_{\rm out}}
 h_{12,e}h_{13,r}h_{14,r}.
\label{eq:V47-visible-envelope}
\end{equation}
Then that block satisfies \(BSC_2\).
\end{proposition}

\begin{proof}
By \(q\)-dominance and
\Cref{lem:V47-dominant-coordinate},
\[
 \Lambda_{\rm out}
 \ge c a_{1,q}
 \ge\frac c3A_1(Q_{e,r,q}).
\]
Insert this in \eqref{eq:V47-visible-envelope} and apply
\Cref{lem:V36-exponential-two-debit}.  The envelope is already paid,
so no further numerator or resolvent is introduced.
\end{proof}

\section{The exact balanced residual after mass visibility}
\label{sec:V47-residual}

\begin{theorem}[V47 balanced mass-location residual]
\label{thm:V47-residual}
After the V36 compiled cases, the V45 aggregate reformulation, and
the unbalanced closure
\Cref{thm:V47-unbalanced-closes}, every unresolved marked-heavy
double-center word lies in exactly one of:
\begin{enumerate}[label=\textup{(\roman*)},leftmargin=4em]
\item an \(e\)-dominant balanced class with the one genuine prefix
      debit \eqref{eq:V47-e-one-debit} and one missing center debit;
\item an \(r\)-dominant balanced protected class requiring retention
      of the two local center coefficients through the operator
      estimate \eqref{eq:V47-r-cancel}; or
\item a \(q\)-dominant balanced low-output class in which generic
      output heat does not see the large input mass.
\end{enumerate}
The first usable joint object remains the aggregate square
\eqref{eq:V45-aggregate-square}, restricted to one of these three
disjoint dominance classes.  No class is declared closed except the
compiled unbalanced and V41 fully local subcells stated above.
\end{theorem}

\begin{proof}
\Cref{lem:V47-dominant-coordinate} gives the disjoint three-role
partition.  \Cref{thm:V47-unbalanced-closes} removes every dominant
coordinate lying in an applicable unbalanced complete bank.
\Cref{prop:V47-e-one-debit,prop:V47-r-stock,
cor:V47-payer-not-input-debit} give the exact status of the three
balanced alternatives.  V45 proves that the aggregate square is a
sufficient weaker interface.  The closure boundary is
\Cref{cor:V47-r-V41,prop:V47-output-visible-closes}.
\end{proof}

\begin{theorem}[Third-series V47 result ledger]
\label{thm:V47-results}
Relative to V1--V46, V47 proves:
\begin{enumerate}[label=\textup{(V47.\arabic*)},leftmargin=6.3em]
\item a disjoint, repeated-coordinate-safe dominant marked-coordinate
      partition;
\item equivalence up to constants of dominant local mass and marked
      mass;
\item closure of every dominant coordinate covered by the complete
      unbalanced exponential compiler;
\item the exact one-debit endpoint estimate on the
      \(e\)-dominant class;
\item algebraic cancellation of the two center coefficients in the
      \(r\)-dominant inverse-mass stock, with the V41 subcell boundary;
\item an \(SU(2)\) Clebsch--Gordan obstruction to converting balanced
      input mass into output heat mass;
\item conditional closure of every output-visible paid block; and
\item the exact three-class balanced mass-location residual.
\end{enumerate}
\end{theorem}

\begin{proof}
These are
\Cref{lem:V47-dominant-coordinate,
cor:V47-local-targets,
thm:V47-unbalanced-closes,
prop:V47-e-one-debit,
prop:V47-r-stock,
cor:V47-r-V41,
prop:V47-SU2-low-output,
prop:V47-output-visible-closes,
thm:V47-residual}.
\end{proof}

\begin{programme}[V48: balanced low-output product-state capacity]
\label{prog:V47-V48}
\begin{enumerate}[leftmargin=3em]
\item Begin with the \(q\)-dominant balanced low-output class, since
      V47 proves that output heat alone cannot close it.
\item Resolve the one-coordinate output projection across the cut
      separating row \(1\) from the other rows, but keep its
      multiplicity space and its coupling to the binary--ternary word.
\item Compute the exact product-state capacity of each fixed low
      output, not only the invariant output, and compare its
      highest-weight scaling with \(a_{1,q}^{-2}\).
\item Test the estimate first on
      \(V_j\otimes V_j\to V_0,V_1\) for \(SU(2)\); a bound which fails
      there cannot be used in the general compact-group proof.
\item If the fixed-output capacity supplies less than two powers,
      retain its product with the actual prefix debit or the two
      ternary replica differences before taking the supremum.
\item Do not use dimension suppression as Casimir suppression without
      an explicit Weyl-growth exponent and a uniform multiplicity
      argument.
\item Leave the \(e\)- and \(r\)-dominant balanced classes unchanged
      until the low-output capacity is decided.  No companion audit is
      due before V50.
\end{enumerate}
\end{programme}

\begin{firewall}[Claim boundary after third-series V47]
\label{fire:V47-claim-boundary}
V47 closes the compiled unbalanced and output-visible subcells and
gives an exact balanced residual.  It does not prove the balanced
low-output capacity, \(ACWQ_2\), \(ABSC_2\), all local-payer
polarizations, the one-occurrence or directed/coupled center cases,
the \(R_1\) paths, or the full \(2+1+1\) source.  The \(3+1\), global
discrete, globally aggregated \([4]\), and \([3,3]\) families,
higher MTA, WUFC, CPM, cutoff removal, reflection positivity,
Osterwalder--Schrader reconstruction, construction of the continuum
Yang--Mills measure, and a uniform positive continuum spectral gap
remain open.  The full Yang--Mills mass gap has not been proved.
\end{firewall}

\paragraph{Parent-version record.}
V47 was copied byte-for-byte from
\texttt{Renewal\_Yang\_Mills\_Mass\_Gap\_Research\_Notes\_V46.tex}
before this chapter was appended.  The V46 parent had \(24\,009\)
source lines, \(873\,052\) bytes, and SHA--256
\[
 \texttt{72f593f485e4593d0886002509684984e07516ecbe48dd938f23b31aaf5b0100}.
\]
The next compulsory companion audit is V50.

% =====================================================================
% V48 APPENDIX: BALANCED LOW-OUTPUT PRODUCT-STATE CAPACITY
% =====================================================================

\chapter{V48: Balanced Low-Output Projectors Do Not Supply Two
Casimir Debits}
\label{ch:V48-low-output}

\section{A bipartite projection lemma}
\label{sec:V48-projection}

For finite-dimensional Hilbert spaces \(H_1,H_2\), define the pure
product-state capacity of a positive operator \(P\) by
\begin{equation}
 \kappa_{\rm prod}(P)
 :=
 \sup_{\substack{\|u\|=\|v\|=1}}
 \langle u\otimes v,P(u\otimes v)\rangle.
\label{eq:V48-product-capacity}
\end{equation}
For a projection, this is the squared norm of the largest product
component in its range.  Mixed product densities give the same
supremum by convexity.

\begin{lemma}[Universal rank lower bound]
\label{lem:V48-rank-lower}
If \(P\) is an orthogonal projection of rank \(r\) on
\(\mathbb C^{d_1}\otimes\mathbb C^{d_2}\), then
\begin{equation}
 \boxed{
 \kappa_{\rm prod}(P)\ge\frac{r}{d_1d_2}.}
\label{eq:V48-rank-lower}
\end{equation}
\end{lemma}

\begin{proof}
Average the product expectation in
\eqref{eq:V48-product-capacity} over independent Haar unit vectors
\(u,v\).  Since
\[
 \int |u\rangle\langle u|\,du=\frac{I_{d_1}}{d_1},
 \qquad
 \int |v\rangle\langle v|\,dv=\frac{I_{d_2}}{d_2},
\]
the average equals
\[
 \operatorname{Tr}\left[
 P\left(\frac{I_{d_1}}{d_1}\otimes
        \frac{I_{d_2}}{d_2}\right)\right]
 =\frac{\operatorname{Tr}P}{d_1d_2}
 =\frac r{d_1d_2}.
\]
The supremum is at least the average.
\end{proof}

\begin{lemma}[Upper bound from a scalar marginal]
\label{lem:V48-scalar-marginal}
Let \(P\ge0\) act on
\(\mathbb C^d\otimes\mathbb C^d\).  If
\begin{equation}
 \operatorname{Tr}_2P=cI_d,
\label{eq:V48-scalar-marginal}
\end{equation}
then
\begin{equation}
 \kappa_{\rm prod}(P)\le c.
\label{eq:V48-marginal-upper}
\end{equation}
If \(P\) is a rank-\(r\) projection, then \(c=r/d\).
\end{lemma}

\begin{proof}
For a unit vector \(u\), let
\[
 B_u:=(\langle u|\otimes I)P(|u\rangle\otimes I).
\]
This is positive and
\[
 \operatorname{Tr}B_u
 =
 \langle u,(\operatorname{Tr}_2P)u\rangle=c.
\]
For every unit \(v\),
\[
 \langle u\otimes v,P(u\otimes v)\rangle
 =\langle v,B_uv\rangle
 \le\|B_u\|\le\operatorname{Tr}B_u=c.
\]
Taking the supremum proves \eqref{eq:V48-marginal-upper}.  When \(P\)
is a projection, tracing \eqref{eq:V48-scalar-marginal} gives
\(cd=\operatorname{Tr}P=r\).
\end{proof}

\section{Application to equal-spin \(SU(2)\) outputs}
\label{sec:V48-SU2}

Let \(d_j=2j+1\) and let \(P_{j,\ell}\) denote the orthogonal
projection from \(V_j\otimes V_j\) onto the multiplicity-one
Clebsch--Gordan summand \(V_\ell\), where
\(0\le\ell\le2j\).

\begin{theorem}[Two-sided product-capacity bounds for a fixed output]
\label{thm:V48-SU2-projector}
For every allowed \(j,\ell\),
\begin{equation}
 \boxed{
 \frac{2\ell+1}{d_j^2}
 \le
 \kappa_{\rm prod}(P_{j,\ell})
 \le
 \frac{2\ell+1}{d_j}.}
\label{eq:V48-SU2-projector}
\end{equation}
In particular, for each fixed \(\ell\), the capacity cannot decay
faster than order \(d_j^{-2}\) as \(j\to\infty\).
\end{theorem}

\begin{proof}
The projection has rank \(2\ell+1\), so the lower bound is
\Cref{lem:V48-rank-lower}.  It is invariant under the diagonal
\(SU(2)\) action:
\[
 (U_j(g)\otimes U_j(g))P_{j,\ell}
 =
 P_{j,\ell}(U_j(g)\otimes U_j(g)).
\]
Therefore \(\operatorname{Tr}_2P_{j,\ell}\) commutes with the
irreducible representation \(U_j\).  Schur's lemma makes it a scalar
multiple of the identity.  Its trace is
\(\operatorname{rank}P_{j,\ell}=2\ell+1\), so
\[
 \operatorname{Tr}_2P_{j,\ell}
 =\frac{2\ell+1}{d_j}I_{d_j}.
\]
Apply \Cref{lem:V48-scalar-marginal}.
\end{proof}

\begin{proposition}[Exact invariant-output capacity]
\label{prop:V48-singlet}
For the singlet projection \(P_{j,0}\),
\begin{equation}
 \kappa_{\rm prod}(P_{j,0})=\frac1{d_j}.
\label{eq:V48-singlet}
\end{equation}
\end{proposition}

\begin{proof}
The normalized singlet vector has Schmidt decomposition with all
\(d_j\) Schmidt coefficients equal to \(d_j^{-1/2}\).  The maximum
overlap of a bipartite unit vector with a product vector is its
largest Schmidt coefficient.  Squaring gives \(d_j^{-1}\), which is
also the upper bound in \eqref{eq:V48-SU2-projector} at \(\ell=0\).
\end{proof}

\begin{corollary}[The fixed spin-one output misses the two-debit
scale]
\label{cor:V48-spin-one-no-two-debit}
For \(j\ge1\), the nontrivial fixed output \(\ell=1\) satisfies
\begin{equation}
 \kappa_{\rm prod}(P_{j,1})
 \ge\frac3{(2j+1)^2}.
\label{eq:V48-spin-one-lower}
\end{equation}
There is no constant \(C\), independent of \(j\), such that
\begin{equation}
 \kappa_{\rm prod}(P_{j,1})
 \le\frac{C}{(1+j(j+1))^2}.
\label{eq:V48-false-two-debit}
\end{equation}
\end{corollary}

\begin{proof}
The lower bound is
\eqref{eq:V48-SU2-projector}.  If
\eqref{eq:V48-false-two-debit} held, then
\[
 \frac3{(2j+1)^2}
 \le\frac{C}{(1+j(j+1))^2}.
\]
Multiplying by the denominator on the right would bound a quantity
asymptotic to \(3j^2/4\), a contradiction.
\end{proof}

\begin{firewall}[Dimension suppression is not two Casimir powers]
\label{fire:V48-dimension-Casimir}
For \(SU(2)\), \(d_j\asymp a_j^{1/2}\).  The universal rank lower
bound for a fixed low output is therefore of order \(a_j^{-1}\), and
the singlet capacity is of order \(a_j^{-1/2}\).  Neither may be
relabeled as \(a_j^{-2}\).  The fact that the output rank is fixed
does not change this conclusion.
\end{firewall}

\section{Insertion of the paid low-output scalar}
\label{sec:V48-paid-output}

For a resolved output Casimir \(\Lambda>0\), write the usual
one-payer spectral scalar abstractly as
\begin{equation}
 g_{s_\alpha}(\Lambda)
 :=
 \frac{\Lambda e^{-s_\alpha\Lambda}}
 {1-e^{-s_\alpha\Lambda}}.
\label{eq:V48-payer-scalar}
\end{equation}
Changing the fixed normalization of the heat parameter only changes
constants.  For the spin-one \(SU(2)\) output,
\(\Lambda=C_2(1)=2\), so
\(g_{s_\alpha}(2)>0\) is independent of the input spin \(j\).

\begin{proposition}[The isolated paid projector still fails]
\label{prop:V48-paid-projector-fails}
Fix \(s_\alpha>0\).  The positive paid spin-one block
\[
 g_{s_\alpha}(2)P_{j,1}
\]
does not satisfy a product-state capacity bound of order \(a_j^{-2}\)
uniformly in \(j\).
\end{proposition}

\begin{proof}
Positive homogeneity and
\eqref{eq:V48-spin-one-lower} give
\[
 \kappa_{\rm prod}(g_{s_\alpha}(2)P_{j,1})
 \ge
 \frac{3g_{s_\alpha}(2)}{(2j+1)^2}.
\]
The scalar prefactor is fixed and positive.  The contradiction
argument in \Cref{cor:V48-spin-one-no-two-debit} applies unchanged.
\end{proof}

\begin{firewall}[The payer is not duplicated]
\label{fire:V48-one-payer}
The scalar in \eqref{eq:V48-payer-scalar} is the one already
allocated numerator/resolvent.  V48 neither squares it nor introduces
a second resolvent.  The obstruction is that this actual payer is
constant on a fixed low-output channel while the balanced input
Casimir diverges.
\end{firewall}

\section{Consequence for the \(q\)-dominant residual}
\label{sec:V48-consequence}

\begin{theorem}[No isolated-output-capacity proof of the
\(q\)-dominant card]
\label{thm:V48-no-isolated-output}
The balanced \(q\)-dominant \(BSC_2\) or \(ABSC_2\) estimate cannot be
deduced from a tensor-product factorization in which the paid
low-output projection is estimated in product-state capacity and all
remaining binary--ternary factors are replaced only by
support-uniform constants.  This failure occurs already in the
physical \(SU(2)\) channel
\[
 V_j\otimes V_j\longrightarrow V_1.
\]
\end{theorem}

\begin{proof}
Such a factorization would require the isolated paid projection to
supply the \(a_{1,q}^{-2}\) branch, because the other factors have
been replaced by mass-independent constants.  Take
\(a_{1,q}=a_j=1+j(j+1)\).  By
\Cref{prop:V48-paid-projector-fails}, the isolated block has capacity
bounded below at order \(a_j^{-1}\), and therefore cannot obey a
uniform \(Ca_j^{-2}\) upper bound.  The Clebsch--Gordan channel is an
actual balanced \(SU(2)\) output by
\eqref{eq:V47-CG}.
\end{proof}

\begin{corollary}[The low-output projection must remain coupled]
\label{cor:V48-must-couple}
Any proof of the balanced \(q\)-dominant class must retain the
low-output projection together with at least one mass-sensitive
factor from the exact binary--ternary carrier before taking
product-state capacity.  The candidates currently present are the
selected binary prefix occurrence, the two center-row replica
differences in the ternary square, and their signed
cross-coordinate aggregate.
\end{corollary}

\begin{proof}
\Cref{thm:V48-no-isolated-output} excludes replacing all those factors
by constants before the output projection is estimated.  The listed
objects are exactly the center-bearing occurrences inventoried in
\Cref{cor:V43-center-factors} and the legal aggregate terms of
\Cref{thm:V45-aggregate-square}.
\end{proof}

\begin{assessment}[What the exact capacity computation changes]
\label{ass:V48-direction}
V47 left two possible mechanisms on the balanced \(q\)-dominant
branch: isolated low-output suppression or coupling to the signed
binary--ternary word.  V48 eliminates the first.  The next smallest
object is therefore not another bound on \(P_{j,\ell}\), but the
compressed sandwich
\[
 P_{j,\ell}\,
 D_r(\boldsymbol\xi)R_{e,\pi}\rho
 R_{f,\pi}^*D_r(\boldsymbol\eta)^*\,
 P_{j,\ell}
\]
inside the product-state duality, with the actual row-\(1\) prefix or
ternary differences retained.  A successful estimate must gain at
least one further inverse power beyond the rank lower-bound scale.
\end{assessment}

\begin{theorem}[Third-series V48 result ledger]
\label{thm:V48-results}
Relative to V1--V47, V48 proves:
\begin{enumerate}[label=\textup{(V48.\arabic*)},leftmargin=6.3em]
\item a universal rank lower bound for the product-state capacity of
      a bipartite projection;
\item a scalar-marginal upper bound;
\item two-sided capacity bounds for every
      \(V_j\otimes V_j\to V_\ell\) \(SU(2)\) projection;
\item the exact singlet product-state capacity;
\item a fixed spin-one physical obstruction to two inverse input
      Casimir powers;
\item persistence of that obstruction after insertion of the unique
      paid low-output scalar;
\item exclusion of every isolated-output-capacity proof of the
      balanced \(q\)-dominant card; and
\item reduction to a low-output projection coupled to an actual
      prefix or ternary response occurrence.
\end{enumerate}
\end{theorem}

\begin{proof}
These are
\Cref{lem:V48-rank-lower,
lem:V48-scalar-marginal,
thm:V48-SU2-projector,
prop:V48-singlet,
cor:V48-spin-one-no-two-debit,
prop:V48-paid-projector-fails,
thm:V48-no-isolated-output,
cor:V48-must-couple}.
\end{proof}

\begin{programme}[V49: compressed low-output response sandwich]
\label{prog:V48-V49}
\begin{enumerate}[leftmargin=3em]
\item Fix the \(SU(2)\) spin-one output first and retain its exact
      Clebsch--Gordan projection inside the V45 aggregate square.
\item Move the product-state vectors through the outer contractions,
      but do not take the capacity of \(P_{j,1}\) separately.
\item Express one actual row-\(1\) replacement difference in the
      Clebsch--Gordan basis and compute the compressed matrix
      coefficients between \(V_j\otimes V_j\) and \(V_1\).
\item Determine whether the difference annihilates the leading
      \(d_j^{-2}\) product-overlap channel or supplies an additional
      factor comparable to \(a_j^{-1}\).
\item Test both the diagonal and cross-prefix terms; do not infer the
      cross estimate from the diagonal.
\item If the spin-one channel still has order larger than
      \(a_j^{-2}\), record a physical counterpacket to the proposed
      coupled card and return upstream to the full five-role
      cancellation.
\item Generalize to fixed compact groups only after the \(SU(2)\)
      regression succeeds.  The next companion audit is V50.
\end{enumerate}
\end{programme}

\begin{firewall}[Claim boundary after third-series V48]
\label{fire:V48-claim-boundary}
V48 proves a physical low-output obstruction and selects the coupled
projector-response sandwich.  It does not estimate that sandwich,
prove \(ACWQ_2\) or \(ABSC_2\), close the balanced \(e\)-,
\(r\)-, or \(q\)-dominant classes, prove all local-payer
polarizations, close the one-occurrence or directed/coupled center
cases or the \(R_1\) paths, or prove the full \(2+1+1\) source.  The
\(3+1\), global discrete, globally aggregated \([4]\), and
\([3,3]\) families, higher MTA, WUFC, CPM, cutoff removal,
reflection positivity, Osterwalder--Schrader reconstruction,
construction of the continuum Yang--Mills measure, and a uniform
positive continuum spectral gap remain open.  The full Yang--Mills
mass gap has not been proved.
\end{firewall}

\paragraph{Parent-version record.}
V48 was copied byte-for-byte from
\texttt{Renewal\_Yang\_Mills\_Mass\_Gap\_Research\_Notes\_V47.tex}
before this chapter was appended.  The V47 parent had \(24\,467\)
source lines, \(890\,364\) bytes, and SHA--256
\[
 \texttt{55abcbe4acb36b1ded5c29006c250e8a9f3fc93f1c8c350fbd20fded94df3593}.
\]
The next compulsory companion audit is V50.

% =====================================================================
% V49 APPENDIX: ARCHIVE CORRECTION AND COMPOSITE POLARIZATION
% =====================================================================

\chapter{V49: Archive Reconciliation and Full Composite-Source
Polarization}
\label{ch:V49-composite}

\section{Correction of the V48 novelty claim}
\label{sec:V49-correction}

Before carrying out the proposed Clebsch--Gordan calculation, the
frozen f2 archive was searched through V84--V91.  That search changes
the correct provenance of V48.

\begin{record}[Archived results that precede V48]
\label{rec:V49-archive-priority}
The frozen file
\texttt{Renewal\_Yang\_Mills\_Mass\_Gap\_Research\_Notes\_V100\_f2.tex}
already contains:
\begin{enumerate}[label=\textup{(\roman*)},leftmargin=4em]
\item the V32 exact \(1/d_j\) one-row invariant-projector capacity,
      including arbitrary multiplicity;
\item the V84 partial-trace upper bound
      \((2J+1)/(2j+1)\) for
      \(V_j\otimes V_j\to V_J\);
\item the exact anti-aligned spin-one product weight
      \(6j/((2j+1)(2j+2))\asymp j^{-1}\);
\item paid and unpaid physical lower bounds on the spin-one joining
      block;
\item the V90 tilted-coherent counterexample to an isolated local
      signed joining card; and
\item the V91 proof that this obstruction survives the complete
      first-position sum and full connected \(2+2\) carrier, while a
      different legal marked path still controls that packet.
\end{enumerate}
\end{record}

\begin{proposition}[Exact status of V48 after archive reconciliation]
\label{prop:V49-V48-status}
The abstract rank lower bound
\eqref{eq:V48-rank-lower} is a valid general lemma not identified in
the archive search.  However, the \(SU(2)\) upper bound and the claim
that a fixed paid spin-one output cannot supply two inverse Casimir
powers are re-derivations of the stronger V84 calculation, not new
programme advances.  Items asking V49 to recompute the isolated
spin-one projection are withdrawn.  All V48 inequalities remain
mathematically valid.
\end{proposition}

\begin{proof}
V84's intermediate-dimension lemma is exactly the upper half of
\eqref{eq:V48-SU2-projector}.  Its exact product vector has spin-one
weight of order \(j^{-1}\), stronger than the universal V48 lower
bound of order \(j^{-2}\).  V84 also inserts the unique payer and
retains a lower bound of order \(j^{-1}\).  Hence
\Cref{prop:V48-paid-projector-fails} follows from the archived result
with a sharper exponent.  The universal averaging proof in
\Cref{lem:V48-rank-lower} applies to arbitrary bipartite projections
and is retained.  This proves the stated division between valid
content and novelty.
\end{proof}

\begin{firewall}[No transfer from the \(2+2\) counterpacket]
\label{fire:V49-no-S22-transfer}
The V84--V91 counterpacket belongs to the global \(2+2\) source.  It
is a mandatory regression against isolated low-output cards, but it
is not automatically a counterexample to the present
\(2+1+1\) source.  Its positive lower bound may be used here only
after an exact embedding into the current composite source is proved.
No such embedding is asserted in V49.
\end{firewall}

\section{Polarization before the five-role split}
\label{sec:V49-polarization}

At the \(2+1+1\) joining coordinate, put
\[
 Y_{12}:=X_1\otimes X_2.
\]
The local source letter before the V37 Leonov--Shiryaev split is
\begin{equation}
 \mathsf S_{12|3|4}
 :=
 \kappa_3(Y_{12},X_3,X_4).
\label{eq:V49-composite-source}
\end{equation}
Let \(\xi_0,\xi_1,\xi_2,\xi_3\) be independent replicas.  Write
\begin{align}
 \delta_iX_i
 &:=
 X_i(\xi_0)-X_i(\xi_i),
 \qquad i=1,2,3,4,
\label{eq:V49-row-differences}\\
 \delta_1Y_{12}
 &:=
 Y_{12}(\xi_0)-Y_{12}(\xi_1).
\label{eq:V49-composite-difference}
\end{align}
Only \(\delta_1Y_{12}\), \(\delta_3X_3\), and
\(\delta_4X_4\) enter the three-variable polarization; the notation
\(\delta_2X_2\) is used in the exact Leibniz split below.

\begin{theorem}[Exact composite four-replica polarization]
\label{thm:V49-composite-polarization}
The full five-role source letter satisfies
\begin{equation}
 \boxed{
 \mathsf S_{12|3|4}
 =
 {\mathbb E}_{\xi_0,\xi_1,\xi_3,\xi_4}
 \left[
 \delta_1Y_{12}
 \otimes\delta_3X_3
 \otimes\delta_4X_4
 \right].}
\label{eq:V49-composite-polarization}
\end{equation}
Moreover,
\begin{equation}
 \delta_1Y_{12}
 =
 \delta_1X_1\otimes X_2(\xi_0)
 +
 X_1(\xi_1)\otimes
 \bigl(X_2(\xi_0)-X_2(\xi_1)\bigr).
\label{eq:V49-Leibniz}
\end{equation}
Thus
\begin{equation}
 \mathsf S_{12|3|4}
 =
 {\mathbb E}\,[D^{(1)}(\boldsymbol\xi)
                 +D^{(2)}(\boldsymbol\xi)],
\label{eq:V49-two-branch}
\end{equation}
where
\begin{align}
 D^{(1)}(\boldsymbol\xi)
 &:=
 \delta_1X_1\otimes X_2(\xi_0)
 \otimes\delta_3X_3\otimes\delta_4X_4,
\label{eq:V49-D1}\\
 D^{(2)}(\boldsymbol\xi)
 &:=
 X_1(\xi_1)\otimes
 \bigl(X_2(\xi_0)-X_2(\xi_1)\bigr)
 \otimes\delta_3X_3\otimes\delta_4X_4.
\label{eq:V49-D2}
\end{align}
\end{theorem}

\begin{proof}
Apply \Cref{thm:V40-four-replica} to the three variables
\((A,B,C)=(Y_{12},X_3,X_4)\).  These act on the disjoint tensor legs
\(\{1,2\},\{3\},\{4\}\), so the theorem applies without reordering
operators.  This gives
\eqref{eq:V49-composite-polarization}.  Add and subtract
\(X_1(\xi_1)\otimes X_2(\xi_0)\) in
\eqref{eq:V49-composite-difference}; the two resulting differences
are exactly \eqref{eq:V49-Leibniz}.  Substitution gives
\eqref{eq:V49-two-branch}--\eqref{eq:V49-D2}.
\end{proof}

\begin{corollary}[The five local roles are reassembled]
\label{cor:V49-five-roles-reassembled}
The right side of \eqref{eq:V49-two-branch} equals the sum of the five
operators \(C_4,C_{31}^{(1)},C_{31}^{(2)},C_{22}^{(a)},C_{22}^{(b)}\)
in \eqref{eq:V37-five-role-identity}.  Consequently the composite
polarization retains every cancellation between the selected
fourth-cumulant, binary--ternary, and covariance-product roles before
an absolute value.
\end{corollary}

\begin{proof}
\Cref{thm:V49-composite-polarization} equals
\(\kappa_3(X_1X_2,X_3,X_4)\).  The V37
Leonov--Shiryaev identity is an exact second expression for the same
operator.  Equality of both expressions with
\eqref{eq:V49-composite-source} proves the claim.
\end{proof}

\begin{firewall}[The two Leibniz branches are not separately
positive]
\label{fire:V49-branches-signed}
Neither expectation of \(D^{(1)}\) nor that of \(D^{(2)}\) is
declared positive, and their cross terms in the square may have
either sign.  The split \eqref{eq:V49-Leibniz} inventories actual row
differences; it is not permission to apply the triangle inequality
before the two branches are recombined.
\end{firewall}

\section{The exact composite-source square}
\label{sec:V49-square}

Fix first an outside-local payer placement and write the complete
source word as
\begin{equation}
 \widehat W_\pi=L_\pi\mathsf S_{12|3|4}R_\pi.
\label{eq:V49-complete-word}
\end{equation}
All binary prefixes, singleton moments, separators, suffix factors,
output compression, and the unique payer remain in
\(L_\pi,R_\pi\).

\begin{theorem}[Exact four-block composite square]
\label{thm:V49-composite-square}
Let \(\boldsymbol\xi,\boldsymbol\eta\) be independent copies of the
replicas in \Cref{thm:V49-composite-polarization}.  Then
\begin{equation}
 \boxed{
 \widehat W_\pi^*\widehat W_\pi
 =
 \sum_{a,b\in\{1,2\}}
 {\mathbb E}_{\boldsymbol\eta,\boldsymbol\xi}
 \left[
 R_\pi^*D^{(a)}(\boldsymbol\eta)^*
 L_\pi^*L_\pi
 D^{(b)}(\boldsymbol\xi)R_\pi
 \right].}
\label{eq:V49-composite-square}
\end{equation}
The complete sum is positive.  Its \((1,1)\) block contains two
actual row-\(1\) replacement differences, while its \((2,2)\) block
contains two actual row-\(2\) replacement differences.  The mixed
blocks contain one of each.
\end{theorem}

\begin{proof}
Insert \eqref{eq:V49-two-branch} and its independent adjoint into
\[
 \widehat W_\pi^*\widehat W_\pi
 =
 R_\pi^*\mathsf S_{12|3|4}^*
 L_\pi^*L_\pi
 \mathsf S_{12|3|4}R_\pi,
\]
then distribute the two finite sums and apply Fubini.  This proves
\eqref{eq:V49-composite-square}.  Positivity follows from the left
side.  The occurrence statement is visible in
\eqref{eq:V49-D1}--\eqref{eq:V49-D2}: each independent replica copy
of \(D^{(1)}\) has one row-\(1\) difference, and each copy of
\(D^{(2)}\) has one row-\(2\) difference.
\end{proof}

\begin{corollary}[The two star centers are present in one carrier]
\label{cor:V49-two-centers}
Before the five-role split, the positive square
\eqref{eq:V49-composite-square} contains a genuine two-occurrence
quadratic carrier for each of the two possible star centers.  This
does not prove two inverse Casimir powers, but it removes the V36
factorization defect in which the scalar arms \(13,14\) were
incorrectly treated as separate operator endpoints.
\end{corollary}

\begin{proof}
The row-\(1\) and row-\(2\) diagonal blocks have the stated actual
differences by \Cref{thm:V49-composite-square}.  They arise from two
copies of the complete composite source, not from the two scalar
marks in a positive ternary envelope.  Quantitative response bounds
on the compressed blocks remain unproved.
\end{proof}

\section{A source-level quadratic interface}
\label{sec:V49-interface}

Let \(\mathfrak A_{211}(\pi)\) denote the complete literal
marked-tree right side for this fixed \(2+1+1\) source packet and
payer placement: it includes the eight lower-role trees of
\Cref{lem:V35-eight-trees} and the selected-four-row currency, with
their existing coefficient normalizations.  This is notation for the
already specified target, not a new enlargement of its tree atlas.

\begin{definition}[Full composite-source quadratic card]
\label{def:V49-FCQ}
The output-oriented \(FCQ_{211}(\pi)\) card is
\begin{equation}
 \boxed{
 \kappa_\otimes(\widehat W_\pi^*\widehat W_\pi)
 \le C\,\mathfrak A_{211}(\pi)^2,}
\label{eq:V49-FCQ}
\end{equation}
or the corresponding left-square estimate when required by output
duality.
\end{definition}

\begin{theorem}[The full composite card bypasses the separated
\(BSC_2\) interface]
\label{thm:V49-FCQ-suffices}
If \(FCQ_{211}(\pi)\) holds for every allowed payer placement, then
the complete \(2+1+1\) joining packet satisfies its literal
marked-tree source estimate.  This implication does not require
separate \(BSC_2\), \(ABSC_2\), or selected lower-role estimates.
\end{theorem}

\begin{proof}
Apply the product-state Cauchy--Schwarz inequality
\Cref{lem:V43-complete-word-CS} to the complete operator
\(\widehat W_\pi\), then take the nonnegative square root of
\eqref{eq:V49-FCQ}.  This gives
\[
 \kappa_\otimes^{\rm abs}(\widehat W_\pi)
 \le C\mathfrak A_{211}(\pi),
\]
which is precisely the fixed-payer literal marked-tree target by the
definition of \(\mathfrak A_{211}(\pi)\).  Since the five roles have
not been separated, no estimate on an individual role is invoked.
\end{proof}

\begin{proposition}[\(BSC_2\) is sufficient but no longer mandatory]
\label{prop:V49-BSC2-not-mandatory}
The earlier chain
\[
 BSC_2\Longrightarrow BTE_2
\Longrightarrow\hbox{control of the two lower-order stars}
\]
remains valid.  It is not a necessary declared interface for the
source-level route: \(FCQ_{211}\) permits cancellation between all
five exact roles in \eqref{eq:V37-five-role-identity}, which an
individual \(BSC_2\) estimate forbids.
\end{proposition}

\begin{proof}
The first implication chain is
\Cref{thm:V42-BSC2-restored,thm:V43-CWQ2-suffices}.  The second
statement follows from
\Cref{cor:V49-five-roles-reassembled,thm:V49-FCQ-suffices}.  This is
a statement about sufficient proof interfaces; it does not assert
that the five-role cancellation is quantitatively strong enough.
\end{proof}

\begin{firewall}[The source-level card remains open]
\label{fire:V49-FCQ-open}
Equation \eqref{eq:V49-composite-square} is an identity, while
\eqref{eq:V49-FCQ} is an unproved estimate.  Counting two actual
row-\(1\) or row-\(2\) differences does not turn them into two
inverse mass powers.  The mixed \((1,2)\) and \((2,1)\) blocks may
not be discarded, because positivity belongs only to their complete
four-block sum.
\end{firewall}

\section{V49 status and next compulsory audit}
\label{sec:V49-status}

\begin{theorem}[Third-series V49 result ledger]
\label{thm:V49-results}
Relative to V1--V48, V49 proves:
\begin{enumerate}[label=\textup{(V49.\arabic*)},leftmargin=6.3em]
\item exact archive reconciliation showing which V48 projector
      statements were already proved more sharply in V32/V84;
\item withdrawal of the redundant proposed spin-one recomputation;
\item a four-replica polarization of the unsplit composite source
      \(\kappa_3(X_1X_2,X_3,X_4)\);
\item the exact two-branch Leibniz decomposition into actual row-\(1\)
      and row-\(2\) differences;
\item reassembly of all five V37 local roles before absolute values;
\item the exact four-block positive square of the complete composite
      word;
\item genuine two-occurrence carriers for both possible star centers;
      and
\item a source-level quadratic interface sufficient without separate
      \(BSC_2\) or \(ABSC_2\) cards.
\end{enumerate}
\end{theorem}

\begin{proof}
These are
\Cref{rec:V49-archive-priority,
prop:V49-V48-status,
thm:V49-composite-polarization,
cor:V49-five-roles-reassembled,
thm:V49-composite-square,
cor:V49-two-centers,
def:V49-FCQ,
thm:V49-FCQ-suffices,
prop:V49-BSC2-not-mandatory}.
\end{proof}

\begin{programme}[V50: companion audit and composite-square
regression]
\label{prog:V49-V50}
\begin{enumerate}[leftmargin=3em]
\item Perform the compulsory audit of the newest active SM and NS
      notes, recording hashes and applying the type firewall.
\item Test \eqref{eq:V49-FCQ} against the archived V37 normalized
      character family and the V84--V91 coherent families, but claim
      a counterexample only after an exact \(2+1+1\) embedding.
\item Expand the \((1,1)\) block of
      \eqref{eq:V49-composite-square} through the actual local Wilson
      heat output, retaining the mixed blocks.
\item Determine whether the five-role recombination cancels the
      leading balanced low-output channel or merely reroutes it to a
      different legal tree.
\item If the channel survives, compare its exact coefficient with the
      full eight-tree atlas \(\mathfrak A_{211}\), not with the star
      currency alone.
\item Keep the unique payer fixed in both square copies and leave
      local-payer words separate until their paid composite
      polarization is written.
\end{enumerate}
\end{programme}

\begin{firewall}[Claim boundary after third-series V49]
\label{fire:V49-claim-boundary}
V49 supplies an exact unsplit composite-source carrier and a weaker
sufficient interface.  It does not prove \(FCQ_{211}\), close the
complete \(2+1+1\) source, prove all local-payer polarizations, close
the remaining one-occurrence or directed/coupled center cases or the
\(R_1\) paths, or close the \(3+1\), global discrete, globally
aggregated \([4]\), and \([3,3]\) families.  Higher MTA, WUFC, CPM,
cutoff removal, reflection positivity, Osterwalder--Schrader
reconstruction, construction of the continuum Yang--Mills measure,
and a uniform positive continuum spectral gap remain open.  The full
Yang--Mills mass gap has not been proved.
\end{firewall}

\paragraph{Parent-version record.}
V49 was copied byte-for-byte from
\texttt{Renewal\_Yang\_Mills\_Mass\_Gap\_Research\_Notes\_V48.tex}
before this chapter was appended.  The V48 parent had \(24\,856\)
source lines, \(903\,555\) bytes, and SHA--256
\[
 \texttt{cb35a016f5ff2a4f02df5b2131405161dfde75c2f08435174846fe4133474d77}.
\]
The next compulsory companion audit is V50.

% =====================================================================
% V50 APPENDIX: COMPANION AUDIT AND REGRESSION ELIGIBILITY
% =====================================================================

\chapter{V50: Companion Audit and the Composite-Source Regression
Boundary}
\label{ch:V50-audit}

\section{Compulsory companion audit}
\label{sec:V50-companion-audit}

At the time of this audit, the newest active companion files were
\begin{align*}
 &\texttt{Renewal\_SM\_Einstein\_Continuum\_Research\_Notes\_V29.tex},\\
 &\texttt{Renewal\_NS\_Stability\_Research\_Notes\_V31.tex}.
\end{align*}
The SM file had \(10\,638\) source lines, \(402\,860\) bytes, and
SHA--256
\[
 \texttt{481d727916e8716611025b0442dc83d317d1b8871e812ad0e1bdd2b5f9103b5f}.
\]
The NS file had \(23\,052\) source lines, \(887\,537\) bytes, and
SHA--256
\[
 \texttt{f1121b62238ad5491bb7cf03635ea9934942edf573ed8faaa9ee79e1d38a6696}.
\]
The NS file is unchanged from V45.  The SM programme advanced from
V23 to V29.  Its new V24--V29 results were inspected, with particular
attention to the V28 exact second density score and the V29
massless-chain blocking obstruction.

\begin{theorem}[V50 companion transfer decision]
\label{thm:V50-companion-decision}
No new companion estimate applies to the Wilson composite square
\eqref{eq:V49-composite-square}.  Two proof-order conclusions do
apply:
\begin{enumerate}[label=\textup{(\roman*)},leftmargin=4em]
\item the SM V28 second density derivative is a genuine centered
      second score, not the square of its first score; analogously,
      the four blocks in \eqref{eq:V49-composite-square} must remain
      the exact square of the composite source and may not be replaced
      by duplicated one-difference estimates; and
\item SM V29 proves that coarse blocking amplifies an already strict
      contraction but does not create one, while NS V31 proves the
      analogous sharpness of heat lifting; therefore neither support
      aggregation nor an additional heat step may be cited as the
      missing two-center estimate without a strict microscopic
      Wilson inequality.
\end{enumerate}
\end{theorem}

\begin{proof}
SM V28 works with a strongly convex quartic onsite measure, hopping
bonds, density derivatives, and a Finner/KP expansion.  SM V29 works
with conditional influence matrices.  NS V31 works with Fourier
shells and the Navier--Stokes Duhamel map.  None of those object types
matches the finite-dimensional Wilson tensor carrier, so no analytic
bound transfers.  The two conclusions stated above are the
type-independent logical content of their exact-score and
no-generated-contraction theorems.
\end{proof}

\begin{firewall}[V50 companion type boundary]
\label{fire:V50-companion-types}
No Higgs convexity, Finner, KP, Dobrushin, Gaussian-chain, Leray,
coarea, or Duhamel estimate is used as a Yang--Mills inequality.
The audit changes only the admissible proof order.
\end{firewall}

\section{The archived coherent packet is not yet a \(2+1+1\)
regression}
\label{sec:V50-regression-boundary}

The V49 archive reconciliation found a strong V90--V91
\(SU(2)\) obstruction.  Its local joining letter is a covariance
between a same-side pair and a third row.  The current source letter
is the third cumulant
\(\kappa_3(X_1X_2,X_3,X_4)\).  These are not the same connected
functional.

\begin{lemma}[A constant fourth row annihilates the composite source]
\label{lem:V50-constant-fourth-row}
Let \(A,B\) be bounded operator-valued random variables on disjoint
tensor legs and let \(I\) be deterministic.  Then
\begin{equation}
 \kappa_3(A,B,I)=0.
\label{eq:V50-constant-cumulant}
\end{equation}
Equivalently, the composite polarization
\eqref{eq:V49-composite-polarization} vanishes when \(X_4\) is
constant.
\end{lemma}

\begin{proof}
In the four-replica formula, the factor
\(I(\xi_0)-I(\xi_4)\) is zero pointwise.  Hence the integrand and its
expectation vanish.  This is also immediate from the moment--cumulant
formula: all terms factor through the deterministic argument and
their Möbius coefficients sum to zero.
\end{proof}

\begin{corollary}[No idle-row embedding of V90--V91]
\label{cor:V50-no-idle-embedding}
Tensoring the archived three-row coherent joining packet with a
trivial or deterministic fourth row does not embed it into
\(\mathsf S_{12|3|4}\); it produces the zero third cumulant.
Therefore the V90--V91 lower bound is not, without a new four-row
construction, a counterexample to \(FCQ_{211}\).
\end{corollary}

\begin{proof}
Take \(A=X_1X_2\), \(B=X_3\), and \(X_4=I\) in
\Cref{lem:V50-constant-fourth-row}.  The resulting local source is
zero, whereas the archived joining covariance is nonzero on its
coherent packet.  Thus the proposed tensoring does not preserve the
operator being tested.
\end{proof}

\begin{proposition}[The first admissible regression must use four
active rows]
\label{prop:V50-four-active}
Any physical regression test for \eqref{eq:V49-FCQ} derived from the
archived coherent family must choose a nonconstant fourth-row Wilson
coefficient and verify all of the following on one common sequence:
\begin{enumerate}[label=\textup{(\alph*)},leftmargin=4em]
\item the \(2+1+1\) first-connectivity word is nonzero;
\item the V90 coherent low-output contribution survives the extra
      replacement difference \(\delta_4X_4\);
\item the other five-role components do not cancel the proposed lower
      bound; and
\item the full eight-tree atlas
      \(\mathfrak A_{211}\), not merely its star term, is smaller than
      that lower bound if a counterexample is claimed.
\end{enumerate}
Absent these four verifications, the archive supplies only a
regression warning, not a disproof of the source-level card.
\end{proposition}

\begin{proof}
Item \textup{(a)} is required because the target concerns the
specified source.  Item \textup{(b)} follows from the exact
polarization \eqref{eq:V49-composite-polarization}; without the fourth
difference the source is zero by
\Cref{lem:V50-constant-fourth-row}.  Item \textup{(c)} is required by
\Cref{cor:V49-five-roles-reassembled}.  Item \textup{(d)} is required
by the definition \eqref{eq:V49-FCQ} and by the archived V91 example,
which violates one proposed currency while satisfying a rerouted
legal tree.  These conditions are therefore necessary for the stated
regression use.
\end{proof}

\begin{firewall}[No counterexample by analogy]
\label{fire:V50-no-analogy-counterexample}
A lower bound for a covariance, a fixed output projection, or one
separated five-role component is not a lower bound for the full
composite third cumulant.  Cancellation must be evaluated after the
four active rows and all exact local roles are assembled.
\end{firewall}

\section{V50 status}
\label{sec:V50-status}

\begin{theorem}[Third-series V50 result ledger]
\label{thm:V50-results}
Relative to V1--V49, V50 proves:
\begin{enumerate}[label=\textup{(V50.\arabic*)},leftmargin=6.3em]
\item the compulsory SM V29 and NS V31 audit with exact file hashes;
\item an explicit no-transfer decision and its proof-order
      consequences;
\item exact annihilation of the composite third cumulant by a
      deterministic fourth row;
\item proof that the archived three-row V90--V91 coherent packet has
      no idle-row embedding into the \(2+1+1\) source; and
\item the four necessary checks for any valid physical regression
      against \(FCQ_{211}\).
\end{enumerate}
\end{theorem}

\begin{proof}
These are
\Cref{thm:V50-companion-decision,
lem:V50-constant-fourth-row,
cor:V50-no-idle-embedding,
prop:V50-four-active}.
\end{proof}

\begin{programme}[V51: genuine four-row composite regression]
\label{prog:V50-V51}
\begin{enumerate}[leftmargin=3em]
\item Choose the smallest nonconstant fourth-row \(SU(2)\) coefficient
      compatible with the \(12|3|4\) source and the V90 coherent
      pair.
\item Evaluate its replacement difference in
      \eqref{eq:V49-composite-polarization} before resolving output
      blocks.
\item Determine whether the V90 heat-window contribution survives in
      the full composite cumulant and compute its sign and scale.
\item Expand the complete five-role identity on the same row-product
      vector to verify cancellation or survival.
\item Compare any surviving lower bound with every legal tree in the
      V35 eight-tree atlas.  If a rerouted tree controls it, retain
      that routing rather than demanding \(BSC_2\).
\item Keep one payer and all prefix coordinates exact.  No companion
      audit is due before V55.
\end{enumerate}
\end{programme}

\begin{firewall}[Claim boundary after third-series V50]
\label{fire:V50-claim-boundary}
V50 validates the audit and regression boundary but does not prove or
disprove \(FCQ_{211}\), close the complete \(2+1+1\) source, or close
any later source, MTA, WUFC, CPM, cutoff-removal, reconstruction, or
spectral-gap gate.  The full Yang--Mills mass gap has not been proved.
\end{firewall}

\paragraph{Parent-version record.}
V50 was copied byte-for-byte from
\texttt{Renewal\_Yang\_Mills\_Mass\_Gap\_Research\_Notes\_V49.tex}
before this chapter was appended.  The V49 parent had \(25\,282\)
source lines, \(919\,153\) bytes, and SHA--256
\[
 \texttt{8dc8683b12ed46b93657795131deb16eb1b1630ddb8d4c6cd0c4f05cca6ceb0d}.
\]
The next compulsory companion audit is V55.

% =====================================================================
% V51 APPENDIX: FINITE-HEAT FOUR-ROW CHARACTER REGRESSION
% =====================================================================

\chapter{V51: A Genuine Four-Active-Row Finite-Heat Regression}
\label{ch:V51-character}

\section{Finite-heat character moments}
\label{sec:V51-moments}

Let \(G=SU(2)\), let \(\chi_j\) be the spin-\(j\) character,
\(d_j=2j+1\), and put \(x_j=\chi_j/d_j\).  For fixed \(s>0\), let
\({\mathbb E}_s\) denote expectation against the normalized central
heat kernel based at the identity, with the standing Fourier
normalization
\begin{equation}
 {\mathbb E}_s\chi_\ell
 =
 d_\ell q_\ell,
 \qquad
 q_\ell:=e^{-s\ell(\ell+1)}.
\label{eq:V51-heat-character}
\end{equation}
Define the finite constant
\begin{equation}
 C_s:=\sum_{\ell=0}^{\infty}(2\ell+1)
 e^{-s\ell(\ell+1)}<\infty.
\label{eq:V51-Cs}
\end{equation}

\begin{lemma}[First, second, and fourth heat moments]
\label{lem:V51-heat-moments}
For every half-integer \(j\ge0\),
\begin{align}
 {\mathbb E}_sx_j
 &=q_j,
\label{eq:V51-first}\\
 0\le{\mathbb E}_sx_j^2
 &=
 \frac1{d_j^2}
 \sum_{\ell=0}^{2j}(2\ell+1)q_\ell
 \le\frac{C_s}{d_j^2},
\label{eq:V51-second}\\
 {\mathbb E}_sx_j^4
 &\ge\frac1{d_j^3}.
\label{eq:V51-fourth}
\end{align}
Also \(|x_j(g)|\le1\) pointwise, and hence
\begin{equation}
 |{\mathbb E}_sx_j^3|\le1.
\label{eq:V51-third-crude}
\end{equation}
\end{lemma}

\begin{proof}
Equation \eqref{eq:V51-first} is
\eqref{eq:V51-heat-character}.  The Clebsch--Gordan identity
\[
 \chi_j^2=\sum_{\ell=0}^{2j}\chi_\ell
\]
and \eqref{eq:V51-heat-character} give
\eqref{eq:V51-second}; extending the positive finite sum gives the
upper bound \(C_s/d_j^2\).

Write
\[
 \chi_j^4=\sum_{\nu\ge0}m_\nu\chi_\nu
\]
with nonnegative tensor-product multiplicities \(m_\nu\).  The
trivial multiplicity is
\[
 m_0
 =
 \dim (V_j^{\otimes4})^{SU(2)}
 =
 \sum_{\ell=0}^{2j}1
 =
 d_j,
\]
because each intermediate spin \(\ell\) in
\(V_j\otimes V_j\) occurs once and
\((V_\ell\otimes V_\ell)^{SU(2)}\) is one-dimensional.  Every term
\(m_\nu d_\nu q_\nu\) in the heat expectation is nonnegative.
Keeping only \(\nu=0\) yields
\[
 {\mathbb E}_s x_j^4
 =
 d_j^{-4}\sum_\nu m_\nu d_\nu q_\nu
 \ge d_j^{-4}m_0
 =d_j^{-3}.
\]
Finally, the character is the trace of a \(d_j\)-dimensional unitary,
so \(|\chi_j|\le d_j\).  Thus \(|x_j|\le1\), proving
\eqref{eq:V51-third-crude}.
\end{proof}

\section{The composite third cumulant survives at finite heat}
\label{sec:V51-cumulant}

Use four active row variables
\[
 X_1=X_2=X_3=X_4=x_j.
\]
Then the unsplit local \(2+1+1\) source of V49 is
\[
 \mathsf S_j
 :=
 \kappa_{3,s}(x_j^2,x_j,x_j).
\]

\begin{theorem}[Finite-heat four-row lower bound]
\label{thm:V51-finite-heat-lower}
For every fixed \(s>0\), there is \(j_0(s)\) such that
\begin{equation}
 \boxed{
 \kappa_{3,s}(x_j^2,x_j,x_j)
 \ge\frac1{2d_j^3}>0
 \qquad(j\ge j_0(s)).}
\label{eq:V51-finite-heat-lower}
\end{equation}
All four row variables are nonconstant and belong to the same
physical one-link heat law.
\end{theorem}

\begin{proof}
For arbitrary scalar random variables \(A,B,C\),
\[
 \kappa_3(A,B,C)
 =
 {\mathbb E}ABC
2{\mathbb E}A\,{\mathbb E}B\,{\mathbb E}C
-{\mathbb E}A\,{\mathbb E}BC
-{\mathbb E}B\,{\mathbb E}AC
-{\mathbb E}C\,{\mathbb E}AB.
\]
Set \(A=x_j^2\) and \(B=C=x_j\).  Then
\begin{align}
 \mathsf S_j
 =
 {\mathbb E}_sx_j^4
 -({\mathbb E}_sx_j^2)^2
 -2({\mathbb E}_sx_j)({\mathbb E}_sx_j^3)
 +2({\mathbb E}_sx_j^2)({\mathbb E}_sx_j)^2.
\label{eq:V51-cumulant-expansion}
\end{align}
The last term is nonnegative.  By
\Cref{lem:V51-heat-moments},
\[
 \mathsf S_j
 \ge
 d_j^{-3}-C_s^2d_j^{-4}-2q_j.
\]
For fixed \(s>0\), \(C_s\) is fixed and
\(q_j=e^{-sj(j+1)}=o(d_j^{-N})\) for every \(N\).
Choose \(j_0(s)\) so that
\(C_s^2d_j^{-4}+2q_j\le\tfrac12d_j^{-3}\).
This proves \eqref{eq:V51-finite-heat-lower}.  Nonconstancy follows
from \(j>0\).
\end{proof}

\begin{corollary}[The isolated composite letter has no two-Casimir
decay]
\label{cor:V51-no-local-two-Casimir}
Let \(a_j=1+j(j+1)\).  There is no constant \(C_s\), independent of
\(j\), such that
\begin{equation}
 |\kappa_{3,s}(x_j^2,x_j,x_j)|
 \le C_s a_j^{-2}
\label{eq:V51-false-local-two-Casimir}
\end{equation}
for all \(j\).
\end{corollary}

\begin{proof}
By \Cref{thm:V51-finite-heat-lower}, the left side is bounded below
by \(c d_j^{-3}\asymp j^{-3}\), whereas
\(a_j^{-2}\asymp j^{-4}\).  Their ratio grows linearly in \(j\).
\end{proof}

\begin{remark}[Genuine extension of the V37 Haar test]
\label{rem:V51-nonduplication}
V37 computed the exact Haar value
\((d_j-1)/d_j^4\).  V51 does not repeat that calculation: it proves a
uniform lower bound for every fixed finite heat time using the heat
Fourier coefficients, including the nonzero first and third moments
which are absent in the Haar limit.
\end{remark}

\begin{firewall}[The local lower bound is not yet a source
counterexample]
\label{fire:V51-not-source-counterexample}
The inequality
\eqref{eq:V51-finite-heat-lower} concerns the unsplit local joining
letter.  It does not include the earlier binary prefix, the unique
payer, the suffix, or the full eight-tree right side.  It therefore
refutes an isolated \(a_j^{-2}\) composite-letter estimate but does
not disprove \(FCQ_{211}\) or quartic MTA.
\end{firewall}

\section{Effect on the V49 route}
\label{sec:V51-effect}

\begin{theorem}[The second debit cannot come from the local composite
letter alone]
\label{thm:V51-route-decision}
On the physical four-equal-spin family above, the complete local
composite cumulant supplies at most the scaling tested in
\eqref{eq:V51-finite-heat-lower}, which is one power of \(j\) larger
than the desired two-Casimir scale.  Consequently any proof of
\(FCQ_{211}\) must obtain the remaining suppression from the actual
binary prefix, the payer/suffix carrier, cancellation with other
coordinate positions, or domination by a different legal tree in
the full atlas.  It cannot be proved by first replacing
\(\mathsf S_{12|3|4}\) with an isolated
\(Ca_{1,r}^{-2}\) norm or capacity bound.
\end{theorem}

\begin{proof}
The impossibility of the isolated bound is
\Cref{cor:V51-no-local-two-Casimir}.  The complete word
\eqref{eq:V49-complete-word} contains exactly the additional
structures listed in the conclusion, and
\Cref{thm:V49-FCQ-suffices} permits them to remain coupled.  No other
source of mass dependence is present after the local letter is
isolated.
\end{proof}

\begin{firewall}[No return to \(TQ_2\)]
\label{fire:V51-no-TQ2}
The finite-heat character family reinforces the V42 withdrawal of
the isolated ternary endpoint target.  It may not be bypassed by
squaring the local lower bound or by calling the two replica copies
two coordinate banks.  The next calculation must be performed on a
complete source word.
\end{firewall}

\section{V51 status and next acquisition}
\label{sec:V51-status}

\begin{theorem}[Third-series V51 result ledger]
\label{thm:V51-results}
Relative to V1--V50, V51 proves:
\begin{enumerate}[label=\textup{(V51.\arabic*)},leftmargin=6.3em]
\item exact first and second normalized-character moments under the
      finite \(SU(2)\) heat law;
\item a fourth-moment lower bound from the exact invariant
      multiplicity \(d_j\);
\item a positive \(cd_j^{-3}\) lower bound for the genuine
      four-active-row composite third cumulant at every fixed heat
      time;
\item failure of an isolated two-inverse-Casimir estimate for that
      composite letter; and
\item reduction of the remaining debit to the complete prefix,
      payer/suffix, coordinate cancellation, or legal-tree rerouting.
\end{enumerate}
\end{theorem}

\begin{proof}
These are
\Cref{lem:V51-heat-moments,
thm:V51-finite-heat-lower,
cor:V51-no-local-two-Casimir,
thm:V51-route-decision}.
\end{proof}

\begin{programme}[V52: prefix transmission of the four-row character
channel]
\label{prog:V51-V52}
\begin{enumerate}[leftmargin=3em]
\item Add the smallest exact \(12\) first-connectivity prefix to the
      four-equal-spin joining coordinate while keeping a row-product
      density realizing the character trace.
\item Compute the prefix covariance expectation at the same finite
      heat time and determine whether it transmits the
      \(d_j^{-3}\) local lower bound.
\item Insert one actual payer and verify its resolved output block
      rather than using a scalar cap.
\item Evaluate all eight legal tree weights on the resulting
      two-coordinate representation assignment.
\item If one rerouted path dominates the lower bound, record the
      routing certificate; if all are smaller, obtain a genuine
      \(2+1+1\) counterpacket to \(FCQ_{211}\) and return upstream of
      the first-connectivity split.
\item No companion audit is due before V55.
\end{enumerate}
\end{programme}

\begin{firewall}[Claim boundary after third-series V51]
\label{fire:V51-claim-boundary}
V51 proves a finite-heat local obstruction, not a complete source
counterexample or estimate.  It does not decide \(FCQ_{211}\), close
the \(2+1+1\) source, or close any subsequent MTA, WUFC, CPM,
cutoff-removal, reconstruction, or spectral-gap gate.  The full
Yang--Mills mass gap has not been proved.
\end{firewall}

\paragraph{Parent-version record.}
V51 was copied byte-for-byte from
\texttt{Renewal\_Yang\_Mills\_Mass\_Gap\_Research\_Notes\_V50.tex}
before this chapter was appended.  The V50 parent had \(25\,507\)
source lines, \(928\,443\) bytes, and SHA--256
\[
 \texttt{5d15b3543c148265a8c41a586628ffa6387e7fdcec249e03773220bb878ffc7f}.
\]
The next compulsory companion audit is V55.

% =====================================================================
% V52 APPENDIX: PREFIX TRANSMISSION AND TREE REROUTING
% =====================================================================

\chapter{V52: Prefix Transmission, Correct Normalization, and a
Two-Coordinate Routing Certificate}
\label{ch:V52-prefix-routing}

\section{Purpose and correction of the proposed comparison}
\label{sec:V52-purpose}

V51 proved that, for four equal spin-\(j\) normalized characters at
one joining coordinate,
\[
 \kappa_{3,s}(x_j^2,x_j,x_j)\gtrsim d_j^{-3}.
\]
The next required test was to insert the smallest legal binary prefix
and one actual payer, then compare the resulting word with all legal
trees.  This chapter performs that test.

There is a necessary normalization correction at the outset.  The
quantities
\[
 h_{uv,e}=a_{u,e}a_{v,e}
\]
are derivative stocks in a positive influence majorant.  They are not
scalar coefficients multiplying the exact character cumulant.
Consequently the V51 lower bound may not be multiplied by
\(a_j^4\) merely because the positive ternary estimate contains two
quotient-edge stocks.  The exact word and the marked-tree currency
must be evaluated separately.

\begin{firewall}[Stocks are not source coefficients]
\label{fire:V52-stock-coefficient}
The implication
\[
 |\kappa_{3,s}(x_j^2,x_j,x_j)|\gtrsim d_j^{-3}
 \quad\Longrightarrow\quad
 |J_r|\gtrsim a_j^4d_j^{-3}
\]
is invalid.  The factor \(a_j^4\) occurs only after differentiating
and positively bounding the complete cumulant.  No identity in
V1--V51 inserts it into the signed multiplier.
\end{firewall}

\section{The smallest legal two-coordinate packet}
\label{sec:V52-packet}

Fix an admissible heat time \(s=s_\alpha>0\) and a nonzero
half-integer \(k\).  Put
\[
 \alpha:=a_k=1+k(k+1).
\]
Use two ordered coordinates \(e<r\).  At \(e\), rows \(1,2\) carry
spin \(k\), while rows \(3,4\) are inactive.  At \(r\), all four rows
carry spin \(j\).  There are no other active coordinates.  Thus
\begin{equation}
 \Xi_1=\Xi_2=A+\alpha,\qquad
 \Xi_3=\Xi_4=A,\qquad
 A:=a_j=1+j(j+1).
\label{eq:V52-row-masses}
\end{equation}

Let
\begin{equation}
 \beta_{s,k}
 :=
 {\mathbb E}_s x_k^2-({\mathbb E}_s x_k)^2
 =
 \frac1{d_k^2}
 \sum_{\ell=0}^{2k}(2\ell+1)e^{-s\ell(\ell+1)}
 -e^{-2sk(k+1)}.
\label{eq:V52-beta}
\end{equation}

\begin{lemma}[Strict binary-prefix transmission coefficient]
\label{lem:V52-prefix-variance}
For every \(s>0\) and \(k>0\),
\begin{equation}
 0<\beta_{s,k}\le1.
\label{eq:V52-beta-positive}
\end{equation}
On the normalized-character scalar functional, the one-coordinate
complete binary prefix covariance
\(T_e^{12}-R_e^{12}\) has value exactly \(\beta_{s,k}\).
\end{lemma}

\begin{proof}
The displayed formula is
\Cref{lem:V51-heat-moments} applied with \(j=k\), followed by the
definition of variance.  The central heat-kernel measure at time
\(s>0\) has a strictly positive density on the connected compact
group \(SU(2)\).  Since \(k>0\), the real character \(x_k\) is not
constant.  Hence its variance is strictly positive.  The pointwise
bound \(|x_k|\le1\) gives
\(\beta_{s,k}\le{\mathbb E}_s x_k^2\le1\).

By \eqref{eq:V10-prefix-covariance}, on a one-coordinate prefix,
\[
 T_e^{12}-R_e^{12}
 =
 {\mathbb E}_s(x_kx_k)
 -({\mathbb E}_sx_k)({\mathbb E}_sx_k),
\]
which is \(\beta_{s,k}\).
\end{proof}

Let
\[
 \mathsf W_{k,j}^{\rm un}
 :=
 (T_e^{12}-R_e^{12})
 \otimes
 \kappa_{3,s}(x_j^2,x_j,x_j).
\]
This is the scalar evaluation of the unique \(r\)-summand of
\eqref{eq:V35-S211} on the present two-coordinate packet; all
singleton prefix and suffix spectators are identities.

\begin{theorem}[Exact survival through the legal prefix]
\label{thm:V52-prefix-transmission}
For every fixed \(s>0\) and \(k>0\), and every
\(j\ge j_0(s)\),
\begin{equation}
 \boxed{
 \mathsf W_{k,j}^{\rm un}
 =
 \beta_{s,k}\,
 \kappa_{3,s}(x_j^2,x_j,x_j)
 \ge
 \frac{\beta_{s,k}}{2d_j^3}>0.}
\label{eq:V52-prefix-transmission}
\end{equation}
Thus the V51 four-row channel is not cancelled by the smallest exact
\(12\) first-connectivity prefix.
\end{theorem}

\begin{proof}
The heat law on the product group at \(e\) and \(r\) is a product
measure.  The exact source word is a tensor product of its complete
binary prefix covariance at \(e\) and its assembled composite
third cumulant at \(r\).  Therefore its scalar evaluation factors.
\Cref{lem:V52-prefix-variance} gives the first factor, and
\Cref{thm:V51-finite-heat-lower} gives the second.
\end{proof}

\begin{corollary}[The normalized signed lower channel is very small]
\label{cor:V52-normalized-signed}
The exact signed scalar packet satisfies
\begin{equation}
 \frac{\mathsf W_{k,j}^{\rm un}}
 {\Xi_1\Xi_2\Xi_3\Xi_4}
 \ge
 \frac{\beta_{s,k}}
 {2d_j^3(A+\alpha)^2A^2},
\label{eq:V52-normalized-lower}
\end{equation}
while its absolute value is at most
\begin{equation}
 \frac6{(A+\alpha)^2A^2}.
\label{eq:V52-normalized-absolute}
\end{equation}
In particular the lower channel in
\eqref{eq:V52-normalized-lower} is of order \(j^{-11}\), not
\(j^{-3}\).
\end{corollary}

\begin{proof}
Use \eqref{eq:V52-row-masses} and
\eqref{eq:V52-prefix-transmission}.  For the upper bound,
\(\beta_{s,k}\le1\), and the five terms in the third-cumulant
moment formula have total absolute coefficient
\(1+2+1+1+1=6\).  Every moment of variables bounded by one has
absolute value at most one.  Finally \(d_j\asymp j\) and
\(A\asymp j^2\).
\end{proof}

\section{Exact evaluation of the legal tree currencies}
\label{sec:V52-tree-evaluation}

The total stocks on this packet are
\begin{align}
 H_{12}&=\alpha^2+A^2,
\label{eq:V52-H12}\\
 H_{13}=H_{14}=H_{23}=H_{24}=H_{34}&=A^2.
\label{eq:V52-other-H}
\end{align}
Define
\begin{equation}
 u:=\frac{\alpha^2+A^2}{(A+\alpha)^2},
 \qquad
 \rho:=\frac{A}{A+\alpha}.
\label{eq:V52-u-rho}
\end{equation}
Then
\begin{equation}
 \theta_{12}=u,\quad
 \theta_{13}=\theta_{14}=\theta_{23}=\theta_{24}=\rho,
 \quad
 \theta_{34}=1.
\label{eq:V52-theta-table}
\end{equation}

\begin{lemma}[All eight total-stock trees]
\label{lem:V52-eight-values}
On the packet above, each of the four trees
\(T_{ij}^{(0)}\) has weight
\begin{equation}
 \prod_{uv\in T_{ij}^{(0)}}\theta_{uv}=u\rho^2,
\label{eq:V52-zero-tree-value}
\end{equation}
and each of the four path trees
\(T_i^{(3)},T_j^{(4)}\) has weight
\begin{equation}
 \prod_{uv\in T_i^{(3)}}\theta_{uv}
 =
 \prod_{uv\in T_j^{(4)}}\theta_{uv}
 =
 u\rho.
\label{eq:V52-path-tree-value}
\end{equation}
If \(A\ge\alpha\), then
\begin{equation}
 \sum_{T\in\mathfrak A_{211}^{(8)}}
 \prod_{uv\in E(T)}\theta_{uv}
 =
 4u\rho^2+4u\rho
 \ge\frac32.
\label{eq:V52-atlas-floor}
\end{equation}
\end{lemma}

\begin{proof}
Substitute \eqref{eq:V52-theta-table} into the exact census
\eqref{eq:V35-eight-trees}.  The four \(T^{(0)}\) trees use the
edge \(12\) and two edges from \(\{1,2\}\) to \(3,4\), proving
\eqref{eq:V52-zero-tree-value}.  Each path tree uses \(12\), one
such crossing, and \(34\), proving
\eqref{eq:V52-path-tree-value}.  The elementary inequality
\(\alpha^2+A^2\ge(A+\alpha)^2/2\) gives \(u\ge1/2\), while
\(A\ge\alpha\) gives \(\rho\ge1/2\).  Hence
\[
 4u\rho^2+4u\rho
 \ge
 4\left(\frac12\right)\left(\frac14\right)
 +4\left(\frac12\right)\left(\frac12\right)
 =\frac32.
\]
\end{proof}

The total-stock expansion permits the \(12\) edge to be taken at the
joining coordinate.  To ensure that this Cartesian enlargement is
not hiding the difficulty, retain instead the actual ordered prefix
mark \(e<r\).  The literal \(T_{11}^{(0)}\) marked weight is
\begin{equation}
 \tau_{11}(e,r)
 :=
 \frac{\alpha^2}{(A+\alpha)^2}
 \frac{A^2}{(A+\alpha)A}
 \frac{A^2}{(A+\alpha)A}
 =
 \frac{\alpha^2A^2}{(A+\alpha)^4}.
\label{eq:V52-ordered-star}
\end{equation}

\begin{lemma}[The strict ordered mark still has spare decay]
\label{lem:V52-ordered-floor}
If \(A\ge\alpha\), then
\begin{equation}
 \tau_{11}(e,r)
 \ge\frac{\alpha^2}{16A^2}.
\label{eq:V52-ordered-floor}
\end{equation}
Consequently
\begin{equation}
 \frac1{(A+\alpha)^2A^2}
 \le
 \frac4{\alpha^2A^2}\tau_{11}(e,r)
 \le
 \frac4{\alpha^2}\tau_{11}(e,r).
\label{eq:V52-normalization-vs-ordered}
\end{equation}
\end{lemma}

\begin{proof}
Since \(A+\alpha\le2A\),
\[
 \tau_{11}(e,r)
 =
 \frac{\alpha^2A^2}{(A+\alpha)^4}
 \ge\frac{\alpha^2}{16A^2}.
\]
For the sharper comparison in
\eqref{eq:V52-normalization-vs-ordered}, divide the left side by
\(\tau_{11}(e,r)\):
\[
 \frac{(A+\alpha)^2}{\alpha^2A^4}
 \le\frac4{\alpha^2A^2}.
\]
The last inequality uses \(A\ge1\).
\end{proof}

\section{Insertion of the actual prefix payer}
\label{sec:V52-payer}

Assign the unique output numerator and resolvent to the complete
binary prefix covariance at \(e\).  Let
\(\mathcal W_{k,j}^{\rm pre}\) denote the resulting physically
compressed two-coordinate word, with the local joining cumulant kept
assembled.

\begin{theorem}[Once-paid ordered routing certificate]
\label{thm:V52-paid-routing}
There is a constant \(C_{s,k}\), independent of \(j\), such that
\begin{equation}
 \boxed{
 \frac{
 \Gamma_\otimes(1,\mathcal W_{k,j}^{\rm pre})}
 {\Xi_1\Xi_2\Xi_3\Xi_4}
 \le
 C_{s,k}\,\tau_{11}(e,r)}
\label{eq:V52-paid-routing}
\end{equation}
whenever \(A\ge\alpha\).  Hence this genuine four-active-row,
two-coordinate family is controlled already by the strictly ordered
star mark.  A fortiori it is controlled by the complete eight-tree
atlas.
\end{theorem}

\begin{proof}
For the one-coordinate prefix, \(H_{12}(<r)=\alpha^2\).
The paid half of \Cref{thm:V10-prefix-response} is the estimate for
the complete physical envelope after the sole output numerator and
resolvent have been assigned; it gives
\[
 \|\mathcal P_e^{12}\|,
 \ \kappa_\otimes(\mathcal P_e^{12})
 \le C\min\{\alpha^2,s^{-1}\}.
\]
This is an actual payer allocation, not a coefficient inserted into
the scalar regression.

The assembled local third cumulant is a signed sum whose total
absolute coefficient is six.  Its three inputs are complete moment
contractions, so its positive physical envelope has norm at most
six.  There are no nontrivial singleton or suffix factors in this
packet.  Tensor submultiplicativity therefore gives
\[
 \Gamma_\otimes(1,\mathcal W_{k,j}^{\rm pre})
 \le
 6C\min\{\alpha^2,s^{-1}\}
 =:C_{s,k}^{(0)}.
\]
Divide by the row-mass product in
\eqref{eq:V52-row-masses}, then apply
\eqref{eq:V52-normalization-vs-ordered}.  This proves
\eqref{eq:V52-paid-routing}, after replacing
\(C_{s,k}^{(0)}\) by \(4C_{s,k}^{(0)}/\alpha^2\).
\end{proof}

\begin{remark}[What was and was not resolved]
\label{rem:V52-payer-resolution}
The exact scalar lower functional
\eqref{eq:V52-prefix-transmission} does not select a unique physical
output constituent of the payer.  Such a selection is unnecessary
for the present decision: the complete, physically compressed paid
prefix theorem controls every resolved output block simultaneously
and therefore rules out this family as a counterpacket.  No paid
scalar lower bound is claimed.
\end{remark}

\begin{corollary}[Resolution of the V51 regression alternative]
\label{cor:V52-regression-resolved}
The V51 equal-spin character family transmits through the legal
prefix, but it does not disprove \(FCQ_{211}\).  The alternative in
\Cref{prog:V51-V52} resolves on this family as a legal-tree routing
certificate, not as a source counterexample.
\end{corollary}

\begin{proof}
Transmission is \Cref{thm:V52-prefix-transmission}.  The genuine paid
word is controlled by a member of the legal atlas by
\Cref{thm:V52-paid-routing}.  Both statements hold on the same
two-coordinate representation assignment.
\end{proof}

\section{A general exclusion of locally saturated regressions}
\label{sec:V52-saturation}

\begin{definition}[Locally \(c\)-saturated four-row joining]
\label{def:V52-local-saturation}
A joining coordinate \(r\) is locally \(c\)-saturated if
\begin{equation}
 a_{i,r}\ge c\Xi_i
 \qquad(1\le i\le4)
\label{eq:V52-local-saturation}
\end{equation}
for some \(0<c\le1\).
\end{definition}

\begin{proposition}[No small-tree regression at a saturated joining]
\label{prop:V52-saturation-no-go}
If \(r\) is locally \(c\)-saturated, then for every pair \(u\ne v\),
\begin{equation}
 \theta_{uv}\ge c^2.
\label{eq:V52-saturated-edge}
\end{equation}
Consequently every three-edge spanning-tree weight is at least
\(c^6\).  In particular, no family with a uniformly saturated
four-row joining coordinate can make all eight \(2+1+1\) tree
currencies tend to zero.
\end{proposition}

\begin{proof}
The joining coordinate is one nonnegative summand of \(H_{uv}\), so
\[
 \theta_{uv}
 =
 \frac{H_{uv}}{\Xi_u\Xi_v}
 \ge
 \frac{a_{u,r}a_{v,r}}{\Xi_u\Xi_v}
 \ge c^2.
\]
Multiply this estimate along the three edges of any spanning tree.
\end{proof}

\begin{assessment}[Upstream consequence]
\label{ass:V52-upstream-consequence}
Any further character regression against the full \(2+1+1\) source
must be genuinely multicoordinate and mass-diluted: after passage to
a subsequence, at least one fixed row must satisfy
\[
 \frac{a_{i,r}}{\Xi_i}\longrightarrow0.
\]
Indeed, if the minimum of the four displayed ratios did not approach
zero along a subsequence, that subsequence would be locally
\(c\)-saturated; then \Cref{prop:V52-saturation-no-go} would supply a
fixed positive legal tree weight.  Finiteness of the row set selects
one fixed diluted row on a further subsequence.  The off-joining mass
must at the same time pass through
complete moment or invariant output channels without exponential heat
suppression.  This is the next concrete construction problem.
\end{assessment}

\section{V52 status and next acquisition}
\label{sec:V52-status}

\begin{theorem}[Third-series V52 result ledger]
\label{thm:V52-results}
Relative to V1--V51, V52 proves:
\begin{enumerate}[label=\textup{(V52.\arabic*)},leftmargin=6.3em]
\item the exact positive finite-heat covariance of the smallest
      nontrivial binary character prefix;
\item exact transmission of the V51 \(d_j^{-3}\) four-row joining
      channel through that legal prefix;
\item the correction separating derivative stocks from coefficients
      of the signed source;
\item exact evaluation of all eight total-stock tree weights on the
      two-coordinate packet;
\item an exact comparison with the stricter ordered
      \(T_{11}^{(0)}\) prefix mark;
\item insertion of the genuine physically compressed binary-prefix
      payer and a uniform once-paid routing certificate;
\item exclusion of the equal-spin packet as a counterexample to
      \(FCQ_{211}\); and
\item a general proof that any future small-tree regression must
      dilute at least one row away from the joining coordinate.
\end{enumerate}
\end{theorem}

\begin{proof}
These are
\Cref{lem:V52-prefix-variance,
thm:V52-prefix-transmission,
fire:V52-stock-coefficient,
lem:V52-eight-values,
lem:V52-ordered-floor,
thm:V52-paid-routing,
cor:V52-regression-resolved,
prop:V52-saturation-no-go}.
\end{proof}

\begin{programme}[V53: the first mass-diluted protected packet]
\label{prog:V52-V53}
\begin{enumerate}[leftmargin=3em]
\item Add the smallest off-joining coordinate bank that makes
      \(a_{i,r}/\Xi_i\) small for a selected row while retaining a
      nonzero four-row joining cumulant.
\item Resolve the off-joining outputs before taking norms and select
      invariant pairings which avoid exponential heat loss.
\item Compute the exact binary prefix and local composite cumulant on
      the same product-state functional; do not multiply either by
      influence stocks.
\item Keep the unique payer in one fixed complete bank and use its
      physically compressed paid theorem before any coordinate
      expansion.
\item Evaluate all eight literal tree currencies, including rerouted
      \(12\) and \(34\) edges at every active coordinate.
\item If every protected dilution creates a rerouted tree, formulate
      the resulting covering lemma; if one avoids every tree, compare
      its complete paid word with \(FCQ_{211}\).
\item No companion audit is due before V55.
\end{enumerate}
\end{programme}

\begin{firewall}[Claim boundary after third-series V52]
\label{fire:V52-claim-boundary}
V52 resolves one complete regression family and proves a necessary
mass-dilution condition for any successor.  It does not prove
\(FCQ_{211}\), close the general \(2+1+1\) source, or close any
remaining quartic source family.  Higher MTA, WUFC, CPM, cutoff
removal, reflection positivity, Osterwalder--Schrader reconstruction,
construction of the continuum Yang--Mills measure, and a uniform
positive continuum spectral gap remain open.  The full Yang--Mills
mass gap has not been proved.
\end{firewall}

\paragraph{Parent-version record.}
V52 was copied byte-for-byte from
\texttt{Renewal\_Yang\_Mills\_Mass\_Gap\_Research\_Notes\_V51.tex}
before this chapter was appended.  The V51 parent had \(25\,813\)
source lines, \(938\,029\) bytes, and SHA--256
\[
 \texttt{85998ad9c785a8da59a10ad87cbd7ec14fa5b2d472620f753fc6ba086dfea59d}.
\]
The next compulsory companion audit is V55.

% =====================================================================
% V53 APPENDIX: THE FIRST MASS-DILUTED PROTECTED PACKET
% =====================================================================

\chapter{V53: The Symmetric Protected-Pair Dilution and Its
Source-Coordinate Tree}
\label{ch:V53-protected-dilution}

\section{The first packet beyond local saturation}
\label{sec:V53-purpose}

V52 proves that a bad sequence must dilute at least one input row away
from the four-row joining coordinate.  The smallest physical way to
add a large row mass without an exponentially small singleton heat
moment is to add a balanced pair and retain its invariant output.
This chapter tests the symmetric realization: a high-spin \(12\)
prefix pair, a fixed-spin four-row joining coordinate, and a
high-spin \(34\) suffix pair.

The packet is genuinely outside the locally saturated class of
\Cref{def:V52-local-saturation}.  Its exact signed character
functional is only polynomially small and remains nonzero.  Even so,
the packet is not a counterexample.  A single source-coordinate path
tree has exactly the missing \(B^{-4}\) normalization and controls
every payer placement without borrowing the large suffix \(34\)
stock.

\begin{assessment}[Nonduplication boundary]
\label{ass:V53-nonduplication}
V46 gave an abstract three-coordinate scalar obstruction after
decoupling the marked mass from all Wilson relations.  V51--V52 gave
a physical equal-spin joining regression without off-joining mass.
The calculation below is different: all three coordinates carry
actual \(SU(2)\) characters, the two large off-joining banks have
protected pair outputs, the exact \(2+1+1\) word is evaluated before
absolute values, and every admissible payer location is checked
against one literal source-coordinate tree.
\end{assessment}

\section{Exact three-coordinate character evaluation}
\label{sec:V53-character}

Fix \(s=s_\alpha>0\).  Choose once and for all a half-integer
\(k\ge j_0(s)\), and put
\begin{equation}
 \alpha:=a_k,\qquad
 \sigma_{s,k}:=
 \kappa_{3,s}(x_k^2,x_k,x_k)>0.
\label{eq:V53-fixed-join}
\end{equation}
The positivity follows from
\Cref{thm:V51-finite-heat-lower}.  For a variable half-integer \(J\),
write
\[
 B:=a_J,\qquad D:=d_J,
\]
and use three ordered coordinates \(e<r<q\):
\begin{center}
\begin{tabular}{@{}c|cccc@{}}
\toprule
coordinate&row \(1\)&row \(2\)&row \(3\)&row \(4\)\\
\midrule
\(e\)&\(J\)&\(J\)&inactive&inactive\\
\(r\)&\(k\)&\(k\)&\(k\)&\(k\)\\
\(q\)&inactive&inactive&\(J\)&\(J\)\\
\bottomrule
\end{tabular}
\end{center}
The coordinate \(e\) is the complete binary \(12\) prefix, \(r\) is
the \(12|3|4\) joining coordinate, and \(q\) lies in the complete
suffix.

Define
\begin{align}
 \mu_{s,J}
 &:={\mathbb E}_s x_J^2
 =
 \frac1{D^2}\sum_{\ell=0}^{2J}
 (2\ell+1)e^{-s\ell(\ell+1)},
\label{eq:V53-mu}\\
 \beta_{s,J}
 &:=
 \mu_{s,J}-e^{-2sJ(J+1)}.
\label{eq:V53-beta}
\end{align}

\begin{lemma}[Two protected-pair factors]
\label{lem:V53-pair-factors}
There is \(J_1(s)\) such that for all \(J\ge J_1(s)\),
\begin{align}
 \frac1{D^2}
 \le\mu_{s,J}
 &\le\frac{C_s}{D^2},
\label{eq:V53-mu-bounds}\\
 \frac1{2D^2}
 \le\beta_{s,J}
 &\le\frac{C_s}{D^2}.
\label{eq:V53-beta-bounds}
\end{align}
\end{lemma}

\begin{proof}
The formula and upper bound for \(\mu_{s,J}\) are
\eqref{eq:V51-second}.  Its \(\ell=0\) summand is \(D^{-2}\), proving
the lower bound.  Variance gives
\(0\le\beta_{s,J}\le\mu_{s,J}\).  Since
\[
 e^{-2sJ(J+1)}=o(D^{-N})
\]
for every \(N\), choose \(J_1(s)\) so that the exponential is at most
\(\tfrac12D^{-2}\).  Then
\(\beta_{s,J}\ge\mu_{s,J}-\tfrac12D^{-2}
\ge\tfrac12D^{-2}\).
\end{proof}

Let \(\mathsf V_{J;k}^{\rm un}\) be the scalar
normalized-character evaluation of the unique \(r\)-summand in the
exact source \eqref{eq:V35-S211} on this packet.  Thus the binary
prefix is kept as its complete covariance, the joining carrier is
the unsplit composite third cumulant, and the suffix is its complete
joint \(34\) moment.

\begin{theorem}[Exact survival under symmetric protected dilution]
\label{thm:V53-exact-survival}
For every \(J\ge J_1(s)\),
\begin{equation}
 \boxed{
 \mathsf V_{J;k}^{\rm un}
 =
 \beta_{s,J}\,\sigma_{s,k}\,\mu_{s,J}.}
\label{eq:V53-exact-factorization}
\end{equation}
In particular,
\begin{equation}
 \frac{\sigma_{s,k}}{2D^4}
 \le
 \mathsf V_{J;k}^{\rm un}
 \le
 \frac{\sigma_{s,k}C_s^2}{D^4}.
\label{eq:V53-exact-two-sided}
\end{equation}
Hence the exact signed source is positive and has polynomial, rather
than exponential, high-spin decay.
\end{theorem}

\begin{proof}
The three coordinate heat laws are independent.  At \(e\), the
complete joint \(12\) moment minus its independent-replica moment is
\(\beta_{s,J}\) by
\Cref{lem:V52-prefix-variance}.  At \(r\), the assembled joining
cumulant is \(\sigma_{s,k}\) by
\eqref{eq:V53-fixed-join}.  At \(q\), rows \(1,2\) are inactive and
the complete suffix moment on rows \(3,4\) is
\({\mathbb E}_sx_J^2=\mu_{s,J}\).  Tensor-product evaluation gives
\eqref{eq:V53-exact-factorization}.  Apply
\Cref{lem:V53-pair-factors}.
\end{proof}

\begin{firewall}[Protected does not mean unsuppressed]
\label{fire:V53-protected-polynomial}
The invariant constituent prevents exponential heat loss, but the
normalized character trace still contributes \(D^{-2}\) for each
protected pair.  It is therefore false to replace either
\(\mu_{s,J}\) or \(\beta_{s,J}\) by a representation-independent
positive constant in the signed scalar functional.
\end{firewall}

\section{Row masses, stocks, and the eight-tree census}
\label{sec:V53-tree-census}

All four rows have the same total mass:
\begin{equation}
 \Xi_1=\Xi_2=\Xi_3=\Xi_4=B+\alpha.
\label{eq:V53-Xi}
\end{equation}
The pair stocks are
\begin{align}
 H_{12}=H_{34}&=B^2+\alpha^2,
\label{eq:V53-pair-stocks}\\
 H_{13}=H_{14}=H_{23}=H_{24}&=\alpha^2.
\label{eq:V53-cross-stocks}
\end{align}
Put
\begin{equation}
 u_B:=\frac{B^2+\alpha^2}{(B+\alpha)^2},
 \qquad
 \varepsilon_B:=\frac{\alpha^2}{(B+\alpha)^2}.
\label{eq:V53-u-epsilon}
\end{equation}
Then
\begin{equation}
 \theta_{12}=\theta_{34}=u_B,
 \qquad
 \theta_{13}=\theta_{14}=\theta_{23}=\theta_{24}
 =\varepsilon_B.
\label{eq:V53-theta}
\end{equation}

\begin{lemma}[Exact total-stock atlas on the diluted packet]
\label{lem:V53-atlas}
The four star-family terms \(T_{ij}^{(0)}\) have weight
\begin{equation}
 u_B\varepsilon_B^2,
\label{eq:V53-star-total}
\end{equation}
and the four path-family terms
\(T_i^{(3)},T_j^{(4)}\) have weight
\begin{equation}
 u_B^2\varepsilon_B.
\label{eq:V53-path-total}
\end{equation}
Therefore
\begin{equation}
 \sum_{T\in\mathfrak A_{211}^{(8)}}
 \prod_{uv\in E(T)}\theta_{uv}
 =
 4u_B\varepsilon_B^2+4u_B^2\varepsilon_B.
\label{eq:V53-atlas-sum}
\end{equation}
For \(B\ge\alpha\), each path-family term is at least
\begin{equation}
 u_B^2\varepsilon_B
 \ge\frac{\alpha^2}{16B^2}.
\label{eq:V53-total-path-floor}
\end{equation}
\end{lemma}

\begin{proof}
Substitute \eqref{eq:V53-theta} into the tree list
\eqref{eq:V35-eight-trees}.  A \(T^{(0)}\) term uses \(12\) and two
crossing edges; a path term uses \(12\), \(34\), and one crossing
edge.  This proves the exact values and their sum.

As in V52, \(u_B\ge1/2\).  If \(B\ge\alpha\), then
\((B+\alpha)^2\le4B^2\), so
\(\varepsilon_B\ge\alpha^2/(4B^2)\).  Hence
\[
 u_B^2\varepsilon_B
 \ge\frac14\frac{\alpha^2}{4B^2}.
\]
\end{proof}

The total-stock path in
\eqref{eq:V53-path-total} can use the large suffix \(34\) stock.
That is a legal Cartesian enlargement, but it is not needed.  Mark
the \(12\) edge at the actual prefix coordinate \(e\), and mark both
quotient edges of \(T_1^{(3)}=\{12,13,34\}\) at the actual joining
coordinate \(r\).  The resulting source-coordinate weight is
\begin{equation}
 \tau_{\rm src}(J;k)
 :=
 \frac{B^2}{(B+\alpha)^2}
 \frac{\alpha^2}{(B+\alpha)^2}
 \frac{\alpha^2}{(B+\alpha)^2}
 =
 \frac{B^2\alpha^4}{(B+\alpha)^6}.
\label{eq:V53-source-tree}
\end{equation}

\begin{lemma}[The source-coordinate path has the exact fourth power]
\label{lem:V53-source-tree}
If \(B\ge\alpha\), then
\begin{equation}
 \tau_{\rm src}(J;k)
 \ge\frac{\alpha^4}{64B^4},
\label{eq:V53-source-tree-floor}
\end{equation}
and
\begin{equation}
 \frac1{(B+\alpha)^4}
 \le
 \frac4{\alpha^4}\,
 \tau_{\rm src}(J;k).
\label{eq:V53-denominator-tree}
\end{equation}
\end{lemma}

\begin{proof}
The first bound follows from \(B+\alpha\le2B\):
\[
 \frac{B^2\alpha^4}{(B+\alpha)^6}
 \ge\frac{B^2\alpha^4}{64B^6}.
\]
For the second, divide the left side by
\(\tau_{\rm src}(J;k)\):
\[
 \frac{(B+\alpha)^2}{B^2\alpha^4}
 \le\frac4{\alpha^4}.
\]
\end{proof}

\begin{remark}[Why the small total-stock path was not essential]
\label{rem:V53-no-suffix-borrow}
The marked edge \(34\) in
\eqref{eq:V53-source-tree} uses \(h_{34,r}=\alpha^2\), not the large
suffix contribution \(h_{34,q}=B^2\).  Thus the control below is
valid even before the Cartesian enlargement
\(h_{34,r}\le H_{34}\).  The same local quotient-tree coordinate
which defines the exact joining carrier already supplies the needed
currency.
\end{remark}

\section{All genuine payer locations}
\label{sec:V53-payers}

For each admissible payer placement \(\pi\), let
\(\mathcal V_{J;k}^{(\pi)}\) denote the complete physically
compressed word on the three-coordinate packet.  The nontrivial
locations are the binary prefix, the assembled joining cumulant, and
the complete suffix.  The prefix singleton banks on rows \(3,4\) are
empty on this packet.

\begin{lemma}[Uniform paid envelope on every branch]
\label{lem:V53-paid-envelope}
There is \(C_{s,k}<\infty\), independent of \(J\) and of the payer
placement \(\pi\), such that
\begin{equation}
 \Gamma_\otimes(1,\mathcal V_{J;k}^{(\pi)})
 \le C_{s,k}.
\label{eq:V53-paid-envelope}
\end{equation}
\end{lemma}

\begin{proof}
If the binary prefix pays, the paid part of
\Cref{thm:V10-prefix-response} gives
\[
 \|\mathcal P_{12}^{<r}\|
 \le C\min\{B^2,s^{-1}\}\le C/s.
\]
The joining positive envelope has norm at most six and the suffix is
a complete moment contraction.

If the joining carrier pays, use the archived f2 local estimate
\[
 \|\mathcal P_r^{12|3|4}\|
 \le Cs\,Q_r^{12|3|4},
\]
recorded there as
\(\texttt{eq:V67-local-ternary-paid}\).  Its quotient stocks at \(r\)
are fixed functions of \(\alpha\), so its norm is bounded by a
constant depending only on \(s,k\).  The unpaid complete prefix
covariance has norm at most two, and the suffix remains a
contraction.

If the complete suffix pays, use the complete-bank payer estimate in
the V67 ledger, equivalently the universal paid complete-bank bound
\eqref{eq:V33-paid-universal}.  It costs at most \(C/s\),
independently of \(J\).  The other two factors have the same bounds
as above.  An empty prefix singleton bank has zero allocated output
mass and produces no additional branch.  These cases exhaust the
V67 payer ledger for the present packet.
\end{proof}

\begin{theorem}[Every payer is controlled by one actual-coordinate
path]
\label{thm:V53-all-payer-routing}
For all \(J\) with \(B\ge\alpha\), and for every admissible payer
placement \(\pi\),
\begin{equation}
 \boxed{
 \frac{
 \Gamma_\otimes(1,\mathcal V_{J;k}^{(\pi)})}
 {\Xi_1\Xi_2\Xi_3\Xi_4}
 \le
 C_{s,k}'\,\tau_{\rm src}(J;k).}
\label{eq:V53-all-payer-routing}
\end{equation}
The constant is independent of \(J\) and of \(\pi\).
\end{theorem}

\begin{proof}
By \eqref{eq:V53-Xi} and
\Cref{lem:V53-paid-envelope}, the left side is at most
\[
 \frac{C_{s,k}}{(B+\alpha)^4}.
\]
Apply \eqref{eq:V53-denominator-tree}.  The finite number of payer
types is already absorbed into \(C_{s,k}\).
\end{proof}

\begin{corollary}[The first protected dilution is not a counterpacket]
\label{cor:V53-not-counterpacket}
The symmetric protected-pair dilution has a positive polynomial
signed character channel, but it satisfies the once-paid
\(2+1+1\) marked-tree target at every payer location.  It therefore
does not disprove \(FCQ_{211}\) or quartic MTA.
\end{corollary}

\begin{proof}
The first assertion is
\Cref{thm:V53-exact-survival}; the second is
\Cref{thm:V53-all-payer-routing}.  The controlling tree belongs to
the exact eight-tree family by
\eqref{eq:V35-tree-three}.
\end{proof}

\section{What the packet teaches about the remaining search}
\label{sec:V53-consequence}

\begin{proposition}[Pair protection creates its own denominator
payment]
\label{prop:V53-self-payment}
On the symmetric packet, the four row denominators contribute
\((B+\alpha)^{-4}\).  The literal source-coordinate path
\(\{12,13,34\}\) recovers exactly the same high-mass power:
\begin{equation}
 \tau_{\rm src}(J;k)\asymp_{k}B^{-4}.
\label{eq:V53-self-payment}
\end{equation}
The large prefix pair supplies the \(B^2\) numerator, while the two
extra incidences at the fixed joining coordinate supply
\(\alpha^4\).  No character-dimension estimate is required for this
positive routing.
\end{proposition}

\begin{proof}
This is the two-sided estimate obtained directly from
\eqref{eq:V53-source-tree} when \(B\ge\alpha\):
\[
 \frac{\alpha^4}{64B^4}
 \le\tau_{\rm src}(J;k)
 \le\frac{\alpha^4}{B^4}.
\]
The factor attribution follows from its three marked numerators
\(h_{12,e}=B^2\),
\(h_{13,r}=\alpha^2\), and
\(h_{34,r}=\alpha^2\).
\end{proof}

\begin{assessment}[Refined shape of a possible bad packet]
\label{ass:V53-refined-bad-packet}
A successor cannot merely move large mass from the joining coordinate
into two protected complementary pairs.  To make the source-coordinate
path in \eqref{eq:V53-source-tree} smaller than the four-row
normalization, it must also deplete at least one of its actual
numerators:
\[
 h_{12,e},\qquad h_{i3,r},\qquad h_{34,r},
\]
for both \(i=1,2\), and must do so without making a star
\(\{12,i3,j4\}\) or the symmetric path
\(\{12,j4,34\}\) large.  This is a finite covering problem on the
eight source trees coupled to the locations carrying the four row
masses.
\end{assessment}

\section{V53 status and next acquisition}
\label{sec:V53-status}

\begin{theorem}[Third-series V53 result ledger]
\label{thm:V53-results}
Relative to V1--V52, V53 proves:
\begin{enumerate}[label=\textup{(V53.\arabic*)},leftmargin=6.3em]
\item an exact physical three-coordinate protected-pair character
      packet outside the locally saturated class;
\item two-sided \(D^{-4}\) bounds for its nonzero signed
      \(2+1+1\) source functional;
\item the exact eight-tree census on its diluted row masses;
\item an actual-coordinate path tree which avoids borrowing the large
      suffix stock;
\item a uniform positive-envelope bound for every genuine payer
      placement;
\item a once-paid routing certificate uniform in the large spin;
\item exclusion of symmetric complementary-pair dilution as an
      \(FCQ_{211}\) counterpacket; and
\item reduction of the next search to a finite simultaneous depletion
      problem for all eight source trees.
\end{enumerate}
\end{theorem}

\begin{proof}
These are
\Cref{lem:V53-pair-factors,
thm:V53-exact-survival,
lem:V53-atlas,
lem:V53-source-tree,
lem:V53-paid-envelope,
thm:V53-all-payer-routing,
cor:V53-not-counterpacket,
prop:V53-self-payment,
ass:V53-refined-bad-packet}.
\end{proof}

\begin{programme}[V54: finite tree-covering theorem for diluted mass]
\label{prog:V53-V54}
\begin{enumerate}[leftmargin=3em]
\item Encode each row mass by its prefix, joining, singleton, and
      suffix locations, with banks already used by the payer kept
      distinct.
\item For each of the eight trees, retain the coordinate-restricted
      numerator before replacing it by a total \(H_{uv}\).
\item Prove or refute the finite alternative: either one literal tree
      pays the four normalized row masses, or a fixed row has two
      depleted incidences in every admissible tree.
\item In the latter branch, return to the V49 complete quadratic
      source and test whether its two actual replacement differences
      yield the missing pair of depleted-row debits.
\item Keep protected pair outputs assembled; do not convert their
      invariant multiplicity into a coordinate or representation
      count.
\item No companion audit is due before V55.
\end{enumerate}
\end{programme}

\begin{firewall}[Claim boundary after third-series V53]
\label{fire:V53-claim-boundary}
V53 closes one mass-diluted physical regression family, not the
general \(2+1+1\) source.  It does not prove \(FCQ_{211}\), the
finite tree-covering alternative proposed for V54, or any remaining
quartic source family.  Higher MTA, WUFC, CPM, cutoff removal,
reflection positivity, Osterwalder--Schrader reconstruction,
construction of the continuum Yang--Mills measure, and a uniform
positive continuum spectral gap remain open.  The full Yang--Mills
mass gap has not been proved.
\end{firewall}

\paragraph{Parent-version record.}
V53 was copied byte-for-byte from
\texttt{Renewal\_Yang\_Mills\_Mass\_Gap\_Research\_Notes\_V52.tex}
before this chapter was appended.  The V52 parent had \(26\,340\)
source lines, \(954\,971\) bytes, and SHA--256
\[
 \texttt{a10dde015d9ffa2b0dadb35032ce1fa601939b34a2d14971e7ea97bc2e7f1cda}.
\]
The next compulsory companion audit is V55.

% =====================================================================
% V54 APPENDIX: THE QUOTIENT-CUT COVERING THEOREM
% =====================================================================

\chapter{V54: An Exact Quotient-Cut Covering Theorem for the
Eight \(2+1+1\) Trees}
\label{ch:V54-cut-covering}

\section{Why the eight-tree list should be compressed}
\label{sec:V54-purpose}

V53 showed directly that symmetric protected-pair dilution creates a
source-coordinate path which pays the four row denominators.  The
next question is whether this was an accident of a symmetric
three-coordinate packet or a consequence of the topology of the
complete eight-tree atlas.

This chapter gives the exact finite answer at the level of normalized
edge currencies.  The eight terms compress to a binary-prefix factor
times the spanning-tree polynomial of the three quotient vertices
\[
 \{12\},\quad\{3\},\quad\{4\}.
\]
The second-largest quotient connection then gives a sharp covering
parameter.  A small full atlas has exactly four possible depletion
faces: the prefix edge is small, the composite block is cut from both
singletons, row \(3\) is cut, or row \(4\) is cut.

\begin{assessment}[Nonduplication boundary]
\label{ass:V54-nonduplication}
\Cref{lem:V35-eight-trees} enumerated the eight labelled trees, and
\Cref{lem:V35-demand-census} counted their row degrees.  Neither result
gave a quantitative converse describing every way in which all eight
weights can vanish simultaneously.  V54 proves that converse and
uses it to replace an unconstrained mass-dilution search by four
exhaustive cut faces.
\end{assessment}

\section{Normalized edge space}
\label{sec:V54-edge-space}

For nonnegative row weights, recall
\[
 \theta_{uv}:=\frac{H_{uv}}{\Xi_u\Xi_v}.
\]

\begin{lemma}[Every normalized edge lies in the unit interval]
\label{lem:V54-edge-unit}
For every \(u\ne v\),
\begin{equation}
 0\le\theta_{uv}\le1.
\label{eq:V54-edge-unit}
\end{equation}
\end{lemma}

\begin{proof}
Nonnegativity is immediate.  Also
\[
 H_{uv}
 =
 \sum_e a_{u,e}a_{v,e}
 \le
 \left(\sum_ea_{u,e}\right)
 \left(\sum_fa_{v,f}\right)
 =
 \Xi_u\Xi_v,
\]
because the product on the middle right contains every diagonal
summand and additional nonnegative off-diagonal summands.
\end{proof}

Introduce the four quotient variables
\begin{align}
 p&:=\theta_{12},
\label{eq:V54-p}\\
 x&:=\theta_{13}+\theta_{23},
\label{eq:V54-x}\\
 y&:=\theta_{14}+\theta_{24},
\label{eq:V54-y}\\
 z&:=\theta_{34}.
\label{eq:V54-z}
\end{align}
Here \(p\) is the normalized binary-prefix stock, while \(x,y,z\)
are the three edge weights on the quotient graph
\(\{12\}|3|4\).  By \Cref{lem:V54-edge-unit},
\[
 0\le p,z\le1,\qquad 0\le x,y\le2.
\]

Let
\begin{equation}
 \mathscr A_{211}
 :=
 \sum_{T\in\mathfrak A_{211}^{(8)}}
 \prod_{uv\in E(T)}\theta_{uv}
\label{eq:V54-atlas-definition}
\end{equation}
be the complete normalized eight-tree polynomial.

\begin{theorem}[Exact quotient-tree factorization]
\label{thm:V54-exact-factorization}
One has the identity
\begin{equation}
 \boxed{
 \mathscr A_{211}
 =
 p(xy+xz+yz).}
\label{eq:V54-exact-factorization}
\end{equation}
The three terms in parentheses are precisely the three spanning trees
of the quotient triangle with edge weights \(x,y,z\).
\end{theorem}

\begin{proof}
The four trees \(T_{ij}^{(0)}=\{12,i3,j4\}\) contribute
\[
 p(\theta_{13}+\theta_{23})
  (\theta_{14}+\theta_{24})
 =pxy.
\]
The two trees \(T_i^{(3)}=\{12,i3,34\}\) contribute \(pxz\), and
the two trees \(T_j^{(4)}=\{12,j4,34\}\) contribute \(pyz\).
These are all eight trees by
\Cref{lem:V35-eight-trees}.
\end{proof}

\section{The second-largest connection}
\label{sec:V54-second-largest}

For three nonnegative numbers \(x,y,z\), let
\[
 m_2(x,y,z)
\]
denote their second-largest value, with multiplicities retained.

\begin{lemma}[Second-largest lower bound]
\label{lem:V54-second-largest}
For all \(x,y,z\ge0\),
\begin{equation}
 xy+xz+yz\ge m_2(x,y,z)^2.
\label{eq:V54-second-largest}
\end{equation}
Consequently
\begin{equation}
 \boxed{
 \mathscr A_{211}
 \ge p\,m_2(x,y,z)^2.}
\label{eq:V54-atlas-cover}
\end{equation}
\end{lemma}

\begin{proof}
Let \(m_1\ge m_2\ge m_3\) be the decreasing rearrangement of
\(x,y,z\).  The sum \(xy+xz+yz\) contains the product
\(m_1m_2\), so it is at least \(m_1m_2\ge m_2^2\).
Use \eqref{eq:V54-exact-factorization}.
\end{proof}

\begin{theorem}[Quantitative four-face alternative]
\label{thm:V54-four-face}
Fix \(0<\delta\le1\).  If
\begin{equation}
 \mathscr A_{211}<\delta^3,
\label{eq:V54-small-atlas}
\end{equation}
then at least one of the following four alternatives holds:
\begin{enumerate}[label=\textup{(\Alph*)},leftmargin=4.5em]
\item \(p<\delta\);
      \label{item:V54-prefix-face}
\item \(x<\delta\) and \(y<\delta\);
      \label{item:V54-composite-face}
\item \(x<\delta\) and \(z<\delta\);
      \label{item:V54-row3-face}
\item \(y<\delta\) and \(z<\delta\).
      \label{item:V54-row4-face}
\end{enumerate}
Equivalently, if \(p\ge\delta\) and no two of \(x,y,z\) are below
\(\delta\), then \(\mathscr A_{211}\ge\delta^3\).
\end{theorem}

\begin{proof}
Suppose item \ref{item:V54-prefix-face} fails, so \(p\ge\delta\).
By \eqref{eq:V54-atlas-cover} and
\eqref{eq:V54-small-atlas},
\[
 \delta\,m_2(x,y,z)^2
 \le p\,m_2(x,y,z)^2
 \le\mathscr A_{211}
 <\delta^3.
\]
Thus \(m_2(x,y,z)<\delta\).  At least two of \(x,y,z\) are then
strictly below \(\delta\).  The three possible pairs give items
\ref{item:V54-composite-face}--%
\ref{item:V54-row4-face}.
\end{proof}

\begin{corollary}[Sequential cut classification]
\label{cor:V54-sequential-cuts}
Let \(\{\theta_{uv}^{(n)}\}\) be a sequence of normalized edge
currencies for which
\(\mathscr A_{211}^{(n)}\to0\).  After passage to a subsequence, one
of the following four statements may be selected:
\begin{align}
 p_n&\longrightarrow0;
\label{eq:V54-prefix-limit}\\
 x_n\longrightarrow0,\quad y_n&\longrightarrow0;
\label{eq:V54-composite-limit}\\
 x_n\longrightarrow0,\quad z_n&\longrightarrow0;
\label{eq:V54-row3-limit}\\
 y_n\longrightarrow0,\quad z_n&\longrightarrow0.
\label{eq:V54-row4-limit}
\end{align}
The cases need not be disjoint, but their union is exhaustive.
\end{corollary}

\begin{proof}
Pass first to a subsequence on which either \(p_n\to0\) or
\(p_n\ge c>0\).  In the first case use
\eqref{eq:V54-prefix-limit}.  In the second,
\eqref{eq:V54-atlas-cover} gives
\(m_2(x_n,y_n,z_n)\to0\).  At least two of the three variables tend
to zero after a further subsequence.  There are only three possible
pairs, giving the remaining alternatives.
\end{proof}

\section{Graph meaning of the four faces}
\label{sec:V54-graph-meaning}

\begin{proposition}[The depletion faces are quotient cuts]
\label{prop:V54-cut-meaning}
The four cases in
\Cref{thm:V54-four-face} have the following exact graph meaning:
\begin{enumerate}[label=\textup{(\Alph*)},leftmargin=4.5em]
\item \(p\) depletes the required binary prefix edge \(12\);
\item \(x,y\) deplete both quotient edges incident to the composite
      vertex \(\{12\}\);
\item \(x,z\) deplete both quotient edges incident to the singleton
      vertex \(3\);
\item \(y,z\) deplete both quotient edges incident to the singleton
      vertex \(4\).
\end{enumerate}
Thus every small-atlas sequence exposes either the prefix edge or one
two-edge cut of the quotient triangle.
\end{proposition}

\begin{proof}
By definition, \(x\) is the total quotient edge between
\(\{12\}\) and \(3\), \(y\) is the edge between \(\{12\}\) and
\(4\), and \(z\) is the edge between \(3\) and \(4\).  Deleting two
edges incident to one quotient vertex isolates precisely that vertex.
The factor \(p\) lies outside the quotient triangle and is required by
all eight lifted trees.
\end{proof}

\begin{remark}[No cancellation or analytic estimate is used]
\label{rem:V54-pure-covering}
The cut theorem is a nonnegative polynomial identity and a
three-number ordering argument.  It does not estimate the source
operator, assume a thermal bound, use a row-mass moment, or take an
absolute value of a signed Wilson word.  It is therefore an exact
routing constraint on any future analytic regression.
\end{remark}

\section{Compatibility with joining-coordinate fractions}
\label{sec:V54-local-fractions}

At a selected joining coordinate \(r\), put
\begin{equation}
 \rho_i(r):=\frac{a_{i,r}}{\Xi_i}\in[0,1].
\label{eq:V54-rho}
\end{equation}

\begin{lemma}[Joining fractions feed the quotient edges]
\label{lem:V54-fraction-edges}
One has
\begin{align}
 p&\ge\rho_1\rho_2,
\label{eq:V54-rho-p}\\
 x&\ge\rho_3(\rho_1+\rho_2),
\label{eq:V54-rho-x}\\
 y&\ge\rho_4(\rho_1+\rho_2),
\label{eq:V54-rho-y}\\
 z&\ge\rho_3\rho_4.
\label{eq:V54-rho-z}
\end{align}
\end{lemma}

\begin{proof}
For every pair \(u,v\), the coordinate \(r\) contributes
\(a_{u,r}a_{v,r}\) to \(H_{uv}\).  Hence
\(\theta_{uv}\ge\rho_u\rho_v\).  The four inequalities follow after
adding the appropriate crossing pairs in \(x\) and \(y\).
\end{proof}

\begin{corollary}[Typed row dilution on each nonprefix face]
\label{cor:V54-typed-dilution}
Consider a sequence on one fixed cut face.
\begin{enumerate}[label=\textup{(\roman*)},leftmargin=4em]
\item On the composite face \(x,y\to0\), if either
      \(\rho_1+\rho_2\) stays bounded below, then both
      \(\rho_3,\rho_4\to0\).
\item On the row-\(3\) face \(x,z\to0\), if \(\rho_3\) stays bounded
      below, then
      \(\rho_1,\rho_2,\rho_4\to0\).
\item On the row-\(4\) face \(y,z\to0\), if \(\rho_4\) stays bounded
      below, then
      \(\rho_1,\rho_2,\rho_3\to0\).
\end{enumerate}
\label{eq:V54-typed-dilution}
\end{corollary}

\begin{proof}
For item \textup{(i)}, use
\eqref{eq:V54-rho-x}--\eqref{eq:V54-rho-y}.
For item \textup{(ii)}, a positive lower bound on \(\rho_3\) and
\eqref{eq:V54-rho-x} force
\(\rho_1+\rho_2\to0\), while
\eqref{eq:V54-rho-z} forces \(\rho_4\to0\).
Item \textup{(iii)} is symmetric.
\end{proof}

\begin{firewall}[One-sided implications only]
\label{fire:V54-no-converse}
The bounds in \Cref{lem:V54-fraction-edges} are lower bounds.
Small joining fractions do not imply small total edges, because
prefix, singleton, or suffix coordinates may contribute to
\(H_{uv}\).  Those off-joining contributions are precisely the
rerouting mechanism seen in V53 and must remain in the full atlas.
\end{firewall}

\section{Consequences for physical regression design}
\label{sec:V54-regression}

\begin{theorem}[Necessary shape of any bounded-envelope
counterpacket]
\label{thm:V54-counterpacket-shape}
Consider any sequence of complete \(2+1+1\) packets at fixed
\(s>0\) for which:
\begin{enumerate}[label=\textup{(\roman*)},leftmargin=4em]
\item the once-paid normalized positive envelopes are uniformly
      bounded; and
\item their ratio to \(\mathscr A_{211}\) is unbounded.
\end{enumerate}
Then \(\mathscr A_{211}\to0\) along a subsequence, and after a further
subsequence the packets lie asymptotically on one of the four cut
faces in \Cref{cor:V54-sequential-cuts}.
\end{theorem}

\begin{proof}
Let \(G_n\) be the normalized paid envelope and suppose
\(G_n/\mathscr A_{211}^{(n)}\to\infty\) along a subsequence.  If the
atlas were bounded below there, uniform boundedness of \(G_n\) would
bound the ratio, a contradiction.  Thus pass to a subsequence with
\(\mathscr A_{211}^{(n)}\to0\), and apply
\Cref{cor:V54-sequential-cuts}.
\end{proof}

\begin{assessment}[The analytic problem after the covering theorem]
\label{ass:V54-analytic-residual}
The finite tree census is no longer the bottleneck.  A proof or
counterexample to \(FCQ_{211}\) can be sought separately on four
typed regions:
\begin{enumerate}[label=\textup{(\Alph*)},leftmargin=4.5em]
\item a depleted complete binary prefix;
\item a two-edge cut around the composite block \(\{12\}\);
\item a two-edge cut around singleton row \(3\); or
\item a two-edge cut around singleton row \(4\).
\end{enumerate}
The V49 complete square has two actual copies of the relevant source
differences, but V49 did not prove that their compressed capacities
pay the two depleted edges.  That cut-specific quadratic response,
with the sole payer fixed across both copies, is the first remaining
analytic acquisition.
\end{assessment}

\section{V54 status and compulsory next audit}
\label{sec:V54-status}

\begin{theorem}[Third-series V54 result ledger]
\label{thm:V54-results}
Relative to V1--V53, V54 proves:
\begin{enumerate}[label=\textup{(V54.\arabic*)},leftmargin=6.3em]
\item the unit-interval bound for every normalized edge currency;
\item exact factorization of the eight-tree atlas into the prefix
      edge and the quotient-triangle tree polynomial;
\item the second-largest-connection lower bound;
\item a quantitative four-face alternative at every threshold;
\item an exhaustive sequential classification of vanishing atlases;
\item the precise quotient-cut meaning of all four faces;
\item lower bounds connecting joining-coordinate row fractions to
      quotient edges;
\item typed row-dilution consequences on the three nonprefix faces;
      and
\item reduction of every bounded-envelope regression to one of four
      cut-specific quadratic response problems.
\end{enumerate}
\end{theorem}

\begin{proof}
These are
\Cref{lem:V54-edge-unit,
thm:V54-exact-factorization,
lem:V54-second-largest,
thm:V54-four-face,
cor:V54-sequential-cuts,
prop:V54-cut-meaning,
lem:V54-fraction-edges,
cor:V54-typed-dilution,
thm:V54-counterpacket-shape}.
\end{proof}

\begin{programme}[V55: compulsory companion audit and first
cut-specific response]
\label{prog:V54-V55}
\begin{enumerate}[leftmargin=3em]
\item Audit the newest active Standard-Model and Navier--Stokes notes,
      record exact hashes, and apply the companion type firewall.
\item Check whether either programme has acquired a genuine
      cut-specific square estimate rather than a covering or
      first-moment analogy.
\item Return first to the prefix face \(p\to0\), where the exact V44
      doubled-prefix identity is already available.
\item Keep the joining cumulant, suffix, output compression, and one
      payer inside the complete V49 square while testing whether the
      prefix response can be factored without an off-diagonal loss.
\item If the prefix face still fails, move upstream to its exact
      output-block commutator rather than repeating the positive
      prefix telescope.
\end{enumerate}
\end{programme}

\begin{firewall}[Claim boundary after third-series V54]
\label{fire:V54-claim-boundary}
V54 proves an exact finite covering theorem, not any of the four
cut-specific analytic response estimates.  It does not prove
\(FCQ_{211}\), close the general \(2+1+1\) source, or close any
remaining quartic source family.  Higher MTA, WUFC, CPM, cutoff
removal, reflection positivity, Osterwalder--Schrader reconstruction,
construction of the continuum Yang--Mills measure, and a uniform
positive continuum spectral gap remain open.  The full Yang--Mills
mass gap has not been proved.
\end{firewall}

\paragraph{Parent-version record.}
V54 was copied byte-for-byte from
\texttt{Renewal\_Yang\_Mills\_Mass\_Gap\_Research\_Notes\_V53.tex}
before this chapter was appended.  The V53 parent had \(26\,867\)
source lines, \(972\,034\) bytes, and SHA--256
\[
 \texttt{e0ec416d1eb23c71f8b7f8d6edb277a7c0193dbfde28de3c00e47258c6ef8c12}.
\]
The next compulsory companion audit is V55.

% =====================================================================
% V55 APPENDIX: COMPANION AUDIT AND PREFIX-CUT CLOSURE
% =====================================================================

\chapter{V55: Companion Audit and Closure of the Pure
Prefix-Depletion Face}
\label{ch:V55-audit-prefix}

\section{Compulsory companion audit}
\label{sec:V55-companion-audit}

The five-version checkpoint required after V50 was performed after
V54.  During the audit the parallel Standard-Model programme advanced
from V33 to V34, so the later complete file was used.  The exact
snapshots inspected were:
\begin{center}
\small
\begin{tabular}{@{}p{0.43\linewidth}rrp{0.29\linewidth}@{}}
\toprule
file&lines&bytes&SHA--256\\
\midrule
\texttt{Renewal\_SM\_Einstein\_Continuum\_Research\_Notes\_V34.tex}
&\(12\,696\)&\(471\,805\)&
\texttt{40410012eaeb74f25c5229d59ecad51f39c398d8c64e2ae46a9b2318cdc2e694}\\
\texttt{Renewal\_NS\_Stability\_Research\_Notes\_V31.tex}
&\(23\,052\)&\(887\,537\)&
\texttt{f1121b62238ad5491bb7cf03635ea9934942edf573ed8faaa9ee79e1d38a6696}\\
\bottomrule
\end{tabular}
\end{center}

The NS file is unchanged from the V50 audit.  Its terminal sharpness
tests still show that a flat stopping mark, heat lift, or
polarization hinge does not create the missing positive first moment.
Any gain must remain in the complete signed nonlinear correlation
before the source norm is taken.

The SM file contains four new stages beyond the V29 snapshot used at
V50:
\begin{enumerate}[label=\textup{(\roman*)},leftmargin=4em]
\item V30 retains a uniformly convex correlated Higgs parent and
      writes its covariance as the exact Helffer--Sjostrand
      one-form resolvent between two actual gradients.  A
      Combes--Thomas conjugation gives exponential spatial decay.
\item V31 obtains gap-free local Yukawa coefficient moments from the
      convex parent, including the quadratic exponential estimate
      later used for rare cells.
\item V32 keeps the antisymmetric Wick determinant intact, uses Gram
      bounds for every forest-interpolated minor, and extracts
      connected cell support without a factorial loss.
\item V33 differentiates the hard covariance through the exact
      resolvent identities.  At second order it retains both ordered
      two-insertion words and the genuine second-derivative contact
      word; marked complementary minors remain Gram controlled.
\item V34 makes an exact smooth good/tail split of the Yukawa
      multiplication operator.  The good determinant is zero-free,
      promotion of its modulus preserves convexity, and every actual
      tail cell owns one simultaneous exponential payment.
\end{enumerate}

\begin{theorem}[V55 companion transfer decision]
\label{thm:V55-transfer-decision}
No analytic SM or NS estimate transfers to a Yang--Mills cut face.
The audit enforces the following proof order:
\begin{enumerate}[label=\textup{(\alph*)},leftmargin=4em]
\item a second cut debit must come from an exact quadratic word,
      including all ordered and contact contributions, not from
      squaring a one-debit positive estimate;
\item a good/tail or hot/thermal split must be an exact partition of
      the original word, and each exceptional occurrence may pay
      only once;
\item distinct connection mechanisms must be assembled into one
      typed spanning tree before exponentiation or absolute values;
      and
\item an assumed strict gap or convexity may be propagated but may
      not be cited as the mechanism which creates the missing
      Yang--Mills response.
\end{enumerate}
\end{theorem}

\begin{proof}
The SM one-form inverse acts on Higgs field gradients under the
hypothesis \(D^2S\ge\kappa>0\); the SM determinant and Gram minors are
fermionic objects; and the NS estimates concern Fourier-shell
Duhamel and coarea operators.  None is the Wilson product-state
capacity of the V49 composite square, so no numerical bound is
type-compatible.

Items \textup{(a)--(d)} are the common logical content of, respectively,
the SM V33 exact second resolvent word, the SM V34 exact good/tail
factorization and one-tail ownership, the SM V32--V33 mixed-tree
assembly, and the SM V30/NS V31 gap and sharpness firewalls.
\end{proof}

\begin{firewall}[Companion type boundary]
\label{fire:V55-companion-types}
No Witten Laplacian, Higgs Hessian, fermion determinant, Gram minor,
Navier--Stokes heat ladder, coarea identity, or helical hinge is used
below as a Yang--Mills inequality.  The mathematical result in this
chapter uses only the YM prefix estimates, one-payer bounds, row
normalization, and the V54 tree identity.
\end{firewall}

\section{Separating the pure prefix face}
\label{sec:V55-pure-prefix}

Retain the V54 notation
\[
 \mathscr A_{211}=p(xy+xz+yz),
 \qquad
 m_2=m_2(x,y,z).
\]
For \(0<\eta\le1\), define the pure prefix sector
\begin{equation}
 \mathfrak P_\eta
 :=
 \{p<\eta,\ m_2(x,y,z)\ge\eta\}.
\label{eq:V55-pure-prefix-sector}
\end{equation}
The numerical cutoff \(p<\eta\) only identifies the face; the proof
below uses the lower bound on \(m_2\).

\begin{lemma}[Atlas lower bound on the pure prefix sector]
\label{lem:V55-pure-prefix-atlas}
On \(\mathfrak P_\eta\),
\begin{equation}
 \mathscr A_{211}\ge\eta^2p.
\label{eq:V55-pure-prefix-atlas}
\end{equation}
\end{lemma}

\begin{proof}
Apply \eqref{eq:V54-atlas-cover} and
\(m_2\ge\eta\).
\end{proof}

Fix one selected joining-coordinate occurrence in
\eqref{eq:V35-S211}, one admissible payer placement \(\pi\), and all
its complete prefix, singleton, and suffix banks.  Write the
resulting once-paid operator as
\(\widehat W_\pi\), as in
\eqref{eq:V49-complete-word}, and put
\begin{equation}
 \widetilde W_\pi
 :=
 \frac{\widehat W_\pi}
 {\Xi_1\Xi_2\Xi_3\Xi_4}.
\label{eq:V55-normalized-word}
\end{equation}
If a row is inactive at the joining coordinate, the corresponding
replacement difference vanishes.  Thus it suffices to consider
nonzero packets, for which \(\Xi_i\ge a_*>0\).

\begin{lemma}[Uniform complete-word prefix response]
\label{lem:V55-complete-prefix-response}
At fixed \(s=s_\alpha>0\), there is \(C_s<\infty\), independent of
support, representations, physical multiplicities, the selected
joining coordinate, and the payer placement, such that
\begin{equation}
 \|\widehat W_\pi\|\le C_sH_{12}.
\label{eq:V55-complete-prefix-response}
\end{equation}
\end{lemma}

\begin{proof}
Keep the complete binary prefix covariance as one factor.  If it is
unpaid, \Cref{thm:V10-prefix-response} gives
\[
 \|{\mathfrak J}_{<r}^{12}\|
 \le C\min\{1,sH_{12}(<r)\}
 \le CsH_{12}.
\]
If it carries the payer, the same theorem gives
\[
 \|{\mathfrak J}_{<r}^{12,\rm pay}\|
 \le C\min\{H_{12}(<r),s^{-1}\}
 \le CH_{12}.
\]

In every other payer branch, the sole output numerator and resolvent
are assigned to the assembled joining factor, a complete singleton
bank, or the complete suffix.  The one-payer estimates bound that
complete factor by \(C_s\), uniformly in its representations and
support.  Every unpaid complete moment is a contraction.  The
assembled third cumulant has norm at most six by its moment formula.
The fixed finite separators and physical compressions are
contractions.  Multiplying these operator-norm bounds proves
\eqref{eq:V55-complete-prefix-response}.  Finally
\(H_{12}(<r)\le H_{12}\).
\end{proof}

\begin{firewall}[Why the proof is not a duplicated debit]
\label{fire:V55-one-prefix-use}
The factor \(H_{12}\) in
\eqref{eq:V55-complete-prefix-response} is acquired once from the
complete signed prefix covariance.  The square below is the actual
operator square \(\widehat W_\pi^*\widehat W_\pi\); it is not a
second use of the prefix bank in the original source word.
\end{firewall}

\begin{theorem}[Quadratic closure of the pure prefix face]
\label{thm:V55-prefix-face-closure}
For every \(0<\eta\le1\), every complete selected-coordinate
\(2+1+1\) packet in \(\mathfrak P_\eta\), and every admissible payer
placement,
\begin{equation}
 \boxed{
 \kappa_\otimes(
 \widetilde W_\pi^*\widetilde W_\pi)
 \le
 C_{s,\eta}\,
 \mathscr A_{211}^{\,2}.}
\label{eq:V55-prefix-face-closure}
\end{equation}
The bound is uniform in support, representations, physical
multiplicities, and the selected joining coordinate.
\end{theorem}

\begin{proof}
Product-state capacity of a positive operator is at most its operator
norm.  Hence \Cref{lem:V55-complete-prefix-response} gives
\begin{align}
 \kappa_\otimes(
 \widetilde W_\pi^*\widetilde W_\pi)
 &\le
 \|\widetilde W_\pi\|^2
\nonumber\\
 &\le
 \frac{C_sH_{12}^2}
 {(\Xi_1\Xi_2\Xi_3\Xi_4)^2}
\nonumber\\
 &=
 \frac{C_sp^2}{\Xi_3^2\Xi_4^2}
 \le
 C_sa_*^{-4}p^2.
\label{eq:V55-prefix-square}
\end{align}
By \Cref{lem:V55-pure-prefix-atlas},
\[
 p^2\le\eta^{-4}\mathscr A_{211}^{\,2}.
\]
Absorb \(a_*^{-4}\eta^{-4}\) into \(C_{s,\eta}\).
\end{proof}

\begin{remark}[Relation to the V44 square]
\label{rem:V55-V44-relation}
V44 proves a sharper support-uniform capacity estimate for the
isolated doubled prefix, including its off-diagonal telescope terms.
The pure-face argument needs only its operator-norm consequence
because the quotient connectivity is bounded below by \(\eta\).
No V44 off-diagonal term is discarded; it remains inside the complete
prefix norm used in
\Cref{lem:V55-complete-prefix-response}.
\end{remark}

\section{The prefix face leaves the residual list}
\label{sec:V55-residual}

\begin{corollary}[Only quotient cuts can support a counterpacket]
\label{cor:V55-only-quotient-cuts}
Let a bounded-envelope counterpacket sequence be as in
\Cref{thm:V54-counterpacket-shape}.  After passage to a subsequence,
it must lie on one of the three quotient-cut faces
\begin{align*}
 x,y&\longrightarrow0,\\
 x,z&\longrightarrow0,\\
 y,z&\longrightarrow0.
\end{align*}
The pure prefix face cannot support such a sequence.
\end{corollary}

\begin{proof}
V54 gives either \(p\to0\) or one of the three quotient cuts.  In the
first case, pass to a further subsequence on which either
\(m_2\to0\) or \(m_2\ge\eta>0\).  If \(m_2\to0\), at least two of
\(x,y,z\) tend to zero, so this is already a quotient-cut face.  If
\(m_2\ge\eta\), the sequence eventually lies in
\(\mathfrak P_\eta\) and is controlled by
\Cref{thm:V55-prefix-face-closure}, contradicting the counterpacket
property.
\end{proof}

\begin{assessment}[Audit-adjusted next face]
\label{ass:V55-next-face}
The row-\(3\) face \(x,z\to0\) is the first remaining target.
In the V49 polarization, the row-\(3\) replacement difference is a
common complete factor of both Leibniz branches.  Its actual operator
square therefore contains two row-\(3\) occurrences before the
mixed row-\(1\)/row-\(2\) blocks are separated.  V56 should expose
that common factor while retaining:
\begin{enumerate}[label=\textup{(\roman*)},leftmargin=4em]
\item both composite Leibniz branches and their cross terms;
\item the row-\(4\) difference;
\item all exterior banks and their exact order; and
\item one fixed payer across the two square copies.
\end{enumerate}
This follows the SM V33 proof-order lesson without importing its
resolvent bound: use the exact quadratic word, not the square of a
previous positive one-edge estimate.
\end{assessment}

\section{V55 status and next acquisition}
\label{sec:V55-status}

\begin{theorem}[Third-series V55 result ledger]
\label{thm:V55-results}
Relative to V1--V54, V55 proves:
\begin{enumerate}[label=\textup{(V55.\arabic*)},leftmargin=6.3em]
\item the compulsory audit of SM V34 and unchanged NS V31, with exact
      file hashes;
\item a type-correct no-transfer decision and four proof-order
      constraints;
\item an exact separation of the pure prefix-depletion sector by the
      second-largest quotient connection;
\item a uniform complete-word operator bound linear in the actual
      prefix stock at every payer location;
\item a quadratic normalized \(FCQ_{211}\) bound on the entire pure
      prefix face;
\item removal of the prefix face from the possible counterpacket
      list; and
\item reduction of the remaining search to the three quotient cuts,
      beginning with the row-\(3\) cut.
\end{enumerate}
\end{theorem}

\begin{proof}
These are
\Cref{thm:V55-transfer-decision,
lem:V55-pure-prefix-atlas,
lem:V55-complete-prefix-response,
thm:V55-prefix-face-closure,
cor:V55-only-quotient-cuts,
ass:V55-next-face}.
\end{proof}

\begin{programme}[V56: row-\(3\) cut square]
\label{prog:V55-V56}
\begin{enumerate}[leftmargin=3em]
\item In the exact V49 four-block square, factor the common
      \(\delta_3X_3\) occurrence on both replica sides without
      separating the two composite Leibniz branches.
\item Retain the complete row-\(3\) off-joining banks and identify the
      exact coordinate support on which the two occurrences overlap.
\item Split the row-\(3\) cut into \(x\)- and \(z\)-dominant
      coordinate cells only after the signed square is formed.
\item Prove the quadratic cut-sum response in \((x+z)^2\), or exhibit
      a physical protected packet on which that response fails.
\item Keep the row-\(4\) difference and all mixed composite blocks;
      no diagonal block may be used as a lower or upper bound for the
      complete square by itself.
\item No companion audit is due before V60.
\end{enumerate}
\end{programme}

\begin{firewall}[Claim boundary after third-series V55]
\label{fire:V55-claim-boundary}
V55 closes only the pure prefix-depletion face of the selected
\(2+1+1\) quadratic source.  The composite, row-\(3\), and row-\(4\)
quotient cuts remain open, as do the unselected/global aggregation
and the other quartic source families.  Higher MTA, WUFC, CPM, cutoff
removal, reflection positivity, Osterwalder--Schrader reconstruction,
construction of the continuum Yang--Mills measure, and a uniform
positive continuum spectral gap remain open.  The full Yang--Mills
mass gap has not been proved.
\end{firewall}

\paragraph{Parent-version record.}
V55 was copied byte-for-byte from
\texttt{Renewal\_Yang\_Mills\_Mass\_Gap\_Research\_Notes\_V54.tex}
before this chapter was appended.  The V54 parent had \(27\,320\)
source lines, \(987\,348\) bytes, and SHA--256
\[
 \texttt{67c38196dec737fb0143fb746386e24fa2e7a4f2fd59d504b7bf40b169dda24f}.
\]
The next compulsory companion audit is V60.

% =====================================================================
% V56 APPENDIX: THE PURE ROW-3 CUT
% =====================================================================

\chapter{V56: The Pure Row-\(3\) Cut, Its Exact Square, and the
Protected-Output Obstruction}
\label{ch:V56-row3-cut}

\section{The correct cut currency}
\label{sec:V56-currency}

V55 selected the row-\(3\) quotient face
\[
 x,z\longrightarrow0.
\]
The first point is to identify its target correctly.  On a pure
row-\(3\) face the prefix edge \(p\) and the opposite quotient edge
\(y\) remain nondepleted.  Hence the eight-tree polynomial is linear
in the cut sum \(x+z\), not quadratic in the product \(xz\).  Since
\(FCQ_{211}\) is a square estimate, the required response is
\((x+z)^2\).

For \(0<\eta\le1\), define
\begin{equation}
 \mathfrak R_{3,\eta}
 :=
 \{p\ge\eta,\ y\ge\eta,\ x+z<\eta\}.
\label{eq:V56-pure-row3}
\end{equation}
The final inequality only identifies a neighborhood of the face and
is not used in the comparison below.

\begin{theorem}[Exact atlas equivalence on the pure row-\(3\) face]
\label{thm:V56-atlas-equivalence}
On \(\mathfrak R_{3,\eta}\),
\begin{equation}
 \boxed{
 \eta^2(x+z)
 \le
 \mathscr A_{211}
 \le
 3(x+z).}
\label{eq:V56-atlas-equivalence}
\end{equation}
\end{theorem}

\begin{proof}
By \eqref{eq:V54-exact-factorization},
\[
 \mathscr A_{211}
 =
 p\{y(x+z)+xz\}.
\]
The first term and \(p,y\ge\eta\) give the lower bound.  For the
upper bound, use \(p\le1\), \(y\le2\), and
\(xz\le x\le x+z\), since \(z\le1\).
\end{proof}

\begin{definition}[Row-\(3\) cut-square card]
\label{def:V56-RCQ3}
For the normalized once-paid word
\(\widetilde W_\pi\) in
\eqref{eq:V55-normalized-word}, the card \(RCQ_3(\eta)\) is
\begin{equation}
 \boxed{
 \kappa_\otimes(
 \widetilde W_\pi^*\widetilde W_\pi)
 \le
 C_{s,\eta}(x+z)^2}
\label{eq:V56-RCQ3}
\end{equation}
for every complete selected-coordinate packet in
\(\mathfrak R_{3,\eta}\) and every admissible payer placement.
\end{definition}

\begin{corollary}[Exact reduction of the pure row-\(3\) face]
\label{cor:V56-RCQ3-equivalence}
On \(\mathfrak R_{3,\eta}\), \(RCQ_3(\eta)\) is equivalent, up to
constants depending on \(\eta\), to the restriction of
\(FCQ_{211}\):
\begin{equation}
 \kappa_\otimes(
 \widetilde W_\pi^*\widetilde W_\pi)
 \le C_{s,\eta}\mathscr A_{211}^{\,2}.
\label{eq:V56-row3-FCQ}
\end{equation}
\end{corollary}

\begin{proof}
By \eqref{eq:V56-atlas-equivalence},
\[
 x+z\le\eta^{-2}\mathscr A_{211},
 \qquad
 \mathscr A_{211}\le3(x+z).
\]
The first implication sends \(RCQ_3\) to \(FCQ_{211}\), and the
second sends the restricted \(FCQ_{211}\) estimate to \(RCQ_3\).
\end{proof}

\section{The two actual row-\(3\) occurrences}
\label{sec:V56-exact-square}

At the joining coordinate, keep the composite difference unsplit:
\begin{equation}
 \Delta_{12}(\boldsymbol\xi)
 :=
 Y_{12}(\xi_0)-Y_{12}(\xi_1),
\label{eq:V56-composite-difference}
\end{equation}
and set
\begin{equation}
 {\cal D}(\boldsymbol\xi)
 :=
 \Delta_{12}(\boldsymbol\xi)
 \otimes\delta_3X_3(\boldsymbol\xi)
 \otimes\delta_4X_4(\boldsymbol\xi).
\label{eq:V56-unsplit-D}
\end{equation}
The row order is \(\{1,2\},3,4\).

\begin{theorem}[Exact common-row square]
\label{thm:V56-common-row-square}
For every outside-local payer placement,
\begin{equation}
 \boxed{
 \widehat W_\pi^*\widehat W_\pi
 =
 {\mathbb E}_{\boldsymbol\eta,\boldsymbol\xi}
 \left[
 R_\pi^*
 {\cal D}(\boldsymbol\eta)^*
 L_\pi^*L_\pi
 {\cal D}(\boldsymbol\xi)
 R_\pi
 \right].}
\label{eq:V56-common-row-square}
\end{equation}
The integrand contains exactly two actual row-\(3\) replacement
differences,
\[
 \delta_3X_3(\boldsymbol\eta)^*,
 \qquad
 \delta_3X_3(\boldsymbol\xi),
\]
while the complete composite and row-\(4\) differences remain
attached.  Expanding \(\Delta_{12}\) by
\eqref{eq:V49-Leibniz} recovers all four blocks of
\eqref{eq:V49-composite-square}.
\end{theorem}

\begin{proof}
\Cref{thm:V49-composite-polarization} states that
\[
 \mathsf S_{12|3|4}
 =
 {\mathbb E}_{\boldsymbol\xi}{\cal D}(\boldsymbol\xi).
\]
Insert this formula and an independent adjoint copy into
\(\widehat W_\pi=L_\pi\mathsf S_{12|3|4}R_\pi\), then use Fubini.
This gives \eqref{eq:V56-common-row-square}.  The two row-\(3\)
occurrences are visible in
\eqref{eq:V56-unsplit-D}.  Finally the two terms of the exact
Leibniz identity on each replica side produce the four pairs
\((a,b)\in\{1,2\}^2\) in
\eqref{eq:V49-composite-square}.
\end{proof}

\begin{corollary}[No composite-branch separation is needed]
\label{cor:V56-no-branch-separation}
The row-\(3\) cut is common to the complete composite-source square.
It may therefore be investigated before selecting a row-\(1\) or
row-\(2\) Leibniz branch.  Any proof which estimates the four
\((a,b)\) blocks separately imposes a stronger condition than
\(RCQ_3(\eta)\).
\end{corollary}

\begin{proof}
Equation \eqref{eq:V56-common-row-square} is formed with the unsplit
\(\Delta_{12}\).  The four-block expression is only its distributive
expansion.  Since individual mixed blocks need not be positive,
separate estimates discard cancellations retained by the exact
square.
\end{proof}

\begin{firewall}[Compression prevents formal extraction]
\label{fire:V56-no-formal-extraction}
The two row-\(3\) differences are genuine occurrences, but they do
not automatically factor outside
\(L_\pi^*L_\pi\), the physical output compression, or the payer.
Equation \eqref{eq:V56-common-row-square} is an occurrence identity,
not the inequality \(RCQ_3(\eta)\).
\end{firewall}

\begin{firewall}[The local-payer square is still separate]
\label{fire:V56-local-payer-separate}
\Cref{thm:V56-common-row-square} uses the outside-local placement
assumed in V49.  If the joining carrier itself pays, its resolved
output mass and resolvent must be included inside both replica copies
before the square is formed.  The occurrence of two row-\(3\)
differences remains suggestive, but no local-payer polarization is
inferred from \eqref{eq:V56-common-row-square}.
\end{firewall}

\section{Why protected output cannot simply be squared}
\label{sec:V56-protected-obstruction}

The exact square suggests applying the one-row invariant-capacity
bound once to each row-\(3\) occurrence.  That argument is false when
both occurrences resolve to the same protected output projection.
The obstruction is idempotence.

\begin{lemma}[Sharp invariant projection and its square]
\label{lem:V56-idempotent-capacity}
Let \(V\) be a nontrivial irreducible unitary representation of
dimension \(d\), let \(P_{\rm sing}\) be the rank-one invariant
projection in \(V\otimes V^*\), and let
\(\kappa_\otimes\) be product-state capacity across
\(V|V^*\).  Then
\begin{equation}
 \boxed{
 \kappa_\otimes(P_{\rm sing})
 =
 \kappa_\otimes(P_{\rm sing}^2)
 =
 \frac1d.}
\label{eq:V56-idempotent-capacity}
\end{equation}
Consequently no representation-independent \(C\) can satisfy
\begin{equation}
 \kappa_\otimes(P_{\rm sing}^2)
 \le C\,\kappa_\otimes(P_{\rm sing})^2
\label{eq:V56-false-capacity-square}
\end{equation}
for all irreducible \(V\).
\end{lemma}

\begin{proof}
With dual orthonormal bases, the invariant unit vector is
\[
 \Omega=d^{-1/2}\sum_{i=1}^d e_i\otimes e_i^*,
 \qquad
 P_{\rm sing}=|\Omega\rangle\langle\Omega|.
\]
\Cref{thm:V32-invariant-capacity} gives the upper bound \(1/d\).
The product vector \(e_1\otimes e_1^*\) gives
\[
 |\langle\Omega,e_1\otimes e_1^*\rangle|^2=\frac1d,
\]
so the bound is attained.  Since \(P_{\rm sing}^2=P_{\rm sing}\),
the same value holds for the square.  If
\eqref{eq:V56-false-capacity-square} held, it would imply
\(d\le C\); dimensions are unbounded already for \(SU(2)\).
\end{proof}

\begin{proposition}[Exact limitation of the two-occurrence argument]
\label{prop:V56-protected-limit}
The occurrence count in
\Cref{thm:V56-common-row-square}, together with the statewise
one-row protected bound
\eqref{eq:V32-protected-capacity}, does not imply
\(RCQ_3(\eta)\).  A valid proof must do at least one of the following
before protected output pinching:
\begin{enumerate}[label=\textup{(\roman*)},leftmargin=4em]
\item retain cancellation with the complete composite or row-\(4\)
      differences;
\item use distinct physical coordinate banks for the two required
      gains;
\item obtain an additional partner-edge response from the exact
      Wilson word; or
\item show that the relevant output is gapped rather than protected.
\end{enumerate}
\end{proposition}

\begin{proof}
After both row-\(3\) occurrences are compressed to one common
protected block, their product contains the same invariant
projection twice.  By
\Cref{lem:V56-idempotent-capacity}, the product-state capacity is
still \(1/d\), not \(1/d^2\).  Therefore two occurrences plus the
one-use protected-capacity estimate cannot furnish the square of a
cut currency.  Each listed alternative introduces information absent
from that invalid inference.
\end{proof}

\begin{remark}[This does not disprove \(RCQ_3\)]
\label{rem:V56-not-RCQ-counterexample}
\Cref{prop:V56-protected-limit} is a no-go theorem for one proof
step, not a physical counterpacket to the complete source card.  The
composite difference, row-\(4\) difference, prefix, suffix, and payer
may remove or reroute the protected channel before the final
capacity is taken.
\end{remark}

\section{Normalized row-mass overlap}
\label{sec:V56-overlap}

For each active row define a probability distribution on physical
coordinates:
\begin{equation}
 w_i(e):=\frac{a_{i,e}}{\Xi_i},
 \qquad
 \sum_e w_i(e)=1.
\label{eq:V56-row-distribution}
\end{equation}

\begin{lemma}[Edges are exact mass-distribution overlaps]
\label{lem:V56-overlap-identity}
For every \(u\ne v\),
\begin{equation}
 \theta_{uv}=\sum_e w_u(e)w_v(e).
\label{eq:V56-overlap-edge}
\end{equation}
In particular, with
\begin{equation}
 g_3(e):=w_1(e)+w_2(e)+w_4(e),
\label{eq:V56-g3}
\end{equation}
the row-\(3\) cut currency is
\begin{equation}
 \boxed{
 x+z
 =
 \sum_e w_3(e)g_3(e).}
\label{eq:V56-cut-overlap}
\end{equation}
\end{lemma}

\begin{proof}
Divide
\(H_{uv}=\sum_ea_{u,e}a_{v,e}\) by
\(\Xi_u\Xi_v\).  Summing the cases
\((u,v)=(3,1),(3,2),(3,4)\) gives
\eqref{eq:V56-cut-overlap}.
\end{proof}

\begin{corollary}[Exact low-partner-density localization]
\label{cor:V56-low-density}
For \(0<\lambda\le1\), set
\[
 S_\lambda:=\{e:g_3(e)\ge\lambda\}.
\]
Then
\begin{equation}
 \sum_{e\in S_\lambda}w_3(e)
 \le\frac{x+z}{\lambda}.
\label{eq:V56-shared-mass}
\end{equation}
Thus, when \(x+z\) is small, almost all normalized row-\(3\) mass
lies at coordinates where the total normalized density of the other
three rows is small.
\end{corollary}

\begin{proof}
On \(S_\lambda\),
\(w_3(e)g_3(e)\ge\lambda w_3(e)\).  Sum and use
\eqref{eq:V56-cut-overlap}.
\end{proof}

\begin{firewall}[Normalized separation is not raw
representation imbalance]
\label{fire:V56-normalized-not-raw}
The inequality \(g_3(e)\ll1\) compares the fractions
\(a_{u,e}/\Xi_u\).  It does not imply
\(\sum_{u\ne3}\|\lambda_{u,e}\|
\ll\|\lambda_{3,e}\|\), because another row may have a much larger
total mass \(\Xi_u\).  Therefore the unbalanced-output theorem
\Cref{lem:V32-unbalanced-output} may not be applied from
\eqref{eq:V56-shared-mass} alone.
\end{firewall}

\section{Balanced atoms force mass escalation}
\label{sec:V56-escalation}

Let \(c_G>0\) be chosen so that at every coordinate balanced at row
\(3\), \Cref{lem:V32-balanced-partner} supplies some
\(u\in\{1,2,4\}\) with
\begin{equation}
 a_{u,e}\ge c_Ga_{3,e}.
\label{eq:V56-balanced-partner}
\end{equation}

\begin{lemma}[Atomic mass-escalation inequality]
\label{lem:V56-mass-escalation}
Suppose \(e\notin S_\lambda\), \(w_3(e)\ge\tau>0\), and the local
representation assignment is balanced at row \(3\).  For a partner
\(u\) satisfying
\eqref{eq:V56-balanced-partner},
\begin{equation}
 \boxed{
 \frac{\Xi_u}{\Xi_3}
 \ge
 c_G\frac{w_3(e)}{w_u(e)}
 \ge
 c_G\frac{\tau}{\lambda}.}
\label{eq:V56-mass-escalation}
\end{equation}
\end{lemma}

\begin{proof}
Since \(e\notin S_\lambda\),
\(w_u(e)\le g_3(e)<\lambda\).  Also
\[
 w_u(e)
 =
 \frac{a_{u,e}}{\Xi_u}
 \ge
 c_G\frac{a_{3,e}}{\Xi_u}
 =
 c_G\frac{\Xi_3}{\Xi_u}w_3(e).
\]
Rearrange and use \(w_3(e)\ge\tau\).
\end{proof}

\begin{corollary}[Four-way row-\(3\) decomposition]
\label{cor:V56-four-way}
Let \(\varepsilon:=x+z\in(0,1)\), and set
\[
 \lambda:=\varepsilon^{1/2},
 \qquad
 \tau:=\varepsilon^{1/4}.
\]
Every coordinate active at row \(3\) belongs to exactly one of:
\begin{enumerate}[label=\textup{(\roman*)},leftmargin=4em]
\item the shared set \(S_\lambda\), carrying at most
      \(\varepsilon^{1/2}\) of the normalized row-\(3\) mass;
\item a low-density balanced atom with \(w_3(e)\ge\tau\), of which
      there are at most \(\varepsilon^{-1/4}\), and whose selected
      partner obeys
      \[
      \Xi_u/\Xi_3\ge c_G\varepsilon^{-1/4};
      \]
\item a low-density balanced diffuse coordinate with
      \(w_3(e)<\varepsilon^{1/4}\); or
\item a low-density coordinate unbalanced at row \(3\), whose every
      output retains a fixed fraction of the local row-\(3\)
      Casimir by
      \Cref{lem:V32-unbalanced-output}.
\end{enumerate}
\label{eq:V56-four-way}
\end{corollary}

\begin{proof}
Split first into \(S_\lambda\) and its complement, then split the
complement by balanced/unbalanced and by
\(w_3(e)\ge\tau\) or \(w_3(e)<\tau\) on the balanced part.
The shared-mass bound is
\eqref{eq:V56-shared-mass}:
\[
 \varepsilon/\lambda=\varepsilon^{1/2}.
\]
There are at most \(1/\tau=\varepsilon^{-1/4}\) atoms because
\(\sum_ew_3(e)=1\).  Their mass escalation follows from
\Cref{lem:V56-mass-escalation}.  The unbalanced output statement is
\Cref{lem:V32-unbalanced-output}.
\end{proof}

\begin{assessment}[The two genuine residual mechanisms]
\label{ass:V56-residual-mechanisms}
The row-\(3\) cut has now been reduced beyond a generic
\emph{hot-row} label.  Shared mass is already small in the exact cut
currency.  Unbalanced low-density coordinates have a gapped output.
Balanced atomic coordinates force a strict escalation to a larger
partner-row mass and can occur only finitely many times at a fixed
threshold.  The unresolved sector is the union of:
\begin{enumerate}[label=\textup{(\alph*)},leftmargin=4em]
\item the mass-escalation chain, which must be routed into the strong
      \(p,y\) edges or terminate in a gapped output; and
\item the diffuse balanced protected bank, where
      \Cref{lem:V56-idempotent-capacity} forbids simply squaring the
      one-use invariant capacity.
\end{enumerate}
These mechanisms require different estimates and may not be combined
by a positive union bound before the complete square is formed.
\end{assessment}

\section{V56 status and next acquisition}
\label{sec:V56-status}

\begin{theorem}[Third-series V56 result ledger]
\label{thm:V56-results}
Relative to V1--V55, V56 proves:
\begin{enumerate}[label=\textup{(V56.\arabic*)},leftmargin=6.3em]
\item exact comparability of the eight-tree atlas with the row-\(3\)
      cut sum \(x+z\) on the pure face;
\item equivalence of pure-face \(FCQ_{211}\) with the quadratic
      cut-square card \(RCQ_3(\eta)\);
\item an unsplit replica identity displaying two actual row-\(3\)
      differences while preserving all composite cross terms;
\item proof that common output compression prevents their formal
      extraction;
\item the sharp idempotent invariant-projection obstruction to
      squaring a protected capacity;
\item the exact representation of normalized edge currencies as
      overlaps of row-mass probability distributions;
\item localization of row-\(3\) mass to low partner-density
      coordinates on the cut face;
\item a quantitative mass-escalation theorem for every balanced
      atomic coordinate; and
\item an exhaustive shared/atomic/diffuse/unbalanced decomposition of
      the remaining row-\(3\) mechanism.
\end{enumerate}
\end{theorem}

\begin{proof}
These are
\Cref{thm:V56-atlas-equivalence,
cor:V56-RCQ3-equivalence,
thm:V56-common-row-square,
fire:V56-no-formal-extraction,
lem:V56-idempotent-capacity,
lem:V56-overlap-identity,
cor:V56-low-density,
lem:V56-mass-escalation,
cor:V56-four-way}.
\end{proof}

\begin{programme}[V57: escalation routing and the diffuse protected
bank]
\label{prog:V56-V57}
\begin{enumerate}[leftmargin=3em]
\item Iterate \eqref{eq:V56-mass-escalation} only across distinct
      row labels and prove that a cycle produces a nondepleted edge,
      while an acyclic chain terminates after at most three steps.
\item Keep the strong pure-face edges \(p\) and \(y\) visible during
      that routing; do not replace them by constants before selecting
      the literal tree.
\item On diffuse balanced banks, apply the V32 protected/gapped
      decomposition to the complete bank before the two square copies
      are pinched together.
\item Test whether the one-failure resolver supplies a second
      independent bank occurrence.  If it resolves both copies to the
      same invariant projector, retain the composite and row-\(4\)
      factors rather than squaring its capacity.
\item Prove \(RCQ_3(\eta)\) on one of the two residual mechanisms or
      construct a complete physical counterpacket including its
      payer and all legal trees.
\item No companion audit is due before V60.
\end{enumerate}
\end{programme}

\begin{firewall}[Claim boundary after third-series V56]
\label{fire:V56-claim-boundary}
V56 identifies the exact pure row-\(3\) cut card and its remaining
protected mechanisms; it does not prove \(RCQ_3(\eta)\) or close that
face.  The composite and row-\(4\) quotient cuts also remain open, as
do the other quartic source families, higher MTA, WUFC, CPM, cutoff
removal, reflection positivity, Osterwalder--Schrader reconstruction,
construction of the continuum Yang--Mills measure, and a uniform
positive continuum spectral gap.  The full Yang--Mills mass gap has
not been proved.
\end{firewall}

\paragraph{Parent-version record.}
V56 was copied byte-for-byte from
\texttt{Renewal\_Yang\_Mills\_Mass\_Gap\_Research\_Notes\_V55.tex}
before this chapter was appended.  The V55 parent had \(27\,696\)
source lines, \(1\,001\,537\) bytes, and SHA--256
\[
 \texttt{a0d6005ebb17648651d24a5db4b0d30770046d21b4ea9a6af5ccff9d9aa3ce52}.
\]
The next compulsory companion audit is V60.

% =====================================================================
% V57 APPENDIX: OWNED COMPARABLE BANKS AND MASS ESCALATION
% =====================================================================

\chapter{V57: Quadratic Currency from an Owned Comparable Bank and
Finite Mass Escalation}
\label{ch:V57-owned-bank}

\section{Purpose and scope}
\label{sec:V57-purpose}

V56 isolated two residual mechanisms on the pure row-\(3\) cut:
mass escalation through a balanced atomic partner, and a diffuse
balanced protected bank on which a single invariant projection is
idempotent.  The latter obstruction disappears when one complete
unmarked bank owns fixed fractions of both participating row masses.
The tensorized protected estimate then contains an additional
inverse bank length, and that factor turns the one-use capacity into
the square of the normalized edge currency.

This chapter proves the exact bank theorem and an exhaustive
alternative: a cofinal balanced row-\(3\) bank either contains such
an owned comparable subbank or forces a prescribed increase of total
mass to one of rows \(1,2,4\).  Strict mass increases cannot form a
cycle and terminate after at most three row changes.

\begin{assessment}[Nonduplication boundary]
\label{ass:V57-nonduplication}
\Cref{cor:V32-V70-currencies} proves protected and gapped currencies
in the internal bank variables \(N,S,H\).  It does not compare those
currencies with the globally normalized edge
\(H/(\Xi_3\Xi_u)\).  V57 adds the two-sided local comparability and
two-row ownership hypotheses which make that comparison exact, then
proves the complementary mass-escalation alternative.  No V32
representation theorem is reproved.
\end{assessment}

\section{Owned comparable banks}
\label{sec:V57-definitions}

Fix \(u\in\{1,2,4\}\), and let \(D\) be a nonempty complete unmarked
coordinate bank assigned to the pair \(3u\).  Put
\begin{equation}
 N:=|D|,\qquad
 S:=\sum_{e\in D}a_{3,e},\qquad
 T:=\sum_{e\in D}a_{u,e},\qquad
 H:=\sum_{e\in D}a_{3,e}a_{u,e}.
\label{eq:V57-bank-data}
\end{equation}
Its globally normalized edge contribution is
\begin{equation}
 \theta_{3u}(D):=\frac{H}{\Xi_3\Xi_u}.
\label{eq:V57-local-theta}
\end{equation}

\begin{definition}[Comparable and owned]
\label{def:V57-owned-comparable}
For constants \(0<c\le C<\infty\) and
\(0<\sigma_3,\sigma_u\le1\), the bank \(D\) is
\((c,C;\sigma_3,\sigma_u)\)-owned comparable if
\begin{align}
 ca_{3,e}\le a_{u,e}&\le Ca_{3,e}
 \qquad(e\in D),
\label{eq:V57-comparable}\\
 S\ge\sigma_3\Xi_3,\qquad
 T&\ge\sigma_u\Xi_u.
\label{eq:V57-ownership}
\end{align}
\end{definition}

\begin{lemma}[Internal currency and global edge are comparable]
\label{lem:V57-internal-global}
On an owned comparable bank,
\begin{align}
 \frac cN
 \le\frac{H}{S^2}
 &\le C,
\label{eq:V57-internal-range}\\
 \theta_{3u}(D)
 &\ge
 \frac{\sigma_3\sigma_u}{C}\,
 \frac{H}{S^2}.
\label{eq:V57-global-lower}
\end{align}
\end{lemma}

\begin{proof}
The lower comparability bound gives
\[
 H\ge c\sum_{e\in D}a_{3,e}^2
 \ge\frac cN\left(\sum_{e\in D}a_{3,e}\right)^2
 =\frac cN S^2.
\]
The upper bound follows from
\[
 H\le C\sum_ea_{3,e}^2\le CS^2.
\]
Also \(T\le CS\).  Ownership gives
\[
 \Xi_3\Xi_u
 \le
 \frac{S}{\sigma_3}\frac{T}{\sigma_u}
 \le
 \frac{CS^2}{\sigma_3\sigma_u}.
\]
Insert this into \eqref{eq:V57-local-theta}.
\end{proof}

\section{Protected and gapped sectors both become quadratic}
\label{sec:V57-quadratic-currency}

Let \(\mathcal P_{3u}(D)\) and
\(\mathcal G_{3u}(D)\) be the complete protected and gapped bank
envelopes of \Cref{cor:V32-V70-currencies}.  The bank is assumed not
to contain the unique payer or any selected tree coordinate; those
finite separators are kept outside \(D\).

\begin{theorem}[Owned comparable bank quadratic currency]
\label{thm:V57-bank-quadratic}
For every fixed group and fixed \(s=s_\alpha>0\), an owned comparable
bank satisfies
\begin{align}
 \kappa_\otimes(\mathcal P_{3u}(D))
 &\le
 C_{G,c,C,\sigma_3,\sigma_u}\,
 \theta_{3u}(D)^2,
\label{eq:V57-protected-quadratic}\\
 \kappa_\otimes(\mathcal G_{3u}(D))
 &\le
 C_{G,s,c,C,\sigma_3,\sigma_u}\,
 \theta_{3u}(D)^2.
\label{eq:V57-gapped-quadratic}
\end{align}
The constants are independent of \(N\), support, representations, and
physical multiplicities.
\end{theorem}

\begin{proof}
The protected currency
\eqref{eq:V32-protected-currency} gives
\[
 \kappa_\otimes(\mathcal P_{3u}(D))
 \le
 \frac{C_G}{N}\frac{H}{S^2}.
\]
By \eqref{eq:V57-internal-range},
\[
 \frac1N\frac{H}{S^2}
 \le
 \frac1c\left(\frac{H}{S^2}\right)^2.
\]
Equation \eqref{eq:V57-global-lower} now proves
\eqref{eq:V57-protected-quadratic}.

For the gapped sector use the exponential entry of
\eqref{eq:V32-gapped-currency}:
\[
 \kappa_\otimes(\mathcal G_{3u}(D))
 \le C_Ge^{-\gamma_GsN}.
\]
The first inequality in
\eqref{eq:V57-internal-range} and
\eqref{eq:V57-global-lower} give
\[
 \theta_{3u}(D)\ge
 \frac{c\sigma_3\sigma_u}{CN}.
\]
Since
\(\sup_{N\ge1}N^2e^{-\gamma_GsN}<\infty\),
\eqref{eq:V57-gapped-quadratic} follows.
\end{proof}

\begin{corollary}[Idempotence is harmless after complete-bank
tensorization]
\label{cor:V57-idempotence-resolved}
If either complete bank envelope is a positive contraction
\(\mathcal M_D\), then
\begin{equation}
 \kappa_\otimes(\mathcal M_D^2)
 \le
 \kappa_\otimes(\mathcal M_D)
 \le
 C\,\theta_{3u}(D)^2.
\label{eq:V57-square-bank}
\end{equation}
Thus the V56 single-projection obstruction does not persist on an
owned comparable complete bank.
\end{corollary}

\begin{proof}
Functional calculus gives
\(0\preccurlyeq\mathcal M_D^2\preccurlyeq\mathcal M_D\).
Product-state capacity is monotone on positive operators.  Apply
\Cref{thm:V57-bank-quadratic}.
\end{proof}

\begin{corollary}[The bank pays the row-\(3\) cut scale]
\label{cor:V57-bank-pays-cut}
For \(u\in\{1,2,4\}\),
\begin{equation}
 \theta_{3u}(D)\le\theta_{3u}\le x+z.
\label{eq:V57-bank-below-cut}
\end{equation}
Consequently every bank sector covered by
\Cref{thm:V57-bank-quadratic} has capacity at most
\(C(x+z)^2\), the exact \(RCQ_3\) scale.
\end{corollary}

\begin{proof}
The local stock \(H(D)\) is a sub-sum of \(H_{3u}\), proving the
first inequality.  The second follows from
\[
 x+z=\theta_{31}+\theta_{32}+\theta_{34}.
\]
Square and apply \Cref{thm:V57-bank-quadratic}.
\end{proof}

\begin{firewall}[What remains to insert the bank into the source]
\label{fire:V57-insertion-boundary}
\Cref{thm:V57-bank-quadratic} is a complete-bank capacity estimate.
To close a full word, the protected/gapped bank must still be exposed
as one unmarked tensor factor of
\eqref{eq:V56-common-row-square} before physical pinching.  If the
bank contains the payer, a selected coordinate, or is interlaced with
the output compression, its separate extraction is not asserted.
\end{firewall}

\section{Every failed bank forces total-mass escalation}
\label{sec:V57-dichotomy}

Let \(E\) be a complete unmarked bank of coordinates balanced at row
\(3\).  At each coordinate choose one partner
\(u(e)\in\{1,2,4\}\) supplied by
\eqref{eq:V56-balanced-partner}, and let
\[
 E_u:=\{e\in E:u(e)=u\}.
\]

\begin{theorem}[Comparable-bank or mass-escalation dichotomy]
\label{thm:V57-bank-or-escalation}
Assume
\begin{equation}
 \sum_{e\in E}a_{3,e}\ge\sigma\Xi_3
\label{eq:V57-cofinal-row3}
\end{equation}
for some \(0<\sigma\le1\).  Fix \(K>c_G\) and
\(0<\omega<1\).  Then at least one of the following holds:
\begin{enumerate}[label=\textup{(\roman*)},leftmargin=4em]
\item there exist \(u\in\{1,2,4\}\) and a complete subbank
      \(D\subseteq E_u\) which is
      \[
      (c_G,K;\sigma/6,\omega)\text{-owned comparable};
      \]
\item for some \(u\in\{1,2,4\}\),
      \begin{equation}
      \frac{\Xi_u}{\Xi_3}
      \ge
      \min\left\{
      \frac{K\sigma}{6},
      \frac{c_G\sigma}{6\omega}
      \right\}.
      \label{eq:V57-escalation-factor}
      \end{equation}
\end{enumerate}
\end{theorem}

\begin{proof}
The three partner banks partition \(E\), so one of them, say
\(E_u\), carries row-\(3\) mass
\[
 S_u:=\sum_{e\in E_u}a_{3,e}
 \ge\frac{\sigma}{3}\Xi_3.
\]
Split it into
\begin{align*}
 D&:=\{e\in E_u:a_{u,e}\le Ka_{3,e}\},\\
 U&:=\{e\in E_u:a_{u,e}>Ka_{3,e}\}.
\end{align*}
Every coordinate in both sets has
\(a_{u,e}\ge c_Ga_{3,e}\) by the partner choice.

If the row-\(3\) mass of \(U\) is at least \(S_u/2\), then
\[
 \Xi_u
 \ge\sum_{e\in U}a_{u,e}
 >
 K\sum_{e\in U}a_{3,e}
 \ge\frac{K\sigma}{6}\Xi_3,
\]
which is item \textup{(ii)}.

Otherwise \(D\) has row-\(3\) mass
\[
 S(D)\ge S_u/2\ge\frac{\sigma}{6}\Xi_3
\]
and is \((c_G,K)\)-comparable.  If
\(T(D)\ge\omega\Xi_u\), it satisfies all conditions in item
\textup{(i)}.  If not, comparability from below gives
\[
 \Xi_u
 >
 \frac{T(D)}{\omega}
 \ge
 \frac{c_GS(D)}{\omega}
 \ge
 \frac{c_G\sigma}{6\omega}\Xi_3,
\]
again item \textup{(ii)}.
\end{proof}

\begin{corollary}[Arbitrarily strong finite alternative]
\label{cor:V57-arbitrary-escalation}
For every prescribed \(R>1\), the parameters \(K,\omega\) may be
chosen, depending only on \(R,\sigma,G\), so that
\Cref{thm:V57-bank-or-escalation} yields either an owned comparable
bank covered by \Cref{thm:V57-bank-quadratic}, or
\begin{equation}
 \Xi_u\ge R\Xi_3
\label{eq:V57-R-escalation}
\end{equation}
for some \(u\in\{1,2,4\}\).
\end{corollary}

\begin{proof}
Take
\[
 K:=\max\{2c_G,6R/\sigma\},
 \qquad
 \omega:=\min\{1/2,c_G\sigma/(6R)\}.
\]
Both entries of the minimum in
\eqref{eq:V57-escalation-factor} are then at least \(R\).
\end{proof}

\section{Escalation cannot cycle}
\label{sec:V57-cycle}

\begin{lemma}[Strict mass arrows are acyclic]
\label{lem:V57-no-cycle}
Fix \(R>1\).  Draw a directed arrow \(i\to u\) only when
\begin{equation}
 \Xi_u\ge R\Xi_i.
\label{eq:V57-mass-arrow}
\end{equation}
The resulting directed graph on the four row labels has no directed
cycle.  Every directed path has at most three arrows.
\end{lemma}

\begin{proof}
Along a directed cycle of length \(m\), multiplying
\eqref{eq:V57-mass-arrow} would give
\[
 \Xi_i\ge R^m\Xi_i,
\]
which is impossible because \(\Xi_i>0\) and \(R^m>1\).
An acyclic directed path on four vertices visits each vertex at most
once and therefore has at most three arrows.
\end{proof}

\begin{theorem}[Finite termination of repeated bank failure]
\label{thm:V57-finite-termination}
Suppose at each visited row a complete balanced bank owns a fixed
fraction \(\sigma\) of that row mass, and apply
\Cref{cor:V57-arbitrary-escalation} with one fixed \(R>1\).
After at most three escalations, one of the following must occur:
\begin{enumerate}[label=\textup{(\alph*)},leftmargin=4em]
\item an owned comparable bank supplies a quadratic normalized edge
      currency;
\item the cofinal row mass lies in a marked or payer bank and leaves
      the unmarked theorem's scope;
\item the cofinal bank is unbalanced and its output is gapped; or
\item the terminal row has no cofinal balanced bank of the stated
      type.
\end{enumerate}
An infinite hierarchy of increasingly massive protected partner rows
is impossible.
\end{theorem}

\begin{proof}
Whenever the first alternative of
\Cref{cor:V57-arbitrary-escalation} holds, item \textup{(a)} occurs.
If its hypotheses fail, the failure is exactly one of
\textup{(b)--(d)}.  Otherwise draw the strict mass arrow supplied by
its second alternative and repeat at the new row.  By
\Cref{lem:V57-no-cycle}, at most three arrows can be drawn.
\end{proof}

\begin{assessment}[Residual after the bank theorem]
\label{ass:V57-residual}
For the pure row-\(3\) cut, a diffuse balanced protected bank is no
longer intrinsically obstructive.  If one comparable partner bank
owns cofinal fractions of both row masses and is unmarked, its
complete protected and gapped sectors already have the
\((x+z)^2\) scale.  The remaining work is now localized to:
\begin{enumerate}[label=\textup{(\roman*)},leftmargin=4em]
\item inserting that bank through the exact V56 square and output
      compression;
\item routing the terminal row of the finite mass-escalation chain
      into the strong pure-face edges \(p,y\);
\item treating the finite marked/payer coordinates without applying
      an unmarked tensor estimate; and
\item proving the local-payer replica square separately.
\end{enumerate}
\end{assessment}

\section{V57 status and next acquisition}
\label{sec:V57-status}

\begin{theorem}[Third-series V57 result ledger]
\label{thm:V57-results}
Relative to V1--V56, V57 proves:
\begin{enumerate}[label=\textup{(V57.\arabic*)},leftmargin=6.3em]
\item exact comparison of internal bank currency with a globally
      normalized row edge under two-sided comparability and ownership;
\item quadratic protected and gapped capacity bounds for every owned
      comparable unmarked bank;
\item removal of the idempotent-projection loss after complete-bank
      tensorization;
\item direct domination of those bank capacities by the
      \(RCQ_3\) target \((x+z)^2\);
\item an exhaustive comparable-bank or total-mass-escalation
      dichotomy;
\item an arbitrarily strong fixed-factor version of that dichotomy;
\item acyclicity and three-step termination of strict row-mass
      escalation; and
\item localization of the residual to bank insertion, terminal-row
      routing, marked/payer coordinates, and the local-payer square.
\end{enumerate}
\end{theorem}

\begin{proof}
These are
\Cref{lem:V57-internal-global,
thm:V57-bank-quadratic,
cor:V57-idempotence-resolved,
cor:V57-bank-pays-cut,
thm:V57-bank-or-escalation,
cor:V57-arbitrary-escalation,
lem:V57-no-cycle,
thm:V57-finite-termination}.
\end{proof}

\begin{programme}[V58: insertion and terminal-row routing]
\label{prog:V57-V58}
\begin{enumerate}[leftmargin=3em]
\item Write the exact tensor placement of one unmarked
      row-\(3\)--partner bank inside
      \eqref{eq:V56-common-row-square}, with every marked coordinate
      removed before applying
      \Cref{thm:V57-bank-quadratic}.
\item Determine whether physical output compression is block-diagonal
      with respect to that protected/gapped bank split.  If it is,
      prove the corresponding \(RCQ_3\) subcase; if not, isolate the
      exact off-diagonal block.
\item Follow every strict mass arrow to its terminal row and use
      \(p\ge\eta\), \(y\ge\eta\) to test whether the terminal mass has
      an already nondepleted literal edge.
\item Keep payer-bank and selected-coordinate terms finite and
      separate; do not absorb them into the unmarked bank length.
\item Begin the local-payer polarization only after the outside-local
      bank insertion has been decided.
\item No companion audit is due before V60.
\end{enumerate}
\end{programme}

\begin{firewall}[Claim boundary after third-series V57]
\label{fire:V57-claim-boundary}
V57 proves a quadratic currency for an owned comparable unmarked
bank and finite termination of mass escalation.  It does not insert
that currency through every physical compression, prove
\(RCQ_3(\eta)\), or close the row-\(3\) face.  The composite and
row-\(4\) cuts, other quartic source families, higher MTA, WUFC, CPM,
cutoff removal, reflection positivity, Osterwalder--Schrader
reconstruction, construction of the continuum Yang--Mills measure,
and a uniform positive continuum spectral gap remain open.  The full
Yang--Mills mass gap has not been proved.
\end{firewall}

\paragraph{Parent-version record.}
V57 was copied byte-for-byte from
\texttt{Renewal\_Yang\_Mills\_Mass\_Gap\_Research\_Notes\_V56.tex}
before this chapter was appended.  The V56 parent had \(28\,258\)
source lines, \(1\,019\,800\) bytes, and SHA--256
\[
 \texttt{c15653d0deb1c2b43465f3f0e335417d4324652c0e8a0ad8a02c9ae702404ba7}.
\]
The next compulsory companion audit is V60.

% =====================================================================
% V58 APPENDIX: CAPACITY-PRESERVING BANK INSERTION
% =====================================================================

\chapter{V58: Capacity-Preserving Separator Norms and Exact Bank
Insertion}
\label{ch:V58-insertion}

\section{Why coordinate factorization is not yet capacity
factorization}
\label{sec:V58-purpose}

V57 proves the required \((x+z)^2\) capacity for an owned comparable
unmarked bank.  A full \(RCQ_3\) estimate still requires inserting
that bank into the exact source square.  Coordinate tensor products
can be exposed by the V14 restriction identity, but the output
dualizer, recoupling maps, and physical compression remain as bounded
separators.  Their ordinary operator norms do not preserve
product-state capacity across the cut
\[
 \mathcal H_3\mid\mathcal H_{\{1,2,4\}}.
\]

This chapter identifies the correct sufficient norm: the projective
operator tensor norm across that row cut.  It proves a dressed-bank
insertion theorem under a uniform projective bound and gives a sharp
norm-one counterexample showing why an ordinary separator norm cannot
replace it.

\begin{assessment}[Nonduplication boundary]
\label{ass:V58-nonduplication}
V14 proves exact subset resolution while retaining finite separators.
V22 studies dressed endpoint capacities through commutators.  V57
proves a new quadratic bank currency before dressing.  V58 does not
repeat these results: it supplies the missing functional-analytic
criterion under which the V57 currency survives a separator, and
proves that operator norm alone cannot supply that criterion.
\end{assessment}

\section{The cut-projective separator norm}
\label{sec:V58-projective}

Let \(\mathcal H=\mathcal H_3\otimes\mathcal H_c\), where
\(\mathcal H_c=\mathcal H_{\{1,2,4\}}\), including all finite
multiplicity fibres.  For an operator \(S\in\mathcal B(\mathcal H)\),
define
\begin{equation}
 \|S\|_{\pi,3|c}
 :=
 \inf\left\{
 \sum_{\nu=1}^m\|A_\nu\|\,\|B_\nu\|:
 S=\sum_{\nu=1}^mA_\nu\otimes B_\nu
 \right\}.
\label{eq:V58-projective-norm}
\end{equation}
The infimum is finite in every physical finite-cutoff block.

\begin{lemma}[Projective separators preserve product capacity]
\label{lem:V58-projective-capacity}
For every positive operator \(M\) on \(\mathcal H\) and every
operator \(S\),
\begin{equation}
 \boxed{
 \kappa_\otimes(S^*MS)
 \le
 \|S\|_{\pi,3|c}^{\,2}\,
 \kappa_\otimes(M).}
\label{eq:V58-projective-capacity}
\end{equation}
\end{lemma}

\begin{proof}
Fix a decomposition
\(S=\sum_\nu A_\nu\otimes B_\nu\) and a unit product vector
\(\zeta=v\otimes w\).  Put
\[
 z_\nu=A_\nu v\otimes B_\nu w.
\]
Positivity of \(M\) makes
\(\langle\cdot,M\cdot\rangle\) a semidefinite inner product, so
Cauchy--Schwarz gives
\[
 |\langle z_\mu,Mz_\nu\rangle|
 \le
 \langle z_\mu,Mz_\mu\rangle^{1/2}
 \langle z_\nu,Mz_\nu\rangle^{1/2}.
\]
Each \(z_\nu\) is a scalar multiple of a product vector.  Therefore
\[
 \langle z_\nu,Mz_\nu\rangle
 \le
 \kappa_\otimes(M)
 \|A_\nu v\|^2\|B_\nu w\|^2.
\]
Summing the cross terms yields
\[
 \langle S\zeta,MS\zeta\rangle
 \le
 \kappa_\otimes(M)
 \left(\sum_\nu\|A_\nu\|\,\|B_\nu\|\right)^2.
\]
Take the supremum over product vectors and then the infimum over
decompositions.
\end{proof}

\begin{corollary}[Product separators are harmless]
\label{cor:V58-product-separator}
If \(S=A\otimes B\), then
\begin{equation}
 \kappa_\otimes(S^*MS)
 \le\|A\|^2\|B\|^2\kappa_\otimes(M).
\label{eq:V58-product-separator}
\end{equation}
A sum of a uniformly bounded number of product separators is harmless
provided the sum of their product norms is uniformly bounded.
\end{corollary}

\begin{proof}
Use the displayed decomposition in
\eqref{eq:V58-projective-norm} and apply
\Cref{lem:V58-projective-capacity}.
\end{proof}

\section{Oriented insertion into an operator square}
\label{sec:V58-oriented}

\begin{theorem}[Dressed positive-bank insertion]
\label{thm:V58-dressed-bank}
Let \(0\preccurlyeq M\preccurlyeq I\), and let
\[
 W=LMS.
\]
If
\begin{equation}
 \|L\|\le L_0,\qquad
 \|S\|_{\pi,3|c}\le K_0,
\label{eq:V58-dressing-hypotheses}
\end{equation}
then
\begin{equation}
 \boxed{
 \kappa_\otimes(W^*W)
 \le
 L_0^2K_0^2\kappa_\otimes(M).}
\label{eq:V58-dressed-bank}
\end{equation}
There is a left-square analogue: if the left dressing has bounded
projective norm and the right dressing has bounded operator norm,
then the same estimate holds for \(WW^*\).
\end{theorem}

\begin{proof}
Since \(L^*L\preccurlyeq L_0^2I\),
\[
 W^*W
 =
 S^*ML^*LMS
 \preccurlyeq
 L_0^2S^*M^2S
 \preccurlyeq
 L_0^2S^*MS,
\]
where the last inequality uses \(0\preccurlyeq M\preccurlyeq I\).
Monotonicity of product-state capacity and
\Cref{lem:V58-projective-capacity} prove
\eqref{eq:V58-dressed-bank}.  Taking adjoints and reversing the
orientation proves the left-square statement.
\end{proof}

\begin{corollary}[Insertion of the V57 bank currency]
\label{cor:V58-V57-insertion}
Let \(M_D\) be a protected or gapped owned comparable bank from
\Cref{thm:V57-bank-quadratic}.  If a restriction-resolved source word
has the form
\[
 W_\nu=L_\nu M_DS_\nu
\]
and one of the two orientations in
\Cref{thm:V58-dressed-bank} has uniform constants, then
\begin{equation}
 \kappa_\otimes(W_\nu^*W_\nu)
 \le C(x+z)^2
\label{eq:V58-bank-RCQ}
\end{equation}
in the right-square orientation, or the corresponding left-square
bound.  The constant is uniform in the bank length and
representations.
\end{corollary}

\begin{proof}
By
\Cref{thm:V57-bank-quadratic,cor:V57-bank-pays-cut},
\[
 \kappa_\otimes(M_D)\le C(x+z)^2.
\]
Apply \Cref{thm:V58-dressed-bank}.
\end{proof}

\section{Exact restriction of the \(2+1+1\) word}
\label{sec:V58-restriction}

Fix a physical coordinate set \(D\) disjoint from the selected
joining coordinate and all marked or payer coordinates.  Apply the
moment and covariance restriction identities
\eqref{eq:V14-moment-split}--%
\eqref{eq:V14-covariance-split} to every complete occurrence in one
fixed \(2+1+1\) source word.

\begin{theorem}[Finite subset resolution of a selected
\(2+1+1\) word]
\label{thm:V58-subset-resolution}
Before any positive envelope, output norm, or capacity is taken, a
fixed selected-coordinate \(2+1+1\) word is an exact sum
\begin{equation}
 \widehat W_\pi=\sum_{\nu=1}^{N_*}W_{\pi,\nu},
 \qquad
 N_*\le2^{J_*},
\label{eq:V58-subset-resolution}
\end{equation}
where \(J_*\) depends only on the fixed four-row source atlas.  In
each resolved word:
\begin{enumerate}[label=\textup{(\roman*)},leftmargin=4em]
\item every complete moment factors between \(D\) and \(D^c\);
\item every covariance is localized to one side and retains a
      complete endpoint moment on the other;
\item the sole payer remains at its original occurrence;
\item coordinate order and the selected joining letter remain
      unchanged; and
\item every output dualizer, recoupling, and physical compression
      remains in its original finite separator position.
\end{enumerate}
The number of resolved words is independent of support,
representations, and physical multiplicities.
\end{theorem}

\begin{proof}
The selected \(2+1+1\) word has a fixed finite number of covariance
occurrences: the complete binary prefix covariance, the assembled
local cumulant written as a fixed finite moment combination, and any
fixed covariance occurrence in the payer representation.  Apply
\eqref{eq:V14-covariance-split} to each covariance and
\eqref{eq:V14-moment-split} to each moment, then distribute the
finite binary sums.  This gives at most \(2^{J_*}\) terms.

All separator maps multiply the resulting identity without being
expanded.  Because \(D\) excludes payer coordinates, the payer label
does not split; the general split
\Cref{lem:V14-payer-split} shows that it would still remain unique if
the exclusion were relaxed.  No step depends on \(|D|\).
\end{proof}

\begin{lemma}[Finite recombination costs only a fixed constant]
\label{lem:V58-finite-recombination}
If every term in \eqref{eq:V58-subset-resolution} satisfies
\[
 \kappa_\otimes(W_{\pi,\nu}^*W_{\pi,\nu})
 \le C_\nu(x+z)^2,
\]
then
\begin{equation}
 \kappa_\otimes(\widehat W_\pi^*\widehat W_\pi)
 \le
 N_*\left(\sum_{\nu=1}^{N_*}C_\nu\right)(x+z)^2.
\label{eq:V58-recombination}
\end{equation}
\end{lemma}

\begin{proof}
For arbitrary operators,
\[
 \left(\sum_\nu W_\nu\right)^*
 \left(\sum_\nu W_\nu\right)
 \preccurlyeq
 N_*\sum_\nu W_\nu^*W_\nu.
\]
This follows by applying
\(\|\sum_\nu z_\nu\|^2\le N_*\sum_\nu\|z_\nu\|^2\) to every vector.
Use positivity and subadditivity of
\(\kappa_\otimes\).
\end{proof}

\begin{assessment}[Exact insertion interface]
\label{ass:V58-insertion-interface}
For an owned comparable bank, the only missing hypothesis in
\Cref{cor:V58-V57-insertion} is now explicit.  In each
restriction-resolved word, one orientation must have:
\begin{enumerate}[label=\textup{(\roman*)},leftmargin=4em]
\item an ordinary uniform norm on the dressing adjacent to the first
      copy of the bank; and
\item a uniform row-\(3|3^c\) projective norm on the dressing through
      which product states enter the other copy.
\end{enumerate}
If these hold term by term, \Cref{lem:V58-finite-recombination}
closes that complete-bank subcase of \(RCQ_3(\eta)\).
\end{assessment}

\section{Ordinary separator norm is insufficient}
\label{sec:V58-no-norm}

\begin{theorem}[Sharp entangling-separator obstruction]
\label{thm:V58-entangler-no-go}
For every \(d\ge2\), there are Hilbert spaces
\(\mathcal H_A=\mathcal H_B=\mathbb C^d\), a rank-one positive
projection \(P_d\), and a unitary \(U_d\) such that
\begin{align}
 \kappa_\otimes(P_d)&=\frac1d,
\label{eq:V58-small-capacity}\\
 \|U_d\|&=1,
\label{eq:V58-unitary-norm}\\
 \kappa_\otimes(U_d^*P_dU_d)&=1.
\label{eq:V58-entangled-amplification}
\end{align}
Therefore no constant \(C\), independent of \(d\), can satisfy
\begin{equation}
 \kappa_\otimes(S^*MS)
 \le C\|S\|^2\kappa_\otimes(M)
\label{eq:V58-false-norm-insertion}
\end{equation}
for all positive \(M\) and bounded \(S\).
\end{theorem}

\begin{proof}
Let \(P_d=|\Omega_d\rangle\langle\Omega_d|\) be the maximally
entangled invariant projection of
\Cref{lem:V56-idempotent-capacity}.  Then
\(\kappa_\otimes(P_d)=1/d\).  Choose a unitary \(U_d\) satisfying
\[
 U_d(e_1\otimes e_1)=\Omega_d;
\]
such a unitary exists because any two unit vectors extend to
orthonormal bases.  The product vector \(e_1\otimes e_1\) then gives
\[
 \langle e_1\otimes e_1,
 U_d^*P_dU_d(e_1\otimes e_1)\rangle=1.
\]
The capacity cannot exceed the norm of the projection, so it equals
one.  Inserting these objects into
\eqref{eq:V58-false-norm-insertion} would give \(1\le C/d\).
\end{proof}

\begin{corollary}[The projective norm detects the entangler]
\label{cor:V58-entangler-projective}
The unitary above satisfies
\begin{equation}
 \|U_d\|_{\pi,A|B}\ge\sqrt d.
\label{eq:V58-entangler-projective}
\end{equation}
Thus the projective criterion excludes exactly the
dimension-dependent amplification hidden by the ordinary norm.
\end{corollary}

\begin{proof}
Apply \Cref{lem:V58-projective-capacity} with
\(M=P_d\) and \(S=U_d\):
\[
 1
 \le
 \|U_d\|_{\pi,A|B}^2/d.
\]
\end{proof}

\begin{firewall}[Orthogonality is not separability]
\label{fire:V58-orthogonal-not-product}
An orthogonal output compression has operator norm one, but this does
not imply a uniform projective norm across the row cut.  Statements
such as the V38 absence of an output multiplicity count concern
operator norm after orthogonal block resolution; they do not by
themselves prove product-capacity preservation through the block
map.
\end{firewall}

\section{The commutator alternative}
\label{sec:V58-commutator}

\begin{proposition}[Two valid insertion routes]
\label{prop:V58-two-routes}
For each physical separator adjacent to an owned comparable bank,
one of the following independent routes would suffice:
\begin{enumerate}[label=\textup{(\alph*)},leftmargin=4em]
\item prove a uniform cut-projective norm and apply
      \Cref{thm:V58-dressed-bank}; or
\item retain the separator inside the square and prove the
      corresponding paid commutator-square estimate required by the
      V22 dressed-endpoint identity.
\end{enumerate}
Ordinary norm control without either input is insufficient.
\end{proposition}

\begin{proof}
Route \textup{(a)} is
\Cref{thm:V58-dressed-bank}.  V22 decomposes a dressed endpoint
square into the undressed quadratic energy plus an exact commutator
square; its compiler closes once both are bounded.  This is route
\textup{(b)}.  The impossibility of using only the ordinary norm is
\Cref{thm:V58-entangler-no-go}.
\end{proof}

\begin{assessment}[Upstream object selected for V59]
\label{ass:V58-upstream}
The next calculation is no longer an abstract bank estimate.  It is
the row-\(3|3^c\) tensor structure of the actual finite separator
formed by:
\begin{enumerate}[label=\textup{(\roman*)},leftmargin=4em]
\item coordinatewise Clebsch--Gordan recoupling at the boundary of
      the unmarked bank;
\item the physical output isotypic compression;
\item the V94 output dualizer; and
\item the unique payer if it lies outside the bank.
\end{enumerate}
V59 must compute its projective norm or its commutator with the
protected/gapped bank.  Treating it merely as a norm-one map would
repeat the obstruction in
\Cref{thm:V58-entangler-no-go}.
\end{assessment}

\section{V58 status and next acquisition}
\label{sec:V58-status}

\begin{theorem}[Third-series V58 result ledger]
\label{thm:V58-results}
Relative to V1--V57, V58 proves:
\begin{enumerate}[label=\textup{(V58.\arabic*)},leftmargin=6.3em]
\item the exact row-\(3|3^c\) projective operator norm relevant to
      product-state capacity;
\item a capacity-preservation theorem for every projectively bounded
      separator;
\item an oriented dressed-bank insertion theorem;
\item conditional insertion of the V57 quadratic bank currency at
      the \(RCQ_3\) scale;
\item finite exact subset resolution of the selected
      \(2+1+1\) word and support-uniform recombination;
\item a sharp dimension-dependent counterexample to insertion using
      ordinary separator norm alone;
\item a lower bound showing that the projective norm detects the
      entangling amplification; and
\item reduction of the remaining insertion problem to either an
      actual separator projective bound or an exact commutator square.
\end{enumerate}
\end{theorem}

\begin{proof}
These are
\Cref{lem:V58-projective-capacity,
thm:V58-dressed-bank,
cor:V58-V57-insertion,
thm:V58-subset-resolution,
lem:V58-finite-recombination,
thm:V58-entangler-no-go,
cor:V58-entangler-projective,
prop:V58-two-routes}.
\end{proof}

\begin{programme}[V59: actual separator tensor audit]
\label{prog:V58-V59}
\begin{enumerate}[leftmargin=3em]
\item Extract from the archived V92--V94 construction the exact
      matrix of the output dualizer and adjacent recouplings on one
      fixed physical output block.
\item Compute or bound its operator Schmidt/projective norm across
      row \(3|3^c\), retaining multiplicity spaces.
\item If that norm grows with a representation dimension, identify
      the exact maximally entangled channel and move to the V22
      commutator route rather than weakening the target.
\item If the norm is uniform after the complete block is assembled,
      insert \Cref{thm:V57-bank-quadratic} and close the owned
      comparable unmarked-bank subcase of \(RCQ_3(\eta)\).
\item Keep local-payer words separate.
\item No companion audit is due before V60.
\end{enumerate}
\end{programme}

\begin{firewall}[Claim boundary after third-series V58]
\label{fire:V58-claim-boundary}
V58 proves an exact sufficient insertion interface and a sharp
ordinary-norm obstruction.  It does not verify that interface for
the physical V94 separator, prove its commutator card, or close
\(RCQ_3(\eta)\).  The other quotient cuts, quartic source families,
higher MTA, WUFC, CPM, cutoff removal, reflection positivity,
Osterwalder--Schrader reconstruction, construction of the continuum
Yang--Mills measure, and a uniform positive continuum spectral gap
remain open.  The full Yang--Mills mass gap has not been proved.
\end{firewall}

\paragraph{Parent-version record.}
V58 was copied byte-for-byte from
\texttt{Renewal\_Yang\_Mills\_Mass\_Gap\_Research\_Notes\_V57.tex}
before this chapter was appended.  The V57 parent had \(28\,746\)
source lines, \(1\,035\,444\) bytes, and SHA--256
\[
 \texttt{697efa02ffe4ca340b17ce87ec5d04e885eacf364f6fc9d9d80226907942bd89}.
\]
The next compulsory companion audit is V60.

% =====================================================================
% V59 APPENDIX: CENTRAL BANK COMMUTATION AND PHYSICAL INSERTION
% =====================================================================

\chapter{V59: Central Bank Commutation and the Exact Physical
Insertion}
\label{ch:V59-central-insertion}

\section{Upstream correction: the V94 dualizer is not a source
factor}
\label{sec:V59-dualizer-audit}

V58 deliberately left two possible insertion routes: a uniform
row-
\(3|3^c\) projective norm for the physical dressing, or a localized
commutator square.  Its V59 programme included the V94 output dualizer
among the operators whose projective norm should be computed.  A
return to the exact V94 construction shows that this item was placed
in the wrong category.  The output dualizer is a test contraction
introduced by Schatten duality after the complete source operator has
been formed.  It is not a fixed operator in
\(\widehat W_\pi\), and hence is not a fixed dressing in
\(\widehat W_\pi^*\widehat W_\pi\).

\begin{proposition}[Exact logical position of the V94 contractions]
\label{prop:V59-dualizer-position}
With the notation of \Cref{thm:V94-exact-block-duality}, the family
\(\boldsymbol D=(D_U)_U\) occurs only in the dual representation
\begin{equation}
 \sum_Uw_U\|Y_U(X)\|_1
 =
 \sup_{\|D_U\|\le1}
 \left|
  \Tr\!\left[X{\cal T}_{w,K}(\boldsymbol D)\right]
 \right|.
\label{eq:V59-dualizer-position}
\end{equation}
For a complete equivariant sandwich
\(K_U=a_Uc_Ub_U\), it is absorbed in the proof into
\begin{equation}
 C_{\boldsymbol D}
 =
 \sum_UW_U(D_U^*\otimes c_U)W_U^*,
 \qquad
 \|C_{\boldsymbol D}\|
 \le\sup_U\|c_U\|.
\label{eq:V59-dualizer-middle}
\end{equation}
Neither \(D_U\) nor \(C_{\boldsymbol D}\) is part of the source word
before the trace-norm duality in
\eqref{eq:V59-dualizer-position}.  Consequently no fixed matrix of an
``output dualizer'' exists whose cut-projective norm could be inserted
into \Cref{thm:V58-dressed-bank}.
\end{proposition}

\begin{proof}
Equation \eqref{eq:V59-dualizer-position} is the identity in
\Cref{thm:V94-exact-block-duality}.  In its proof, \(D_U\) is chosen
as a polar dual contraction for the already formed trace-class block
\(Y_U(X)\).  Thus its choice depends on the block being tested and
occurs under the supremum.  When
\(K_U=a_Uc_Ub_U\), the block calculation in
\Cref{thm:V94-dual-sandwich} gives
\eqref{eq:V59-dualizer-middle} and
\({\cal T}_{1,K}(\boldsymbol D)=AC_{\boldsymbol D}B\).
This is a dual proof device, not a factorization of the original
source operator.  The norm estimate follows from orthogonality of the
output blocks and \(\|D_U\|\le1\).
\end{proof}

\begin{remark}[Why equivariance cannot be assigned to the dualizer]
\label{rem:V59-dualizer-not-equivariant}
On an output block,
\begin{equation}
 [D_U^*\otimes c_U,\,\pi_U(g)\otimes I]
 =
 [D_U^*,\pi_U(g)]\otimes c_U.
\label{eq:V59-dualizer-commutator}
\end{equation}
The contraction \(D_U\) is arbitrary.  If \(U\) is irreducible, it
commutes with every \(\pi_U(g)\) only when it is scalar, by Schur's
lemma.  Therefore the central-commutation argument below applies to
the physical equivariant source factors, not to
\(C_{\boldsymbol D}\).  This distinction is consistent with V94:
that chapter uses only the ordinary norm of the dual middle operator
inside a linear trace estimate.
\end{remark}

\begin{assessment}[Correction to the V58 projective audit]
\label{ass:V59-correction}
The projective criterion of \Cref{lem:V58-projective-capacity} remains
valid.  What is corrected is its proposed input list.  The V94
dualizer must be removed from the fixed physical separator in
\Cref{ass:V58-upstream}.  The actual RCQ square must instead be audited
at the level of its source factors: complete moments and cumulants,
canonical recouplings, physical isotypic compressions, off-bank
coordinate factors, and the unique payer at its genuine location.
\end{assessment}

\section{Centrality of the protected and gapped bank sectors}
\label{sec:V59-centrality}

The source factors in the preceding list have more structure than a
generic norm-one separator.  They are coordinatewise intertwiners.
The V32 protected/gapped split is spectral for the same local group
action and is therefore central in the corresponding intertwiner
algebra.

\begin{lemma}[Central isotypic functional calculus]
\label{lem:V59-central-isotypic}
Let a finite-dimensional unitary \(G\)-module have isotypic
decomposition
\begin{equation}
 \mathcal H
 =
 \bigoplus_{\lambda\in\Lambda}
 V_\lambda\otimes M_\lambda.
\label{eq:V59-isotypic-decomposition}
\end{equation}
If \(T\) commutes with the \(G\)-action, then
\begin{equation}
 T=\bigoplus_{\lambda\in\Lambda}
 I_{V_\lambda}\otimes t_\lambda.
\label{eq:V59-commutant-form}
\end{equation}
For every bounded scalar function \(f\) on \(\Lambda\), the operator
\begin{equation}
 Z_f
 :=
 \bigoplus_{\lambda\in\Lambda}
 f(\lambda)I_{V_\lambda}\otimes I_{M_\lambda}
\label{eq:V59-central-spectral}
\end{equation}
commutes with every such \(T\).  If \(0\le f\le1\), then
\(0\preccurlyeq Z_f\preccurlyeq I\).
\end{lemma}

\begin{proof}
The commutant form \eqref{eq:V59-commutant-form} follows from Schur's
lemma, including arbitrary multiplicity spaces.  Both products
\(TZ_f\) and \(Z_fT\) restrict on the \(\lambda\)-block to
\(f(\lambda)I_{V_\lambda}\otimes t_\lambda\), proving commutation.
The last assertion follows blockwise.
\end{proof}

\begin{lemma}[Complete-bank commutation]
\label{lem:V59-bank-commutation}
Let \(D\) be a finite coordinate bank.  At coordinate \(e\in D\),
let \(Z_e\) be either an isotypic projection or a normalized positive
heat multiplier which is scalar on each local isotypic block.  Put
\begin{equation}
 \mathcal M_D:=\bigotimes_{e\in D}Z_e.
\label{eq:V59-bank-central-operator}
\end{equation}
After the canonical replica tensor factors have been identified, let
\(Q\) be a finite sum of finite products of:
\begin{enumerate}[label=\textup{(\roman*)},leftmargin=4em]
\item coordinatewise complete moment or cumulant intertwiners;
\item canonical equivariant recouplings;
\item local physical isotypic compressions;
\item operators supported in \(D^c\).
\end{enumerate}
Then
\begin{equation}
 [\mathcal M_D,Q]=0.
\label{eq:V59-bank-commutation}
\end{equation}
Moreover \(0\preccurlyeq\mathcal M_D\preccurlyeq I\).
\end{lemma}

\begin{proof}
For each \(e\in D\), the first three kinds of local factors are
intertwiners for the local action.  For moments and cumulants this is
\Cref{lem:V94-separator-equivariant}; canonical recoupling merely
changes an equivariant realization, and an isotypic compression is
itself a central spectral projection.  After the canonical
identification, \Cref{lem:V59-central-isotypic} shows that every such
factor commutes with \(Z_e\).  A factor supported in \(D^c\) commutes
tensorially.  Commutation is preserved under products and finite
sums, giving \eqref{eq:V59-bank-commutation}.  Positivity and the
contraction bound tensorize from the final assertion of
\Cref{lem:V59-central-isotypic}.
\end{proof}

\begin{corollary}[The V32 bank envelopes are central contractions]
\label{cor:V59-V32-central}
For every complete unmarked bank sector used in
\Cref{cor:V32-V70-currencies}, the protected envelope
\(\mathcal P_{3u}(D)\) and the positive gapped heat envelope
\(\mathcal G_{3u}(D)\) are positive contractions of the form
\eqref{eq:V59-bank-central-operator}.  They commute with every physical
source separator in \Cref{lem:V59-bank-commutation}.
\end{corollary}

\begin{proof}
The protected local factor is the projection onto the all-trivial
isotypic sector from \Cref{cor:V32-protected-capacity}.  The gapped
factor is the complementary isotypic sum with normalized multiplier
\(e^{-s_\alpha C_2(\lambda)}\), as in
\Cref{lem:V32-output-gap}.  Both are scalar functions of the local
isotypic label with values in \([0,1]\).  Tensor over the complete
unmarked bank and apply \Cref{lem:V59-bank-commutation}.
\end{proof}

\begin{firewall}[The exact symmetry being used]
\label{fire:V59-symmetry-scope}
Centrality is invoked only for a source factor which is an intertwiner
for every local action meeting the bank.  Mere equivariance for a
coarser diagonal action would not imply
\eqref{eq:V59-bank-commutation}.  In the restriction-resolved source,
the required local actions are retained because \(D\) is a complete
coordinate bank, while marked and payer coordinates have been
removed before the bank is selected.  No arbitrary output dual
contraction is included; see \Cref{prop:V59-dualizer-position}.
\end{firewall}

\section{Commuting insertion into the quadratic source}
\label{sec:V59-insertion}

\begin{theorem}[Central positive-bank insertion]
\label{thm:V59-central-insertion}
Let \(0\preccurlyeq M\preccurlyeq I\), let \([M,Q]=0\), and set
\begin{equation}
 W=LMQ.
\label{eq:V59-central-word}
\end{equation}
If \(\|L\|\le L_0\) and \(\|Q\|\le Q_0\), then
\begin{equation}
 \boxed{
 \kappa_\otimes(W^*W)
 \le
 L_0^2Q_0^2\kappa_\otimes(M^2)
 \le
 L_0^2Q_0^2\kappa_\otimes(M).}
\label{eq:V59-central-insertion}
\end{equation}
The adjoint-oriented assertion holds for \(WW^*\) when the central
bank lies on the opposite side.
\end{theorem}

\begin{proof}
Because \(M=M^*\), the relation \([M,Q]=0\) also gives
\([M,Q^*]=0\).  Hence
\begin{align*}
 W^*W
 &=Q^*ML^*LMQ\\
 &\preccurlyeq L_0^2Q^*M^2Q\\
 &=L_0^2M Q^*Q M\\
 &\preccurlyeq L_0^2Q_0^2M^2.
\end{align*}
Take product-state capacity.  Finally
\(M^2\preccurlyeq M\) by positive functional calculus.  The
left-square assertion follows after taking adjoints.
\end{proof}

\begin{remark}[Relation to V22 and V58]
\label{rem:V59-relation}
Theorem \ref{thm:V59-central-insertion} is the zero-commutator branch
of \Cref{cor:V22-commuting-safe}, specialized to the V57 positive bank
currency and written in the exact orientation needed here.  It is not
a replacement for \Cref{thm:V58-dressed-bank}: the latter applies to
noncommuting projectively bounded dressings.  Centrality makes a
projective estimate unnecessary for the present physical subcase.
\end{remark}

\section{Closure of the owned comparable unmarked-bank subcase}
\label{sec:V59-physical-closure}

Fix an outside-local payer placement and apply
\Cref{thm:V58-subset-resolution} to a complete selected-coordinate
\(2+1+1\) word.  Suppose one resolved term contains a nonempty owned
comparable unmarked row-\(3\)--\(u\) bank \(D\), and resolve that bank
into one of its protected or gapped positive envelopes
\(\mathcal M_D\).

\begin{theorem}[Exact physical insertion of an owned bank]
\label{thm:V59-physical-owned-bank}
For every such resolved term, the normalized word can be written,
after commuting only central bank factors through physical source
intertwiners, as
\begin{equation}
 \widetilde W_{\pi,\nu}
 =L_{\pi,\nu}\mathcal M_DQ_{\pi,\nu},
 \qquad
 [\mathcal M_D,Q_{\pi,\nu}]=0,
\label{eq:V59-physical-factorization}
\end{equation}
where
\begin{equation}
 \|L_{\pi,\nu}\|+\|Q_{\pi,\nu}\|
 \le C_s
\label{eq:V59-physical-norms}
\end{equation}
uniformly in the bank length, support, representations, physical
multiplicities, and selected coordinate.  Therefore
\begin{equation}
 \boxed{
 \kappa_\otimes(
  \widetilde W_{\pi,\nu}^*\widetilde W_{\pi,\nu})
 \le C_s\theta_{3u}(D)^2
 \le C_s(x+z)^2.}
\label{eq:V59-owned-term-RCQ}
\end{equation}
The same conclusion holds in the left-square orientation.
\end{theorem}

\begin{proof}
The exact V14 restriction used in
\Cref{thm:V58-subset-resolution} factors every complete moment between
\(D\) and \(D^c\) before a positive envelope or capacity is taken.
The protected/gapped resolution on \(D\) produces
\(\mathcal M_D\).  All marked coordinates and the unique payer were
excluded from \(D\).  Factors supported in \(D^c\) commute
tensorially with \(\mathcal M_D\); local complete moments, cumulants,
recouplings, and physical isotypic compressions meeting \(D\) commute
by \Cref{cor:V59-V32-central}.  Moving only these commuting factors
gives \eqref{eq:V59-physical-factorization} without changing their
order relative to one another.

Every unpaid complete moment and physical compression is a
contraction.  The fixed local cumulant and joining separators have the
uniform norms in \Cref{lem:V92-separator-size}.  The unique
outside-local payer has its complete one-payer norm, and the binary
prefix is kept as one complete factor exactly as in
\Cref{lem:V55-complete-prefix-response}.  After division by
\(\Xi_1\Xi_2\Xi_3\Xi_4\), the active-row lower bound
\(\Xi_i\ge a_*\) absorbs the unused denominators.  This proves
\eqref{eq:V59-physical-norms}.

Apply \Cref{thm:V59-central-insertion}, then
\Cref{thm:V57-bank-quadratic,cor:V57-idempotence-resolved}:
\[
 \kappa_\otimes(\mathcal M_D^2)
 \le C_s\theta_{3u}(D)^2.
\]
Finally use \Cref{cor:V57-bank-pays-cut}.  The adjoint orientation is
identical.
\end{proof}

\begin{corollary}[Finite recombination closes the owned-bank packet]
\label{cor:V59-owned-packet}
If every nonzero restriction-resolved term of a fixed outside-local
packet contains an owned comparable unmarked bank sector as in
\Cref{thm:V59-physical-owned-bank}, then that packet satisfies
\(RCQ_3(\eta)\).
\end{corollary}

\begin{proof}
Each term satisfies \eqref{eq:V59-owned-term-RCQ}.  The number of terms
is the fixed \(N_*\) of \Cref{thm:V58-subset-resolution}.  Apply
\Cref{lem:V58-finite-recombination}.
\end{proof}

\begin{assessment}[What was actually closed]
\label{ass:V59-closed-scope}
The bank insertion obstruction in
\Cref{fire:V57-insertion-boundary} is removed for complete owned
comparable unmarked protected/gapped banks and outside-local payer
placements.  The result uses neither a dimension-dependent
cut-projective norm nor a false capacity invariance under arbitrary
unitaries.  It uses centrality of the actual physical bank split.

This does not close a restriction-resolved term whose cofinal mass is
confined to marked or payer coordinates, whose balanced bank fails
two-row ownership, or whose mass-escalation chain terminates at a row
without an owned bank.  It also does not treat a joining-coordinate
local payer.
\end{assessment}

\section{V59 result ledger and V60 audit mandate}
\label{sec:V59-status}

\begin{theorem}[Third-series V59 result ledger]
\label{thm:V59-results}
Relative to V1--V58, V59 proves:
\begin{enumerate}[label=\textup{(V59.\arabic*)},leftmargin=6.3em]
\item the V94 output dualizer is a variable trace-duality test and not
      a fixed factor of the quadratic source word;
\item its arbitrary block contractions cannot be assigned physical
      equivariance;
\item local isotypic spectral contractions are central in the full
      intertwiner algebra, with multiplicities retained;
\item complete protected and gapped bank envelopes commute with the
      actual coordinatewise equivariant source separators;
\item a sharp central positive-bank insertion inequality;
\item exact physical insertion of the V57 quadratic bank currency for
      every outside-local owned-bank resolved term; and
\item finite recombination and closure of every packet all of whose
      resolved terms possess that bank.
\end{enumerate}
\end{theorem}

\begin{proof}
These are
\Cref{prop:V59-dualizer-position,
rem:V59-dualizer-not-equivariant,
lem:V59-central-isotypic,
lem:V59-bank-commutation,
cor:V59-V32-central,
thm:V59-central-insertion,
thm:V59-physical-owned-bank,
cor:V59-owned-packet}.
\end{proof}

\begin{programme}[V60: compulsory audit and terminal mass routing]
\label{prog:V59-V60}
\begin{enumerate}[leftmargin=3em]
\item Perform the compulsory read-only audit of the newest active SM
      and NS notes, recording exact filenames, sizes, and hashes.
\item Import only type-correct proof-order lessons; do not transfer an
      SM or NS estimate into Yang--Mills without a proved interface.
\item Return to the finite mass-escalation chain of
      \Cref{thm:V57-finite-termination}.  At its terminal row, classify
      the alternatives using the two nondepleted pure-face edges
      \(p\ge\eta\) and \(y\ge\eta\).
\item Route every terminal owned bank through
      \Cref{thm:V59-physical-owned-bank}; isolate the exact scalar
      deficit in the remaining unowned terminal alternatives.
\item Keep finite marked/payer banks and the local-payer square as
      separate branches.
\end{enumerate}
\end{programme}

\begin{firewall}[Claim boundary after third-series V59]
\label{fire:V59-claim-boundary}
V59 closes the exact physical insertion of an owned comparable
unmarked protected/gapped bank for outside-local payer placements.  It
does not prove that every pure row-\(3\) packet contains such a bank,
close the terminal mass-escalation alternatives, prove the
local-payer square, or establish full \(RCQ_3(\eta)\).  The other
quotient cuts, quartic source families, higher MTA, WUFC, CPM, cutoff
removal, reflection positivity, Osterwalder--Schrader reconstruction,
construction of the continuum Yang--Mills measure, and a uniform
positive continuum spectral gap remain open.  The full Yang--Mills
mass gap has not been proved.
\end{firewall}

\paragraph{Parent-version record.}
V59 was copied byte-for-byte from
\texttt{Renewal\_Yang\_Mills\_Mass\_Gap\_Research\_Notes\_V58.tex}
before this chapter was appended.  The V58 parent had \(29\,240\)
source lines, \(1\,051\,969\) bytes, and SHA--256
\[
 \texttt{ef9bfbc3142c38b4a0ea64f8274d3c32a6960fcb791eebd5359f027b94b5b536}.
\]
The compulsory companion audit is due in V60.

% =====================================================================
% V60 APPENDIX: COMPANION AUDIT AND OWNERSHIP COMPENSATION
% =====================================================================

\chapter{V60: Companion Audit, Terminal Rows, and the Missing
Ownership Factor}
\label{ch:V60-compensated-bank}

\section{Compulsory companion audit}
\label{sec:V60-audit}

The active companion files inspected before choosing the V60
direction were:
\begin{enumerate}[leftmargin=3em]
\item
\texttt{Renewal\_SM\_Einstein\_Continuum\_Research\_Notes\_V41.tex},
with \(15\,091\) source lines and \(556\,448\) bytes, SHA--256
\[
 \texttt{faa34f77cb2346d9c117d4a37603ef402542fd78ab1aa21066304088d06d86a0}.
\]
Its newest result obtains a strict nonlinear two-color conditional
angle from two Poincare inequalities, a conditional variance identity,
and a bounded mixed Hessian.  Individual norm-one conditional
projections would not give the strict gain.
\item
\texttt{Renewal\_NS\_Stability\_Research\_Notes\_V31.tex}, with
\(23\,053\) source lines and \(887\,537\) bytes, SHA--256
\[
 \texttt{f1121b62238ad5491bb7cf03635ea9934942edf573ed8faaa9ee79e1d38a6696}.
\]
Its terminal result identifies the missing critical first moment and
proves that heat lifting or a polarization hinge cannot manufacture
that moment after the correlated nonlinear source has been replaced
by separate positive bounds.
\end{enumerate}

\begin{assessment}[Type-correct transfer]
\label{ass:V60-audit-transfer}
No Higgs conditional-angle inequality and no Navier--Stokes heat or
helicity estimate is imported into Yang--Mills.  The common
proof-order lesson is used instead.  A strict gain must be attached to
the relative quantity which actually carries it before coarse
positive estimates are taken.  In the present bank dichotomy, that
relative quantity is the partner-ownership fraction
\(T/\Xi_u\).  Replacing it only by the total-mass arrow
\(\Xi_u\gg\Xi_i\) discards the normalization which can pay the
quadratic edge scale.
\end{assessment}

\section{The terminal row has no fourth combinatorial alternative}
\label{sec:V60-terminal}

Fix a row \(i\).  Let \(F_i\) be the finite set of its marked, selected,
or payer coordinates.  Among the remaining active coordinates, let
\(B_i\) be those balanced in the sense of
\eqref{eq:V32-unbalanced}, and let \(U_i\) be its complement.  Thus
\begin{equation}
 \Xi_i
 =
 \sum_{e\in F_i}a_{i,e}
 +
 \sum_{e\in B_i}a_{i,e}
 +
 \sum_{e\in U_i}a_{i,e}.
\label{eq:V60-terminal-partition}
\end{equation}

\begin{lemma}[Terminal mass trichotomy]
\label{lem:V60-terminal-trichotomy}
Let \(t\) maximize the four total row masses, and fix
\(0<\sigma\le1/3\).  At least one of the following holds:
\begin{enumerate}[label=\textup{(\roman*)},leftmargin=4em]
\item \(F_t\) carries at least \(\sigma\Xi_t\);
\item \(U_t\) carries at least \(\sigma\Xi_t\); or
\item the complete unmarked balanced set \(B_t\) carries at least
      \(\sigma\Xi_t\), and it contains an owned comparable partner
      subbank.
\end{enumerate}
In item \textup{(iii)}, the ownership and comparability constants
depend only on \(\sigma\), the fixed group, and one fixed choice of
\(R>1\), not on support or representations.
\end{lemma}

\begin{proof}
If neither \textup{(i)} nor \textup{(ii)} holds, then
\eqref{eq:V60-terminal-partition} gives
\[
 \sum_{e\in B_t}a_{t,e}
 >(1-2\sigma)\Xi_t
 \ge\sigma\Xi_t.
\]
Assign every coordinate of \(B_t\) to a maximal partner supplied by
\Cref{lem:V32-balanced-partner}.  Apply
\Cref{cor:V57-arbitrary-escalation} to this complete bank with any
fixed \(R>1\).  Its escalation alternative would give a row \(u\)
with
\(\Xi_u\ge R\Xi_t>\Xi_t\), contradicting maximality of \(t\).
The owned comparable alternative must therefore hold.
\end{proof}

\begin{corollary}[The V57 terminal failure is classified]
\label{cor:V60-terminal-classified}
At the terminal row of \Cref{thm:V57-finite-termination}, the
previously unspecified absence of a cofinal balanced bank reduces to
one of two explicit cases: cofinal finite exceptional mass or cofinal
unbalanced mass.  Otherwise there is an owned comparable unmarked
bank, whose physical insertion is covered by
\Cref{thm:V59-physical-owned-bank} after relabelling the root row.
\end{corollary}

\begin{proof}
Apply \Cref{lem:V60-terminal-trichotomy}.  The representation and
centrality arguments used in V57 and V59 are symmetric in the four
row labels.
\end{proof}

\begin{firewall}[Classification is not yet the row-3 estimate]
\label{fire:V60-terminal-not-RCQ}
The terminal owned bank has a quadratic capacity in its own normalized
edge \(\theta_{tu}\).  That edge need not be small when the original
row-\(3\) cut \(x+z\) is small.  Therefore
\Cref{cor:V60-terminal-classified} removes a combinatorial ambiguity,
but does not by itself prove \(RCQ_3(\eta)\).
\end{firewall}

\section{The exact scalar discarded by mass escalation}
\label{sec:V60-beta}

Let \(D\) be a nonempty complete unmarked bank assigned from a root
row \(i\) to a partner row \(u\).  Put
\begin{equation}
 N=|D|,
 \quad
 S=\sum_{e\in D}a_{i,e},
 \quad
 T=\sum_{e\in D}a_{u,e},
 \quad
 H=\sum_{e\in D}a_{i,e}a_{u,e},
\label{eq:V60-bank-data}
\end{equation}
and define
\begin{equation}
 q_D:=\frac{H}{S^2},
 \qquad
 \alpha_D:=\frac{S}{\Xi_i},
 \qquad
 \beta_D:=\frac{T}{\Xi_u},
 \qquad
 \theta_{iu}(D):=\frac{H}{\Xi_i\Xi_u}.
\label{eq:V60-alpha-beta}
\end{equation}

\begin{lemma}[Exact ownership factorization]
\label{lem:V60-ownership-factorization}
Assume only the pointwise comparability
\begin{equation}
 c_0a_{i,e}\le a_{u,e}\le C_0a_{i,e}
 \qquad(e\in D),
\label{eq:V60-comparability}
\end{equation}
where \(0<c_0\le C_0<\infty\).  Then
\begin{align}
 \frac{c_0}{N}\le q_D&\le C_0,
\label{eq:V60-q-range}\\
 \frac{q_D\alpha_D\beta_D}{C_0}
 \le\theta_{iu}(D)
 &\le
 \frac{q_D\alpha_D\beta_D}{c_0}.
\label{eq:V60-theta-factorization}
\end{align}
Thus partner ownership enters the global edge through exactly one
factor \(\beta_D\).
\end{lemma}

\begin{proof}
The proof of \eqref{eq:V60-q-range} is the first part of
\Cref{lem:V57-internal-global}: Cauchy--Schwarz gives the lower bound,
and \(\sum_ea_{i,e}^2\le S^2\) gives the upper bound.  Put
\(r_D=T/S\).  Comparability gives
\(c_0\le r_D\le C_0\).  Directly from the definitions,
\[
 \theta_{iu}(D)
 =
 q_D\frac{S}{\Xi_i}\frac{S}{\Xi_u}
 =
 q_D\alpha_D\frac{\beta_D}{r_D}.
\]
The bounds on \(r_D\) prove
\eqref{eq:V60-theta-factorization}.
\end{proof}

\begin{theorem}[Partner-ownership compensation]
\label{thm:V60-ownership-compensation}
Assume \eqref{eq:V60-comparability} and root ownership
\begin{equation}
 \alpha_D\ge\sigma>0.
\label{eq:V60-root-ownership}
\end{equation}
Let \(\mathcal P_{iu}(D)\) and
\(\mathcal G_{iu}(D)\) be the protected and positive gapped bank
envelopes of \Cref{cor:V32-V70-currencies}.  Then
\begin{align}
 \boxed{
 \beta_D^2
 \kappa_\otimes(\mathcal P_{iu}(D))
 \le C_{G,c_0,C_0,\sigma}
 \theta_{iu}(D)^2,}
\label{eq:V60-protected-compensation}\\
 \boxed{
 \beta_D^2
 \kappa_\otimes(\mathcal G_{iu}(D))
 \le C_{G,s,c_0,C_0,\sigma}
 \theta_{iu}(D)^2.}
\label{eq:V60-gapped-compensation}
\end{align}
No lower bound on \(\beta_D\) is assumed.  The constants are uniform
in \(N\), support, representations, and physical multiplicities.
\end{theorem}

\begin{proof}
For the protected sector,
\Cref{cor:V32-V70-currencies} and
\eqref{eq:V60-q-range} give
\[
 \kappa_\otimes(\mathcal P_{iu}(D))
 \le\frac{C_G}{N}q_D
 \le\frac{C_G}{c_0}q_D^2.
\]
By the lower bound in \eqref{eq:V60-theta-factorization} and
\(\alpha_D\ge\sigma\),
\[
 \beta_D^2q_D^2
 \le\frac{C_0^2}{\sigma^2}\theta_{iu}(D)^2.
\]
This proves \eqref{eq:V60-protected-compensation}.

For the gapped sector, use the exponential entry of
\eqref{eq:V32-gapped-currency}:
\[
 \kappa_\otimes(\mathcal G_{iu}(D))
 \le C_Ge^{-\gamma_GsN}.
\]
Equations \eqref{eq:V60-q-range} and
\eqref{eq:V60-theta-factorization} yield
\[
 \theta_{iu}(D)
 \ge
 \frac{c_0\sigma}{C_0N}\beta_D.
\]
Consequently
\[
 \beta_D^2e^{-\gamma_GsN}
 \le
 \frac{C_0^2}{c_0^2\sigma^2}
 N^2e^{-\gamma_GsN}\theta_{iu}(D)^2.
\]
The supremum of \(N^2e^{-\gamma_GsN}\) over \(N\ge1\) is finite.
This proves \eqref{eq:V60-gapped-compensation}.
\end{proof}

\begin{corollary}[Compensated central insertion]
\label{cor:V60-compensated-insertion}
Suppose a restriction-resolved normalized source word has the exact
form
\begin{equation}
 \widetilde W_{\pi,\nu}
 =
 \beta_D L_{\pi,\nu}\mathcal M_DQ_{\pi,\nu},
\label{eq:V60-compensated-word}
\end{equation}
where \(\mathcal M_D\) is either bank envelope in
\Cref{thm:V60-ownership-compensation}, the root-ownership and
comparability hypotheses hold,
\([\mathcal M_D,Q_{\pi,\nu}]=0\), and the two dressing norms are
uniform.  Then
\begin{equation}
 \kappa_\otimes(
  \widetilde W_{\pi,\nu}^*\widetilde W_{\pi,\nu})
 \le C_s\theta_{iu}(D)^2.
\label{eq:V60-compensated-RCQ}
\end{equation}
In particular, for \(i=3\) this is at most \(C_s(x+z)^2\).
\end{corollary}

\begin{proof}
Apply \Cref{thm:V59-central-insertion} to
\eqref{eq:V60-compensated-word}.  The scalar contributes
\(\beta_D^2\).  Then use
\Cref{thm:V60-ownership-compensation} and, for \(i=3\),
\(\theta_{3u}(D)\le x+z\).
\end{proof}

\begin{assessment}[The revised bank dichotomy]
\label{ass:V60-revised-dichotomy}
On a root-owned comparable bank, there are now two analytically valid
ways to close the quadratic normalized edge:
\begin{enumerate}[label=\textup{(\alph*)},leftmargin=4em]
\item if \(\beta_D\) is bounded below, use the owned-bank theorem
      \Cref{thm:V57-bank-quadratic} and its V59 physical insertion;
\item if \(\beta_D\) is small, retain one exact factor
      \(\beta_D\) in each replica occurrence and use
      \Cref{cor:V60-compensated-insertion}.
\end{enumerate}
The total-mass escalation
\(\Xi_u/\Xi_i\gg1\) is only a coarse certificate that the second
branch may be relevant.  It is not itself the analytic payment.
\end{assessment}

\section{Why terminal strong edges cannot replace ancestry}
\label{sec:V60-no-terminal-substitution}

\begin{proposition}[Sharp scalar separation of the initial and
terminal edges]
\label{prop:V60-terminal-counterexample}
For every \(M\ge1\), consider two coordinates \(a,b\) with nonnegative
row weights
\begin{equation}
 (a_{1,a},a_{2,a},a_{3,a},a_{4,a})=(1,0,1,0),
 \qquad
 (a_{1,b},a_{2,b},a_{3,b},a_{4,b})=(M,M,0,M).
\label{eq:V60-scalar-example}
\end{equation}
Then
\begin{align}
 \Xi_1&=M+1,&
 \Xi_2&=M,&
 \Xi_3&=1,&
 \Xi_4&=M,
\label{eq:V60-example-masses}\\
 p&=\frac{M}{M+1},&
 x+z&=\frac1{M+1},&
 y&=1+\frac{M}{M+1}.
\label{eq:V60-example-edges}
\end{align}
Thus the pure row-\(3\) cut tends to zero, row \(1\) is a terminal
maximal-mass row, and its edges to rows \(2\) and \(4\) tend to one.
No bound of a terminal strong edge by \(C(x+z)\) follows from the
total masses and the pure-face inequalities.
\end{proposition}

\begin{proof}
The displayed masses are immediate.  The only row-\(3\) overlap is
\(H_{13}=1\), so
\(x+z=\theta_{13}=1/(M+1)\).  At coordinate \(b\),
\(H_{12}=H_{14}=H_{24}=M^2\).  Division by the corresponding total
mass products gives
\[
 \theta_{12}=\theta_{14}=\frac{M}{M+1},
 \qquad
 \theta_{24}=1,
\]
which is \eqref{eq:V60-example-edges}.  The ratio
\(\Xi_1/\Xi_3=M+1\) is an arbitrarily strong mass escalation, while
the terminal incident edges stay order one.
\end{proof}

\begin{firewall}[Scope of the scalar example]
\label{fire:V60-example-scope}
The example is a nonnegative normalized-edge configuration, not a
claim that every such table is realized by a fixed compact-group
source word.  It rigorously rules out any proof which uses only the
mass-arrow inequalities and the scalar pure-face bounds.  A valid
physical proof must retain additional coordinate/operator ancestry.
In \Cref{thm:V60-ownership-compensation}, that ancestry is precisely
\(\beta_D\).
\end{firewall}

\section{V60 result ledger and the exact V61 acquisition}
\label{sec:V60-status}

\begin{theorem}[Third-series V60 result ledger]
\label{thm:V60-results}
Relative to V1--V59, V60 proves:
\begin{enumerate}[label=\textup{(V60.\arabic*)},leftmargin=6.3em]
\item the compulsory active SM/NS audit with exact file snapshots;
\item a type-correct proof-order transfer retaining the relative
      ownership factor before positive relaxation;
\item a terminal mass trichotomy which removes the unspecified fourth
      combinatorial failure in V57;
\item the exact factorization of a global normalized edge into
      internal bank currency, root ownership, and partner ownership;
\item protected and gapped partner-ownership compensation without a
      lower bound on the partner fraction;
\item a compensated central-insertion theorem at the quadratic edge
      scale; and
\item a sharp scalar counterexample forbidding replacement of the
      initial small edge by a terminal strong edge.
\end{enumerate}
\end{theorem}

\begin{proof}
These are
\Cref{ass:V60-audit-transfer,
lem:V60-terminal-trichotomy,
cor:V60-terminal-classified,
lem:V60-ownership-factorization,
thm:V60-ownership-compensation,
cor:V60-compensated-insertion,
prop:V60-terminal-counterexample}.
\end{proof}

\begin{programme}[V61: exact extraction of the partner fraction]
\label{prog:V60-V61}
\begin{enumerate}[leftmargin=3em]
\item Return to the V14 restriction identity before the row-
      \(3\)--partner bank is replaced by a protected/gapped envelope.
\item In the partner-nonowned branch, track the complete partner-row
      moment on \(D\) together with the global normalization
      \(\Xi_u^{-1}\).  Determine whether their exact one-use
      contribution is \(T/\Xi_u=\beta_D\).
\item Perform this calculation in both independent replica
      occurrences before applying a square or a triangle inequality.
\item If the factor is exact, combine it with
      \Cref{cor:V60-compensated-insertion} and close the diffuse
      comparable row-\(3\) branch without mass escalation.
\item If it is not exact, display the precise remainder operator and
      move upstream to its complete moment/covariance identity.
\item Keep cofinal finite exceptional mass, unbalanced mass, and the
      joining-coordinate local payer separate.  The next compulsory
      companion audit is V65.
\end{enumerate}
\end{programme}

\begin{firewall}[Claim boundary after third-series V60]
\label{fire:V60-claim-boundary}
V60 proves that the missing partner-ownership factor would exactly
compensate both protected and gapped bank capacities.  It does not yet
prove that this scalar factor occurs in every physical
partner-nonowned source word, close the exceptional or unbalanced
terminal branches, prove the local-payer square, or establish full
\(RCQ_3(\eta)\).  The other quotient cuts, quartic source families,
higher MTA, WUFC, CPM, cutoff removal, reflection positivity,
Osterwalder--Schrader reconstruction, construction of the continuum
Yang--Mills measure, and a uniform positive continuum spectral gap
remain open.  The full Yang--Mills mass gap has not been proved.
\end{firewall}

\paragraph{Parent-version record.}
V60 was copied byte-for-byte from
\texttt{Renewal\_Yang\_Mills\_Mass\_Gap\_Research\_Notes\_V59.tex}
before this chapter was appended.  The V59 parent had \(29\,687\)
source lines, \(1\,069\,418\) bytes, and SHA--256
\[
 \texttt{a5554ebac47f64580b8b271b26cc3d28cae5da36c5bd27b8e5d4f7bff15b7760}.
\]
The next compulsory companion audit is V65.

% =====================================================================
% V61 APPENDIX: THE MOMENT-SPLIT NO-GO AND RESPONSE-CARRIER TARGET
% =====================================================================

\chapter{V61: The Moment-Split No-Go and the Exact Response Carrier
Required for Ownership Compensation}
\label{ch:V61-beta-no-go}

\section{Return to the coefficient-one V14 identity}
\label{sec:V61-V14-audit}

V60 proved that one factor
\(\beta_D=T(D)/\Xi_u\) in each replica occurrence would compensate
the failure of partner ownership.  The first item of its V61 programme
asked whether this factor is already present in the V14 subset
restriction.  It is not.  The V14 identity separates coordinate
operators with coefficient one; it does not insert a scalar row-mass
fraction.

\begin{proposition}[No ownership coefficient in subset restriction]
\label{prop:V61-no-beta-in-V14}
For a complete moment and covariance on a support \(S\), restriction
to \(D\subseteq S\) gives
\begin{align}
 M_S(\rho)
 &=M_{S\cap D}(\rho)\otimes M_{S\cap D^c}(\rho),
\label{eq:V61-moment-one}\\
 \Delta_S^{\rho,\sigma}
 &=\Delta_{S\cap D}^{\rho,\sigma}
   \otimes M_{S\cap D^c}(\rho)
 +M_{S\cap D}(\sigma)
   \otimes\Delta_{S\cap D^c}^{\rho,\sigma}.
\label{eq:V61-covariance-one}
\end{align}
Every displayed coefficient is one.  In the first covariance
summand, the partner moment on \(D^c\) is an endpoint contraction; in
the second, the bank factor on \(D\) is an endpoint moment.  Neither
summand contains the scalar
\(T(D)/\Xi_u\) merely by virtue of this identity.
\end{proposition}

\begin{proof}
Equations \eqref{eq:V61-moment-one} and
\eqref{eq:V61-covariance-one} are
\eqref{eq:V14-moment-split} and
\eqref{eq:V14-covariance-split} with \(E=D\).  They are multilinear
operator identities before normalization or estimates.  The row
weights \(a_{u,e}\), their bank sum \(T(D)\), and the total
\(\Xi_u\) do not occur in either formula.  An estimate on an endpoint
moment may introduce a contraction constant, but a contraction bound
alone is not the vanishing fraction \(T(D)/\Xi_u\).
\end{proof}

\begin{corollary}[V14 resolution preserves but does not create a
response]
\label{cor:V61-resolution-not-response}
The finite word resolution of \Cref{thm:V58-subset-resolution}
preserves every genuine signed response occurrence already present in
the source word.  It cannot turn a complete moment occurrence into a
new response carrying \(T(D)/\Xi_u\).
\end{corollary}

\begin{proof}
Apply \Cref{prop:V61-no-beta-in-V14} to every occurrence and distribute
the fixed finite sum.  Each localized covariance remains the one
localized covariance selected by
\eqref{eq:V61-covariance-one}; every other factor is a complete
moment or an unchanged finite separator.  No scalar coefficient is
created by distribution.
\end{proof}

\section{A spectator-extension obstruction}
\label{sec:V61-spectator}

The absence of \(\beta_D\) is not only notational.  It is detected by
the same off-bank mass test which corrected the overstrong endpoint
factorization in V41--V42.

\begin{theorem}[A restricted moment bank cannot pay partner ownership]
\label{thm:V61-spectator-no-go}
Fix a nonzero restricted bank word \(Z_D\) made from complete moment
factors on \(D\) and fixed finite source intertwiners.  Consider a
spectator extension outside \(D\) for which
\begin{equation}
 Z_D\longmapsto Z_D\otimes I,
 \qquad
 \Xi_u\longmapsto\Xi_u+A,
 \qquad A>0,
\label{eq:V61-spectator-extension}
\end{equation}
while the weights on \(D\), and hence \(T(D)\), stay fixed.  Then
\begin{equation}
 \|Z_D\otimes I\|=\|Z_D\|,
 \qquad
 \beta_D(A)=\frac{T(D)}{\Xi_u+A}\longrightarrow0.
\label{eq:V61-spectator-limits}
\end{equation}
Consequently no support-uniform estimate
\begin{equation}
 \|Z_D\otimes I\|\le C\beta_D(A)
\label{eq:V61-false-moment-beta}
\end{equation}
can follow from complete-moment factorization alone.
\end{theorem}

\begin{proof}
The spatial tensor norm is multiplicative and \(\|I\|=1\), proving
the first equality in \eqref{eq:V61-spectator-limits}.  The second is
elementary because \(T(D)\) is fixed.  If
\eqref{eq:V61-false-moment-beta} held with \(C\) independent of the
spectator mass, letting \(A\to\infty\) would imply
\(\|Z_D\|=0\), contrary to the hypothesis.
\end{proof}

\begin{firewall}[Scope of the spectator test]
\label{fire:V61-spectator-scope}
The theorem does not say that every physical spectator contributes an
identity.  It says that the V14 algebra and the ordinary complete
moment contraction estimate permit such a tensor extension and hence
cannot prove a vanishing ownership factor.  A physical proof may
still obtain \(\beta_D\) from a signed covariance, a derivative, a
heat gap depending on the spectator representation, or a larger
source-level cancellation.  That additional occurrence must be
displayed and estimated explicitly.
\end{firewall}

\begin{assessment}[Correction to the V60 acquisition]
\label{ass:V61-V60-correction}
The conditional implication
\Cref{cor:V60-compensated-insertion} is valid, but its hypothesis
\eqref{eq:V60-compensated-word} is not supplied by
\eqref{eq:V14-moment-split}.  In particular, a complete partner-row
moment on \(D\) must not be relabelled as the scalar
\(T(D)/\Xi_u\).  The route now moves upstream to an actual normalized
response occurrence.
\end{assessment}

\section{Scalar response attachment has the wrong inequality
direction}
\label{sec:V61-attachment}

There is an existing scalar response on a comparable bank.  It is
important to distinguish its legal use in a tree ledger from the
operator smallness needed for ownership compensation.

\begin{lemma}[The reverse directed response dominates the ownership
fraction]
\label{lem:V61-reverse-dominates-beta}
Assume every root weight on \(D\) is active and hence
\(a_{i,e}\ge a_*>0\).  Define the reverse directed response
\begin{equation}
 d_{u\to i}(D):=\frac{H_{iu}(D)}{\Xi_u}.
\label{eq:V61-reverse-response}
\end{equation}
Then
\begin{equation}
 d_{u\to i}(D)\ge a_*\beta_D.
\label{eq:V61-response-lower-beta}
\end{equation}
\end{lemma}

\begin{proof}
Since \(a_{i,e}\ge a_*\) on \(D\),
\[
 H_{iu}(D)
 =\sum_{e\in D}a_{i,e}a_{u,e}
 \ge a_*\sum_{e\in D}a_{u,e}
 =a_*T(D).
\]
Divide by \(\Xi_u\).
\end{proof}

\begin{proposition}[A lower response is not an operator damping
factor]
\label{prop:V61-lower-not-damping}
The lower bound \eqref{eq:V61-response-lower-beta} permits legal
attachment of a scalar tree edge after a complete operator estimate,
as in \Cref{lem:V26-common-moment-attachment}.  It does not imply that
an operator occurrence \(Z_D\) satisfies
\(\|Z_D\|\le C\beta_D\), and therefore does not verify
\eqref{eq:V60-compensated-word}.
\end{proposition}

\begin{proof}
Scalar response attachment uses
\[
 \beta_D\le a_*^{-1}d_{u\to i}(D)
\]
to enlarge a proved right side by a legal marked-tree response.  The
needed damping estimate has the opposite form: it must show that the
left-side operator becomes small with \(\beta_D\).  Taking
\(Z_D=I\) and \(\beta_D\downarrow0\) already disproves any logical
implication from the displayed lower bound to such an operator norm.
The spectator family in \Cref{thm:V61-spectator-no-go} realizes the
same distinction at the bank level.
\end{proof}

\section{The exact response-carrier interface}
\label{sec:V61-interface}

\begin{definition}[Normalized partner response carrier]
\label{def:V61-NPRC}
For a root-owned comparable bank \(D\) assigned from \(i\) to \(u\),
an \(NPRC_{u|D}\) is an actual complete source factor
\(\mathcal Z_{u|D}\), present before positive envelopes, such that:
\begin{enumerate}[label=\textup{(\roman*)},leftmargin=4em]
\item it is coordinatewise equivariant on every coordinate meeting
      \(D\);
\item it occurs exactly once in each independent replica occurrence;
\item in the unpaid placement,
\begin{equation}
 \|\mathcal Z_{u|D}\|
 \le C_s\frac{T(D)}{\Xi_u}=C_s\beta_D;
\label{eq:V61-NPRC-unpaid}
\end{equation}
\item if it carries the unique payer, its corresponding complete paid
      norm has the prescribed single \(s_\alpha^{-1}\) debit and no
      second payer is introduced.
\end{enumerate}
\end{definition}

\begin{theorem}[An NPRC closes the partner-nonowned bank]
\label{thm:V61-NPRC-closes}
Suppose a restriction-resolved normalized source word can be ordered,
using only central commutations proved in V59, as
\begin{equation}
 \widetilde W_{\pi,\nu}
 =L_{\pi,\nu}\mathcal M_D
   \mathcal Z_{u|D}Q_{\pi,\nu},
\label{eq:V61-NPRC-word}
\end{equation}
where \(\mathcal M_D\) is a protected or positive gapped bank
envelope, \(\mathcal Z_{u|D}\) is an \(NPRC_{u|D}\), and the remaining
dressings have uniform norms and commute with \(\mathcal M_D\) on
the bank.  Then
\begin{equation}
 \boxed{
 \kappa_\otimes(
  \widetilde W_{\pi,\nu}^*\widetilde W_{\pi,\nu})
 \le C_s\theta_{iu}(D)^2.}
\label{eq:V61-NPRC-quadratic}
\end{equation}
For \(i=3\), the right side is at most \(C_s(x+z)^2\).
\end{theorem}

\begin{proof}
Equivariance of \(\mathcal Z_{u|D}\) and
\Cref{lem:V59-bank-commutation} imply
\([\mathcal M_D,\mathcal Z_{u|D}Q_{\pi,\nu}]=0\).  Its norm is at most
\(C_s\beta_D\) by \eqref{eq:V61-NPRC-unpaid} and the bound on
\(Q_{\pi,\nu}\).  Apply
\Cref{thm:V59-central-insertion}; the resulting scalar factor is
\(C_s\beta_D^2\).  The two estimates in
\Cref{thm:V60-ownership-compensation} prove
\eqref{eq:V61-NPRC-quadratic}.  Finally use
\(\theta_{3u}(D)\le x+z\).
\end{proof}

\begin{firewall}[NPRC is a target, not a renamed moment]
\label{fire:V61-NPRC-open}
\Cref{thm:V61-NPRC-closes} is a sufficient compiler.  V61 does not
assert that every partner-nonowned \(2+1+1\) word contains an NPRC.
In particular, \(\mathcal Z_{u|D}\) may not be chosen to be the
ordinary moment \(M_D\), and the scalar lower response
\eqref{eq:V61-response-lower-beta} may not be substituted for the
operator norm hypothesis \eqref{eq:V61-NPRC-unpaid}.
\end{firewall}

\section{V61 result ledger and the upstream V62 target}
\label{sec:V61-status}

\begin{theorem}[Third-series V61 result ledger]
\label{thm:V61-results}
Relative to V1--V60, V61 proves:
\begin{enumerate}[label=\textup{(V61.\arabic*)},leftmargin=6.3em]
\item the V14 moment/covariance restriction has coefficient one and
      creates no partner-ownership fraction;
\item finite subset resolution preserves genuine responses but does
      not turn moments into responses;
\item a spectator-extension no-go for extracting
      \(T(D)/\Xi_u\) from an ordinary restricted moment bank;
\item the exact lower comparison between reverse directed response and
      partner ownership, together with its inequality-direction
      firewall;
\item the normalized partner response-carrier interface; and
\item quadratic closure of every partner-nonowned bank for which that
      actual carrier is present in both replica occurrences.
\end{enumerate}
\end{theorem}

\begin{proof}
These are
\Cref{prop:V61-no-beta-in-V14,
cor:V61-resolution-not-response,
thm:V61-spectator-no-go,
lem:V61-reverse-dominates-beta,
prop:V61-lower-not-damping,
def:V61-NPRC,
thm:V61-NPRC-closes}.
\end{proof}

\begin{programme}[V62: locate or refute the NPRC in the exact row-3
word]
\label{prog:V61-V62}
\begin{enumerate}[leftmargin=3em]
\item Restore the complete row-\(u\) prefix, joining, singleton, and
      suffix occurrences suppressed in
      \(L_\pi\) and \(R_\pi\) of
      \eqref{eq:V49-complete-word}.
\item For a fixed partner-nonowned bank \(D\), apply the V14
      covariance telescope occurrence by occurrence and tabulate which
      terms contain a genuine row-\(u\) difference on \(D\).
\item Test those differences against the established complete
      one-row heat/covariance response bounds.  The required upper
      scale is exactly \(T(D)/\Xi_u\), not
      \(H_{iu}(D)/\Xi_u\) and not an attached lower response.
\item If no such occurrence exists, construct the complete
      spectator-invariant source term and move upstream to the unsplit
      composite polarization rather than adding a fictitious factor.
\item Keep the local payer and cofinal finite exceptional coordinates
      outside this audit.  The next compulsory companion audit remains
      V65.
\end{enumerate}
\end{programme}

\begin{firewall}[Claim boundary after third-series V61]
\label{fire:V61-claim-boundary}
V61 proves that the V14 moment split alone cannot supply the missing
partner-ownership factor and gives an exact sufficient response-carrier
interface.  It does not prove that the physical row-\(3\) word contains
that carrier, close all partner-nonowned, unbalanced, exceptional, or
local-payer branches, or establish full \(RCQ_3(\eta)\).  The other
quotient cuts, quartic source families, higher MTA, WUFC, CPM, cutoff
removal, reflection positivity, Osterwalder--Schrader reconstruction,
construction of the continuum Yang--Mills measure, and a uniform
positive continuum spectral gap remain open.  The full Yang--Mills
mass gap has not been proved.
\end{firewall}

\paragraph{Parent-version record.}
V61 was copied byte-for-byte from
\texttt{Renewal\_Yang\_Mills\_Mass\_Gap\_Research\_Notes\_V60.tex}
before this chapter was appended.  The V60 parent had \(30\,134\)
source lines, \(1\,084\,793\) bytes, and SHA--256
\[
 \texttt{ef70b09ad3d450e1a38b2a29899cfa97bd66ecd17ac0435d5c493a2766d23c0b}.
\]
The next compulsory companion audit is V65.

% =====================================================================
% V62 APPENDIX: RESIDUAL NORMALIZATION AND THE ROW-4 BANK
% =====================================================================

\chapter{V62: Residual Normalization Closes the Comparable Row-3--Row-4
Bank}
\label{ch:V62-row4-bank}

\section{The denominator ledger after the binary prefix}
\label{sec:V62-ledger}

V61 proves that a complete moment does not manufacture the missing
partner fraction.  The normalized source nevertheless contains four
global row denominators.  The complete binary prefix uses the row-
\(1\) and row-\(2\) denominators through its genuine response
\(H_{12}\).  The row-\(3\) and row-\(4\) denominators remain visible.

\begin{lemma}[Exact residual scalar]
\label{lem:V62-residual-scalar}
On the pure row-\(3\) face, where \(p=\theta_{12}>0\),
\begin{equation}
 \frac{H_{12}}
 {\Xi_1\Xi_2\Xi_3\Xi_4}
 =
 \frac{p}{\Xi_3\Xi_4}.
\label{eq:V62-residual-scalar}
\end{equation}
For every outside-local restriction-resolved word containing a central
unmarked bank envelope \(\mathcal M_D\), the proof of
\Cref{lem:V55-complete-prefix-response} and the commutations in V59
give an exact scalar--operator presentation
\begin{equation}
 \widetilde W_{\pi,\nu}
 =
 \frac{p}{\Xi_3\Xi_4}
 \overline L_{\pi,\nu}\mathcal M_DQ_{\pi,\nu},
\label{eq:V62-normalized-factorization}
\end{equation}
with
\begin{equation}
 \|\overline L_{\pi,\nu}\|
 +\|Q_{\pi,\nu}\|
 \le C_s,
 \qquad
 [\mathcal M_D,Q_{\pi,\nu}]=0.
\label{eq:V62-normalized-dressings}
\end{equation}
The scalar in \eqref{eq:V62-normalized-factorization} is not a new
response; it is the original normalization after the one actual
binary-prefix estimate.
\end{lemma}

\begin{proof}
The identity \eqref{eq:V62-residual-scalar} is the definition
\(p=H_{12}/(\Xi_1\Xi_2)\).  In the unpaid prefix placement,
\Cref{thm:V10-prefix-response} bounds the complete binary prefix by
\(C_sH_{12}\).  In the prefix-payer placement, its paid bound is
\(CH_{12}\).  At every other payer location, the binary prefix is
still bounded by \(C_sH_{12}\), while the unique complete paid factor
has its support-uniform norm.  All other complete moments are
contractions and all fixed local separators have uniform norm, as in
the proof of \Cref{lem:V55-complete-prefix-response}.

Since \(H_{12}>0\), divide the assembled prefix-side operator by
\(H_{12}\) and absorb its uniform normalized norm into
\(\overline L_{\pi,\nu}\).  Divide the complete word by
\(\Xi_1\Xi_2\Xi_3\Xi_4\), then use
\eqref{eq:V62-residual-scalar}.  The V59 centrality theorem moves only
the bank envelope through coordinatewise equivariant source factors,
giving the displayed order and commutator.  No row denominator is
discarded in this calculation.
\end{proof}

\begin{firewall}[No duplicated prefix debit]
\label{fire:V62-one-prefix}
The factor \(p\) in \eqref{eq:V62-normalized-factorization} is obtained
once from the complete signed binary prefix.  Its appearance in the
positive square is the square of that one already estimated source
operator, exactly as in \Cref{fire:V55-one-prefix-use}; it is not a
second use of the prefix bank.
\end{firewall}

\section{The unused row-4 denominator is the missing partner factor}
\label{sec:V62-row4}

Let \(D\) be a nonempty comparable bank assigned from row \(3\) to row
\(4\), with V60 notation
\begin{equation}
 T(D)=\sum_{e\in D}a_{4,e},
 \qquad
 \beta_D=\frac{T(D)}{\Xi_4}.
\label{eq:V62-row4-beta}
\end{equation}
Every coordinate in \(D\) is active at row \(4\), so
\begin{equation}
 T(D)\ge a_*.
\label{eq:V62-T-floor}
\end{equation}

\begin{lemma}[Residual denominator supplies row-4 ownership]
\label{lem:V62-row4-denominator}
For every nonempty row-\(3\)--row-\(4\) bank,
\begin{equation}
 \frac{p}{\Xi_3\Xi_4}
 \le
 \frac{\beta_D}{a_*^2}.
\label{eq:V62-row4-denominator}
\end{equation}
\end{lemma}

\begin{proof}
By \Cref{lem:V54-edge-unit}, \(p\le1\).  Activity gives
\(\Xi_3\ge a_*\), while
\[
 \frac1{\Xi_4}
 =\frac{\beta_D}{T(D)}
 \le\frac{\beta_D}{a_*}
\]
by \eqref{eq:V62-T-floor}.  Multiply the two bounds.
\end{proof}

\begin{theorem}[Comparable row-3--row-4 bank closure]
\label{thm:V62-row4-bank-closure}
Fix root ownership
\begin{equation}
 \sum_{e\in D}a_{3,e}\ge\sigma\Xi_3
\label{eq:V62-root-owned}
\end{equation}
and pointwise comparability
\begin{equation}
 c_0a_{3,e}\le a_{4,e}\le C_0a_{3,e}
 \qquad(e\in D).
\label{eq:V62-row4-comparable}
\end{equation}
No lower bound on \(T(D)/\Xi_4\) is required.  For either the protected
or positive gapped bank envelope \(\mathcal M_D\), every
outside-local restriction-resolved word of the form
\eqref{eq:V62-normalized-factorization} obeys
\begin{equation}
 \boxed{
 \kappa_\otimes(
  \widetilde W_{\pi,\nu}^*\widetilde W_{\pi,\nu})
 \le C_s\theta_{34}(D)^2
 \le C_s(x+z)^2.}
\label{eq:V62-row4-RCQ}
\end{equation}
The assertion is uniform in bank length, support, representations,
physical multiplicities, and selected coordinate, and has the
corresponding left-square form.
\end{theorem}

\begin{proof}
Apply \Cref{thm:V59-central-insertion} to
\eqref{eq:V62-normalized-factorization}.  Equations
\eqref{eq:V62-normalized-dressings} and
\eqref{eq:V62-row4-denominator} give
\begin{align*}
 \kappa_\otimes(
  \widetilde W_{\pi,\nu}^*\widetilde W_{\pi,\nu})
 &\le
 C_s\left(\frac{p}{\Xi_3\Xi_4}\right)^2
 \kappa_\otimes(\mathcal M_D)\\
 &\le
 C_sa_*^{-4}\beta_D^2
 \kappa_\otimes(\mathcal M_D).
\end{align*}
Use \Cref{thm:V60-ownership-compensation} with \(i=3,u=4\),
then \(\theta_{34}(D)\le z\le x+z\).  The left-square proof is the
adjoint-oriented part of \Cref{thm:V59-central-insertion}.
\end{proof}

\begin{corollary}[Both ownership cells of the row-4 bank are closed]
\label{cor:V62-all-row4-ownership}
For a root-owned comparable unmarked row-\(3\)--row-\(4\) bank:
\begin{enumerate}[label=\textup{(\roman*)},leftmargin=4em]
\item if \(T(D)/\Xi_4\) is bounded below, V57 and V59 close the
      physical bank insertion;
\item if it is arbitrarily small, \Cref{thm:V62-row4-bank-closure}
      closes it using the residual row-\(4\) denominator.
\end{enumerate}
Thus partner ownership is no longer a hypothesis anywhere in the
comparable unmarked row-\(3\)--row-\(4\) branch.
\end{corollary}

\begin{proof}
The two alternatives exhaust the value of the nonnegative fraction
\(T(D)/\Xi_4\).  Apply the stated theorems in the two cells.
\end{proof}

\begin{corollary}[Finite row-4 packet recombination]
\label{cor:V62-row4-recombination}
If every nonzero term in the V58 finite restriction resolution of a
fixed outside-local packet contains a root-owned comparable row-
\(3\)--row-\(4\) bank, then that packet satisfies \(RCQ_3(\eta)\).
\end{corollary}

\begin{proof}
Use \Cref{cor:V62-all-row4-ownership} term by term and then
\Cref{lem:V58-finite-recombination}.
\end{proof}

\section{Why the same normalization does not close partners 1 and 2}
\label{sec:V62-lateral}

\begin{proposition}[Exact denominator limitation]
\label{prop:V62-denominator-limitation}
After the complete binary-prefix response has been estimated as in
\Cref{lem:V62-residual-scalar}, the residual scalar contains
\(\Xi_3^{-1}\Xi_4^{-1}\) and contains neither
\(\Xi_1^{-1}\) nor \(\Xi_2^{-1}\).  Therefore this normalization
argument supplies the partner factor \(T(D)/\Xi_u\) directly only for
\(u=4\).  It supplies no such conclusion for \(u=1,2\).
\end{proposition}

\begin{proof}
This is the exact identity \eqref{eq:V62-residual-scalar}.  The two
prefix denominators have already formed the normalized prefix response
\(p=H_{12}/(\Xi_1\Xi_2)\).  Reusing either one would charge the same
source occurrence twice.
\end{proof}

\begin{proposition}[Strong pure-face edges do not restore a lateral
denominator]
\label{prop:V62-lateral-counterexample}
For \(A\ge1\), take two-coordinate weights
\begin{equation}
 (a_{1,a},a_{2,a},a_{3,a},a_{4,a})=(1,0,1,0),
 \qquad
 (a_{1,b},a_{2,b},a_{3,b},a_{4,b})=(A,A,0,1).
\label{eq:V62-lateral-example}
\end{equation}
Then
\begin{equation}
 p=\frac{A}{A+1},
 \qquad
 y=1+\frac{A}{A+1},
 \qquad
 x+z=\frac1{A+1},
\label{eq:V62-lateral-edges}
\end{equation}
but for the row-\(3\)--row-\(1\) bank \(D=\{a\}\),
\begin{equation}
 \beta_D=\frac1{A+1}\longrightarrow0,
 \qquad
 \frac{p}{\Xi_3\Xi_4}=\frac{A}{A+1}\longrightarrow1.
\label{eq:V62-lateral-separation}
\end{equation}
Hence no estimate of the residual scalar by \(C\beta_D\) follows for
the lateral partner from \(p,y\ge\eta\).
\end{proposition}

\begin{proof}
The total masses are
\(\Xi_1=A+1\), \(\Xi_2=A\), \(\Xi_3=1\), and \(\Xi_4=1\).
The nonzero overlaps are
\(H_{13}=1\), \(H_{12}=A^2\),
\(H_{14}=A\), and \(H_{24}=A\).  Division by the corresponding
products of total masses gives
\eqref{eq:V62-lateral-edges}.  On \(D\), \(T(D)=a_{1,a}=1\), proving
the first part of \eqref{eq:V62-lateral-separation}; the residual
scalar follows from \eqref{eq:V62-residual-scalar}.
\end{proof}

\begin{firewall}[Scope of the lateral example]
\label{fire:V62-lateral-scope}
As in \Cref{fire:V60-example-scope}, the example rules out an argument
from the nonnegative mass and edge ledger alone; it is not asserted to
be a complete representation-realizable source packet.  It shows that
the row-\(1\) or row-\(2\) partner fraction must come from an additional
physical response occurrence or a larger signed cancellation, not
from the already spent prefix denominator.
\end{firewall}

\section{V62 result ledger and V63 acquisition order}
\label{sec:V62-status}

\begin{theorem}[Third-series V62 result ledger]
\label{thm:V62-results}
Relative to V1--V61, V62 proves:
\begin{enumerate}[label=\textup{(V62.\arabic*)},leftmargin=6.3em]
\item the exact residual normalization
      \(p/(\Xi_3\Xi_4)\) after the one binary-prefix response;
\item a nonduplication ledger for that prefix debit;
\item conversion of the unused row-\(4\) denominator into the missing
      partner-ownership fraction;
\item quadratic physical closure of every root-owned comparable
      unmarked row-\(3\)--row-\(4\) protected or gapped bank, without
      partner ownership;
\item finite packet recombination for the all-row-4-bank cell;
\item the exact limitation of the residual denominator for partners
      \(1,2\); and
\item a sharp scalar separation showing that the strong edges
      \(p,y\) do not repair that lateral deficit.
\end{enumerate}
\end{theorem}

\begin{proof}
These are
\Cref{lem:V62-residual-scalar,
fire:V62-one-prefix,
lem:V62-row4-denominator,
thm:V62-row4-bank-closure,
cor:V62-all-row4-ownership,
cor:V62-row4-recombination,
prop:V62-denominator-limitation,
prop:V62-lateral-counterexample}.
\end{proof}

\begin{programme}[V63: lateral-partner response audit]
\label{prog:V62-V63}
\begin{enumerate}[leftmargin=3em]
\item Restrict the remaining partner set to \(u=1,2\); the comparable
      row-\(4\) bank is closed by V62.
\item Restore the two exact lateral occurrences which involve row
      \(u\): the binary prefix covariance and the composite joining
      difference.  Do not estimate either separately before their
      shared coordinate bank is exposed.
\item Apply the V14 covariance telescope to the lateral bank and test
      whether a normalized commutator or covariance remainder satisfies
      the NPRC norm \eqref{eq:V61-NPRC-unpaid}.
\item Preserve the unsplit composite difference of V56.  If neither
      lateral occurrence supplies the factor, test their exact signed
      combination before declaring the NPRC absent.
\item Keep row-\(4\), finite exceptional, unbalanced, and local-payer
      branches separate.  The next compulsory companion audit remains
      V65.
\end{enumerate}
\end{programme}

\begin{firewall}[Claim boundary after third-series V62]
\label{fire:V62-claim-boundary}
V62 closes the comparable unmarked row-\(3\)--row-\(4\) bank branch
for outside-local payers.  It does not close comparable banks with
partners \(1,2\), the unbalanced or finite exceptional branches, the
local-payer square, or full \(RCQ_3(\eta)\).  The other quotient cuts,
quartic source families, higher MTA, WUFC, CPM, cutoff removal,
reflection positivity, Osterwalder--Schrader reconstruction,
construction of the continuum Yang--Mills measure, and a uniform
positive continuum spectral gap remain open.  The full Yang--Mills
mass gap has not been proved.
\end{firewall}

\paragraph{Parent-version record.}
V62 was copied byte-for-byte from
\texttt{Renewal\_Yang\_Mills\_Mass\_Gap\_Research\_Notes\_V61.tex}
before this chapter was appended.  The V61 parent had \(30\,483\)
source lines, \(1\,098\,122\) bytes, and SHA--256
\[
 \texttt{53ecbd3348df95fde8ca89b7e80581132821842fb82ddcb4ffa1082161f16e66}.
\]
The next compulsory companion audit is V65.

% =====================================================================
% V63 APPENDIX: LATERAL MASS COMPARISON AND BACKBONE IMBALANCE
% =====================================================================

\chapter{V63: Lateral Mass Comparison and Strong-Backbone Imbalance
Localization}
\label{ch:V63-lateral-mass}

\section{A larger closure cell from the residual row-4 denominator}
\label{sec:V63-moderate-lateral}

V62 used the identity
\[
 \frac{H_{12}}
 {\Xi_1\Xi_2\Xi_3\Xi_4}
 =\frac{p}{\Xi_3\Xi_4}
\]
to close a row-\(3\)--row-\(4\) bank.  The same residual denominator
also closes a lateral partner \(u\in\{1,2\}\) whenever its total mass
is controlled by \(\Xi_4\).

\begin{lemma}[Transport of the unused denominator]
\label{lem:V63-denominator-transport}
Let \(D\) be a nonempty row-\(3\)--row-\(u\) bank with
\(u\in\{1,2\}\), and put
\begin{equation}
 \beta_D:=\frac{T(D)}{\Xi_u},
 \qquad
 T(D)=\sum_{e\in D}a_{u,e}.
\label{eq:V63-beta}
\end{equation}
If
\begin{equation}
 \Xi_u\le K\Xi_4
\label{eq:V63-moderate-mass}
\end{equation}
for a fixed \(K\ge1\), then
\begin{equation}
 \frac{p}{\Xi_3\Xi_4}
 \le
 \frac{K}{a_*^2}\beta_D.
\label{eq:V63-denominator-transport}
\end{equation}
\end{lemma}

\begin{proof}
The bank is nonempty and row \(u\) is active on every coordinate
assigned to it, so \(T(D)\ge a_*\).  Hence
\[
 \frac1{\Xi_4}
 \le\frac K{\Xi_u}
 =\frac{K\beta_D}{T(D)}
 \le\frac K{a_*}\beta_D.
\]
Also \(p\le1\) by \Cref{lem:V54-edge-unit} and
\(\Xi_3\ge a_*\).  Multiply the inequalities.
\end{proof}

\begin{theorem}[Moderate lateral-partner closure]
\label{thm:V63-moderate-lateral-closure}
Let \(u\in\{1,2\}\), and let \(D\) be a root-owned comparable
unmarked row-\(3\)--row-\(u\) bank satisfying
\eqref{eq:V63-moderate-mass}.  For either its protected or positive
gapped central bank envelope, every outside-local
restriction-resolved normalized word satisfies
\begin{equation}
 \boxed{
 \kappa_\otimes(
  \widetilde W_{\pi,\nu}^*\widetilde W_{\pi,\nu})
 \le C_{s,K}\theta_{3u}(D)^2
 \le C_{s,K}(x+z)^2.}
\label{eq:V63-moderate-lateral-RCQ}
\end{equation}
The estimate is uniform in support, bank length, representations,
physical multiplicities, and selected coordinate.
\end{theorem}

\begin{proof}
Use the exact normalized factorization
\eqref{eq:V62-normalized-factorization}.  By
\Cref{lem:V63-denominator-transport}, its scalar prefactor is at most
\(K a_*^{-2}\beta_D\).  Apply the central insertion theorem
\Cref{thm:V59-central-insertion}, followed by partner-ownership
compensation in \Cref{thm:V60-ownership-compensation}.  The last
inequality is \(\theta_{3u}(D)\le x+z\).
\end{proof}

\begin{corollary}[Current comparable-bank residual]
\label{cor:V63-comparable-residual}
Fix \(K\ge1\).  Among outside-local root-owned comparable unmarked
row-\(3\) banks:
\begin{enumerate}[label=\textup{(\roman*)},leftmargin=4em]
\item partner \(u=4\) is closed for all mass ratios by V62;
\item partner \(u=1,2\) is closed when \(\Xi_u\le K\Xi_4\); and
\item the only remaining comparable-bank cells have
      \(u\in\{1,2\}\) and \(\Xi_u>K\Xi_4\).
\end{enumerate}
\end{corollary}

\begin{proof}
Use \Cref{cor:V62-all-row4-ownership} for item \textup{(i)} and
\Cref{thm:V63-moderate-lateral-closure} for item \textup{(ii)}.  The
three partner labels exhaust \(\{1,2,4\}\).
\end{proof}

\begin{firewall}[The constant \(K\) is fixed]
\label{fire:V63-fixed-K}
The constant in \eqref{eq:V63-moderate-lateral-RCQ} depends on \(K\).
One may not choose \(K\) as a function of the vanishing cut currency
\(x+z\) and then call the resulting constant uniform.  The
large-ratio complement is treated structurally below.
\end{firewall}

\section{The strong pure-face backbone}
\label{sec:V63-backbone}

On \(\mathfrak R_{3,\eta}\),
\begin{equation}
 \theta_{12}=p\ge\eta,
 \qquad
 \theta_{14}+\theta_{24}=y\ge\eta.
\label{eq:V63-strong-backbone}
\end{equation}
Choose \(v\in\{1,2\}\) with
\begin{equation}
 \theta_{v4}\ge\eta/2.
\label{eq:V63-strong-v4}
\end{equation}
The edges \(12\) and \(v4\) form a path connecting both lateral rows
to row \(4\).

\begin{lemma}[A lateral mass ratio appears on a strong edge]
\label{lem:V63-path-ratio}
Let \(u\in\{1,2\}\) and assume
\begin{equation}
 \frac{\Xi_u}{\Xi_4}>K.
\label{eq:V63-large-lateral}
\end{equation}
Along the path from \(u\) to \(4\) using the edges in
\eqref{eq:V63-strong-backbone}--\eqref{eq:V63-strong-v4}, there is an
oriented edge \(a\to b\) such that
\begin{equation}
 \theta_{ab}\ge\eta/2,
 \qquad
 \frac{\Xi_a}{\Xi_b}>K^{1/2}.
\label{eq:V63-strong-imbalanced-edge}
\end{equation}
If \(u=v\), the exponent \(1/2\) may be replaced by \(1\).
\end{lemma}

\begin{proof}
If \(u=v\), the path is the single edge \(u4\), which has weight at
least \(\eta/2\), and
\eqref{eq:V63-large-lateral} gives the mass ratio.

If \(u\ne v\), the path is \(u-v-4\).  Both edge weights are at
least \(\eta/2\), since \(\theta_{uv}=\theta_{12}\ge\eta\).  The
product of the two oriented mass ratios is
\[
 \frac{\Xi_u}{\Xi_v}
 \frac{\Xi_v}{\Xi_4}
 =\frac{\Xi_u}{\Xi_4}>K.
\]
At least one factor exceeds \(K^{1/2}\).  Orient that edge from its
larger-mass endpoint to its smaller-mass endpoint.
\end{proof}

\section{A strong imbalanced edge owns a pointwise-imbalanced bank}
\label{sec:V63-localization}

The next lemma is purely nonnegative and retains the coordinate set
which carries the imbalance.  It is the relative quantity needed
before a representation-theoretic gapped/balanced split.

\begin{lemma}[Quantitative strong-edge imbalance localization]
\label{lem:V63-imbalance-localization}
Let nonnegative arrays \((A_e)_e,(B_e)_e\) have totals
\(A=\sum_eA_e\), \(B=\sum_eB_e>0\), and suppose
\begin{equation}
 \frac{\sum_eA_eB_e}{AB}\ge\delta>0.
\label{eq:V63-strong-overlap}
\end{equation}
Fix \(C_0>0\) and define
\begin{equation}
 E_{\rm bal}:=\{e:A_e\le C_0B_e\},
 \qquad
 E_{\rm im}:=\{e:A_e>C_0B_e\}.
\label{eq:V63-pair-split}
\end{equation}
If
\begin{equation}
 \frac AB\ge\frac{2C_0}{\delta},
\label{eq:V63-large-ratio-threshold}
\end{equation}
then
\begin{align}
 \sum_{e\in E_{\rm im}}A_eB_e
 &\ge\frac\delta2AB,
\label{eq:V63-imbalanced-overlap}\\
 \sum_{e\in E_{\rm im}}A_e
 &\ge\frac\delta2A.
\label{eq:V63-imbalanced-mass}
\end{align}
Thus a fixed fraction of the heavy endpoint mass lies on coordinates
where it exceeds the light endpoint pointwise by the prescribed
factor \(C_0\).
\end{lemma}

\begin{proof}
On \(E_{\rm bal}\),
\[
 \sum_{e\in E_{\rm bal}}A_eB_e
 \le C_0\sum_eB_e^2
 \le C_0B^2
 \le\frac\delta2AB
\]
by \eqref{eq:V63-large-ratio-threshold}.  Subtract this from
\eqref{eq:V63-strong-overlap} to obtain
\eqref{eq:V63-imbalanced-overlap}.  Since \(B_e\le B\),
\[
 B\sum_{e\in E_{\rm im}}A_e
 \ge
 \sum_{e\in E_{\rm im}}A_eB_e,
\]
which gives \eqref{eq:V63-imbalanced-mass}.
\end{proof}

\begin{corollary}[Localization on the Yang--Mills backbone]
\label{cor:V63-YM-localization}
Choose \(K\) so that
\begin{equation}
 K^{1/2}\ge\frac{4C_0}{\eta}.
\label{eq:V63-K-choice}
\end{equation}
In every remaining cell of
\Cref{cor:V63-comparable-residual}, there is an actual strong edge
\(a b\in\{12,14,24\}\), oriented from its heavier endpoint, and a
complete coordinate bank \(E_{\rm im}\) such that
\begin{align}
 \theta_{ab}&\ge\eta/2,
\label{eq:V63-YM-edge}\\
 a_{a,e}&>C_0a_{b,e}
 \quad(e\in E_{\rm im}),
\label{eq:V63-YM-pointwise}\\
 \sum_{e\in E_{\rm im}}a_{a,e}
 &\ge\frac\eta4\Xi_a.
\label{eq:V63-YM-heavy-mass}
\end{align}
\end{corollary}

\begin{proof}
Use \Cref{lem:V63-path-ratio} with
\(\delta=\eta/2\).  Condition \eqref{eq:V63-K-choice} is exactly
\(K^{1/2}\ge2C_0/\delta\).  Apply
\Cref{lem:V63-imbalance-localization} with
\(A_e=a_{a,e}\), \(B_e=a_{b,e}\).
\end{proof}

\begin{assessment}[What the large-ratio branch has become]
\label{ass:V63-large-ratio-result}
The remaining lateral branch is no longer merely
\(\Xi_u\gg\Xi_4\).  It contains an actual strong prefix-or-quotient
edge and a cofinal bank on which the heavy endpoint dominates the
light endpoint pointwise.  At each coordinate of that bank, either
the heavy row is unbalanced against all other active rows, giving the
V32 retained-output gap, or some third row balances it and supplies a
new physical partner.  This last finite representation-theoretic split
is the V64 target.  The strong edge and its coordinate bank must remain
attached during that split.
\end{assessment}

\section{V63 result ledger and V64 acquisition order}
\label{sec:V63-status}

\begin{theorem}[Third-series V63 result ledger]
\label{thm:V63-results}
Relative to V1--V62, V63 proves:
\begin{enumerate}[label=\textup{(V63.\arabic*)},leftmargin=6.3em]
\item transport of the residual row-\(4\) denominator to every
      lateral partner with \(\Xi_u\le K\Xi_4\);
\item quadratic closure of all such moderate-mass comparable banks;
\item reduction of the comparable-bank residual to large lateral mass
      ratios;
\item a strong-backbone path lemma locating a large mass ratio on an
      actual edge of weight at least \(\eta/2\);
\item quantitative localization of strong overlap to a cofinal
      pointwise-imbalanced bank; and
\item the corresponding Yang--Mills backbone bank with its exact
      strong edge retained.
\end{enumerate}
\end{theorem}

\begin{proof}
These are
\Cref{lem:V63-denominator-transport,
thm:V63-moderate-lateral-closure,
cor:V63-comparable-residual,
lem:V63-path-ratio,
lem:V63-imbalance-localization,
cor:V63-YM-localization}.
\end{proof}

\begin{programme}[V64: full-row resolution of the imbalanced backbone
bank]
\label{prog:V63-V64}
\begin{enumerate}[leftmargin=3em]
\item On \(E_{\rm im}\), split coordinates according to whether the
      heavy row is V32-unbalanced against all other active rows.
\item Tensor the sharp Casimir-dependent heat loss on that subbank;
      retain dependence on the total heavy-row mass rather than only
      the fixed spectral floor.
\item On the complementary balanced subbank, assign the maximal
      partner supplied by \Cref{lem:V32-balanced-partner}.  Keep the
      original strong edge \(ab\) in the same word.
\item Prove that either a root-owned comparable bank is reached or a
      strictly larger total-mass row is reached.  Use acyclicity only
      after recording the coordinate bank and strong edge.
\item Determine whether the resulting gapped or owned-bank currency,
      together with the original small row-\(3\) edge, supplies the
      lateral NPRC scale.  Do not replace the retained edge by its
      lower bound before the operator estimate.
\item The next compulsory companion audit remains V65.
\end{enumerate}
\end{programme}

\begin{firewall}[Claim boundary after third-series V63]
\label{fire:V63-claim-boundary}
V63 closes comparable lateral banks whose partner mass is controlled
by \(\Xi_4\) and localizes the complementary large-ratio branch to an
actual pointwise-imbalanced strong-edge bank.  It does not yet compile
that bank through the full V32 balanced/unbalanced split, close all
lateral partners, the finite exceptional or local-payer branches, or
establish full \(RCQ_3(\eta)\).  The other quotient cuts, quartic
source families, higher MTA, WUFC, CPM, cutoff removal, reflection
positivity, Osterwalder--Schrader reconstruction, construction of the
continuum Yang--Mills measure, and a uniform positive continuum
spectral gap remain open.  The full Yang--Mills mass gap has not been
proved.
\end{firewall}

\paragraph{Parent-version record.}
V63 was copied byte-for-byte from
\texttt{Renewal\_Yang\_Mills\_Mass\_Gap\_Research\_Notes\_V62.tex}
before this chapter was appended.  The V62 parent had \(30\,835\)
source lines, \(1\,110\,792\) bytes, and SHA--256
\[
 \texttt{3dfc24dbcb3a7e241deabc610884eda3de05fbfa2d1e003f0ed11f62303bb660}.
\]
The next compulsory companion audit is V65.

% =====================================================================
% V64 APPENDIX: FULL-ROW IMBALANCE RESOLUTION
% =====================================================================

\chapter{V64: Full-Row Resolution of a Strong Imbalanced Backbone
Bank}
\label{ch:V64-full-row-resolution}

\section{Pair imbalance versus full local balance}
\label{sec:V64-local-split}

V63 produces a complete coordinate bank \(E\), an oriented strong
edge \(a\to b\), and constants \(\delta,\sigma_0>0\) such that
\begin{equation}
 \theta_{ab}\ge\delta,
 \qquad
 \sum_{e\in E}a_{a,e}\ge\sigma_0\Xi_a,
 \qquad
 a_{a,e}>C_0a_{b,e}\qquad(e\in E).
\label{eq:V64-V63-data}
\end{equation}
Pairwise dominance over row \(b\) does not by itself make row \(a\)
V32-unbalanced: rows outside the edge may balance it.  The following
lemma makes that alternative exhaustive.

\begin{lemma}[Every nongapped coordinate supplies another physical
partner]
\label{lem:V64-nongapped-partner}
There is \(c_G^{\rm p}>0\), depending only on the fixed group, with
the following property.  At a coordinate active in row \(a\), exactly
one of the following alternatives may be selected:
\begin{enumerate}[label=\textup{(\roman*)},leftmargin=4em]
\item row \(a\) is unbalanced in the sense of
      \eqref{eq:V32-unbalanced};
\item some row \(c\ne a\) satisfies
\begin{equation}
 a_{c,e}\ge c_G^{\rm p}a_{a,e}.
\label{eq:V64-new-partner}
\end{equation}
\end{enumerate}
The partner in \textup{(ii)} may be chosen from the at most three
other active rows.
\end{lemma}

\begin{proof}
If the entire local assignment is balanced, apply
\Cref{lem:V32-balanced-partner}.  It gives
\(a_{a,e}\le C_Ga_{c,e}\), so
\eqref{eq:V64-new-partner} holds with \(C_G^{-1}\).

Suppose instead that the assignment is unbalanced, but not at row
\(a\).  Then some \(c\ne a\) obeys
\[
 \sum_{j\ne c}\|\lambda_{j,e}\|
 \le\frac12\|\lambda_{c,e}\|.
\]
In particular
\(\|\lambda_{c,e}\|\ge2\|\lambda_{a,e}\|\).  The two-sided
Casimir comparison \eqref{eq:V32-Casimir-comparison} gives
\[
 a_{c,e}
 \ge c_G(1+\|\lambda_{c,e}\|^2)
 \ge c_G(1+\|\lambda_{a,e}\|^2)
 \ge\frac{c_G}{C_G}a_{a,e}.
\]
Take \(c_G^{\rm p}\) to be the smaller of the two fixed constants.
The remaining alternative is precisely that row \(a\) is
unbalanced.
\end{proof}

\begin{definition}[Heavy-row gapped and partner banks]
\label{def:V64-bank-split}
For the bank \(E\) in \eqref{eq:V64-V63-data}, let
\(E_{\rm gap}\) be the coordinates at which row \(a\) is unbalanced.
On its complement choose a partner by
\Cref{lem:V64-nongapped-partner} and write
\begin{equation}
 E\setminus E_{\rm gap}
 =\dot\bigcup_{c\ne a}E_c.
\label{eq:V64-partner-partition}
\end{equation}
\end{definition}

\section{The gapped bank remembers total heavy-row mass}
\label{sec:V64-sharp-gap}

V32 used the fixed spectral floor to obtain a factor exponential in
the number of bank coordinates.  On a heavy-row-unbalanced bank, the
stronger Casimir estimate in the same theorem retains the sum of the
actual row weights.

\begin{lemma}[Casimir-mass heat tensorization]
\label{lem:V64-Casimir-mass-gap}
Let \(F\) be a complete bank on which row \(a\) is unbalanced at every
coordinate, and put
\begin{equation}
 S_a(F):=\sum_{e\in F}a_{a,e}.
\label{eq:V64-heavy-bank-mass}
\end{equation}
For fixed \(s=s_\alpha>0\), the normalized positive gapped output
envelope \(\mathcal G_a(F)\) satisfies
\begin{equation}
 \boxed{
 \|\mathcal G_a(F)\|
 \le e^{-c_{G}'sS_a(F)}.}
\label{eq:V64-mass-gap}
\end{equation}
The same bounds hold for product-state capacity.
\end{lemma}

\begin{proof}
At coordinate \(e\),
\Cref{lem:V32-unbalanced-output} gives for every output constituent
\(\nu_e\)
\[
 1+C_2(\nu_e)\ge c_Ga_{a,e}.
\]
If \(a_{a,e}\ge2/c_G\), this implies
\[
 C_2(\nu_e)\ge\frac{c_G}{2}a_{a,e}.
\]
If \(a_{a,e}<2/c_G\), the same V32 lemma says that \(\nu_e\ne0\).
The fixed output spectral floor therefore gives
\[
 C_2(\nu_e)
 \ge c_{\mathrm{gap},G}
 \ge \frac{c_Gc_{\mathrm{gap},G}}2a_{a,e}.
\]
Thus in both cases
\[
 C_2(\nu_e)\ge c_G'a_{a,e},
 \qquad
 c_G':=\frac{c_G}{2}\min\{1,c_{\mathrm{gap},G}\}>0.
\]
Orthogonality of output blocks gives the local norm bound
\(e^{-c_G'sa_{a,e}}\).  Tensoring over \(F\) proves
\eqref{eq:V64-mass-gap} without a support-dependent prefactor.
Capacity is at most operator norm for a positive operator by
\Cref{lem:V41-capacity-norm}.
\end{proof}

\begin{remark}[Normalization convention in the exponential]
\label{rem:V64-gap-convention}
If the standing heat multiplier is written as
\(e^{-s(1+C_2)}\), the estimate is only stronger.  For the convention
\(e^{-sC_2}\), the nontrivial-output spectral floor is essential on
the finite small-input range.  Only the positive rate \(c_G's\) in
\eqref{eq:V64-mass-gap} is used below.
\end{remark}

\begin{lemma}[Cofinal heat decay pays every finite mass-ratio power]
\label{lem:V64-gap-pays-ratio}
Suppose
\begin{equation}
 S_a(F)\ge\sigma\Xi_a
\label{eq:V64-gap-cofinal}
\end{equation}
for some \(\sigma>0\).  For every fixed integer \(m\ge1\) and every
active row \(b\),
\begin{equation}
 \boxed{
 \kappa_\otimes(\mathcal G_a(F))
 \le C_{G,s,\sigma,m}
 \left(\frac{\Xi_b}{\Xi_a}\right)^m.}
\label{eq:V64-gap-ratio}
\end{equation}
No assumption on the order of \(\Xi_a,\Xi_b\) is required.
\end{lemma}

\begin{proof}
By \Cref{lem:V64-Casimir-mass-gap}, the left side is at most
\(e^{-c_G's\sigma\Xi_a}\).  If
\(\Xi_b\ge\Xi_a\), the right-hand ratio is at least one and the
claim follows from the contraction bound after increasing the
constant.  If \(\Xi_b<\Xi_a\), activity gives \(\Xi_b\ge a_*\), so
\[
 \left(\frac{\Xi_b}{\Xi_a}\right)^m
 \ge\frac{a_*^m}{\Xi_a^m}.
\]
The function \(x^me^{-c_G's\sigma x}\) is bounded on
\([a_*,\infty)\).  Rearrangement proves
\eqref{eq:V64-gap-ratio}.
\end{proof}

\section{Owned bank, mass escalation, or cofinal heat gap}
\label{sec:V64-trichotomy}

\begin{theorem}[Full-row resolution of the V63 bank]
\label{thm:V64-full-row-trichotomy}
Let \(E\) satisfy \eqref{eq:V64-V63-data}.  Fix \(R>1\).  Then at
least one of the following occurs:
\begin{enumerate}[label=\textup{(\alph*)},leftmargin=4em]
\item a complete subbank \(F\subseteq E_{\rm gap}\) satisfies
\begin{equation}
 S_a(F)\ge\frac{\sigma_0}{2}\Xi_a
\label{eq:V64-cofinal-gap}
\end{equation}
and hence the mass-ratio estimates
\eqref{eq:V64-gap-ratio};
\item for some \(c\ne a\), a complete subbank \(D\subseteq E_c\)
is root-owned and pointwise comparable in both directions, and owns a
fixed fraction of \(\Xi_c\);
\item for some \(c\ne a\),
\begin{equation}
 \Xi_c\ge R\Xi_a.
\label{eq:V64-new-escalation}
\end{equation}
\end{enumerate}
All fractions and comparability constants depend only on
\(G,\sigma_0,R\), not on support or representations.
\end{theorem}

\begin{proof}
If
\(S_a(E_{\rm gap})\ge\sigma_0\Xi_a/2\), take
\(F=E_{\rm gap}\) and use
\Cref{lem:V64-gap-pays-ratio}.

Otherwise the partner part of
\eqref{eq:V64-partner-partition} carries more than
\(\sigma_0\Xi_a/2\).  One of its at most three banks, say \(E_c\),
therefore satisfies
\[
 S_a(E_c)\ge\frac{\sigma_0}{6}\Xi_a,
 \qquad
 a_{c,e}\ge c_G^{\rm p}a_{a,e}\qquad(e\in E_c).
\]
Apply the proof of
\Cref{thm:V57-bank-or-escalation} to this bank, with root \(a\),
lower comparability constant \(c_G^{\rm p}\), and parameters chosen
as in \Cref{cor:V57-arbitrary-escalation}.  The proof uses only the
displayed lower pointwise comparison and root ownership: split at an
upper comparison threshold, then test partner ownership.  It yields
either a two-sided comparable bank owning fixed fractions of both row
masses, which is item \textup{(b)}, or
\eqref{eq:V64-new-escalation}, which is item \textup{(c)}.
\end{proof}

\begin{corollary}[Maximal heavy rows have only two outcomes]
\label{cor:V64-maximal-row}
If row \(a\) has maximal total mass among the four rows, then item
\textup{(c)} of \Cref{thm:V64-full-row-trichotomy} is impossible.
The cofinal V63 bank therefore yields either a cofinal
Casimir-mass-gapped subbank or an owned comparable subbank.
\end{corollary}

\begin{proof}
For \(R>1\), \(\Xi_c\ge R\Xi_a\) contradicts maximality.
\end{proof}

\begin{corollary}[Escalation still terminates with its bank retained]
\label{cor:V64-finite-chain}
Starting from the oriented strong-edge bank in
\Cref{cor:V63-YM-localization}, repeat item \textup{(c)} only when it
occurs.  After at most three row changes, the chain terminates in a
cofinal gapped bank or an owned comparable bank.  At every step the
coordinate bank and the incoming strong edge remain recorded.
\end{corollary}

\begin{proof}
Every escalation multiplies the total row mass by the same fixed
factor \(R>1\).  Apply \Cref{lem:V57-no-cycle}.  The other two
alternatives terminate the construction.  No bank or edge is removed
by this scalar graph argument, so their labels may be retained.
\end{proof}

\begin{assessment}[Analytic and incidence work now separate]
\label{ass:V64-separation}
The terminal local analytic alternatives are closed at their natural
scale: an owned comparable bank has the V57 quadratic currency and
V59 central insertion, while a cofinal heavy-row-unbalanced bank has
exponential decay strong enough to pay every finite mass-ratio power.
What is not yet proved is that these terminal gains and the original
small row-\(3\) edge occur in one exact physical square with the
required orientation.  The V65 audit must therefore preserve the
entire labelled chain and test its source-word ancestry before any
product of scalar estimates is taken.
\end{assessment}

\section{V64 result ledger and V65 mandate}
\label{sec:V64-status}

\begin{theorem}[Third-series V64 result ledger]
\label{thm:V64-results}
Relative to V1--V63, V64 proves:
\begin{enumerate}[label=\textup{(V64.\arabic*)},leftmargin=6.3em]
\item every coordinate not unbalanced at the selected heavy row has
      another physical partner with a uniform Casimir comparison;
\item the exact gapped heat bound exponential in total heavy-row
      Casimir mass;
\item conversion of cofinal heat decay into every fixed inverse
      total-mass-ratio power;
\item a full-row trichotomy: cofinal heat gap, owned comparable bank,
      or strict total-mass escalation;
\item removal of escalation at a maximal row; and
\item finite termination while retaining the coordinate banks and
      incoming strong edges.
\end{enumerate}
\end{theorem}

\begin{proof}
These are
\Cref{lem:V64-nongapped-partner,
lem:V64-Casimir-mass-gap,
lem:V64-gap-pays-ratio,
thm:V64-full-row-trichotomy,
cor:V64-maximal-row,
cor:V64-finite-chain}.
\end{proof}

\begin{programme}[V65: compulsory audit and labelled-chain insertion]
\label{prog:V64-V65}
\begin{enumerate}[leftmargin=3em]
\item Perform the compulsory read-only audit of the newest active SM
      and NS notes and record exact file metadata.
\item Import only a proof-order lesson which is type-correct for the
      labelled strong-edge chain.
\item Write the complete chain data
      \((E_j,a_j\to b_j,D_j)\) before physical compression and identify
      which factors are actual occurrences of the fixed
      \(2+1+1\) word.
\item For a gapped terminal bank, commute its central positive
      envelope to the safe side of the square and use
      \eqref{eq:V64-gap-ratio} without dropping the original row-
      \(3\) bank.
\item For an owned terminal bank, test whether V59 insertion and the
      intervening strong-edge occurrences give the product of
      ownership fractions required to return to \((x+z)^2\).
\item If the labelled chain contains an edge which is only a scalar
      lower response and not an operator occurrence, stop at that edge
      and return upstream to its exact covariance identity.
\end{enumerate}
\end{programme}

\begin{firewall}[Claim boundary after third-series V64]
\label{fire:V64-claim-boundary}
V64 closes the local representation-theoretic classification and the
cofinal gapped mass-ratio estimate on the imbalanced strong-edge bank.
It does not yet insert an entire labelled escalation chain into the
physical source square, close all lateral, finite exceptional, or
local-payer branches, or establish full \(RCQ_3(\eta)\).  The other
quotient cuts, quartic source families, higher MTA, WUFC, CPM, cutoff
removal, reflection positivity, Osterwalder--Schrader reconstruction,
construction of the continuum Yang--Mills measure, and a uniform
positive continuum spectral gap remain open.  The full Yang--Mills
mass gap has not been proved.
\end{firewall}

\paragraph{Parent-version record.}
V64 was copied byte-for-byte from
\texttt{Renewal\_Yang\_Mills\_Mass\_Gap\_Research\_Notes\_V63.tex}
before this chapter was appended.  The V63 parent had \(31\,185\)
source lines, \(1\,122\,550\) bytes, and SHA--256
\[
 \texttt{644b3a6d7634ab7dfcd4a7f6a53763e74939e2ffd6de17c67dd9cdb811a2e525}.
\]
The compulsory companion audit is due in V65.

% =====================================================================
% V65 APPENDIX: AUDIT AND LABELLED-CHAIN INSERTION
% =====================================================================

\chapter{V65: Companion Audit and the Exact Labelled-Chain
Insertion Test}
\label{ch:V65-chain-insertion}

\section{Compulsory read-only companion audit}
\label{sec:V65-audit}

The companion audit was performed before the direction of this
chapter was fixed.  The active files at the completed read were:
\begin{enumerate}[leftmargin=3em]
\item
\texttt{Renewal\_SM\_Einstein\_Continuum\_Research\_Notes\_V46.tex},
with \(16\,867\) logical source lines and \(620\,928\) bytes,
SHA--256
\[
 \texttt{4befE83cfbf34c813a56a88cc77e4b69523e269077da7a7b2b5531594d25dcc5}.
\]
Case is immaterial in this hexadecimal digest.  Its newest exact
calculation differentiates a two-point Witten representation before
estimating it.  One derivative is a genuine two-resolvent insertion;
the other differentiates the normalized carrier and contains an
additional field insertion.  Pair decay alone therefore does not
manufacture a three-point carrier.
\item
\texttt{Renewal\_NS\_Stability\_Research\_Notes\_V31.tex},
with \(23\,052\) logical source lines and \(887\,537\) bytes,
SHA--256
\[
 \texttt{f1121b62238ad5491bb7cf03635ea9934942edf573ed8faaa9ee79e1d38a6696}.
\]
This file is unchanged from the V60 audit.  Its terminal ledger keeps
an exact stopped quantity until signed incidence has been assembled,
and forbids charging the persistent branch a second time.
\end{enumerate}
The line counts above are the logical-line counts returned by the
same read used for the hashes.

\begin{assessment}[Type-correct transfer from the companion audit]
\label{ass:V65-audit-transfer}
Only proof order is transferred.  A chain selected from nonnegative
Yang--Mills row weights is not thereby a product of source
operators, just as a pair estimate in the SM note is not an
unwritten third-order carrier.  Likewise the original prefix or
bank currency may be charged only once.  V65 will first expose all
claimed bank factors in one exact source word, and only then combine
their estimates.
\end{assessment}

\section{Simultaneous restriction of finitely many banks}
\label{sec:V65-simultaneous-resolution}

We first separate an algebraic fact from an occurrence hypothesis.
Let \(D_1,\dots,D_k\) be complete coordinate banks.  Their Boolean
atoms are the nonempty sets
\begin{equation}
 A_\varepsilon
 :=
 \bigcap_{j=1}^k D_j^{\varepsilon_j},
 \qquad
 D_j^1=D_j,\quad D_j^0=D_j^c,
 \qquad
 \varepsilon\in\{0,1\}^k.
\label{eq:V65-Boolean-atoms}
\end{equation}
There are at most \(2^k\) such atoms.

\begin{lemma}[Exact simultaneous bank restriction]
\label{lem:V65-simultaneous-restriction}
Consider a fixed complete source word containing \(J\) covariance
or difference occurrences and any fixed number of complete moment
occurrences.  Restriction to all banks \(D_1,\dots,D_k\) has an exact
expansion into at most
\begin{equation}
 (2^k)^J
\label{eq:V65-resolution-count}
\end{equation}
terms.  Every moment tensorizes over the atoms
\eqref{eq:V65-Boolean-atoms}; in each resulting term every covariance
is localized to exactly one atom.  All coefficients in the
telescopes are one.  In particular, the expansion preserves actual
response occurrences but creates no row-mass fraction.
\end{lemma}

\begin{proof}
Order the nonempty atoms as \(A_1,\dots,A_N\), where \(N\le2^k\).
For a complete moment,
\[
 M_{\cup_rA_r}(\rho)=\bigotimes_{r=1}^NM_{A_r}(\rho).
\]
For a complete covariance
\(\Delta^{\rho,\sigma}=M(\rho)-M(\sigma)\), the elementary tensor
telescope is
\begin{align}
 \Delta_{\cup_rA_r}^{\rho,\sigma}
 &=
 \sum_{q=1}^N
 \left(\bigotimes_{r<q}M_{A_r}(\rho)\right)
 \otimes\Delta_{A_q}^{\rho,\sigma}
 \otimes
 \left(\bigotimes_{r>q}M_{A_r}(\sigma)\right).
\label{eq:V65-atom-telescope}
\end{align}
Thus a covariance occurrence contributes at most \(N\) choices and
a moment contributes no sum.  Distributing the \(J\) covariance
sums gives at most \(N^J\le(2^k)^J\) terms.  Formula
\eqref{eq:V65-atom-telescope} has coefficient one, so no quantity
depending on \(a_{i,e}\), \(\Xi_i\), or the selected bank is created.
\end{proof}

\begin{definition}[Source-occurring and jointly source-occurring
banks]
\label{def:V65-source-occurring}
A protected or gapped bank factor is \emph{source-occurring} in a
resolved word if:
\begin{enumerate}[label=\textup{(\roman*)},leftmargin=4em]
\item its coordinate bank is exposed by
      \Cref{lem:V65-simultaneous-restriction} inside an actual
      complete source occurrence;
\item the local \(G\)-module on that occurrence is the module to
      which the protected/gapped isotypic resolution is applied; and
\item the resulting positive central factor is retained before any
      norm or product-state capacity is taken.
\end{enumerate}
A finite family is \emph{jointly source-occurring} if one common
exact expansion of the original word contains all its factors in
each term under consideration.  This is an incidence property of
the source word, not a property of the scalar row-weight ledger.
\end{definition}

\begin{corollary}[Finite joint resolution has no support cost]
\label{cor:V65-joint-resolution}
For a V64 chain of length at most three, simultaneous resolution of
the original bank and every chain bank costs only a fixed finite
number of terms, provided the asserted factors are
source-occurring.  The number is independent of support,
representations, physical multiplicities, and bank sizes.
\end{corollary}

\begin{proof}
There are at most four recorded banks and the number \(J\) of
covariance occurrences in the fixed \(2+1+1\) source word is fixed.
Apply \Cref{lem:V65-simultaneous-restriction}.  Any subsequent
protected/gapped choice is a fixed finite isotypic resolution on an
already exposed physical occurrence; centrality is V59.
\end{proof}

\begin{firewall}[Restriction does not prove occurrence]
\label{fire:V65-resolution-not-occurrence}
The conditional phrase in
\Cref{cor:V65-joint-resolution} is essential.  A bank selected by
V64 from the arrays \(a_{i,e}\) can be used to restrict every
existing occurrence, but restriction does not prove that a
particular occurrence carries the local heat/isotypic operator whose
gapped envelope is wanted.  That incidence must be checked in the
uncompressed \(2+1+1\) word.
\end{firewall}

\section{The endpoint mass retained by the V63--V64 chain}
\label{sec:V65-endpoint-mass}

Fix a remaining V63 lateral cell: \(u\in\{1,2\}\),
\(\Xi_u>K\Xi_4\), where \(K\) is the fixed constant in
\eqref{eq:V63-K-choice}.  Let \(a_0\) be the heavy endpoint of the
strong imbalanced edge selected in
\Cref{lem:V63-path-ratio}, and let
\(a_0,a_1,\dots,a_\ell\), \(\ell\le3\), be the successive heavy rows
in the V64 escalation, so
\begin{equation}
 \Xi_{a_{j+1}}\ge R\Xi_{a_j}.
\label{eq:V65-escalation-masses}
\end{equation}

\begin{lemma}[The terminal row controls the original lateral mass]
\label{lem:V65-terminal-controls-lateral}
For the retained chain,
\begin{equation}
 \boxed{\Xi_u\le K^{1/2}\Xi_{a_\ell}.}
\label{eq:V65-terminal-controls-lateral}
\end{equation}
\end{lemma}

\begin{proof}
If the V63 path from \(u\) to \(4\) has one edge, then \(a_0=u\).
If it has two edges \(u-v-4\), V63 chooses an edge whose oriented
mass ratio exceeds \(K^{1/2}\).  If this is the first edge, its heavy
endpoint is \(u\).  If it is the second, failure to choose the first
ratio means
\(\Xi_u/\Xi_v\le K^{1/2}\), while its heavy endpoint is \(v\).
Thus in both cases
\(\Xi_u\le K^{1/2}\Xi_{a_0}\).  The escalation inequalities
\eqref{eq:V65-escalation-masses} give
\(\Xi_{a_\ell}\ge\Xi_{a_0}\).
\end{proof}

\begin{lemma}[Cofinal gap gives an ordinary norm mass ratio]
\label{lem:V65-gap-norm-ratio}
Let \(F\) be the terminal V64 gapped bank in row \(a_\ell\), so
\begin{equation}
 S_{a_\ell}(F)\ge\sigma\Xi_{a_\ell}
\label{eq:V65-terminal-gap-cofinal}
\end{equation}
for a fixed \(\sigma>0\).  For every active row \(b\) and every fixed
integer \(m\ge1\),
\begin{equation}
 \boxed{
 \|\mathcal G_{a_\ell}(F)\|
 \le C_{G,s,\sigma,m}
 \left(\frac{\Xi_b}{\Xi_{a_\ell}}\right)^m.}
\label{eq:V65-gap-operator-ratio}
\end{equation}
In particular,
\begin{equation}
 \|\mathcal G_{a_\ell}(F)\|
 \le C_{G,s,\sigma}K^{1/2}
 \frac{\Xi_4}{\Xi_u}.
\label{eq:V65-gap-lateral-ratio}
\end{equation}
\end{lemma}

\begin{proof}
\Cref{lem:V64-Casimir-mass-gap} gives the ordinary operator bound
\[
 \|\mathcal G_{a_\ell}(F)\|
 \le e^{-c_G's\sigma\Xi_{a_\ell}}.
\]
The two cases in the proof of
\Cref{lem:V64-gap-pays-ratio} apply verbatim to operator norm:
if \(\Xi_b\ge\Xi_{a_\ell}\), use the contraction bound; otherwise
use \(\Xi_b\ge a_*\) and boundedness of
\(x^me^{-c_G's\sigma x}\).  This proves
\eqref{eq:V65-gap-operator-ratio}.  Take \(b=4,m=1\) and use
\eqref{eq:V65-terminal-controls-lateral}.
\end{proof}

\section{A jointly occurring terminal gap closes the lateral term}
\label{sec:V65-gapped-insertion}

Let \(D_0\) be the original root-owned comparable unmarked
row-\(3\)--row-\(u\) bank and put
\begin{equation}
 T_0=\sum_{e\in D_0}a_{u,e},
 \qquad
 \beta_0=\frac{T_0}{\Xi_u}.
\label{eq:V65-original-beta}
\end{equation}
Let \(\mathcal M_{D_0}\) be either its protected or positive gapped
central envelope.

\begin{theorem}[Exact terminal-gap insertion]
\label{thm:V65-terminal-gap-insertion}
Suppose the original factor \(\mathcal M_{D_0}\) and the terminal
cofinal gapped factor \(\mathcal G_{a_\ell}(F)\) are jointly
source-occurring in an outside-local restriction-resolved word.
After the simultaneous resolution, suppose the word has the exact
central presentation
\begin{equation}
 \widetilde W_{\pi,\nu}
 =
 \frac{p}{\Xi_3\Xi_4}\,
 L_{\pi,\nu}\,
 \mathcal M_{D_0}\,
 \mathcal G_{a_\ell}(F)\,
 Q_{\pi,\nu},
\label{eq:V65-joint-gapped-word}
\end{equation}
where the two displayed positive contractions commute with one
another and with \(Q_{\pi,\nu}\), and
\begin{equation}
 \|L_{\pi,\nu}\|+\|Q_{\pi,\nu}\|\le C_s.
\label{eq:V65-joint-dressings}
\end{equation}
Then
\begin{equation}
 \boxed{
 \kappa_\otimes(
  \widetilde W_{\pi,\nu}^*\widetilde W_{\pi,\nu})
 \le C_{s,K}\theta_{3u}(D_0)^2
 \le C_{s,K}(x+z)^2.}
\label{eq:V65-terminal-gap-RCQ}
\end{equation}
The same conclusion holds in the opposite square orientation.
\end{theorem}

\begin{proof}
Write \(M=\mathcal M_{D_0}\), \(G=\mathcal G_{a_\ell}(F)\), and
\(\lambda=p/(\Xi_3\Xi_4)\).  From
\eqref{eq:V65-joint-gapped-word} and
\eqref{eq:V65-joint-dressings},
\begin{align*}
 \widetilde W_{\pi,\nu}^*\widetilde W_{\pi,\nu}
 &=
 \lambda^2Q^*GM L^*LMGQ\\
 &\preccurlyeq
 C_s\lambda^2Q^*M^2G^2Q\\
 &\preccurlyeq
 C_s\lambda^2 M^2G^2\\
 &\preccurlyeq
 C_s\lambda^2\|G\|^2M^2.
\end{align*}
Here only the stated central commutations were used.  Product-state
capacity is monotone on positive operators, so
\[
 \kappa_\otimes(
  \widetilde W_{\pi,\nu}^*\widetilde W_{\pi,\nu})
 \le
 C_s
 \left(\frac{p}{\Xi_3\Xi_4}\right)^2
 \|G\|^2\kappa_\otimes(M^2).
\]
Apply \eqref{eq:V65-gap-lateral-ratio}; the factors \(\Xi_4^2\)
cancel:
\begin{equation}
 \kappa_\otimes(
  \widetilde W_{\pi,\nu}^*\widetilde W_{\pi,\nu})
 \le
 C_{s,K}
 \frac{p^2}{\Xi_3^2\Xi_u^2}
 \kappa_\otimes(M^2).
\label{eq:V65-after-gap}
\end{equation}
The original bank is nonempty, so \(T_0\ge a_*\); also
\(\Xi_3\ge a_*\), \(p\le1\), and
\(M^2\preccurlyeq M\).  Consequently
\[
 \frac{p^2}{\Xi_3^2\Xi_u^2}
 \kappa_\otimes(M^2)
 \le
 a_*^{-4}\beta_0^2\kappa_\otimes(M).
\]
The protected and gapped ownership-compensation bounds in
\Cref{thm:V60-ownership-compensation} now give the first inequality
in \eqref{eq:V65-terminal-gap-RCQ}.  The second is the row-\(3\)
cut estimate \(\theta_{3u}(D_0)\le x+z\).  Taking adjoints proves the
other orientation.
\end{proof}

\begin{corollary}[Finite recombination of the joint-gap cell]
\label{cor:V65-gap-recombination}
Every outside-local packet for which each nonzero simultaneous
resolution term satisfies the hypotheses of
\Cref{thm:V65-terminal-gap-insertion} obeys \(RCQ_3(\eta)\) on this
cell.
\end{corollary}

\begin{proof}
The number of simultaneous terms is fixed by
\Cref{cor:V65-joint-resolution}.  Apply
\Cref{lem:V58-finite-recombination} to
\eqref{eq:V65-terminal-gap-RCQ}.
\end{proof}

\begin{firewall}[The terminal heat is not a second prefix debit]
\label{fire:V65-no-double-charge}
The factor \(p\) in \eqref{eq:V65-joint-gapped-word} is the single
binary-prefix response already present in V62.  The terminal gain is
the ordinary norm of a different, actually occurring cofinal heat
factor.  The proof squares the one source word; it neither inserts a
second prefix response nor charges the original bank twice.
\end{firewall}

\section{Why the terminal-owned alternative does not yet multiply}
\label{sec:V65-owned-no-go}

\begin{proposition}[A scalar escalation chain creates no product
carrier]
\label{prop:V65-chain-no-product}
Apply any finite sequence of V14 restrictions to a fixed source word
along the banks selected by the V64 escalation algorithm.  If none
of the original source occurrences carries an upper norm response
factor \(\mathcal Z_j\) with
\begin{equation}
 \|\mathcal Z_j\|\le C\beta_j,
\label{eq:V65-chain-NPRC}
\end{equation}
then the coefficient-one restriction algebra does not produce the
scalar product \(\prod_j\beta_j\).  In particular, retaining the
labels of strong edges and owned terminal banks is not by itself a
proof of a multiplicative damping factor.
\end{proposition}

\begin{proof}
For one bank this is
\Cref{prop:V61-no-beta-in-V14}.  For finitely many banks use the
single Boolean-atom expansion
\eqref{eq:V65-atom-telescope}.  Every coefficient is one, and every
factor in every term is the restriction of an occurrence already in
the original word.  Induction over its fixed covariance occurrences
therefore produces only tensor factors and localized differences,
not a scalar depending on \(T_j/\Xi_{b_j}\).  Such a scalar can enter
only through an estimate of an actual response occurrence, such as
\eqref{eq:V65-chain-NPRC}, or through a normalization already
present before restriction.  The V64 algorithm changes neither.
\end{proof}

\begin{proposition}[Commuting capacity currencies need not
multiply]
\label{prop:V65-capacity-not-multiplicative}
There is no general inequality
\begin{equation}
 \kappa_\otimes(MN)
 \le
 \kappa_\otimes(M)\kappa_\otimes(N)
\label{eq:V65-false-capacity-product}
\end{equation}
for commuting positive contractions on a bipartite space.
\end{proposition}

\begin{proof}
On \(\mathbb C^d\otimes\mathbb C^d\), let
\[
 \Omega=d^{-1/2}\sum_{r=1}^de_r\otimes e_r,
 \qquad
 P=|\Omega\rangle\langle\Omega|,
 \qquad d\ge2.
\]
For product unit vectors \(x\otimes y\),
\[
 \langle x\otimes y,P(x\otimes y)\rangle
 =|\langle\Omega,x\otimes y\rangle|^2
 \le\frac1d,
\]
with equality for suitable basis vectors.  Hence
\(\kappa_\otimes(P)=1/d\).  Taking \(M=N=P\), the factors are
commuting positive contractions and \(MN=P^2=P\).  Therefore
\[
 \kappa_\otimes(MN)=\frac1d
 >
 \frac1{d^2}
 =
 \kappa_\otimes(M)\kappa_\otimes(N).
\]
\end{proof}

\begin{assessment}[Exact residual after V65]
\label{ass:V65-residual}
The cofinal-gapped terminal alternative is now analytically closed
whenever the original bank and terminal heat factor are jointly
source-occurring.  Ordinary heat norm, rather than a second capacity
currency, transports the unused row-\(4\) denominator to the large
lateral mass.

The terminal-owned alternative is different.  Its gain is presently
a product-state capacity, not a small ordinary norm.  By
\Cref{prop:V65-capacity-not-multiplicative}, it cannot simply be
multiplied by the original bank capacity.  By
\Cref{prop:V65-chain-no-product}, the scalar chain cannot replace the
missing operator carrier.  The remaining work is therefore a finite
incidence audit of the uncompressed \(2+1+1\) source word: determine
which chain banks are jointly source-occurring and, at the first
failure, return to the exact covariance or composite joining
difference that generated that edge.
\end{assessment}

\section{V65 result ledger and V66 acquisition order}
\label{sec:V65-status}

\begin{theorem}[Third-series V65 result ledger]
\label{thm:V65-results}
Relative to V1--V64, V65 proves:
\begin{enumerate}[label=\textup{(V65.\arabic*)},leftmargin=6.3em]
\item the compulsory active SM V46 and NS V31 audit, with exact
      metadata and a type-correct proof-order transfer;
\item an exact simultaneous restriction theorem for finitely many
      banks, with a support-independent term count;
\item a precise distinction between a source-occurring bank and a
      bank selected only from scalar weights;
\item control of the original lateral mass by the terminal V64 row;
\item an ordinary operator-norm mass-ratio estimate for a cofinal
      gapped bank;
\item quadratic closure of every jointly source-occurring terminal
      gap cell, including finite recombination;
\item a proof that coefficient-one chain resolution creates no
      product of ownership fractions; and
\item an explicit counterexample to multiplicative product-state
      capacity for commuting positive contractions.
\end{enumerate}
\end{theorem}

\begin{proof}
These are
\Cref{sec:V65-audit,
lem:V65-simultaneous-restriction,
def:V65-source-occurring,
cor:V65-joint-resolution,
lem:V65-terminal-controls-lateral,
lem:V65-gap-norm-ratio,
thm:V65-terminal-gap-insertion,
cor:V65-gap-recombination,
prop:V65-chain-no-product,
prop:V65-capacity-not-multiplicative}.
\end{proof}

\begin{programme}[V66: exact source-occurrence incidence table]
\label{prog:V65-V66}
\begin{enumerate}[leftmargin=3em]
\item Restore the full uncompressed word from
      \eqref{eq:V49-complete-word}, including prefix covariance,
      both singleton occurrences, the composite joining difference,
      suffix, payer, and replica orientation.
\item For the original bank \(D_0\) and every V64 chain bank, record
      the exact factor slot and local \(G\)-module on which its
      protected/gapped resolution would act.
\item Apply \Cref{lem:V65-simultaneous-restriction} only after this
      occurrence table is complete.  Route every genuine terminal
      gap through \Cref{thm:V65-terminal-gap-insertion}.
\item In the terminal-owned branch, stop at the first bank whose
      only evidence is a scalar strong edge.  Expand the exact
      covariance or unsplit composite joining occurrence at that
      slot and test it for an NPRC upper bound.
\item Do not multiply capacity currencies without an additional
      tensor-independence or ordinary-norm theorem.  The next
      compulsory companion audit is V70.
\end{enumerate}
\end{programme}

\begin{firewall}[Claim boundary after third-series V65]
\label{fire:V65-claim-boundary}
V65 closes the jointly source-occurring cofinal-gapped terminal cell.
It does not prove that every V64 terminal bank is an occurrence in
the fixed \(2+1+1\) word, close the terminal-owned, finite
exceptional, or local-payer branches, or establish full
\(RCQ_3(\eta)\).  The other quotient cuts, quartic source families,
higher MTA, WUFC, CPM, cutoff removal, reflection positivity,
Osterwalder--Schrader reconstruction, construction of the continuum
Yang--Mills measure, and a uniform positive continuum spectral gap
remain open.  The full Yang--Mills mass gap has not been proved.
\end{firewall}

\paragraph{Parent-version and correction record.}
V65 began from
\texttt{Renewal\_Yang\_Mills\_Mass\_Gap\_Research\_Notes\_V64.tex}.
The V64 parent had \(31\,542\) logical source lines,
\(1\,135\,364\) bytes, and SHA--256
\[
 \texttt{ffb7fe75ff3c3014c2c032630b34e66b0836066bbe0cb7c5d452463b09ceb2e6}.
\]
Before appending V65, two missing backslashes in V64 displays were
corrected: the literal source text \texttt{quad} in
\eqref{eq:V64-V63-data} and in the proof of
\Cref{thm:V64-full-row-trichotomy} now reads \verb|\qquad|.  These
were typographical corrections only.  The next compulsory companion
audit is V70.

% =====================================================================
% V66 APPENDIX: EXACT SOURCE-SLOT INCIDENCE
% =====================================================================

\chapter{V66: Exact Source-Slot Incidence and Closure of Every
Outside-Local Terminal-Gap Branch}
\label{ch:V66-source-incidence}

\section{Return to the uncompressed first-connectivity summand}
\label{sec:V66-uncompressed}

V65 deliberately left joint source occurrence as a hypothesis.  The
fixed joining-coordinate summand of the already established
first-connectivity identity \eqref{eq:V35-S211} is
\begin{equation}
 \mathcal W_r
 =
 K_{12}^{<r,\mathrm{conn}}
 \otimes M_{\{3\}}^{<r}
 \otimes M_{\{4\}}^{<r}
 \otimes J_r(12|3|4)
 \otimes M_{>r}.
\label{eq:V66-uncompressed-word}
\end{equation}
Separators, physical compression, and one fixed payer placement are
understood in their original order.  No norm has been taken in
\eqref{eq:V66-uncompressed-word}.

The coordinate slots are:
\begin{enumerate}[label=\textup{(\roman*)},leftmargin=4em]
\item \(e<r\): the complete signed binary prefix covariance on rows
      \(1,2\), together with the two complete singleton moments on
      rows \(3,4\);
\item \(e=r\): the assembled local composite third cumulant
      \(J_r(12|3|4)\); and
\item \(e>r\): the complete four-row suffix moment \(M_{>r}\).
\end{enumerate}
Every V63--V64 chain bank is unmarked and excludes the selected
joining coordinate and every payer coordinate.  Thus it meets only
the first and third slots.

\begin{definition}[Complete prefix carrier]
\label{def:V66-prefix-carrier}
Write
\begin{equation}
 \mathcal C_{<r}
 :=
 K_{12}^{<r,\mathrm{conn}}
 \otimes M_{\{3\}}^{<r}
 \otimes M_{\{4\}}^{<r}.
\label{eq:V66-prefix-carrier}
\end{equation}
An \emph{endpoint expansion} of \(\mathcal C_{<r}\) means: first
apply the coefficient-one covariance telescope to the binary
prefix, and then write each remaining localized covariance
\(\Delta_A^{\rho,\sigma}\) as
\(M_A(\rho)-M_A(\sigma)\).  The singleton moments are left complete.
\end{definition}

\begin{lemma}[Every endpoint-expanded prefix atom is a complete
local moment state]
\label{lem:V66-prefix-endpoints}
Fix finitely many banks and form their Boolean atoms as in
\eqref{eq:V65-Boolean-atoms}.  The complete prefix carrier
\(\mathcal C_{<r}\) is an exact finite signed sum such that, on every
atom \(A\subset\{e:e<r\}\), every summand is a tensor product of:
\begin{enumerate}[label=\textup{(\alph*)},leftmargin=4em]
\item one complete endpoint moment on the row-\(1,2\) block; and
\item the complete row-\(3\) and row-\(4\) singleton moments.
\end{enumerate}
Consequently the full local tensor on \(A\) is an actual complete
moment-partition state to which the V32 protected/gapped resolution
applies.  The number of signed endpoint terms is fixed independently
of support and representations.
\end{lemma}

\begin{proof}
Apply \Cref{lem:V65-simultaneous-restriction} to the complete binary
prefix covariance.  On an atom to which that covariance is
localized, use its definition
\[
 \Delta_A^{\rho,\sigma}=M_A(\rho)-M_A(\sigma).
\]
This gives two complete endpoint moments.  On every other atom, the
tensor telescope already supplies a complete endpoint moment.  The
row-\(3\) and row-\(4\) prefix factors tensorize as moments and
introduce no sum.  There is a fixed number of covariance occurrences
and finitely many banks, so the total expansion is fixed.  Before
physical compression, each endpoint factor and the singleton
factors are precisely the local moment blocks used in the
moment-partition resolution of V32.
\end{proof}

\begin{lemma}[Every suffix atom is already a complete moment state]
\label{lem:V66-suffix-endpoints}
On every Boolean atom
\(A\subset\{e:e>r\}\), restriction of \(M_{>r}\) is a complete
four-row moment.  Its protected/gapped resolution is therefore an
actual source occurrence, with no endpoint expansion.
\end{lemma}

\begin{proof}
This is the complete-moment identity
\eqref{eq:V61-moment-one} applied to the suffix support.  It is also
the contraction structure used in the V67 payer ledger: every
unpaid suffix factor is a complete moment, while the sole paid
suffix factor is formed only after the complete product is kept.
\end{proof}

\begin{theorem}[Exact chain-bank incidence in the \(2+1+1\) word]
\label{thm:V66-chain-incidence}
Fix an outside-local payer placement and one selected joining
coordinate \(r\).  Let \(D_1,\dots,D_k\), \(k\le4\), be the original
bank and the retained banks of a V64 chain.  After a single exact
finite expansion of \eqref{eq:V66-uncompressed-word}:
\begin{enumerate}[label=\textup{(\roman*)},leftmargin=4em]
\item every \(D_j\cap\{e<r\}\) is source-occurring in the endpoint
      expansion of \(\mathcal C_{<r}\);
\item every \(D_j\cap\{e>r\}\) is source-occurring in \(M_{>r}\);
\item all those occurrences are joint in the sense of
      \Cref{def:V65-source-occurring}; and
\item their positive spectral factors commute with each other and
      with every source separator allowed by
      \Cref{lem:V59-bank-commutation}.
\end{enumerate}
The expansion count is independent of support, representations,
physical multiplicities, and bank sizes.
\end{theorem}

\begin{proof}
Use the common Boolean atoms for all \(D_j\).
\Cref{lem:V66-prefix-endpoints} treats every prefix atom and
\Cref{lem:V66-suffix-endpoints} treats every suffix atom.  The
selected coordinate \(r\) and the payer coordinates do not belong
to the banks, so there is no remaining coordinate slot.  Perform all
endpoint expansions before inserting protected/gapped isotypic
resolutions.  This produces one common exact expansion rather than
separate incompatible expansions for the different banks.

Each positive spectral factor is a scalar function of the local
isotypic label.  Pairwise commutation and commutation with the
remaining coordinatewise intertwiners follow from
\Cref{lem:V59-central-isotypic,lem:V59-bank-commutation}.  The term
count is bounded by \Cref{lem:V65-simultaneous-restriction} times a
fixed endpoint factor for each of the fixed number of binary
covariance occurrences.
\end{proof}

\begin{firewall}[The joining cumulant was not split]
\label{fire:V66-joining-unsplit}
The incidence expansion does not polarize or separately majorize
\(J_r(12|3|4)\).  Chain banks exclude \(r\), so the local composite
cumulant remains one assembled occurrence.  In particular, V66
does not replace it by its five partition roles or discard the mixed
Leibniz blocks of V49.
\end{firewall}

\section{Cofinality survives the prefix--suffix split}
\label{sec:V66-half-bank}

\begin{lemma}[One side retains half the terminal mass]
\label{lem:V66-half-cofinal}
Let \(F\) be a terminal V64 gapped bank in row \(a\), with
\begin{equation}
 S_a(F)\ge\sigma\Xi_a.
\label{eq:V66-F-cofinal}
\end{equation}
Put
\[
 F^-:=F\cap\{e:e<r\},
 \qquad
 F^+:=F\cap\{e:e>r\}.
\]
Then one may choose
\(\widehat F\in\{F^-,F^+\}\) such that
\begin{equation}
 S_a(\widehat F)\ge\frac{\sigma}{2}\Xi_a.
\label{eq:V66-half-cofinal}
\end{equation}
The chosen bank remains unbalanced at row \(a\) at every coordinate.
\end{lemma}

\begin{proof}
Since \(r\notin F\), the two displayed banks partition \(F\), and
\[
 S_a(F)=S_a(F^-)+S_a(F^+).
\]
At least one summand is at least half the total.  Unbalance is a
coordinatewise property inherited by subsets.
\end{proof}

\begin{corollary}[The chosen half-bank is jointly source-occurring]
\label{cor:V66-half-source}
The original row-\(3\)--row-\(u\) bank \(D_0\) and the terminal
gapped half-bank \(\widehat F\) in
\Cref{lem:V66-half-cofinal} are jointly source-occurring after the
single expansion of \Cref{thm:V66-chain-incidence}.
\end{corollary}

\begin{proof}
Include \(D_0,F\) among the banks in
\Cref{thm:V66-chain-incidence}.  The half-bank is a union of Boolean
atoms on one side of \(r\); retain the corresponding local spectral
factors.  Both factors occur in the same resolved terms, including
when their coordinate sets overlap, because central spectral
functions on a common atom commute.  A sector with incompatible
spectral choices is simply the zero term.
\end{proof}

\section{Unconditional closure of the terminal-gap branch}
\label{sec:V66-terminal-gap}

\begin{theorem}[Every outside-local terminal-gap chain closes]
\label{thm:V66-all-terminal-gaps}
Fix \(0<\eta\le1\) and the fixed \(K\) chosen in
\eqref{eq:V63-K-choice}.  In every remaining large-lateral cell of
\Cref{cor:V63-comparable-residual}, suppose the V64 finite chain
terminates in alternative \textup{(a)} of
\Cref{thm:V64-full-row-trichotomy}.  Then every outside-local
restriction-resolved normalized word satisfies
\begin{equation}
 \boxed{
 \kappa_\otimes(
  \widetilde W_{\pi,\nu}^*\widetilde W_{\pi,\nu})
 \le C_{s,K}\theta_{3u}(D_0)^2
 \le C_{s,K}(x+z)^2.}
\label{eq:V66-all-terminal-gaps}
\end{equation}
The opposite square orientation obeys the same bound.  After finite
recombination, the complete outside-local terminal-gap packet
satisfies \(RCQ_3(\eta)\).
\end{theorem}

\begin{proof}
Choose the half-bank \(\widehat F\) by
\Cref{lem:V66-half-cofinal}.  It has cofinality constant
\(\sigma/2\).  By \Cref{cor:V66-half-source}, the original positive
bank factor and
\(\mathcal G_a(\widehat F)\) are jointly source-occurring.

The scalar normalization and uniform dressing bounds are exactly
\eqref{eq:V62-normalized-factorization}--%
\eqref{eq:V62-normalized-dressings}.  All additional endpoint
moments are contractions, and the fixed endpoint expansion changes
only the uniform constant.  Centrality permits the order
\eqref{eq:V65-joint-gapped-word}.  Apply
\Cref{thm:V65-terminal-gap-insertion}, with \(\sigma\) replaced by
\(\sigma/2\), to obtain \eqref{eq:V66-all-terminal-gaps}.

The simultaneous restriction and endpoint expansion have a fixed
number of terms by \Cref{thm:V66-chain-incidence}.  Apply
\Cref{lem:V58-finite-recombination}.  The adjoint orientation is the
left-square clause of
\Cref{thm:V65-terminal-gap-insertion}.
\end{proof}

\begin{assessment}[A V65 hypothesis has been discharged]
\label{ass:V66-gap-discharged}
Joint occurrence is no longer an open condition for an unmarked
outside-local terminal gap in the fixed \(2+1+1\) word.  The exact
first-connectivity slots, endpoint expansion, and half-bank lemma
supply it.  Therefore a V64 chain ending in a cofinal heavy-row gap
cannot remain in an \(RCQ_3\) counterpacket.
\end{assessment}

\section{The row-\(3\)-incident terminal-owned branch also closes}
\label{sec:V66-owned-row3}

\begin{corollary}[Terminal ownership incident to the cut row]
\label{cor:V66-row3-owned}
Suppose the V64 chain terminates in an owned comparable bank
\(D_\ell\) whose two row labels are \(3\) and
\(c\in\{1,2,4\}\).  Then every outside-local
restriction-resolved term satisfies
\begin{equation}
 \kappa_\otimes(
  \widetilde W_{\pi,\nu}^*\widetilde W_{\pi,\nu})
 \le C_s\theta_{3c}(D_\ell)^2
 \le C_s(x+z)^2.
\label{eq:V66-row3-owned}
\end{equation}
The packet closes after finite recombination.
\end{corollary}

\begin{proof}
\Cref{thm:V66-chain-incidence} makes \(D_\ell\)
source-occurring in every relevant prefix or suffix sector.  Resolve
it into its protected and gapped positive central envelopes.  Since
the bank is owned at both endpoints and comparable,
\Cref{thm:V59-physical-owned-bank} applies to each sector and gives
the first inequality.  The edge is incident to row \(3\), so its
normalized bank overlap is a substock of
\(\theta_{31}+\theta_{32}+\theta_{34}=x+z\).  Fixed finite
recombination gives the packet statement.
\end{proof}

\begin{corollary}[Exact terminal-owned residual graph]
\label{cor:V66-owned-residual}
After V66, the only terminal-owned V64 chains not closed by incidence
and central insertion have their terminal edge entirely in
\begin{equation}
 \{12,14,24\}.
\label{eq:V66-lateral-owned-edges}
\end{equation}
These edges lie in the graph obtained after deleting the cut row
\(3\).
\end{corollary}

\begin{proof}
There are six edges on four rows.  The three edges incident to row
\(3\) close by \Cref{cor:V66-row3-owned}.  The other three are
exactly \eqref{eq:V66-lateral-owned-edges}.
\end{proof}

\begin{firewall}[Why a lateral owned edge is not yet the cut
currency]
\label{fire:V66-lateral-not-cut}
For \(ab\in\{12,14,24\}\), the scalar
\(\theta_{ab}(D_\ell)\) need not tend to zero with \(x+z\).
The central bank estimate at that edge therefore cannot replace the
missing row-\(3\) partner fraction.  Multiplying its capacity by the
capacity of the original row-\(3\)--row-\(u\) bank is forbidden by
\Cref{prop:V65-capacity-not-multiplicative}.  A new estimate must
retain an actual signed occurrence connecting the lateral terminal
edge back to row \(3\).
\end{firewall}

\section{V66 result ledger and V67 acquisition order}
\label{sec:V66-status}

\begin{theorem}[Third-series V66 result ledger]
\label{thm:V66-results}
Relative to V1--V65, V66 proves:
\begin{enumerate}[label=\textup{(V66.\arabic*)},leftmargin=6.3em]
\item the exact prefix/joining/suffix occurrence table for a fixed
      \(2+1+1\) first-connectivity summand;
\item a fixed endpoint expansion making every prefix Boolean atom a
      complete local moment-partition state;
\item source occurrence of every suffix bank without expansion;
\item joint source occurrence and central commutation of all
      unmarked banks in a V64 chain;
\item preservation of half the terminal cofinal mass on one side of
      the joining coordinate;
\item unconditional \(RCQ_3\) closure of every outside-local chain
      ending in a cofinal gapped bank;
\item closure of every terminal-owned edge incident to row \(3\);
      and
\item reduction of the terminal-owned residual to the lateral graph
      \(\{12,14,24\}\).
\end{enumerate}
\end{theorem}

\begin{proof}
These are
\Cref{eq:V66-uncompressed-word,
lem:V66-prefix-endpoints,
lem:V66-suffix-endpoints,
thm:V66-chain-incidence,
lem:V66-half-cofinal,
cor:V66-half-source,
thm:V66-all-terminal-gaps,
cor:V66-row3-owned,
cor:V66-owned-residual}.
\end{proof}

\begin{programme}[V67: the first signed bridge back to row \(3\)]
\label{prog:V66-V67}
\begin{enumerate}[leftmargin=3em]
\item Keep the residual terminal edge
      \(ab\in\{12,14,24\}\), the original row-\(3\)--row-\(u\)
      bank \(D_0\), and the exact source slot of each in one
      simultaneous expansion.
\item If \(ab=12\), return to the complete signed binary prefix
      covariance before its endpoint expansion.  If
      \(ab\in\{14,24\}\), return to the unsplit composite difference
      in \eqref{eq:V56-composite-difference} together with the
      row-\(4\) difference.
\item Derive the exact conditional square or covariance identity
      which contains both the lateral bank and the row-\(3\)
      occurrence.  Do not multiply their separate capacities.
\item Test that actual bridge for an NPRC upper norm
      \(O(T(D_0)/\Xi_u)\).  A scalar lower response is not
      sufficient.
\item Keep the local-payer, marked-heavy, and finite exceptional
      branches separate.  The next compulsory companion audit is
      V70.
\end{enumerate}
\end{programme}

\begin{firewall}[Claim boundary after third-series V66]
\label{fire:V66-claim-boundary}
V66 closes every outside-local cofinal terminal-gap chain and every
terminal-owned chain whose final edge meets row \(3\).  It does not
close lateral terminal-owned edges, finite exceptional or
local-payer branches, or full \(RCQ_3(\eta)\).  The other quotient
cuts, quartic source families, higher MTA, WUFC, CPM, cutoff removal,
reflection positivity, Osterwalder--Schrader reconstruction,
construction of the continuum Yang--Mills measure, and a uniform
positive continuum spectral gap remain open.  The full Yang--Mills
mass gap has not been proved.
\end{firewall}

\paragraph{Parent-version record.}
V66 was formed from
\texttt{Renewal\_Yang\_Mills\_Mass\_Gap\_Research\_Notes\_V65.tex}.
The V65 parent had \(32\,089\) logical source lines,
\(1\,155\,330\) bytes, and SHA--256
\[
 \texttt{de57a93f281798d4c80ae16c68fe8086ac7b1712561610fbc21bca752390080d}.
\]
The next compulsory companion audit is V70.

% =====================================================================
% V67 APPENDIX: CUT ALIGNMENT AND THE CONDITIONAL BRIDGE
% =====================================================================

\chapter{V67: Cut Alignment, Slot-Level No-Go Theorems, and the
Exact Conditional Bridge}
\label{ch:V67-cut-alignment}

\section{Source occurrence is not cut alignment}
\label{sec:V67-cut-mismatch}

V66 proves that every retained bank is an actual central spectral
factor in one common expansion.  It does not say that a capacity
estimate rooted at a terminal lateral row is an estimate across the
row-\(3\) cut.  We now write the cut on the capacity:
\[
 \kappa_{3|c}:=\kappa_\otimes^{\,3|\{1,2,4\}},
 \qquad
 \kappa_{a|a^c}:=\kappa_\otimes^{\,a|[4]\setminus\{a\}}.
\]

\begin{proposition}[Small capacity on a lateral cut need not be
small on the row-\(3\) cut]
\label{prop:V67-cut-transfer-no-go}
For every \(d\ge2\), there is a positive contraction \(P_d\) on a
three-row Hilbert space such that
\begin{equation}
 \kappa_{1|\{2,3\}}(P_d)=\frac1d,
 \qquad
 \kappa_{3|\{1,2\}}(P_d)=1.
\label{eq:V67-cut-transfer-no-go}
\end{equation}
Consequently no representation-independent bound can transfer a
terminal capacity currency on a lateral cut to the row-\(3\) cut
without additional source incidence.
\end{proposition}

\begin{proof}
Take
\[
 \mathcal H_1=\mathcal H_2=\mathbb C^d,
 \qquad
 \mathcal H_3\ne\{0\},
 \qquad
 \Omega_d=d^{-1/2}\sum_{j=1}^de_j\otimes e_j,
\]
and set
\[
 P_d=|\Omega_d\rangle\langle\Omega_d|_{12}
 \otimes I_3.
\]
Across the cut \(1|\{2,3\}\), a product unit vector has the form
\(x_1\otimes z_{23}\).  Expanding \(z_{23}\) in an orthonormal basis
of \(\mathcal H_3\) and applying Cauchy--Schwarz to the maximally
entangled projection gives expectation at most \(1/d\); equality is
attained by
\(e_1\otimes e_1\otimes v_3\).  Hence the first equality holds.

Across \(3|\{1,2\}\), choose the product vector
\(v_3\otimes\Omega_d\), with the tensor factors reordered to
\(\mathcal H_3\otimes(\mathcal H_1\otimes\mathcal H_2)\).
Its expectation is one.  Since \(P_d\) is a contraction, the second
capacity equals one.
\end{proof}

\begin{corollary}[What V57 does not imply at a lateral terminal
edge]
\label{cor:V67-V57-cut-scope}
If the terminal owned edge
\(ab\in\{12,14,24\}\), the V57 estimate
\(\kappa_{a|a^c}(\mathcal M_{ab})\lesssim
\theta_{ab}(D)^2\) does not by itself bound
\(\kappa_{3|c}(\mathcal M_{ab})\).  Source occurrence from V66 and
central commutation from V59 do not change this conclusion.
\end{corollary}

\begin{proof}
The two capacities optimize over different families of product
states.  \Cref{prop:V67-cut-transfer-no-go} rules out a general
comparison even for a central-looking rank-one projection tensored
with an identity spectator.  V66 proves where the operator occurs,
not a relation between these two state families.
\end{proof}

\begin{firewall}[No relabelling of the terminal root]
\label{fire:V67-no-cut-relabel}
The terminal root \(a\) was selected by a scalar mass-escalation
algorithm.  It may not be relabelled as row \(3\) inside a fixed
physical source word.  Row labels determine the product-state cut,
the prefix partition, and the composite joining occurrence.
\end{firewall}

\section{Exact conditional slices of the row-\(3\) capacity}
\label{sec:V67-slices}

Let
\(\mathcal H=\mathcal H_3\otimes\mathcal H_c\), where
\(\mathcal H_c=\mathcal H_{\{1,2,4\}}\) includes all multiplicity
spaces.  For \(A\in\mathcal B(\mathcal H)\) and a unit vector
\(w\in\mathcal H_c\), define the row-\(3\) slice
\begin{equation}
 A[w]
 :=
 (I_3\otimes\langle w|)\,A\,
 (I_3\otimes|w\rangle)
 \in\mathcal B(\mathcal H_3).
\label{eq:V67-slice}
\end{equation}

\begin{lemma}[Exact slice formula for positive capacity]
\label{lem:V67-slice-capacity}
If \(A\succcurlyeq0\), then
\begin{equation}
 \boxed{
 \kappa_{3|c}(A)
 =
 \sup_{\|w\|=1}\|A[w]\|.}
\label{eq:V67-slice-capacity}
\end{equation}
\end{lemma}

\begin{proof}
For unit vectors \(v\in\mathcal H_3\) and
\(w\in\mathcal H_c\),
\[
 \langle v\otimes w,A(v\otimes w)\rangle
 =
 \langle v,A[w]v\rangle.
\]
The slice is positive.  For fixed \(w\), the supremum over \(v\) is
therefore its operator norm.  Taking the supremum over \(w\) gives
\eqref{eq:V67-slice-capacity}.  Mixed product states do not enlarge
the supremum because they are convex combinations of pure product
states.
\end{proof}

Let \(D_0\) be the original row-\(3\)--row-\(u\) bank,
\(u\in\{1,2\}\), let \(M=\mathcal M_{D_0}\), and let \(N\) be a
positive central factor coming from a terminal lateral owned bank.
Both are contractions and commute after V66 resolution.

\begin{definition}[Lateral conditional bridge card]
\label{def:V67-LCB}
The card \(LCB_{3,u}(M,N)\) is the estimate
\begin{equation}
 \boxed{
 \sup_{\|w\|=1}
 \|(M^2N^2)[w]\|
 \le
 C_s
 \left(\frac{\Xi_4}{\Xi_u}\right)^2
 \kappa_{3|c}(M).}
\label{eq:V67-LCB}
\end{equation}
The constant must be uniform in support, representations, physical
multiplicities, and the selected coordinate.
\end{definition}

\begin{theorem}[The exact conditional bridge is sufficient]
\label{thm:V67-LCB-suffices}
Suppose a residual outside-local resolved word has the central
presentation
\begin{equation}
 \widetilde W_{\pi,\nu}
 =
 \frac{p}{\Xi_3\Xi_4}\,
 L_{\pi,\nu}MNQ_{\pi,\nu},
\qquad
 \|L_{\pi,\nu}\|+\|Q_{\pi,\nu}\|\le C_s,
\label{eq:V67-lateral-word}
\end{equation}
with \(M,N\) commuting with \(Q_{\pi,\nu}\) and with each other.  If
\eqref{eq:V67-LCB} holds, then
\begin{equation}
 \boxed{
 \kappa_{3|c}(
 \widetilde W_{\pi,\nu}^*\widetilde W_{\pi,\nu})
 \le C_s\theta_{3u}(D_0)^2
 \le C_s(x+z)^2.}
\label{eq:V67-LCB-RCQ}
\end{equation}
\end{theorem}

\begin{proof}
As in the proof of
\Cref{thm:V65-terminal-gap-insertion}, centrality gives
\[
 \widetilde W_{\pi,\nu}^*\widetilde W_{\pi,\nu}
 \preccurlyeq
 C_s
 \left(\frac{p}{\Xi_3\Xi_4}\right)^2
 M^2N^2.
\]
Use monotonicity, \Cref{lem:V67-slice-capacity}, and
\eqref{eq:V67-LCB}:
\[
 \kappa_{3|c}(
 \widetilde W_{\pi,\nu}^*\widetilde W_{\pi,\nu})
 \le
 C_s
 \frac{p^2}{\Xi_3^2\Xi_u^2}
 \kappa_{3|c}(M).
\]
The original bank is nonempty, so
\(T_0\ge a_*\), while \(p\le1\) and \(\Xi_3\ge a_*\).
Thus
\[
 \frac{p^2}{\Xi_3^2\Xi_u^2}\kappa_{3|c}(M)
 \le
 a_*^{-4}\beta_0^2\kappa_{3|c}(M).
\]
Apply the two protected/gapped estimates in
\Cref{thm:V60-ownership-compensation}, then
\(\theta_{3u}(D_0)\le x+z\).
\end{proof}

\begin{corollary}[V66 terminal heat proves the bridge card]
\label{cor:V67-gap-is-LCB}
If \(N\) is a terminal cofinal gapped factor satisfying
\(\|N\|\le C_s\Xi_4/\Xi_u\), then
\eqref{eq:V67-LCB} holds.
\end{corollary}

\begin{proof}
Because \(M,N\) commute and are positive,
\[
 0\preccurlyeq M^2N^2
 \preccurlyeq\|N\|^2M^2
 \preccurlyeq\|N\|^2M.
\]
Take row-\(3\) capacity, use
\(\kappa_{3|c}(M^2)\le\kappa_{3|c}(M)\), and insert the norm
hypothesis.  The slice formulation is
\Cref{lem:V67-slice-capacity}.
\end{proof}

\begin{assessment}[The remaining owned branch is one precise
conditional inequality]
\label{ass:V67-LCB-position}
V65--V66 did not merely prove one isolated heat estimate: they
verified \(LCB_{3,u}\) for every cofinal terminal gap.  A protected
lateral terminal factor has norm one, so the same proof gives no
mass ratio.  The residual analytic question is exactly whether its
conditional slices against the original row-\(3\) bank satisfy
\eqref{eq:V67-LCB}.
\end{assessment}

\section{The isolated prefix and joining slots cannot be the NPRC}
\label{sec:V67-slot-no-go}

\begin{lemma}[Slot invariance under later spectator extension]
\label{lem:V67-slot-invariance}
Fix a joining coordinate \(r\).  Adding coordinates strictly after
\(r\) changes neither the complete prefix covariance
\(K_{12}^{<r,\mathrm{conn}}\) nor the local joining carrier
\(J_r(12|3|4)\).  It changes only the complete suffix and the global
row masses.
\end{lemma}

\begin{proof}
This is immediate from the ordered supports in
\eqref{eq:V66-uncompressed-word}: the prefix uses only \(e<r\), the
joining carrier uses \(e=r\), and every \(e>r\) belongs to
\(M_{>r}\).
\end{proof}

\begin{proposition}[No isolated prefix NPRC]
\label{prop:V67-prefix-NPRC-no-go}
Fix a nonzero prefix restriction \(Z_{<r}\) and an original bank
\(D_0\subset\{e:e<r\}\) with fixed \(T_0>0\).  The prefix algebra
alone cannot prove a support-uniform estimate
\begin{equation}
 \|Z_{<r}\|
 \le C\frac{T_0}{\Xi_u}.
\label{eq:V67-false-prefix-NPRC}
\end{equation}
\end{proposition}

\begin{proof}
Adjoin after \(r\) a row-\(u\) spectator stock of mass \(A>0\).
By \Cref{lem:V67-slot-invariance}, \(Z_{<r}\) and \(T_0\) are
unchanged, whereas \(\Xi_u\) becomes \(\Xi_u+A\).  The right side of
\eqref{eq:V67-false-prefix-NPRC} tends to zero as \(A\to\infty\),
while the nonzero left side is fixed.  This is impossible.  The
argument concerns an isolated prefix factor; the complete suffix is
allowed to respond to the added stock.
\end{proof}

\begin{proposition}[No isolated local-joining NPRC]
\label{prop:V67-joining-NPRC-no-go}
Fix a nonzero local restriction \(Z_r\) of
\(J_r(12|3|4)\) and a bank \(D_0\) away from \(r\) with fixed
\(T_0>0\).  The local joining algebra alone cannot prove
\begin{equation}
 \|Z_r\|
 \le C\frac{T_0}{\Xi_u}
\label{eq:V67-false-joining-NPRC}
\end{equation}
uniformly under spectator extensions.
\end{proposition}

\begin{proof}
Use the same later spectator extension.  The local joining letter is
unchanged by \Cref{lem:V67-slot-invariance}, while the proposed right
side tends to zero.
\end{proof}

\begin{firewall}[Scope of the slot no-go]
\label{fire:V67-slot-scope}
\Cref{prop:V67-prefix-NPRC-no-go,prop:V67-joining-NPRC-no-go} do not
rule out a complete-word cancellation or a suffix heat response.
They prove that neither the prefix nor the selected local joining
factor, estimated in isolation, is the missing NPRC.  The next
argument must retain the suffix or sum the exact first-connectivity
positions before taking a norm.
\end{firewall}

\section{The first legitimate upstream carrier}
\label{sec:V67-upstream-carrier}

For fixed payer type \(\pi\), let
\(\mathfrak I_\pi^{\mathrm{lat}}\) be the set of joining coordinates
and endpoint-expansion labels whose V64 chain terminates in the
lateral owned residual.  Define the signed aggregate before positive
majorization:
\begin{equation}
 \mathbf W_\pi^{\mathrm{lat}}
 :=
 \sum_{(r,\nu)\in\mathfrak I_\pi^{\mathrm{lat}}}
 \widetilde W_{r,\pi,\nu}.
\label{eq:V67-lateral-aggregate}
\end{equation}

\begin{theorem}[Exact aggregate square retains every possible
bridge]
\label{thm:V67-aggregate-square}
One has the operator identity
\begin{equation}
 \boxed{
 (\mathbf W_\pi^{\mathrm{lat}})^*
 \mathbf W_\pi^{\mathrm{lat}}
 =
 \sum_{\substack{(r,\nu)\in\mathfrak I_\pi^{\mathrm{lat}}\\
                 (s,\mu)\in\mathfrak I_\pi^{\mathrm{lat}}}}
 \widetilde W_{r,\pi,\nu}^*
 \widetilde W_{s,\pi,\mu}.}
\label{eq:V67-aggregate-square}
\end{equation}
The diagonal terms contain the conditional bridge target
\eqref{eq:V67-LCB}; the off-diagonal terms retain the only possible
cancellation between different first-connectivity positions.  The
same payer label \(\pi\) occurs in both factors.
\end{theorem}

\begin{proof}
Multiply the finite signed sum
\eqref{eq:V67-lateral-aggregate} by its adjoint and distribute.
Nothing is estimated.  The diagonal is the fixed-word square in
\eqref{eq:V67-lateral-word}.  The common label \(\pi\) follows from
forming the aggregate only after fixing the unique payer placement.
\end{proof}

\begin{firewall}[No entrywise absolute value]
\label{fire:V67-no-entrywise}
Applying capacity or absolute value to each summand in
\eqref{eq:V67-aggregate-square} returns to the diagonal
\(LCB_{3,u}\) obligation and discards every cross-position
cancellation.  Positivity belongs to the complete left side, not to
an individual off-diagonal term.
\end{firewall}

\begin{assessment}[Upstream direction after the no-go]
\label{ass:V67-direction}
The residual cannot be closed by relabelling the terminal capacity
cut, by attaching a scalar lower edge, or by estimating the prefix
or joining slot alone.  There are now exactly two legitimate routes:
\begin{enumerate}[label=\textup{(\roman*)},leftmargin=4em]
\item prove the conditional slice inequality
      \eqref{eq:V67-LCB} from the exact common local
      moment-partition incidence of \(M\) and a protected terminal
      factor \(N\); or
\item prove a stronger signed estimate directly on the aggregate
      square \eqref{eq:V67-aggregate-square}, where later spectators
      are no longer invisible.
\end{enumerate}
V68 should test the first route by splitting each protected local
block according to whether the invariant component containing the
terminal root also contains row \(3\).  The row-\(3\)-disconnected
part must then be returned to the aggregate rather than assigned a
false cut capacity.
\end{assessment}

\section{V67 result ledger and V68 acquisition order}
\label{sec:V67-status}

\begin{theorem}[Third-series V67 result ledger]
\label{thm:V67-results}
Relative to V1--V66, V67 proves:
\begin{enumerate}[label=\textup{(V67.\arabic*)},leftmargin=6.3em]
\item a sharp counterexample to transferring product-state capacity
      between the terminal lateral cut and the row-\(3\) cut;
\item the exact conditional-slice formula for positive
      row-\(3\) capacity;
\item the precise lateral conditional bridge card and its sufficiency
      for \(RCQ_3(\eta)\);
\item verification that every V66 cofinal terminal gap satisfies
      that bridge card;
\item slot invariance under later spectator extension;
\item no-go theorems excluding the isolated prefix and isolated
      local joining factors as NPRCs;
\item the exact signed aggregate square over the remaining
      first-connectivity positions; and
\item reduction of the lateral protected residual to either a
      conditional local-block slice estimate or an aggregate signed
      estimate.
\end{enumerate}
\end{theorem}

\begin{proof}
These are
\Cref{prop:V67-cut-transfer-no-go,
cor:V67-V57-cut-scope,
lem:V67-slice-capacity,
def:V67-LCB,
thm:V67-LCB-suffices,
cor:V67-gap-is-LCB,
lem:V67-slot-invariance,
prop:V67-prefix-NPRC-no-go,
prop:V67-joining-NPRC-no-go,
thm:V67-aggregate-square}.
\end{proof}

\begin{programme}[V68: protected-block incidence across row \(3\)]
\label{prog:V67-V68}
\begin{enumerate}[leftmargin=3em]
\item In one endpoint-expanded prefix or suffix moment state, write
      the exact local partition block containing the terminal root.
\item Split according to whether that block contains row \(3\).
      Preserve multiplicity spaces and the actual row-\(3\)
      representation.
\item On the row-\(3\)-connected part, compute the invariant
      projection's Schmidt bound across \(3|c\) and tensor it over
      the complete common bank before estimating dressings.
\item On the row-\(3\)-disconnected part, prove the exact tensor
      separation and test whether the row-\(3\) replacement
      difference annihilates it.  If it does not, retain it only
      inside \eqref{eq:V67-aggregate-square}.
\item Compare any obtained slice bound with the precise mass-ratio
      scale in \eqref{eq:V67-LCB}; do not infer a square from
      idempotence or from a different cut.
\item The next compulsory companion audit is V70.
\end{enumerate}
\end{programme}

\begin{firewall}[Claim boundary after third-series V67]
\label{fire:V67-claim-boundary}
V67 identifies the exact cut-aligned conditional bridge but does not
prove it for protected lateral terminal banks or close their signed
aggregate.  Finite exceptional and local-payer branches and full
\(RCQ_3(\eta)\) remain open.  The other quotient cuts, quartic
source families, higher MTA, WUFC, CPM, cutoff removal, reflection
positivity, Osterwalder--Schrader reconstruction, construction of
the continuum Yang--Mills measure, and a uniform positive continuum
spectral gap remain open.  The full Yang--Mills mass gap has not
been proved.
\end{firewall}

\paragraph{Parent-version record.}
V67 was formed from
\texttt{Renewal\_Yang\_Mills\_Mass\_Gap\_Research\_Notes\_V66.tex}.
The V66 parent had \(32\,501\) logical source lines,
\(1\,171\,436\) bytes, and SHA--256
\[
 \texttt{6c32338beb3c96faf36c03bc01c17e9c3dd81b0e02cb976fdf76d529af97b9d3}.
\]
The next compulsory companion audit is V70.

% =====================================================================
% V68 APPENDIX: RIGOROUS PROTECTED-BANK TENSORIZATION
% =====================================================================

\chapter{V68: Partial-Transpose Tensorization Repairs the
Protected-Bank Resolver}
\label{ch:V68-protected-tensorization}

\section{Upstream audit of the V32--V33 tensor step}
\label{sec:V68-upstream-audit}

V67 asks for a row-\(3\)-cut estimate on a protected bank.  Before
using the older protected-bank currency, one must audit its
tensorization.  A row density matrix may be entangled among all
coordinates in that row.  Therefore a one-coordinate product-state
capacity bound does not, merely by tensoring, imply the product of
the one-coordinate bounds.

The proof of
\Cref{lem:V33-abstract-one-failure} invokes tensor
submultiplicativity of product-state capacity for arbitrary
projections satisfying only
\(\kappa_\otimes(P_e)\le\rho\).  That inference is not among the
proved abstract hypotheses.  The physical Yang--Mills projectors
have additional structure: they are invariant projections for
independent coordinate copies of \(G\).  A partial-transpose
calculation supplies the stronger tensor-stable statement actually
needed.

\begin{firewall}[Correction of scope, not deletion of the physical
result]
\label{fire:V68-abstract-scope}
From V68 onward,
\Cref{lem:V33-abstract-one-failure} is not used as a theorem about
arbitrary positive projections with a scalar capacity bound.  Its
Yang--Mills applications are replaced by the invariant-projector
theorems below.  The local V32 invariant-capacity theorem remains
valid, and the physical V33 resolver estimates will be recovered
with all coordinate entanglement retained.
\end{firewall}

\section{Partial transpose of an invariant projection}
\label{sec:V68-partial-transpose}

Let \(V\) be a nontrivial irreducible \(G\)-module of dimension
\(d\), and let \(W\) be an arbitrary finite-dimensional unitary
\(G\)-module.  On the only contributing isotypic component write
\[
 W[V^*]\simeq V^*\otimes\mathcal M.
\]
Choose the basis used in
\Cref{thm:V32-invariant-capacity}, and let \(\Gamma_W\) denote matrix
transpose on the \(W\)-side of \(V\otimes W\).

\begin{lemma}[Dimension-sharp partial-transpose norm]
\label{lem:V68-invariant-PT}
For the invariant projection \(P_{\rm inv}\) onto
\((V\otimes W)^G\),
\begin{equation}
 \boxed{
 \|P_{\rm inv}^{\Gamma_W}\|=\frac1d}
\label{eq:V68-invariant-PT}
\end{equation}
when the invariant subspace is nonzero; if it is zero, the left side
is zero.  The equality is independent of the multiplicity space
\(\mathcal M\).
\end{lemma}

\begin{proof}
On \(V\otimes V^*\otimes\mathcal M\),
\[
 P_{\rm inv}
 =
 |\Omega\rangle\langle\Omega|
 \otimes I_{\mathcal M},
 \qquad
 \Omega=d^{-1/2}\sum_{j=1}^de_j\otimes e_j^*.
\]
It is zero on the orthogonal isotypic complement.  Partial transpose
on \(V^*\) gives
\[
 (|\Omega\rangle\langle\Omega|)^{\Gamma}
 =\frac1d\,\mathsf F,
\]
where \(\mathsf F\) is the flip between \(V\) and \(V^*\) after the
chosen conjugate-linear identification.  The flip is unitary and
self-adjoint.  Tensoring with \(I_{\mathcal M}\) and adding a zero
orthogonal block does not change the norm, proving
\eqref{eq:V68-invariant-PT}.
\end{proof}

\begin{corollary}[Ancilla-stable invariant capacity]
\label{cor:V68-ancilla-stable}
Let \(\mathcal A,\mathcal B\) be arbitrary finite-dimensional local
ancillas.  For all density matrices
\(\rho_{V\mathcal A}\) and \(\sigma_{W\mathcal B}\),
\begin{equation}
 \operatorname{Tr}\!\left[
 (P_{\rm inv}\otimes I_{\mathcal A\mathcal B})
 (\rho_{V\mathcal A}\otimes\sigma_{W\mathcal B})
 \right]
 \le\frac1d.
\label{eq:V68-ancilla-stable}
\end{equation}
\end{corollary}

\begin{proof}
Trace out the two ancillas.  The reduced matrices
\(\rho_V\) and \(\sigma_W\) remain density matrices and remain a
product across \(V|W\).  Apply
\Cref{thm:V32-invariant-capacity}.  Equivalently, transpose the
\(W\mathcal B\)-side and use
\Cref{lem:V68-invariant-PT}.
\end{proof}

\section{Exact product-group protected capacity}
\label{sec:V68-product-group}

Let \(E\) be a finite coordinate bank.  At each \(e\in E\), let
\(V_e\) be the active irreducible module in one fixed root row \(k\),
let \(d_e=\dim V_e\), let \(W_e\) contain the remaining row factors
and multiplicities, and let \(P_e\) be the local all-trivial
invariant projection.

\begin{theorem}[Protected projectors tensorize across entangled row
states]
\label{thm:V68-protected-product}
Across the grouped cut
\[
 \left(\bigotimes_{e\in E}V_e\right)
 \ \bigg|\ 
 \left(\bigotimes_{e\in E}W_e\right),
\]
one has
\begin{equation}
 \boxed{
 \kappa_{k|k^c}\!\left(\bigotimes_{e\in E}P_e\right)
 \le
 \prod_{e\in E}\frac1{d_e}
 \le\rho_G^{|E|}.}
\label{eq:V68-protected-product}
\end{equation}
The row density matrices may be arbitrarily entangled among the
coordinates on their own side of the cut.
\end{theorem}

\begin{proof}
Take partial transpose on the complete complementary-row side.
It commutes with coordinate tensor products, so
\[
 \left\|
 \left(\bigotimes_{e\in E}P_e\right)^\Gamma
 \right\|
 =
 \prod_{e\in E}\|P_e^\Gamma\|
 \le
 \prod_{e\in E}\frac1{d_e}
\]
by \Cref{lem:V68-invariant-PT}.  For arbitrary density matrices
\(\rho_k,\sigma_{k^c}\), including coordinate-entangled ones,
\begin{align*}
 &\operatorname{Tr}\!\left[
 \left(\bigotimes_eP_e\right)
 (\rho_k\otimes\sigma_{k^c})\right]\\
 &\qquad=
 \operatorname{Tr}\!\left[
 \left(\bigotimes_eP_e\right)^\Gamma
 (\rho_k\otimes\sigma_{k^c}^{\,T})\right]
 \le
 \left\|
 \left(\bigotimes_eP_e\right)^\Gamma
 \right\|.
\end{align*}
The transposed density matrix is positive with trace one.  Taking
the supremum proves the first inequality.  The second follows from
\(d_e\ge d_{\min,G}\) and
\(\rho_G=d_{\min,G}^{-1}\).
\end{proof}

\begin{remark}[Why this proof is stronger than a scalar capacity
iteration]
\label{rem:V68-not-scalar-iteration}
The proof never conditions on one coordinate and then asserts that
the remaining row state is product.  It evaluates the complete
coordinate product at once through the norm of its partial
transpose.  This is precisely what permits entanglement within each
row.
\end{remark}

\section{The physical protected-plus-gapped heat product}
\label{sec:V68-heat-product}

For \(0\le r\le1\), put
\begin{equation}
 B_e(r):=P_e+r(I-P_e),
 \qquad
 A_\rho(r):=\rho+(1-\rho)r.
\label{eq:V68-isotropic-majorant}
\end{equation}

\begin{lemma}[Tensor-stable heat majorant]
\label{lem:V68-heat-majorant}
For every coordinate-entangled product density across \(k|k^c\),
\begin{equation}
 \boxed{
 \kappa_{k|k^c}\!\left(
 \bigotimes_{e\in E}B_e(r)\right)
 \le A_{\rho_G}(r)^{|E|}.}
\label{eq:V68-heat-majorant}
\end{equation}
\end{lemma}

\begin{proof}
By \Cref{lem:V68-invariant-PT},
\[
 \|B_e(r)^\Gamma\|
 \le
 r+(1-r)\|P_e^\Gamma\|
 \le
 r+(1-r)\rho_G
 =
 A_{\rho_G}(r).
\]
Partial transpose of the coordinate product is the tensor product of
the partial transposes.  Repeat the trace calculation in
\Cref{thm:V68-protected-product} and multiply their ordinary
operator norms.
\end{proof}

The physical local heat block has the exact orthogonal spectral
form
\begin{equation}
 M_e(t)=P_e+G_e(t),
\qquad
 P_eG_e(t)=0,
\qquad
 0\preccurlyeq G_e(t)
 \preccurlyeq r_t(I-P_e),
\label{eq:V68-physical-heat-split}
\end{equation}
where \(r_t=(1-\delta)^t\) and
\(\delta=\gamma_Gs_\alpha\).  Put
\[
 M_E(t)=\bigotimes_{e\in E}M_e(t),
 \qquad
 P_*=\bigotimes_{e\in E}P_e.
\]

\begin{theorem}[Entanglement-safe one-failure bound]
\label{thm:V68-one-failure}
For \(N=|E|\ge1\),
\begin{equation}
 \boxed{
 \kappa_{k|k^c}\bigl(M_E(t)-P_*\bigr)
 \le
 \Psi_{N,\rho_G}(r_t)
 :=
 \min\!\left\{
 r_t,\,
 A_{\rho_G}(r_t)^N
 \right\}.}
\label{eq:V68-one-failure}
\end{equation}
\end{theorem}

\begin{proof}
The orthogonal decomposition
\eqref{eq:V68-physical-heat-split} shows that
\(M_E(t)-P_*\succcurlyeq0\).  Every nonprotected joint spectral
sector contains at least one \(G_e(t)\); all other factors are
contractions.  Hence
\[
 \|M_E(t)-P_*\|\le r_t.
\]
Also \(M_e(t)\preccurlyeq B_e(r_t)\), so tensor positivity gives
\[
 M_E(t)\preccurlyeq\bigotimes_eB_e(r_t).
\]
Since
\(0\preccurlyeq M_E(t)-P_*\preccurlyeq M_E(t)\), monotonicity and
\Cref{lem:V68-heat-majorant} give
\[
 \kappa_{k|k^c}(M_E(t)-P_*)
 \le A_{\rho_G}(r_t)^N.
\]
Combine the two upper bounds.
\end{proof}

\begin{firewall}[The exact spectral split is essential]
\label{fire:V68-physical-split}
The lower relation \(M_e(t)\succcurlyeq P_e\), the orthogonality
\(P_eG_e(t)=0\), and the heat bound on \(G_e(t)\) are physical
spectral facts.  They are stronger than the upper majorant alone in
the old abstract V33 statement.  They are what makes
\(M_E(t)-P_*\) positive and gives the ordinary norm \(r_t\).
\end{firewall}

\section{Power sums for the repaired majorant}
\label{sec:V68-power-sums}

\begin{lemma}[Geometric power sum]
\label{lem:V68-Psi-sum}
Let \(0<\rho<1\), \(0<\delta\le\tfrac12\), and
\(r_t=(1-\delta)^t\).  There are constants \(C_\rho,c_\rho>0\)
such that for every \(N\ge1\),
\begin{align}
 \sum_{n\ge1}\Psi_{N,\rho}(r_n)
 &\le
 C_\rho e^{-c_\rho\delta N}
 \left(1+\frac1{\delta N}\right),
\label{eq:V68-integer-Psi}\\
 \sum_{n\ge0}\Psi_{N,\rho}(r_{n+1/2})
 &\le
 C_\rho e^{-c_\rho\delta N}
 \left(1+\frac1{\delta N}\right).
\label{eq:V68-half-Psi}
\end{align}
\end{lemma}

\begin{proof}
Put \(\lambda=-\log(1-\delta)\), so
\(\delta\le\lambda\le2\delta\), and define
\[
 f_N(r):=\min\{r,[\rho+(1-\rho)r]^N\}.
\]
This is increasing.  Geometric interval comparison gives
\[
 \sum_{n\ge1}f_N(e^{-\lambda n})
 \le
 f_N(e^{-\lambda})
 +\lambda^{-1}
 \int_0^{e^{-\lambda}}\frac{f_N(r)}r\,dr.
\label{eq:V68-geometric-comparison}
\]
Let \(r_1=e^{-\lambda}=1-\delta\).  Since
\[
 \rho+(1-\rho)r
 \le e^{-(1-\rho)(1-r)},
\]
the endpoint term is at most
\(e^{-(1-\rho)\delta N}\).

For the integral use
\[
 f_N(r)
 \le
 r^{1/2}[\rho+(1-\rho)r]^{N/2}.
\]
On \(r_1/2\le r\le r_1\), one has \(r^{-1/2}\le2\) and
\[
 1-r\ge\delta+(r_1-r).
\]
Integration therefore gives
\[
 \int_{r_1/2}^{r_1}\frac{f_N(r)}r\,dr
 \le
 \frac{C_\rho}{N}e^{-c_\rho\delta N}.
\]
On \(0<r<r_1/2\), the extra lower bound
\(1-r\ge\delta+r_1/2\), together with
\(\int_0^{r_1/2}r^{-1/2}\,dr<\infty\), gives the stronger estimate
\(C_\rho e^{-c_\rho N}\).  Insert both estimates into
\eqref{eq:V68-geometric-comparison} and use
\(\lambda\asymp\delta\).  This proves
\eqref{eq:V68-integer-Psi}, after reducing \(c_\rho\) and enlarging
\(C_\rho\).

For half powers start the same comparison at
\(r_{1/2}=(1-\delta)^{1/2}\).  Since
\[
 1-(1-\delta)^{1/2}\ge\delta/2,
\]
the proof is identical with changed constants.
\end{proof}

\begin{corollary}[Repaired physical V33 resolvers]
\label{cor:V68-physical-resolvers}
For the actual invariant projectors and heat blocks
\eqref{eq:V68-physical-heat-split}, the geometric and
Casimir-weighted complete-bank resolvers obey
\begin{align}
 \kappa_{k|k^c}(\mathscr G_E)
 &\le
 C_Ge^{-c_Gs_\alpha N}
 \left(1+\frac1{s_\alpha N}\right),
\label{eq:V68-geometric-resolver}\\
 \kappa_{k|k^c}(\mathscr C_E)
 &\le
 \frac{C_G}{s_\alpha}e^{-c_Gs_\alpha N}.
\label{eq:V68-Casimir-resolver}
\end{align}
These are the physical conclusions
\eqref{eq:V33-geometric-resolver}--%
\eqref{eq:V33-Casimir-resolver}, now proved without an abstract
capacity-product assumption.
\end{corollary}

\begin{proof}
For the geometric resolver use
\[
 \frac{Q}{1-Q}=\sum_{n\ge1}Q^n.
\]
On the non-all-protected sector, the \(n\)-th complete heat power is
\(M_E(n)-P_*\).  Apply
\Cref{thm:V68-one-failure,lem:V68-Psi-sum} and subadditivity.

For the Casimir resolver, the scalar half-heat estimate
\eqref{eq:V33-half-heat} bounds its block coefficient by
\(C_Gs_\alpha^{-1}\) times the half-powered heat block.  The
all-protected sector has zero Casimir and is absent.  Apply
\eqref{eq:V68-half-Psi} and absorb its exponentially smaller
polynomial factor into the exponential after changing constants.
\end{proof}

\begin{corollary}[The downstream physical bank currencies are
restored]
\label{cor:V68-downstream-restored}
The protected and gapped physical bank currencies in
\Cref{cor:V32-V70-currencies}, the paid square tail in
\Cref{thm:V33-paid-square-tail}, and their uses in V57--V67 remain
valid for the actual Yang--Mills invariant/heat sectors, with
possibly changed group constants.
\end{corollary}

\begin{proof}
The protected currency uses
\(\rho_G^N\), now supplied by
\Cref{thm:V68-protected-product}.  The non-all-protected geometric
and Casimir resolvers are
\Cref{cor:V68-physical-resolvers}.  The row-power inequality,
coordinate-additive payer bound, and scalar comparisons used in
V32--V33 are unchanged.  Repeating those displayed downstream
proofs with \eqref{eq:V68-geometric-resolver}--%
\eqref{eq:V68-Casimir-resolver} gives the same conclusions.
\end{proof}

\begin{assessment}[Effect on the V67 bridge]
\label{ass:V68-bridge-effect}
The older bank estimates needed by the V67 conditional bridge now
have a tensorization proof which permits arbitrary entanglement
inside each row.  This repair does not itself prove
\(LCB_{3,u}\): the terminal protected currency may still be rooted
at a lateral cut.  It does ensure that the next local-block
incidence analysis starts from valid protected and one-failure
currencies rather than an unproved generic crossnorm rule.
\end{assessment}

\section{V68 result ledger and V69 acquisition order}
\label{sec:V68-status}

\begin{theorem}[Third-series V68 result ledger]
\label{thm:V68-results}
Relative to V1--V67, V68 proves:
\begin{enumerate}[label=\textup{(V68.\arabic*)},leftmargin=6.3em]
\item identification of the missing hypothesis in the abstract V33
      tensor-capacity argument;
\item the exact partial-transpose norm \(1/\dim V\) of an invariant
      projection with arbitrary multiplicity;
\item ancilla stability and full coordinate tensorization for the
      physical all-protected sector;
\item a partial-transpose heat-majorant estimate valid for
      coordinate-entangled row states;
\item an entanglement-safe one-failure bound combining ordinary heat
      norm and protected capacity;
\item exponentially localized integer and half-integer power sums
      for the repaired majorant;
\item recovery of the physical geometric and Casimir resolvers; and
\item restoration of the downstream V32--V33 and V57--V67 physical
      bank currencies with corrected proof.
\end{enumerate}
\end{theorem}

\begin{proof}
These are
\Cref{fire:V68-abstract-scope,
lem:V68-invariant-PT,
cor:V68-ancilla-stable,
thm:V68-protected-product,
lem:V68-heat-majorant,
thm:V68-one-failure,
lem:V68-Psi-sum,
cor:V68-physical-resolvers,
cor:V68-downstream-restored}.
\end{proof}

\begin{programme}[V69: row-\(3\) incidence inside a protected
terminal bank]
\label{prog:V68-V69}
\begin{enumerate}[leftmargin=3em]
\item In every local moment-partition state of the common original
      and terminal banks, record the block containing row \(3\) and
      whether the row-\(3\) representation is active.
\item Apply \Cref{lem:V68-invariant-PT} across the actual row-\(3\)
      cut, with all other bank coordinates treated simultaneously
      by \Cref{thm:V68-protected-product}.
\item Split the terminal bank into row-\(3\)-active and
      row-\(3\)-inactive atoms before comparing with
      \eqref{eq:V67-LCB}.
\item On an active atom, retain the exact row-\(3\) representation
      dimension rather than replacing it immediately by
      \(\rho_G\).  On an inactive atom, record the exact tensor
      separation and test the aggregate square.
\item Prove a cut-aligned mass-ratio bound or isolate a finite
      high-representation concentration cell where Weyl-dimension
      growth is sharp.
\item The next compulsory companion audit is V70.
\end{enumerate}
\end{programme}

\begin{firewall}[Claim boundary after third-series V68]
\label{fire:V68-claim-boundary}
V68 repairs the protected/gapped tensor resolver for physical
Yang--Mills invariant sectors.  It does not prove the lateral
conditional bridge, close protected terminal-owned, finite
exceptional, or local-payer branches, or establish full
\(RCQ_3(\eta)\).  The other quotient cuts, quartic source families,
higher MTA, WUFC, CPM, cutoff removal, reflection positivity,
Osterwalder--Schrader reconstruction, construction of the continuum
Yang--Mills measure, and a uniform positive continuum spectral gap
remain open.  The full Yang--Mills mass gap has not been proved.
\end{firewall}

\paragraph{Parent-version record.}
V68 was formed from
\texttt{Renewal\_Yang\_Mills\_Mass\_Gap\_Research\_Notes\_V67.tex}.
The V67 parent had \(32\,978\) logical source lines,
\(1\,187\,747\) bytes, and SHA--256
\[
 \texttt{fe7c5f78ffea0e0c7b754ff6093fd51a082272b57abcc231ea05702fabcab4a0}.
\]
The next compulsory companion audit is V70.

% =====================================================================
% V69 APPENDIX: CUT-ALIGNED PROTECTED MASS
% =====================================================================

\chapter{V69: Cut-Aligned Protected Mass, Weyl-Dimension Decay, and
a New Closed Lateral Cell}
\label{ch:V69-protected-mass}

\section{The row-\(3\) block in a local protected state}
\label{sec:V69-local-block}

Fix one endpoint-expanded local moment-partition state at a bank
coordinate \(e\).  Its row partition is \(\pi_e\), and its
all-trivial output projection factors as
\begin{equation}
 P_e^{\mathrm{prot}}
 =
 \bigotimes_{B\in\pi_e}P_{B,e}^{\mathrm{inv}},
\label{eq:V69-local-protected}
\end{equation}
where \(P_{B,e}^{\mathrm{inv}}\) projects
\(\bigotimes_{i\in B}V_{\lambda_{i,e}}\) onto its invariant
subspace.  Inactive rows carry the trivial representation.

\begin{lemma}[Exact row-\(3\) partial-transpose norm]
\label{lem:V69-local-row3-PT}
If row \(3\) is active at \(e\), then
\begin{equation}
 \boxed{
 \left\|
 (P_e^{\mathrm{prot}})^{\Gamma_{3^c}}
 \right\|
 \le
 \frac1{d_{\lambda_{3,e}}}.}
\label{eq:V69-local-row3-PT}
\end{equation}
If row \(3\) is inactive, the corresponding general upper bound is
one.  The conclusion is independent of whether the partition block
containing row \(3\) also contains the terminal lateral root.
\end{lemma}

\begin{proof}
Let \(B_3\in\pi_e\) be the block containing row \(3\).  Regroup its
module as
\[
 V_{\lambda_{3,e}}\otimes
 \bigotimes_{i\in B_3\setminus\{3\}}V_{\lambda_{i,e}}.
\]
Apply \Cref{lem:V68-invariant-PT} to its invariant projection,
treating any dual multiplicity as part of the complementary factor.
This gives norm at most \(d_{\lambda_{3,e}}^{-1}\) after partial
transpose on the non-row-\(3\) factors of \(B_3\).

Every other block \(B\ne B_3\) lies wholly on the complementary
side.  Full transpose preserves the operator norm of its orthogonal
projection, which is one when the block is nonzero.  Multiplying the
block norms in \eqref{eq:V69-local-protected} proves
\eqref{eq:V69-local-row3-PT}.  If row \(3\) is inactive, its factor
is one-dimensional and the remaining projection can have norm one;
no strict bound follows.
\end{proof}

\begin{assessment}[The correct incidence split]
\label{ass:V69-incidence-split}
The relevant dichotomy is row-\(3\) activity, not merely whether row
\(3\) shares the terminal root's partition block.  Even in a
different invariant block, an active row-\(3\) representation gives
the cut-aligned factor
\eqref{eq:V69-local-row3-PT}.  An inactive row gives none.
\end{assessment}

\section{Weyl dimension versus Casimir mass}
\label{sec:V69-Weyl}

\begin{lemma}[Uniform dimension--Casimir lower bound]
\label{lem:V69-Weyl-dimension}
For a fixed compact connected semisimple group \(G\), there is
\(c_G>0\) such that every nontrivial irreducible highest-weight
module \(V_\lambda\) satisfies
\begin{equation}
 \boxed{
 \dim V_\lambda
 \ge
 c_G\sqrt{1+C_2(\lambda)}.}
\label{eq:V69-Weyl-dimension}
\end{equation}
\end{lemma}

\begin{proof}
It suffices to treat one simple factor; for a semisimple product,
dimensions multiply and the Casimir is a fixed positive sum over
the finitely many factors.

Write the dominant integral highest weight as
\(\lambda=\sum_{i=1}^{r_G}n_i\omega_i\), \(n_i\in\mathbb N_0\), and
put \(n_*=\max_i n_i\).  Weyl's dimension formula contains, for a
simple root \(\alpha_i\) with \(n_i=n_*\), the factor
\[
 \frac{\langle\lambda+\rho,\alpha_i^\vee\rangle}
      {\langle\rho,\alpha_i^\vee\rangle}
 =n_*+1.
\]
Every other positive-root factor is at least one.  Hence
\(\dim V_\lambda\ge n_*+1\).

Because the fundamental weights form a fixed finite basis,
\(\|\lambda\|\le C_G n_*\).  The Casimir identity
\[
 C_2(\lambda)=\langle\lambda,\lambda+2\rho\rangle
\]
and norm equivalence therefore give
\[
 1+C_2(\lambda)\le C_G(1+n_*)^2.
\]
Combine the two inequalities and adjust the group constant.
\end{proof}

\begin{corollary}[One active protected coordinate]
\label{cor:V69-one-coordinate-decay}
At a row-\(3\)-active coordinate,
\begin{equation}
 \left\|
 (P_e^{\mathrm{prot}})^{\Gamma_{3^c}}
 \right\|
 \le
 \min\left\{
 \rho_G,\,
 C_Ga_{3,e}^{-1/2}
 \right\},
\qquad
a_{3,e}=1+C_2(\lambda_{3,e}).
\label{eq:V69-one-coordinate-decay}
\end{equation}
\end{corollary}

\begin{proof}
The inverse-dimension estimate is
\Cref{lem:V69-local-row3-PT}.  Use
\(d_{\lambda_{3,e}}\ge d_{\min,G}=\rho_G^{-1}\) and
\Cref{lem:V69-Weyl-dimension}.
\end{proof}

\section{A complete row-\(3\)-active protected subbank}
\label{sec:V69-active-bank}

Let \(A\) be a nonempty complete bank on which row \(3\) is active
at every coordinate, and put
\begin{equation}
 S_3(A):=\sum_{e\in A}a_{3,e}.
\label{eq:V69-active-mass}
\end{equation}
Let \(P_A^{\mathrm{prot}}\) be the tensor product of the local
all-trivial projections in one fixed moment-partition state.

\begin{theorem}[Universal half-power protected-mass decay]
\label{thm:V69-half-power}
There is \(C_G<\infty\) such that
\begin{equation}
 \boxed{
 \left\|
 (P_A^{\mathrm{prot}})^{\Gamma_{3^c}}
 \right\|
 \le
 C_GS_3(A)^{-1/2}.}
\label{eq:V69-half-power}
\end{equation}
Consequently the same bound holds for
\(\kappa_{3|c}(P_A^{\mathrm{prot}})\).
\end{theorem}

\begin{proof}
Let \(N=|A|\) and choose \(e_*\) with maximal \(a_{3,e}\).  By
\Cref{cor:V69-one-coordinate-decay} and tensor
multiplicativity of ordinary norm after partial transpose,
\[
 \left\|
 (P_A^{\mathrm{prot}})^{\Gamma_{3^c}}
 \right\|
 \le
 C_Ga_{3,e_*}^{-1/2}\rho_G^{N-1}.
\]
Since \(a_{3,e_*}\ge S_3(A)/N\), the right side is at most
\[
 C_G
 \rho_G^{N-1}\sqrt N\,S_3(A)^{-1/2}.
\]
The scalar sequence
\(\rho_G^{N-1}\sqrt N\) is bounded for \(N\ge1\), proving
\eqref{eq:V69-half-power}.  Capacity is bounded by the partial-
transpose norm by the trace calculation in
\Cref{thm:V68-protected-product}.
\end{proof}

\begin{proposition}[The half power is sharp for \(SU(2)\)]
\label{prop:V69-half-power-sharp}
No universal estimate
\[
 \kappa_{3|c}(P_A^{\mathrm{prot}})
 \le C S_3(A)^{-\gamma}
\]
with \(\gamma>1/2\) follows for all protected banks from
representation dimension alone, already for \(G=SU(2)\).
\end{proposition}

\begin{proof}
Take a one-coordinate bank and put spin \(j\) in row \(3\), with a
dual spin \(j\) on the complementary side.  The invariant projection
has capacity
\[
 (2j+1)^{-1}
\]
by \Cref{thm:V32-invariant-capacity}, and this value is attained by a
basis product vector.  Meanwhile
\[
 S_3(A)=1+j(j+1)\asymp j^2.
\]
Thus the capacity is comparable to \(S_3(A)^{-1/2}\), and any
\(\gamma>1/2\) fails as \(j\to\infty\).
\end{proof}

\begin{firewall}[Large representation is not arbitrary polynomial
heat]
\label{fire:V69-no-arbitrary-power}
The cofinal unbalanced heat sector of V64 pays every fixed mass
power because it has ordinary exponential norm decay in Casimir.
A protected invariant projection has no such heat norm: its
universal dimension decay is only the sharp half power of
\Cref{prop:V69-half-power-sharp}.  The two mechanisms must not be
interchanged.
\end{firewall}

\section{Partial-transpose currency of the original protected bank}
\label{sec:V69-original-bank}

Return to the original root-owned comparable row-\(3\)--row-\(u\)
bank \(D_0\).  Put
\[
 N_0=|D_0|,
 \qquad
 q_0=\frac{H_0}{S_0^2},
 \qquad
 \beta_0=\frac{T_0}{\Xi_u}.
\]

\begin{lemma}[The V60 protected currency holds in partial-transpose
norm]
\label{lem:V69-original-PT}
For the all-protected envelope \(M_0\) on \(D_0\),
\begin{equation}
 \boxed{
 \|M_0^{\Gamma_{3^c}}\|
 \le C_{G,c_0}q_0^2.}
\label{eq:V69-original-PT}
\end{equation}
\end{lemma}

\begin{proof}
\Cref{thm:V68-protected-product} gives the stronger intermediate
quantity
\[
 \|M_0^{\Gamma_{3^c}}\|
 \le\rho_G^{N_0}.
\]
Pointwise comparability gives
\(q_0\ge c_0/N_0\) by
\eqref{eq:V60-q-range}.  Since
\(\sup_{N\ge1}N^2\rho_G^N<\infty\),
\[
 \rho_G^{N_0}
 \le C_GN_0^{-2}
 \le C_{G,c_0}q_0^2.
\]
\end{proof}

\section{A new closed disjoint protected cell}
\label{sec:V69-disjoint-closure}

Let \(A\) be the row-\(3\)-active part of a terminal lateral
all-protected bank.  Suppose \(A\cap D_0=\varnothing\), and denote
its protected factor by \(N_A\).  Factors on terminal coordinates
where row \(3\) is inactive may be included in a contraction
\(N_{\rm in}\); after complementary-side transpose their norm is at
most one.

\begin{theorem}[Disjoint protected-mass bridge]
\label{thm:V69-disjoint-bridge}
Assume
\begin{equation}
 S_3(A)
 \ge
 L\left(\frac{\Xi_u}{\Xi_4}\right)^4
\label{eq:V69-mass-threshold}
\end{equation}
for a fixed sufficiently large group-dependent \(L\).  If a
restriction-resolved outside-local word has its original protected
factor \(M_0\) on \(D_0\) and terminal protected factor
\(N_A N_{\rm in}\) on the disjoint terminal bank, then
\begin{equation}
 \boxed{
 \kappa_{3|c}(
 \widetilde W_{\pi,\nu}^*\widetilde W_{\pi,\nu})
 \le
 C_s\theta_{3u}(D_0)^2
 \le C_s(x+z)^2.}
\label{eq:V69-disjoint-RCQ}
\end{equation}
\end{theorem}

\begin{proof}
By \Cref{thm:V69-half-power} and
\eqref{eq:V69-mass-threshold},
\begin{equation}
 \|(N_AN_{\rm in})^{\Gamma_{3^c}}\|
 \le
 C_GS_3(A)^{-1/2}
 \le
 C_{G,L}\left(\frac{\Xi_4}{\Xi_u}\right)^2.
\label{eq:V69-terminal-PT-ratio}
\end{equation}
Because \(A\) and \(D_0\) are disjoint coordinate banks, partial
transpose and ordinary norm tensorize between their factors.
\Cref{lem:V69-original-PT} therefore gives
\[
 \|
 (M_0N_AN_{\rm in})^{\Gamma_{3^c}}
 \|
 \le
 C_s q_0^2
 \left(\frac{\Xi_4}{\Xi_u}\right)^2.
\]
The factors are projections and commute with the resolved
dressings by V59 and V66.  Repeating the central-square inequality
in \Cref{thm:V65-terminal-gap-insertion}, but bounding final
row-\(3\) capacity by partial-transpose norm, yields
\begin{align*}
 \kappa_{3|c}(
 \widetilde W_{\pi,\nu}^*\widetilde W_{\pi,\nu})
 &\le
 C_s
 \left(\frac{p}{\Xi_3\Xi_4}\right)^2
 q_0^2
 \left(\frac{\Xi_4}{\Xi_u}\right)^2\\
 &\le
 C_sa_*^{-4}q_0^2\beta_0^2.
\end{align*}
For the last line use \(p\le1\), \(\Xi_3\ge a_*\), and
\(T_0\ge a_*\).  The lower side of the exact factorization
\eqref{eq:V60-theta-factorization}, together with root ownership
\(\alpha_0\ge\sigma\), gives
\[
 q_0^2\beta_0^2\le C_s\theta_{3u}(D_0)^2.
\]
Finally \(\theta_{3u}(D_0)\le x+z\).
\end{proof}

\begin{corollary}[The disjoint diffuse protected cell leaves the
counterpacket list]
\label{cor:V69-disjoint-cell}
After fixed finite simultaneous resolution, every term satisfying
\eqref{eq:V69-mass-threshold} and the disjointness hypothesis of
\Cref{thm:V69-disjoint-bridge} closes under recombination.
\end{corollary}

\begin{proof}
Apply \Cref{thm:V69-disjoint-bridge} term by term and use
\Cref{lem:V58-finite-recombination}.  The Boolean-atom and endpoint
term count is fixed by \Cref{thm:V66-chain-incidence}.
\end{proof}

\begin{assessment}[Exact residual before the V70 audit]
\label{ass:V69-residual}
The protected lateral terminal branch has gained a genuinely
cut-aligned closed cell.  The unresolved protected terms satisfy at
least one of:
\begin{enumerate}[label=\textup{(\roman*)},leftmargin=4em]
\item the terminal bank's row-\(3\)-active part overlaps the original
      bank \(D_0\);
\item its active mass is below the fourth-power threshold
      \eqref{eq:V69-mass-threshold};
\item row \(3\) is inactive on the terminal protected mass; or
\item the original factor lies in the non-all-protected gapped
      sector rather than \(M_0\).
\end{enumerate}
The \(SU(2)\) one-coordinate example proves that representation
dimension alone cannot remove item \textup{(ii)}.  These branches
must be tested in the exact aggregate square or by an additional
signed response.
\end{assessment}

\section{V69 result ledger and V70 audit mandate}
\label{sec:V69-status}

\begin{theorem}[Third-series V69 result ledger]
\label{thm:V69-results}
Relative to V1--V68, V69 proves:
\begin{enumerate}[label=\textup{(V69.\arabic*)},leftmargin=6.3em]
\item the exact row-\(3\) partial-transpose norm of an arbitrary
      local all-protected moment-partition state;
\item correction of the incidence split from same-block incidence to
      row-\(3\) activity;
\item a group-uniform Weyl dimension--Casimir lower bound;
\item universal \(S_3(A)^{-1/2}\) decay of a complete
      row-\(3\)-active protected bank;
\item sharpness of that exponent for one-coordinate \(SU(2)\)
      banks;
\item the original V60 protected currency in partial-transpose norm;
\item a disjoint protected-mass bridge at the exact fourth-power
      threshold; and
\item \(RCQ_3\) closure and finite recombination of that new lateral
      cell.
\end{enumerate}
\end{theorem}

\begin{proof}
These are
\Cref{lem:V69-local-row3-PT,
ass:V69-incidence-split,
lem:V69-Weyl-dimension,
cor:V69-one-coordinate-decay,
thm:V69-half-power,
prop:V69-half-power-sharp,
lem:V69-original-PT,
thm:V69-disjoint-bridge,
cor:V69-disjoint-cell}.
\end{proof}

\begin{programme}[V70: compulsory audit and concentrated protected
residual]
\label{prog:V69-V70}
\begin{enumerate}[leftmargin=3em]
\item Perform the compulsory read-only audit of the newest active SM
      and NS notes and record their exact metadata.
\item Import only proof-order information relevant to overlapping or
      concentrated protected banks.
\item Partition the row-\(3\)-active terminal mass into its overlap
      with \(D_0\) and its disjoint complement.  Route the latter
      through \Cref{thm:V69-disjoint-bridge} whenever the threshold
      holds.
\item On the overlap, compute the partial transpose of the product
      of the two actual commuting invariant projections; do not
      multiply their separate capacities.
\item On the subthreshold or row-\(3\)-inactive cell, restore the
      signed aggregate \eqref{eq:V67-aggregate-square} and test its
      first cross-position cancellation.
\item Keep original gapped, finite exceptional, and local-payer
      branches separate.
\end{enumerate}
\end{programme}

\begin{firewall}[Claim boundary after third-series V69]
\label{fire:V69-claim-boundary}
V69 closes the disjoint high-row-\(3\)-mass protected lateral cell.
It does not close overlapping, subthreshold, row-\(3\)-inactive, or
original-gapped lateral protected cells, finite exceptional or
local-payer branches, or full \(RCQ_3(\eta)\).  The other quotient
cuts, quartic source families, higher MTA, WUFC, CPM, cutoff removal,
reflection positivity, Osterwalder--Schrader reconstruction,
construction of the continuum Yang--Mills measure, and a uniform
positive continuum spectral gap remain open.  The full Yang--Mills
mass gap has not been proved.
\end{firewall}

\paragraph{Parent-version record.}
V69 was formed from
\texttt{Renewal\_Yang\_Mills\_Mass\_Gap\_Research\_Notes\_V68.tex}.
The V68 parent had \(33\,510\) logical source lines,
\(1\,204\,561\) bytes, and SHA--256
\[
 \texttt{c1df5db6079318eb5a5a20f3ee28304d70d7ea71a1e151cdf3e71e7bfc4cb6f1}.
\]
The compulsory companion audit is due in V70.

% =====================================================================
% V70 APPENDIX: AUDIT AND OVERLAPPING PROTECTED BANKS
% =====================================================================

\chapter{V70: Companion Audit, Overlap Idempotence, and the
New-Mass Protected Bridge}
\label{ch:V70-overlap}

\section{Compulsory read-only companion audit}
\label{sec:V70-audit}

The scheduled audit was performed before choosing the V70
mathematical direction.  The completed reads were:
\begin{enumerate}[leftmargin=3em]
\item
\texttt{Renewal\_SM\_Einstein\_Continuum\_Research\_Notes\_V52.tex},
with \(19\,005\) logical source lines and \(698\,889\) bytes,
SHA--256
\[
 \texttt{bd14be468074dd88f1d5f4c1b31d04d80c36c7ea97561b5de016cbfdf7be498a}.
\]
The parallel SM programme advanced through a transient V51 while
the audit was under way; V52 was the file actually read to
completion.  Its new sequence first relocates exact conditional
cumulants to collision components, then groups a complete product
before applying a covariance across a cut.  It proves that
factorial local merger weights can be summed without duplication,
but also proves that one \(L^2\) contraction cannot manufacture the
\(k-1\) spatial factors of a simultaneous \(k\)-way merger.
\item
\texttt{Renewal\_NS\_Stability\_Research\_Notes\_V31.tex},
with \(23\,052\) logical source lines and \(887\,537\) bytes,
SHA--256
\[
 \texttt{f1121b62238ad5491bb7cf03635ea9934942edf573ed8faaa9ee79e1d38a6696}.
\]
It is unchanged.  Its persistent branch still carries the explicit
no-double-charge rule, while its heat-formation ladder retains the
complete stopped quantity until the actual signed incidence is
assembled.
\end{enumerate}

\begin{assessment}[Type-correct V70 transfer]
\label{ass:V70-audit-transfer}
No SM conditional-cumulant estimate and no NS heat estimate is a
Yang--Mills operator inequality.  The applicable proof-order lesson
is exact incidence with one charge: when two bank labels select the
same local protected output, their product is one idempotent
projection.  Only protected coordinates genuinely outside the
original bank can supply a new tensor factor.  V70 proves that
statement directly in the YM moment-partition algebra.
\end{assessment}

\section{One all-trivial projector per local state}
\label{sec:V70-local-idempotence}

Fix one endpoint-expanded moment-partition state at coordinate \(e\).
The all-trivial output projection \(P_e^{\mathrm{prot}}\) is
\eqref{eq:V69-local-protected}.  It is determined by that local
state, not by the root/partner labels subsequently used to estimate
it.

\begin{lemma}[Bank labels do not duplicate the local protected
projector]
\label{lem:V70-same-local-projector}
Suppose two retained bank labels both contain \(e\) and both select
the all-trivial sector of the same resolved local
moment-partition state.  Their two local factors are equal to
\(P_e^{\mathrm{prot}}\), and hence
\begin{equation}
 P_e^{\mathrm{prot}}P_e^{\mathrm{prot}}
 =
 P_e^{\mathrm{prot}}.
\label{eq:V70-local-idempotence}
\end{equation}
This remains true if the banks were produced at different steps of
the scalar mass-escalation chain.
\end{lemma}

\begin{proof}
For the fixed state, the protected condition says that every output
block of the partition is trivial.  The corresponding spectral
projection is uniquely
\[
 \bigotimes_{B\in\pi_e}P_{B,e}^{\mathrm{inv}}.
\]
Changing the chosen root, comparison partner, or the step at which
the coordinate was recorded does not change \(\pi_e\) or this
spectral event.  Both factors are therefore the same orthogonal
projection, whose square is itself.
\end{proof}

\begin{theorem}[Exact union law for overlapping protected banks]
\label{thm:V70-protected-union}
Let \(D,F\) be two complete banks in one common resolved
moment-partition state, and let
\[
 P_D:=\bigotimes_{e\in D}P_e^{\mathrm{prot}},
 \qquad
 P_F:=\bigotimes_{e\in F}P_e^{\mathrm{prot}},
\]
with identities off the indicated supports.  Then
\begin{equation}
 \boxed{
 P_DP_F=P_FP_D=P_{D\cup F}.}
\label{eq:V70-protected-union}
\end{equation}
In particular, the overlap \(D\cap F\) is counted once.
\end{theorem}

\begin{proof}
All local spectral projections commute.  On \(D\cap F\), apply
\eqref{eq:V70-local-idempotence}; on the symmetric difference only
one projection occurs; off the union both factors are identities.
Tensoring these coordinatewise identities gives
\eqref{eq:V70-protected-union}.
\end{proof}

\begin{corollary}[No second protected currency on a contained bank]
\label{cor:V70-contained-no-gain}
If \(F\subseteq D\), then \(P_DP_F=P_D\).  Thus a second bank label
supported inside an already protected bank supplies no additional
capacity or partial-transpose factor.
\end{corollary}

\begin{proof}
Put \(D\cup F=D\) in
\Cref{thm:V70-protected-union}.
\end{proof}

\begin{firewall}[This is the overlap form of the V56 obstruction]
\label{fire:V70-overlap-no-double-charge}
The exact identity \(P^2=P\), not a loose inequality, is why overlap
mass cannot be charged twice.  Assigning one protected currency to
the original bank and another to the same coordinates under a
terminal label would reproduce the false capacity squaring excluded
by \Cref{lem:V56-idempotent-capacity}.
\end{firewall}

\section{Partial transpose sees only genuinely new coordinates}
\label{sec:V70-new-coordinates}

Let \(D_0\) be the original row-\(3\)--row-\(u\) bank and \(F\) a
terminal lateral protected bank.  Split
\begin{align}
 F_{\rm old}&:=F\cap D_0,\nonumber\\
 F_{\rm new}&:=F\setminus D_0,\nonumber\\
 A_{\rm new}&:=\{e\in F_{\rm new}:\text{row \(3\) is active at }e\}.
\label{eq:V70-new-split}
\end{align}
Put
\begin{equation}
 S_3^{\rm new}(F):=\sum_{e\in A_{\rm new}}a_{3,e}.
\label{eq:V70-new-active-mass}
\end{equation}

\begin{lemma}[Exact partial-transpose factorization after overlap
removal]
\label{lem:V70-new-PT}
Let \(M_0=P_{D_0}\) and \(N_F=P_F\).  Then
\begin{equation}
 M_0N_F
 =
 P_{D_0}P_{F_{\rm new}},
\label{eq:V70-remove-overlap}
\end{equation}
and
\begin{equation}
 \boxed{
 \|(M_0N_F)^{\Gamma_{3^c}}\|
 \le
 \|M_0^{\Gamma_{3^c}}\|\,
 C_G\bigl(S_3^{\rm new}(F)\bigr)^{-1/2}}
\label{eq:V70-new-PT}
\end{equation}
when \(A_{\rm new}\ne\varnothing\).  New coordinates on which row
\(3\) is inactive contribute at most one.
\end{lemma}

\begin{proof}
The union law gives
\[
 M_0N_F=P_{D_0\cup F}=P_{D_0}P_{F\setminus D_0},
\]
which is \eqref{eq:V70-remove-overlap}.  The two displayed factors
now have disjoint coordinate supports.  Partial transpose and
ordinary norm therefore tensorize between them.

On \(A_{\rm new}\), apply
\Cref{thm:V69-half-power}.  On
\(F_{\rm new}\setminus A_{\rm new}\), row \(3\) is inactive; full
transpose on the complementary side preserves the norm-one
projection bound.  Multiplying the disjoint coordinate factors
proves \eqref{eq:V70-new-PT}.
\end{proof}

\begin{remark}[Zero new active mass]
\label{rem:V70-zero-new}
If \(A_{\rm new}=\varnothing\), the valid conclusion is only
\[
 \|(M_0N_F)^{\Gamma_{3^c}}\|
 \le\|M_0^{\Gamma_{3^c}}\|.
\]
The expression
\((S_3^{\rm new})^{-1/2}\) is not used at zero.
\end{remark}

\section{The overlapping protected bridge}
\label{sec:V70-overlap-bridge}

\begin{theorem}[New-mass protected bridge]
\label{thm:V70-new-mass-bridge}
Suppose the original and terminal factors are all-protected in one
common outside-local resolved word.  If
\begin{equation}
 S_3^{\rm new}(F)
 \ge
 L\left(\frac{\Xi_u}{\Xi_4}\right)^4
\label{eq:V70-new-threshold}
\end{equation}
for a fixed sufficiently large group-dependent \(L\), then
\begin{equation}
 \boxed{
 \kappa_{3|c}(
 \widetilde W_{\pi,\nu}^*\widetilde W_{\pi,\nu})
 \le
 C_s\theta_{3u}(D_0)^2
 \le C_s(x+z)^2.}
\label{eq:V70-new-mass-RCQ}
\end{equation}
No disjointness of the full banks \(D_0,F\) is required.
\end{theorem}

\begin{proof}
\Cref{lem:V69-original-PT} and
\Cref{lem:V70-new-PT} give
\[
 \|(M_0N_F)^{\Gamma_{3^c}}\|
 \le
 C_sq_0^2\bigl(S_3^{\rm new}(F)\bigr)^{-1/2}.
\]
Under \eqref{eq:V70-new-threshold}, this is at most
\[
 C_sq_0^2\left(\frac{\Xi_4}{\Xi_u}\right)^2.
\]
The remainder of the proof is exactly the scalar normalization
calculation in \Cref{thm:V69-disjoint-bridge}: the factor
\((\Xi_4/\Xi_u)^2\) cancels the residual \(\Xi_4^{-2}\), the active
floors turn \(\Xi_u^{-2}\) into at most
\(a_*^{-2}\beta_0^2\), and
\eqref{eq:V60-theta-factorization} with root ownership turns
\(q_0^2\beta_0^2\) into
\(C_s\theta_{3u}(D_0)^2\).  The cut estimate gives the final
inequality.
\end{proof}

\begin{corollary}[V69 disjointness is removed]
\label{cor:V70-V69-extended}
The closure theorem
\Cref{thm:V69-disjoint-bridge} remains true for overlapping original
and terminal protected banks after replacing \(S_3(A)\) by the
genuinely new mass \(S_3^{\rm new}(F)\) in
\eqref{eq:V70-new-threshold}.
\end{corollary}

\begin{proof}
This is \Cref{thm:V70-new-mass-bridge}.  When the banks are
disjoint, \(F_{\rm new}=F\), so it contains the V69 result.
\end{proof}

\begin{assessment}[Exact concentrated residual after V70]
\label{ass:V70-residual}
For the all-protected original/terminal sector, overlap itself is no
longer an ambiguity.  It is removed by the exact union law.  Every
remaining term satisfies
\begin{equation}
 S_3^{\rm new}(F)
 <
 L\left(\frac{\Xi_u}{\Xi_4}\right)^4,
\label{eq:V70-subthreshold}
\end{equation}
possibly with \(S_3^{\rm new}(F)=0\).  Because the half-power
Weyl estimate is sharp, this subthreshold region requires signed
source information or a different cut-aligned response.  It cannot
be closed by assigning another protected capacity to the overlap.
\end{assessment}

\section{V70 result ledger and V71 acquisition order}
\label{sec:V70-status}

\begin{theorem}[Third-series V70 result ledger]
\label{thm:V70-results}
Relative to V1--V69, V70 proves:
\begin{enumerate}[label=\textup{(V70.\arabic*)},leftmargin=6.3em]
\item the compulsory active SM V52 and NS V31 audit with exact
      fingerprints;
\item a type-correct transfer of the one-collision/one-charge proof
      order;
\item uniqueness and idempotence of the local all-trivial projection
      under multiple bank labels;
\item the exact union law for overlapping protected banks;
\item removal of every duplicate protected currency on the overlap;
\item partial-transpose factorization through only the genuinely new
      row-\(3\)-active coordinates;
\item the new-mass protected bridge; and
\item removal of the full-bank disjointness hypothesis from the V69
      closed cell.
\end{enumerate}
\end{theorem}

\begin{proof}
These are
\Cref{sec:V70-audit,
ass:V70-audit-transfer,
lem:V70-same-local-projector,
thm:V70-protected-union,
cor:V70-contained-no-gain,
lem:V70-new-PT,
thm:V70-new-mass-bridge,
cor:V70-V69-extended}.
\end{proof}

\begin{programme}[V71: subthreshold and row-\(3\)-inactive terminal
mass]
\label{prog:V70-V71}
\begin{enumerate}[leftmargin=3em]
\item On the cell \eqref{eq:V70-subthreshold}, split the new terminal
      coordinates into row-\(3\)-active and inactive parts in the
      exact aggregate \eqref{eq:V67-aggregate-square}.
\item Prove that a protected factor on a row-\(3\)-inactive
      coordinate is a complementary-side contraction across
      \(3|c\), and determine whether it can be removed without
      conditioning the remaining product density.
\item On the active subthreshold part, retain the exact local
      representation dimensions and the signed first-connectivity
      coefficient; the sharp \(SU(2)\) one-coordinate test must be
      applied.
\item Treat an original non-all-protected factor separately using
      the repaired V68 one-failure resolver; do not transfer a
      terminal protected capacity across cuts.
\item If neither diagonal route supplies
      \eqref{eq:V67-LCB}, analyze the first off-diagonal
      \((r,s)\) pair in \eqref{eq:V67-aggregate-square} before taking
      absolute values.
\item The next compulsory companion audit is V75.
\end{enumerate}
\end{programme}

\begin{firewall}[Claim boundary after third-series V70]
\label{fire:V70-claim-boundary}
V70 closes the all-protected lateral terminal cell with sufficient
genuinely new row-\(3\)-active mass, including overlapping banks.  It
does not close the subthreshold, row-\(3\)-inactive,
original-gapped, finite exceptional, or local-payer branches, or
full \(RCQ_3(\eta)\).  The other quotient cuts, quartic source
families, higher MTA, WUFC, CPM, cutoff removal, reflection
positivity, Osterwalder--Schrader reconstruction, construction of
the continuum Yang--Mills measure, and a uniform positive continuum
spectral gap remain open.  The full Yang--Mills mass gap has not
been proved.
\end{firewall}

\paragraph{Parent-version record.}
V70 was formed from
\texttt{Renewal\_Yang\_Mills\_Mass\_Gap\_Research\_Notes\_V69.tex}.
The V69 parent had \(33\,982\) logical source lines,
\(1\,219\,490\) bytes, and SHA--256
\[
 \texttt{b18d637cc8a20e99033eb9673dbb8c54d2cd1b9f663804cd8673bb09c0442edd}.
\]
The next compulsory companion audit is V75.

% =====================================================================
% V71 APPENDIX: EXACT DIMENSION CURRENCY AND DIAGONAL EXHAUSTION
% =====================================================================

\chapter{V71: Exact Dimension Currency, One-Sided Neutrality, and
Diagonal Exhaustion}
\label{ch:V71-dimension-currency}

\section{What the subthreshold question actually asks}
\label{sec:V71-question}

The residual \eqref{eq:V70-subthreshold} was obtained only after
replacing the exact representation dimensions on the genuinely new
row-\(3\)-active coordinates by the universal Weyl half-power.  This
replacement is useful for a mass criterion, but it is not lossless.
The present chapter first restores the exact dimension product.

There is a second issue.  A new protected coordinate on which row
\(3\) is inactive is wholly on the complementary side of the cut.
Such a factor can be deleted from an upper bound without conditioning
the remaining row density.  Nevertheless, if the factor is a nonzero
projection, this deletion is exact: it supplies no smallness at all.
That distinction decides the diagonal branch.

\section{One-sided ancillary factors are exactly neutral}
\label{sec:V71-one-sided}

\begin{lemma}[Exact one-sided ancillary formula]
\label{lem:V71-one-sided}
Let \(A\succcurlyeq0\) act on
\(\mathcal H_3\otimes\mathcal H_c\), and let
\(Q\succcurlyeq0\) act on an ancillary complementary space
\(\mathcal K_c\).  Then
\begin{equation}
 \boxed{
 \kappa_{\mathcal H_3\,|\,
 \mathcal H_c\otimes\mathcal K_c}(A\otimes Q)
 =
 \kappa_{\mathcal H_3\,|\,\mathcal H_c}(A)\,\|Q\|.}
\label{eq:V71-one-sided}
\end{equation}
No product assumption among the tensor coordinates inside
\(\mathcal H_c\otimes\mathcal K_c\) is made.
\end{lemma}

\begin{proof}
For a unit vector \(v\in\mathcal H_3\), put
\[
 B_v:=(\langle v|\otimes I_c)A(|v\rangle\otimes I_c)
 \succcurlyeq0.
\]
After \(v\) is fixed, the complementary vector is allowed to be an
arbitrary unit vector of
\(\mathcal H_c\otimes\mathcal K_c\).  Hence
\[
 \sup_{\|w\|=1}
 \langle v\otimes w,(A\otimes Q)(v\otimes w)\rangle
 =
 \|B_v\otimes Q\|
 =
 \|B_v\|\,\|Q\|.
\]
Taking the supremum over \(v\) and using
\Cref{lem:V67-slice-capacity} proves
\eqref{eq:V71-one-sided}.
\end{proof}

\begin{corollary}[Row-\(3\)-inactive protected coordinates are
neutral]
\label{cor:V71-inactive-neutral}
Let \(I\) be any set of new terminal coordinates on which row \(3\)
is inactive, and let
\[
 Q_I=\bigotimes_{e\in I}P_e^{\mathrm{prot}}
\]
be a nonzero resolved protected factor.  For every positive operator
\(A\) on the remaining coordinates,
\begin{equation}
 \boxed{\kappa_{3|c}(A\otimes Q_I)=\kappa_{3|c}(A).}
\label{eq:V71-inactive-neutral}
\end{equation}
In particular, \(Q_I\) may be removed from an upper bound without
conditioning a product state, but it cannot be charged as a strict
row-\(3\)-cut currency.
\end{corollary}

\begin{proof}
Because row \(3\) is inactive, its local representation is
one-dimensional and
\(P_e^{\mathrm{prot}}\) acts only on the complementary-row factors.
The resolved atom is nonzero, so every projection has norm one and
\(\|Q_I\|=1\).  Apply \Cref{lem:V71-one-sided}.
\end{proof}

\begin{firewall}[Safe removal is not decay]
\label{fire:V71-removal-not-decay}
The inequality obtained by discarding a complementary-side
contraction is valid even when complementary row states are entangled
among coordinates.  What is invalid is to count the discarded factor
as another protected capacity.  Equation
\eqref{eq:V71-inactive-neutral} shows that the latter gain can be
exactly zero.
\end{firewall}

\section{The exact representation-dimension product}
\label{sec:V71-exact-product}

Retain the V70 notation
\[
 A_{\rm new}
 =
 \{e\in F\setminus D_0:\text{row \(3\) is active at }e\}.
\]
Define the empty product to be one and put
\begin{equation}
 {\mathfrak D}_3^{\rm new}(F)
 :=
 \prod_{e\in A_{\rm new}}d_{\lambda_{3,e}}.
\label{eq:V71-dimension-product}
\end{equation}

\begin{lemma}[Exact discrete refinement of the V70 partial-transpose
bound]
\label{lem:V71-exact-dimension-PT}
In the all-protected original/terminal sector,
\begin{equation}
 \boxed{
 \|(M_0N_F)^{\Gamma_{3^c}}\|
 \le
 C_{G,c_0}q_0^2
 \bigl({\mathfrak D}_3^{\rm new}(F)\bigr)^{-1}.}
\label{eq:V71-exact-dimension-PT}
\end{equation}
The coordinates of \(F\setminus D_0\) on which row \(3\) is inactive
do not occur in the product.
\end{lemma}

\begin{proof}
The union identity \eqref{eq:V70-remove-overlap} leaves disjoint
coordinate factors on \(D_0\) and \(F\setminus D_0\).  For each
\(e\in A_{\rm new}\),
\Cref{lem:V68-invariant-PT,lem:V69-local-row3-PT} gives the factor
\(d_{\lambda_{3,e}}^{-1}\).  Every remaining new factor is wholly on
the complementary side and has transposed norm one.  Ordinary norm
is multiplicative on the resulting coordinate tensor product.
Finally
\(\|M_0^{\Gamma_{3^c}}\|\le C_{G,c_0}q_0^2\) by
\Cref{lem:V69-original-PT}.
\end{proof}

\begin{theorem}[Exact dimension-product bridge]
\label{thm:V71-dimension-bridge}
There is a constant \(L_D<\infty\), depending only on the fixed
structural and group data, such that the following implication holds:
\begin{equation}
 {\mathfrak D}_3^{\rm new}(F)
 \ge
 L_D\left(\frac{\Xi_u}{\Xi_4}\right)^2
\quad\Longrightarrow\quad
 \boxed{
 \kappa_{3|c}(
 \widetilde W_{\pi,\nu}^*\widetilde W_{\pi,\nu})
 \le C_s(x+z)^2.}
\label{eq:V71-dimension-bridge}
\end{equation}
The full terminal bank may overlap \(D_0\), and no lower bound on
\(S_3^{\rm new}(F)\) is required.
\end{theorem}

\begin{proof}
Under the displayed hypothesis,
\Cref{lem:V71-exact-dimension-PT} gives
\[
 \|(M_0N_F)^{\Gamma_{3^c}}\|
 \le
 C_sq_0^2\left(\frac{\Xi_4}{\Xi_u}\right)^2.
\]
Insert this estimate into the central-square argument of
\Cref{thm:V69-disjoint-bridge}.  Explicitly,
\begin{align*}
 \kappa_{3|c}(
 \widetilde W_{\pi,\nu}^*\widetilde W_{\pi,\nu})
 &\le
 C_s
 \left(\frac{p}{\Xi_3\Xi_4}\right)^2
 q_0^2
 \left(\frac{\Xi_4}{\Xi_u}\right)^2\\
 &\le
 C_sa_*^{-4}q_0^2\beta_0^2\\
 &\le C_s\theta_{3u}(D_0)^2
 \le C_s(x+z)^2.
\end{align*}
The second line uses \(p\le1\), \(\Xi_3\ge a_*\), and
\(T_0\ge a_*\); the third uses
\eqref{eq:V60-theta-factorization} and root ownership.
\end{proof}

\begin{corollary}[The exact criterion strictly refines the mass
criterion]
\label{cor:V71-strict-refinement}
The V70 fourth-power mass cell is contained, after adjustment of
constants, in the cell of
\Cref{thm:V71-dimension-bridge}.  The containment is strict along
families with many bounded nontrivial row-\(3\) representations.
\end{corollary}

\begin{proof}
The proof of \Cref{thm:V69-half-power} gives directly
\[
 \bigl({\mathfrak D}_3^{\rm new}(F)\bigr)^{-1}
 \le C_G\bigl(S_3^{\rm new}(F)\bigr)^{-1/2}.
\]
Thus \(S_3^{\rm new}(F)\ge L(\Xi_u/\Xi_4)^4\) implies the dimension
criterion when \(L\) is sufficiently large.

For strictness, take \(G=SU(2)\) and \(N\) new active coordinates,
each with spin \(1/2\) in row \(3\).  Then
\[
 {\mathfrak D}_3^{\rm new}=2^N,
 \qquad
 S_3^{\rm new}=\frac74N.
\]
Given a large ratio \(R=\Xi_u/\Xi_4\), choose
\(N=\lceil2\log_2 R+\log_2L_D\rceil\).  The dimension criterion
holds, whereas \(S_3^{\rm new}=O(\log R)\), far below every fixed
multiple of \(R^4\).
\end{proof}

\begin{assessment}[Exact all-protected diagonal residual]
\label{ass:V71-diagonal-residual}
After V71, every unresolved all-protected diagonal term lies in
\begin{equation}
 \boxed{
 {\mathfrak D}_3^{\rm new}(F)
 <
 L_D\left(\frac{\Xi_u}{\Xi_4}\right)^2.}
\label{eq:V71-discrete-residual}
\end{equation}
Row-\(3\)-inactive coordinates have been removed exactly and do not
define an additional case.  The older mass subthreshold
\eqref{eq:V70-subthreshold} is therefore replaced by the sharper
discrete residual \eqref{eq:V71-discrete-residual}.
\end{assessment}

\section{A physical spectator test exhausts the diagonal route}
\label{sec:V71-spectator-test}

\begin{proposition}[Complementary protected spectators destroy every
purely diagonal ratio bound]
\label{prop:V71-diagonal-no-go}
Already for \(G=SU(2)\), there is no support-uniform implication from
the resolved diagonal protected factors alone of the form
\begin{equation}
 \|(M_0N_F)^{\Gamma_{3^c}}\|
 \le
 C q_0^2\left(\frac{\Xi_4}{\Xi_u}\right)^2
\label{eq:V71-false-diagonal-ratio}
\end{equation}
throughout the residual \eqref{eq:V71-discrete-residual}.
The obstruction persists with all local factors genuine invariant
projections.
\end{proposition}

\begin{proof}
Fix any nonzero resolved protected packet containing the original
factor and one new row-\(3\)-active singlet projection.  Thus its
left side in \eqref{eq:V71-false-diagonal-ratio} is a fixed positive
number; for a spin-\(j\) singlet the new factor has partial-transpose
norm \((2j+1)^{-1}\).

Now append \(m\) coordinates after the selected joining position.
At each added coordinate activate only rows \(1\) and \(2\), put
dual spin-\(1/2\) representations on them, and select their invariant
singlet.  These are genuine all-trivial local output projections.
They are row-\(3\)-inactive and wholly on the complementary side of
the cut.  By \Cref{cor:V71-inactive-neutral}, tensoring all \(m\) of
them leaves the row-\(3\) capacity, and equivalently the relevant
partial-transpose norm, unchanged.

Choose \(u\in\{1,2\}\).  The added bank increases \(\Xi_u\) linearly
in \(m\), while it changes neither \(\Xi_4\) nor
\({\mathfrak D}_3^{\rm new}(F)\).  Consequently
\[
 \left(\frac{\Xi_4}{\Xi_u}\right)^2\longrightarrow0,
 \qquad
 {\mathfrak D}_3^{\rm new}(F)
 <
 L_D\left(\frac{\Xi_u}{\Xi_4}\right)^2
\]
for all sufficiently large \(m\), whereas the left side of
\eqref{eq:V71-false-diagonal-ratio} stays fixed and positive.  This
contradicts any uniform constant \(C\).
\end{proof}

\begin{remark}[Scope of the spectator test]
\label{rem:V71-spectator-scope}
The proposition does not construct a counterexample to the complete
signed first-connectivity source.  Appending the coordinates changes
its suffix and its sum over possible first-connectivity positions.
It proves the narrower statement needed here: no estimate of one
diagonal resolved factor, based only on its protected projections,
can supply the missing global mass ratio.  The complete source is
precisely the information which remains available upstream.
\end{remark}

\section{The original one-failure sector does not change the verdict}
\label{sec:V71-one-failure}

\begin{lemma}[Repaired one-failure decay is not a denominator
response]
\label{lem:V71-one-failure-not-ratio}
Let \(\mathcal G_{D_0}(t)=M_{D_0}(t)-P_{D_0}\) be a nonzero physical
one-failure factor from \Cref{thm:V68-one-failure}.  Its valid bound
\[
 \kappa_{3|c}(\mathcal G_{D_0}(t))
 \le
 \Psi_{N_0,\rho_G}(r_t)
\]
contains heat time and original-bank cardinality, but it does not,
without a source-occurring response, imply a factor
\((\Xi_4/\Xi_u)^2\).
\end{lemma}

\begin{proof}
The displayed estimate is exactly
\Cref{thm:V68-one-failure} with root \(3\).  To test whether it can
also yield the mass ratio, keep the nonzero finite original block
fixed and append the complementary \(12\)-singlet spectators used in
\Cref{prop:V71-diagonal-no-go}.  By
\Cref{lem:V71-one-sided}, tensoring those projections leaves its
row-\(3\) capacity unchanged, while
\(\Xi_4/\Xi_u\to0\).  Thus the one-failure bound is a legitimate heat
or cardinality currency, but it is not the missing normalization
response.  If an already proved scalar factor \(\beta_0\) occurs in
the same word, \Cref{thm:V60-ownership-compensation} applies; this
lemma concerns exactly the residual in which that source occurrence
has not been obtained.
\end{proof}

\begin{firewall}[No transfer between the two gapped mechanisms]
\label{fire:V71-gap-separation}
A terminal cofinal gap supplies the mass ratio by
\Cref{cor:V67-gap-is-LCB}; an original one-failure factor supplies
heat/cardinality decay by \Cref{thm:V68-one-failure}.  These are
different occurrences.  The second cannot be renamed as the first,
and multiplying them is allowed only when both occur in the same
resolved word.
\end{firewall}

\section{What can and cannot be inferred from one off-diagonal pair}
\label{sec:V71-off-diagonal}

\begin{lemma}[Sharp two-term Gram interval]
\label{lem:V71-Gram-interval}
For arbitrary operators \(X,Y\),
\begin{equation}
 -\bigl(X^*X+Y^*Y\bigr)
 \preccurlyeq
 X^*Y+Y^*X
 \preccurlyeq
 X^*X+Y^*Y.
\label{eq:V71-Gram-interval}
\end{equation}
Both endpoints are attained: the upper one when \(Y=X\), and the
lower one when \(Y=-X\).
\end{lemma}

\begin{proof}
Expand the positive operators
\[
 (X-Y)^*(X-Y)\succcurlyeq0,
 \qquad
 (X+Y)^*(X+Y)\succcurlyeq0.
\]
Rearranging gives the two inequalities.  Substitution of
\(Y=\pm X\) proves sharpness.
\end{proof}

\begin{proposition}[First-position labels alone impose no cancelling
sign]
\label{prop:V71-no-abstract-cross-cancellation}
The algebraic facts that two terms have different first-connectivity
positions, share a payer label, and occur in the positive square
\eqref{eq:V67-aggregate-square} do not imply a negative
off-diagonal contribution.  A useful gain must come from an exact
source-specific relation between their signed local cumulants and
complete intervening moments.
\end{proposition}

\begin{proof}
The contribution of a pair is
\[
 X^*X+Y^*Y+X^*Y+Y^*X=(X+Y)^*(X+Y).
\]
\Cref{lem:V71-Gram-interval} permits both complete reinforcement and
complete cancellation with identical diagonal data.  Distinct
position labels do not enter this operator identity.  Therefore a
sign cannot be assigned before the actual first-connectivity
coefficients and the intervening suffix/prefix factors are written.
\end{proof}

\begin{assessment}[Upstream conclusion of V71]
\label{ass:V71-upstream}
The diagonal protected analysis is now exhausted without ambiguity:
\begin{enumerate}[label=\textup{(\roman*)},leftmargin=4em]
\item overlap is removed by the V70 union law;
\item row-\(3\)-inactive factors are exactly neutral;
\item the exact dimension-product cell
      \eqref{eq:V71-dimension-bridge} is closed;
\item the residual is the discrete concentration cell
      \eqref{eq:V71-discrete-residual};
\item original one-failure decay does not create the missing mass
      ratio; and
\item abstract positivity supplies no cross-position sign.
\end{enumerate}
Hence the next valid move is upstream to the exact two-position
first-connectivity algebra, retaining every coordinate between the
two joining positions.  Repeating a diagonal partial-transpose
estimate cannot close the residual.
\end{assessment}

\section{V71 result ledger and V72 acquisition order}
\label{sec:V71-status}

\begin{theorem}[Third-series V71 result ledger]
\label{thm:V71-results}
Relative to V1--V70, V71 proves:
\begin{enumerate}[label=\textup{(V71.\arabic*)},leftmargin=6.3em]
\item the exact capacity formula for a one-sided ancillary
      contraction;
\item exact neutrality of every row-\(3\)-inactive protected factor;
\item the representation-dimension product refinement of the V70
      partial-transpose estimate;
\item the exact dimension-product bridge and its strict improvement
      over the fourth-power mass threshold;
\item reduction of the all-protected diagonal residual to
      \eqref{eq:V71-discrete-residual};
\item a physical \(SU(2)\) spectator no-go for every purely diagonal
      mass-ratio argument;
\item separation of repaired one-failure decay from a terminal
      denominator response; and
\item the sharp two-term Gram interval excluding any abstract
      off-diagonal sign claim.
\end{enumerate}
\end{theorem}

\begin{proof}
These are
\Cref{lem:V71-one-sided,
cor:V71-inactive-neutral,
lem:V71-exact-dimension-PT,
thm:V71-dimension-bridge,
cor:V71-strict-refinement,
ass:V71-diagonal-residual,
prop:V71-diagonal-no-go,
lem:V71-one-failure-not-ratio,
lem:V71-Gram-interval,
prop:V71-no-abstract-cross-cancellation}.
\end{proof}

\begin{programme}[V72: the exact first off-diagonal incidence]
\label{prog:V71-V72}
\begin{enumerate}[leftmargin=3em]
\item Choose \(r<s\) in the same payer aggregate
      \eqref{eq:V67-lateral-aggregate} and write both
      first-connectivity words with the three ordered regions
      \(e<r\), \(r<e<s\), and \(e>s\).
\item Expand only the complete moments in the intervening region;
      do not resolve either joining cumulant by an absolute envelope.
\item Determine the exact local row-partition joins required for a
      word first connected at \(r\) and one first connected at \(s\).
\item Compute whether their cross term vanishes under any genuine
      orthogonal output compression.  Do not infer orthogonality from
      distinct coordinate labels.
\item If no cancellation holds, locate the earliest intervening
      response occurrence whose deletion changes the later joining
      position, and formulate its normalized response card.
\item Keep finite-exceptional and local-payer branches separate.  The
      next compulsory companion audit remains V75.
\end{enumerate}
\end{programme}

\begin{firewall}[Claim boundary after third-series V71]
\label{fire:V71-claim-boundary}
V71 closes the exact large-dimension-product protected lateral cell
and removes row-\(3\)-inactive factors from the case list.  It does
not close the discrete concentration residual, original-gapped
residual without a source-occurring response, finite exceptional or
local-payer branches, or full \(RCQ_3(\eta)\).  The other quotient
cuts, quartic source families, higher MTA, WUFC, CPM, cutoff removal,
reflection positivity, Osterwalder--Schrader reconstruction,
construction of the continuum Yang--Mills measure, and a uniform
positive continuum spectral gap remain open.  The full Yang--Mills
mass gap has not been proved.
\end{firewall}

\paragraph{Parent-version record.}
V71 was formed from
\texttt{Renewal\_Yang\_Mills\_Mass\_Gap\_Research\_Notes\_V70.tex}.
The V70 parent had \(34\,354\) logical source lines,
\(1\,232\,463\) bytes, and SHA--256
\[
 \texttt{8c66442af54b83c0d812f3ba0e1ee84cfee6e5935ea9abebbff781d2b7690794}.
\]
The next compulsory companion audit is V75.

% =====================================================================
% V72 APPENDIX: TWO-POSITION FIRST-CONNECTIVITY ALGEBRA
% =====================================================================

\chapter{V72: Exact Two-Position Incidence and the First-Escape
Telescope}
\label{ch:V72-two-position}

\section{Local notation at the \(2+1+1\) prefix}
\label{sec:V72-notation}

Fix
\[
 \rho:=12|3|4,
 \qquad
 \widehat1:=1234.
\]
The exact local moment
\(M_e(\rho)=\sum_{\pi\le\rho}L_{e,\pi}\) was defined in
\eqref{eq:V2-local-partition-moment}.  In this chapter write
\begin{equation}
 R_e^\rho:=M_e(\rho),
 \qquad
 U_e:=M_e(\widehat1)=\sum_{\pi\in\Pi([4])}L_{e,\pi}.
\label{eq:V72-local-RU}
\end{equation}
For an ordered interval \(I\), put
\[
 R_I^\rho:=\bigotimes_{e\in I}R_e^\rho,
 \qquad
 U_I:=\bigotimes_{e\in I}U_e.
\]
Empty interval products are identities.  The joining carrier is
\[
 J_e(\rho)
 =
 \sum_{\rho\vee\pi=\widehat1}L_{e,\pi}.
\]

\begin{lemma}[Refinement characterization of a delayed join]
\label{lem:V72-delayed-refinement}
Suppose the cumulative row partition immediately before coordinate
\(r\) is \(\rho\), and the first global join occurs at \(s>r\) with
pre-join partition still equal to \(\rho\).  Then every local
partition \(\pi_e\) with \(r\le e<s\) refines \(\rho\), and the exact
sum over all such intervening letters is \(R_{[r,s)}^\rho\).
\end{lemma}

\begin{proof}
Cumulative joins are monotone in the partition lattice.  If some
\(\pi_e\nleq\rho\), then
\(\rho\vee\pi_e>\rho\), and no later join can return the cumulative
partition to \(\rho\).  Hence every \(\pi_e\le\rho\).  Conversely,
any sequence of refinements of \(\rho\) leaves the cumulative join
equal to \(\rho\).  Summing independently at each coordinate gives
\[
 \bigotimes_{r\le e<s}\sum_{\pi_e\le\rho}L_{e,\pi_e}
 =
 R_{[r,s)}^\rho.
\]
\end{proof}

\section{Exact splitting of the later prefix}
\label{sec:V72-prefix-split}

For a coordinate set \(A\), write \(K_{12}^{A,\mathrm{conn}}\) for
the complete connected binary carrier and \(M_i^A\) for the complete
one-row moment.  The following elementary identity is the exact
binary first-connection split which is needed before the two
four-row positions are compared.

\begin{lemma}[Binary covariance concatenation]
\label{lem:V72-binary-concatenation}
For disjoint ordered coordinate sets \(A<B\),
\begin{align}
 K_{12}^{A\cup B,\mathrm{conn}}
 ={}&
 K_{12}^{A,\mathrm{conn}}\otimes M_{12}^{B}
 \nonumber\\
 &+
 (M_1^A\otimes M_2^A)
 \otimes K_{12}^{B,\mathrm{conn}}.
\label{eq:V72-binary-concatenation}
\end{align}
\end{lemma}

\begin{proof}
Use
\[
 K_{12}^{A,\mathrm{conn}}
 =
 M_{12}^A-M_1^A\otimes M_2^A
\]
and the coordinate factorizations
\[
 M_{12}^{A\cup B}=M_{12}^A\otimes M_{12}^B,
 \qquad
 M_i^{A\cup B}=M_i^A\otimes M_i^B.
\]
Adding and subtracting
\((M_1^A\otimes M_2^A)\otimes M_{12}^B\) proves the displayed
identity.  Equivalently, row \(1\) and row \(2\) first connect in
\(A\) or remain separate throughout \(A\) and connect in \(B\).
\end{proof}

For \(r<s\), define the late-binary carrier
\begin{align}
 \mathcal D_{r,s}:={}&
 M_1^{<r}\otimes M_2^{<r}
 \otimes M_3^{<r}\otimes M_4^{<r}
 \nonumber\\
 &\otimes
 K_{12}^{[r,s),\mathrm{conn}}
 \otimes M_3^{[r,s)}
 \otimes M_4^{[r,s)}.
\label{eq:V72-late-binary}
\end{align}

\begin{theorem}[Exact \(r<s\) decomposition of the later
\(2+1+1\) word]
\label{thm:V72-later-word-split}
Let
\(\mathcal C_{<r}\) be \eqref{eq:V66-prefix-carrier}.  Then
\begin{equation}
 \boxed{
 \mathcal C_{<s}
 =
 \mathcal C_{<r}\otimes R_{[r,s)}^\rho
 +
 \mathcal D_{r,s}.}
\label{eq:V72-prefix-split}
\end{equation}
Consequently the uncompressed first-connectivity word at \(s\)
splits exactly as
\begin{align}
 \mathcal W_s
 ={}&
 \underbrace{
 \mathcal C_{<r}\otimes R_{[r,s)}^\rho
 \otimes J_s(\rho)\otimes M_{>s}}_
 {\mathcal W_{s\leftarrow r}^{\mathrm{old}}}
 \nonumber\\
 &+
 \underbrace{
 \mathcal D_{r,s}\otimes J_s(\rho)\otimes M_{>s}}_
 {\mathcal W_{s\leftarrow r}^{\mathrm{new}}}.
\label{eq:V72-later-word-split}
\end{align}
The first term already has its binary \(12\) connection before \(r\);
the second obtains that binary connection in \([r,s)\).
\end{theorem}

\begin{proof}
Apply \Cref{lem:V72-binary-concatenation} with
\(A=\{e:e<r\}\) and \(B=[r,s)\).  Tensor the first summand with the
complete row-\(3\) and row-\(4\) singleton moments.  Its interval
factor is
\[
 M_{12}^{[r,s)}\otimes M_3^{[r,s)}
 \otimes M_4^{[r,s)}
 =
 R_{[r,s)}^\rho
\]
by block factorization.  The second summand is
\(\mathcal D_{r,s}\).  This proves
\eqref{eq:V72-prefix-split}.  Tensoring with the joining carrier at
\(s\) and the complete suffix gives
\eqref{eq:V72-later-word-split}.
\end{proof}

\begin{remark}[The split is not an estimate]
\label{rem:V72-no-estimate}
Both summands in \eqref{eq:V72-later-word-split} retain their signs.
No endpoint, Fourier-output, payer, or norm resolution has been
applied.  A fixed outside-local payer may be reinserted in its actual
slot after this algebraic split; if it lies in \([r,s]\), it remains
inside the corresponding interval factor.
\end{remark}

\section{The common-prefix cross kernel}
\label{sec:V72-cross-kernel}

The word first connected at \(r\) is
\begin{equation}
 \mathcal W_r
 =
 \mathcal C_{<r}\otimes J_r(\rho)
 \otimes U_{(r,s]}\otimes M_{>s}.
\label{eq:V72-early-word}
\end{equation}

\begin{theorem}[Exact factorization of the first common-prefix cross
term]
\label{thm:V72-cross-factorization}
Before physical compression and payer insertion,
\begin{align}
 \mathcal W_r^*
 \mathcal W_{s\leftarrow r}^{\mathrm{old}}
 ={}&
 (\mathcal C_{<r}^*\mathcal C_{<r})
 \otimes
 \bigl(J_r(\rho)^*R_r^\rho\bigr)
 \nonumber\\
 &\otimes
 \bigotimes_{r<e<s}(U_e^*R_e^\rho)
 \otimes
 \bigl(U_s^*J_s(\rho)\bigr)
 \nonumber\\
 &\otimes
 (M_{>s}^*M_{>s}).
\label{eq:V72-cross-factorization}
\end{align}
The adjoint cross term is the adjoint of the displayed tensor.
\end{theorem}

\begin{proof}
Insert \eqref{eq:V72-later-word-split} and
\eqref{eq:V72-early-word}.  Operators belonging to different
physical coordinates tensorize.  Multiplication at coordinate \(r\)
gives \(J_r^*R_r^\rho\); at \(r<e<s\) it gives
\(U_e^*R_e^\rho\); at \(s\) it gives \(U_s^*J_s\); the prefix and
suffix factors are common.  This is exactly
\eqref{eq:V72-cross-factorization}.
\end{proof}

\begin{corollary}[Different first-connectivity positions are not an
orthogonality certificate]
\label{cor:V72-no-position-orthogonality}
The labels \(r\ne s\) alone do not imply
\(\mathcal W_r^*\mathcal W_s=0\).  On the common-prefix branch,
orthogonality must be proved from an actual vanishing local compressed
kernel in \eqref{eq:V72-cross-factorization}, or from a cancellation
after the late-binary branch is restored.
\end{corollary}

\begin{proof}
The tensor product in
\eqref{eq:V72-cross-factorization} contains no factor that vanishes
as a formal consequence of \(r<s\).  In particular, in the universal
partition-letter algebra one may evaluate one refining letter and
one joining letter as the scalar \(1\), all unused letters as zero,
and all common moments as \(1\).  Then the displayed cross kernel is
\(1\).  Thus no identity in the partition lattice makes it zero.
Physical orthogonality, if present, must come from the actual output
maps rather than the position labels.
\end{proof}

\begin{firewall}[Orthogonal spectral blocks are not orthogonal
partition histories]
\label{fire:V72-output-not-history}
The output resolution is orthogonal in representation labels at one
fixed complete moment-partition state.  The two histories in
\eqref{eq:V72-cross-factorization} use different local partition
letters at \(r\) and \(s\).  One may not infer that their ranges are
orthogonal without computing the compressed products
\(J_r^*R_r^\rho\) and \(U_s^*J_s\).  This is the cross-position
version of \Cref{fire:V58-orthogonal-not-product}.
\end{firewall}

\section{A coefficient-one first-escape telescope}
\label{sec:V72-first-escape}

Define the local escape letter
\begin{equation}
 E_e^\rho
 :=
 U_e-R_e^\rho
 =
 \sum_{\pi\nleq\rho}L_{e,\pi}.
\label{eq:V72-escape-letter}
\end{equation}

\begin{theorem}[Exact first-escape telescope]
\label{thm:V72-first-escape}
For every \(r<s\),
\begin{equation}
 \boxed{
 U_{[r,s)}-R_{[r,s)}^\rho
 =
 \sum_{t=r}^{s-1}
 R_{[r,t)}^\rho
 \otimes E_t^\rho
 \otimes U_{(t,s)}.}
\label{eq:V72-first-escape}
\end{equation}
Every word on the right occurs exactly once, at its first coordinate
whose local partition does not refine \(\rho\).
\end{theorem}

\begin{proof}
This is the coefficient-one tensor telescope
\[
 \bigotimes_{e=r}^{s-1}U_e
 -
 \bigotimes_{e=r}^{s-1}R_e^\rho
 =
 \sum_{t=r}^{s-1}
 \left(\bigotimes_{r\le e<t}R_e^\rho\right)
 \otimes(U_t-R_t^\rho)
 \otimes
 \left(\bigotimes_{t<e<s}U_e\right).
\]
Substitute \eqref{eq:V72-escape-letter}.  In a partition word not
entirely refining \(\rho\), the first nonrefining coordinate \(t\) is
unique, which proves the final assertion.
\end{proof}

Let the three quotient vertices be
\[
 A=\{1,2\},\qquad B=\{3\},\qquad C=\{4\}.
\]
For \(\pi\nleq\rho\), the coarser partition
\(\rho\vee\pi\) is one of
\[
 AB|C,\qquad AC|B,\qquad BC|A,\qquad ABC.
\]
Accordingly define
\begin{align}
 E_e^{AB}
 &:=
 \sum_{\rho\vee\pi=AB|C}L_{e,\pi},&
 E_e^{AC}
 &:=
 \sum_{\rho\vee\pi=AC|B}L_{e,\pi},
\nonumber\\
 E_e^{BC}
 &:=
 \sum_{\rho\vee\pi=BC|A}L_{e,\pi},&
 E_e^{ABC}
 &:=
 \sum_{\rho\vee\pi=ABC}L_{e,\pi}.
\label{eq:V72-four-escape-types}
\end{align}

\begin{lemma}[Exhaustive quotient-edge split]
\label{lem:V72-escape-types}
One has the exact disjoint sum
\begin{equation}
 \boxed{
 E_e^\rho
 =
 E_e^{AB}+E_e^{AC}+E_e^{BC}+E_e^{ABC},
 \qquad
 E_e^{ABC}=J_e(\rho).}
\label{eq:V72-escape-split}
\end{equation}
The first three terms create exactly one edge of the quotient
triangle; the fourth connects all three quotient vertices.
\end{lemma}

\begin{proof}
The interval \([\rho,\widehat1]\) in the partition lattice is
isomorphic to the partition lattice on the three blocks \(A,B,C\).
Its elements strictly above \(\rho\) are exactly the three two-block
partitions and \(\widehat1\).  Group the summands in
\eqref{eq:V72-escape-letter} by the value of
\(\rho\vee\pi\).  The final group is the definition of
\(J_e(\rho)\).
\end{proof}

\begin{corollary}[Every delay response has an actual quotient
incidence]
\label{cor:V72-response-incidence}
Every summand of
\(U_{[r,s)}-R_{[r,s)}^\rho\) contains a uniquely marked first
coordinate \(t\) and at least one of the quotient incidences
\[
 AB,\qquad AC,\qquad BC.
\]
There is no response summand supported only inside the composite
block \(A=\{1,2\}\).
\end{corollary}

\begin{proof}
Apply \Cref{thm:V72-first-escape,lem:V72-escape-types}.  A local
partition supported only inside \(A\), together with singleton blocks
\(B,C\), refines \(\rho\), so it belongs to \(R_t^\rho\) and not to
\(E_t^\rho\).
\end{proof}

\section{The no-escape obstruction and the next response card}
\label{sec:V72-no-escape}

\begin{proposition}[Internal \(12\) spectators survive every
first-escape telescope]
\label{prop:V72-internal-spectators}
Append any collection of coordinates at which all local partition
letters refine \(\rho\), in particular the row-\(3\)-inactive
protected \(12\)-singlet spectators of
\Cref{prop:V71-diagonal-no-go}.  These coordinates occur entirely in
the \(R^\rho\) factors of
\eqref{eq:V72-first-escape}.  They neither create an escape letter nor
produce a coefficient from the difference
\(U-R^\rho\).
\end{proposition}

\begin{proof}
At such a coordinate \(U_e=R_e^\rho\) on the selected local
partition word.  Hence its contribution to
\(E_e^\rho=U_e-R_e^\rho\) is zero.  In every nonzero telescope
summand it is an unmarked factor before or after the unique escape
coordinate.  The special \(12\)-singlet spectators are internal to
the quotient block \(A\), so the conclusion also follows from
\Cref{cor:V72-response-incidence}.
\end{proof}

\begin{assessment}[Exact location of a possible cross-position gain]
\label{ass:V72-cross-location}
The first off-diagonal analysis has two sharply separated pieces:
\begin{enumerate}[label=\textup{(\roman*)},leftmargin=4em]
\item the common-prefix/no-escape kernel
      \eqref{eq:V72-cross-factorization}, which has no formal
      orthogonality and retains all internal \(12\) spectators; and
\item the coefficient-one escape sum
      \eqref{eq:V72-first-escape}, in which every marked term has an
      actual quotient incidence.
\end{enumerate}
Thus an off-diagonal proof of \(RCQ_3(\eta)\) must either establish a
physical compressed cancellation in the no-escape kernel or estimate
the escape sum using its \(AB,AC,BC\) incidence.  The position labels
and partition lattice alone do neither.
\end{assessment}

For later use, define the compressed no-escape kernel for the fixed
outside-local payer label \(\pi\) by
\begin{equation}
 {\cal N}_{r,s}^{(\pi)}
 :=
 \operatorname{Comp}_{\pi}\!\left[
 \mathcal W_r^*
 \mathcal W_{s\leftarrow r}^{\mathrm{old}}
 \right],
\label{eq:V72-no-escape-kernel}
\end{equation}
where \(\operatorname{Comp}_{\pi}\) denotes the already fixed physical
output compression, separators, and payer insertion in their source
order.

\begin{definition}[No-escape response card]
\label{def:V72-NERC}
The card \(NERC_3(\eta)\) is the paired estimate
\begin{equation}
 \boxed{
 \kappa_{3|c}\!\left(
 \sum_{r<s}
 \bigl({\cal N}_{r,s}^{(\pi)}
 +({\cal N}_{r,s}^{(\pi)})^*\bigr)
 +
 \sum_r
 \widetilde W_{r,\pi}^*\widetilde W_{r,\pi}
 \right)
 \le C_{s,\eta}(x+z)^2,}
\label{eq:V72-NERC}
\end{equation}
with the sums restricted to one finite resolved lateral aggregate.
The diagonal is included because the off-diagonal part alone need not
be positive.
\end{definition}

\begin{firewall}[The response card is not yet a theorem]
\label{fire:V72-NERC-open}
\Cref{def:V72-NERC} records the smallest positive combination in
which no-escape cross-position cancellation could close the discrete
residual.  V72 proves its exact source ancestry but does not prove the
inequality.  The escape terms and the late-binary terms from
\eqref{eq:V72-later-word-split} must be added back before the complete
aggregate is recovered.
\end{firewall}

\section{V72 result ledger and V73 acquisition order}
\label{sec:V72-status}

\begin{theorem}[Third-series V72 result ledger]
\label{thm:V72-results}
Relative to V1--V71, V72 proves:
\begin{enumerate}[label=\textup{(V72.\arabic*)},leftmargin=6.3em]
\item the refinement characterization of a delayed \(2+1+1\) join;
\item exact binary covariance concatenation across \(r<s\);
\item the old-binary/late-binary decomposition of the word first
      connected at \(s\);
\item coordinatewise factorization of the first common-prefix cross
      kernel;
\item failure of position labels alone to provide orthogonality;
\item the coefficient-one first-escape telescope;
\item its exhaustive \(AB,AC,BC,ABC\) quotient-incidence split;
\item proof that internal protected \(12\) spectators are invisible
      to every escape response; and
\item the exact no-escape response card on the smallest positive
      paired aggregate.
\end{enumerate}
\end{theorem}

\begin{proof}
These are
\Cref{lem:V72-delayed-refinement,
lem:V72-binary-concatenation,
thm:V72-later-word-split,
thm:V72-cross-factorization,
cor:V72-no-position-orthogonality,
thm:V72-first-escape,
lem:V72-escape-types,
cor:V72-response-incidence,
prop:V72-internal-spectators,
def:V72-NERC}.
\end{proof}

\begin{programme}[V73: resolve the no-escape kernel before estimating
it]
\label{prog:V72-V73}
\begin{enumerate}[leftmargin=3em]
\item Insert the physical Fourier-output projectors at coordinates
      \(r\) and \(s\) in
      \eqref{eq:V72-cross-factorization}; compute the common output
      labels of \(J_r^*R_r^\rho\) and \(U_s^*J_s\).
\item Determine whether the all-protected discrete residual makes
      either local kernel zero.  If it does not, construct the
      smallest \(SU(2)\) recoupling block on which it is nonzero.
\item Split the escape telescope by
      \eqref{eq:V72-escape-split}; route the \(AB\) and \(BC\)
      incidences directly to the row-\(3\) cut and keep the lateral
      \(AC\) incidence separate.
\item Test whether the late-binary term in
      \eqref{eq:V72-later-word-split} carries the missing
      partner-ownership fraction.  Do not reuse the already charged
      prefix response \(p\).
\item If the no-escape \(SU(2)\) block is nonzero and neutral, return
      upstream to the unsplit two-row prefix covariance rather than
      asserting cancellation.
\item The next compulsory companion audit remains V75.
\end{enumerate}
\end{programme}

\begin{firewall}[Claim boundary after third-series V72]
\label{fire:V72-claim-boundary}
V72 identifies the exact two-position cross kernel and every
first-escape incidence.  It does not prove \(NERC_3(\eta)\), close the
discrete concentration, original-gapped, finite exceptional, or
local-payer residuals, or establish full \(RCQ_3(\eta)\).  The other
quotient cuts, quartic source families, higher MTA, WUFC, CPM, cutoff
removal, reflection positivity, Osterwalder--Schrader reconstruction,
construction of the continuum Yang--Mills measure, and a uniform
positive continuum spectral gap remain open.  The full Yang--Mills
mass gap has not been proved.
\end{firewall}

\paragraph{Parent-version record.}
V72 was formed from
\texttt{Renewal\_Yang\_Mills\_Mass\_Gap\_Research\_Notes\_V71.tex}.
The V71 parent had \(34\,860\) logical source lines,
\(1\,250\,532\) bytes, and SHA--256
\[
 \texttt{8f57ca5fef72b9e470965cdc1e600bed15017ae51490d8678e02c20989f13cb7}.
\]
The next compulsory companion audit is V75.

% =====================================================================
% V73 APPENDIX: LOCALIZED PREFIX NORMALIZATION
% =====================================================================

\chapter{V73: Localized Prefix Normalization and Closure of the
Low-Prefix-Mass Lateral Cell}
\label{ch:V73-localized-prefix}

\section{The upstream loss in the global prefix fraction}
\label{sec:V73-upstream-loss}

V72 was designed to test cross-position cancellation.  Its
first-escape telescope also reveals why protected spectators after a
joining coordinate appeared to defeat every diagonal mass-ratio
argument: the scalar factor in
\eqref{eq:V62-normalized-factorization} used the global stock
\[
 H_{12}=\sum_e h_{12,e},
\]
although the actual binary prefix at position \(r\) uses only
coordinates \(e<r\).  The replacement
\(H_{12}^{<r}\le H_{12}\) is harmless for a global tree majorant, but
it is not harmless when a missing lateral denominator must respond to
mass located after \(r\).

Define
\begin{equation}
 H_{12}^{<r}:=\sum_{e<r}h_{12,e},
 \qquad
 p_r^-:=\frac{H_{12}^{<r}}{\Xi_1\Xi_2}.
\label{eq:V73-local-prefix-fraction}
\end{equation}
Then \(0\le p_r^-\le p\le1\), but the first inequality can be
arbitrarily strict.

\begin{assessment}[Why V73 changes direction]
\label{ass:V73-direction-change}
The V72 no-escape kernel need not be orthogonal, and V51 supplies a
physical scalar channel in which the local joining cumulant is
strictly positive.  Computing a recoupling matrix cannot repair a
scalar overestimate made before that matrix is reached.  V73
therefore goes upstream to the exact prefix stock.  This is not a new
use of the prefix response: it replaces the coarse numerator
\(H_{12}\) by the smaller numerator that actually occurs in the same
source word.
\end{assessment}

\section{Exact localized normalization of a fixed word}
\label{sec:V73-word-normalization}

\begin{lemma}[The binary prefix uses only \(H_{12}^{<r}\)]
\label{lem:V73-prefix-stock}
For the complete connected binary prefix in
\eqref{eq:V66-uncompressed-word},
\begin{equation}
 \|K_{12}^{<r,\mathrm{conn}}\|
 \le C_s H_{12}^{<r},
\label{eq:V73-prefix-stock}
\end{equation}
with the corresponding once-paid estimate when the unique payer lies
in that prefix.  No coordinate \(e\ge r\) occurs on the right.
\end{lemma}

\begin{proof}
Apply the complete binary-prefix response theorem used in
\Cref{lem:V55-complete-prefix-response,lem:V62-residual-scalar} to
the actual support \(\{e:e<r\}\), rather than enlarging its pair stock
to the full coordinate set.  That theorem is support-local: its
coefficient-one covariance telescope sums one \(12\) influence
\(h_{12,e}\) over coordinates belonging to the covariance support.
The sum is therefore precisely \(H_{12}^{<r}\).  The payer version
has the same support because the payer is retained at its actual
prefix occurrence.
\end{proof}

\begin{lemma}[Localized resolved-word factorization]
\label{lem:V73-localized-factorization}
Fix \(r\), an outside-local payer placement, and one common
restriction/output resolution containing a positive original bank
factor \(\mathcal M_{D_0}\) and a commuting positive terminal factor
\(\mathcal N_F\).  The normalized word has a presentation
\begin{equation}
 \boxed{
 \widetilde W_{r,\pi,\nu}
 =
 \frac{p_r^-}{\Xi_3\Xi_4}\,
 L_{r,\pi,\nu}\,
 \mathcal M_{D_0}\mathcal N_F\,
 Q_{r,\pi,\nu},}
\label{eq:V73-localized-factorization}
\end{equation}
where
\begin{equation}
 \|L_{r,\pi,\nu}\|+\|Q_{r,\pi,\nu}\|\le C_s,
 \qquad
 [\mathcal M_{D_0},\mathcal N_F]
 =
 [\mathcal M_{D_0}\mathcal N_F,Q_{r,\pi,\nu}]
 =0.
\label{eq:V73-localized-dressings}
\end{equation}
If \(H_{12}^{<r}=0\), the word is zero.
\end{lemma}

\begin{proof}
Repeat the proof of
\Cref{lem:V62-residual-scalar} with
\Cref{lem:V73-prefix-stock} in place of the enlarged bound
\(C_sH_{12}\).  Division by the four input masses gives the exact
scalar identity
\[
 \frac{H_{12}^{<r}}
 {\Xi_1\Xi_2\Xi_3\Xi_4}
 =
 \frac{p_r^-}{\Xi_3\Xi_4}.
\]
When the numerator is nonzero, divide the bounded prefix-side
operator by \(H_{12}^{<r}\) and absorb it into \(L_{r,\pi,\nu}\).
The V66 joint source resolution and V59 centrality place the two
positive bank factors in the displayed order.  All remaining moments,
compressions, separators, and the fixed payer have the same uniform
bounds as in \eqref{eq:V62-normalized-dressings}.  If the numerator
vanishes, \eqref{eq:V73-prefix-stock} forces the prefix, hence the
word, to vanish.
\end{proof}

\begin{firewall}[Localized replacement, not a second prefix charge]
\label{fire:V73-one-prefix-charge}
Equation \eqref{eq:V73-localized-factorization} uses the binary
prefix once.  It does not multiply the old factor \(p\) by
\(p_r^-\).  The former was obtained by enlarging the latter's
numerator; V73 keeps the un-enlarged numerator from the start.
\end{firewall}

\section{A deterministic prefix-fraction inequality}
\label{sec:V73-prefix-fraction}

For \(i\in[4]\), put
\begin{equation}
 \Xi_i^{<r}:=\sum_{e<r}a_{i,e},
 \qquad
 \alpha_i^-(r):=\frac{\Xi_i^{<r}}{\Xi_i}.
\label{eq:V73-prefix-mass}
\end{equation}

\begin{lemma}[The localized prefix edge is bounded by either endpoint
fraction]
\label{lem:V73-prefix-endpoint-fraction}
For every \(r\),
\begin{equation}
 \boxed{
 p_r^-
 \le
 \min\{\alpha_1^-(r),\alpha_2^-(r)\}.}
\label{eq:V73-prefix-endpoint-fraction}
\end{equation}
\end{lemma}

\begin{proof}
For \(u\in\{1,2\}\), let \(v\) be the other endpoint.  Positivity
gives
\[
 H_{12}^{<r}
 =
 \sum_{e<r}a_{u,e}a_{v,e}
 \le
 \left(\sum_{e<r}a_{u,e}\right)
 \left(\sum_{e<r}a_{v,e}\right)
 \le
 \Xi_u^{<r}\Xi_v.
\]
Divide by \(\Xi_u\Xi_v=\Xi_1\Xi_2\).  Doing this for both choices of
\(u\) proves the minimum.
\end{proof}

\begin{corollary}[The missing lateral ratio is already present on the
low-prefix cell]
\label{cor:V73-prefix-ratio}
Let \(u\in\{1,2\}\).  If
\begin{equation}
 \Xi_u^{<r}\le K\Xi_4,
\label{eq:V73-low-prefix-cell}
\end{equation}
then
\begin{equation}
 \boxed{
 p_r^-\le K\frac{\Xi_4}{\Xi_u}.}
\label{eq:V73-prefix-ratio}
\end{equation}
\end{corollary}

\begin{proof}
Use the \(u\)-entry of
\eqref{eq:V73-prefix-endpoint-fraction}:
\[
 p_r^-\le\frac{\Xi_u^{<r}}{\Xi_u}
 \le K\frac{\Xi_4}{\Xi_u}.
\]
\end{proof}

\section{Closure of the low-prefix-mass lateral cell}
\label{sec:V73-closure}

\begin{theorem}[Localized-prefix bridge]
\label{thm:V73-localized-prefix-bridge}
Assume that \(D_0\) is the original root-owned comparable
row-\(3\)--row-\(u\) bank, \(u\in\{1,2\}\), and that
\(\mathcal M_{D_0}\) is either its physical protected envelope or its
positive non-all-protected envelope from
\Cref{thm:V60-ownership-compensation}.  Let
\(\mathcal N_F\) be any commuting positive contraction occurring in
the same outside-local resolved word.  On the cell
\eqref{eq:V73-low-prefix-cell},
\begin{equation}
 \boxed{
 \kappa_{3|c}\!\left(
 \widetilde W_{r,\pi,\nu}^*
 \widetilde W_{r,\pi,\nu}\right)
 \le
 C_{s,K}\theta_{3u}(D_0)^2
 \le C_{s,K}(x+z)^2.}
\label{eq:V73-low-prefix-RCQ}
\end{equation}
No terminal protected capacity and no terminal gap are required.
\end{theorem}

\begin{proof}
All positive factors commute with the resolved dressing by
\eqref{eq:V73-localized-dressings}.  Since
\(0\preccurlyeq\mathcal N_F\preccurlyeq I\) and
\(0\preccurlyeq\mathcal M_{D_0}\preccurlyeq I\), the central-square
argument gives
\begin{align*}
 \widetilde W_{r,\pi,\nu}^*
 \widetilde W_{r,\pi,\nu}
 &\preccurlyeq
 C_s
 \left(\frac{p_r^-}{\Xi_3\Xi_4}\right)^2
 \mathcal M_{D_0}^2\mathcal N_F^2\\
 &\preccurlyeq
 C_s
 \left(\frac{p_r^-}{\Xi_3\Xi_4}\right)^2
 \mathcal M_{D_0}.
\end{align*}
Take row-\(3\) product capacity and apply
\eqref{eq:V73-prefix-ratio}:
\[
 \kappa_{3|c}(
 \widetilde W_{r,\pi,\nu}^*
 \widetilde W_{r,\pi,\nu})
 \le
 C_{s,K}
 \frac1{\Xi_3^2\Xi_u^2}
 \kappa_{3|c}(\mathcal M_{D_0}).
\]
The original bank is nonempty, so
\(T_0\ge a_*\), and a nonzero word has
\(\Xi_3\ge a_*\).  With
\(\beta_0=T_0/\Xi_u\),
\[
 \frac1{\Xi_3^2\Xi_u^2}
 \kappa_{3|c}(\mathcal M_{D_0})
 \le
 a_*^{-4}\beta_0^2
 \kappa_{3|c}(\mathcal M_{D_0}).
\]
Apply respectively the protected or gapped line of
\Cref{thm:V60-ownership-compensation}.  This gives the first
inequality in \eqref{eq:V73-low-prefix-RCQ}; the second is
\(\theta_{3u}(D_0)\le x+z\).
\end{proof}

\begin{corollary}[The original one-failure branch is closed on the
same cell]
\label{cor:V73-original-gapped-cell}
The original-gapped residual left in
\Cref{fire:V70-claim-boundary,fire:V71-claim-boundary} satisfies
\(RCQ_3(\eta)\) whenever
\(\Xi_u^{<r}\le K\Xi_4\).
\end{corollary}

\begin{proof}
Take the positive non-all-protected envelope as
\(\mathcal M_{D_0}\) in
\Cref{thm:V73-localized-prefix-bridge}.  Its entanglement-safe
construction is supplied by V68, and its ownership-compensation
estimate is the gapped line of V60.
\end{proof}

\begin{corollary}[The all-protected discrete concentration branch is
closed on the same cell]
\label{cor:V73-protected-cell}
Even when
\({\mathfrak D}_3^{\rm new}(F)\) fails
\eqref{eq:V71-dimension-bridge}, the all-protected original/terminal
word satisfies \(RCQ_3(\eta)\) on
\(\Xi_u^{<r}\le K\Xi_4\).
\end{corollary}

\begin{proof}
Use the protected envelope for \(\mathcal M_{D_0}\) and the terminal
protected projection for \(\mathcal N_F\) in
\Cref{thm:V73-localized-prefix-bridge}.  The proof uses only
\(\|\mathcal N_F\|\le1\), so no dimension threshold is needed.
\end{proof}

\section{The spectator obstruction is removed at its source}
\label{sec:V73-spectators}

\begin{proposition}[Later \(12\)-spectators shrink the actual prefix
fraction]
\label{prop:V73-spectator-repair}
Keep a fixed nonzero source word with joining coordinate \(r\), and
append after \(r\) a protected \(12\)-spectator bank of row-\(1\) and
row-\(2\) mass \(A\).  Then \(H_{12}^{<r}\) is unchanged, whereas
\begin{equation}
 p_r^-
 =
 O(A^{-2})
\label{eq:V73-spectator-prefix-decay}
\end{equation}
when both appended row masses are comparable to \(A\).  By contrast,
the enlarged global fraction \(p=H_{12}/(\Xi_1\Xi_2)\) can remain of
order one.
\end{proposition}

\begin{proof}
The support of \(H_{12}^{<r}\) is fixed.  The two denominators in
\eqref{eq:V73-local-prefix-fraction} each grow comparably to \(A\),
which proves \eqref{eq:V73-spectator-prefix-decay}.  If a single
appended coordinate has comparable row masses \(A\), its contribution
to the global stock \(H_{12}\) is comparable to \(A^2\), so the
global fraction need not decay.
\end{proof}

\begin{assessment}[Resolution of the V71--V72 spectator no-go]
\label{ass:V73-spectator-resolution}
The physical spectators in
\Cref{prop:V71-diagonal-no-go,prop:V72-internal-spectators} correctly
show that neither a diagonal protected factor nor the first-escape
telescope sees internal \(12\) mass.  They do not obstruct the
complete normalized word.  That word sees the mass through the two
global input denominators while its actual prefix numerator stays
fixed, exactly as
\eqref{eq:V73-spectator-prefix-decay} records.  The apparent
obstruction was created by the premature enlargement
\(H_{12}^{<r}\le H_{12}\).
\end{assessment}

\section{The exact residual after localization}
\label{sec:V73-residual}

\begin{corollary}[Cofinal prefix mass is necessary in every remaining
lateral word]
\label{cor:V73-cofinal-prefix-residual}
Fix a sufficiently large constant \(K\) used in the V63
mass-ratio split.  Every outside-local lateral word not already closed
by V66, V70--V71, or
\Cref{thm:V73-localized-prefix-bridge} satisfies
\begin{equation}
 \boxed{
 \Xi_u^{<r}>K\Xi_4,
 \qquad u\in\{1,2\}.}
\label{eq:V73-cofinal-prefix-residual}
\end{equation}
In particular, the mass forcing the lateral escalation is already
present before the joining coordinate and cannot be placed entirely
in the suffix.
\end{corollary}

\begin{proof}
The complement of
\eqref{eq:V73-cofinal-prefix-residual} is
\eqref{eq:V73-low-prefix-cell}, which is closed by
\Cref{thm:V73-localized-prefix-bridge}.  The cited earlier theorems
remove terminal gaps, sufficiently large terminal protected
dimension products, and the already closed ownership cells.
\end{proof}

\begin{assessment}[New upstream object]
\label{ass:V73-new-upstream-object}
The unresolved mass is no longer an arbitrary terminal spectator
stock.  It is cofinal row-\(u\) mass inside the actual binary prefix
\[
 K_{12}^{<r,\mathrm{conn}}.
\]
The next proof must keep the localized scalar \(p_r^-\) and resolve
that prefix bank before replacing it by its global stock.  A
protected/gapped split internal to the complete binary covariance is
the first legitimate candidate.  Cross-position cancellation should
be revisited only for the prefix subcell which survives that split.
\end{assessment}

\section{V73 result ledger and V74 acquisition order}
\label{sec:V73-status}

\begin{theorem}[Third-series V73 result ledger]
\label{thm:V73-results}
Relative to V1--V72, V73 proves:
\begin{enumerate}[label=\textup{(V73.\arabic*)},leftmargin=6.3em]
\item identification of the global-prefix-stock enlargement as the
      source of the later-spectator obstruction;
\item the support-local binary-prefix estimate
      \eqref{eq:V73-prefix-stock};
\item exact localized resolved-word factorization with scalar
      \(p_r^-/(\Xi_3\Xi_4)\);
\item the deterministic bound of \(p_r^-\) by either prefix endpoint
      mass fraction;
\item the localized-prefix bridge for both protected and physical
      one-failure original bank sectors;
\item closure of the low-prefix-mass part of the discrete
      concentration residual without a terminal capacity;
\item exact repair of the V71--V72 internal-spectator no-go; and
\item reduction of every remaining lateral word to cofinal
      row-\(u\) mass inside the actual binary prefix.
\end{enumerate}
\end{theorem}

\begin{proof}
These are
\Cref{ass:V73-direction-change,
lem:V73-prefix-stock,
lem:V73-localized-factorization,
lem:V73-prefix-endpoint-fraction,
cor:V73-prefix-ratio,
thm:V73-localized-prefix-bridge,
cor:V73-original-gapped-cell,
cor:V73-protected-cell,
prop:V73-spectator-repair,
ass:V73-spectator-resolution,
cor:V73-cofinal-prefix-residual}.
\end{proof}

\begin{programme}[V74: protected/gapped resolution inside the actual
binary prefix]
\label{prog:V73-V74}
\begin{enumerate}[leftmargin=3em]
\item On the residual
      \eqref{eq:V73-cofinal-prefix-residual}, select a complete
      unmarked comparable subbank of the row-\(u\) prefix mass using
      the V57 finite escalation, but do so inside
      \(\{e:e<r\}\).
\item Apply the simultaneous restriction identity to the complete
      binary covariance before taking its norm; do not treat it as a
      positive moment.
\item Determine whether the selected bank occurs in both covariance
      endpoints.  If it does, use the exact endpoint difference to
      prevent duplicate protected charging.
\item Close any prefix-gapped subbank by ordinary heat norm, retaining
      \(p_r^-\) rather than restoring \(p\).
\item On the all-protected prefix subbank, compute its cut alignment
      across row \(3|3^c\).  Since the prefix acts on rows \(1,2\)
      only, expect neutrality and seek the missing gain in the
      endpoint covariance difference, not in its capacity.
\item Keep the V72 first-escape and no-escape kernels dormant until
      this upstream prefix split is complete.  The compulsory
      companion audit is due in V75.
\end{enumerate}
\end{programme}

\begin{firewall}[Claim boundary after third-series V73]
\label{fire:V73-claim-boundary}
V73 closes every outside-local lateral word whose terminal lateral
row has at most \(K\Xi_4\) mass before its joining coordinate,
including protected and original-gapped cells.  It does not close the
cofinal-prefix residual, finite exceptional or local-payer branches,
or full \(RCQ_3(\eta)\).  The V72 aggregate response card, other
quotient cuts, quartic source families, higher MTA, WUFC, CPM, cutoff
removal, reflection positivity, Osterwalder--Schrader reconstruction,
construction of the continuum Yang--Mills measure, and a uniform
positive continuum spectral gap remain open.  The full Yang--Mills
mass gap has not been proved.
\end{firewall}

\paragraph{Parent-version record.}
V73 was formed from
\texttt{Renewal\_Yang\_Mills\_Mass\_Gap\_Research\_Notes\_V72.tex}.
The V72 parent had \(35\,415\) logical source lines,
\(1\,268\,179\) bytes, and SHA--256
\[
 \texttt{6e11f34cb48522e73cac9f841ebcf7d8875126ee29f7e5bcb21d923d6f756a69}.
\]
The next compulsory companion audit is V75.

% =====================================================================
% V74 APPENDIX: UNIFORM WORD CAP AND PURE-CUT CLOSURE
% =====================================================================

\chapter{V74: The Uniform Complete-Word Cap Closes Every Pure
Quotient-Cut Face}
\label{ch:V74-uniform-cap}

\section{A stronger estimate already contained in V10 and V55}
\label{sec:V74-cap}

The V73 localized response is valid, but the residual it leaves asks
for an estimate which is already implicit in the minimum bounds used
in V55.  At fixed heat time, the complete binary covariance has a
support-uniform cap in both its unpaid and paid forms.  The V55 proof
also gives a support-uniform cap for every other possible payer
location.  Keeping the cap branch, rather than the response branch,
eliminates the entire pure row-\(3\) face.

\begin{lemma}[Uniform complete selected-word norm]
\label{lem:V74-uniform-word}
Fix \(s=s_\alpha>0\), a joining coordinate \(r\), and any admissible
payer placement \(\pi\).  The unnormalized, physically compressed,
complete selected-coordinate word of
\eqref{eq:V49-complete-word} satisfies
\begin{equation}
 \boxed{\|\widehat W_{r,\pi}\|\le C_s.}
\label{eq:V74-uniform-word}
\end{equation}
The constant is independent of support, representations, physical
multiplicities, and \(r\).  The same bound holds for the adjoint
orientation.
\end{lemma}

\begin{proof}
Keep the complete binary prefix covariance as one factor.  If it is
unpaid, \eqref{eq:V10-prefix-unpaid} gives
\[
 \|{\cal N}_{<r}^{12}\|
 \le C\min\{1,sH_{12}(<r)\}\le C.
\]
If it carries the unique payer,
\eqref{eq:V10-prefix-paid} gives
\[
 \|{\cal P}_{<r}^{12}\|
 \le C\min\{H_{12}(<r),s^{-1}\}\le C_s.
\]

In every other payer placement the prefix is unpaid and has the first
uniform bound.  The complete paid joining factor, singleton bank, or
suffix has norm at most \(C_s\), exactly as recorded in the proof of
\Cref{lem:V55-complete-prefix-response}.  Every other complete moment
is a contraction, the assembled third cumulant has norm at most six,
and the fixed separators and physical compressions are contractions.
Multiplication gives \eqref{eq:V74-uniform-word}.  Taking adjoints
does not change operator norm.
\end{proof}

\begin{remark}[Relation to the V55 response bound]
\label{rem:V74-two-branches}
The same proof yields both
\[
 \|\widehat W_{r,\pi}\|
 \le C_sH_{12}(<r)
 \quad\hbox{and}\quad
 \|\widehat W_{r,\pi}\|\le C_s.
\]
V55 kept the first branch to match a depleted prefix edge.  V74 keeps
the second branch on a quotient-cut face.  They are two outputs of one
prefix theorem, not two charges of the covariance.
\end{remark}

\section{A nonzero joining word supplies a local cut floor}
\label{sec:V74-local-floor}

\begin{lemma}[All four input rows are active at a nonzero joining
coordinate]
\label{lem:V74-four-active}
If
\(\widehat W_{r,\pi}\ne0\), then every row \(i\in[4]\) is active at
\(r\), and hence
\begin{equation}
 \Xi_i\ge a_{i,r}\ge a_*.
\label{eq:V74-active-floor}
\end{equation}
Moreover, for every \(i\ne j\),
\begin{equation}
 H_{ij}\ge a_{i,r}a_{j,r}\ge a_*^2.
\label{eq:V74-pair-floor}
\end{equation}
\end{lemma}

\begin{proof}
As observed immediately after
\eqref{eq:V55-normalized-word}, the replacement difference belonging
to an inactive row vanishes.  The assembled joining cumulant is
multilinear in the four row occurrences, so a nonzero joining word
has all four active.  Apply the nontrivial mass floor
\eqref{eq:V1-floor}.  The selected coordinate \(r\) occurs in each
total mass and in each pair stock, giving the two displayed
inequalities.
\end{proof}

\begin{lemma}[Inverse-volume domination by any quotient cut]
\label{lem:V74-inverse-volume-cut}
For a nonzero selected word,
\begin{align}
 \frac1{\Xi_1\Xi_2\Xi_3\Xi_4}
 &\le a_*^{-4}(x+z),
\label{eq:V74-row3-volume}\\
 \frac1{\Xi_1\Xi_2\Xi_3\Xi_4}
 &\le a_*^{-4}(y+z),
\label{eq:V74-row4-volume}\\
 \frac1{\Xi_1\Xi_2\Xi_3\Xi_4}
 &\le a_*^{-4}(x+y).
\label{eq:V74-composite-volume}
\end{align}
\end{lemma}

\begin{proof}
For the row-\(3\) cut,
\[
 x+z\ge\theta_{13}
 =\frac{H_{13}}{\Xi_1\Xi_3}.
\]
Therefore
\[
 \frac1{\Xi_1\Xi_2\Xi_3\Xi_4}
 =
 \frac{\theta_{13}}{\Xi_2\Xi_4H_{13}}
 \le a_*^{-4}\theta_{13}
 \le a_*^{-4}(x+z),
\]
where \Cref{lem:V74-four-active} gives
\(\Xi_2\Xi_4H_{13}\ge a_*^4\).
For the row-\(4\) cut use
\(\theta_{14}\le y+z\) and
\(\Xi_2\Xi_3H_{14}\ge a_*^4\).  For the composite cut use again
\(\theta_{13}\le x+y\) and the first calculation.
\end{proof}

\section{Unconditional closure of the fixed-position row-\(3\) card}
\label{sec:V74-RCQ3}

\begin{theorem}[The full selected-coordinate \(RCQ_3\) estimate]
\label{thm:V74-RCQ3}
For every fixed \(s>0\), every selected joining coordinate \(r\),
every physical representation assignment, and every admissible payer
placement,
\begin{equation}
 \boxed{
 \kappa_{3|c}\!\left(
 \widetilde W_{r,\pi}^*\widetilde W_{r,\pi}\right)
 \le C_s(x+z)^2.}
\label{eq:V74-RCQ3}
\end{equation}
The estimate holds without the pure-face assumptions and with the
same conclusion for
\(\widetilde W_{r,\pi}\widetilde W_{r,\pi}^*\).
\end{theorem}

\begin{proof}
If the word is zero there is nothing to prove.  Otherwise
\Cref{lem:V74-uniform-word,lem:V74-inverse-volume-cut} give
\[
 \|\widetilde W_{r,\pi}\|
 =
 \frac{\|\widehat W_{r,\pi}\|}
 {\Xi_1\Xi_2\Xi_3\Xi_4}
 \le C_sa_*^{-4}(x+z).
\]
Product-state capacity of a positive operator is no larger than its
ordinary norm, so
\[
 \kappa_{3|c}(
 \widetilde W_{r,\pi}^*\widetilde W_{r,\pi})
 \le
 \|\widetilde W_{r,\pi}\|^2
 \le C_s(x+z)^2.
\]
The left-square orientation has the same ordinary norm.
\end{proof}

\begin{corollary}[The pure row-\(3\) face of \(FCQ_{211}\) is closed]
\label{cor:V74-row3-FCQ}
On \(\mathfrak R_{3,\eta}\), every complete selected-coordinate
\(2+1+1\) packet and every payer placement satisfies
\begin{equation}
 \boxed{
 \kappa_\otimes(
 \widetilde W_{r,\pi}^*\widetilde W_{r,\pi})
 \le C_{s,\eta}\mathscr A_{211}^{\,2}.}
\label{eq:V74-row3-FCQ}
\end{equation}
Thus the pure row-\(3\) face is removed from the \(FCQ_{211}\)
counterpacket list.
\end{corollary}

\begin{proof}
This is \Cref{thm:V74-RCQ3} followed by the upper comparison
\(x+z\le\eta^{-2}\mathscr A_{211}\) in
\Cref{thm:V56-atlas-equivalence}.
\end{proof}

\begin{assessment}[The V57--V73 bank route is bypassed on this face]
\label{ass:V74-bank-bypass}
The protected/gapped bank theorems V57--V73 remain valid and record
useful cut-aligned information.  They are not needed to prove the
fixed-position card \eqref{eq:V74-RCQ3}.  The decisive observation is
that a pure quotient-cut target asks for one small cut factor at the
operator level; its square is supplied only after taking the genuine
operator square.  The uniform complete-word cap supplies that one
factor through the four input denominators and one actual joining
edge.  No invariant projection is charged twice.
\end{assessment}

\begin{corollary}[The local-payer row-\(3\) square is also closed]
\label{cor:V74-local-payer}
The conclusion \eqref{eq:V74-RCQ3} includes the placement in which
the joining carrier itself contains the unique payer.
\end{corollary}

\begin{proof}
That placement is one of the cases in
\Cref{lem:V74-uniform-word}.  Ordinary norm controls both square
orientations, so no separate extraction of the payer from the local
replica polarization is required.
\end{proof}

\section{The two other pure quotient cuts}
\label{sec:V74-other-cuts}

Define
\begin{align}
 \mathfrak R_{4,\eta}
 &:=
 \{p\ge\eta,\ x\ge\eta,\ y+z<\eta\},
\label{eq:V74-row4-face}\\
 \mathfrak R_{12,\eta}
 &:=
 \{p\ge\eta,\ z\ge\eta,\ x+y<\eta\}.
\label{eq:V74-composite-face}
\end{align}

\begin{theorem}[Uniform closure of all three pure quotient-cut faces]
\label{thm:V74-three-pure-cuts}
For every selected-coordinate packet and every payer placement,
\begin{align}
 \kappa_\otimes(
 \widetilde W_{r,\pi}^*\widetilde W_{r,\pi})
 &\le C_s(y+z)^2
 &&\text{on all data},
\label{eq:V74-row4-card}\\
 \kappa_\otimes(
 \widetilde W_{r,\pi}^*\widetilde W_{r,\pi})
 &\le C_s(x+y)^2
 &&\text{on all data}.
\label{eq:V74-composite-card}
\end{align}
Consequently,
\begin{equation}
 \boxed{
 \kappa_\otimes(
 \widetilde W_{r,\pi}^*\widetilde W_{r,\pi})
 \le C_{s,\eta}\mathscr A_{211}^{\,2}}
\label{eq:V74-other-pure-FCQ}
\end{equation}
on each of
\(\mathfrak R_{4,\eta}\) and
\(\mathfrak R_{12,\eta}\).
\end{theorem}

\begin{proof}
Use respectively
\eqref{eq:V74-row4-volume} and
\eqref{eq:V74-composite-volume} in the proof of
\Cref{thm:V74-RCQ3}.

On \(\mathfrak R_{4,\eta}\), the exact atlas factorization gives
\[
 \mathscr A_{211}
 =
 p\{x(y+z)+yz\}
 \ge\eta^2(y+z).
\]
On \(\mathfrak R_{12,\eta}\),
\[
 \mathscr A_{211}
 =
 p\{z(x+y)+xy\}
 \ge\eta^2(x+y).
\]
Insert these lower bounds into
\eqref{eq:V74-row4-card} and
\eqref{eq:V74-composite-card}.
\end{proof}

\begin{corollary}[Only multiple-depletion intersections remain among
the quotient faces]
\label{cor:V74-intersections-only}
The three regions in which exactly one quotient vertex is separated
from two nondepleted complementary connections are closed.  Any
remaining small-atlas quotient packet must approach an intersection
where at least two of \(x,y,z\) vanish and no complementary quotient
edge stays uniformly positive.
\end{corollary}

\begin{proof}
The pure row-\(3\), row-\(4\), and composite regions are
\Cref{cor:V74-row3-FCQ,thm:V74-three-pure-cuts}.  If a pair of
quotient variables vanishes while the third remains bounded below,
one of those regions applies after choosing a fixed smaller
\(\eta\).  Thus failure of all three pure-region hypotheses along a
small-atlas sequence forces the third variable to vanish as well,
after a subsequence.
\end{proof}

\begin{firewall}[Why V74 does not yet prove all of \(FCQ_{211}\)]
\label{fire:V74-not-full-FCQ}
When \(x,y,z\) all tend to zero, the atlas
\[
 p(xy+xz+yz)
\]
is quadratic in the small quotient connections.  The uniform word
cap plus one inverse-volume comparison supplies only one such factor
before the operator is squared, whereas \(FCQ_{211}\) requires the
operator itself to be controlled by the quadratic atlas.  The local
third-cumulant two-edge response, not another use of the uniform cap,
must treat this intersection.  Likewise, intersections involving
\(p\to0\) must retain the binary prefix response already used in V55.
\end{firewall}

\section{V74 result ledger and the V75 audit mandate}
\label{sec:V74-status}

\begin{theorem}[Third-series V74 result ledger]
\label{thm:V74-results}
Relative to V1--V73, V74 proves:
\begin{enumerate}[label=\textup{(V74.\arabic*)},leftmargin=6.3em]
\item a support-, representation-, multiplicity-, coordinate-, and
      payer-uniform norm cap for every complete selected
      \(2+1+1\) word;
\item the four-row activity and pair-stock floors at every nonzero
      joining coordinate;
\item inverse-volume domination by each of the three quotient-cut
      sums;
\item the unconditional fixed-position \(RCQ_3\) estimate in both
      square orientations;
\item closure of the full pure row-\(3\) face, including its
      local-payer placement;
\item symmetric closure of the pure row-\(4\) and pure composite
      quotient-cut faces; and
\item reduction of the quotient residual to multiple-depletion
      intersections.
\end{enumerate}
\end{theorem}

\begin{proof}
These are
\Cref{lem:V74-uniform-word,
lem:V74-four-active,
lem:V74-inverse-volume-cut,
thm:V74-RCQ3,
cor:V74-row3-FCQ,
cor:V74-local-payer,
thm:V74-three-pure-cuts,
cor:V74-intersections-only}.
\end{proof}

\begin{programme}[V75: compulsory audit and the triple-depleted
quotient corner]
\label{prog:V74-V75}
\begin{enumerate}[leftmargin=3em]
\item Perform the compulsory full read-only audit of the newest active
      SM and NS notes and record filenames, lines, bytes, and SHA--256
      fingerprints.
\item Import only proof-order information relevant to a simultaneous
      two-edge response; no SM or NS inequality transfers by analogy.
\item Return to the unsplit local carrier
      \(J_r(12|3|4)\) and retain its already proved ternary
      two-edge minimum before normalization.
\item On \(x,y,z\to0\), compare the local joining stocks with the
      global quotient variables without tensorizing two separately
      enveloped arms.
\item Combine the local quadratic quotient response with the uniform
      binary-prefix cap.  If \(p\to0\) simultaneously, use the
      localized prefix response branch exactly once.
\item Decide the complete triple-depleted and
      prefix-plus-quotient intersection before returning to the V72
      cross-position aggregate.
\end{enumerate}
\end{programme}

\begin{firewall}[Claim boundary after third-series V74]
\label{fire:V74-claim-boundary}
V74 proves the complete selected-coordinate \(RCQ_3\) card and closes
all three pure quotient-cut faces of \(FCQ_{211}\), at every payer
placement.  It does not close the multiple-depletion intersections,
the complete \(2+1+1\) family, other quartic source families, higher
MTA, WUFC, CPM, cutoff removal, reflection positivity,
Osterwalder--Schrader reconstruction, construction of the continuum
Yang--Mills measure, or a uniform positive continuum spectral gap.
The full Yang--Mills mass gap has not been proved.
\end{firewall}

\paragraph{Parent-version record.}
V74 was formed from
\texttt{Renewal\_Yang\_Mills\_Mass\_Gap\_Research\_Notes\_V73.tex}.
The V73 parent had \(35\,899\) logical source lines,
\(1\,284\,653\) bytes, and SHA--256
\[
 \texttt{ee164db3b61bfe1e5a838b86bdf68395f8add448a39f7d0de3b1ee8a84391ada}.
\]
The compulsory companion audit is due in V75.

% =====================================================================
% V75 APPENDIX: COMPANION AUDIT AND NONDEPLETED-PREFIX CLOSURE
% =====================================================================

\chapter{V75: Companion Audit, the Crossed Joining Floor, and the
Simultaneous-Depletion Residual}
\label{ch:V75-audit-crossed-floor}

\section{Compulsory read-only companion audit}
\label{sec:V75-audit}

The scheduled audit was completed before the V75 proof direction was
chosen.  The newest active companion files read to completion were:
\begin{enumerate}[leftmargin=3em]
\item
\texttt{Renewal\_SM\_Einstein\_Continuum\_Research\_Notes\_V57.tex},
with \(20\,455\) logical source lines and \(752\,895\) bytes,
SHA--256
\[
 \texttt{f924b518bda667bb5b7adc644433d27c72529e6ba5ab060cedc74434f3bbdf8c}.
\]
Its new V56--V57 chain retains the exact six trace-log word types
through a weighted third-derivative estimate, then sums hierarchical
star mergers.  At a \(k\)-child merger the unordered-child
\(1/k!\) cancels the local \(k!\) before block positions are summed;
one or two marks leave only polynomial branching factors and do not
change the convergence radius.  The resulting positive good-field
pressure remains conditional on the same collision interface at
every rescaling stage, and the complex phase is explicitly open.
\item
\texttt{Renewal\_NS\_Stability\_Research\_Notes\_V31.tex},
with \(23\,052\) logical source lines and \(887\,537\) bytes,
SHA--256
\[
 \texttt{f1121b62238ad5491bb7cf03635ea9934942edf573ed8faaa9ee79e1d38a6696}.
\]
The file is unchanged since the V70 audit.  Its exact dyadic
formation ladder still shows that removing a flat probability mark
restores one first-moment factor; weaker heat norms merely hide that
factor in a source weight.  Its persistent branch retains the
no-double-charge rule.
\end{enumerate}

\begin{assessment}[Type-correct V75 transfer]
\label{ass:V75-audit-transfer}
Neither a Standard-Model collision kernel nor a Navier--Stokes heat
source is a Yang--Mills operator.  The only transferred information
is proof order:
\begin{enumerate}[label=\textup{(\roman*)},leftmargin=4em]
\item preserve the exact local branching or partition hierarchy until
      its combinatorial factor has cancelled;
\item do not infer a product of two small factors from two separate
      minimum bounds; and
\item charge a prefix response and a quotient response only if both
      occur in the same complete YM word.
\end{enumerate}
The mathematical closure below is proved solely from the YM uniform
word cap, selected-coordinate activity, and the exact quotient-tree
identity.
\end{assessment}

\section{The crossed matching hidden in \(xy\)}
\label{sec:V75-crossed-floor}

\begin{lemma}[Crossed joining-coordinate product floor]
\label{lem:V75-crossed-floor}
For every nonzero selected \(2+1+1\) joining word,
\begin{equation}
 \boxed{
 xy
 \ge
 \theta_{13}\theta_{24}
 \ge
 \frac{a_*^4}
 {\Xi_1\Xi_2\Xi_3\Xi_4}.}
\label{eq:V75-crossed-floor}
\end{equation}
The alternative crossed product
\(\theta_{23}\theta_{14}\) obeys the same lower bound.
\end{lemma}

\begin{proof}
By definition,
\[
 x=(\theta_{13}+\theta_{23}),
 \qquad
 y=(\theta_{14}+\theta_{24}),
\]
so their product contains
\(\theta_{13}\theta_{24}\).  By
\Cref{lem:V74-four-active}, the selected joining coordinate \(r\)
contributes
\[
 H_{13}\ge a_{1,r}a_{3,r}\ge a_*^2,
 \qquad
 H_{24}\ge a_{2,r}a_{4,r}\ge a_*^2.
\]
Therefore
\[
 \theta_{13}\theta_{24}
 =
 \frac{H_{13}H_{24}}
 {\Xi_1\Xi_2\Xi_3\Xi_4}
 \ge
 \frac{a_*^4}
 {\Xi_1\Xi_2\Xi_3\Xi_4}.
\]
Interchange rows \(1,2\) for the alternative product.
\end{proof}

\begin{corollary}[Inverse volume is bounded by the full quotient-tree
polynomial]
\label{cor:V75-volume-quotient}
Put
\begin{equation}
 q:=xy+xz+yz.
\label{eq:V75-q}
\end{equation}
Then every nonzero selected word satisfies
\begin{equation}
 \boxed{
 \frac1{\Xi_1\Xi_2\Xi_3\Xi_4}
 \le a_*^{-4}xy
 \le a_*^{-4}q.}
\label{eq:V75-volume-quotient}
\end{equation}
\end{corollary}

\begin{proof}
Rearrange \eqref{eq:V75-crossed-floor} and use nonnegativity of
\(xz,yz\).
\end{proof}

\section{All quotient geometry closes when the prefix is
nondepleted}
\label{sec:V75-p-nondepleted}

\begin{theorem}[Complete \(p\)-nondepleted \(FCQ_{211}\) closure]
\label{thm:V75-p-nondepleted}
Fix \(0<\eta\le1\).  For every complete selected-coordinate
\(2+1+1\) packet satisfying
\begin{equation}
 p\ge\eta,
\label{eq:V75-p-lower}
\end{equation}
and every admissible payer placement,
\begin{equation}
 \boxed{
 \kappa_\otimes(
 \widetilde W_{r,\pi}^*\widetilde W_{r,\pi})
 \le
 C_{s,\eta}\mathscr A_{211}^{\,2}.}
\label{eq:V75-p-FCQ}
\end{equation}
The result includes the region \(x,y,z\to0\) and both square
orientations.
\end{theorem}

\begin{proof}
\Cref{lem:V74-uniform-word,cor:V75-volume-quotient} give
\[
 \|\widetilde W_{r,\pi}\|
 \le
 \frac{C_s}{\Xi_1\Xi_2\Xi_3\Xi_4}
 \le C_sa_*^{-4}q.
\]
The exact quotient-tree factorization is
\[
 \mathscr A_{211}=pq.
\]
Under \eqref{eq:V75-p-lower},
\(q\le\eta^{-1}\mathscr A_{211}\).  Hence
\[
 \|\widetilde W_{r,\pi}\|
 \le C_{s,\eta}\mathscr A_{211}.
\]
Square the ordinary norm and use capacity bounded by norm.  The
adjoint orientation is identical.
\end{proof}

\begin{remark}[Why the triple-depleted quotient corner is included]
\label{rem:V75-triple-corner}
No quotient variable is assumed bounded below in
\Cref{thm:V75-p-nondepleted}.  The crossed product \(xy\) supplies
the entire inverse input volume even when all of \(x,y,z\) are small.
The only nondepletion used is the separate binary-prefix factor
\(p\ge\eta\), which converts \(q\) into the full atlas \(pq\).
\end{remark}

\begin{corollary}[The V74 pure-cut hypotheses are unnecessary when
\(p\) is nondepleted]
\label{cor:V75-V74-strengthened}
The pure row-\(3\), row-\(4\), and composite-cut conclusions of V74
extend to their full quotient-cut regions, with no lower bound on the
opposite quotient edge, provided \(p\ge\eta\).
\end{corollary}

\begin{proof}
All such packets are contained in the hypothesis of
\Cref{thm:V75-p-nondepleted}.
\end{proof}

\section{The exact remaining \(2+1+1\) corner}
\label{sec:V75-residual}

\begin{theorem}[Simultaneous prefix/quotient depletion is necessary]
\label{thm:V75-simultaneous-depletion}
Let a sequence of complete selected-coordinate \(2+1+1\) packets,
with fixed payer type and fixed \(s>0\), violate every uniform
\(FCQ_{211}\) bound.  After passage to a subsequence,
\begin{equation}
 \boxed{
 p_n\longrightarrow0,
 \qquad
 m_2(x_n,y_n,z_n)\longrightarrow0.}
\label{eq:V75-simultaneous-depletion}
\end{equation}
Consequently at least two of \(x_n,y_n,z_n\) tend to zero along a
further subsequence.
\end{theorem}

\begin{proof}
If \(p_n\not\to0\), pass to a subsequence with
\(p_n\ge\eta>0\).  This contradicts
\Cref{thm:V75-p-nondepleted}.  Hence \(p_n\to0\).

If \(m_2(x_n,y_n,z_n)\not\to0\), pass to a further subsequence with
\(m_2\ge\eta>0\).  The packets then lie eventually in the V55 pure
prefix sector and are closed by
\Cref{thm:V55-prefix-face-closure}, again a contradiction.  Therefore
\(m_2\to0\).  By definition of the second-largest value, at least two
quotient variables tend to zero after selecting their labels.
\end{proof}

\begin{corollary}[Quantitative residual cell]
\label{cor:V75-quantitative-cell}
For every fixed \(0<\delta\le1\), all packets outside
\begin{equation}
 \mathfrak D_\delta
 :=
 \{p<\delta,\ m_2(x,y,z)<\delta\}
\label{eq:V75-double-cell}
\end{equation}
are covered by the union of
\Cref{thm:V75-p-nondepleted} and
\Cref{thm:V55-prefix-face-closure}, with constants depending only on
\(s,\delta\) and the fixed structural data.
\end{corollary}

\begin{proof}
If \(p\ge\delta\), apply
\Cref{thm:V75-p-nondepleted}.  Otherwise, if
\(m_2\ge\delta\), apply V55.
\end{proof}

\section{The missing estimate is a product-coupling card}
\label{sec:V75-product-card}

The two presently valid operator bounds on
\(\mathfrak D_\delta\) have different ancestry.  The V55
prefix-response branch gives
\begin{equation}
 \|\widetilde W_{r,\pi}\|
 \le C_s p,
\label{eq:V75-prefix-only}
\end{equation}
after using the active lower floors for \(\Xi_3,\Xi_4\).  The V75
uniform-volume branch gives
\begin{equation}
 \|\widetilde W_{r,\pi}\|
 \le C_s q.
\label{eq:V75-quotient-only}
\end{equation}
Thus the separate bounds give only
\begin{equation}
 \|\widetilde W_{r,\pi}\|
 \le C_s\min\{p,q\}.
\label{eq:V75-two-minima}
\end{equation}

\begin{definition}[Prefix--quotient product-coupling card]
\label{def:V75-PQC}
The selected-word card \(PQC_{211}\) is
\begin{equation}
 \boxed{
 \|\widetilde W_{r,\pi}\|
 \le C_s\,p\,q
 =
 C_s\mathscr A_{211}.}
\label{eq:V75-PQC}
\end{equation}
\end{definition}

\begin{lemma}[The product-coupling card is sufficient and genuinely
stronger than the two separate minima]
\label{lem:V75-PQC-sufficient}
If \eqref{eq:V75-PQC} holds for every payer placement on
\(\mathfrak D_\delta\), then the complete selected-coordinate
\(2+1+1\) family satisfies \(FCQ_{211}\).  Conversely,
\eqref{eq:V75-two-minima} alone does not imply
\eqref{eq:V75-PQC}.
\end{lemma}

\begin{proof}
Squaring \eqref{eq:V75-PQC} and bounding product-state capacity by
ordinary norm gives the quadratic card.  The complement of
\(\mathfrak D_\delta\) is already closed by
\Cref{cor:V75-quantitative-cell}.

For strictness, take the scalar values \(p=q=\varepsilon\).  Then
\(\min\{p,q\}=\varepsilon\), whereas \(pq=\varepsilon^2\).  No
constant independent of \(\varepsilon\) bounds the first by the
second.  Hence the two valid source bounds cannot merely be written
next to each other and called a product.
\end{proof}

\begin{assessment}[The correct post-audit bottleneck]
\label{ass:V75-bottleneck}
The \(2+1+1\) problem is no longer any single quotient face, protected
capacity, terminal mass ratio, or cross-position sign.  It is the
simultaneous-depletion product card
\eqref{eq:V75-PQC}.  A proof must retain in one complete word:
\begin{enumerate}[label=\textup{(\roman*)},leftmargin=4em]
\item the exact binary prefix covariance which supplies \(p\);
\item the unsplit composite joining cumulant which supplies two
      quotient edges; and
\item the sole payer and physical compression.
\end{enumerate}
This conclusion agrees with the proof-order part of the SM and NS
audits, but no companion estimate has been imported.
\end{assessment}

\section{V75 result ledger and V76 acquisition order}
\label{sec:V75-status}

\begin{theorem}[Third-series V75 result ledger]
\label{thm:V75-results}
Relative to V1--V74, V75 proves:
\begin{enumerate}[label=\textup{(V75.\arabic*)},leftmargin=6.3em]
\item the compulsory full SM V57 and NS V31 audit with exact
      fingerprints;
\item the type-correct hierarchical/one-charge transfer rule;
\item the crossed joining-coordinate product floor
      \eqref{eq:V75-crossed-floor};
\item domination of inverse input volume by the full quotient-tree
      polynomial \(q\);
\item complete closure of every \(p\)-nondepleted quotient geometry,
      including the triple-depleted quotient corner;
\item reduction of every possible \(2+1+1\) counterpacket to
      simultaneous \(p\)- and second-quotient depletion;
\item the quantitative residual cell \(\mathfrak D_\delta\); and
\item the exact prefix--quotient product-coupling card, together with
      proof that the two separate minimum bounds do not imply it.
\end{enumerate}
\end{theorem}

\begin{proof}
These are
\Cref{sec:V75-audit,
ass:V75-audit-transfer,
lem:V75-crossed-floor,
cor:V75-volume-quotient,
thm:V75-p-nondepleted,
thm:V75-simultaneous-depletion,
cor:V75-quantitative-cell,
def:V75-PQC,
lem:V75-PQC-sufficient}.
\end{proof}

\begin{programme}[V76: one-word prefix--quotient coupling]
\label{prog:V75-V76}
\begin{enumerate}[leftmargin=3em]
\item Return to
      \(K_{12}^{<r,\mathrm{conn}}\otimes
      J_r(12|3|4)\) before physical compression.
\item Insert the exact prefix telescope
      \eqref{eq:V10-prefix-telescope} and the unsplit centered or
      replica formula for the joining third cumulant in the same
      word.  Do not multiply their positive envelopes after separate
      compression.
\item Normalize the prefix occurrence by \(\Xi_1\Xi_2\) and the
      quotient occurrence by \(\Xi_3\Xi_4\), keeping all local
      \(a_{1,r},a_{2,r}\) factors which may compensate a diluted
      joining coordinate.
\item Split only after this assembly by whether the prefix response
      coordinate and the joining coordinate share a dominant row.
\item Test the resulting card on the V51 equal-spin and V53 protected
      dilution packets and on a new simultaneous-depletion scalar
      family.
\item If a same-row center still loses a denominator, return to the
      exact V44 doubled-prefix square; do not claim a product from
      \eqref{eq:V75-two-minima}.
\item The next compulsory companion audit is V80.
\end{enumerate}
\end{programme}

\begin{firewall}[Claim boundary after third-series V75]
\label{fire:V75-claim-boundary}
V75 closes the entire \(p\)-nondepleted part of the selected
\(2+1+1\) source and reduces the remainder to the explicit
prefix--quotient product-coupling card.  It does not prove that card,
close the complete \(2+1+1\) source, its unselected/global
aggregation, other quartic source families, higher MTA, WUFC, CPM,
cutoff removal, reflection positivity, Osterwalder--Schrader
reconstruction, construction of the continuum Yang--Mills measure,
or a uniform positive continuum spectral gap.  The full Yang--Mills
mass gap has not been proved.
\end{firewall}

\paragraph{Parent-version record.}
V75 was formed from
\texttt{Renewal\_Yang\_Mills\_Mass\_Gap\_Research\_Notes\_V74.tex}.
The V74 parent had \(36\,314\) logical source lines,
\(1\,298\,260\) bytes, and SHA--256
\[
 \texttt{11dc1fa35ffe760c90c6d05f6f23afff880115e5035495a19d5a7a2fb7979313}.
\]
The next compulsory companion audit is V80.

% =====================================================================
% V76 APPENDIX: PREFIX--BANK COUPLING ON A DOMINANT QUOTIENT EDGE
% =====================================================================

\chapter{V76: A Genuine Prefix--Bank Product and the Dominant-Edge
Reduction}
\label{ch:V76-prefix-bank}

\section{Why the next split is by the largest quotient edge}
\label{sec:V76-largest-edge}

Retain
\[
 p=\theta_{12},\qquad
 q=xy+xz+yz,\qquad
 \mathscr A_{211}=pq.
\]
Put
\begin{equation}
 m:=\max\{x,y,z\}.
\label{eq:V76-m}
\end{equation}
Choose a dominant label \(\ell\in\{x,y,z\}\), using the order
\(x\prec y\prec z\) to break ties, and define its complementary cut
sum by
\begin{equation}
 c_\ell:=
 \begin{cases}
  y+z,&\ell=x,\\
  x+z,&\ell=y,\\
  x+y,&\ell=z.
 \end{cases}
\label{eq:V76-complementary-cut}
\end{equation}
Thus \(c_x\) is the row-\(4\) cut, \(c_y\) is the row-\(3\) cut,
and \(c_z\) is the cut between the composite block \(\{1,2\}\) and
the two singleton blocks.

\begin{lemma}[Largest-edge factorization of the quotient tree]
\label{lem:V76-largest-edge}
For all \(x,y,z\ge0\),
\begin{equation}
 \boxed{
 mc_\ell\le q\le\frac32mc_\ell.}
\label{eq:V76-largest-edge}
\end{equation}
Consequently, on \(m\ge\eta>0\),
\begin{equation}
 pc_\ell
 \le\eta^{-1}\mathscr A_{211}.
\label{eq:V76-p-cut-by-atlas}
\end{equation}
\end{lemma}

\begin{proof}
Write the two nondominant variables as \(a,b\).  Then
\[
 q=m(a+b)+ab,
 \qquad
 c_\ell=a+b.
\]
The lower bound is immediate.  Since \(a,b\le m\),
\[
 2ab\le m(a+b);
\]
hence \(q\le\frac32m(a+b)\).  If \(m\ge\eta\), then
\(\mathscr A_{211}=pq\ge\eta pc_\ell\).
\end{proof}

\begin{remark}[This is the correct two-factor target on the
dominant cell]
\label{rem:V76-correct-target}
When \(m\ge\eta\), the three-factor atlas \(pq\) is, up to the fixed
factor \(\eta^{-1}\), the product of the binary-prefix edge \(p\)
and the cut opposite the dominant quotient edge.  Thus one need not
manufacture the whole polynomial \(q\) from one operator occurrence:
it is enough to retain \(p\) and a genuine quadratic currency for
\(c_\ell\) in the same complete square.
\end{remark}

\section{Row-symmetric owned banks and compatible cuts}
\label{sec:V76-compatible-bank}

For an unordered row pair \(a b\), let
\[
 H_{ab}(D):=\sum_{e\in D}a_{a,e}a_{b,e},
 \qquad
 \theta_{ab}(D):=
 \frac{H_{ab}(D)}{\Xi_a\Xi_b}.
\]
The V57 definition of an owned comparable bank and its protected and
gapped envelopes are transported from \(3u\) to \(ab\) by permuting
the four row labels.

\begin{lemma}[Row-permutation form of the V57--V59 bank theorem]
\label{lem:V76-row-symmetric-bank}
Fix ownership and comparability constants.  If \(D\) is an unmarked
owned comparable \(a\)--\(b\) bank, disjoint from the selected and
payer coordinates, and \(\mathcal M_D\) is one of its complete
protected or positive gapped envelopes, then
\begin{equation}
 0\preccurlyeq\mathcal M_D\preccurlyeq I,
 \qquad
 \kappa_\otimes(\mathcal M_D)
 \le C_s\theta_{ab}(D)^2.
\label{eq:V76-symmetric-bank}
\end{equation}
Moreover \(\mathcal M_D\) commutes with all coordinatewise
equivariant physical source separators meeting \(D\).
\end{lemma}

\begin{proof}
Relabel rows \(a,b\) as \(3,u\) in
\Cref{thm:V57-bank-quadratic,cor:V59-V32-central}.  The proofs of
those results use only pointwise comparison of the two row weights,
ownership of the two row masses, complete-bank tensorization,
isotypic centrality, and product-state capacity.  All these
properties are invariant under a permutation of the four tensor
factors.  Conjugation by the corresponding tensor-factor permutation
therefore gives \eqref{eq:V76-symmetric-bank} and the commutation
statement with the same structural constants.
\end{proof}

Define the compatible edge sets
\begin{equation}
 \mathcal E_x=\{14,24,34\},\qquad
 \mathcal E_y=\{13,23,34\},\qquad
 \mathcal E_z=\{13,23,14,24\}.
\label{eq:V76-compatible-edges}
\end{equation}

\begin{lemma}[Every compatible bank lies below the complementary
cut]
\label{lem:V76-bank-below-cut}
If \(ab\in\mathcal E_\ell\), then
\begin{equation}
 \theta_{ab}(D)\le\theta_{ab}\le c_\ell.
\label{eq:V76-bank-below-cut}
\end{equation}
\end{lemma}

\begin{proof}
The bank stock is a sub-sum of \(H_{ab}\), which proves the first
inequality.  For \(\ell=x\),
\[
 c_x=y+z=\theta_{14}+\theta_{24}+\theta_{34}.
\]
For \(\ell=y\),
\[
 c_y=x+z=\theta_{13}+\theta_{23}+\theta_{34}.
\]
For \(\ell=z\),
\[
 c_z=x+y
 =\theta_{13}+\theta_{23}+\theta_{14}+\theta_{24}.
\]
Nonnegativity proves the second inequality in all cases.
\end{proof}

\section{Keeping the prefix response during central insertion}
\label{sec:V76-prefix-preserving}

The V59 physical insertion theorem deliberately absorbed every
unused row denominator into the active-row floor because it was
targeting one quotient cut.  For the full atlas this is too early:
the genuine binary-prefix response must remain visible.

\begin{lemma}[All-payer prefix-preserving bank factorization]
\label{lem:V76-prefix-bank-factorization}
Fix a joining coordinate \(r\), a payer placement \(\pi\), and one
term \(\nu\) in the exact V58 restriction resolution.  Suppose the
term exposes a complete unmarked bank envelope \(\mathcal M_D\)
which is disjoint from \(r\) and from the payer coordinate.  Then,
after commuting only the central factor \(\mathcal M_D\), the
normalized term has the exact presentation
\begin{equation}
 \boxed{
 \widetilde W_{r,\pi,\nu}
 =
 \frac{p_r^-}{\Xi_3\Xi_4}\,
 L_{r,\pi,\nu}\mathcal M_DQ_{r,\pi,\nu},}
\label{eq:V76-prefix-bank-factorization}
\end{equation}
where
\[
 p_r^-=\frac{H_{12}^{<r}}{\Xi_1\Xi_2},
 \qquad
 \|L_{r,\pi,\nu}\|+\|Q_{r,\pi,\nu}\|\le C_s,
 \qquad
 [\mathcal M_D,Q_{r,\pi,\nu}]=0.
\]
The constants are uniform in the support, representations,
physical multiplicities, \(r\), and \(\pi\).  If
\(H_{12}^{<r}=0\), the resolved word vanishes.  The statement holds
for both square orientations.
\end{lemma}

\begin{proof}
The coefficient-one V10 covariance telescope on the actual prefix
support gives
\[
 \|K_{12}^{<r,\mathrm{conn}}\|
 \le C_sH_{12}^{<r}.
\]
If the prefix carries the payer, its paid estimate is instead
\[
 \|K_{12}^{<r,\mathrm{conn,pay}}\|
 \le C\min\{H_{12}^{<r},s^{-1}\}
 \le CH_{12}^{<r}.
\]
For every other payer placement, including the selected joining
coordinate, the prefix is unpaid and the complete once-paid factor
has the support-uniform \(C_s\) norm used in
\Cref{lem:V55-complete-prefix-response,lem:V74-uniform-word}.
All other complete moments are contractions, while the fixed local
cumulant and physical separators have their uniform norms.

The V58 resolution is an exact finite moment/covariance identity.
Each resolved prefix covariance is supported either in \(D\) or in
\(D^c\), and its response stock is a sub-sum of
\(H_{12}^{<r}\).  Hence the preceding bound applies term by term,
not merely after recombination.  When \(H_{12}^{<r}>0\), divide the
resolved prefix-side operator by that scalar and absorb it, together
with all uniformly bounded factors to the left of the exposed bank,
into \(L_{r,\pi,\nu}\).

By \Cref{lem:V76-row-symmetric-bank}, the complete bank envelope is
central for every physical source intertwiner meeting \(D\); factors
on \(D^c\) commute tensorially.  It may therefore be moved to the
displayed position, leaving the order of all noncentral factors
unchanged.  Division by the four input masses gives
\[
 \frac{H_{12}^{<r}}
 {\Xi_1\Xi_2\Xi_3\Xi_4}
 =
 \frac{p_r^-}{\Xi_3\Xi_4},
\]
which proves \eqref{eq:V76-prefix-bank-factorization}.  If the
prefix stock vanishes, its norm bound forces the resolved word to
vanish.  Taking adjoints gives the opposite orientation.
\end{proof}

\begin{firewall}[Two actual currencies, not two uses of the prefix]
\label{fire:V76-two-currencies}
The scalar \(p_r^-\) comes from the unique complete binary-prefix
covariance.  The second small factor below comes from the independent
complete-bank capacity theorem
\eqref{eq:V76-symmetric-bank}.  The adjoint occurrence in the square
does not create a second payer or a second prefix response.  Thus the
argument respects
\Cref{thm:V44-one-debit,fire:V55-one-prefix-use}.
\end{firewall}

\begin{theorem}[One-word prefix--bank product coupling]
\label{thm:V76-prefix-bank-product}
Under the hypotheses of
\Cref{lem:V76-prefix-bank-factorization}, suppose
\(\mathcal M_D\) is an owned comparable \(a\)--\(b\) bank with
\(ab\in\mathcal E_\ell\).  Then
\begin{equation}
 \boxed{
 \kappa_\otimes\!\left(
 \widetilde W_{r,\pi,\nu}^*
 \widetilde W_{r,\pi,\nu}\right)
 \le C_s p^2\theta_{ab}(D)^2
 \le C_s p^2c_\ell^2.}
\label{eq:V76-prefix-bank-product}
\end{equation}
The same estimate holds for the left-square orientation.
\end{theorem}

\begin{proof}
Apply the central insertion theorem
\Cref{thm:V59-central-insertion} to
\eqref{eq:V76-prefix-bank-factorization}:
\[
 \kappa_\otimes(
 \widetilde W_{r,\pi,\nu}^*
 \widetilde W_{r,\pi,\nu})
 \le
 C_s
 \left(\frac{p_r^-}{\Xi_3\Xi_4}\right)^2
 \kappa_\otimes(\mathcal M_D).
\]
For a nonzero word,
\(\Xi_3,\Xi_4\ge a_*\) by
\Cref{lem:V74-four-active}; also \(p_r^-\le p\).  Therefore
\[
 \kappa_\otimes(
 \widetilde W_{r,\pi,\nu}^*
 \widetilde W_{r,\pi,\nu})
 \le
 C_sa_*^{-4}p^2\theta_{ab}(D)^2
\]
by \Cref{lem:V76-row-symmetric-bank}.  Absorb \(a_*^{-4}\) into
the fixed constant and use \Cref{lem:V76-bank-below-cut}.  The
left-square proof is the adjoint version of central insertion.
\end{proof}

\section{A new closed dominant-edge cell}
\label{sec:V76-dominant-closure}

\begin{definition}[Compatible-bank-covered selected word]
\label{def:V76-compatible-covered}
Fix ownership/comparability constants and \(\eta>0\).  A selected
word with \(m\ge\eta\) is \emph{compatible-bank-covered} if, after
one common V58 restriction resolution, every nonzero resolved term
contains a protected or positive gapped owned comparable unmarked
bank on an edge in \(\mathcal E_\ell\), disjoint from the selected
and payer coordinates.
\end{definition}

\begin{theorem}[Full \(FCQ_{211}\) bound on the
dominant-edge/compatible-bank cell]
\label{thm:V76-dominant-bank-FCQ}
For fixed \(s>0\), \(\eta>0\), and fixed bank structural constants,
every compatible-bank-covered selected word satisfies
\begin{equation}
 \boxed{
 \kappa_\otimes(
 \widetilde W_{r,\pi}^*\widetilde W_{r,\pi})
 \le
 C_{s,\eta}\mathscr A_{211}^{\,2}.}
\label{eq:V76-dominant-bank-FCQ}
\end{equation}
The conclusion is uniform in support, representations, physical
multiplicities, joining coordinate, and payer placement, and holds
in both square orientations.
\end{theorem}

\begin{proof}
For each resolved term,
\Cref{thm:V76-prefix-bank-product} gives
\[
 \kappa_\otimes(
 \widetilde W_{r,\pi,\nu}^*
 \widetilde W_{r,\pi,\nu})
 \le C_sp^2c_\ell^2.
\]
By \eqref{eq:V76-p-cut-by-atlas},
\[
 p^2c_\ell^2
 \le\eta^{-2}\mathscr A_{211}^{\,2}.
\]
The number of exact resolved terms is bounded by the fixed
\(N_*\) in \Cref{thm:V58-subset-resolution}.  Apply
\Cref{lem:V58-finite-recombination}; its constant is independent of
support.  The argument is unchanged after taking adjoints.
\end{proof}

\begin{corollary}[All three dominant labels are covered by the same
compiler]
\label{cor:V76-three-labels}
The theorem applies as follows:
\begin{enumerate}[label=\textup{(\roman*)},leftmargin=4em]
\item if \(x=m\ge\eta\), any compatible owned bank across the
      row-\(4\) cut supplies \(p(y+z)\);
\item if \(y=m\ge\eta\), any compatible owned bank across the
      row-\(3\) cut supplies \(p(x+z)\);
\item if \(z=m\ge\eta\), any compatible owned bank across the
      composite cut supplies \(p(x+y)\).
\end{enumerate}
No quotient face is privileged by the analytic insertion theorem.
\end{corollary}

\begin{proof}
This is the list
\eqref{eq:V76-compatible-edges} inserted into
\Cref{thm:V76-dominant-bank-FCQ}.
\end{proof}

\section{What a remaining counterpacket must now do}
\label{sec:V76-residual}

\begin{theorem}[Dominant-bank avoidance or total quotient
depletion]
\label{thm:V76-counterpacket}
Let a sequence of complete selected \(2+1+1\) words violate every
uniform \(FCQ_{211}\) bound.  After passage to a subsequence,
\[
 p_n\to0,\qquad m_{2,n}\to0,
\]
and precisely one of the following alternatives remains:
\begin{enumerate}[label=\textup{(\alph*)},leftmargin=4em]
\item \(m_n\to0\), hence
      \(p_n,x_n,y_n,z_n\to0\);
\item \(m_n\ge\eta>0\), one fixed label
      \(\ell\in\{x,y,z\}\) is dominant, and the word is not
      compatible-bank-covered: at least one nonzero term of every
      admissible common finite resolution lacks a compatible bank
      with the fixed structural constants.
\end{enumerate}
\label{eq:V76-counterpacket-alternatives}
\end{theorem}

\begin{proof}
The first two limits are
\Cref{thm:V75-simultaneous-depletion}.  Pass to a subsequence on
which either \(m_n\to0\) or \(m_n\ge\eta>0\).  In the latter case,
pass again so that the tie-broken dominant label is constant.  If
the words were compatible-bank-covered, then
\Cref{thm:V76-dominant-bank-FCQ} would contradict the assumed
failure.  This proves the dichotomy.
\end{proof}

\begin{corollary}[Mass geometry of the totally depleted branch]
\label{cor:V76-total-depletion-masses}
In alternative \textup{(a)}, every normalized pair edge tends to
zero:
\begin{equation}
 \theta_{ij}^{(n)}\longrightarrow0
 \qquad(1\le i<j\le4).
\label{eq:V76-all-edges-zero}
\end{equation}
For nonzero joining words,
\begin{equation}
 \Xi_i^{(n)}\Xi_j^{(n)}\longrightarrow\infty
 \qquad(1\le i<j\le4).
\label{eq:V76-all-pair-masses}
\end{equation}
After a further subsequence, at least three of the four row masses
tend to infinity.
\end{corollary}

\begin{proof}
The edge \(12\) is \(p_n\).  Each of the other five edges is a
nonnegative summand of one of \(x_n,y_n,z_n\), so
\eqref{eq:V76-all-edges-zero} follows.  At a nonzero joining
coordinate,
\Cref{lem:V74-four-active} gives
\[
 \theta_{ij}^{(n)}
 =
 \frac{H_{ij}^{(n)}}
 {\Xi_i^{(n)}\Xi_j^{(n)}}
 \ge
 \frac{a_*^2}
 {\Xi_i^{(n)}\Xi_j^{(n)}}.
\]
Thus every pair product must diverge.  Pass to a subsequence on
which each \(\Xi_i^{(n)}\) converges in the compactified interval
\([a_*,\infty]\).  Two finite limits would contradict the
corresponding pair-product divergence, so at most one is finite.
\end{proof}

\begin{firewall}[Compatible-bank occurrence is not automatic]
\label{fire:V76-bank-not-automatic}
A small cut sum does not by itself create an owned comparable bank
inside the source word.  V57 proves a bank-or-mass-escalation
dichotomy for a specified complete bank, while
\Cref{thm:V76-dominant-bank-FCQ} assumes that a compatible bank is
actually exposed in every resolved term.  Inferring this occurrence
from the scalar inequality
\(\theta_{ab}(D)\le c_\ell\) would reverse the direction of the
argument and is not made here.
\end{firewall}

\section{Regression tests and next acquisition}
\label{sec:V76-regression}

\begin{assessment}[Regression tests]
\label{ass:V76-regression}
\begin{enumerate}[label=\textup{(\roman*)},leftmargin=4em]
\item The V51 equal-spin scalar channel is not declared cancelling:
      the new estimate uses a separate complete bank currency and
      remains silent when no such bank occurs.
\item V53 spectator dilution cannot create a false gain.  A bank
      contributes only when it satisfies the two-sided ownership and
      comparability hypotheses; arbitrary suffix spectators are not
      inserted into \(\mathcal M_D\).
\item The scalar family \(p=q=\varepsilon\) still disproves any
      attempt to infer \(pq\) from two separate minimum bounds.  V76
      instead displays the two factors in one exact resolved word.
\item A local payer is counted once.  Its complete paid norm belongs
      to \(L\) or \(Q\); neither the bank square nor the adjoint copy
      is relabelled as an additional payer.
\end{enumerate}
\end{assessment}

\begin{theorem}[Third-series V76 result ledger]
\label{thm:V76-results}
Relative to V1--V75, V76 proves:
\begin{enumerate}[label=\textup{(V76.\arabic*)},leftmargin=6.3em]
\item the sharp two-sided largest-edge comparison
      \(mc_\ell\le q\le\frac32mc_\ell\);
\item the row-permutation form of the owned-bank quadratic and
      centrality theorems;
\item the exact compatible edge sets for all three quotient labels;
\item an all-payer prefix-preserving central-bank factorization;
\item a genuine one-word product of the binary-prefix response and
      a compatible bank cut currency;
\item full \(FCQ_{211}\) closure on every nondepleted dominant-edge
      cell whose resolved words are compatible-bank-covered;
\item reduction of a remaining counterpacket to compatible-bank
      avoidance or total quotient depletion; and
\item divergence of every pair of row masses, with at least three
      individually divergent rows, on the totally depleted branch.
\end{enumerate}
\end{theorem}

\begin{proof}
These are
\Cref{lem:V76-largest-edge,
lem:V76-row-symmetric-bank,
lem:V76-bank-below-cut,
lem:V76-prefix-bank-factorization,
thm:V76-prefix-bank-product,
thm:V76-dominant-bank-FCQ,
thm:V76-counterpacket,
cor:V76-total-depletion-masses}.
\end{proof}

\begin{programme}[V77: source occurrence on the two residual
branches]
\label{prog:V76-V77}
\begin{enumerate}[leftmargin=3em]
\item On the nondepleted-dominant branch, return to the exact V35
      first-connectivity word and list the complete unmarked banks
      crossing \(\mathcal E_\ell\) before any envelope is taken.
\item Apply the V57 bank-or-escalation theorem only to banks which
      actually occur in that list.  Determine whether failure of a
      compatible owned bank forces a cofinal mass arrow in one of
      the two depleted quotient directions.
\item Keep local-payer and selected coordinates outside every bank,
      and retain the finite terms on which their removal destroys
      ownership.
\item On the totally depleted branch use
      \eqref{eq:V76-all-pair-masses}, not the weaker statement that
      the total input volume diverges.  Split by the unique possibly
      bounded row and test whether the other three rows contain a
      cofinal gapped bank or an all-protected three-row carrier.
\item If the mass-arrow routing revisits a previously used bank,
      return upstream to the unsplit joining cumulant rather than
      charging that bank twice.
\item The next compulsory companion audit remains V80.
\end{enumerate}
\end{programme}

\begin{firewall}[Claim boundary after third-series V76]
\label{fire:V76-claim-boundary}
V76 proves a genuine prefix--quotient product only when a compatible
owned complete bank occurs in the same resolved word, and closes the
corresponding nondepleted-dominant cell.  It does not prove universal
compatible-bank occurrence, close the totally depleted branch, close
the complete \(2+1+1\) source or its global aggregation, or close
other quartic source families, higher MTA, WUFC, CPM, cutoff removal,
reflection positivity, Osterwalder--Schrader reconstruction,
construction of the continuum Yang--Mills measure, or a uniform
positive continuum spectral gap.  The full Yang--Mills mass gap has
not been proved.
\end{firewall}

\paragraph{Parent-version record.}
V76 was formed from
\texttt{Renewal\_Yang\_Mills\_Mass\_Gap\_Research\_Notes\_V75.tex}.
The V75 parent had \(36\,733\) logical source lines,
\(1\,312\,134\) bytes, and SHA--256
\[
 \texttt{e741dc0ebed6b41047ed6873111e7eeb80781f104296b1e055eb5301c4a1df8b}.
\]
The next compulsory companion audit is V80.

% =====================================================================
% V77 APPENDIX: PERFECT-MATCHING CLOSURE AND THE Z-DOMINANT RESIDUAL
% =====================================================================

\chapter{V77: Perfect-Matching Closure and the Exact
\(z\)-Dominant Residual}
\label{ch:V77-matching-z}

\section{The disjoint matching missed by the quotient-cut split}
\label{sec:V77-matching}

V76 splits the quotient polynomial by its largest edge.  Before
searching for compatible banks on all three branches, one should use
the fact that the binary-prefix edge \(12\) has a disjoint quotient
partner: the singleton edge \(34\).  Their product uses every row
mass exactly once.

\begin{lemma}[Perfect-matching inverse-volume floor]
\label{lem:V77-matching-floor}
For every nonzero selected joining word,
\begin{equation}
 \boxed{
 pz=\theta_{12}\theta_{34}
 \ge
 \frac{a_*^4}
 {\Xi_1\Xi_2\Xi_3\Xi_4}.}
\label{eq:V77-matching-floor}
\end{equation}
\end{lemma}

\begin{proof}
At the selected joining coordinate \(r\), all four rows are active
by \Cref{lem:V74-four-active}.  Therefore
\[
 H_{12}\ge a_{1,r}a_{2,r}\ge a_*^2,
 \qquad
 H_{34}\ge a_{3,r}a_{4,r}\ge a_*^2.
\]
Using \(p=H_{12}/(\Xi_1\Xi_2)\) and
\(z=H_{34}/(\Xi_3\Xi_4)\) gives the result.
\end{proof}

\begin{remark}[Why this differs from the V75 crossed floor]
\label{rem:V77-two-matchings}
V75 used the quotient-only matching
\(\theta_{13}\theta_{24}\le xy\).  It bounds inverse volume by
\(q\), leaving the prefix factor \(p\) unpaid.  The matching
\(\theta_{12}\theta_{34}=pz\) already includes the actual prefix
edge.  It becomes decisive precisely when \(x\) or \(y\) supplies a
separate nondepleted factor in \(q\).
\end{remark}

\section{Unconditional closure whenever \(x\) or \(y\) is
nondepleted}
\label{sec:V77-xy-closure}

\begin{theorem}[Perfect-matching \(FCQ_{211}\) closure]
\label{thm:V77-xy-FCQ}
Fix \(0<\eta\le1\).  Every complete selected-coordinate
\(2+1+1\) word satisfying
\begin{equation}
 \max\{x,y\}\ge\eta
\label{eq:V77-xy-nondepleted}
\end{equation}
obeys, for every payer placement and both square orientations,
\begin{equation}
 \boxed{
 \kappa_\otimes(
 \widetilde W_{r,\pi}^*\widetilde W_{r,\pi})
 \le
 C_{s,\eta}\mathscr A_{211}^{\,2}.}
\label{eq:V77-xy-FCQ}
\end{equation}
No bank occurrence, ownership, gap, or protected-sector hypothesis
is required.
\end{theorem}

\begin{proof}
If \(x\ge\eta\), then
\[
 q=xy+xz+yz\ge xz\ge\eta z.
\]
If \(y\ge\eta\), the same conclusion follows from
\(q\ge yz\).  In either case,
\[
 \mathscr A_{211}=pq\ge\eta pz.
\]
By \Cref{lem:V77-matching-floor},
\begin{equation}
 \frac1{\Xi_1\Xi_2\Xi_3\Xi_4}
 \le
 a_*^{-4}pz
 \le
 a_*^{-4}\eta^{-1}\mathscr A_{211}.
\label{eq:V77-volume-by-atlas}
\end{equation}
The all-payer uniform complete-word cap
\Cref{lem:V74-uniform-word} now gives
\[
 \|\widetilde W_{r,\pi}\|
 \le
 C_{s,\eta}\mathscr A_{211}.
\]
Square the ordinary norm and use product-state capacity bounded by
ordinary norm.  Taking adjoints proves the other orientation.
\end{proof}

\begin{corollary}[Two V76 dominant-label branches disappear]
\label{cor:V77-two-dominant-branches}
The \(x\)-dominant and \(y\)-dominant nondepleted branches of
\Cref{thm:V76-counterpacket} are closed unconditionally.  The V76
compatible-bank hypothesis is needed on a nondepleted dominant
branch only when \(z\) is dominant and both \(x,y\) are depleted.
\end{corollary}

\begin{proof}
If \(x=m\ge\eta\) or \(y=m\ge\eta\), then
\eqref{eq:V77-xy-nondepleted} holds.  Apply
\Cref{thm:V77-xy-FCQ}.
\end{proof}

\begin{corollary}[Sharper quantitative residual cell]
\label{cor:V77-residual-cell}
For every fixed \(0<\delta\le1\), all selected words outside
\begin{equation}
 \mathfrak Z_\delta
 :=
 \{p<\delta,\ x<\delta,\ y<\delta\}
\label{eq:V77-residual-cell}
\end{equation}
satisfy \(FCQ_{211}\), with constants depending only on
\(s,\delta\) and the fixed structural data.
\end{corollary}

\begin{proof}
If \(p\ge\delta\), use
\Cref{thm:V75-p-nondepleted}.  If \(x\ge\delta\) or
\(y\ge\delta\), use \Cref{thm:V77-xy-FCQ}.
\end{proof}

\section{A second unconditional cell inside the \(z\)-dominant
sector}
\label{sec:V77-mass-ratio}

The perfect matching leaves the factor \(x+y\) when \(z\) is the
only nondepleted quotient edge.  A local cross edge can pay this
factor and the remaining mismatch is exactly one row-mass ratio.

\begin{lemma}[Repeated-row volume comparison]
\label{lem:V77-repeated-row}
Let \(u\in\{1,2\}\) and let \(v,w\in\{3,4\}\) be distinct.  Then
\begin{equation}
 \frac1{\Xi_1\Xi_2\Xi_3\Xi_4}
 \le
 a_*^{-4}\frac{\Xi_u}{\Xi_v}\,
 p\,\theta_{uw}.
\label{eq:V77-repeated-row}
\end{equation}
\end{lemma}

\begin{proof}
At the selected coordinate,
\[
 p\ge\frac{a_*^2}{\Xi_1\Xi_2},
 \qquad
 \theta_{uw}\ge\frac{a_*^2}{\Xi_u\Xi_w}.
\]
Since \(\{v,w\}=\{3,4\}\), multiplication gives
\[
 p\theta_{uw}
 \ge
 \frac{a_*^4}
 {\Xi_1\Xi_2\Xi_u\Xi_w}
 =
 a_*^4\frac{\Xi_v}{\Xi_u}\,
 \frac1{\Xi_1\Xi_2\Xi_3\Xi_4}.
\]
Rearrange.
\end{proof}

\begin{theorem}[Mass-comparable \(z\)-dominant closure]
\label{thm:V77-z-mass-comparable}
Fix \(\eta>0\) and \(K\ge1\).  Suppose
\begin{equation}
 z\ge\eta,
 \qquad
 \min\{\Xi_1,\Xi_2\}
 \le K\max\{\Xi_3,\Xi_4\}.
\label{eq:V77-z-mass-cell}
\end{equation}
Then every selected word and every payer placement satisfies
\begin{equation}
 \boxed{
 \kappa_\otimes(
 \widetilde W_{r,\pi}^*\widetilde W_{r,\pi})
 \le C_{s,\eta,K}\mathscr A_{211}^{\,2}.}
\label{eq:V77-z-mass-FCQ}
\end{equation}
The left-square orientation obeys the same estimate.
\end{theorem}

\begin{proof}
Choose \(u\in\{1,2\}\) with
\(\Xi_u=\min\{\Xi_1,\Xi_2\}\), and
\(v\in\{3,4\}\) with
\(\Xi_v=\max\{\Xi_3,\Xi_4\}\).  Let \(w\) be the other singleton
row.  Then \(\Xi_u/\Xi_v\le K\) and
\(\theta_{uw}\le x+y\).  By
\Cref{lem:V77-repeated-row},
\[
 \frac1{\Xi_1\Xi_2\Xi_3\Xi_4}
 \le
 K a_*^{-4}p(x+y).
\]
Also
\[
 q=xy+z(x+y)\ge\eta(x+y),
\]
so
\[
 p(x+y)\le\eta^{-1}\mathscr A_{211}.
\]
The uniform word cap and the ordinary-norm domination of capacity
finish the proof exactly as in
\Cref{thm:V77-xy-FCQ}.
\end{proof}

\begin{remark}[The new condition is exhaustive over the four local
cross edges]
\label{rem:V77-four-cross-edges}
Choosing \(uw=13,14,23,24\) in
\Cref{lem:V77-repeated-row} yields respectively the ratios
\[
 \frac{\Xi_1}{\Xi_4},\quad
 \frac{\Xi_1}{\Xi_3},\quad
 \frac{\Xi_2}{\Xi_4},\quad
 \frac{\Xi_2}{\Xi_3}.
\]
At least one is bounded by \(K\) exactly when the second condition
in \eqref{eq:V77-z-mass-cell} holds.  Thus the proof has not hidden
a preferred cross edge.
\end{remark}

\section{Probability-overlap form of the remaining geometry}
\label{sec:V77-overlap}

For each row define its normalized coordinate weights
\begin{equation}
 w_i(e):=\frac{a_{i,e}}{\Xi_i}.
\label{eq:V77-w}
\end{equation}
They are probability vectors, and
\begin{equation}
 \theta_{ij}=\sum_e w_i(e)w_j(e).
\label{eq:V77-overlap}
\end{equation}
When \(z>0\), define the normalized \(34\)-overlap law
\begin{equation}
 \nu_{34}(e):=\frac{w_3(e)w_4(e)}{z}.
\label{eq:V77-nu34}
\end{equation}

\begin{lemma}[The \(34\) overlap avoids composite-row mass]
\label{lem:V77-overlap-avoidance}
For \(z>0\),
\begin{equation}
 \boxed{
 \mathbb E_{\nu_{34}}[w_1+w_2]
 \le
 \frac{\min\{x,y\}}{z}.}
\label{eq:V77-overlap-mean}
\end{equation}
Consequently, for every \(\tau>0\),
\begin{equation}
 \nu_{34}\{e:w_1(e)+w_2(e)>\tau\}
 \le
 \frac{\min\{x,y\}}{\tau z}.
\label{eq:V77-overlap-Markov}
\end{equation}
\end{lemma}

\begin{proof}
Because every probability weight is at most one,
\begin{align*}
 \sum_e(w_1+w_2)w_3w_4
 &\le
 \sum_e(w_1+w_2)w_3
 =x,\\
 \sum_e(w_1+w_2)w_3w_4
 &\le
 \sum_e(w_1+w_2)w_4
 =y.
\end{align*}
Divide the better inequality by
\(\sum_e w_3w_4=z\), proving
\eqref{eq:V77-overlap-mean}.  Markov's inequality gives
\eqref{eq:V77-overlap-Markov}.
\end{proof}

\begin{corollary}[Quantitative three-cluster separation]
\label{cor:V77-three-cluster}
If \(z_n\ge\eta>0\) while \(x_n,y_n\to0\), then for every fixed
\(\tau>0\),
\begin{equation}
 \nu_{34}^{(n)}
 \{w_1^{(n)}+w_2^{(n)}>\tau\}
 \longrightarrow0.
\label{eq:V77-three-cluster}
\end{equation}
Thus almost all normalized \(34\)-overlap mass lies on coordinates
which are simultaneously light for rows \(1\) and \(2\).
\end{corollary}

\begin{proof}
Insert \(z_n\ge\eta\) into
\eqref{eq:V77-overlap-Markov}.
\end{proof}

\begin{firewall}[Overlap localization is not yet operator damping]
\label{fire:V77-overlap-not-operator}
Equation \eqref{eq:V77-three-cluster} is an exact scalar
localization statement.  It does not say that the joining cumulant,
a protected projector, or a physical compression has norm
\(O(x+y)\).  Such a conclusion would require the localized
coordinates to occur in a signed response factor of the same source
word.  V77 does not make that conversion.
\end{firewall}

\section{The sharpened counterpacket theorem}
\label{sec:V77-counterpacket}

\begin{theorem}[Only a heavy-composite \(z\)-sector or total
depletion remains]
\label{thm:V77-counterpacket}
Let a sequence of complete selected-coordinate words violate every
uniform \(FCQ_{211}\) bound.  After passage to a subsequence,
\begin{equation}
 p_n\to0,\qquad x_n\to0,\qquad y_n\to0,
\label{eq:V77-pxy-zero}
\end{equation}
and one of the following holds:
\begin{enumerate}[label=\textup{(\alph*)},leftmargin=4em]
\item \(z_n\to0\); then all six normalized row-pair edges tend to
      zero and the pair-mass divergence conclusion
      \eqref{eq:V76-all-pair-masses} holds;
\item \(z_n\ge\eta>0\), and
\begin{equation}
 \frac{\min\{\Xi_1^{(n)},\Xi_2^{(n)}\}}
 {\max\{\Xi_3^{(n)},\Xi_4^{(n)}\}}
 \longrightarrow\infty.
\label{eq:V77-heavy-composite}
\end{equation}
Moreover the \(34\)-overlap law satisfies the three-cluster
localization \eqref{eq:V77-three-cluster}.
\end{enumerate}
\end{theorem}

\begin{proof}
\Cref{thm:V75-simultaneous-depletion} gives \(p_n\to0\).
If \(x_n\not\to0\), pass to a subsequence with
\(x_n\ge\eta>0\), contradicting
\Cref{thm:V77-xy-FCQ}; hence \(x_n\to0\).  The same argument applies
to \(y_n\).

Pass to a subsequence on which either \(z_n\to0\) or
\(z_n\ge\eta>0\).  The first alternative is
\Cref{cor:V76-total-depletion-masses}.  In the second, if the ratio
in \eqref{eq:V77-heavy-composite} did not tend to infinity, a further
subsequence would satisfy
\(\min\{\Xi_1,\Xi_2\}\le K\max\{\Xi_3,\Xi_4\}\) for some fixed
\(K\), contradicting
\Cref{thm:V77-z-mass-comparable}.  Finally apply
\Cref{cor:V77-three-cluster}.
\end{proof}

\begin{corollary}[The role of the V76 compatible-bank theorem is now
confined]
\label{cor:V77-bank-confined}
On a counterpacket with \(z_n\ge\eta\), the only remaining use of
\Cref{thm:V76-dominant-bank-FCQ} is the composite-cut edge set
\[
 \mathcal E_z=\{13,23,14,24\}
\]
inside the heavy-composite regime
\eqref{eq:V77-heavy-composite}.  If every resolved term is
compatible-bank-covered there, the counterpacket is impossible.
\end{corollary}

\begin{proof}
The \(x\)- and \(y\)-nondepleted sectors are closed by V77, and the
mass-comparable part of the remaining \(z\)-sector is closed by
\Cref{thm:V77-z-mass-comparable}.  Apply the \(\ell=z\) case of
\Cref{thm:V76-dominant-bank-FCQ}.
\end{proof}

\section{V77 result ledger and V78 acquisition order}
\label{sec:V77-status}

\begin{theorem}[Third-series V77 result ledger]
\label{thm:V77-results}
Relative to V1--V76, V77 proves:
\begin{enumerate}[label=\textup{(V77.\arabic*)},leftmargin=6.3em]
\item the disjoint perfect-matching floor
      \(pz\ge a_*^4/(\Xi_1\Xi_2\Xi_3\Xi_4)\);
\item unconditional all-payer \(FCQ_{211}\) closure whenever
      \(x\) or \(y\) is nondepleted;
\item elimination of the \(x\)- and \(y\)-dominant residuals;
\item the sharper residual cell
      \(\{p<\delta,x<\delta,y<\delta\}\);
\item a repeated-row inverse-volume comparison for each of the four
      composite-cut edges;
\item unconditional closure of the mass-comparable part of the
      \(z\)-dominant sector;
\item exact concentration of normalized \(34\)-overlap on
      composite-light coordinates; and
\item reduction of every counterpacket to either total edge
      depletion or a \(z\)-nondepleted sector in which both composite
      row masses dominate both singleton row masses without bound.
\end{enumerate}
\end{theorem}

\begin{proof}
These are
\Cref{lem:V77-matching-floor,
thm:V77-xy-FCQ,
cor:V77-two-dominant-branches,
cor:V77-residual-cell,
lem:V77-repeated-row,
thm:V77-z-mass-comparable,
lem:V77-overlap-avoidance,
thm:V77-counterpacket}.
\end{proof}

\begin{programme}[V78: ordered-source analysis in the
heavy-composite \(z\)-sector]
\label{prog:V77-V78}
\begin{enumerate}[leftmargin=3em]
\item Fix \(z\ge\eta\), the heavy-composite ratio
      \eqref{eq:V77-heavy-composite}, and the exact first-connectivity
      word \eqref{eq:V35-S211}.
\item Split the \(34\)-overlap coordinates by
      \(w_1+w_2\le\tau\) and its complement using
      \eqref{eq:V77-overlap-Markov}.  Keep this as an exact coordinate
      restriction before inserting any protected/gapped envelope.
\item Determine whether the composite-light \(34\) bank occurs
      before or after the selected joining coordinate.  Use the V73
      localized prefix \(p_r^-\) on the latter branch rather than
      restoring the global \(p\).
\item On the former branch, test the coefficient-one first-escape
      telescope of V72 for an actual \(13,23,14,\) or \(24\)
      response.  Do not replace the scalar localization by an
      operator norm.
\item Route a genuine compatible response through the V76
      prefix--bank compiler; if none occurs, record the exact
      no-escape source factor before taking a positive envelope.
\item Treat total edge depletion only after the \(z\)-nondepleted
      branch is exhausted.  The next compulsory companion audit is
      V80.
\end{enumerate}
\end{programme}

\begin{firewall}[Claim boundary after third-series V77]
\label{fire:V77-claim-boundary}
V77 unconditionally closes all selected words with \(x\) or \(y\)
nondepleted and the mass-comparable part of the \(z\)-nondepleted
sector.  It does not turn the \(34\)-overlap localization into an
operator response, close the heavy-composite or totally depleted
branches, prove compatible-bank occurrence, close the complete
\(2+1+1\) source or its global aggregation, or close other quartic
source families, higher MTA, WUFC, CPM, cutoff removal, reflection
positivity, Osterwalder--Schrader reconstruction, construction of
the continuum Yang--Mills measure, or a uniform positive continuum
spectral gap.  The full Yang--Mills mass gap has not been proved.
\end{firewall}

\paragraph{Parent-version record.}
V77 was formed from
\texttt{Renewal\_Yang\_Mills\_Mass\_Gap\_Research\_Notes\_V76.tex}.
The V76 parent had \(37\,295\) logical source lines,
\(1\,331\,571\) bytes, and SHA--256
\[
 \texttt{3b096609baf7012a1db99f27e4b8ea6759d13cf49aa0b933845e55783e62af45}.
\]
The next compulsory companion audit is V80.

% =====================================================================
% V78 APPENDIX: THE HEAVY-COMPOSITE DEFECT AND ITS MISSING RESPONSE
% =====================================================================

\chapter{V78: The Exact Heavy-Composite Defect and a Sharp
Insufficiency Test}
\label{ch:V78-heavy-defect}

\section{The one remaining mass ratio}
\label{sec:V78-ratio}

On the \(z\)-nondepleted residual define
\begin{equation}
 R:=
 \frac{\min\{\Xi_1,\Xi_2\}}
 {\max\{\Xi_3,\Xi_4\}},
 \qquad
 \chi:=R^{-1}.
\label{eq:V78-R-chi}
\end{equation}
The counterpacket theorem
\Cref{thm:V77-counterpacket} says that the still-unclosed branch has
\(R\to\infty\), or equivalently \(\chi\to0\).

\begin{lemma}[The residual denominator costs exactly \(R\)]
\label{lem:V78-denominator-defect}
Assume \(z\ge\eta>0\).  Then
\begin{align}
 \frac1{\Xi_1\Xi_2\Xi_3\Xi_4}
 &\le
 a_*^{-4}R\,p(x+y),
\label{eq:V78-volume-defect}\\
 \frac{p_r^-}{\Xi_3\Xi_4}
 &\le
 a_*^{-2}R\,p(x+y).
\label{eq:V78-prefix-defect}
\end{align}
Consequently
\begin{equation}
 \boxed{
 \|\widetilde W_{r,\pi}\|
 \le
 C_{s,\eta}R\,\mathscr A_{211}}
\label{eq:V78-R-loss}
\end{equation}
for every payer placement.  The same \(R\)-loss is obtained from the
uniform-volume branch and from the localized-prefix branch.
\end{lemma}

\begin{proof}
Choose \(u\in\{1,2\}\) with
\(\Xi_u=\min\{\Xi_1,\Xi_2\}\), and
\(v\in\{3,4\}\) with
\(\Xi_v=\max\{\Xi_3,\Xi_4\}\).  Let \(w\) be the other member of
\(\{3,4\}\).  Thus \(\Xi_u/\Xi_v=R\), and
\(\theta_{uw}\le x+y\).

Equation \eqref{eq:V78-volume-defect} is
\Cref{lem:V77-repeated-row}.  For the second estimate, the selected
coordinate gives
\[
 x+y\ge\theta_{uw}
 \ge\frac{a_*^2}{\Xi_u\Xi_w}.
\]
Since \(\Xi_3\Xi_4=\Xi_v\Xi_w\),
\[
 \frac1{\Xi_3\Xi_4}
 \le
 a_*^{-2}\frac{\Xi_u}{\Xi_v}(x+y)
 =
 a_*^{-2}R(x+y).
\]
Multiply by \(p_r^-\le p\).

The uniform word cap combined with
\eqref{eq:V78-volume-defect} gives
\[
 \|\widetilde W_{r,\pi}\|
 \le C_sR\,p(x+y).
\]
Alternatively use the exact localized prefix norm implicit in
\Cref{lem:V76-prefix-bank-factorization} after omitting the
contraction bank, followed by
\eqref{eq:V78-prefix-defect}.  Finally,
\[
 \mathscr A_{211}
 =
 p\{xy+z(x+y)\}
 \ge\eta p(x+y).
\]
This proves \eqref{eq:V78-R-loss}.
\end{proof}

\begin{corollary}[Subpolynomial ratio closure]
\label{cor:V78-bounded-ratio}
Every fixed cell \(R\le K\) is closed with
\[
 \kappa_\otimes(
 \widetilde W_{r,\pi}^*\widetilde W_{r,\pi})
 \le
 C_{s,\eta,K}\mathscr A_{211}^{\,2}.
\]
Thus a counterpacket on \(z\ge\eta\) must use the unbounded factor
\(R\), not merely large total input volume.
\end{corollary}

\begin{proof}
Square \eqref{eq:V78-R-loss} and use \(R\le K\).  This recovers
\Cref{thm:V77-z-mass-comparable} in the ratio notation.
\end{proof}

\section{A scalar support family saturating the defect}
\label{sec:V78-saturation}

The next example is not asserted to be the matrix value of a
Yang--Mills source word.  It is a sharp logical test of the scalar
weight information and of every positive norm estimate used up to
V77.

\begin{proposition}[Heavy-composite support family]
\label{prop:V78-support-family}
Let \(N\ge2\), take five ordered coordinates
\(e_0<r<e_1<e_2<e_3\), and assign the row weights
\begin{center}
\begin{tabular}{c|ccccc}
 &\(e_0\)&\(r\)&\(e_1\)&\(e_2\)&\(e_3\)\\
\hline
\(a_1\)&\(1\)&\(1\)&\(N^2\)&\(0\)&\(0\)\\
\(a_2\)&\(1\)&\(1\)&\(0\)&\(N^2\)&\(0\)\\
\(a_3\)&\(0\)&\(1\)&\(0\)&\(0\)&\(N\)\\
\(a_4\)&\(0\)&\(1\)&\(0\)&\(0\)&\(N\)
\end{tabular}
\end{center}
Then \(r\) is four-row active and
\begin{align}
 \Xi_1=\Xi_2&=A_N:=N^2+2,
 &
 \Xi_3=\Xi_4&=B_N:=N+1,
\label{eq:V78-family-masses}\\
 p&=\frac2{A_N^2},
 &
 x=y&=\frac2{A_NB_N},
\label{eq:V78-family-pxy}\\
 z&=\frac{N^2+1}{B_N^2}.
\label{eq:V78-family-z}
\end{align}
In particular,
\begin{equation}
 R_N=\frac{A_N}{B_N}\asymp N,
 \qquad
 \mathscr A_{211}^{(N)}\asymp N^{-7},
 \qquad
 \frac1{\Xi_1\Xi_2\Xi_3\Xi_4}\asymp N^{-6}.
\label{eq:V78-family-asymptotics}
\end{equation}
Moreover \(H_{12}^{<r}=1\), so the localized-prefix envelope has
the same \(N^{-6}\) scale, and
\begin{equation}
 \frac{
 (\Xi_1\Xi_2\Xi_3\Xi_4)^{-1}}
 {\mathscr A_{211}^{(N)}}
 \asymp R_N.
\label{eq:V78-family-saturates}
\end{equation}
\end{proposition}

\begin{proof}
Summing the table gives \eqref{eq:V78-family-masses}.  The only
\(12\) overlaps are \(e_0,r\), the only cross-composite overlaps are
at \(r\), and the \(34\) overlaps are \(r,e_3\).  Hence
\[
 H_{12}=2,\qquad
 H_{13}=H_{14}=H_{23}=H_{24}=1,\qquad
 H_{34}=N^2+1,
\]
which gives \eqref{eq:V78-family-pxy}--%
\eqref{eq:V78-family-z}.  Therefore
\[
 q
 =
 \frac4{A_N^2B_N^2}
 +
 \frac{4(N^2+1)}{A_NB_N^3}
 \asymp N^{-3},
\]
and \(pq\asymp N^{-7}\).  The inverse input volume is
\(A_N^{-2}B_N^{-2}\asymp N^{-6}\), while
\(R_N=A_N/B_N\asymp N\).  Since \(e_0\) is the only \(12\)-overlap
strictly before \(r\), \(H_{12}^{<r}=1\), and its normalized
prefix scale is
\[
 \frac{H_{12}^{<r}}
 {\Xi_1\Xi_2\Xi_3\Xi_4}
 =
 A_N^{-2}B_N^{-2}.
\]
These identities prove the last two assertions.
\end{proof}

\begin{corollary}[Scalar overlap localization alone cannot remove
the defect]
\label{cor:V78-localization-sharp}
For the family in \Cref{prop:V78-support-family},
\(\nu_{34}^{(N)}\) concentrates at \(e_3\), where
\(w_1+w_2=0\), while the ratio loss in
\eqref{eq:V78-family-saturates} diverges.  Thus even perfect
composite-light localization of the \(34\) overlap is insufficient
without an operator response connecting it to the prefix/joining
word.
\end{corollary}

\begin{proof}
The \(34\) overlap at \(e_3\) is \(N^2\), while that at \(r\) is
one, so
\[
 \nu_{34}^{(N)}(e_3)=\frac{N^2}{N^2+1}\to1.
\]
Rows \(1,2\) are inactive at \(e_3\), hence
\(w_1(e_3)+w_2(e_3)=0\).  The divergent defect is
\eqref{eq:V78-family-saturates}.
\end{proof}

\begin{corollary}[Compatible composite-cut banks need not occur]
\label{cor:V78-no-compatible-bank}
In the same family, after excluding the selected coordinate \(r\),
there is no nonempty two-row bank on any edge of
\(\mathcal E_z=\{13,23,14,24\}\).  Consequently the V76
compatible-bank theorem cannot by itself close the full
heavy-composite sector.
\end{corollary}

\begin{proof}
Outside \(r\), coordinate \(e_0\) carries only rows \(1,2\);
\(e_1,e_2\) carry only one composite row each; and \(e_3\) carries
only rows \(3,4\).  No remaining coordinate has both endpoints of a
composite-cut edge active.
\end{proof}

\section{A no-go theorem for the currently separated positive data}
\label{sec:V78-no-go}

\begin{theorem}[The V55, V74, and V77 envelopes do not imply the
product card]
\label{thm:V78-positive-data-no-go}
There is no constant \(C\), independent of the support weights, for
which the following data alone imply
\(\|W\|\le C\mathscr A_{211}\):
\begin{enumerate}[label=\textup{(\roman*)},leftmargin=4em]
\item the uniform word envelope
      \(\|W\|\le C_0/(\Xi_1\Xi_2\Xi_3\Xi_4)\);
\item the localized-prefix envelope
      \(\|W\|\le
      C_0p_r^-/(\Xi_3\Xi_4)\);
\item the composite-cut envelope
      \(\|W\|\le C_0(x+y)\);
\item four-row activity, the nontrivial mass floor, and
      \(z\ge\eta\); and
\item the scalar \(34\)-overlap localization of
      \Cref{lem:V77-overlap-avoidance}.
\end{enumerate}
\end{theorem}

\begin{proof}
Use the weights of
\Cref{prop:V78-support-family} and take the scalar one-dimensional
operator
\[
 W_N:=\frac1{A_N^2B_N^2}.
\]
It satisfies the uniform and localized-prefix envelopes with
constant one.  Since \(x+y\asymp N^{-3}\), it also satisfies the
composite-cut envelope for all \(N\).  The four-row activity, mass
floor, \(z\ge1/3\) for \(N\ge2\), and overlap localization were
verified above.  But
\[
 \frac{\|W_N\|}{\mathscr A_{211}^{(N)}}
 \asymp R_N\longrightarrow\infty.
\]
Thus no implication using only items \textup{(i)--(v)} is possible.
The scalar \(W_N\) is a logical countermodel to an implication
between estimates; it is not claimed to be a physical YM word.
\end{proof}

\begin{assessment}[Mandatory upstream move]
\label{ass:V78-upstream}
The missing \(R^{-1}\) cannot be recovered by recombining the
already separated V55 prefix norm, V74 cut norm, V76 bank
conditional, and V77 scalar localization.  The next proof must
inspect an operator occurrence suppressed by at least one of those
envelopes.  In the exact V35 word the earliest candidate is the
unsplit composite joining cumulant
\[
 J_r(12|3|4)
 =
 \kappa_3(Z_{12,r},X_{3,r},X_{4,r}),
\]
kept next to the complete binary prefix before the row-\(1\) or
row-\(2\) heavy demand is estimated.
\end{assessment}

\section{The exact response card now required}
\label{sec:V78-card}

\begin{definition}[Heavy-composite response card]
\label{def:V78-HCRC}
On \(z\ge\eta\) and \(R\ge1\), a restriction-resolved word satisfies
\(HCRC_{12|34}\) if
\begin{equation}
 \boxed{
 \|\widetilde W_{r,\pi,\nu}\|
 \le
 C_s\chi\,
 \frac{p_r^-}{\Xi_3\Xi_4},
 \qquad
 \chi=\frac{\max\{\Xi_3,\Xi_4\}}
 {\min\{\Xi_1,\Xi_2\}}.}
\label{eq:V78-HCRC}
\end{equation}
The card must hold for the actual once-paid word, with the same
payer placement in its adjoint square.
\end{definition}

\begin{theorem}[\(HCRC_{12|34}\) is sufficient on the
heavy-composite sector]
\label{thm:V78-HCRC-suffices}
If every term of one fixed finite exact resolution satisfies
\eqref{eq:V78-HCRC}, then
\begin{equation}
 \kappa_\otimes(
 \widetilde W_{r,\pi}^*\widetilde W_{r,\pi})
 \le
 C_{s,\eta}\mathscr A_{211}^{\,2}.
\label{eq:V78-HCRC-FCQ}
\end{equation}
\end{theorem}

\begin{proof}
The proof of \eqref{eq:V78-prefix-defect} gives
\[
 \frac{p_r^-}{\Xi_3\Xi_4}
 \le a_*^{-2}R\,p(x+y).
\]
Multiplication by \(\chi=R^{-1}\) and
\eqref{eq:V78-HCRC} yield
\[
 \|\widetilde W_{r,\pi,\nu}\|
 \le C_sp(x+y)
 \le C_{s,\eta}\mathscr A_{211}.
\]
Square, take capacity, and recombine the fixed number of exact terms
using \Cref{lem:V58-finite-recombination}.
\end{proof}

\begin{firewall}[The card is not proved by inserting a scalar ratio]
\label{fire:V78-card-open}
The factor \(\chi\) in \eqref{eq:V78-HCRC} must be the norm or
capacity of an actual source response.  The numerical fact
\(\chi\to0\), a lower response inequality, or the concentration
\eqref{eq:V77-three-cluster} cannot be multiplied into the word.
V78 proves sufficiency and necessity relative to the present
envelopes, not the \(HCRC_{12|34}\) estimate itself.
\end{firewall}

\section{V78 result ledger and V79 acquisition order}
\label{sec:V78-status}

\begin{theorem}[Third-series V78 result ledger]
\label{thm:V78-results}
Relative to V1--V77, V78 proves:
\begin{enumerate}[label=\textup{(V78.\arabic*)},leftmargin=6.3em]
\item the exact heavy-composite ratio
      \(R=\min(\Xi_1,\Xi_2)/\max(\Xi_3,\Xi_4)\);
\item matching \(R\)-loss bounds from both the uniform-volume and
      localized-prefix routes;
\item a five-coordinate family saturating that loss while
      \(z\) remains nondepleted;
\item failure of compatible composite-cut bank occurrence on the
      same family;
\item a sharp no-go theorem showing that the presently separated
      positive envelopes and scalar localization do not imply
      \(PQC_{211}\);
\item identification of the mandatory upstream object as the
      unsplit composite joining cumulant adjacent to the prefix;
\item the exact \(HCRC_{12|34}\) response card carrying the missing
      \(R^{-1}\); and
\item proof that this card closes the complete heavy-composite
      selected-word sector after finite resolution.
\end{enumerate}
\end{theorem}

\begin{proof}
These are
\Cref{lem:V78-denominator-defect,
prop:V78-support-family,
cor:V78-no-compatible-bank,
thm:V78-positive-data-no-go,
ass:V78-upstream,
def:V78-HCRC,
thm:V78-HCRC-suffices}.
\end{proof}

\begin{programme}[V79: unsplit composite joining response]
\label{prog:V78-V79}
\begin{enumerate}[leftmargin=3em]
\item Write
      \(\kappa_3(Z_{12},X_3,X_4)\) as one centered composite
      covariance, before expanding it into its five moment-partition
      roles.
\item Apply the coefficient-one heat/covariance telescope to the
      composite entry \(Z_{12}\), retaining the order of its row-\(1\)
      and row-\(2\) differences and the unique payer.
\item Normalize only after the composite covariance is adjacent to
      \(K_{12}^{<r,\mathrm{conn}}\).  Test whether one heavy
      composite-row denominator appears as
      \(\chi\), rather than as the lossy scalar \(R\).
\item Separate the two Leibniz branches only after their common
      centered \(34\) factor is retained.  Do not infer a product
      from their separate norms.
\item Test the exact carrier on
      \Cref{prop:V78-support-family}; a valid estimate must gain
      \(N^{-1}\) beyond the localized-prefix envelope there.
\item If the centered covariance still gives only \(R^0\), inspect
      its compressed output-block commutator before returning to
      bank routing.
\item The next compulsory companion audit remains V80.
\end{enumerate}
\end{programme}

\begin{firewall}[Claim boundary after third-series V78]
\label{fire:V78-claim-boundary}
V78 quantifies the sole mass-ratio defect on the \(z\)-nondepleted
heavy-composite sector, proves that the existing separated estimates
cannot remove it, and gives a sufficient exact response card.  It
does not prove \(HCRC_{12|34}\), close the heavy-composite or totally
depleted branches, close the complete \(2+1+1\) source or its global
aggregation, or close other quartic source families, higher MTA,
WUFC, CPM, cutoff removal, reflection positivity,
Osterwalder--Schrader reconstruction, construction of the continuum
Yang--Mills measure, or a uniform positive continuum spectral gap.
The full Yang--Mills mass gap has not been proved.
\end{firewall}

\paragraph{Parent-version record.}
V78 was formed from
\texttt{Renewal\_Yang\_Mills\_Mass\_Gap\_Research\_Notes\_V77.tex}.
The V77 parent had \(37\,780\) logical source lines,
\(1\,346\,503\) bytes, and SHA--256
\[
 \texttt{d3579299cb4ff9b16ac805554948dd09df3581f10076f871ab9c27d825e7421f}.
\]
The next compulsory companion audit is V80.

% =====================================================================
% V79 APPENDIX: PRIVATE HEAVY-MASS CLOSURE AND INCIDENCE ROUTING
% =====================================================================

\chapter{V79: Physical-Private Damping Closes the Scalar Saturation
Mechanism}
\label{ch:V79-private-heavy}

\section{The V78 countermodel is scalar, not physically private}
\label{sec:V79-correction}

The V78 support family correctly proves that scalar floors, overlap
localization, and the already separated word envelopes do not imply
\(HCRC_{12|34}\).  Its large row-\(1\) and row-\(2\) masses,
however, were placed on coordinates at which only one physical row
is active.  Such coordinates are not neutral in the complete
Yang--Mills word: the packetwise physical-private theorem factors
their Wilson expectations exactly, and heat gives arbitrarily many
inverse powers of their total mass.  V79 retains the V78 no-go as a
logical statement and removes its physically private realization
from the counterpacket list.

For \(i\in\{1,2\}\), define the complete physical-private bank
\begin{equation}
 \mathcal P_i^{\rm only}
 :=
 \{e:a_{i,e}>0,\ a_{j,e}=0\text{ for every }j\ne i\},
 \qquad
 P_i:=\sum_{e\in\mathcal P_i^{\rm only}}a_{i,e}.
\label{eq:V79-private-bank}
\end{equation}
Let \(q_{P_i}\) be its complete normalized Wilson multiplier, as in
\eqref{eq:V27-private-product}.

\begin{lemma}[Cofinal private damping supplies the heavy-composite
ratio]
\label{lem:V79-private-chi}
Assume
\begin{equation}
 P_i\ge\beta\Xi_i
\label{eq:V79-private-cofinal}
\end{equation}
for some \(i\in\{1,2\}\) and fixed \(\beta>0\).  On \(R\ge1\),
\begin{equation}
 q_{P_i}\le C_{s,\beta}\chi.
\label{eq:V79-private-chi}
\end{equation}
If the unique payer lies in this private bank, its complete paid
scalar \(\mathfrak p_{P_i}\) satisfies the same estimate with a
different fixed-\(s\) constant:
\begin{equation}
 \mathfrak p_{P_i}\le C_{s,\beta}\chi.
\label{eq:V79-paid-private-chi}
\end{equation}
\end{lemma}

\begin{proof}
The \(m=1\) line of
\Cref{lem:V27-private-denominator}, whose proof uses only the
cofinal bound \eqref{eq:V79-private-cofinal}, gives
\[
 q_{P_i}\le\frac{C_{s,\beta}}{\Xi_i}.
\]
For a private payer,
\Cref{lem:V29-paid-private-moments} with \(m=1\) gives
\[
 \mathfrak p_{P_i}
 \le\frac{C_s}{P_i}
 \le\frac{C_{s,\beta}}{\Xi_i}.
\]
Now \(\Xi_i\ge\min\{\Xi_1,\Xi_2\}\), while activity gives
\(\max\{\Xi_3,\Xi_4\}\ge a_*\).  Hence
\[
 \frac1{\Xi_i}
 \le
 \frac1{\min\{\Xi_1,\Xi_2\}}
 \le
 a_*^{-1}
 \frac{\max\{\Xi_3,\Xi_4\}}
 {\min\{\Xi_1,\Xi_2\}}
 =
 a_*^{-1}\chi.
\]
Absorb \(a_*^{-1}\).
\end{proof}

\section{Exact insertion into the heavy-composite word}
\label{sec:V79-insertion}

\begin{lemma}[Private-factorized prefix word]
\label{lem:V79-private-prefix-word}
Fix a selected coordinate \(r\) and a payer placement.  If
\(\mathcal P_i^{\rm only}\) is disjoint from \(r\), then the complete
first-connectivity family, after its exact local-role reassembly, has
one of the presentations
\begin{align}
 \widetilde W_{r,\pi}
 &=
 q_{P_i}\,
 \frac{p_r^-}{\Xi_3\Xi_4}\,
 L_{r,\pi}Q_{r,\pi},
\label{eq:V79-private-prefix-unpaid}\\
 \widetilde W_{r,\pi}
 &=
 \mathfrak p_{P_i}\,
 \frac{p_r^-}{\Xi_3\Xi_4}\,
 L_{r,\pi}Q_{r,\pi}
\label{eq:V79-private-prefix-paid}
\end{align}
according as the payer lies outside or inside the private bank, with
\[
 \|L_{r,\pi}\|+\|Q_{r,\pi}\|\le C_s.
\]
If exact subset resolution is required, the same statement holds
term by term with a fixed number of terms.
\end{lemma}

\begin{proof}
At a physical-private coordinate the local letter is the discrete
partition letter and its complete normalized expectation is scalar.
The packetwise identity
\Cref{lem:V27-exact-private-factor} therefore factors
\(q_{P_i}\) from the complete first-connectivity family before any
positive envelope.  Such a coordinate cannot be the four-row
joining coordinate.  Deleting it leaves the relative order of all
other coordinates unchanged and removes no \(12\) covariance stock,
because the other endpoint row is inactive there.

If the payer lies in the private bank, the complete paid-bank
calculation of
\Cref{lem:V29-paid-private-moments,fire:V29-private-payer}
replaces \(q_{P_i}\) by the single paid scalar
\(\mathfrak p_{P_i}\); it does not multiply an outside payer by a
second private factor.  In either case, apply the support-local
prefix estimate \Cref{lem:V73-prefix-stock} to the remaining word.
Every other paid or unpaid complete factor has the uniform norm used
in \Cref{lem:V55-complete-prefix-response}.  Division by the four
input masses gives
\[
 \frac{H_{12}^{<r}}
 {\Xi_1\Xi_2\Xi_3\Xi_4}
 =
 \frac{p_r^-}{\Xi_3\Xi_4}.
\]
This proves the displayed presentations.  Exact restriction
commutes with a scalar common factor and has the fixed finite term
count of \Cref{thm:V58-subset-resolution}.
\end{proof}

\begin{theorem}[Physical-private heavy-composite closure]
\label{thm:V79-private-HCRC}
Fix \(z\ge\eta>0\), \(R\ge1\), and \(\beta>0\).  If
\eqref{eq:V79-private-cofinal} holds for either composite row
\(i=1\) or \(i=2\), then \(HCRC_{12|34}\) holds for every payer
placement.  Consequently
\begin{equation}
 \boxed{
 \kappa_\otimes(
 \widetilde W_{r,\pi}^*\widetilde W_{r,\pi})
 \le
 C_{s,\eta,\beta}\mathscr A_{211}^{\,2}.}
\label{eq:V79-private-FCQ}
\end{equation}
The same conclusion holds in the other square orientation.
\end{theorem}

\begin{proof}
Use \Cref{lem:V79-private-prefix-word}.  In the outside-bank payer
case, \eqref{eq:V79-private-chi} gives
\[
 \|\widetilde W_{r,\pi}\|
 \le
 C_{s,\beta}\chi
 \frac{p_r^-}{\Xi_3\Xi_4}.
\]
In the private-payer case use
\eqref{eq:V79-paid-private-chi}; all other payer placements are
already included in the first line.  This is exactly
\eqref{eq:V78-HCRC}.  Apply
\Cref{thm:V78-HCRC-suffices}.  Taking adjoints changes no scalar or
ordinary norm.
\end{proof}

\begin{corollary}[The V78 saturation family is physically damped]
\label{cor:V79-V78-damped}
In the support family
\Cref{prop:V78-support-family}, the large \(e_1\) and \(e_2\)
coordinates form physical-private banks with
\[
 \frac{P_i}{\Xi_i}
 =
 \frac{N^2}{N^2+2}\longrightarrow1.
\]
If those scalar weights arise from actual nontrivial Wilson input
blocks, their complete private multipliers satisfy
\[
 q_{P_i}\le C_sN^{-2},
\]
and any one such factor supplies more than the missing
\(\chi\asymp N^{-1}\).  Therefore the family is not a physical
counterexample to \(HCRC_{12|34}\).
\end{corollary}

\begin{proof}
The cofinal fraction is immediate from the table in V78.
\Cref{lem:V27-private-denominator} gives
\[
 q_{P_i}\le C_s/\Xi_i\le C_sN^{-2}.
\]
Compare with \(\chi=B_N/A_N\asymp N^{-1}\), and apply
\Cref{thm:V79-private-HCRC}.
\end{proof}

\begin{firewall}[Correction does not weaken the V78 no-go]
\label{fire:V79-no-go-retained}
\Cref{thm:V78-positive-data-no-go} remains valid: its list of
hypotheses deliberately omitted the exact physical-private
multiplier.  V79 adds new operator ancestry not present in that
logical implication.  The distinction is precisely why the
programme returned upstream from scalar weights to the complete
source word.
\end{firewall}

\section{Where heavy mass must go after private closure}
\label{sec:V79-incidence}

For \(i\in\{1,2\}\), partition the nonprivate coordinates at which
row \(i\) is active by the least other physically active row.  Write
\begin{equation}
 \mathcal S_{ij}^{(i)}
 :=
 \{e:a_{i,e}>0,\ j=\min\{k\ne i:a_{k,e}>0\}\},
 \qquad
 S_{ij}^{(i)}
 :=
 \sum_{e\in\mathcal S_{ij}^{(i)}}a_{i,e}.
\label{eq:V79-shared-cells}
\end{equation}

\begin{lemma}[Exact private/shared mass partition]
\label{lem:V79-private-shared}
For \(i=1,2\),
\begin{equation}
 \boxed{
 \Xi_i=P_i+\sum_{j\ne i}S_{ij}^{(i)}.}
\label{eq:V79-private-shared}
\end{equation}
If \(P_i\le\beta\Xi_i\), some \(j\ne i\) satisfies
\begin{equation}
 S_{ij}^{(i)}
 \ge
 \frac{1-\beta}{3}\Xi_i.
\label{eq:V79-cofinal-shared}
\end{equation}
\end{lemma}

\begin{proof}
Every coordinate active in row \(i\) is either physically private or
has a unique least other active row.  The sets in
\eqref{eq:V79-private-bank} and
\eqref{eq:V79-shared-cells} are therefore disjoint and exhaustive,
which proves \eqref{eq:V79-private-shared}.  If the private summand
is at most \(\beta\Xi_i\), the three shared summands total at least
\((1-\beta)\Xi_i\); pigeonhole gives
\eqref{eq:V79-cofinal-shared}.
\end{proof}

\begin{lemma}[A singleton partner becomes strongly imbalanced]
\label{lem:V79-cross-imbalance}
Consider a heavy-composite sequence with
\(\chi_n\to0\).  Suppose that for a fixed
\(i\in\{1,2\}\), fixed \(j\in\{3,4\}\), and fixed \(\sigma>0\),
\begin{equation}
 S_{ij}^{(i,n)}\ge\sigma\Xi_i^{(n)}.
\label{eq:V79-cross-cofinal}
\end{equation}
For any fixed \(L>0\), define
\[
 C_L^{(n)}
 :=
 \{e\in\mathcal S_{ij}^{(i,n)}:
       a_{i,e}\le L a_{j,e}\},
 \qquad
 U_L^{(n)}
 :=
 \mathcal S_{ij}^{(i,n)}\setminus C_L^{(n)}.
\]
Then
\begin{align}
 \sum_{e\in C_L^{(n)}}a_{i,e}
 &\le L\Xi_j^{(n)}
 \le L\chi_n\Xi_i^{(n)},
\label{eq:V79-comparable-small}\\
 \sum_{e\in U_L^{(n)}}a_{i,e}
 &\ge(\sigma-L\chi_n)\Xi_i^{(n)},
\label{eq:V79-unbalanced-cofinal}\\
 a_{i,e}&>La_{j,e}\qquad(e\in U_L^{(n)}).
\label{eq:V79-pointwise-imbalance}
\end{align}
Thus every cofinal composite-to-singleton cell becomes a cofinal
oriented strong-imbalance bank.
\end{lemma}

\begin{proof}
On \(C_L^{(n)}\),
\[
 \sum_{e\in C_L^{(n)}}a_{i,e}
 \le
 L\sum_{e\in C_L^{(n)}}a_{j,e}
 \le L\Xi_j^{(n)}.
\]
Since
\[
 \frac{\Xi_j^{(n)}}{\Xi_i^{(n)}}
 \le
 \frac{\max\{\Xi_3^{(n)},\Xi_4^{(n)}\}}
 {\min\{\Xi_1^{(n)},\Xi_2^{(n)}\}}
 =\chi_n,
\]
this proves \eqref{eq:V79-comparable-small}.  Subtract it from
\eqref{eq:V79-cross-cofinal} to obtain
\eqref{eq:V79-unbalanced-cofinal}; the last line is the definition
of \(U_L^{(n)}\).
\end{proof}

\begin{remark}[Pair imbalance is not yet a heat gap]
\label{rem:V79-pair-not-full}
The inequality \(a_{i,e}>La_{j,e}\) does not say that row \(i\) is
V32-unbalanced against all other rows.  Another row can balance it.
The exact next tool is therefore the V64 full-row split: each
coordinate in \(U_L\) either has a gapped heavy-row output or owns a
new uniformly comparable physical partner.  Applying heat decay
without that split would repeat the pair-versus-full-balance error
corrected in V64.
\end{remark}

\section{The new residual incidence table}
\label{sec:V79-residual}

\begin{theorem}[Private-free heavy-composite routing]
\label{thm:V79-routing}
Let a \(z\)-nondepleted heavy-composite counterpacket sequence be as
in \Cref{thm:V77-counterpacket}.  After passage to a subsequence:
\begin{enumerate}[label=\textup{(\roman*)},leftmargin=4em]
\item for both \(i=1,2\),
\begin{equation}
 \frac{P_i^{(n)}}{\Xi_i^{(n)}}\longrightarrow0;
\label{eq:V79-private-vanishes}
\end{equation}
\item for each \(i\), one fixed partner \(j_i\ne i\) satisfies
\begin{equation}
 S_{ij_i}^{(i,n)}
 \ge\frac14\Xi_i^{(n)}
\label{eq:V79-stable-partner}
\end{equation}
for all sufficiently large \(n\);
\item if \(j_i\in\{3,4\}\), the corresponding cell contains, for
      every fixed \(L\), a cofinal oriented bank satisfying
      \eqref{eq:V79-unbalanced-cofinal}--%
      \eqref{eq:V79-pointwise-imbalance};
\item if \(j_i\) is the other composite row, a fixed fraction of
      row \(i\)'s mass lies on physical \(12\)-shared coordinates.
\end{enumerate}
\end{theorem}

\begin{proof}
If \eqref{eq:V79-private-vanishes} failed for some \(i\), a
subsequence would satisfy \(P_i\ge\beta\Xi_i\), contradicting
\Cref{thm:V79-private-HCRC}.  Hence for large \(n\),
\(P_i\le\Xi_i/4\).  Apply
\Cref{lem:V79-private-shared}: one of three partners carries at least
\(\Xi_i/4\).  Pass to a subsequence so that this partner is fixed.
If it is a singleton row, apply
\Cref{lem:V79-cross-imbalance}.  If it is the other composite row,
\eqref{eq:V79-stable-partner} is exactly the last assertion.
\end{proof}

\begin{assessment}[What has genuinely changed]
\label{ass:V79-change}
The heavy-composite defect is no longer supported by arbitrary
one-row spectator mass.  Both heavy rows must place asymptotically
all their mass into physical shared-incidence cells.  A cofinal
cross-cut cell is automatically strongly imbalanced and enters the
already established V64 full-row trichotomy.  The only alternative
keeps a cofinal \(12\)-shared bank adjacent to the actual binary
prefix.  These are operator-bearing coordinate banks, not scalar
mass labels.
\end{assessment}

\section{V79 result ledger and V80 audit mandate}
\label{sec:V79-status}

\begin{theorem}[Third-series V79 result ledger]
\label{thm:V79-results}
Relative to V1--V78, V79 proves:
\begin{enumerate}[label=\textup{(V79.\arabic*)},leftmargin=6.3em]
\item cofinal physical-private heat damping at the missing
      \(\chi=R^{-1}\) scale, including a payer inside the bank;
\item exact insertion of that scalar multiplier next to the
      localized binary-prefix response;
\item \(HCRC_{12|34}\) and full selected-word closure of every
      cofinal-private heavy-composite packet;
\item physical damping, rather than physical obstruction, of the
      V78 scalar saturation family;
\item an exact private/shared physical-incidence partition for each
      heavy composite row;
\item extraction of a cofinal shared partner when private mass is
      subcofinal;
\item conversion of every cofinal composite-to-singleton cell into a
      cofinal pointwise strong-imbalance bank; and
\item reduction of the remaining heavy-composite counterpacket to
      V64 full-row routing or a cofinal internal \(12\)-shared bank.
\end{enumerate}
\end{theorem}

\begin{proof}
These are
\Cref{lem:V79-private-chi,
lem:V79-private-prefix-word,
thm:V79-private-HCRC,
cor:V79-V78-damped,
lem:V79-private-shared,
lem:V79-cross-imbalance,
thm:V79-routing}.
\end{proof}

\begin{programme}[V80: compulsory audit and physical shared-bank
termination]
\label{prog:V79-V80}
\begin{enumerate}[leftmargin=3em]
\item Perform the compulsory full read-only audit of the newest
      active Standard-Model and Navier--Stokes notes, record exact
      metadata, and transfer only type-correct proof order.
\item For a cross-cut partner in
      \Cref{thm:V79-routing}, apply the V64 coordinatewise
      full-balance split to \(U_L\), preserving the source occurrence
      and the unique payer.
\item Close the cofinal gapped alternative with one power of the V64
      mass-ratio decay, and test whether its central insertion proves
      \(HCRC_{12|34}\) for every payer placement.
\item If V64 selects a new comparable partner, distinguish a
      singleton partner, which is a compatible V76 bank, from the
      other composite row, which moves the mass to a \(12\)-shared
      bank.
\item On the internal \(12\)-shared branch retain that bank inside
      the doubled prefix square.  Do not charge its overlap once as
      \(p\) and again as an independent bank response.
\item Remove marked/payer coordinates only by an exact finite
      exceptional split; if one is mass-cofinal, route it to the
      selected or paid hot-output estimate instead of declaring the
      unmarked bank cofinal.
\end{enumerate}
\end{programme}

\begin{firewall}[Claim boundary after third-series V79]
\label{fire:V79-claim-boundary}
V79 closes the cofinal physical-private part of the
\(z\)-nondepleted heavy-composite sector and proves an exhaustive
shared-incidence routing for the remainder.  It does not yet insert
the resulting V64 trichotomy at the \(HCRC_{12|34}\) scale, close the
cofinal internal \(12\)-bank, close total edge depletion, close the
complete \(2+1+1\) source or its global aggregation, or close other
quartic source families, higher MTA, WUFC, CPM, cutoff removal,
reflection positivity, Osterwalder--Schrader reconstruction,
construction of the continuum Yang--Mills measure, or a uniform
positive continuum spectral gap.  The full Yang--Mills mass gap has
not been proved.
\end{firewall}

\paragraph{Parent-version record.}
V79 was formed from
\texttt{Renewal\_Yang\_Mills\_Mass\_Gap\_Research\_Notes\_V78.tex}.
The V78 parent had \(38\,232\) logical source lines,
\(1\,360\,498\) bytes, and SHA--256
\[
 \texttt{540ac8cab8165c892e95884d22ca782b1b516a8e761a2f3ed54cbe56818604a7}.
\]
The next compulsory companion audit is V80.

% =====================================================================
% V80 APPENDIX: COMPANION AUDIT AND COFINAL SHARED-BANK TERMINATION
% =====================================================================

\chapter{V80: Companion Audit and Maximal-Composite Shared-Bank
Termination}
\label{ch:V80-audit-shared}

\section{Compulsory read-only companion audit}
\label{sec:V80-audit}

The scheduled five-version audit was completed before selecting the
V80 proof direction.  The newest active companion files were:
\begin{enumerate}[leftmargin=3em]
\item
\texttt{Renewal\_SM\_Einstein\_Continuum\_Research\_Notes\_V62.tex},
with \(22\,350\) logical source lines and \(823\,395\) bytes,
SHA--256
\[
 \texttt{21a12e67f8673e52e29ffef50fe6c19bf6f46d61e88b0956a08afc142a1d83b5}.
\]
Since the V75 snapshot, SM V58 proves a conditional two-species
rooted tail composition; V59 acquires the first averaged rare-tail
collision and gives a Gaussian counterexample to a conditional
supremum; V60 replaces that false supremum by a guarded/first-owner
partition; V61 proves the all-valence sector with exactly one
primitive tail; and V62 consolidates all repetitions at one physical
tail cell before assigning activities.  Sibling histories then have
disjoint owner sets, so their activities tensor projectively.
Arbitrary multi-tail local cumulants are proved, but their joint
pre-collision spatial response remains explicitly open.
\item
\texttt{Renewal\_NS\_Stability\_Research\_Notes\_V31.tex},
with \(23\,052\) logical source lines and \(887\,537\) bytes,
SHA--256
\[
 \texttt{f1121b62238ad5491bb7cf03635ea9934942edf573ed8faaa9ee79e1d38a6696}.
\]
This file is byte-for-byte unchanged from the V75 audit.  Its exact
formation ladder still permits one payment for one actual removed
flat mark and forbids a second charge on a persistent occurrence.
\end{enumerate}

\begin{assessment}[Type-correct V80 transfer]
\label{ass:V80-transfer}
No SM tail activity, Witten response, or NS heat estimate is a
Yang--Mills operator bound.  The audit changes only the YM proof
order:
\begin{enumerate}[label=\textup{(\roman*)},leftmargin=4em]
\item consolidate repeated uses of one physical coordinate before
      assigning bank currencies;
\item multiply capacities only for banks made disjoint by one common
      exact restriction;
\item keep a unique payer or owner label through recombination; and
\item do not promote separate local estimates to a joint response.
\end{enumerate}
In particular, a physical \(12\)-shared bank below cannot be charged
once through the prefix edge \(p\) and again as a nominally
independent bank.  The terminal factor must either be disjoint and
source-occurring, or remain in the unresolved joint carrier.
\end{assessment}

\begin{firewall}[Companion type boundary]
\label{fire:V80-companion-types}
The proofs below use only the YM physical-incidence partition,
V32/V57/V64 representation dichotomies, V59 centrality, V65 exact
restriction, and the V76--V79 response cards.  No companion analytic
inequality is imported.
\end{firewall}

\section{A cofinal-row trichotomy without a pre-existing strong edge}
\label{sec:V80-cofinal-trichotomy}

V64 stated its full-row trichotomy for a bank inherited from a strong
oriented edge.  Inspection of the proof shows that the terminal
classification itself needs only a complete bank owning a fixed
fraction of one row.  The strong edge was used to construct that
bank upstream, not inside the classification.

\begin{theorem}[Cofinal complete-bank trichotomy]
\label{thm:V80-cofinal-trichotomy}
Let \(E\) be a complete unmarked coordinate bank and let row \(a\)
satisfy
\begin{equation}
 S_a(E):=\sum_{e\in E}a_{a,e}\ge\sigma\Xi_a
\label{eq:V80-root-cofinal}
\end{equation}
for fixed \(\sigma>0\).  For every fixed \(R_0>1\), at least one of
the following occurs:
\begin{enumerate}[label=\textup{(\alph*)},leftmargin=4em]
\item a subbank \(F\subseteq E\) is V32-unbalanced at row \(a\) at
      every coordinate and
\begin{equation}
 S_a(F)\ge\frac{\sigma}{2}\Xi_a;
\label{eq:V80-gap-cofinal}
\end{equation}
\item for some \(c\ne a\), a subbank \(D\subseteq E\) is two-sided
      pointwise comparable between \(a,c\) and owns fixed fractions
      of both \(\Xi_a,\Xi_c\);
\item for some \(c\ne a\),
\begin{equation}
 \Xi_c\ge R_0\Xi_a.
\label{eq:V80-escalation}
\end{equation}
\end{enumerate}
All fractions and comparison constants depend only on
\(G,\sigma,R_0\).
\end{theorem}

\begin{proof}
At every coordinate in \(E\), apply
\Cref{lem:V64-nongapped-partner}.  Let \(E_{\rm gap}\) contain the
coordinates which are unbalanced at row \(a\), and partition the
remainder into at most three partner banks \(E_c\) on which
\[
 a_{c,e}\ge c_G^{\rm p}a_{a,e}.
\]
If \(S_a(E_{\rm gap})\ge\sigma\Xi_a/2\), take
\(F=E_{\rm gap}\).

Otherwise the partner banks carry more than
\(\sigma\Xi_a/2\), so one \(E_c\) satisfies
\[
 S_a(E_c)\ge\frac{\sigma}{6}\Xi_a.
\]
Apply the elementary split used in the proof of
\Cref{thm:V57-bank-or-escalation}: first restrict to
\(a_{c,e}\le C a_{a,e}\) or its complement, then test whether the
retained partner mass owns a fixed fraction of \(\Xi_c\).  The lower
comparison \(a_{c,e}\ge c_G^{\rm p}a_{a,e}\) and root ownership are
the only hypotheses used there.  With the upper threshold chosen
from \(R_0\), the result is either a two-sided comparable subbank
owning both rows or \(\Xi_c\ge R_0\Xi_a\).  These are
\textup{(b)} and \textup{(c)}.
\end{proof}

\begin{corollary}[A maximal row has only two terminal outcomes]
\label{cor:V80-maximal-trichotomy}
If \(\Xi_a=\max_i\Xi_i\), alternative \textup{(c)} is impossible.
Every complete cofinal unmarked row-\(a\) bank therefore contains a
cofinal gapped subbank or an owned comparable subbank.
\end{corollary}

\begin{proof}
For \(R_0>1\), \eqref{eq:V80-escalation} contradicts maximality.
\end{proof}

\section{The gapped terminal factor proves the missing
heavy-composite card}
\label{sec:V80-gap-insertion}

\begin{lemma}[Cofinal heavy-composite gap has \(\chi\)-norm]
\label{lem:V80-gap-chi}
Let \(a\in\{1,2\}\), assume \(R\ge1\), and let \(F\) satisfy
\eqref{eq:V80-gap-cofinal}.  Its complete positive gapped envelope
\(\mathcal G_a(F)\) obeys
\begin{equation}
 \boxed{
 \|\mathcal G_a(F)\|
 \le C_{G,s,\sigma}\chi.}
\label{eq:V80-gap-chi}
\end{equation}
\end{lemma}

\begin{proof}
By \Cref{lem:V64-Casimir-mass-gap},
\[
 \|\mathcal G_a(F)\|
 \le e^{-c_G's\sigma\Xi_a/2}.
\]
Let \(b\in\{3,4\}\) have maximal singleton mass.  The proof of
\Cref{lem:V64-gap-pays-ratio}, applied with \(m=1\), is an
ordinary-norm proof before capacity is taken and gives
\[
 \|\mathcal G_a(F)\|
 \le C_{G,s,\sigma}\frac{\Xi_b}{\Xi_a}.
\]
Since \(\Xi_a\ge\min\{\Xi_1,\Xi_2\}\),
\[
 \frac{\Xi_b}{\Xi_a}
 \le
 \frac{\max\{\Xi_3,\Xi_4\}}
 {\min\{\Xi_1,\Xi_2\}}
 =\chi.
\]
\end{proof}

\begin{theorem}[Source-occurring cofinal gap closes the
heavy-composite word]
\label{thm:V80-gap-HCRC}
Assume \(z\ge\eta>0\), \(R\ge1\), and that a common exact V65
resolution exposes the gapped factor
\(\mathcal G_a(F)\) of
\Cref{lem:V80-gap-chi} as a source-occurring central factor in every
term under consideration.  Then those terms satisfy
\(HCRC_{12|34}\), at every payer placement, and hence obey
\begin{equation}
 \kappa_\otimes(
 \widetilde W_{r,\pi}^*\widetilde W_{r,\pi})
 \le
 C_{s,\eta,\sigma}\mathscr A_{211}^{\,2}.
\label{eq:V80-gap-FCQ}
\end{equation}
\end{theorem}

\begin{proof}
The bank is unmarked and excludes the payer and selected coordinates.
By V59 centrality and the all-payer prefix factorization of
\Cref{lem:V76-prefix-bank-factorization}, each resolved term can be
written
\[
 \widetilde W_{r,\pi,\nu}
 =
 \frac{p_r^-}{\Xi_3\Xi_4}\,
 L_\nu\mathcal G_a(F)Q_\nu,
 \qquad
 \|L_\nu\|+\|Q_\nu\|\le C_s.
\]
Insert \eqref{eq:V80-gap-chi} and take the ordinary norm.  This is
exactly \eqref{eq:V78-HCRC}.  Apply
\Cref{thm:V78-HCRC-suffices} and the fixed finite recombination
\Cref{lem:V58-finite-recombination}.
\end{proof}

\section{Removing the finite exceptional coordinates}
\label{sec:V80-exceptional}

For a fixed selected word let \(Q_{r,\pi}\) contain the selected
joining coordinate, the payer coordinate if distinct, and the
finitely many marked separator coordinates in its literal tree
packet.  The fixed source arity gives
\begin{equation}
 |Q_{r,\pi}|\le q_*.
\label{eq:V80-qstar}
\end{equation}

\begin{lemma}[Unmarked-cofinal or exceptional-heavy]
\label{lem:V80-unmarked-or-exceptional}
Suppose a physical shared-incidence cell \(E_0\) satisfies
\[
 S_a(E_0)\ge\sigma\Xi_a.
\]
Put \(E=E_0\setminus Q_{r,\pi}\).  Then either
\begin{equation}
 S_a(E)\ge\frac{\sigma}{2}\Xi_a,
\label{eq:V80-unmarked-cofinal}
\end{equation}
or some \(e\in E_0\cap Q_{r,\pi}\) satisfies
\begin{equation}
 a_{a,e}\ge\frac{\sigma}{2q_*}\Xi_a.
\label{eq:V80-exceptional-heavy}
\end{equation}
\end{lemma}

\begin{proof}
If \eqref{eq:V80-unmarked-cofinal} fails, the removed coordinates
carry more than \(\sigma\Xi_a/2\).  There are at most \(q_*\) of
them, so one carries at least the average.
\end{proof}

\begin{firewall}[An exceptional coordinate is not an unmarked bank]
\label{fire:V80-exceptional-not-bank}
When \eqref{eq:V80-exceptional-heavy} occurs, the selected or paid
local factor must be estimated with its actual hot-output theorem.
It may not be placed inside the V57/V64 complete unmarked bank or
given a second payer.  V80 records this branch but does not close it.
\end{firewall}

\section{Termination of the unmarked maximal-composite branch}
\label{sec:V80-termination}

\begin{theorem}[Maximal-composite shared-bank reduction]
\label{thm:V80-maximal-reduction}
Consider a \(z\)-nondepleted heavy-composite counterpacket sequence
after the private closure of V79.  Choose
\[
 a\in\{1,2\},
 \qquad
 \Xi_a=\max\{\Xi_1,\Xi_2\}.
\]
After passage to a subsequence, one of the following occurs:
\begin{enumerate}[label=\textup{(\roman*)},leftmargin=4em]
\item a fixed exceptional coordinate satisfies
      \eqref{eq:V80-exceptional-heavy};
\item a common exact resolution exposes a source-occurring cofinal
      row-\(a\) gapped bank, and the packet closes by
      \Cref{thm:V80-gap-HCRC};
\item a common exact resolution exposes an owned comparable bank
      between \(a\) and a singleton row \(3\) or \(4\), and the
      packet closes by the \(\ell=z\) case of
      \Cref{thm:V76-dominant-bank-FCQ};
\item the terminal owned comparable bank joins rows \(1,2\);
\item the gapped or owned terminal bank supplied by
      \Cref{thm:V80-cofinal-trichotomy} is not source-occurring in
      at least one signed restriction term.
\end{enumerate}
\label{eq:V80-maximal-alternatives}
\end{theorem}

\begin{proof}
By \Cref{thm:V79-routing}, row \(a\) has a fixed shared partner cell
\(E_0\) carrying at least \(\Xi_a/4\).  Apply
\Cref{lem:V80-unmarked-or-exceptional} with \(\sigma=1/4\).
The exceptional alternative is \textup{(i)}.

Otherwise the unmarked bank \(E\) owns at least \(\Xi_a/8\).
On the heavy-composite sequence,
\[
 \Xi_a
 =
 \max\{\Xi_1,\Xi_2\}
 >
 \max\{\Xi_3,\Xi_4\}
\]
eventually, so row \(a\) is maximal among all four rows.  Apply
\Cref{cor:V80-maximal-trichotomy}.  If the terminal bank fails the
source-occurrence test of
\Cref{def:V65-source-occurring}, record \textup{(v)}.

For a source-occurring gapped bank use
\Cref{thm:V80-gap-HCRC}, giving \textup{(ii)}.  For a
source-occurring owned bank, its partner is either a singleton or
the other composite row.  In the first case its edge belongs to
\(\mathcal E_z\), and
\Cref{thm:V76-dominant-bank-FCQ} gives \textup{(iii)}.  The second
case is precisely the internal \(12\) alternative \textup{(iv)}.
\end{proof}

\begin{firewall}[The internal \(12\) bank is consolidated with the
prefix]
\label{fire:V80-internal-no-double}
Alternative \textup{(iv)} is not closed by inserting the bank's
quadratic \(\theta_{12}(D)^2\) currency next to the prefix factor
\(p\).  Both arise from the same physical \(12\)-incidence set and a
common restriction may place the prefix covariance on that bank.
The V80 companion audit requires consolidation of this repeated
cell set before any product is taken.  The correct next object is
the exact prefix covariance with its internal protected/gapped
resolution retained in the same square.
\end{firewall}

\begin{corollary}[Sharpened heavy-composite residual]
\label{cor:V80-heavy-residual}
After V80, a \(z\)-nondepleted heavy-composite counterpacket can
persist only through:
\begin{enumerate}[label=\textup{(\alph*)},leftmargin=4em]
\item a mass-cofinal selected/payer exceptional coordinate;
\item a cofinal internal \(12\)-shared bank consolidated with the
      binary prefix; or
\item a precisely identified source-occurrence failure in a signed
      restriction term.
\end{enumerate}
\label{eq:V80-heavy-residual}
\end{corollary}

\begin{proof}
Physical-private mass is excluded by
\Cref{thm:V79-private-HCRC}.  The other two closed alternatives in
\Cref{thm:V80-maximal-reduction} cannot support a counterpacket.
\end{proof}

\section{V80 result ledger and V81 acquisition order}
\label{sec:V80-status}

\begin{theorem}[Third-series V80 result ledger]
\label{thm:V80-results}
Relative to V1--V79, V80 proves:
\begin{enumerate}[label=\textup{(V80.\arabic*)},leftmargin=6.3em]
\item the compulsory SM V62 and NS V31 audit with exact
      fingerprints;
\item the no-duplicate-cell/no-separate-response proof-order rule;
\item a cofinal complete-row bank trichotomy which does not require
      an incoming strong normalized edge;
\item elimination of mass escalation at a maximal row;
\item the missing \(\chi\)-norm from a cofinal heavy-composite
      Casimir-gapped bank;
\item exact central insertion of that gain and proof of
      \(HCRC_{12|34}\);
\item an unmarked-cofinal versus finite-exceptional-heavy dichotomy;
\item closure of every source-occurring gapped or singleton-partner
      owned terminal bank at the full atlas scale; and
\item reduction of the remaining \(z\)-nondepleted
      heavy-composite sector to three explicit operator-bearing
      branches.
\end{enumerate}
\end{theorem}

\begin{proof}
These are
\Cref{sec:V80-audit,
ass:V80-transfer,
thm:V80-cofinal-trichotomy,
cor:V80-maximal-trichotomy,
lem:V80-gap-chi,
thm:V80-gap-HCRC,
lem:V80-unmarked-or-exceptional,
thm:V80-maximal-reduction,
cor:V80-heavy-residual}.
\end{proof}

\begin{programme}[V81: consolidated internal-prefix square]
\label{prog:V80-V81}
\begin{enumerate}[leftmargin=3em]
\item Fix alternative \textup{(b)} of
      \eqref{eq:V80-heavy-residual} and restrict the exact V10
      binary covariance telescope to the internal \(12\) bank and
      its complement before taking norms.
\item Resolve the internal bank into protected/gapped sectors inside
      both occurrences of the V44 prefix square.  Consolidate equal
      physical coordinate atoms before assigning any capacity.
\item Determine whether the cofinal gapped sector supplies
      \(\chi^2\) before the final square root, or whether its heat
      factor has already been spent in the one-prefix response.
\item For the protected sector compute the product-state capacity of
      the restricted covariance, not of the isolated invariant
      projection.  Retain off-diagonal prefix-coordinate terms.
\item Treat the exceptional branch with the V38 selected hot-output
      moments or the V29 paid-private moment according to the actual
      coordinate role; do not move an exceptional coordinate into an
      unmarked bank.
\item In source-occurrence-failure terms, return to the signed V14
      covariance restriction identity and identify the first term
      in which the desired positive factor cannot be exposed.
\item The next companion audit is V85.
\end{enumerate}
\end{programme}

\begin{firewall}[Claim boundary after third-series V80]
\label{fire:V80-claim-boundary}
V80 closes the source-occurring cofinal gapped and compatible owned
branches of the heavy-composite sector and isolates its only three
remaining mechanisms.  It does not close the internal \(12\)
consolidated square, exceptional-heavy or source-occurrence-failure
branches, total edge depletion, the complete \(2+1+1\) source or its
global aggregation, other quartic source families, higher MTA, WUFC,
CPM, cutoff removal, reflection positivity,
Osterwalder--Schrader reconstruction, construction of the continuum
Yang--Mills measure, or a uniform positive continuum spectral gap.
The full Yang--Mills mass gap has not been proved.
\end{firewall}

\paragraph{Parent-version record.}
V80 was formed from
\texttt{Renewal\_Yang\_Mills\_Mass\_Gap\_Research\_Notes\_V79.tex}.
The V79 parent had \(38\,711\) logical source lines,
\(1\,376\,773\) bytes, and SHA--256
\[
 \texttt{e4b1e5be0624ba23e5ddc6435497942fcf31e80a691346ecfe7cd62317b930f1}.
\]
The next compulsory companion audit is V85.

% =====================================================================
% V81 APPENDIX: DIFFUSE PROTECTED INTERNAL-PREFIX CLOSURE
% =====================================================================

\chapter{V81: Prefix Depletion Forces Protected Internal Banks to
Tensorize Strongly}
\label{ch:V81-protected-internal}

\section{Effective diffuseness of an owned internal \(12\) bank}
\label{sec:V81-diffuse}

Let \(D\) be an owned comparable unmarked bank between rows \(1,2\).
Write
\begin{equation}
 N:=|D|,\qquad
 S:=\sum_{e\in D}a_{1,e},\qquad
 T:=\sum_{e\in D}a_{2,e},\qquad
 H:=\sum_{e\in D}a_{1,e}a_{2,e}.
\label{eq:V81-bank-data}
\end{equation}
Thus, for fixed \(c,C,\sigma_1,\sigma_2>0\),
\begin{equation}
 ca_{1,e}\le a_{2,e}\le Ca_{1,e},
 \qquad
 S\ge\sigma_1\Xi_1,
 \qquad
 T\ge\sigma_2\Xi_2.
\label{eq:V81-owned-comparable}
\end{equation}

\begin{lemma}[Prefix depletion bounds the effective concentration]
\label{lem:V81-effective-concentration}
Put
\begin{equation}
 q_D^{(2)}
 :=
 \frac{\sum_{e\in D}a_{1,e}^2}{S^2}.
\label{eq:V81-q2}
\end{equation}
Then
\begin{equation}
 \boxed{
 q_D^{(2)}\le C_0p,}
\label{eq:V81-q2-by-p}
\end{equation}
where \(C_0\) depends only on the constants in
\eqref{eq:V81-owned-comparable}.  In particular,
\begin{equation}
 \max_{e\in D}\frac{a_{1,e}}S
 \le\sqrt{C_0p}.
\label{eq:V81-max-fraction}
\end{equation}
\end{lemma}

\begin{proof}
Pointwise lower comparability gives
\[
 H\ge c\sum_{e\in D}a_{1,e}^2.
\]
Also \(T\le CS\), while ownership gives
\[
 \Xi_1\Xi_2
 \le
 \frac S{\sigma_1}\frac T{\sigma_2}
 \le
 \frac{C}{\sigma_1\sigma_2}S^2.
\]
Since \(H\le H_{12}\),
\[
 p=\frac{H_{12}}{\Xi_1\Xi_2}
 \ge
 \frac{c\sigma_1\sigma_2}{C}
 \frac{\sum_ea_{1,e}^2}{S^2}.
\]
This is \eqref{eq:V81-q2-by-p}.  Finally
\[
 \left(\max_e\frac{a_{1,e}}S\right)^2
 \le q_D^{(2)}.
\]
\end{proof}

\begin{corollary}[No finite-coordinate protected obstruction on a
depleted prefix]
\label{cor:V81-many-effective-coordinates}
There is a structural \(p_0>0\) such that \(p\le p_0\) implies
\[
 \max_{e\in D}a_{1,e}\le\frac18S
\]
and hence \(N\ge8\).
\end{corollary}

\begin{proof}
Choose \(p_0\le(64C_0)^{-1}\) in
\eqref{eq:V81-max-fraction}.  At least eight nonzero summands are
then required to sum to \(S\).
\end{proof}

\section{A joint cardinality--dimension optimization}
\label{sec:V81-optimization}

\begin{lemma}[Four-largest-coordinate product lemma]
\label{lem:V81-four-largest}
Let \(b_1,\ldots,b_N>0\), put \(S_b=\sum_eb_e\), and assume
\begin{equation}
 \max_eb_e\le\varepsilon S_b,
 \qquad
 0<\varepsilon\le\frac18.
\label{eq:V81-diffuse-sequence}
\end{equation}
If \(0<\rho<1\) and
\[
 0\le r_e\le\min\{\rho,Cb_e^{-1/2}\},
\]
then
\begin{equation}
 \boxed{
 \prod_{e=1}^Nr_e
 \le C_{\rho,C}S_b^{-2}.}
\label{eq:V81-product-S2}
\end{equation}
\end{lemma}

\begin{proof}
Reorder so that \(b_1\ge\cdots\ge b_N\).  Assumption
\eqref{eq:V81-diffuse-sequence} implies \(N\ge8\).  For
\(1\le k\le4\), the first \(k-1\) entries carry at most
\(3\varepsilon S_b\le3S_b/8\).  Hence
\[
 (N-k+1)b_k
 \ge
 \sum_{j=k}^Nb_j
 \ge\frac58S_b,
\]
and therefore
\[
 b_k\ge\frac{5S_b}{8N}.
\]
Use the dimension branch of the bound on the four largest entries
and the uniform branch on the others:
\[
 \prod_{e=1}^Nr_e
 \le
 C^4(b_1b_2b_3b_4)^{-1/2}\rho^{N-4}
 \le
 C'\frac{N^2\rho^{N-4}}{S_b^2}.
\]
Since \(\sup_{N\ge8}N^2\rho^{N-4}<\infty\), this proves
\eqref{eq:V81-product-S2}.
\end{proof}

\begin{remark}[Why both inputs are necessary]
\label{rem:V81-two-inputs}
The inverse-dimension estimate alone has the sharp one-coordinate
power \(S_b^{-1/2}\), as shown in V69.  The uniform
\(\rho^{N}\) estimate alone does not see arbitrarily large weights
on a fixed number of coordinates.  Prefix depletion supplies the
missing diffuseness: it prevents a fixed number of weights from
carrying \(S_b\), allowing four dimension factors and the remaining
uniform factors to be used simultaneously.
\end{remark}

\section{Protected capacity at two inverse heavy-row powers}
\label{sec:V81-protected-capacity}

Let
\[
 P_D^{\rm prot}:=\bigotimes_{e\in D}P_e^{\rm prot}
\]
be the complete all-protected invariant projector, viewed across
the row-\(1|1^c\) product cut.

\begin{theorem}[Diffuse owned protected bank pays two row-mass
powers]
\label{thm:V81-protected-two-powers}
For \(p\le p_0\),
\begin{equation}
 \boxed{
 \kappa_{1|1^c}(P_D^{\rm prot})
 \le
 \|(P_D^{\rm prot})^{\Gamma_{1^c}}\|
 \le
 \frac{C}{\Xi_1^2}.}
\label{eq:V81-protected-Xi2}
\end{equation}
The symmetric row-\(2\) statement also holds.  On the
heavy-composite cell \(R\ge1\),
\begin{equation}
 \kappa_\otimes(P_D^{\rm prot})
 \le C\chi^2.
\label{eq:V81-protected-chi2}
\end{equation}
\end{theorem}

\begin{proof}
At every coordinate, the partial-transpose calculation and Weyl
dimension bound
\Cref{cor:V69-one-coordinate-decay} give
\[
 \|(P_e^{\rm prot})^{\Gamma_{1^c}}\|
 \le
 \min\{\rho_G,C_Ga_{1,e}^{-1/2}\}.
\]
Partial transpose tensorizes exactly over \(D\).  By
\Cref{lem:V81-effective-concentration,
cor:V81-many-effective-coordinates}, the sequence
\((a_{1,e})_{e\in D}\) satisfies
\eqref{eq:V81-diffuse-sequence}.  Apply
\Cref{lem:V81-four-largest}:
\[
 \|(P_D^{\rm prot})^{\Gamma_{1^c}}\|
 \le C_GS^{-2}
 \le C_{G,\sigma_1}\Xi_1^{-2}.
\]
The partial-transpose trace estimate used in
\Cref{thm:V68-protected-product} bounds product-state capacity by
this norm.  This proves \eqref{eq:V81-protected-Xi2}; exchange rows
\(1,2\) for the symmetric assertion.

Let \(a\in\{1,2\}\) have maximal composite mass.  Use its version of
\eqref{eq:V81-protected-Xi2}.  Since
\(\Xi_a\ge\min\{\Xi_1,\Xi_2\}\) and
\(\max\{\Xi_3,\Xi_4\}\ge a_*\),
\[
 \Xi_a^{-2}
 \le
 a_*^{-2}
 \left(
 \frac{\max\{\Xi_3,\Xi_4\}}
 {\min\{\Xi_1,\Xi_2\}}
 \right)^2
 =
 a_*^{-2}\chi^2.
\]
Full four-row product states are a subset of the product states
across the coarser cut \(a|a^c\), giving
\eqref{eq:V81-protected-chi2}.
\end{proof}

\begin{corollary}[The V69 half-power sharpness test is respected]
\label{cor:V81-V69-respected}
The one-coordinate \(SU(2)\) family in
\Cref{prop:V69-half-power-sharp} does not contradict
\eqref{eq:V81-protected-Xi2}: an owned internal bank with one
effective coordinate has \(q_D^{(2)}=1\), and therefore cannot lie
on \(p\le p_0\).
\end{corollary}

\begin{proof}
For one nonzero weight,
\(\sum_ea_e^2/(\sum_ea_e)^2=1\).  Use
\eqref{eq:V81-q2-by-p} and the choice of \(p_0\).
\end{proof}

\section{The consolidated protected prefix term}
\label{sec:V81-consolidated}

\begin{lemma}[Protected restriction is supported on the complete
invariant projector]
\label{lem:V81-protected-support}
Restrict the exact binary prefix covariance simultaneously to \(D\)
and \(D^c\), and resolve every local physical output on \(D\).
Every term in which all coordinates of \(D\) choose their protected
output has a factor \(Z_D\) satisfying
\begin{equation}
 Z_D=P_D^{\rm prot}Z_DP_D^{\rm prot},
 \qquad
 \|Z_D\|\le C_s.
\label{eq:V81-protected-support}
\end{equation}
Consequently
\begin{equation}
 Z_D^*Z_D\preccurlyeq C_sP_D^{\rm prot}.
\label{eq:V81-protected-square}
\end{equation}
This statement includes both possibilities that the prefix
covariance is localized to \(D\) or to \(D^c\).
\end{lemma}

\begin{proof}
The exact V14 restriction identity writes the prefix covariance as
a coefficient-one sum in which one covariance is localized to one
side and the other side carries a complete endpoint moment.  Resolve
the local physical outputs on \(D\) before taking a norm.  On the
all-protected term every local range is the invariant range
\(P_e^{\rm prot}\); their intersection over independent coordinates
is the tensor product \(P_D^{\rm prot}\).  Thus every local moment,
covariance endpoint, recoupling, and compression meeting \(D\) is
supported between two copies of \(P_D^{\rm prot}\).

All complete moments and physical compressions are contractions.
The binary covariance has two moment terms and the number of fixed
separators is support-independent, so the operator between the two
projectors has norm at most \(C_s\).  This gives
\eqref{eq:V81-protected-support}.  Then
\[
 Z_D^*Z_D
 =
 P_D^{\rm prot}Z_D^*Z_DP_D^{\rm prot}
 \preccurlyeq
 \|Z_D\|^2P_D^{\rm prot},
\]
proving \eqref{eq:V81-protected-square}.
\end{proof}

\begin{firewall}[Consolidation, not a second \(p\)-charge]
\label{fire:V81-no-double-p}
The proof of \Cref{lem:V81-protected-support} does not extract a
separate protected bank beside the binary prefix.  It resolves the
prefix covariance itself on the common physical coordinate set.
The scalar \(p\) is used only in
\Cref{lem:V81-effective-concentration} to prove diffuseness; no
factor \(p\,\theta_{12}(D)\) is formed.  This is consistent with the
V80 companion audit and the V44 one-debit ceiling.
\end{firewall}

\begin{theorem}[All-protected internal \(12\) branch satisfies full
\(FCQ_{211}\)]
\label{thm:V81-protected-internal-FCQ}
Fix \(z\ge\eta>0\), \(R\ge1\), and \(p\le p_0\).  Suppose a common
exact resolution of the complete selected word exposes an owned
comparable internal \(12\) bank \(D\), and consider the terms which
are all-protected on \(D\).  Then, for every payer placement and
both square orientations,
\begin{equation}
 \boxed{
 \kappa_\otimes(
 \widetilde W_{r,\pi}^{\,*}
 \widetilde W_{r,\pi})
 \le
 C_{s,\eta}\mathscr A_{211}^{\,2}.}
\label{eq:V81-protected-internal-FCQ}
\end{equation}
The assertion holds after the fixed finite recombination of all such
protected terms.
\end{theorem}

\begin{proof}
For each resolved term, keep \(Z_D\) from
\Cref{lem:V81-protected-support} inside the complete physical word.
V59 centrality moves only the outer protected projectors through the
other physical source intertwiners.  All remaining complete factors,
including the actual once-paid factor, have norm at most \(C_s\).
With
\[
 V:=\frac1{\Xi_1\Xi_2\Xi_3\Xi_4},
\]
the appropriate square is therefore bounded in positive order by
\[
 C_sV^2P_D^{\rm prot}
\]
using \eqref{eq:V81-protected-square}.  Take product-state capacity
and apply \eqref{eq:V81-protected-chi2}:
\[
 \kappa_\otimes(
 \widetilde W_{r,\pi}^{\,*}
 \widetilde W_{r,\pi})
 \le C_sV^2\chi^2.
\]
The heavy-composite defect estimate
\eqref{eq:V78-volume-defect}, together with \(z\ge\eta\), gives
\[
 V\le C_{s,\eta}R\mathscr A_{211}.
\]
Since \(R\chi=1\), the last two displays prove
\eqref{eq:V81-protected-internal-FCQ}.  The adjoint orientation is
identical.  Apply
\Cref{lem:V58-finite-recombination} to the fixed exact resolution.
\end{proof}

\section{The residual internal sector after protected closure}
\label{sec:V81-residual}

\begin{corollary}[Only non-all-protected internal terms remain]
\label{cor:V81-internal-residual}
On a heavy-composite counterpacket, alternative
\textup{(b)} of \eqref{eq:V80-heavy-residual} can persist only in
restriction terms whose internal \(12\) bank has at least one
nonprotected local output, or in terms where the exact
source-occurrence/support statement
\eqref{eq:V81-protected-support} fails.
\end{corollary}

\begin{proof}
The counterpacket has \(p\to0\), so eventually \(p\le p_0\).
Every all-protected term satisfying the exact support statement is
closed by \Cref{thm:V81-protected-internal-FCQ}.
\end{proof}

\begin{assessment}[Why the remaining gapped term is not yet
automatic]
\label{ass:V81-gapped-open}
The V68 one-failure majorant gives exponential decay in the number
of bank coordinates, while V81 uses both that cardinality decay and
four representation-dimension factors.  A non-all-protected term
may have a low nontrivial output at one or more coordinates.
Replacing every such output complement by a norm-one projection
loses its input-dimension suppression; retaining only the uniform
heat gap may be insufficient when the comparable input weights grow
faster than exponentially in the bank length.  The next estimate
must therefore keep the input/output isotypic dimensions of the
nontrivial sectors, rather than citing the coarse V57
\(\theta_{12}(D)^2\) currency.
\end{assessment}

\section{V81 result ledger and V82 acquisition order}
\label{sec:V81-status}

\begin{theorem}[Third-series V81 result ledger]
\label{thm:V81-results}
Relative to V1--V80, V81 proves:
\begin{enumerate}[label=\textup{(V81.\arabic*)},leftmargin=6.3em]
\item effective \(\ell^2\)-diffuseness of every owned comparable
      internal \(12\) bank from prefix depletion;
\item exclusion of fixed-coordinate concentration for sufficiently
      small \(p\);
\item a joint cardinality--dimension product inequality using the
      four largest weights;
\item two inverse heavy-row mass powers for the complete protected
      projector;
\item compatibility of that result with the sharp one-coordinate
      \(SU(2)\) half-power obstruction;
\item exact protected support of the consolidated restricted prefix
      covariance;
\item full atlas closure of every all-protected internal \(12\)
      restriction term without a duplicated prefix charge; and
\item reduction of the internal-bank residual to non-all-protected
      output or an explicit source-support failure.
\end{enumerate}
\end{theorem}

\begin{proof}
These are
\Cref{lem:V81-effective-concentration,
cor:V81-many-effective-coordinates,
lem:V81-four-largest,
thm:V81-protected-two-powers,
cor:V81-V69-respected,
lem:V81-protected-support,
thm:V81-protected-internal-FCQ,
cor:V81-internal-residual}.
\end{proof}

\begin{programme}[V82: variable-dimension nonprotected outputs]
\label{prog:V81-V82}
\begin{enumerate}[leftmargin=3em]
\item On each coordinate of the internal \(12\) bank, retain the
      complete output isotypic projection \(P_{U_e}\) rather than
      replacing \(I-P_e^{\rm prot}\) by the identity.
\item Prove an input/output partial-transpose bound depending on
      \(\dim U_e/\dim V_{1,e}\), uniformly in multiplicity, for a
      fixed compact group.  Test first the exact \(SU(2)\) formula
      from V48.
\item Combine low-output dimension suppression with high-output heat
      decay in one local minimum, then tensorize that minimum over
      the diffuse bank.
\item Reuse the four-largest-coordinate optimization only after the
      nonprotected local factor has been proved.  Do not infer it
      from the all-protected projector.
\item Treat a payer inside the bank by keeping the complete paid
      spectral coefficient before summing output sectors.
\item If a fixed low-output sector saturates only the V48
      half-power, determine whether diffuseness over four effective
      coordinates still supplies \(\chi^2\); otherwise isolate the
      exact output-rank obstruction.
\item The next companion audit remains V85.
\end{enumerate}
\end{programme}

\begin{firewall}[Claim boundary after third-series V81]
\label{fire:V81-claim-boundary}
V81 closes the all-protected part of the consolidated internal
\(12\) bank on every depleted-prefix heavy-composite counterpacket.
It does not close its non-all-protected spectral sectors,
exceptional-heavy or source-occurrence-failure branches, total edge
depletion, the complete \(2+1+1\) source or its global aggregation,
other quartic source families, higher MTA, WUFC, CPM, cutoff removal,
reflection positivity, Osterwalder--Schrader reconstruction,
construction of the continuum Yang--Mills measure, or a uniform
positive continuum spectral gap.  The full Yang--Mills mass gap has
not been proved.
\end{firewall}

\paragraph{Parent-version record.}
V81 was formed from
\texttt{Renewal\_Yang\_Mills\_Mass\_Gap\_Research\_Notes\_V80.tex}.
The V80 parent had \(39\,162\) logical source lines,
\(1\,393\,575\) bytes, and SHA--256
\[
 \texttt{0bc8392720c63f4d51b2b0761fecb5c2ec48d19c20d4b2553a2b3f514b92cb39}.
\]
The next compulsory companion audit is V85.

% =====================================================================
% V82 APPENDIX: LOW-OUTPUT MARGINALS AND THE PT HEAT INTERFACE
% =====================================================================

\chapter{V82: Multiplicity-Uniform Low-Output Suppression and the
Remaining Partial-Transpose Interface}
\label{ch:V82-output-PT}

\section{One local isotypic output across a two-row cut}
\label{sec:V82-isotypic}

Let \(V,W\) be nontrivial irreducible unitary \(G\)-modules.  Resolve
\[
 V\otimes W
 =
 \bigoplus_U U\otimes\mathcal M_U,
 \qquad
 m_U:=\dim\mathcal M_U,
\]
and let \(P_U\) denote the orthogonal projection onto the complete
\(U\)-isotypic summand, including its multiplicity space.  Write
\(d_V=\dim V\), \(d_W=\dim W\), and \(d_U=\dim U\).

\begin{lemma}[Exact scalar marginal of an isotypic projection]
\label{lem:V82-isotypic-marginal}
\begin{equation}
 \boxed{
 \operatorname{Tr}_W P_U
 =
 \frac{d_Um_U}{d_V}I_V.}
\label{eq:V82-isotypic-marginal}
\end{equation}
Consequently
\begin{equation}
 \kappa_{V|W}(P_U)
 \le
 \min\left\{1,\frac{d_Um_U}{d_V}\right\}.
\label{eq:V82-isotypic-capacity}
\end{equation}
\end{lemma}

\begin{proof}
The isotypic projection commutes with
\(V(g)\otimes W(g)\) for every \(g\in G\).  Invariance of partial
trace under conjugation on \(W\) therefore gives
\[
 V(g)(\operatorname{Tr}_WP_U)V(g)^*
 =
 \operatorname{Tr}_WP_U.
\]
Schur's lemma makes the marginal a scalar multiple of \(I_V\).
Its trace is
\(\operatorname{Tr}P_U=d_Um_U\), proving
\eqref{eq:V82-isotypic-marginal}.

For density matrices \(\rho,\sigma\), one has
\(0\preccurlyeq\sigma\preccurlyeq I_W\).  Hence
\[
 \operatorname{Tr}[P_U(\rho\otimes\sigma)]
 \le
 \operatorname{Tr}[P_U(\rho\otimes I_W)]
 =
 \operatorname{Tr}[\rho\,\operatorname{Tr}_WP_U]
 =
 \frac{d_Um_U}{d_V}.
\]
The projection norm supplies the additional upper bound one.
\end{proof}

\begin{lemma}[Multiplicity is absorbed by the larger input
dimension]
\label{lem:V82-multiplicity}
If \(d_V\ge d_W\), then
\begin{equation}
 m_Ud_V\le d_Ud_W
\label{eq:V82-multiplicity-dimension}
\end{equation}
and
\begin{equation}
 \boxed{
 \kappa_{V|W}(P_U)
 \le
 \min\left\{1,\frac{d_U^2}{d_V}\right\}.}
\label{eq:V82-low-output}
\end{equation}
After exchanging \(V,W\), the denominator can always be taken to be
\(\max\{d_V,d_W\}\).
\end{lemma}

\begin{proof}
Frobenius reciprocity gives
\[
 m_U
 =
 \dim\operatorname{Hom}_G(U,V\otimes W)
 =
 \dim\operatorname{Hom}_G(V,U\otimes W^*).
\]
Thus \(U\otimes W^*\) contains \(m_U\) orthogonal copies of \(V\).
Comparing dimensions yields
\[
 m_Ud_V\le d_Ud_W.
\]
Insert this into \eqref{eq:V82-isotypic-capacity}:
\[
 \frac{d_Um_U}{d_V}
 \le
 \frac{d_U^2d_W}{d_V^2}
 \le
 \frac{d_U^2}{d_V}.
\]
The symmetric statement follows by interchanging the tensor legs.
\end{proof}

\begin{remark}[Recovery of the \(SU(2)\) fixed-output scale]
\label{rem:V82-SU2}
For \(V=W=V_j\) and fixed output \(U=V_\ell\),
\eqref{eq:V82-isotypic-capacity} gives
\[
 \kappa(P_{j,\ell})
 \le\frac{2\ell+1}{2j+1},
\]
because the \(SU(2)\) tensor product is multiplicity free.  This is
the V48 upper bound.  The general estimate
\eqref{eq:V82-low-output} is weaker by one factor \(d_U\), but is
uniform in arbitrary tensor-product multiplicity.
\end{remark}

\section{Heat summation preserves the inverse input dimension}
\label{sec:V82-heat-sum}

For fixed \(s>0\), put
\begin{equation}
 M_{V,W}(s)
 :=
 \sum_{U\subset V\otimes W}
 e^{-sC_2(U)}P_U.
\label{eq:V82-local-heat}
\end{equation}

\begin{theorem}[Multiplicity-uniform local heat capacity]
\label{thm:V82-heat-capacity}
There is a finite constant
\begin{equation}
 K_G(s):=\sum_{U\in\widehat G}d_U^2e^{-sC_2(U)}<\infty
\label{eq:V82-heat-trace}
\end{equation}
such that
\begin{equation}
 \boxed{
 \kappa_{V|W}(M_{V,W}(s))
 \le
 \frac{K_G(s)}{\max\{d_V,d_W\}}.}
\label{eq:V82-heat-capacity}
\end{equation}
\end{theorem}

\begin{proof}
Assume \(d_V\ge d_W\); the other case is symmetric.  Positivity and
subadditivity of product-state capacity, followed by
\eqref{eq:V82-low-output}, give
\[
 \kappa_{V|W}(M_{V,W}(s))
 \le
 \sum_{U\subset V\otimes W}
 e^{-sC_2(U)}\frac{d_U^2}{d_V}
 \le
 \frac1{d_V}
 \sum_{U\in\widehat G}d_U^2e^{-sC_2(U)}.
\]
The last series is the heat trace of the compact-group Laplacian at
positive time and is finite.  This proves the assertion.
\end{proof}

\begin{corollary}[The desired Casimir half-power is present before
tensorization]
\label{cor:V82-local-half-power}
If the larger-dimensional input has Casimir mass \(a_{\max}\), then
\begin{equation}
 \kappa_{V|W}(M_{V,W}(s))
 \le C_{G,s}a_{\max}^{-1/2}.
\label{eq:V82-local-half-power}
\end{equation}
\end{corollary}

\begin{proof}
Apply \Cref{lem:V69-Weyl-dimension} to the larger-dimensional input
module.  Its dimension is at least the dimension of the other input
by definition and at least \(c_Ga_{\max}^{1/2}\) when
\(a_{\max}\) denotes its own Casimir mass.
\end{proof}

\begin{firewall}[Dimension order and Casimir order need not agree]
\label{fire:V82-dimension-Casimir-order}
The symbol \(a_{\max}\) in
\eqref{eq:V82-local-half-power} refers to the Casimir mass of the
input selected by maximal \emph{dimension}; it is not automatically
the larger of the two Casimir masses for a higher-rank group.
Pointwise comparable Casimirs alone do not give a uniform ratio of
Weyl dimensions in different highest-weight directions.  V82 uses
only the lower bound for the selected module.
\end{firewall}

\section{Why the local capacity theorem does not yet tensorize}
\label{sec:V82-tensorization}

The bank product is tested against states which may be arbitrarily
entangled among coordinates within the row-\(1\) side and within
its complement.  A product of one-coordinate product capacities is
therefore not a proved upper bound.  V68 solved this for invariant
projections by computing a partial-transpose norm exactly.

\begin{definition}[Output-isotypic partial-transpose heat card]
\label{def:V82-OIPT}
The local \(OIPT_G(s)\) card is the assertion that, after choosing
the input tensor leg of larger dimension and transposing the other
leg,
\begin{equation}
 \boxed{
 \|M_{V,W}(s)^\Gamma\|
 \le
 \frac{C_{G,s}}{\max\{d_V,d_W\}}.}
\label{eq:V82-OIPT}
\end{equation}
The bound must include the entire multiplicity space of every output
isotype.
\end{definition}

\begin{proposition}[The proved marginal bound has the correct scale
but not the required tensor stability]
\label{prop:V82-marginal-not-PT}
\Cref{thm:V82-heat-capacity} proves the right side of
\eqref{eq:V82-OIPT} after replacing partial-transpose norm by
one-coordinate product-state capacity.  No result through V81
implies that replacement in the opposite direction for a general
nonprotected output sum.
\end{proposition}

\begin{proof}
The first assertion is exactly
\eqref{eq:V82-heat-capacity}.  Product-state capacity tests only
positive product densities.  Partial-transpose norm tests the
operator against all trace-class directions after transpose and is
strictly stronger in general.  The V68 tensorization proof computed
that stronger norm only for the invariant projection; its
\Cref{fire:V68-abstract-scope} explicitly withdraws tensorization
from scalar one-coordinate capacity data alone.  The V58 entangling
separator example gives the same general obstruction.  Hence no
logical implication to \eqref{eq:V82-OIPT} has been proved.
\end{proof}

\section{What the \(OIPT\) card would close}
\label{sec:V82-compiler}

At a coordinate \(e\), let
\[
 B_e(s):=P_e^{\rm prot}
 +r_s(I-P_e^{\rm prot}),
 \qquad 0<r_s<1,
\]
be the V68 tensor-stable uniform heat majorant.  Its
partial-transpose norm is at most
\[
 \rho_s:=r_s+(1-r_s)\rho_G<1.
\]

\begin{theorem}[\(OIPT_G(s)\) closes the complete internal heat
bank]
\label{thm:V82-OIPT-compiler}
Assume \eqref{eq:V82-OIPT}.  Let \(D\) be the diffuse owned
comparable internal \(12\) bank of V81, and suppose its physical
local module at each coordinate is a two-row tensor product to which
\eqref{eq:V82-OIPT} applies.  Then, for \(p\le p_0\),
\begin{equation}
 \kappa_\otimes\!\left(
 \bigotimes_{e\in D}M_e(s)\right)
 \le
 \frac{C_{G,s}}{\Xi_a^2}
 \le C_{G,s}\chi^2,
\label{eq:V82-bank-chi2}
\end{equation}
where \(a\in\{1,2\}\) is a row whose input dimension is maximal at
each of the four selected largest-weight coordinates.  Consequently
the non-all-protected internal sector satisfies
\(FCQ_{211}\) by the V81 consolidated-square proof.
\end{theorem}

\begin{proof}
Order the row-\(a\) weights on \(D\) and select the four largest.
There are only two possible larger-dimension input legs, so after a
fixed finite split one row supplies at least two, but the stated
hypothesis selects a common row \(a\) on all four; this is the exact
orientation needed for the displayed conclusion.

On each of the four selected coordinates use
\eqref{eq:V82-OIPT}; by
\Cref{lem:V69-Weyl-dimension} the corresponding factor is at most
\(C_{G,s}a_{a,e}^{-1/2}\).  On every other coordinate use the V68
majorant \(B_e(s)\), whose partial-transpose norm is at most
\(\rho_s\).  Positivity gives
\[
 \bigotimes_{e\in D}M_e(s)
 \preccurlyeq
 \left(\bigotimes_{e\ {\rm selected}}M_e(s)\right)
 \otimes
 \left(\bigotimes_{e\ {\rm other}}B_e(s)\right).
\]
The product-state expectation of the right side is bounded by the
norm of its full partial transpose.  That norm tensorizes and is at
most
\[
 C_{G,s}^4
 (a_{a,1}a_{a,2}a_{a,3}a_{a,4})^{-1/2}
 \rho_s^{N-4}.
\]
The proof of \Cref{lem:V81-four-largest}, together with ownership,
bounds this by \(C_{G,s}\Xi_a^{-2}\).  The \(\chi^2\) comparison is
the last step of
\Cref{thm:V81-protected-two-powers}.  Since the non-all-protected
positive sector is bounded by the complete heat product, repeat the
central consolidated-square argument in
\Cref{thm:V81-protected-internal-FCQ}.
\end{proof}

\begin{firewall}[The common orientation is also an open detail]
\label{fire:V82-orientation}
For a higher-rank group, the larger-dimensional member of a
pointwise Casimir-comparable pair may switch between rows across
coordinates.  A pigeonhole argument supplies only two of the four
largest weights to one row, which is insufficient for the
\(\Xi^{-2}\) exponent in the present proof.  Thus V82 records two
distinct missing inputs: the \(OIPT_G(s)\) estimate itself and a
four-coordinate common-orientation argument, or a symmetric
replacement using dimensions from both rows.
\end{firewall}

\section{V82 result ledger and V83 acquisition order}
\label{sec:V82-status}

\begin{theorem}[Third-series V82 result ledger]
\label{thm:V82-results}
Relative to V1--V81, V82 proves:
\begin{enumerate}[label=\textup{(V82.\arabic*)},leftmargin=6.3em]
\item the exact scalar marginal of an arbitrary complete isotypic
      projection, including multiplicity;
\item a Frobenius-reciprocity multiplicity bound absorbed by the
      larger input dimension;
\item a multiplicity-uniform fixed-output product-capacity estimate;
\item summation of all heat-weighted outputs at the inverse larger
      input-dimension scale;
\item recovery of the expected Casimir half-power before
      coordinate tensorization;
\item identification of the exact partial-transpose heat card needed
      for entanglement-safe tensorization;
\item proof that the established marginal estimate does not imply
      that card; and
\item a compiler showing that \(OIPT_G(s)\), together with a common
      four-coordinate orientation, closes the non-all-protected
      internal heat bank.
\end{enumerate}
\end{theorem}

\begin{proof}
These are
\Cref{lem:V82-isotypic-marginal,
lem:V82-multiplicity,
thm:V82-heat-capacity,
cor:V82-local-half-power,
def:V82-OIPT,
prop:V82-marginal-not-PT,
thm:V82-OIPT-compiler,
fire:V82-orientation}.
\end{proof}

\begin{programme}[V83: partial transpose of a fixed output isotype]
\label{prog:V82-V83}
\begin{enumerate}[leftmargin=3em]
\item Write \(P_U=\sum_{\alpha=1}^{m_U}T_\alpha T_\alpha^*\) for
      orthogonal Clebsch--Gordan isometries and compute
      \((T_\alpha T_\alpha^*)^\Gamma\) before applying a triangle
      inequality.
\item Determine whether the orthogonality of multiplicity channels
      survives partial transpose or whether their operator norms add.
\item For \(SU(2)\), evaluate the fixed-\(\ell\) block through the
      recoupling/\(6j\) matrix and test the scale
      \((2\ell+1)/(2j+1)\).
\item Sum the output heat only after the fixed-output
      partial-transpose estimate is established; do not infer
      \eqref{eq:V82-OIPT} from
      \eqref{eq:V82-heat-capacity}.
\item If multiplicities force a loss, seek a symmetric two-row
      Schatten bound which can use two large weights from each row
      in the V81 optimization.
\item Keep payer coefficients inside the output sum and test the
      low-output sector separately from the exponentially damped
      high-output sector.
\item The next companion audit remains V85.
\end{enumerate}
\end{programme}

\begin{firewall}[Claim boundary after third-series V82]
\label{fire:V82-claim-boundary}
V82 proves the complete one-coordinate low-output and heat-summed
product-capacity estimates, including arbitrary multiplicity.  It
does not prove their partial-transpose analogue, the common
orientation needed by the displayed compiler, or therefore close
all non-all-protected internal \(12\) sectors.  Exceptional-heavy,
source-occurrence-failure, total edge-depletion, the complete
\(2+1+1\) source and its global aggregation, other quartic source
families, higher MTA, WUFC, CPM, cutoff removal, reflection
positivity, Osterwalder--Schrader reconstruction, construction of
the continuum Yang--Mills measure, and a uniform positive continuum
spectral gap remain open.  The full Yang--Mills mass gap has not
been proved.
\end{firewall}

\paragraph{Parent-version record.}
V82 was formed from
\texttt{Renewal\_Yang\_Mills\_Mass\_Gap\_Research\_Notes\_V81.tex}.
The V81 parent had \(39\,654\) logical source lines,
\(1\,409\,259\) bytes, and SHA--256
\[
 \texttt{38303aff81ef6d453ee1b344d0389e703da351dcd551e855a9b879f6c4eccf17}.
\]
The next compulsory companion audit is V85.

\clearpage
\part{Third-series V83: a scalar-marginal partial-transpose
sandwich and unconditional closure of the internal heat bank}
\label{part:V83}

\section{Purpose and upstream change of direction}
\label{sec:V83-purpose}

V82 isolated two apparent losses between the complete local heat
operator and the diffuse-bank estimate: the unproved
\(OIPT_G(s)\) card and the possibility that the larger-dimensional
input leg changes from coordinate to coordinate.  Direct expansion
of Clebsch--Gordan coefficients is unnecessary.  The decisive fact
is a general positive-block-matrix inequality.  It bounds a partial
transpose by the corresponding partial trace, and the norms of the
two possible partial transposes are equal.  Consequently one may
transpose one fixed global row while taking the better of the two
local marginal bounds at every coordinate.  This section proves
that statement from first principles and then re-runs the V82
compiler without either conditional hypothesis.

\section{The positive-block partial-transpose sandwich}
\label{sec:V83-block-sandwich}

\begin{theorem}[Scalar-marginal partial-transpose sandwich]
\label{thm:V83-PT-sandwich}
Let \(\mathcal H_A,\mathcal H_B\) be finite-dimensional Hilbert
spaces, choose an orthonormal basis of \(\mathcal H_B\), and let
\(X\ge0\) on \(\mathcal H_A\otimes\mathcal H_B\).  Put
\[
 S_A:=\operatorname{Tr}_{\mathcal H_B}X.
\]
Then, after the evident identification of tensor factors,
\begin{equation}
 -\,S_A\otimes I_B
 \preccurlyeq X^{\Gamma_B}
 \preccurlyeq S_A\otimes I_B.
\label{eq:V83-PT-sandwich}
\end{equation}
In particular,
\begin{equation}
 \|X^{\Gamma_B}\|\le \|S_A\|.
\label{eq:V83-one-marginal}
\end{equation}
\end{theorem}

\begin{proof}
Write \(m=\dim\mathcal H_B\), fix the chosen basis
\((e_i)_{i=1}^m\), and regard \(X\) as the positive operator-valued
block matrix
\[
 X=[X_{ij}]_{i,j=1}^m,\qquad
 X_{ij}:=(I_A\otimes\langle e_i|)X
              (I_A\otimes|e_j\rangle).
\]
Thus
\[
 X^{\Gamma_B}=[X_{ji}]_{i,j=1}^m,
 \qquad
 S_A=\sum_{i=1}^mX_{ii}.
\]
For every \(i<j\), the principal compression
\[
 P_{ij}:=
 \begin{pmatrix}
  X_{ii}&X_{ij}\\ X_{ji}&X_{jj}
 \end{pmatrix}
\]
is positive on \(\mathcal H_A\oplus\mathcal H_A\).  Conjugating
first by the block swap and, when necessary, by
\(\operatorname{diag}(I_A,-I_A)\), shows that both
\[
 Q_{ij}^{+}:=
 \begin{pmatrix}
  X_{jj}&X_{ji}\\ X_{ij}&X_{ii}
 \end{pmatrix},
 \qquad
 Q_{ij}^{-}:=
 \begin{pmatrix}
  X_{jj}&-X_{ji}\\ -X_{ij}&X_{ii}
 \end{pmatrix}
\]
are positive.

Embed these two-by-two blocks into
\(\bigoplus_{i=1}^m\mathcal H_A\).  A block-by-block comparison gives
\begin{align}
 I_B\otimes S_A-X^{\Gamma_B}
   &=\sum_{1\le i<j\le m}Q_{ij}^{-},\label{eq:V83-upper-sum}\\
 I_B\otimes S_A+X^{\Gamma_B}
   &=2\bigoplus_{i=1}^mX_{ii}
     +\sum_{1\le i<j\le m}Q_{ij}^{+}.
\label{eq:V83-lower-sum}
\end{align}
Every summand on the right is positive, proving
\eqref{eq:V83-PT-sandwich}.  Since \(X^{\Gamma_B}\) is self-adjoint,
for every unit vector \(\zeta\),
\[
 |\langle\zeta,X^{\Gamma_B}\zeta\rangle|
 \le
 \langle\zeta,(I_B\otimes S_A)\zeta\rangle
 \le \|S_A\|.
\]
Taking the supremum proves \eqref{eq:V83-one-marginal}.
\end{proof}

\begin{corollary}[Symmetric marginal bound with one fixed
transpose]
\label{cor:V83-symmetric-PT}
For every positive \(X\) as above,
\begin{equation}
 \boxed{\;
 \|X^{\Gamma_B}\|
 \le
 \min\bigl\{
 \|\operatorname{Tr}_{\mathcal B}X\|,
 \|\operatorname{Tr}_{\mathcal A}X\|
 \bigr\}.\;}
\label{eq:V83-symmetric-PT}
\end{equation}
The transpose on the left remains the transpose of the single fixed
factor \(\mathcal H_B\).
\end{corollary}

\begin{proof}
The first marginal bound is
\Cref{thm:V83-PT-sandwich}.  In the chosen product bases,
\[
 X^{\Gamma_A}=(X^{\Gamma_B})^{\mathsf T}.
\]
Full matrix transpose preserves singular values and hence operator
norm.  Applying \Cref{thm:V83-PT-sandwich} with \(A\) and \(B\)
interchanged therefore gives
\[
 \|X^{\Gamma_B}\|
 =\|X^{\Gamma_A}\|
 \le\|\operatorname{Tr}_{\mathcal H_A}X\|.
\]
Taking the better of the two estimates proves the claim.
\end{proof}

\begin{firewall}[What positivity supplies]
\label{fire:V83-positivity}
The order sandwich \eqref{eq:V83-PT-sandwich} uses \(X\ge0\).
It does not assert that partial transpose preserves arbitrary
operator inequalities.  Later domination arguments are therefore
made before transpose, at the level of positive product-state
expectations; the transpose is used only on the chosen dominating
tensor product.
\end{firewall}

\section{The \(OIPT_G(s)\) card is a theorem}
\label{sec:V83-OIPT}

\begin{theorem}[Complete output-isotypic heat card]
\label{thm:V83-OIPT}
Let \(V,W\) be irreducible unitary \(G\)-modules and let
\[
 M_{V,W}(s)
 =
 \sum_{U\in\widehat G}e^{-sC_2(U)}P_U
\]
be the complete local output-isotypic heat operator of V82,
including all multiplicity spaces.  For the partial transpose of
either fixed input leg,
\begin{equation}
 \boxed{\;
 \|M_{V,W}(s)^\Gamma\|
 \le
 \frac{K_G(s)}{\max\{d_V,d_W\}},
 \qquad
 K_G(s):=\sum_{U\in\widehat G}d_U^2e^{-sC_2(U)}<\infty.\;}
\label{eq:V83-OIPT}
\end{equation}
Thus the card postulated in \Cref{def:V82-OIPT} holds with its full
multipity content and without choosing a coordinate-dependent
transpose.
\end{theorem}

\begin{proof}
The operator \(M_{V,W}(s)\) is positive.  It commutes with the
diagonal \(G\)-action, so the two partial traces are scalar by
Schur's lemma.  With
\[
 T(s):=\operatorname{Tr}M_{V,W}(s)
      =\sum_{U\in\widehat G}
        e^{-sC_2(U)}d_Um_U,
\]
the exact identities are
\[
 \operatorname{Tr}_{W}M_{V,W}(s)
   ={T(s)\over d_V}I_V,
 \qquad
 \operatorname{Tr}_{V}M_{V,W}(s)
   ={T(s)\over d_W}I_W.
\]
Hence \Cref{cor:V83-symmetric-PT} gives, for either fixed local
transpose,
\begin{equation}
 \|M_{V,W}(s)^\Gamma\|
 \le {T(s)\over\max\{d_V,d_W\}}.
\label{eq:V83-trace-over-max}
\end{equation}

Assume first that \(d_V\ge d_W\).  The Frobenius-reciprocity
estimate of \Cref{lem:V82-multiplicity} states
\(m_Ud_V\le d_Ud_W\); therefore
\[
 d_Um_U\le d_U^2{d_W\over d_V}\le d_U^2.
\]
Summing with the heat weights gives \(T(s)\le K_G(s)\).
The case \(d_W\ge d_V\) is identical after interchanging the input
modules.  Insert this estimate into
\eqref{eq:V83-trace-over-max}.
\end{proof}

\begin{corollary}[Fixed-row Casimir half-power]
\label{cor:V83-fixed-row-half}
Fix either input row \(a\), let its irreducible module have Casimir
mass \(a_a\), and transpose the other row.  Then
\begin{equation}
 \|M_{V,W}(s)^\Gamma\|
 \le C_{G,s}(1+a_a)^{-1/2}.
\label{eq:V83-fixed-row-half}
\end{equation}
No comparison between the two input dimensions is required.
\end{corollary}

\begin{proof}
The denominator in \eqref{eq:V83-OIPT} is at least the dimension
of the fixed row-\(a\) module.  The Weyl-dimension lower bound of
\Cref{lem:V69-Weyl-dimension}, with the harmless \(1+a_a\)
normalization for the trivial module, finishes the proof.
\end{proof}

\section{Unconditional entanglement-safe internal-bank closure}
\label{sec:V83-bank}

\begin{theorem}[Complete diffuse internal heat bank]
\label{thm:V83-complete-bank}
Let \(D\) be the owned, pointwise-comparable, diffuse internal
\(12\) bank of V81, and let \(M_e(s)\) be the complete physical
two-row heat operator at \(e\).  For either fixed physical row
\(a\in\{1,2\}\), and for \(p\le p_0\),
\begin{equation}
 \kappa_\otimes\!\left(\bigotimes_{e\in D}M_e(s)\right)
 \le {C_{G,s}\over\Xi_a^2}
 \le C_{G,s}\chi^2.
\label{eq:V83-bank-chi2}
\end{equation}
Consequently the non-all-protected part of the consolidated
internal \(12\) sector satisfies \(FCQ_{211}\).  The common
dimension-orientation assumption in
\Cref{thm:V82-OIPT-compiler} is unnecessary.
\end{theorem}

\begin{proof}
Order the row-\(a\) Casimir weights on \(D\), and denote the four
largest by \(a_{a,1},\ldots,a_{a,4}\).  At each of these four
coordinates there are two positive dominants for the physical heat
factor: the factor \(M_e(s)\) itself and the V68 majorant
\[
 B_e(s)=P_e^{\rm prot}+r_s(I-P_e^{\rm prot}).
\]
For the first choice,
\Cref{cor:V83-fixed-row-half} bounds the fixed-row partial-transpose
norm by \(C_{G,s}(1+a_{a,e})^{-1/2}\).  For the second choice, V68
bounds that norm by the uniform constant
\(\rho_s<1\).  Choose at each of the four coordinates whichever
positive dominant has the smaller proved partial-transpose norm.
At every other coordinate choose \(B_e(s)\).

Let \(H_D\) be the tensor product of these chosen dominants.
Positivity and coordinatewise domination give
\[
 0\preccurlyeq\bigotimes_{e\in D}M_e(s)
 \preccurlyeq H_D.
\]
If \(\varrho_a\otimes\varrho_{\bar a}\) is an arbitrary product
density across the two full rows, then
\[
 \operatorname{Tr}\!\left[
  (\varrho_a\otimes\varrho_{\bar a})
  \bigotimes_{e\in D}M_e(s)\right]
 \le
 \operatorname{Tr}[
  (\varrho_a\otimes\varrho_{\bar a})H_D]
 \le \|H_D^{\Gamma_{\bar a}}\|.
\]
The last transpose is one fixed full-row transpose.  It factorizes
coordinatewise, so its norm is at most
\[
 \left[
 \prod_{j=1}^4
  \min\{\rho_s,C_{G,s}(1+a_{a,j})^{-1/2}\}
 \right]\rho_s^{|D|-4}.
\]
The four-largest-coordinate estimate
\Cref{lem:V81-four-largest}, with the harmless replacement of
\(a\) by \(1+a\), bounds this expression by
\(C_{G,s}\Xi_a^{-2}\).  Ownership and the heavy-row comparison used
in \Cref{thm:V81-protected-two-powers} then give the second
inequality in \eqref{eq:V83-bank-chi2}.

The non-all-protected positive output sector is dominated by the
complete heat product.  Substitution of
\eqref{eq:V83-bank-chi2} into the consolidated-square argument of
\Cref{thm:V81-protected-internal-FCQ} proves \(FCQ_{211}\).
\end{proof}

\begin{corollary}[Resolution of the two V82 cards]
\label{cor:V83-V82-resolved}
The \(OIPT_G(s)\) card of \Cref{def:V82-OIPT} is proved, and the
orientation obstruction of \Cref{fire:V82-orientation} is removed.
Every source-supported owned diffuse comparable internal \(12\)
heat bank, protected or nonprotected, now has the two heavy
inverse-mass powers required by the V81 consolidated-square
compiler.
\end{corollary}

\begin{proof}
\Cref{thm:V83-OIPT} proves the card.  The fixed-row argument in
\Cref{thm:V83-complete-bank} removes the orientation hypothesis.
The protected sector was already closed by
\Cref{thm:V81-protected-internal-FCQ}; the complementary positive
sector is closed by \Cref{thm:V83-complete-bank}.
\end{proof}

\begin{firewall}[Scope of the new closure]
\label{fire:V83-scope}
The conclusion is an entanglement-safe capacity estimate for the
complete local heat product on the already constructed owned
diffuse comparable internal \(12\) bank.  It does not manufacture
that bank when source occurrence fails, does not treat the
exceptional-heavy alternative of the V80 trichotomy, and does not
resolve total edge depletion.  It also does not by itself aggregate
all quartic source families or construct the continuum measure.
\end{firewall}

\section{V83 result ledger and V84 acquisition order}
\label{sec:V83-status}

\begin{theorem}[Third-series V83 result ledger]
\label{thm:V83-results}
Relative to V1--V82, V83 proves:
\begin{enumerate}[label=\textup{(V83.\arabic*)},leftmargin=6.3em]
\item the two-sided partial-transpose order sandwich
      \eqref{eq:V83-PT-sandwich} for every positive bipartite
      operator;
\item the symmetric minimum-of-marginals norm bound while keeping
      one fixed transposed side;
\item the full multiplicity-uniform \(OIPT_G(s)\) heat card;
\item a Casimir half-power from either preassigned row, independent
      of local dimension ordering;
\item entanglement-safe tensorization of the complete diffuse
      internal heat bank;
\item removal of the common-orientation hypothesis in the V82
      compiler; and
\item closure of both protected and nonprotected source-supported
      owned diffuse comparable internal \(12\) sectors.
\end{enumerate}
\end{theorem}

\begin{proof}
These are
\Cref{thm:V83-PT-sandwich,
cor:V83-symmetric-PT,
thm:V83-OIPT,
cor:V83-fixed-row-half,
thm:V83-complete-bank,
cor:V83-V82-resolved}.
\end{proof}

\begin{programme}[V84: eliminate the source-occurrence-failure
alternative]
\label{prog:V83-V84}
\begin{enumerate}[leftmargin=3em]
\item Return to the maximal shared-bank trichotomy of V80 and write
      source-occurrence failure as an explicit incidence statement,
      not as a verbal complement.
\item Determine whether a maximal internal \(12\) bank which avoids
      every source-occurring cofinal row forces the quartic source
      into the \(2+1+1\) or \(1+1+1+1\) thread shapes.
\item Consolidate repeated visits to one physical cell before
      assigning any further payer, as required by the V80
      SM/NS audit.
\item Use the V83 complete-bank theorem immediately whenever the
      extracted bank is source-supported, rather than reopening
      output-sector analysis.
\item If incidence alone does not force a new thread shape, isolate
      the weakest quantitative hitting-number lemma that would do
      so and test it on a finite hypergraph model.
\item The next compulsory SM/NS companion audit remains V85.
\end{enumerate}
\end{programme}

\begin{firewall}[Claim boundary after third-series V83]
\label{fire:V83-claim-boundary}
V83 rigorously proves the previously missing partial-transpose heat
card and closes every source-supported owned diffuse comparable
internal \(12\) heat bank, including arbitrary output
multiplicities and row-internal entanglement.  Source-occurrence
failure, the exceptional-heavy alternative, total edge depletion,
the complete \(2+1+1\) source and its global aggregation, other
quartic source families, higher MTA, WUFC, CPM, cutoff removal,
reflection positivity, Osterwalder--Schrader reconstruction,
construction of the continuum Yang--Mills measure, and a uniform
positive continuum spectral gap remain open.  The full Yang--Mills
mass gap has not been proved.
\end{firewall}

\paragraph{Parent-version record.}
V83 was formed from
\texttt{Renewal\_Yang\_Mills\_Mass\_Gap\_Research\_Notes\_V82.tex}.
The V82 parent had \(40\,080\) logical source lines,
\(1\,423\,323\) bytes, and SHA--256
\[
 \texttt{57f4ab279b64fbff4b735d6eedf2d3878e45e43af0e9729fe62673e4457d2fe9}.
\]
The next compulsory companion audit is V85.

\clearpage
\part{Third-series V84: exact incidence for arbitrary unmarked
terminal banks and elimination of source-occurrence failure}
\label{part:V84}

\section{The apparent failure is an order-of-expansion issue}
\label{sec:V84-purpose}

Alternative \textup{(v)} of
\Cref{thm:V80-maximal-reduction} was retained because the V80
terminal bank was selected after a scalar cofinal-row trichotomy,
whereas V66 had stated its incidence theorem for the earlier
V64 chain banks.  The distinction is not physical.  The terminal
bank is still a literal, unmarked subset of the coordinates, chosen
from Casimir weights before any signed endpoint expansion.  It can
therefore be included in the same Boolean refinement as the earlier
banks.  The only extra point is that a bank crossing the joining
coordinate must be split into its prefix and suffix pieces without
losing cofinality or two-row ownership.  The next lemmas make this
argument explicit.

\section{Universal unmarked-bank incidence}
\label{sec:V84-incidence}

\begin{definition}[Admissible unmarked bank in a fixed word]
\label{def:V84-admissible-bank}
Fix a joining coordinate \(r\), a payer placement \(\pi\), and its
finite marked set \(Q_{r,\pi}\) from
\eqref{eq:V80-qstar}.  An \emph{admissible unmarked bank} is a
coordinate set \(D\) such that:
\begin{enumerate}[label=\textup{(\roman*)},leftmargin=4em]
\item \(D\cap Q_{r,\pi}=\varnothing\);
\item membership in \(D\) is determined by the representation and
      Casimir-weight data, not by a choice made after expanding a
      signed covariance endpoint; and
\item the local resolution requested on \(D\) is the complete
      protected/gapped output resolution of the physical moment
      module already present in the word.
\end{enumerate}
\end{definition}

\begin{lemma}[V80 terminal banks are admissible]
\label{lem:V84-terminal-admissible}
Every gapped or owned comparable bank produced from
\(E=E_0\setminus Q_{r,\pi}\) by
\Cref{thm:V80-cofinal-trichotomy} is admissible in the sense of
\Cref{def:V84-admissible-bank}.
\end{lemma}

\begin{proof}
The construction in \Cref{thm:V80-cofinal-trichotomy} partitions
\(E\) by coordinatewise inequalities among the numerical Casimir
weights \(a_{i,e}\), followed by numerical mass tests.  Hence the
output banks are subsets of \(E\), are disjoint from \(Q_{r,\pi}\),
and are fixed before any endpoint sign is chosen.  The trichotomy
asks precisely for the V32 complete physical output resolution.
These are the three clauses of
\Cref{def:V84-admissible-bank}.
\end{proof}

\begin{theorem}[Exact incidence for any finite admissible family]
\label{thm:V84-universal-incidence}
Let \(D_1,\ldots,D_k\) be a fixed finite family of admissible
unmarked banks in the \(2+1+1\) word
\eqref{eq:V66-uncompressed-word}.  There is one exact signed
expansion, with a number of terms depending only on \(k\) and the
fixed source arity, such that in every term:
\begin{enumerate}[label=\textup{(\alph*)},leftmargin=4em]
\item each prefix piece \(D_j^-:=D_j\cap\{e<r\}\) is a complete
      source occurrence in an endpoint moment;
\item each suffix piece \(D_j^+:=D_j\cap\{e>r\}\) is a complete
      source occurrence in the four-row suffix moment;
\item all requested positive spectral factors on all nonempty
      pieces occur jointly and commute; and
\item the joining coordinate, payer, and marked separator packet
      are unchanged and occur exactly once.
\end{enumerate}
The statement is independent of how the banks were found.
\end{theorem}

\begin{proof}
Form the finite Boolean algebra generated by
\(D_1,\ldots,D_k\) separately on the two sides of \(r\).  On the
prefix side, apply the coefficient-one simultaneous covariance
restriction of \Cref{lem:V65-simultaneous-restriction}, and only
then expand each localized covariance into its two complete
endpoint moments as in \Cref{def:V66-prefix-carrier}.  By
\Cref{lem:V66-prefix-endpoints}, every resulting Boolean atom is a
complete physical moment-partition state.  Insert the requested
positive central output resolution there.

On the suffix side no endpoint expansion is needed:
\Cref{lem:V66-suffix-endpoints} says that every Boolean atom is
already a complete four-row moment.  Insert the corresponding
central output resolution.  The banks avoid \(Q_{r,\pi}\), so no
operation touches the joining, payer, or marked separator slots.
Centrality and commutation are exactly
\Cref{lem:V59-central-isotypic,lem:V59-bank-commutation}.

All restrictions and endpoint expansions are performed
simultaneously, before any norm or capacity is taken.  Thus this is
one common exact expansion, rather than a multiplication of
estimates from different signed terms.  A fixed finite family has
at most \(2^k\) Boolean atoms and a source-arity-bounded number of
covariance endpoints, proving the uniform term-count assertion.
\end{proof}

\begin{firewall}[No reconstructed or duplicate occurrence]
\label{fire:V84-no-duplicate}
\Cref{thm:V84-universal-incidence} does not infer a source factor
from a scalar mass ledger.  It exhibits the factor inside the
original prefix or suffix slot of
\eqref{eq:V66-uncompressed-word}.  If the same physical cell is
visited by several named banks, it is one Boolean atom carrying
commuting central resolutions; it is not assigned several
independent payers.
\end{firewall}

\section{Cofinality and ownership survive the source-slot split}
\label{sec:V84-splitting}

\begin{lemma}[Cofinal gapped bank has a source-occurring half]
\label{lem:V84-gapped-half}
If an admissible gapped bank \(F\) satisfies
\[
 S_a(F)\ge\sigma\Xi_a,
\]
then one of \(F^-\) or \(F^+\) is source-occurring and satisfies
\[
 S_a(F^\pm)\ge{\sigma\over2}\Xi_a.
\]
It remains pointwise gapped at row \(a\).
\end{lemma}

\begin{proof}
Because \(r\in Q_{r,\pi}\) and \(F\cap Q_{r,\pi}=\varnothing\),
\(F=F^-\dot\cup F^+\).  Additivity gives the half-mass conclusion,
and a coordinatewise gap is inherited by a subset.
\Cref{thm:V84-universal-incidence} supplies the source occurrence.
\end{proof}

\begin{lemma}[An owned comparable bank has a source-occurring owned
half]
\label{lem:V84-owned-half}
Suppose \(D\) is admissible, pointwise two-sided comparable between
rows \(a,c\),
\begin{equation}
 \lambda a_{a,e}\le a_{c,e}\le\Lambda a_{a,e}
 \quad(e\in D),
\label{eq:V84-comparable}
\end{equation}
and owns both rows:
\[
 S_a(D)\ge\sigma_a\Xi_a,
 \qquad
 S_c(D)\ge\sigma_c\Xi_c.
\]
Then one of \(D^-\) or \(D^+\), denoted \(\widehat D\), is
source-occurring, remains pointwise comparable, and owns fixed
fractions of both \(\Xi_a\) and \(\Xi_c\):
\begin{equation}
 S_a(\widehat D)\ge{\sigma_a\over2}\Xi_a,
 \qquad
 S_c(\widehat D)\ge
 {\lambda\sigma_a\sigma_c\over2\Lambda}\Xi_c.
\label{eq:V84-owned-half}
\end{equation}
\end{lemma}

\begin{proof}
Choose \(\widehat D\in\{D^-,D^+\}\) with
\[
 S_a(\widehat D)\ge\tfrac12S_a(D)
 \ge{\sigma_a\over2}\Xi_a.
\]
The lower pointwise comparison gives
\[
 S_c(\widehat D)\ge\lambda S_a(\widehat D)
 \ge{\lambda\sigma_a\over2}\Xi_a.
\]
On the other hand, ownership of row \(c\) and the upper comparison
on the full bank imply
\[
 \sigma_c\Xi_c\le S_c(D)\le\Lambda S_a(D)
 \le\Lambda\Xi_a,
\]
and hence \(\Xi_a\ge(\sigma_c/\Lambda)\Xi_c\).  Substitution proves
the second inequality in \eqref{eq:V84-owned-half}.  Comparability
is inherited by subsets, and source occurrence follows from
\Cref{thm:V84-universal-incidence}.
\end{proof}

\begin{assessment}[Why the split does not spend a second currency]
\label{ass:V84-split-currency}
The half-bank is used because it is a literal subfactor of the full
terminal bank in one complete source slot.  Its mass estimates are
monotone consequences of the original ownership ledger.  No prefix
response, payer, or second copy of the bank is introduced.
\end{assessment}

\section{Elimination of the V80 occurrence-failure branch}
\label{sec:V84-elimination}

\begin{theorem}[No source-occurrence failure for an unmarked V80
terminal bank]
\label{thm:V84-no-occurrence-failure}
In a fixed \(2+1+1\) first-connectivity word, alternative
\textup{(v)} of \Cref{thm:V80-maximal-reduction} is empty.  More
precisely:
\begin{enumerate}[label=\textup{(\roman*)},leftmargin=4em]
\item a cofinal gapped V80 terminal bank contains a source-occurring
      cofinal gapped half-bank;
\item an owned comparable V80 terminal bank contains a
      source-occurring half-bank owning fixed fractions of both
      endpoint rows; and
\item all earlier banks needed in the same proof occur jointly with
      that half-bank in one fixed finite exact expansion.
\end{enumerate}
\end{theorem}

\begin{proof}
By \Cref{lem:V84-terminal-admissible}, the terminal bank and the
fixed finite list of earlier banks form an admissible family.
Apply \Cref{thm:V84-universal-incidence} to the whole family.
For a gapped terminal outcome, use
\Cref{lem:V84-gapped-half}; for an owned comparable outcome, use
\Cref{lem:V84-owned-half}.  These lemmas retain exactly the
cofinality or ownership hypotheses required by the downstream
cards, with only fixed losses in their constants.  Thus every
signed term produced by the common expansion contains the needed
positive central factor in an actual source slot.  There is no
remaining term satisfying the definition of alternative
\textup{(v)}.
\end{proof}

\begin{theorem}[Heavy-composite counterpackets reduce to the marked
exceptional branch]
\label{thm:V84-heavy-to-exceptional}
Assume \(z\ge\eta>0\), \(R\ge1\), and the heavy-composite
counterpacket hypotheses of \Cref{thm:V80-maximal-reduction}.  After
V83--V84, a persistent counterpacket can occur only if a selected,
paid, or marked separator coordinate \(e\in Q_{r,\pi}\) satisfies
\begin{equation}
 a_{a,e}\ge c_*\Xi_a
\label{eq:V84-only-exceptional}
\end{equation}
for a fixed heavy row \(a\in\{1,2\}\) and fixed \(c_*>0\).
\end{theorem}

\begin{proof}
The five alternatives of
\Cref{thm:V80-maximal-reduction} are exhaustive.  Its cofinal gapped
alternative closes by \Cref{thm:V80-gap-HCRC}; its singleton-partner
owned alternative closes by
\Cref{thm:V76-dominant-bank-FCQ}; and its internal \(12\) owned
alternative closes by \Cref{cor:V83-V82-resolved}.  The
source-occurrence-failure alternative is empty by
\Cref{thm:V84-no-occurrence-failure}.  Only the exceptional
alternative remains.  Equation \eqref{eq:V80-exceptional-heavy}
gives \eqref{eq:V84-only-exceptional} with
\(c_*=\sigma/(2q_*)\).
\end{proof}

\begin{firewall}[The exceptional coordinate still carries its
literal role]
\label{fire:V84-exceptional-role}
\Cref{thm:V84-heavy-to-exceptional} is a reduction, not an estimate
of the surviving word.  The heavy coordinate lies in
\(Q_{r,\pi}\), hence may be the joining coordinate, the payer, or a
marked separator.  It must be treated by the hot-output theorem
appropriate to that actual role.  Moving it into the unmarked bank
would undo the incidence proof.
\end{firewall}

\section{V84 result ledger and V85 acquisition order}
\label{sec:V84-status}

\begin{theorem}[Third-series V84 result ledger]
\label{thm:V84-results}
Relative to V1--V83, V84 proves:
\begin{enumerate}[label=\textup{(V84.\arabic*)},leftmargin=6.3em]
\item an intrinsic admissibility criterion for unmarked banks in a
      fixed first-connectivity word;
\item admissibility of every V80 terminal bank;
\item a common exact source-slot expansion for any fixed finite
      admissible family;
\item preservation of a cofinal gap on one source-occurring side of
      the joining coordinate;
\item preservation of two-row ownership and comparability on one
      source-occurring side;
\item elimination of the V80 source-occurrence-failure alternative;
      and
\item reduction of every heavy-composite counterpacket to a
      mass-cofinal coordinate in the finite marked packet.
\end{enumerate}
\end{theorem}

\begin{proof}
These are
\Cref{def:V84-admissible-bank,
lem:V84-terminal-admissible,
thm:V84-universal-incidence,
lem:V84-gapped-half,
lem:V84-owned-half,
thm:V84-no-occurrence-failure,
thm:V84-heavy-to-exceptional}.
\end{proof}

\begin{programme}[V85: compulsory audit and literal-role
exceptional analysis]
\label{prog:V84-V85}
\begin{enumerate}[leftmargin=3em]
\item Perform the compulsory read-only audit of the newest active SM
      and NS notes before selecting the proof direction.
\item Partition \(Q_{r,\pi}\) into joining, payer, and separator
      roles, consolidating coincidences before estimating anything.
\item For a heavy joining coordinate, apply the complete assembled
      \(J_r(12|3|4)\) hot-output estimate, retaining every mixed
      Leibniz block.
\item For a distinct heavy payer, apply its actual paid local
      moment theorem exactly once.
\item For a heavy separator coordinate, use its selected
      hot-output theorem without reclassifying it as a bank.
\item Compare the resulting inverse heavy-mass power with the
      \(\chi\)-response required by \(HCRC_{12|34}\); if one power is
      missing, isolate whether the retained prefix square supplies
      it without duplicate charging.
\end{enumerate}
\end{programme}

\begin{firewall}[Claim boundary after third-series V84]
\label{fire:V84-claim-boundary}
V84 eliminates source-occurrence failure for all unmarked terminal
banks in the fixed \(2+1+1\) first-connectivity word and reduces the
entire \(z\)-nondepleted heavy-composite counterpacket to its finite
marked exceptional branch.  It does not yet close that exceptional
branch, total edge depletion, the complete \(2+1+1\) source and its
global aggregation, other quartic source families, higher MTA,
WUFC, CPM, cutoff removal, reflection positivity,
Osterwalder--Schrader reconstruction, construction of the continuum
Yang--Mills measure, or a uniform positive continuum spectral gap.
The full Yang--Mills mass gap has not been proved.
\end{firewall}

\paragraph{Parent-version record.}
V84 was formed from
\texttt{Renewal\_Yang\_Mills\_Mass\_Gap\_Research\_Notes\_V83.tex}.
The V83 parent had \(40\,501\) logical source lines,
\(1\,437\,531\) bytes, and SHA--256
\[
 \texttt{23295eebe27bdb288b67e94f82d92a30beacd6102b92da0d5843596289dfa0c1}.
\]
The next compulsory companion audit is V85.

\clearpage
\part{Third-series V85: companion audit and unbalanced closure of
the finite marked exceptional packet}
\label{part:V85}

\section{Compulsory five-version companion audit}
\label{sec:V85-audit}

The required read-only audit was performed before the V85 proof was
chosen.  The newest active companion files were:
\begin{enumerate}[leftmargin=3em]
\item
\texttt{Renewal\_SM\_Einstein\_Continuum\_Research\_Notes\_V64.tex},
with \(23\,157\) logical source lines and \(852\,521\) bytes,
SHA--256
\[
 \texttt{5127cf0774eb4c7d2b1fd0a273a2f9f690ab84cb1360d560dacf503a035e5642}.
\]
Relative to the V80 audit of SM V62, SM V63 proves that a guarded
conditional integration turns one tail carrier into a small
quasi-local Lipschitz load, while its complementary first-owner
branch remains only in a skeleton \(L^p\) class.  SM V64 treats
simultaneous owner confluence by forming the complete distinct
tail-cell set first and applying one group guard.  A weighted
boundary-energy alternative then preserves one rare activity for
each distinct physical cell without multiplying payments for a
shared boundary fluctuation.  Its remaining gate is an all-depth
stopping forest for the bad group outputs.
\item
\texttt{Renewal\_NS\_Stability\_Research\_Notes\_V31.tex},
with \(23\,052\) logical source lines and \(887\,537\) bytes,
SHA--256
\[
 \texttt{f1121b62238ad5491bb7cf03635ea9934942edf573ed8faaa9ee79e1d38a6696}.
\]
This file is byte-for-byte unchanged from the V75 and V80 audits.
Its exact formation ladder still permits one payment for one actual
removed flat mark and forbids charging a persistent occurrence
again at a later stage.
\end{enumerate}

\begin{assessment}[Type-correct V85 transfer]
\label{ass:V85-transfer}
Neither the SM group guard nor the NS formation estimate is a
Yang--Mills heat-kernel bound.  The audit transfers only proof order:
first consolidate all role names belonging to one physical marked
coordinate; next classify the complete local YM tuple; and only then
apply one role-appropriate estimate.  A coordinate which is
simultaneously joining and paying is not two sources of decay.
\end{assessment}

\begin{firewall}[No companion analytic import]
\label{fire:V85-no-import}
Every estimate below uses only the YM overlap identities, the
counterpacket limits, compact-group Casimir comparison, and the
V34--V38 unbalanced compilers.  No SM tail probability, Witten
kernel, group guard, or NS formation inequality is used.
\end{firewall}

\section{Consolidating the finite marked packet}
\label{sec:V85-role-consolidation}

For a fixed selected word, let \(\mathcal R_{r,\pi}\) be its finite
set of marked \emph{role names}: joining, payer when present, binary
prefix/tree marks, and marked separators.  Let
\[
 c:\mathcal R_{r,\pi}\longrightarrow Q_{r,\pi}
\]
send a role to its physical coordinate.  Several roles may have the
same image.

\begin{lemma}[Physical-coordinate owner partition]
\label{lem:V85-role-owner}
Fix a deterministic order on \(\mathcal R_{r,\pi}\).  For every
physical coordinate \(d\in Q_{r,\pi}\), assign all roles in
\(c^{-1}(d)\) to the first role in that fibre.  Then:
\begin{enumerate}[label=\textup{(\roman*)},leftmargin=4em]
\item the owner classes are disjoint and cover every marked role;
\item every physical coordinate occurs in exactly one owner class;
\item the original complete word and its unique payer are unchanged;
      and
\item if \(S_a(Q_{r,\pi})\ge\sigma\Xi_a\), the owner of some
      coordinate \(d\) satisfies
\begin{equation}
 a_{a,d}\ge{\sigma\over q_*}\Xi_a.
\label{eq:V85-owner-heavy}
\end{equation}
\end{enumerate}
\end{lemma}

\begin{proof}
The first three statements are properties of fibres of the map
\(c\); ownership changes only the bookkeeping label, not the
operator word.  For the last statement, \(Q_{r,\pi}\) has at most
\(q_*\) distinct coordinates by \eqref{eq:V80-qstar}.  One summand
of \(S_a(Q_{r,\pi})\) is at least its average.
\end{proof}

\begin{firewall}[Coincident roles receive one estimate]
\label{fire:V85-one-coordinate-one-estimate}
After \Cref{lem:V85-role-owner}, the complete local factor at \(d\)
is assembled with all of its joining, paying, or separating
operations in their original order.  It is never replaced by the
product of separate estimates attached to the names in
\(c^{-1}(d)\).
\end{firewall}

\section{Counterpacket asymmetry forces local unbalance}
\label{sec:V85-unbalance}

\begin{lemma}[Depletion of the opposite composite input at a
cofinal marked cell]
\label{lem:V85-opposite-depletion}
Let \(a\in\{1,2\}\), let \(b\) be the other composite row, assume
\(\Xi_a=\max\{\Xi_1,\Xi_2\}\), and suppose
\begin{equation}
 a_{a,d}\ge c_*\Xi_a
\label{eq:V85-cofinal-cell}
\end{equation}
at one coordinate.  Then
\begin{equation}
 {a_{b,d}\over a_{a,d}}
 \le {p\over c_*^2},
 \qquad p=\theta_{12}.
\label{eq:V85-opposite-ratio}
\end{equation}
\end{lemma}

\begin{proof}
The overlap stock contains the \(d\)-summand, so
\[
 p={H_{12}\over\Xi_1\Xi_2}
 \ge
 {a_{a,d}a_{b,d}\over\Xi_a\Xi_b}
 \ge c_*{a_{b,d}\over\Xi_b}.
\]
Because \(\Xi_b\le\Xi_a\) and
\(a_{a,d}\ge c_*\Xi_a\),
\[
 {a_{b,d}\over a_{a,d}}
 \le {a_{b,d}\over c_*\Xi_a}
 \le {a_{b,d}\over c_*\Xi_b}
 \le {p\over c_*^2}.
\]
\end{proof}

\begin{lemma}[Singleton inputs are locally negligible]
\label{lem:V85-singleton-depletion}
Under \eqref{eq:V85-cofinal-cell}, for \(j\in\{3,4\}\),
\begin{equation}
 {a_{j,d}\over a_{a,d}}
 \le{\chi\over c_*},
 \qquad
 \chi=
 {\max\{\Xi_3,\Xi_4\}\over
  \min\{\Xi_1,\Xi_2\}}.
\label{eq:V85-singleton-ratio}
\end{equation}
\end{lemma}

\begin{proof}
Since \(a_{j,d}\le\Xi_j\) and
\(\min\{\Xi_1,\Xi_2\}\le\Xi_a\),
\[
 {a_{j,d}\over a_{a,d}}
 \le
 {\max\{\Xi_3,\Xi_4\}\over c_*\Xi_a}
 \le
 {\max\{\Xi_3,\Xi_4\}\over
  c_*\min\{\Xi_1,\Xi_2\}}
 ={\chi\over c_*}.
\]
\end{proof}

\begin{theorem}[Every exceptional counterpacket cell is eventually
V32-unbalanced]
\label{thm:V85-exceptional-unbalanced}
Let a heavy-composite counterpacket satisfy
\[
 p_n\longrightarrow0,\qquad
 \chi_n\longrightarrow0,
\]
as in \Cref{thm:V77-counterpacket} and
\Cref{eq:V78-R-chi}.  Suppose that, after passage to a subsequence,
a marked physical coordinate \(d_n\) and a fixed row
\(a\in\{1,2\}\) satisfy
\[
 a_{a,d_n}\ge c_*\Xi_a^{(n)}
\]
with fixed \(c_*>0\).  Then for all sufficiently large \(n\),
\begin{equation}
 \sum_{i\ne a}\|\lambda_{i,d_n}\|
 \le\frac12\|\lambda_{a,d_n}\|.
\label{eq:V85-local-unbalance}
\end{equation}
Thus the complete local representation tuple at \(d_n\) is
unbalanced at row \(a\) in the exact sense of
\eqref{eq:V32-unbalanced}.
\end{theorem}

\begin{proof}
The nonzero selected joining word has the fixed active Casimir
floor used in \Cref{lem:V77-matching-floor}.  Since
\(R_n=\chi_n^{-1}\to\infty\), both composite row masses diverge;
hence \(\Xi_a^{(n)}\to\infty\), and the cofinality hypothesis gives
\(a_{a,d_n}\to\infty\).

\Cref{lem:V85-opposite-depletion,
lem:V85-singleton-depletion} give
\begin{equation}
 \max_{i\ne a}
 {a_{i,d_n}\over a_{a,d_n}}
 \le
 \max\left\{
 {p_n\over c_*^2},{\chi_n\over c_*}
 \right\}
 \longrightarrow0.
\label{eq:V85-all-mass-ratios}
\end{equation}
The compact-group Casimir comparison
\eqref{eq:V32-Casimir-comparison} says that
\(a_{\lambda}=1+C_2(\lambda)\) is uniformly comparable to
\(1+\|\lambda\|^2\).  Because the denominator in
\eqref{eq:V85-all-mass-ratios} diverges, it follows for each
\(i\ne a\) that
\[
 {\|\lambda_{i,d_n}\|\over
  \|\lambda_{a,d_n}\|}
 \longrightarrow0.
\]
There are at most three other active rows.  Their sum is therefore
at most half of the dominant norm for all sufficiently large \(n\),
which is \eqref{eq:V85-local-unbalance}.
\end{proof}

\begin{remark}[Why V47's balanced obstruction does not survive this
counterpacket]
\label{rem:V85-V47-balance}
V47 correctly exhibited balanced low-output marked cells in the
unrestricted source.  V85 adds the simultaneous asymptotic
information \(p\to0\), \(\chi\to0\), and one cofinal marked
composite input.  These three facts force the marked tuple into the
unbalanced, not balanced, side of the exhaustive V32 split.  No
claim of generic input/output Casimir comparison is made.
\end{remark}

\section{Literal-role closure of the exceptional branch}
\label{sec:V85-exceptional-closure}

\begin{theorem}[The finite marked exceptional branch closes]
\label{thm:V85-exceptional-closes}
No \(z\)-nondepleted heavy-composite \(FCQ_{211}\) counterpacket can
satisfy the exceptional alternative
\eqref{eq:V84-only-exceptional}.  The conclusion holds for all
coincidences among joining, payer, prefix-mark, and separator roles.
\end{theorem}

\begin{proof}
Consolidate coincident roles by
\Cref{lem:V85-role-owner}; this leaves one complete local physical
factor and the unique payer at the exceptional coordinate.  By
\Cref{thm:V85-exceptional-unbalanced}, its representation tuple is
eventually V32-unbalanced at the heavy composite row.

The unbalanced output theorem
\Cref{lem:V32-unbalanced-output} retains a fixed fraction of the
dominant input Casimir in every resolved local output.  The
all-fixed-order moment and payer bounds of
\Cref{thm:V34-unbalanced-bank} therefore apply to the complete
one-coordinate bank \(\{d\}\), with the output resolution made only
after the complete local product is formed.  For a one- or
two-demand marked-tree role, this is exactly the immediate
unbalanced compiler in
\Cref{prop:V35-immediate-compilers}.  For the selected
fourth-cumulant role it is also the unbalanced selected branch of
\Cref{thm:V38-GSEL4}; for the balanced-marked interface notation,
the same conclusion is recorded in
\Cref{thm:V47-unbalanced-closes}.

These compilers explicitly include a payer at the same coordinate
and all other payer placements.  Hence role coincidence does not
require, and by \Cref{fire:V85-one-coordinate-one-estimate} does not
permit, a second payment.  The number of role-owner classes and
local moment-cumulant terms is bounded by the fixed source arity, so
finite recombination changes only the uniform constant.  The
exceptional words satisfy their literal marked-tree, hence
\(FCQ_{211}\), bound and cannot form a counterpacket.
\end{proof}

\begin{corollary}[The \(z\)-nondepleted heavy-composite sector is
closed]
\label{cor:V85-heavy-composite-closed}
Every selected \(2+1+1\) word with \(z\ge\eta>0\) in the
heavy-composite regime satisfies the uniform \(FCQ_{211}\) estimate.
Equivalently, the heavy-composite alternative of
\Cref{thm:V77-counterpacket} is empty.
\end{corollary}

\begin{proof}
\Cref{thm:V84-heavy-to-exceptional} reduces a persistent
heavy-composite counterpacket to the exceptional alternative.
\Cref{thm:V85-exceptional-closes} eliminates that alternative.
\end{proof}

\begin{corollary}[Only total overlap depletion remains for the
selected \(2+1+1\) source]
\label{cor:V85-only-total-depletion}
If a sequence of complete selected-coordinate \(2+1+1\) words
violates every uniform \(FCQ_{211}\) bound, then, after passage to a
subsequence,
\begin{equation}
 p_n\to0,\qquad x_n\to0,\qquad y_n\to0,\qquad z_n\to0.
\label{eq:V85-total-depletion}
\end{equation}
All pairwise normalized overlaps are therefore depleted.
\end{corollary}

\begin{proof}
\Cref{thm:V77-counterpacket} gives
\(p_n,x_n,y_n\to0\) and says that either \(z_n\to0\) or a
\(z\)-nondepleted heavy-composite subsequence remains.  The latter
is impossible by \Cref{cor:V85-heavy-composite-closed}.
\end{proof}

\begin{firewall}[Source closure has not yet been inferred]
\label{fire:V85-not-source-closure}
\Cref{cor:V85-only-total-depletion} is a sharp counterpacket
reduction, not yet a contradiction.  V76 proves that total depletion
forces divergent row masses, but a family of four large rows can
still have all normalized pair overlaps small.  A new complete-word
estimate or combinatorial concentration theorem is required there.
\end{firewall}

\section{V85 result ledger and V86 acquisition order}
\label{sec:V85-status}

\begin{theorem}[Third-series V85 result ledger]
\label{thm:V85-results}
Relative to V1--V84, V85 proves:
\begin{enumerate}[label=\textup{(V85.\arabic*)},leftmargin=6.3em]
\item the compulsory SM V64 and NS V31 audit with exact
      fingerprints;
\item a type-correct one-coordinate/one-estimate role-owner
      partition;
\item depletion of the opposite composite input at every
      mass-cofinal marked heavy-row cell;
\item depletion of both singleton inputs at that cell in the
      heavy-composite limit;
\item eventual V32 unbalance of every exceptional marked tuple;
\item literal-role closure of the complete finite exceptional
      packet, including all payer coincidences;
\item closure of the entire \(z\)-nondepleted heavy-composite sector;
      and
\item reduction of every remaining selected \(2+1+1\)
      counterpacket to total pair-overlap depletion.
\end{enumerate}
\end{theorem}

\begin{proof}
These are
\Cref{sec:V85-audit,
ass:V85-transfer,
lem:V85-role-owner,
lem:V85-opposite-depletion,
lem:V85-singleton-depletion,
thm:V85-exceptional-unbalanced,
thm:V85-exceptional-closes,
cor:V85-heavy-composite-closed,
cor:V85-only-total-depletion}.
\end{proof}

\begin{programme}[V86: total-depletion geometry]
\label{prog:V85-V86}
\begin{enumerate}[leftmargin=3em]
\item Work on the exact counterpacket
      \eqref{eq:V85-total-depletion}; do not reopen the
      \(z\)-nondepleted branches.
\item Use \Cref{cor:V76-total-depletion-masses} to normalize each
      row into a probability distribution on physical coordinates
      and translate \(\theta_{ij}\to0\) into quantitative pairwise
      avoidance.
\item Partition coordinates by their active-row incidence set.
      Determine whether pairwise avoidance forces each row mass into
      physical-private cells, which V79 closes, or permits a cyclic
      family of small shared cells.
\item If cyclic sharing persists, extract its finite incidence
      hypergraph and identify the first complete source slot in which
      two of its row distributions must meet.
\item Keep marked coordinates separate.  A cofinal marked cell is
      already closed by the V85 unbalance argument.
\item Seek a four-row concentration inequality at the level of the
      exact source word, not a false lower bound on one pair overlap.
\end{enumerate}
\end{programme}

\begin{firewall}[Claim boundary after third-series V85]
\label{fire:V85-claim-boundary}
V85 closes the complete marked exceptional branch and hence every
\(z\)-nondepleted heavy-composite counterpacket.  It leaves total
pair-overlap depletion for the selected \(2+1+1\) source, global
aggregation of that source, other quartic source families, higher
MTA, WUFC, CPM, cutoff removal, reflection positivity,
Osterwalder--Schrader reconstruction, construction of the continuum
Yang--Mills measure, and a uniform positive continuum spectral gap
open.  The full Yang--Mills mass gap has not been proved.
\end{firewall}

\paragraph{Parent-version record.}
V85 was formed from
\texttt{Renewal\_Yang\_Mills\_Mass\_Gap\_Research\_Notes\_V84.tex}.
The V84 parent had \(40\,855\) logical source lines,
\(1\,451\,653\) bytes, and SHA--256
\[
 \texttt{515b2960b1a22bb59687e858f90b2eb0079463ab4523604a2b517837828cee73}.
\]
The next compulsory companion audit is V90.

\clearpage
\part{Third-series V86: maximal-row routing in the totally depleted
sector}
\label{part:V86}

\section{Purpose}
\label{sec:V86-purpose}

After V85, every hypothetical selected \(2+1+1\) counterpacket has
all six normalized pair overlaps tending to zero.  This does not
force physical privacy: two rows may share many coordinates while
their normalized weights are diffuse.  The correct upstream move is
to root the physical incidence partition at a globally maximal row.
Marked concentration is already closed by V85.  On the remaining
cofinal unmarked bank, maximality removes the mass-escalation branch
of the V80 trichotomy.  A cofinal gap closes exponentially; the only
surviving geometry is an owned comparable bank whose maximal-row
weights are quantitatively diffuse.  This chapter proves each step.

\section{A polynomial floor for the full atlas}
\label{sec:V86-atlas-floor}

\begin{lemma}[Uniform joining-coordinate overlap floor]
\label{lem:V86-edge-floor}
Let \(r\) be the selected joining coordinate of a nonzero
\(2+1+1\) word and put
\[
 M:=\max_{1\le i\le4}\Xi_i.
\]
There is a fixed \(c_0>0\) such that
\begin{equation}
 \theta_{ij}\ge {c_0\over M^2}
 \qquad(1\le i<j\le4).
\label{eq:V86-edge-floor}
\end{equation}
\end{lemma}

\begin{proof}
All four rows are active at \(r\) by
\Cref{lem:V74-four-active}; hence
\(a_{i,r}\ge a_*>0\).  Since the \(r\)-summand occurs in every
pair stock,
\[
 H_{ij}\ge a_{i,r}a_{j,r}\ge a_*^2.
\]
Also \(\Xi_i\Xi_j\le M^2\).  Thus
\(\theta_{ij}=H_{ij}/(\Xi_i\Xi_j)\ge a_*^2/M^2\).
\end{proof}

\begin{corollary}[Polynomial lower bound for the \(2+1+1\) atlas]
\label{cor:V86-atlas-floor}
With
\[
 p=\theta_{12},\qquad
 q=xy+xz+yz,\qquad
 \mathscr A_{211}=pq,
\]
one has
\begin{equation}
 \boxed{\;
 \mathscr A_{211}\ge {c_1\over M^6},
 \qquad
 \mathscr A_{211}^{\,2}\ge {c_1^2\over M^{12}}.\;}
\label{eq:V86-atlas-floor}
\end{equation}
\end{corollary}

\begin{proof}
Each of \(x,y,z\) is a sum of two normalized pair edges and is
therefore at least \(2c_0M^{-2}\) by
\Cref{lem:V86-edge-floor}.  Hence
\[
 q=xy+xz+yz\ge12c_0^2M^{-4},
 \qquad
 p\ge c_0M^{-2}.
\]
Multiplication proves the assertion with \(c_1=12c_0^3\).
\end{proof}

\begin{firewall}[The floor is used only against exponential decay]
\label{fire:V86-floor-scope}
\eqref{eq:V86-atlas-floor} is intentionally coarse.  It does not
replace the exact eight-tree stock and cannot close a merely
polynomial response with the wrong incidence.  Its role below is
only to show that a cofinal heat gap at the maximal row dominates
the full atlas target.
\end{firewall}

\section{A complete unmarked maximal-row bank}
\label{sec:V86-maximal-bank}

For each word in a counterpacket, choose the least row label
\(a\) attaining \(M=\max_i\Xi_i\).  Pass to a subsequence so that
\(a\) is fixed.  Let
\[
 E_a:=\{e:e\notin Q_{r,\pi},\ a_{a,e}>0\}.
\]
This is the full unmarked physical support of row \(a\).

\begin{lemma}[Marked-small or already closed]
\label{lem:V86-marked-small}
On a persistent counterpacket,
\begin{equation}
 S_a(Q_{r,\pi})=o(\Xi_a).
\label{eq:V86-marked-small}
\end{equation}
Consequently, for all sufficiently large indices,
\begin{equation}
 S_a(E_a)\ge\frac12\Xi_a.
\label{eq:V86-unmarked-cofinal}
\end{equation}
\end{lemma}

\begin{proof}
If \eqref{eq:V86-marked-small} failed, a subsequence would satisfy
\(S_a(Q_{r,\pi})\ge\sigma\Xi_a\) for some fixed \(\sigma>0\).
\Cref{lem:V85-role-owner} would then produce a fixed-fraction
marked coordinate, and the proof of
\Cref{thm:V85-exceptional-closes} would close that subsequence.
This contradicts persistence.  The complement of
\eqref{eq:V86-marked-small} in the full row support is \(E_a\), so
\eqref{eq:V86-unmarked-cofinal} follows.
\end{proof}

\begin{lemma}[The maximal-row bank is a complete admissible source
bank]
\label{lem:V86-bank-admissible}
The bank \(E_a\) is admissible in the sense of
\Cref{def:V84-admissible-bank}.  Every fixed finite subbank selected
from it by numerical Casimir comparisons is exposed in one common
exact prefix/suffix source resolution.
\end{lemma}

\begin{proof}
The bank is disjoint from \(Q_{r,\pi}\) by definition and is
selected from physical activity and numerical weights before any
signed endpoint expansion.  Its requested V32 resolution is the
complete local physical output resolution.  Thus it is admissible.
Apply \Cref{thm:V84-universal-incidence} to \(E_a\) and to the fixed
finite family of subbanks produced by a later numerical
trichotomy.
\end{proof}

\section{Maximality leaves a gap or a diffuse owned bank}
\label{sec:V86-trichotomy}

\begin{theorem}[Totally depleted maximal-row dichotomy]
\label{thm:V86-maximal-dichotomy}
On a persistent counterpacket satisfying
\eqref{eq:V85-total-depletion}, after passage to a subsequence at
least one of the following occurs:
\begin{enumerate}[label=\textup{(\alph*)},leftmargin=4em]
\item a source-occurring row-\(a\) gapped bank \(F\subset E_a\)
      satisfies
\begin{equation}
 S_a(F)\ge\sigma\Xi_a;
\label{eq:V86-gapped-cofinal}
\end{equation}
\item for one fixed \(c\ne a\), a source-occurring bank
      \(D\subset E_a\) is pointwise two-sided comparable between
      \(a,c\) and owns fixed fractions of both \(\Xi_a,\Xi_c\).
\end{enumerate}
\label{eq:V86-dichotomy}
\end{theorem}

\begin{proof}
Apply \Cref{thm:V80-cofinal-trichotomy} to the complete unmarked bank
\(E_a\), using \eqref{eq:V86-unmarked-cofinal}.  The mass-escalation
alternative would give \(\Xi_c\ge R_0\Xi_a\) with \(R_0>1\), contrary
to the definition \(\Xi_a=M\).  The other two alternatives are the
gapped and owned comparable banks above.  Their exact source
occurrence, including a cofinal or owned prefix/suffix half, follows
from
\Cref{lem:V86-bank-admissible,
lem:V84-gapped-half,
lem:V84-owned-half}.  There are only three possible partner labels,
so one is fixed after taking a subsequence.
\end{proof}

\begin{theorem}[The cofinal gapped alternative is impossible]
\label{thm:V86-gap-closes}
Alternative \textup{(a)} of
\eqref{eq:V86-dichotomy} satisfies the full \(FCQ_{211}\) estimate
and cannot support a counterpacket.
\end{theorem}

\begin{proof}
The V64 Casimir gap estimate used in
\Cref{lem:V80-gap-chi} gives
\[
 \|\mathcal G_a(F)\|
 \le e^{-c_{G,s,\sigma}\Xi_a}
 =e^{-c_{G,s,\sigma}M}.
\]
By the exact source occurrence from
\Cref{thm:V86-maximal-dichotomy}, this is a central positive factor
inside each term of one fixed finite expansion.  All complete
spectators are contractions; the once-paid factor and the fixed
local cumulant arity have the uniform bounds used in
\Cref{thm:V66-all-terminal-gaps,thm:V80-gap-HCRC}.  Therefore each
oriented complete-word square has product-state capacity at most
\[
 C_{G,s}e^{-2c_{G,s,\sigma}M}.
\]
For \(M\ge1\),
\[
 e^{-2cM}\le C_cM^{-12}.
\]
\Cref{cor:V86-atlas-floor} then bounds the last display by
\(C_{G,s,\sigma}\mathscr A_{211}^{\,2}\).  Fixed finite
recombination proves \(FCQ_{211}\).
\end{proof}

\begin{lemma}[The owned comparable alternative is effectively
diffuse]
\label{lem:V86-owned-diffuse}
Let \(D\) be alternative \textup{(b)} of
\eqref{eq:V86-dichotomy}, with
\[
 \lambda a_{a,e}\le a_{c,e}\le\Lambda a_{a,e}
 \quad(e\in D),
 \qquad
 S_a(D)\ge\sigma_a\Xi_a.
\]
Then
\begin{equation}
 {1\over\Xi_a^2}\sum_{e\in D}a_{a,e}^2
 \le {1\over\lambda}\theta_{ac},
\label{eq:V86-l2-diffuse}
\end{equation}
and consequently
\begin{equation}
 \max_{e\in D}{a_{a,e}\over\Xi_a}
 \le\lambda^{-1/2}\theta_{ac}^{1/2}
 \longrightarrow0.
\label{eq:V86-linf-diffuse}
\end{equation}
\end{lemma}

\begin{proof}
Pointwise lower comparability gives
\[
 H_{ac}\ge H_{ac}(D)
 =\sum_{e\in D}a_{a,e}a_{c,e}
 \ge\lambda\sum_{e\in D}a_{a,e}^2.
\]
Since the root row is globally maximal,
\(\Xi_c\le\Xi_a\).  Hence
\[
 \sum_{e\in D}a_{a,e}^2
 \le{\theta_{ac}\Xi_a\Xi_c\over\lambda}
 \le{\theta_{ac}\Xi_a^2\over\lambda},
\]
which is \eqref{eq:V86-l2-diffuse}.  The maximum square is at most
the sum of squares.  Finally
\(\theta_{ac}\to0\) by
\eqref{eq:V76-all-edges-zero}.
\end{proof}

\begin{theorem}[Exact residual after maximal-row routing]
\label{thm:V86-residual}
Every persistent totally depleted \(2+1+1\) counterpacket has, after
passing to a subsequence, a fixed maximal row \(a\), a fixed partner
\(c\ne a\), and a source-occurring pointwise-comparable bank \(D\)
which:
\begin{enumerate}[label=\textup{(\roman*)},leftmargin=4em]
\item owns fixed fractions of both endpoint row masses;
\item avoids all marked and paid coordinates;
\item is jointly exposed with every earlier required bank in one
      exact finite resolution; and
\item obeys the effective diffuseness estimates
      \eqref{eq:V86-l2-diffuse}--%
      \eqref{eq:V86-linf-diffuse}.
\end{enumerate}
No cofinal gapped or marked-concentration alternative remains.
\end{theorem}

\begin{proof}
\Cref{thm:V86-maximal-dichotomy} is exhaustive.  Its gapped branch
is removed by \Cref{thm:V86-gap-closes}; marked concentration was
removed by \Cref{lem:V86-marked-small}.  The properties of the
remaining owned bank follow from
\Cref{thm:V86-maximal-dichotomy,
lem:V86-bank-admissible,
thm:V84-universal-incidence,
lem:V86-owned-diffuse}.
\end{proof}

\begin{firewall}[Diffuseness alone is not yet the full tree
currency]
\label{fire:V86-diffuse-not-closed}
V83 turns a diffuse owned \emph{two-row internal heat bank} into two
inverse powers of one endpoint mass.  The bank in
\Cref{thm:V86-residual} may lie in a complete three- or four-row
moment slot, and its pair \(ac\) need not be one of the two edges
required by the same legal \(2+1+1\) tree term.  Applying V83 without
first preserving the physical local block and the tree incidence
would repeat the incompatible-bank error of V76.
\end{firewall}

\section{V86 result ledger and V87 acquisition order}
\label{sec:V86-status}

\begin{theorem}[Third-series V86 result ledger]
\label{thm:V86-results}
Relative to V1--V85, V86 proves:
\begin{enumerate}[label=\textup{(V86.\arabic*)},leftmargin=6.3em]
\item a uniform \(M^{-6}\) lower bound for the full
      \(2+1+1\) atlas at a nonzero joining word;
\item subcofinality of the entire finite marked packet on any
      persistent totally depleted counterpacket;
\item exact source admissibility of the complete unmarked support of
      a globally maximal row;
\item a maximal-row dichotomy between a cofinal gap and an owned
      comparable bank;
\item exponential closure of the cofinal-gap branch at the full
      atlas scale;
\item quantitative \(\ell^2\) and atomwise diffuseness of the owned
      comparable branch; and
\item reduction of total depletion to a source-occurring, marked-free,
      owned diffuse comparable bank with a fixed endpoint pair.
\end{enumerate}
\end{theorem}

\begin{proof}
These are
\Cref{lem:V86-edge-floor,
cor:V86-atlas-floor,
lem:V86-marked-small,
lem:V86-bank-admissible,
thm:V86-maximal-dichotomy,
thm:V86-gap-closes,
lem:V86-owned-diffuse,
thm:V86-residual}.
\end{proof}

\begin{programme}[V87: complete-slot partial transpose for the
diffuse pair bank]
\label{prog:V86-V87}
\begin{enumerate}[leftmargin=3em]
\item Fix the partner pair \(ac\) of
      \Cref{thm:V86-residual} and return to the exact local
      moment-partition state on every coordinate of \(D\).
\item Decompose the complementary active rows into irreducible
      blocks before inserting the heat output, keeping all
      multiplicities.
\item Extend the V83 positive-block marginal sandwich from one
      irreducible \(V\otimes W\) block to the complete fixed-arity
      physical local sum without paying the number of complementary
      output labels.
\item If the extension holds, use the four-largest-coordinate lemma
      with \eqref{eq:V86-l2-diffuse} and a fixed full-row transpose.
\item Only after obtaining the complete bank capacity, test its pair
      \(ac\) against the eight legal tree incidences.  Reroute the
      tree rather than declaring an arbitrary pair compatible.
\item If complementary multiplicities obstruct the heat sum, isolate
      their exact trace growth and test whether the second owned row
      absorbs it.
\end{enumerate}
\end{programme}

\begin{firewall}[Claim boundary after third-series V86]
\label{fire:V86-claim-boundary}
V86 closes the cofinal-gap part of total depletion and reduces every
remaining selected \(2+1+1\) counterpacket to one exact
source-occurring diffuse owned comparable bank.  It does not yet
prove the complete-slot partial-transpose estimate or its legal-tree
routing, close the selected source or its global aggregation, close
other quartic source families, higher MTA, WUFC, CPM, cutoff removal,
reflection positivity, Osterwalder--Schrader reconstruction,
construction of the continuum Yang--Mills measure, or a uniform
positive continuum spectral gap.  The full Yang--Mills mass gap has
not been proved.
\end{firewall}

\paragraph{Parent-version record.}
V86 was formed from
\texttt{Renewal\_Yang\_Mills\_Mass\_Gap\_Research\_Notes\_V85.tex}.
The V85 parent had \(41\,272\) logical source lines,
\(1\,466\,858\) bytes, and SHA--256
\[
 \texttt{0b6f7ccc251e1b8b3457dd391a0b470b5d96e65937875a9ba4236862881f4c22}.
\]
The next compulsory companion audit is V90.

\clearpage
\part{Third-series V87: reducible-complement heat tensorization and
closure of total depletion}
\label{part:V87}

\section{Purpose and exponent choice}
\label{sec:V87-purpose}

V86 reduces total depletion to a marked-free source-occurring bank
which owns the globally maximal row and is diffuse on that row.
There are two remaining questions.  First, the block containing the
maximal row may contain several other rows, so its complementary
module is reducible.  Second, the full atlas square has only the
coarse lower bound \(M^{-12}\), whereas V83 selected four
coordinates and obtained \(M^{-2}\).

Both points have direct resolutions.  Decomposing the complementary
module into irreducibles makes the heat operator a direct sum, and
partial-transpose norm takes the maximum, not the sum, over those
blocks.  Then the \(k\)-largest-coordinate argument may be run for
any fixed \(k\).  Choosing \(k=24\) gives \(M^{-12}\), exactly the
power required by \eqref{eq:V86-atlas-floor}.

\section{The complete-slot heat card}
\label{sec:V87-complete-slot}

\begin{theorem}[Reducible-complement output heat card]
\label{thm:V87-reducible-OIPT}
Let \(V\) be an irreducible unitary \(G\)-module and let \(W\) be an
arbitrary finite-dimensional unitary \(G\)-module.  Decompose
\[
 W\simeq
 \bigoplus_{T\in\widehat G}
   \mathcal N_T\otimes T
\]
orthogonally.  Let the complete output heat multiplier on
\(V\otimes W\) be
\begin{equation}
 M_{V,W}(s):=
 \bigoplus_{T\in\widehat G}
 I_{\mathcal N_T}\otimes M_{V,T}(s),
\label{eq:V87-reducible-heat}
\end{equation}
where \(M_{V,T}(s)\) is the full multiplicity-preserving operator of
\Cref{thm:V83-OIPT}.  Then, with partial transpose on either one
fixed side of \(V|W\),
\begin{equation}
 \boxed{\;
 \|M_{V,W}(s)^\Gamma\|
 \le {K_G(s)\over\dim V}.\;}
\label{eq:V87-reducible-OIPT}
\end{equation}
The constant is independent of the irreducible decomposition and
all multiplicities \(\dim\mathcal N_T\).
\end{theorem}

\begin{proof}
Choose an orthonormal basis adapted to the displayed decomposition
of \(W\).  Partial transpose of
\eqref{eq:V87-reducible-heat} preserves its orthogonal direct-sum
form and transposes the identity on each multiplicity space.  Hence
\[
 \|M_{V,W}(s)^\Gamma\|
 =
 \max_{T:\mathcal N_T\ne0}
 \|M_{V,T}(s)^\Gamma\|.
\]
By \Cref{thm:V83-OIPT}, every term on the right is at most
\[
 {K_G(s)\over\max\{\dim V,\dim T\}}
 \le {K_G(s)\over\dim V}.
\]
This proves the estimate.  A different orthonormal basis of \(W\)
conjugates the partially transposed operator by the conjugate and
transpose of a unitary, so its norm is unchanged.
\end{proof}

\begin{corollary}[Complete moment-partition local factor]
\label{cor:V87-local-slot}
Fix one endpoint-expanded local moment-partition state at a
coordinate \(e\), and let row \(a\) be active there with irreducible
module \(V_{\lambda_{a,e}}\).  Let \(\mathbb M_e(s)\) be the product
of the complete output heat multipliers on all blocks of that local
partition.  Across the cut \(a|a^c\),
\begin{equation}
 \boxed{\;
 \|\mathbb M_e(s)^{\Gamma_{a^c}}\|
 \le
 {K_G(s)\over d_{\lambda_{a,e}}}
 \le C_{G,s}a_{a,e}^{-1/2}.\;}
\label{eq:V87-local-slot}
\end{equation}
\end{corollary}

\begin{proof}
Let \(B_a\) be the partition block containing row \(a\), and regroup
its module as
\[
 V_{\lambda_{a,e}}\otimes
 W_e,\qquad
 W_e=\bigotimes_{i\in B_a\setminus\{a\}}
 V_{\lambda_{i,e}}.
\]
The first inequality for the \(B_a\)-factor is
\Cref{thm:V87-reducible-OIPT}.  Every other partition block lies
wholly on the complementary side; full transpose preserves the
norm of its positive heat contraction.  Multiplication over the
partition blocks proves the first inequality.  The second is the
Weyl dimension bound \Cref{lem:V69-Weyl-dimension}.
\end{proof}

\begin{firewall}[Why complementary multiplicities do not add]
\label{fire:V87-no-multiplicity-sum}
The direct-sum argument is performed before any triangle inequality.
The multiplicity spaces are orthogonal blocks on which the heat
operator is the identity tensored with the corresponding complete
\(V\otimes T\) heat multiplier.  Replacing the maximum by a sum
would create a false representation-count loss and is not done.
\end{firewall}

\section{A fixed-\(k\) diffuse-coordinate compiler}
\label{sec:V87-k-largest}

\begin{lemma}[\(k\)-largest-coordinate optimization]
\label{lem:V87-k-largest}
Fix \(k\ge1\), \(0<\rho<1\), \(C_0,\sigma>0\).  Let
\(b_1\ge\cdots\ge b_N>0\) satisfy
\begin{equation}
 \sum_{j=1}^N b_j\ge\sigma S,
 \qquad
 b_1\le{\sigma\over2k}S.
\label{eq:V87-k-hyp}
\end{equation}
If numbers \(r_j\) obey
\[
 r_j\le\min\{\rho,C_0b_j^{-1/2}\}
 \quad(1\le j\le k),
 \qquad
 r_j\le\rho\quad(j>k),
\]
then
\begin{equation}
 \prod_{j=1}^Nr_j\le
 C_{k,\rho,C_0,\sigma}S^{-k/2}.
\label{eq:V87-k-product}
\end{equation}
\end{lemma}

\begin{proof}
The second condition in \eqref{eq:V87-k-hyp} and the mass lower
bound imply \(N\ge2k\).  For \(1\le j\le k\), the first \(j-1\)
weights carry at most
\[
 (j-1)b_1\le{k-1\over2k}\sigma S<{\sigma S\over2}.
\]
The remaining weights therefore carry at least \(\sigma S/2\).
Since \(b_j\) is the largest remaining weight,
\[
 b_j\ge{\sigma S\over2(N-j+1)}
 \ge{\sigma S\over2N}.
\]
Consequently
\[
 \prod_{j=1}^Nr_j
 \le
 C_0^k
 \left({2N\over\sigma S}\right)^{k/2}
 \rho^{N-k}.
\]
The supremum of \(N^{k/2}\rho^{N-k}\) over integers \(N\ge k\) is
finite.  Absorb it into the stated constant.
\end{proof}

\begin{theorem}[Arbitrary fixed inverse-mass power on a diffuse
complete heat bank]
\label{thm:V87-diffuse-bank-power}
Let \(D\) be a complete source-occurring bank on which row \(a\) is
active at every coordinate.  Suppose
\begin{equation}
 S_a(D)\ge\sigma\Xi_a,
 \qquad
 \max_{e\in D}{a_{a,e}\over\Xi_a}
 \le{\sigma\over2k}.
\label{eq:V87-bank-diffuse-hyp}
\end{equation}
Let
\[
 \mathbb M_D(s):=\bigotimes_{e\in D}\mathbb M_e(s)
\]
be the complete product of local moment-partition heat factors.
For every fixed \(k\),
\begin{equation}
 \boxed{\;
 \kappa_{a|a^c}(\mathbb M_D(s))
 \le C_{G,s,k,\sigma}\Xi_a^{-k/2}.\;}
\label{eq:V87-diffuse-bank-power}
\end{equation}
The row states may be arbitrarily entangled among coordinates on
each side of the cut.
\end{theorem}

\begin{proof}
Order \(b_e=a_{a,e}\) decreasingly.  On each of the first \(k\)
coordinates there are two positive dominants of the physical heat
factor.  The factor \(\mathbb M_e(s)\) itself has fixed-row
partial-transpose norm at most
\(C_{G,s}b_e^{-1/2}\) by
\Cref{cor:V87-local-slot}.  The V68 isotropic majorant
\[
 B_e(s)=P_e^{\rm prot}+r_s(I-P_e^{\rm prot})
\]
dominates \(\mathbb M_e(s)\) and, by
\Cref{lem:V68-heat-majorant,lem:V69-local-row3-PT} with the row
label relabelled, has partial-transpose norm at most a fixed
\(\rho_s<1\).  Choose the dominant giving the smaller proved norm.
At all remaining coordinates choose \(B_e(s)\).

Their tensor product \(H_D\) positively dominates
\(\mathbb M_D(s)\).  For every product density across the complete
row cut, first use this positive order and then transpose the fixed
full complementary side.  The norm tensorizes coordinatewise, so
\[
 \kappa_{a|a^c}(\mathbb M_D(s))
 \le
 \prod_{j=1}^k
 \min\{\rho_s,C_{G,s}b_j^{-1/2}\}
 \rho_s^{N-k}.
\]
Apply \Cref{lem:V87-k-largest} with \(S=\Xi_a\).
\end{proof}

\begin{corollary}[The V86 residual bank supplies twelve powers]
\label{cor:V87-twelve-powers}
For the bank \(D\) in \Cref{thm:V86-residual}, along every
persistent totally depleted counterpacket and for all sufficiently
large indices,
\begin{equation}
 \boxed{\;
 \kappa_{a|a^c}(\mathbb M_D(s))
 \le C_{G,s}\,M^{-12},
 \qquad M=\Xi_a=\max_i\Xi_i.\;}
\label{eq:V87-twelve-powers}
\end{equation}
The same bound holds for every positive protected/gapped bank sector
\(\mathcal M_D\preccurlyeq\mathbb M_D(s)\), and for
\(\mathcal M_D^2\).
\end{corollary}

\begin{proof}
Take \(k=24\).  Ownership supplies the fixed \(\sigma\) in
\eqref{eq:V87-bank-diffuse-hyp}.  By
\eqref{eq:V86-linf-diffuse}, its atomwise upper bound tends to zero,
so the second hypothesis of
\eqref{eq:V87-bank-diffuse-hyp} holds eventually.  Apply
\Cref{thm:V87-diffuse-bank-power}.

Every resolved positive bank sector is bounded by the complete heat
product.  Also \(0\preccurlyeq\mathcal M_D\preccurlyeq I\), so
\(\mathcal M_D^2\preccurlyeq\mathcal M_D\).  Monotonicity of
product-state capacity proves the remaining assertions.
\end{proof}

\section{Closure of total depletion and of \(FCQ_{211}\)}
\label{sec:V87-closure}

\begin{theorem}[No totally depleted counterpacket]
\label{thm:V87-no-total-depletion}
The totally depleted alternative
\eqref{eq:V85-total-depletion} cannot support a sequence violating
every uniform \(FCQ_{211}\) estimate.
\end{theorem}

\begin{proof}
Use the fixed exact resolution and source-occurring bank \(D\) of
\Cref{thm:V86-residual}.  It is disjoint from the selected joining
coordinate and from the payer coordinate.  Thus the all-payer
factorization \Cref{lem:V76-prefix-bank-factorization} applies for
every payer placement and both square orientations.  The central
insertion inequality used in
\Cref{thm:V76-prefix-bank-product} gives
\begin{align}
 \kappa_\otimes(
 \widetilde W_{r,\pi,\nu}^*
 \widetilde W_{r,\pi,\nu})
 &\le
 C_s
 \left({p_r^-\over\Xi_3\Xi_4}\right)^2
 \kappa_{a|a^c}(\mathcal M_D^2)
\notag\\
 &\le C_{G,s}M^{-12}.
\label{eq:V87-word-M12}
\end{align}
Here \(p_r^-\le1\) and the active row floor absorbs
\(\Xi_3^{-2}\Xi_4^{-2}\); the second line is
\Cref{cor:V87-twelve-powers}.  No compatibility between the pair
\(ac\) and a preselected quotient label is required, because the
last line is compared with the complete atlas rather than with one
named edge.

In this insertion, the fully row-product capacity is no larger than
the bipartite capacity across \(a|a^c\): the latter permits an
arbitrary, possibly row-entangled, density on the complementary
side.  Thus using \(\kappa_{a|a^c}(\mathcal M_D^2)\) in the first
line is a valid enlargement of the original test class.

By \Cref{cor:V86-atlas-floor},
\[
 M^{-12}\le C\mathscr A_{211}^{\,2}.
\]
Thus every term of the fixed resolution satisfies the required
oriented \(FCQ_{211}\) square estimate.  The resolution has fixed
cardinality, so \Cref{lem:V58-finite-recombination} gives the same
bound for the complete word, contradicting the counterpacket
hypothesis.
\end{proof}

\begin{theorem}[Uniform complete selected-word \(FCQ_{211}\)]
\label{thm:V87-FCQ211}
For every fixed compact connected semisimple \(G\), admissible heat
parameter, selected joining coordinate, and payer placement, the
complete \(2+1+1\) word satisfies
\begin{equation}
 \boxed{\;
 \kappa_\otimes(
 \widetilde W_{r,\pi}^{\,*}\widetilde W_{r,\pi})
 \le C_{G,s}\mathscr A_{211}^{\,2},\;}
\label{eq:V87-FCQ211}
\end{equation}
or the corresponding left-square estimate dictated by output
duality.  The constant is uniform in support, representations,
physical multiplicities, and the selected coordinate.
\end{theorem}

\begin{proof}
If no such uniform constant existed, choose a sequence for which
the ratio of the left side to
\(\mathscr A_{211}^{\,2}\) diverges.  The counterpacket reductions
\Cref{thm:V76-counterpacket,thm:V77-counterpacket} give, after
subsequences, either a \(z\)-nondepleted heavy-composite packet or
total depletion.  The first is impossible by
\Cref{cor:V85-heavy-composite-closed}; the second is impossible by
\Cref{thm:V87-no-total-depletion}.  This contradiction proves the
uniform estimate.
\end{proof}

\begin{corollary}[The full composite \(2+1+1\) card is acquired]
\label{cor:V87-composite-card}
The source-level card \(FCQ_{211}(\pi)\) of
\Cref{def:V49-FCQ} holds for every payer placement.  Consequently
the complete \(2+1+1\) joining packet satisfies its literal
marked-tree source estimate by
\Cref{thm:V49-FCQ-suffices}.
\end{corollary}

\begin{proof}
\Cref{thm:V87-FCQ211} is exactly the required oriented complete-word
square estimate.  Apply \Cref{thm:V49-FCQ-suffices}.
\end{proof}

\begin{firewall}[Boundary of the acquired source card]
\label{fire:V87-source-boundary}
The theorem closes the complete fixed \(2+1+1\)
first-connectivity joining packet, including every payer placement.
It does not by itself prove the subsequent global aggregation over
all source families, nor does it treat the distinct \(3+1\),
\(2+2\), discrete, globally aggregated \([4]\), or \([3,3]\)
quartic families.  Those require their own exact source
identities and routing.
\end{firewall}

\section{V87 result ledger and V88 acquisition order}
\label{sec:V87-status}

\begin{theorem}[Third-series V87 result ledger]
\label{thm:V87-results}
Relative to V1--V86, V87 proves:
\begin{enumerate}[label=\textup{(V87.\arabic*)},leftmargin=6.3em]
\item the multiplicity-free reducible-complement extension of the
      V83 partial-transpose heat card;
\item the corresponding complete moment-partition local estimate;
\item a \(k\)-largest-coordinate optimization for arbitrary fixed
      \(k\);
\item arbitrary fixed inverse maximal-row mass powers on a diffuse
      source-occurring heat bank;
\item the \(M^{-12}\) bank capacity required by the full atlas
      square;
\item closure of the totally depleted counterpacket for all payer
      placements;
\item the uniform complete selected-word \(FCQ_{211}\) theorem; and
\item acquisition of the complete composite \(2+1+1\) source card.
\end{enumerate}
\end{theorem}

\begin{proof}
These are
\Cref{thm:V87-reducible-OIPT,
cor:V87-local-slot,
lem:V87-k-largest,
thm:V87-diffuse-bank-power,
cor:V87-twelve-powers,
thm:V87-no-total-depletion,
thm:V87-FCQ211,
cor:V87-composite-card}.
\end{proof}

\begin{programme}[V88: install the \(2+1+1\) card in the quartic
source census]
\label{prog:V87-V88}
\begin{enumerate}[leftmargin=3em]
\item Return to the exact proper-quartic partition census of V31 and
      mark every labelled \(2+1+1\) contribution now covered by
      \Cref{cor:V87-composite-card}.
\item Check the first-connectivity summation over joining
      coordinates before declaring the entire labelled family
      closed; retain the fixed finite payer ledger.
\item Identify which remaining \(3+1\), \(2+2\), and discrete
      source families are related by a row permutation and which
      require a genuinely different connected-prefix identity.
\item Apply the V87 \(k\)-coordinate diffuse-bank compiler only to a
      bank exposed in the same exact source word.
\item Isolate the next minimal quartic carrier after removing all
      \(2+1+1\) labels, and go upstream to its complete source
      identity if its existing estimate is only conditional.
\item The next compulsory companion audit remains V90.
\end{enumerate}
\end{programme}

\begin{firewall}[Claim boundary after third-series V87]
\label{fire:V87-claim-boundary}
V87 proves the full complete-word \(FCQ_{211}\) card and closes the
fixed \(2+1+1\) joining packet.  Global aggregation of that packet,
other quartic source families, higher MTA, WUFC, CPM, cutoff removal,
reflection positivity, Osterwalder--Schrader reconstruction,
construction of the continuum Yang--Mills measure, and a uniform
positive continuum spectral gap remain open.  The full Yang--Mills
mass gap has not been proved.
\end{firewall}

\paragraph{Parent-version record.}
V87 was formed from
\texttt{Renewal\_Yang\_Mills\_Mass\_Gap\_Research\_Notes\_V86.tex}.
The V86 parent had \(41\,657\) logical source lines,
\(1\,480\,000\) bytes, and SHA--256
\[
 \texttt{6b39e1e685218c95475ee93b34b35c624f8ba7d580d2a06280cb5e70a07f2fb6}.
\]
The next compulsory companion audit is V90.

\clearpage
\part{Third-series V88: quartic source census after \(FCQ_{211}\)
and the exact \(3+1\) first-connectivity word}
\label{part:V88}

\section{Census update without overclaiming aggregation}
\label{sec:V88-census}

\begin{theorem}[Nine of the fourteen proper-prefix labels have
complete fixed-packet cards]
\label{thm:V88-nine-labels}
In the proper quartic prefix census of
\Cref{lem:V31-prefix-census}:
\begin{enumerate}[label=\textup{(\roman*)},leftmargin=4em]
\item all three labelled \(2+2\) families satisfy their complete
      global source theorem by \Cref{cor:V34-S22-no-gap}; and
\item all six labelled \(2+1+1\) prefixes satisfy the complete
      fixed-joining-packet card of
      \Cref{cor:V87-composite-card}.
\end{enumerate}
The proper-prefix labels not covered by either statement are the
four \(3+1\) labels and the one discrete label.
\end{theorem}

\begin{proof}
The first count and conclusion are exactly
\Cref{cor:V34-S22-no-gap}.  A \(2+1+1\) prefix is determined by its
unique pair, hence there are \(\binom42=6\).  The proofs of
\Cref{thm:V87-FCQ211,cor:V87-composite-card} use only a relabelling
of the pair as \(12\) and the singleton rows as \(3,4\); all
constants are invariant under the finite row permutation.  Thus all
six labels are covered.  Subtracting \(3+6\) from the fourteen
labels in \Cref{lem:V31-prefix-census} leaves the stated four plus
one labels.
\end{proof}

\begin{firewall}[Fixed packet versus joining-coordinate sum]
\label{fire:V88-fixed-vs-global}
\Cref{thm:V87-FCQ211} is uniform in the selected coordinate \(r\),
but its right side is the complete global atlas
\(\mathscr A_{211}\).  The estimate
\[
 \kappa_\otimes^{\rm abs}(W_r)
 \le C\mathscr A_{211}
\]
for each \(r\) does not by itself imply the same estimate for
\(\sum_rW_r\) with a support-independent constant.  The \(2+2\)
statement in \Cref{thm:V88-nine-labels} is global; the \(2+1+1\)
statement is deliberately recorded only at fixed joining packet
until the weighted summation below is acquired.
\end{firewall}

\begin{definition}[Weighted \(2+1+1\) joining card]
\label{def:V88-WFCQ211}
For a fixed payer type \(\pi\), let
\(\mathfrak a_{r,\pi}\ge0\) be the literal normalized marked-tree
stock of the \(r\)-summand in
\eqref{eq:V35-S211}, with all coordinate marks retained.  The
\(WFCQ_{211}(\pi)\) card is
\begin{align}
 \kappa_\otimes^{\rm abs}(W_{r,\pi})
 &\le C\mathfrak a_{r,\pi},
\label{eq:V88-weighted-local}\\
 \sum_r\mathfrak a_{r,\pi}
 &\le C\mathscr A_{211}.
\label{eq:V88-weighted-sum}
\end{align}
\end{definition}

\begin{lemma}[The scalar summation line is already available]
\label{lem:V88-tree-summation}
For each of the eight lower-role trees in
\Cref{lem:V35-eight-trees}, summing its exact ordered coordinate
marks over the joining position is bounded by the corresponding
product of total pair stocks.  Consequently
\eqref{eq:V88-weighted-sum} holds for the sum of the eight literal
tree stocks and the already established selected-junction stock.
\end{lemma}

\begin{proof}
For a star term with prefix mark \(e<r\), nonnegativity and
Cartesian enlargement give, for \(i,j\in\{1,2\}\),
\[
 \sum_r\sum_{e<r}
 h_{12,e}h_{i3,r}h_{j4,r}
 \le
 \left(\sum_eh_{12,e}\right)
 \left(\sum_rh_{i3,r}\right)
 \left(\sum_rh_{j4,r}\right)
 =
 H_{12}H_{i3}H_{j4}.
\]
The four path terms are identical after replacing one of the two
joining arms by its literal \(34\) mark; their sums are bounded by
\(H_{12}H_{i3}H_{34}\) or
\(H_{12}H_{j4}H_{34}\).  These are exactly the eight products in
\eqref{eq:V35-eight-trees}.  Division by the fixed row
normalizations gives \eqref{eq:V88-weighted-sum}.  The selected
four-row term has its own disjoint first-selected-coordinate
summation in \Cref{thm:V38-GSEL4}; it is not duplicated among the
eight lower-role terms.
\end{proof}

\begin{proposition}[The missing \(2+1+1\) global step is precisely
the weighted local line]
\label{prop:V88-weighted-gate}
If \eqref{eq:V88-weighted-local} holds for every payer type, then the
entire first-connectivity sum \(\mathcal S_{211}\) in
\eqref{eq:V35-S211} satisfies its global literal marked-tree source
estimate.  The unweighted V87 card alone does not logically imply
\eqref{eq:V88-weighted-local}.
\end{proposition}

\begin{proof}
By subadditivity of absolute product-state capacity,
\[
 \kappa_\otimes^{\rm abs}\!\left(\sum_rW_{r,\pi}\right)
 \le
 \sum_r\kappa_\otimes^{\rm abs}(W_{r,\pi})
 \le C\sum_r\mathfrak a_{r,\pi}
 \le C\mathscr A_{211},
\]
where the last line is \Cref{lem:V88-tree-summation}.  There are
only finitely many payer types.

For nonimplication, take \(N\) positive scalar summands \(W_r=1\)
and a global scalar majorant \(\mathscr A=1\).  Every individual
bound \(|W_r|\le\mathscr A\) holds, while
\(|\sum_rW_r|=N\).  Thus uniformity in \(r\) is not a summable
weight.  This logical example does not claim that the physical
first-connectivity terms have equal sign.
\end{proof}

\section{The next untouched proper-prefix family}
\label{sec:V88-S31}

Fix the prefix partition \(123|4\); the other three \(3+1\) labels
are row permutations.  Let
\(K_{123}^{<r,\mathrm{conn}}\) be the complete connected prefix
carrier on rows \(1,2,3\), let \(M_{\{4\}}^{<r}\) be the singleton
prefix moment, and let \(M_{>r}\) be the complete four-row suffix.
At the joining coordinate put
\[
 Z_{123,r}:=X_{1,r}\otimes X_{2,r}\otimes X_{3,r},
 \qquad
 J_r(123|4):=\kappa_2(Z_{123,r},X_{4,r}).
\]

\begin{theorem}[Exact \(3+1\) first-connectivity identity]
\label{thm:V88-S31-identity}
Before any norm, positive envelope, output split, or payer
allocation,
\begin{equation}
 \boxed{\;
 \mathcal S_{31}^{123|4}
 =
 \sum_{r=1}^n
 K_{123}^{<r,\mathrm{conn}}
 \otimes M_{\{4\}}^{<r}
 \otimes J_r(123|4)
 \otimes M_{>r}.\;}
\label{eq:V88-S31}
\end{equation}
The summands are indexed by the first coordinate at which the
already connected triple block joins the singleton block.
\end{theorem}

\begin{proof}
Apply the coefficient-one first-connectivity telescope to the
ordered product of local four-row moment letters, starting from the
proper prefix partition \(123|4\).  Immediately before the first
global join, the \(123\) block is connected and row \(4\) is a
singleton, giving the first two prefix factors in
\eqref{eq:V88-S31}.  At the first joining coordinate, the quotient
partition has two blocks.  The connected local contribution between
two quotient blocks is their complete covariance
\(\kappa_2(Z_{123,r},X_{4,r})\), with all local moment-partition
roles assembled.  After the join, the suffix is the complete
four-row moment.  The first-join positions are disjoint and
exhaustive, so their coefficient-one sum is the original source.
\end{proof}

\begin{lemma}[Exact three-branch polarization of the joining
covariance]
\label{lem:V88-S31-polarization}
Let \(\xi_0,\xi_1,\xi_4\) be independent replicas.  Then
\begin{equation}
 J_r(123|4)
 =
 \mathbb E\!\left[
  \bigl(Z_{123,r}(\xi_0)-Z_{123,r}(\xi_1)\bigr)
  \otimes
  \bigl(X_{4,r}(\xi_0)-X_{4,r}(\xi_4)\bigr)
 \right].
\label{eq:V88-S31-polarization}
\end{equation}
Moreover,
\begin{align}
 Z_{123}(\xi_0)-Z_{123}(\xi_1)
 ={}&
 (X_1^0-X_1^1)\otimes X_2^0\otimes X_3^0
\notag\\
 &+X_1^1\otimes(X_2^0-X_2^1)\otimes X_3^0
\notag\\
 &+X_1^1\otimes X_2^1\otimes(X_3^0-X_3^1).
\label{eq:V88-S31-Leibniz}
\end{align}
The three branches must be reassembled before an absolute value.
\end{lemma}

\begin{proof}
Expansion of the expectation in
\eqref{eq:V88-S31-polarization} gives
\[
 \mathbb E(ZX_4)-\mathbb EZ\,\mathbb EX_4
 -\mathbb EZ\,\mathbb EX_4
 +\mathbb EZ\,\mathbb EX_4
 =
 \kappa_2(Z,X_4).
\]
The second formula follows by adding and subtracting
\(X_1^1X_2^0X_3^0\) and
\(X_1^1X_2^1X_3^0\).  It is an exact noncommutative tensor
Leibniz telescope in the fixed row order.
\end{proof}

\begin{firewall}[The three branches are not three independent
responses]
\label{fire:V88-S31-signed}
The right side of \eqref{eq:V88-S31-Leibniz} is a signed
decomposition of one composite difference.  Its square contains
nine diagonal and cross terms.  Applying three separate positive
response bounds before reassembly would discard those cross terms
and may charge the same joining occurrence more than once.
\end{firewall}

\section{The literal \(3+1\) tree atlas}
\label{sec:V88-S31-atlas}

Define the connected triple-prefix stock
\begin{equation}
 \mathscr T_{123}
 :=
 H_{12}H_{13}+H_{12}H_{23}+H_{13}H_{23},
\label{eq:V88-triple-tree}
\end{equation}
and the joining bridge stock
\begin{equation}
 \mathscr B_{123|4}:=H_{14}+H_{24}+H_{34}.
\label{eq:V88-S31-bridge}
\end{equation}

\begin{lemma}[The \(3+1\) packet has nine literal spanning trees]
\label{lem:V88-S31-nine}
The product
\begin{equation}
 \mathscr T_{123}\mathscr B_{123|4}
\label{eq:V88-S31-nine-product}
\end{equation}
expands into exactly nine labelled spanning-tree monomials on
\([4]\).  Each consists of one of the three trees on
\(\{1,2,3\}\) and one edge joining row \(4\) to that tree.
\end{lemma}

\begin{proof}
There are three spanning trees on three labelled vertices and three
choices of the edge incident to vertex \(4\).  Every resulting
three-edge graph is connected and acyclic: adding a new vertex by
one edge to a tree preserves both properties.  Conversely, a
spanning tree whose restriction to \(\{1,2,3\}\) is connected has
exactly two edges in that restriction and exactly one edge incident
to \(4\); hence it occurs uniquely in the displayed product.
\end{proof}

\begin{definition}[Normalized \(3+1\) target and square card]
\label{def:V88-FCQ31}
Put
\begin{equation}
 \mathscr A_{31}^{123|4}
 :=
 {\mathscr T_{123}\mathscr B_{123|4}
  \over\Xi_1\Xi_2\Xi_3\Xi_4}.
\label{eq:V88-A31}
\end{equation}
For a fixed joining position and payer placement, \(FCQ_{31}\) is
the oriented complete-word square estimate
\begin{equation}
 \kappa_\otimes(W_{r,\pi}^*W_{r,\pi})
 \le C\bigl(\mathscr A_{31}^{123|4}\bigr)^2,
\label{eq:V88-FCQ31}
\end{equation}
or its left-square analogue.
\end{definition}

\begin{proposition}[The exact next analytic carrier]
\label{prop:V88-S31-next}
A proof of \(FCQ_{31}\) for every payer placement, followed by the
coordinate-weighted analogue of
\Cref{prop:V88-weighted-gate}, closes all four labelled \(3+1\)
families.  The first new local object is the positive square of the
assembled covariance \eqref{eq:V88-S31-polarization}, with all nine
Leibniz cross terms retained.
\end{proposition}

\begin{proof}
Product-state Cauchy--Schwarz turns the square estimate into the
fixed-word absolute estimate.  The ordered first-connectivity marks
sum by the same nonnegative Cartesian enlargement used in
\Cref{lem:V88-tree-summation}, now to the nine monomials of
\Cref{lem:V88-S31-nine}.  Row permutation covers the four choices
of singleton.  The exact polarization has three branches, so its
positive square has \(3^2=9\) blocks; by
\Cref{fire:V88-S31-signed}, that complete square is the first
legitimate analytic carrier.
\end{proof}

\section{V88 result ledger and V89 acquisition order}
\label{sec:V88-status}

\begin{theorem}[Third-series V88 result ledger]
\label{thm:V88-results}
Relative to V1--V87, V88 proves:
\begin{enumerate}[label=\textup{(V88.\arabic*)},leftmargin=6.3em]
\item the exact post-V87 proper-prefix census: three global \(2+2\)
      labels, six fixed-packet \(2+1+1\) labels, and five untouched
      labels;
\item the distinction between a uniform fixed-\(r\) card and a
      support-uniform joining-coordinate sum;
\item summability of the literal \(2+1+1\) marked-tree stocks;
\item the precise weighted local card sufficient for global
      \(2+1+1\) aggregation;
\item the exact \(3+1\) first-connectivity identity;
\item a three-branch replica polarization of its assembled joining
      covariance;
\item the exact nine-tree \(3+1\) atlas; and
\item the complete positive-square card which is the next analytic
      target.
\end{enumerate}
\end{theorem}

\begin{proof}
These are
\Cref{thm:V88-nine-labels,
fire:V88-fixed-vs-global,
lem:V88-tree-summation,
prop:V88-weighted-gate,
thm:V88-S31-identity,
lem:V88-S31-polarization,
lem:V88-S31-nine,
def:V88-FCQ31,
prop:V88-S31-next}.
\end{proof}

\begin{programme}[V89: first \(3+1\) source estimate]
\label{prog:V88-V89}
\begin{enumerate}[leftmargin=3em]
\item Form the exact nine-block positive square of
      \eqref{eq:V88-S31-polarization}, with one fixed payer in both
      copies.
\item Prove the thermal-volume branch directly from
      \eqref{eq:V88-A31}.
\item On a hot counterpacket, choose a maximal demanded row and
      repeat the marked/unmarked trichotomy only after identifying
      the demanded vertices of each of the nine trees.
\item Use the V87 arbitrary-\(k\) diffuse-bank theorem when a
      source-occurring bank survives; choose \(k\) from the
      polynomial floor of the \(3+1\) atlas rather than copying
      \(k=24\).
\item Keep the joining bridge covariance assembled.  Do not treat
      its three Leibniz branches as separate banks.
\item In parallel, test whether the V87 counterpacket proof can be
      localized to the weight
      \(\mathfrak a_{r,\pi}\) required by
      \eqref{eq:V88-weighted-local}.
\item The next compulsory companion audit remains V90.
\end{enumerate}
\end{programme}

\begin{firewall}[Claim boundary after third-series V88]
\label{fire:V88-claim-boundary}
V88 updates the quartic census, isolates the exact weighted
aggregation gate for \(2+1+1\), and constructs the full \(3+1\)
first-connectivity carrier and target.  It does not prove the
weighted \(2+1+1\) card, \(FCQ_{31}\), the global discrete source,
globally aggregated \([4]\), \([3,3]\), higher MTA, WUFC, CPM,
cutoff removal, reflection positivity, Osterwalder--Schrader
reconstruction, construction of the continuum Yang--Mills measure,
or a uniform positive continuum spectral gap.  The full
Yang--Mills mass gap has not been proved.
\end{firewall}

\paragraph{Parent-version record.}
V88 was formed from
\texttt{Renewal\_Yang\_Mills\_Mass\_Gap\_Research\_Notes\_V87.tex}.
The V87 parent had \(42\,105\) logical source lines,
\(1\,495\,528\) bytes, and SHA--256
\[
 \texttt{a26293157dcaf190be729b6351ba76f9211278a71162a8945b1b6d2c2ed9cee1}.
\]
The next compulsory companion audit is V90.

\clearpage
\part{Third-series V89: normalization repair for the
\texorpdfstring{\(3+1\)}{3+1} source and a reusable diffuse-bank
terminal}
\label{part:V89}

\section{Why the V88 target must be repaired upstream}
\label{sec:V89-purpose}

V88 correctly returned to the assembled \(3+1\) covariance and
correctly kept its three Leibniz branches together.  Two points in
that chapter require correction before its proposed square estimate
can be attacked.

First, the exact \(3+1\) first-connectivity word is not a previously
untouched algebraic family.  It is already the source
\(\texttt{eq:V66-S31}\) in the frozen f2 archive.  F2 V66 proves its
fully thermal DRQC estimate, and f2 V77 proves a conditional
two-block-bank resummation.  What remains open is the uniform hot
normalization and the coordinate-weighted global card, not the source
identity itself.

Second, \eqref{eq:V88-A31} divides a product of three raw pair stocks
by only one copy of each row mass.  The normalized MTA monomial divides
each raw edge stock by both endpoint masses.  The V88 expression is
therefore larger than the MTA target by the exact excess-degree
factors.  Proving its square card would not prove quartic MTA.

\begin{assessment}[Nonduplication and correction boundary]
\label{ass:V89-nonduplication}
This chapter does not reprove frozen V66 or V77.  It:
\begin{enumerate}[label=\textup{(\roman*)},leftmargin=4em]
\item replaces the erroneous V88 normalized target by the literal
      normalized tree atlas;
\item records the exact Gram-square organization of the three
      covariance branches;
\item extracts from frozen V66 a coordinate-weighted thermal square
      card; and
\item isolates an incidence-neutral diffuse-bank terminal which may
      be reused in the remaining hot counterpacket.
\end{enumerate}
\end{assessment}

\section{The corrected literal atlas}
\label{sec:V89-correct-atlas}

Recall
\[
 \theta_{uv}:={H_{uv}\over\Xi_u\Xi_v}.
\]
For the prefix \(123|4\), define
\begin{align}
 \widehat{\mathscr T}_{123}
 &:=
 \theta_{12}\theta_{13}
 +\theta_{12}\theta_{23}
 +\theta_{13}\theta_{23},
\label{eq:V89-normalized-triple}\\
 \widehat{\mathscr B}_{123|4}
 &:=
 \theta_{14}+\theta_{24}+\theta_{34},
\label{eq:V89-normalized-bridge}\\
 \widehat{\mathscr A}_{31}^{123|4}
 &:=
 \widehat{\mathscr T}_{123}
 \widehat{\mathscr B}_{123|4}.
\label{eq:V89-correct-A31}
\end{align}

\begin{lemma}[Exact excess-degree correction]
\label{lem:V89-degree-correction}
Let \(T\) be a tree on \(\{1,2,3\}\), let
\(i\in\{1,2,3\}\), and put
\(\tau=T\cup\{i4\}\).  Then
\begin{equation}
 { \left(\prod_{uv\in E(T)}H_{uv}\right)H_{i4}
   \over \Xi_1\Xi_2\Xi_3\Xi_4}
 =
 \left(\prod_{k=1}^4
       \Xi_k^{\deg_\tau(k)-1}\right)
 \prod_{uv\in E(\tau)}\theta_{uv}.
\label{eq:V89-degree-correction}
\end{equation}
Consequently
\begin{equation}
 \mathscr A_{31}^{123|4}\ \text{from \eqref{eq:V88-A31}}
 =
 \sum_{\substack{T\in\mathfrak T(\{1,2,3\})\\i\in\{1,2,3\}}}
 \left(\prod_{k=1}^4
       \Xi_k^{\deg_{T\cup\{i4\}}(k)-1}\right)
 \prod_{uv\in E(T\cup\{i4\})}\theta_{uv}.
\label{eq:V89-old-A31-expanded}
\end{equation}
The MTA target is instead
\(\widehat{\mathscr A}_{31}^{123|4}\), the same sum with every
excess-degree factor deleted.
\end{lemma}

\begin{proof}
For every edge \(uv\),
\(H_{uv}=\Xi_u\Xi_v\theta_{uv}\).  In the product over the
three edges of \(\tau\), row \(k\) therefore contributes
\(\Xi_k^{\deg_\tau(k)}\).  Division by
\(\prod_k\Xi_k\) proves \eqref{eq:V89-degree-correction}.
Summing over the three triple trees and the three choices of bridge
gives \eqref{eq:V89-old-A31-expanded}.  On the other hand,
expanding \eqref{eq:V89-correct-A31} gives precisely the nine
products \(\prod_{uv\in E(\tau)}\theta_{uv}\).
\end{proof}

\begin{corollary}[The V88 square card is insufficient]
\label{cor:V89-old-card-insufficient}
Even a uniform proof of \eqref{eq:V88-FCQ31} with the V88 target
\(\mathscr A_{31}^{123|4}\) would not imply the normalized
\(3+1\) MTA card.  For a star centred at a row \(k\), the gap is
\(\Xi_k^2\); for a path it is the product of the masses at its two
internal vertices.
\end{corollary}

\begin{proof}
In a four-vertex star the centre has degree three and every leaf
degree one, so the excess-degree product in
\eqref{eq:V89-degree-correction} is \(\Xi_k^2\).
In a path the two internal vertices have degree two and the leaves
degree one, giving the stated product.  These factors are unbounded
on hot representation sequences.  Thus a bound by the larger V88
quantity cannot be converted to a bound by the literal normalized
tree sum with a uniform constant.
\end{proof}

\begin{definition}[Correct \(3+1\) fixed-word and weighted cards]
\label{def:V89-correct-cards}
Let \(\widetilde W_{r,\pi}\) denote the complete \(r\)-th
\(3+1\) word with the four input-row normalization and the fixed
sole-payer placement \(\pi\) included.  The corrected fixed-word
card is
\begin{equation}
 \boxed{\;
 \kappa_\otimes(
   \widetilde W_{r,\pi}^{\,*}\widetilde W_{r,\pi})
 \le
 C_{G,s}
 \bigl(\widehat{\mathscr A}_{31}^{123|4}\bigr)^2,\;}
\label{eq:V89-correct-FCQ31}
\end{equation}
or the left-square analogue selected by output duality.

For
\begin{align}
 \theta_{uv}^{<r}
 &:={H_{uv}^{<r}\over\Xi_u\Xi_v},
 &
 \theta_{i4,r}
 &:={h_{i4,r}\over\Xi_i\Xi_4},
\label{eq:V89-local-theta}\\
 \mathfrak a_{31,r}
 &:=
 \sum_{\substack{T\in\mathfrak T(\{1,2,3\})\\i\in\{1,2,3\}}}
 \left(\prod_{uv\in E(T)}\theta_{uv}^{<r}\right)
 \theta_{i4,r},
\label{eq:V89-local-a31}
\end{align}
the weighted card is
\begin{equation}
 \opnorm{\widetilde W_{r,\pi}}
 \le C_{G,s}\mathfrak a_{31,r}.
\label{eq:V89-weighted-card}
\end{equation}
\end{definition}

\begin{lemma}[The corrected local stocks are summable]
\label{lem:V89-correct-summation}
\begin{equation}
 \boxed{\;
 \sum_r\mathfrak a_{31,r}
 \le
 \widehat{\mathscr A}_{31}^{123|4}.\;}
\label{eq:V89-correct-summation}
\end{equation}
\end{lemma}

\begin{proof}
Fix \(T\) and \(i\).  Expanding the two prefix stocks into their
nonnegative coordinate sums and then discarding their restrictions
to coordinates \(<r\) gives
\[
 \sum_r
 \left(\prod_{uv\in E(T)}\theta_{uv}^{<r}\right)
 \theta_{i4,r}
 \le
 \left(\prod_{uv\in E(T)}\theta_{uv}\right)\theta_{i4}.
\]
Summing over \(T\) and \(i\) gives
\eqref{eq:V89-correct-summation}.  This is a Cartesian enlargement;
it neither counts a support size nor asserts independence of the two
prefix marks.
\end{proof}

\begin{firewall}[Superseded notation]
\label{fire:V89-superseded}
From this point onward, every reference to \(FCQ_{31}\) means
\eqref{eq:V89-correct-FCQ31}, and every \(3+1\) atlas means
\eqref{eq:V89-correct-A31}.  Equations
\eqref{eq:V88-A31} and \eqref{eq:V88-FCQ31} are retained only as an
append-only audit trail and are not used as proof interfaces.
\end{firewall}

\section{The assembled covariance square}
\label{sec:V89-Gram-square}

Use the three independent replicas in
\eqref{eq:V88-S31-polarization}.  Let \(D_i\) be the \(i\)-th
summand of the ordered product difference
\eqref{eq:V88-S31-Leibniz}, let
\(\Delta_4=X_4(\xi_0)-X_4(\xi_4)\), and put
\begin{equation}
 J_{i,r}:=\mathbb E[D_i\otimes\Delta_4],
 \qquad
 J_r=\sum_{i=1}^3J_{i,r}.
\label{eq:V89-joining-branches}
\end{equation}
Fix the payer location \(\pi\) before making this decomposition, and
write
\(\widetilde W_{r,\pi}=\sum_i\widetilde W_{i,r,\pi}\)
for the resulting linear branch decomposition of the complete word.

\begin{lemma}[Exact nine-block Gram square]
\label{lem:V89-nine-Gram}
For either square orientation,
\begin{align}
 \widetilde W_{r,\pi}^{\,*}\widetilde W_{r,\pi}
 &=
 \sum_{i,j=1}^3
 \widetilde W_{i,r,\pi}^{\,*}
 \widetilde W_{j,r,\pi},
\label{eq:V89-nine-square}\\
 \widetilde W_{r,\pi}\widetilde W_{r,\pi}^{\,*}
 &=
 \sum_{i,j=1}^3
 \widetilde W_{i,r,\pi}
 \widetilde W_{j,r,\pi}^{\,*}.
\label{eq:V89-nine-left-square}
\end{align}
The block matrix
\[
 \left(
 \widetilde W_{i,r,\pi}^{\,*}
 \widetilde W_{j,r,\pi}
 \right)_{i,j=1}^3
\]
is positive on the three-fold Hilbert-space direct sum.  The six
off-diagonal summands in \eqref{eq:V89-nine-square} need not be
positive individually.
\end{lemma}

\begin{proof}
The two identities are distributivity applied after the one payer
has been fixed.  For vectors \(v_1,v_2,v_3\),
\[
 \sum_{i,j}
 \left\langle
 v_i,\widetilde W_i^*\widetilde W_jv_j
 \right\rangle
 =
 \left\|\sum_j\widetilde W_jv_j\right\|^2
 \ge0.
\]
This proves block positivity.  No corresponding statement holds for
an individual cross term, which is generally not even self-adjoint.
\end{proof}

\begin{corollary}[The assembled norm controls the complete square]
\label{cor:V89-norm-controls-square}
If
\[
 \opnorm{\widetilde W_{r,\pi}}\le C\mathfrak a_{31,r},
\]
then both oriented square cards hold with right side
\(C^2\mathfrak a_{31,r}^2\).  It is neither necessary nor legitimate
to give the nine terms of \eqref{eq:V89-nine-square} nine separate
positive envelopes.
\end{corollary}

\begin{proof}
For every product state \(\rho\),
\[
 \left|\Tr\!\left(
  \rho\,\widetilde W^*\widetilde W\right)\right|
 \le\opnorm{\widetilde W^*\widetilde W}
 =\opnorm{\widetilde W}^2.
\]
Take the supremum over product states.  The left-square proof is
identical.
\end{proof}

\section{The already-available thermal weighted card}
\label{sec:V89-thermal}

\begin{theorem}[Weighted thermal \(3+1\) square card]
\label{thm:V89-thermal-weighted}
Fix \(M<\infty\) and assume
\begin{equation}
 s_\alpha\Xi_k\le M
 \qquad(1\le k\le4).
\label{eq:V89-thermal}
\end{equation}
Then every payer placement \(\pi\) satisfies
\begin{align}
 \opnorm{\widetilde W_{r,\pi}}
 &\le C_{G,s,M}\mathfrak a_{31,r},
\label{eq:V89-thermal-weighted}\\
 \kappa_\otimes(
  \widetilde W_{r,\pi}^{\,*}\widetilde W_{r,\pi})
 &\le C_{G,s,M}\mathfrak a_{31,r}^{\,2},
\label{eq:V89-thermal-square}
\end{align}
and likewise for the left square.  Consequently the entire thermal
\(3+1\) source obeys
\begin{equation}
 \kappa_\otimes^{\rm abs}\!\left(
  \sum_r\widetilde W_{r,\pi}\right)
 \le
 C_{G,s,M}\widehat{\mathscr A}_{31}^{123|4}.
\label{eq:V89-thermal-global}
\end{equation}
\end{theorem}

\begin{proof}
Use the complete ternary-prefix numerator and response bounds
\(\texttt{eq:V66-prefix-unpaid}\) and
\(\texttt{eq:V66-prefix-paid}\) from frozen f2.  Use the assembled
joining-covariance bounds
\(\texttt{eq:V66-unpaid-crossing}\) and
\(\texttt{eq:V66-paid-crossing}\); the latter are applied to
\(J_r\), not to the three \(J_{i,r}\) separately.

If the prefix pays, multiply its paid bound by the unpaid joining
bound.  If the joining covariance pays, use the reverse choice.  If
the singleton prefix or complete suffix pays, use the bankwise
one-payer debit \(\texttt{lem:V67-bank-payer}\) with the two unpaid
factors.  In all four cases, after division by
\(\Xi_1\Xi_2\Xi_3\Xi_4\), the term indexed by \(T,i\) is bounded by
\[
 C_{G,s}\,(s_\alpha\Xi_i)
 \left(\prod_{uv\in E(T)}\theta_{uv}^{<r}\right)
 \theta_{i4,r};
\]
a harmless additional fixed power of \(s_\alpha\) is absorbed in
\(C_{G,s}\).  Hypothesis \eqref{eq:V89-thermal} removes the displayed
row demand.  Summation over the fixed nine pairs \((T,i)\) proves
\eqref{eq:V89-thermal-weighted}.

\Cref{cor:V89-norm-controls-square} proves
\eqref{eq:V89-thermal-square}.  Finally, subadditivity of absolute
product-state capacity followed by
\Cref{lem:V89-correct-summation} gives
\eqref{eq:V89-thermal-global}.
\end{proof}

\begin{assessment}[What is new relative to frozen V66]
\label{ass:V89-thermal-new}
Frozen \(\texttt{thm:V66-thermal-31}\) already closes the assembled
thermal DRQC source.  The new information in
\Cref{thm:V89-thermal-weighted} is only the explicit local weight
\(\mathfrak a_{31,r}\), its oriented complete-word square
consequence, and the correction of the target used in V88.
\end{assessment}

\section{A normalization-correct polynomial floor}
\label{sec:V89-floor}

\begin{lemma}[Polynomial floor for the corrected \(3+1\) atlas]
\label{lem:V89-correct-floor}
Let \(r\) be a selected nonzero four-row joining coordinate and put
\(M_*:=\max_k\Xi_k\).  Then
\begin{equation}
 \boxed{\;
 \widehat{\mathscr A}_{31}^{123|4}
 \ge {c_{31}\over M_*^6},
 \qquad
 \bigl(\widehat{\mathscr A}_{31}^{123|4}\bigr)^2
 \ge {c_{31}^2\over M_*^{12}}.\;}
\label{eq:V89-correct-floor}
\end{equation}
\end{lemma}

\begin{proof}
The selected-coordinate argument of
\Cref{lem:V86-edge-floor} gives
\(\theta_{uv}\ge c_0M_*^{-2}\) for every pair \(u<v\).
Therefore
\[
 \widehat{\mathscr T}_{123}
 \ge3c_0^2M_*^{-4},
 \qquad
 \widehat{\mathscr B}_{123|4}
 \ge3c_0M_*^{-2}.
\]
Multiply and take \(c_{31}=9c_0^3\).
\end{proof}

\begin{remark}[The exponent choice is again \(k=24\)]
\label{rem:V89-k24}
The corrected atlas, unlike the V88 raw-stock expression, has the
same worst polynomial floor \(M_*^{-6}\) as the \(2+1+1\) atlas.
Its square therefore requires \(M_*^{-12}\).  Since
\Cref{thm:V87-diffuse-bank-power} supplies \(M_*^{-k/2}\), the
normalization-correct choice is \(k=24\), not a smaller exponent
suggested by the superseded target.
\end{remark}

\section{An incidence-neutral diffuse-bank terminal}
\label{sec:V89-diffuse-terminal}

\begin{definition}[Central diffuse-bank insertion]
\label{def:V89-central-insertion}
A resolved normalized word \(\widetilde W\) has a central
\((a,D)\)-insertion with constant \(C_*\) if:
\begin{enumerate}[label=\textup{(\roman*)},leftmargin=4em]
\item \(D\) is a complete source-occurring coordinate bank, disjoint
      from the selected and payer coordinates;
\item row \(a\) is active on \(D\), and the complete heat factor on
      \(D\) is \(\mathcal M_D\);
\item the resolution has at most \(N_*\) terms, independent of
      support and representations; and
\item in the square orientation selected by output duality, every
      resolved term \(W_\nu\) satisfies
\begin{equation}
 \kappa_\otimes(W_\nu^*W_\nu)
 \le C_*\kappa_{a|a^c}(\mathcal M_D^2),
\label{eq:V89-central-insertion}
\end{equation}
or the corresponding left-square inequality.
\end{enumerate}
\end{definition}

\begin{theorem}[Diffuse-bank terminal for any corrected quartic
packet]
\label{thm:V89-diffuse-terminal}
Suppose a sequence of selected four-row words has
\(M_*=\max_k\Xi_k\to\infty\), a central \((a,D)\)-insertion with
uniform \(C_*,N_*\), and
\begin{equation}
 \Xi_a=M_*,
 \qquad
 S_a(D)\ge\sigma\Xi_a,
 \qquad
 \max_{e\in D}{a_{a,e}\over\Xi_a}\longrightarrow0.
\label{eq:V89-diffuse-terminal-hyp}
\end{equation}
Then, eventually,
\begin{equation}
 \boxed{\;
 \kappa_\otimes(\widetilde W^*\widetilde W)
 \le
 C_{G,s,\sigma,C_*,N_*}
 \bigl(\widehat{\mathscr A}_{31}^{123|4}\bigr)^2,\;}
\label{eq:V89-diffuse-terminal-result}
\end{equation}
with the analogous conclusion in the other orientation.
\end{theorem}

\begin{proof}
Apply \Cref{thm:V87-diffuse-bank-power} with \(k=24\).
The last hypothesis in \eqref{eq:V89-diffuse-terminal-hyp} eventually
implies its atom-size condition, and hence
\[
 \kappa_{a|a^c}(\mathcal M_D^2)
 \le C_{G,s,\sigma}M_*^{-12}.
\]
Insert this in \eqref{eq:V89-central-insertion}.  Recombination of
the fixed \(N_*\) terms costs at most a constant depending on \(N_*\)
by the finite-recombination lemma
\Cref{lem:V58-finite-recombination}.  Finally
\Cref{lem:V89-correct-floor} gives
\[
 M_*^{-12}
 \le c_{31}^{-2}
 \bigl(\widehat{\mathscr A}_{31}^{123|4}\bigr)^2.
\]
This proves the claim.  The bipartite capacity across \(a|a^c\)
enlarges the fully row-product test class, exactly as in
\Cref{thm:V87-no-total-depletion}.
\end{proof}

\begin{firewall}[What the terminal does not supply]
\label{fire:V89-terminal-scope}
\Cref{thm:V89-diffuse-terminal} closes a counterpacket only after
the exact \(3+1\) word has exposed a source-occurring owned diffuse
bank with the central insertion property.  It does not prove that
every hot \(3+1\) counterpacket has such a bank.  The missing step is
now a bank-extraction theorem compatible with the two-block cut
\(123|4\), not a new local covariance estimate and not an estimate
of nine independent Leibniz squares.
\end{firewall}

\section{V89 result ledger and V90 acquisition order}
\label{sec:V89-status}

\begin{theorem}[Third-series V89 result ledger]
\label{thm:V89-results}
Relative to V1--V88 and the frozen archives, V89 proves:
\begin{enumerate}[label=\textup{(V89.\arabic*)},leftmargin=6.3em]
\item the exact excess-degree error in the V88 \(3+1\) target;
\item the corrected literal nine-tree atlas and corrected fixed and
      weighted cards;
\item support-uniform summability of the corrected local weights;
\item the exact positive \(3\times3\) Gram organization of the
      assembled covariance square;
\item the coordinate-weighted thermal square card;
\item the \(M_*^{-6}\) floor for the corrected atlas and the resulting
      \(k=24\) exponent choice; and
\item an incidence-neutral diffuse-bank terminal for the hot
      counterpacket.
\end{enumerate}
\end{theorem}

\begin{proof}
These are
\Cref{lem:V89-degree-correction,
def:V89-correct-cards,
lem:V89-correct-summation,
lem:V89-nine-Gram,
thm:V89-thermal-weighted,
lem:V89-correct-floor,
thm:V89-diffuse-terminal}.
\end{proof}

\begin{programme}[V90: audit and hot-bank extraction]
\label{prog:V89-V90}
\begin{enumerate}[leftmargin=3em]
\item Perform the compulsory read-only SM/NS audit and transfer only
      type-correct proof-order information.
\item Assume failure of the corrected
      \eqref{eq:V89-correct-FCQ31}; remove the thermal branch using
      \Cref{thm:V89-thermal-weighted}.
\item Apply the complete two-block covariance minimum across
      \(123|4\) before selecting a bridge label.
\item Split a maximal demanded row into marked concentration, a
      cut-crossing bank, and cut-internal/private mass.  Use frozen
      V77 immediately on the cut-crossing dominant alternative.
\item On the remaining internal/private alternative, prove exact
      source occurrence and either a cofinal gap or an owned
      comparable diffuse bank.
\item Invoke \Cref{thm:V89-diffuse-terminal} only after its central
      insertion hypotheses have been verified for the same resolved
      source word.
\end{enumerate}
\end{programme}

\begin{firewall}[Claim boundary after third-series V89]
\label{fire:V89-claim-boundary}
V89 corrects the \(3+1\) normalization, records the exact square
algebra, and closes the weighted thermal branch.  It does not yet
prove hot-bank extraction, the corrected uniform \(FCQ_{31}\), the
weighted \(2+1+1\) card, the global discrete source, globally
aggregated \([4]\), \([3,3]\), higher MTA, WUFC, CPM, cutoff removal,
reflection positivity, Osterwalder--Schrader reconstruction,
construction of the continuum Yang--Mills measure, or a uniform
positive continuum spectral gap.  The full Yang--Mills mass gap has
not been proved.
\end{firewall}

\paragraph{Parent-version record.}
V89 was formed from
\texttt{Renewal\_Yang\_Mills\_Mass\_Gap\_Research\_Notes\_V88.tex}.
The V88 parent had \(42\,506\) logical source lines,
\(1\,509\,729\) bytes, and SHA--256
\[
 \texttt{83e59835cb6d67d094aaea84ad430f906ab97c7ca33f2278760d1e9be3bef755}.
\]
The next compulsory companion audit is V90.

\clearpage
\part{Third-series V90: companion audit and reduction of the hot
\texorpdfstring{\(3+1\)}{3+1} packet to one internal concentrated
triple block}
\label{part:V90}

\section{Compulsory SM/NS companion audit}
\label{sec:V90-audit}

The fifth-version protocol requires a read-only audit at V90.  The
newest active companion files are
\begin{align}
 {\cal C}_{\rm SM}
 &:=
 \texttt{Renewal\_SM\_Einstein\_Continuum\_Research\_Notes\_V68.tex},
\label{eq:V90-SM-file}\\
 {\cal C}_{\rm NS}
 &:=
 \texttt{Renewal\_NS\_Stability\_Research\_Notes\_V31.tex}.
\label{eq:V90-NS-file}
\end{align}
Their exact identities at this audit are:
\begin{center}
\begin{tabular}{lrrl}
\toprule
file & lines & bytes & SHA--256\\
\midrule
SM V68 & \(25\,190\) & \(932\,109\) &
\texttt{ffc71685c66ab743ac405db0134afe79b}\\
&&&\texttt{bda4ec00bc1c856a1baae23edc04ee9}\\
NS V31 & \(23\,052\) & \(887\,537\) &
\texttt{f1121b62238ad5491bb7cf03635ea993}\\
&&&\texttt{4942edf573ed8faaa9ee79e1d38a6696}\\
\bottomrule
\end{tabular}
\end{center}

\begin{assessment}[Type-correct content of the SM audit]
\label{ass:V90-SM-audit}
SM V68 withdraws an analytic use of its V65 theorem after detecting
that the theorem was proved under a finite-range master parent but had
been inserted after a nonlocal determinant modulus was promoted into
the parent law.  Its repair restarts every collision under one common
master parent and keeps all determinant descendants external.

The type-correct YM consequence is a proof-order restriction:
a source-occurrence theorem, a heat-bank capacity, and an all-payer
factorization may be combined only inside the same exact
first-connectivity resolution.  A bank exposed after changing the
moment-partition role cannot be inserted into a factorization proved
for the original role unless an exact tower or role-change identity is
proved.  No SM determinant, Witten constant, or Higgs estimate is
imported.
\end{assessment}

\begin{assessment}[Type-correct content of the NS audit]
\label{ass:V90-NS-audit}
NS V31 is unchanged since the V85 audit.  It proves that flat dyadic
probability marks encode the desired unmarked critical quantity only
after multiplication by their exact first-moment weight.  A uniform
marked estimate does not remove that weight, and a persistent time
interval may not be charged twice.

The YM analogue is already the distinction in
\Cref{fire:V88-fixed-vs-global}: a uniform fixed-\(r\) card does not
imply the coordinate-weighted card.  The literal
\(\mathfrak a_{31,r}\) must remain attached to the same source word
through the bank extraction.  No Navier--Stokes coarea, heat-formation,
or helicity estimate is imported.
\end{assessment}

\begin{firewall}[Companion audit boundary]
\label{fire:V90-audit-boundary}
The companion programmes supply only falsification tests for proof
organization: retain one parent/source resolution and retain the
literal first-moment weight.  Every estimate below is proved from the
YM source identity, archived YM covariance/MTA interfaces, and the
V87 heat-bank theorem.
\end{firewall}

\section{Finite tree-branch resolution of a failed square card}
\label{sec:V90-branch-resolution}

For \(T\in\mathfrak T(\{1,2,3\})\) and
\(i\in\{1,2,3\}\), write
\begin{equation}
 \mathfrak a_{T,i,r}
 :=
 \left(\prod_{uv\in E(T)}\theta_{uv}^{<r}\right)
 \theta_{i4,r}.
\label{eq:V90-branch-stock}
\end{equation}
Thus
\(\mathfrak a_{31,r}=\sum_{T,i}\mathfrak a_{T,i,r}\).
Use the fixed complete ternary-tree resolution underlying frozen
V66 and then the three-branch decomposition
\eqref{eq:V89-joining-branches}.  Denote the corresponding normalized
word by \(\widetilde W_{T,i,r,\pi}\).

\begin{lemma}[Nine-branch square recombination]
\label{lem:V90-nine-recombination}
For either output-duality orientation,
\begin{equation}
 \kappa_\otimes(
  \widetilde W_{r,\pi}^{\,*}\widetilde W_{r,\pi})
 \le
 9\sum_{T,i}
 \kappa_\otimes(
  \widetilde W_{T,i,r,\pi}^{\,*}
  \widetilde W_{T,i,r,\pi}),
\label{eq:V90-nine-recombination}
\end{equation}
after enlarging the fixed constant to absorb the finite internal
ternary resolution.  Hence if every branch obeys
\begin{equation}
 \kappa_\otimes(
  \widetilde W_{T,i,r,\pi}^{\,*}
  \widetilde W_{T,i,r,\pi})
 \le C\mathfrak a_{T,i,r}^{\,2},
\label{eq:V90-branch-card}
\end{equation}
then the corrected card \eqref{eq:V89-correct-FCQ31} holds.
\end{lemma}

\begin{proof}
For operators \(A_1,\ldots,A_N\),
\[
 \left(\sum_{\nu=1}^NA_\nu\right)^*
 \left(\sum_{\nu=1}^NA_\nu\right)
 \preccurlyeq
 N\sum_{\nu=1}^NA_\nu^*A_\nu,
\]
because
\[
 N\sum_\nu A_\nu^*A_\nu-
 \left(\sum_\nu A_\nu\right)^*
 \left(\sum_\nu A_\nu\right)
 =
 \sum_{\mu<\nu}(A_\mu-A_\nu)^*(A_\mu-A_\nu).
\]
Apply this first to the three joining branches and then to the fixed
ternary-tree resolution.  Positive product-state capacity is
monotone and subadditive.  Finally
\[
 \sum_{T,i}\mathfrak a_{T,i,r}^{\,2}
 \le
 \left(\sum_{T,i}\mathfrak a_{T,i,r}\right)^2
 \le
 \bigl(\widehat{\mathscr A}_{31}^{123|4}\bigr)^2
\]
by \Cref{lem:V89-correct-summation}.  The left-square identity follows
by applying the same argument to adjoints.
\end{proof}

\begin{firewall}[A finite branch sum is not repeated bank payment]
\label{fire:V90-no-repeated-bank}
The proof uses one common payer placement and one common
source-occurring bank in each resolved word.  It bounds the Gram cross
terms by a fixed finite operator inequality; it does not assign a new
resolvent or a new copy of the bank to each scalar edge in the
positive tree polynomial.
\end{firewall}

\section{The unique residual row demand}
\label{sec:V90-one-demand}

\begin{lemma}[One-demand ordinary bound]
\label{lem:V90-one-demand}
For every \(T,i,r,\pi\),
\begin{equation}
 \boxed{\;
 \opnorm{\widetilde W_{T,i,r,\pi}}
 \le
 C_{G,s}\Xi_i\,\mathfrak a_{T,i,r}.\;}
\label{eq:V90-one-demand}
\end{equation}
The sole unnormalized row demand is the endpoint \(i\) of the joining
bridge \(i4\).
\end{lemma}

\begin{proof}
Frozen
\(\texttt{eq:V66-prefix-unpaid}\)--%
\(\texttt{eq:V66-prefix-paid}\)
give the complete ternary prefix the normalized tree factor
\(\prod_{uv\in E(T)}\theta_{uv}^{<r}\), with the sole payer removing
one heat power when it lies there.  Frozen
\(\texttt{eq:V66-unpaid-crossing}\)--%
\(\texttt{eq:V66-paid-crossing}\)
give the assembled joining covariance its actual bridge
\(h_{i4,r}=\Xi_i\Xi_4\theta_{i4,r}\).
The singleton and suffix payer alternatives use
\(\texttt{lem:V67-bank-payer}\).  In every location division by the
four input masses cancels all row masses except the second occurrence
of \(\Xi_i\).  Fixed powers of \(s_\alpha\) are absorbed into
\(C_{G,s}\), proving \eqref{eq:V90-one-demand}.
\end{proof}

\begin{corollary}[A failed branch is hot at its demanded row]
\label{cor:V90-demand-hot}
If \eqref{eq:V90-branch-card} fails along a sequence for one fixed
\((T,i,\pi)\), then, after passing to a subsequence,
\begin{equation}
 \Xi_i\longrightarrow\infty.
\label{eq:V90-demand-hot}
\end{equation}
\end{corollary}

\begin{proof}
If \(\Xi_i\) were bounded, then
\[
 \kappa_\otimes(\widetilde W^*\widetilde W)
 \le\opnorm{\widetilde W}^2
 \le C_{G,s}\Xi_i^2\mathfrak a_{T,i,r}^2
\]
would give \eqref{eq:V90-branch-card}.
\end{proof}

\section{A four-coordinate diffuse compiler for one demand}
\label{sec:V90-k4}

\begin{definition}[Branch-compatible demanded-row bank]
\label{def:V90-compatible-bank}
A bank \(D\) is compatible with the resolved branch
\(\widetilde W_{T,i,r,\pi}\) if it is exposed by one exact
prefix/suffix incidence resolution of that same word, is disjoint
from the selected and payer coordinates, and its complete heat
factor \(\mathcal M_D\) satisfies the central insertion
\begin{equation}
 \kappa_\otimes(
  \widetilde W_{T,i,r,\pi}^{\,*}
  \widetilde W_{T,i,r,\pi})
 \le
 C_{G,s}\Xi_i^2\mathfrak a_{T,i,r}^2
 \kappa_{i|i^c}(\mathcal M_D^2).
\label{eq:V90-branch-insertion}
\end{equation}
The other orientation is defined analogously.
\end{definition}

\begin{theorem}[Diffuse ownership closes one \(3+1\) branch]
\label{thm:V90-k4-terminal}
Suppose a compatible bank \(D\) satisfies, for fixed \(\sigma>0\),
\begin{equation}
 S_i(D)\ge\sigma\Xi_i,
 \qquad
 \max_{e\in D}{a_{i,e}\over\Xi_i}
 \longrightarrow0.
\label{eq:V90-k4-hyp}
\end{equation}
Then \eqref{eq:V90-branch-card} holds eventually.  Only four selected
coordinates of \(D\) are required.
\end{theorem}

\begin{proof}
Use \Cref{thm:V87-diffuse-bank-power} with \(k=4\):
\[
 \kappa_{i|i^c}(\mathcal M_D^2)
 \le C_{G,s,\sigma}\Xi_i^{-2}.
\]
Insert this in \eqref{eq:V90-branch-insertion}; the two inverse powers
cancel the square of the unique demand in
\eqref{eq:V90-one-demand}.  This gives
\eqref{eq:V90-branch-card}.
\end{proof}

\begin{lemma}[Diffuse full support produces a compatible bank]
\label{lem:V90-diffuse-compatible}
Let a failed branch satisfy \eqref{eq:V90-demand-hot} and suppose
\begin{equation}
 \max_e{a_{i,e}\over\Xi_i}\longrightarrow0.
\label{eq:V90-full-diffuse}
\end{equation}
Then, after an exact finite incidence refinement and passage to a
subsequence, there is a branch-compatible bank \(D\) satisfying
\eqref{eq:V90-k4-hyp}.
\end{lemma}

\begin{proof}
Delete the selected coordinate, the payer coordinate when distinct,
and the fixed finite marked packet of the branch resolution.
By \eqref{eq:V90-full-diffuse}, the deleted mass is
\(o(\Xi_i)\), so the remaining support owns
\((1-o(1))\Xi_i\).

Apply the universal prefix/suffix incidence theorem
\Cref{thm:V84-universal-incidence} to this entire unmarked bank in
the exact \(3+1\) word.  Its finite Boolean endpoint resolution
places every coordinate in an actual prefix or suffix moment role;
one cell retains a fixed positive fraction of the row-\(i\) mass.
Call that cell \(D\).  The all-payer central factorization used in
\Cref{thm:V87-no-total-depletion} applies because \(D\) is disjoint
from the selected and payer coordinates and has not changed its
moment-partition role.  Keep the already proved numerator factors
in \eqref{eq:V90-one-demand}; inserting the complete \(D\)-heat
factor gives \eqref{eq:V90-branch-insertion}.  The maximum in
\eqref{eq:V90-k4-hyp} is no larger than the maximum in
\eqref{eq:V90-full-diffuse}.
\end{proof}

\begin{corollary}[Every persistent failure has a cofinal atom]
\label{cor:V90-cofinal-atom}
If a corrected fixed \(3+1\) square card fails, then after fixed
finite branch selection there exist \(\sigma>0\), a demanded row
\(i\in\{1,2,3\}\), and coordinates \(e_n\) such that
\begin{equation}
 a_{i,e_n}\ge\sigma\Xi_i
\label{eq:V90-cofinal-atom}
\end{equation}
along the counterpacket.
\end{corollary}

\begin{proof}
Use \Cref{lem:V90-nine-recombination} to select a failed
\((T,i,\pi)\) branch.  If no such \(\sigma,e_n\) existed, the maximum
in \eqref{eq:V90-full-diffuse} would tend to zero after passage to a
subsequence.  Then
\Cref{lem:V90-diffuse-compatible,thm:V90-k4-terminal} would close
the branch, a contradiction.
\end{proof}

\section{Classification of the cofinal atom}
\label{sec:V90-atom-classification}

At \(e=e_n\), compare local Casimir weights to \(a_{i,e}\).
The finite row set permits passage to a subsequence on which each
ratio \(a_{k,e}/a_{i,e}\) tends in
\([0,\infty]\).

\begin{lemma}[Unbalanced and cut-crossing atoms are already closed]
\label{lem:V90-known-atom-branches}
For the cofinal atom of \Cref{cor:V90-cofinal-atom}:
\begin{enumerate}[label=\textup{(\roman*)},leftmargin=4em]
\item if the local highest-weight tuple is eventually V32-unbalanced,
      the branch is closed by the unbalanced local compilers;
\item if it is not unbalanced and
      \(a_{4,e}\ge c\,a_{i,e}\) for some fixed \(c>0\), the branch is
      closed by the complete two-block bank theorem across
      \(123|4\).
\end{enumerate}
\end{lemma}

\begin{proof}
Item \textup{(i)} is precisely the scope assembled in
\Cref{cor:V85-heavy-composite-closed}: the V34, V35, V38, and V47
unbalanced compilers cover joining, payer, separator, and coincidence
roles after one exact finite resolution.

For \textup{(ii)}, the one-coordinate bank \(B=\{e\}\) is
cut-exposed from row \(i\) to row \(4\), and
\[
 X_{i\to4}(B)=a_{i,e}\ge\sigma\Xi_i.
\]
Thus row \(i\) is \(\sigma\)-bank-dominant across the two-block cut.
Frozen \(\texttt{cor:V77-bank-pays-row}\), inserted by
\(\texttt{prop:V77-tensor-insertion}\), supplies exactly the missing
factor \(\Xi_i^{-1}\) in the word and its adjoint.  The covariance is
kept assembled across \(123|4\), and the resulting branch obeys
\eqref{eq:V90-branch-card}.
\end{proof}

\begin{theorem}[Sharp residual form of a hot \(3+1\) counterpacket]
\label{thm:V90-sharp-residual}
If the corrected \(FCQ_{31}\) card fails, then after subsequences and
fixed finite choices there are:
\[
 T\in\mathfrak T(\{1,2,3\}),\quad
 i\in\{1,2,3\},\quad
 c\in\{1,2,3\}\setminus\{i\},\quad
 \sigma,\lambda,\Lambda>0,
\]
and a cofinal coordinate \(e_n\) such that
\begin{align}
 \Xi_i&\longrightarrow\infty,
\label{eq:V90-residual-hot}\\
 a_{i,e_n}&\ge\sigma\Xi_i,
\label{eq:V90-residual-own}\\
 \lambda a_{i,e_n}
 \le a_{c,e_n}
 &\le\Lambda a_{i,e_n},
\label{eq:V90-residual-internal}\\
 {a_{4,e_n}\over a_{i,e_n}}
 &\longrightarrow0,
\label{eq:V90-residual-no-bridge}
\end{align}
and the local tuple is not V32-unbalanced.  Every diffuse,
unbalanced, or cut-crossing alternative has been removed.
\end{theorem}

\begin{proof}
\Cref{lem:V90-nine-recombination,cor:V90-demand-hot,
cor:V90-cofinal-atom} give \(T,i,e_n\) and the first two
properties.  Remove the unbalanced branch by
\Cref{lem:V90-known-atom-branches}.  In the remaining local
highest-weight geometry, no row weight can dominate
\(a_{i,e_n}\) by an unbounded factor; otherwise the tuple becomes
unbalanced after relabelling the dominant row.  Hence all finite
nonzero limiting ratios are bounded above.

If row \(4\) has a positive limiting ratio, item \textup{(ii)} of
\Cref{lem:V90-known-atom-branches} closes the branch.  Thus
\eqref{eq:V90-residual-no-bridge} holds.  Since the tuple is not
unbalanced while row \(i\) is cofinally large and row \(4\) is
negligible, at least one of the other two triple rows must retain a
positive limiting ratio to row \(i\).  Select it as \(c\) and pass
once more to obtain fixed \(\lambda,\Lambda\), proving
\eqref{eq:V90-residual-internal}.
\end{proof}

\begin{assessment}[The exact remaining local object]
\label{ass:V90-remaining-object}
The obstruction is no longer a three-branch covariance square, a
support-size sum, or an arbitrary hot row.  It is one coordinate
which owns a fixed fraction of the demanded bridge endpoint mass,
has a comparable partner inside the already connected triple, has
negligible row-\(4\) mass, and lies in a locally balanced/protected
output sector.  The prefix ternary carrier and the joining covariance
share this high internal representation through their common row
labels, but the bridge \(i4\) itself is weak there.

The next estimate must therefore be a \emph{rooted triple-block
partial-transpose card}: retain the protected coupling of rows
\(i,c\) inside the complete three-row output, and show that the row-4
difference or the earlier connected-prefix tree supplies the missing
\(\Xi_i^{-1}\) without separating the protected multiplicity space.
\end{assessment}

\section{V90 result ledger and V91 acquisition order}
\label{sec:V90-status}

\begin{theorem}[Third-series V90 result ledger]
\label{thm:V90-results}
Relative to V1--V89, V90 proves:
\begin{enumerate}[label=\textup{(V90.\arabic*)},leftmargin=6.3em]
\item the compulsory SM V68/NS V31 audit with exact file identities
      and no cross-program estimate import;
\item a fixed finite Gram recombination reducing \(FCQ_{31}\) to
      nine normalized tree branches;
\item the exact single demanded row in each branch;
\item a \(k=4\) diffuse-bank compiler preserving the literal local
      tree weight;
\item source-compatible bank extraction on every fully diffuse
      demanded-row support;
\item necessity of a cofinal demanded-row atom on any persistent
      failure;
\item closure of unbalanced and cut-crossing cofinal atoms; and
\item reduction of the entire hot failure to the internally balanced
      concentrated triple block
      \eqref{eq:V90-residual-hot}--%
      \eqref{eq:V90-residual-no-bridge}.
\end{enumerate}
\end{theorem}

\begin{proof}
These are
\Cref{ass:V90-SM-audit,
ass:V90-NS-audit,
lem:V90-nine-recombination,
lem:V90-one-demand,
thm:V90-k4-terminal,
lem:V90-diffuse-compatible,
cor:V90-cofinal-atom,
lem:V90-known-atom-branches,
thm:V90-sharp-residual}.
\end{proof}

\begin{programme}[V91: rooted triple-block response]
\label{prog:V90-V91}
\begin{enumerate}[leftmargin=3em]
\item Fix the residual rows \(i,c\) and cofinal coordinate \(e\) of
      \Cref{thm:V90-sharp-residual}; do not average over bridge labels.
\item Decompose the complete local module at \(e\) by first coupling
      \(V_i\otimes V_c\), retaining all multiplicities and the third
      triple row as part of the complementary module.
\item Apply the V83 positive-block marginal sandwich only after this
      recoupling.  Determine whether partial transpose on the demanded
      row gives \(d_{V_i}^{-1}\), \(d_{V_i}^{-1/2}\), or no gain on
      the protected low-output block.
\item Test the sharp result on the \(SU(2)\) channel
      \(V_j\otimes V_j\supset V_0\), where internal cancellation is
      maximal and row \(4\) is trivial or fixed.
\item If one local coordinate cannot supply the full inverse
      \(\Xi_i\), retain the earlier triple-prefix tree occurrence and
      seek a two-occurrence Schatten sandwich.  Do not manufacture a
      second copy of the joining bridge.
\item Keep the literal weight
      \(\mathfrak a_{T,i,r}\) attached throughout, as required by the
      NS audit and the global summation gate.
\end{enumerate}
\end{programme}

\begin{firewall}[Claim boundary after third-series V90]
\label{fire:V90-claim-boundary}
V90 performs the compulsory audit and reduces every possible failure
of the corrected \(3+1\) fixed-word card to one internally balanced
cofinal triple-block atom.  It does not prove the rooted
triple-block response, corrected uniform \(FCQ_{31}\), the weighted
global \(3+1\) or \(2+1+1\) cards, the global discrete source,
globally aggregated \([4]\), \([3,3]\), higher MTA, WUFC, CPM,
cutoff removal, reflection positivity, Osterwalder--Schrader
reconstruction, construction of the continuum Yang--Mills measure,
or a uniform positive continuum spectral gap.  The full Yang--Mills
mass gap has not been proved.
\end{firewall}

\paragraph{Parent-version record.}
V90 was formed from
\texttt{Renewal\_Yang\_Mills\_Mass\_Gap\_Research\_Notes\_V89.tex}.
The V89 parent had \(43\,079\) logical source lines,
\(1\,528\,815\) bytes, and SHA--256
\[
 \texttt{4c566035a65ba9582c7dbe84b68682ea68277e664e997cb5cee09fa1d24b9c04}.
\]
The next compulsory companion audit is V95.

\clearpage
\part{Third-series V91: sharp singlet obstruction and
mass-ascent rerouting of the concentrated
\texorpdfstring{\(3+1\)}{3+1} residual}
\label{part:V91}

\section{The proposed one-coordinate gain is already falsified}
\label{sec:V91-purpose}

V90 reduced a failed corrected \(3+1\) card to a cofinal coordinate
where the demanded triple row \(i\) has a comparable partner \(c\)
inside the triple and row \(4\) is locally negligible.  Its V91
programme asked whether recoupling that one local block could supply
the full inverse demanded-row mass.

The frozen and active record already contains the decisive sharp
test.  \Cref{prop:V48-singlet} computes the exact product-state
capacity \(1/d_j\) of the invariant channel in
\(V_j\otimes V_j\).  The complete heat multiplier has eigenvalue one
on that channel, even after squaring.  Thus one balanced local
coordinate supplies only one inverse representation dimension,
which is one half Casimir power for \(SU(2)\).  The two inverse
Casimir powers required in the square insertion
\eqref{eq:V90-branch-insertion} are impossible from that coordinate
alone.

\begin{assessment}[Nonduplication boundary]
\label{ass:V91-nonduplication}
V91 does not repeat the V48 projection-capacity proof.  It proves the
new implication for the complete heat square appearing in the V90
\(3+1\) residual, shows sharpness of the V83 half-power, and replaces
the failed one-coordinate route by a finite total-mass ascent.  The
ascent either finds a larger diffuse row, which pays the original
demand, or ends at a uniformly strong normalized internal edge.
\end{assessment}

\section{The complete heat square retains the singlet}
\label{sec:V91-singlet}

For \(G=SU(2)\), let \(V_j\) be the spin-\(j\) module,
\(d_j=2j+1\), and
\[
 a_j:=1+C_2(j)=1+j(j+1).
\]
On \(V_j\otimes V_j\), write
\[
 M_j(s):=
 \sum_{\ell=0}^{2j}e^{-s\ell(\ell+1)}P_{j,\ell}.
\]

\begin{theorem}[Sharp one-coordinate heat-square obstruction]
\label{thm:V91-heat-square-obstruction}
For every \(s>0\),
\begin{equation}
 M_j(s)^2
 =
 \sum_{\ell=0}^{2j}e^{-2s\ell(\ell+1)}P_{j,\ell}
 \succcurlyeq P_{j,0},
\label{eq:V91-heat-square-dominates}
\end{equation}
and hence
\begin{equation}
 \boxed{\;
 \kappa_{\rm prod}(M_j(s)^2)
 \ge {1\over d_j}
 \asymp a_j^{-1/2}.\;}
\label{eq:V91-heat-square-lower}
\end{equation}
In particular, there is no \(j\)-independent \(C\) such that
\begin{equation}
 \kappa_{\rm prod}(M_j(s)^2)
 \le Ca_j^{-1},
 \quad\text{or}\quad
 \kappa_{\rm prod}(M_j(s)^2)
 \le Ca_j^{-2}.
\label{eq:V91-false-powers}
\end{equation}
\end{theorem}

\begin{proof}
The Clebsch--Gordan projections are mutually orthogonal and the
singlet eigenvalue is \(e^{-sC_2(0)}=1\).  This proves
\eqref{eq:V91-heat-square-dominates}.  Positive product-state
capacity is monotone, and \Cref{prop:V48-singlet} gives
\(\kappa_{\rm prod}(P_{j,0})=d_j^{-1}\).  Finally
\(d_j\asymp j\) and \(a_j\asymp j^2\).  Multiplying either proposed
bound in \eqref{eq:V91-false-powers} by \(d_j\) gives a right side
tending to zero, contradicting \eqref{eq:V91-heat-square-lower}.
\end{proof}

\begin{lemma}[Exact partial-transpose sharpness on the singlet]
\label{lem:V91-singlet-PT}
Up to local unitaries determined by the self-duality of \(V_j\),
\begin{equation}
 P_{j,0}^{\Gamma_2}={1\over d_j}F,
\label{eq:V91-singlet-PT}
\end{equation}
where \(F(u\otimes v)=v\otimes u\) is the flip.  Consequently
\begin{equation}
 \|P_{j,0}^{\Gamma_2}\|={1\over d_j}.
\label{eq:V91-singlet-PT-norm}
\end{equation}
Thus the \(d_j^{-1}\) scale in the V83/V87 local
partial-transpose card is sharp.
\end{lemma}

\begin{proof}
After a local unitary change of basis, the normalized invariant
vector is
\[
 \Omega_j={1\over\sqrt{d_j}}\sum_{m=1}^{d_j}e_m\otimes e_m.
\]
Then
\[
 (|\Omega_j\rangle\langle\Omega_j|)^{\Gamma_2}
 ={1\over d_j}\sum_{m,n}
 |e_m\otimes e_n\rangle\langle e_n\otimes e_m|
 ={F\over d_j}.
\]
The flip is unitary and self-adjoint, so its norm is one.  Local
unitaries do not change the norm.
\end{proof}

\begin{corollary}[The V90 single-atom insertion cannot close by
heat alone]
\label{cor:V91-no-single-atom-close}
There is no universal argument which closes
\eqref{eq:V90-branch-insertion} on the internally balanced cofinal
atom of \Cref{thm:V90-sharp-residual} by replacing every companion
factor by a contraction and estimating only that atom's complete heat
square.  This remains false if the third triple row and row \(4\) are
put in fixed trivial or bounded representations.
\end{corollary}

\begin{proof}
Take rows \(i,c\) in equal spin \(j\), put their coupled output in
the singlet channel, and tensor fixed unit vectors on the other
bounded factors.  The demanded-row mass is comparable to \(a_j\).
After the companions have been replaced by contractions, closing
\eqref{eq:V90-branch-insertion} would require the local heat square
to contribute \(O(a_j^{-2})\).  This contradicts
\eqref{eq:V91-heat-square-lower}.  Tensoring fixed one-dimensional
factors changes neither the product overlap nor its \(j\)-scaling.
\end{proof}

\begin{firewall}[Scope of the obstruction]
\label{fire:V91-obstruction-scope}
The theorem rules out an isolated local-heat proof, not the complete
\(3+1\) source card.  As in V48, an actual prefix response, a second
distinct heat coordinate, or a global tree rerouting may supply the
missing dimension.  The signed assembled covariance has not been
replaced by the singlet projector.
\end{firewall}

\section{A larger row may pay the original demand}
\label{sec:V91-larger-row}

\begin{lemma}[Diffuse closure is monotone under mass ascent]
\label{lem:V91-larger-diffuse}
Fix a failed branch whose unique ordinary demand is \(\Xi_i\).
Suppose another row \(a\) satisfies
\(\Xi_a\ge\Xi_i\) and has a branch-compatible bank \(D\) with
\begin{equation}
 S_a(D)\ge\sigma\Xi_a,
 \qquad
 \max_{e\in D}{a_{a,e}\over\Xi_a}\longrightarrow0.
\label{eq:V91-larger-diffuse}
\end{equation}
Then the branch closes.
\end{lemma}

\begin{proof}
Apply \Cref{thm:V87-diffuse-bank-power} to row \(a\) with \(k=4\):
\[
 \kappa_{a|a^c}(\mathcal M_D^2)
 \le C\Xi_a^{-2}
 \le C\Xi_i^{-2}.
\]
The central insertion proof of
\Cref{thm:V90-k4-terminal} does not require the transposed row to be
the row whose scalar demand is being paid; it requires only that the
bipartite test class enlarge the fully row-product class.  The last
display cancels the factor \(\Xi_i^2\) in
\eqref{eq:V90-branch-insertion}.
\end{proof}

\begin{lemma}[Cofinal comparable atoms give an edge or an ascent]
\label{lem:V91-edge-or-ascent}
Suppose
\begin{equation}
 a_{a,e}\ge\sigma\Xi_a,
 \qquad
 a_{c,e}\ge\lambda a_{a,e}.
\label{eq:V91-comparable-atom}
\end{equation}
For every fixed \(L>1\), exactly one of the following alternatives
holds:
\begin{enumerate}[label=\textup{(\alph*)},leftmargin=4em]
\item \(\Xi_c\le L\Xi_a\), in which case
\begin{equation}
 \boxed{\;\theta_{ac}\ge{\lambda\sigma^2\over L};\;}
\label{eq:V91-strong-edge}
\end{equation}
\item \(\Xi_c>L\Xi_a\), a strict total-mass ascent.
\end{enumerate}
\end{lemma}

\begin{proof}
The alternatives partition the possibilities.  In the first,
\[
 H_{ac}\ge a_{a,e}a_{c,e}
 \ge\lambda a_{a,e}^2
 \ge\lambda\sigma^2\Xi_a^2.
\]
Divide by
\(\Xi_a\Xi_c\le L\Xi_a^2\).
\end{proof}

\section{Finite ascent inside the triple}
\label{sec:V91-ascent}

Call a row \(a\) \emph{diffuse} if
\(\max_e a_{a,e}/\Xi_a\to0\), and \emph{concentrated} otherwise.
After passage to a subsequence, a concentrated row has a cofinal atom
with a fixed positive ownership fraction.

\begin{theorem}[Mass-ascent/strong-edge alternative]
\label{thm:V91-ascent-alternative}
Starting from the residual branch of
\Cref{thm:V90-sharp-residual}, fix \(L>1\).
After subsequences, one of the following occurs:
\begin{enumerate}[label=\textup{(\roman*)},leftmargin=4em]
\item a triple row \(a\) with \(\Xi_a\ge\Xi_i\) has a
      source-compatible diffuse owned bank, and the branch closes by
      \Cref{lem:V91-larger-diffuse};
\item a cofinal atom is V32-unbalanced or cut-crossing to row \(4\),
      and the branch closes by
      \Cref{lem:V90-known-atom-branches}; or
\item there are distinct triple rows \(a,c\), a cofinal atom \(e\),
      and fixed \(\sigma,\lambda,\Lambda,\eta>0\) such that
\begin{align}
 \Xi_a&\ge\Xi_i,
 &
 a_{a,e}&\ge\sigma\Xi_a,
\label{eq:V91-terminal-own}\\
 \lambda a_{a,e}
 \le a_{c,e}&\le\Lambda a_{a,e},
 &
 \theta_{ac}&\ge\eta,
\label{eq:V91-terminal-edge}\\
 a_{4,e}/a_{a,e}&\longrightarrow0,
\label{eq:V91-terminal-row4}
\end{align}
and the local tuple is not unbalanced.
\end{enumerate}
The ascent uses at most two strict steps.
\end{theorem}

\begin{proof}
Begin with \(a=i\) and the internal comparable partner \(c\) from
\Cref{thm:V90-sharp-residual}.  Apply
\Cref{lem:V91-edge-or-ascent}.  If the strong-edge alternative holds,
we are in \textup{(iii)}.

Otherwise \(\Xi_c>L\Xi_a\).  If row \(c\) is diffuse, delete the fixed
selected, payer, and marked coordinates; the remainder still owns
\((1-o(1))\Xi_c\).  The exact incidence argument of
\Cref{lem:V90-diffuse-compatible} produces a source-compatible bank,
and \textup{(i)} holds.

If row \(c\) is concentrated, select a cofinal atom.  The local
unbalanced and row-\(4\)-crossing cases are \textup{(ii)}.  Otherwise
the proof of \Cref{thm:V90-sharp-residual}, with \(c\) as root row,
selects a comparable partner inside the triple.  Apply
\Cref{lem:V91-edge-or-ascent} again.

Every strict step multiplies total mass by more than \(L\).  A row
cannot repeat in an ascent chain, since repetition would imply
\(\Xi_a>L^m\Xi_a\) for some \(m\ge1\).  There are only three triple
rows, so at most two strict steps occur.  If the process does not
exit through \textup{(i)} or \textup{(ii)}, it terminates in the
strong-edge alternative.  The fixed ratio bounds and the local
row-\(4\) and balance properties are retained by the same finite
subsequence selection as in \Cref{thm:V90-sharp-residual}.
\end{proof}

\begin{corollary}[Every persistent failure has a strong internal
normalized edge]
\label{cor:V91-strong-edge-residual}
If the corrected \(FCQ_{31}\) card fails, then alternative
\textup{(iii)} of \Cref{thm:V91-ascent-alternative} holds.
\end{corollary}

\begin{proof}
Alternatives \textup{(i)} and \textup{(ii)} close the selected failed
branch, contradicting its selection in
\Cref{lem:V90-nine-recombination}.
\end{proof}

\section{Improved atlas floor on the strong-edge residual}
\label{sec:V91-improved-floor}

\begin{lemma}[A strong internal edge saves two floor powers]
\label{lem:V91-improved-floor}
Let
\[
 M_*:=\max_{1\le k\le4}\Xi_k.
\]
On the residual of
\Cref{cor:V91-strong-edge-residual},
\begin{equation}
 \boxed{\;
 \widehat{\mathscr A}_{31}^{123|4}
 \ge {c_\eta\over M_*^4},
 \qquad
 \bigl(\widehat{\mathscr A}_{31}^{123|4}\bigr)^2
 \ge {c_\eta^2\over M_*^8}.\;}
\label{eq:V91-improved-floor}
\end{equation}
\end{lemma}

\begin{proof}
Let \(d\) be the third row in \(\{1,2,3\}\setminus\{a,c\}\).
By \Cref{lem:V86-edge-floor},
\[
 \theta_{ad},\theta_{cd}\ge c_0M_*^{-2},
 \qquad
 \widehat{\mathscr B}_{123|4}\ge3c_0M_*^{-2}.
\]
Since \(\theta_{ac}\ge\eta\), the two triple-tree terms containing
the edge \(ac\) give
\[
 \widehat{\mathscr T}_{123}
 \ge
 \theta_{ac}(\theta_{ad}+\theta_{cd})
 \ge2\eta c_0M_*^{-2}.
\]
Multiplication proves \eqref{eq:V91-improved-floor} with
\(c_\eta=6\eta c_0^2\).
\end{proof}

\begin{corollary}[Incidence-neutral fallback exponent]
\label{cor:V91-k16}
On the strong-edge residual, any source-compatible diffuse bank
owning a globally maximal row closes the complete square by the
incidence-neutral argument of \Cref{thm:V89-diffuse-terminal} with
\(k=16\) in place of \(k=24\).
\end{corollary}

\begin{proof}
\Cref{thm:V87-diffuse-bank-power} with \(k=16\) gives
\(CM_*^{-8}\).  Compare with
\eqref{eq:V91-improved-floor}.
\end{proof}

\begin{assessment}[The next irreducible gate]
\label{ass:V91-next-gate}
The remaining packet has both a cofinal internally balanced atom and
a strong normalized edge \(ac\) inside the triple.  The isolated
heat projector at that atom cannot provide the missing power, by
\Cref{thm:V91-heat-square-obstruction}.  Therefore the next proof
must retain one additional complete occurrence.

There are two genuinely different placements to classify:
\begin{enumerate}[label=\textup{(\alph*)},leftmargin=4em]
\item the strong atom lies before the first \(123|4\) join and is
      part of the connected triple-prefix carrier; or
\item it lies after that join in the complete suffix.
\end{enumerate}
In the prefix case, the correct candidate is a coupled
prefix-response/low-output sandwich.  In the suffix case, the strong
edge is global tree currency but not a literal prefix mark; this tests
whether the coordinate-weighted local card
\eqref{eq:V89-weighted-card} is too strong and global tree rerouting is
necessary.
\end{assessment}

\section{V91 result ledger and V92 acquisition order}
\label{sec:V91-status}

\begin{theorem}[Third-series V91 result ledger]
\label{thm:V91-results}
Relative to V1--V90, V91 proves:
\begin{enumerate}[label=\textup{(V91.\arabic*)},leftmargin=6.3em]
\item the complete \(SU(2)\) heat square dominates the singlet and
      has capacity at least \(a_j^{-1/2}\);
\item the exact \(d_j^{-1}\) partial-transpose norm of the singlet,
      showing sharpness of the V83 half-power;
\item impossibility of closing the V90 residual from one isolated
      balanced heat coordinate;
\item closure by any source-compatible diffuse row whose total mass
      dominates the demanded row;
\item the exact strong-edge-or-total-mass-ascent dichotomy;
\item termination of the triple-row ascent after at most two strict
      steps;
\item reduction of every persistent failure to a cofinal balanced
      atom carrying a uniformly strong internal normalized edge; and
\item improvement of the corrected atlas-square floor from
      \(M_*^{-12}\) to \(M_*^{-8}\) on that residual.
\end{enumerate}
\end{theorem}

\begin{proof}
These are
\Cref{thm:V91-heat-square-obstruction,
lem:V91-singlet-PT,
cor:V91-no-single-atom-close,
lem:V91-larger-diffuse,
lem:V91-edge-or-ascent,
thm:V91-ascent-alternative,
cor:V91-strong-edge-residual,
lem:V91-improved-floor}.
\end{proof}

\begin{programme}[V92: prefix versus suffix placement of the strong
atom]
\label{prog:V91-V92}
\begin{enumerate}[leftmargin=3em]
\item Apply the exact universal incidence resolution to the selected
      strong atom without changing its moment-partition role.
\item If it lies in the triple prefix, retain the complete prefix
      cumulant on both sides of the square and compute its compression
      to the internally protected \(ac\) output.  Compare with the V48
      rank lower bound before claiming a second dimension gain.
\item If it lies in the suffix, test
      \eqref{eq:V89-weighted-card} on an explicit \(SU(2)\) family with
      a high \(ac\) singlet after the joining coordinate and fixed
      nontrivial prefix/joining data.
\item If that family violates the local weighted card, withdraw that
      sufficient interface and return to a global rerouted tree bound;
      do not infer failure of global quartic MTA.
\item If the local card survives, identify the exact non-suffix factor
      supplying the additional \(d_j^{-1}\).
\item Keep the sole payer fixed and test both square orientations.
\end{enumerate}
\end{programme}

\begin{firewall}[Claim boundary after third-series V91]
\label{fire:V91-claim-boundary}
V91 proves a sharp single-coordinate obstruction and reduces every
persistent corrected \(3+1\) failure to a strong-edge internally
balanced atom, split next by prefix/suffix placement.  It does not
prove the required coupled response or decide the weighted local
card, corrected uniform \(FCQ_{31}\), the weighted global
\(3+1\) or \(2+1+1\) cards, the global discrete source, globally
aggregated \([4]\), \([3,3]\), higher MTA, WUFC, CPM, cutoff removal,
reflection positivity, Osterwalder--Schrader reconstruction,
construction of the continuum Yang--Mills measure, or a uniform
positive continuum spectral gap.  The full Yang--Mills mass gap has
not been proved.
\end{firewall}

\paragraph{Parent-version record.}
V91 was formed from
\texttt{Renewal\_Yang\_Mills\_Mass\_Gap\_Research\_Notes\_V90.tex}.
The V90 parent had \(43\,594\) logical source lines,
\(1\,547\,662\) bytes, and SHA--256
\[
 \texttt{47aeb1db15213e86152e1e5699d85ba108617714234004495135b7d9f807d266}.
\]
The next compulsory companion audit is V95.

\clearpage
\part{Third-series V92: a physical suffix-rescue counterpacket to
the literal local card and an ordered rerouted replacement}
\label{part:V92}

\section{Purpose}
\label{sec:V92-purpose}

V91 separated a strong internal atom according to whether it occurs
before or after the first \(123|4\) join.  If the atom lies in the
suffix, its normalized edge is present in the global quartic atlas
but absent from the literal prefix weight
\(\mathfrak a_{31,r}\).  The proposed weighted local estimate
\eqref{eq:V89-weighted-card} therefore asks the suffix heat capacity
to reproduce an edge which its right side has discarded.

This chapter gives an explicit \(SU(2)\) ordered-word family showing
that the literal local card is false on the suffix branch.  The family
does not violate the global \(3+1\) atlas: its strong suffix edge
enlarges the global target by exactly the missing scale.  We then
define a one-rescue ordered weight which retains that edge, prove its
support-uniform summability, and make it the next analytic interface.

\begin{assessment}[Nonduplication boundary]
\label{ass:V92-nonduplication}
The singlet capacity is not recomputed; V92 uses
\Cref{thm:V91-heat-square-obstruction}.  The new content is the
placement of that singlet after a fixed nonzero \(3+1\) source, the
comparison of local and global normalized tree stocks, and the
ordered rescue weight.
\end{assessment}

\section{A fixed nonzero \(3+1\) base word}
\label{sec:V92-base-word}

Let \(\phi\) be a real nonconstant normalized matrix coefficient of
the spin-\(\tfrac12\) representation of \(SU(2)\), evaluated under a
fixed nondegenerate Wilson heat law.  Then
\begin{equation}
 v_\phi:=\mathbb E(\phi^2)-(\mathbb E\phi)^2>0.
\label{eq:V92-positive-variance}
\end{equation}
Take three ordered coordinates \(p<q<r\):
\begin{enumerate}[label=\textup{(\roman*)},leftmargin=4em]
\item at \(p\), only rows \(1,2\) carry \(\phi\);
\item at \(q\), only rows \(1,3\) carry \(\phi\);
\item at \(r\), only rows \(1,4\) carry \(\phi\).
\end{enumerate}
Inactive row factors are identities and are omitted from the active
Casimir support.

\begin{lemma}[Existence of a fixed nonzero physical base block]
\label{lem:V92-base-nonzero}
There is a fixed physical Fourier-output resolution and a fixed payer
placement among \(p,q,r\) for which the \(123|4\) word with first
global join at \(r\) contains a nonzero block \(A_0\).  Its operator
norm and the product-state capacity of \(A_0^*A_0\) are positive
constants independent of every later representation.
\end{lemma}

\begin{proof}
The local covariance at \(p\) is
\(\kappa_2(\phi_1,\phi_2)=v_\phi\), and the covariance at \(q\) is
\(\kappa_2(\phi_1,\phi_3)=v_\phi\).  Their row graphs
\(\{12\}\) and \(\{13\}\) have connected join on
\(\{1,2,3\}\), so the coefficient-one linked-word formula
\Cref{prop:V7-linked-word} places their product in
\(K_{123}^{<r,\mathrm{conn}}\).  At \(r\), rows \(2,3\) are
identities, and the assembled covariance
\(\kappa_2(Z_{123,r},X_{4,r})\) contains
\(\kappa_2(\phi_1,\phi_4)=v_\phi\).  Thus the indicated source
coefficient is \(v_\phi^3>0\) before the fixed physical pinching.

Choose a nonzero finite output block in this coefficient.  If the
payer is placed at \(r\), choose a fixed nontrivial output, whose
one-payer scalar is strictly positive; alternatively place it at
\(p\) or \(q\).  All data are now fixed and finite-dimensional, so
the resulting block \(A_0\) is nonzero.  The positive operator
\(A_0^*A_0\) has positive expectation on some product vector because
product vectors span the finite tensor product.  This gives the two
fixed positive constants.
\end{proof}

\section{Appending a high protected suffix}
\label{sec:V92-high-suffix}

Add one coordinate \(e>r\).  At \(e\), rows \(1,2\) carry equal
spin \(j\), rows \(3,4\) are inactive, and the complete suffix moment
is the heat operator \(M_j(s)\) of V91.  Choose normalized rank-one
Fourier coefficients on rows \(1,2\) which approach the maximizing
product vectors for the singlet projection.  Let \(W_j\) be the
resulting complete source block with the fixed payer of
\Cref{lem:V92-base-nonzero}.

\begin{lemma}[Exact tensor factor and mass scales]
\label{lem:V92-tensor-scales}
For fixed constants \(c_k,C_k>0\),
\begin{align}
 W_j&=A_0\otimes M_j(s),
\label{eq:V92-tensor-word}\\
 c_1a_j\le\Xi_1,\Xi_2&\le C_1a_j,
 &
 c_2\le\Xi_3,\Xi_4&\le C_2,
\label{eq:V92-mass-scales}
\end{align}
where \(a_j=1+j(j+1)\).  If
\[
 \widetilde W_j:=
 {W_j\over\Xi_1\Xi_2\Xi_3\Xi_4},
\]
then
\begin{equation}
 \boxed{\;
 \kappa_\otimes(\widetilde W_j^*\widetilde W_j)
 \ge c\,a_j^{-4}d_j^{-1}
 \asymp a_j^{-9/2}.\;}
\label{eq:V92-word-lower}
\end{equation}
\end{lemma}

\begin{proof}
The coordinate variables are independent and \(e\) lies in the
complete suffix, so exact first-connectivity gives the tensor
factorization \eqref{eq:V92-tensor-word}.  Only the two spin-\(j\)
entries grow, proving \eqref{eq:V92-mass-scales}.

Choose a product state attaining a fixed positive expectation of
\(A_0^*A_0\), and independently choose the row-\(1|2\) product state
from \Cref{thm:V91-heat-square-obstruction}.  Product-state
expectations multiply over the coordinate tensor factors.  Hence
\[
 \kappa_\otimes(W_j^*W_j)
 \ge c_0\kappa_{\rm prod}(M_j(s)^2)
 \ge {c_0\over d_j}.
\]
Division by
\((\Xi_1\Xi_2\Xi_3\Xi_4)^2\asymp a_j^4\)
proves \eqref{eq:V92-word-lower}.  Since
\(d_j\asymp a_j^{1/2}\), the last asymptotic follows.
\end{proof}

\section{Failure of the literal local weight}
\label{sec:V92-local-failure}

\begin{theorem}[Physical counterpacket to
\(\mathfrak a_{31,r}\)-locality]
\label{thm:V92-local-card-false}
For the preceding family,
\begin{equation}
 \mathfrak a_{31,r}\asymp a_j^{-4},
\label{eq:V92-local-weight-scale}
\end{equation}
and therefore
\begin{equation}
 {\kappa_\otimes(\widetilde W_j^*\widetilde W_j)
  \over\mathfrak a_{31,r}^{\,2}}
 \ge c\,a_j^{7/2}\longrightarrow\infty.
\label{eq:V92-local-card-diverges}
\end{equation}
Thus the hot weighted local card
\eqref{eq:V89-weighted-card}, and even its oriented square analogue,
is false for an admissible physical \(SU(2)\) source family.
\end{theorem}

\begin{proof}
Before \(r\), the only internal triple edges are the fixed edges
\(12\) at \(p\) and \(13\) at \(q\).  At \(r\), the only joining
bridge is \(14\).  Their raw stocks are fixed positive constants.
Using \eqref{eq:V92-mass-scales},
\begin{align}
 \theta_{12}^{<r}&\asymp a_j^{-2},
 &
 \theta_{13}^{<r}&\asymp a_j^{-1},
 &
 \theta_{14,r}&\asymp a_j^{-1}.
\label{eq:V92-local-edge-scales}
\end{align}
Every other term of \eqref{eq:V89-local-a31} vanishes for this
support pattern.  Their product proves
\eqref{eq:V92-local-weight-scale}.  Combine its square with
\eqref{eq:V92-word-lower}.
\end{proof}

\begin{firewall}[The counterpacket does not refute global MTA]
\label{fire:V92-not-global-counterexample}
The theorem refutes a sufficient termwise localization introduced in
V89.  It does not refute the corrected global fixed-word card
\eqref{eq:V89-correct-FCQ31}, the global \(3+1\) source, or quartic
MTA.  The rescuing suffix edge is a legitimate global marked-tree
edge and must not be deleted from the target.
\end{firewall}

\section{The global atlas controls the same family}
\label{sec:V92-global-rescue}

\begin{proposition}[Exact suffix-edge rescue]
\label{prop:V92-global-rescue}
For the V92 family,
\begin{align}
 \theta_{12}&\ge c>0,
\label{eq:V92-global-strong-edge}\\
 \theta_{13}&\asymp a_j^{-1},
 &
 \theta_{14}&\asymp a_j^{-1},
\label{eq:V92-global-weak-edges}\\
 \widehat{\mathscr A}_{31}^{123|4}
 &\ge c\,a_j^{-2}.
\label{eq:V92-global-atlas-scale}
\end{align}
Moreover
\begin{equation}
 \kappa_\otimes(\widetilde W_j^*\widetilde W_j)
 \le C a_j^{-4}
 \le C'
 \bigl(\widehat{\mathscr A}_{31}^{123|4}\bigr)^2.
\label{eq:V92-global-controls}
\end{equation}
\end{proposition}

\begin{proof}
At the suffix coordinate,
\(H_{12}\ge a_j^2\), while
\(\Xi_1\Xi_2\asymp a_j^2\), proving
\eqref{eq:V92-global-strong-edge}.  The only \(13\) and \(14\)
stocks are fixed and their denominators contain one growing row mass,
giving \eqref{eq:V92-global-weak-edges}.  The global tree
\(\{12,13,14\}\) is a term of
\eqref{eq:V89-correct-A31}; its product proves
\eqref{eq:V92-global-atlas-scale}.

All heat moments are contractions and \(A_0\) is fixed, so
\(\|W_j\|\le C\).  After the four-row normalization,
\[
 \kappa_\otimes(\widetilde W_j^*\widetilde W_j)
 \le\|\widetilde W_j\|^2
 \le {C\over
  (\Xi_1\Xi_2\Xi_3\Xi_4)^2}
 \le Ca_j^{-4}.
\]
Compare with the square of
\eqref{eq:V92-global-atlas-scale}.
\end{proof}

\begin{assessment}[Correction of the V89--V91 route]
\label{ass:V92-route-correction}
The thermal theorem \Cref{thm:V89-thermal-weighted} remains valid.
The assertion that \eqref{eq:V89-weighted-card} is the universal hot
interface is withdrawn.  V90--V91 remain valid as conditional
reductions showing why that stronger local interface fails: their
terminal strong atom is precisely the suffix rescue constructed here.
They may not be cited as a proof that the global \(3+1\) card fails.
\end{assessment}

\section{An ordered one-rescue weight}
\label{sec:V92-rescue-weight}

For a pair edge \(uv\), put
\[
 \theta_{uv}^{>r}:=
 {H_{uv}^{>r}\over\Xi_u\Xi_v}.
\]
If \(T\) has its two edges written as the unordered set
\(\{f,g\}\), define
\begin{equation}
 \mathfrak b_{T,i,r}
 :=
 \left(
 \theta_f^{<r}\theta_g^{<r}
 +\theta_f^{<r}\theta_g^{>r}
 +\theta_g^{<r}\theta_f^{>r}
 \right)\theta_{i4,r},
\label{eq:V92-rescue-branch}
\end{equation}
and
\begin{equation}
 \mathfrak b_{31,r}
 :=
 \sum_{\substack{T\in\mathfrak T(\{1,2,3\})\\
                  i\in\{1,2,3\}}}
 \mathfrak b_{T,i,r}.
\label{eq:V92-rescue-weight}
\end{equation}
The first term is literal.  The other two retain one actual prefix
edge and permit the second internal-tree edge to be paid by one
complete suffix bank.

\begin{lemma}[Ordered rescue weights sum without support loss]
\label{lem:V92-rescue-summation}
\begin{equation}
 \boxed{\;
 \sum_r\mathfrak b_{31,r}
 \le3\widehat{\mathscr A}_{31}^{123|4}.\;}
\label{eq:V92-rescue-summation}
\end{equation}
\end{lemma}

\begin{proof}
Fix \(T=\{f,g\}\) and \(i\).  Nonnegativity and Cartesian enlargement
give
\begin{align*}
 \sum_r\theta_f^{<r}\theta_g^{<r}\theta_{i4,r}
 &\le\theta_f\theta_g\theta_{i4},\\
 \sum_r\theta_f^{<r}\theta_g^{>r}\theta_{i4,r}
 &\le\theta_f\theta_g\theta_{i4},\\
 \sum_r\theta_g^{<r}\theta_f^{>r}\theta_{i4,r}
 &\le\theta_f\theta_g\theta_{i4}.
\end{align*}
The order restrictions only delete nonnegative triples of coordinate
indices.  Sum the three inequalities over the nine pairs \((T,i)\).
\end{proof}

\begin{proposition}[The V92 counterpacket is exactly rescued]
\label{prop:V92-rescue-exact}
For the V92 family,
\begin{equation}
 \mathfrak b_{31,r}\ge
 \theta_{13}^{<r}\theta_{12}^{>r}\theta_{14,r}
 \asymp a_j^{-2}.
\label{eq:V92-rescue-scale}
\end{equation}
Consequently its complete square is bounded by
\(C\mathfrak b_{31,r}^2\).
\end{proposition}

\begin{proof}
The displayed term occurs in
\eqref{eq:V92-rescue-branch} for
\(T=\{12,13\}\) and \(i=1\).
The high suffix gives \(\theta_{12}^{>r}\ge c\), while
\eqref{eq:V92-local-edge-scales} gives the other two scales.
The final assertion is \eqref{eq:V92-global-controls}.
\end{proof}

\begin{definition}[Rerouted weighted \(3+1\) card]
\label{def:V92-RWFCQ31}
The card \(RWFCQ_{31}\) is the pair of estimates
\begin{align}
 \kappa_\otimes(
  \widetilde W_{r,\pi}^*\widetilde W_{r,\pi})
 &\le C_{G,s}\mathfrak b_{31,r}^{\,2},
\label{eq:V92-RWFCQ-square}\\
 \kappa_\otimes(
  \widetilde W_{r,\pi}\widetilde W_{r,\pi}^*)
 &\le C_{G,s}\mathfrak b_{31,r}^{\,2},
\label{eq:V92-RWFCQ-left-square}
\end{align}
for every payer placement.  Product-state Cauchy--Schwarz and
\Cref{lem:V92-rescue-summation} then give the global \(3+1\) source
bound with the corrected atlas.
\end{definition}

\begin{proposition}[The rerouted square card implies the global
\(3+1\) source bound]
\label{prop:V92-RWFCQ-suffices}
If \eqref{eq:V92-RWFCQ-square}--%
\eqref{eq:V92-RWFCQ-left-square} hold at the orientation required
by each output-dual summand and for every payer placement, then the
complete \(3+1\) first-connectivity source is bounded by
\(C_{G,s}\widehat{\mathscr A}_{31}^{123|4}\).
\end{proposition}

\begin{proof}
Product-state Cauchy--Schwarz gives
\[
 \kappa_\otimes^{\rm abs}(\widetilde W_{r,\pi})
 \le C\mathfrak b_{31,r}.
\]
Subadditivity and \eqref{eq:V92-rescue-summation} bound the sum over
\(r\).  The payer list is finite.
\end{proof}

\begin{firewall}[One rescue edge, not an arbitrary global tree]
\label{fire:V92-one-rescue}
The weight \(\mathfrak b_{31,r}\) does not attach an unrelated
global atlas term to every joining coordinate.  Each monomial keeps
the actual joining bridge at \(r\), one actual edge in the connected
triple prefix, and one actual edge in the complete suffix.  The three
coordinates are ordered and are summed once by
\Cref{lem:V92-rescue-summation}.
\end{firewall}

\section{V92 result ledger and V93 acquisition order}
\label{sec:V92-status}

\begin{theorem}[Third-series V92 result ledger]
\label{thm:V92-results}
Relative to V1--V91, V92 proves:
\begin{enumerate}[label=\textup{(V92.\arabic*)},leftmargin=6.3em]
\item a fixed nonzero physical \(SU(2)\) \(3+1\) base word;
\item exact tensorization after appending an equal-spin protected
      suffix;
\item the lower bound \(a_j^{-9/2}\) for its normalized complete
      square capacity;
\item divergence by \(a_j^{7/2}\) of the proposed literal local
      square ratio;
\item falsity of the universal hot
      \(\mathfrak a_{31,r}\)-weighted card;
\item control of the same family by the legitimate strong global
      suffix edge;
\item an ordered one-rescue weight summable by three copies of the
      corrected global atlas; and
\item exact rescue of the counterpacket by that new weight, together
      with a proof that its square card implies the global source
      estimate.
\end{enumerate}
\end{theorem}

\begin{proof}
These are
\Cref{lem:V92-base-nonzero,
lem:V92-tensor-scales,
thm:V92-local-card-false,
prop:V92-global-rescue,
ass:V92-route-correction,
lem:V92-rescue-summation,
prop:V92-rescue-exact,
prop:V92-RWFCQ-suffices}.
\end{proof}

\begin{programme}[V93: prove or falsify the one-rescue card]
\label{prog:V92-V93}
\begin{enumerate}[leftmargin=3em]
\item Resolve whether the strong atom of
      \Cref{cor:V91-strong-edge-residual} lies in the prefix or
      suffix without changing its exact source role.
\item On the suffix branch, retain the high internal edge and the
      complete suffix heat factor together.  Use the V83 marginal
      sandwich for one dimension gain and the enlarged rescue weight
      for the remaining normalization; do not ask the heat factor to
      recreate the suffix edge.
\item On the prefix branch, use the literal term of
      \eqref{eq:V92-rescue-branch} and retain the connected ternary
      response before output pinching.
\item Test \(RWFCQ_{31}\) on equal-spin singlets with arbitrary payer
      location and on the spin-one output obstruction of V48.
\item If one suffix rescue is insufficient, construct the minimal
      two-rescue ordered atlas and check its summability before any
      analytic estimate.
\item Revisit the \(2+1+1\) weighted gate only after deciding whether
      the same suffix-rescue defect occurs there.
\end{enumerate}
\end{programme}

\begin{firewall}[Claim boundary after third-series V92]
\label{fire:V92-claim-boundary}
V92 disproves the hot literal local card but not the global
\(3+1\) card.  It supplies a summable one-rescue replacement and
closes the explicit counterpacket under that replacement.  It does
not prove \(RWFCQ_{31}\), the corrected global \(3+1\) source, the
weighted global \(2+1+1\) source, the global discrete source,
globally aggregated \([4]\), \([3,3]\), higher MTA, WUFC, CPM,
cutoff removal, reflection positivity, Osterwalder--Schrader
reconstruction, construction of the continuum Yang--Mills measure,
or a uniform positive continuum spectral gap.  The full Yang--Mills
mass gap has not been proved.
\end{firewall}

\paragraph{Parent-version record.}
V92 was formed from
\texttt{Renewal\_Yang\_Mills\_Mass\_Gap\_Research\_Notes\_V91.tex}.
The V91 parent had \(44\,056\) logical source lines,
\(1\,564\,123\) bytes, and SHA--256
\[
 \texttt{168cb5de62e2cf908865f8d9fbadab31af7c70c8566f0000f0917ca972f9751c}.
\]
The next compulsory companion audit is V95.

\clearpage
\part{Third-series V93: globally marked contraction rescue and
closure of all \texorpdfstring{\(3+1\)}{3+1}
first-connectivity sources}
\label{part:V93}

\section{From the V92 counterpacket to the correct local weight}
\label{sec:V93-purpose}

The V92 ordered rescue weight uses a suffix edge explicitly.  For a
uniform theorem, the strong internal edge found by V91 may lie either
before or after the joining coordinate.  There is no need to split
those placements: keep that edge as a complete global pair stock,
retain one actual prefix connector for the triple, and retain the
actual local bridge at the joining coordinate.  These three edges form
a spanning tree and their ordered sum is still support-uniform.

The terminal analytic estimate is then elementary but decisive.  A
complete fixed-payer Wilson word has a representation-uniform raw
norm.  On the strong-edge counterpacket, its four input-mass
normalization is bounded by the rescued tree monomial.  Thus the
terminal packet closes by ordinary norm, without asking one protected
heat coordinate for a false second Casimir power.

\begin{assessment}[Nonduplication boundary]
\label{ass:V93-nonduplication}
V93 uses the V90--V91 counterpacket reduction and the V92 warning.  It
does not redo their heat or ascent arguments.  Its new ingredients are
the global-prefix-local ordered weight, a uniform raw complete-word
bound, the mass-cancellation comparison on the strong-edge residual,
and the contradiction argument upgrading those facts to the full
weighted \(3+1\) source theorem.
\end{assessment}

\section{The total-edge rescue weight}
\label{sec:V93-weight}

If \(T=\{f,g\}\) is one of the three trees on
\(\{1,2,3\}\), define
\begin{equation}
 \mathfrak c_{T,i,r}
 :=
 \left(
   \theta_f^{<r}\theta_g^{<r}
  +\theta_f\theta_g^{<r}
  +\theta_g\theta_f^{<r}
 \right)\theta_{i4,r},
\qquad i\in\{1,2,3\},
\label{eq:V93-rescue-branch}
\end{equation}
and
\begin{equation}
 \mathfrak c_{31,r}
 :=
 \sum_{\substack{T\in\mathfrak T(\{1,2,3\})\\
                  i\in\{1,2,3\}}}
 \mathfrak c_{T,i,r}.
\label{eq:V93-rescue-weight}
\end{equation}
The second and third terms designate one internal edge as global and
retain the other as an actual prefix connector.

\begin{lemma}[The total-edge rescue remains summable]
\label{lem:V93-rescue-summation}
\begin{equation}
 \boxed{\;
 \sum_r\mathfrak c_{31,r}
 \le3\widehat{\mathscr A}_{31}^{123|4}.\;}
\label{eq:V93-rescue-summation}
\end{equation}
\end{lemma}

\begin{proof}
Fix \(T=\{f,g\}\) and \(i\).  Since all stocks are nonnegative,
\begin{align*}
 \sum_r\theta_f^{<r}\theta_g^{<r}\theta_{i4,r}
 &\le\theta_f\theta_g\theta_{i4},\\
 \sum_r\theta_f\theta_g^{<r}\theta_{i4,r}
 &\le\theta_f\theta_g\theta_{i4},\\
 \sum_r\theta_g\theta_f^{<r}\theta_{i4,r}
 &\le\theta_f\theta_g\theta_{i4}.
\end{align*}
In the last two lines the global factor is taken outside the sum and
the ordered prefix/local pair is enlarged to its Cartesian product.
Sum over \(T,i\).  Each right side is one of the nine monomials of
\eqref{eq:V89-correct-A31}.
\end{proof}

\begin{corollary}[V93 contains both earlier weights]
\label{cor:V93-weight-dominance}
Pointwise in \(r\),
\begin{equation}
 \mathfrak c_{31,r}\ge
 \mathfrak b_{31,r}\ge
 \mathfrak a_{31,r}.
\label{eq:V93-weight-dominance}
\end{equation}
\end{corollary}

\begin{proof}
Use
\(\theta_f=\theta_f^{<r}+\theta_{f,r}+\theta_f^{>r}\)
and nonnegativity in
\eqref{eq:V93-rescue-branch}.  The literal first term is exactly the
corresponding term of \(\mathfrak a_{31,r}\); replacing the global
factor by its suffix part gives the two rescue terms of
\eqref{eq:V92-rescue-branch}.
\end{proof}

\section{A uniform raw bound for a fixed-payer word}
\label{sec:V93-raw-bound}

\begin{lemma}[Complete \(3+1\) words are raw contractions up to a
fixed constant]
\label{lem:V93-raw-bound}
Fix the joining coordinate \(r\), one ternary-tree/bridge resolution
\((T,i)\), one payer placement \(\pi\), and either output-duality
orientation.  Let \(W_{T,i,r,\pi}^{\rm raw}\) be the complete
un-normalized physical word.  Then
\begin{equation}
 \boxed{\;
 \opnorm{W_{T,i,r,\pi}^{\rm raw}}\le C_{G,s},\;}
\label{eq:V93-raw-bound}
\end{equation}
uniformly in support, input representations, and physical
multiplicities.  Consequently
\begin{equation}
 \opnorm{\widetilde W_{T,i,r,\pi}}
 \le {C_{G,s}\over\Xi_1\Xi_2\Xi_3\Xi_4}.
\label{eq:V93-normalized-contraction}
\end{equation}
\end{lemma}

\begin{proof}
Every complete Wilson moment is a contraction.  A binary covariance
is the difference of two contractions, and a ternary cumulant is the
fixed seven-term moment--cumulant polynomial; hence their raw norms
are bounded by constants depending only on the fixed order.  Physical
isotypic pinching is contractive when kept complete and does not sum
multiplicity copies.

If the payer lies in the prefix or joining covariance, use the paid
minimum in the frozen V66/V77 interfaces and discard its stock entry,
leaving the uniform \(C_{G,s}\) branch.  If it lies in a singleton or
suffix bank, use the complete bankwise debit
\(\texttt{lem:V67-bank-payer}\), or the arbitrary-suffix
one-resolvent theorem recorded in f2 V81/V96.  The payer occurs once,
so these alternatives are not multiplied.  All remaining factors are
complete contractions.  This proves \eqref{eq:V93-raw-bound}.
Division by the four input masses proves
\eqref{eq:V93-normalized-contraction}.
\end{proof}

\begin{firewall}[Why the raw bound was not enough before V91]
\label{fire:V93-raw-not-general}
Without a strong global edge, a normalized tree can be as small as
order \(M_*^{-6}\), while the square of the four-row input
normalization need be only order \(M_*^{-8}\) when all rows are
comparable.  The raw contraction then misses four powers.  V93 uses it
only after the strong internal edge and its cofinal mass comparison
have been proved.
\end{firewall}

\section{The strong-edge contraction comparison}
\label{sec:V93-strong-comparison}

Retain alternative \textup{(iii)} of
\Cref{thm:V91-ascent-alternative}.  Thus \(a,c\) are distinct
triple rows, \(d\) is the remaining triple row,
\[
 \Xi_a\ge\Xi_i,\qquad
 a_{a,e}\ge\sigma\Xi_a,\qquad
 a_{c,e}\ge\lambda a_{a,e},\qquad
 \theta_{ac}\ge\eta.
\]
In a nonzero resolved ternary prefix tree \(T\), at least one edge
\[
 f\in\{ad,cd\}
\]
connects \(d\) to \(\{a,c\}\).  Retain such an edge.  The resolved
joining branch has one actual bridge \(i4\) at \(r\).

\begin{lemma}[Actual prefix and bridge floors]
\label{lem:V93-actual-floors}
There is a fixed \(a_*>0\) such that
\begin{align}
 \theta_f^{<r}
 &\ge {a_*^2\over\Xi_u\Xi_v},
 \qquad f=uv,
\label{eq:V93-prefix-floor}\\
 \theta_{i4,r}
 &\ge {a_*^2\over\Xi_i\Xi_4}.
\label{eq:V93-bridge-floor}
\end{align}
\end{lemma}

\begin{proof}
The selected resolved prefix-tree term is nonzero, so its edge stock
\(H_f^{<r}\) contains at least one active coordinate summand
\(a_{u,q}a_{v,q}\ge a_*^2\).  Likewise the nonzero \(i\)-th branch
of the joining covariance contains the active local product
\(h_{i4,r}\ge a_*^2\).  Divide by the corresponding total row masses.
\end{proof}

\begin{theorem}[Strong-edge words are bounded by one rescued tree]
\label{thm:V93-strong-word}
On the strong-edge residual, every selected nonzero branch satisfies
\begin{equation}
 \boxed{\;
 \opnorm{\widetilde W_{T,i,r,\pi}}
 \le
 C_{G,s,\sigma,\lambda,\eta}
 \theta_{ac}\theta_f^{<r}\theta_{i4,r}
 \le
 C\mathfrak c_{31,r}.\;}
\label{eq:V93-strong-word}
\end{equation}
The estimate holds for every payer placement and both square
orientations.
\end{theorem}

\begin{proof}
First,
\[
 \Xi_c\ge a_{c,e}
 \ge\lambda a_{a,e}
 \ge\lambda\sigma\Xi_a.
\]
If \(f=ad\), then
\Cref{lem:V93-actual-floors} and \(\theta_{ac}\ge\eta\)
give
\[
 \theta_{ac}\theta_{ad}^{<r}\theta_{i4,r}
 \ge
 {\eta a_*^4\over
  \Xi_a\Xi_d\Xi_i\Xi_4}.
\]
Compared with the right side of
\eqref{eq:V93-normalized-contraction}, the required scalar ratio is
\[
 {\Xi_a\Xi_d\Xi_i\Xi_4
  \over\Xi_a\Xi_c\Xi_d\Xi_4}
 =
 {\Xi_i\over\Xi_c}
 \le{\Xi_a\over\Xi_c}
 \le {1\over\lambda\sigma}.
\]

If \(f=cd\), the rescued monomial is at least
\[
 {\eta a_*^4\over
  \Xi_c\Xi_d\Xi_i\Xi_4}.
\]
The corresponding ratio is
\[
 {\Xi_c\Xi_d\Xi_i\Xi_4
  \over\Xi_a\Xi_c\Xi_d\Xi_4}
 =
 {\Xi_i\over\Xi_a}\le1.
\]
Insert these comparisons in
\eqref{eq:V93-normalized-contraction}.  In either case,
\(\{ac,f,i4\}\) is a spanning tree: \(ac\) joins two triple rows,
\(f\) attaches the third, and \(i4\) attaches row \(4\).  Therefore
its ordered monomial occurs in one of the global-prefix-local terms
of \eqref{eq:V93-rescue-weight}.  Ordinary norm is invariant under
adjoint, so both orientations are covered.
\end{proof}

\section{The rerouted weighted square card}
\label{sec:V93-card}

\begin{theorem}[Uniform total-edge-rescued \(FCQ_{31}\)]
\label{thm:V93-rescued-FCQ31}
For every compact connected semisimple \(G\), fixed admissible heat
parameter, joining coordinate \(r\), and payer placement \(\pi\),
\begin{align}
 \kappa_\otimes(
  \widetilde W_{r,\pi}^*\widetilde W_{r,\pi})
 &\le C_{G,s}\mathfrak c_{31,r}^{\,2},
\label{eq:V93-rescued-square}\\
 \kappa_\otimes(
  \widetilde W_{r,\pi}\widetilde W_{r,\pi}^*)
 &\le C_{G,s}\mathfrak c_{31,r}^{\,2}.
\label{eq:V93-rescued-left-square}
\end{align}
The constant is uniform in support, representations, physical
multiplicities, and \(r\).
\end{theorem}

\begin{proof}
Suppose no uniform constant exists.  Select a sequence for which one
of the two ratios diverges.  The finite Gram recombination
\Cref{lem:V90-nine-recombination} selects fixed
\((T,i,\pi)\) with
\[
 {\kappa_\otimes(
  \widetilde W_{T,i,r,\pi}^*
  \widetilde W_{T,i,r,\pi})\over
  \mathfrak c_{31,r}^{\,2}}
 \longrightarrow\infty
\]
in the relevant orientation.  Since
\(\mathfrak c_{31,r}\ge\mathfrak a_{T,i,r}\) by
\Cref{cor:V93-weight-dominance}, this is also a counterpacket to the
literal branch card.

Apply the exhaustive V90--V91 reduction.  A bounded demanded row is
closed by \Cref{cor:V90-demand-hot}; a source-compatible diffuse row
by \Cref{thm:V90-k4-terminal,lem:V91-larger-diffuse}; an unbalanced
or cut-crossing cofinal atom by
\Cref{lem:V90-known-atom-branches}; and a strict mass ascent
terminates by \Cref{thm:V91-ascent-alternative}.  Therefore every
persistent counterpacket reaches the strong-edge residual of
\Cref{cor:V91-strong-edge-residual}.

On that residual,
\Cref{thm:V93-strong-word} gives
\(\|\widetilde W_{T,i,r,\pi}\|\le C\mathfrak c_{31,r}\).
Positive capacity is bounded by ordinary square norm, contradicting
the selected divergence.  The same proof applied to adjoints covers
the other orientation.
\end{proof}

\begin{corollary}[Corrected fixed-word \(FCQ_{31}\)]
\label{cor:V93-correct-FCQ31}
The corrected global-atlas card
\eqref{eq:V89-correct-FCQ31} holds for every payer placement.
\end{corollary}

\begin{proof}
Every local or total normalized pair stock is at most its global
counterpart.  Hence each term of
\(\mathfrak c_{31,r}\) is one of the nine global tree monomials, up
to the fixed threefold duplication in
\eqref{eq:V93-rescue-branch}.  Thus
\(\mathfrak c_{31,r}\le3
\widehat{\mathscr A}_{31}^{123|4}\).  Apply
\Cref{thm:V93-rescued-FCQ31}.
\end{proof}

\section{Global \(3+1\) source closure}
\label{sec:V93-global}

\begin{theorem}[Complete global \(3+1\) first-connectivity source]
\label{thm:V93-global-S31}
For the source \(\mathcal S_{31}^{123|4}\) in
\eqref{eq:V88-S31}, with the sole final payer allocated exactly once,
\begin{equation}
 \boxed{\;
 \kappa_\otimes^{\rm abs}
 \bigl((\mathcal S_{31}^{123|4})^{\rm pay}\bigr)
 \le
 C_{G,s}\widehat{\mathscr A}_{31}^{123|4}.\;}
\label{eq:V93-global-S31}
\end{equation}
The same theorem holds for all four choices of singleton row.
\end{theorem}

\begin{proof}
Product-state Cauchy--Schwarz and
\Cref{thm:V93-rescued-FCQ31} give, in the orientation selected for
each word,
\[
 \kappa_\otimes^{\rm abs}(\widetilde W_{r,\pi})
 \le C_{G,s}\mathfrak c_{31,r}.
\]
Sum first over \(r\), use
\Cref{lem:V93-rescue-summation}, and then sum over the fixed finite
payer and physical-resolution types.  Exact first-connectivity
\Cref{thm:V88-S31-identity} identifies the resulting sum with the
complete source.  Row permutation preserves every constant and gives
the other three labels.
\end{proof}

\begin{corollary}[Updated proper-prefix census]
\label{cor:V93-census}
All three \(2+2\) labels and all four \(3+1\) labels now satisfy
complete global source theorems.  All six \(2+1+1\) labels retain
their fixed-word card from V87 but still require a valid global
aggregation.  The discrete prefix remains untreated.  Thus seven of
the fourteen proper-prefix labels are globally closed and six more
are fixed-word closed.
\end{corollary}

\begin{proof}
Use \Cref{cor:V34-S22-no-gap,thm:V93-global-S31} and the count in
\Cref{thm:V88-nine-labels}.  The local \(2+1+1\) weighted interface
from V88 has not been repaired after the suffix-rescue defect
identified in V92.
\end{proof}

\begin{firewall}[What the contraction rescue does not prove]
\label{fire:V93-scope}
The raw contraction works only after a strong global edge is tied to
a cofinal mass comparison.  It does not license attaching an
arbitrary global edge to every source word.  The actual prefix
connector and joining bridge remain ordered in
\(\mathfrak c_{31,r}\), and their sum is proved in
\Cref{lem:V93-rescue-summation}.
\end{firewall}

\section{V93 result ledger and V94 acquisition order}
\label{sec:V93-status}

\begin{theorem}[Third-series V93 result ledger]
\label{thm:V93-results}
Relative to V1--V92, V93 proves:
\begin{enumerate}[label=\textup{(V93.\arabic*)},leftmargin=6.3em]
\item the total-edge rescue weight and its three-atlas summation;
\item a representation- and support-uniform raw norm for every
      fixed-payer complete \(3+1\) word;
\item exact cancellation of the four input masses by one strong
      global edge, one actual prefix connector, and the actual joining
      bridge;
\item the uniform rerouted oriented square card;
\item the corrected fixed-word \(FCQ_{31}\);
\item the complete global \(3+1\) source estimate for all four
      singleton labels; and
\item the updated proper-prefix census.
\end{enumerate}
\end{theorem}

\begin{proof}
These are
\Cref{lem:V93-rescue-summation,
lem:V93-raw-bound,
thm:V93-strong-word,
thm:V93-rescued-FCQ31,
cor:V93-correct-FCQ31,
thm:V93-global-S31,
cor:V93-census}.
\end{proof}

\begin{programme}[V94: test and repair \(2+1+1\) aggregation]
\label{prog:V93-V94}
\begin{enumerate}[leftmargin=3em]
\item Apply the V92 suffix construction to the
      \(2+1+1\) word before attempting
      \eqref{eq:V88-weighted-local}.  Decide whether that literal local
      card is false.
\item If it fails, retain the actual prefix edge \(12\), the actual
      joining-coordinate quotient edges, and one global rescuing edge
      in a finite rerouted atlas.
\item Prove scalar summability of the rerouted weights before using
      V87's fixed-word card.
\item Revisit the V76--V87 counterpacket reduction with the rerouted
      target; use the raw-contraction terminal whenever a strong edge
      and cofinal mass comparison appear.
\item Only after global \(2+1+1\) aggregation is valid, proceed to the
      discrete prefix.
\item The next compulsory companion audit remains V95.
\end{enumerate}
\end{programme}

\begin{firewall}[Claim boundary after third-series V93]
\label{fire:V93-claim-boundary}
V93 closes every \(3+1\) first-connectivity source globally.  It does
not close global \(2+1+1\) aggregation, the global discrete source,
globally aggregated \([4]\), \([3,3]\), higher MTA, WUFC, CPM,
cutoff removal, reflection positivity, Osterwalder--Schrader
reconstruction, construction of the continuum Yang--Mills measure,
or a uniform positive continuum spectral gap.  The full Yang--Mills
mass gap has not been proved.
\end{firewall}

\paragraph{Parent-version record.}
V93 was formed from
\texttt{Renewal\_Yang\_Mills\_Mass\_Gap\_Research\_Notes\_V92.tex}.
The V92 parent had \(44\,523\) logical source lines,
\(1\,580\,603\) bytes, and SHA--256
\[
 \texttt{9026d1a880622f74a905b61fc9ab3d2789305a57fb158f0c98e9d0304a95a3d2}.
\]
The next compulsory companion audit is V95.

\clearpage
\part{Third-series V94: failure of literal
\texorpdfstring{\(2+1+1\)}{2+1+1} localization and a summable
two-rescue atlas}
\label{part:V94}

\section{Purpose}
\label{sec:V94-purpose}

V88 correctly observed that the V87 fixed-\(r\) card did not by itself
sum over joining coordinates, but proposed the literal local weight
\(\mathfrak a_{r,\pi}\) of
\eqref{eq:V88-weighted-local}.  V92 showed that the analogous
\(3+1\) demand is false: a protected suffix can carry a strong global
edge absent from the literal prefix weight.  V93 repaired that family
by allowing a globally marked rescue edge while retaining enough
ordered source marks to sum.

The same defect occurs for \(2+1+1\).  This chapter constructs a
physical \(SU(2)\) counterpacket.  It then gives a finite ordered
two-rescue weight: among the three edges of each legal tree, the
prefix edge and at most one of the two joining edges may be replaced
by their total stocks, while at least one actual joining-coordinate
edge remains.  Every resulting term sums into the same global tree
monomial.

\begin{assessment}[Nonduplication boundary]
\label{ass:V94-nonduplication}
V94 does not redo the V35 eight-tree census or V87 fixed-word card.
It tests the stronger V88 localization, withdraws it after a physical
counterpacket, and proves the scalar summability of its replacement.
The analytic two-rescue square card is left to V95.
\end{assessment}

\section{A nonzero local ternary joining block}
\label{sec:V94-base}

Let \(\mu_s\) be a nondegenerate \(SU(2)\) heat law.  We first record
that its Peter--Weyl algebra contains a scalar with nonzero third
cumulant.

\begin{lemma}[A finite Fourier polynomial with nonzero third
cumulant]
\label{lem:V94-nonzero-third}
There is a bounded real finite Peter--Weyl polynomial \(f\) such that
\begin{equation}
 \kappa_3(f,f,f)
 =
 \mathbb E_{\mu_s}\bigl[(f-\mathbb E_{\mu_s}f)^3\bigr]\ne0.
\label{eq:V94-nonzero-third}
\end{equation}
\end{lemma}

\begin{proof}
The heat law has a positive smooth density and is non-atomic.  Choose
a Borel set \(E\) with
\(p:=\mu_s(E)\in(0,1)\setminus\{\tfrac12\}\).  For
\(g=1_E-p\),
\[
 \mathbb E_{\mu_s}g^3=p(1-p)(1-2p)\ne0.
\]
Finite real Peter--Weyl polynomials are dense in
\(L^3(\mu_s)\).  The centered third moment is continuous in
\(L^3\), by H\"older's inequality and the uniform \(L^3\) control of
a convergent sequence.  Hence a sufficiently close finite polynomial
\(f\) retains a nonzero centered third moment.
\end{proof}

Choose a second bounded real finite Peter--Weyl polynomial \(\phi\)
with
\[
 \operatorname{Var}_{\mu_s}(\phi)>0.
\]
Take ordered coordinates \(p<r\).  At \(p\), rows \(1,2\) carry
\(\phi\), producing the connected pair prefix.  At \(r\), row \(1\)
inside the composite block \(12\), row \(3\), and row \(4\) carry
\(f\), while row \(2\)'s local factor is the identity.  Thus
\[
 J_r(12|3|4)
 =
 \kappa_3(f,f,f)\ne0
\]
on a fixed scalar coefficient.

\begin{lemma}[Fixed nonzero \(2+1+1\) base block]
\label{lem:V94-base-nonzero}
There is a fixed physical output block and a fixed payer placement in
\(\{p,r\}\) for which the exact \(2+1+1\) source word is a nonzero
finite-dimensional operator \(B_0\).  Both \(\|B_0\|\) and
\(\kappa_\otimes(B_0^*B_0)\) are positive fixed constants.
\end{lemma}

\begin{proof}
At \(p\), the prefix coefficient is
\(\operatorname{Var}_{\mu_s}(\phi)>0\).  At \(r\), the assembled
joining coefficient is nonzero by
\Cref{lem:V94-nonzero-third}.  Their row partitions have cumulative
join \(\widehat1\), first reached at \(r\), so the linked-word
criterion \Cref{prop:V7-linked-word} assigns the product coefficient
to \eqref{eq:V35-S211}.  Resolve a nonzero physical output block and
place the payer on a fixed nontrivial output in \(p\) or \(r\).
The resulting operator \(B_0\) is nonzero.  As in
\Cref{lem:V92-base-nonzero}, a nonzero positive operator
\(B_0^*B_0\) has positive expectation on some product vector.
\end{proof}

\section{A protected suffix counterpacket}
\label{sec:V94-suffix}

Add a coordinate \(e>r\) at which rows \(1,2\) carry equal spin \(j\)
and rows \(3,4\) are inactive.  Its complete suffix heat factor is
\(M_j(s)\).  Let
\[
 a_j=1+j(j+1),\qquad d_j=2j+1.
\]
The exact block is
\begin{equation}
 W_j=B_0\otimes M_j(s).
\label{eq:V94-tensor-word}
\end{equation}
Moreover
\begin{equation}
 \Xi_1,\Xi_2\asymp a_j,
 \qquad
 \Xi_3,\Xi_4\asymp1.
\label{eq:V94-mass-scales}
\end{equation}

\begin{lemma}[Normalized square lower bound]
\label{lem:V94-square-lower}
For
\[
 \widetilde W_j:=
 {W_j\over\Xi_1\Xi_2\Xi_3\Xi_4},
\]
one has
\begin{equation}
 \boxed{\;
 \kappa_\otimes(\widetilde W_j^*\widetilde W_j)
 \ge c\,a_j^{-4}d_j^{-1}
 \asymp a_j^{-9/2}.\;}
\label{eq:V94-square-lower}
\end{equation}
\end{lemma}

\begin{proof}
Use a product vector with fixed positive
\(B_0^*B_0\) expectation and tensor it with a maximizing singlet
product vector at \(e\).  The exact tensor factorization
\eqref{eq:V94-tensor-word} and
\Cref{thm:V91-heat-square-obstruction} give
\(\kappa_\otimes(W_j^*W_j)\ge c/d_j\).
Now use \eqref{eq:V94-mass-scales}.
\end{proof}

\section{The V88 literal \(2+1+1\) weight is false}
\label{sec:V94-literal-false}

Put
\begin{align}
 p_r^-&:=\theta_{12}^{<r},
\label{eq:V94-local-p}\\
 x_r&:=\theta_{13,r}+\theta_{23,r},
 &
 y_r&:=\theta_{14,r}+\theta_{24,r},
 &
 z_r&:=\theta_{34,r},
\label{eq:V94-local-xyz}\\
 q_r&:=x_ry_r+x_rz_r+y_rz_r.
\label{eq:V94-local-q}
\end{align}
The literal normalized joining stock of the V35 eight trees is
\begin{equation}
 \mathfrak a_{211,r}:=p_r^-q_r.
\label{eq:V94-literal-weight}
\end{equation}

\begin{theorem}[Physical failure of the literal \(2+1+1\) card]
\label{thm:V94-literal-false}
For the V94 family,
\begin{equation}
 p_r^-\asymp a_j^{-2},
 \qquad
 x_r,y_r\asymp a_j^{-1},
 \qquad
 z_r\asymp1,
\label{eq:V94-local-scales}
\end{equation}
so
\begin{equation}
 \mathfrak a_{211,r}\asymp a_j^{-3}.
\label{eq:V94-literal-scale}
\end{equation}
Consequently
\begin{equation}
 {\kappa_\otimes(\widetilde W_j^*\widetilde W_j)
  \over\mathfrak a_{211,r}^{\,2}}
 \ge c\,a_j^{3/2}\longrightarrow\infty.
\label{eq:V94-literal-divergence}
\end{equation}
Thus the universal weighted local line
\eqref{eq:V88-weighted-local} is false, including its oriented-square
strengthening.
\end{theorem}

\begin{proof}
The only prefix \(12\) stock before \(r\) is fixed, while its
denominator contains both growing row masses, giving
\(p_r^-\asymp a_j^{-2}\).  At \(r\), only row \(1\) in the composite
block grows in total mass; hence each \(A\)-to-\(3\) or
\(A\)-to-\(4\) normalized edge is of order \(a_j^{-1}\).
Rows \(3,4\) have fixed masses and a fixed nonzero common local
coefficient, so \(z_r\asymp1\).  Therefore
\[
 q_r=x_ry_r+x_rz_r+y_rz_r\asymp a_j^{-1},
\]
which proves \eqref{eq:V94-literal-scale}.  Combine its square with
\eqref{eq:V94-square-lower}.
\end{proof}

\begin{firewall}[Fixed-word V87 and global MTA survive]
\label{fire:V94-not-global}
The counterpacket refutes only the stronger literal localization
introduced in V88.  It does not refute
\Cref{thm:V87-FCQ211}, the global \(2+1+1\) source estimate, or
quartic MTA.  Indeed the high suffix makes the total edge
\(\theta_{12}\) strong, and that edge is legitimate global tree
currency.
\end{firewall}

\begin{proposition}[The global atlas controls the V94 family]
\label{prop:V94-global-controls}
Along the V94 family,
\begin{equation}
 p=\theta_{12}\ge c>0,
 \qquad
 q=xy+xz+yz\asymp a_j^{-1},
 \qquad
 \mathscr A_{211}=pq\asymp a_j^{-1}.
\label{eq:V94-global-scales}
\end{equation}
Moreover
\begin{equation}
 \kappa_\otimes(\widetilde W_j^*\widetilde W_j)
 \le Ca_j^{-4}
 \le C'\mathscr A_{211}^{\,2}.
\label{eq:V94-global-controls}
\end{equation}
\end{proposition}

\begin{proof}
The suffix contribution \(a_j^2\) to \(H_{12}\) is comparable to
\(\Xi_1\Xi_2\), so \(p\ge c\).  The remaining totals have the same
orders as the local quantities in \eqref{eq:V94-local-scales}, giving
the stated \(q\).  Finally \(W_j\) has uniformly bounded raw norm,
so division by the squared four-row normalization gives
\(Ca_j^{-4}\).  Since
\(\mathscr A_{211}^2\asymp a_j^{-2}\), the last inequality follows.
\end{proof}

\section{A finite two-rescue atlas}
\label{sec:V94-two-rescue}

Let \(\mathcal E_{211}\) be the eight labelled trees of
\Cref{lem:V35-eight-trees}.  Each
\(\tau\in\mathcal E_{211}\) has the distinguished prefix edge
\(f_0=12\) and two distinguished joining edges \(f_1,f_2\).
For an edge \(f\), write \(\theta_{f,r}\) for its local normalized
stock and \(\theta_f\) for its total stock.

\begin{definition}[Two-rescue weight for one tree]
\label{def:V94-tree-rescue}
Set
\begin{equation}
\begin{split}
 \mathfrak d_{\tau,r}:={}&
 \theta_{12}^{<r}\theta_{f_1,r}\theta_{f_2,r}
 +\theta_{12}\theta_{f_1,r}\theta_{f_2,r}\\
 &+\theta_{12}^{<r}\theta_{f_1}\theta_{f_2,r}
 +\theta_{12}^{<r}\theta_{f_1,r}\theta_{f_2}\\
 &+\theta_{12}\theta_{f_1}\theta_{f_2,r}
 +\theta_{12}\theta_{f_1,r}\theta_{f_2}.
\end{split}
\label{eq:V94-tree-rescue}
\end{equation}
Define
\begin{equation}
 \mathfrak d_{211,r}
 :=
 \sum_{\tau\in\mathcal E_{211}}\mathfrak d_{\tau,r}.
\label{eq:V94-two-rescue-weight}
\end{equation}
\end{definition}

\begin{remark}[Admissible placements]
\label{rem:V94-rescue-placements}
The first line of \eqref{eq:V94-tree-rescue} contains the literal
term and the term with only the prefix edge globalized.  The second
line globalizes one joining edge while retaining the prefix order.
The third line may globalize both the prefix and one joining edge.
At least one of \(f_1,f_2\) remains at the actual joining coordinate
in every term.
\end{remark}

\begin{lemma}[Two-rescue summability]
\label{lem:V94-rescue-summation}
\begin{equation}
 \boxed{\;
 \sum_r\mathfrak d_{211,r}
 \le6\mathscr A_{211}.\;}
\label{eq:V94-rescue-summation}
\end{equation}
\end{lemma}

\begin{proof}
Fix \(\tau=\{12,f_1,f_2\}\).  For the literal term,
\[
 \sum_r
 \theta_{12}^{<r}\theta_{f_1,r}\theta_{f_2,r}
 \le\theta_{12}\theta_{f_1}\theta_{f_2}.
\]
For a term with exactly one local joining edge, take the one or two
global factors outside the sum and use, for example,
\[
 \sum_r\theta_{12}^{<r}\theta_{f_2,r}
 \le\theta_{12}\theta_{f_2}.
\]
For a term with both joining edges local, use
\(\sum_r\theta_{f_1,r}\theta_{f_2,r}
\le\theta_{f_1}\theta_{f_2}\).
Thus each of the six terms is at most the global tree monomial
\(\prod_{f\in\tau}\theta_f\).  Sum over the eight trees; their sum is
\(\mathscr A_{211}=p(xy+xz+yz)\).
\end{proof}

\begin{proposition}[The V94 counterpacket is rescued by the global
prefix edge]
\label{prop:V94-counterpacket-rescued}
For the V94 family,
\begin{equation}
 \mathfrak d_{211,r}
 \ge
 \theta_{12}\theta_{13,r}\theta_{34,r}
 \asymp a_j^{-1}.
\label{eq:V94-rescue-scale}
\end{equation}
Hence
\[
 \kappa_\otimes(\widetilde W_j^*\widetilde W_j)
 \le C\mathfrak d_{211,r}^{\,2}.
\]
\end{proposition}

\begin{proof}
The tree \(\{12,13,34\}\) is \(T_1^{(3)}\) in
\eqref{eq:V35-tree-three}.  Its second term in
\eqref{eq:V94-tree-rescue} is the displayed monomial.  Use
\eqref{eq:V94-global-scales,eq:V94-local-scales}, followed by
\eqref{eq:V94-global-controls}.
\end{proof}

\begin{definition}[Two-rescue \(2+1+1\) square card]
\label{def:V94-TRFCQ211}
The card \(TRFCQ_{211}\) is
\begin{align}
 \kappa_\otimes(
  \widetilde W_{r,\pi}^*\widetilde W_{r,\pi})
 &\le C_{G,s}\mathfrak d_{211,r}^{\,2},
\label{eq:V94-TRFCQ-square}\\
 \kappa_\otimes(
  \widetilde W_{r,\pi}\widetilde W_{r,\pi}^*)
 &\le C_{G,s}\mathfrak d_{211,r}^{\,2},
\label{eq:V94-TRFCQ-left-square}
\end{align}
in the output-dual orientation required by every resolved term and
for every payer placement.
\end{definition}

\begin{proposition}[\(TRFCQ_{211}\) implies global aggregation]
\label{prop:V94-TRFCQ-suffices}
If \eqref{eq:V94-TRFCQ-square}--%
\eqref{eq:V94-TRFCQ-left-square} hold, then the complete
\(2+1+1\) source \eqref{eq:V35-S211} satisfies its global literal
marked-tree estimate.
\end{proposition}

\begin{proof}
Product-state Cauchy--Schwarz gives
\[
 \kappa_\otimes^{\rm abs}(\widetilde W_{r,\pi})
 \le C\mathfrak d_{211,r}.
\]
Sum over \(r\), apply
\Cref{lem:V94-rescue-summation}, and sum over the fixed finite payer
types.  Row permutation treats all six choices of prefix pair.
\end{proof}

\begin{firewall}[The two-rescue weight is still source ordered]
\label{fire:V94-source-ordered}
No term of \(\mathfrak d_{211,r}\) consists of three global edges.
At least one actual joining edge retains the coordinate \(r\), and
any retained prefix edge is restricted to \(<r\).  This is why
\Cref{lem:V94-rescue-summation} has no support-size factor.
\end{firewall}

\section{V94 result ledger and V95 acquisition order}
\label{sec:V94-status}

\begin{theorem}[Third-series V94 result ledger]
\label{thm:V94-results}
Relative to V1--V93, V94 proves:
\begin{enumerate}[label=\textup{(V94.\arabic*)},leftmargin=6.3em]
\item existence of a finite Peter--Weyl scalar with nonzero third
      cumulant under the Wilson heat law;
\item a fixed nonzero physical \(2+1+1\) base word;
\item an equal-spin protected-suffix square lower bound;
\item divergence of the V88 literal local-card ratio by
      \(a_j^{3/2}\);
\item falsity of the proposed universal literal weighted card;
\item control of that family by the legitimate strong global prefix
      edge;
\item a six-term-per-tree two-rescue atlas and its support-uniform
      summability; and
\item exact rescue of the counterpacket plus a sufficient
      \(TRFCQ_{211}\) interface.
\end{enumerate}
\end{theorem}

\begin{proof}
These are
\Cref{lem:V94-nonzero-third,
lem:V94-base-nonzero,
lem:V94-square-lower,
thm:V94-literal-false,
prop:V94-global-controls,
lem:V94-rescue-summation,
prop:V94-counterpacket-rescued,
prop:V94-TRFCQ-suffices}.
\end{proof}

\begin{programme}[V95: audit and analytic two-rescue card]
\label{prog:V94-V95}
\begin{enumerate}[leftmargin=3em]
\item Perform the compulsory SM/NS audit before importing any new
      proof-order heuristic.
\item Assume failure of \(TRFCQ_{211}\) and select one of the fixed
      eight tree labels, payer types, and square orientations.
\item Reuse the V76--V87 counterpacket reduction only after checking
      each step against the enlarged target
      \(\mathfrak d_{211,r}\).
\item Close every terminal strong-edge packet by a V93-style raw
      contraction, choosing a rescued tree which retains one actual
      joining edge.
\item Use V87 diffuse-bank powers only on a bank exposed in the same
      exact source role.
\item If a terminal packet needs all three edges global, record that
      as a failure of \(TRFCQ_{211}\) and move upstream to the
      aggregate source rather than adding a nonsummable term.
\end{enumerate}
\end{programme}

\begin{firewall}[Claim boundary after third-series V94]
\label{fire:V94-claim-boundary}
V94 disproves the literal \(2+1+1\) weighted card and supplies a
summable two-rescue replacement.  It does not prove
\(TRFCQ_{211}\), global \(2+1+1\) aggregation, the global discrete
source, globally aggregated \([4]\), \([3,3]\), higher MTA, WUFC,
CPM, cutoff removal, reflection positivity,
Osterwalder--Schrader reconstruction, construction of the continuum
Yang--Mills measure, or a uniform positive continuum spectral gap.
The full Yang--Mills mass gap has not been proved.
\end{firewall}

\paragraph{Parent-version record.}
V94 was formed from
\texttt{Renewal\_Yang\_Mills\_Mass\_Gap\_Research\_Notes\_V93.tex}.
The V93 parent had \(44\,999\) logical source lines,
\(1\,596\,796\) bytes, and SHA--256
\[
 \texttt{6eea30786d3c483fe5b8b1ddf7a713927b280fc7293927de836aabeed34e3891}.
\]
The next compulsory companion audit is V95.

\clearpage
\part{Third-series V95: companion audit, maximal-row dichotomy,
and global closure of the
\texorpdfstring{\(2+1+1\)}{2+1+1} sources}
\label{part:V95}

\section{Compulsory SM/NS companion audit}
\label{sec:V95-audit}

The newest active companion files at this fifth-version checkpoint
are:
\begin{align}
 {\cal C}_{\rm SM}
 &:=
 \texttt{Renewal\_SM\_Einstein\_Continuum\_Research\_Notes\_V72.tex},
\label{eq:V95-SM-file}\\
 {\cal C}_{\rm NS}
 &:=
 \texttt{Renewal\_NS\_Stability\_Research\_Notes\_V31.tex}.
\label{eq:V95-NS-file}
\end{align}
Their exact identities are:
\begin{center}
\begin{tabular}{lrrl}
\toprule
file & lines & bytes & SHA--256\\
\midrule
SM V72 & \(27\,155\) & \(1\,001\,966\) &
\texttt{3a1c69717c12d5daf8221ced8bc3bcd8}\\
&&&\texttt{e4ad873ce2bb98f3e08d717c76ef1744}\\
NS V31 & \(23\,052\) & \(887\,537\) &
\texttt{f1121b62238ad5491bb7cf03635ea993}\\
&&&\texttt{4942edf573ed8faaa9ee79e1d38a6696}\\
\bottomrule
\end{tabular}
\end{center}

\begin{assessment}[Type-correct content of SM V69--V72]
\label{ass:V95-SM-audit}
SM V69 restores lattice units and proves that several quantities
previously hidden in fixed-cutoff constants grow at the ultraviolet
scale.  V70 moves the complete marginal packet into its running
parent rather than charging it as a small residual.  V72 assigns each
potentially singular left/right spectral cluster to one soft owner
before a hard logarithm is formed; transfer prevents a false inverse
but does not solve the soft sector.

The YM proof-order consequences are:
\begin{enumerate}[label=\textup{(\roman*)},leftmargin=4em]
\item the sum over joining coordinates must be proved before its
      constant is called support-uniform; and
\item a rescuing global edge must be assigned once to one displayed
      tree monomial, not counted both as heat suppression and as an
      additional edge.
\end{enumerate}
The V94 two-rescue atlas meets both requirements through
\Cref{lem:V94-rescue-summation}.  No determinant, singular-value, or
renormalization-group estimate is imported.
\end{assessment}

\begin{assessment}[Type-correct content of NS V31]
\label{ass:V95-NS-audit}
NS V31 remains unchanged.  Its flat dyadic mark removes a stopping
family multiplicity only when the exact first shell moment remains
visible; weaker data do not recover that weight.  The YM counterpart
is the actual local joining edge retained in every term of
\(\mathfrak d_{211,r}\).  It absorbs the \(r\)-sum and cannot be
replaced by an unweighted count.  No NS estimate is imported.
\end{assessment}

\begin{firewall}[Companion type boundary]
\label{fire:V95-audit-boundary}
The audit supplies consistency tests only.  All estimates below use
the exact YM first-connectivity word, normalized pair stocks, the
complete heat-bank card, and finite compact-group representation
algebra.
\end{firewall}

\section{A local polynomial floor for the two-rescue target}
\label{sec:V95-floor}

Fix one nonzero resolved \(2+1+1\) word at joining coordinate \(r\).
Its prefix contains an actual \(12\) edge before \(r\), and its
joining cumulant contains a quotient-tree pair of actual row edges
\(f_1,f_2\) at \(r\).  The lifted row graph
\(\tau=\{12,f_1,f_2\}\) is one of
\(\mathcal E_{211}\).

\begin{lemma}[Literal local floor]
\label{lem:V95-local-floor}
Let \(M_*:=\max_k\Xi_k\).  There is a fixed \(c_*>0\) such that
\begin{equation}
 \boxed{\;
 \mathfrak d_{211,r}\ge {c_*\over M_*^6},
 \qquad
 \mathfrak d_{211,r}^{\,2}\ge {c_*^2\over M_*^{12}}.\;}
\label{eq:V95-local-floor}
\end{equation}
\end{lemma}

\begin{proof}
The nonzero prefix edge has
\(H_{12}^{<r}\ge a_*^2\).  Each of the two selected local joining
edges has \(h_{f_\nu,r}\ge a_*^2\).  Therefore the literal first term
of \(\mathfrak d_{\tau,r}\) is at least
\[
 {a_*^6\over
  \prod_{uv\in\tau}\Xi_u\Xi_v}
 \ge {a_*^6\over M_*^6}.
\]
It is a term of \(\mathfrak d_{211,r}\).  Square the inequality.
\end{proof}

\begin{firewall}[The floor is not the final tree estimate]
\label{fire:V95-floor-scope}
The floor is used only against heat decay on a diffuse or gapped
maximal-row bank.  On the concentrated balanced branch, the exact
strong edge and exact rescued tree are retained; the floor is not
substituted for their incidence.
\end{firewall}

\section{Uniform raw norm of the complete word}
\label{sec:V95-raw}

\begin{lemma}[Raw \(2+1+1\) contraction bound]
\label{lem:V95-raw-bound}
For every fixed resolved quotient tree, payer placement, and
output-duality orientation,
\begin{equation}
 \opnorm{W_{r,\pi}^{\rm raw}}\le C_{G,s},
\qquad
 \opnorm{\widetilde W_{r,\pi}}
 \le {C_{G,s}\over\Xi_1\Xi_2\Xi_3\Xi_4}.
\label{eq:V95-raw-bound}
\end{equation}
The constant is uniform in support, representations, and physical
multiplicities.
\end{lemma}

\begin{proof}
The binary prefix cumulant is a difference of two complete
contractions.  The assembled joining third cumulant is a fixed
moment--cumulant polynomial and has a representation-independent raw
norm.  Singleton prefixes, the complete suffix, and physical
compressions are contractions.  The one payer is bounded uniformly by
the paid V67/V77 minimum when it is in a connected factor and by the
bankwise one-resolvent debit when it is in a singleton or suffix.
Complete pinching is performed before any triangle inequality, so no
multiplicity count appears.  Divide by the four row masses.
\end{proof}

\section{Diffuse maximal-row closure}
\label{sec:V95-diffuse}

\begin{lemma}[A diffuse maximal row closes against the local floor]
\label{lem:V95-diffuse-close}
Suppose row \(a\) has \(\Xi_a=M_*\to\infty\) and
\begin{equation}
 \max_e{a_{a,e}\over\Xi_a}\longrightarrow0.
\label{eq:V95-max-diffuse}
\end{equation}
Then every resolved oriented square is bounded by
\(C_{G,s}\mathfrak d_{211,r}^{\,2}\).
\end{lemma}

\begin{proof}
Delete the selected and payer coordinates and the fixed finite marked
packet.  Their total row-\(a\) mass is \(o(M_*)\).  Apply
\Cref{thm:V84-universal-incidence} to the remaining support in the
same exact word.  One cell of its fixed finite endpoint/suffix
resolution is a source-occurring bank \(D\) owning
\(S_a(D)\ge\sigma M_*\) and retaining
\(\max_{e\in D}a_{a,e}/M_*\to0\).

The all-payer central factorization of
\Cref{thm:V87-no-total-depletion} inserts the complete heat square on
\(D\) without changing its moment-partition role.  By
\Cref{thm:V87-diffuse-bank-power} with \(k=24\),
\[
 \kappa_{a|a^c}(\mathcal M_D^2)
 \le C_{G,s,\sigma}M_*^{-12}.
\]
All companion factors have the raw bounds of
\Cref{lem:V95-raw-bound}; the four-row normalization is at most one.
Thus the complete oriented word square is at most \(CM_*^{-12}\).
Use \Cref{lem:V95-local-floor}.
\end{proof}

\section{A cofinal balanced atom forces a strong edge}
\label{sec:V95-cofinal}

\begin{lemma}[Maximal-row atom dichotomy]
\label{lem:V95-atom-dichotomy}
Suppose instead that
\eqref{eq:V95-max-diffuse} fails.  After a subsequence there are
\(\sigma>0\) and coordinates \(e\) such that
\begin{equation}
 a_{a,e}\ge\sigma M_*.
\label{eq:V95-cofinal-max}
\end{equation}
Then either:
\begin{enumerate}[label=\textup{(\alph*)},leftmargin=4em]
\item the local tuple at \(e\) is V32-unbalanced and the word closes
      by the existing gapped/unbalanced compiler; or
\item there are a fixed row \(c\ne a\) and \(\lambda>0\) with
\begin{equation}
 a_{c,e}\ge\lambda a_{a,e},
 \qquad
 \Xi_c\ge\lambda\sigma M_*,
 \qquad
 \theta_{ac}\ge\lambda\sigma^2.
\label{eq:V95-strong-edge}
\end{equation}
\end{enumerate}
\end{lemma}

\begin{proof}
Failure of diffuseness gives
\eqref{eq:V95-cofinal-max} after a subsequence.  If every other local
row weight is \(o(a_{a,e})\), compact-group Casimir comparison makes
the largest highest weight exceed the sum of the remaining highest
weights; this is the V32-unbalanced alternative used in
\Cref{cor:V85-heavy-composite-closed}.

Otherwise finite row selection gives
\(a_{c,e}\ge\lambda a_{a,e}\) for some fixed \(c,\lambda\).
The first two inequalities in \eqref{eq:V95-strong-edge} follow
immediately.  Also
\[
 H_{ac}\ge a_{a,e}a_{c,e}
 \ge\lambda\sigma^2M_*^2.
\]
Since \(\Xi_a=M_*\) and \(\Xi_c\le M_*\), division by
\(\Xi_a\Xi_c\le M_*^2\) gives the last inequality.
\end{proof}

\begin{corollary}[The unbalanced atom closes at the local-floor
scale]
\label{cor:V95-unbalanced-close}
Alternative \textup{(a)} of
\Cref{lem:V95-atom-dichotomy} satisfies the
\(\mathfrak d_{211,r}^{\,2}\) square card for every payer placement.
\end{corollary}

\begin{proof}
The V34/V35/V38/V47 compilers, assembled in
\Cref{cor:V85-heavy-composite-closed}, apply to joining, payer,
separator, and coincidence roles.  They either give the literal
marked-tree bound, which is a term of
\(\mathfrak d_{211,r}\), or a cofinal heat-gap bound
\(C_{G,s}e^{-cM_*}\).  In the latter case
\(e^{-cM_*}\le C M_*^{-12}\), and
\Cref{lem:V95-local-floor} applies.
\end{proof}

\section{A tree containing the strong edge}
\label{sec:V95-tree-routing}

Let \(Q_r\) be the selected two-edge quotient tree on the three
blocks
\[
 A=\{1,2\},\qquad B=\{3\},\qquad C=\{4\}.
\]
Its two row representatives, together with \(12\), form the source
tree \(\tau_r\).

\begin{lemma}[One actual joining edge survives strong-edge rerouting]
\label{lem:V95-tree-routing}
For every row pair \(ac\):
\begin{enumerate}[label=\textup{(\roman*)},leftmargin=4em]
\item if \(ac=12\), the tree consisting of the global edge \(12\)
      and both actual joining edges of \(Q_r\) is one of
      \(\mathcal E_{211}\);
\item if \(ac\ne12\), there is an actual joining edge
      \(h\in Q_r\) such that
\begin{equation}
 \tau'=\{12,ac,h\}
\label{eq:V95-rerouted-tree}
\end{equation}
is a spanning tree on \([4]\).
\end{enumerate}
The corresponding monomial occurs in
\(\mathfrak d_{211,r}\), with \(ac\) global, \(h\) local at \(r\),
and \(12\) either prefix-local or global.
\end{lemma}

\begin{proof}
Contract the prefix block \(\{1,2\}\).  If \(ac=12\), the statement
is the forest--quotient lift proved in frozen
\(\texttt{thm:V63-forest-quotient}\).

Otherwise \(ac\) represents one edge of the three-vertex quotient
graph on \(A,B,C\).  If that quotient edge belongs to \(Q_r\), take
the other edge of \(Q_r\) as \(h\).  If it does not, take either edge
of \(Q_r\): any two distinct edges of a triangle form a quotient
tree.  Lifting the two quotient edges and adjoining the internal edge
\(12\) gives a four-row tree by
frozen \(\texttt{thm:V63-forest-quotient}\).  The monomial is
respectively the
second, third/fourth, or fifth/sixth term of
\eqref{eq:V94-tree-rescue}.
\end{proof}

\begin{lemma}[Two actual edge floors]
\label{lem:V95-two-floors}
In case \textup{(i)} of \Cref{lem:V95-tree-routing}, the two local
joining edges \(f,g\) satisfy
\[
 \theta_{f,r}\ge {a_*^2\over\Xi_u\Xi_v},
 \qquad
 \theta_{g,r}\ge {a_*^2\over\Xi_x\Xi_y}.
\]
In case \textup{(ii)}, the prefix edge and retained joining edge
satisfy
\[
 \theta_{12}^{<r}\ge {a_*^2\over\Xi_1\Xi_2},
 \qquad
 \theta_{h,r}\ge {a_*^2\over\Xi_x\Xi_y}.
\]
\end{lemma}

\begin{proof}
Every displayed source edge is nonzero in the selected resolved word,
so its raw stock contains an active coordinate product at least
\(a_*^2\).  Divide by the endpoint row masses.
\end{proof}

\begin{lemma}[Forest ratio bound]
\label{lem:V95-forest-ratio}
Let \(\tau'\) be a four-vertex tree containing a distinguished edge
\(ac\), and suppose
\begin{equation}
 \Xi_a,\Xi_c\ge m_0M_*,
 \qquad
 \Xi_k\le M_*\quad(1\le k\le4).
\label{eq:V95-two-maximal-endpoints}
\end{equation}
If \(f,g\) are the other two edges, then
\begin{equation}
 {\displaystyle\prod_{uv\in\{f,g\}}\Xi_u\Xi_v
  \over\Xi_1\Xi_2\Xi_3\Xi_4}
 \le m_0^{-1}.
\label{eq:V95-forest-ratio}
\end{equation}
\end{lemma}

\begin{proof}
Delete \(ac\) from \(\tau'\).  If the remaining two-edge forest is a
matching, the numerator is exactly the denominator.  Otherwise it is
a two-edge path plus one isolated vertex.  The ratio is the mass of
the path centre divided by the mass of the isolated vertex.  The
isolated vertex is one endpoint of the deleted edge \(ac\), hence its
mass is at least \(m_0M_*\), while the path-centre mass is at most
\(M_*\).
\end{proof}

\begin{theorem}[A strong maximal-row edge closes by raw contraction]
\label{thm:V95-strong-close}
Alternative \textup{(b)} of
\Cref{lem:V95-atom-dichotomy} satisfies, for every resolved word,
\begin{equation}
 \boxed{\;
 \opnorm{\widetilde W_{r,\pi}}
 \le C_{G,s,\sigma,\lambda}\mathfrak d_{211,r}.\;}
\label{eq:V95-strong-close}
\end{equation}
Hence both oriented square cards hold.
\end{theorem}

\begin{proof}
By \eqref{eq:V95-strong-edge}, both endpoints \(a,c\) have total mass
at least \(m_0M_*\), with
\(m_0=\min\{1,\lambda\sigma\}\), and
\(\theta_{ac}\ge\lambda\sigma^2\).

Choose the tree \(\tau'\) from
\Cref{lem:V95-tree-routing}.  Its two nondistinguished edges have
the actual lower bounds of \Cref{lem:V95-two-floors}.  Therefore the
corresponding term of \(\mathfrak d_{211,r}\) is at least
\[
 (\lambda\sigma^2)a_*^4
 {\vphantom{\prod}\over
  \displaystyle\prod_{uv\in\tau'\setminus\{ac\}}\Xi_u\Xi_v}.
\]
\Cref{lem:V95-forest-ratio} shows that this is at least
\[
 {c_{\sigma,\lambda}\over
  \Xi_1\Xi_2\Xi_3\Xi_4}.
\]
Now apply the normalized raw bound
\eqref{eq:V95-raw-bound}.  Squaring ordinary norm bounds either
oriented positive capacity.
\end{proof}

\section{The analytic two-rescue card}
\label{sec:V95-card}

\begin{theorem}[Uniform \(TRFCQ_{211}\)]
\label{thm:V95-TRFCQ211}
For every compact connected semisimple \(G\), admissible heat
parameter, joining coordinate, payer placement, and physical output
resolution,
\begin{align}
 \kappa_\otimes(
  \widetilde W_{r,\pi}^*\widetilde W_{r,\pi})
 &\le C_{G,s}\mathfrak d_{211,r}^{\,2},
\label{eq:V95-TRFCQ-square}\\
 \kappa_\otimes(
  \widetilde W_{r,\pi}\widetilde W_{r,\pi}^*)
 &\le C_{G,s}\mathfrak d_{211,r}^{\,2}.
\label{eq:V95-TRFCQ-left-square}
\end{align}
\end{theorem}

\begin{proof}
If no uniform constant exists, select a violating sequence and then
a fixed row \(a\) with \(\Xi_a=M_*\).  If \(M_*\) stays bounded,
\Cref{lem:V95-raw-bound,lem:V95-local-floor} close the sequence.
Thus \(M_*\to\infty\).

After a subsequence, either
\eqref{eq:V95-max-diffuse} holds or there is a cofinal atom.  The
diffuse case is \Cref{lem:V95-diffuse-close}.  In the atom case,
\Cref{lem:V95-atom-dichotomy} is exhaustive.  Its unbalanced branch
is \Cref{cor:V95-unbalanced-close}; its balanced branch is
\Cref{thm:V95-strong-close}.  Every alternative contradicts the
selected violation.  The argument is invariant under adjoint and
therefore proves both orientations.
\end{proof}

\begin{theorem}[Complete global \(2+1+1\) source]
\label{thm:V95-global-S211}
The exact source \(\mathcal S_{211}\) in
\eqref{eq:V35-S211}, with the sole payer allocated once, satisfies
\begin{equation}
 \boxed{\;
 \kappa_\otimes^{\rm abs}
  (\mathcal S_{211}^{\rm pay})
 \le C_{G,s}\mathscr A_{211}.\;}
\label{eq:V95-global-S211}
\end{equation}
The same holds for all six labelled choices of prefix pair.
\end{theorem}

\begin{proof}
\Cref{thm:V95-TRFCQ211} is the card
\Cref{def:V94-TRFCQ211}.  Apply
\Cref{prop:V94-TRFCQ-suffices} and
\Cref{lem:V94-rescue-summation}.  Row permutation preserves the
finite list and every constant.
\end{proof}

\begin{corollary}[Thirteen proper-prefix labels are globally closed]
\label{cor:V95-thirteen}
All three \(2+2\), all four \(3+1\), and all six \(2+1+1\)
first-connectivity sources satisfy their complete global normalized
tree estimates.  The sole untreated proper-prefix label is the
discrete partition \(1|2|3|4\).
\end{corollary}

\begin{proof}
Combine
\Cref{cor:V34-S22-no-gap,thm:V93-global-S31,
thm:V95-global-S211}
with the prefix census \Cref{lem:V31-prefix-census}.
\end{proof}

\begin{firewall}[One owner for the rescuing edge]
\label{fire:V95-one-owner}
In the balanced atom branch, the strong edge \(ac\) is used only in
the rescued tree monomial.  The proof then bounds the remaining
operator by its raw contraction; it does not also claim heat decay
from the same protected atom.  This is the one-owner rule required by
the SM audit.
\end{firewall}

\section{V95 result ledger and V96 acquisition order}
\label{sec:V95-status}

\begin{theorem}[Third-series V95 result ledger]
\label{thm:V95-results}
Relative to V1--V94, V95 proves:
\begin{enumerate}[label=\textup{(V95.\arabic*)},leftmargin=6.3em]
\item the compulsory SM V72/NS V31 audit with exact file identities;
\item an \(M_*^{-6}\) local floor for the two-rescue atlas;
\item a uniform raw norm for complete fixed-payer \(2+1+1\) words;
\item closure of a diffuse maximal row by the V87 \(k=24\) bank
      theorem;
\item a cofinal-atom dichotomy between known unbalance and a strong
      normalized edge with two maximal-scale endpoints;
\item an exact tree rerouting retaining an actual joining edge;
\item the forest mass-ratio lemma and raw-contraction closure of the
      strong-edge branch;
\item the uniform two-rescue oriented square card;
\item the complete global \(2+1+1\) source theorem for all six
      labels; and
\item global closure of thirteen of the fourteen proper-prefix labels.
\end{enumerate}
\end{theorem}

\begin{proof}
These are
\Cref{ass:V95-SM-audit,
ass:V95-NS-audit,
lem:V95-local-floor,
lem:V95-raw-bound,
lem:V95-diffuse-close,
lem:V95-atom-dichotomy,
lem:V95-tree-routing,
lem:V95-forest-ratio,
thm:V95-strong-close,
thm:V95-TRFCQ211,
thm:V95-global-S211,
cor:V95-thirteen}.
\end{proof}

\begin{programme}[V96: the discrete-prefix source]
\label{prog:V95-V96}
\begin{enumerate}[leftmargin=3em]
\item Return to the exact first-connectivity identity for prefix
      \(1|2|3|4\).  Its joining carrier is the complete local fourth
      cumulant and must remain assembled.
\item Enumerate the local quotient-tree branches and their payer
      placements without expanding the cumulant into independent
      partition-letter norms.
\item Define the minimal globally rescued weight which retains at
      least one actual joining-coordinate edge and prove its scalar
      summability.
\item Establish a local floor for that weight.  Then repeat the
      maximal-row diffuse/cofinal-atom dichotomy, using raw
      contraction only after a strong edge has been assigned to a
      legal spanning tree.
\item Check whether one local fourth cumulant can create all three
      tree edges at the same coordinate without a new positive-square
      interface.
\item After the discrete source is settled, audit whether these
      fourteen proper-prefix theorems plus the already closed
      selected-\([4]\) theorem assemble the complete quartic MTA.
\end{enumerate}
\end{programme}

\begin{firewall}[Claim boundary after third-series V95]
\label{fire:V95-claim-boundary}
V95 closes all six \(2+1+1\) sources globally and leaves only the
discrete proper prefix.  It does not close that source, globally
aggregate every \([4]\) term, close \([3,3]\), higher MTA, WUFC, CPM,
cutoff removal, reflection positivity, Osterwalder--Schrader
reconstruction, construct the continuum Yang--Mills measure, or
prove a uniform positive continuum spectral gap.  The full
Yang--Mills mass gap has not been proved.
\end{firewall}

\paragraph{Parent-version record.}
V95 was formed from
\texttt{Renewal\_Yang\_Mills\_Mass\_Gap\_Research\_Notes\_V94.tex}.
The V94 parent had \(45\,482\) logical source lines,
\(1\,612\,453\) bytes, and SHA--256
\[
 \texttt{56ac567788f12f6c5062825f5de7c69669fdbfe3e5fef962be0e1892f2432a7a}.
\]
The next compulsory companion audit is V100.

\clearpage
\part{Third-series V96: the discrete first-join source and
restoration of global quartic MTA}
\label{part:V96}

\section{Why the last prefix needs no new local heat theorem}
\label{sec:V96-purpose}

The only proper four-row prefix not covered after V95 is the discrete
partition
\[
 \widehat0=1|2|3|4.
\]
Immediately before its first global join there are no earlier
connected blocks.  The joining coordinate must therefore carry the
complete local fourth cumulant.  This is exactly the selected local
object covered for every compact connected semisimple group by
\Cref{thm:V38-GSEL4}.

The V38 firewall correctly declined to infer a multiple-selection
theorem merely from uniformity in the selected coordinate.  The
missing summation is simple once the exact first-join decomposition
is written: each local fourth-cumulant tree carries three marks at the
same coordinate, and its diagonal coordinate sum is bounded by the
Cartesian product of the three global edge stocks.

\begin{assessment}[Nonduplication boundary]
\label{ass:V96-nonduplication}
V96 does not repeat the selected-junction analysis of V38.  It proves
the discrete-prefix identity, performs the previously missing
coordinate-weighted summation, and combines that family with the
thirteen source theorems already proved.  The resulting assembly is a
new proof of the complete global quartic marked-tree atlas.
\end{assessment}

\section{Exact discrete-prefix identity}
\label{sec:V96-identity}

Let \(M_{\{i\}}^{<r}\) be the complete singleton moment of row \(i\)
on coordinates before \(r\), let
\[
 K_{4,r}:=\kappa_4(X_{1,r},X_{2,r},X_{3,r},X_{4,r}),
\]
with every local moment-partition letter assembled, and let
\(M_{>r}\) be the complete four-row suffix.

\begin{theorem}[Exact discrete first-connectivity source]
\label{thm:V96-discrete-identity}
Before any norm, positive envelope, Fourier-output resolution, or
payer allocation, the discrete-prefix source is
\begin{equation}
 \boxed{\;
 \mathcal S_{\widehat0}
 =
 \sum_r
 \left(\bigotimes_{i=1}^4M_{\{i\}}^{<r}\right)
 \otimes K_{4,r}\otimes M_{>r}.\;}
\label{eq:V96-discrete-source}
\end{equation}
Every connected coordinate word whose cumulative partition is
discrete immediately before its first full join occurs exactly once.
\end{theorem}

\begin{proof}
Apply the coefficient-one linked-word and first-connectivity
telescope \Cref{prop:V7-linked-word} to a connected four-row word.
In the present family, the cumulative join of every coordinate
strictly before \(r\) is \(\widehat0\).  Hence no two rows share a
prefix replica block, and the prefix factors into the four complete
singleton moments.

At the first joining coordinate,
\(\widehat0\vee\pi_r=\widehat1\) implies
\(\pi_r=\widehat1\).  M\"obius inversion on the local partition
lattice identifies its connected coefficient with the complete
fourth cumulant \(K_{4,r}\).  After that coordinate the word carries
the unrestricted complete moment suffix.  Uniqueness of the first
full-join coordinate gives disjointness and coefficient one.
\end{proof}

\begin{firewall}[The local cumulant remains assembled]
\label{fire:V96-assembled}
Equation \eqref{eq:V96-discrete-source} does not expand \(K_{4,r}\)
into absolute moment-partition letters.  The V38 selected-junction
theorem is applied to the complete signed cumulant with its physical
multiplicity spaces and one payer intact.
\end{firewall}

\section{The local sixteen-tree weight}
\label{sec:V96-local-weight}

Put
\begin{equation}
 \theta_{ij,r}:={h_{ij,r}\over\Xi_i\Xi_j},
\label{eq:V96-local-theta}
\end{equation}
and let \(\mathfrak T_4\) be the sixteen labelled spanning trees on
\([4]\).  Define
\begin{align}
 \mathfrak t_{4,r}
 &:=
 \sum_{\tau\in\mathfrak T_4}
 \prod_{ij\in E(\tau)}\theta_{ij,r},
\label{eq:V96-local-tree-weight}\\
 \mathscr T_4
 &:=
 \sum_{\tau\in\mathfrak T_4}
 \prod_{ij\in E(\tau)}\theta_{ij}.
\label{eq:V96-global-tree-weight}
\end{align}

\begin{lemma}[Diagonal three-mark summation]
\label{lem:V96-diagonal-sum}
\begin{equation}
 \boxed{\;
 \sum_r\mathfrak t_{4,r}\le\mathscr T_4.\;}
\label{eq:V96-diagonal-sum}
\end{equation}
\end{lemma}

\begin{proof}
Fix \(\tau\in\mathfrak T_4\), with edges \(f_1,f_2,f_3\).
Nonnegativity gives
\[
 \sum_r
 \theta_{f_1,r}\theta_{f_2,r}\theta_{f_3,r}
 \le
 \left(\sum_{r_1}\theta_{f_1,r_1}\right)
 \left(\sum_{r_2}\theta_{f_2,r_2}\right)
 \left(\sum_{r_3}\theta_{f_3,r_3}\right)
 =
 \prod_{f\in\tau}\theta_f.
\]
Sum over the fixed sixteen trees.  No coordinate count appears.
\end{proof}

\begin{remark}[Why sixteen is fixed]
\label{rem:V96-sixteen}
Cayley's formula gives \(4^{4-2}=16\).  This is an arity constant,
not a support or representation multiplicity.
\end{remark}

\section{Insertion of the selected-junction theorem}
\label{sec:V96-GSEL}

For a fixed payer placement \(\pi\), let
\(\widetilde W_{r,\pi}^{\widehat0}\) be the normalized \(r\)-summand
of \eqref{eq:V96-discrete-source}.

\begin{theorem}[Coordinate-weighted selected fourth-cumulant card]
\label{thm:V96-weighted-GSEL}
For every fixed compact connected semisimple \(G\) and admissible heat
parameter,
\begin{equation}
 \boxed{\;
 \kappa_\otimes^{\rm abs}
  (\widetilde W_{r,\pi}^{\widehat0})
 \le C_{G,s}\mathfrak t_{4,r},\;}
\label{eq:V96-weighted-GSEL}
\end{equation}
uniformly in \(r\), support, representations, physical
multiplicities, and payer placement.
\end{theorem}

\begin{proof}
The \(r\)-summand has one first-connectivity-selected complete
four-row local cumulant, complete singleton prefix spectators,
a complete positive moment suffix, exact physical output
compression, and one allocated final payer.  It is therefore an
admissible selected-junction packet in the sense of
\Cref{def:V38-admissible-selected}.

\Cref{thm:V38-GSEL4} gives its literal coefficient-sensitive
marked-tree bound.  In the discrete-prefix case all three tree edges
belong to the selected local fourth cumulant, so its literal
normalized right side is exactly
\(\mathfrak t_{4,r}\).  The theorem already includes every payer
location and retains complete multiplicities.
\end{proof}

\begin{firewall}[What resolves the V38 aggregation warning]
\label{fire:V96-V38-warning}
Uniformity of \eqref{eq:V96-weighted-GSEL} in \(r\) is not used by
itself.  The coordinate-dependent right side is summed by
\Cref{lem:V96-diagonal-sum}.  This is the weighted input absent from
the multiple-selection inference forbidden by
\Cref{fire:V38-no-global-selection}.
\end{firewall}

\begin{theorem}[Complete global discrete-prefix source]
\label{thm:V96-global-discrete}
With the sole final payer allocated exactly once,
\begin{equation}
 \boxed{\;
 \kappa_\otimes^{\rm abs}
  (\mathcal S_{\widehat0}^{\rm pay})
 \le C_{G,s}\mathscr T_4.\;}
\label{eq:V96-global-discrete}
\end{equation}
\end{theorem}

\begin{proof}
Use the exact decomposition
\eqref{eq:V96-discrete-source}, subadditivity of absolute
product-state capacity, and
\Cref{thm:V96-weighted-GSEL}:
\[
 \kappa_\otimes^{\rm abs}
  (\mathcal S_{\widehat0}^{\rm pay})
 \le
 C_{G,s}\sum_r\mathfrak t_{4,r}.
\]
Apply \Cref{lem:V96-diagonal-sum} and sum over the fixed finite payer
types.
\end{proof}

\section{Exhaustive assembly of the quartic carrier}
\label{sec:V96-assembly}

Let \(\Pi_4^\circ:=\Pi([4])\setminus\{\widehat1\}\).
Its fourteen partitions have block-size types
\begin{center}
\begin{tabular}{c|cccc}
type & \(3+1\) & \(2+2\) & \(2+1+1\) & \(1+1+1+1\)\\
\hline
number & \(4\) & \(3\) & \(6\) & \(1\).
\end{tabular}
\end{center}

\begin{lemma}[Exact fourteen-family partition]
\label{lem:V96-fourteen-family}
Every connected four-row coordinate word has a unique first full-join
coordinate \(r\) and a unique cumulative partition
\(\rho\in\Pi_4^\circ\) immediately before \(r\).  The resulting
fourteen source families are disjoint and exhaustive, with coefficient
one.
\end{lemma}

\begin{proof}
This is the first-connectivity uniqueness in
\Cref{thm:V96-discrete-identity}, without fixing \(\rho\).
The cumulative join is monotone in coordinate order.  Immediately
before its first value \(\widehat1\), it is a unique proper partition.
Conversely every connected word reaches \(\widehat1\) once for the
first time.  The block-size census is the elementary set-partition
count displayed above.
\end{proof}

\begin{theorem}[Restored global quartic marked-tree atlas]
\label{thm:V96-global-quartic-MTA}
For every fixed compact connected semisimple \(G\), admissible heat
parameter, finite support, and trace-class row-factorized Fourier
coefficients, the complete once-paid connected four-row physical
carrier \(\mathcal C_4^{\rm pay}\) satisfies
\begin{equation}
 \boxed{\;
 \kappa_\otimes^{\rm abs}(\mathcal C_4^{\rm pay})
 \le
 C_{G,s}
 \sum_{\tau\in\mathfrak T_4}
 \prod_{ij\in E(\tau)}\theta_{ij}.\;}
\label{eq:V96-global-quartic-MTA}
\end{equation}
The constant is uniform in support, input representations, physical
output multiplicities, and payer placement.
\end{theorem}

\begin{proof}
Decompose the connected carrier by
\Cref{lem:V96-fourteen-family}.  For the three \(2+2\) partitions use
\Cref{cor:V34-S22-no-gap}.  For the four \(3+1\) partitions use
\Cref{thm:V93-global-S31}.  For the six \(2+1+1\) partitions use
\Cref{thm:V95-global-S211}.  For the discrete partition use
\Cref{thm:V96-global-discrete}.

Each family theorem has a finite sum of literal spanning-tree
monomials on the right.  Enlarge it to the complete sixteen-tree sum
\(\mathscr T_4\).  There are fourteen families and finitely many
payer/output-duality types, so their sum changes only
\(C_{G,s}\).  Exact coefficient-one disjointness prevents omission or
multiple counting on the left.
\end{proof}

\begin{corollary}[All quartic coordinate topologies are included]
\label{cor:V96-all-quartic-topologies}
The global theorem
\eqref{eq:V96-global-quartic-MTA} includes the minimal coordinate
topologies
\[
 [4],\qquad[3,2],\qquad[3,3],\qquad[2,2,2],
\]
including collisions, repeated joining coordinates, and every payer
coincidence.  No separate unaggregated \([4]\) or \([3,3]\) family
remains at order four.
\end{corollary}

\begin{proof}
Every connected four-row word, irrespective of its minimal
hypergraph topology, belongs to one first-connectivity family by
\Cref{lem:V96-fourteen-family}.  The family theorems use complete
local cumulants and complete moment suffixes, so coincident local
partition letters and multiplicity spaces are retained.  The four
listed topologies exhaust the connected four-row skeletons recorded
in the archive map, hence all are included.
\end{proof}

\begin{assessment}[Relation to the withdrawn earlier quartic claim]
\label{ass:V96-restoration}
Frozen f2 V52 stated a global quartic MTA theorem before later
chapters found that its separated source proof lost attachment-row
normalization and weighted coordinate aggregation.  V89--V95 repair
those precise defects, including physical counterpackets to the two
overstrong local cards.  \Cref{thm:V96-global-quartic-MTA} is therefore
a restored theorem with a different proof: exact first-connectivity,
globally rescued but summable weights, and the already valid selected
fourth-cumulant interface.
\end{assessment}

\begin{firewall}[Quartic MTA is not the mass gap]
\label{fire:V96-not-gap}
The theorem closes the normalized connected carrier only at fixed
row order four.  The Yang--Mills mass-gap route still requires
higher/all-order MTA with controlled arity constants, WUFC, CPM,
uniform cutoff estimates, continuum construction, reflection
positivity, Osterwalder--Schrader reconstruction, and a positive
continuum spectral gap.
\end{firewall}

\section{V96 result ledger and V97 acquisition order}
\label{sec:V96-status}

\begin{theorem}[Third-series V96 result ledger]
\label{thm:V96-results}
Relative to V1--V95, V96 proves:
\begin{enumerate}[label=\textup{(V96.\arabic*)},leftmargin=6.3em]
\item the exact discrete first-connectivity identity;
\item the local sixteen-tree normalized weight;
\item support-uniform diagonal summation of its three local marks;
\item direct insertion of the general-\(G\) selected-junction theorem;
\item the complete global discrete-prefix source;
\item the disjoint and exhaustive fourteen-family partition;
\item restoration of the complete global quartic marked-tree atlas;
      and
\item inclusion of every quartic coordinate topology and coincidence
      sector.
\end{enumerate}
\end{theorem}

\begin{proof}
These are
\Cref{thm:V96-discrete-identity,
lem:V96-diagonal-sum,
thm:V96-weighted-GSEL,
thm:V96-global-discrete,
lem:V96-fourteen-family,
thm:V96-global-quartic-MTA,
cor:V96-all-quartic-topologies}.
\end{proof}

\begin{programme}[V97: exact quintic first-connectivity census]
\label{prog:V96-V97}
\begin{enumerate}[leftmargin=3em]
\item Return to the frozen quintic record and separate proved local
      fifth-cumulant interfaces from source-level claims which used an
      unweighted coordinate sum.
\item Enumerate the \(B_5-1=51\) proper prefix partitions by block
      sizes:
      \(4+1\) \((5)\), \(3+2\) \((10)\),
      \(3+1+1\) \((10)\), \(2+2+1\) \((15)\),
      \(2+1+1+1\) \((10)\), and the discrete partition \((1)\).
\item Write the exact quotient-cumulant first-join source for each
      block-size type, using the now-proved lower connected carriers
      through order four.
\item Determine which types need zero, one, or two global rescue edges
      while retaining at least one actual joining mark.
\item Prove scalar summability before attempting any fifth-order
      capacity estimate.
\item Keep arity dependence explicit; a fixed order-five constant is
      not yet an all-order compiler.
\end{enumerate}
\end{programme}

\begin{firewall}[Claim boundary after third-series V96]
\label{fire:V96-claim-boundary}
V96 restores complete global quartic MTA.  It does not prove quintic
or all-order MTA, WUFC, CPM, cutoff removal, reflection positivity,
Osterwalder--Schrader reconstruction, construction of the continuum
Yang--Mills measure, or a uniform positive continuum spectral gap.
The full Yang--Mills mass gap has not been proved.
\end{firewall}

\paragraph{Parent-version record.}
V96 was formed from
\texttt{Renewal\_Yang\_Mills\_Mass\_Gap\_Research\_Notes\_V95.tex}.
The V95 parent had \(46\,044\) logical source lines,
\(1\,631\,937\) bytes, and SHA--256
\[
 \texttt{220d56bc574d0a89bc0b04b1cfa42ce3f6fd4a309c0487ce8fbe567bc6760641}.
\]
The next compulsory companion audit is V100.

% =============================================================================
% BEGIN APPENDED THIRD-SERIES V97 MATERIAL
% =============================================================================

\clearpage
\part{Third-series V97: exact quintic first-connectivity census,
summable rescue atlas, and the one-strong-edge deficit}
\label{part:V97}

\section{Purpose and archive reconciliation}
\label{sec:V97-purpose}

V96 restored the complete global marked-tree atlas at row order four.
The next order is not obtained by copying the old quintic
graph-difference argument.  The frozen file
\texttt{Renewal\_Yang\_Mills\_Mass\_Gap\_Research\_Notes\_V100\_f2.tex}
contains both valid quintic inputs and a route which was subsequently
shown to leak.  This chapter first separates them and then restarts the
global quintic source analysis from exact first connectivity.

\begin{assessment}[Status of the frozen quintic record]
\label{ass:V97-frozen-status}
The following frozen results remain valid inputs.
\begin{enumerate}[label=\textup{(\roman*)},leftmargin=4em]
\item V53 computes the leading fifth Wilson cumulant and proves that a
      one-Casimir-per-input estimate is false.
\item \texttt{thm:V56-uniform-one-link-quintic} decomposes the complete
      one-coordinate fifth cumulant into its \(125\) labelled tree
      fibres with representation-uniform tree-valence bounds.
\item V57 proves the one-payer forest-plus-quotient-tree compiler,
      conditional on obtaining the quotient marks without splitting an
      assembled cancellation.
\item V58 gives a correct six-orbit partition-lattice identity, including
      three genuinely unicyclic orbits.
\end{enumerate}
However, \texttt{prop:V59-path-leakage} proves that an individual path
difference contains the endpoint covariance times the means of the
internal vertices.  Thus a graph edge of the V58 identity does not, by
itself, deliver a Wilson edge mark.  The conditional statewise
edge-marking step after V58 is not available.  None of the graphwise
orbit capacities is used below.
\end{assessment}

\begin{firewall}[The route being retired]
\label{fire:V97-no-graphwise-route}
The exact V58 identity is not false; the attempted analytic
interpretation of every one of its displayed merge differences as an
independent normalized pair mark is false.  Taking norms of the six
orbits separately would destroy the cancellation exposed by V59.  The
replacement carrier is the complete quotient cumulant at the unique
first globally connecting coordinate.
\end{firewall}

\section{The exact fifty-one-family formula}
\label{sec:V97-exact-source}

Let \(I=[5]\), let \(E=\{1,\ldots,n\}\) be the ordered finite coordinate
set, and let \(\Pi_5\) be the partition lattice of \(I\).  As in V96, a
local cumulant letter at coordinate \(e\) is indexed by
\(\pi_e\in\Pi_5\), and the cumulative join through \(r\) is
\(j_r:=\bigvee_{e\le r}\pi_e\).

Fix a proper partition \(\rho<\widehat1\), write
\(k:=|\rho|\), and let \(\mathcal K_B^{<r,\mathrm{conn}}\) denote the
complete connected prefix carrier on a block \(B\in\rho\), using only
coordinates below \(r\).  Singleton blocks carry their complete
one-row prefix moments.  At coordinate \(r\), set
\begin{equation}
 Y_{B,r}:=\prod_{i\in B}F_{i,r},
 \qquad
 \mathcal J_r^\rho
 :=\kappa_k\bigl(Y_{B,r}:B\in\rho\bigr),
\label{eq:V97-quotient-cumulant}
\end{equation}
where the row order inside each displayed product is fixed once and for
all.  Let \(\mathcal M_{>r}\) be the complete local-moment suffix.

\begin{theorem}[Exact quintic first-connectivity source]
\label{thm:V97-exact-quintic-source}
The complete connected five-row carrier has the coefficient-one
decomposition
\begin{equation}
 \boxed{\;
 \mathcal C_5
 =
 \sum_{r=1}^{n}\ \sum_{\rho\in\Pi_5\setminus\{\widehat1\}}
 \left(\bigotimes_{B\in\rho}
       \mathcal K_B^{<r,\mathrm{conn}}\right)
 \otimes\mathcal J_r^\rho
 \otimes\mathcal M_{>r}. \;}
\label{eq:V97-exact-quintic-source}
\end{equation}
Every connected coordinate word occurs exactly once.  No graph-orbit
difference and no accepted-state propagator occurs in this formula.
\end{theorem}

\begin{proof}
Expand the complete coordinate product into local cumulant letters.
The row cumulant keeps exactly those words for which
\(\bigvee_e\pi_e=\widehat1\), by M\"obius inversion on \(\Pi_5\).
Such a word has a unique first index \(r\) with
\(j_r=\widehat1\), and its preceding join
\(\rho=j_{r-1}\) is proper.  This partitions the connected words
disjointly and with coefficient one.

For fixed \((r,\rho)\), a prefix word with join exactly \(\rho\)
contains no letter crossing two blocks of \(\rho\).  Its restrictions
to the blocks have full join in each block and vary independently.
M\"obius inversion inside each block therefore gives
\(\bigotimes_{B\in\rho}\mathcal K_B^{<r,\mathrm{conn}}\).

It remains to identify the joining letter.  Expand the right side of
\eqref{eq:V97-quotient-cumulant} by the moment--cumulant formula on the
\(k\) quotient labels \(B\in\rho\), and then expand each local moment
into row-partition cumulant letters.  A row partition \(\pi\) has
coefficient
\[
 \sum_{\sigma\in\Pi(\rho):\,q_\rho(\pi)\le\sigma}
       \mu(\sigma,\widehat1),
\]
where \(q_\rho(\pi)\) is the partition induced on the blocks of
\(\rho\) by \(\rho\vee\pi\).  The M\"obius cancellation identity makes
this coefficient one precisely when
\(q_\rho(\pi)=\widehat1\), equivalently
\(\rho\vee\pi=\widehat1\), and zero otherwise.  Hence
\(\mathcal J_r^\rho\) is exactly the sum of all admissible joining
letters, with their cancellation still assembled.  Later letters are
unrestricted and sum to \(\mathcal M_{>r}\), proving the formula.
\end{proof}

\begin{lemma}[Exact proper-prefix census]
\label{lem:V97-prefix-census}
The \(B_5-1=51\) proper partitions in
\eqref{eq:V97-exact-quintic-source} have the following block-size
types:
\begin{center}
\begin{tabular}{c|rrrrrr|r}
type
 & \(4+1\)&\(3+2\)&\(3+1+1\)&\(2+2+1\)
 &\(2+1+1+1\)&\(1+1+1+1+1\)&total\\
\hline
number&5&10&10&15&10&1&51.
\end{tabular}
\end{center}
\end{lemma}

\begin{proof}
Choose the singleton in type \(4+1\), the two-block in type \(3+2\),
the triple in type \(3+1+1\), the pair in type \(2+1+1+1\), and
nothing in the discrete type.  This gives respectively
\(5,\binom52,\binom53,\binom52,1\).  For type \(2+2+1\), choose the
singleton and then one of the three pairings of the remaining four
labels, giving \(5\cdot3=15\).  Their sum is \(51\).
\end{proof}

\begin{corollary}[Exact analytic interface by quotient order]
\label{cor:V97-interface-table}
The source families in \eqref{eq:V97-exact-quintic-source} require
exactly the following already identified interfaces:
\begin{center}
\begin{tabular}{c|c|c|c}
prefix type&\(k=|\rho|\)&prefix connected blocks&joining carrier\\
\hline
\(4+1\)&2&one quartic&complete covariance\\
\(3+2\)&2&one ternary, one binary&complete covariance\\
\(3+1+1\)&3&one ternary&complete third cumulant\\
\(2+2+1\)&3&two binary&complete third cumulant\\
\(2+1+1+1\)&4&one binary&complete fourth cumulant\\
\(1+1+1+1+1\)&5&none&complete fifth cumulant.
\end{tabular}
\end{center}
Thus V96 supplies the only newly required prefix theorem, namely the
complete quartic block in the \(4+1\) family.  The frozen uniform local
fifth-cumulant card applies only in the last row of the table.
\end{corollary}

\begin{proof}
This is immediate from
\(\mathcal J_r^\rho=\kappa_k(Y_{B,r}:B\in\rho)\) and the sizes of the
blocks of \(\rho\).  Every proper block has size at most four, and the
display lists all six types from \Cref{lem:V97-prefix-census}.
\end{proof}

\section{Forest--quotient lift without graph differences}
\label{sec:V97-lift}

For each nonsingleton \(B\in\rho\), choose a labelled spanning tree
\(T_B\) on \(B\).  Let \(\zeta\) be a tree on the \(k\) quotient
vertices \(B\in\rho\), and, for every \(BC\in E(\zeta)\), choose one
lift \(e_{BC}=ij\) with \(i\in B\), \(j\in C\).  Put
\begin{equation}
 P:=\bigcup_{B\in\rho}E(T_B),
 \qquad
 Q:=\{e_{BC}:BC\in E(\zeta)\},
 \qquad
 T:=P\sqcup Q .
\label{eq:V97-PQT}
\end{equation}

\begin{lemma}[Exact five-row tree lift]
\label{lem:V97-tree-lift}
The graph \(T\) in \eqref{eq:V97-PQT} is a labelled spanning tree on
\([5]\), and
\begin{equation}
 |P|=5-k,\qquad |Q|=k-1,\qquad |T|=4.
\label{eq:V97-edge-count}
\end{equation}
Conversely, contracting every block of \(\rho\) in any spanning tree
whose internal restriction on each block is connected produces such a
quotient tree.
\end{lemma}

\begin{proof}
The internal tree in \(B\) has \(|B|-1\) edges, hence
\[
 |P|=\sum_{B\in\rho}(|B|-1)=5-k.
\]
The quotient tree has \(k-1\) lifted edges.  Internal trees connect
each block, and the lifted quotient edges connect the blocks, so \(T\)
is connected.  It has five vertices and four edges and is therefore a
tree.  Conversely, contraction preserves connectedness.  Since the
original graph has four edges and the connected internal restrictions
already use \(5-k\), exactly \(k-1\) cross-block edges remain; the
connected quotient graph is therefore a tree.
\end{proof}

\begin{firewall}[What the lift does and does not prove]
\label{fire:V97-lift-scope}
\Cref{lem:V97-tree-lift} is applied to tree fibres of the complete
prefix carriers and the complete quotient cumulant.  It does not assign
marks to V58 merge differences.  It is a graph conversion after the
valid connected-cumulant tree cards have been invoked, not an
independent norm estimate.
\end{firewall}

\section{A universal finite rescue atlas}
\label{sec:V97-rescue}

For an edge \(e=ij\), retain the established nonnegative stocks
\(\theta_{e,r}\), \(\theta_e^{<r}\), and
\(\theta_e=\sum_s\theta_{e,s}\).  For fixed data
\((\rho,(T_B),\zeta,(e_{BC}))\), define its literal local weight
\begin{equation}
 \mathfrak l_{\rho,T,r}
 :=
 \prod_{e\in P}\theta_e^{<r}
 \prod_{e\in Q}\theta_{e,r}.
\label{eq:V97-literal-weight}
\end{equation}
The quartic counterpackets of V92 and V94 show why this literal weight
cannot be the final local square target: a heat suffix may supply a
strong copy of an edge already used by the source tree.

\begin{definition}[Quintic rescue weight]
\label{def:V97-rescue-weight}
The finite rescue weight attached to the lifted tree \(T=P\sqcup Q\)
at \(r\) is
\begin{equation}
 \boxed{\;
 \mathfrak r_{\rho,T,r}
 :=
 \sum_{A\subseteq P}
 \ \sum_{\varnothing\ne L\subseteq Q}
 \left(\prod_{e\in A}\theta_e\right)
 \left(\prod_{e\in P\setminus A}\theta_e^{<r}\right)
 \left(\prod_{e\in L}\theta_{e,r}\right)
 \left(\prod_{e\in Q\setminus L}\theta_e\right).\;}
\label{eq:V97-rescue-weight}
\end{equation}
Thus any prefix edge may be rescued globally and any joining edge may
be rescued globally, but at least one actual joining edge remains at
the first-connectivity coordinate.  The literal term is
\((A,L)=(\varnothing,Q)\).
\end{definition}

\begin{lemma}[Support-uniform rescue summation]
\label{lem:V97-rescue-summation}
For every fixed lifted tree,
\begin{equation}
 \boxed{\;
 \sum_r\mathfrak r_{\rho,T,r}
 \le
 2^{\,5-k}\bigl(2^{\,k-1}-1\bigr)
 \prod_{e\in T}\theta_e .\;}
\label{eq:V97-rescue-summation}
\end{equation}
The multipliers for \(k=2,3,4,5\) are respectively
\(8,12,14,15\).
\end{lemma}

\begin{proof}
Fix \(A\subseteq P\) and nonempty \(L\subseteq Q\).
All global factors may be taken outside the \(r\)-sum.  Since the
stocks are nonnegative,
\(\theta_e^{<r}\le\theta_e\), and diagonal summation is bounded by
Cartesian enlargement:
\[
 \sum_r
 \left(\prod_{e\in P\setminus A}\theta_e^{<r}\right)
 \left(\prod_{e\in L}\theta_{e,r}\right)
 \le
 \left(\prod_{e\in P\setminus A}\theta_e\right)
 \left(\prod_{e\in L}\theta_e\right).
\]
After restoring the factors indexed by \(A\) and \(Q\setminus L\),
the right side is \(\prod_{e\in T}\theta_e\).  There are
\(2^{|P|}(2^{|Q|}-1)\) choices.  Use
\eqref{eq:V97-edge-count}; the four displayed values follow.
\end{proof}

\begin{remark}[Why one local joining edge is compulsory]
\label{rem:V97-local-anchor}
If \(L=\varnothing\) were admitted, a summand of
\(\mathfrak r_{\rho,T,r}\) could be independent of \(r\); summing it
would cost the support size.  The local joining anchor is exactly what
makes \eqref{eq:V97-rescue-summation} support-uniform.
\end{remark}

\begin{definition}[Quintic tree-rescue square card]
\label{def:V97-QTRSC5}
Resolve one summand of
\eqref{eq:V97-exact-quintic-source} into fixed prefix-tree fibres,
one quotient-cumulant tree fibre, chosen endpoint lifts, physical
output blocks, a single payer placement, and write the normalized word
as \(\widetilde W_{\rho,T,r,\pi}\).  The card \(QTRSC_5\) is the pair
of estimates
\begin{align}
 \kappa_\otimes(
  \widetilde W_{\rho,T,r,\pi}^{*}
  \widetilde W_{\rho,T,r,\pi})
 &\le C_{G,s}\mathfrak r_{\rho,T,r}^{\,2},
\label{eq:V97-QTRSC5-right}\\
 \kappa_\otimes(
  \widetilde W_{\rho,T,r,\pi}
  \widetilde W_{\rho,T,r,\pi}^{*})
 &\le C_{G,s}\mathfrak r_{\rho,T,r}^{\,2},
\label{eq:V97-QTRSC5-left}
\end{align}
in the output-dual orientation required by that resolved word.
\end{definition}

\begin{proposition}[\(QTRSC_5\) is sufficient for global quintic MTA]
\label{prop:V97-QTRSC5-suffices}
If \eqref{eq:V97-QTRSC5-right}--%
\eqref{eq:V97-QTRSC5-left} hold for every member of the fixed finite
resolution, then the complete once-paid five-row carrier satisfies
\begin{equation}
 \boxed{\;
 \kappa_\otimes^{\rm abs}(\mathcal C_5^{\rm pay})
 \le C_{G,s}
 \sum_{\tau\in\mathfrak T_5}
 \prod_{e\in E(\tau)}\theta_e .\;}
\label{eq:V97-conditional-quintic-MTA}
\end{equation}
\end{proposition}

\begin{proof}
Product-state Cauchy--Schwarz applied in the required orientation gives
\[
 \kappa_\otimes^{\rm abs}
   (\widetilde W_{\rho,T,r,\pi})
 \le C_{G,s}\mathfrak r_{\rho,T,r}.
\]
Sum first in \(r\) and apply
\Cref{lem:V97-rescue-summation}.  There are only \(51\) prefix
partitions, finitely many prefix and quotient tree fibres at fixed order
five, finitely many endpoint lifts, and finitely many payer/output
duality types.  By \Cref{lem:V97-tree-lift}, every resulting monomial is
a five-row spanning-tree monomial.  Absorb the finite preimage count
into \(C_{G,s}\) and enlarge to the complete \(125\)-tree sum.
\end{proof}

\section{Polynomial floor and the diffuse terminal}
\label{sec:V97-diffuse}

Put \(\Xi_i\ge1\) for the five row masses and
\(M:=\max_i\Xi_i\).  On a fixed nonzero resolved source fibre, every
edge in \(P\) has an actual occurrence below \(r\), and every edge in
\(Q\) has an actual joining occurrence at \(r\).  The established
edge-occurrence floor gives, for a fixed \(a_*>0\),
\begin{equation}
 \theta_e^{<r}\ge {a_*^2\over\Xi_i\Xi_j}
 \quad(e=ij\in P),
 \qquad
 \theta_{e,r}\ge {a_*^2\over\Xi_i\Xi_j}
 \quad(e=ij\in Q).
\label{eq:V97-actual-edge-floor}
\end{equation}

\begin{lemma}[Fifth-order local rescue floor]
\label{lem:V97-rescue-floor}
Every nonzero lifted-tree fibre satisfies
\begin{equation}
 \boxed{\;
 \mathfrak r_{\rho,T,r}
 \ge\mathfrak l_{\rho,T,r}
 \ge {a_*^8\over M^8},
 \qquad
 \mathfrak r_{\rho,T,r}^{\,2}
 \ge {a_*^{16}\over M^{16}}.\;}
\label{eq:V97-rescue-floor}
\end{equation}
\end{lemma}

\begin{proof}
The literal term belongs to
\eqref{eq:V97-rescue-weight}.  Multiplying
\eqref{eq:V97-actual-edge-floor} over the four edges of \(T\) gives
\[
 \mathfrak l_{\rho,T,r}
 \ge
 {a_*^8\over
   \prod_{i=1}^5\Xi_i^{\deg_T(i)}}.
\]
The degrees sum to eight and every \(\Xi_i\le M\), giving the first
display.  Square it.
\end{proof}

\begin{proposition}[Diffuse maximal-row terminal at order five]
\label{prop:V97-diffuse-terminal}
Suppose a resolved quintic word contains, centrally and in the same
exact source role, a complete source-occurring coordinate bank \(D\)
on which a row \(a\) with \(\Xi_a=M\) is active.  Suppose for fixed
\(\sigma>0\)
\[
 S_a(D)\ge\sigma M,
 \qquad
 \max_{e\in D}{a_{a,e}\over M}\longrightarrow0,
\]
and that the bank is disjoint from the selected and payer coordinates.
Then along that packet both orientations in \(QTRSC_5\) hold
eventually.
\end{proposition}

\begin{proof}
For all sufficiently large packet indices, the second hypothesis is
at most \(\sigma/(2\cdot32)\).  Apply
\Cref{thm:V87-diffuse-bank-power} with \(k=32\) to the complete heat
bank, in the orientation required by the resolved square.  Its
product-state capacity is at most \(C_{G,s,\sigma}M^{-16}\).
The remaining selected fixed-arity local carrier, complete spectators,
physical compression, and sole payer have the established uniform raw
bounds.  Central positive insertion therefore gives the same upper
bound for either complete-word square.  Compare with
\eqref{eq:V97-rescue-floor}.
\end{proof}

\begin{firewall}[Source occurrence remains part of the hypothesis]
\label{fire:V97-diffuse-occurrence}
The arbitrary-power heat theorem may be used only on a complete bank
which occurs in the same exact word and survives the selected/payer
deletions.  \Cref{prop:V97-diffuse-terminal} does not infer such a bank
from a scalar row-mass inequality.
\end{firewall}

\section{Why the quartic one-strong-edge terminal stops at five rows}
\label{sec:V97-strong-defect}

For a tree \(T\), distinguish one edge \(e_0=ac\) which is globally
strong, and put \(F:=T\setminus\{e_0\}\).  The raw fixed-arity
normalization scale is
\begin{equation}
 \mathfrak n_5(\Xi):={1\over\Xi_1\Xi_2\Xi_3\Xi_4\Xi_5}.
\label{eq:V97-raw-scale}
\end{equation}

\begin{lemma}[Exact forest mass defect]
\label{lem:V97-forest-defect}
If \(\theta_{ac}\ge\eta>0\) and the other three edges are actual
source edges obeying \eqref{eq:V97-actual-edge-floor}, then
\begin{equation}
 \theta_{ac}\prod_{e\in F}\theta_e^{\rm act}
 \ge
 {\eta a_*^6\over
   \prod_i\Xi_i^{\deg_F(i)}}
 =
 {\eta a_*^6\over
   \displaystyle\prod_i\Xi_i^{\deg_F(i)-1}}\,
 \mathfrak n_5(\Xi).
\label{eq:V97-forest-defect}
\end{equation}
Here a negative exponent in the displayed defect product is retained
as a denominator; no signwise deletion is permitted.
\end{lemma}

\begin{proof}
Multiply the three actual-edge floors.  Dividing and multiplying by
\(\prod_i\Xi_i\) gives the second equality.
\end{proof}

\begin{proposition}[One strong edge is insufficient from scale data
alone]
\label{prop:V97-one-strong-insufficient}
The hypotheses of \Cref{lem:V97-forest-defect} do not give a uniform
lower bound by \(c\,\mathfrak n_5(\Xi)\).  This remains false when both
endpoints of the strong edge have maximal-scale mass.
\end{proposition}

\begin{proof}
Take \(\Xi_1=\cdots=\Xi_5=M\), let \(T\) be any five-row spanning
tree, choose \(e_0\in E(T)\), set \(\theta_{e_0}=\eta\), and let each
remaining actual edge have its minimal scale \(a_*^2M^{-2}\).
Then
\[
 \theta_{e_0}\prod_{e\in F}\theta_e^{\rm act}
 =\eta a_*^6M^{-6},
 \qquad
 \mathfrak n_5(\Xi)=M^{-5}.
\]
Their ratio is \(\eta a_*^6M^{-1}\to0\).  Both endpoints of
\(e_0\) have mass \(M\).  Equivalently,
\(\sum_i(\deg_F(i)-1)=2|F|-5=1\), which is the uncancelled mass demand
in \eqref{eq:V97-forest-defect}.
\end{proof}

\begin{assessment}[The first genuinely new quintic obstruction]
\label{ass:V97-new-obstruction}
At four rows, deleting one strong edge leaves two actual edges and
\(\sum_i(\deg_F(i)-1)=0\); the V95 forest ratio can therefore close the
raw normalized word after appropriate routing.  At five rows, deleting
one strong edge leaves three actual edges and the exponent sum is one.
Thus the V93--V95 terminal cannot simply be relabelled at order five.
The scalar rescue atlas and the diffuse branch are complete, but a
concentrated balanced packet with only one certified strong edge still
needs one additional inverse row-mass response, or a second independent
strong edge that can be placed in the same lifted source tree.
\end{assessment}

\begin{corollary}[Precise V97 reduction]
\label{cor:V97-reduction}
For global quintic MTA it now suffices to prove \(QTRSC_5\).
Every source-compatible diffuse maximal-row packet satisfies that card.
Any failure sequence may therefore be reduced, after the established
fixed-fibre and fixed-payer selection, to a concentrated packet.  A
terminal argument which produces only one strong normalized edge does
not close that packet without an additional analytic input.
\end{corollary}

\begin{proof}
Sufficiency is \Cref{prop:V97-QTRSC5-suffices}; the diffuse terminal is
\Cref{prop:V97-diffuse-terminal}; and the last assertion is
\Cref{prop:V97-one-strong-insufficient}.  The finite selections do not
cost support multiplicity because \eqref{eq:V97-exact-quintic-source}
has a unique first-connectivity coordinate and only fixed-order fibre
sets.
\end{proof}

\section{V97 result ledger and V98 acquisition order}
\label{sec:V97-status}

\begin{theorem}[Third-series V97 result ledger]
\label{thm:V97-results}
Relative to V1--V96, V97 proves:
\begin{enumerate}[label=\textup{(V97.\arabic*)},leftmargin=6.3em]
\item the archive reconciliation separating the valid local fifth
      cumulant from the invalid graphwise mark route;
\item the exact coefficient-one quintic first-connectivity formula;
\item the exhaustive \(51\)-prefix census and its six analytic
      interfaces;
\item the forest-plus-quotient lift to an ordinary five-row tree;
\item a universal finite rescue atlas retaining a local joining anchor;
\item support-uniform summation of that atlas with explicit multipliers;
\item the sufficient \(QTRSC_5\) square interface for global quintic
      MTA;
\item the \(M^{-8}\) local rescue floor and the \(k=32\) diffuse-bank
      terminal; and
\item the exact one-power forest defect showing why a single strong
      edge no longer closes the balanced concentrated branch.
\end{enumerate}
\end{theorem}

\begin{proof}
These are
\Cref{ass:V97-frozen-status,
thm:V97-exact-quintic-source,
lem:V97-prefix-census,
cor:V97-interface-table,
lem:V97-tree-lift,
lem:V97-rescue-summation,
prop:V97-QTRSC5-suffices,
lem:V97-rescue-floor,
prop:V97-diffuse-terminal,
lem:V97-forest-defect,
prop:V97-one-strong-insufficient,
cor:V97-reduction}.
\end{proof}

\begin{programme}[V98: acquire the missing concentrated response]
\label{prog:V97-V98}
\begin{enumerate}[leftmargin=3em]
\item Run the maximal-row concentration dichotomy inside one fixed
      five-row first-connectivity fibre, without changing its prefix
      block or payer role.
\item Determine when two independent comparable cofinal edges can be
      routed into the same lifted spanning tree while retaining a local
      joining anchor.
\item Prove the exact two-strong-edge forest ratio, including unequal
      row masses and every placement of the two edges.
\item In the residual exactly-one-strong-edge sector, test the complete
      \(SU(2)\) equal-spin singlet heat square.  V91 shows that one
      partial-transpose slot is sharp, so any extra inverse mass must
      come from an independent slot or an assembled fifth-cumulant
      response.
\item If neither occurs, construct a physical five-row counterpacket to
      \(QTRSC_5\) and move upstream to a larger but still summable rescue
      carrier.
\item Keep the \(125\) local fifth-cumulant tree fibres assembled until
      their valid tree-valence decomposition is invoked; do not revive
      the V58 graphwise edge assignment.
\end{enumerate}
\end{programme}

\begin{firewall}[Claim boundary after third-series V97]
\label{fire:V97-claim-boundary}
V97 gives the exact quintic source census, a summable finite rescue
atlas, and closes its source-compatible diffuse branch.  It does not
prove \(QTRSC_5\) in the concentrated one-strong-edge sector, global
quintic or all-order MTA, WUFC, CPM, cutoff removal, reflection
positivity, Osterwalder--Schrader reconstruction, construction of the
continuum Yang--Mills measure, or a uniform positive continuum spectral
gap.  The full Yang--Mills mass gap has not been proved.
\end{firewall}

\paragraph{Parent-version record.}
V97 was formed from
\texttt{Renewal\_Yang\_Mills\_Mass\_Gap\_Research\_Notes\_V96.tex}.
The V96 parent had \(46\,445\) logical source lines,
\(1\,646\,335\) bytes, and SHA--256
\[
 \texttt{8b8fbd64e4001ce5e758221fbf9b236cc10fdb06711386a2ea8777ef92090d54}.
\]
The next compulsory companion audit is V100.

% =============================================================================
% END APPENDED THIRD-SERIES V97 MATERIAL
% =============================================================================

% =============================================================================
% BEGIN APPENDED THIRD-SERIES V98 MATERIAL
% =============================================================================

\clearpage
\part{Third-series V98: exchange rescue and the exact
two-strong-edge geometry}
\label{part:V98}

\section{Purpose}
\label{sec:V98-purpose}

V97 showed that a single globally strong normalized edge pays exactly
enough raw row mass at order four and one power too little at order
five.  The next possible concentrated terminal is therefore a pair of
strong edges.  There are two separate questions:
\begin{enumerate}[label=\textup{(\roman*)},leftmargin=4em]
\item can the two strong edges be exchanged into a spanning tree while
      retaining an actual first-connectivity joining edge; and
\item once that is possible, does the remaining two-edge forest pay the
      five row normalizations for unequal row masses?
\end{enumerate}
This chapter answers both questions exactly.  It also isolates the
only combinatorial failure of the exchange: the local joining anchor is
already one of the strong edge labels, or it is the chord completing
their triangle.  That residual case still has the one-power deficit and
is not hidden inside a support-dependent constant.

\section{The exchange-rescue atlas}
\label{sec:V98-exchange}

Fix one literal source tree \(T_0=P\sqcup Q\) from
\eqref{eq:V97-PQT}.  Thus \(P\) consists of actual prefix-tree edges
and \(Q\) of actual quotient-tree edges at the first joining coordinate
\(r\).  The V97 rescue weight changes the coordinate role of an edge
but keeps the edge set \(T_0\).  A concentrated atom can instead give a
strong pair label outside \(T_0\), so a legal terminal must also allow
finite graphic-basis exchange.

\begin{definition}[Exchange-rescue weight]
\label{def:V98-exchange-weight}
For the fixed source tree \(T_0=P\sqcup Q\), define
\begin{equation}
\begin{split}
 \mathfrak x_{T_0,r}
 :=\sum_{\substack{\tau\in\mathfrak T_5\\
                   E(\tau)\cap Q\ne\varnothing}}
 \ \sum_{A\subseteq E(\tau)\cap P}
 \ \sum_{\varnothing\ne L\subseteq E(\tau)\cap Q}
 &\left(\prod_{e\in A}\theta_e\right)
  \left(\prod_{e\in(E(\tau)\cap P)\setminus A}
                   \theta_e^{<r}\right)\\
 {}\times&
  \left(\prod_{e\in L}\theta_{e,r}\right)
  \left(\prod_{e\in(E(\tau)\cap Q)\setminus L}
                   \theta_e\right)
  \left(\prod_{e\in E(\tau)\setminus(P\cup Q)}
                   \theta_e\right).
\end{split}
\label{eq:V98-exchange-weight}
\end{equation}
Edges of \(\tau\) outside the literal source tree are necessarily
global rescue edges.  At least one edge of the actual joining set \(Q\)
remains local at \(r\).
\end{definition}

\begin{lemma}[Exchange-rescue summation]
\label{lem:V98-exchange-summation}
The enlarged weight is still support-uniform:
\begin{equation}
 \boxed{\;
 \sum_r\mathfrak x_{T_0,r}
 \le 15
 \sum_{\tau\in\mathfrak T_5}
 \prod_{e\in E(\tau)}\theta_e .\;}
\label{eq:V98-exchange-summation}
\end{equation}
Moreover,
\(\mathfrak x_{T_0,r}\ge\mathfrak r_{\rho,T_0,r}\).
\end{lemma}

\begin{proof}
Fix \(\tau,A,L\).  Take every global factor outside the \(r\)-sum,
bound each remaining prefix stock by its global stock, and enlarge the
diagonal local sum to a Cartesian product exactly as in
\Cref{lem:V97-rescue-summation}.  The result is
\(\prod_{e\in E(\tau)}\theta_e\).

Put \(p=|E(\tau)\cap P|\) and \(q=|E(\tau)\cap Q|\).  When \(q\ge1\),
the number of role choices is
\(2^p(2^q-1)\).  Since \(p+q\le4\), this is at most \(15\).
Sum over the \(125\) labelled trees.  Finally, the subfamily
\(\tau=T_0\) in \eqref{eq:V98-exchange-weight} is exactly
\eqref{eq:V97-rescue-weight}.
\end{proof}

\begin{definition}[Exchange-rescue quintic square card]
\label{def:V98-XQTRSC5}
For the resolved word of \Cref{def:V97-QTRSC5}, the card
\(XQTRSC_5\) is
\begin{align}
 \kappa_\otimes(
  \widetilde W_{\rho,T_0,r,\pi}^{*}
  \widetilde W_{\rho,T_0,r,\pi})
 &\le C_{G,s}\mathfrak x_{T_0,r}^{\,2},
\label{eq:V98-XQTRSC5-right}\\
 \kappa_\otimes(
  \widetilde W_{\rho,T_0,r,\pi}
  \widetilde W_{\rho,T_0,r,\pi}^{*})
 &\le C_{G,s}\mathfrak x_{T_0,r}^{\,2}.
\label{eq:V98-XQTRSC5-left}
\end{align}
\end{definition}

\begin{proposition}[\(XQTRSC_5\) is sufficient]
\label{prop:V98-XQTRSC5-suffices}
If \(XQTRSC_5\) holds for every fixed resolved fibre, payer, and
orientation, then the global quintic MTA bound
\eqref{eq:V97-conditional-quintic-MTA} holds.
\end{proposition}

\begin{proof}
Product-state Cauchy--Schwarz gives
\(\kappa_\otimes^{\rm abs}(\widetilde W)
 \le C\mathfrak x_{T_0,r}\).
Apply \Cref{lem:V98-exchange-summation}, then sum the fixed finite
prefix, tree-fibre, endpoint-lift, payer, and physical-output lists.
The right side is already the complete labelled five-row tree atlas.
\end{proof}

\begin{remark}[This enlargement does not revive V58]
\label{rem:V98-not-V58}
Every term of \(\mathfrak x_{T_0,r}\) is a nonnegative scalar
spanning-tree monomial.  No V58 merge difference is normed or assigned
an edge.  The analytic left side remains the complete assembled
first-connectivity word of V97.
\end{remark}

\section{The uniform raw quintic word}
\label{sec:V98-raw}

\begin{lemma}[Raw fixed-payer five-row contraction]
\label{lem:V98-raw-quintic}
For every fixed source type, connected-tree fibre, endpoint lift,
physical output block, payer placement, and dual orientation,
\begin{equation}
 \boxed{\;
 \opnorm{W_{\rho,T_0,r,\pi}^{\rm raw}}\le C_{G,s},
 \qquad
 \opnorm{\widetilde W_{\rho,T_0,r,\pi}}
 \le {C_{G,s}\over\Xi_1\Xi_2\Xi_3\Xi_4\Xi_5}.\;}
\label{eq:V98-raw-quintic}
\end{equation}
The constant is independent of support, representation heights, and
physical multiplicities.
\end{lemma}

\begin{proof}
For every \(B\in\rho\), its complete prefix connected carrier is a
fixed-order moment--cumulant polynomial in complete moment
contractions; its order is at most four.  The joining carrier
\(\mathcal J_r^\rho\) is one complete cumulant of order
\(|\rho|\le5\), again a fixed Bell-number polynomial of contractions.
The complete suffix, singleton moments, and physical compression are
contractions.  The sole final payer is bounded by the established
one-resolvent debit in its assigned factor.  Complete isotypic pinching
precedes the finite triangle inequality, so multiplicities give a
maximum rather than a sum.  This proves the raw bound.  Dividing once
by each of the five row masses gives the normalized bound.
\end{proof}

\section{Two strong edges and a two-edge forest}
\label{sec:V98-two-strong}

\begin{definition}[Maximal-scale strong pair]
\label{def:V98-strong-pair}
Two distinct row edges \(S=\{s_1,s_2\}\) form an
\((\eta,m_0,M)\)-strong pair if
\begin{equation}
 \theta_{s_\ell}\ge\eta,
 \qquad
 \Xi_v\ge m_0M
 \quad
 (v\hbox{ an endpoint of }s_\ell,\ \ell=1,2),
 \qquad
 M=\max_i\Xi_i,
\label{eq:V98-strong-pair}
\end{equation}
where \(\eta,m_0>0\) are fixed along the packet.
\end{definition}

\begin{lemma}[Two-strong-edge forest ratio]
\label{lem:V98-two-strong-ratio}
Let \(\tau\) be a five-vertex spanning tree containing the strong pair
\(S\), and write \(F:=E(\tau)\setminus S\).  Then
\begin{equation}
 { \displaystyle\prod_{uv\in F}\Xi_u\Xi_v
   \over
   \Xi_1\Xi_2\Xi_3\Xi_4\Xi_5}
 \le \max\{1,m_0^{-2}\}.
\label{eq:V98-two-strong-ratio}
\end{equation}
Consequently, if both edges of \(F\) are actual source edges,
\begin{equation}
 \boxed{\;
 \left(\prod_{s\in S}\theta_s\right)
 \left(\prod_{e\in F}\theta_e^{\rm act}\right)
 \ge
 {c_{\eta,m_0,a_*}\over
   \Xi_1\Xi_2\Xi_3\Xi_4\Xi_5}.\;}
\label{eq:V98-two-strong-lower}
\end{equation}
\end{lemma}

\begin{proof}
Deleting two edges from a tree leaves a two-edge forest with three
components.  If its two edges are disjoint, four vertices have degree
one and the fifth vertex \(p\) is isolated.  The ratio in
\eqref{eq:V98-two-strong-ratio} is \(1/\Xi_p\le1\).

Otherwise \(F\) is a two-edge path with centre \(c\) and two isolated
vertices \(p,q\).  Since adding the two strong edges makes the full
tree connected, each isolated component is incident to a strong edge.
Thus \(\Xi_p,\Xi_q\ge m_0M\), while \(\Xi_c\le M\).  The ratio is
\[
 {\Xi_c\over\Xi_p\Xi_q}
 \le {1\over m_0^2M}
 \le m_0^{-2}.
\]
This proves the ratio bound.  The actual edge floors give
\(\prod_{e\in F}\theta_e^{\rm act}\ge
a_*^4/\prod_{uv\in F}\Xi_u\Xi_v\); multiply by
\(\theta_{s_1}\theta_{s_2}\ge\eta^2\) and use the ratio bound.
\end{proof}

\begin{definition}[Anchor-compatible strong pair]
\label{def:V98-anchor-compatible}
The strong pair \(S\) is compatible with the fixed source tree
\(T_0=P\sqcup Q\) if there is
\begin{equation}
 q\in Q\setminus S
 \quad\hbox{such that}\quad
 S\cup\{q\}\ \hbox{is acyclic}.
\label{eq:V98-anchor-compatible}
\end{equation}
\end{definition}

\begin{lemma}[Exact anchored basis exchange]
\label{lem:V98-anchored-exchange}
If \eqref{eq:V98-anchor-compatible} holds, there is an actual edge
\(f\in E(T_0)\setminus(S\cup\{q\})\) such that
\begin{equation}
 \tau:=S\cup\{q,f\}
\label{eq:V98-exchanged-tree}
\end{equation}
is a five-row spanning tree.  Its monomial occurs in
\(\mathfrak x_{T_0,r}\), with both edges of \(S\) global, \(q\) local
at \(r\), and \(f\) in its actual prefix or local joining role.
\end{lemma}

\begin{proof}
The graph \(S\cup\{q\}\) is a three-edge forest on five vertices, so it
has two components.  The source tree \(T_0\) is connected; hence some
edge \(f\in E(T_0)\) crosses those two components.  Such an edge is not
already in \(S\cup\{q\}\).  Adding it gives a connected five-vertex
graph with four edges, hence the tree \(\tau\).

In \eqref{eq:V98-exchange-weight}, select this \(\tau\), put \(q\) in
\(L\), put \(f\) in its literal role, and choose the global role for
each strong edge which happens also to belong to \(P\cup Q\).
Every edge outside \(T_0\) is global by definition.  This is precisely
the stated monomial.
\end{proof}

\begin{theorem}[Anchor-compatible two-strong-edge closure]
\label{thm:V98-two-strong-close}
Every resolved word possessing an
\((\eta,m_0,M)\)-strong pair compatible with its literal source tree
satisfies both \(XQTRSC_5\) orientations.
\end{theorem}

\begin{proof}
Choose \(\tau=S\cup\{q,f\}\) by
\Cref{lem:V98-anchored-exchange}.  The two nondistinguished edges
\(q,f\) are actual source edges and therefore obey
\eqref{eq:V97-actual-edge-floor}.  By
\Cref{lem:V98-two-strong-ratio}, the corresponding term of
\(\mathfrak x_{T_0,r}\) is at least
\[
 {c_{\eta,m_0,a_*}\over\prod_i\Xi_i}.
\]
Compare with \eqref{eq:V98-raw-quintic}.  The ordinary operator norm
bound, squared, bounds either oriented positive product-state capacity.
\end{proof}

\section{Exact failure of anchored exchange}
\label{sec:V98-anchor-failure}

\begin{lemma}[Chord-or-coverage classification]
\label{lem:V98-anchor-classification}
Let \(S\) be two distinct edges and let \(Q\ne\varnothing\).
There is no \(q\in Q\setminus S\) satisfying
\eqref{eq:V98-anchor-compatible} if and only if
\begin{equation}
 Q\subseteq S\cup\{q_\triangle\},
\label{eq:V98-obstruction-set}
\end{equation}
where either \(Q\subseteq S\), or the two edges of \(S\) form a
length-two path and \(q_\triangle\) is its unique closing chord.
\end{lemma}

\begin{proof}
A two-edge graph is always a forest.  Adding a third distinct edge
creates a cycle precisely when the two original edges are adjacent and
the third edge joins the two free endpoints; that edge is unique.
Thus every edge of \(Q\setminus S\) fails acyclicity exactly when it is
the unique chord \(q_\triangle\).  This is equivalent to
\eqref{eq:V98-obstruction-set}, including the vacuous case
\(Q\setminus S=\varnothing\).
\end{proof}

\begin{corollary}[Automatic compatibility for a five-block junction]
\label{cor:V98-five-block-compatible}
For the discrete prefix, \(|Q|=4\), so every two-edge strong pair is
anchor-compatible.  More generally compatibility is automatic whenever
\(|Q\setminus S|\ge2\).
\end{corollary}

\begin{proof}
At most one edge outside \(S\) is its closing chord.  If two candidates
are available, at least one is acyclic.  For the discrete prefix,
\(|Q|=4>|S|+1\), so \(|Q\setminus S|\ge2\).
\end{proof}

\begin{assessment}[Affected quotient orders]
\label{ass:V98-affected-orders}
The residual obstruction is finite and explicit.
\begin{enumerate}[label=\textup{(\roman*)},leftmargin=4em]
\item For \(k=5\), it cannot occur.
\item For \(k=4\), it can occur only when the three local quotient-tree
      edges are the two strong labels and their closing chord.
\item For \(k=3\), it can occur only when the two local quotient-tree
      edges lie in the three-edge set of
      \eqref{eq:V98-obstruction-set}.
\item For \(k=2\), the sole covariance edge is either a strong label or
      the closing chord.
\end{enumerate}
This classification concerns labels and coordinate roles, not the
validity of the complete local cumulant card.
\end{assessment}

\begin{proposition}[The anchor obstruction still has a one-power
deficit]
\label{prop:V98-anchor-deficit}
Assume equal row masses \(\Xi_i=M\), two global strong factors of fixed
size, and no additional heat response.  In either obstruction of
\Cref{lem:V98-anchor-classification}, retaining both strong factors,
an actual local joining anchor, and connectedness requires at least
three factors at the actual-edge scale \(M^{-2}\).  The resulting
positive carrier can therefore be only \(O(M^{-6})\) from these data,
whereas the raw normalized word has scale \(M^{-5}\).
\end{proposition}

\begin{proof}
First suppose \(q_\triangle\in Q\setminus S\).  The three edges
\(S\cup\{q_\triangle\}\) form a triangle on three vertices.  The two
remaining vertices are separate components, so at least two further
actual source edges are required for a connected spanning carrier.
Besides the two global strong factors, there are therefore the local
anchor and two further actual factors: three factors of scale
\(M^{-2}\).

If instead every joining anchor lies in \(S\), using its global strong
stock does not retain a local \(r\)-summable factor.  Retaining both
roles requires an additional local copy of that pair stock.  A
spanning tree containing the two strong labels still requires two
further actual edges, so again there are at least three actual-scale
factors.  In both cases their product is \(O(M^{-6})\), while
\(\mathfrak n_5(\Xi)=M^{-5}\).
\end{proof}

\begin{firewall}[This is an insufficiency theorem, not a counterexample
to quintic MTA]
\label{fire:V98-deficit-scope}
\Cref{prop:V98-anchor-deficit} says that two global strong-edge lower
bounds plus minimal actual-edge floors do not prove the raw contraction
estimate in the obstructed label geometry.  The complete Wilson word
may contain additional damping or cancellation.  The proposition does
not assert that \(XQTRSC_5\), global quintic MTA, or Yang--Mills theory
is false.
\end{firewall}

\section{V98 result ledger and V99 acquisition order}
\label{sec:V98-status}

\begin{theorem}[Third-series V98 result ledger]
\label{thm:V98-results}
Relative to V1--V97, V98 proves:
\begin{enumerate}[label=\textup{(V98.\arabic*)},leftmargin=6.3em]
\item the exchange-rescue atlas allowing strong labels outside the
      literal source tree;
\item its support-uniform reduction to the \(125\)-tree global atlas;
\item the sufficient \(XQTRSC_5\) analytic interface;
\item the uniform raw fixed-payer quintic contraction;
\item the exact unequal-mass two-strong-edge forest ratio;
\item anchored graphic-basis exchange retaining an actual local joining
      edge;
\item closure of every anchor-compatible two-strong-edge word;
\item the chord-or-coverage classification of every failed exchange;
      and
\item the surviving one-power deficit in precisely that obstruction.
\end{enumerate}
\end{theorem}

\begin{proof}
These are
\Cref{lem:V98-exchange-summation,
prop:V98-XQTRSC5-suffices,
lem:V98-raw-quintic,
lem:V98-two-strong-ratio,
lem:V98-anchored-exchange,
thm:V98-two-strong-close,
lem:V98-anchor-classification,
cor:V98-five-block-compatible,
prop:V98-anchor-deficit}.
\end{proof}

\begin{programme}[V99: concentration multiplicity and the obstructed
anchor]
\label{prog:V98-V99}
\begin{enumerate}[leftmargin=3em]
\item Starting with a maximal row, refine the V95 concentration
      dichotomy by the number of distinct comparable partner labels,
      while retaining the same V97 source word.
\item If two distinct strong labels are obtained, apply
      \Cref{thm:V98-two-strong-close} immediately unless the exact
      chord-or-coverage obstruction occurs.
\item In the obstructed case, determine whether the local joining
      occurrence itself is cofinal.  If so, use its local strength and
      summability simultaneously rather than duplicating a global mark.
\item If the joining occurrence is not cofinal, seek one additional
      inverse mass from a complete heat square on a source-occurring
      bank.  The required gain is \(M^{-1}\), not the full
      \(M^{-16}\) diffuse terminal.
\item Treat repeated cofinal atoms with the same partner as one strong
      pair label; do not count coordinate multiplicity as graph
      independence.
\item If a single partner owns all comparable concentration, construct
      the corresponding two-row cofinal core and test it against the
      exact \(SU(2)\) singlet lower bound before proposing a response
      estimate.
\end{enumerate}
\end{programme}

\begin{firewall}[Claim boundary after third-series V98]
\label{fire:V98-claim-boundary}
V98 closes every anchor-compatible two-strong-edge quintic word and
completely classifies failed anchored exchanges.  It does not prove
that every concentrated packet supplies such a pair, close the
chord-or-coverage sector, prove global quintic or all-order MTA, WUFC,
CPM, cutoff removal, reflection positivity, Osterwalder--Schrader
reconstruction, construction of the continuum Yang--Mills measure, or
a uniform positive continuum spectral gap.  The full Yang--Mills mass
gap has not been proved.
\end{firewall}

\paragraph{Parent-version record.}
V98 was formed from
\texttt{Renewal\_Yang\_Mills\_Mass\_Gap\_Research\_Notes\_V97.tex}.
The V97 parent had \(47\,082\) logical source lines,
\(1\,670\,602\) bytes, and SHA--256
\[
 \texttt{07c470e1b729cd68f25ca047d007d6d3efe32fd748b0f3703c98e1a85e0b51e2}.
\]
The next compulsory companion audit is V100.

% =============================================================================
% END APPENDED THIRD-SERIES V98 MATERIAL
% =============================================================================

% =============================================================================
% BEGIN APPENDED THIRD-SERIES V99 MATERIAL
% =============================================================================

\clearpage
\part{Third-series V99: the sharp one-power supplement and
cofinal-partner reduction}
\label{part:V99}

\section{Purpose and exponent correction}
\label{sec:V99-purpose}

The V98 anchor obstruction leaves a deficit of \(M^{-1}\) at the
normalized word level:
\[
 \opnorm{\widetilde W}=O(M^{-5}),
 \qquad
 \hbox{available one-strong tree weight}=O(M^{-6})
\]
when all five row masses are comparable.  Since the analytic cards are
oriented square estimates, the missing heat capacity is \(M^{-2}\),
not \(M^{-1}\).  This chapter proves that four diffuse complete heat
slots supply exactly that square-level supplement and that fewer than
four independent \(SU(2)\) singlet slots cannot do so uniformly.

It then organizes the concentrated alternatives by the number of
distinct comparable partners of a maximal row.  Three partners always
contain an anchor-compatible strong pair.  Two partners reduce to the
finite V98 obstruction and one partner gives a genuine two-row cofinal
core.  Repeated coordinates with the same partner are not miscounted as
distinct tree edges.

\section{One strong edge inside the exchange atlas}
\label{sec:V99-one-strong}

\begin{lemma}[Anchored one-strong exchange]
\label{lem:V99-one-strong-exchange}
Let \(T_0=P\sqcup Q\) be a literal source tree and let
\(S=\{s_1,s_2\}\) be two distinct row-edge labels.  For every
\(q\in Q\), one can choose \(s\in S\setminus\{q\}\) and actual edges
\(f,g\in E(T_0)\) such that
\begin{equation}
 \tau=\{s,q,f,g\}
\label{eq:V99-one-strong-tree}
\end{equation}
is a five-row spanning tree.  The exchange weight
\(\mathfrak x_{T_0,r}\) contains its monomial with \(s\) global,
\(q\) local at \(r\), and \(f,g\) in their actual source roles.
\end{lemma}

\begin{proof}
Because \(S\) has two distinct edges, at least one of them differs from
the fixed \(q\); choose it as \(s\).  The two-edge graph
\(\{s,q\}\) is a forest.  Starting from its three components, use the
connected source tree \(T_0\) to choose successively an edge crossing
two current components.  Two such edges \(f,g\) produce a connected
five-vertex graph with four edges, hence a spanning tree.  Neither
chosen edge is already in the growing forest.

Select this tree in \eqref{eq:V98-exchange-weight}, retain \(q\) in
the local set \(L\), place \(f,g\) in their literal roles, and choose
the global role for \(s\) if \(s\in P\cup Q\).  If \(s\notin T_0\),
its role is global by definition.
\end{proof}

\begin{lemma}[Sharp one-strong forest loss]
\label{lem:V99-one-strong-loss}
Let \(\tau\) be a five-row tree containing a distinguished edge
\(s=ac\) such that
\[
 \theta_{ac}\ge\eta,
 \qquad
 \Xi_a,\Xi_c\ge m_0M,
 \qquad M=\max_i\Xi_i.
\]
If the three edges \(F=E(\tau)\setminus\{ac\}\) are actual source
edges, then
\begin{align}
 { \displaystyle\prod_{uv\in F}\Xi_u\Xi_v
   \over \displaystyle\prod_{i=1}^5\Xi_i}
 &\le C_{m_0}M,
\label{eq:V99-one-strong-ratio}\\
 \theta_{ac}\prod_{e\in F}\theta_e^{\rm act}
 &\ge
 {c_{\eta,m_0,a_*}\over
  M\Xi_1\Xi_2\Xi_3\Xi_4\Xi_5}.
\label{eq:V99-one-strong-floor}
\end{align}
The power \(M\) in \eqref{eq:V99-one-strong-ratio} is sharp when all
row masses are \(M\).
\end{lemma}

\begin{proof}
Deleting \(ac\) leaves a three-edge forest with two components.
Their vertex-size types are \(2+3\) or \(1+4\).

For type \(2+3\), the two-vertex component is an edge and the
three-vertex component is a path.  Cancelling one copy of every row
mass leaves only the mass of the path centre, which is at most \(M\).

For type \(1+4\), let \(p\) be the isolated vertex.  It is an endpoint
of the deleted edge \(ac\), since that edge is the only connection
between the two components of the original tree.  Hence
\(\Xi_p\ge m_0M\).  In the four-vertex tree the sum of the positive
exponents \(\deg(v)-1\) is two, so its excess product is at most
\(M^2\).  Division by the isolated factor gives at most \(M/m_0\).
This proves \eqref{eq:V99-one-strong-ratio}.

Multiply the three actual-edge floors and the strong lower bound, then
use the ratio estimate to obtain
\eqref{eq:V99-one-strong-floor}.  If every \(\Xi_i=M\), the ratio on
the left of \eqref{eq:V99-one-strong-ratio} is exactly \(M\), because
\(\sum_i(\deg_F(i)-1)=1\).
\end{proof}

\begin{corollary}[Uniform floor available in every V98 obstruction]
\label{cor:V99-obstructed-floor}
In a chord-or-coverage obstruction with an
\((\eta,m_0,M)\)-strong pair, the exchange atlas satisfies
\begin{equation}
 \boxed{\;
 \mathfrak x_{T_0,r}
 \ge
 {c_{\eta,m_0,a_*}\over
  M\prod_i\Xi_i}.\;}
\label{eq:V99-obstructed-floor}
\end{equation}
\end{corollary}

\begin{proof}
Choose any \(q\in Q\) and then the strong label \(s\ne q\) supplied by
\Cref{lem:V99-one-strong-exchange}.  Apply
\Cref{lem:V99-one-strong-loss} to the resulting tree.  Its monomial is
one of the nonnegative terms of \(\mathfrak x_{T_0,r}\).
\end{proof}

\section{The exact square-level supplement}
\label{sec:V99-supplement}

\begin{definition}[One-power supplement bank]
\label{def:V99-supplement-bank}
A resolved oriented quintic square has a one-power supplement if it
contains centrally, in the same exact source word, a complete positive
bank factor \(\mathcal M_D^2\) such that
\begin{equation}
 \kappa_{a|a^c}(\mathcal M_D^2)
 \le C_{G,s}M^{-2}
\label{eq:V99-supplement}
\end{equation}
for a maximal row \(a\).  The bank is disjoint from the selected,
payer, and fixed strong coordinates.
\end{definition}

\begin{theorem}[A one-power supplement closes every obstructed strong
pair]
\label{thm:V99-supplement-close}
If a chord-or-coverage word has an
\((\eta,m_0,M)\)-strong pair and the supplement
\eqref{eq:V99-supplement}, then it satisfies both orientations of
\(XQTRSC_5\).
\end{theorem}

\begin{proof}
The normalized raw contraction
\eqref{eq:V98-raw-quintic}, with central positive insertion of the
bank, gives either oriented square capacity at most
\[
 {C_{G,s}M^{-2}\over(\prod_i\Xi_i)^2}.
\]
By \eqref{eq:V99-obstructed-floor}, this is at most
\(C\mathfrak x_{T_0,r}^{\,2}\).  All comparisons occur inside one
fixed source word and one payer placement.
\end{proof}

\begin{proposition}[Four diffuse slots supply the supplement]
\label{prop:V99-four-slot-supplement}
Let \(D\) be a complete source-occurring bank on a maximal row \(a\)
such that, for fixed \(\sigma>0\),
\begin{equation}
 S_a(D)\ge\sigma M,
 \qquad
 \max_{e\in D}{a_{a,e}\over M}
 \le{\sigma\over8}.
\label{eq:V99-four-slot-hyp}
\end{equation}
Then \(D\) is a one-power supplement bank.
\end{proposition}

\begin{proof}
Apply \Cref{thm:V87-diffuse-bank-power} with \(k=4\).
It gives
\[
 \kappa_{a|a^c}(\mathcal M_D^2)
 \le C_{G,s,\sigma}M^{-4/2}
 =C_{G,s,\sigma}M^{-2}.
\]
The theorem applies to the square because
\(0\preccurlyeq\mathcal M_D^2\preccurlyeq\mathcal M_D\).
\end{proof}

\begin{proposition}[\(SU(2)\) four-slot sharpness]
\label{prop:V99-four-slot-sharp}
Within the class of tensor products of independent complete
equal-spin \(SU(2)\) heat-square slots, no product of \(q<4\) slots can
obey a uniform \(CM^{-2}\) upper bound when their total row mass is
comparable to \(M\).
\end{proposition}

\begin{proof}
For spin \(j\), V91 proves that the complete heat square retains the
singlet and has product-state capacity at least
\[
 {1\over 2j+1}\asymp a_j^{-1/2},
\]
where \(a_j=1+C_2(j)\).  Take \(q\) independent equal-spin slots.
Testing the tensor product on the tensor product of their optimizing
product states gives a lower bound
\[
 c_q a_j^{-q/2}.
\]
For fixed \(q\), their total row mass is \(M=qa_j\asymp a_j\), so this
is \(c'_qM^{-q/2}\).  If \(q<4\), its ratio to \(M^{-2}\) diverges as
\(M^{2-q/2}\).  Hence fewer than four such slots cannot supply
\eqref{eq:V99-supplement} uniformly.
\end{proof}

\begin{firewall}[The word-level and square-level powers]
\label{fire:V99-power-ledger}
The obstructed tree is smaller than the raw word by one factor
\(M^{-1}\).  Squaring makes the required capacity gain \(M^{-2}\).
One complete equal-spin heat-square slot can be only
\(M^{-1/2}\) by \Cref{prop:V99-four-slot-sharp}; four independent
slots are the sharp tensor-product threshold.  A proof that silently
uses one slot four times is invalid.
\end{firewall}

\section{The maximal-row partner graph}
\label{sec:V99-partner}

Fix a counterpacket with maximal row \(a\), \(\Xi_a=M\to\infty\).
Call a selected row-\(a\) coordinate \(\sigma\)-cofinal if
\(a_{a,e}\ge\sigma M\).  In the balanced alternative of the
established cofinal-atom dichotomy, choose a row
\(c(e)\ne a\) with
\begin{equation}
 a_{c(e),e}\ge\lambda a_{a,e},
 \qquad
 \Xi_{c(e)}\ge\lambda\sigma M,
 \qquad
 \theta_{a\,c(e)}\ge\lambda\sigma^2.
\label{eq:V99-partner-edge}
\end{equation}
The partner graph is the simple star whose edge labels are the
distinct pairs \(ac(e)\).  Several coordinates with the same partner
give one edge label.

\begin{lemma}[Cofinal atoms are finite but repeated labels do not add]
\label{lem:V99-cofinal-count}
For fixed \(\sigma\), row \(a\) has at most
\(\lfloor\sigma^{-1}\rfloor\) pairwise distinct
\(\sigma\)-cofinal coordinates.  Nevertheless, if all their balanced
partners equal one row \(c\), the partner graph contains only the
single edge \(ac\).
\end{lemma}

\begin{proof}
The coordinate masses are nonnegative and sum to \(\Xi_a=M\), so
\(N\sigma M\le M\) for \(N\) such coordinates.  The second assertion
is the definition of a simple pair-label graph.  The global stock
\(\theta_{ac}\) sums repeated occurrences; it does not create parallel
edges in a spanning-tree monomial.
\end{proof}

\begin{lemma}[Three distinct partners force compatible exchange]
\label{lem:V99-three-partners}
Suppose the partner graph contains three distinct edges
\(ac,ad,ae\).  For every nonempty joining set \(Q\), two of those
strong edges form an anchor-compatible pair in the sense of
\Cref{def:V98-anchor-compatible}.
\end{lemma}

\begin{proof}
Choose any \(q\in Q\).
If \(q\) is one of \(ac,ad,ae\), choose the other two; adjoining \(q\)
gives a three-edge star and is acyclic.

If \(q\) joins two of the partner vertices, say \(q=cd\), choose
\(ac\) and \(ae\).  The edges \(ac,ae,cd\) form the path
\(d-c-a-e\), hence are acyclic.  In every other case, choose two star
edges so that \(q\) is not the chord between their leaves; this is
possible because there are three partner leaves.  Then \(q\) is
outside the chosen pair and its addition creates no cycle.
\end{proof}

\begin{corollary}[Three-partner concentrated closure]
\label{cor:V99-three-partner-close}
Every balanced concentrated resolved word whose maximal row has at
least three distinct comparable partners satisfies \(XQTRSC_5\).
\end{corollary}

\begin{proof}
Each chosen partner edge obeys
\eqref{eq:V98-strong-pair}, with constants determined by
\(\sigma,\lambda\).  Use \Cref{lem:V99-three-partners} and then
\Cref{thm:V98-two-strong-close}.
\end{proof}

\begin{theorem}[Exact concentrated-partner reduction]
\label{thm:V99-partner-reduction}
After fixing one quintic first-connectivity fibre, payer, and
orientation, every persistent maximal-row failure not already closed by
the unbalanced/gapped compiler or the V97 diffuse terminal belongs to
one of the following disjoint analytic classes:
\begin{enumerate}[label=\textup{(\roman*)},leftmargin=4em]
\item at least three distinct comparable partners, closed by
      \Cref{cor:V99-three-partner-close};
\item exactly two distinct comparable partners and an
      anchor-compatible pair, closed by
      \Cref{thm:V98-two-strong-close};
\item exactly two distinct partners in the chord-or-coverage geometry
      with a four-slot source-occurring bank, closed by
      \Cref{thm:V99-supplement-close};
\item exactly two distinct partners in the chord-or-coverage geometry
      with no proved supplement bank; or
\item one distinct partner label, giving a two-row cofinal core, with no
      proved supplement bank.
\end{enumerate}
Only \textup{(iv)} and \textup{(v)} remain open after V99.
\end{theorem}

\begin{proof}
The fixed maximal row is either diffuse or has a cofinal atom.  The
former is V97.  In the latter case the established atom dichotomy is
unbalanced, which is already closed, or balanced, which gives
\eqref{eq:V99-partner-edge}.  Count distinct partner labels.
Three or more reduce to three and use
\Cref{cor:V99-three-partner-close}.  Two labels are either compatible
or belong to the exhaustive obstruction
\Cref{lem:V98-anchor-classification}; in the latter case the supplement
is either available or not.  One label is the final class.  Zero labels
is exactly the unbalanced alternative.  The alternatives are disjoint
by construction.
\end{proof}

\begin{assessment}[The remaining core is analytic, not combinatorial]
\label{ass:V99-remaining-core}
The \(51\)-family source census, global scalar summation, graphic
exchange, unequal-mass forest ratios, and diffuse heat powers have now
all been performed.  The remaining quintic issue is narrower: in a
fixed chord-obstructed two-partner packet or a fixed one-partner
two-row cofinal core, prove a source-compatible \(M^{-2}\) square
response, or exhibit a physical counterpacket to \(XQTRSC_5\).
Repeated large coordinates with the same row-pair label cannot be
declared independent tree edges, and a single heat slot cannot be
declared to supply four-slot decay.
\end{assessment}

\section{V99 result ledger and V100 acquisition order}
\label{sec:V99-status}

\begin{theorem}[Third-series V99 result ledger]
\label{thm:V99-results}
Relative to V1--V98, V99 proves:
\begin{enumerate}[label=\textup{(V99.\arabic*)},leftmargin=6.3em]
\item the exact distinction between the missing word-level
      \(M^{-1}\) and square-level \(M^{-2}\) powers;
\item anchored exchange with one retained strong label;
\item the sharp unequal-mass one-strong forest loss;
\item a uniform obstructed-sector exchange-atlas floor;
\item closure from a source-compatible \(M^{-2}\) supplement;
\item acquisition of that supplement from four diffuse slots;
\item \(SU(2)\) sharpness showing that fewer independent singlet slots
      cannot supply it;
\item the simple maximal-row partner graph and its no-parallel-edge
      rule;
\item automatic closure from three distinct comparable partners; and
\item reduction of the open concentrated sector to the exact classes
      \textup{(iv)}--\textup{(v)} of
      \Cref{thm:V99-partner-reduction}.
\end{enumerate}
\end{theorem}

\begin{proof}
These are
\Cref{lem:V99-one-strong-exchange,
lem:V99-one-strong-loss,
cor:V99-obstructed-floor,
thm:V99-supplement-close,
prop:V99-four-slot-supplement,
prop:V99-four-slot-sharp,
lem:V99-cofinal-count,
lem:V99-three-partners,
cor:V99-three-partner-close,
thm:V99-partner-reduction}.
\end{proof}

\begin{programme}[V100: compulsory audit and the two-row cofinal core]
\label{prog:V99-V100}
\begin{enumerate}[leftmargin=3em]
\item Perform the compulsory audit of the newest active SM and NS
      notes before changing the analytic route.
\item Freeze the validated V100 note as
      \texttt{\_V100\_f3.tex}, then restart the active series at V1
      with an updated major-result summary, as required by the
      hundred-version rollover rule.
\item In the one-partner core, condition on the complete pair
      representation rather than on individual cofinal coordinates.
\item Determine whether four independent complete heat slots survive
      inside the same first-connectivity word.  If not, do not apply
      \Cref{prop:V99-four-slot-supplement}.
\item Test the exact \(SU(2)\) dual-singlet sector: its one-slot
      \(M^{-1/2}\) lower bound is binding and rules out any claimed
      single-slot \(M^{-2}\) estimate.
\item For the chord-obstructed two-partner class, keep the assembled
      quotient cumulant and test whether its local joining block
      supplies a genuine second response rather than a duplicate edge
      mark.
\end{enumerate}
\end{programme}

\begin{firewall}[Claim boundary after third-series V99]
\label{fire:V99-claim-boundary}
V99 closes the three-partner sector and every obstructed two-partner
sector possessing a source-compatible four-slot supplement.  It does
not close the supplement-free chord obstruction or the one-partner
two-row cofinal core, prove global quintic or all-order MTA, WUFC, CPM,
cutoff removal, reflection positivity, Osterwalder--Schrader
reconstruction, construction of the continuum Yang--Mills measure, or
a uniform positive continuum spectral gap.  The full Yang--Mills mass
gap has not been proved.
\end{firewall}

\paragraph{Parent-version record.}
V99 was formed from
\texttt{Renewal\_Yang\_Mills\_Mass\_Gap\_Research\_Notes\_V98.tex}.
The V98 parent had \(47\,562\) logical source lines,
\(1\,689\,126\) bytes, and SHA--256
\[
 \texttt{8dc4c10106a62629f34a308629e9f70a3c7128e9344becb9889457b969d51192}.
\]
The compulsory companion audit is due at V100.

% =============================================================================
% END APPENDED THIRD-SERIES V99 MATERIAL
% =============================================================================

% =============================================================================
% BEGIN APPENDED THIRD-SERIES V100 MATERIAL
% =============================================================================

\clearpage
\part{Third-series V100: compulsory companion audit and the
finite cofinal-slot core}
\label{part:V100}

\section{Compulsory SM/NS audit}
\label{sec:V100-audit}

The active companion files inspected before this continuation were:
\begin{center}
\begin{tabular}{p{0.57\textwidth}rr}
\toprule
file & source lines & bytes\\
\midrule
\texttt{Renewal\_SM\_Einstein\_Continuum\_Research\_Notes\_V77.tex}
& \(26\,855\) & \(1\,095\,259\)\\
\texttt{Renewal\_NS\_Stability\_Research\_Notes\_V31.tex}
& \(21\,557\) & \(887\,537\)\\
\bottomrule
\end{tabular}
\end{center}
Their SHA--256 digests are respectively
\[
\begin{split}
 &\texttt{4b71e6138094b927d757684f27ccd8eca6798265044970aca8b75bba25954ed5},\\
 &\texttt{f1121b62238ad5491bb7cf03635ea9934942edf573ed8faaa9ee79e1d38a6696}.
\end{split}
\]

\begin{assessment}[SM V77 audit]
\label{ass:V100-SM-audit}
SM V77 proves that a smooth zeroth-order Yukawa perturbation preserves
the regulated chiral Fredholm index.  It therefore refuses to replace
an unmatched electroweak exterior-power sector by a square mass
determinant, while allowing the exact compressed mass minor in the
index-zero vectorlike color sector.

The type-correct lesson for the present quintic core is that a scalar
lower bound cannot repair a mismatch of operator role.  A global strong
edge, a local joining anchor, and a complete heat slot are different
objects even when they carry the same row-pair label.  They may be
combined only when the exact first-connectivity word contains those
roles.  No SM determinant, index theorem, or lower bound is imported
into Yang--Mills.
\end{assessment}

\begin{assessment}[NS V31 audit]
\label{ass:V100-NS-audit}
NS V31 proves that flat dyadic marks remove stopping-time multiplicity,
but unmarking costs their exact first radial moment.  Its
weak/intermediate/critical heat ladder is scale matched, and a forced
mode proves the heat exponent sharp.  A balanced-helicity plane wave
also shows that the signed hinge has a nullspace on the total-tail
branch.

For the Yang--Mills core, the applicable proof-order rule is to retain
the exact missing \(M^{-2}\) square response.  Replacing it by an
unweighted count of cofinal coordinates merely hides the same power,
and heat response alone cannot be strengthened past its sharp slotwise
exponent.  Any improvement in the residual \(q<4\) core must use the
assembled local quotient cumulant before norms, just as NS must use its
complete signed nonlinear source before separate positive parts.  No
Navier--Stokes estimate is imported.
\end{assessment}

\begin{record}[Audit direction decision]
\label{rec:V100-audit-decision}
The audit changes no Yang--Mills theorem or constant.  It confirms the
V99 direction: first count complete source-occurring heat slots with
their coordinate roles; then isolate the precise response still owed by
the assembled quotient cumulant.  Further scalar graph enlargement is
not the next task.
\end{record}

\section{Cofinal complete slots}
\label{sec:V100-cofinal-slots}

\begin{definition}[Usable cofinal heat slot]
\label{def:V100-usable-slot}
For a fixed resolved quintic square and a maximal row \(a\), a
coordinate \(e\) is a usable \((\sigma,M)\)-cofinal slot if:
\begin{enumerate}[label=\textup{(\roman*)},leftmargin=4em]
\item its complete moment-partition heat factor occurs centrally in
      that same exact source word;
\item row \(a\) is active there and \(a_{a,e}\ge\sigma M\);
\item it is disjoint from the first-connectivity, payer, and fixed
      strong-coordinate roles; and
\item its full complementary module and every physical multiplicity
      remain inside the complete heat factor.
\end{enumerate}
\end{definition}

\begin{theorem}[Four cofinal complete slots supply \(M^{-2}\)]
\label{thm:V100-four-cofinal}
If a resolved oriented square contains four distinct usable
\((\sigma,M)\)-cofinal slots, their complete product is a one-power
supplement bank:
\begin{equation}
 \boxed{\;
 \kappa_{a|a^c}
 \left(\bigotimes_{\ell=1}^{4}\mathbb M_{e_\ell}(s)\right)
 \le C_{G,s,\sigma}M^{-2}.\;}
\label{eq:V100-four-cofinal}
\end{equation}
The same estimate holds for every positive protected sector and for
the square of such a sector.
\end{theorem}

\begin{proof}
At slot \(e_\ell\), the reducible-complement heat card
\Cref{cor:V87-local-slot} gives
\[
 \|\mathbb M_{e_\ell}(s)^{\Gamma_{a^c}}\|
 \le C_{G,s}a_{a,e_\ell}^{-1/2}
 \le C_{G,s,\sigma}M^{-1/2}.
\]
The slots are distinct tensor coordinates.  Partial transpose on the
fixed row cut and operator norm tensorize, even when the row state is
entangled among the four coordinates.  Multiplication gives \(CM^{-2}\).
Positive domination and
\(0\preccurlyeq\mathcal M\preccurlyeq I\) imply the assertions for a
protected sector and its square.
\end{proof}

\begin{corollary}[Four-cofinal closure of every V98 obstruction]
\label{cor:V100-four-cofinal-close}
A chord-or-coverage word with a maximal-scale strong pair and four
usable cofinal slots satisfies both \(XQTRSC_5\) orientations.
\end{corollary}

\begin{proof}
\Cref{thm:V100-four-cofinal} gives
\eqref{eq:V99-supplement}.  Apply
\Cref{thm:V99-supplement-close}.
\end{proof}

\begin{remark}[Diffuse and cofinal four-slot routes are distinct]
\label{rem:V100-two-four-slot-routes}
\Cref{prop:V99-four-slot-supplement} extracts four effective powers
from a diffuse bank of many small coordinates by optimization.
\Cref{thm:V100-four-cofinal} uses four individually large coordinates.
Neither theorem implies the other, but both produce the identical
\(M^{-2}\) square interface required by
\Cref{def:V99-supplement-bank}.
\end{remark}

\section{The residual slot budget}
\label{sec:V100-budget}

After deleting the selected, payer, and fixed strong coordinates, let
\(q\in\{0,1,2,3\}\) be the number of usable cofinal complete slots in a
supplement-free core.  Their slotwise heat response is
\(M^{-q/2}\).  Define the response still owed by the assembled remainder
to be
\begin{equation}
 \delta_q:={4-q\over2}.
\label{eq:V100-deltaq}
\end{equation}

\begin{lemma}[Exact residual response table]
\label{lem:V100-response-table}
The square-level response required beyond \(q\) usable cofinal slots is:
\begin{center}
\begin{tabular}{c|cccc}
\(q\)&0&1&2&3\\
\hline
available heat power&\(1\)&\(M^{-1/2}\)&\(M^{-1}\)&\(M^{-3/2}\)\\
still required&\(M^{-2}\)&\(M^{-3/2}\)&\(M^{-1}\)&\(M^{-1/2}\).
\end{tabular}
\end{center}
Equivalently the remainder must supply \(M^{-\delta_q}\).
\end{lemma}

\begin{proof}
The four-slot target is \(M^{-2}=M^{-4/2}\).  Tensorization of the
\(q\) complete-slot bounds gives \(M^{-q/2}\).  Division leaves
\(M^{-(4-q)/2}\), producing the table.
\end{proof}

\begin{definition}[Assembled \(q\)-slot core response]
\label{def:V100-qcore}
Let \(\mathcal R_{q,\rho,r,\pi}\) be the complete factor remaining in
the oriented first-connectivity square after extracting exactly \(q\)
usable complete heat slots.  It includes the complete quotient
cumulant \(\mathcal J_r^\rho\), physical output compression, and the
single payer in its fixed role.  The \(q\)-slot core card is
\begin{equation}
 \kappa_{a|a^c}
  (\mathcal R_{q,\rho,r,\pi})
 \le C_{G,s}M^{-\delta_q}.
\label{eq:V100-qcore}
\end{equation}
No moment-partition summand of \(\mathcal J_r^\rho\) is normed
separately in this definition.
\end{definition}

\begin{proposition}[The \(q\)-slot card is sufficient for the missing
supplement]
\label{prop:V100-qcore-suffices}
If \eqref{eq:V100-qcore} holds in the orientation of a fixed residual
core, then that core has the supplement
\eqref{eq:V99-supplement} and hence satisfies \(XQTRSC_5\).
\end{proposition}

\begin{proof}
The \(q\) extracted heat slots contribute
\(CM^{-q/2}\).  They are tensor-coordinate factors disjoint from the
remainder, so their partially transposed norm multiplies the product
state capacity of \(\mathcal R_{q,\rho,r,\pi}\).  By
\eqref{eq:V100-deltaq},
\[
 M^{-q/2}M^{-\delta_q}=M^{-2}.
\]
This is the supplement.  Apply
\Cref{thm:V99-supplement-close}.
\end{proof}

\begin{proposition}[Sharp heat-only obstruction for \(q<4\)]
\label{prop:V100-heat-only-obstruction}
For every \(q<4\), the bound
\eqref{eq:V100-qcore} cannot be obtained by replacing
\(\mathcal R_{q,\rho,r,\pi}\) with a bounded contraction and using only
the \(q\) extracted complete heat slots.  The failure already occurs on
the \(SU(2)\) equal-spin singlet family.
\end{proposition}

\begin{proof}
\Cref{prop:V99-four-slot-sharp} gives a lower bound
\(c_qM^{-q/2}\) for the tensor product of \(q\) complete heat-square
slots.  A bounded remainder supplies no additional negative power.
For \(q<4\), this is asymptotically larger than \(M^{-2}\).  Equivalently,
the factor \(M^{-\delta_q}\) in
\eqref{eq:V100-qcore} is genuinely missing.  Thus it must come from
the assembled quotient/payer response or another source-occurring
factor, not from recounting the same heat slots.
\end{proof}

\begin{firewall}[What the \(q\)-slot definition does not assume]
\label{fire:V100-qcore-scope}
The card \eqref{eq:V100-qcore} is an exact remaining interface, not a
proved estimate.  V100 proves its sufficiency and the sharpness of its
exponent.  It does not assume that a complete quotient cumulant has
partial-transpose decay merely because one of its pair labels is
strong.
\end{firewall}

\section{The finite two-row cofinal core}
\label{sec:V100-two-row-core}

\begin{theorem}[Finite residual core after heat extraction]
\label{thm:V100-finite-core}
Within the one-partner class \textup{(v)} of
\Cref{thm:V99-partner-reduction}, either:
\begin{enumerate}[label=\textup{(\roman*)},leftmargin=4em]
\item a source-compatible diffuse bank exists and the word closes by
      \Cref{prop:V99-four-slot-supplement};
\item at least four usable cofinal complete slots exist and the word
      closes by \Cref{cor:V100-four-cofinal-close}; or
\item after their removal the unresolved core contains at most three
      usable cofinal coordinates, all comparable interactions of the
      maximal row use the same partner label \(ac\), and its exact
      missing target is one of the four cards
      \eqref{eq:V100-qcore}.
\end{enumerate}
The same trichotomy applies to the supplement-free two-partner
chord-or-coverage class, with two fixed partner labels.
\end{theorem}

\begin{proof}
The diffuse alternative is exactly the source-compatible hypothesis of
\Cref{prop:V99-four-slot-supplement}.  If it fails, count usable
cofinal slots.  Four or more contain a selected four-slot subbank and
close by \Cref{cor:V100-four-cofinal-close}.  Otherwise their number is
\(q\le3\); \Cref{lem:V100-response-table} and
\Cref{prop:V100-qcore-suffices} identify the exact remaining card.
The partner-label statements are the definitions of the two residual
classes in \Cref{thm:V99-partner-reduction}.
\end{proof}

\begin{assessment}[Direction after the audit]
\label{ass:V100-direction}
The open quintic problem is no longer an infinite support census.  It
is a finite family indexed by
\[
 q\in\{0,1,2,3\},\quad
 \rho\in\Pi_5\setminus\{\widehat1\},\quad
 \hbox{payer side},\quad\hbox{square orientation},
\]
with one or two fixed cofinal partner labels.  The next proof must keep
\(\mathcal J_r^\rho\) assembled and obtain exactly
\(M^{-\delta_q}\).  SM V77 and NS V31 both reinforce that proof order:
role or scale deficits cannot be erased after taking an incompatible
determinant or norm.
\end{assessment}

\section{V100 result ledger and rollover}
\label{sec:V100-status}

\begin{theorem}[Third-series V100 result ledger]
\label{thm:V100-results}
Relative to V1--V99, V100 proves:
\begin{enumerate}[label=\textup{(V100.\arabic*)},leftmargin=7em]
\item the compulsory hashed SM V77 and NS V31 audit with no
      cross-program estimate import;
\item the four-cofinal-complete-slot \(M^{-2}\) theorem;
\item closure of every obstructed word carrying such four slots;
\item the exact \(q=0,1,2,3\) residual response table;
\item sufficiency of the assembled \(q\)-slot core card;
\item the sharp \(SU(2)\) heat-only obstruction at every \(q<4\); and
\item reduction of the supplement-free one- and two-partner sectors to
      a finite assembled quotient-cumulant core.
\end{enumerate}
\end{theorem}

\begin{proof}
These are
\Cref{ass:V100-SM-audit,
ass:V100-NS-audit,
thm:V100-four-cofinal,
cor:V100-four-cofinal-close,
lem:V100-response-table,
prop:V100-qcore-suffices,
prop:V100-heat-only-obstruction,
thm:V100-finite-core}.
\end{proof}

\begin{programme}[Post-rollover V1 acquisition]
\label{prog:V100-rollover}
\begin{enumerate}[leftmargin=3em]
\item Freeze this complete validated file with suffix
      \texttt{\_V100\_f3.tex} and begin a new contextual V1.
\item Carry forward the exact quartet of response exponents
      \(\delta_q=(4-q)/2\), not the full historical case tree.
\item Split the \(51\) prefix partitions by whether the sole cofinal
      pair lies inside one prefix block or crosses two prefix blocks.
\item In the crossing class, test local cofinal alignment at the joining
      coordinate before using a global rescue copy.
\item In the internal class, keep the lower-order prefix carrier
      assembled and seek the missing response in the quotient cumulant.
\item Test every proposed \(q\)-slot response on the V91
      \(SU(2)\) singlet family before claiming uniform decay.
\end{enumerate}
\end{programme}

\begin{firewall}[Claim boundary after third-series V100]
\label{fire:V100-claim-boundary}
V100 reduces the open concentrated quintic sector to finitely many
assembled \(q\)-slot response cards with sharp exponents.  It does not
prove those cards, global quintic or all-order MTA, WUFC, CPM, cutoff
removal, reflection positivity, Osterwalder--Schrader reconstruction,
construction of the continuum Yang--Mills measure, or a uniform
positive continuum spectral gap.  The full Yang--Mills mass gap has not
been proved.
\end{firewall}

\paragraph{Parent-version record.}
V100 was formed from
\texttt{Renewal\_Yang\_Mills\_Mass\_Gap\_Research\_Notes\_V99.tex}.
The V99 parent had \(48\,008\) logical source lines,
\(1\,706\,328\) bytes, and SHA--256
\[
 \texttt{9d6284ad1f21bab8d684b75f9290830ca030aa9aa447f7fa18d92d379762f359}.
\]
This file is the third hundred-version rollover point.

% =============================================================================
% END APPENDED THIRD-SERIES V100 MATERIAL
% =============================================================================

\end{document}
